% Auto-generated from data/segments.json + layout.json (+ transforms) — do not edit by hand.
% volume: 10Ma02
% mode: reading
% segments: 1683
% transform_rules: 0
% toc_marks: 122

% Volume layout defaults (layout.json ``layout``).
\csromanlayoutapply{1}{1.5}{21.6pt}{5pt}{6.3pt}{65pt}{2.5em}%

% Reading physical-page layout (layout.json ``page_layout_reading_mode``).
\csromanreadingpagelayoutvolume{1}{1.5}{21.6pt}{5pt}{6.3pt}{65pt}{2.5em}%
\csromanreadingpagelayoutenable

\csromanheader{2}{\textbf{มชฺฌิมนิกาย}}
\csromanmatikaanchor{mtk.0}
\csromantocmarknopage{chapter}{มชฺฌิมปณฺณาสปาฬิ}{mtk.0}
\bookmark[dest={mtk.0},level=0]{มชฺฌิมปณฺณาสปาฬิ}
\gambhira{\textbf{มชฺฌิมปณฺณาสปาฬิ}}
\namakkaram{นโม ตสฺส ภควโต อรหโต สมฺมาสมฺพุทฺธสฺส.}
\csromanmatikaanchor{mtk.1}
\csromantocmarknopage{chapter}{1. คหปติวคฺค}{mtk.1}
\bookmark[dest={mtk.1},level=0]{1. คหปติวคฺค}
\chapterhead{1. คหปติวคฺค}
\csromanmatikaanchor{mtk.2}
\csromantocmarkref{subsection}{1. กนฺทรกสุตฺต}{mtk.2}
\bookmark[dest={mtk.2},level=2]{1. กนฺทรกสุตฺต}
\csromanheader{2}{1. กนฺทรกสุตฺต}
\proseitem{1}{\csromanfolio{1}เอวํ เม สุตํ\csromandash{}เอกํ สมยํ ภควา จมฺปายํ วิหรติ คคฺคราย โปกฺขรณิยา ตีเร มหตา ภิกฺขุสํเฆน สทฺธิํ.\csromanspacer{} อถ โข เปสฺโส\footnote{เปโย (ก)} จ หตฺถาโรหปุตฺโต กนฺทรโก จ ปริพฺพาชโก เยน ภควา เตนุปสงฺกมิํสุ, อุปสงฺกมิตฺวา เปสฺโส หตฺถาโรหปุตฺโต ภควนฺตํ อภิวาเทตฺวา เอกมนฺตํ นิสีทิ.\csromanspacer{} กนฺทรโก ปน ปริพฺพาชโก ภควตา สทฺธิํ สมฺโมทิ, สมฺโมทนียํ กถํ สารณียํ\footnote{สาราณียํ (สี, สฺยา, กํ, อิ)} วีติสาเรตฺวา เอกมนฺตํ อฏฺฐาสิ, เอกมนฺตํ ฐิโต โข กนฺทรโก ปริพฺพาชโก ตุณฺหีภูตํ ตุณฺหีภูตํ ภิกฺขุสํฆํ อนุวิโลเกตฺวา ภควนฺตํ เอตทโวจ, “อจฺฉริยํ โภ โคตม, อพฺภุตํ โภ โคตม, ยาวญฺจิทํ โภตา โคตเมน สมฺมา ภิกฺขุสํโฆ ปฏิปาทิโต.\csromanspacer{} เยปิ เต โภ โคตม อเหสุํ อตีตมทฺธานํ อรหนฺโต สมฺมาสมฺพุทฺธา, เตปิ ภควนฺโต เอตปรมํ เยว สมฺมา ภิกฺขุสํฆํ ปฏิปาเทสุํ, เสยฺยถาปิ เอตรหิ โภตา โคตเมน สมฺมา ภิกฺขุสํโฆ ปฏิปาทิโต.\csromanspacer{} เยปิ เต โภ โคตม ภวิสฺสนฺติ อนาคตมทฺธานํ อรหนฺโต สมฺมาสมฺพุทฺธา, เตปิ ภควนฺโต เอตปรมํ เยว สมฺมา ภิกฺขุสํฆํ ปฏิปาเทสฺสนฺติ, เสยฺยถาปิ เอตรหิ โภตา โคตเมน สมฺมา ภิกฺขุสํโฆ ปฏิปาทิโต”ติ.}

\proseitem{2}{\csromanfolio{2}เอวเมตํ กนฺทรก, เอวเมตํ กนฺทรก, เยปิ เต กนฺทรก อเหสุํ อตีตมทฺธานํ อรหนฺโต สมฺมาสมฺพุทฺธา, เตปิ ภควนฺโต เอตปรมํเยว สมฺมา ภิกฺขุสํฆํ ปฏิปาเทสุํ, เสยฺยถาปิ เอตรหิ มยา สมฺมา ภิกฺขุสํโฆ ปฏิปาทิโต.\csromanspacer{} เยปิ เต กนฺทรก ภวิสฺสนฺติ อนาคตมทฺธานํ อรหนฺโต สมฺมาสมฺพุทฺธา, เตปิ ภควนฺโต เอตปรมํเยว สมฺมา ภิกฺขุสํฆํ ปฏิปาเทสฺสนฺติ, เสยฺยถาปิ เอตรหิ มยา สมฺมา ภิกฺขุสํโฆ ปฏิปาทิโต.}

\prose{สนฺติ หิ กนฺทรก ภิกฺขู อิมสฺมิํ ภิกฺขุสํเฆ อรหนฺโต ขีณาสวา วุสิตวนฺโต กตกรณียา โอหิตภารา อนุปฺปตฺตสทตฺถา ปริกฺขีณภวสํโยชนา สมฺมทญฺญา วิมุตฺตา, สนฺติ หิ กนฺทรก ภิกฺขู อิมสฺมิํ ภิกฺขุสํเฆ เสกฺขา สนฺตตสีลา สนฺตตวุตฺติโน นิปกา นิปกวุตฺติโน, เต จตูสุ\footnote{นิปกวุตฺติโน จตูสุ (สี)} สติปฏฺฐาเนสุ สุปฺปติฏฺฐิตจิตฺตา\footnote{สุปฏฺฐิตจิตฺตา (สี, อิ, ก)} วิหรนฺติ.\csromanspacer{} กตเมสุ จตูสุ, อิธ กนฺทรก ภิกฺขุ กาเย กายานุปสฺสี วิหรติ อาตาปี สมฺปชาโน สติมา วิเนยฺย โลเก อภิชฺฌาโทมนสฺสํ, เวทนาสุ เวทนานุปสฺสี วิหรติ อาตาปี สมฺปชาโน สติมา วิเนยฺย โลเก อภิชฺฌาโทมนสฺสํ, จิตฺเต จิตฺตานุปสฺสี วิหรติ อาตาปี สมฺปชาโน สติมา วิเนยฺย โลเก อภิชฺฌาโทมนสฺสํ, ธมฺเมสุ ธมฺมานุปสฺสี วิหรติ อาตาปี สมฺปชาโน สติมา วิเนยฺย โลเก อภิชฺฌาโทมนสฺสนฺติ.}

\proseitem{3}{เอวํ วุตฺเต เปสฺโส หตฺถาโรหปุตฺโต ภควนฺตํ เอตทโวจ “อจฺฉริยํ ภนฺเต, อพฺภุตํ ภนฺเต, ยาว สุปญฺญตฺตา จิเม ภนฺเต ภควตา จตฺตาโร สติปฏฺฐานา สตฺตานํ วิสุทฺธิยา โสกปริเทวานํ\footnote{โสกปริทฺทวานํ (สี, อิ)} สมติกฺกมาย ทุกฺขโทมนสฺสานํ อตฺถงฺคมาย ญายสฺส อธิคมาย นิพฺพานสฺส สจฺฉิกิริยาย.\csromanspacer{} มยมฺปิ หิ ภนฺเต คิหี โอทาตวสนา กาเลน กาลํ อิเมสุ จตูสุ สติปฏฺฐาเนสุ สุปฺปติฏฺฐิตจิตฺตา วิหราม.\csromanspacer{} อิธ มยํ ภนฺเต กาเย กายานุปสฺสิโน วิหราม อาตปิโน สมฺปชานา สติมนฺโต วิเนยฺย โลเก อภิชฺฌาโทมนสฺสํ, เวทนาสุ เวทนานุปสฺสิโน วิหราม อาตาปิโน สมฺปชานา สติมนฺโต วิเนยฺย โลเก อภิชฺฌาโทมนสฺสํ, \csromanfolio{3}จิตฺเต จิตฺตานุปสฺสิโน วิหราม อาตาปิโน สมฺปชานา สติมนฺโต วิเนยฺย โลเก อภิชฺฌาโทมนสฺสํ, ธมฺเมสุ ธมฺมานุปสฺสิโน วิหราม อาตาปิโน สมฺปชานา สติมนฺโต วิเนยฺย โลเก อภิชฺฌาโทมนสฺสํ.\csromanspacer{} อจฺฉริยํ ภนฺเต, อพฺภุตํ ภนฺเต, ยาวญฺจิทํ ภนฺเต ภควา เอวํ มนุสฺสคหเน เอวํ มนุสฺสกสเฏ เอวํ มนุสฺสสาเฐยฺเย วตฺตมาเน สตฺตานํ หิตาหิตํ ชานาติ.\csromanspacer{} คหนเญฺหตํ ภนฺเต ยทิทํ มนุสฺสา, อุตฺตานกเญฺหตํ ภนฺเต ยทิทํ ปสโว.\csromanspacer{} อหํ หิ ภนฺเต ปโหมิ หตฺถิทมฺมํ สาเรตุํ, ยาวตเกน อนฺตเรน จมฺปํ คตาคตํ กริสฺสติ, สพฺพานิ ตานิ สาเฐยฺยานิ กูเฏยฺยานิ วงฺเกยฺยานิ ชิเมฺหยฺยานิ ปาตุกริสฺสติ.\csromanspacer{} อมฺหากํ ปน ภนฺเต ทาสาติ วา เปสฺสาติ วา กมฺมกราติ วา อญฺญถาว กาเยน สมุทาจรนฺติ, อญฺญถาว วาจาย, อญฺญถาว เนสํ จิตฺตํ โหติ.\csromanspacer{} อจฺฉริยํ ภนฺเต, อพฺภุตํ ภนฺเต, ยาวญฺจิทํ ภนฺเต ภควา เอวํ มนุสฺสคหเน เอวํ มนุสฺสกสเฏ เอวํ มนุสฺสสาเฐยฺเย วตฺตมาเน สตฺตานํ หิตาหิตํ ชานาติ.\csromanspacer{} คหนเญฺหตํ ภนฺเต ยทิทํ มนุสฺสา, อุตฺตานกเญฺหตํ ภนฺเต ยทิทํ ปสโว”ติ.}

\proseitem{4}{เอวเมตํ เปสฺส, เอวเมตํ เปสฺส, คหนเญฺหตํ เปสฺส ยทิทํ มนุสฺสา, อุตฺตานกเญฺหตํ เปสฺส ยทิทํ ปสโว.\csromanspacer{} จตฺตาโรเม เปสฺส ปุคฺคลา สนฺโต สํวิชฺชมานา โลกสฺมิํ.\csromanspacer{} กตเม จตฺตาโร, อิธ เปสฺส เอกจฺโจ ปุคฺคโล อตฺตนฺตโป โหติ อตฺตปริตาปนานุโยคมนุยุตฺโต, อิธ ปน เปสฺส เอกจฺโจ ปุคฺคโล ปรนฺตโป โหติ ปรปริตาปนานุโยคมนุยุตฺโต, อิธ ปน เปสฺส เอกจฺโจ ปุคฺคโล อตฺตนฺตโป จ โหติ อตฺตปริตาปนานุโยคมนุยุตฺโต ปรนฺตโป จ ปรปริตาปนานุโยคมนุยุตฺโต, อิธ ปน เปสฺส เอกจฺโจ ปุคฺคโล เนวตฺตนฺตโป โหติ นาตฺตปริตาปนานุโยคมนุยุตฺโต น ปรนฺตโป น ปรปริตาปนานุโยคมนุยุตฺโต, โส อนตฺตนฺตโป อปรนฺตโป ทิฏฺเฐว ธมฺเม นิจฺฉาโต นิพฺพุโต สีตีภูโต\footnote{สีติภูโต (สี, อิ, ก)} สุขปฺปฏิสํเวที พฺรหฺมภูเตน อตฺตนา วิหรติ.\csromanspacer{} อิเมสํ เปสฺส จตุนฺนํ ปุคฺคลานํ กตโม เต ปุคฺคโล จิตฺตํ อาราเธตีติ.}

\prose{ยฺวายํ ภนฺเต ปุคฺคโล อตฺตนฺตโป อตฺตปริตาปนานุโยคมนุยุตฺโต, อยํ เม ปุคฺคโล จิตฺตํ นาราเธติ.\csromanspacer{} โยปายํ ภนฺเต ปุคฺคโล ปรนฺตโป \csromanfolio{4}ปรปริตาปนานุโยคมนุยุตฺโต, อยมฺปิ เม ปุคฺคโล จิตฺตํ นาราเธติ.\csromanspacer{} โยปายํ ภนฺเต ปุคฺคโล อตฺตนฺตโป จ อตฺตปริตาปนานุโยคมนุยุตฺโต ปรนฺตโป จ ปรปริตาปนานุโยคมนุยุตฺโต, อยมฺปิ เม ปุคฺคโล จิตฺตํ นาราเธติ.\csromanspacer{} โย จ โข อยํ ภนฺเต ปุคฺคโล เนวตฺตนฺตโป นาตฺตปริตาปนานุโยคมนุยุตฺโต น ปรนฺตโป น ปรปริตาปนานุโยคมนุยุตฺโต, โส อนตฺตนฺตโป อปรนฺตโป ทิฏฺเฐว ธมฺเม นิจฺฉาโต นิพฺพุโต สีตีภูโต สุขปฺปฏิสํเวที พฺรหฺมภูเตน อตฺตนา วิหรติ, อยเมว\footnote{อยํ (สี, สฺยา, กํ, อิ)} เม ปุคฺคโล จิตฺตํ อาราเธตีติ.}

\proseitem{5}{กสฺมา ปน เต เปสฺส อิเม ตโย ปุคฺคลา จิตฺตํ นาราเธนฺตีติ.\csromanspacer{} ยฺวายํ ภนฺเต ปุคฺคโล อตฺตนฺตโป อตฺตปริตาปนานุโยคมนุยุตฺโต, โส อตฺตานํ สุขกามํ ทุกฺขปฏิกฺกูลํ อาตาเปติ ปริตาเปติ, อิมินา เม อยํ ปุคฺคโล จิตฺตํ นาราเธติ.\csromanspacer{} โยปายํ ภนฺเต ปุคฺคโล ปรนฺตโป ปรปริตาปนานุโยคมนุยุตฺโต, โส ปรํ สุขกามํ ทุกฺขปฏิกฺกูลํ อาตาเปติ ปริตาเปติ, อิมินา เม อยํ ปุคฺคโล จิตฺตํ นาราเธติ.\csromanspacer{} โยปายํ ภนฺเต ปุคฺคโล อตฺตนฺตโป จ อตฺตปริตาปนานุโยคมนุยุตฺโต ปรนฺตโป จ ปรปริตาปนานุโยคมนุยุตฺโต, โส อตฺตานญฺจ ปรญฺจ สุขกามํ ทุกฺขปฏิกฺกูลํ\footnote{สุขกาเม ทุกฺขปฏิกฺกูเล (สี. อิ) 3-3. วิหรติ. อิมินา (สี, สฺยา, กํ, อิ)} อาตาเปติ ปริตาเปติ, อิมินา เม อยํ ปุคฺคโล จิตฺตํ นาราเธติ.\csromanspacer{} โย จ โข อยํ ภนฺเต ปุคฺคโล เนวตฺตนฺตโป นาตฺตปริตาปนานุโยคมนุยุตฺโต น ปรนฺตโป น ปรปริตาปนานุโยคมนุยุตฺโต, โส อนตฺตนฺตโป อปรนฺตโป ทิฏฺเฐว ธมฺเม นิจฺฉาโต นิพฺพุโต สีตีภูโต สุขปฺปฏิสํเวที พฺรหฺมภูเตน อตฺตนา3 วิหรติ.\csromanspacer{} โส อตฺตานญฺจ ปรญฺจ สุขกามํ ทุกฺขปฏิกฺกูลํ เนว อาตาเปติ น ปริตาเปหิ, อิมินา3 เม อยํ ปุคฺคโล จิตฺตํ อาราเธติ.\csromanspacer{} หนฺท จ ทานิ มยํ ภนฺเต คจฺฉาม พหุกิจฺจา มยํ พหุกรณียาติ.\csromanspacer{} ยสฺสทานิ ตฺวํ เปสฺส กาลํ มญฺญสีติ.\csromanspacer{} อถ โข เปสฺโส หตฺถาโรหปุตฺโต ภควโต ภาสิตํ อภินนฺทิตฺวา อนุโมทิตฺวา อุฏฺฐายาสนา ภควนฺตํ อภิวาเทตฺวา ปทกฺขิณํ กตฺวา ปกฺกามิ.}

\proseitem{6}{\csromanfolio{5}อถ โข ภควา อจิรปกฺกนฺเต เปสฺเส หตฺถาโรหปุตฺเต ภิกฺขู อามนฺเตสิ “ปณฺฑิโต ภิกฺขเว เปสฺโส หตฺถาโรหปุตฺโต, มหาปญฺโญ ภิกฺขเว เปสฺโส หตฺถาโรหปุตฺโต, สเจ ภิกฺขเว เปสฺโส หตฺถาโรหปุตฺโต มุหุตฺตํ นิสีเทยฺย, ยาวสฺสาหํ อิเม จตฺตาโร ปุคฺคเล วิตฺถาเรน วิภชิสฺสามิ\footnote{วิภชามิ (สี, อิ)}, มหตา อตฺเถน สํยุตฺโต อภวิสฺส, อปิ จ ภิกฺขเว เอตฺตาวตาปิ เปสฺโส หตฺถาโรหปุตฺโต มหตา อตฺเถน สํยุตฺโต”ติ.\csromanspacer{} เอตสฺส ภควา กาโล, เอตสฺส สุคต กาโล, ยํ ภควา อิเม จตฺตาโร ปุคฺคเล วิตฺถาเรน วิภเชยฺย, ภควโต สุตฺวา ภิกฺขู ธาเรสฺสนฺตีติ.\csromanspacer{} เตน หิ ภิกฺขเว สุณาถสาธุกํ มนสิ กโรถ ภาสิสฺสามีติ. “เอวํ ภนฺเต”ติ โข เต ภิกฺขู ภควโต ปจฺจสฺโสสุํ.\csromanspacer{} ภควา เอตทโวจ\csromandash{}}

\proseitem{7}{กตโม จ ภิกฺขเว ปุคฺคโล อตฺตนฺตโป อตฺตปริตาปนานุโยคมนุยุตฺโต.\csromanspacer{} อิธ ภิกฺขเว เอกจฺโจ ปุคฺคโล อเจลโก โหติ, มุตฺตาจาโร, หตฺถาปเลขโน\footnote{หตฺถาวเลขโน (สฺยา, กํ)}, น-เอหิภทฺทนฺติโก, นติฏฺฐภทฺทนฺติโก\footnote{น-เอหิภทนฺติโก, นติฏฺฐภทนฺติโก (สี, สฺยา, กํ, อิ)}, นาภิหฏํ, น อุทฺทิสฺสกตํ, น นิมนฺตนํ สาทิยติ, โส น กุมฺภิมุขา ปฏิคฺคณฺหาติ, น กโฬปิมุขา\footnote{ขโฬปิมุขา (สี)} ปฏิคฺคณฺหาติ, น เอฬกมนฺตรํ, น ทณฺฑมนฺตรํ, น มุสลมนฺตรํ, น ทฺวินฺนํ ภุญฺชมานานํ, น คพฺภินิยา, น ปายมานาย, น ปุริสนฺตรคตาย, น สงฺกิตฺตีสุ, น ยตฺถ สา อุปฏฺฐิโต โหติ, น ยตฺถ มกฺขิกา สณฺฑสณฺฑจารินี, น มจฺฉํ, น มํสํ, น สุรํ, น เมรยํ, น ถุโสทกํ ปิวติ.\csromanspacer{} โส เอกาคาริโก วา โหติ เอกาโลปิโก, ทฺวาคาริโก วา โหติ ทฺวาโลปิโก ฯเปฯ สตฺตาคาริโก วา โหติ สตฺตาโลปิโก, เอกิสฺสาปิ ทตฺติยา ยาเปติ, ทฺวีหิปิ ทตฺตีหิ ยาเปติ ฯเปฯ สตฺตหิปิ ทตฺตีหิ ยาเปติ, เอกาหิกมฺปิ อาหารํ อาหาเรติ, ทฺวีหิกมฺปิ อาหารํ อาหาเรติ ฯเปฯ สตฺตาหิกมฺปิ อาหารํ อาหาเรติ, อิติ เอวรูปํ อฑฺฒมาสิกํ ปริยายภตฺตโภชนานุโยคมนุยุตฺโต วิหรติ.\csromanspacer{} โส สากภกฺโข วา โหติ, สามากภกฺโข วา โหติ, นีวารภกฺโข วา โหติ, ททฺทุลภกฺโข วา โหติ, หฏภกฺโข วา โหติ, กณภกฺโข วา โหติ, อาจามภกฺโข วา โหติ, ปิญฺญากภกฺโข วา โหติ, ติณภกฺโข วา โหติ, โคมยภกฺโข \csromanfolio{6}วา โหติ, วนมูลผลาหาโร ยาเปติ ปวตฺตผลโภชี.\csromanspacer{} โส สาณานิปิ ธาเรติ, มสาณานิปิ ธาเรติ, ฉวทุสฺสานิปิ ธาเรติ, ปํสุกูลานิปิ ธาเรติ, ติรีฏานิปิ ธาเรติ, อชินมฺปิ ธาเรติ, อชินกฺขิปมฺปิ ธาเรติ, กุสจีรมฺปิ ธาเรติ, วากจีรมฺปิ ธาเรติ, ผลกจีรมฺปิ ธาเรติ, เกสกมฺพลมฺปิ ธาเรติ, วาฬกมฺพลมฺปิ ธาเรติ, อุลูกปกฺขมฺปิ ธาเรติ, เกสมสฺสุโลจโกปิ โหติ เกสมสฺสุโลจนานุโยคมนุยุตฺโต, อุพฺภฏฺฐโกปิ โหติ อาสนปฏิกฺขิตฺโต, อุกฺกุฏิโกปิ โหติ อุกฺกุฏิกปฺปธานมนุยุตฺโต, กณฺฏกาปสฺสยิโกปิ โหติ กณฺฏกาปสฺสเย เสยฺยํ กปฺเปติ\footnote{ปสฺส ม 1 มหาสีหนาทสุตฺเต (111) ปิฏฺเฐ.}, สายตติยกํปิ อุทโกโรหนานุโยคมนุยุตฺโต วิหรติ.\csromanspacer{} อิติ เอวรูปํ อเนกวิหิตํ กายสฺส อาตาปนปริตาปนานุโยคมนุยุตฺโต วิหรติ.\csromanspacer{} อยํ วุจฺจติ ภิกฺขเว ปุคฺคโล อตฺตนฺตโป อตฺตปริตาปนานุโยคมนุยุตฺโต.}

\proseitem{8}{กตโม จ ภิกฺขเว ปุคฺคโล ปรนฺตโป ปรปริตาปนานุโยคมนุยุตฺโต.\csromanspacer{} อิธ ภิกฺขเว เอกจฺโจ ปุคฺคโล โอรพฺภิโก โหติ สูกริโก สากุณิโก มาควิโก ลุทฺโท มจฺฉฆาตโก โจโร โจรฆาตโก โคฆาตโก พนฺธนาคาริโก, เย วา ปนญฺเญปิ เกจิ กุรูรกมฺมนฺตา.\csromanspacer{} อยํ วุจฺจติ ภิกฺขเว ปุคฺคโล ปรนฺตโป ปรปริตาปนานุโยคมนุยุตฺโต.}

\proseitem{9}{กตโม จ ภิกฺขเว ปุคฺคโล อตฺตนฺตโป จ อตฺตปริตาปนานุโยคมนุยุตฺโต ปรนฺตโป จ ปรปริตาปนานุโยคมนุยุตฺโต.\csromanspacer{} อิธ ภิกฺขเว เอกจฺโจ ปุคฺคโล ราชา วา โหติ ขตฺติโย มุทฺธาวสิตฺโต พฺราหฺมโณ วา มหาสาโล.\csromanspacer{} โส ปุรตฺถิเมน นครสฺส นวํ สนฺถาคารํ\footnote{สนฺธาคารํ (ฏีกา)} การาเปตฺวา เกสมสฺสุํ โอหาเรตฺวา ขราชินํ นิวาเสตฺวา สปฺปิเตเลน กายํ อพฺภญฺชิตฺวา มควิสาเณน ปิฏฺฐิํ กณฺฑุวมาโน นวํ สนฺถาคารํ ปวิสติ สทฺธิํ มเหสิยา พฺราหฺมเณน จ ปุโรหิเตน.\csromanspacer{} โส ตตฺถ อนนฺตรหิตาย ภูมิยา หริตุปลิตฺตาย เสยฺยํ กปฺเปติ.\csromanspacer{} เอกิสฺสาย คาวิยา สรูปวจฺฉาย ยํ เอกสฺมิํ ถเน ขีรํ โหติ, เตน ราชา ยาเปติ.\csromanspacer{} ยํ ทุติยสฺมิํ ถเน ขีรํ โหติ, เตน มเหสี ยาเปติ.\csromanspacer{} ยํ ตติยสฺมิํ ถเน ขีรํ โหติ, เตน พฺราหฺมโณ ปุโรหิโต \csromanfolio{7}ยาเปติ.\csromanspacer{} ยํ จตุตฺถสฺมิํ ถเน ขีรํ โหติ, เตน อคฺคิํ ชุหติ.\csromanspacer{} อวเสเสน วจฺฉโก ยาเปติ.\csromanspacer{} โส เอวมาห “เอตฺตกา อุสภา หญฺญนฺตุ ยญฺญตฺถาย, เอตฺตกา วจฺฉตรา หญฺญนฺตุ ยญฺญตฺถาย, เอตฺตกา วจฺฉตริโย หญฺญนฺตุ ยญฺญตฺถาย, เอตฺตกา อชา หญฺญนฺตุ ยญฺญตฺถาย, เอตฺตกา อุรพฺภา หญฺญนฺตุ ยญฺญตฺถาย, (เอตฺตกา อสฺสา หญฺญนฺตุ ยญฺญตฺถาย)\footnote{( ) นตฺถิ สี-อิ-โปตฺถเกสุ.} เอตฺตกา รุกฺขา ฉิชฺชนฺตุ ยูปตฺถาย, เอตฺตกา ทพฺภา ลูยนฺตุ พริหิสตฺถายา”ติ\footnote{ปริหิํสตฺถาย (ก)}.\csromanspacer{} เยปิสฺส เต โหนฺติ ทาสาติ วา เปสฺสาติ วา กมฺมกราติ วา, เตปิ ทณฺฑตชฺชิตา ภยตชฺชิตา อสฺสุมุขา รุทมานา ปริกมฺมานิ กโรนฺติ.\csromanspacer{} อยํ วุจฺจติ ภิกฺขเว ปุคฺคโล อตฺตนฺตโป จ อตฺตปริตาปนานุโยคมนุยุตฺโต ปรนฺตโป จ ปรปริตาปนานุโยคมนุยุตฺโต.}

\proseitem{10}{กตโม จ ภิกฺขเว ปุคฺคโล เนวตฺตนฺตโป นาตฺตปริตาปนานุโยคมนุยุตฺโต น ปรนฺตโป น ปรปริตาปนานุโยคมนุยุตฺโต.\csromanspacer{} โส อนตฺตนฺตโป อปรนฺตโป ทิฏฺเฐว ธมฺเม นิจฺฉาโต นิพฺพุโต สีตีภูโต สุขปฺปฏิสํเวที พฺรหฺมภูเตน อตฺตนา วิหรติ.\csromanspacer{} อิธ ภิกฺขเว ตถาคโต โลเก อุปฺปชฺชติ อรหํ สมฺมาสมฺพุทฺโธ วิชฺชาจรณสมฺป สุคโต โลกวิทู อนุตฺตโต ปุริสทมฺมสารถิ สตฺถา เทวมนุสฺสานํ พุทฺโธ ภควา.\csromanspacer{} โส อิมํ โลกํ สเทวกํ สมารกํ สพฺรหฺมกํ สสฺสมณพฺราหฺมณิํ ปชํ สเทวมนุสฺสํ สยํ อภิญฺญา สจฺฉิกตฺวา ปเวเทติ, โส ธมฺมํ เทเสติ อาทิกลฺยาณํ มชฺเฌกลฺยาณํ ปริโยสานกลฺยาณํ สาตฺถํ สพฺยญฺชนํ เกวลปริปุณฺณํ ปริสุทฺธํ พฺรหฺมจริยํ ปกาเสติ, ตํ ธมฺมํ สุณาติ คหปติ วา คหปติปุตฺโต วา อญฺญตรสฺมิํ วา กุเล ปจฺจาชาโต.\csromanspacer{} โส ตํ ธมฺมํ สุตฺวา ตถาคเต สทฺธํ ปฏิลภติ, โส เตน สทฺธาปฏิลาเภน สมนฺนาคโต อิติ ปฏิสญฺจิกฺขติ “สมฺพาโธ ฆราวาโส รชาปโถ, อพฺโภกาโส ปพฺพชฺชา, นยิทํ สุกรํ อคารํ อชฺฌาวสตา เอกนฺตปริปุณฺณํ เอกนฺตปริสุทฺธํ สงฺขลิขิตํ พฺรหฺมจริยํ จริตุํ, ยํนูนาตํ เกสมสฺสุํ โอหาเรตฺวา กาสายานิ วตฺถานิ อจฺฉาเทตฺวา อคารสฺมา อนคาริยํ ปพฺพเชยฺยนฺ”ติ.\csromanspacer{} โส อปเรน สมเยน อปฺปํ วา โภคกฺขนฺธํ ปหาย มหนฺตํ วา โภคกฺขนฺธํ ปหาย อปฺปํ วา ญาติปริวฏฺฏํ \csromanfolio{8}ปหาย มหนฺตํ วา ญาติปริวฏฺฏํ ปหาย เกสมสฺสุํ โอหาเรตฺวา กาสายานิ วตฺถานิ อจฺฉาเทตฺวา อคารสฺมา อนคาริยํ ปพฺพชติ.}

\proseitem{11}{โส เอวํ ปพฺพชิโต สมาโน ภิกฺขูนํ สิกฺขาสาชีวสมาปนฺโน ปาณาติปาตํ ปหาย ปาณาติปาตา ปฏิวิรโต โหติ นิหิตทณฺโฑ นิหิตสตฺโถ ลชฺชี ทยาปนฺโน สพฺพปาณภูตหิตานุกมฺปี วิหรติ.\csromanspacer{} อทินฺนาทานํ ปหาย อทินฺนาทานา ปฏิวิรโต โหติ ทินฺนาทายี ทินฺนปาฏิกงฺขี, อเถเนน สุจิภูเตน อตฺตนา วิหรติ.\csromanspacer{} อพฺรหฺมจริยํ ปหาย พฺรหฺมจารี โหติ อาราจารี วิรโต เมถุนา คามธมฺมา.\csromanspacer{} มุสาวาทํ ปหาย มุสาวาทา ปฏิวิรโต โหติ สจฺจวาที สจฺจสนฺโธ เถโต ปจฺจยิโก อวิสํวาทโก โลกสฺส.\csromanspacer{} ปิสุณํ วาจํ ปหาย ปิสุณาย วาจาย ปฏิวิรโต โหติ, อิโต สุตฺวา น อมุตฺร อกฺขาตา อิเมสํ เภทาย, อมุตฺร วา สุตฺวา น อิเมสํ อกฺขาตา อมูสํ เภทาย, อิติ ภินฺนานํ วา สนฺธาตา สหิตานํ วา อนุปฺปทาตา สมคฺคาราโม สมคฺครโต สมคฺคนนฺที สมคฺคกรณิํ วาจํ ภาสิตา โหติ.\csromanspacer{} ผรุสํ วาจํ ปหาย ผรุสาย วาจาย ปฏิวิรโต โหติ, ยา สา วาจา เนลา กณฺณสุขา เปมนียา หทยงฺคมา โปรี พหุชนกนฺตา พหุชนมนาปา, ตถารูปิํ วาจํ ภาสิตา โหติ.\csromanspacer{} สมฺผปฺปลาปํ ปหาย สมฺผปฺปลาปา ปฏิวิรโต โหติ กาลวาที ภูตวาที อตฺถวาที ธมฺมวาที วินยวาที, นิธานวติํ วาจํ ภาสิตา กาเลน สาปเทสํ ปริยนฺตวติํ อตฺถสํหิตํ.\csromanspacer{} โส พีชคามภูตคามสมารมฺภา ปฏิวิรโต โหติ.\csromanspacer{} เอกภตฺติโก โหติ รตฺตูปรโต, วิรโต วิกาลโภชนา.\csromanspacer{} นจฺจคีตวาทิตวิสูกทสฺสนา ปฏิวิรโต โหติ.\csromanspacer{} มาลาคนฺธ วิเลปน ธารณ มณฺฑน วิภูสน ฏฺฐานา ปฏิวิรโต โหติ.\csromanspacer{} อุจฺจาสยนมหาสยนา ปฏิวิรโต โหติ.\csromanspacer{} ชาตรูปรชตปฏิคฺคหณา ปฏิวิรโต โหติ.\csromanspacer{} อามกธญฺญปฏิคฺคหณา ปฏิวิรโต โหติ.\csromanspacer{} อามกมํสปฏิคฺคหณา ปฏิวิรโต โหติ.\csromanspacer{} อิตฺถิกุมาริกปฏิคฺคหณา ปฏิวิรโต โหติ.\csromanspacer{} ทาสิทาสปฏิคฺคหณา ปฏิวิรโต โหติ.\csromanspacer{} อเชฬกปฏิคฺคหณา ปฏิวิรโต โหติ.\csromanspacer{} กุกฺกุฏสูกรปฏิคฺคหณา ปฏิวิรโต โหติ.\csromanspacer{} หตฺถิควสฺสวฬวปฏิคฺคหณา ปฏิวิรโต โหติ.\csromanspacer{} เขตฺตวตฺถุปฏิคฺคหณา ปฏิวิรโต โหติ.\csromanspacer{} ทูเตยฺยปหิณคมนานุโยคา ปฏิวิรโต โหติ.\csromanspacer{} กยวิกฺกยา \csromanfolio{9}ปฏิวิรโต โหติ.\csromanspacer{} ตุลากูฏกํสกูฏมานกูฏา ปฏิวิรโต โหติ.\csromanspacer{} อุกฺโกฏนวญฺจนนิกติสาจิโยคา\footnote{สาวิโยคา (สฺยา, กํ, ก) สาจิ กุฏิลปริยาโย.} ปฏิวิรโต โหติ.\csromanspacer{} เฉทนวธพนฺธนวิปราโมส-อาโลปสหสาการา ปฏิวิรโต โหติ 2.}

\prose{โส สนฺตุฏฺโฐ โหติ กายปริหาริเกน จีวเรน กุจฺฉิปริหาริเกน ปิณฺฑปาเตน, โส เยน เยเนว ปกฺกมติ, สมาทาเยว ปกฺกมติ.\csromanspacer{} เสยฺยถาปิ นาม ปกฺขี สกุโณ เยน เยเนว เฑติ, สปตฺตภาโรว เฑติ.\csromanspacer{} เอวเมว ภิกฺขุ สนฺตุฏฺโฐ โหติ กายปริหาริเกน จีวเรน กุจฺฉิปริหาริเกน ปิณฺฑปาเตน.\csromanspacer{} โส เยน เยเนว ปกฺกมติ, สมาทาเยว ปกฺกมติ.\csromanspacer{} โส อิมินา อริเยน สีลกฺขนฺเธน สมนฺนาคโต อชฺฌตฺตํ อนวชฺชสุขํ ปฏิสํเวเทติ.}

\proseitem{12}{โส จกฺขุนา รูปํ ทิสฺวา น นิมิตฺตคฺคาหี โหติ นานุพฺยญฺชนคฺคาหี, ยตฺวาธิกรณเมนํ จกฺขุนฺทฺริยํ อสํวุตํ วิหรนฺตํ อภิชฺฌาโทมนสฺสา ปาปกา อกุสลา ธมฺมา อนฺวาสฺสเวยฺยุํ, ตสฺส สํวราย ปฏิปชฺชติ, รกฺขติ จกฺขุนฺทฺริยํ, จกฺขุนฺทฺริเย สํวรํ อาปชฺชติ.\csromanspacer{} โสเตน สทฺทํ สุตฺวา ฯเปฯ.\csromanspacer{} ฆาเนน คนฺธํ ฆายิตฺวา ฯเปฯ.\csromanspacer{} ชิวฺหาย รสํ สายิตฺวา ฯเปฯ.\csromanspacer{} กาเยน โผฏฺฐพฺพํ ผุสิตฺวา ฯเปฯ.\csromanspacer{} มนสา ธมฺมํ วิญฺญาย น นิมิตฺตคฺคาหี โหติ นานุพฺยญฺชนคฺคาหี, ยตฺวาธิกรณเมนํ มนินฺทฺริยํ อสํวุตํ วิหรนฺตํ อภิชฺฌาโทมนสฺสา ปาปกา อกุสลา ธมฺมา อนฺวาสฺสเวยฺยุํ, ตสฺส สํวราย ปฏิปชฺชติ, รกฺขติ มนินฺทฺริยํ, มนินฺทฺริเย สํวรํ อาปชฺชติ.\csromanspacer{} โส อิมินา อริเยน อินฺทฺริยสํวเรน สมนฺนาคโต อชฺฌตฺตํ อพฺยาเสกสุขํ ปฏิสํเวเทติ.}

\prose{โส อภิกฺกนฺเต ปฏิกฺกนฺเต สมฺปชานการี โหติ, อาโลกิเต วิโลกิเต สมฺปชานการี โหติ, สมิญฺชิเต ปสาริเต สมฺปชานการี โหติ, สํฆาฏิปตฺตจีวรธารเณ สมฺปชานการี โหติ, อสิเต ปีเต ขายิเต สายิเต สมฺปชานการี โหติ, อุจฺจารปสฺสาวกมฺเม สมฺปชานการี โหติ, คเต ฐิเต นิสินฺเน สุตฺเต ชาคริเต ภาสิเต ตุณฺหีภาเว สมฺปชานการี โหติ.}

\proseitem{13}{\csromanfolio{10}โส อิมินา จ อริเยน สีลกฺขนฺเธน สมนฺนาคโต (อิมาย จ อริยาย สนฺตุฏฺฐิยา สมนฺนาคโต)\footnote{ปสฺส ม 1 จูฬหตฺถิปโทปเม (239) ปิฏฺเฐ.} อิมินา จ อริเยน อินฺทฺริยสํวเรน สมนฺนาคโต อิมินา จ อริเยน สติสมฺปชญฺเญน สมนฺนาคโต วิวิตฺตํ เสนาสนํ ภชติ อรญฺญํ รุกฺขมูลํ ปพฺพตํ กนฺทรํ คิริคุหํ สุสานํ วนปตฺถํ อพฺโภกาสํ ปลาลปุญฺชํ.\csromanspacer{} โส ปจฺฉาภตฺตํ ปิณฺฑปาตปฏิกฺกนฺโต นิสีทติ ปลฺลงฺกํ อาภุชิตฺวา อุชุํ กายํ ปณิธาย ปริมุขํ สติํ อุปฏฺฐเปตฺวา, โส อภิชฺฌํ โลเก ปหาย วิคตาภิชฺเฌน เจตสา วิหรติ, อภิชฺฌาย จิตฺตํ ปริโสเธติ.\csromanspacer{} พฺยาปาทปโทสํ ปหาย อพฺยาปนฺนจิตฺโต วิหรติ สพฺพปาณภูตหิตานุกมฺปี, พฺยาปาทปโทสา จิตฺตํ ปริโสเธติ.\csromanspacer{} ถินมิทฺธํ ปหาย วิคตถินมิทฺโธ วิหรติ อาโลกสญฺญี สโต สมฺปชาโน, ถินมิทฺธา จิตฺตํ ปริโสเธติ.\csromanspacer{} อุทฺธจฺจกุกฺกุจฺจํ ปหาย อนุทฺธโต วิหรติ อชฺฌตฺตํ วูปสนฺตจิตฺโต, อุทฺธจฺจกุกฺกุจฺจา จิตฺตํ ปริโสเธติ.\csromanspacer{} วิจิกิจฺฉํ ปหาย ติณฺณวิจิกิจฺโฉ วิหรติ อกถํกถี กุสเลสุ ธมฺเมสุ, วิจิกิจฺฉาย จิตฺตํ ปริโสเธติ.}

\prose{โส อิเม ปญฺจ นีวรเณ ปหาย เจตโส อุปกฺกิเลเส ปญฺญาย ทุพฺพลีกรเณ วิวิจฺเจว กาเมหิ วิวิจฺจ อกุสเลหิ ธมฺเมหิ สวิตกฺกํ สวิจารํ วิเวกชํ ปีติสุขํ ปฐมํ ฌานํ อุปสมฺปชฺช วิหรติ.\csromanspacer{} วิตกฺกวิจารานํ วูปสมา อชฺฌตฺตํ สมฺปสาทนํ เจตโส เอโกทิภาวํ อวิตกฺกํ อวิจารํ สมาธิชํ ปีติสุขํ ทุติยํ ฌานํ อุปสมฺปชฺช วิหรติ.\csromanspacer{} ปีติยา จ วิราคา อุเปกฺขโก จ วิหรติ สโต จ สมฺปชาโน สุขญฺจ กาเยน ปฏิสํเวเทติ, ยํ ตํ อริยา อาจิกฺขนฺติ “อุเปกฺขโก สติมา สุขวิหารี”ติ, ตติยํ ฌานํ อุปสมฺปชฺช วิหรติ.\csromanspacer{} สุขสฺส จ ปหานา ทุกฺขสฺส จ ปหานา ปุพฺเพว โสมนสฺสโทมนสฺสานํ อตฺถงฺคมา อทุกฺขมสุขํ อุเปกฺขาสติปาริสุทฺธิํ จตุตฺถํ ฌานํ อุปสมฺปชฺช วิหรติ.}

\proseitem{14}{โส เอวํ สมาหิเต จิตฺเต ปริสุทฺเธ ปริโยทาเต อนงฺคเณ วิคตูปกฺกิเลเส มุทุภูเต กมฺมนิเย ฐิเต อาเนญฺชปฺปตฺเต ปุพฺเพนิวาสา นุสฺสติญาณาย จิตฺตํ อภินินฺนาเมติ.\csromanspacer{} โส อเนกวิหิตํ ปุพฺเพนิวาสํ อนุสฺสรติ, เสยฺยถิทํ, เอกมฺปิ ชาติํ เทฺวปิ ชาติโย ติสฺโสปิ ชาติโย จตสฺโสปิ ชาติโย ปญฺจปิ ชาติโย ทสปิ ชาติโย \csromanfolio{11}วีสมฺปิ ชาติโย ติํสมฺปิ ชาติโย จตฺตาลีสมฺปี ชาติโย ปญฺญาสมฺปิ ชาติโย ชาติสตมฺปิ ชาติสหสฺสมฺปิ ชาติสตสหสฺสมฺปิ อเนเกปิ สํวฏฺฏกปฺเป อเนเกปิ วิวฏฺฏกปฺเป อเนเกปิ สํวฏฺฏวิวฏฺฏกปฺเป, “อมุตฺราสิํ เอวํนาโม เอวํโคตฺโต เอวํวณฺโณ เอวมาหาโร เอวํสุขทุกฺขปฺปฏิสํเวที เอวมายุปริยนฺโต, โส ตโต จุโต อมุตฺร อุทปาทิํ, ตตฺราปาสิํ เอวํนาโม เอวํโคตฺโต เอวํวณฺโณ เอวมาหาโร เอวํสุขทุกฺขปฺปฏิสํเวที เอวมายุปริยนฺโต, โส ตโต จุโต อิธูปปนฺโน”ติ, อิติ สาการํ ส-อุทฺเทสํ อเนกวิหิตํ ปุพฺเพนิวาสํ อนุสฺสรติ.}

\proseitem{15}{โส เอวํ สมาหิเต จิตฺเต ปริสุทฺเธ ปริโยทาเต อนงฺคเณ วิคตูปกฺกิเลเส มุทุภูเต กมฺมนิเย ฐิเต อาเนญฺชปฺปตฺเต สตฺตานํ จุตูปปาตญาณาย จิตฺตํ อภินินฺนาเมติ.\csromanspacer{} โส ทิพฺเพน จกฺขุนา วิสุทฺเธน อติกฺกนฺตมานุสเกน สตฺเต ปสฺสติ จวมาเน อุปปชฺชมาเน หีเน ปณีเต สุวณฺเณ ทุพฺพณฺเณ สุคเต ทุคฺคเต ยถากมฺมูปเค สตฺเต ปชานาติ.\csromanspacer{} อิเม วต โภนฺโต สตฺตา กายทุจฺจริเตน สมนฺนาคตา วิจีทุจฺจริเตน สมนฺนาคตา มโนทุจฺจริเตน สมนฺนาคตา, อริยานํ อุปาวาทกา มิจฺฉาทิฏฺฐิกา มิจฺฉาทิฏฺฐิกมฺมสมาทานา, เต กายสฺส เภทา ปรํ มรณา อปายํ ทุคฺคติํ วินิปาตํ นิรยํ อุปปนฺนา.\csromanspacer{} อิเม วา ปน โภนฺโต สตฺตา กายสุจริเตน สมนฺนาคตา วจีสุจริเตน สมนฺนาคตา มโนสุจริเตน สมนฺนาคตา อริยานํ อนุปวาทกา สมฺมาทิฏฺฐิกา สมฺมาทิฏฺฐิกมฺมสมาทานา, เต กายสฺส เภทา ปรํ มรณา สุคติํ สคฺคํ โลกํ อุปปนฺนาติ.\csromanspacer{} อิติ ทิพฺเพน จกฺขุนา วิสุทฺเธน อติกฺกนฺตมานุสเกน สตฺเต ปสฺสติ จวมาเน อุปปชฺชมาเน หีเน ปณีเต สุวณฺเณ ทุพฺพณฺเณ สุคเต ทุคฺคเต ยถากมฺมูปเค สตฺเต ปชานาติ.}

\proseitem{16}{โส เอวํ สมาหิเต จิตฺเต ปริสุทฺเธ ปริโยทาเต อนงฺคเณ วิคตูปกฺกิเลเส มุทุภูเต กมฺมนิเย ฐิเต อาเนญฺชปฺปตฺเต อาสวานํ ขยญาณาย จิตฺตํ อภินินฺนาเมติ.\csromanspacer{} โส อิทํ ทุกฺขนฺติ ยถาภูตํ ปชานาติ, อยํ ทุกฺขสมุทโยติ ยถาภูตํ ปชานาติ, อยํ ทุกฺขนิโรโธติ ยถาภูตํ ปชานาติ, อยํ ทุกฺขนิโรธคามินี ปฏิปทาติ ยถาภูตํ ปชานาติ.\csromanspacer{} อิเม อาสวาติ ยถาภูตํ ปชานาติ, อยํ อาสวสมุทโยติ ยถาภูตํ ปชานาติ, อยํ อาสวนิโรโธติ ยถาภูตํ \csromanfolio{12}ปชานาติ, อยํ อาสวนิโรธคามินี ปฏิปทาติ ยถาภูตํ ปชานาติ.\csromanspacer{} ตสฺส เอวํ ชานโต เอวํ ปสฺสโต กามาสวาปิ จิตฺตํ วิมุจฺจติ, ภวาสวาปิ จิตฺตํ วิมุจฺจติ, อวิชฺชาสวาปิ จิตฺตํ วิมุจฺจติ, วิมุตฺตสฺมิํ วิมุตฺตมิติ ญาณํ โหติ, “ขีณา ชาติ วุสิตํ พฺรหฺมจริยํ กตํ กรณียํ นาปรํ อิตฺถตฺตายา”ติ ปชานาติ.\csromanspacer{} อยํ วุจฺจติ ภิกฺขเว ปุคฺคโล เนวตฺตนฺตโป นาตฺตปริตาปนานุโยคมนุยุตฺโต น ปรนฺตโป น ปรปริตาปนานุโยคมนุยุตฺโต, โส อนตฺตนฺตโป อปรนฺตโป ทิฏฺเฐว ธมฺเม นิจฺฉาโต นิพฺพุโต สีตีภูโต สุขปฺปฏิสํเวที พฺรหฺมภูเตน อตฺตนา วิหรตีติ.}

\prose{อิทมโวจ ภควา.\csromanspacer{} อตฺตมนา เต ภิกฺขู ภควโต ภาสิตํ อภินนฺทุนฺติ.}

\nitthitamruled{กนฺทรกสุตฺตํ นิฏฺฐิตํ ปฐมํ.}
\csromanmatikaanchor{mtk.3}
\csromantocmarkref{subsection}{2. อฏฺฐกนาครสุตฺต}{mtk.3}
\bookmark[dest={mtk.3},level=2]{2. อฏฺฐกนาครสุตฺต}
\csromanheader{2}{2. อฏฺฐกนาครสุตฺต}
\proseitem{17}{เอวํ เม สุตํ\csromandash{}เอกํ สมยํ อายสฺมา อานนฺโท เวสาลิยํ วิหรติ เพลุวคามเก\footnote{เวฬุวคามเก (สฺยา, กํ, ก)}.\csromanspacer{} เตน โข ปน สมเยน ทสโม คหปติ อฏฺฐกนาคโร ปาฏลิปุตฺตํ อนุปฺปตฺโต โหติ เกนจิเทว กรณีเยน.\csromanspacer{} อถ โข ทสโม คหปติ อฏฺฐกนาคโร เยน กุกฺกุฏาราโม, เยน อญฺญตโร ภิกฺขุ เตนุปสงฺกมิ, อุปสงฺกมิตฺวา ตํ ภิกฺขุํ อภิวาเทตฺวา เอกมนฺตํ นิสีทิ, เอกมนฺตํ นิสินฺโน โข ทสโม คหปติ อฏฺฐกนาคโร ตํ ภิกฺขุํ เอตทโวจ “กหํ นุ โข ภนฺเต อายสฺมา อานนฺโท เอตรหิ วิหรติ, ทสฺสนกามา หิ มยํ ตํ อายสฺมนฺตํ อานนฺทนฺ”ติ.\csromanspacer{} เอโส คหปติ อายสฺมา อานนฺโท เวสาลิยํ วิหรติ เพลุวคามเกติ.\csromanspacer{} อถ โข ทสโม คหปติ อฏฺฐกนาคโร ปาฏลิปุตฺเต ตํ กรณียํ ตีเรตฺวา เยน เวสาลี, เยน เพลุวคามโก, เยนายสฺมา อานนฺโท เตนุปสงฺกมิ, อุปสงฺกมิตฺวา อายสฺมนฺตํ อานนฺทํ อภิวาเทตฺวา เอกมนฺตํ นิสีทิ.}

\proseitem{18}{เอกมนฺตํ นิสินฺโน โข ทสโม คหปติ อฏฺฐกนาคโร อายสฺมนฺตํ อานนฺทํ เอตทโวจ “อตฺถิ นุ โข ภนฺเต อานนฺท เตน}

\vspace{\tipitakaparskip}
\csromancenter{อฏฺฐกนาครสุตฺต (52) ภควตา ชานตา ปสฺสตา อรหตา สมฺมาสมฺพุทฺเธน เอกธมฺโม อกฺขาโต, ยตฺถ ภิกฺขุโน อปฺปมตฺตส อาตาปิโน ปหิตตฺตสฺส วิหรโต อวิมุตฺตญฺเจว จิตฺตํ วิมุจฺจติ, อปริกฺขีณา จ อาสวา ปริกฺขยํ คจฺฉนฺติ, อนนุปฺปตฺตญฺจ อนุตฺตรํ โยคกฺเขมํ อนุปาปุณาตี”ติ.}
\prose{\csromanfolio{13}อตฺถิ โข คหปติ เตน ภควตา ชานตา ปสฺสตา อรหตา สมฺมาสมฺพุทฺเธน เอกธมฺโม อกฺขาโต, ยตฺถ ภิกฺขุโน อปฺปมตฺตสฺส อาตาปิโน ปหิตตฺตสฺส วิหรโต อวิมุตฺตญฺเจว จิตฺตํ วิมุจฺจติ, อปริกฺขีณา จ อาสวา ปริกฺขยํ คจฺฉนฺติ, อนนุปฺปตฺตญฺจ อนุตฺตรํ โยคกฺเขมํ อนุปาปุณาตีติ.}

\prose{กตโม ปน ภนฺเต อานนฺท เตน ภควตา ชานตา ปสฺสตา อรหตา สมฺมาสมฺพุทฺเธน เอกธมฺโม อกฺขาโต, ยตฺถ ภิกฺขุโน อปฺปมตฺตสฺส อาตาปิโน ปหิตตฺตสฺส วิหรโต อวิมุตฺตญฺเจว จิตฺตํ วิมุจฺจติ, อปริกฺขีณา จ อาสวา ปริกฺขยํ คจฺฉนฺติ, อนนุปฺปตฺตญฺจ อนุตฺตรํ โยคกฺเขมํ อนุปาปุณาตีติ.}

\proseitem{19}{อิธ คหปติ ภิกฺขุ วิวิจฺเจว กาเมหิ วิวิจฺจ อกุสเลหิ ธมฺเมหิ สวิตกฺกํ สวิจารํ วิเวกชํ ปีติสุขํ ปฐมํ ฌานํ อุปสมฺปชฺช วิหรติ.\csromanspacer{} โส อิติ ปฏิสญฺจิกฺขติ “อิทมฺปิ โข ปฐมํ ฌานํ อภิสงฺขตํ อภิสญฺเจตยิตํ.\csromanspacer{} ยํ โข ปน กิญฺจิ อภิสงฺขตํ อภิสญฺเจตยิตํ, ตทนิจฺจํ นิโรธธมฺมนฺ”ติ ปชานาติ.\csromanspacer{} โส ตตฺถ ฐิโต อาสวานํ ขยํ ปาปุณาติ.\csromanspacer{} โน เจ อาสวานํ ขยํ ปาปุณาติ, เตเนว ธมฺมราเคน ตาย ธมฺมนนฺทิยา ปญฺจนฺนํ โอรมฺภาคิยานํ สํโยชนานํ ปริกฺขยา โอปปาติโก โหติ ตตฺถ ปรินิพฺพายี อนาวตฺติธมฺโม ตสฺมา โลกา.\csromanspacer{} อยมฺปิ โข คหปติ เตน ภควตา ชานตา ปสฺสตา อรหตา สมฺมาสมฺพุทฺเธน เอกธมฺโม อกฺขาโต, ยตฺถ ภิกฺขุโน อปฺปมตฺตสฺส อาตาปิโน ปหิตตฺตสฺส วิหรโต อวิมุตฺตญฺเจว จิตฺตํ วิมุจฺจติ, อปริกฺขีณา จ อาสวา ปริกฺขยํ คจฺฉนฺติ, อนนุปฺปตฺตญฺจ อนุตฺตรํ โยคกฺเขมํ อนุปาปุณาติ. (1)}

\proseitem{20}{ปุน จปรํ คหปติ ภิกฺขุ วิตกฺกวิจารานํ วูปสมา อชฺฌตฺตํ สมฺปสาทนํ ฯเปฯ ทุติยํ ฌานํ อุปสมฺปชฺช วิหรติ.\csromanspacer{} โส อิติ ปฏิสญฺจิกฺขติ อิทมฺปิ โข ทุติยํ ฌานํ อภิสงฺขตํ อภิสญฺเจตยิตํ ฯเปฯ อนุตฺตรํ โยคกฺเขมํ อนุปาปุณาติ. (2)}

\prose{\csromanfolio{14}ปุน จปรํ คหปติ ภิกฺขุ ปีติยา จ วิราคา ฯเปฯ ตติยํ ฌานํ อุปสมฺปชฺช วิหรติ.\csromanspacer{} โส อิติ ปฏิสญฺจิกฺขติ อิทมฺปิ โข ตติยํ ฌานํ อภิสงฺขตํ อภิสญฺเจตยิตํ ฯเปฯ อนุตฺตรํ โยคกฺเขมํ อนุปาปุณาติ. (3)}

\prose{ปุน จปรํ คหปติ ภิกฺขุ สุขสฺส จ ปหานา ฯเปฯ จตุตฺถํ ฌานํ อุปสมฺปชฺช วิหรติ.\csromanspacer{} โส อิติ ปฏิสญฺจิกฺขติ อิทมฺปิ โข จตุตฺถํ ฌานํ อภิสงฺขตํ อภิสญฺเจตยิตํ ฯเปฯ อนุตฺตรํ โยคกฺเขมํ อนุปาปุณาติ. (4)}

\prose{ปุน จปรํ คหปติ ภิกฺขุ เมตฺตาสหคเตน เจตสา เอกํ ทิสํ ผริตฺวา วิหรติ.\csromanspacer{} ตถา ทุติยํ.\csromanspacer{} ตถา ตติยํ.\csromanspacer{} ตถา จตุตฺถํ\footnote{จตุตฺถิํ (สี, อิ)}.\csromanspacer{} อิติ อุทฺธมโธ ติริยํ สพฺพธิ สพฺพตฺตตาย สพฺพาวนฺตํ โลกํ เมตฺตาสหคเตน เจตสา วิปุเลน มหคฺคเตน อปฺปมาเณน อเวเรน อพฺยาพชฺเฌน\footnote{อพฺยาปชฺเฌน (สี, สฺยา, อิ), อพฺยาปชฺเชน (ก) องฺคุตฺตรติกนิปาตฏีกา โอโลเกตพฺพา.} ผริตฺวา วิหรติ.\csromanspacer{} โส อิติ ปฏิสญฺจิกฺขติ “อยมฺปิ โข เมตฺตาเจโตวิมุตฺติ อภิสงฺขตา อภิสญฺเจตยิตา.\csromanspacer{} ยํ โข ปน กิญฺจิ อภิสงฺขตํ อภิสญฺเจตยิตํ, ตทนิจฺจํ นิโรธธมฺมนฺ”ติ ปชานาติ.\csromanspacer{} โส ตตฺถ ฐิโต ฯเปฯ อนุตฺตรํ โยคกฺเขมํ อนุปาปุณาติ. (5)}

\prose{ปุน จปรํ คหปติ ภิกฺขุ กรุณาสหคเตน เจตสา ฯเปฯ มุทิตาสหคเตน เจตสา ฯเปฯ อุเปกฺขาสหคเตน เจตสา เอกํ ทิสํ ผุริตฺวา วิหรติ.\csromanspacer{} ตถา ทุติยํ.\csromanspacer{} ตถา ตติยํ.\csromanspacer{} ตถา จตุตฺถํ.\csromanspacer{} อิติ อุทฺธมโธ ติริตํ สพฺพธิ สพฺพตฺตตาย สพฺพาวนฺตํ โลกํ อุเปกฺขาสหคเตน เจตสา วิปุเลน มหคฺคเตน อปฺปมาเณน อเวเรน อพฺยาพชฺเฌน ผริตฺวา วิหรติ.\csromanspacer{} โส อิติ ปฏิสญฺจิกฺขติ “อยมฺปิ โข อุเปกฺขาเจโตวิมุตฺติ อภิสงฺขตา อภิสญฺเจตยิตา.\csromanspacer{} ยํ โข ปน กิญฺจิ อภิสงฺขตํ อภิสญฺเจตยิตํ, ตทนิจฺจํ นิโรธธมฺมนฺ”ติ ปชานาติ.\csromanspacer{} โส ตตฺถ ฐิโต ฯเปฯ อนุตฺตรํ โยคกฺเขมํ อนุปาปุณาติ. (6-7-8)}

\prose{ปุน จปรํ คหปติ ภิกฺขุ สพฺพโส รูปสญฺญานํ สมติกฺกมา ปฏิฆสญฺญานํ อตฺถงฺคมา นานตฺตสญฺญานํ อมนสิการา “อนนฺโต อากาโส”ติ อากาสานญฺจายตนํ อุปสมฺปชฺช วิหรติ.\csromanspacer{} โส อิติ ปฏิสญฺจิกฺขติ “อยมฺปิ โข อากาสานญฺจายตนสมาปตฺติ อภิสงฺขตา อภิสญฺเจตยิตา.\csromanspacer{} ยํ โข ปน กิญฺจิ อภิสงฺขตํ อภิสญฺเจตยิตํ,}

\vspace{\tipitakaparskip}
\csromancenter{อฏฺฐกนาครสุตฺต (52) ตทนิจฺจํ นิโรธธมฺมนฺติ ปชานาติ.\csromanspacer{} โส ตตฺถ ฐิโต ฯเปฯ อนุตฺตรํ โยคกฺเขมํ อนุปาปุณาติ. (9)}
\prose{\csromanfolio{15}ปุน จปรํ คหปติ ภิกฺขุ สพฺพโส อากาสานญฺจายตนํ สมติกฺกมฺม “อนนฺตํ วิญฺญาณนฺ”ติ วิญฺญาณญฺจายตนํ อุปสมฺปชฺช วิหรติ.\csromanspacer{} โส อิติ ปฏิสญฺจิกฺขติ “อยมฺปิ โข วิญฺญาณญฺจายตนสมาปตฺติ อภิสงฺขตา อภิสญฺเจตยิตา.\csromanspacer{} ยํ โข ปน กิญฺจิ อภิสงฺขตํ อภิสญฺเจตยิตํ, ตทนิจฺจํ นิโรธธมฺมนฺติ ปชานาติ.\csromanspacer{} โส ตตฺถ ฐิโต ฯเปฯ อนุตฺตรํ โยคกฺเขมํ อนุปาปุณาติ. (10)}

\prose{ปุน จปรํ คหปติ ภิกฺขุ สพฺพโส วิญฺญาณญฺจายตนํ สมติกฺกมฺม “นตฺถิ กิญฺจี”ติ อากิญฺจญฺญายตนํ อุปสมฺปชฺช วิหรติ.\csromanspacer{} โส อิติ ปฏิสญฺจิกฺขติ “อยมฺปิ โข อากิญฺจญฺญายตนสมาปตฺติ อภิสงฺขตา อภิสญฺเจตยิตา.\csromanspacer{} ยํ โข ปน กิญฺจิ อภิสงฺขตํ อภิสญฺเจตยิตํ, ตทนิจฺจํ นิโรธธมฺมนฺ”ติ ปชานาติ.\csromanspacer{} โส ตตฺถ ฐิโต อาสวานํ ขยํ ปาปุณาติ.\csromanspacer{} โน เจ อาสวานํ ขยํ ปาปุณาติ, เตเนว ธมฺมราเคน ตาย ธมฺมนนฺทิยา ปญฺจนฺนํ โอรมฺภาคิยานํ สํโยชนานํ ปริกฺขยา โอปปาติโก โหติ ตตฺถ ปรินิพฺพายี อนาวตฺติธมฺโม ตสฺมา โลกา.\csromanspacer{} อยมฺปิ โข คหปติ เตน ภควตา ชานตา ปสฺสตา อรหตา สมฺมาสมฺพุทฺเธน เอกธมฺโม อกฺขาโต, ยตฺถ ภิกฺขุโน อปฺปมตฺตสฺส อาตาปิโน ปหิตตฺตสฺส วิหรโต อวิมุตฺตญฺเจว จิตฺตํ วิมุจฺจติ, อปริกฺขีณา จ อาสวา ปริกฺขยํ คจฺฉนฺติ, อนนุปฺปตฺตญฺจ อนุตฺตรํ โยคกฺเขมํ อนุปาปุณาตีติ. (11)}

\proseitem{21}{เอวํ วุตฺเต ทสโม คหปติ อฏฺฐกนาคโร อายสฺมนฺตํ อานนฺทํ เอตทโวจ “เสยฺยถาปิ ภนฺเต อานนฺท ปุริโส เอกํว นิธิมุขํ คเวสนฺโต สกิเทว เอกาทส นิธิมุขานิ อธิคจฺเฉยฺย, เอวเมว โข อหํ ภนฺเต เอกํ อมตทฺวารํ คเวสนฺโต สกิเทว\footnote{สกิํเทว (ก)} เอกาทส อมตทฺวารานิ อลตฺถํ ภาวนาย.\csromanspacer{} เสยฺยถาปิ ภนฺเต ปุริสสฺส อคารํ เอกาทสทฺวารํ, โส ตสฺมิํ อคาเร อาทิตฺเต เอกเมเกนปิ ทฺวาเรน สกฺกุเณยฺย อตฺตานํ โสตฺถิํ กาตุํ, เอวเมว โข อหํ ภนฺเต อิเมสํ เอกาทสนฺนํ อมตทฺวารานํ เอกเมเกนปิ อมตทฺวาเรน สกฺกุณิสฺสามิ อตฺตานํ โสตฺถิํ กาตุํ.\csromanspacer{} อิเมหิ นาม ภนฺเต อญฺญาติตฺถิยา อาจริยสฺส อาจริยธนํ ปริเยสิสฺสนฺติ, กิมงฺคํ\footnote{กิํ (สี, อิ)} ปนาหํ อายสฺมโต อานนฺทสฺส ปูชํ น \csromanfolio{16}กริสฺสามี”ติ.\csromanspacer{} อถ โข ทสโม คหปติ อฏฺฐกนาคโร ปาฏลิปุตฺตกญฺจ เวสาลิกญฺจ ภิกฺขุสํฆํ สนฺนิปาเตตฺวา ปณีเตน ขาทนีเยน โภชนีเยน สหตฺถา สนฺตปฺเปสิ สมฺปวาเรสิ, เอกเมกญฺจ ภิกฺขุํ ปจฺเจกํ ทุสฺสยุเคน อจฺฉาเทสิ, อายสฺมนฺตญฺจ อานนฺทํ ติจีวเรน อจฺฉาเทสิ, อายสฺมโต จ อานนฺทสฺส ปญฺจสตวิหารํ การาเปสีติ.}

\nitthitamruled{อฏฺฐกนาครสุตฺตํ นิฏฺฐิตํ ทุติยํ.}
\csromanmatikaanchor{mtk.4}
\csromantocmarkref{subsection}{3. เสขสุตฺต}{mtk.4}
\bookmark[dest={mtk.4},level=2]{3. เสขสุตฺต}
\csromanheader{2}{3. เสขสุตฺต}
\proseitem{22}{เอวํ เม สุตํ\csromandash{}เอกํ สมยํ ภควา สกฺเกสุ วิหรติ กปิลวตฺถุสฺมิํ นิโคฺรธาราเม.\csromanspacer{} เตน โข ปน สมเยน กาปิลวตฺถวานํ\footnote{กปิลวตฺถุวาสีนํ (ก)} สกฺยานํ นวํ สนฺถาคารํ อจิรการิตํ โหติ อนชฺฌาวุฏฺฐํ\footnote{อนชฺฌาวุตฺถํ (สี, สฺยา, กํ, อิ)} สมเณน วา พฺราหฺมเณน วา เกนจิ วา มนุสฺสภูเตน.\csromanspacer{} อถ โข กาปิลวตฺถวา สกฺยา เยน ภควา เตนุปสงฺกมิํสุ, อุปสงฺกมิตฺวา ภควนฺตํ อภิวาเทตฺวา เอกมนฺตํ นิสีทิํสุ.\csromanspacer{} เอกมนฺตํ นิสินฺนา โข กาปิลวตฺถวา สกฺยา ภควนฺตํ เอตทโวจุํ “อิธ ภนฺเต กาปิลวตฺถวานํ สกฺยานํ นวํ สนฺถาคารํ อจิรการิตํ\footnote{อจิรการิตํ โหติ (สฺยา, กํ, ก)} อนชฺฌาวุฏฺฐํ สมเณน วา พฺราหฺมเณน วา เกนจิ วา มนุสฺสภูเตน, ตํ ภนฺเต ภควา ปฐมํ ปริภุญฺชตุ, ภควตา ปฐมํ ปริภุตฺตํ ปจฺฉา กาปิลวตฺถวา สกฺยา ปริภุญฺชิสฺสนฺติ, ตทสฺส กาปิลวตฺถวานํ สกฺยานํ ทีฆรตฺตํ หิตาย สุขายา”ติ.\csromanspacer{} อธิวาเสสิ ภควา ตุณฺหีภาเวน อถ โข กาปิลวตฺถวา สกฺยา ภควโต อธิวาสนํ วิทิตฺวา อุฏฺฐายาสนา ภควนฺตํ อภิวาเทตฺวา ปทกฺขิณํ กตฺวา เยน นวํ สนฺถาคารํ เตนุปสงฺกมิํสุ, อุปสงฺกมิตฺวา สพฺพสนฺถริํ สนฺถาคารํ\footnote{สพฺพสนฺถริํ สนฺถตํ (ก)} สนฺถริตฺวา อาสนานิ ปญฺญเปตฺวา อุทกมณิกํ อุปฏฺฐเปตฺวา เตลปฺปทีปํ อาโรเปตฺวา เยน ภควา เตนุปสงฺกมิํสุ, อุปสงฺกมิตฺวา ภควนฺตํ อภิวาเทตฺวา เอกมนฺตํ อฏฺฐํสุ.\csromanspacer{} เอกมนฺตํ ฐิตา โข กาปิลวตฺถวา สกฺยา ภควนฺตํ เอตทโวจุํ “สพฺพสนฺถริํ สนฺถตํ ภนฺเต สนฺถาคารํ อาสนานิ ปญฺญตฺตานิ, อุทกมณิโก อุปฏฺฐาปิโต, เตลปฺปทีโป อาโรปิโต, ยสฺสทานิ}

\vspace{\tipitakaparskip}
\csromancenter{เสขสุตฺต (53) ภนฺเต ภควา กาลํ มญฺญตี”ติ.\csromanspacer{} อถ โข ภควา นิวาเสตฺวา ปตฺตจีวรมาทาย สทฺธิํ ภิกฺขุสํเฆน เยน สนฺถาคารํ เตนุปสงฺกมิ, อุปสงฺกมิตฺวา ปาเท ปกฺขาเลตฺวา สนฺถาคารํ ปวิสิตฺวา มชฺฌิมํ ถมฺภํ นิสฺสาย ปุรตฺถาภิมุโข นิสีทิ.\csromanspacer{} ภิกฺขุสํโฆปิ โข ปาเท ปกฺขาเลตฺวา สนฺถาคารํ ปวิสิตฺวา ปจฺฉิมํ ภิตฺติํ นิสฺสาย ปุรตฺถาภิมุโข นิสีทิ ภควนฺตํเยว ปุรกฺขตฺวา.\csromanspacer{} กาปิลวตฺถวาปิ โข สกฺยา ปาเท ปกฺขาเลตฺวา สนฺถาคารํ ปวิสิตฺวา ปุรตฺถิมํ ภิตฺติํ นิสฺสาย ปจฺฉิมาภิมุขา นิสีทิํสุ ภควนฺตํเยว ปุรกฺขตฺวา.\csromanspacer{} อถ โข ภควา กาปิลวตฺถเว สเกฺย พหุเทว รตฺติํ ธมฺมิยา กถาย สนฺทสฺเสตฺวา สมาทเปตฺวา สมุตฺเตเชตฺวา สมฺปหํเสตฺวา อายสฺมนฺตํ อานนฺทํ อามนฺเตสิ “ปฏิภาตุ ตํ อานนฺท กาปิลวตฺถวานํ สกฺยานํ เสโข ปาฏิปโท\footnote{ปฏิปโท (สฺยา, กํ, ก)}, ปิฏฺฐิ เม อาคิลายติ, ตมหํ อายมิสฺสามี”ติ. “เอวํ ภนฺเต”ติ โข อายสฺมา อานนฺโท ภควโต ปจฺจสฺโสสิ.\csromanspacer{} อต โข ภควา จตุคฺคุณํ สงฺฆาฏิํ ปญฺญาเปตฺวา ทกฺขิเณน ปสฺเสน สีหเสยฺยํ กปฺเปสิ ปาเท ปาทํ อจฺจาธาย สโต สมฺปชาโน อุฏฺฐานสญฺญํ มนสิ กริตฺวา.}
\proseitem{23}{\csromanfolio{17}อถ โข อายสฺมา อานนฺโท มหานามํ สกฺกํ อามนฺเตสิ\csromandash{}อิธ มหานาม อริยสาวโก สีลสมฺปนฺโน โหติ, อินฺทฺริเยสุ คุตฺตทฺวาโร โหติ, โภชเน มตฺตญฺญู โหติ, ชาคริยํ อนุยุตฺโต โหติ, สตฺตหิ สทฺธมฺเมหิ สมนฺนาคโต โหติ, จตุนฺนํ ฌานานํ อาภิเจตสิกานํ ทิฏฺฐธมฺมสุขวิหารานํ นิกามลาภี โหติ อกิจฺฉลาภี อกสิรลาภี.}

\proseitem{24}{กถญฺจ มหานาม อริยสาวโก สีลสมฺปนฺโน โหติ.\csromanspacer{} อิธ มหานาม อริยสาวโก สีลวา โหติ, ปาติโมกฺขสํวรสํวุโต วิหรติ อาจรโคจรสมฺปนฺโน, อณุมตฺเตสุ วชฺเชสุ ภยทสฺสาวี สมาทาย สิกฺขติ สิกฺขาปเทสุ.\csromanspacer{} เอวํ โข มหานาม อริยสาวโก สีลสมฺปนฺโน โหติ. (1)}

\prose{กถญฺจ มหานาม อริยสาวโก อินฺทฺริเยสุ คุตฺตทฺวาโร โหติ.\csromanspacer{} อิธ มหานาม อริยสาวโก จกฺขุนา รูปํ ทิสฺวา น นิมิตฺตคฺคาหี โหติ นานุพฺยญฺชนคฺคาหี, ยตฺวาธิกรณเมนํ จกฺขุนฺทฺริยํ อสํวุตํ วิหรนฺตํ \csromanfolio{18}อภิชฺฌาโทมนสฺสา ปาปกา อกุสลา ธมฺมา อนฺวาสฺสเวยฺยุํ, ตสฺส สํวราย ปฏิปชฺชติ, รกฺขติ จกฺขุนฺทฺริยํ, จกฺขุนฺทฺริเย สํวรํ อาปชฺชติ.\csromanspacer{} โสเตน สทฺทํ สุตฺวา ฯเปฯ.\csromanspacer{} ฆาเนน คนฺธํ ฆายิตฺวา ฯเปฯ.\csromanspacer{} ชิวฺหาย รสํ สายิตฺวา ฯเปฯ.\csromanspacer{} กาเยน โผฏฺฐพฺพํ ผุสิตฺวา ฯเปฯ.\csromanspacer{} มนสา ธมฺมํ วิญฺญาย น นิมิตฺตคฺคาหี โหติ นานุพฺยญฺชนคฺคาหี, ยตฺวาธิกรณเมนํ มนินฺทฺริยํ อสํวุตํ วิหรนฺตํ อภิชฺฌาโทมนสฺสา ปาปกา อกุสลา ธมฺมา อนฺวาสฺสเวยฺยุํ, ตสฺส สํวราย ปฏิปชฺชติ, รกฺขติ มนินฺทฺริยํ, มนินฺทฺริเย สํวรํ อาปชฺชติ.\csromanspacer{} เอวํ โข มหานาม อริยสาวโก อินฺทฺริเยสุ คุตฺตทฺวาโร โหติ. (2)}

\prose{กถญฺจ มหานาม อริยสาวโก โภชเน มตฺตญฺญู โหติ.\csromanspacer{} อิธ มหานาม อริยสาวโก ปฏิสงฺขา โยนิโส อาหารํ อาหาเรติ เนว ทวาย น มทาย น มณฺฑนาย น วิภูสนาย, ยาวเทว อิมสฺส กายสฺส ฐิติยา ยาปนาย วิหิํสูปรติยา พฺรหฺมจริยานุคฺคหาย, อิติ ปุราณญฺจ เวทนํ ปฏิหงฺขามิ, นวญฺจ เวทนํ น อุปฺปาเทสฺสามิ, ยาตฺรา จ เม ภวิสฺสติ อนวชฺชตา จ ผาสุวิหาโร จาติ.\csromanspacer{} เอวํ โข มหานาม อริยสาวโก โภชเน มตฺตญฺญู โหติ. (3)}

\prose{กถญฺจ มหานาม อริยสาวโก ชาคริยํ อนุยุตฺโต โหติ.\csromanspacer{} อิธ มหานาม อริยสาวโก ทิวสํ จงฺกเมน นิสชฺชาย อาวรณีเยหิ ธมฺเมหิ จิตฺตํ ปริโสเธติ, รตฺติยา ปฐมํ ยามํ จงฺกเมน นิสชฺชาย อาวรณีเยหิ ธมฺเมหิ จิตฺตํ ปริโสเธติ, รตฺติยา มชฺฌิมํ ยามํ ทกฺขิเณน ปสฺเสน สีหเสยฺยํ กปฺเปติ ปาเท ปาทํ อจฺจาธาย สโต สมฺปชาโน อุฏฺฐานสญฺญํ มนสิ กริตฺวา, รตฺถิยา ปจฺฉิมํ ยามํ ปจฺจุฏฺฐาย จงฺกเมน นิสชฺชาย อาวรณีเยหิ ธมฺเมหิ จิตฺตํ ปริโสเธติ.\csromanspacer{} เอวํ โข มหานาม อริยสาวโก ชาคริยํ อนุยุตฺโต โหติ. (4)}

\proseitem{25}{กถญฺจ มหานาม อริยสาวโก สตฺตหิ สทฺธมฺเมหิ สมนฺนาคโต โหติ.\csromanspacer{} อิธ มหานาม อริยสาวโก สทฺโธ โหติ, สทฺทหติ ตถาคตสฺส โพธิํ “อิติปิ โส ภควา อรหํ สมฺมาสมฺพุทฺโธ วิชฺชาจรณสมฺปนฺโน สุคโต โลกวิทู อนุตฺตโร ปุริสทมฺมสารถิ สตฺถา เทวมนุสฺสานํ พุทฺโธ ภควา”ติ.\csromanspacer{} หิริมา โหติ, หิรียติ กายทุจฺจริเตน วจีทุจฺจริเตน มโนทุจฺจริเตน, หิรียติ ปาปกานํ อกุสลานํ ธมฺมานํ สมาปตฺติยา.\csromanspacer{} โอตฺตปฺปี โหติ, โอตฺตปฺปติ}

\vspace{\tipitakaparskip}
\csromancenter{เสขสุตฺต (53) กายทุจฺจริเตน วจีทุจฺจริเตน มโนทุจฺจริเตน, โอตฺตปฺปติ ปาปกานํ อกุสลานํ ธมฺมานํ สมาปตฺติยา.\csromanspacer{} พหุสฺสุโต โหติ สุตธโร สุตสนฺนิจโย, เย เต ธมฺมา อาทิกลฺยาณา มชฺเฌกลฺยาณา ปริโยสานกลฺยาณา สาตฺถา สพฺยญฺชนา เกวลปริปุณฺณํ ปริสุทฺธํ พฺรหฺมจริยํ อภิวทนฺติ, ตถารูปาสฺส ธมฺมา พหุสฺสุตา\footnote{พหู สุตา (?)} โหนฺติ ธาตา\footnote{ธตา (สี, สฺยา, กํ, อิ)} วจสา ปริจิตา มนสานุเปกฺขิตา ทิฏฺฐิยา สุปฺปฏิวิทฺธา.\csromanspacer{} อารทฺธวีริโย วิหรติ อกุสลานํ ธมฺมานํ ปหานาย กุสลานํ ธมฺมานํ อุปสมฺปทาย, ถามวา ทฬฺหปรกฺกโม อนิกฺขิตฺตธุโร กุสเลสุ ธมฺเมสุ.\csromanspacer{} สติมา โหติ ปรเมน สติเนปกฺเกน สมนฺนาคโต, จิรกตํปิ จิรภาสิตํปิ สริตา อนุสฺสริตา.\csromanspacer{} ปญฺญวา โหติ อุทยตฺถคามินิยา ปญฺญาย สมนฺนาคโต อริยาย นิพฺเพธิกาย สมฺมา ทุกฺขกฺขยคามินิยา.\csromanspacer{} เอวํ โข มหานาม อริยสาวโก สตฺตหิ สทฺธมฺเมหิ สมนฺนาคโต โหติ. (5-11)}
\proseitem{26}{\csromanfolio{19}กถญฺจ มหานาม อริยสาวโก จตุนฺนํ ฌานานํ อาภิเจตสิกานํ ทิฏฺฐธมฺมสุขวิหารานํ นิกามลาภี โหติ อกิจฺฉลาภี อกสิรลาภี.\csromanspacer{} อิธ มหานาม อริยสาวโก วิวิจฺเจว กาเมหิ ฯเปฯ ปฐมํ ฌานํ อุปสมฺปชฺช วิหรติ.\csromanspacer{} วิตกฺกวิจารานํ วูปสมา อชฺฌตฺตํ สมฺปสาทนํ ฯเปฯ ทุติยํ ฌานํ อุปสมฺปชฺช วิหรติ.\csromanspacer{} ปีติยา จ วิราคา ฯเปฯ ตติยํ ฌานํ อุปสมฺปชฺช วิหรติ.\csromanspacer{} สุขสฺส จ ปหานา ทุกฺขสฺส จ ปหานา ปุพฺเพว โสมนสฺสโทมนสฺสานํ อตฺถงฺคมา ฯเปฯ จตุตฺถํ ฌานํ อุปสมฺปชฺช วิหรติ.\csromanspacer{} เอวํ โข มหานาม อริยสาวโก จตุนฺนํ ฌานานํ อาภิเจตสิกานํ ทิฏฺฐธมฺมสุขวิหารานํ นิกามลาภี โหติ อกิจฺฉลาภี อกสิรลาภี. (12-15)}

\proseitem{27}{ยโต โข มหานาม อริยสาวโก เอวํ สีลสมฺปนฺโน โหติ, เอวํ อินฺทฺริเยสุ คุตฺตทฺวาโร โหติ, เอวํ โภชเน มตฺตญฺญู โหติ, เอวํ ชาคริยํ อนุยุตฺโต โหติ, เอวํ สตฺตหิ สทฺธมฺเมหิ สมนฺนาคโต โหติ, เอวํ จตุนฺนํ ฌานานํ อาภิเจตสิกานํ ทิฏฺฐธมฺมสุขวิหารานํ นิกามลาภี โหติ อกิจฺฉลาภี อกสิรลาภี.\csromanspacer{} อยํ วุจฺจติ มหานาม อริยสาวโก เสโข ปาฏิปโท. \csromanfolio{20}อปุจฺจณฺฑตาย สมาปนฺโน ภพฺโพ อภินิพฺภิทาย, ภพฺโพ สมฺโพธาย, ภพฺโพ อนุตฺตรสฺส โยคกฺเขมสฺส อธิคมาย.\csromanspacer{} เสยฺยถาปิ มหานาม กุกฺกุฏิยา อณฺฑานิ อฏฺฐ วา ทส วา ทฺวาทส วา, ตานาสฺสุ กุกฺกุฏิยา สมฺมา อธิสยิตานิ สมฺมา ปริเสทิตานิ สมฺมา ปริภาวิตานิ.\csromanspacer{} กิญฺจาปิ ตสฺสา กุกฺกุฏิยา น เอวํ อิจฺฉา อุปฺปชฺเชยฺย “อโห วติเม กุกฺกุฏโปตกา ปาทนขสิขาย วา มุขตุณฺฑเกน วา อณฺฑโกสํ ปทาเลตฺวา โสตฺถินา อภินิพฺภิชฺเชยฺยุนฺ”ติ.\csromanspacer{} อถ โข ภพฺพาว เต กุกฺกุฏโปตกา ปาทนขสิขาย วา มุขตุณฺฑเกน วา อณฺฑโกสํ ปทาเลตฺวา โสตฺถินา อภินิพฺภิชฺชิตุํ, เอวเมว โข มหานาม ยโต อริยสาวโก เอวํ สีลสมฺปนฺโน โหติ, เอวํ อินฺทฺริเยสุ คุตฺตทฺวาโร โหติ, เอวํ โภชเน มตฺตญฺญู โหติ, เอวํ ชาคริยํ อนุยุตฺโต โหติ, เอวํ สตฺตหิ สทฺธมฺเมหิ สมนฺนาคโต โหติ, เอวํ จตุนฺนํ ฌานานํ อาภิเจตสิกานํ ทิฏฺฐธมฺมสุขวิหารานํ นิกามลาภี โหติ อกิจฺฉลาภี อกสิรลาภี.\csromanspacer{} อยํ วุจฺจติ มหานาม อริยสาวโก เสโข ปาฏิปโท.\csromanspacer{} อปุจฺจณฺฑตาย สมาปนฺโน ภพฺโพ อภินิพฺภิทาย, ภพฺโพ สมฺโพธาย, ภพฺโพ อนุตฺตรสฺส โยคกฺเขมสฺส อธิคมาย.}

\proseitem{28}{ส โข โส มหานาม อริยสาวโก อิมํเยว อนุตฺตรํ อุเปกฺขาสติปาริสุทฺธิํ อาคมฺม อเนกวิหิตํ ปุพฺเพนิวาสํ อนุสฺสรติ.\csromanspacer{} เสยฺยถิทํ, เอกมฺปิ ชาติํ เทฺวปิ ชาติโย ฯเปฯ อิติ สาการํ ส-อุทฺเทสํ อเนกวิหิตํ ปุพฺเพนิวาสํ อนุสฺสรติ.\csromanspacer{} อยมสฺส ปฐมาภินิพฺภิทา โหติ กุกฺกุฏจฺฉาปกสฺเสว อณฺฑโกสมฺหา. (1)}

\prose{ส โข โส มหานาม อริยสาวโก อิมํเยว อนุตฺตรํ อุเปกฺขาสติปาริสุทฺธิํ อาคมฺม ทิพฺเพน จกฺขุนา วิสุทฺเธน อติกฺกนฺตมานุสเกน สตฺเต ปสฺสติ จวมาเน อุปปชฺชมาเน หีเน ปณีเต สุวณฺเณ ทุพฺพณฺเณ สุคเต ทุคฺคเต ฯเปฯ ยถากมฺมูปเค สตฺเต ปชานาติ.\csromanspacer{} อยมสฺส ทุติยาภินิพฺภิทา โหติ กุกฺกุฏจฺฉาปกสฺเสว อณฺฑโกสมฺหา. (2)}

\prose{ส โข โส มหานาม อริยสาวโก อิมํเยว อนุตฺตรํ อุเปกฺขาสติปาริสุทฺธิํ อาคมฺม อาสวานํ ขยา อนาสวํ เจโตวิมุตฺติํ ปญฺญาวิมุตฺติํ ทิฏฺเฐว ธมฺเม สยํ อภิญฺญา สจฺฉิกตฺวา อุปสมฺปชฺช วิหรติ.\csromanspacer{} อยมสฺส ตติยาภินิพฺภิทา โหติ กุกฺกุฏจฺฉาปกสฺเสว อณฺฑโกสมฺหา. (3)}

\csromanheader{2}{3. เสขสุตฺต (53)}
\proseitem{29}{\csromanfolio{21}ยมฺปิ\footnote{ยมฺปิ โข (ก)} มหานาม อริยสาวโก สีลสมฺปนฺโน โหติ, อิทมฺปิสฺส โหติ จรณสฺมิํ.\csromanspacer{} ยมฺปิ มหานาม อริยสาวโก อินฺทฺริเยสุ คุตฺตทฺวาโร โหติ, อิทมฺปิสฺส โหติ จรณสฺมิํ.\csromanspacer{} ยมฺปิ มหานาม อริยสาวโก โภชเน มตฺตญฺญู โหติ, อิทมฺปิสฺส โหติ จรณสฺมิํ.\csromanspacer{} ยมฺปิ มหานาม อริยสาวโก ชาคริยํ อนุยุตฺโต โหติ, อิทมฺปิสฺส โหติ จรณสฺมิํ.\csromanspacer{} ยมฺปิ มหานาม อริยสาวโก สตฺตหิ สทฺธมฺเมหิ สมนฺนาคโต โหติ, อิทมฺปิสฺส โหติ จรณสฺมิํ.\csromanspacer{} ยมฺปิ มหานาม อริยสาวโก จตุนฺนํ ฌานานํ อาภิเจตสิกานํ ทิฏฺฐธมฺมสุขวิหารานํ นิกามลาภี โหติ อกิจฺฉลาภี อกสิรลาภี, อิทมฺปิสฺส โหติ จรณสฺมิํ.}

\prose{ยญฺจ โข มหานาม อริยสาวโก อเนกวิหิตํ ปุพฺเพนิวาสํ อนุสฺสรติ.\csromanspacer{} เสยฺยถิทํ, เอกมฺปิ ชาติํ เทฺวปิ ชาติโย ฯเปฯ อิติ สาการํ ส-อุทฺเทสํ อเนกวิหิตํ ปุพฺเพนิวาสํ อนุสฺสรติ, อิทมฺปิสฺส โหติ วิชฺชาย.\csromanspacer{} ยมฺปิ มหานาม อริยสาวโก ทิพฺเพน จกฺขุนา วิสุทฺเธน อติกฺกนฺตมานุสเกน สตฺเต ปสฺสติ จวมาเน อุปปชฺชมาเน หีเน ปณีเต สุวณฺเณ ทุพฺพณฺเณ สุคเต ทุคฺคเต ฯเปฯ ยถากมฺมูปเค สตฺเต ปชานาติ, อิทมฺปิสฺส โหติ วิชฺชาย.\csromanspacer{} ยมฺปิ มหานาม อริยสาวโก อาสวานํ ขยา อนาสวํ เจโตวิมุตฺติํ ปญฺญาวิมุตฺติํ ทิฏฺเฐว ธมฺเม สยํ อภิญฺญาสจฺฉิกตฺวา อุปสมฺปชฺช วิหรติ, อิทมฺปิสฺส โหติ วิชฺชาย.}

\prose{อยํ วุจฺจติ มหานาม อริยสาวโก วิชฺชาสมฺปนฺโน อิติปิ, จรณสมฺปนฺโน อิติปิ, วิชฺชาจรณสมฺปนฺโน อิติปิ.}

\csromanheader{2}{30. พฺรหฺมุนาเปสา มหานาม สนงฺกุมาเรน คาถา ภาสิตา\csromandash{}}
\csromangathameasure{ขตฺติโย เสฏฺโฐ ชเนตสฺมิํ, เย โคตฺตปฏิสาริโน. \\ วิชฺชาจรณสมฺปนฺโน, โส เสฏฺโฐ เทวมานุเสติ. \\ สา โข ปเนสา มหานาม พฺรหฺมุนา สนงฺกุมาเรน คาถา สุคีตา โน ทุคฺคีตา, สุภาสิตา โน ทุพฺภาสิตา, อตฺถสํหิตา โน อนตฺถสํหิตา, อนุมตา ภควตาติ. \\ อถ โข ภควา อุฏฺฐหิตฺวา อายสฺมนฺตํ อานนฺทํ อามนฺเตสิ “สาธุ สาธุ อานนฺท, สาธุ โข ตฺวํ อานนฺท กาปิลวตฺถวานํ สกฺยานํ เสขํ ปาฏิปทํ อภาสี”ติ. \\ อิทมโวจายสฺมา อานนฺโท, สมนุญฺโญ สตฺตา อโหสิ. อตฺตมนา กาปิลวตฺถวา สกฺยา อายสฺมโต อานนฺทสฺส ภาสิตํ อภินนฺทุนฺติ.}
\csromangathagroup{\csromangathastanza{ขตฺติโย เสฏฺโฐ ชเนตสฺมิํ, เย โคตฺตปฏิสาริโน. \\ วิชฺชาจรณสมฺปนฺโน, โส เสฏฺโฐ เทวมานุเสติ.}\csromangathastanza{สา โข ปเนสา มหานาม พฺรหฺมุนา สนงฺกุมาเรน คาถา สุคีตา โน ทุคฺคีตา, สุภาสิตา โน ทุพฺภาสิตา, อตฺถสํหิตา โน อนตฺถสํหิตา, อนุมตา ภควตาติ. \\ อถ โข ภควา อุฏฺฐหิตฺวา อายสฺมนฺตํ อานนฺทํ อามนฺเตสิ “สาธุ สาธุ อานนฺท, สาธุ โข ตฺวํ อานนฺท กาปิลวตฺถวานํ สกฺยานํ เสขํ ปาฏิปทํ อภาสี”ติ. \\ อิทมโวจายสฺมา อานนฺโท, สมนุญฺโญ สตฺตา อโหสิ.\csromanspacer{} อตฺตมนา กาปิลวตฺถวา สกฺยา อายสฺมโต อานนฺทสฺส ภาสิตํ อภินนฺทุนฺติ.}}

\nitthitamruled{เสขสุตฺตํ นิฏฺฐิตํ ตติยํ.}
\csromanmatikaanchor{mtk.5}
\csromantocmarkref{subsection}{4. โปตลิยสุตฺต}{mtk.5}
\bookmark[dest={mtk.5},level=2]{4. โปตลิยสุตฺต}
\csromanheader{2}{4. โปตลิยสุตฺต}
\proseitem{31}{\csromanfolio{22}เอวํ เม สุตํ\csromandash{}เอกํ สมยํ ภควา องฺคุตฺตราเปสุ วิหรติ, อาปณํ นาม องฺคุตฺตราปานํ นิคโม.\csromanspacer{} อถ โข ภควา ปุพฺพณฺหสมยํ นิวาเสตฺวา ปตฺตจีวรมาทาย อาปณํ ปิณฺฑาย ปาวิสิ, อาปเณ ปิณฺฑาย จริตฺวา ปจฺฉาภตฺตํ ปิณฺฑปาตปฏิกฺกนฺโต เยนญฺญตโร วนสณฺโฑ เตนุปสงฺกมิ ทิวาวิหาราย, ตํ วนสณฺฑํ อชฺโฌคาเหตฺวา\footnote{อชฺโฌคเหตฺวา (สี, สฺยา, กํ), อชฺโฌคาหิตฺวา (อิ, ก)} อญฺญตรสฺมิํ รุกฺขมูเล ทิวาวิหารํ นิสีทิ.\csromanspacer{} โปตลิโยปิ โข คหปติ สมฺปนฺนนิวาสนปาวุรโณ\footnote{ปาปุรโณ (สี, สฺยา, กํ) 3. ฉตฺตุปาหโน (ก)} ฉตฺตุปาหนาหิ3 ชงฺฆาวิหารํ อนุจงฺกมมาโน อนุวิจรมาโน เยน โส วนสณฺโฑ เตนุปสงฺกมิ, อุปสงฺกมิตฺวา ตํ วนสณฺฑํ อชฺโฌคาเหตฺวา เยน ภควา เตนุปสงฺกมิ, อุปสงฺกมิตฺวา ภควตา สทฺธิํ สมฺโมทิ, สมฺโมทนียํ กถํ สารณียํ วีติสาเรตฺวา เอกมนฺตํ อฏฺฐาสิ, เอกมนฺตํ ฐิตํ โข โปตลิยํ คหปติํ ภควา เอตทโวจ “สํวิชฺชนฺติ โข คหปติ อาสนานิ, สเจ อากงฺขสิ นิสีทา”ติ.\csromanspacer{} เอวํ วุตฺเต โปตลิโย คหปติ “คหปติวาเทน มํ สมโณ โคตโม สมุทาจรตี”ติ กุปิโต อนตฺตมโน ตุณฺหี อโหสิ.\csromanspacer{} ทุติยมฺปิ โข ภควา ฯเปฯ.\csromanspacer{} ตติยมฺปิ โข ภควา โปตลิยํ คหปติํ เอตทโวจ “สํวิชฺชนฺติ โข คหปติ อาสนานิ, สเจ อากงฺขสิ นิสีทา”ติ.\csromanspacer{} เอวํ วุตฺเต โปตลิโย คหปติ “คหปติวาเทน มํ สมโณ โคตโม สมุทาจรตี”ติ กุปิโต อนตฺตมโน ภควนฺตํ เอตทโวจ “ตยิทํ โภ โคตม นจฺฉนฺนํ ตยิทํ นปฺปติรูปํ, ยํ มํ ตฺวํ คหปติวาเทน สมุทาจรสี”ติ.\csromanspacer{} เต หิ เต คหปติ อาการา เต}

\vspace{\tipitakaparskip}
\csromancenter{โปตลิยสุตฺต (54) ลิงฺคา เต นิมิตฺตา ยถา ตํ คหปติสฺสาติ.\csromanspacer{} ตถา หิ ปน เม โภ โคตม สพฺเพ กมฺมนฺตา ปฏิกฺขิตฺตา, สพฺเพ โวหารา สมุจฺฉินฺนาติ.\csromanspacer{} ยถา กถํ ปน เต คหปติ สพฺเพ กมฺมนฺตา ปฏิกฺขิตฺตา, สพฺเพ โวหารา สมุจฺฉินฺนาติ.\csromanspacer{} อิธ เม โภ โคตม ยํ อโหสิ ธนํ วา ธญฺญํ วา รชตํ วา ชาตรูปํ วา, สพฺพํ ตํ ปุตฺตานํ ทายชฺชํ นิยฺยาตํ.\csromanspacer{} ตตฺถาหํ อโนวาที อนุปวาที ฆาสจฺฉาทนปรโม วิหรามิ.\csromanspacer{} เอวํ โข เม\footnote{เอวญฺจ เม (สฺยา), เอวํ เม (ก)} โภ โคตม สพฺเพ กมฺมนฺตา ปฏิกฺขิตฺตา, สพฺเพ โวหารา สมุจฺฉินฺนาติ.\csromanspacer{} อญฺญถา โข ตฺวํ คหปติ โวหารสมุจฺเฉทํ วทสิ, อญฺญถา จ ปน อริยสฺส วินเย โวหารสมุจฺเฉโท โหตีติ.\csromanspacer{} ยถา กถํ ปน ภนฺเต อริยสฺส วินเย โวหารสมุจฺเฉโท โหติ, สาธุ เม ภนฺเต ภควา ตถา ธมฺมํ เทเสตุ, ยถา อริยสฺส วินเย โวหารสมุจฺเฉโท โหตีติ.\csromanspacer{} เตน หิ คหปติ สุณาหิ สาธุกํ มนสิ กโรหิ ภาสิสฺสามีติ. “เอวํ ภนฺเต”ติ โข โปตลิโย คหปติ ภควโต ปจฺจสฺโสสิ.}
\proseitem{32}{\csromanfolio{23}ภควา เอตทโวจ\csromandash{}อฏฺฐ โข อิเม คหปติ ธมฺมา อริยสฺส วินเย โวหารสมุจฺเฉทาย สํวตฺตนฺติ.\csromanspacer{} กตเม อฏฺฐ.\csromanspacer{} อปาณาติปาตํ นิสฺสาย ปาณาติปาโต ปหาตพฺโพ, ทินฺนาทานํ นิสฺสาย อทินฺนาทานํ ปหาตพฺพํ, สจฺจวาจํ\footnote{สจฺจํ วาจํ (สฺยา)} นิสฺสาย มุสาวาโท ปหาตพฺโพ, อปิสุณํ วาจํ นิสฺสาย ปิสุณา วาจา ปหาตพฺพา, อคิทฺธิโลภํ นิสฺสาย คิทฺธิโลโภ ปหาตพฺโพ, อนินฺทาโรสํ นิสฺสาย นินฺทาโรโส ปหาตพฺโพ, อกฺโกธูปายาสํ นิสฺสาย โกธูปายาโส ปหาตพฺโพ, อนติมานํ นิสฺสาย อติมาโน ปหาตพฺโพ.\csromanspacer{} อิเม โข คหปติ อฏฺฐ ธมฺมา สํขิตฺเตน วุตฺตา, วิตฺถาเรน อวิภตฺตา, อริยสฺส วินเย โวหารสมุจฺเฉทาย สํวตฺตนฺตีติ.\csromanspacer{} เย เม\footnote{เย เม ปน (สฺยา, ก)} ภนฺเต ภควตา อฏฺฐ ธมฺมา สํขิตฺเตน วุตฺตา, วิตฺถาเรน อวิภตฺตา, อริยสฺส วินเย โวหารสมุจฺเฉทาย สํวตฺตนฺติ, สาธุ เม ภนฺเต ภควา อิเม อฏฺฐ ธมฺเม วิตฺถาเรน\footnote{วิตฺถาเรตฺวา (ก)} วิภชตุ อนุกมฺปํ อุปาทายาติ.\csromanspacer{} เตน หิ คหปติ สุณาหิ สาธุกํ มนสิ กโรหิ ภาสิสฺสามีติ. “เอวํ ภนฺเต”ติ โข โปตลิโย คหปติ ภควโต ปจฺจสฺโสสิ.\csromanspacer{} ภควา เอตทโวจ\csromandash{}}

\proseitem{33}{\csromanfolio{24}“อปาณาติปาตํ นิสฺสาย ปาณาติปาโต ปหาตพฺโพ”ติ อิติ โข ปเนตํ วุตฺตํ, กิญฺเจตํ ปฏิจฺจ วุตฺตํ.\csromanspacer{} อิธ คหปติ อริยสาวโก อิติ ปฏิสญฺจิกฺขติ\csromandash{}เยสํ โข อหํ สํโยชนานํ เหตุ ปาณาติปาตี อสฺสํ, เตสาหํ สํโยชนานํ ปหานาย สมุจฺเฉทาย ปฏิปนฺโน.\csromanspacer{} อหญฺเจว\footnote{อหญฺเจ (?)} โข ปน ปาณาติปาตี อสฺสํ, อตฺตาปิ มํ อุปวเทยฺย ปาณาติปาตปจฺจยา, อนุวิจฺจาปิ มํ วิญฺญู\footnote{อนุวิจฺจ วิญฺญู (สี, สฺยา, อิ)} ครเหยฺยุํ ปาณาติปาตปจฺจยา, กายสฺส เภทา ปรํ มรณา ทุคฺคติ ปาฏิกงฺขา ปาณาติปาตปจฺจยา.\csromanspacer{} เอตเทว โข ปน สํโยชนํ เอตํ นีวรณํ, ยทิทํ ปาณาติปาโต.\csromanspacer{} เย จ ปาณาติปาตปจฺจยา อุปฺปชฺเชยฺยุํ อาสวา วิฆาตปริฬาหา, ปาณาติปาตา ปฏิวิรตสฺส เอวํส เต อาสวา วิฆาตปริฬาหา น โหนฺติ. “อปาณาติปาตํ นิสฺสาย ปาณาติปาโต ปหาตพฺโพ”ติ อิติ ยนฺตํ วุตฺตํ, อิทเมตํ ปฏิจฺจ วุตฺตํ.}

\proseitem{34}{“ทินฺนาทานํ นิสฺสาย อทินฺนาทานํ ปหาตพฺพนฺ”ติ อิติ โข ปเนตํ วุตฺตํ, กิญฺเจตํ ปฏิจฺจ วุตฺตํ.\csromanspacer{} อิธ คหปติ อริยสาวโก อิติ ปฏิสญฺจิกฺขติ\csromandash{}เยสํ โข อหํ สํโยชนานํ เหตุ อทินฺนาทายี อสฺสํ, เตสาหํ สํโยชนานํ ปหานาย สมุจฺเฉทาย ปฏิปนฺโน.\csromanspacer{} อหญฺเจว โข ปน อทินฺนาทายี อสฺสํ, อตฺตาปิ มํ อุปวเทยฺย อทินฺนาทานปจฺจยา, อนุวิจฺจาปิ มํ วิญฺญู ครเหยฺยุํ อทินฺนาทานปจฺจยา, กายสฺส เภทา ปรํ มรณา ทุคฺคติ ปาฏิกงฺขา อทินฺนาทานปจฺจยา.\csromanspacer{} เอตเทว โข ปน สํโยชนํ เอตํ นีวรณํ, ยทิทํ อทินฺนาทานํ.\csromanspacer{} เย จ อทินฺนาทานปจฺจยา อุปฺปชฺเชยฺยุํ อาสวา วิฆาตปริฬาหา, อทินฺนาทานา ปฏิวิรตสฺส เอวํส เต อาสวา วิฆาตปริฬาหา น โหนฺติ. “ทินฺนาทานํ นิสฺสาย อทินฺนาทานํ ปหาตพฺพนฺ”ติ อิติ ยนฺตํ วุตฺตํ, อิทเมตํ ปฏิจฺจ วุตฺตํ.}

\proseitem{35}{“สจฺจวาจํ นิสฺสาย มุสาวาโท ปหาตพฺโพ”ติ อิติ โข ปเนตํ วุตฺตํ, กิญฺเจตํ ปฏิจฺจ วุตฺตํ.\csromanspacer{} อิธ คหปติ อริยสาวโก อิติ ปฏิสญฺจิกฺขติ\csromandash{}เยสํ โข อหํ สํ โยชนานํ เหตุ มุสาวาที อสฺสํ, เตสาหํ สํโยชนานํ ปหานาย สมุจฺเฉทาย ปฏิปนฺโน.\csromanspacer{} อหญฺเจว โข ปน มุสาวาที อสฺสํ, อตฺตาปิ มํ อุปวเทยฺย มุสาวาทปจฺจยา, อนุวิจฺจาปิ มํ วิญฺญู ครเหยฺยุํ มุสาวาทปจฺจยา, กายสฺส เภทา ปรํ มรณา ทุคฺคติ ปาฏิกงฺขา มุสาวาทปจฺจยา.\csromanspacer{} เอตเทว โข ปน สํโยชนํ เอตํ นีวรณํ,}

\vspace{\tipitakaparskip}
\csromancenter{โปตลิยสุตฺต (54) ยทิทํ มุสาวาโท.\csromanspacer{} เย จ มุสาวาทปจฺจยา อุปฺปชฺเชยฺยุํ อาสวา วิฆาตปริฬาหา, มุสาวาทา ปฏิวิรตสฺส เอวํส เต อาสวา วิฆาตปริฬาหา น โหนฺติ. “สจฺจวาจํ นิสฺสาย มุสาวาโท ปหาตพฺโพ”ติ อิติ ยนฺตํ วุตฺตํ, อิทเมตํ ปฏิจฺจ วุตฺตํ.}
\proseitem{36}{\csromanfolio{25}“อปิสุณํ วาจํ นิสฺสาย ปิสุณา วาจา ปหาตพฺพา”ติ อิติ โข ปเนตํ วุตฺตํ, กิญฺเจตํ ปฏิจฺจ วุตฺตํ.\csromanspacer{} อิธ คหปติ อริยสาวโก อิติ ปฏิสญฺจิกฺขติ\csromandash{}เยสํ โข อหํ สํโยชนานํ เหตุ ปิสุณวาโจ อสฺสํ, เตสาหํ สํโยชนานํ ปหานาย สมุจฺเฉทาย ปฏิปนฺโน.\csromanspacer{} อหญฺเจว โข ปน ปิสุณวาโจ อสฺสํ, อตฺตาปิ มํ อุปวเทยฺย ปิสุณาวาจาปจฺจยา, อนุวิจฺจาปิ มํ วิญฺญู ครเหยฺยุํ ปิสุณวาจาปจฺจยา, กายสฺส เภทา ปรํ มรณา ทุคฺคติ ปาฏิกงฺขา ปิสุณวาจาปจฺจยา.\csromanspacer{} เอตเทว โข ปน สํโยชนํ เอตํ นีวรณํ, ยทิทํ ปิสุณา วาจา.\csromanspacer{} เย จ ปิสุณวาจาปจฺจยา อุปฺปชฺเชยฺยุํ อาสวา วิฆาตปริฬาหา, ปิสุณาย วาจาย ปฏิวิรตสฺส เอวํส เต อาสวา วิฆาตปริฬาหา น โหนฺติ. “อปิสุณํ วาจํ นิสฺสาย ปิสุณา วาจา ปหาตพฺพา”ติ อิติ ยนฺตํ วุตฺตํ, อิทเมตํ ปฏิจฺจ วุตฺตํ.}

\proseitem{37}{“อคิทฺธิโลภํ นิสฺสาย คิทฺธิโลโภ ปหาตพฺโพ”ติ อิติ โข ปเนตํ วุตฺตํ, กิญฺเจตํ ปฏิจฺจ วุตฺตํ.\csromanspacer{} อิธ คหปติ อริยสาวโก อิติ ปฏิสญฺจิกฺขติ\csromandash{}เยสํ โข อหํ สํโยชนานํ เหตุ คิทฺธิโลภี อสฺสํ, เตสาหํ สํโยชนานํ ปหานาย สมุจฺเฉทาย ปฏิปนฺโน.\csromanspacer{} อหญฺเจว โข ปน คิทฺธิโลภี อสฺสํ, อตฺตาปิ มํ อุปวเทยฺย คิทฺธิโลภปจฺจยา, อนุวิจฺจาปิ มํ วิญฺญู ครเหยฺยุํ คิทฺธิโลภปจฺจยา, กายสฺส เภทา ปรํ มรณา ทุคฺคติ ปาฏิกงฺขา คิทฺธิโลภปจฺจยา.\csromanspacer{} เอตเทว โข ปน สํโยชนํ เอตํ นีวรณํ, ยทิทํ คิทฺธิโลโภ.\csromanspacer{} เย จ คิทฺธิโลภปจฺจยา อุปฺปชฺเชยฺยุํ อาสวา วิฆาตปริฬาหา, คิทฺธิโลภา ปฏิวิรตสฺส เอวํส เต อาสวา วิฆาตปริฬาหา น โหนฺติ. “อคิทฺธิโลภํ นิสฺสาย คิทฺธิโลโภ ปหาตพฺโพ”ติ อิติ ยนฺตํ วุตฺตํ, อิทเมตํ ปฏิจฺจ วุตฺตํ.}

\proseitem{38}{“อนินฺทาโรสํ นิสฺสาย นินฺทาโรโส ปหาตพฺโพ”ติ อิติ โข ปเนตํ วุตฺตํ, กิญฺเจตํ ปฏิจฺจ วุตฺตํ.\csromanspacer{} อิธ คหปติ อริยสาวโก อิติ ปฏิสญฺจิกฺขติ\csromandash{}เยสํ โข อหํ สํโยชนานํ เหตุ นินฺทาโรสี อสฺสํ, \csromanfolio{26}เตสาหํ สํโยชนานํ ปหานาย สมุจฺเฉทาย ปฏิปนฺโน.\csromanspacer{} อหญฺเจว โข ปน นินฺทาโรสี อสฺสํ, อตฺตาปิ มํ อุปวเทยฺย นินฺทาโรสปจฺจยา, อนุวิจฺจาปิ มํ วิญฺญู ครเหยฺยุํ นินฺทาโรสปจฺจยา, กายสฺส เภทา ปรํ มรณา ทุคฺคติ ปาฏิกงฺขา นินฺทาโรสปจฺจยา.\csromanspacer{} เอตเทว โข ปน สํโยชนํ เอตํ นีวรณํ, ยทิทํ นินฺทาโรโส.\csromanspacer{} เย จ นินฺทาโรสปจฺจยา อุปฺปชฺเชยฺยุํ อาสวา วิฆาตปริฬาหา, อนินฺทาโรสิสฺส เอวํส เต อาสวา วิฆาตปริฬาหา น โหนฺติ. “อนินฺทาโรสํ นิสฺสาย นินฺทาโรโส ปหาตพฺโพ”ติ อิติ ยนฺตํ วุตฺตํ, อิทเมตํ ปฏิจฺจ วุตฺตํ.}

\proseitem{39}{“อกฺโกธูปายาสํ นิสฺสาย โกธูปายาโส ปหาตพฺโพ”ติ อิติ โข ปเนตํ วุตฺตํ, กิญฺเจตํ ปฏิจฺจ วุตฺตํ.\csromanspacer{} อิธ คหปติ อริยสาวโก อิติ ปฏิสญฺจิกฺขติ\csromandash{}เยสํ โข อหํ สํโยชนานํ เหตุ โกธูปายาสี อสฺสํ, เตสาหํ สํโยชนานํ ปหานาย สมุจฺเฉทาย ปฏิปนฺโน.\csromanspacer{} อหญฺเจว โข ปน โกธูปายาสี อสฺสํ, อตฺตาปิ มํ อุปวเทยฺย โกธูปายาสปจฺจยา, อนุวิจฺจาปิ มํ วิญฺญู ครเหยฺยุํ โกธูปายาสปจฺจยา, กายสฺส เภทา ปรํ มรณา ทุคฺคติ ปาฏิกงฺขา โกธูปายาสปจฺจยา.\csromanspacer{} เอตเทว โข ปน สํโยชนํ เอตํ นีวรณํ, ยทิทํ โกธูปายาโส.\csromanspacer{} เย จ โกธูปายาสปจฺจยา อุปฺปชฺเชยฺยุํ อาสวา วิฆาตปริฬาหา, อกฺโกธูปายาสิสฺส เอวํส เต อาสวา วิฆาตปริฬาหา น โหนฺติ. “อกฺโกธูปายาสํ นิสฺสาย โกธูปายาโส ปหาตพฺโพ”ติ อิติ ยนฺตํ วุตฺตํ, อิทเมตํ ปฏิจฺจ วุตฺตํ.}

\proseitem{40}{“อนติมานํ นิสฺสาย อติมาโน ปหาตพฺโพ”ติ อิติ โข ปเนตํ วุตฺตํ, กิญฺเจตํ ปฏิจฺจ วุตฺตํ.\csromanspacer{} อิธ คหปติ อริยสาวโก อิติ ปฏิสญฺจิกฺขติ\csromandash{} เยสํ โข อหํ สํโยชนานํ เหตุ อติมานี อสฺสํ, เตสาหํ สํโยชนานํ ปหานาย สมุจฺเฉทาย ปฏิปนฺโน.\csromanspacer{} อหญฺเจว โข ปน อติมานี อสฺสํ, อตฺตาปิ มํ อุปวเทยฺย อติมานปจฺจยา, อนุวิจฺจาปิ มํ วิญฺญู ครเหยฺยุํ อติมานปจฺจยา, กายสฺส เภทา ปรํ มรณา ทุคฺคติ ปาฏิกงฺขา อติมานปจฺจยา.\csromanspacer{} เอตเทว โข ปน สํโยชนํ เอตํ นีวรณํ, ยทิทํ อติมาโน.\csromanspacer{} เย จ อติมานปจฺจยา อุปฺปชฺเชยฺยุํ อาสวา วิฆาตปริฬาหา, อนติมานิสฺส เอวํส เต อาสวา วิฆาตปริฬาหา น โหนฺติ. “อนติมานํ นิสฺสาย อติมาโน ปหาตพฺโพ”ติ อิติ ยนฺตํ วุตฺตํ, อิทเมตํ ปฏิจฺจ วุตฺตํ.}

\csromanheader{2}{4. โปตลิยสุตฺต (54)}
\proseitem{41}{\csromanfolio{27}อิเม โข คหปติ อฏฺฐ ธมฺมา สํขิตฺเตน วุตฺตา, วิตฺถาเรน วิภตฺตา\footnote{อวิภตฺตา (สฺยา, ก)}, เย อริยสฺส วินเย โวหารสมุจฺเฉทาย สํวตฺตนฺติ.\csromanspacer{} น เตฺวว ตาว อริยสฺส วินเย สพฺเพน สพฺพํ สพฺพถา สพฺพํ โวหารสมุจฺเฉโท โหตีติ.}

\prose{ยถา กถํ ปน ภนฺเต อริยสฺส วินเย สพฺเพน สพฺพํ สพฺพถา สพฺพํ โวหารสมุจฺเฉโท โหติ.\csromanspacer{} สาธุ เม ภนฺเต ภควา ตถา ธมฺมํ เทเสตุ, ยถา อริยสฺส วินเย สพฺเพน สพฺพํ สพฺพถา สพฺพํ โวหารสมุจฺเฉโท โหตีติ.\csromanspacer{} เตน หิ คหปติ สุณาหิ สาธุกํ มนสิ กโรหิ ภาสิสฺสามีติ. “เอวํ ภนฺเต”ติ โข โปตลิโย คหปติ ภควโต ปจฺจสฺโสสิ.\csromanspacer{} ภควา เอตทโวจ\csromandash{}}

\csromanheader{2}{\textbf{กามาทีนวกถา}}
\proseitem{42}{เสยฺยถาปิ คหปติ กุกฺกุโร ชิฆจฺฉาทุพฺพลฺยปเรโต โคฆาตกสูนํ ปจฺจุปฏฺฐิโต อสฺส, ตเมนํ ทกฺโข โคฆาตโก วา โคฆาตกนฺเตวาสี วา อฏฺฐิกงฺกลํ สุนิกฺกนฺตํ นิกฺกนฺตํ นิมฺมํสํ โลหิตมกฺขิตํ อุปสุมฺเภยฺย\footnote{อุปจฺฉุเภยฺย (สี, อิ), อุปจฺฉูเภยฺย (สฺยา, กํ), อุปจฺจุมฺเภยฺย (ก)}.\csromanspacer{} ตํ กิํ มญฺญสิ คหปติ, อปิ นุ โข โส กุกฺกุโร อมุํ อฏฺฐิกงฺกลํ สุนิกฺกนฺตํ นิกฺกนฺตํ นิมฺมํสํ โลหิตมกฺขิตํ ปเลหนฺโต ชิฆจฺฉาทุพฺพลฺยํ ปฏิวิเนยฺยาติ.\csromanspacer{} โน เหตํ ภนฺเต.\csromanspacer{} ตํ กิสฺส เหตุ, อทุํ หิ ภนฺเต อฏฺฐิกงฺกลํ สุนิกฺกนฺตํ นิกฺกนฺตํ นิมฺมํสํ โลหิตมกฺขิตํ, ยาวเทว ปน โส กุกฺกุโร กิลมถสฺส วิฆาตสฺส ภาคี อสฺสาติ.\csromanspacer{} เอวเมว โข คหปติ อริยสาวโก อิติ ปฏิสญฺจิกฺขติ “อฏฺฐิกงฺกลูปมา กามา วุตฺตา ภควตา พหุทุกฺขา พหุปายาสา\footnote{พหูปายาสา (สี, สฺยา, กํ, อิ)}, อาทีนโว เอตฺถ ภิยฺโย”ติ.\csromanspacer{} เอวเมตํ ยถาภูตํ สมฺมปฺปญฺญาย ทิสฺวา ยายํ อุเปกฺขา นานตฺตา นานตฺตสิตา, ตํ อภินิวชฺเชตฺวา ยายํ อุเปกฺขา เอกตฺตา เอกตฺตสิตา, ยตฺถ สพฺพโส โลกามิสูปาทานา อปริเสสา นิรุชฺฌนฺติ, ตเมวูเปกฺขํ ภาเวติ.}

\proseitem{43}{เสยฺยถาปิ คหปติ คิชฺโฌ วา กงฺโก วา กุลโล วา มํสเปสิํ อาทาย อุฑฺฑีเยยฺย\footnote{อุฑฺฑเยยฺย (สฺยา, อิ)}, ตเมนํ คิชฺฌาปิ กงฺกาปิ กุลลาปิ \csromanfolio{28}อนุปติตฺวา อนุปติตฺวา วิตจฺเฉยฺยุํ วิสฺสชฺเชยฺยุํ\footnote{วิราเชยฺยุํ (สี, สฺยา, กํ, อิ)}.\csromanspacer{} ตํ กิํ มญฺญสิ คหปติ, สเจ โส คิชฺโฌ วา กงฺโก วา กุลโล วา ตํ มํสเปสิํ น ขิปฺปเมว ปฏินิสฺสชฺเชยฺย.\csromanspacer{} โส ตโตนิทานํ มรณํ วา นิคจฺเฉยฺย, มรณมตฺตํ วา ทุกฺขนฺติ.\csromanspacer{} เอวํ ภนฺเต.\csromanspacer{} เอวเมว โข คหปติ อริยสาวโก อิติ ปฏิสญฺจิกฺขติ “มํสเปสูปมา กามา วุตฺตา ภควตา พหุทุกฺขา พหุปายาสา, อาทีนโว เอตฺถ ภิยฺโย”ติ.\csromanspacer{} เอวเมตํ ยถาภูตํ สมฺมปฺปญฺญาย ทิสฺวา ยายํ อุเปกฺขา นานตฺตา นานตฺตสิตา, ตํ อภินิวชฺเชตฺวา ยายํ อุเปกฺขา เอกตฺตา เอกตฺตสิตา, ยตฺถ สพฺพโส โลกามิสูปาทานา อปริเสสา นิรุชฺฌนฺติ, ตเมวูเปกฺขํ ภาเวติ.}

\proseitem{44}{เสยฺยถาปิ คหปติ ปุริโส อาทิตฺตํ ติณุกฺกํ อาทาย ปฏิวาตํ คจฺฉยฺย.\csromanspacer{} ตํ กิํ มญฺญสิ คหปติ, สเจ โส ปุริโส ตํ อาทิตฺตํ ติณุกฺกํ น ขิปฺปเมว ปฏินิสฺสชฺเชยฺย.\csromanspacer{} ตสฺส สา อาทิตฺตา ติณุกฺกา หตฺถํ วา ทเหยฺย, พาหุํ วา ทเหยฺย, อญฺญตรํ วา อญฺญตรํ วา องฺคปจฺจงฺคํ\footnote{ทเหยฺย. อญฺญตรํ วา องฺคปจฺจงฺคํ (สี, อิ)} ทเหยฺย.\csromanspacer{} โส ตโตนิทานํ มรณํ วา นิคจฺเฉยฺย, มรณมตฺตํ วา ทุกฺขนฺติ เอวํ ภนฺเต.\csromanspacer{} เอวเมว โข คหปติ อริยสาวโก อิติ ปฏิสญฺจิกฺขติ “ติณุกฺกูปมา กามา วุตฺตา ภควตา พหุทุกฺขา พหุปายาสา, อาทีนโว เอตฺถ ภิยฺโย”ติ.\csromanspacer{} เอวเมตํ ยถาภูตํ สมฺมปฺปญฺญาย ทิสฺวา ฯเปฯ ตเมวูเปกฺขํ ภาเวติ.}

\proseitem{45}{เสยฺยถาปิ คหปติ องฺคารกาสุ สาธิกโปริสา ปูรา องฺคารานํ วีตจฺฉิกานํ วีตธูมานํ.\csromanspacer{} อถ ปุริโส อาคจฺเฉยฺย ชีวิตุกาโม อมริตุกาโม สุขกาโม ทุกฺขปฺปฏิกฺกูโล, ตเมนํ เทฺว พลวนฺโต ปุริสา นานาพาหาสุ คเหตฺวา องฺคารกาสุํ อุปกฑฺเฒยฺยุํ.\csromanspacer{} ตํ กิํ มญฺญสิ คหปติ, อปิ นุ โส ปุริโส อิติจิติเจว กายํ สนฺนาเมยฺยาติ.\csromanspacer{} เอวํ ภนฺเต.\csromanspacer{} ตํ กิสฺส เหตุ, วิทิตญฺหิ ภนฺเต ตสฺส ปุริสสฺส “อิมํ จาหํ องฺคารกาสุํ ปปติสฺสามิ, ตโตนิทานํ มรณํ วา นิคจฺฉิสฺสามิ, มรณมตฺตํ วา ทุกฺขนฺ”ติ.\csromanspacer{} เอวเมว โข คหปติ อริยสาวโก อิติ ปฏิสญฺจิกฺขติ “องฺคารกาสูปมา กามา วุตฺตา ภควตา พหุทุกฺขา พหุปายาสา, อาทีนโว เอตฺถ ภิยฺโย”ติ.\csromanspacer{} เอวเมตํ ยถาภูตํ สมฺมปฺปญฺญาย ทิสฺวา ฯเปฯ ตเมวูเปกฺขํ ภาเวติ.}

\csromanheader{2}{4. โปตลิยสุตฺต (54)}
\proseitem{46}{\csromanfolio{29}เสยฺยถาปิ คหปติ ปุริโส สุปินกํ ปสฺเสยฺย อารามรามเณยฺยกํ วนรามเณยฺยกํ ภูมิรามเณยฺยกํ โปกฺขรณิรามเณยฺยกํ, โส ปฏิพุทฺโธ น กิญฺจิ ปฏิปสฺเสยฺย\footnote{ปสฺเสยฺย (สี, สฺยา, กํ, อิ)}.\csromanspacer{} เอวเมว โข คหปติ อริยสาวโก อิติ ปฏิสญฺจิกฺขติ “สุปินกูปมา กามา วุตฺตา ภควตา พหุทุกฺขา พหุปายาสา, อาทีนโว เอตฺถ ภิยฺโย”ติ ฯเปฯ ตเมวูเปกฺขํ ภาเวติ.}

\proseitem{47}{เสยฺยถาปิ คหปติ ปุริโส ยาจิตกํ โภคํ ยาจิตฺวา ยานํ วา\footnote{ยานํ (สฺยา, กํ, อิ)} โปริเสยฺยํ\footnote{โปโรเสยฺยํ (สี, อิ, ก), โอโรเปยฺย (สฺยา,กํ)} ปวรมณิกุณฺฑลํ, โส เตหิ ยาจิตเกหิ โภเคหิ ปุรกฺขโต ปริวุโต อนฺตราปณํ ปฏิปชฺเชยฺย.\csromanspacer{} ตเมนํ ชโน ทิสฺวา เอวํ วเทยฺย “โภคี วต โภ ปุริโส, เอวํ กิร โภคิโน โภคานิ ภุญฺชนฺตี”ติ.\csromanspacer{} ตเมนํ สามิกา ยตฺต ยตฺเถว ปสฺเสยฺยุํ, ตตฺถ ตตฺเถว สานิ หเรยฺยุํ.\csromanspacer{} ตํ กิํ มญฺญสิ คหปติ, อลํ นุ โข ตสฺส ปุริสสฺส อญฺญถตฺตายาติ.\csromanspacer{} เอวํ ภนฺเต.\csromanspacer{} ตํ กิสฺส เหตุ, สามิโน หิ ภนฺเต สานิ หรนฺตีติ.\csromanspacer{} เอวเมว โข คหปติ อริยสาวโก อิติ ปฏิสญฺจิกฺขติ “ยาจิตกูปมา กามา วุตฺตา ภควตา พหุทุกฺขา พหุปายาสา, อาทีนโว เอตฺถ ภิยฺโย”ติ ฯเปฯ ตเมวูเปกฺขํ ภาเวติ.}

\proseitem{48}{เสยฺยถาปิ คหปติ คามสฺส วา นิคมสฺส วา อวิทูเร ติพฺโพ วนสณฺโฑ, ตตฺรสฺส รุกฺโข สมฺปนฺนผโล จ อุปปนฺนผโล\footnote{อุปฺปนฺนผโล (สฺยา)} จ, น จสฺสุ กานิจิ ผลานิ ภูมิยํ ปติตานิ.\csromanspacer{} อถ ปุริโส อาคจฺเฉยฺย ผลตฺถิโก ผลคเวสี ผลปริเยสนํ จรมาโน.\csromanspacer{} โส ตํ วนสณฺฑํ อชฺโฌคาเหตฺวา ตํ รุกฺขํ ปสฺเสยฺย สมฺปนฺนผลญฺจ อุปปนฺนผลญฺจ.\csromanspacer{} ตสฺส เอวมสฺส “อยํ โข รุกฺโข สมฺปนฺนผโล จ อุปปนฺนผโล จ, นตฺถิ จ กานิจิ ผลานิ ภูมิยํ ปติตานิ, ชานามิ โข ปนาหํ รุกฺขํ อาโรหิตุํ\footnote{อารุหิตุํ (สี)}.\csromanspacer{} ยํนูนาหํ อิมํ รุกฺขํ อาโรหิตฺวา ยาวทตฺถญฺจ ขาเทยฺยํ, อุจฺฉงฺคญฺจ ปูเรยฺยนฺ”ติ.\csromanspacer{} โส ตํ รุกฺขํ อาโรหิตฺวา ยาวทตฺถญฺจ ขาเทยฺย, อุจฺฉงฺคญฺจ ปูเรยฺย.\csromanspacer{} อถ ทุติโย ปุริโส อาคจฺเฉยฺย ผลตฺถิโก ผลคเวสี ผลปริเยสนํ จรมาโน ติณฺหํ กุฐาริํ\footnote{กุธาริํ (สฺยา, กํ, ก)} อาทาย, โส ตํ วนสณฺฑํ \csromanfolio{30}อชฺโฌคาเหตฺวา ตํ รุกฺขํ ปสฺเสยฺย สมฺปนฺนผลญฺจ อุปปนฺนผลญฺจ.\csromanspacer{} ตสฺส เอวมสฺส “อยํ โข รุกฺโข สมฺปนฺนผโล จ อุปปนฺนผโล จ, นตฺถิ จ กานิจิ ผลานิ ภูมิยํ ปติตานิ, น โข ปนาหํ ชานามิ รุกฺขํ อาโรหิตุํ.\csromanspacer{} ยํนูนาหํ อิมํ รุกฺขํ มูลโต เฉตฺวา ยาวทตฺถญฺจ ขาเทยฺยํ, อุจฺฉงฺคญฺจ ปูเรยฺยนฺ”ติ.\csromanspacer{} โส ตํ รุกฺขํ มูลโตว ฉินฺเทยฺย.\csromanspacer{} ตํ กิํ มญฺญสิ คหปติ, อมุโก\footnote{อสุ (สี, อิ)} โย โส ปุริโส ปฐมํ รุกฺขํ อารูโฬ, สเจ โส น ขิปฺปเมว โอโรเหยฺย.\csromanspacer{} ตสฺส โส รุกฺโข ปปตนฺโต หตฺถํ วา ภญฺเชยฺย, ปาทํ วา ภญฺเชยฺย, อญฺญตรํ วา อญฺญตรํ วา องฺคปจฺจงฺคํ ภญฺเชยฺย.\csromanspacer{} โส ตโตนิทานํ มรณํ วา นิคจฺเฉยฺย, มรณมตฺตํ วา ทุกฺขนฺติ.\csromanspacer{} เอวํ ภนฺเต.\csromanspacer{} เอวเมว โข คหปติ อริยสาวโก อิติ ปฏิสญฺจิกฺขติ “รุกฺขผลูปมา กามา วุตฺตา ภควตา พหุทุกฺขา พหุปายาสา, อาทีนโว เอตฺถ ภิยฺโย”ติ.\csromanspacer{} เอวเมตํ ยถาภูตํ สมฺมปฺปญฺญาย ทิสฺวา ยายํ อุเปกฺขา นานตฺตา นานตฺตสิตา, ตํ อภินิวชฺเชตฺวา ยายํ อุเปกฺขา เอกตฺตา เอกตฺตสิตา, ยตฺถ สพฺพโส โลกามิสูปาทานา อปริเสสา นิรุชฺฌนฺติ, ตเมวูเปกฺขํ ภาเวติ.}

\proseitem{49}{ส โข โส คหปติ อริยสาวโก อิมํเยว อนุตฺตรํ อุเปกฺขาสติปาริสุทฺธิํ อาคมฺม อเนกวิหิตํ ปุพฺเพนิวาสํ อนุสฺสรติ.\csromanspacer{} เสยฺยถิทํ เอกมฺปิ ชาติํ เทฺวปิ ชาติโย ฯเปฯ อิติ สาการํ ส-อุทฺเทสํ อเนกวิหิตํ ปุพฺเพนิวาสํ อนุสฺสรติ.}

\prose{ส โข โส คหปติ อริยสาวโก อิมํเยว อนุตฺตรํ อุเปกฺขาสติปาริสุทฺธิํ อาคมฺม ทิพฺเพน จกฺขุนา วิสุทฺเธน อติกฺกนฺตมานุสเกน สตฺเต ปสฺสติ จวมาเน อุปปชฺชมาเน หีเน ปณีเต สุวณฺเณ ทุพฺพณฺเณ สุคเต ทุคฺคเต ฯเปฯ ยถากมฺมูปเค สตฺเต ปชานาติ.}

\prose{ส โข โส คหปติ อริยสาวโก อิมํเยว อนุตฺตรํ อุเปกฺขาสติปาริสุทฺธิํ อาคมฺม อาสวานํ ขยา อนาสวํ เจโตวิมุตฺติํ ปญฺญาวิมุตฺติํ ทิฏฺเฐว ธมฺเม สยํ อภิญฺญา สจฺฉิกตฺวา อุปสมฺปชฺช วิหรติ.\csromanspacer{} เอตฺตาวตา โข คหปติ อริยสฺส วินเย สพฺเพน สพฺพํ สพฺพถา สพฺพํ โวหารสมุจฺเฉโท โหติ.}

\csromanheader{2}{5. ชีวกสุตฺต (55)}
\proseitem{50}{\csromanfolio{31}ตํ กิํ มญฺญสิ คหปติ, ยถา อริยสฺส วินเย สพฺเพน สพฺพํ สพฺพถา สพฺพํ โวหารสมุจฺเฉโท โหติ, อปิ นุ ตฺวํ เอวรูปํ โวหารสมุจฺเฉทํ อตฺตนิ สมนุปสฺสสีติ.\csromanspacer{} โก จาหํ ภนฺเต, โก จ อริยสฺส วินเย สพฺเพน สพฺพํ สพฺพถา สพฺพํ โวหารสมุจฺเฉโท.\csromanspacer{} อารกา อหํ ภนฺเต อริยสฺส วินเย สพฺเพน สพฺพํ สพฺพถา สพฺพํ โวหารสมุจฺเฉทา.\csromanspacer{} มยํ หิ ภนฺเต ปุพฺเพ อญฺญติตฺถิเย ปริพฺพาชเก อนาชานีเยว สมาเน อาชานียาติ อมญฺญิมฺหา, อนาชานีเยว สมาเน อาชานียโภชนํ โภชิมฺห, อนาชานีเยว สมาเน อาชานียฐาเน ฐปิมฺห.\csromanspacer{} ภิกฺขู ปน มยํ ภนฺเต อาชานีเยว สมาเน อนาชานียาติ อมญฺญิมฺห, อาชานีเยว สมาเน อนาชานียโภชนํ โภชิมฺห, อาชานีเยว สมาเน อนาชานียฐาเน ฐปิมฺห.\csromanspacer{} อิทานิ ปน มยํ ภนฺเต อญฺญติตฺถิเย ปริพฺพาชเก อนาชานีเยว สมาเน อนาชานียาติ ชานิสฺสาม, อนาชานีเยว สมาเน อนาชานียโภชนํ โภเชสฺสาม, อนาชานีเยว สมาเน อนาชานียฐาเน ฐเปสฺสาม.\csromanspacer{} ภิกฺขู ปน มยํ ภนฺเต อาชานีเยว สมาเน อาชานียาติ ชานิสฺสาม, อาชานีเยว สมาเน อาชานียโภชนํ โภเชสฺสาม, อาชานีเยว สมาเน อาชานิยฐาเน ฐเปสฺสาม.\csromanspacer{} อชเนสิ วต เม ภนฺเต ภควา สมเณสุ สมณปฺเปมํ สมเณสุ สมณปฺปสาทํ สมเณสุ สมณคารวํ.\csromanspacer{} อภิกฺกนฺตํ ภนฺเต, อภิกฺกนฺตํ ภนฺเต, เสยฺยถาปิ ภนฺเต นิกฺกุชฺชิตํ วา อุกฺกุชฺเชยฺย, ปฏิจฺฉนฺนํ วา วิวเรยฺย, มูฬฺหสฺส วา มคฺคํ อาจิกฺเขยฺย, อนฺธกาเร วา เตลปชฺโชตํ ธาเรยฺย ‘จกฺขุมนฺโต รูปานิ ทกฺขนฺตี’ติ.\csromanspacer{} เอวเมวํ โข ภนฺเต ภควตา อเนกปริยาเยน ธมฺโม ปกาสิโต, เอสาหํ ภนฺเต ภควนฺตํ สรณํ คจฺฉามิ ธมฺมญฺจ ภิกฺขุสํฆญฺจ, อุปาสกํ มํ ภควา ธาเรตุ อชฺชตคฺเค ปาณุเปตํ สรณํ คตนฺติ.}

\nitthitamruled{โปตลิยสุตฺตํ นิฏฺฐิตํ จตุตฺถํ.}
\csromanmatikaanchor{mtk.6}
\csromantocmarkref{subsection}{5. ชีวกสุตฺต}{mtk.6}
\bookmark[dest={mtk.6},level=2]{5. ชีวกสุตฺต}
\csromanheader{2}{5. ชีวกสุตฺต}
\proseitem{51}{เอวํ เม สุตํ\csromandash{}เอกํ สมยํ ภควา ราชคเห วิหรติ ชีวกสฺส โกมารภจฺจสฺส อมฺพวเน.\csromanspacer{} อถ โข ชีวโก โกมารภจฺโจ เยน ภควา เตนุปสงฺกมิ, อุปสงฺกมิตฺวา ภควนฺตํ อภิวาเทตฺวา เอกมนฺตํ \csromanfolio{32}นิสีทิ, เอกมนฺตํ นิสินฺโน โข ชีวโก โกมารภจฺโจ ภควนฺตํ เอตทโวจ “สุตํ เมตํ ภนฺเต ‘สมณํ โคตมํ อุทฺทิสฺส ปาณํ อารภนฺติ\footnote{อารมฺภนฺติ (ก)}, ตํ สมโณ โคตโม ชานํ อุทฺทิสฺสกตํ\footnote{อุทฺทิสฺสกฏํ (สี, อิ)} มํสํ ปริภุญฺชติ ปฏิจฺจกมฺมนฺ’ติ.\csromanspacer{} เย เต ภนฺเต เอวมาหํสุ ‘สมณํ โคตมํ อุทฺทิสฺส ปาณํ อารภนฺติ, ตํ สมโณ โคตโม ชานํ อุทฺทิสฺสกตํ มํสํ ปริภุญฺชติ ปฏิจฺจกมฺมนฺ’ติ.\csromanspacer{} กจฺจิ เต ภนฺเต ภควโต วุตฺตวาทิโน, น จ ภควนฺตํ อภูเตน อพฺภาจิกฺขนฺติ, ธมฺมสฺส จานุธมฺมํ พฺยากโรนฺติ, น จ โกจิ สหธมฺมิโก วาทานุวาโท คารยฺหํ ฐานํ อาคจฺฉตี”ติ.}

\proseitem{52}{เย เต ชีวก เอวมาหํสุ “สมณํ โคตมํ อุทฺทิสฺส ปาณํ อารภนฺติ, ตํ สมโณ โคตโม ชานํ อุทฺทิสฺสกตํ มํสํ ปริภุญฺชติ ปฏิจฺจกมฺมนฺ”ติ.\csromanspacer{} น เม เต วุตฺตวาทิโน, อพฺภาจิกฺขนฺติ จ มํ เต อสตา อภูเตน.\csromanspacer{} ตีหิ โข อหํ ชีวก ฐาเนหิ มํสํ อปริโภคนฺติ วทามิ ทิฏฺฐํ สุตํ ปริสงฺกิตํ, อิเมหิ โข อหํ ชีวก ตีหิ ฐาเนหิ มํสํ อปริโภคนฺติ วทามิ.\csromanspacer{} ตีหิ โข อหํ ชีวก ฐาเนหิ มํสํ ปริโภคนฺติ วทามิ อทิฏฺฐํ อสุตํ อปริสงฺกิตํ, อิเมหิ โข อหํ ชีวก ตีหิ ฐาเนหิ มํสํ ปริโภคนฺติ วทามิ.}

\proseitem{53}{อิธ ชีวก ภิกฺขุ อญฺญตรํ คามํ วา นิคมํ วา อุปนิสฺสาย วิหรติ, โส เมตฺตาสหคเตน เจตสา เอกํ ทิสํ ผริตฺวา วิหรติ.\csromanspacer{} ตถา ทุติยํ.\csromanspacer{} ตถา ตติยํ.\csromanspacer{} ตถา จตุตฺถํ.\csromanspacer{} อิติ อุทฺธมโธ ติริยํ สพฺพธิ สพฺพตฺตตาย สพฺพาวนฺตํ โลกํ เมตฺตาสหคเตน เจตสา วิปุเลน มหคฺคเตน อปฺปมาเณน อเวเรน อพฺยาพชฺเฌน ผริตฺวา วิหรติ, ตเมนํ คหปติ วา คหปติปุตฺโต วา อุปสงฺกมิตฺวา สฺวาตนาย ภตฺเตน นิมนฺเตติ, อากงฺขมาโนว\footnote{อากงฺขมาโน (สฺยา, กํ)} ชีวก ภิกฺขุ อธิวาเสติ.\csromanspacer{} โส ตสฺสา รตฺติยา อจฺจเยน ปุพฺพณฺหสมยํ นิวาเสตฺวา ปตฺตจีวรมาทาย เยน ตสฺส คหปติสฺส วา คหปติปุตฺตสฺส วา นิเวสนํ เตนุปสงฺกมติ, อุปสงฺกมิตฺวา ปญฺญตฺเต อาสเน นิสีทติ.\csromanspacer{} ตเมนํ โส คหปติ วา คหปติปุตฺโต วา ปณีเตน ปิณฺฑปาเตน ปริวิสติ, ตสฺส น เอวํ โหติ “สาธุ วต มายํ\footnote{มํ + อยํ = มายํ.} คหปติ วา คหปติปุตฺโต วา ปณีเตน}

\vspace{\tipitakaparskip}
\csromancenter{ชีวกสุตฺต (55) ปิณฺฑปาเตน ปริวิเสยฺยา”ติ. “อโห วต มายํ คหปติ วา คหปติปุตฺโต วา อายติํปิ เอวรูเปน ปณีเตน ปิณฺฑปาเตน ปริวิเสยฺยา”ติ เอวมฺปิสฺส น โหติ.\csromanspacer{} โส ตํ ปิณฺฑปาตํ อคถิโต\footnote{อคธิโต (สฺยา, กํ, ก)} อมุจฺฉิโต อนชฺโฌปนฺโน\footnote{อนชฺฌาปนฺโน (สฺยา, กํ, ก)} อาทีนวทสฺสาวี นิสฺสรณปญฺโญ ปริภุญฺชติ.\csromanspacer{} ตํ กิํ มญฺญสิ ชีวก, อปิ นุ โส ภิกฺขุ ตสฺมิํ สมเย อตฺตพฺยาพาธาย วา เจเตติ, ปรพฺยาพาธาย วา เจเตติ, อุภยพฺยาพาธาย วา เจเตตีติ.\csromanspacer{} โน เหตํ ภนฺเต.\csromanspacer{} นนุ โส ชีวก ภิกฺขุ ตสฺมิํ สมเย อนวชฺชํเยว อาหารํ อาหาเรตีติ.\csromanspacer{} เอวํ ภนฺเต, สุตํ เมตํ ภนฺเต “พฺรหฺมา เมตฺตาวิหารี”ติ, ตํ เม อิทํ ภนฺเต ภควา สกฺขิทิฏฺโฐ, ภควา หิ ภนฺเต เมตฺตาวิหารีติ.\csromanspacer{} เยน โข ชีวก ราเคน เยน โทเสน เยน โมเหน พฺยาปาทวา อสฺส, โส ราโค โส โทโส โส โมโห ตถาคตสฺส ปหีโน อุจฺฉินฺนมูโล ตาลาวตฺถุกโต อนภาวํกโต\footnote{อนภาวกโต (สี, อิ), อนภาวํคโต (สฺยา, กํ)}, อายติํ อนุปฺปาทธมฺโม.\csromanspacer{} สเจ โข เต ชีวก อิทํ สนฺธาย ภาสิตํ “อนุชานามิ เต เอตนฺ”ติ.\csromanspacer{} เอตเทว โข ปน เม ภนฺเต สนฺธาย ภาสิตํ\footnote{ภาสิตนฺติ (สฺยา)}.}
\proseitem{54}{\csromanfolio{33}อิธ ชีวก ภิกฺขุ อญฺญตรํ คามํ วา นิคมํ วา อุปนิสฺสาย วิหรติ, โส กรุณาสหคเตน เจตสา ฯเปฯ มุทิตาสหคเตน เจตสา ฯเปฯ อุเปกฺขาสหคเตน เจตสา เอกํ ทิสํ ผริตฺวา วิหรติ.\csromanspacer{} ตถา ทุติยํ.\csromanspacer{} ตถา ตติยํ.\csromanspacer{} ตถา จตุตฺถํ.\csromanspacer{} อิติ อุทฺธมโธ ติริยํ สพฺพธิ สพฺพตฺตตาย สพฺพาวนฺตํ โลกํ อุเปกฺขาสหคเตน เจตสา วิปุเลน มหคฺคเตน อปฺปมาเณน อเวเรน อพฺยาพชฺเฌน ผริตฺวา วิหรติ.\csromanspacer{} ตเมนํ คหปติ วา คหปติปุตฺโต วา อุปสงฺกมิตฺวา สฺวาตนาย ภตฺเตน นิมนฺเตติ, อากงฺขมาโนว ชีวก ภิกฺขุ อธิวาเสติ.\csromanspacer{} โส ตสฺสา รตฺติยา อจฺจเยน ปุพฺพณฺหสมยํ นิวาเสตฺวา ปตฺตจีวรมาทาย เยน คหปติสฺส วา คหปติปุตฺตสฺส วา นิเวสนํ เตนุปสงฺกมติ, อุปสงฺกมิตฺวา ปญฺญตฺเต อาสเน นิสีทติ.\csromanspacer{} ตเมนํ โส คหปติ วา คหปติปุตฺโต วา ปณีเตน ปิณฺฑปาเตน ปริวิสติ, ตสฺส น เอวํ โหติ “สาธุ วต มายํ คหปติ วา คหปติปุตฺโต วา ปณีเตน ปิณฺฑปาเตน ปริวิเสยฺยา”ติ. “อโห วต มายํ คหปติ วา คหปติปุตฺโต วา \csromanfolio{34}อายติํปิ เอวรูเปน ปณีเตน ปิณฺฑปาเตน ปริวิเสยฺยา”ติ เอวมฺปิสฺส น โหติ.\csromanspacer{} โส ตํ ปิณฺฑปาตํ อคถิโต อมุจฺฉิโต อนชฺโฌปนฺโน อาทีนวทสฺสาวี นิสฺสรณปญฺโญ ปริภุญฺชติ.\csromanspacer{} ตํ กิํ มญฺญสิ ชีวก, อปิ นุ โส ภิกฺขุ ตสฺมิํ สมเย อตฺตพฺยาพาธาย วา เจเตติ, ปรพฺยาพาธาย วา เจเตติ, อุภยพฺยาพาธาย วา เจเตตีติ.\csromanspacer{} โน เหตํ ภนฺเต.\csromanspacer{} นนุ โส ชีวก ภิกฺขุ ตสฺมิํ สมเย อนวชฺชํเยว อาหารํ อาหาเรตีติ.\csromanspacer{} เอวํ ภนฺเต, สุตํ เมตํ ภนฺเต “พฺรหฺมา อุเปกฺขาวิหารี”ติ, ตํ เม อิทํ ภนฺเต ภควา สกฺขิทิฏฺโฐ, ภควา หิ ภนฺเต อุเปกฺขาวิหารีติ.\csromanspacer{} เยน โข ชีวก ราเคน เยน โทเสน เยน โมเหน วิเหสวา อสฺส, อรติวา อสฺส, ปฏิฆวา อสฺส.\csromanspacer{} โส ราโค โส โทโส โส โมโห ตถาคหสฺส ปหีโน อุจฺฉินฺนมูโล ตาลาวตฺถุกโต อนภาวํกโต อายติํ อนุปฺปาทธมฺโม, สเจ โข เต ชีวก อิทํ สนฺธาย ภาสิตํ “อนุชานามิ เต เอตนฺ”ติ.\csromanspacer{} เอตเทว โข ปน เม ภนฺเต สนฺธาย ภาสิตํ.}

\proseitem{55}{โย โข ชีวก ตถาคตํ วา ตถาคตสาวกํ วา อุทฺทิสฺส ปาณํ อารภติ, โส ปญฺจหิ ฐาเนหิ พหุํ อปุญฺญํ ปสวติ.\csromanspacer{} ยมฺปิ โส คหปติ เอวมาห “คจฺฉถ อมุกํ นาม ปาณํ อาเนถา”ติ, อิมินา ปฐเมน ฐาเนน พหุํ อปุญฺญํ ปสวติ.\csromanspacer{} ยมฺปิ โส ปาโณ คลปฺปเวฐเกน\footnote{คลปฺปเวธเกน (พหูสุ)} อานียมาโน ทุกฺขํ โทมนสฺสํ ปฏิสํเวเทติ, อิมินา ทุติเยน ฐาเนน พหุํ อปุญฺญํ ปสวติ.\csromanspacer{} ยมฺปิ โส เอวมาห “คจฺฉถ อิมํ ปาณํ อารภถา”ติ, อิมินา ตติเยน ฐาเนน พหุํ อปุญฺญํ ปสวติ.\csromanspacer{} ยมฺปิ โส ปาโณ อารภิยมาโน ทุกฺขํ โทมนสฺสํ ปฏิสํเวเทติ, อิมินา จตุตฺเถน ฐาเนน พหุํ อปุญฺญํ ปสวติ.\csromanspacer{} ยมฺปิ โส ตถาคตํ วา ตถาคต- สาวกํ วา อกปฺปิเยน อาสาเทติ, อิมินา ปญฺจเมน ฐาเนน พหุํ อปุญฺญํ ปสวติ.\csromanspacer{} โย โข ชีวก ตถาคตํ วา ตถาคตสาวกํ วา อุทฺทิสฺส ปาณํ อารภติ, โส อิเมหิ ปญฺจหิ ฐาเนหิ พหุํ อปุญฺญํ ปสวตีติ.}

\prose{เอวํ วุตฺเต ชีวโก โกมารภจฺโจ ภควนฺตํ เอตทโวจ “อจฺฉริยํ ภนฺเต, อพฺภุตํ ภนฺเต, กปฺปิยํ วต ภนฺเต ภิกฺขู อาหารํ}

\vspace{\tipitakaparskip}
\csromancenter{อุปาลิสุตฺต (56) อาหาเรนฺติ, อนวชฺชํ วต ภนฺเต ภิกฺขู อาหารํ อาหาเรนฺติ, อภิกฺกนฺตํ ภนฺเต, อภิกฺกนฺตํ ภนฺเต ฯเปฯ อุปาสกํ มํ ภควา ธาเรตุ, อชฺชตคฺเค ปาณุเปตํ สรณํ คตนฺ”ติ.}
\nitthitamruled{ชีวกสุตฺตํ นิฏฺฐิตํ ปญฺจมํ.}
\csromanmatikaanchor{mtk.7}
\csromantocmarkref{subsection}{6. อุปาลิสุตฺต}{mtk.7}
\bookmark[dest={mtk.7},level=2]{6. อุปาลิสุตฺต}
\csromanheader{2}{6. อุปาลิสุตฺต}
\proseitem{56}{\csromanfolio{35}เอวํ เม สุตํ\csromandash{}เอกํ สมยํ ภควา นาฬนฺทายํ วิหรติ ปาวาริกมฺพวเน.\csromanspacer{} เตน โข ปน สมเยน นิคณฺโฐ นาฏปุตฺโต\footnote{นาถปุตฺโต (สี), นาตปุตฺโต (อิ)} นาฬนฺทายํ ปฏิวสติ มหติยา นิคณฺฐปริสาย สทฺธิํ.\csromanspacer{} อถ โข ทีฆตปสฺสี นิคณฺโฐ นาฬนฺทายํ ปิณฺฑาย จริตฺวา ปจฺฉาภตฺตํ ปิณฺฑปาตปฏิกฺกนฺโต เยน ปาวาริกมฺพวนํ เยน ภควา เตนุปสงฺกมิ, อุปสงฺกมิตฺวา ภควตา สทฺธิํ สมฺโมทิ, สมฺโมทนียํ กถํ สารณียํ วีติสาเรตฺวา เอกมนฺตํ อฏฺฐาสิ, เอกมนฺตํ ฐิตํ โข ทีฆตปสฺสิํ นิคณฺฐํ ภควา เอตทโวจ “สํวิชฺชนฺติ โข ตปสฺสิ\footnote{ทีฆตปสฺสิ (สฺยา, กํ, ก)} อาสนานิ, สเจ อากงฺขสิ นิสีทา”ติ.\csromanspacer{} เอวํ วุตฺเต ทีฆตปสฺสี นิคณฺโฐ อญฺญตรํ นีจํ อาสนํ คเหตฺวา เอกมนฺตํ นิสีทิ.\csromanspacer{} เอกมนฺตํ นิสินฺนํ โข ทีฆตปสฺสิํ นิคณฺฐํ ภควา เอตทโวจ “กติ ปน ตปสฺสิ นิคณฺโฐ นาฏปุตฺโต กมฺมานิ ปญฺญเปติ ปาปสฺส กมฺมสฺส กิริยาย ปาปสฺส กมฺมสฺส ปวตฺติยา”ติ.\csromanspacer{} น โข อาวุโส โคตม อาจิณฺณํ นิคณฺฐสฺส นาฏปุตฺตสฺส “กมฺมํ กมฺมนฺ”ติ ปญฺญเปตุํ, “ทณฺฑํ ทณฺฑนฺ”ติ โข อาวุโส โคตม อาจิณฺณํ นิคณฺฐสฺส นาฏปุตฺตสฺส ปญฺญเปตุนฺติ.\csromanspacer{} กติ ปน ตปสฺสิ นิคณฺโฐ นาฏปุตฺโต ทณฺฑานิ ปญฺญเปติ ปาปสฺส กมฺมสฺส กิริยาย ปาปสฺส กมฺมสฺส ปวตฺติยาติ.\csromanspacer{} ตีณิ โข อาวุโส โคตม นิคณฺโฐ นาฏปุตฺโต ทณฺฑานิ ปญฺญเปติ ปาปสฺส กมฺมสฺส กิริยาย ปาปสฺส กมฺมสฺส ปวตฺติยาติ.\csromanspacer{} เสยฺยถิทํ, กายทณฺฑํ วจีทณฺฑํ มโนทณฺฑนฺติ.\csromanspacer{} กิํ ปน ตปสฺสิ อญฺญเทว กายทณฺฑํ อญฺญํ วจีทณฺฑํ อญฺญํ มโนทณฺฑนฺติ.\csromanspacer{} อญฺญเทว อาวุโส โคตม กายทณฺฑํ อญฺญํ วจีทณฺฑํ อญฺญํ มโนทณฺฑนฺติ.\csromanspacer{} อิเมสํ ปน ตปสฺสิ ติณฺณํ ทณฺฑานํ เอวํ ปฏิวิภตฺตานํ เอวํ ปฏิวิสิฏฺฐานํ กตมํ ทณฺฑํ นิคณฺโฐ นาฏปุตฺโต มหาสาวชฺชตรํ ปญฺญเปติ ปาปสฺส กมฺมสฺส กิริยาย ปาปสฺส กมฺมสฺส \csromanfolio{36}ปวตฺติยา, ยทิ วา กายทณฺฑํ ยทิ วา วจีทณฺฑํ ยทิ วา มโนทณฺฑนฺติ.\csromanspacer{} อิเมสํ โข อาวุโส โคตม ติณฺณํ ทณฺฑานํ เอวํ ปฏิวิภตฺตานํ เอวํ ปฏิวิสิฏฺฐานํ กายทณฺฑํ นิคณฺโฐ นาฏปุตฺโต มหาสาวชฺชตรํ ปญฺญเปติ ปาปสฺส กมฺมสฺส กิริยาย ปาปสฺส กมฺมสฺส ปวตฺติยา, โน ตถา วจีทณฺฑํ โน ตถา มโนทณฺฑนฺติ. “กายทณฺฑนฺ”ติ ตปสฺสิ วเทสิ? “กายทณฺฑนฺ”ติ อาวุโส โคตม วทามิ. “กายทณฺฑนฺ”ติ ตปสฺสิ วเทสิ? “กายทณฺฑนฺ”ติ อาวุโส โคตม วทามิ. “กายทณฺฑนฺ”ติ ตปสฺสิ วเทสิ? “กายทณฺฑนฺ”ติ อาวุโส โคตม วทามีติ.\csromanspacer{} อิติห ภควา ทีฆตปสฺสิํ นิคณฺฐํ อิมสฺมิํ กถาวตฺถุสฺมิํ ยาวตติยกํ ปติฏฺฐาเปสิ.}

\proseitem{57}{เอวํ วุตฺเต ทีฆตปสฺสี นิคณฺโฐ ภควนฺตํ เอตทโวจ “ตฺวํ ปนาวุโส โคตม กติ ทณฺฑานิ ปญฺญเปสิ ปาปสฺส กมฺมสฺส กิริยาย ปาปสฺส กมฺมสฺส ปวตฺติยา”ติ.\csromanspacer{} น โข ตปสฺสิ อาจิณฺณํ ตถาคตสฺส “ทณฺฑํ ทณฺฑนฺ”ติ ปญฺญเปตุํ, “กมฺมํ กมฺมนฺ”ติ โข ตปสฺสิ อาจิณฺณํ ตถาคตสฺส ปญฺญเปตุนฺติ.\csromanspacer{} ตฺวํ ปนาวุโส โคตม กติ กมฺมานิ ปญฺญเปสิ ปาปสฺส กมฺมสฺส กิริยาย ปาปสฺส กมฺมสฺส ปวตฺติยาติ.\csromanspacer{} ตีณิ โข อหํ ตปสฺสิ กมฺมานิ ปญฺญเปมิ ปาปสฺส กมฺมสฺส กิริยาย ปาปสฺส กมฺมสฺส ปวตฺติยา, เสยฺยถิทํ, กายกมฺมํ วจีกมฺมํ มโนกมฺมนฺติ.\csromanspacer{} กิํ ปนาวุโส โคตม อญฺญเทว กายกมฺมํ อญฺญํ วจีกมฺมํ อญฺญํ มโนกมฺมนฺติ.\csromanspacer{} อญฺญเทว ตปสฺสิ กายกมฺมํ อญฺญํ วจีกมฺมํ อญฺญํ มโนกมฺมนฺติ.\csromanspacer{} อิเมสํ ปนาวุโส โคตม ติณฺณํ กมฺมานํ เอวํ ปฏิวิภตฺตานํ เอวํ ปฏิวิสิฏฺฐานํ กตมํ กมฺมํ มหาสาวชฺชตรํ ปญฺญเปสิ ปาปสฺส กมฺมสฺส กิริยาย ปาปสฺส กมฺมสฺส ปวตฺติยา, ยทิ วา กายกมฺมํ ยทิ วา วจีกมฺมํ ยทิ วา มโนกมฺมนฺติ.\csromanspacer{} อิเมสํ โข อหํ ตปสฺสิ ติณฺณํ กมฺมานํ เอวํ ปฏิวิภตฺตานํ เอวํ ปฏิวิสิฏฺฐานํ มโนกมฺมํ มหาสาวชฺชตรํ ปญฺญเปมิ ปาปสฺส กมฺมสฺส กิริยาย ปาปสฺส กมฺมสฺส ปวตฺติยา, โน ตถา กายกมฺมํ, โน ตถา วจีกมฺมนฺติ. “มโนกมฺมนฺ”ติ อาวุโส โคตม วเทสิ? “มโนกมฺมนฺ”ติ ตปสฺสิ วทามิ. “มโนกมฺมนฺ”ติ อาวุโส โคตม วเทสิ? “มโนกมฺมนฺ”ติ ตปสฺสิ วทามิ. “มโนกมฺมนฺ”ติ อาวุโส โคตม วเทสิ? “มโนกมฺมนฺ”ติ ตปสฺสิ วทามีติ.\csromanspacer{} อิติห ทีฆตปสฺสี นิคณฺโฐ ภควนฺตํ อิมสฺมิํ กถาวตฺถุสฺมิํ ยาวตติยกํ ปติฏฺฐาเปตฺวา อุฏฺฐายาสนา เยน นิคณฺโฐ นาฏปุตฺโต เตนุปสงฺกมิ.}

\csromanheader{2}{6. อุปาลิสุตฺต (56)}
\proseitem{58}{\csromanfolio{37}เตน โข ปน สมเยน นิคณฺโฐ นาฏปุตฺโต มหติยา คิหิปริสาย สทฺธิํ นิสินฺโน โหติ พาลกินิยา ปริสาย อุปาลิปมุขาย.\csromanspacer{} อทฺทสา โข นิคณฺโฐ นาฏปุตฺโต ทีฆตปสฺสิํ นิคณฺฐํ ทูรโตว อาคจฺฉนฺตํ, ทิสฺวาน ทีฆตปสฺสิํ นิคณฺฐํ เอตทโวจ “หนฺท กุโต นุ ตฺวํ ตปสฺสิ อาคจฺฉสิ ทิวา ทิวสฺสา”ติ.\csromanspacer{} อิโต หิ โข อหํ ภนฺเต อาคจฺฉามิ สมณสฺส โคตมสฺส สนฺติกาติ.\csromanspacer{} อหุ ปน เต ตปสฺสิ สมเณน โคตเมน สทฺธิํ โกจิเทว กถาสลฺลาโปติ.\csromanspacer{} อหุ โข เม ภนฺเต สมเณน โคตเมน สทฺธิํ โกจิเทว กถาสลฺลาโปติ.\csromanspacer{} ยถา กถํ ปน เต ตปสฺสิ อหุ สมเณน โคตเมน สทฺธิํ โกจิเทว กถาสลฺลาโปติ.\csromanspacer{} อถ โข ทีฆตปสฺสี นิคณฺโฐ ยาวตโก อโหสิ ภควตา สทฺธิํ กถาสลฺลาโป, ตํ สพฺพํ นิคณฺฐสฺส นาฏปุตฺตสฺส อาโรเจสิ.\csromanspacer{} เอวํ วุตฺเต นิคณฺโฐ นาฏปุตฺโต ทีฆตปสฺสิํ นิคณฺฐํ เอตทโวจ “สาธุ สาธุ ตปสฺสิ, ยถา ตํ สุตวตา สาวเกน สมฺมเทว สตฺถุสาสนํ อาชานนฺเตน, เอวเมว ทีฆตปสฺสินา นิคณฺเฐน สมณสฺส โคตมสฺส พฺยากตํ, กิํ หิ โสภติ ฉโว มโนทณฺโฑ อิมสฺส เอวํ โอฬาริกสฺส กายทณฺฑสฺส อุปนิธาย, อถ โข กายทณฺโฑว มหาสาวชฺชตโร ปาปสฺส กมฺมสฺส กิริยาย ปาปสฺส กมฺมสฺส ปวตฺติยา, โน ตถา วจีทณฺโฑ, โน ตถา มโนทณฺโฑ”ติ.}

\proseitem{59}{เอวํ วุตฺเต อุปาลิ คหปติ นิคณฺฐํ นาฏปุตฺตํ เอตทโวจ “สาธุ สาธุ ภนฺเต ทีฆตปสฺสี\footnote{ตปสฺสี (สี, อิ)}, ยถา ตํ สุตวตา สาวเกน สมฺมเทว สตฺถุสาสนํ อาชานนฺเตน, เอวเมวํ ภทนฺเตน ตปสฺสินา สมณสฺส โคตมสฺส พฺยากตํ, กิํ หิ โสภติ ฉโว มโนทณฺโฑ อิมสฺส เอวํ โอฬาริกสฺส กายทณฺฑสฺส อุปนิธาย, อถ โข กายทณฺโฑว มหาสาวชฺชตโร ปาปสฺส กมฺมสฺส กิริยาย ปาปสฺส กมฺมสฺส ปวตฺติยา, โน ตถา วจีทณฺโฑ, โน ตถา มโนทณฺโฑ.\csromanspacer{} หนฺท จาหํ ภนฺเต คจฺฉามิ สมณสฺส โคตมสฺส อิมสฺมิํ กถาวตฺถุสฺมิํ วาทํ อาโรเปสฺสามิ.\csromanspacer{} สเจ เม สมโณ โคตโม ตถา ปติฏฺฐหิสฺสติ, ยถา ภทนฺเตน ตปสฺสินา ปติฏฺฐาปิตํ.\csromanspacer{} เสยฺยถาปิ นาม พลวา ปุริโส ทีฆโลมิกํ เอฬกํ โลเมสุ คเหตฺวา อากฑฺเฒยฺย ปริกฑฺเฒยฺย สมฺปริกฑฺเฒยฺย.\csromanspacer{} เอวเมวาหํ สมณํ โคตมํ วาเทน วาทํ อากฑฺฒิสฺสามิ ปริกฑฺฒิสฺสามิ \csromanfolio{38}สมฺปริกฑฺฒิสฺสามิ.\csromanspacer{} เสยฺยถาปิ นาม พลวา โสณฺฑิกากมฺมกาโร มหนฺตํ โสณฺฑิกากิลญฺชํ คมฺภีเร อุทกรหเท ปกฺขิปิตฺวา กณฺเณ คเหตฺวา อากฑฺเฒยฺย ปริกฑฺเฒยฺย สมฺปริกฑฺเฒยฺย.\csromanspacer{} เอวเมวาหํ สมณํ โคตมํ วาเทน วาทํ อากฑฺฒิสฺสามิ ปริกฑฺฒิสฺสามิ สมฺปริกฑฺฒิสฺสามิ.\csromanspacer{} เสยฺยถาปิ นาม พลวา โสณฺฑิกาธุตฺโต วาลํ\footnote{ถาลํ (ก)} กณฺเณ คเหตฺวา โอธุเนยฺย นิทฺธุเนยฺย นิปฺโผเฏยฺย\footnote{นิจฺฉาเทยฺย (สี, อิ, ก), นิจฺฉาเฏยฺย (ก), นิปฺโปเฐยฺย (สฺยา, กํ)}.\csromanspacer{} เอวเมวาหํ สมณํ โคตมํ วาเทน วาทํ โอธุนิสฺสามิ นิทฺทุนิสฺสามิ นิปฺโผเฏสฺสามิ.\csromanspacer{} เสยฺยถาปิ นาม กุญฺชโร สฏฺฐิหายโน คมฺภีรํ โปกฺขรณิํ โอคาเหตฺวา สาณโธวิกํ นาม กีฬิตชาตํ กีฬติ.\csromanspacer{} เอวเมวาหํ สมณํ โคตมํ สาณโธวิกํ มญฺเญ กีฬิตชาตํ กีฬิสฺสามิ.\csromanspacer{} หนฺท จาหํ ภนฺเต คจฺฉามิ สมณสฺส โคตมสฺส อิมสฺมิํ กถาวตฺถุสฺมิํ วาทํ อาโรเปสฺสามี”ติ.\csromanspacer{} คจฺฉ ตฺวํ คหปติ สมณสฺส โคตมสฺส อิมสฺมิํ กถาวตฺถุสฺมิํ วาทํ อาโรเปหิ, อหํ วา หิ คหปติ สมณสฺส โคตมสฺส วาทํ อาโรเปยฺยํ, ทีฆตปสฺสี วา นิคณฺโฐ, ตฺวํ วาติ.}

\proseitem{60}{เอวํ วุตฺเต ทีฆตปสฺสี นิคณฺโฐ นิคณฺฐํ นาฏปุตฺตํ เอตทโวจ “น โข เมตํ ภนฺเต รุจฺจติ, ยํ อุปาลิ คหปติ สมณสฺส โคตมสฺส วาทํ อาโรเปยฺย, สมโณ หิ ภนฺเต โคตโม มายาวี อาวฏฺฏนิํ.\csromanspacer{} มายํ ชานาติ, ยาย อญฺญติตฺถิยานํ สาวเก อาวฏฺเฏตี”ติ.\csromanspacer{} อฏฺฐานํ โข เอตํ ตปสฺสิ อนวกาโส, ยํ อุปาลิ คหปติ สมณสฺส โคตมสฺส สาวกตฺตํ อุปคจฺเฉยฺย.\csromanspacer{} ฐานํ จ โข เอตํ วิชฺชติ, ยํ สมโณ โคตโม อุปาลิสฺส คหปติสฺส สาวกตฺตํ อุปคจฺเฉยฺย.\csromanspacer{} คจฺฉ ตฺวํ คหปติ สมณสฺส โคตมสฺส อิมสฺมิํ กถาวตฺถุสฺมิํ วาทํ อาโรเปหิ, อหํ วา หิ คหปติ สมณสฺส โคตมสฺส วาทํ อาโรเปยฺยํ, ทีฆตปสฺสี วา นิคณฺโฐ, ตฺวํ วาติ.\csromanspacer{} ทุติยมฺปิ โข ทีฆตปสฺสี ฯเปฯ.\csromanspacer{} ตติยมฺปิ โข ทีฆตปสฺสี นิคณฺโฐ นิคณฺฐํ นาฏปุตฺตํ เอตทโวจ “น โข เมตํ ภนฺเต รุจฺจติ, ยํ อุปาลิ คหปติ สมณสฺส โคตมสฺส วาทํ อาโรเปยฺย, สมโณ หิ ภนฺเต โคตโม มายาวี อาวฏฺฏนิํ มายํ ชานาติ, ยาย อญฺญติตฺถิยานํ สาวเก อาวฏฺเฏตี”ติ.\csromanspacer{} อฏฺฐานํ โข เอตํ ตปสฺสิ อนวกาโส, ยํ อุปาลิ คหปติ สมณสฺส โคตมสฺส}

\vspace{\tipitakaparskip}
\csromancenter{อุปาลิสุตฺต (56) สาวกตฺตํ อุปคจฺเฉยฺย.\csromanspacer{} ฐานํ จ โข เอตํ วิชฺชติ, ยํ สมโณ โคตโม อุปาลิสฺส คหปติสฺส สาวกตฺตํ อุปคจฺเฉยฺย.\csromanspacer{} คจฺฉ ตฺวํ คหปติ สมณสฺส โคตมสฺส อิมสฺมิํ กถาวตฺถุสฺมิํ วาทํ อาโรเปหิ, อหํ วา หิ คหปติ สมณสฺส โคตมสฺส วาทํ อาโรเปยฺยํ, ทีฆตปสฺสี วา นิคณฺโฐ, ตฺวํ วาติ. “เอวํ ภนฺเต”ติ โข อุปาลิ คหปติ นิคณฺฐสฺส นาฏปุตฺตสฺส ปฏิสฺสุตฺวา อุฏฺฐายาสนา นิคณฺฐํ นาฏปุตฺตํ อภิวาเทตฺวา ปทกฺขิณํ กตฺวา เยน ปาวาริกมฺพวนํ เยน ภควา เตนุปสงฺกมิ, อุปสงฺกมิตฺวา ภควนฺตํ อภิวาเทตฺวา เอกมนฺตํ นิสีทิ, เอกมนฺตํ นิสินฺโน โข อุปาลิ คหปติ ภควนฺตํ เอตทโวจ “อาคมา นุ ขฺวิธ ภนฺเต ทีฆตปสฺสี นิคณฺโฐ”ติ.\csromanspacer{} อาคมา ขฺวิธ คหปติ ทีฆตปสฺสี นิคณฺโฐติ.\csromanspacer{} อหุ โข ปน เต ภนฺเต ทีฆตปสฺสินา นิคณฺเฐน สทฺธิํ โกจิเทว กถาสลฺลาโปติ.\csromanspacer{} อหุ โข เม คหปติ ทีฆตปสฺสินา นิคณฺเฐน สทฺธิํ โกจิเทว กถาสลฺลาโปติ.\csromanspacer{} ยถา กถํ ปน เต ภนฺเต อหุ ทีฆตปสฺสินา นิคณฺเฐน สทฺธิํ โกจิเทว กถาสลฺลาโปติ.\csromanspacer{} อถ โข ภควา ยาวตโก อโหสิ ทีฆตปสฺสินา นิคณฺเฐน สทฺธิํ กถาสลฺลาโป, ตํ สพฺพํ อุปาลิสฺส คหปติสฺส อาโรเจสิ.}
\proseitem{61}{\csromanfolio{39}เอวํ วุตฺเต อุปาลิ คหปติ ภควนฺตํ เอตทโวจ “สาธุ สาธุ ภนฺเต ตปสฺสี, ยถา ตํ สุตวตา สาวเกน สมฺมเทว สตฺถุสาสนํ อาชานนฺเตน, เอวเมวํ ทีฆตปสฺสินา นิคณฺเฐน ภควโต พฺยากตํ, กิํ หิ โสภติ ฉโว มโนทณฺโฑ อิมสฺส เอวํ โอฬาริกสฺส กายทณฺฑสฺส อุปนิธาย, อถ โข กายทณฺโฑว มหาสาวชฺชตโร ปาปสฺส กมฺมสฺส กิริยาย ปาปสฺส กมฺมสฺส ปวตฺติยา, โน ตถา วจีทณฺโฑ, โน ตถา มโนทณฺโฑติ.\csromanspacer{} สเจ โข ตฺวํ คหปติ สจฺเจ ปติฏฺฐาย มนฺเตยฺยาสิ, สิยา โน เอตฺถ กถาสลฺลาโปติ.\csromanspacer{} สจฺเจ อหํ ภนฺเต ปติฏฺฐาย มนฺเตสฺสามิ, โหตุ โน เอตฺถ กถาสลฺลาโปติ.}

\proseitem{62}{ตํ กิํ มญฺญสิ คหปติ, อิธสฺส นิคณฺโฐ อาพาธิโก ทุกฺขิโต พาฬฺหคิลาโน สีโตทกปฏิกฺขิตฺโต อุโณฺหทกปฏิเสวี, โส สีโตทกํ อลภมาโน กาลงฺกเรยฺย, อิมสฺส ปน คหปติ นิคณฺโฐ นาฏปุตฺโต กตฺถูปปตฺติํ ปญฺญเปตีติ.\csromanspacer{} อตฺถิ ภนฺเต มโนสตฺตา นาม เทวา, ตตฺถ โส อุปปชฺชติ.\csromanspacer{} ตํ กิสฺส เหตุ.\csromanspacer{} อสุ หิ \csromanfolio{40}ภนฺเต มโนปฏิพทฺโธ กาลงฺกโรตีติ.\csromanspacer{} มนสิ กโรหิ คหปติ\footnote{คหปติ คหปติ มนสิ กโรหิ (สี, สฺยา, กํ), คหปติ มนสิ กโรหิ (ก), คหปติ คหปติ (อิ)}, มนสิ กริตฺวา โข คหปติ พฺยากโรหิ, น โข เต สนฺธิยติ ปุริเมน วา ปจฺฉิมํ, ปจฺฉิเมน วา ปุริมํ, ภาสิตา โข ปน เต คหปติ เอสา วาจา “สจฺเจ อหํ ภนฺเต ปติฏฺฐาย มนฺเตสฺสามิ, โหตุ โน เอตฺถ กถาสลฺลาโป”ติ.\csromanspacer{} กิญฺจาปิ ภนฺเต ภควา เอวมาห, อถ โข กายทณฺโฑว มหาสาวชฺชตโร ปาปสฺส กมฺมสฺส กิริยาย ปาปสฺส กมฺมสฺส ปวตฺติยา, โน ตถา วจีทณฺโฑ, โน ตถา มโนทณฺโฑติ.}

\proseitem{63}{ตํ กิํ มญฺญสิ คหปติ, อิธสฺส นิคณฺโฐ นาฏปุตฺโต จาตุยามสํวรสํวุโต สพฺพวาริวาริโต สพฺพวาริยุตฺโต สพฺพวาริธุโต สพฺพวาริผุโฏ, โส อภิกฺกมนฺโต ปฏิกฺกมนฺโต พหู ขุทฺทเก ปาเณ สํฆาตํ อาปาเทติ, อิมสฺส ปน คหปติ นิคณฺโฐ นาฏปุตฺโต กํ วิปากํ ปญฺญเปตีติ.\csromanspacer{} อสญฺเจตนิกํ ภนฺเต นิคณฺโฐ นาฏปุตฺโต โน มหาสาวชฺชํ ปญฺญเปตีติ.\csromanspacer{} สเจ ปน คหปติ เจเตตีติ.\csromanspacer{} มหาสาวชฺชํ ภนฺเต โหตีติ.\csromanspacer{} เจตนํ ปน คหปติ นิคณฺโฐ นาฏปุตฺโต กิสฺมิํ ปญฺญเปตีติ.\csromanspacer{} มโนทณฺฑสฺมิํ ภนฺเตติ.\csromanspacer{} มนสิ กโรหิ คหปติ, มนสิ กริตฺวา โข คหปติ พฺยากโรหิ, น โข เต สนฺธิยติ ปุริเมน วา ปจฺฉิมํ, ปจฺฉิเมน วา ปุริมํ, ภาสิตา โข ปน เต คหปติ เอสา วาจา “สจฺเจ อหํ ภนฺเต ปติฏฺฐาย มนฺเตสฺสามิ, โหตุ โน เอตฺถ กถาสลฺลาโป”ติ.\csromanspacer{} กิญฺจาปิ ภนฺเต ภควา เอวมาห, อถ โข กายทณฺโฑว มหาสาวชฺชตโร ปาปสฺส กมฺมสฺส กิริยาย ปาปสฺส กมฺมสฺส ปวตฺติยา, โน ตถา วจีทณฺโฑ, โน ตถา มโนทณฺโฑติ.}

\proseitem{64}{ตํ กิํ มญฺญสิ คหปติ, อยํ นาฬนฺทา อิทฺธา เจว ผีตา จ พหุชนา อากิณฺณมนุสฺสาติ.\csromanspacer{} เอวํ ภนฺเต อยํ นาฬนฺทา อิทฺธา เจว ผีตา จ พหุชนา อากิณฺณมนุสฺสาติ.\csromanspacer{} ตํ กิํ มญฺญสิ คหปติ, อิธ ปุริโส อาคจฺเฉยฺย อุกฺขิตฺตาสิโก, โส เอวํ วเทยฺย “อหํ ยาวติกา อิมิสฺสา นาฬนฺทาย ปาณา, เต เอเกน ขเณน เอเกน มุหุตฺเตน เอกํ มํสขลํ เอกํ มํสปุญฺชํ กริสฺสามี”ติ, ตํ กิํ มญฺญสิ คหปติ, ปโหติ นุ โข โส ปุริโส ยาวติกา อิมิสฺสา นาฬนฺทาย ปาณา, เต เอเกน ขเณน เอเกน}

\vspace{\tipitakaparskip}
\csromancenter{อุปาลิสุตฺต (56) มุหุตฺเตน เอกํ มํสขลํ เอกํ มํสปุญฺชํ กาตุนฺติ.\csromanspacer{} ทสปิ ภนฺเต ปุริสา วีสํปิ ภนฺเต ปุริสา ติํสํปิ ภนฺเต ปุริสา จตฺตารีสํปิ ภนฺเต ปุริสา ปญฺญาสํปิ ภนฺเต ปุริสา นปฺปโหนฺติ ยาวติกา อิมิสฺสา นาฬนฺทาย ปาณา, เต เอเกน ขเณน เอเกน มุหุตฺเตน เอกํ มํสขลํ เอกํ มํสปุญฺชํ กาตุํ, กิํ หิ โสภติ เอโก ฉโว ปุริโสติ.\csromanspacer{} ตํ กิํ มญฺญสิ คหปติ, อิธ อาคจฺเฉยฺย สมโณ วา พฺราหฺมโณ วา อิทฺธิมา เจโตวสิปฺปตฺโต, โส เอวํ วเทยฺย “อหํ อิมํ นาฬนฺทํ เอเกน มโนปโทเสน ภสฺมํ กริสฺสามี”ติ, ตํ กิํ มญฺญสิ คหปติ, ปโหติ นุ โข โส สมโณ วาพฺราหฺมโณ วา อิทฺธิมา เจโตวสิปฺปตฺโต อิมํ นาฬนฺทํ เอเกน มโนปโทเสน ภสฺมํ กาตุนฺติ.\csromanspacer{} ทสปิ ภนฺเต นาฬนฺทา วีสํปิ นาฬนฺทา ติํสํปิ นาฬานฺทา จตฺตารีสํปิ นาฬนฺทา ปญฺญาสํปิ นาฬนฺทา ปโหติ โส สมโณ วา พฺราหฺมโณ วา อิทฺธิมา เจโตวสิปฺปตฺโต เอเกน มโนปโทเสน ภสฺมํ กาตุํ, กิํ หิ โสภติ เอกา ฉวา นาฬนฺทาติ.\csromanspacer{} มนสิ กโรหิ คหปติ, มนสิ กริตฺวา โข คหปติ พฺยากโรหิ, น โข เต สนฺธิยติ ปุริเมน วา ปจฺฉิมํ, ปจฺฉิเมน วา ปุริมํ, ภาสิตา โข ปน เต คหปติ เอสา วาจา “สจฺเจ อหํ ภนฺเต ปติฏฺฐาย มนฺเตสฺสามิ, โหตุ โน เอตฺถ กถาสลฺลาโป”ติ.\csromanspacer{} กิญฺจาปิ ภนฺเต ภควา เอวมาห, อถ โข กายทณฺโฑว มหาสาวชฺชตโร ปาปสฺส กมฺมสฺส กิริยาย ปาปสฺส กมฺมสฺส ปวตฺติยา, โน ตถา วจีทณฺโฑ, โน ตถา มโนทณฺโฑติ.}
\proseitem{65}{\csromanfolio{41}ตํ กิํ มญฺญสิ คหปติ, สุตํ เต ทณฺฑกีรญฺญํ\footnote{ทณฺฑการญฺญํ (สี, อิ)} กาลิงฺคารญฺญํ มชฺฌารญฺญํ\footnote{เมชฺฌารญฺญํ (สี, สฺยา, กํ, อิ)} มาตงฺคารญฺญํ อรญฺญํ อรญฺญภูตนฺติ.\csromanspacer{} เอวํ ภนฺเต, สุตํ เม ทณฺฑกีรญฺญํ กาลิงฺคารญฺญํ มชฺฌารญฺญํ มาตงฺคารญฺญํ อรญฺญํ อรญฺญภูตนฺติ.\csromanspacer{} ตํ กิํ มญฺญสิ คหปติ, กินฺติ เต สุตํ, เกน ตํ ทณฺฑกีรญฺญํ กาลิงฺคารญฺญํ มชฺฌารญฺญํ มาตงฺคารญฺญํ อรญฺญํ อรญฺญภูตนฺติ.\csromanspacer{} สุตํ เมตํ ภนฺเต อิสีนํ มโนปโทเสน ตํ ทณฺฑกีรญฺญํ กาลิงฺคารญฺญํ มชฺฌารญฺญํ มาตงฺคารญฺญํ อรญฺญํ อรญฺญภูตนฺติ.\csromanspacer{} มนสิ กโรหิ คหปติ, มนสิ กริตฺวา โข คหปติ พฺยากโรหิ, น โข เต สนฺธิยติ ปุริเมน วา ปจฺฉิมํ, ปจฺฉิเมน วา ปุริมํ, ภาสิตา โข ปน เต คหปติ เอสา วาจา “สจฺเจ อหํ ภนฺเต ปติฏฺฐาย มนฺเตสฺสามิ, โหตุ โน เอตฺถ กถาสลฺลาโป”ติ.}

\proseitem{66}{\csromanfolio{42}ปุริเมเนวาหํ ภนฺเต โอปมฺเมน ภควโต อตฺตมโน อภิรทฺโธ, อปิ จาหํ อิมานิ ภควโต วิจิตฺรานิ ปญฺหปฏิภานานิ โสตุกาโม เอวาหํ ภควนฺตํ ปจฺจนีกํ กาตพฺพํ อมญฺญิสฺสํ.\csromanspacer{} อภิกฺกนฺตํ ภนฺเต, อภิกฺกนฺตํ ภนฺเต, เสยฺยถาปิ ภนฺเต นิกฺกุชฺชิตํ วา อุกฺกุชฺเชยฺย, ปฏิจฺฉนฺนํ วา วิวเรยฺย, มูฬฺหสฺส วา มคฺคํ อาจิกฺเขยฺย, อนฺธกาเร วา เตลปชฺโชตํ ธาเรยฺย “จกฺขุมนฺโต รูปานิ ทกฺขนฺตี”ติ, เอวเมวํ ภควตา อเนกปริยาเยน ธมฺโม ปกาสิโต, เอสาหํ ภนฺเต ภควนฺตํ สรณํ คจฺฉามิ ธมฺมญฺจ ภิกฺขุสํฆญฺจ, อุปาสกํ มํ ภควา ธาเรตุ อชฺชตคฺเค ปาณุเปตํ สรณํ คตนฺติ.}

\proseitem{67}{อนุวิจฺจการํ โข คหปติ กโรหิ, อนุวิจฺจกาโร ตุมฺหาทิสานํ ญาตมนุสฺสานํ สาธุ โหตีติ.\csromanspacer{} อิมินาปาหํ ภนฺเต ภควโต ภิยฺโยโส มตฺตาย อตฺตมโน อภิรทฺโธ, ยํ มํ ภควา เอวมาห “อนุวิจฺจการํ โข คหปติ กโรหิ, อนุวิจฺจกาโร ตุมฺหาทิสานํ ญาตมนุสฺสานํ สาธุ โหตี”ติ, มํ หิ ภนฺเต อญฺญติตฺถิยา สาวกํ ลภิตฺวา เกวลกปฺปํ นาฬนฺทํ ปฏากํ ปริหเรยฺยุํ “อุปาลิ อมฺหากํ คหปติ สาวกตฺตํ อุปคโต”ติ.\csromanspacer{} อถ จ ปน มํ ภควา เอวมาห “อนุวิจฺจการํ โข คหปติ กโรหิ, อนุวิจฺจกาโร ตุมฺหาทิสานํ ญาตมนุสฺสานํ สาธุ โหตี”ติ.\csromanspacer{} เอสาหํ ภนฺเต ทุติยมฺปิ ภควนฺตํ สรณํ คจฺฉามิ ธมฺมญฺจ ภิกฺขุสํฆญฺจ, อุปาสกํ มํ ภควา ธาเรตุ อชฺชตคฺเค ปาณุเปตํ สรณํ คตนฺติ.}

\proseitem{68}{ทีฆรตฺตํ โข เต คหปติ นิคณฺฐานํ โอปานภูตํ กุลํ, เยน เนสํ อุปคตานํ ปิณฺฑกํ ทาตพฺพํ มญฺเญยฺยาสีติ.\csromanspacer{} อิมินาปาหํ ภนฺเต ภควโต ภิยฺโยโส มตฺตาย อตฺตมโน อภิรทฺโธ, ยํ มํ ภควา เอวมาห “ทีฆรตฺตํ โข เต คหปติ นิคณฺฐานํ โอปานภูตํ กุลํ, เยน เนสํ อุปคตานํ ปิณฺฑกํ ทาตพฺพํ มญฺเญยฺยาสี”ติ.\csromanspacer{} สุตํ เมตํ ภนฺเต สมโณ โคตโม เอวมาห “มยฺหเมว ทานํ ทาตพฺพํ, นาญฺเญสํ ทานํ ทาตพฺพํ, มยฺหเมว สาวกานํ ทานํ ทาตพฺพํ, นาญฺเญสํ สาวกานํ ทานํ ทาตพฺพํ, มยฺหเมว ทินฺนํ มหปฺผลํ, นาญฺเญสํ ทินฺนํ มหปฺผลํ, มยฺหเมว สาวกานํ ทินฺนํ มหปฺผลํ, นาญฺเญสํ สาวกานํ ทินฺนํ มหปฺผลนฺ”ติ, อถ จ ปน มํ ภควา นิคณฺเฐสุปิ ทาเน สมาทเปติ, อปิ จ ภนฺเต มยเมตฺถ กาลํ ชานิสฺสาม, เอสาหํ ภนฺเต ตติยมฺปิ ภควนฺตํ สรณํ คจฺฉามิ ธมฺมญฺจ ภิกฺขุสํฆญฺจ, อุปาสกํ มํ ภควา ธาเรตุ อชฺชตคฺเค ปาณุเปตํ สรณํ คตนฺติ.}

\csromanheader{2}{6. อุปาลิสุตฺต (56)}
\proseitem{69}{\csromanfolio{43}อถ โข ภควา อุปาลิสฺส คหปติสฺส อนุปุพฺพิํ กถํ\footnote{อานุปุพฺพีกถํ (สี), อานุปุพฺพิกถํ (อิ), อนุปุพฺพิกถํ (สฺยา, กํ, ก)} กเถสิ, เสยฺยถิทํ, ทานกถํ, สีลกถํ สคฺคกถํ กามานํ อาทีนวํ โอการํ สํกิเลสํ เนกฺขมฺเม อานิสํสํ ปกาเสสิ.\csromanspacer{} ยทา ภควา อญฺญาสิ อุปาลิํ คหปติํ กลฺลจิตฺตํ มุทุจิตฺตํ วินีวรณจิตฺตํ อุทคฺคจิตฺตํ ปสนฺนจิตฺตํ, อถ ยา พุทฺธานํ สามุกฺกํสิกา ธมฺมเทสนา, ตํ ปกาเสสิ ทุกฺขํ สมุทยํ นิโรธํ มคฺคํ, เสยฺยถาปิ นาม สุทฺธํ วตฺถํ อปคตกาฬกํ สมฺมเทว รชนํ ปฏิคฺคเณฺหยฺย, เอวเมว อุปาลิสฺส คหปติสฺส ตสฺมิํ เยว อาสเน วิรชํ วีตมลํ ธมฺมจกฺขุํ อุทปาทิ “ยํ กิญฺจิ สมุทยธมฺมํ, สพฺพํ ตํ นิโรธธมฺมนฺ”ติ.\csromanspacer{} อถ โข อุปาลิ คหปติ ทิฏฺฐธมฺโม ปตฺตธมฺโม วิทิตธมฺโม ปริโยคาฬฺหธมฺโม ติณฺณวิจิกิจฺโฉ วิคตกถํกโถ เวสารชฺชปฺปตฺโต อปรปฺปจฺจโย สตฺถุสาสเน ภควนฺตํ เอตทโวจ “หนฺท จ ทานิ มยํ ภนฺเต คจฺฉาม พหุกิจฺจา มยํ พหุกรณียา”ติ.\csromanspacer{} ยสฺสทานิ ตฺวํ คหปติ กาลํ มญฺญสีติ.}

\proseitem{70}{อถ โข อุปาลิ คหปติ ภควโต ภาสิตํ อภินนฺทิตฺวา อนุโมทิตฺวา อุฏฺฐายาสนา ภควนฺตํ อภิวาเทตฺวา ปทกฺขิณํ กตฺวา เยน สกํ นิเวสนํ เตนุปสงฺกมิ, อุปสงฺกมิตฺวา โทวาริกํ อามนฺเตสิ “อชฺชตคฺเค สมฺม โทวาริก อาวรามิ ทฺวารํ นิคณฺฐานํ นิคณฺฐีนํ, อนาวฏํ ทฺวารํ ภควโต ภิกฺขูนํ ภิกฺขุนีนํ อุปาสกานํ อุปาสิกานํ.\csromanspacer{} สเจ โกจิ นิคณฺโฐ อาคจฺฉติ, ตเมนํ ตฺวํ เอวํ วเทยฺยาสิ ‘ติฏฺฐ ภนฺเต มา ปาวิสิ, อชฺชตคฺเค อุปาลิ คหปติ สมณสฺส โคตมสฺส สาวกตฺตํ อุปคโต, อาวฏํ ทฺวารํ นิคณฺฐานํ นิคณฺฐีนํ, อนาวฏํ ทฺวารํ ภควโต ภิกฺขูนํ ภิกฺขุนีนํ อุปาสกานํ อุปาสิกานํ, สเจ เต ภนฺเต ปิณฺฑเกน อตฺโถ, เอตฺเถว ติฏฺฐ, เอตฺเถว เต อาหริสฺสนฺตี’ติ”. “เอวํ ภนฺเต”ติ โข โทวาริโก อุปาลิสฺส คหปติสฺส ปจฺจสฺโสสิ.}

\proseitem{71}{อสฺโสสิ โข ทีฆตปสฺสี นิคณฺโฐ “อุปาลิ กิร คหปติ สมณสฺส โคตมสฺส สาวกตฺตํ อุปคโต”ติ.\csromanspacer{} อถ โข ทีฆตปสฺสี นิคณฺโฐ เยน นิคณฺโฐ นาฏปุตฺโต เตนุปสงฺกมิ, อุปสงฺกมิตฺวา นิคณฺฐํ นาฏปุตฺตํ เอตทโวจ “สุตํ เมตํ ภนฺเต อุปาลิ กิร คหปติ สมณสฺส โคตมสฺส สาวกตฺตํ อุปคโต”ติ.\csromanspacer{} อฏฺฐานํ โข เอตํ \csromanfolio{44}ตปสฺสิ อนวกาโส, ยํ อุปาลิ คหปติ สมณสฺส โคตมสฺส สาวกตฺตํ อุปคจฺเฉยฺย.\csromanspacer{} ฐานํ จ โข เอตํ วิชฺชติ, ยํ สมโณ โคตโม อุปาลิสฺส คหปติสฺส สาวกตฺตํ อุปคจฺเฉยฺยาติ.\csromanspacer{} ทุติยมฺปิ โข ทีฆตปสฺสี นิคณฺโฐ ฯเปฯ.\csromanspacer{} ตติยมฺปิ โข ทีฆตปสฺสี นิคณฺโฐ นิคณฺฐํ นาฏปุตฺตํ เอตทโวจ, สุตํ เมตํ ภนฺเต ฯเปฯ อุปาลิสฺส คหปติสฺส สาวกตฺตํ อุปคจฺเฉยฺยาติ.\csromanspacer{} หนฺทาหํ ภนฺเต คจฺฉามิ, ยาว ชานามิ ยทิ วา อุปาลิ คหปติ สมณสฺส โคตมสฺส สาวกตฺตํ อุปคโต ยทิ วา โนติ.\csromanspacer{} คจฺฉ ตฺวํ ตปสฺสิ ชานาหิ ยทิ วา อุปาลิ คหปติ สมณสฺส โคตมสฺส สาวกตฺตํ อุปคโต ยทิ วา โนติ.}

\proseitem{72}{อถ โข ทีฆตปสฺสี นิคณฺโฐ เยน อุปาลิสฺส คหปติสฺส นิเวสนํ เตนุปสงฺกมิ.\csromanspacer{} อทฺทสา โข โทวาริโก ทีฆตปสฺสิํ นิคณฺฐํ ทูรโตว อาคจฺฉนฺตํ, ทิสฺวาน ทีฆตปสฺสิํ นิคณฺฐํ เอตทโวจ “ติฏฺฐ ภนฺเต มา ปาวิสิ, อชฺชตคฺเค อุปาลิ คหปติ สมณสฺส โคตมสฺส สาวกตฺตํ อุปคโต, อาวฏํ ทฺวารํ นิคณฺฐานํ นิคณฺฐีนํ, อนาวฏํ ทฺวารํ ภควโต ภิกฺขูนํ ภิกฺขุนีนํ อุปาสกานํ อุปาสิกานํ, สเจ เต ภนฺเต ปิณฺฑเกน อตฺโถ, เอตฺเถว ติฏฺฐ, เอตฺเถว เต อาหริสฺสนฺตี”ติ. “น เม อาวุโส ปิณฺฑเกน อตฺโถ”ติ วตฺวา ตโต ปฏินิวตฺติตฺวา เยน นิคณฺโฐ นาฏปุตฺโต เตนุปสงฺกมิ, อุปสงฺกมิตฺวา นิคณฺฐํ นาฏปุตฺตํ เอตทโวจ “สจฺจํเยว โข ภนฺเต ยํ อุปาลิ คหปติ สมณสฺส โคตมสฺส สาวกตฺตํ อุปคโต, เอตํ โข เต อหํ ภนฺเต นาลตฺถํ, ‘น โข เม ภนฺเต รุจฺจติ, ยํ อุปาลิ คหปติ สมณสฺส โคตมสฺส วาทํ อาโรเปยฺย.\csromanspacer{} สมโณ หิ ภนฺเต โคตโม มายาวี อาวฏฺฏนิํ มายํ ชานาติ, ยาย อญฺญติตฺถิยานํ สาวเก อาวฏฺเฏตี’ติ.\csromanspacer{} อาวฏฺโฏ โข เต ภนฺเต อุปาลิ คหปติ สมเณน โคตเมน อาวฏฺฏนิยา มายายา”ติ.\csromanspacer{} อฏฺฐานํ โข เอตํ ตปสฺสิ อนวกาโส, ยํ อุปาลิ คหปติ สมณสฺส โคตมสฺส สาวกตฺตํ อุปคจฺเฉยฺย.\csromanspacer{} ฐานํ จ โข เอตํ วิชฺชติ, ยํ สมโณ โคตโม อุปาลิสฺส คหปติสฺส สาวกตฺตํ อุปคจฺเฉยฺยาติ.\csromanspacer{} ทุติยมฺปิ โข ทีฆตปสฺสี นิคณฺโฐ นิคณฺฐํ นาฏปุตฺตํ เอตทโวจ “สจฺจํเยว ภนฺเต ฯเปฯ อุปาลิสฺส คหปติสฺส สาวกตฺตํ อุปคจฺเฉยฺยาติ.\csromanspacer{} ตติยมฺปิ โข ทีฆตปสฺสี นิคณฺโฐ นิคณฺฐํ นาฏปุตฺตํ เอตทโวจ “สจฺจํเยว โข ภนฺเต ฯเปฯ อุปาลิสฺส คหปติสฺส สาวกตฺตํ อุปคจฺเฉยฺยาติ.\csromanspacer{} หนฺท}

\vspace{\tipitakaparskip}
\csromancenter{อุปาลิสุตฺต (56) จาหํ ตปสฺสิ คจฺฉามิ, ยาว จาหํ สามํเยว ชานามิ ยทิ วา อุปาลิ คหปติ สมณสฺส โคตมสฺส สาวกตฺตํ อุปคโต ยทิ วา โนติ.}
\prose{\csromanfolio{45}อถ โข นิคณฺโฐ นาฏปุตฺโต มหติยา นิคณฺฐปริสาย สทฺธิํ เยน อุปาลิสฺส คหปติสฺส นิเวสนํ เตนุปสงฺกมิ.\csromanspacer{} อทฺทสา โข โทวาริโก นิคณฺฐํ นาฏปุตฺตํ ทูรโตว อาคจฺฉนฺตํ, ทิสฺวาน นิคณฺฐํ นาฏปุตฺตํ เอตทโวจ “ติฏฺฐ ภนฺเต มา ปาวิสิ, อชฺชตคฺเค อุปาลิ คหปติ สมณสฺส โคตมสฺส สาวกตฺตํ อุปคโต, อาวฏํ ทฺวารํ นิคณฺฐานํ นิคณฺฐีนํ, อนาวฏํ ทฺวารํ ภควโต ภิกฺขูนํ ภิกฺขุนีนํ อุปาสกานํ อุปาสิกานํ, สเจ เต ภนฺเต ปิณฺฑเกน อตฺโถ, เอตฺเถว ติฏฺฐ, เอตฺเถว เต อาหริสฺสนฺตี”ติ.\csromanspacer{} เตน หิ สมฺม โทวาริก เยน อุปาลิ คหปติ เตนุปสงฺกม, อุปสงฺกมิตฺวา อุปาลิํ คหปติํ เอวํ วเทหิ “นิคณฺโฐ ภนฺเต นาฏปุตฺโต มหติยา นิคณฺฐปริสาย สทฺธิํ พหิทฺวารโกฏฺฐเก ฐิโต, โส เต ทสฺสนกาโม”ติ. “เอวํ ภนฺเต”ติ โข โทวาริโก นิคณฺฐสฺส นาฏปุตฺตสฺส ปฏิสฺสุตฺวา เยน อุปาลิ คหปติ เตนุปสงฺกมิ, อุปสงฺกมิตฺวา อุปาลิํ คหปติํ เอตทโวจ “นิคณฺโฐ ภนฺเต นาฏปุตฺโต มหติยา นิคณฺฐปริสาย สทฺธิํ พหิทฺวารโกฏฺฐเก ฐิโต, โส เต ทสฺสนกาโม”ติ.\csromanspacer{} เตน หิ สมฺม โทวาริก มชฺฌิมาย ทฺวารสาลาย อาสนานิ ปญฺญเปหีติ. “เอวํ ภนฺเต”ติ โข โทวาริโก อุปาลิสฺส คหปติสฺส ปฏิสฺสุตฺวา มชฺฌิมาย ทฺวารสาลาย อาสนานิ ปญฺญเปตฺวา เยน อุปาลิ คหปติ เตนุปสงฺกมิ, อุปสงฺกมิตฺวา อุปาลิํ คหปติํ เอตทโวจ “ปญฺญตฺตานิ โข ภนฺเต มชฺฌิมาย ทฺวารสาลาย อาสนานิ, ยสฺสทานิ กาลํ มญฺญสี”ติ.}

\proseitem{73}{อถ โข อุปาลิ คหปติ เยน มชฺฌิมา ทฺวารสาลา เตนุปสงฺกมิ, อุปสงฺกมิตฺวา ยํ ตตฺถ อาสนํ อคฺคญฺจ เสฏฺฐญฺจ อุตฺตมญฺจ ปณีตญฺจ, ตตฺถ สามํ นิสีทิตฺวา โทวาริกํ อามนฺเตสิ\csromandash{}เตน หิ สมฺม โทวาริก เยน นิคณฺโฐ นาฏปุตฺโต เตนุปสงฺกม, อุปสงฺกมิตฺวา นิคณฺฐํ นาฏปุตฺตํ เอวํ วเทหิ “อุปาลิ ภนฺเต คหปติ เอวมาห ‘ปวิส กิร ภนฺเต สเจ อากงฺขสี’ติ”. “เอวํ ภนฺเต”ติ โข โทวาริโก อุปาลิสฺส คหปติสฺส ปฏิสฺสุตฺวา เยน นิคณฺโฐ นาฏปุตฺโต เตนุปสงฺกมิ, อุปสงฺกมิตฺวา นิคณฺฐํ นาฏปุตฺตํ เอตทโวจ “อุปาลิ ภนฺเต คหปติ เอวมาห ‘ปวิส กิร ภนฺเต สเจ อากงฺขสี’ติ”. \csromanfolio{46}อถ โข นิคณฺโฐ นาฏปุตฺโต มหติยา นิคณฺฐปริสาย สทฺธิํ เยน มชฺฌิมา ทฺวารสาลา เตนุปสงฺกมิ.\csromanspacer{} อถ โข อุปาลิ คหปติ ยํ สุทํ ปุพฺเพ ยโต ปสฺสติ นิคณฺฐํ นาฏปุตฺตํ ทูรโตว อาคจฺฉนฺตํ, ทิสฺวาน ตโต ปจฺจุคฺคนฺตฺวา ยํ ตตฺถ อาสนํ อคฺคญฺจ เสฏฺฐญฺจ อุตฺตมญฺจ ปณีตญฺจ, ตํ อุตฺตราสงฺเคน สมฺมชฺชิตฺวา\footnote{ปมชฺชิตฺวา (สี, อิ)} ปริคฺคเหตฺวา นิสีทาเปติ.\csromanspacer{} โส ทานิ ยํ ตตฺถ อาสนํ อคฺคญฺจ เสฏฺฐญฺจ อุตฺตมญฺจ ปณีตญฺจ, ตตฺถ สามํ นิสีทิตฺวา นิคณฺฐํ นาฏปุตฺตํ เอตทโวจ “สํวิชฺชนฺติ โข ภนฺเต อาสนานิ, สเจ อากงฺขสิ นิสีทา”ติ.\csromanspacer{} เอวํ วุตฺเต นิคณฺโฐ นาฏปุตฺโต อุปาลิํ คหปติํ เอตทโวจ “อุมฺมตฺโตสิ ตฺวํ คหปติ, ทตฺโตสิ ตฺวํ คหปติ ‘คจฺฉามหํ ภนฺเต สมณสฺส โคตมสฺส วาทํ อาโรเปสฺสามี’ติ คนฺตฺวา มหตาสิ วาทสงฺฆาเฏน ปฏิมุกฺโก อาคโต.\csromanspacer{} เสยฺยถาปิ คหปติ ปุริโส อณฺฑหารโก คนฺตฺวา อุพฺภเตหิ อณฺเฑหิ อาคจฺเฉยฺย, เสยฺยถา วา, ปน คหปติ ปุริโส อกฺขิกหารโก คนฺตฺวา อุพฺภเตหิ อกฺขีหิ อาคจฺเฉยฺย, เอวเมว โข ตฺวํ คหปติ ‘คจฺฉามหํ ภนฺเต สมณสฺส โคตมสฺส วาทํ อาโรเปสฺสามี’ติ คนฺตฺวา มหตาสิ วาทสงฺฆาเฏน ปฏิมุกฺโก อาคโต, อาวฏฺโฏสิ โข ตฺวํ คหปติ สมเณน โคตเมน อาวฏฺฏนิยา มายายา”ติ.}

\proseitem{74}{ภทฺทิกา ภนฺเต อาวฏฺฏนี มายา, กลฺยาณี ภนฺเต อาวฏฺฏนี มายา, ปิยา เม ภนฺเต ญาติสาโลหิตา อิมาย อาวฏฺฏนิยา อาวฏฺเฏยฺยุํ, ปิยานํปิ เม อสฺส ญาติสาโลหิตานํ ทีฆรตฺตํ หิตาย สุขาย.\csromanspacer{} สพฺเพ เจปิ ภนฺเต ขตฺติยา อิมาย อาวฏฺฏนิยา อาวฏฺเฏยฺยุํ, สพฺเพสานํปิสฺส ขตฺติยานํ ทีฆรตฺตํ หิตาย สุขาย.\csromanspacer{} สพฺเพ เจปิ ภนฺเต พฺราหฺมณา ฯเปฯ เวสฺสา ฯเปฯ สุทฺทา อิมาย อาวฏฺฏนิยา อาวฏฺเฏยฺยุํ, สพฺเพสานํปิสฺส สุทฺทานํ ทีฆรตฺตํ หิตาย สุขาย.\csromanspacer{} สเทวโก เจปิ ภนฺเต โลโก สมารโก สพฺรหฺมโก สสฺสมณพฺราหฺมณี ปชา สเทวมนุสฺสา อิมาย อาวฏฺฏนิยา อาวฏฺเฏยฺยุํ, สเทวกสฺสปิสฺส โลกสฺส สมารกสฺส สพฺรหฺมกสฺส สสฺสมณพฺราหฺมณิยา ปชาย สเทวมนุสฺสาย ทีฆรตฺตํ หิตาย สุขายาติ.\csromanspacer{} เตน หิ ภนฺเต อุปมํ เต กริสฺสามิ, อุปมายปิเธกจฺเจ วิญฺญู ปุริสา ภาสิตสฺส อตฺถํ อาชานนฺติ.}

\csromanheader{2}{6. อุปาลิสุตฺต (56)}
\proseitem{75}{\csromanfolio{47}ภูตปุพฺพํ ภนฺเต อญฺญตรสฺส พฺราหฺมณสฺส ชิณฺณสฺส วุฑฺฒสฺส มหลฺลกสฺส ทหรา มาณวิกา ปชาปตี อโหสิ คพฺภินี อุปวิชญฺญา, อถ โข ภนฺเต สา มาณวิกา ตํ พฺราหฺมณํ เอตทโวจ “คจฺฉ ตฺวํ พฺราหฺมณ, อาปณา มกฺกฏจฺฉาปกํ กิณิตฺวา อาเนหิ, โย เม กุมารกสฺส กีฬาปนโก ภวิสฺสตี”ติ.\csromanspacer{} เอวํ วุตฺเต โส พฺราหฺมโณ ตํ มาณวิกํ เอตทโวจ “อาคเมหิ ตาว โภติ, ยาว วิชายติ, สเจ ตฺวํ โภติ กุมารกํ วิชายิสฺสสิ, ตสฺสา เต อหํ อาปณา มกฺกฏจฺฉาปกํ กิณิตฺวา อาเนสฺสามิ, โย เต กุมารกสฺส กีฬาปนโก ภวิสฺสติ.\csromanspacer{} สเจ ปน ตฺวํ โภติ กุมาริกํ วิชายิสฺสสิ, ตสฺสา เต อหํ อาปณา มกฺกฏจฺฉาปิกํ กิณิตฺวา อาเนสฺสามิ, ยา เต กุมาริกาย กีฬาปนิกา ภวิสฺสตี”ติ.\csromanspacer{} ทุติยมฺปิ โข ภนฺเต สา มาณวิกา ฯเปฯ.\csromanspacer{} ตติยมฺปิ โข ภนฺเต สา มาณวิกา ตํ พฺราหฺมณํ เอตทโวจ “คจฺฉ ตฺวํ พฺราหฺมณ, อาปณา มกฺกฏจฺฉาปกํ กิณิตฺวา อาเนหิ, โย เม กุมารกสฺส กีฬาปนโก ภวิสฺสตี”ติ.\csromanspacer{} อถ โข ภนฺเต โส พฺราหฺมโณ ตสฺสา มาณวิกาย สารตฺโต ปฏิพทฺธจิตฺโต อาปณา มกฺกฏจฺฉาปกํ กิณิตฺวา อาเนตฺวา ตํ มาณวิกํ เอตทโวจ “อยํ เต โภติ อาปณา มกฺกฏจฺฉาปโก กิณิตฺวา อานีโต, โย เต กุมารกสฺส กีฬาปนโก ภวิสฺสตี”ติ.\csromanspacer{} เอวํ วุตฺเต ภนฺเต สา มาณวิกา ตํ พฺราหฺมณํ เอตทโวจ “คจฺฉ ตฺวํ พฺราหฺมณ อิมํ มกฺกฏจฺฉาปกํ อาทาย เยน รตฺตปาณิ รชกปุตฺโต เตนุปสงฺกม, อุปสงฺกมิตฺวา รตฺตปาณิํ รชกปุตฺตํ เอวํ วเทหิ, อิจฺฉามหํ สมฺม รตฺตปาณิ อิมํ มกฺกฏจฺฉาปกํ ปีตาวเลปนํ นาม รงฺคชาตํ รชิตํ อาโกฏิตปจฺจาโกฏิตํ อุภโตภาควิมฏฺฐนฺ”ติ.}

\prose{อถ โข ภนฺเต โส พฺราหฺมโณ ตสฺสา มาณวิกาย สารตฺโต ปฏิพทฺธจิตฺโต ตํ มกฺกฏจฺฉาปกํ อาทาย เยน รตฺตปาณิ รชกปุตฺโต เตนุปสงฺกมิ, อุปสงฺกมิตฺวา รตฺตปาณิํ รชกปุตฺตํ เอตทโวจ “อิจฺฉามหํ สมฺม รตฺตปาณิ อิมํ มกฺกฏจฺฉาปกํ ปีตาวเลปนํ นาม รงฺคชาตํ รชิตํ อาโกฏิตปจฺจาโกฏิตํ อุภโตภาควิมฏฺฐนฺ”ติ.\csromanspacer{} เอวํ วุตฺเต ภนฺเต รตฺตปาณิ รชกปุตฺโต ตํ พฺราหฺมณํ เอตทโวจ “อยํ โข เต ภนฺเต มกฺกฏจฺฉาปโก รงฺคกฺขโม หิ โข, โน อาโกฏนกฺขโม, โน วิมชฺชนกฺขโม”ติ.\csromanspacer{} เอวเมว โข ภนฺเต พาลานํ นิคณฺฐานํ วาโท, \csromanfolio{48}รงฺคกฺขโม หิ โข พาลานํ โน ปณฺฑิตานํ, โน อนุโยคกฺขโม, โน วิมชฺชนกฺขโม.\csromanspacer{} อถ โข ภนฺเต โส พฺราหฺมโณ อปเรน สมเยน นวํ ทุสฺสยุคํ อาทาย เยน รตฺตปาณิ รชกปุตฺโต เตนุปสงฺกมิ, อุปสงฺกมิตฺวา รตฺตปาณิํ รชกปุตฺตํ เอตทโวจ “อิจฺฉามหํ สมฺม รตฺตปาณิ อิมํ นวํ ทุสฺสยุคํ ปีตาวเลปนํ นาม รงฺคชาตํ รชิตํ อาโกฏิตปจฺจาโกฏิตํ อุภโตภาควิมฏฺฐนฺ”ติ.\csromanspacer{} เอวํ วุตฺเต ภนฺเต รตฺตปาณิ รชกปุตฺโต ตํ พฺราหฺมณํ เอตทโวจ “อิทํ โข เต ภนฺเต นวํ ทุสฺสยุคํ รงฺคกฺขมํ เจว อาโกฏนกฺขมํ จ วิมชฺชนกฺขมํ จา”ติ.\csromanspacer{} เอวเมว โข ภนฺเต ตสฺส ภควโต วาโท อรหโต สมฺมาสมฺพุทฺธสฺส รงฺคกฺขโม เจว ปณฺฑิตานํ โน พาลานํ, อนุโยคกฺขโม จ วิมชฺชนกฺขโม จาติ.}

\prose{สราชิกา โข คหปติ ปริสา เอวํ ชานาติ “อุปาลิ คหปติ นิคณฺฐสฺส นาฏปุตฺตสฺส สาวโก”ติ, กสฺส ตํ คหปติ สาวกํ ธาเรมาติ.\csromanspacer{} เอวํ วุตฺเต อุปาลิ คหปติ อุฏฺฐายาสนา เอกํสํ อุตฺตราสงฺคํ กริตฺวา เยน ภควา เตนญฺชลิํ ปณาเมตฺวา นิคณฺฐํ นาฏปุตฺตํ เอตทโวจ “เตน หิ ภนฺเต สุโณหิ ยสฺสาหํ สาวโก”ติ\csromandash{}}

\csromanhangingbegin
\hangingheaditem{76}{ธีรสฺส วิคตโมหสฺส, ปภินฺนขีลสฺส วิชิตวิชยสฺส.}
\hangingbody{อนีฆสฺส สุสมจิตฺตสฺส, วุทฺธสีลสฺส สาธุปญฺญสฺส.}
\hangingbody{เวสมนฺตรสฺส1 วิมลสฺส, ภควโต ตสฺส สาวโกหมสฺมิ.}
\hangingbody{อกถํกถิสฺส ตุสิตสฺส, วนฺตโลกามิสสฺส มุทิตสฺส.}
\hangingbody{กตสมณสฺส มนุชสฺส, อนฺติมสารีรสฺส นรสฺส.}
\hangingbody{อโนปมสฺส วิรชสฺส, ภควโต ตสฺส สาวโกหมสฺมิ.}
\hangingbody{อสํสยสฺส กุสลสฺส, เวนยิกสฺส สารถิวรสฺส.}
\hangingbody{อนุตฺตรสฺส รุจิรธมฺมสฺส, นิกฺกงฺขสฺส ปภาสกสฺส2.}
\hangingbody{มานจฺฉิทสฺส วีรสฺส, ภควโต ตสฺส สาวโกหมสฺมิ.}
\hangingbody{นิสภสฺส อปฺปเมยฺยสฺส, คมฺภีรสฺส โมนปตฺตสฺส.}
\hangingbody{เขมงฺกรสฺส เวทสฺส, ธมฺมฏฺฐสฺส สํวุตตฺตสฺส.}
\hangingbody{สงฺคาติคสฺส มุตฺตสฺส, ภควโต ตสฺส สาวโกหมสฺมิ.}
\csromanhangingend

\csromanheader{2}{6. อุปาลิสุตฺต (56)}
\csromangathameasure{นาคสฺส ปนฺตเสนสฺส, ขีณสํโยชนสฺส มุตฺตสฺส. \\ ปฏิมนฺตกสฺส โธนสฺส, ปนฺนธชสฺส วีตราคสฺส. \\ ทนฺตสฺส นิปฺปปญฺจสฺส, ภควโต ตสฺส สาวโกหมสฺมิ. \\ อิสิสตฺตมสฺส อกุหสฺส, เตวิชฺชสฺส พฺรหฺมปตฺตสฺส. \\ นฺหาตกสฺส ปทกสฺส, ปสฺสทฺธสฺส วิทิตเวทสฺส. \\ ปุรินฺททสฺส อ, ภควโต ตสฺส สาวโกหมสฺมิ. \\ อริยสฺส ภาวิตตฺตสฺส, ปตฺติปตฺตสฺส เวยฺยากรณสฺส. \\ สติมโต วิปสฺสิสฺส, อนภินตสฺส โน อปนตสฺส. \\ อเนชสฺส วสิปฺปตฺตสฺส, ภควโต ตสฺส สาวโกหสฺมิ. \\ สมุคฺคตสฺส ฌายิสฺส, อนนุคตนฺตรสฺส สุทฺธสฺส. \\ อสิตสฺส หิตสฺส, ปวิวิตฺตสฺส อคฺคปฺปตฺตสฺส. \\ ติณฺณสฺส ตารยนฺตสฺส, ภควโต ตสฺส สาวโกหมสฺมิ. \\ สนฺตสฺส ภูริปญฺญสฺส, มหาปญฺญสฺส วีตโลภสฺส. \\ ตถาคตสฺส สุคตสฺส, อปฺปฏิปุคฺคลสฺส อสมสฺส. \\ วิสารทสฺส นิปุณสฺส, ภควโต ตสฺส สาวโกมหสฺมิ. \\ ตณฺหจฺฉิทสฺส พุทฺธสฺส, วีตธูมสฺส อนุปลิตฺตสฺส. \\ อาหุเนยฺยสฺส ยกฺขสฺส, อุตฺตมปุคฺคลสฺส อตุลสฺส. \\ มหโต ยสคฺคปตฺตสฺส, ภควโต ตสฺส สาวโกหมสฺมีติ.}
\csromangathagroup{\csromangathastanza{\csromanfolio{49}นาคสฺส ปนฺตเสนสฺส, ขีณสํโยชนสฺส มุตฺตสฺส. \\ ปฏิมนฺตกสฺส\footnote{ปฏิมนฺตสฺส (ก)} โธนสฺส, ปนฺนธชสฺส วีตราคสฺส.}\csromangathastanza{ทนฺตสฺส นิปฺปปญฺจสฺส, ภควโต ตสฺส สาวโกหมสฺมิ. \\ อิสิสตฺตมสฺส อกุหสฺส, เตวิชฺชสฺส พฺรหฺมปตฺตสฺส.}\csromangathastanza{นฺหาตกสฺส\footnote{นหาตกสฺส (สี, สฺยา, อิ)} ปทกสฺส, ปสฺสทฺธสฺส วิทิตเวทสฺส. \\ ปุรินฺททสฺส อ, ภควโต ตสฺส สาวโกหมสฺมิ.}\csromangathastanza{อริยสฺส ภาวิตตฺตสฺส, ปตฺติปตฺตสฺส เวยฺยากรณสฺส. \\ สติมโต วิปสฺสิสฺส, อนภินตสฺส โน อปนตสฺส.}\csromangathastanza{อเนชสฺส วสิปฺปตฺตสฺส, ภควโต ตสฺส สาวโกหสฺมิ. \\ สมุคฺคตสฺส\footnote{สมฺมคฺคตสฺส (สี, สฺยา, อิ)} ฌายิสฺส, อนนุคตนฺตรสฺส สุทฺธสฺส.}\csromangathastanza{อสิตสฺส หิตสฺส\footnote{อปฺปหีนสฺส (สี, อิ), อปฺปภีตสฺส(สฺยา)}, ปวิวิตฺตสฺส อคฺคปฺปตฺตสฺส. \\ ติณฺณสฺส ตารยนฺตสฺส, ภควโต ตสฺส สาวโกหมสฺมิ.}\csromangathastanza{สนฺตสฺส ภูริปญฺญสฺส, มหาปญฺญสฺส วีตโลภสฺส. \\ ตถาคตสฺส สุคตสฺส, อปฺปฏิปุคฺคลสฺส อสมสฺส.}\csromangathastanza{วิสารทสฺส นิปุณสฺส, ภควโต ตสฺส สาวโกมหสฺมิ. \\ ตณฺหจฺฉิทสฺส พุทฺธสฺส, วีตธูมสฺส อนุปลิตฺตสฺส.}\csromangathastanza{อาหุเนยฺยสฺส ยกฺขสฺส, อุตฺตมปุคฺคลสฺส อตุลสฺส. \\ มหโต ยสคฺคปตฺตสฺส, ภควโต ตสฺส สาวโกหมสฺมีติ.}}

\proseitem{77}{กทา สญฺญูฬฺหา ปน เต คหปติ อิเม สมณสฺส โคตมสฺส วณฺณาติ.\csromanspacer{} เสยฺยถาปิ ภนฺเต นานาปุปฺผานํ มหาปุปฺผราสิ, ตเมนํ ทกฺโข มาลากาโร วา มาลาการนฺเตวาสี วา วิจิตฺตํ มาลํ คนฺเถยฺย, เอวเมว โข ภนฺเต โส ภควา อเนกวณฺโณ อเนกสตวณฺโณ, โก หิ ภนฺเต วณฺณารหสฺส วณฺณํ น กริสฺสตีติ.\csromanspacer{} อถ โข นิคณฺฐสฺส นาฏปุตฺตสฺส ภควโต สกฺการํ อสหมานสฺส ตตฺเถว อุณฺหํ โลหิตํ มุขโต อุคฺคจฺฉีติ\footnote{อุคฺคญฺฉิ (สี, สฺยา, อิ)}.}

\nitthitam{อุปาลิสุตฺตํ นิฏฺฐิตํ ฉฏฺฐํ.}
\csromanmatikaanchor{mtk.8}
\csromantocmarkref{subsection}{7. กุกฺกุรวติกสุตฺต}{mtk.8}
\bookmark[dest={mtk.8},level=2]{7. กุกฺกุรวติกสุตฺต}
\csromanheader{2}{7. กุกฺกุรวติกสุตฺต}
\proseitem{78}{\csromanfolio{50}เอวํ เม สุตํ\csromandash{}เอกํ สมยํ ภควา โกลิเยสุ วิหรติ หลิทฺทวสนํ นาม โกลิยานํ นิคโม.\csromanspacer{} อถ โข ปุณฺโณ จ โกลิยปุตฺโต โควติโก อเจโล จ เสนิโย กุกฺกุรวติโก เยน ภควา เตนุปสงฺกมิํสุ, อุปสงฺกมิตฺวา ปุณฺโณ โกลิยปุตฺโต โควติโก ภควนฺตํ อภิวาเทตฺวา เอกมนฺตํ นิสีทิ.\csromanspacer{} อเจโล ปน เสนิโย กุกฺกุรวติโก ภควตา สทฺธิํ สมฺโมทิ, สมฺโมทนียํ กถํ สารณียํ วีติสาเรตฺวา กุกฺกุโรว ปลิกุชฺชิตฺวา\footnote{ปลิกุณฺฐิตฺวา (สฺยา, กํ), ปลิคุณฺฐิตฺวา (ก)} เอกมนฺตํ นิสีทิ.\csromanspacer{} เอกมนฺตํ นิสินฺโน โข ปุณฺโณ โกลิยปุตฺโต โควติโก ภควนฺตํ เอตทโวจ “อยํ ภนฺเต อเจโล เสนิโย กุกฺกุรวติโก ทุกฺกรการโก ฉมานิกฺขิตฺตํ โภชนํ ภุญฺชติ, ตสฺส ตํ กุกฺกุรวตํ ทีฆรตฺตํ สมตฺตํ สมาทินฺนํ, ตสฺส กา คติ, โก อภิสมฺปราโย”ติ.\csromanspacer{} อลํ ปุณฺณ, ติฏฺฐเตตํ, มา มํ เอตํ ปุจฺฉีติ.\csromanspacer{} ทุติยมฺปิ โข ปุณฺโณ โกลิยปุตฺโต โควติโก ฯเปฯ.\csromanspacer{} ตติยมฺปิ โข ปุณฺโณ โกลิยปุตฺโต โควติโก ภควนฺตํ เอตทโวจ “อยํ ภนฺเต อเจโล เสนิโย กุกฺกุรวติโก ทุกฺกรการโก ฉมานิกฺขิตฺตํ โภชนํ ภุญฺชติ, ตสฺส ตํ กุกฺกุรวตํ ทีฆรตฺตํ สมตฺตํ สมาทินฺนํ, ตสฺส กา คติ, โก อภิสมฺปราโย”ติ.}

\proseitem{79}{อทฺธา โข เต อหํ ปุณฺณ น ลภามิ, “อลํ ปุณฺณ ติฏฺฐเตตํ, มา มํ เอตํ ปุจฺฉี”ติ, อปิ จ ตฺยาหํ พฺยากริสฺสามิ.\csromanspacer{} อิธ ปุณฺณ เอกจฺโจ กุกฺกุรวตํ ภาเวติ ปริปุณฺณํ อพฺโพกิณฺณํ, กุกฺกุรสีลํ ภาเวติ ปริปุณฺณํ อพฺโพกิณฺณํ, กุกฺกุรจิตฺตํ ภาเวติ ปริปุณฺณํ อพฺโพกิณฺณํ, กุกฺกุรากปฺปํ ภาเวติ ปริปุณฺณํ อพฺโพกิณฺณํ.\csromanspacer{} โส กุกฺกุรวตํ ภาเวตฺวา ปริปุณฺณํ อพฺโพกิณฺณํ กุกฺกุรสีลํ ภเวตฺวา ปริปุณฺณํ อพฺโพกิณฺณํ กุกฺกุรจิตฺตํ ภาเวตฺวา ปริปุณฺณํ อพฺโพกิณฺณํ กุกฺกุรากปฺปํ ภาเวตฺวา ปริปุณฺณํ อพฺโพกิณฺณํ กายสฺส เภทา ปรํ มรณา กุกฺกุรานํ สหพฺยตํ อุปปชฺชติ.\csromanspacer{} สเจ โข ปนสฺส เอวํทิฏฺฐิ โหติ “อิมินาหํ สีเลน วา วเตน วา ตเปน วา พฺรหฺมจริเยน วา เทโว วา ภวิสฺสามิ, เทวญฺญตโร วา”ติ, สาสฺส\footnote{สายํ (ก)} โหติ มิจฺฉาทิฏฺฐิ.\csromanspacer{} มิจฺฉาทิฏฺฐิสฺส\footnote{มิจฺฉาทิฏฺฐิกสฺส (สี)} โข อหํ ปุณฺณ ทฺวินฺนํ คตีนํ อญฺญตรํ คติํ วทามิ นิรยํ วา ติรจฺฉานโยนิํ วา.\csromanspacer{} อิติ โข ปุณฺณ}

\vspace{\tipitakaparskip}
\csromancenter{กุกฺกรวติกสุตฺต (57) สมฺปชฺชมานํ กุกฺกุรวตํ กุกฺกุรานํ สหพฺยตํ อุปเนติ วิปชฺชมานํ นิรยนฺติ.\csromanspacer{} เอวํ วุตฺเต อเจโล เสนิโย กุกฺกุรวติโก ปโรทิ, อสฺสูนิ ปวตฺเตสิ.}
\prose{\csromanfolio{51}อถ โข ภควา ปุณฺณํ โกลิยปุตฺตํ โควติกํ เอตทโวจ “เอตํ โข เต อหํ ปุณฺณ นาลตฺถํ, ‘อลํ ปุณฺณ ติฏฺฐเตตํ, มา มํ เอตํ ปุจฺฉี’ติ”.\csromanspacer{} นาหํ ภนฺเต เอตํ โรทามิ, ยํ มํ ภควา เอวมาห “อปิ จ เม อิทํ ภนฺเต กุกฺกุรวตํ ทีฆรตฺตํ สมตฺตํ สมาทินฺนํ.\csromanspacer{} อยํ ภนฺเต ปุณฺโณ โกลิยปุตฺโต โควติโก, ตสฺส ตํ โควตํ ทีฆรตฺตํ สมตฺตํ สมาทินฺนํ, ตสฺส กา คติ, โก อภิสมฺปราโย”ติ.\csromanspacer{} อลํ เสนิย ติฏฺฐเตตํ, มา มํ เอตํ ปุจฺฉีติ.\csromanspacer{} ทุติยมฺปิ โข อเจโล เสนิโย ฯเปฯ.\csromanspacer{} ตติยมฺปิ โข อเจโล เสนิโย กุกฺกุรวติโก ภควนฺตํ เอตทโวจ “อยํ ภนฺเต ปุณฺโณ โกลิยปุตฺโต โควติโก, ตสฺส ตํ โควตํ ทีฆรตฺตํ สมตฺตํ สมาทินฺนํ, ตสฺส กา คติ, โก อภิสมฺปราโย”ติ.}

\proseitem{80}{อทฺธา โข เต อหํ เสนิย น ลภามิ, “อลํ เสนิย ติฏฺฐเตตํ, มา มํ เอตํ ปุจฺฉี”ติ, อปิ จ ตฺยาหํ พฺยากริสฺสามิ.\csromanspacer{} อิธ เสนิย เอกจฺโจ โควตํ ภาเวติ ปริปุณฺณํ อพฺโพกิณฺณํ, โคสีลํ ภาเวติ ปริปุณฺณํ อพฺโพกิณฺณํ, โคจิตฺตํ ภาเวติ ปริปุณฺณํ อพฺโพกิณฺณํ, ควากปฺปํ\footnote{คฺวากปฺปํ (ก)} ภาเวติ ปริปุณฺณํ อพฺโพกิณฺณํ.\csromanspacer{} โส โควตํ ภาเวตฺวา ปริปุณฺณํ อพฺโพกิณฺณํ โคสีลํ ภาเวตฺวา ปริปุณฺณํ อพฺโพกิณฺณํ โคจิตฺตํ ภาเวตฺวา ปริปุณฺณํ อพฺโพกิณฺณํ ควากปฺปํ ภาเวตฺวา ปริปุณฺณํ อพฺโพกิณฺณํ กายสฺส เภทา ปรํ มรณา คุนฺนํ สหพฺยตํ อุปปชฺชติ.\csromanspacer{} สเจ โข ปนสฺส เอวํ ทิฏฺฐิ โหติ “อิมินาหํ สีเลน วา วเตน วา ตเปน วา พฺรหฺมจริเยน วา เทโว วา ภวิสฺสามิ เทวญฺญตโร วา”ติ, สาสฺส โหติ มิจฺฉาทิฏฺฐิ.\csromanspacer{} มิจฺฉทิฏฺฐิสฺส โข อหํ เสนิย ทฺวินฺนํ คตีนํ อญฺญตรํ คติํ วทามิ นิรยํ วา ติรจฺฉานโยนิํ วา.\csromanspacer{} อิติ โข เสนิย สมฺปชฺชมานํ โควตํ คุนฺนํ สหพฺยตํ อุปเนติ วิปชฺชมานํ นิรยนฺติ.\csromanspacer{} เอวํ วุตฺเต ปุณฺโณ โกลิยปุตฺโต โควติโก ปโรทิ, อสฺสูนิ ปวตฺเตสิ.}

\prose{อถ โข ภควา อเจลํ เสนิยํ กุกฺกุรวติกํ เอตทโวจ “เอตํ โข เต อหํ เสนิย นาลตฺถํ, อลํ เสนิย ติฏฺฐเตลํ, มา \csromanfolio{52}มํ เอตํ ปุจฺฉี”ติ.\csromanspacer{} นาหํ ภนฺเต เอตํ โรทามิ, ยํ มํ ภควา เอวมาห, อปิ จ เม อิทํ ภนฺเต โควตํ ทีฆรตฺตํ สมตฺตํ สมาทินฺนํ.\csromanspacer{} เอวํ ปสนฺโน อหํ ภนฺเต ภควติ, ปโหติ ภควา ตถา ธมฺมํ เทเสตุํ, ยถา อหํ เจวิมํ โควตํ ปชเหยฺยํ, อยํ เจว อเจโล เสนิโย กุกฺกุรวติโก ตํ กุกฺกุรวตํ ปชเหยฺยาติ.\csromanspacer{} เตน หิ ปุณฺณ สุณาหิ สาธุกํ มนสิ กโรหิ ภาสิสฺสามีติ. “เอวํ ภนฺเต”ติ โข ปุณฺโณ โกลิยปุตฺโต โควติโก ภควโต ปจฺจสฺโสสิ.\csromanspacer{} ภควา เอตทโวจ\csromandash{}}

\proseitem{81}{จตฺตาริมานิ ปุณฺณ กมฺมานิ มยา สยํ อภิญฺญา สจฺฉิกตฺวา ปเวทิตานิ.\csromanspacer{} กตมานิ จตฺตาริ, อตฺถิ ปุณฺณ กมฺมํ กณฺหํ กณฺหวิปากํ, อตฺถิ ปุณฺณ กมฺมํ สุกฺกํ สุกฺกวิปากํ, อตฺถิ ปุณฺณ กมฺมํ กณฺหสุกฺกํ กณฺหสุกฺกวิปากํ, อตฺถิ ปุณฺณ กมฺมํ อกณฺหํ อสุกฺกํ อกณฺห-อสุกฺกวิปากํ กมฺมกฺขยาย สํวตฺตติ.}

\prose{กตมญฺจ ปุณฺณ กมฺมํ กณฺหํ กณฺหวิปากํ.\csromanspacer{} อิธ ปุณฺณ เอกจฺโจ สพฺยาพชฺฌํ\footnote{สพฺยาปชฺฌํ (สี, สฺยา, กํ)} กายสงฺขารํ อภิสงฺขโรติ, สพฺยาพชฺฌํ วจีสงฺขารํ อภิสงฺขโรติ, สพฺยาพชฺฌํ มโนสงฺขารํ อภิสงฺขโรติ.\csromanspacer{} โส สพฺยาพชฺฌํ กายสงฺขารํ อภิสงฺขริตฺวา สพฺยาพชฺฌํ วจีสงฺขารํ อภิสงฺขริตฺวา สพฺยาพชฺฌํ มโนสงฺขารํ อภิสงฺขริตฺวา สพฺยาพชฺฌํ โลกํ อุปปชฺชติ.\csromanspacer{} ตเมนํ สพฺยาพชฺฌํ โลกํ อุปปนฺนํ สมานํ สพฺยาพชฺฌา ผสฺสา ผุสนฺติ, โส สพฺยาพชฺเฌหิ ผสฺเสหิ ผุฏฺโฐ สมาโน สพฺยาพชฺฌํ เวทนํ เวทติ เอกนฺตทุกฺขํ, เสยฺยถาปิ สตฺตา เนรยิกา.\csromanspacer{} อิติ โข ปุณฺณ ภูตา ภูตสฺส อุปปตฺติ โหติ.\csromanspacer{} ยํ กโรติ เตน อุปปชฺชติ, อุปปนฺนเมนํ ผสฺสา ผุสนฺติ.\csromanspacer{} เอวํปาหํ ปุณฺณ ‘กมฺมทายาทา สตฺตา’ติ วทามิ.\csromanspacer{} อิทํ วุจฺจติ ปุณฺณ กมฺมํ กณฺหํ ปณฺหวิปากํ. (1)}

\prose{กตมญฺจ ปุณฺณ กมฺมํ สุกฺกํ สุกฺกวิปากํ.\csromanspacer{} อิธ ปุณฺณ เอกจฺโจ อพฺยาพชฺฌํ กายสงฺขารํ อภิสงฺขโรติ, อพฺยาพชฺฌํ วจีสงฺขารํ อภิสงฺขโรติ, อพฺยาพชฺฌํ มโนสงฺขารํ อภิสงฺขโรติ.\csromanspacer{} โส อพฺยาพชฺฌํ กายสงฺขารํ อภิสงฺขริตฺวา อพฺยาพชฺฌํ วจีสงฺขารํ อภิสงฺขริตฺวา อพฺยาพชฺฌํ มโนสงฺขารํ อภิสงฺขริตฺวา อพฺยาพชฺฌํ โลกํ อุปปชฺชติ.\csromanspacer{} ตเมนํ อพฺยาพชฺฌํ โลกํ}

\vspace{\tipitakaparskip}
\csromancenter{กุกฺกรวติกสุตฺต (57) อุปปนฺนํ สมานํ อพฺยาพชฺฌา ผสฺสา ผุสนฺติ, โส อพฺยาพชฺเฌหิ ผสฺเสหิ ผุฏฺโฐ สมาโน อพฺยาพชฺฌํ เวทนํ เวเทติ เอกนฺตสุขํ, เสยฺยถาปิ เทวา สุภกิณฺหา.\csromanspacer{} อิติ โข ปุณฺณ ภูตา ภูตสฺส อุปปตฺติ โหติ.\csromanspacer{} ยํ กโรติ เตน อุปปชฺชติ, อุปปนฺนเมนํ ผสฺสา ผุสนฺติ.\csromanspacer{} เอวํปาหํ ปุณฺณ ‘กมฺมทายาทา สตฺตา’ติ วทามิ.\csromanspacer{} อิทํ วุจฺจติ ปุณฺณ กมฺมํ สุกฺกํ สุกฺกวิปากํ. (2)}
\prose{\csromanfolio{53}กตมญฺจ ปุณฺณ กมฺมํ กณฺหสุกฺกํ กณฺหสุกฺกวิปากํ.\csromanspacer{} อิธ ปุณฺณ เอกจฺโจ สพฺยาพชฺชมฺปิ อพฺยาพชฺชมฺปิ กายสงฺขารํ อภิสงฺขโรติ, สพฺยาพชฺฌมฺปิ อพฺยาพชฺฌมฺปิ วจีสงฺขารํ อภิสงฺขโรติ, สพฺยาพชฺฌมฺปิ อพฺยาพชฺฌมฺปิ มโนสงฺขารํ อภิสงฺขโรติ.\csromanspacer{} โส สพฺยาพชฺฌมฺปิ อพฺยาพชฺฌมฺปิ กายสงฺขารํ อภิสงฺขริตฺวา สพฺยาพชฺฌมฺปิ อพฺยาพชฺฌมฺปิ วจีสงฺขารํ อภิสงฺขริตฺวา สพฺยาพชฺฌมฺปิ อพฺยาพชฺฌมฺปิ มโนสงฺขารํ อภิสงฺขริตฺวา สพฺยาพชฺฌมฺปิ อพฺยาพชฺฌมฺปิ โลกํ อุปปชฺชติ.\csromanspacer{} ตเมนํ สพฺยาพชฺฌมฺปิ อพฺยาพชฺฌมฺปิ โลกํ อุปปนฺนํ สมานํ สพฺยาพชฺฌาปิ อพฺยาพชฺฌาปิ ผสฺสา ผุสนฺติ, โส สพฺยาพชฺเฌหิปิ อพฺยาพชฺเฌหิปิ ผสฺเสหิ ผุฏฺโฐ สมาโน สพฺยาพชฺฌมฺปิ อพฺยาพชฺฌมฺปิ เวทนํ เวเทติ โวกิณฺณสุขทุกฺขํ, เสยฺยถาปิ มนุสฺสา เอกจฺเจ จ เทวา เอกจฺเจ จ วินิปาติกา.\csromanspacer{} อิติ โข ปุณฺณ ภูตา ภูตสฺส อุปปตฺติ โหติ.\csromanspacer{} ยํ กโรติ เตน อุปปชฺชติ, อุปปนฺนเมนํ ผสฺสา ผุสนฺติ.\csromanspacer{} เอวํปาหํ ปุณฺณ ‘กมฺมทายาทา สตฺตา’ติ วทามิ.\csromanspacer{} อิทํ วุจฺจติ ปุณฺณ กมฺมํ กณฺหสุกฺกํ กณฺหสุกฺกวิปากํ. (3)}

\prose{กตมญฺจ ปุณฺณ กมฺมํ อกณฺหํ อสุกฺกํ อกณฺห-อสุกฺกวิปากํ กมฺมกฺขยาย สํวตฺตติ.\csromanspacer{} ตตฺร ปุณฺณ ยมิทํ กมฺมํ กณฺหํ กณฺหวิปากํ ตสฺส ปหานาย ยา เจตนา, ยมิทํ\footnote{ยมฺปิทํ (สี, อิ)} กมฺมํ สุกฺกํ สุกฺกวิปากํ ตสฺส ปหานาย ยา เจตนา, ยมิทํ1 กมฺมํ กณฺหสุกฺกํ กณฺหสุกฺกวิปากํ ตสฺส ปหานาย ยา เจตนา.\csromanspacer{} อิทํ วุจฺจติ ปุณฺณ กมฺมํ อกณฺหํ อสุกฺกํ อกณฺห-อสุกฺกวิปากํ กมฺมกฺขยาย สํวตฺตตีติ.\csromanspacer{} อิมานิ โข ปุณฺณ จตฺตาริ กมฺมานิ มยา สยํ อภิญฺญา สจฺฉิกตฺวา ปเวทิตานีติ. (4)}

\csromanmatikaanchor{mtk.9}
\csromantocmarkref{subsection}{8. อภยราชกุมารสุตฺต}{mtk.9}
\bookmark[dest={mtk.9},level=2]{8. อภยราชกุมารสุตฺต}
\proseitem{82}{เอวํ วุตฺเต ปุณฺโณ โกลิยปุตฺโต โควติโก ภควนฺตํ เอตทโวจ “อภิกฺกนฺตํ ภนฺเต, อภิกฺกนฺตํ ภนฺเต, เสยฺยถาปิ ภนฺเต ฯเปฯ อุปาสกํ มํ ภควา ธาเรตุ อชฺชตคฺเค ปาณุเปตํ สรณํ คตนฺ”ติ. \csromanfolio{54}อเจโล ปน เสนิโย กุกฺกุรวติโก ภควนฺตํ เอตทโวจ “อภิกฺกนฺตํ ภนฺเต, อภิกฺกนฺตํ ภนฺเต, เสยฺยถาปิ ภนฺเต ฯเปฯ ปกาสิโต, เอสาหํ ภนฺเต ภควนฺตํ สรณํ คจฺฉามิ ธมฺมญฺจ ภิกฺขุสํฆญฺจ, ลเภยฺยาหํ ภนฺเต ภควโต สนฺติเก ปพฺพชฺชํ, ลเภยฺยํ อุปสมฺปทนฺ”ติ.\csromanspacer{} โย โข เสนิย อญฺญติตฺถิยปุพฺโพ อิมสฺมิํ ธมฺมวินเย อากงฺขติ ปพฺพชฺชํ อากงฺขติ อุปสมฺปทํ, โส จตฺตาโร มาเส ปริวสติ, จตุนฺนํ มาสานํ อจฺเจเยน อารทฺธจิตฺตา ภิกฺขู ปพฺพาเชนฺติ อุปสมฺปาเทนฺติ ภิกฺขุภาวาย, อปิ จ เมตฺถ ปุคฺคลเวมตฺตตา วิทิตาติ.}

\prose{สเจ ภนฺเต อญฺญติตฺถิยปุพฺพา อิมสฺมิํ ธมฺมวินเย อากงฺขนฺตา ปพฺพชฺชํ อากงฺขนฺตา อุปสมฺปทํ, เต จตฺตาโร มาเส ปริวสนฺติ, จตุนฺนํ มาสานํ อจฺจเยน อารทฺธจิตฺตา ภิกฺขู ปพฺพาเชนฺติ อุปสมฺปาเทนฺติ ภิกฺขุภาวาย, อหํ จตฺตาริ วสฺสานิ ปริวสิสฺสามิ จตุนฺนํ วสฺสานํ อจฺจเยน อารทฺธจิตฺตา ภิกฺขู ปพฺพาเชนฺตุ อุปสมฺปาเทนฺตุ ภิกฺขุ ภาวายาติ.\csromanspacer{} อลตฺถ โข อเจโล เสนิโย กุกฺกุรวติโก ภควโต สนฺติเก ปพฺพชฺชํ, อลตฺถ อุปสมฺปทํ, อจิรูปสมฺปนฺโน โข ปนาÿอสฺมา เสนิโย เอโก วูปกฏฺโฐ อปฺปมตฺโต อาตาปี ปหิตตฺโต วิหรนฺโต นจิรสฺเสว, ยสฺสตฺถาย กุลปุตฺตา สมฺมเทว อคารสฺมา อนคาริยํ ปพฺพชนฺติ.\csromanspacer{} ตทนุตฺตรํ พฺรหฺมจริยปริโยสานํ ทิฏฺเฐว ธมฺเม สยํ อภิญฺญา สจฺฉิกตฺวา อุปสมฺปชฺช วิหาสิ “ขีณา ชาติ วุสิตํ พฺรหฺมจริยํ กตํ กรณียํ นาปรํ อิตฺถตฺตายา”ติ อพฺภญฺญาสิ.\csromanspacer{} อญฺญตโร โข ปนายสฺมา เสนิโย อรหตํ อโหสีติ.}

\nitthitamruled{กุกฺกุรวติกสุตฺตํ นิฏฺฐิตํ สตฺตมํ.}
\csromanheader{2}{8. อภยราชากุมารสุตฺต}
\proseitem{83}{เอวํ เม สุตํ\csromandash{}เอกํ สมยํ ภควา ราชคเห วิหรติ เวฬุวเน กลนฺทกนิวาเป.\csromanspacer{} อถ โข อภโย ราชกุมาโร เยน นิคณฺโฐ นาฏปุตฺโต เตนุปสงฺกมิ, อุปสงฺกมิตฺวา นิคณฺฐํ นาฏปุตฺตํ อภิวาเทตฺวา เอกมนฺตํ นิสีทิ.\csromanspacer{} เอกมนฺตํ นิสินฺนํ โข อภยํ ราชกุมารํ นิคณฺโฐ นาฏปุตฺโต เอตทโวจ “เอหิ ตฺวํ ราชกุมาร สมณสฺส โคตมสฺส}

\vspace{\tipitakaparskip}
\csromancenter{อภยราชกุมารสุตฺต (58) วาทํ อาโรเปหิ, เอวํ เต กลฺยาโณ กิตฺติสทฺโท อพฺภุคฺคจฺฉิสฺสติ ‘อภเยน ราชมาเรน สมณสฺส โคตมสฺส เอวํ มหิทฺธิกสฺส เอวํ มหานุภาวสฺส วาโท อาโรปิโต’ติ”.\csromanspacer{} ยถา กถํ ปนาหํ ภนฺเต สมณสฺส โคตมสฺส เอวํ มหิทฺธิกสฺส เอวํ มหานุภาวสฺส วาทํ อาโรเปสฺสามีติ.\csromanspacer{} เอหิ ตฺวํ ราชกุมาร เยน สมโณ โคตโม เตนุปสงฺกม, อุปสงฺกมิตฺวา สมณํ โคตมํ เอวํ วเทหิ “ภาเสยฺย นุ โข ภนฺเต ตถาวโต ตํ วาจํ, ยา สา วาจา ปเรสํ อปฺปิยา อมนาปา”ติ.\csromanspacer{} สเจ เต สมโณ โคตโม เอวํ ปุฏฺโฐ เอวํ พฺยากโรติ “ภาเสยฺย ราชกุมาร ตถาคโต ตํ วาจํ, ยา สา วาจา ปเรสํ อปฺปิยา อมนาปา”ติ, ตเมนํ ตฺวํ เอวํ วเทยฺยาสิ “อถ กิญฺจรหิ เต ภนฺเต ปุถุชฺชเนน นานากรณํ.\csromanspacer{} ปุถุชฺชโนปิ หิ ตํ วาจํ ภาเสยฺย, ยา สา วาจา ปเรสํ อปฺปิยา อมนาปา”ติ.\csromanspacer{} สเจ ปน เต สมโณ โคตโม เอวํ ปุฏฺโฐ เอวํ พฺยากโรติ “น ราชกุมาร ตถาคโต ตํ วาจํ ภาเสยฺย, ยา สา วาจา ปเรสํ อปฺปิยา อมนาปา”ติ, ตเมนํ ตฺวํ เอวํ วเทยฺยาสิ “อถ กิญฺจรหิ เต ภนฺเต เทวทตฺโต พฺยากโต ‘อาปายิโก เทวทตฺโต เนรยิโก เทวทตฺโต กปฺปฏฺโฐ เทวทตฺโต อเตกิจฺโฉ เทวทตฺโต’ติ, ตาย จ ปน เต วาจาย เทวทตฺโต กุปิโต อโหสิ อนตฺตมโน”ติ.\csromanspacer{} อิมํ โข เต ราชกุมาร สมโณ โคตโม อุภโตโกฏิกํ ปญฺหํ ปุฏฺโฐ สมาโน เนว สกฺขิติ อุคฺคิลิตุํ, น สกฺขิติ โอคิลิตุํ.\csromanspacer{} เสยฺยถาปิ นาม ปุริสสฺส อโยสิงฺฆาฏกํ กณฺเฐ วิลคฺคํ โส เนว สกฺกุเณยฺย อุคฺคิลิตุํ, น สกฺกุเณยฺย โอคิลิตุํ, เอวเมว โข เต ราชกุมาร สมโณ โคตโม อิมํ อุภโตโกฏิกํ ปญฺหํ ปุฏฺโฐ สมาโน เนว สกฺขิติ อุคฺคิลิตุํ, น สกฺขิติ โอคิลิตุนฺติ. “เอวํ ภนฺเต”ติ โข อภโย ราชกุมาโร นิคณฺฐสฺส นาฏปุตฺตสฺส ปฏิสฺสุตฺวา อุฏฺฐายาสนา นิคณฺฐํ นาฏปุตฺตํ อภิวาเทตฺวา ปทกฺขิณํ กตฺวา เยน ภควา เตนุปสงฺกมิ, อุปสงฺกมิตฺวา ภควนฺตํ อภิวาเทตฺวา เอกมนฺตํ นิสีทิ.}
\proseitem{84}{\csromanfolio{55}เอกมนฺตํ นิสินฺนสฺส โข อภยสฺส ราชกุมารสฺส สูริยํ\footnote{สุริยํ (สี, สฺยา, กํ, อิ)} อุลฺโลเกตฺวา เอตทโหสิ “อกาโล โข อชฺช ภควโต วาทํ \csromanfolio{56}อาโรเปตุํ, เสฺว ทานาหํ สเก นิเวสเน ภควโต วาทํ อาโรเปสฺสามี”ติ ภควนฺตํ เอตทโวจ “อธิวาเสตุ เม ภนฺเต ภควา สฺวาตนาย อตฺตจตุตฺโถ ภตฺตนฺ”ติ.\csromanspacer{} อธิวาเสสิ ภควา ตุณฺหีภาเวน.\csromanspacer{} อถ โข อภโย ราชกุมาโร ภควโต อธิวาสนํ วิทิตฺวา อุฏฺฐายาสนา ภควนฺตํ อภิวาเทตฺวา ปทกฺขิณํ กตฺวา ปกฺกามิ.\csromanspacer{} อถ โข ภควา ตสฺสา รตฺติยา อจฺจเยน ปุพฺพณฺหสมยํ นิวาเสตฺวา ปตฺตจีวรมาทาย เยน อภยสฺส ราชกุมารสฺส นิเวสนํ เตนุปสงฺกมิ, อุปสงฺกมิตฺวา ปญฺญตฺเต อาสเน นิสีทิ.\csromanspacer{} อถ โข อภโย ราชกุมาโร ภควนฺตํ ปณีเตน ขาทนีเยน โภชนีเยน สหตฺถา สนฺตปฺเปสิ สมฺปวาเรสิ.\csromanspacer{} อถ โข อภโย ราชกุมาโร ภควนฺตํ ภุตฺตาวิํ โอนีตปตฺตปาณิํ อญฺญตรํ นีจํ อาสนํ คเหตฺวา เอกมนฺตํ นิสีทิ.}

\proseitem{85}{เอกมนฺตํ นิสินฺโน โข อภโย ราชกุมาโร ภควนฺตํ เอตทโวจ “ภาเสยฺย นุ โข ภนฺเต ตถาคโต ตํ วาจํ, ยา สา วาจา ปเรสํ อปฺปิยา อมนาปา”ติ.\csromanspacer{} น เขฺวตฺถ ราชกุมาร เอกํเสนาติ.\csromanspacer{} เอตฺถ ภนฺเต อนสฺสุํ นิคณฺฐาติ.\csromanspacer{} กิํ ปน ตฺวํ ราชกุมาร เอวํ วเทสิ “เอตฺถ ภนฺเต อนสฺสุํ นิคณฺฐา”ติ.\csromanspacer{} อิธาหํ ภนฺเต เยน นิคณฺโฐ นาฏปุตฺโต เตนุปสงฺกมิ, อุปสงฺกมิตฺวา นิคณฺฐํ นาฏปุตฺตํ อภิวาเทตฺวา เอกมนฺตํ นิสีทิํ, เอกมนฺตํ นิสินฺนํ โข มํ ภนฺเต นิคณฺโฐ นาฏปุตฺโต เอตทโวจ\csromandash{}เอหิ ตฺวํ ราชกุมาร สมณสฺส โคตมสฺส วาทํ อาโรเปหิ, เอวํ เต กลฺยาโณ กิตฺติสทฺโท อพฺภุคฺคจฺฉิสฺสติ “อภเยน ราชกุมาเรน สมณสฺส โคตมสฺส เอวํ มหิทฺธิกสฺส เอวํ มหานุภาวสฺส วาโท อาโรปิโต”ติ.\csromanspacer{} เอวํ วุตฺเต อหํ ภนฺเต นิคณฺฐํ นาฏปุตฺตํ เอตทโวจํ “ยถา กถํ ปนาหํ ภนฺเต สมณสฺส โคตมสฺส เอวํ มหิทฺธิกสฺส เอวํ มหานุภาวสฺส วาทํ อาโรเปสฺสามี”ติ.\csromanspacer{} เอหิ ตฺวํ ราชกุมาร เยน สมโณ โคตโม เตนุปสงฺกม, อุปสงฺกมิตฺวา สมณํ โคตมํ เอวํ วเทหิ “ภาเสยฺย นุ โข ภนฺเต ตถาคโต ตํ วาจํ, ยา สา วาจา ปเรสํ อปฺปิยา อมนาปา”ติ.\csromanspacer{} สเจ เต สมโณ โคตโม เอวํ ปุฏฺโฐ เอวํ พฺยากโรติ “ภาเสยฺย ราชกุมาร ตถาคโต ตํ วาจํ, ยา สา วาจา ปเรสํ อปฺปิยา อมนาปา”ติ, ตเมนํ ตฺวํ เอวํ วเทยฺยาสิ “อถ กิญฺจรหิ เต ภนฺเต ปุถุชฺชเนน นานากรณํ, ปุถุชฺชโนปิ หิ ตํ วาจํ}

\vspace{\tipitakaparskip}
\csromancenter{อภยราชกุมารสุตฺต (58) ภาเสยฺย, ยา สา วาจา ปเรสํ อปฺปิยา อมนาปา”ติ.\csromanspacer{} สเจ ปน เต สมโณ โคตโม เอวํ ปุฏฺโฐ เอวํ พฺยากโรติ “น ราชกุมาร ตถาคโต ตํ วาจํ ภาเสยฺย, ยา สา วาจา ปเรสํ อปฺปิยา อมนาปา”ติ, ตเมนํ ตฺวํ เอวํ วเทยฺยาสิ “อถ กิญฺจรหิ เต ภนฺเต เทวทตฺโต พฺยากโต ‘อาปายิโก เทวทตฺโต เนรยิโก เทวทตฺโต กปฺปฏฺโฐ เทวทตฺโต อเตกิจฺโฉ เทวทตฺโต’ติ, ตาย จ ปน เต วาจาย เทวทตฺโต กุปิโต อโหสิ อนตฺตมโน”ติ.\csromanspacer{} อิมํ โข เต ราชกุมาร สมโณ โคตโม อุภโตโกฏิกํ ปญฺหํ ปุฏฺโฐ สมาโน เนว สกฺขิติ อุคฺคิลิตุํ, น สกฺขิติ โอคิลิตุํ.\csromanspacer{} เสยฺยถาปิ นาม ปุริสสฺส อโยสิงฺฆาฏกํ กณฺเฐ วิลคฺคํ, โส เนว สกฺกุเณยฺย อุคฺคิลิตุํ, น สกฺกุเณยฺย โอคิลิตุํ.\csromanspacer{} เอวเมว โข เต ราชกุมาร สมโณ โคตโม อิมํ อุภโตโกฏิกํ ปญฺหํ ปุโฐ สมาโน เนว สกฺขิติ อุคฺคิลิตุํ, น สกฺขิติ โอคิลิตุนฺติ.}
\proseitem{86}{\csromanfolio{57}เตน โข ปน สมเยน ทหโร กุมาโร มนฺโท อุตฺตานเสยฺยโก อภยสฺส ราชกุมารสฺส องฺเก นิสินฺโน โหติ.\csromanspacer{} อถ โข ภควา อภยํ ราชกุมารํ เอตทโวจ “ตํ กิํ มญฺญสิ ราชกุมาร, สจายํ กุมาโร ตุยฺหํ วา ปมาทมนฺวาย ธาติยา วา ปมาทมนฺวาย กฏฺฐํ วา กฐลํ\footnote{กถลํ (ก)} วา มุเข อาหเรยฺย, กินฺติ นํ กเรยฺยาสี”ติ.\csromanspacer{} อาหเรยฺยสฺสาหํ ภนฺเต, สเจ ภนฺเต น สกฺกุเณยฺยํ อาทิเกเนว อาหตฺตุํ\footnote{อาหริตุํ (สฺยา, กํ)} วาเมน หตฺเถน สีสํ ปริคฺคเหตฺวา\footnote{ปคฺคเหตฺวา (สี)} ทกฺขิเณน หตฺเถน วงฺกงฺคุลิํ กริตฺวา สโลหิตํปิ อาหเรยฺยํ.\csromanspacer{} ตํ กิสฺส เหตุ, อตฺถิ เม ภนฺเต กุมาเร อนุกมฺปาติ.\csromanspacer{} เอวเมว โข ราชกุมาร ยํ ตถาคโต วาจํ ชานาติ อภูตํ อตจฺฉํ อนตฺถสํหิตํ, สา จ ปเรสํ อปฺปิยา อมนาปา, น ตํ ตถาคโต วาจํ ภาสติ.\csromanspacer{} ยมฺปิ ตถาคโต วาจํ ชานาติ ภูตํ ตจฺฉํ อนตฺถสํหิตํ, สา จ ปเรสํ อปฺปิยา อมนาปา, ตมฺปิ ตถาคโต วาจํ น ภาสติ.\csromanspacer{} ยญฺจ โข ตถาคโต วาจํ ชานาติ ภูตํ ตจฺฉํ อตฺถสํหิตํ, สา จ ปเรสํ อปฺปิยา อมนาปา, ตตฺร กาลญฺญู ตถาคโต โหติ ตสฺสา วาจาย เวยฺยากรณาย.\csromanspacer{} ยํ ตถาคโต วาจํ ชานาติ อภูตํ \csromanfolio{58}อตจฺฉํ อนตฺถสํหิตํ, สา จ ปเรสํ ปิยา มนาปา, น ตํ ตถาคโต วาจํ ภาสติ.\csromanspacer{} ยมฺปิ ตถาคโต วาจํ ชานาติ ภูตํ ตจฺฉํ อนตฺถสํหิตํ, สา จ ปเรสํ ปิยา มนาปา, ตมฺปิ ตถาคโต วาจํ น ภาสติ.\csromanspacer{} ยญฺจ ตถาคโต วาจํ ชานาติ ภูตํ ตจฺฉํ อตฺถสํหิตํ, สา จ ปเรสํ ปิยา มนาปา, ตตฺร กาลญฺญู ตถาคโต โหติ ตสฺสา วาจาย เวยฺยากรณาย.\csromanspacer{} ตํ กิสฺส เหตุ, อตฺถิ ราชกุมาร ตถาคตสฺส สตฺเตสุ อนุกมฺปาติ.}

\proseitem{87}{เยเม ภนฺเต ขตฺติยปณฺฑิตาปิ พฺราหฺมณปณฺฑิตาปิ คหปติปณฺฑิตาปิ สมณปณฺฑิตาปิ ปญฺหํ อภิสงฺขริตฺวา ตถาคตํ อุปสงฺกมิตฺวา ปุจฺฉนฺติ “ปุพฺเพว นุ โข เอตํ ภนฺเต ภควโต เจตโส ปริวิตกฺกิตํ โหติ, เย มํ อุปสงฺกมิตฺวา เอวํ ปุจฺฉิสฺสนฺติ, เตสาหํ เอวํ ปุฏฺโฐ เอวํ พฺยากริสฺสามี”ติ, อุทาหุ ฐานโสเวตํ ตถาคตํ ปฏิภาตีติ.\csromanspacer{} เตน หิ ราชกุมาร ตญฺเญเวตฺถ ปฏิปุจฺฉิสฺสามิ, ยถา เต ขเมยฺย, ตถา นํ พฺยากเรยฺยาสิ, ตํ กิํ มญฺญสิ ราชกุมาร, กุสโล ตฺวํ รถสฺส องฺคปจฺจงฺคานนฺติ.\csromanspacer{} เอวํ ภนฺเต กุสโล อหํ รถสฺส องฺคปจฺจงฺคานนฺติ, ตํ กิํ มญฺญสิ ราชกุมาร, เย ตํ อุปสงฺกมิตฺวา เอวํ ปุจฺเฉยฺยุํ “กิํ นามิทํ รถสฺส องฺคปจฺจงฺคนฺ”ติ.\csromanspacer{} ปุพฺเพว นุ โข เต เอตํ เจตโส ปริวิตกฺกิตํ อสฺส, เย มํ อุปสงฺกมิตฺวา เอวํ ปุจฺฉิสฺสนฺติ, เตสาหํ เอวํ ปุฏฺโฐ เอวํ พฺยากริสฺสามีติ, อุทาหุ ฐานโสเวตํ ปฏิภาเสยฺยาติ.\csromanspacer{} อหํ หิ ภนฺเต รถิโก สญฺญาโต กุสโล รถสฺส องฺคปจฺจงฺคานํ สพฺพานิ เม รถสฺส องฺคปจฺจงฺคานิ สุวิทิตานิ, ฐานโสเวตํ มํ ปฏิภาเสยฺยาติ.\csromanspacer{} เอวเมว โข ราชกุมาร เย เต ขตฺติยปณฺฑิตาปิ พฺราหฺมณปณฺฑิตาปิ คหปติปณฺฑิตาปิ สมณปณฺฑิตาปิ ปญฺหํ อภิสงฺขริตฺวา ตถาคตํ อุปสงฺกมิตฺวา ปุจฺฉนฺติ.\csromanspacer{} ฐานโสเวตํ ตถาคตํ ปฏิภาติ.\csromanspacer{} ตํ กิสฺส เหตุ, สา หิ ราชกุมาร ตถาคตสฺส ธมฺมธาตุ สุปฺปฏิวิทฺธา, ยสฺสา ธมฺมธาตุยา สุปฺปฏิวิทฺธตฺตา ฐานโสเวตํ ตถาคตํ ปฏิภาตีติ.}

\prose{เอวํ วุตฺเต อภโย ราชกุมาโร ภควนฺตํ เอตทโวจ “อภิกฺกนฺตํ ภนฺเต, อภิกฺกนฺตํ ภนฺเต ฯเปฯ อชฺชตคฺเค ปาณุเปตํ สรณํ คตนฺ”ติ.}

\nitthitam{อภยราชกุมารสุตฺตํ นิฏฺฐิตํ อฏฺฐมํ.}
\csromanmatikaanchor{mtk.10}
\csromantocmarkref{subsection}{9. พหุเวทนียสุตฺต}{mtk.10}
\bookmark[dest={mtk.10},level=2]{9. พหุเวทนียสุตฺต}
\csromanheader{2}{9. พหุเวทนิยสุตฺต (59) \textbf{9. พหุเวทนียสุตฺต}}
\proseitem{88}{\csromanfolio{59}เอวํ เม สุตํ\csromandash{}เอกํ สมยํ ภควา สาวตฺถิยํ วิหรติ เชตวเน อนาถปิณฺฑิกสฺส อาราเม.\csromanspacer{} อถ โข ปญฺจกงฺโค ถปติ เยนายสฺมา อุทายี เตนุปสงฺกมิ, อุปสงฺกมิตฺวา อายสฺมนฺตํ อุทายิํ อภิวาเทตฺวา เอกมนฺตํ นิสีทิ, เอกมนฺตํ นิสินฺโน โข ปญฺจกงฺโค ถปติ อายสฺมนฺตํ อุทายิํ เอตทโวจ “กติ นุ โข ภนฺเต อุทายิ เวทนา วุตฺตา ภควตา”ติ. “ติสฺโส โข ถปติ\footnote{คหปติ (สฺยา, กํ, อิ)} เวทนา วุตฺตา ภควตา สุขา เวทนา ทุกฺขา เวทนา อทุกฺขมสุขา เวทนา.\csromanspacer{} อิมา โข ถปติ ติสฺโส เวทนา วุตฺตา ภควตา”ติ.\csromanspacer{} เอวํ วุตฺเต ปญฺจกงฺโค ถปติ อายสฺมนฺตํ อุทายิํ เอตทโวจ “น โข ภนฺเต อุทายิ ติสฺโส เวทนา วุตฺตา ภควตา, เทฺว เวทนา วุตฺตา ภควตา สุขา เวทนา ทุกฺขา เวทนา.\csromanspacer{} ยายํ ภนฺเต อทุกฺขมสุขา เวทนา, สนฺตสฺมิํ เอสา ปณีเต สุเข วุตฺตา ภควตา”ติ.\csromanspacer{} ทุติยมฺปิ โข อายสฺมา อุทายี ปญฺจกงฺคํ ถปติํ เอตทโวจ “น โข คหปติ เทฺว เวทนา วุตฺตา ภควตา, ติสฺโส เวทนา วุตฺตา ภควตา สุขา เวทนา ทุกฺขา เวทนา อทุกฺขมสุขา เวทนา.\csromanspacer{} อิมา โข ถปติ ติสฺโส เวทนา วุตฺตา ภควตา”ติ.\csromanspacer{} ทุติยมฺปิ โข ปญฺจกงฺโค ถปติ อายสฺมนฺตํ อุทายิํ เอตทโวจ “น โข ภนฺเต อุทายิ ติสฺโส เวทนา วุตฺตา ภควตา, เทฺว เวทนา วุตฺตา ภควตา สุขา เวทนา ทุกฺขา เวทนา.\csromanspacer{} ยายํ ภนฺเต อทุกฺขมสุขา เวทนา, สนฺตสฺมิํ เอสา ปณีเต สุเข วุตฺตา ภควตา”ติ.\csromanspacer{} ตติยมฺปิ โข อายสฺมา อุทายี ปญฺจกงฺคํ ถปติํ เอตทโวจ “น โข ถปติ เทฺว เวทนา วุตฺตา ภควตา, ติสฺโส เวทนา วุตฺตา ภควตา สุขา เวทนา ทุกฺขา เวทนา อทุกฺขมสุขา เวทนา.\csromanspacer{} อิมา โข ถปติ ติสฺโส เวทนา วุตฺตา ภควตา”ติ.\csromanspacer{} ตติยมฺปิ โข ปญฺจกงฺโค ถปติ อายสฺมนฺตํ อุทายิํ เอตทโวจ “น โข ภนฺเต อุทายิ ติสฺโส เวทนา วุตฺตา ภควตา, เทฺว เวทนา วุตฺตา ภควตา สุขา เวทนา ทุกฺขา เวทนา.\csromanspacer{} ยายํ ภนฺเต อทุกฺขมสุขา เวทนา, สนฺตสฺมิํ เอสา ปณีเต สุเข วุตฺตา ภควตา”ติ.\csromanspacer{} เนว โข สกฺขิ อายสฺมา อุทายี ปญฺจกงฺคํ ถปติํ สญฺญาเปตุํ, น ปนาสกฺขิ ปญฺจกงฺโค ถปติ อายสฺมนฺตํ อุทายิํ สญฺญาเปตุํ.}

\proseitem{89}{อสฺโสสิ โข อายสฺมา อานนฺโท อายสฺมโต อุทายิสฺส ปญฺจกงฺเคน ถปตินา สทฺธิํ อิมํ กถาสลฺลาปํ.\csromanspacer{} อถ โข อายสฺมา \csromanfolio{60}อานนฺโท เยน ภควา เตนุปสงฺกมิ, อุปสงฺกมิตฺวา ภควนฺตํ อภิวาเทตฺวา เอกมนฺตํ นิสีทิ, เอกมนฺตํ นิสินฺโน โข อายสฺมา อานนฺโท ยาวตโก อโหสิ อายสฺมโต อุทายิสฺส ปญฺจกงฺเคน ถปตินา สทฺธิํ กถาสลฺลาโป, ตํ สพฺพํ ภควโต อาโรเจสิ.\csromanspacer{} เอวํ วุตฺเต ภควา อายสฺมนฺตํ อานนฺทํ เอตทโวจ “สนฺตญฺเญว โข อานนฺท ปริยายํ ปญฺจกงฺโค ถปติ อุทายิสฺส นาพฺภนุโมทิ, สนฺตญฺเญว จ ปน ปริยายํ อุทายี ปญฺจกงฺคสฺส ถปติสฺส นาพฺภนุโมทิ.\csromanspacer{} เทฺวปานนฺท เวทนา วุตฺตา มยา ปริยาเยน, ติสฺโสปิ เวทนา วุตฺตา มยา ปริยาเยน, ปญฺจปิ เวทนา วุตฺตา มยา ปริยาเยน, ฉปิ เวทนา วุตฺตา มยา ปริยาเยน, อฏฺฐารสปิ เวทนา วุตฺตา มยา ปริยาเยน, ฉตฺติํสปิ เวทนา วุตฺตา มยา ปริยาเยน, อฏฺฐสตํปิ เวทนา วุตฺตา มยา ปริยาเยน.\csromanspacer{} เอวํ ปริยายเทสิโต โข อานนฺท มยา ธมฺโม, เอวํ ปริยายเทสิเต โข อานนฺท มยา ธมฺเม เย อญฺญมญฺญสฺส สุภาสิตํ สุลปิตํ น สมนุชานิสฺสนฺติ น สมนุมญฺญิสฺสนฺติ น สมนุโมทิสฺสนฺติ, เตสเมตํ ปาฏิกงฺขํ, ภณฺฑนชาตา กลหชาตา วิวาทาปนฺนา อญฺญมญฺญํ มุขสตฺตีหิ วิตุทนฺตา วิหริสฺสนฺติ.\csromanspacer{} เอวํ ปริยายเทสิโต โข อานนฺท มยา ธมฺโม, เอวํ ปริยายเทสิเต โข อานนฺท มยา ธมฺเม เย อญฺญมญฺญสฺส สุภาสิตํ สุลปิตํ สมนุชานิสฺสนฺติ สมนุมญฺญิสฺสนฺติ สมนุโมทิสฺสนฺติ, เตสเมตํ ปาฏิกงฺขํ, สมคฺคา สมฺโมทมานา อวิวทมานา ขีโรทกีภูตา อญฺญมญฺญํ ปิยจกฺขูหิ สมฺปสฺสนฺตา วิหริสฺสนฺติ.}

\proseitem{90}{ปญฺจ โข อิเม อานนฺท กามคุณา.\csromanspacer{} กตเม ปญฺจ, จกฺขุวิญฺเญยฺยา รูปา อิฏฺฐา กนฺตา มนาปา ปิยรูปา กามูปสํหิตา รชนียา, โสตวิญฺเญยฺยา สทฺทา ฯเปฯ ฆานวิญฺเญยฺยา คนฺธา ฯเปฯ ชิวฺหาวิญฺเญยฺยา รสา ฯเปฯ กายวิญฺเญยฺยา โผฏฺฐพฺพา อิฏฺฐา กนฺตา มนาปา ปิยรูปา กามูปสํหิตา รชนียา, อิเม โข อานนฺท ปญฺจ กามคุณา.\csromanspacer{} ยํ โข อานนฺท อิเม ปญฺจ กามคุเณ ปฏิจฺจ อุปฺปชฺชติ สุขํ โสมนสฺสํ, อิทํ วุจฺจติ กามสุขํ.}

\prose{โย โข อานนฺท เอวํ วเทยฺย “เอตปรมํ สตฺตา สุขํ โสมนสฺสํ ปฏิสํเวเทนฺตี”ติ, อิทมสฺส นานุชานามิ, ตํ กิสฺส เหตุ, อตฺถานนฺท เอตมฺหา สุขา อญฺญํ สุขํ อภิกฺกนฺตรญฺจ ปณีตตรญฺจ.\csromanspacer{} กตมญฺจานนฺท เอตมฺหา สุขา อญฺญํ สุขํ อภิกฺกนฺตตรญฺจ ปณีตตรญฺจ.\csromanspacer{} อิธานนฺท ภิกฺขุ วิวิจฺเจว กาเมหิ วิวิจฺจ อกุสเลหิ ธมฺเมหิ ฯเปฯ ปฐมํ ฌานํ อุปสมฺปชฺช}

\vspace{\tipitakaparskip}
\csromancenter{พหุเวทนิยสุตฺต (59) วิหรติ.\csromanspacer{} อิทํ โข อานนฺท เอตมฺหา สุขา อญฺญํ สุขํ อภิกฺกนฺตตรญฺจ ปณีตตรญฺจ.}
\prose{\csromanfolio{61}โย โข อานนฺท เอวํ วเทยฺย “เอตปรมํ สตฺตา สุขํ โสมนสฺสํ ปฏิสํเวเทนฺตี”ติ, อิทมสฺส นานุชานามิ, ตํ กิสฺส เหตุ, อตฺถานนฺท เอตมฺหา สุขา อญฺญํ สุขํ อภิกฺกนฺตตรญฺจ ปณีตตรญฺจ.\csromanspacer{} กตมญฺจานนฺท เอตมฺหา สุขา อญฺญํ สุขํ อภิกฺกนฺตตรญฺจ ปณีตตรญฺจ.\csromanspacer{} อิธานนฺท ภิกฺขุ วิตกฺกวิจารานํ วูปสมา ฯเปฯ ทุติยํ ฌานํ อุปสมฺปชฺช วิหรติ.\csromanspacer{} อิทํ โข อานนฺท เอตมฺหา สุขา อญฺญํ สุขํ อภิกฺกนฺตตรญฺจ ปณีตตรญฺจ.}

\prose{โย โข อานนฺท เอวํ วเทยฺย ฯเปฯ.\csromanspacer{} กตมญฺจานนฺท เอตมฺหา สุขา อญฺญํ สุขํ อภิกฺขนฺตตรญฺจ ปณีตตรญฺจ.\csromanspacer{} อิธานนฺท ภิกฺขุ ปีติยา จ วิราคา ฯเปฯ ตติยํ ฌานํ อุปสมฺปชฺช วิหรติ.\csromanspacer{} อิทํ โข อานนฺท เอตมฺหา สุขา อญฺญํ สุขํ อภิกฺกนฺตตรญฺจ ปณีตตรญฺจ.}

\prose{โย โข อานนฺท เอวํ วเทยฺย ฯเปฯ.\csromanspacer{} กตมญฺจานนฺท เอตมฺหา สุขา อญฺญํ สุขํ อภิกฺกนฺตตรญฺจ ปณีตตรญฺจ.\csromanspacer{} อิธานนฺท ภิกฺขุ สุขสฺส จ ปหานา ฯเปฯ จตุตฺถํ ฌานํ อุปสมฺปชฺช วิหรติ.\csromanspacer{} อิทํ โข อานนฺท เอตมฺหา สุขา อญฺญํ สุขํ อภิกฺกนฺตตรญฺจ ปณีตตรญฺจ.}

\prose{โย โข อานนฺท เอวํ วเทยฺย ฯเปฯ.\csromanspacer{} กตมญฺจานนฺท เอตมฺหา สุขา อญฺญํ สุขํ อภิกฺกนฺตตรญฺจ ปณีตตรญฺจ.\csromanspacer{} อิธานนฺท ภิกฺขุ สพฺพโส รูปสญฺญานํ สมติกฺกมา ปฏิฆสญฺญานํ อตฺถงฺคมา นานตฺตสญฺญานํ อมนสิการา “อนนฺโต อากาโส”ติ อากาสานญฺจายตนํ อุปสมฺปชฺช วิหรติ.\csromanspacer{} อิทํ โข อานนฺท เอตมฺหา สุขา อญฺญํ สุขํ อภิกฺกนฺตตรญฺจ ปณีตตรญฺจ.}

\prose{โย โข อานนฺท เอวํ วเทยฺย ฯเปฯ.\csromanspacer{} กตมญฺจานนฺท เอตมฺหา สุขา อญฺญํ สุขํ อภิกฺกนฺตตรญฺจ ปณีตตรญฺจ.\csromanspacer{} อิธานนฺท ภิกฺขุ สพฺพโส อากาสานญฺจายตนํ สมติกฺกมฺม “อนนฺตํ วิญฺญาณนฺ”ติ วิญฺญาณญฺจายตนํ อุปสมฺปชฺช วิหรติ.\csromanspacer{} อิทํ โข อานนฺท เอตมฺหา สุขา อญฺญํ สุขํ อภิกฺกนฺตตรญฺจ ปณีตตรญฺจ.}

\prose{โย โข อานนฺท เอวํ วเทยฺย ฯเปฯ.\csromanspacer{} กตมญฺจานนฺท เอตมฺหา สุขา อญฺญํ สุขํ อภิกฺกนฺตตรญฺจ ปณีตตรญฺจ.\csromanspacer{} อิธานนฺท ภิกฺขุ สพฺพโส วิญฺญาณญฺจายตนํ สมติกฺกมฺม “นตฺถิ กิญฺจี”ติ อากิญฺจญฺญายตนํ อุปสมฺปชฺช วิหรติ.\csromanspacer{} อิทํ โข อานนฺท เอตมฺหา สุขา อญฺญํ สุขํ อภิกฺกนฺตตรญฺจ ปณีตตรญฺจ.}

\prose{\csromanfolio{62}โย โข อานนฺท เอวํ วเทยฺย ฯเปฯ.\csromanspacer{} กตมญฺจานนฺท เอตมฺหา สุขา อญฺญํ สุขํ อภิกฺกนฺตตรญฺจ ปณีตตรญฺจ.\csromanspacer{} อิธานนฺท ภิกฺขุ สพฺพโส อากิญฺจญฺญายตนํ สมติกฺกมฺม เนวสญฺญานาสญฺญายตนํ อุปสมฺปชฺช วิหรติ.\csromanspacer{} อิทํ โข อานนฺท เอตมฺหา สุขา อญฺญํ สุขํ อภิกฺกนฺตตรญฺจ ปณีตตรญฺจ.}

\prose{โย โข อานนฺท เอวํ วเทยฺย “เอตปรมํ สตฺตา สุขํ โสมนสฺสํ ปฏิสํเวเทนฺตี”ติ, อิทมสฺส นานุชานามิ, ตํ กิสฺส เหตุ, อตฺถานนฺท เอตมฺหา สุขา อญฺญํ สุขํ อภิกฺกนฺตตรญฺจ ปณีตตรญฺจ.\csromanspacer{} กตมาญฺจานนฺท เอตมฺหา สุขา อญฺญํ สุขํ อภิกฺกนฺตตรญฺจ ปณีตตรญฺจ.\csromanspacer{} อิธานนฺท ภิกฺขุ สพฺพโส เนวสญฺญานาสญฺญายตนํ สมติกฺกมฺม สญฺญาเวทยิตนิโรธํ อุปสมฺปชฺช วิหรติ.\csromanspacer{} อิทํ โข อานนฺท เอตมฺหา สุขา อญฺญํ สุขํ อภิกฺกนฺตตรญฺจ ปณีตตรญฺจ.}

\proseitem{91}{ฐานํ โข ปเนตํ อานนฺท วิชฺชติ, ยํ อญฺญติตฺถิยา ปริพฺพาชกา เอวํ วเทยฺยุํ “สญฺญาเวทยิตนิโรธํ สมโณ โคตโม อาห, ตญฺจ สุขสฺมิํ ปญฺญเปติ, ตยิทํ กิํ สุ ตยิทํ กถํสู”ติ.\csromanspacer{} เอวํวาทิโน อานนฺท อญฺญติตฺถิยา ปริพฺพาชกา เอวมสฺสุ วจนียา “น โข อาวุโส ภควา สุขํเยว เวทนํ สนฺธาย สุขสฺมิํ ปญฺญเปติ, อปิ จ อาวุโส ยตฺถ ยตฺถ สุขํ อุปลพฺภติ ยหิํ ยหิํ, ตํ ตํ ตถาคโต สุขสฺมิํ ปญฺญเปตี”ติ.}

\prose{อิทมโวจ ภควา.\csromanspacer{} อตฺตมโน อายสฺมา อานนฺโท ภควโต ภาสิตํ อภินนฺทีติ.}

\nitthitamruled{พหุเวทนียสุตฺตํ นิฏฺฐิตํ นวมํ.}
\csromanmatikaanchor{mtk.11}
\csromantocmarkref{subsection}{10. อปณฺณกสุตฺต}{mtk.11}
\bookmark[dest={mtk.11},level=2]{10. อปณฺณกสุตฺต}
\csromanheader{2}{10. อปณฺณกสุตฺต}
\proseitem{92}{เอวํ เม สุตํ\csromandash{}เอกํ สมยํ ภควา โกสเลสุ จาริกํ จรมาโน มหตา ภิกฺขุสํเฆน สทฺธิํ เยน สาลา นาม โกสลานํ พฺราหฺมณคาโม ตทวสริ.\csromanspacer{} อสฺโสสุํ โข สาเลยฺยกา พฺราหฺมณคหปติกา “สมโณ ขลุ โภ โคตโม สกฺยปุตฺโต สกฺยกุลา ปพฺพชิโต โกสเลสุ จาริกํ จรมาโน มหตา ภิกฺขุสํเฆน สทฺธิํ}

\vspace{\tipitakaparskip}
\csromancenter{อปณฺณกสุตฺต (60) สาลํ อนุปฺปตฺโต, ตํ โข ปน ภวนฺตํ โคตมํ เอวํ กลฺยาโณ กิตฺติสทฺโท อพฺภุคฺคโต ‘อิติปิ โส ภควา อรหํ สมฺมาสมฺพุทฺโธ วิชฺชาจรณสมฺปนฺโน สุคโต โลกวิทู อนุตฺตโร ปุริสทมฺมสารถิ สตฺถา เทวมนุสฺสานํ พุทฺโธ ภควา’ติ, โส อิมํ โลกํ สเทวกํ สมารกํ สพฺรหฺมกํ สสฺสมณพฺราหฺมณิํ ปชํ สเทวมนุสฺสํ สยํ อภิญฺญา สจฺฉิกตฺวา ปเวเทติ, โส ธมฺมํ เทเสติ อาทิกลฺยาณํ มชฺเฌกลฺยาณํ ปริโยสานกลฺยาณํ สาตฺถํ สพฺยญฺชนํ เกวลปริปุณฺณํ ปริสุทฺธํ พฺรหฺมจริยํ ปกาเสติ.\csromanspacer{} สาธุ โข ปน ตถารูปานํ อรหตํ ทสฺสนํ โหตี”ติ.\csromanspacer{} อถ โข สาเลยฺยกา พฺราหฺมณคหปติกา เยน ภควา เตนุปสงฺกมิํสุ, อุปสงฺกมิตฺวา อปฺเปกจฺเจ ภควนฺตํ อภิวาเทตฺวา เอกมนฺตํ นิสีทิํสุ, อปฺเปกจฺเจ ภควตา สทฺธิํ สมฺโมทิํสุ, สมฺโมทนียํ กถํ สารณียํ วีติสาเรตฺวา เอกมนฺตํ นิสีทิํสุ, อปฺเปกจฺเจ เยน ภควา เตนญฺชลิํ ปณาเมตฺวา เอกมนฺตํ นิสีทิํสุ, อปฺเปกจฺเจ ภควโต สนฺติเก นามโคตฺตํ สาเวตฺวา เอกมนฺตํ นิสีทิํสุ, อปฺเปกจฺเจ ตุณฺหีภูตา เอกมนฺตํ นิสีทิํสุ.}
\proseitem{93}{\csromanfolio{63}เอกมนฺตํ นิสินฺเน โข สาเลยฺยเก พฺราหฺมณคหปติเก ภควา เอตทโวจ “อตฺถิ ปน โว คหปตโย โกจิ มนาโป สตฺถา, ยสฺมิํ โว อาการวตี สทฺธา ปฏิลทฺธา”ติ.\csromanspacer{} นตฺถิ โข โน ภนฺเต โกจิ มนาโป สตฺถา, ยสฺมิํ โน อาการวตี สทฺธา ปฏิลทฺธาติ.\csromanspacer{} มนาปํ โว คหปตโย สตฺถารํ อลภนฺเตหิ อยํ อปณฺณโก ธมฺโม สมาทาย วตฺติตพฺโพ.\csromanspacer{} อปณฺณโก หิ คหปตโย ธมฺโม สมตฺโต สมาทินฺโน โส โว ภวิสฺสติ ทีฆรตฺตํ หิตาย สุขาย.\csromanspacer{} กตโม จ คหปตโย อปณฺณโก ธมฺโม.}

\proseitem{94}{สนฺติ คหปตโย เอเก สมณพฺราหฺมณา เอวํวาทิโน เอวํทิฏฺฐิโน “นตฺถิ ทินฺนํ, นตฺถิ ยิฏฺฐํ, นตฺถิ หุตํ, นตฺถิ สุกตทุกฺกฏานํ\footnote{สุกฏทุกฺกฏานํ (สี, สฺยา, กํ, อิ)} กมฺมานํ ผลํ วิปาโก, นตฺถิ อยํ โลโก, นตฺถิ ปโร โลโก, นตฺถิ มาตา, นตฺถิ ปิตา, นตฺถิ สตฺตา โอปปาติกา, นตฺถิ โลเก สมณพฺราหฺมณา สมฺมคฺคตา\footnote{สมคฺคตา (ก)} สมฺมา ปฏิปนฺนา, เย อิมญฺจ โลกํ ปรญฺจ โลกํ สยํ อภิญฺญา สจฺฉิกตฺวา ปเวเทนฺตี”ติ.\csromanspacer{} เตสํเยว โข คหปตโย \csromanfolio{64}สมณพฺราหฺมณานํ เอเก สมณพฺราหฺมณา อุชุวิปจฺจนีกวาทา.\csromanspacer{} เต เอวมาหํสุ “อตฺถิ ทินฺนํ, อตฺถิ ยิฏฺฐํ, อตฺถิ หุตํ, อตฺถิ สุกตทุกฺกฏานํ กมฺมานํ ผลํ วิปาโก, อตฺถิ อยํ โลโก, อตฺถิ ปโร โลโก, อตฺถิ มาตา, อตฺถิ ปิตา, อตฺถิ สตฺตา โอปปาติกา, อตฺถิ โลเก สมณพฺราหฺมณา สมฺมคฺคตา สมฺมา ปฏิปนฺนา, เย อิมญฺจ โลกํ ปรญฺจ โลกํ สยํ อภิญฺญา สจฺฉิกตฺวา ปเวเทนฺตี”ติ.\csromanspacer{} ตํ กิํ มญฺญถ คหปตโย, นนุเม สมณพฺราหฺมณา อญฺญมญฺญสฺส อุชุวิปจฺจนีกวาทาติ.\csromanspacer{} เอวํ ภนฺเต.}

\proseitem{95}{ตตฺร คหปตโย เย เต สมณพฺราหฺมณา เอวํวาทิโน เอวํทิฏฺฐิโน “นตฺถิ ทินฺนํ, นตฺถิ ยิฏฺฐํ ฯเปฯ เย อิมญฺจ โลกํ ปรญฺจ โลกํ สยํ อภิญฺญา สจฺฉิกตฺวา ปเวเทนฺตี”ติ.\csromanspacer{} เตสเมตํ ปาฏิกงฺขํ, ยมิทํ\footnote{ยทิทํ (ก)} กายสุจริตํ วจีสุจริตํ มโนสุจริตํ, อิเม ตโย กุสเล ธมฺเม อภินิวชฺเชตฺวา\footnote{อภินิพฺพชฺเชตฺวา (สฺยา, กํ), อภินิพฺพิชฺชิตฺวา (ก) 3. ปญฺญาเปติ (ก)} ยมิทํ1 กายทุจฺจริตํ วจีทุจฺจริตํ มโนทุจฺจริตํ, อิเม ตโย อกุสเล ธมฺเม สมาทาย วตฺติสฺสนฺติ.\csromanspacer{} ตํ กิสฺส เหตุ, น หิ เต โภนฺโต สมณพฺราหฺมณา ปสฺสนฺติ อกุสลานํ ธมฺมานํ อาทีนวํ โอการํ สํกิเลสํ, กุสลานํ ธมฺมานํ เนกฺขมฺเม อานิสํสํ โวทานปกฺขํ.\csromanspacer{} สนฺตํเยว ปน ปรํ โลกํ “นตฺถิ ปโร โลโก”ติสฺส ทิฏฺฐิ โหติ, สาสฺส โหติ มิจฺฉาทิฏฺฐิ.\csromanspacer{} สนฺตํเยว โข ปน ปรํ โลกํ “นตฺถิ ปโร โลโก”ติ สงฺกปฺเปติ, สฺวาสฺส โหติ มิจฺฉาสงฺกปฺโป.\csromanspacer{} สนฺตํเยว โข ปน ปรํ โลกํ “นตฺถิ ปโร โลโก”ติ วาจํ ภาสติ, สาสฺส โหติ มิจฺฉาวาจา.\csromanspacer{} สนฺตํเยว โข ปน ปรํ โลกํ “นตฺถิ ปโร โลโก”ติ อาห.\csromanspacer{} เย เต อรหนฺโต ปรโลกวิทุโน, เตสมยํ ปจฺจนีกํ กโรติ.\csromanspacer{} สนฺตํเยว โข ปน ปรํ โลกํ “นตฺถิ ปโร โลโก”ติ ปรํ สญฺญาเปติ3, สาสฺส โหติ อสทฺธมฺมสญฺญตฺติ\footnote{อสฺสทฺธมฺมปญฺญตฺติ (ก)}.\csromanspacer{} ตาย จ ปน อสทฺธมฺมสญฺญตฺติยา อตฺตานุกฺกํเสติ, ปรํ วมฺเภติ.\csromanspacer{} อิติ ปุพฺเพว โข ปนสฺส สุสีลฺยํ ปหีนํ โหติ ทุสฺสีลฺยํ ปจฺจุปฏฺฐิตํ, อยญฺจ มิจฺฉาทิฏฺฐิ มิจฺฉาสงฺกปฺโป มิจฺฉาวาจา อริยานํ ปจฺจนีกตา อสทฺธมฺมสญฺญตฺติ อตฺตุกฺกํสนา ปรวมฺภนา, เอวมสฺสิเม\footnote{เอวํ’สิ’เม (สี, สฺยา, กํ, อิ)} อเนเก ปาปกา อกุสลา ธมฺมา สมฺภวนฺติ มิจฺฉาทิฏฺฐิปจฺจยา.}

\csromanheader{2}{10. อปณฺณกสุตฺต (60)}
\prose{\csromanfolio{65}ตตฺร คหปตโย วิญฺญู ปุริโส อิติ ปฏิสญฺจิกฺขติ\csromandash{}สเจ โข นตฺถิ ปโร โลโก, เอวมยํ ภวํ ปุริสปุคฺคโล กายสฺส เภทา โสตฺถิมตฺตานํ กริสฺสติ.\csromanspacer{} สเจ โข อตฺถิ ปโร โลโก, เอวมยํ ภวํ ปุริสปุคฺคโล กายสฺส เภทา ปรํ มรณา อปายํ ทุคฺคติํ วินิปาตํ นิรยํ อุปปชฺชิสฺสติ.\csromanspacer{} กามํ โข ปน มาหุ ปโร โลโก, โหตุ เนสํ ภวตํ สมณพฺราหฺมณานํ สจฺจํ วจนํ, อถ จ ปนายํ ภวํ ปุริสปุคฺคโล ทิฏฺเฐว ธมฺเม วิญฺญูนํ คารโยฺห “ทุสฺสีโล ปุริสปุคฺคโล มิจฺฉาทิฏฺฐิ นตฺถิกวาโท”ติ.\csromanspacer{} สเจ โข อตฺเถว ปโร โลโก, เอวํ อิมสฺส โภโต ปุริสปุคฺคลสฺส อุภยตฺถ กลิคฺคโห, ยญฺจ ทิฏฺเฐว ธมฺเม วิญฺญูนํ คารโยฺห, ยญฺจ กายสฺส เภทา ปรํ มรณา อปายํ ทุคฺคติํ วินิปาตํ นิรยํ อุปปชฺชิสฺสติ, เอวมสฺสายํ อปณฺณโก ธมฺโม ทุสฺสมตฺโต สมาทินฺโน เอกํสํ ผริตฺวา ติฏฺฐติ, ริญฺจติ กุสลํ ฐานํ.}

\proseitem{96}{ตตฺร คหปตโย เย เต สมณพฺราหฺมณา เอวํวาทิโน เอวํทิฏฺฐิโน “อตฺถิ ทินฺนํ ฯเปฯ เย อิมญฺจ โลกํ ปรญฺจ โลกํ สยํ อภิญฺญา สจฺฉิกตฺวา ปเวเทนฺตี”ติ.\csromanspacer{} เตสเมตํ ปาฏิกงฺขํ, ยมิทํ กายทุจฺจริตํ วจีทุจฺจริตํ มโนทุจฺจริตํ, อิเม ตโย อกุสเล ธมฺเม อภินิวชฺเชตฺวา ยมิทํ กายสุจริตํ วจีสุจริตํ มโนสุจริตํ อิเม ตโย กุสเล ธมฺเม สมาทาย วตฺติสฺสนฺติ.\csromanspacer{} ตํ กิสฺส เหตุ, ปสฺสนฺติ หิ เต โภนฺโต สมณพฺราหฺมณา อกุสลานํ ธมฺมานํ อาทีนวํ โอการํ สํกิเลสํ, กุสลานํ ธมฺมานํ เนกฺขมฺเม อานิสํสํ โวทานปกฺขํ.\csromanspacer{} สนฺตํเยว โข ปน ปรํ โลกํ “อตฺถิ ปโร โลโก”ติสฺส ทิฏฺฐิ โหติ, สาสฺส โหติ สมฺมาทิฏฺฐิ.\csromanspacer{} สนฺตํเยว โข ปน ปรํ โลกํ “อตฺถิ ปโร โลโก”ติ สงฺกปฺเปติ, สฺวาสฺส โหติ สมฺมาสงฺกปฺโป.\csromanspacer{} สนฺตํเยว โข ปน ปรํ โลกํ “อตฺถิ ปโร โลโก”ติ วาจํ ภาสติ สาสฺส โหติ สมฺมาวาจา.\csromanspacer{} สนฺตํเยว โข ปน ปรํ โลกํ “อตฺถิ ปโร โลโก”ติ อาห, เย เต อรหนฺโต ปรโลกวิทุโน เตสมยํ น ปจฺจนีกํ กโรติ.\csromanspacer{} สนฺตํเยว โข ปน ปรํ โลกํ “อตฺถิ ปโร โลโก”ติ ปรํ สญฺญาเปติ, สาสฺส โหติ สทฺธมฺมสญฺญตฺติ.\csromanspacer{} ตาย จ ปน สทฺธมฺมสญฺญตฺติยา เนวตฺตานุกฺกํเสติ, น ปรํ วมฺเภติ.\csromanspacer{} อิติ ปุพฺเพว โข ปนสฺส ทุสฺสีลฺยํ ปหีนํ โหติ, สุสีลฺยํ ปจฺจุปฏฺฐิตํ, อยญฺจ สมฺมาทิฏฺฐิ สมฺมาสงฺกปฺโป สมฺมาวาจา อริยานํ อปจฺจนีกตา สทฺธมฺมสญฺญตฺติ \csromanfolio{66}อนตฺตุกฺกํสนา อปรวมฺภนา, เอวมสฺสิเม อเนเก กุสลา ธมฺมา สมฺภวนฺติ สมฺมาทิฏฺฐิปจฺจยา.}

\prose{ตตฺร คหปตโย วิญฺญู ปุริโส อิติ ปฏิสญฺจิกฺขติ\csromandash{}สเจ โข อตฺถิ ปโร โลโก, เอวมยํ ภวํ ปุริสปุคฺคโล กายสฺส เภทา ปรํ มรณา สุคติํ สคฺคํ โลกํ อุปปชฺชิสฺสติ, กามํ โข ปน มาหุ ปโร โลโก, โหตุ เนสํ ภวตํ สมณพฺราหฺมณานํ สจฺจํ วจนํ, อถ จ ปนายํ ภวํ ปุริสปุคฺคโล ทิฏฺเฐว ธมฺเม วิญฺญูนํ ปาสํโส “สีลวา ปุริสปุคฺคโล สมฺมาทิฏฺฐิ อตฺถิกวาโท”ติ.\csromanspacer{} สเจ โข อตฺเถว ปโร โลโก, เอวํ อิมสฺส โภโต ปุริสปุคฺคลสฺส อุภยตฺถ กฏคฺคโห, ยญฺจ ทิฏฺเฐว ธมฺเม วิญฺญูนํ ปาสํโส, ยญฺจ กายสฺส เภทา ปรํ มรณา สุคติํ สคฺคํ โลกํ อุปฺปชฺชิสฺสติ.\csromanspacer{} เอวมสฺสายํ อปณฺณโก ธมฺโม สุสมตฺโต สมาทินฺโน อุภยํสํ ผริตฺวา ติฏฺฐติ, ริญฺจติ อกุสลํ ฐานํ.}

\proseitem{97}{สนฺติ คหปตโย เอเก สมณพฺราหฺมณา เอวํวาทิโน เอวํ ทิฏฺฐิโน “กโรโต การยโต, ฉินฺทโต เฉทาปยโต, ปจโต ปาจาปยโต, โสจยโต โสจาปยโต, กิลมโต กิลมาปยโต, ผนฺทโต ผนฺทาปยโต, ปาณมติปาตยโต\footnote{ปาณมติมาปยโต (สี, อิ), ปาณมติปาตาปยโต (สฺยา, กํ), ปาณมติปาปยโต (ก)}, อทินฺนํ อาทิยโต, สนฺธิํ ฉินฺทโต, นิลฺโลปํ หรโต, เอกาคาริกํ กโรโต, ปริปนฺเถ ติฏฺฐโต, ปรทารํ คจฺฉโต, มุสา ภณโต, กโรโต น กรียติ ปาปํ.\csromanspacer{} ขุรปริยนฺเตน เจปิ จกฺเกน โย อิมิสฺสา ปถวิยา ปาเณ เอกํ มํสขลํ เอกํ มํสปุญฺชํ กเรยฺย, นตฺถิ ตโตนิทานํ ปาปํ, นตฺถิ ปาปสฺส อาคโม.\csromanspacer{} ทกฺขิณญฺเจปิ คงฺคาย ตีรํ คจฺเฉยฺย หนนฺโต ฆาเตนฺโต ฉินฺทนฺโต เฉทาเปนฺโต ปจนฺโต ปาเจนฺโต, นตฺถิ ตโตนิทานํ ปาปํ, นตฺถิ ปาปสฺส อาคโม.\csromanspacer{} อุตฺตรญฺเจปิ คงฺคาย ตีรํ คจฺเฉยฺย ททนฺโต ทาปนฺโต ยชนฺโต ยชาเปนฺโต, นตฺถิ ตโตนิทานํ ปุญฺญํ, นตฺถิ ปุญฺญสฺส อาคโม.\csromanspacer{} ทาเนน ทเมน สํยเมน สจฺจวชฺเชน\footnote{สจฺจวาเจน (ก)} นตฺถิ ปุญฺญํ, นตฺถิ ปุญฺญสฺส อาคโม”ติ.\csromanspacer{} เตสํเยว โข คหปตโย สมณพฺราหฺมณานํ เอเก สมณพฺราหฺมณา อุชุวิปจฺจนีกวาทา, เต เอวมาหํสุ “กโรโต การยโต, ฉินฺทโต เฉทาปยโต, ปจโต ปาจาปยโต,}

\vspace{\tipitakaparskip}
\csromancenter{อปณฺณกสุตฺต (60) โสจยโต โสจาปยโต, กิลมโต กิลมาปยโต, ผนฺทโต ผนฺทาปยโต, ปาณมติปาตยโต, อทินฺนํ อาทิยโต, สนฺธิํ ฉินฺทโต, นิลฺโลปํ หรโต, เอกาคาริกํ กโรโต, ปริปนฺเถ ติฏฺฐโต, ปรทารํ คจฺฉโต, มุสา ภณโต, กโรโต กรียติ ปาปํ.\csromanspacer{} ขุรปริยนฺเตน เจปิ จกฺเกน โย อิมิสฺสา ปถวิยา ปาเณ เอกํ มํสขลํ เอกํ มํสปุญฺชํ กเรยฺย, อตฺถิ ตโตนิทานํ ปาปํ, อตฺถิ ปาปสฺส อาคโม.\csromanspacer{} ทกฺขิณญฺเจปิ คงฺคาย ตีรํ คจฺเฉยฺย, หนนฺโต ฆาเตนฺโต ฉินฺทนฺโต เฉทาเปนฺโต ปจนฺโต ปาเจนฺโต, อตฺถิ ตโตนิทานํ ปาปํ, อตฺถิ ปาปสฺส อาคโม.\csromanspacer{} อุตฺตรญฺเจปิ คงฺคาย ตีรํ คจฺเฉยฺย, ททนฺโต ทาเปนฺโต ยชนฺโต ยชาเปนฺโต, อตฺถิ ตโตนิทานํ ปุญฺญํ, อตฺถิ ปุญฺญสฺส อาคโม.\csromanspacer{} ทาเนน ทเมน สํยเมน สจฺจวชฺเชน อตฺถิ ปุญฺญํ, อตฺถิ ปุญฺญสฺส อาคโม”ติ.\csromanspacer{} ตํ กิํ มญฺญถ คหปตโย, นนุเม สมณพฺราหฺมณา อญฺญมญฺญสฺส อุชุวิปจฺจนีกวาทาติ.\csromanspacer{} เอวํ ภนฺเต.}
\proseitem{98}{\csromanfolio{67}ตตฺร คหปตโย เย เต สมณพฺราหฺมณา เอวํวาทิโน เอวํทิฏฺฐิโน “กโรโต การยโต, ฉินฺทโต เฉทาปยโต, ปจโต ปาจาปยโต, โสจยโต โสจาปยโต, กิลมโต กิลมาปยโต, ผนฺทโต ผนฺทาปยโต, ปาณมติปาตยโต, อทินฺนํ อาทิยโต, สนฺธิํ ฉินฺทโต, นิลฺโลปํ หรโต, เอกาคาริกํ กโรโต, ปริปนฺเถ ติฏฺฐโต, ปรทารํ คจฺฉโต, มุสา ภณโต, กโรโต น กรียติ ปาปํ.\csromanspacer{} ขุรปริยนฺเตน เจปิ จกฺเกน โย อิมิสฺสา ปถวิยา ปาเณ เอกํ มํสขลํ เอกํ มํสปุญฺชํ กเรยฺย, นตฺถิ ตโตนิทานํ ปาปํ, นตฺถิ ปาปสฺส อาคโม.\csromanspacer{} ทกฺขิณญฺเจปิ คงฺคาย ตีรํ คจฺเฉยฺย, หนนฺโต ฆาเตนฺโต ฯเปฯ.\csromanspacer{} ทาเนน ทเมน สํยเมน สจฺจวชฺเชน นตฺถิ ปุญฺญํ, นตฺถิ ปุญฺญสฺส อาคโม”ติ.\csromanspacer{} เตสเมตํ ปาฏิกงฺขํ, ยมิทํ กายสุจริตํ วจีสุจริตํ มโนสุจริตํ, อิเม ตโย กุสเล ธมฺเม อภินิวชฺเชตฺวา ยมิทํ กายทุจฺจริตํ วจีทุจฺจริตํ มโนทุจฺจริตํ, อิเม ตโย อกุสเล ธมฺเม สมาทาย วตฺติสฺสนฺติ, ตํ กิสฺส เหตุ, น หิ เต โภนฺโต สมณพฺราหฺมณา ปสฺสนฺติ อกุสลานํ ธมฺมานํ อาทีนวํ โอการํ สํกิเลสํ, กุสลานํ ธมฺมานํ เนกฺขมฺเม อานิสํสํ โวทานปกฺขํ.\csromanspacer{} สนฺตํเยว โข ปน กิริยํ “นตฺถิ กิริยา”ติสฺส ทิฏฺฐิ โหติ, สาสฺส โหติ มิจฺฉาทิฏฺฐิ.\csromanspacer{} สนฺตํเยว โข ปน กิริยํ “นตฺถิ กิริยา”ติ สงฺกปฺเปติ, สฺวาสฺส \csromanfolio{68}โหติ มิจฺฉาสงฺกปฺโป.\csromanspacer{} สนฺตํเยว โข ปน กิริยํ “นตฺถิ กิริยา”ติ วาจํ ภาสติ, สาสฺส โหติ มิจฺฉาวาจา.\csromanspacer{} สนฺตํเยว โข ปน กิริยํ “นตฺถิ กิริยา”ติ อาห.\csromanspacer{} เย เต อรหนฺโต กิริยวาทา, เตสมยํ ปจฺจนีกํ กโรติ.\csromanspacer{} สนฺตํเยว โข ปน กิริยํ “นตฺถิ กิริยา”ติ ปรํ สญฺญาเปติ.\csromanspacer{} สาสฺส โหติ อสทฺธมฺมสญฺญตฺติ, ตาย จ ปน อสทฺธมฺมสญฺญตฺติยา อตฺตานุกฺกํเสติ, ปรํ วมฺเภติ.\csromanspacer{} อิติ ปุพฺเพว โข ปนสฺส สุสิลฺยํ ปหีนํ โหติ, ทุสฺสิลฺยํ ปจฺจุปฏฺฐิตํ, อยญฺจ มิจฺฉาทิฏฺฐิ มิจฺฉาสงฺกปฺโป มิจฺฉาวาจา อริยานํ ปจฺจนีกตา อสทฺธมฺมสญฺญตฺติ อตฺตุกฺกํสนา ปรวมฺภนา, เอวมสฺสิเม อเนเก ปาปกา อกุสลา ธมฺมา สมฺภวนฺติ มิจฺฉาทิฏฺฐิปจฺจยา.}

\prose{ตตฺร คหปตโย วิญฺญู ปุริโส อิติ ปฏิสญฺจิกฺขติ\csromandash{}สเจ โข นตฺถิ กิริยา, เอวมยํ ภวํ ปุริสปุคฺคโล กายสฺส เภทา โสตฺถิมตฺตานํ กริสฺสติ.\csromanspacer{} สเจ โข อตฺถิ กิริยา, เอวมยํ ภวํ ปุริสปุคฺคโล กายสฺส เภทา ปรํ มรณา อปายํ ทุคฺคติํ วินิปาตํ นิรยํ อุปปชฺชิสฺสติ.\csromanspacer{} กามํ โข ปน มาหุ กิริยา, โหตุ เนสํ ภวตํ สมณพฺราหฺมณานํ สจฺจํ วจนํ, อถ จ ปนายํ ภวํ ปุริสปุคฺคโล ทิฏฺเฐว ธมฺเม วิญฺญูนํ คารโยฺห “ทุสฺสีโล ปุริสปุคฺคโล มิจฺฉาทิฏฺฐิ อกิริยวาโท”ติ.\csromanspacer{} สเจ โข อตฺเถว กิริยา, เอวํ อิมสฺส โภโต ปุริสปุคฺคลสฺส อุภยตฺถ กลิคฺคโห, ยญฺจ ทิฏฺเฐว ธมฺเม วิญฺญูนํ คารโยฺห, ยญฺจ กายสฺส เภทา ปรํ มรณา อปายํ ทุคฺคติํ วินิปาตํ นิรยํ อุปปชฺชิสฺสติ.\csromanspacer{} เอวมสฺสายํ อปณฺณโก ธมฺโม ทุสฺสมตฺโต สมาทินฺโน เอกํสํ ผริตฺวา ติฏฺฐติ, ริญฺจติ กุสลํ ฐานํ.}

\proseitem{99}{ตตฺร คหปตโย เย เต สมณพฺราหฺมณา เอวํวาทิโน เอวํทิฏฺฐิโน “กโรโต การยโต, ฉินฺทโต เฉทาปยโต, ปจโต ปาจาปยโต, โสจยโต โสจาปยโต, กิลมโต กิลมาปยโต, ผนฺทโต ผนฺทาปยโต, ปาณมติปาตยโต, อทินฺนํ อาทิยโต, สนฺธิํ ฉินฺทโต, นิลฺโลปํ หรโต, เอกาคาริกํ กโรโต, ปริปนฺเถ ติฏฺฐโต, ปรทารํ คจฺฉโต, มุสา ภณโต, กโรโต กรียติ ปาปํ.\csromanspacer{} ขุรปริยนฺเตน เจปิ จกฺเกน โย อิมิสฺสา ปถวิยา ปาเณ เอกํ มํสขลํ เอกํ มํสปุญฺชํ กเรยฺย, อตฺถิ ตโตนิทานํ ปาปํ, อตฺถิ ปาปสฺส อาคโม.\csromanspacer{} ทกฺขิณญฺเจปิ คงฺคาย ตีรํ คจฺเฉยฺย หนนฺโต ฆาเตนฺโต}

\vspace{\tipitakaparskip}
\csromancenter{อปณฺณกสุตฺต (60) ฉินฺทนฺโต เฉทาเปนฺโต ปจนฺโต ปาเจนฺโต, อตฺถิ ตโตนิทานํ ปาปํ, อตฺถิ ปาปสฺส อาคโม.\csromanspacer{} อุตฺตรญฺเจปิ คงฺคาย ตีรํ คจฺเฉยฺย ททนฺโต ทาเปนฺโต ยชนฺโต ยชาเปนฺโต, อตฺถิ ตโตนิทานํ ปุญฺญํ, อตฺถิ ปุญฺญสฺส อาคโม.\csromanspacer{} ทาเนน ทเมน สํยเมน สจฺจวชฺเชน อตฺถิ ปุญฺญํ, อตฺถิ ปุญฺญสฺส อาคโม”ติ.\csromanspacer{} เตสเมตํ ปาฏิกงฺขํ, ยมิทํ กายทุจฺจริตํ วจีทุจฺจริตํ มโนทุจฺจริตํ, อิเม ตโย อกุสเล ธมฺเม อภินิวชฺเชตฺวา ยมิทํ กายสุจริตํ วจีสุจริตํ มโนสุจริตํ, อิเม ตโย กุสเล ธมฺเม สมาทาย วตฺติสฺสนฺติ, ตํ กิสฺส เหตุ, ปสฺสนฺติ หิ เต โภนฺโต สมณพฺราหฺมณา อกุสลานํ ธมฺมานํ อาทีนวํ โอการํ สํกิเลสํ, กุสลานํ ธมฺมานํ เนกฺขมฺเม อานิสํสํ โวทานปกฺขํ.\csromanspacer{} สนฺตํเยว โข ปน กิริยํ “อตฺถิ กิริยา”ติสฺส ทิฏฺฐิ โหติ, สาสฺส โหติ สมฺมาทิฏฺฐิ.\csromanspacer{} สนฺตํเยว โข ปน กิริยํ “อตฺถิ กิริยา”ติ สงฺกปฺเปติ, สฺวาสฺส โหติ สมฺมาสงฺกปฺโป.\csromanspacer{} สนฺตํเยว โข ปน กิริยํ “อตฺถิ กิริยา”ติ วาจํ ภาสติ, สาสฺส โหติ สมฺมาวาจา.\csromanspacer{} สนฺตํเยว โข ปน กิริยํ “อตฺถิ กิริยา”ติ อาห.\csromanspacer{} เย เต อรหนฺโต กิริยวาทา, เตสมยํ น ปจฺจนีกํ กโรติ.\csromanspacer{} สนฺตํเยว โข ปน กิริยํ “อตฺถิ กิริยา”ติ ปรํ สญฺญาเปติ, สาสฺส โหติ สทฺธมฺมสญฺญตฺติ.\csromanspacer{} ตาย จ ปน สทฺธมฺมสญฺญตฺติยา เนวตฺตานุกฺกํเสติ, น ปรํ วมฺเภติ.\csromanspacer{} อิติ ปุพฺเพว โข ปนสฺส ทุสฺสีลฺยํ ปหีนํ โหติ, สุสีลฺยํ ปจฺจุปฏฺฐิตํ, อยญฺจ สมฺมาทิฏฺฐิ สมฺมาสงฺกปฺโป สมฺมาวาจา อริยานํ อปจฺจนีกตา สทฺธมฺมสญฺญตฺติ อนตฺตุกฺกํสนา อปรวมฺภนา, เอวมสฺสิเม อนเก กุสลา ธมฺมา สมฺภวนฺติ สมฺมาทิฏฺฐิปจฺจยา.}
\prose{\csromanfolio{69}ตตฺร คหปตโย วิญฺญู ปุริโส อิติ ปฏิสญฺจิกฺขติ\csromandash{}สเจ โข อตฺถิ กิริยา, เอวมยํ ภวํ ปุริสปุคฺคโล กายสฺส เภทา ปรํ มรณา สุคติํ สคฺคํ โลกํ อุปปชฺชิสฺสติ, กามํ โข ปน มาหุ กิริยา, โหตุ เนสํ ภวตํ สมณพฺราหฺมณานํ สจฺจํ วจนํ, อถ จ ปนายํ ภวํ ปุริสปุคฺคโล ทิฏฺเฐว ธมฺเม วิญฺญูนํ ปาสํโส “สีลวา ปุริสปุคฺคโล สมฺมาทิฏฺฐิ กิริยวาโท”ติ.\csromanspacer{} สเจ โข อตฺเถว กิริยา, เอวํ อิมสฺส โภโต ปุริสปุคฺคลสฺส อุภยตฺถ กฏคฺคโห, ยญฺจ ทิฏฺเฐว ธมฺเม วิญฺญูนํ ปาสํโส, ยญฺจ กายสฺส เภทา ปรํ มรณา สุคติํ สคฺคํ โลกํ อุปปชฺชิสฺสติ.\csromanspacer{} เอวมสฺสายํ อปณฺณโก ธมฺโม สุสมโต สมาทินฺโน อุภยํสํ ผริตฺวา ติฏฺฐติ, ริญฺจติ อกุสลํ ฐานํ.}

\proseitem{100}{\csromanfolio{70}สนฺติ คหปตโย เอเก สมณพฺราหฺมณา เอวํวาทิโน เอวํทิฏฺฐิโน “นตฺถิ เหตุ นตฺถิ ปจฺจโย สตฺตานํ สํกิเลสาย, อเหตู อปฺปจฺจยา สตฺตา สํกิลิสฺสนฺติ.\csromanspacer{} นตฺถิ เหตุ นตฺถิ ปจฺจโย สตฺตานํ วิสุทฺธิยา, อเหตู อปฺปจฺจยา สตฺตา วิสุชฺฌนฺติ.\csromanspacer{} นตฺถิ พลํ, นตฺถิ วีริยํ\footnote{วิริยํ (สี, สฺยา, กํ, อิ)}, นตฺถิ ปุริสถาโม, นตฺถิ ปุริสปรกฺกโม, สพฺเพ สตฺตา สพฺเพ ปาณา สพฺเพ ภูตา สพฺเพ ชีวา อวสา อพลา อวีริยา นิยติสํคติภาวปริณตา ฉเสฺว วาภิชาตีสุ สุขทุกฺขํ ปฏิสํเวเทนฺตี”ติ.\csromanspacer{} เตสํเยว โข คหปตโย สมณพฺราหฺมณานํ เอเก สมณพฺราหฺมณา อุชุวิปจฺจนีกวาทา, เต เอวมาหํสุ “อตฺถิ เหตุ อตฺถิ ปจฺจโย สตฺตานํ สํกิเลสาย, สเหตู สปฺปจฺจยา สตฺตา สํกิลิสฺสนฺติ.\csromanspacer{} อตฺถิ เหตุ อตฺถิ ปจฺจโย สตฺตานํ วิสุทฺธิยา, สเหตู สปฺปจฺจยา สตฺตา วิสุชฺฌนฺติ.\csromanspacer{} อตฺถิ พลํ, อตฺถิ วีริยํ, อตฺถิ ปุริสถาโม, อตฺถิ ปุริสปรกฺกโม, น สพฺเพ สตฺตา สพฺเพ ปาณา สพฺเพ ภูตา สพฺเพ ชีวา อวสา อพลา อวีริยา\footnote{อตฺถิ ปุริสปรกฺกโม, สพฺเพ สตฺตา...สวลา สพลา สวีริยา (สฺยา, กํ, ก)} นิยติสํคติภาวปริณตา ฉเสฺววาภิชาตีสุ สุขทุกฺขํ ปฏิสํเวเทนฺตี”ติ.\csromanspacer{} ตํ กิํ มญฺญถ คหปตโย, นนุเม สมณพฺราหฺมณา อญฺญมญฺญสฺส อุชุวิปจฺจนีกวาทาติ.\csromanspacer{} เอวํ ภนฺเต.}

\proseitem{101}{ตตฺร คหปตโย เย เต สมณพฺราหฺมณา เอวํวาทิโน เอวํทิฏฺฐิโน “นตฺถิ เหตุ นตฺถิ ปจฺจโย สตฺตานํ สํกิเลสาย, อเหตู อปฺปจฺจยา สตฺตา สํกิลิสฺสนฺติ.\csromanspacer{} นตฺถิ เหตุ นตฺถิ ปจฺจโย สตฺตานํ วิสุทฺธิยา, อเหตู อปฺปจฺจยา สตฺตา วิสุชฺฌนฺติ.\csromanspacer{} นตฺถิ พลํ, นตฺถิ วีริยํ, นตฺถิ ปุริสถาโม, นตฺถิ ปุริสปรกฺกโม, สพฺเพ สตฺตา สพฺเพ ปาณา สพฺเพ ภูตา สพฺเพ ชีวา อวสา อพลา อวีริยา นิยติสํคติภาวปริณตา, ฉเสฺววาภิชาตีสุ สุขทุกฺขํ ปฏิสํเวเทนฺตี”ติ.\csromanspacer{} เตสเมตํ ปาฏิกงฺขํ, ยมิทํ กายสุจริตํ วจีสุจริตํ มโนสุจริตํ, อิเม ตโย กุสเล ธมฺเม อภินิวชฺเชตฺวา ยมิทํ กายทุจฺจริตํ วจีทุจฺจริตํ มโนทุจฺจริตํ, อิเม ตโย อกุสเล ธมฺเม สมาทาย วตฺติสฺสนฺติ.\csromanspacer{} ตํ กิสฺส เหตุ, น หิ เต โภนฺโต สมณพฺราหฺมณา ปสฺสนฺติ อกุสลานํ ธมฺมานํ อาทีนวํ โอการํ สํกิเลสํ, กุสลานํ ธมฺมานํ เนกฺขมฺเม อานิสํสํ โวทานปกฺขํ.\csromanspacer{} สนฺตํเยว โข ปน เหตุํ}

\vspace{\tipitakaparskip}
\csromancenter{อปณฺณกสุตฺต (60) “นตฺถิ เหตู”ติสฺส ทิฏฺฐิ โหติ, สาสฺส โหติ มิจฺฉาทิฏฺฐิ.\csromanspacer{} สนฺตํเยว โข ปน เหตุํ “นตฺถิ เหตู”ติ สงฺกปฺเปติ, สฺวาสฺส โหติ มิจฺฉาสงฺกปฺโป.\csromanspacer{} สนฺตํเยว โข ปน เหตุํ “นตฺถิ เหตู”ติ วาจํ ภาสติ, สาสฺส โหติ มิจฺฉาวาจา.\csromanspacer{} สนฺตํเยว โข ปน เหตุํ “นตฺถิ เหตู”ติ อาห.\csromanspacer{} เย เต อรหนฺโต เหตุวาทา เตสมยํ ปจฺจนีกํ กโรติ.\csromanspacer{} สนฺตํเยว โข ปน เหตุํ “นตฺถิ เหตู”ติ ปรํ สญฺญาเปติ, สาสฺส โหติ อสทฺธมฺมสญฺญตฺติ.\csromanspacer{} ตาย จ ปน อสทฺธมฺมสญฺญตฺติยา อตฺตานุกฺกํเสติ, ปรํ วมฺเภติ.\csromanspacer{} อิติ ปุพฺเพว โข ปนสฺส สุสีลฺยํ ปหีนํ โหติ, ทุสฺสีลฺยํ ปจฺจุปฏฺฐิตํ, อยญฺจ มิจฺฉาทิฏฺฐิ มิจฺฉาสงฺกปฺโป มิจฺฉาวาจา อริยานํ ปจฺจนีกตา อสทฺธมฺมสญฺญตฺติ อตฺตานุกฺกํสนา ปรวมฺภนา, เอวมสฺสิเม อเนเก ปาปกา อกุสลา ธมฺมา สมฺภวนฺติ มิจฺฉาทิฏฺฐิปจฺจยา.}
\prose{\csromanfolio{71}ตตฺร คหปตโย วิญฺญู ปุริโส อิติ ปฏิสญฺจิกฺขติ\csromandash{}สเจ โข นตฺถิ เหตุ, เอวมยํ ภวํ ปุริสปุคฺคโล กายสฺส เภทา ปรํ มรณา โสตฺถิมตฺตานํ กริสฺสติ.\csromanspacer{} สเจ โข อตฺถิ เหตุ, เอวมยํ ภวํ ปุริสปุคฺคโล กายสฺส เภทา ปรํ มรณา อปายํ ทุคฺคติํ วินิปาตํ นิรยํ อุปปชฺชิสฺสติ.\csromanspacer{} กามํ โข ปน มาหุ เหตุ, โหตุ เนสํ ภวตํ สมณพฺราหฺมณานํ สจฺจํ วจนํ, อถ จ ปนายํ ภวํ ปุริสปุคฺคโล ทิฏฺเฐว ธมฺเม วิญฺญูนํ คารโยฺห “ทุสฺสีโล ปุริสปุคฺคโล มิจฺฉาทิฏฺฐิ อเหตุกวาโท”ติ.\csromanspacer{} สเจ โข อตฺเถว เหตุ, เอวํ อิมสฺส โภโต ปุริสปุคฺคลสฺส อุภยตฺถ กลิคฺคโห, ยญฺจ ทิฏฺเฐว ธมฺเม วิญฺญูนํ คารโยฺห, ยญฺจ กายสฺส เภทา ปรํ มรณา อปายํ ทุคฺคติํ วินิปาตํ นิรยํ อุปปชฺชิสฺสติ.\csromanspacer{} เอวมสฺสายํ อปณฺณโก ธมฺโม ทุสฺสมตฺโต สมาทินฺโน เอกํสํ ผริตฺวา ติฏฺฐติ, ริญฺจติ กุสลํ ฐานํ.}

\proseitem{102}{ตตฺร คหปตโย เย เต สมณพฺราหฺมณา เอวํวาทิโน เอวํทิฏฺฐิโน “อตฺถิ เหตุ อตฺถิ ปจฺจโย สตฺตานํ สํกิเลสาย, สเหตู สปฺปจฺจยา สตฺตา สํกิลิสฺสนฺติ.\csromanspacer{} อตฺถิ เหตุ อตฺถิ ปจฺจโย สตฺตานํ วิสุทฺธิยา, สเหตู สปฺปจฺจยา สตฺตา วิสุชฺฌนฺติ.\csromanspacer{} อตฺถิ พลํ, อตฺถิ วีริยํ, อตฺถิ ปุริสถาโม, อตฺถิ ปุริสปรกฺกโม, น สพฺเพ สตฺตา สพฺเพ ปาณา สพฺเพ ภูตา สพฺเพ ชีวา อวสา อพลา อวีริยา นิยติสํคติภาวปริณตา ฉเสฺววาภิชาตีสุ สุขทุกฺขํ ปฏิสํเวเทนฺตี”ติ.\csromanspacer{} เตสเมตํ ปาฏิกงฺขํ, ยมิทํ กายทุจฺจริตํ วจีทุจฺจริตํ มโนทุจฺจริตํ, อิเม \csromanfolio{72}ตโย อกุสเล ธมฺเม อภินิวชฺเชตฺวา, ยมิทํ กายสุจริตํ วจีสุจริตํ มโนสุจริตํ, อิเม ตโย กุสเล ธมฺเม สมาทาย วตฺติสฺสนฺติ, ตํ กิสฺส เหตุ, ปสฺสนฺติ หิ เต โภนฺโต สมณพฺราหฺมณา อกุสลานํ ธมฺมานํ อาทีนวํ โอการํ สํกิเลสํ, กุสลานํ ธมฺมานํ เนกฺขมฺเม อานิสํสํ โวทานปกฺขํ.\csromanspacer{} สนฺตํเยว โข ปน เหตุํ “อตฺถิ เหตู”ติสฺส ทิฏฺฐิ โหติ, สาสฺส โหติ สมฺมาทิฏฺฐิ.\csromanspacer{} สนฺตํเยว โข ปน เหตุํ “อตฺถิ เหตู”ติ สงฺกปฺเปติ, สฺวาสฺส โหติ สมฺมาสงฺกปฺโป.\csromanspacer{} สนฺตํเยว โข ปน เหตุํ “อตฺถิ เหตู”ติ วาจํ ภาสติ, สาสฺส โหติ สมฺมาวาจา.\csromanspacer{} สนฺตํเยว โข ปน เหตุํ “อตฺถิ เหตู”ติ อาห.\csromanspacer{} เย เต อรหนฺโต เหตุวาทา, เตสมยํ น ปจฺจนีกํ กโรติ.\csromanspacer{} สนฺตํเยว โข ปน เหตุํ “อตฺถิ เหตู”ติ ปรํ สญฺญาเปติ, สาสฺส โหติ สทฺธมฺมสญฺญตฺติ.\csromanspacer{} ตาย จ ปน สทฺธมฺมสญฺญตฺติยา เนวตฺตานุกฺกํเสติ, น ปรํ วมฺเภติ.\csromanspacer{} อิติ ปุพฺเพว โข ปนสฺส ทุสฺสีลฺยํ ปหีนํ โหติ, สุสีลฺยํ ปจฺจุปฏฺฐิตํ, อยญฺจ สมฺมาทิฏฺฐิ สมฺมาสงฺกปฺโป สมฺมาวาจา อริยานํ อปจฺจนีกตา สทฺธมฺมสญฺญตฺติ อนตฺตุกฺกํสนา อปรวมฺภนา, เอวมสฺสิเม อเนเก กุสลา ธมฺมา สมฺภวนฺติ สมฺมาทิฏฺฐิปจฺจยา.}

\prose{ตตฺร คหปตโย วิญฺญู ปุริโส อิติ ปฏิสญฺจิกฺขติ\csromandash{}สเจ โข อตฺถิ เหตุ, เอวมยํ ภวํ ปุริสปุคฺคโล กายสฺส เภทา ปรํ มรณา สุคติํ สคฺคํ โลกํ อุปปชฺชิสฺสติ.\csromanspacer{} กามํ โข ปน มาหุ เหตุ, โหตุ เนสํ ภวตํ สมณพฺราหฺมณานํ สจฺจํ วจนํ, อถ จ ปนายํ ภวํ ปุริสปุคฺคโล ทิฏฺเฐว ธมฺเม วิญฺญูนํ ปาสํโส “สีลวา ปุริสปุคฺคโล สมฺมาทิฏฺฐิ เหตุวาโท”ติ.\csromanspacer{} สเจ โข อตฺถิ เหตุ, เอวํ อิมสฺส โภโต ปุริสปุคฺคลสฺส อุภยตฺถ กฏคฺคโห, ยญฺจ ทิฏฺเฐว ธมฺเม วิญฺญูนํ ปาสํโส, ยญฺจ กายสฺส เภทา ปรํ มรณา สุคติํ สคฺคํ โลกํ อุปปชฺชิสฺสติ.\csromanspacer{} เอวมสฺสายํ อปณฺณโก ธมฺโม สุสมตฺโต สมาทินฺโน อุภยํสํ ผริตฺวา ติฏฺฐติ, ริญฺจติ อกุสลํ ฐานํ.}

\proseitem{103}{สนฺติ คหปตโย เอเก สมณพฺราหฺมณา เอวํวาทิโน เอวํทิฏฺฐิโน “นตฺถิ สพฺพโส อารุปฺปา”ติ.\csromanspacer{} เตสํเยว โข คหปตโย สมณพฺราหฺมณานํ เอเก สมณพฺราหฺมณา อุชุวิปจฺจนีกวาทา, เต เอวมาหํสุ “อตฺถิ สพฺพโส อารุปฺปา”ติ.\csromanspacer{} ตํ กิํ มญฺญถ คหปตโย, นนุเม สมณพฺราหฺมณา อญฺญมญฺญสฺส อุชุวิปจฺจนีกวาทาติ.\csromanspacer{} เอวํ ภนฺเต.}

\vspace{\tipitakaparskip}
\csromancenter{อปณฺณกสุตฺต (60) ตตฺร คหปตโย วิญฺญู ปุริโส อิติ ปฏิสญฺจิกฺขติ “เย โข เต โภนฺโต สมณพฺราหฺมณา เอวํวาทิโน เอวํทิฏฺฐิโน ‘นตฺถิ สพฺพโส อารุปฺปา’ติ, อิทํ เม อทิฏฺฐํ.\csromanspacer{} เยปิ เต โภนฺโต สมณพฺราหฺมณา เอวํวาทิโน เอวํทิฏฺฐิโน ‘อตฺถิ สพฺพโส อารุปฺปา’ติ, อิทํ เม อวิทิตํ.\csromanspacer{} อหํ เจว\footnote{อหญฺเจ (?)} โข ปน อชานนฺโต อปสฺสนฺโต เอกํเสน อาทาย โวหเรยฺยํ ‘อิทเมว สจฺจํ โมฆมญฺญนฺ’ติ, น เมตํ อสฺส ปติรูปํ.\csromanspacer{} เย โข เต โภนฺโต สมณพฺราหฺมณา เอวํวาทิโน เอวํทิฏฺฐิโน ‘นตฺถิ สพฺพโส อารุปฺปา’ติ, สเจ เตสํ ภวตํ สมณพฺราหฺมณานํ สจฺจํ วจนํ, ฐานเมตํ วิชฺชติ, เย เต เทวา รูปิโน มโนมยา, อปณฺณกํ เม ตตฺรูปปตฺติ ภวิสฺสติ.\csromanspacer{} เย ปน เต โภนฺโต สมณพฺราหฺมณา เอวํวาทิโน เอวํทิฏฺฐิโน ‘อตฺถิ สพฺพโส อารุปฺปา’ติ, สเจ เตสํ ภวตํ สมณพฺราหฺมณานํ สจฺจํ วจนํ, ฐานเมตํ วิชฺชติ, เย เต เทวา อรูปิโน สญฺญามยา, อปณฺณกํ เม ตตฺรูปปตฺติ ภวิสฺสติ.\csromanspacer{} ทิสฺสนฺติ โข ปน รูปาธิกรณํ\footnote{รูปการณา (ก)} ทณฺฑาทาน สตฺถาทาน กลห วิคฺคห วิวาทตุวํตุวํ เปสุญฺญ มุสาวาทา.\csromanspacer{} นตฺถิ โข ปเนตํ สพฺพโส อรูเป”ติ.\csromanspacer{} โส อิติ ปฏิสงฺขาย รูปานํเยว นิพฺพิทาย วิราคาย นิโรธาย ปฏิปนฺโน โหติ.}
\proseitem{104}{\csromanfolio{73}สนฺติ คหปตโย เอเก สมณพฺราหฺมณา เอวํวาทิโน เอวํทิฏฺฐิโน “นตฺถิ สพฺพโส ภวนิโรโธ”ติ.\csromanspacer{} เตสํเยว โข คหปตโย สมณพฺราหฺมณานํ เอเก สมณพฺราหฺมณา อุชุวิปจฺจนีกวาทา, เต เอวมาหํสุ “อตฺถิ สพฺพโส ภวนิโรโธ”ติ.\csromanspacer{} ตํ กิํ มญฺญถ คหปตโย, นนุเม สมณพฺราหฺมณา อญฺญมญฺญสฺส อุชุวิปจฺจนีกวาทาติ.\csromanspacer{} เอวํ ภนฺเต.\csromanspacer{} ตตฺร คหปตโย วิญฺญู ปุริโส อิติ ปฏิสญฺจิกฺขติ “เย โข เต โภนฺโต สมณพฺราหฺมณา เอวํวาทิโน เอวํทิฏฺฐิโน ‘นตฺถิ สพฺพโส ภวนิโรโธ’ติ, อิทํ เม อทิฏฺฐํ.\csromanspacer{} เยปิ เต โภนฺโต สมณพฺราหฺมณา เอวํวาทิโน เอวํทิฏฺฐิโน ‘อตฺถิ สพฺพโส ภวนิโรโธ’ติ, อิทํ เม อวิทิตํ.\csromanspacer{} อหํ เจว โข ปน อชานนฺโต อปสฺสนฺโต เอกํเสน อาทาย โวหเรยฺยํ ‘อิทเมว สจฺจํ โมฆมญฺญนฺ’ติ, น เมตํ อสฺส ปติรูปํ.\csromanspacer{} เย โข เต โภนฺโต สมณพฺราหฺมณา เอวํวาทิโน เอวํทิฏฺฐิโน ‘นตฺถิ สพฺพโส ภวนิโรโธ’ติ, สเจ เตสํ \csromanfolio{74}ภวตํ สมณพฺราหฺมณานํ สจฺจํ วจนํ, ฐานเมตํ วิชฺชติ, เย เต เทวา อรูปิโน สญฺญามยา, อปณฺณกํ เม ตตฺรูปปตฺติ ภวิสฺสติ.\csromanspacer{} เย ปน เต โภนฺโต สมณพฺราหฺมณา เอวํวาทิโน เอวํทิฏฺฐิโน ‘อตฺถิ สพฺพโส ภวนิโรโธ’ติ, สเจ เตสํ ภวตํ สมณพฺราหฺมณานํ สจฺจํ วจนํ, ฐานเมตํ วิชฺชติ, ยํ ทิฏฺเฐว ธมฺเม ปรินิพฺพายิสฺสามิ.\csromanspacer{} เย โข เต โภนฺโต สมณพฺราหฺมณา เอวํวาทิโน เอวํทิฏฺฐิโน ‘นตฺถิ สพฺพโส ภวนิโรโธ’ติ, เตสมยํ ทิฏฺฐิ สาราคาย\footnote{สราคาย (สฺยา, กํ)} สนฺติเก สํโยคาย สนฺติเก อภินนฺทนาย สนฺติเก อชฺโฌสานาย สนฺติเก อุปาทานาย สนฺติเก.\csromanspacer{} เย ปน เต โภนฺโต สมณพฺราหฺมณา เอวํวาทิโน เอวํทิฏฺฐิโน ‘อตฺถิ สพฺพโส ภวนิโรโธ’ติ, เตสมยํ ทิฏฺฐิ อสาราคาย สนฺติเก อสํโยคาย สนฺติเก อนภินนฺทนาย สนฺติเก อนชฺโฌสานาย สนฺติเก อนุปาทานาย สนฺติเก”ติ.\csromanspacer{} โส อิติ ปฏิสงฺขาย ภวานํเยว นิพฺพิทาย วิราคาย นิโรธาย ปฏิปนฺโน โหติ.}

\proseitem{105}{จตฺตาโรเม คหปตโย ปุคฺคลา สนฺโต สํวิชฺชมานา โลกสฺมิํ.\csromanspacer{} กตเม จตฺตาโร.\csromanspacer{} อิธ คหปตโย เอกจฺโจ ปุคฺคโล อตฺตนฺตโป โหติ อตฺตปริตาปนานุโยคมนุยุตฺโต, อิธ คหปตโย เอกจฺโจ ปุคฺคโล ปรนฺตโป โหติ ปรปริตาปนานุโยคมนุยุตฺโต, อิธ คหปตโย เอกจฺโจ ปุคฺคโล อตฺตนฺตโป จ โหติ อตฺตปริตาปนานุโยคมนุยุตฺโต ปรนฺตโป จ ปรปริตาปนานุโยคมนุยุตฺโต, อิธ คหปตโย เอกจฺโจ ปุคฺคโล เนวตฺตนฺตโป โหติ นาตฺตปริตาปนานุโยคมนุยุตฺโต น ปรนฺตโป น ปรปริตาปนานุโยคมนุยุตฺโต, โส อนตฺตนฺตโป อปรนฺตโป ทิฏฺเฐว ธมฺเม นิจฺฉาโต นิพฺพุโต สีตีภูโต สุขปฺปฏิสํเวที พฺรหฺมภูเตน อตฺตนา วิหรติ.}

\proseitem{106}{กตโม จ คหปตโย ปุคฺคโล อตฺตนฺตโป อตฺตปริตาปนานุโยคมนุยุตฺโต.\csromanspacer{} อิธ คหปตโย เอกจฺโจ ปุคฺคโล อเจลโก โหติ, มุตฺตาจาโร, หตฺถาปเลขโน ฯเปฯ.\footnote{วิตฺถาโร กนฺทรกสุตฺเต (5) ปิฏฺเฐ.} อิติ เอวรูปํ อเนกวิหิตํ กายสฺส อาตาปนปริตาปนานุโยคมนุยุตฺโต วิหรติ.\csromanspacer{} อยํ วุจฺจติ คหปตโย ปุคฺคโล อตฺตนฺตโป อตฺตปริตาปนานุโยคมนุยุตฺโต.}

\vspace{\tipitakaparskip}
\csromancenter{อปณฺณกสุตฺต (60) กตโม จ คหปตโย ปุคฺคโล ปรนฺตโป ปรปริตาปนานุโยคมนุยุตฺโต.\csromanspacer{} อิธ คหปตโย เอกจฺโจ ปุคฺคโล โอรพฺภิโก โหติ สูกริโก ฯเปฯ เย วา ปนญฺเญปิ เกจิ กุรูรกมฺมนฺตา.\csromanspacer{} อยํ วุจฺจติ คหปตโย ปุคฺคโล ปรนฺตโป ปรปริตาปนานุโยคมนุยุตฺโต.}
\prose{\csromanfolio{75}กตโม จ คหปตโย ปุคฺคโล อตฺตนฺตโป จ อตฺตปริตาปนานุโยคมนุยุตฺโต, ปรนฺตโป จ ปรปริตาปนานุโยคมนุยุตฺโต.\csromanspacer{} อิธ คหปตโย เอกจฺโจ ปุคฺคโล ราชา วา โหติ ขตฺติโย มุทฺธาวสิตฺโต ฯเปฯ เตปิ ทณฺฑตชฺชิตา ภยตชฺชิตา อสฺสุมุขา รุทมานา ปริกมฺมานิ กโรนฺติ.\csromanspacer{} อยํ วุจฺจติ คหปตโย ปุคฺคโล อตฺตนฺตโป จ อตฺตปริตาปนานุโยคมนุยุตฺโต ปรนฺตโป จ ปรปริตาปนานุโยคมนุยุตฺโต.}

\prose{กตโม จ คหปตโย ปุคฺคโล เนวตฺตนฺตโป นาตฺตปริตาปนานุโยคมนุยุตฺโต น ปรนฺตโป น ปรปริตาปนานุโยคมนุยุตฺโต, โส อนตฺตนฺตโป อปรนฺตโป ทิฏฺเฐว ธมฺเม นิจฺฉาโต นิพฺพุโต สีตีภูโต สุขปฺปฏิสํเวที พฺรหฺมภูเตน อตฺตนา วิหรติ.\csromanspacer{} อิธ คหปตโย ตถาคโต โลเก อุปฺปชฺชติ อรหํ สมฺมาสมฺพุทฺโธ ฯเปฯ.\csromanspacer{} โส อิเม ปญฺจ นีวรเณ ปหาย เจตโส อุปกฺกิเลเส ปญฺญาย ทุพฺพลีกรเณ, วิวิจฺเจว กาเมหิ วิวิจฺจ อกุสเลหิ ธมฺเมหิ สวิตกฺกํ สวิจารํ วิเวกชํ ปีติสุขํ ปฐมํ ฌานํ อุปสมฺปชฺช วิหรติ.\csromanspacer{} วิตกฺกวิจารานํ วูปสมา อชฺฌตฺตํ สมฺปสาทนํ เจตโส เอโกทิภาวํ อวิตกฺกํ อวิจารํ สมาธิชํ ปีติสุขํ ทุติยํ ฌานํ ฯเปฯ ตติยํ ฌานํ ฯเปฯ จตุตฺถํ ฌานํ อุปสมฺปชฺช วิหรติ.}

\csromanmatikaanchor{mtk.12}
\csromantocmarkref{subsection}{อุทฺทานคาถาโย}{mtk.12}
\bookmark[dest={mtk.12},level=2]{อุทฺทานคาถาโย}
\prose{โส เอวํ สมาหิเต จิตฺเต ปริสุทฺเธ ปริโยทาเต อนงฺคเณ วิคตูปกฺกิเลเส มุทุภูเต กมฺมนิเย ฐิเต อาเนญฺชปฺปตฺเต ปุพฺเพนิวาสานุสฺสติญาณาย จิตฺตํ อภินินฺนาเมติ, โส อเนกวิหิตํ ปุพฺเพนิวาสํ อนุสฺสรติ.\csromanspacer{} เสยฺยถิทํ, เอกมฺปิ ชาติํ เทฺวปิ ชาติโย ฯเปฯ อิติ สาการํ ส-อุทฺเทสํ อเนกวิหิตํ ปุพฺเพนิวาสํ อนุสฺสรติ.\csromanspacer{} โส เอวํ สมาหิเต จิตฺเต ปริสุทฺเธ ปริโยทาเต อนงฺคเณ วิคตูปกฺกิเลเส มุทุภูเต กมฺมนิเย ฐิเต อาเนญฺชปฺปตฺเต สตฺตานํ จุตูปปาตญาณาย จิตฺตํ อภินินฺนาเมติ.\csromanspacer{} โส ทิพฺเพน จกฺขุนา วิสุทฺเธน อติกฺกนฺตมานุสเกน สตฺเต ปสฺสติ จวมาเน อุปปชฺชมาเน หีเน ปณีเต สุวณฺเณ ทุพฺพณฺเณ สุคเต ทุคฺคเต ฯเปฯ ยถากมฺมูปเค สตฺเต ปชานาติ.\csromanspacer{} โส เอวํ \csromanfolio{76}สมาหิเต จิตฺเต ปริสุทฺเธ ปริโยทาเต อนงฺคเณ วิคตูปกฺกิเลเส มุทุภูเต กมฺมนิเย ฐิเต อาเนญฺชปฺปตฺเต อาสวานํ ขยญาณาย จิตฺตํ อภินินฺนาเมติ.\csromanspacer{} โส อิทํ ทุกฺขนฺติ ยถาภูตํ ปชานาติ ฯเปฯ อยํ อาสวนิโรธคามินี ปฏิปทาติ ยถาภูตํ ปชานาติ.\csromanspacer{} ตสฺส เอวํ ชานโต เอวํ ปสฺสโต กามาสวาปิ จิตฺตํ วิมุจฺจติ, ภวาสวาปิ จิตฺตํ วิมุจฺจติ, อวิชฺชาสวาปิ จิตฺตํ วิมุจฺจติ, วิมุตฺตสฺมิํ วิมุตฺตมิติ ญาณํ โหติ, “ขีณา ชาติ, วุสิตํ พฺรหฺมจริยํ, กตํ กรณียํ, นาปรํ อิตฺถตฺตายา”ติ ปชานาติ.\csromanspacer{} อยํ วุจฺจติ คหปตโย ปุคฺคโล เนวตฺตนฺตโป นาตฺตปริตาปนานุโยคมนุยุตฺโต น ปรนฺตโป น ปรปริตาปนานุโยคมนุยุตฺโต, โส อนตฺตนฺตโป อปรนฺตโป ทิฏฺเฐว ธมฺเม นิจฺฉาโต นิพฺพุโต สีตีภูโต สุขปฺปฏิสํเวที พฺรหฺมภูเตน อตฺตนา วิหรตีติ.}

\prose{เอวํ วุตฺเต สาเลยฺยกา พฺราหฺมณคหปติกา ภควนฺตํ เอตทโวจุํ “อภิกฺกนฺตํ โภ โคตม, อภิกฺกนฺตํ โภ โคตม, เสยฺยถาปิ โภ โคตม นิกฺกุชฺชิตํ วา อุกฺกุชฺเชยฺย, ปฏิจฺฉนฺนํ วา วิวเรยฺย, มูฬฺหสฺส วา มคฺคํ อาจิกฺเขยฺย, อนฺธกาเร วา เตลปชฺโชตํ ธาเรยฺย ‘จกฺขุมนฺโต รูปานิ ทกฺขนฺตี’ติ.\csromanspacer{} เอวเมวํ โภตา โคตเมน อเนกปริยาเยน ธมฺโม ปกาสิโต, เอเต มยํ ภวนฺตํ โคตมํ สรณํ คจฺฉาม ธมฺมญฺจ ภิกฺขุสํฆญฺจ, อุปาสเก โน ภวํ โคตโม ธาเรตุ อชฺชตคฺเค ปาณุเปตํ สรณํ คเต”ติ.}

\nitthitam{อปณฺณกสุตฺตํ นิฏฺฐิตํ ทสมํ.}
\nitthitam{\textbf{คหปติวคฺโค นิฏฺฐิโต ปฐโม.}}
\titlehead{\textbf{ตสฺสุทฺทานํ}}
\csromangathameasure{กนฺทรนาครเสขวโต จ, \\ โปตลิโย ปุน ชีวกภจฺโจ. \\ อุปาลิทมโถ กุกฺกุร-อภโย, \\ พหุเวทนียาปณฺณกโต ทสโม.}
\csromangathagroup{\csromangathastanza{กนฺทรนาครเสขวโต จ, \\ โปตลิโย ปุน ชีวกภจฺโจ. \\ อุปาลิทมโถ กุกฺกุร-อภโย, \\ พหุเวทนียาปณฺณกโต ทสโม.}}

\csromanmatikaanchor{mtk.13}
\csromantocmarknopage{chapter}{2. ภิกฺขุวคฺค}{mtk.13}
\bookmark[dest={mtk.13},level=0]{2. ภิกฺขุวคฺค}
\chapterheadpageread{2. ภิกฺขุวคฺค}
\csromanmatikaanchor{mtk.14}
\csromantocmarkref{subsection}{1. อมฺพลฏฺฐิกราหุโลวาทสุตฺต}{mtk.14}
\bookmark[dest={mtk.14},level=2]{1. อมฺพลฏฺฐิกราหุโลวาทสุตฺต}
\csromanheader{2}{1. อมฺพลฏฺฐิกราหุโลวาทสุตฺต}
\proseitem{107}{\csromanfolio{77}เอวํ เม สุตํ\csromandash{}เอกํ สมยํ ภควา ราชคเห วิหรติ เวฬุวเน กลนฺทกนิวาเป.\csromanspacer{} เตน โข ปน สมเยน อายสฺมา ราหุโล อมฺพลฏฺฐิกายํ วิหรติ.\csromanspacer{} อถ โข ภควา สายนฺหสมยํ ปฏิสลฺลานา วุฏฺฐิโต เยน อมฺพลฏฺฐิกา, เยนายสฺมา ราหุโล เตนุปสงฺกมิ, อทฺทสา โข อายสฺมา ราหุโล ภควนฺตํ ทูรโตว อาคจฺฉนฺตํ, ทิสฺวาน อาสนํ ปญฺญาเปสิ อุทกญฺจ ปาทานํ.\csromanspacer{} นิสีทิ ภควา ปญฺญตฺเต อาสเน, นิสชฺช ปาเท ปกฺขาเลสิ.\csromanspacer{} อายสฺมาปิ โข ราหุโล ภควนฺตํ อภิวาเทตฺวา เอกมนฺตํ นิสีทิ.}

\proseitem{108}{อถ โข ภควา ปริตฺตํ อุทกาวเสสํ อุทกาธาเน ฐเปตฺวา อายสฺมนฺตํ ราหุลํ อามนฺเตสิ “ปสฺสสิ โน ตฺวํ ราหุล อิมํ ปริตฺตํ อุทกาวเสสํ อุทกาธาเน ฐปิตนฺ”ติ.\csromanspacer{} เอวํ ภนฺเต.\csromanspacer{} เอวํ ปริตฺตกํ โข ราหุล เตสํ สามญฺญํ, เยสํ นตฺถิ สมฺปชานมุสาวาเท ลชฺชาติ.\csromanspacer{} อถ โข ภควา ปริตฺตํ อุทกาวเสสํ ฉฑฺเฑตฺวา อายสฺมนฺตํ ราหุลํ อามนฺเตสิ “ปสฺสสิ โน ตฺวํ ราหุล ปริตฺตํ อุทกาวเสสํ ฉฑฺฑิตนฺ”ติ.\csromanspacer{} เอวํ ภนฺเต.\csromanspacer{} เอวํ ฉฑฺฑิตํ โข ราหุล เตสํ สามญฺญํ, เยสํ นตฺถิ สมฺปชานมุสาวาเท ลชฺชาติ.\csromanspacer{} อถ โข ภควา ตํ อุทกาธานํ นิกฺกุชฺชิตฺวา อายสฺมนฺตํ ราหุลํ อามนฺเตสิ “ปสฺสสิ โน ตฺวํ ราหุล อิมํ อุทกาธานํ นิกฺกุชฺชิตนฺ”ติ.\csromanspacer{} เอวํ ภนฺเต.\csromanspacer{} เอวํ นิกฺกุชฺชิตํ โข ราหุล เตสํ สามญฺญํ, เยสํ นตฺถิ สมฺปชานมุสาวาเท ลชฺชาติ.\csromanspacer{} อถ โข ภควา ตํ อุทกาธานํ อุกฺกุชฺชิตฺวา อายสฺมนฺตํ ราหุลํ อามนฺเตสิ “ปสฺสสิ โน ตฺวํ ราหุล อิมํ อุทกาธานํ ริตฺตํ ตุจฺฉนฺ”ติ.\csromanspacer{} เอวํ ภนฺเต.\csromanspacer{} เอวํ ริตฺตํ ตุจฺฉํ โข ราหุล เตสํ สามญฺญํ, เยสํ นตฺถิ สมฺปชานมุสาวาเท ลชฺชาติ.\csromanspacer{} เสยฺยถาปิ ราหุล รญฺโญ นาโค อีสาทนฺโต อุรูฬฺหวา\footnote{อุพฺพูฬฺหวา (สี, อิ)} อภิชาโต สงฺคามาวจโร สงฺคามคโต ปุริเมหิปิ ปาเทหิ กมฺมํ กโรติ, ปจฺฉิเมหิปิ ปาเทหิ กมฺมํ กโรติ, ปุริเมนปิ กาเยน กมฺมํ กโรติ, ปจฺฉิเมนปิ กาเยน กมฺมํ กโรติ, สีเสนปิ กมฺมํ กโรติ, \csromanfolio{78}กณฺเณหิปิ กมฺมํ กโรติ, ทนฺเตหิปิ กมฺมํ กโรติ, นงฺคุฏฺเฐนปิ กมฺมํ กโรติ, รกฺขเตว โสณฺฑํ.\csromanspacer{} ตตฺถ หตฺถาโรหสฺส เอวํ โหติ “อยํ โข รญฺโญ นาโค อีสาทนฺโต อุรูฬฺหวา อภิชาโต สงฺคามาวจโร สงฺคามคโต ปุริเมหิปิ ปาเทหิ กมฺมํ กโรติ, ปจฺฉิเมหิปิ ปาเทหิ กมฺมํ กโรติ ฯเปฯ นงฺคุฏฺเฐนปิ กมฺมํ กโรติ, รกฺขเตว โสณฺฑํ.\csromanspacer{} อปริจฺจตฺตํ โข รญฺโญ นาคสฺส ชีวิตนฺ”ติ.\csromanspacer{} ยโต โข ราหุล รญฺโญ นาโค อีสาทนฺโต อุรูฬฺหวา อภิชาโต สงฺคามาวจโร สงฺคามคโต ปุริเมหิปิ ปาเทหิ กมฺมํ กโรติ, ปจฺฉิเมหิปิ ปาเทหิ กมฺมํ กโรติ ฯเปฯ นงฺคุฏฺเฐนปิ กมฺมํ กโรติ, โสณฺฑายปิ กมฺมํ กโรติ.\csromanspacer{} ตตฺถ หตฺถาโรหสฺส เอวํ โหติ “อยํ โข รญฺโญ นาโค อีสาทนฺโต อุรูฬฺหวา อภิชาโต สงฺคามาวจโร สงฺคามคโต ปุริเมหิปิ ปาเทหิ กมฺมํ กโรติ, ปจฺฉิเมหิปิ ปาเทหิ กมฺมํ กโรติ, ปุริเมนปิ กาเยน กมฺมํ กโรติ, ปจฺฉิเมนปิ กาเยน กมฺมํ กโรติ, สีเสนปิ กมฺมํ กโรติ, กณฺเณหิปิ กมฺมํ กโรติ, ทนฺเตหิปิ กมฺมํ กโรติ, นงฺคุฏฺเฐนปิ กมฺมํ กโรติ, โสณฺฑายปิ กมฺมํ กโรติ, ปริจฺจตฺตํ โข รญฺโญ นาคสฺส ชีวิตํ, นตฺถิ ทานิ กิญฺจิ รญฺโญ นาคสฺส อกรณียนฺ”ติ.\csromanspacer{} เอวเมว โข ราหุล ยสฺส กสฺสจิ สมฺปชานมุสาวาเท นตฺถิ ลชฺชา.\csromanspacer{} นาหํ ตสฺส กิญฺจิ ปาปํ อกรณียนฺติ วทามิ.\csromanspacer{} ตสฺมาติห เต ราหุล “หสฺสาปิ น มุสา ภณิสฺสามี”ติ เอวํ หิ เต ราหุล สิกฺขิตพฺพํ.}

\proseitem{109}{ตํ กิํ มญฺญสิ ราหุล, กิมตฺถิโย อาทาโสติ.\csromanspacer{} ปจฺจเวกฺขณตฺโถ ภนฺเตติ.\csromanspacer{} เอวเมว โข ราหุล ปจฺจเวกฺขิตฺวา ปจฺจเวกฺขิตฺวา กาเยน กมฺมํ กตฺตพฺพํ, ปจฺจเวกฺขิตฺวา ปจฺจเวกฺขิตฺวา วาจาย กมฺมํ กตฺตพฺพํ, ปจฺจเวกฺขิตฺวา ปจฺจเวกฺขิตฺวา มนสา กมฺมํ กตฺตพฺพํ.\csromanspacer{} ยเทว ตฺวํ ราหุล กาเยน กมฺมํ กตฺตุกาโม อโหสิ, ตเทว เต กายกมฺมํ ปจฺจเวกฺขิตพฺพํ “ยํ นุ โข อหํ อิทํ กาเยน กมฺมํ กตฺตุกาโม, อิทํ เม กายกมฺมํ อตฺตพฺยาพาธายปิ สํวตฺเตยฺย, ปรพฺยาพาธายปิ สํวตฺเตยฺย, อุภยพฺยาพาธายปิ สํวตฺเตยฺย.\csromanspacer{} อกุสลํ อิทํ กายกมฺมํ ทุกฺขุทฺรยํ\footnote{ทุกฺขุนฺทฺรยํ, ทุกฺขุทยํ (ก)} ทุกฺขวิปากนฺ”ติ.\csromanspacer{} สเจ ตฺวํ ราหุล ปจฺจเวกฺขมาโน เอวํ ชาเนยฺยาสิ “ยํ โข อหํ อิทํ กาเยน กมฺมํ กตฺตุกาโม, อิทํ เม กายกมฺมํ อตฺตพฺยาพาธายปิ สํวตฺเตยฺย, ปรพฺยาพาธายปิ สํวตฺเตยฺย,}

\vspace{\tipitakaparskip}
\csromancenter{อมฺพลฏฺฐิกราหุโลวาทสุตฺต (61) อุภยพฺยาพาธายปิ สํวตฺเตยฺย, อกุสลํ อิทํ กายกมฺมํ ทุกฺขุทฺรยํ ทุกฺขปิปากนฺ”ติ, เอวรูปํ เต ราหุล กาเยน กมฺมํ สสกฺกํ น กรณียํ\footnote{สํสกฺกํ น จ กรณียํ (ก)}.\csromanspacer{} สเจ ปน ตฺวํ ราหุล ปจฺจเวกฺขมาโน เอวํ ชาเนยฺยาสิ “ยํ โข อหํ อิทํ กาเยน กมฺมํ กตฺตุกาโม, อิทํ เม กายกมฺมํ เนวตฺตพฺยาพาธายปิ สํวตฺเตยฺย, น ปรพฺยาพาธายปิ สํวตฺเตยฺย, น อุภยพฺยาพาธายปิ สํวตฺเตยฺย, กุสลํ อิทํ กายกมฺมํ สุขุทฺรยํ สุขวิปากนฺ”ติ, เอวรูปํ เต ราหุล กาเยน กมฺมํ กรณียํ.}
\prose{\csromanfolio{79}กโรนฺเตนปิ เต ราหุล กาเยน กมฺมํ ตเทว เต กายกมฺมํ ปจฺจเวกฺขิตพฺพํ “ยํ นุ โข อหํ อิทํ กาเยน กมฺมํ กโรมิ, อิทํ เม กายกมฺมํ อตฺตพฺยาพาธายปิ สํวตฺตติ, ปรพฺยาพาธายปิ สํวตฺตติ, อุภยพฺยาพาธายปิ สํวตฺตติ, อกุสลํ อิทํ กายกมฺมํ ทุกฺขุทฺรยํ ทุกฺขวิปากนฺ”ติ.\csromanspacer{} สเจ ปน ตฺวํ ราหุล ปจฺจเวกฺขมาโน เอวํ ชาเนยฺยาสิ “ยํ โข อหํ อิทํ กาเยน กมฺมํ กโรมิ, อิทํ เม กายกมฺมํ อตฺตพฺยาพาธายปิ สํวตฺตติ, ปรพฺยาพาธายปิ สํวตฺตติ, อุภยพฺยาพาธายปิ สํวตฺตติ, อกุสลํ อิทํ กายกมฺมํ ทุกฺขุทฺรยํ ทุกฺขวิปากนฺ”ติ, ปฏิสํหเรยฺยาสิ ตฺวํ ราหุล เอวรูปํ กายกมฺมํ.\csromanspacer{} สเจ ปน ตฺวํ ราหุล ปจฺจเวกฺขมาโน เอวํ ชาเนยฺยาสิ “ยํ โข อหํ อิทํ กาเยน กมฺมํ กโรมิ, อิทํ เม กายกมฺมํ เนวตฺตพฺยาพาธายปิ สํวตฺตติ, น ปรพฺยาพาธายปิ สํวตฺตติ น อุภยพฺยาพาธายปิ สํวตฺตติ, กุสลํ อิทํ กายกมฺมํ สุขุทฺรยํ สุขวิปากนฺ”ติ, อนุปทชฺเชยฺยาสิ ตฺวํ ราหุล เอวรูปํ กายกมฺมํ.}

\prose{กตฺวาปิ เต ราหุล กาเยน กมฺมํ ตเทว เต กายกมฺมํ ปจฺจเวกฺขิตพฺพํ “ยํ นุ โข อหํ อิทํ กาเยน กมฺมํ อกาสิํ, อิทํ เม กายกมฺมํ อตฺตพฺยาพาธายปิ สํวตฺตติ\footnote{สํวตฺติ (อิ)}, ปรพฺยาพาธายปิ สํวตฺตติ, อุภยพฺยาพาธายปิ สํวตฺตติ, อกุสลํ อิทํ กายกมฺมํ ทุกฺขุทฺรยํ ทุกฺขวิปากนฺ”ติ.\csromanspacer{} สเจ โข ตฺวํ ราหุล ปจฺจเวกฺขมาโน เอวํ ชาเนยฺยาสิ “ยํ โข อหํ อิทํ กาเยน กมฺมํ อกาสิํ, อิทํ เม กายกมฺมํ อตฺตพฺยาพาธายปิ สํวตฺตติ, ปรพฺยาพาธายปิ สํวตฺตติ, อุภยพฺยาพาธายปิ สํวตฺตติ, อกุสลํ อิทํ กายกมฺมํ ทุกฺขุทฺรยํ ทุกฺขวิปากนฺ”ติ, เอวรูปํ เต ราหุล กายกมฺมํ สตฺถริ วา วิญฺญูสุ วา \csromanfolio{80}สพฺรหฺมจารีสุ เทเสตพฺพํ วิวริตพฺพํ อุตฺตานีกาตพฺพํ, เทเสตฺวา วิวริตฺวา อุตฺตานีกตฺวา อายติํ สํวรํ อาปชฺชิตพฺพํ.\csromanspacer{} สเจ ปน ตฺวํ ราหุล ปจฺจเวกฺขมาโน เอวํ ชาเนยฺยาสิ “ยํ โข อหํ อิทํ กาเยน กมฺมํ อกาสิํ, อิทํ เม กายกมฺมํ เนวตฺตพฺยาพาธายปิ สํวตฺตติ, น ปรพฺยาพาธายปิ สํวตฺตติ, น อุภยพฺยาพาธายปิ สํวตฺตติ.\csromanspacer{} กุสลํ อิทํ กายกมฺมํ สุขุทฺรยํ สุขวิปากนฺ”ติ, เตเนว ตฺวํ ราหุล ปีติปาโมชฺเชน วิหเรยฺยาสิ อโหรตฺตานุสิกฺขี กุสเลสุ ธมฺเมสุ.}

\proseitem{110}{ยเทว ตฺวํ ราหุล วาจาย กมฺมํ กตฺตุกาโม อโหสิ, ตเทว เต วจีกมฺมํ ปจฺจเวกฺขิตพฺพํ “ยํ นุ โข อหํ อิทํ วาจาย กมฺมํ กตฺตุกาโม, อิทํ เม วจีกมฺมํ อตฺตพฺยาพาธายปิ สํวตฺเตยฺย, ปรพฺยาพาธายปิ สํวตฺเตยฺย, อุภยพฺยาพาธายปิ สํวตฺเตยฺย, อกุสลํ อิทํ วจีกมฺมํ ทุกฺขุทฺรยํ ทุกฺขวิปากนฺ”ติ.\csromanspacer{} สเจ ตฺวํ ราหุล ปจฺจเวกฺขมาโน เอวํ ชาเนยฺยาสิ “ยํ โข อหํ อิทํ วาจาย กมฺมํ กตฺตุกาโม, อิทํ เม วจีกมฺมํ อตฺตพฺยาพาธายปิ สํวตฺเตยฺย, ปรพฺยาพาธายปิ สํวตฺเตยฺย, อุภยพฺยาพาธายปิ สํวตฺเตยฺย, อกุสลํ อิทํ วจีกมฺมํ ทุกฺขุทฺรยํ ทุกฺขวิปากนฺ”ติ, เอวรูปํ เต ราหุล วาจาย กมฺมํ สสกฺกํ น กรณียํ.\csromanspacer{} สเจ ปน ตฺวํ ราหุล ปจฺจเวกฺขมาโน เอวํ ชาเนยฺยาสิ “ยํ โข อหํ อิทํ วาจาย กมฺมํ กตฺตุกาโม, อิทํ เม วจีกมฺมํ เนวตฺตพฺยาพาธายปิ สํวตฺเตยฺย, น ปรพฺยาพาธายปิ สํวตฺเตยฺย, กุสลํ อิทํ วจีกมฺมํ สุขุทฺรยํ สุขวิปากนฺ”ติ, เอวรูปํ เต ราหุล วาจาย กมฺมํ กรณียํ.}

\prose{กโรนฺเตนปิ เต ราหุล วาจาย กมฺมํ ตเทว เต วจีกมฺมํ ปจฺจเวกฺขิตพฺพํ “ยํ นุ โข อหํ อิทํ วาจาย กมฺมํ กโรมิ, อิทํ เม วจีกมฺมํ อตฺตพฺยาพาธายปิ สํวตฺตติ, ปรพฺยาพาธายปิ สํวตฺตติ, อุภยพฺยาพาธายปิ สํวตฺตติ, อกุสลํ อิทํ วจีกมฺมํ ทุกฺขุทฺรยํ ทุกฺขวิปากนฺ”ติ.\csromanspacer{} สเจ ปน ตฺวํ ราหุล ปจฺจเวกฺขมาโน เอวํ ชาเนยฺยาสิ “ยํ โข อหํ อิทํ วาจาย กมฺมํ กโรมิ, อิทํ เม วจีกมฺมํ อตฺตพฺยาพาธายปิ สํวตฺตติ, ปรพฺยาพาธายปิ สํวตฺตติ, อุภยพฺยาพาธายปิ สํวตฺตติ, อกุสลํ อิทํ วจีกมฺมํ ทุกฺขุทฺรยํ ทุกฺขวิปากนฺ”ติ, ปฏิสํหเรยฺยาสิ ตฺวํ ราหุล เอวรูปํ วจีกมฺมํ.\csromanspacer{} สเจ ปน ตฺวํ ราหุล ปจฺจเวกฺขมาโน เอวํ ชาเนยฺยาสิ “ยํ โข อหํ อิทํ วาจาย กมฺมํ กโรมิ, อิทํ เม}

\vspace{\tipitakaparskip}
\csromancenter{อมฺพลฏฺฐิกราหุโลวาทสุตฺต (61) วจีกมฺมํ เนวตฺตพฺยาพาธายปิ สํวตฺตติ, น ปรพฺยาพาธายปิ สํวตฺตติ, น อุภยพฺยาพาธายปิ สํวตฺตติ, กุสลํ อิทํ วจีกมฺมํ สุขุทฺรยํ สุขวิปากนฺ”ติ, อนุปทชฺเชยฺยาสิ ตฺวํ ราหุล เอวรูปํ วจีกมฺมํ.}
\prose{\csromanfolio{81}กตฺวาปิ เต ราหุล วาจาย กมฺมํ, ตเทว เต วจีกมฺมํ ปจฺจเวกฺขิตพฺพํ “ยํ นุ โข อหํ อิทํ วาจาย กมฺมํ อกาสิํ, อิทํ เม วจีกมฺมํ อตฺตพฺยาพาธายปิ สํวตฺตติ\footnote{สํวตฺติ (สี, อิ)}, ปรพฺยาพาธายปิ สํวตฺตติ, อุภยพฺยาพาธายปิ สํวตฺตติ, อกุสลํ อิทํ วจีกมฺมํ ทุกฺขุทฺรยํ ทุกฺขวิปากนฺ”ติ.\csromanspacer{} สเจ โข ตฺวํ ราหุล ปจฺจเวกฺขมาโน เอวํ ชาเนยฺยาสิ “ยํ โข อหํ อิทํ วาจาย กมฺมํ อกาสิํ, อิทํ เม วจีกมฺมํ อตฺตพฺยาพาธายปิ สํวตฺตติ, ปรพฺยาพาธายปิ สํวตฺตติ, อุภยพฺยาพาธายปิ สํวตฺตติ, อกุสลํ อิทํ วจีกมฺมํ ทุกฺขุทฺรยํ ทุกฺขวิปากนฺ”ติ, เอวรูปํ เต ราหุล วจีกมฺมํ สตฺถริ วา วิญฺญูสุ วา สพฺรหฺมจารีสุ เทเสตพฺพํ วิวริตพฺพํ อุตฺตานีกตฺตพฺพํ, เทเสตฺวา วิวริตฺวา อุตฺตานีกตฺวา อายติํ สํวรํ อาปชฺชิตพฺพํ.\csromanspacer{} สเจ ปน ตฺวํ ราหุล ปจฺจเวกฺขมาโน เอวํ ชาเนยฺยาสิ “ยํ โข อหํ อิทํ วาจาย กมฺมํ อกาสิํ, อิทํ เม วจีกมฺมํ เนวตฺตพฺยาพาธายปิ สํวตฺตติ, น ปรพฺยาพาธายปิ สํวตฺตติ, น อุภยพฺยาพาธายปิ สํวตฺตติ, กุสลํ อิทํ วจีกมฺมํ สุขุทฺรยํ สุขวิปากนฺ”ติ, เตเนว ตฺวํ ราหุล ปีติปาโมชฺเชน วิหเรยฺยาสิ อโหรตฺตานุสิกฺขี กุสเลสุ ธมฺเมสุ.}

\proseitem{111}{ยเทว ตฺวํ ราหุล มนสา กมฺมํ กตฺตุกาโม อโหสิ, ตเทว เต มโนกมฺมํ ปจฺจเวกฺขิตพฺพํ “ยํ นุ โข อหํ อิทํ มนสา กมฺมํ กตฺตุกาโม, อิทํ เม มโนกมฺมํ อตฺตพฺยาพาธยปิ สํวตฺเตยฺย, ปรพฺยาพาธายปิ สํวตฺเตยฺย, อุภยพฺยาพาธายปิ สํวตฺเตยฺย, อกุสลํ อิทํ มโนกมฺมํ ทุกฺขุทฺรยํ ทุกฺขวิปากนฺ”ติ.\csromanspacer{} สเจ ตฺวํ ราหุล ปจฺจเวกฺขมาโน เอวํ ชาเนยฺยาสิ “ยํ โข อหํ อิทํ มนสา กมฺมํ กตฺตุกาโม, อิทํ เม มโนกมฺมํ อตฺตพฺยาพาธายปิ สํวตฺเตยฺย, ปรพฺยาพาธายปิ สํวตฺเตยฺย, อุภยพฺยาพาธายปิ สํวตฺเตยฺย, อกุสลํ อิทํ มโนกมฺมํ ทุกฺขุทฺรยํ ทุกฺขวิปากนฺ”ติ, เอวรูปํ เต ราหุล มนสา กมฺมํ สสกฺกํ น กรณียํ.\csromanspacer{} สเจ ปน ตฺวํ ราหุล ปจฺจเวกฺขมาโน เอวํ ชาเนยฺยาสิ “ยํ โข \csromanfolio{82}อหํ อิทํ มนสา กมฺมํ กตฺตุกาโม, อิทํ เม มโนกมฺมํ เนวตฺตพฺยาพาธายปิ สํวตฺเตยฺย, น ปรพฺยาพาธายปิ สํวตฺเตยฺย, น อุภยพฺยาพาธายปิ สํวตฺเตยฺย, กุสลํ อิทํ มโนกมฺมํ สุขุทฺรยํ สุขวิปากนฺ”ติ, เอวรูปํ เต ราหุล มนสา กมฺมํ กรณียํ.}

\prose{กโรนฺเตนปิ เต ราหุล มนสา กมฺมํ ตเทว เต มโนกมฺมํ ปจฺจเวกฺขิตพฺพํ “ยํ นุ โข อหํ อิทํ มนสา กมฺมํ กโรมิ, อิทํ เม มโนกมฺมํ อตฺตพฺยาพาธายปิ สํวตฺตติ, ปรพฺยาพาธายปิ สํวตฺตติ, อุภยพฺยาพาธายปิ สํวตฺตติ, อกุสลํ อิทํ มโนกมฺมํ ทุกฺขุทฺรยํ ทุกฺขวิปากนฺ”ติ.\csromanspacer{} สเจ ปน ตฺวํ ราหุล ปจฺจเวกฺขมาโน เอวํ ชาเนยฺยาสิ “ยํ โข อหํ อิทํ มนสา กมฺมํ กโรมิ, อิทํ เม มโนกมฺมํ อตฺตพฺยาพาธายปิ สํวตฺตติ, ปรพฺยาพาธายปิ สํวตฺตติ, อุภยพฺยาพาธายปิ สํวตฺตติ, อกุสลํ อิทํ มโนกมฺมํ ทุกฺขุทฺรยํ ทุกฺขวิปากนฺ”ติ, ปฏิสํหเรยฺยาสิ ตฺวํ ราหุล เอวรูปํ มโนกมฺมํ.\csromanspacer{} สเจ ปน ตฺวํ ราหุล ปจฺจเวกฺขมาโน เอวํ ชาเนยฺยาสิ “ยํ โข อหํ อิทํ มนสา กมฺมํ กโรมิ, อิทํ เม มโนกมฺมํ เนวตฺตพฺยาพาธายปิ สํวตฺตติ, น ปรพฺยาพาธายปิ สํวตฺตติ, น อุภยพฺยาพาธายปิ สํวตฺตติ, กุสลํ อิทํ มโนกมฺมํ สุขุทฺรยํ สุขวิปากนฺ”ติ, อนุปทชฺเชยฺยาสิ ตฺวํ ราหุล เอวรูปํ มโนกมฺมํ.}

\prose{กตฺวาปิ เย ราหุล มนสา กมฺมํ, ตเทว เต มโนกมฺมํ ปจฺจเวกฺขิตพฺพํ “ยํ นุ โข อหํ อิทํ มนสา กมฺมํ อกาสิํ, อิทํ เม มโนกมฺมํ อตฺตพฺยาพาธายปิ สํวตฺตติ\footnote{สํวตฺติ (สี, อิ)}, ปรพฺยาพาธายปิ สํวตฺตติ, อุภยพฺยาพาธายปิ สํวตฺตติ, อกุสลํ อิทํ มโนกมฺมํ ทุกฺขุทฺรยํ ทุกฺขวิปากนฺ”ติ.\csromanspacer{} สเจ โข ตฺวํ ราหุล ปจฺจเวกฺขมาโน เอวํ ชาเนยฺยาสิ “ยํ โข อหํ อิทํ มนสา กมฺมํ อกาสิํ, อิทํ เม มโนกมฺมํ อตฺตพฺยาพาธายปิ สํวตฺตติ, ปรพฺยาพาธายปิ สํวตฺตติ, อุภยพฺยาพาธายปิ สํวตฺตติ, อกุสลํ อิทํ มโนกมฺมํ ทุกฺขุทฺรยํ ทุกฺขวิปากนฺ”ติ, เอวรูปํ ปน\footnote{เอวรูเป (สี, อิ), เอวรูเป ปน (สฺยา, กํ)} เต ราหุล มโน- กมฺมํ\footnote{มโนกมฺเม (สี, สฺยา, กํ, อิ)} อฏฺฏียิตพฺพํ หรายิตพฺพํ ชิคุจฺฉิตพฺพํ, อฏฺฏียิตฺวา หรายิตฺวา ชิคุจฺฉิตฺวา อายติํ สํวรํ อาปชฺชิตพฺพํ.\csromanspacer{} สเจ ปน ตฺวํ ราหุล ปจฺจเวกฺขมาโน เอวํ ชาเนยฺยาสิ “ยํ โข อหํ อิทํ มนสา กมฺมํ}

\vspace{\tipitakaparskip}
\csromancenter{มหาราหุโลวาทสุตฺต (62) อกาสิํ, อิทํ เม มโนกมฺมํ เนวตฺตพฺยาพาธายปิ สํวตฺตติ, น ปรพฺยาพาธายปิ สํวตฺตปิ, น อุภยพฺยาพาธายปิ สํวตฺตติ, กุสลํ อิทํ มโนกมฺมํ สุขุทฺรยํ สุขวิปากนฺ”ติ, เตเนว ตฺวํ ราหุล ปีติปาโมชฺเชน วิหเรยฺยาสิ อโหรตฺตานุสิกฺขี กุสเลสุ ธมฺเมสุ.}
\proseitem{112}{\csromanfolio{83}เย หิ เกจิ ราหุล อตีตมทฺธานํ สมณา วา พฺราหฺมณา วา กายกมฺมํ ปริโสเธสุํ, วจีกมฺมํ ปริโสเธสุํ, มโนกมฺมํ ปริโสเธสุํ, สพฺเพ เต เอวเมตํ ปจฺจเวกฺขิตฺวา ปจฺจเวกฺขิตฺวา กายกมฺมํ ปริโสเธสุํ, ปจฺจเวกฺขิตฺวา ปจฺจเวกฺขิตฺวา วจีกมฺมํ ปริโสเธสุํ, ปจฺจเวกฺขิตฺวา ปจฺจเวกฺขิตฺวา มโนกมฺมํ ปริโสเธสุํ.\csromanspacer{} เยปิ หิ เกจิ ราหุล อนาคตมทฺธานํ สมณา วา พฺราหฺมณา วา กายกมฺมํ ปริโสเธสฺสนฺติ, วจีกมฺมํ ปริโสเธสฺสนฺติ, มโนกมฺมํ ปริโสเธสฺสนฺติ, สพฺเพ เต เอวเมวํ ปจฺจเวกฺขิตฺวา ปจฺจเวกฺขิตฺวา กายกมฺมํ ปริโสเธสฺสนฺติ, ปจฺจเวกฺขิตฺวา ปจฺจเวกฺขิตฺวา วจีกมฺมํ ปริโสเธสฺสนฺติ, ปจฺจเวกฺขิตฺวา ปจฺจเวกฺขิตฺวา มโนกมฺมํ ปริโสเธสฺสนฺติ.\csromanspacer{} เยปิ หิ เกจิ ราหุล เอตรหิ สมณา วา พฺราหฺมณา วา กายกมฺมํ ปริโสเธนฺติ, วจีกมฺมํ ปริโสเธนฺติ, มโนกมฺมํ ปริโสเธนฺติ, สพฺเพ เต เอวเมวํ ปจฺจเวกฺขิตฺวา ปจฺจเวกฺขิตฺวา กายกมฺมํ ปริโสเธนฺติ, ปจฺจเวกฺขิตฺวา ปจฺจเวกฺขิตฺวา วจีกมฺมํ ปริโสเธนฺติ, ปจฺจเวกฺขิตฺวา ปจฺจเวกฺขิตฺวา มโนกมฺมํ ปริโสเธนฺติ.\csromanspacer{} ตสฺมาติห ราหุล “ปจฺจเวกฺขิตฺวา ปจฺจเวกฺขิตฺวา กายกมฺมํ ปริโสเธสฺสามิ, ปจฺจเวกฺขิตฺวา ปจฺจเวกฺขิตฺวา วจีกมฺมํ ปริโสเธสฺสามิ, ปจฺจเวกฺขิตฺวา ปจฺจเวกฺขิตฺวา มโนกมฺมํ ปริโสเธสฺสามี”ติ เอวํ หิ เต ราหุล สิกฺขิตพฺพนฺติ.}

\prose{อิทมโวจ ภควา.\csromanspacer{} อตฺตมโน อายสฺมา ราหุโล ภควโต ภาสิตํ อภินนฺทีติ.}

\nitthitamruled{อมฺพลฏฺฐิกราหุโลวาทสุตฺตํ นิฏฺฐิตํ ปฐมํ.}
\csromanmatikaanchor{mtk.15}
\csromantocmarkref{subsection}{2. มหาราหุโลวาทสุตฺต}{mtk.15}
\bookmark[dest={mtk.15},level=2]{2. มหาราหุโลวาทสุตฺต}
\csromanheader{2}{2. มหาราหุโลวาทสุตฺต}
\proseitem{113}{เอวํ เม สุตํ\csromandash{}เอกํ สมยํ ภควา สาวตฺถิยํ วิหรติ เชตวเน อนาถปิณฺฑิกสฺส อาราเม.\csromanspacer{} อถ โข ภควา ปุพฺพณฺหสมยํ นิวาเสตฺวา ปตฺตจีวรมาทาย สาวตฺถิํ ปิณฺฑาย ปาวิสิ.\csromanspacer{} อายสฺมาปิ \csromanfolio{84}โข ราหุโล ปุพฺพณฺหสมยํ นิวาเสตฺวา ปตฺตจีวรมาทาย ภควนฺตํ ปิฏฺฐิโต ปิฏฺฐิโต อนุพนฺธิ.\csromanspacer{} อถ โข ภควา อปโลเกตฺวา อายสฺมนฺตํ ราหุลํ อามนฺเตสิ “ยํ กิญฺจิ ราหุล รูปํ อตีตานาคตปจฺจุปฺปนฺนํ อชฺฌตฺตํ วา พหิทฺธา วา โอฬาริกํ วา สุขุมํ วา หีนํ วา ปณีตํ วา ยํ ทูเร สนฺติเก วา, สพฺพํ รูปํ ‘เนตํ มม, เนโสหมสฺมิ, น เมโส อตฺตา’ติ, เอวเมตํ ยถาภูตํ สมฺมปฺปญฺญาย ทฏฺฐพฺพนฺ”ติ.\csromanspacer{} รูปเมว นุ โข ภควา, รูปเมว นุ โข สุคตาติ.\csromanspacer{} รูปมฺปิ ราหุล, เวทนาปิ ราหุล, สญฺญาปิ ราหุล, สงฺขาราปิ ราหุล, วิญฺญาณมฺปิ ราหุลาติ.\csromanspacer{} อถ โข อายสฺมา ราหุโล “โก นชฺช\footnote{โก นุชฺช (สฺยา, กํ)} ภควตา สมฺมุขา โอวาเทน โอวทิโต คามํ ปิณฺฑาย ปวิสิสฺสตี”ติ ตโต ปฏินิวตฺติตฺวา อญฺญตรสฺมิํ รุกฺขมูเล นิสีทิ ปลฺลงฺกํ อาภุชิตฺวา อุชุํ กายํ ปณิธาย ปริมุขํ สติํ อุปฏฺฐเปตฺวา.\csromanspacer{} อทฺทสา โข อายสฺมา สาริปุตฺโต อายสฺมนฺตํ ราหุลํ อญฺญตรสฺมิํ รุกฺขมูเล นิสินฺนํ ปลฺลงฺกํ อาภุชิตฺวา อุชุํ กายํ ปณิธาย ปริมุขํ สติํ อุปฏฺฐเปตฺวา, ทิสฺวาน อายสฺมนฺตํ ราหุลํ อามนฺเตสิ “อานาปานสฺสติํ ราหุล ภาวนํ ภาเวหิ, อานาปานสฺสติ ราหุล ภาวนา ภาวิตา พหุลีกตา มหปฺผลา โหติ มหานิสํสา”ติ.}

\proseitem{114}{อถ โข อายสฺมา ราหุโล สายนฺหสมยํ ปฏิสลฺลานา วุฏฺฐิโต เยน ภควา เตนุปสงฺกมิ, อุปสงฺกมิตฺวา ภควนฺตํ อภิวาเทตฺวา เอกมนฺตํ นิสีทิ, เอกมนฺตํ นิสินฺโน โข อายสฺมา ราหุโล ภควนฺตํ เอตทโวจ “กถํ ภาวิตา นุ โข ภนฺเต อานาปานสฺสติ กถํ พหุลีกตา มหปฺผลา โหติ มหานิสํสา”ติ.\csromanspacer{} ยํ กิญฺจิ ราหุล อชฺฌตฺตํ ปจฺจตฺตํ กกฺขฬํ ขริคตํ อุปาทินฺนํ.\csromanspacer{} เสยฺยถิทํ, เกสา โลมา นขาทนฺตา ตโจ, มํสํ นฺหารุ\footnote{นหารุ (สี, สฺยา, กํ, อิ)} อฏฺฐิ อฏฺฐิมิญฺชํ วกฺกํ, หทยํ ยกนํ กิโลมกํ ปิหกํ ปปฺผาสํ, อนฺตํ อนฺตคุณํ อุทริยํ กรีสํ, ยํ วา ปนญฺญมฺปิ กิญฺจิ อชฺฌตฺตํ ปจฺจตฺตํ กกฺขฬํ ขริคตํ อุปาทินฺนํ.\csromanspacer{} อยํ วุจฺจติ ราหุล อชฺฌตฺติกา ปถวีธาตุ\footnote{ปฐวีธาตุ (สี, สฺยา, กํ, อิ)}.\csromanspacer{} ยา เจว โข ปน อชฺฌตฺติกา ปถวีธาตุ, ยา จ พาหิรา ปถวีธาตุ, ปถวีธาตุเรเวสา.\csromanspacer{} ตํ “เนตํ มม, เนโสหมสฺมิ, น เมโส อตฺตา”ติ เอวเมตํ ยถาภูตํ สมฺมปฺปญฺญาย ทฏฺฐพฺพํ.\csromanspacer{} เอวเมตํ ยถาภูตํ สมฺมปฺปญฺญาย ทิสฺวา ปถวีธาตุยา นิพฺพินฺทติ, ปถวีธาตุยา จิตฺตํ วิราเชติ.}

\csromanheader{2}{2. มหาราหุโลวาทสุตฺต (62)}
\proseitem{115}{\csromanfolio{85}กตมา จ ราหุล อาโปธาตุ.\csromanspacer{} อาโปธาตุ สิยา อชฺฌตฺติกา สิยา พาหิรา.\csromanspacer{} กตมา จ ราหุล อชฺฌตฺติกา อาโปธาตุ.\csromanspacer{} ยํ อชฺฌตฺตํ ปจฺจตฺตํ อาโป อาโปคตํ อุปาทินฺนํ.\csromanspacer{} เสยฺยถิทํ, ปิตฺตํ เสมฺหํ ปุพฺโพ โลหิตํ เสโท เมโท อสฺสุ วสา เขโฬ สิงฺฆาณิกา ลสิกา มุตฺตํ, ยํ วา ปนญฺญมฺปิ กิญฺจิ อชฺฌตฺตํ ปจฺจตฺตํ อาโป อาโปคตํ อุปาทินฺนํ.\csromanspacer{} อยํ วุจฺจติ ราหุล อชฺฌตฺติกา อาโปธาตุ.\csromanspacer{} ยา เจว โข ปน อชฺฌตฺติกา อาโปธาตุ, ยา จ พาหิรา อาโปธาตุ, อาโปธาตุเรเวสา.\csromanspacer{} ตํ “เนตํ มม, เนโสหมสฺมิ, น เมโส อตฺตา”ติ เอวเมตํ ยถาภูตํ สมฺมปฺปญฺญาย ทฏฺฐพฺพํ.\csromanspacer{} เอวเมตํ ยถาภูตํ สมฺมปฺปญฺญาย ทิสฺวา อาโปธาตุยา นิพฺพินฺทติ, อาโปธาตุยา จิตฺตํ วิราเชติ.}

\proseitem{116}{กตมา จ ราหุล เตโชธาตุ.\csromanspacer{} เตโชธาตุ สิยา อชฺฌตฺติกา สิยา พาหิรา.\csromanspacer{} กตมา จ ราหุล อชฺฌตฺติกา เตโชธาตุ.\csromanspacer{} ยํ อชฺฌตฺตํ ปจฺจตฺตํ เตโช เตโชคตํ อุปาทินฺนํ.\csromanspacer{} เสยฺยถิทํ, เยน จ สนฺตปฺปติ, เยน จ ชีรียติ, เยน จ ปริฑยฺหติ, เยน จ อสิตปีตขายิตสายิตํ สมฺมา ปริณามํ คจฺฉติ, ยํ วา ปนญฺญมฺปิ กิญฺจิ อชฺฌตฺตํ ปจฺจตฺตํ เตโช เตโชคตํ อุปาทินฺนํ.\csromanspacer{} อยํ วุจฺจติ ราหุล อชฺฌตฺติกา เตโชธาตุ.\csromanspacer{} ยา เจว โข ปน อชฺฌตฺติกา เตโชธาตุ, ยา จ พาหิรา เตโชธาตุ, เตโชธาตุเรเวสา.\csromanspacer{} ตํ “เนตํ มม, เนโสหมสฺมิ, น เมโส อตฺตา”ติ เอวเมตํ ยถาภูตํ สมฺมปฺปญฺญาย ทฏฺฐพฺพํ.\csromanspacer{} เอวเมตํ ยถาภูตํ สมฺมปฺปญฺญาย ทิสฺวา เตโชธาตุยา นิพฺพินฺทติ, เตโชธาตุยา จิตฺตํ วิราเชติ.}

\proseitem{117}{กตมา จ ราหุล วาโยธาตุ.\csromanspacer{} วาโยธาตุ สิยา อชฺฌตฺติกา สิยา พาหิรา.\csromanspacer{} กตมา จ ราหุล อชฺฌตฺติกา วาโยธาตุ.\csromanspacer{} ยํ อชฺฌตฺตํ ปจฺจตฺตํ วาโย วาโยคตํ อุปาทินฺนํ.\csromanspacer{} เสยฺยถิทํ, อุทฺธงฺคมา วาตา, อโธคมา วาตา, กุจฺฉิสยา วาตา, โกฏฺฐาสยา\footnote{โกฏฺฐสยา (สี, อิ)} วาตา, องฺคมงฺคานุสาริโน วาตา, อสฺสาโส ปสฺสาโส-อิติ, ยํ วา ปนญฺญมฺปิ กิญฺจิ อชฺฌตฺตํ ปจฺจตฺตํ วาโย วาโยคตํ อุปาทินฺนํ.\csromanspacer{} อยํ วุจฺจติ ราหุล อชฺฌตฺติกา วาโยธาตุ.\csromanspacer{} ยา เจว โข ปน อชฺฌตฺติกา วาโยธาตุ, ยา จ พาหิรา วาโยธาตุ, วาโยธาตุเรเวสา.\csromanspacer{} ตํ “เนตํ มม, \csromanfolio{86}เนโสหมสฺมิ, น เมโส อตฺตา”ติ เอวเมตํ ยถาภูตํ สมฺมปฺปญฺญาย ทฏฺฐพฺพํ.\csromanspacer{} เอวเมตํ ยถาภูตํ สมฺมปฺปญฺญาย ทิสฺวา วาโยธาตุยา นิพฺพินฺทติ, วาโยธาตุยา จิตฺตํ วิราเชติ.}

\proseitem{118}{กตมา จ ราหุล อากาสธาตุ.\csromanspacer{} อากาสธาตุ สิยา อชฺฌตฺติกา สิยา พาหิรา.\csromanspacer{} กตมา จ ราหุล อชฺฌตฺติกา อากาสธาตุ.\csromanspacer{} ยํ อชฺฌตฺตํ ปจฺจตฺตํ อากาสํ อากาสคตํ อุปาทินฺนํ.\csromanspacer{} เสยฺยถิทํ, กณฺณจฺฉิทฺทํ นาสจฺฉิทฺทํ มุขทฺวารํ, เยน จ อสิตปีตขายิตสายิตํ อชฺโฌหรติ, ยตฺถ จ อสิตปีตขายิตสายิตํ สนฺติฏฺฐติ, เยน จ อสิตปีตขายิตสายิตํ อโธภาคํ\footnote{อโธภาคา (สี, สฺยา, กํ, อิ)} นิกฺขมติ, ยํ วา ปนญฺญมฺปิ กิญฺจิ อชฺฌตฺตํ ปจฺจตฺตํ อากาสํ อากาสคตํ อฆํ อฆคตํ วิวรํ วิวรคตํ อสมฺผุฏฺฐํ มํสโลหิเตหิ อุปาทินฺนํ\footnote{อากาสคตํ อุปาทินฺนํ (สี, อิ)}.\csromanspacer{} อยํ วุจฺจติ ราหุล อชฺฌตฺติกา อากาสธาตุ.\csromanspacer{} ยา เจว โข ปน อชฺฌตฺติกา อากาสธาตุ, ยา จ พาหิรา อากาสธาตุ, อากาสธาตุ เรเวสา.\csromanspacer{} ตํ “เนตํ มม, เนโสหมสฺมิ, น เมโส อตฺตา”ติ เอวเมตํ ยถาภูตํ สมฺมปฺปญฺญาย ทฏฺฐพฺพํ.\csromanspacer{} เอวเมตํ ยถาภูตํ สมฺมปฺปญฺญาย ทิสฺวา อากาสธาตุยา จิตฺตํ นิพฺพินฺทติ, อากาสธาตุยา จิตฺตํ วิราเชติ.}

\proseitem{119}{ปถวีสมํ ราหุล ภาวนํ ภาเวหิ, ปถวีสมํ หิ เต ราหุล ภาวนํ ภาวยโต อุปฺปนฺนา มนาปามนาปา ผสฺสา จิตฺตํ น ปริยาทาย ฐสฺสนฺติ.\csromanspacer{} เสยฺยถาปิ ราหุล ปถวิยา สุจิํปิ นิกฺขิปนฺติ, อสุจิํปิ นิกฺขิปนฺติ, คูถคตํปิ นิกฺขิปนฺติ, มุตฺตคตํปิ นิกฺขิปนฺติ, เขฬคตํปิ นิกฺขิปนฺติ, ปุพฺพคตํปิ นิกฺขิปนฺติ, โลหิตคตํปิ นิกฺขิปนฺติ.\csromanspacer{} น จ เตน ปถวี อฏฺฏียติ วา หรายติ วา ชิคุจฺฉติ วา.\csromanspacer{} เอวเมว โข ตฺวํ ราหุล ปถวีสมํ ภาวนํ ภาเวหิ, ปถวีสมํ หิ เต ราหุล ภาวนํ ภาวยโต อุปฺปนฺนา มนาปามนาปา ผสฺสา จิตฺตํ น ปริยาทาย ฐสฺสนฺติ.}

\prose{อาโปสมํ ราหุล ภาวนํ ภาเวหิ, อาโปสมํ หิ เต ราหุล ภาวนํ ภาวยโต อุปฺปนฺนา มนาปามนาปา ผสฺสา จิตฺตํ น ปริยาทาย ฐสฺสนฺติ.\csromanspacer{} เสยฺยถาปิ ราหุล อาปสฺมิํ สุจิํปิ โธวนฺติ, อสุจิํปิ โธวนฺติ, คูถคตํปิ โธวนฺติ, มุตฺตคตํปิ โธวนฺติ, เขฬคตํปิ โธวนฺติ, ปุพฺพคตํปิ โธวนฺติ, โลหิตคตํปิ โธวนฺติ.\csromanspacer{} น จ เตน อาโป}

\vspace{\tipitakaparskip}
\csromancenter{มหาราหุโลวาทสุตฺต (62) อฏฺฏียติ วา หรายติ วา ชิคุจฺฉติ วา.\csromanspacer{} เอวเมว โข ตฺวํ ราหุล อาโปสมํ ภาวนํ ภาเวหิ, อาโปสมํ หิ เต ราหุล ภาวนํ ภาวยโต อุปฺปนฺนา มนาปามนาปา ผสฺสา จิตฺตํ น ปริยาทาย ฐสฺสนฺติ.}
\prose{\csromanfolio{87}เตโชสมํ ราหุล ภาวนํ ภาเวหิ, เตโชสมํ หิเต ราหุล ภาวนํ ภาวยโต อุปฺปนฺนา มนาปามนาปา ผสฺสา จิตฺตํ น ปริยาทาย ฐสฺสนฺติ.\csromanspacer{} เสยฺยถาปิ ราหุล เตโช สุจิํปิ ทหติ, อสุจิํปิ ทหติ, คูถคตํปิ ทหติ, มุตฺตคตํปิ ทหติ, เขฬคตํปิ ทหติ, ปุพฺพคตํปิ ทหติ, โลหิตคตํปิ ทหติ.\csromanspacer{} น จ เตน เตโช อฏฺฏียติ วา หรายติ วา ชิคุจฺฉติ วา.\csromanspacer{} เอวเมว โข ตฺวํ ราหุล เตโชสมํ ภาวนํ ภาเวหิ, เตโชสมํ หิ เต ราหุล ภาวนํ ภาวยโต อุปฺปนฺนา มนาปามนาปา ผสฺสา จิตฺตํ น ปริยาทาย ฐสฺสนฺติ.}

\prose{วาโยสมํ ราหุล ภาวนํ ภาเวหิ, วาโยสมํ หิ เต ราหุล ภาวนํ ภวยโต อุปฺปนฺนา มนาปามนาปา ผสฺสา จิตฺตํ น ปริยาทาย ฐสฺสนฺติ.\csromanspacer{} เสยฺยถาปิ ราหุล วาโย สุจิํปิ อุปวายติ, อสุจิํปิ อุปวายติ, คูถคตํปิ อุปวายติ, มุตฺตคตํปิ อุปวายติ, เขฬคตํปิ อุปวายติ, ปุพฺพคตํปิ อุปวายติ, โลหิตคตํปิ อุปวายติ.\csromanspacer{} น จ เตน วาโย อฏฺฏียติ วา หรายติ วา ชิคุจฺฉติ วา.\csromanspacer{} เอวเมว โข ตฺวํ ราหุล วาโยสมํ ภาวนํ ภาเวหิ, วาโยสมํ หิ เต ราหุล ภาวนํ ภาวยโต อุปฺปนฺนา มนาปามนาปา ผสฺสา จิตฺตํ น ปริยาทาย ฐสฺสนฺติ.}

\prose{อากาสสมํ ราหุล ภาวนํ ภาเวหิ, อากาสสมํ หิ เต ราหุล ภาวนํ ภาวยโต อุปฺปนฺนา มนาปามนาปา ผสฺสา จิตฺตํ น ปริยาทาย ฐสฺสนฺติ.\csromanspacer{} เสยฺยถาปิ ราหุล อากาโส น กตฺถจิ ปติฏฺฐิโต.\csromanspacer{} เอวเมว โข ตฺวํ ราหุล อากาสสมํ ภาวนํ ภาเวหิ, อากาสสมํ หิ เต ราหุล ภาวนํ ภาวยโต อุปฺปนฺนา มนาปามนาปา ผสฺสา จิตฺตํ น ปริยาทาย ฐสฺสนฺติ.}

\proseitem{120}{เมตฺตํ ราหุล ภาวนํ ภาเวหิ, เมตฺตํ หิ เต ราหุล ภาวนํ ภาวยโต โย พฺยาปาโท, โส ปหียิสฺสติ.\csromanspacer{} กรุณํ ราหุล ภาวนํ ภาเวหิ, กรุณํ หิ เต ราหุล ภาวนํ ภาวยโต ยา วิเหสา, สา ปหียิสฺสติ.\csromanspacer{} มุทิตํ ราหุล ภาวนํ ภาเวหิ, มุทิตํ หิ เต ราหุล ภาวนํ ภาวยโต ยา อรติ, สา ปหียิสฺสติ. \csromanfolio{88}อุเปกฺขํ ราหุล ภาวนํ ภาเวหิ, อุเปกฺขํ หิ เต ราหุล ภาวนํ ภาวยโต โย ปฏิโฆ, โส ปหียิสฺสติ.\csromanspacer{} อสุภํ ราหุล ภาวนํ ภาเวหิ, อสุภํ หิ เต ราหุล ภาวนํ ภาวยโต โย ราโค, โส ปหียิสฺสติ.\csromanspacer{} อนิจฺจสญฺญํ ราหุล ภาวนํ ภาเวหิ, อนิจฺจสญฺญํ หิ เต ราหุล ภาวนํ ภาวยโต โย อสฺมิมาโน, โส ปหียิสฺสติ.}

\proseitem{121}{อานาปานสฺสติํ ราหุล ภาวนํ ภาเวหิ, อานาปานสฺสติ หิ เต ราหุล ภาวิตา พหุลีกตา มหปฺผลา โหติ มหานิสํสา.\csromanspacer{} กถํ ภาวิตา จ ราหุล อานาปานสฺสติ กถํ พหุลีกตา มหปฺผลา โหติ มหานิสํสา.\csromanspacer{} อิธ ราหุล ภิกฺขุ อรญฺญคโต วา รุกฺขมูลคโต วา สุญฺญาคารคโต วา นิสีทติ ปลฺลงฺกํ อาภุชิตฺวา อุชุํ กายํ ปณิธาย ปริมุขํ สติํ อุปฏฺฐเปตฺวา.\csromanspacer{} โส สโตว อสฺสสติ, สโตว\footnote{สโต (สี, สฺยา, กํ, อิ)} ปสฺสสติ.}

\prose{ทีฆํ วา อสฺสสนฺโต ทีฆํ อสฺสสามีติ ปชานาติ, ทีฆํ วา ปสฺสสนฺโต ทีฆํ ปสฺสสามีติ ปชานาติ.\csromanspacer{} รสฺสํ วา อสฺสสนฺโต รสฺสํ อสฺสสามีติ ปชานาติ, รสฺสํ วา ปสฺสสนฺโต รสฺสํ ปสฺสสามีติ ปชานาติ.\csromanspacer{} สพฺพกายปฺปฏิสํเวที อสฺสสิสฺสามีติ สิกฺขติ, สพฺพกายปฺปฏิสํเวที ปสฺสสิสฺสามีติ สิกฺขติ.\csromanspacer{} ปสฺสมฺภยํ กายสงฺขารํ อสฺสติสฺสามีติ สิกฺขติ, ปสฺสมฺภยํ กายสงฺขารํ ปสฺสสิสฺสามีติ สิกฺขติ.}

\prose{ปีติปฺปฏิสํเวที อสฺสสิสฺสามีติ สิกฺขติ, ปีติปฺปฏิสํเวที ปสฺสสิสฺสามีติ สิกฺขติ.\csromanspacer{} สุขปฺปฏิสํเวที อสฺสสิสฺสามีติ สิกฺขติ, สุขปฺปฏิสํเวที ปสฺสสิสฺสามีติ สิกฺขติ.\csromanspacer{} จิตฺตสงฺขารปฺปฏิสํเวที อสฺสสิสฺสามีติ สิกฺขติ, จิตฺตสงฺขารปฺปฏิสํเวที ปสฺสสิสฺสามีติ สิกฺขติ.\csromanspacer{} ปสฺสมฺภยํ จิตฺตสงฺขารํ อสฺสสิสฺสามีติ สิกฺขติ, ปสฺสมฺภยํ จิตฺตสงฺขารํ ปสฺสสิสฺสามีติ สิกฺขติ.}

\prose{จิตฺตปฺปฏิสํเวสี อสฺสสิสฺสามีติ สิกฺขติ, จิตฺตปฺปฏิสํเวสี ปสฺสสิสฺสามีติ สิกฺขติ.\csromanspacer{} อภิปฺปโมทยํ จิตฺตํ อสฺสสิสฺสามีติ สิกฺขติ, อภิปฺปโมทยํ จิตฺตํ ปสฺสสิสฺสามีติ สิกฺขติ.\csromanspacer{} สมาทหํ จิตฺตํ อสฺสสิสฺสามีติ สิกฺขติ, สมาทหํ จิตฺตํ ปสฺสสิสฺสามีติ สิกฺขติ.\csromanspacer{} วิโมจยํ จิตฺตํ อสฺสสิสฺสามีติ สิกฺขติ, วิโมจยํ จิตฺตํ ปสฺสสิสฺสามีติ สิกฺขติ.}

\csromanheader{2}{3. จูฬมาลุกฺยสุตฺต (63)}
\prose{\csromanfolio{89}อนิจฺจานุปสฺสี อสฺสสิสฺสามีติ สิกฺขติ, อนิจฺจานุปสฺสี ปสฺสสิสฺสามีติ สิกฺขติ.\csromanspacer{} วิราคานุปสฺสี อสฺสสิสฺสามีติ สิกฺขติ, วิราคานุปสฺสี ปสฺสสิสฺสามีติ สิกฺขติ.\csromanspacer{} นิโรธานุปสฺสี อสฺสสิสฺสามีติ สิกฺขติ, นิโรธานุปสฺสี ปสฺสสิสฺสามีติ สิกฺขติ.\csromanspacer{} ปฏินิสฺสคฺคานุปสฺสี อสฺสสิสฺสามีติ สิกฺขติ, ปฏินิสฺสคฺคานุปสฺสี ปสฺสสิสฺสามีติ สิกฺขติ.}

\prose{เอวํ ภาวิตา โข ราหุล อานาปานสฺสติ เอวํ พหุลีกตา มปฺผลา โหติ มหานิสํสา.\csromanspacer{} เอวํ ภาวิตาย ราหุล อานาปานสฺสติยา เอวํ พหุลีกตาย เยปิ เต จริมกา อสฺสาสา, เตปิ วิทิตาว นิรุชฺฌนฺติ, โน อวิทิตาติ.}

\prose{อิทมโวจ ภควา.\csromanspacer{} อตฺตมโน อายสฺมา ราหุโล ภควโต ภาสิตํ อภินนฺทีติ.}

\nitthitamruled{มหาราหุโลวาทสุตฺตํ นิฏฺฐิตํ ทุติยํ.}
\csromanmatikaanchor{mtk.16}
\csromantocmarkref{subsection}{3. จูฬมาลุกฺยสุตฺต}{mtk.16}
\bookmark[dest={mtk.16},level=2]{3. จูฬมาลุกฺยสุตฺต}
\csromanheader{2}{3. จูฬมาลุกฺยสุตฺต}
\proseitem{122}{เอวํ เม สุตํ\csromandash{}เอกํ สมยํ ภควา สาวตฺถิยํ วิหรติ เชตวเน อนาถปิณฺฑิกสฺส อาราเม.\csromanspacer{} อถ โข อายสฺมโต มาลุกฺยปุตฺตสฺส\footnote{มาลุงฺกฺยปุตฺตสฺส (สี, สฺยา, กํ, อิ)} รโหคตสฺส ปฏิสลฺลีนสฺส เอวํ เจตโส ปริวิตกฺโก อุทปาทิ “ยานิมานิ ทิฏฺฐิคตานิ ภควตา อพฺยากตานิ ฐปิตานิ ปฏิกฺขิตฺตานิ ‘สสฺสโต โลโกติปิ, อสสฺสโต โลโกติปิ, อนฺตวา โลโกติปิ, อนนฺตวา โลโกติปิ, ตํ ชีวํ ตํ สรีรนฺติปิ, อญฺญํ ชีวํ อญฺญํ สรีรนฺติปิ, โหติ ตถาคโต ปรํ มรณาติปิ, น โหติ ตถาคโต ปรํ มรณาติปิ, โหติ จ น จ โหติ ตถาคโต ปรํ มรณาติปิ, เนว โหติ น น โหติ ตถาคโต ปรํ มรณาติปิ’, ตานิ เม ภควา น พฺยากโรติ.\csromanspacer{} ยานิ เม ภควา น พฺยากโรติ, ตํ เม น รุจฺจติ, ตํ เม นกฺขมติ.\csromanspacer{} โสหํ ภควนฺตํ อุปสงฺกมิตฺวา เอตมตฺถํ ปุจฺฉิสฺสามิ, สเจ เม ภควา พฺยากริสฺสติ ‘สสฺสโต โลโกติ วา, อสสฺสโต โลโกติ วา ฯเปฯ เนว โหติ น น โหติ ตถาคโต ปรํ มรณาติ วา’.\csromanspacer{} เอวาหํ ภควติ \csromanfolio{90}พฺรหฺมจริยํ จริสฺสามิ.\csromanspacer{} โน เจ เม ภควา พฺยากริสฺสติ ‘สสฺสโต โลโกติ วา, อสสฺสโต โลโกติ วา ฯเปฯ เนว โหติ น น โหติ ตถาคโต ปรํ มรณาติ วา’.\csromanspacer{} เอวาหํ สิกฺขํ ปจฺจกฺขาย หีนายาวตฺติสฺสามี”ติ.}

\proseitem{123}{อถ โข อายสฺมา มาลุกฺยปุตฺโต สายนฺหสมยํ ปฏิสลฺลานา วุฏฺฐิโต เยน ภควา เตนุปสงฺกมิ, อุปสงฺกมิตฺวา ภควนฺตํ อภิวาเทตฺวา เอกมนฺตํ นิสีทิ, เอกมนฺตํ นิสินฺโน โข อายสฺมา มาลุกฺยปุตฺโต ภควนฺตํ เอตทโวจ\csromandash{}}

\proseitem{124}{อิธ มยฺหํ ภนฺเต รโหคตสฺส ปฏิสลฺลีนสฺส เอวํ เจตโส ปริวิตกฺโก อุทปาทิ “ยานิมานิ ทิฏฺฐิคตานิ ภควตา อพฺยากตานิ ฐปิตานิ ปฏิกฺขิตฺตานิ ‘สสฺสโต โลโกติปิ, อสสฺสโต โลโกติปิ ฯเปฯ เนว โหติ น น โหติ ตถาคโต ปรํ มรณาติปิ’, ตานิ เม ภควา น พฺยากโรติ.\csromanspacer{} ยานิ เม ภควา น พฺยากโรติ, ตํ เม น รุจฺจติ, ตํ เม นกฺขมติ.\csromanspacer{} โสหํ ภควนฺตํ อุปสงฺกมิตฺวา เอตมตฺถํ ปุจฺฉิสฺสามิ, สเจ เม ภควา พฺยากริสฺสติ ‘สสฺสโต โลโกติ วา, อสสฺสโต โลโกติ วา ฯเปฯ เนว โหติ น น โหติ ตถาคโต ปรํ มรณาติ วา’.\csromanspacer{} เอวาหํ ภควติ พฺรหฺมจริยํ จริสฺสามิ.\csromanspacer{} โน เจ เม ภควา พฺยากริสฺสติ ‘สสฺสโต โลโกติ วา, อสสฺสโต โลโกติ วา ฯเปฯ เนว โหติ น น โหติ ตถาคโต ปรํ มรณาติ วา’.\csromanspacer{} เอวาหํ สิกฺขํ ปจฺจกฺขาย หีนายาวตฺติสฺสามี”ติ.\csromanspacer{} สเจ ภควา ชานาติ “สสฺสโต โลโก”ติ, “สสฺสโต โลโก”ติ เม ภควา พฺยากโรตุ.\csromanspacer{} สเจ ภควา ชานาติ “อสสฺสโต โลโก”ติ, “อสสฺสโต โลโก”ติ เม ภควา พฺยากโรตุ.\csromanspacer{} โน เจ ภควา ชานาติ “สสฺสโต โลโก”ติ วา “อสสฺสโต โลโก”ติ วา, อชานโต โข ปน อปสฺสโต เอตเทว อุชุกํ โหติ, ยทิทํ น ชานามิ น ปสฺสามีติ.\csromanspacer{} สเจ ภควา ชานาติ “อนฺตวา โลโก”ติ, “อนฺตวา โลโก”ติ เม ภควา พฺยากโรตุ.\csromanspacer{} สเจ ภควา ชานาติ “อนนฺตวา โลโก”ติ, “อนนฺตวา โลโก”ติ เม ภควา พฺยากโรตุ.\csromanspacer{} โน เจ ภควา ชานาติ “อนฺตวา โลโก”ติ วา “อนนฺตวา โลโก”ติ วา, อชานโต โข ปน อปสฺสโต}

\vspace{\tipitakaparskip}
\csromancenter{จูฬมาลุกฺยสุตฺต (63) เอตเทว อุชุกํ โหติ, ยทิทํ น ชานามิ น ปสฺสามีติ.\csromanspacer{} สเจ ภควา ชานาติ “ตํ ชีวํ ตํ สรีรนฺ”ติ, “ตํ ชีวํ ตํ สรีรนฺ”ติ เม ภควา พฺยากโรตุ.\csromanspacer{} สเจ ภควา ชานาติ “อญฺญํ ชีวํ อญฺญํ สรีรนฺ”ติ, “อญฺญํ ชีวํ อญฺญํ สรีรนฺ”ติ เม ภควา พฺยากโรตุ.\csromanspacer{} โน เจ ภควา ชานาติ “ตํ ชีวํ ตํ สรีรนฺ”ติ วา “อญฺญํ ชีวํ อญฺญํ สรีรนฺ”ติ วา, อชานโต โข ปน อปสฺสโต เอตเทว อุชุกํ โหติ, ยทิทํ น ชานามิ น ปสฺสามีติ.\csromanspacer{} สเจ ภควา ชานาติ “โหติ ตถาคโต ปรํ มรณา”ติ, “โหติ ตถาคโต ปรํ มรณา”ติ เม ภควา พฺยากโรตุ.\csromanspacer{} สเจ ภควา ชานาติ “น โหติ ตถาคโต ปรํ มรณา”ติ, “น โหติ ตถาคโต ปรํ มรณา”ติ เม ภควา พฺยากโรตุ.\csromanspacer{} โน เจ ภควา ชานาติ “โหติ ตถาคโต ปรํ มรณา”ติ วา “น โหติ ตถาคโต ปรํ มรณา”ติ วา, อชานโต โข ปน อปสฺสโต เอตเทว อุชุกํ โหติ, ยทิทํ น ชานามิ น ปสฺสามีติ.\csromanspacer{} สเจ ภควา ชานาติ “โหติ จ น จ โหติ ตถาคโต ปรํ มรณา”ติ, “โหติ จ น จ โหติ ตถาคโต ปรํ มรณา”ติ เม ภควา พฺยากโรตุ.\csromanspacer{} สเจ ภควา ชานาติ “เนว โหติ น น โหติ ตถาคโต ปรํ มรณา”ติ, “เนว โหติ น น โหติ ตถาคโต ปรํ มรณา”ติ เม ภควา พฺยากโรตุ.\csromanspacer{} โน เจ ภควา ชานาติ “โหติ จ น จ โหติ ตถาคโต ปรํ มรณา”ติ วา “เนว โหติ น น โหติ ตถาคโต ปรํ มรณา”ติ วา, อชานโต โข ปน อปสฺสโต เอตเทว อุชุกํ โหติ, ยทิทํ น ชานามิ น ปสฺสามีติ.}
\proseitem{125}{\csromanfolio{91}กิํ นุ\footnote{กิํ นุ โข (สฺยา, กํ, ก)} ตาหํ มาลุกฺยปุตฺต เอวํ อวจํ “เอหิ ตฺวํ มาลุกฺยปุตฺต มยิ พฺรหฺมจริยํ จร, อหํ เต พฺยากริสฺสามิ ‘สสฺสโต โลโกติ วา, อสสฺสโต โลโกติ วา, อนฺตวา โลโกติ วา, อนนฺตวา โลโกติ วา, ตํ ชีวํ ตํ สรีรนฺติ วา, อญฺญํ ชีวํ อญฺญํ สรีรนฺติ วา, โหติ ตถาคโต ปรํ มรณาติ วา, น โหติ ตถาคโต ปรํ มรณาติ วา, โหติ จ น จ โหติ ตถาคโต ปรํ มรณาติ วา, เนว โหติ น น โหติ ตถาคโต ปรํ มรณาติ วา’ติ”.\csromanspacer{} โน เหตํ ภนฺเต.\csromanspacer{} ตฺวํ วา ปน มํ เอวํ อวจ “อหํ ภนฺเต ภควติ พฺรหฺมจริยํ \csromanfolio{92}จริสฺสามิ, ภควา เม พฺยากริสฺสติ ‘สสฺสโต โลโกติ วา, อสสฺสโต โลโกติ วา, อนฺตวา โลโกติ วา, อนนฺตวา โลโกติ วา, ตํ ชีวํ ตํ สรีรนฺติ วา, อญฺญํ ชีวํ อญฺญํ สรีรนฺติ วา, โหติ ตถาคโต ปรํ มรณาติ วา, น โหติ ตถาคโต ปรํ มรณาติ วา, โหติ จ น จ โหติ ตถาคโต ปรํ มรณาติ วา, เนว โหติ น น โหติ ตถาคโต ปรํ มรณาติ วา’ติ”.\csromanspacer{} โน เหตํ ภนฺเต.\csromanspacer{} อิติ กิร มาลุกฺยปุตฺต เนวาหํ ตํ วทามิ “เอหิ ตฺวํ มาลุกฺยปุตฺต มยิ พฺรหฺมจริยํ จร, อหํ เต พฺยากริสฺสามิ ‘สสฺสโต โลโกติ วา, อสสฺสโต โลโกติ วา ฯเปฯ เนว โหติ น น โหติ ตถาคโต ปรํ มรณาติ วา’ติ”.\csromanspacer{} นปิ กิร มํ ตฺวํ วเทสิ “อหํ ภนฺเต ภควติ พฺรหฺมจริยํ จริสฺสามิ, ภควา เม พฺยากริสฺสติ ‘สสฺสโตโลโกติ วา, อสสฺสโต โลโกติ วา ฯเปฯ เนว โหติ น น โหติ ตถาคโต ปรํ มรณาติ วา’ติ”.\csromanspacer{} เอวํ สนฺเต โมฆปุริส โก สนฺโต กํ ปจฺจาจิกฺขสิ.}

\proseitem{126}{โย โข มาลุกฺยปุตฺต เอวํ วเทยฺย “น ตาวาหํ ภควติ พฺรหฺมจริยํ จริสฺสามิ, ยาว เม ภควา น พฺยากริสฺสติ ‘สสฺสโต โลโกติ วา, อสสฺสโต โลโกติ วา ฯเปฯ เนว โหติ น น โหติ ตถาคโต ปรํ มรณาติ วา’ติ”.\csromanspacer{} อพฺยากตเมว ตํ มาลุกฺยปุตฺต ตถาคเตน อสฺส, อถ โส ปุคฺคโล กาลํ กเรยฺย.\csromanspacer{} เสยฺยถาปิ มาลุกฺยปุตฺต ปุริโส สลฺเลน วิทฺโธ อสฺส สวิเสน คาฬฺหปเลปเนน, ตสฺส มิตฺตามจฺจา ญาติสาโลหิตา ภิสกฺกํ สลฺลกตฺตํ อุปฏฺฐเปยฺยุํ.\csromanspacer{} โส เอวํ วเทยฺย “น ตาวาหํ อิมํ สลฺลํ อาหริสฺสามิ, ยาว น ตํ ปุริสํ ชานามิ, เยนมฺหิ วิทฺโธ ‘ขตฺติโย วา พฺราหฺมโณ วา เวสฺโส วา สุทฺโท วา’ติ”.\csromanspacer{} โส เอวํ วเทยฺย “น ตาวาหํ อิมํ สลฺลํ อาหริสฺสามิ, ยาว น ตํ ปุริสํ ชานามิ, เยนมฺหิ วิทฺโธ ‘เอวํนาโม เอวํโคตฺโต อิติวา’ติ”.\csromanspacer{} โส เอวํ วเทยฺย “น ตาวาหํ อิมํ สลฺลํ อาหริสฺสามิ, ยาว น ตํ ปุริสํ ชานามิ, เยนมฺหิ วิทฺโธ ‘ทีโฆ วา รสฺโส วา มชฺฌิโม วา’ติ”.\csromanspacer{} โส เอวํ วเทยฺย “น ตาวาตํ อิมํ สลฺลํ อาหริสฺสามิ, ยาว น ตํ ปุริสํ ชานามิ, เยนมฺหิ วิทฺโธ ‘กาโฬ วา สาโม วา มงฺคุรจฺฉวี วา’ติ”.\csromanspacer{} โส เอวํ วเทยฺย “น ตาวาหํ อิมํ สลฺลํ อาหริสฺสามิ, ยาว น ตํ ปุริสํ ชานามิ, เยนมฺหิ วิทฺโธ ‘อมุกสฺมิํ}

\vspace{\tipitakaparskip}
\csromancenter{จูฬมาลุกฺยสุตฺต (63) คาเม วา นิคเม วา นคเร วา’ติ”.\csromanspacer{} โส เอวํ วเทยฺย “น ตาวาหํ อิมํ สลฺลํ อาหริสฺสามิ, ยาว น ตํ ธนุํ ชานามิ, เยนมฺหิ วิทฺโธ “ยทิ วา จาโป ยทิ วา โกทณฺโฑ’ติ”.\csromanspacer{} โส เอวํ วเทยฺย “น ตาวาตํ อิมํ สลฺลํ อาหริสฺสามิ, ยาว น ตํ ชิยํ ชานามิ, ยายมฺหิ วิทฺโธ ‘ยทิ วา อกฺกสฺส ยทิ วา สณฺหสฺส\footnote{สณฺฐสฺส ((สี, สฺยา, กํ, อิ)} ยทิ วา นฺหารุสฺส ยทิ วา มรุวาย ยทิ วา ขีรปณฺณิโน’ติ”.\csromanspacer{} โส เอวํ วเทยฺย “น ตาวาหํ อิมํ สลฺลํ อาหริสฺสามิ, ยาว น ตํ กณฺฑํ ชานามิ, เยนมฺหิ วิทฺโธ ‘ยทิ วา คจฺฉํ ยทิ วา โรปิมนฺ’ติ”.\csromanspacer{} โส เอวํ วเทยฺย “น ตาวาหํ อิมํ สลฺลํ อาหริสฺสามิ, ยาว น ตํ กณฺฑํ ชานามิ, เยนมฺหิ วิทฺโธ ยสฺส ปตฺเตหิ วาชิตํ\footnote{วาขิตฺตํ (ก)} ‘ยทิ วา คิชฺฌสฺส ยทิ วา กงฺกสฺส ยทิ วา กุลลสฺส ยทิ วา โมรสฺส ยทิ วา สิถิลหนุโน’ติ”.\csromanspacer{} โส เอวํ วเทยฺย “น ตาวาหํ อิมํ สลฺลํ อาหริสฺสามิ, ยาว น ตํ กณฺฑํ ชานามิ, เยนมฺหิ วิทฺโธ ยสฺส นฺหารุนา ปริกฺขิตฺตํ ‘ยทิ วา ควสฺส ยทิ วา มหิํสสฺส ยทิ วา เภรวสฺส\footnote{โรรุวสฺส (สี, สฺยา, กํ, อิ)} ยทิ วา เสมฺหารสฺสา’ติ”.\csromanspacer{} โส เอวํ วเทยฺย “น ตาวาหํ อิมํ สลฺลํ อาหริสฺสามิ, ยาว น ตํ สลฺลํ ชานามิ, เยนมฺหิ วิทฺโธ ‘ยทิ วา สลฺลํ ยทิ วา ขุรปฺปํ ยทิ วา เวกณฺฑํ ยทิ วา นาราจํ ยทิ วา วจฺฉทนฺตํ ยทิ วา กรวีรปตฺตนฺ’ติ”.\csromanspacer{} อญฺญาตเมว ตํ มาลุกฺยปุตฺต เตน ปุริเสน อสฺส, อถ โส ปุริโส กาลํ กเรยฺย.\csromanspacer{} เอวเมว โข มาลุกฺยปุตฺต โย เอวํ วเทยฺย “น ตาวาหํ ภควติ พฺรหฺมจริยํ จริสฺสามิ, ยาว เม ภควา น พฺยากริสฺสติ ‘สสฺสโต โลโกติ วา, อสสฺสโต โลโกติ วา ฯเปฯ เนว โหติ น น โหติ ตถาคโต ปรํ มรณาติ วา’ติ”.\csromanspacer{} อพฺยากตเมว ตํ มาลุกฺยปุตฺต ตถาคเตน อสฺส, อถ โส ปุคฺคโล กาลํ กเรยฺย.}
\proseitem{127}{\csromanfolio{93}สสฺสโต โลโกติ มาลุกฺยปุตฺต ทิฏฺฐิยา สติ พฺรหฺมจริยวาโส อภวิสฺสาติ เอวํ โน.\csromanspacer{} อสสฺสโต โลโกติ มาลุกฺยปุตฺต ทิฏฺฐิยา สติ พฺรหฺมจริยวาโส อภวิสฺสาติ เอวมฺปิ โน.\csromanspacer{} สสฺสโต โลโกติ วา มาลุกฺยปุตฺต ทิฏฺฐิยา สติ อสสฺสโต โลโกติ วา ทิฏฺฐิยา สติ อตฺเถว ชาติ, อตฺถิ ชรา, อตฺถิ มรณํ, สนฺติ โสกปริเทวทุกฺขโทมนสฺสุปายาสา, เยสาหํ ทิฏฺเฐว ธมฺเม นิฆาตํ \csromanfolio{94}ปญฺญเปมิ.\csromanspacer{} อนฺตวา โลโกติ มาลุกฺยปุตฺต ทิฏฺฐิยา สติ พฺรหฺมจริยวาโส อภวิสฺสาติ เอวํ โน.\csromanspacer{} อนนฺตวา โลโกติ มาลุกฺยปุตฺต ทิฏฺฐิยา สติ พฺรหมจริยวาโส อภวิสฺสาติ เอวมฺปิ โน.\csromanspacer{} อนฺตวา โลโกติ วา มาลุกฺยปุตฺต ทิฏฺฐิยา สติ อนนฺตวา โลโกติ วา ทิฏฺฐิยา สติ อตฺเถว ชาติ, อตฺถิ ชรา, อตฺถิ มรณํ, สนฺติ โสกปริเทวทุกฺขโทมนสฺสุปายาสา, เยสาหํ ทิฏฺเฐว ธมฺเม นิฆาตํ ปญฺญเปมิ.\csromanspacer{} ตํ ชีวํ ตํ สรีรนฺติ มาลุกฺยปุตฺต ทิฏฺฐิยา สติ พฺรหฺมจริยวาโส อภวิสฺสาติ เอวํ โน.\csromanspacer{} อญฺญํ ชีวํ อญฺญํ สรีรนฺติ มาลุกฺยปุตฺต ทิฏฺฐิยา สติ พฺรหฺมจริยวาโส อภวิสฺสาติ เอวมฺปิ โน.\csromanspacer{} ตํ ชีวํ ตํ สรีรนฺติ วา มาลุกฺยปุตฺต ทิฏฺฐิยา สติ อญฺญํ ชีวํ อญฺญํ สรีรนฺติ วา ทิฏฺฐิยา สติ อตฺเถว ชาติ ฯเปฯ นิฆาตํ ปญฺญเปมิ.\csromanspacer{} โหติ ตถาคโต ปรํ มรณาติ มาลุกฺยปุตฺต ทิฏฺฐิยา สติ พฺรหฺมจริยวาโส อภวิสฺสาติ เอวํ โน.\csromanspacer{} น โหติ ตถาคโต ปรํ มรณาติ มาลุกฺยปุตฺต ทิฏฺฐิยา สติ พฺรหฺมจริยวาโส อภวิสฺสาติ เอวมฺปิ โน.\csromanspacer{} โหติ ตถาคโต ปรํ มรณาติ วา มาลุกฺยปุตฺต ทิฏฺฐิยา สติ น โหติ ตถาคโต ปรํ มรณาติ วา ทิฏฺฐิยา สติ อตฺเถว ชาติ ฯเปฯ เยสาหํ ทิฏฺเฐว ธมฺเม นิฆาตํ ปญฺญเปมิ.\csromanspacer{} โหติ จ น จ โหติ ตถาคโต ปรํ มรณาติ มาลุกฺยปุตฺต ทิฏฺฐิยา สติ พฺรหฺมจริยวาโส อภวิสฺสาติ เอวํ โน.\csromanspacer{} เนว โหติ น น โหติ ตถาคโต ปรํ มรณาติ มาลุกฺยปุตฺต ทิฏฺฐิยา สติ พฺรหฺมจริยวาโส อภวิสฺสาติ เอวมฺปิ โน.\csromanspacer{} โหติ จ น จ โหติ ตถาคโต ปรํ มรณาติ วา มาลุกฺยปุตฺต ทิฏฺฐิยา สติ เนว โหติ น น โหติ ตถาคโต ปรํ มรณาติ วา ทิฏฺฐิยา สติ อตฺเถว ชาติ ฯเปฯ เยสาหํ ทิฏฺเฐว ธมฺเม นิฆาตํ ปญฺญเปมิ.}

\proseitem{128}{ตสฺมาติห มาลุกฺยปุตฺต อพฺยากตญฺจ เม อพฺยากตโต ธาเรถ, พฺยากตญฺจ เม พฺยากตโต ธาเรถ.\csromanspacer{} กิญฺจ มาลุกฺยปุตฺต มยา อพฺยากตํ.\csromanspacer{} สสฺสโต โลโกติ มาลุกฺยปุตฺต มยา อพฺยากตํ, อสสฺสโต โลโกติ มยา อพฺยากตํ, อนฺตวา โลโกติ มยา อพฺยากตํ, อนนฺตวา โลโกติ มยา อพฺยากตํ, ตํ ชีวํ ตํ สรีรนฺติ มยา อพฺยากตํ, อญฺญํ ชีวํ อญฺญํ สรีรนฺติ มยา อพฺยากตํ, โหติ ตถาคโต ปรํ มรณาติ มยา อพฺยากตํ, น โหติ ตถาคโต}

\vspace{\tipitakaparskip}
\csromancenter{มหามาลุกฺยสุตฺต (64) ปรํ มรณาติ มยา อพฺยากตํ, โหติ จ น จ โหติ ตถาคโต ปรํ มรณาติ มยา อพฺยากตํ, เนว โหติ น น โหติ ตถาคโต ปรํ มรณาติ มยา อพฺยากตํ.\csromanspacer{} กสฺมา เจตํ มาลุกฺยปุตฺต มยา อพฺยากตํ.\csromanspacer{} น เหตํ มาลุกฺยปุตฺต อตฺถสํหิตํ, น อาทิพฺรหฺมจริยกํ, น\footnote{เนตํ (สี)} นิพฺพิทาย น วิราคาย น นิโรธาย น อุปสมาย น อภิญฺญาย น สมฺโพธาย น นิพฺพานาย สํวตฺตติ, ตสฺมา ตํ มยา อพฺยากตํ.\csromanspacer{} กิญฺจ มาลุกฺยปุตฺต มยา พฺยากตํ.\csromanspacer{} อิทํ ทุกฺขนฺติ มาลุกฺยปุตฺต มยา พฺยากตํ, อยํ ทุกฺขสมุทโยติ มยา พฺยากตํ, อยํ ทุกฺขนิโรโธติ มยา พฺยากตํ, อยํ ทุกฺขนิโรธคามินี ปฏิปทาติ มยา พฺยากตํ.\csromanspacer{} กสฺมา เจตํ มาลุกฺยปุตฺต มยา พฺยากตํ.\csromanspacer{} เอตํ หิ มาลุกฺยปุตฺต อตฺถสํหิตํ, เอตํ อาทิพฺรหฺมจริยกํ, นิพฺพิทาย วิราคาย นิโรธาย อุปสมาย อภิญฺญาย สมฺโพธาย นิพฺพานาย สํวตฺตติ, ตสฺมา ตํ มยา พฺยากตํ.\csromanspacer{} ตสฺมาติห มาลุกฺยปุตฺต อพฺยากตญฺจ เม อพฺยากตโต ธาเรถ, พฺยากตญฺจ เม พฺยากตโต ธาเรถาติ.}
\csromancenter{อิทมโวจ ภควา.\csromanspacer{} อตฺตมโน อายสฺมา มาลุกฺยปุตฺโต ภควโต ภาสิตํ อภินนฺทีติ.}
\nitthitamruled{จูฬมาลุกฺยสุตฺตํ นิฏฺฐิตํ ตติยํ.}
\csromanmatikaanchor{mtk.17}
\csromantocmarkref{subsection}{4. มหามาลุกฺยสุตฺต}{mtk.17}
\bookmark[dest={mtk.17},level=2]{4. มหามาลุกฺยสุตฺต}
\csromanheader{2}{4. มหามาลุกฺยสุตฺต}
\proseitem{129}{\csromanfolio{95}เอวํ เม สุตํ\csromandash{}เอกํ สมยํ ภควา สาวตฺถิยํ วิหรติ เชตวเน อนาถปิณฺฑิกสฺส อาราเม.\csromanspacer{} ตตฺร โข ภควา ภิกฺขู อามนฺเตสิ ภิกฺขโวติ.\csromanspacer{} ภทนฺเตติ เต ภิกฺขู ภควโต ปจฺจสฺโสสุํ.\csromanspacer{} ภควา เอตทโวจ “ธาเรถ โน ตุเมฺห ภิกฺขเว มยา เทสิตานิ ปญฺโจรมฺภาคิยานิ สํโยชนานี”ติ.}

\prose{เอวํ วุตฺเต อายสฺมา มาลุกฺยปุตฺโต ภควนฺตํ เอตทโวจ “อหํ โข ภนฺเต ธาเรมิ ภควตา เทสิตานิ ปญฺโจรมฺภาคิยานิ สํโยชนานี”ติ.\csromanspacer{} ยถา กถํ ปน ตฺวํ มาลุกฺยปุตฺต ธาเรสิ มยา เทสิตานิ ปญฺโจรมฺภาคิยานิ สํโยชนานีติ.\csromanspacer{} สกฺกายทิฏฺฐิํ โข อหํ ภนฺเต ภควตา \csromanfolio{96}โอรมฺภาคิยํ สํโยชนํ เทสิตํ ธาเรมิ, วิจิกิจฺฉํ โข อหํ ภนฺเต ภควตา โอรมฺภาคิยํ สํโยชนํ เทสิตํ ธาเรมิ, สีลพฺพตปรามาสํ โข อหํ ภนฺเต ภควตา โอรมฺภาคิยํ สํโยชนํ เทสิตํ ธาเรมิ, กามจฺฉนฺทํ โข อหํ ภนฺเต ภควตา โอรมฺภาคิยํ สํโยชนํ เทสิตํ ธาเรมิ, พฺยาปาทํ โข อหํ ภนฺเต ภควตา โอรมฺภาคิยํ สํโยชนํ เทสิตํ ธาเรมิ, เอวํ โข อหํ ภนฺเต ธาเรมิ ภควตา เทสิตานิ ปญฺโจรมฺภาคิยานิ สํโยชนานีติ.}

\prose{กสฺส โข นาม ตฺวํ มาลุกฺยปุตฺต อิมานิ เอวํ ปญฺโจรมฺภาคิยานิ สํโยชนานิ เทสิตานิ ธาเรสิ, นนุ มาลุกฺยปุตฺต อญฺญติตฺถิยา ปริพฺพาชกา อิมินา ตรุณูปเมน อุปารมฺเภน อุปารมฺภิสฺสนฺติ.\csromanspacer{} ทหรสฺส หิ มาลุกฺยปุตฺต กุมารสฺส มนฺทสฺส อุตฺตานเสยฺยกสฺส สกฺกาโยติปิ น โหติ, กุโต ปนสฺส อุปฺปชฺชิสฺสติ สกฺกายทิฏฺฐิ, อนุเสเตฺววสฺส\footnote{อนุเสติ เตฺววสฺส (สี, อิ)} สกฺกายทิฏฺฐานุสโย.\csromanspacer{} ทหรสฺส หิ มาลุกฺยปุตฺต กุมารสฺส มนฺทสฺส อุตฺตานเสยฺยกสฺส ธมฺมาติปิ น โหติ, กุโต ปนสฺส อุปฺปชฺชิสฺสติ ธมฺเมสุ วิจิกิจฺฉา, อนุเสเตฺววสฺส วิจิกิจฺฉานุสโย.\csromanspacer{} ทหรสฺส หิ มาลุกฺยปุตฺต กุมารสฺส มนฺทสฺส อุตฺตานเสยฺยกสฺส สีลาติปิ น โหติ, กุโต ปนสฺส อุปฺปชฺชิสฺสติ สีเลสุ สีลพฺพตปรามาโส, อนุเสเตฺววสฺส สีลพฺพตปรามาสานุสโย.\csromanspacer{} ทหรสฺส หิ มาลุกฺยปุตฺต กุมารสฺส มนฺทสฺส อุตฺตานเสยฺยกสฺส กามาติปิ น โหติ, กุโต ปนสฺส อุปฺปชฺชิสฺสติ กาเมสุ กามจฺฉนฺโท, อนุเสเตฺววสฺส กามราคานุสโย.\csromanspacer{} ทหรสฺส หิ มาลุกฺยปุตฺต กุมารสฺส มนฺทสฺส อุตฺตานเสยฺยกสฺส สตฺตาติปิ น โหติ, กุโต ปนสฺส อุปฺปชฺชิสฺสติ สตฺเตสุ พฺยาปาโท.\csromanspacer{} อนุเสเตฺววสฺส พฺยาปาทานุสโย.\csromanspacer{} นนุ มาลุกฺยปุตฺต อญฺญติตฺถิยา ปริพฺพาชกา อิมินา ตรุณูปเมน อุปารมฺเภน อุปารมฺภิสฺสนฺตีติ.\csromanspacer{} เอวํ วุตฺเต อายสฺมา อานนฺโท ภควนฺตํ เอตทโวจ “เอตสฺส ภควา กาโล, เอตสฺส สุคต กาโล, ยํ ภควา ปญฺโจรมฺภาคิยานิ สํโยชนานิ เทเสยฺย, ภควโต สุตฺวา ภิกฺขู ธาเรสฺสนฺตี”ติ.\csromanspacer{} เตน หานนฺท สุณาหิ สาธุกํ มนสิ กโรหิ ภาสิสฺสามีติ.\csromanspacer{} เอวํ ภนฺเตติ โข อายสฺมา อานนฺโท ภควโต ปจฺจสฺโสสิ.\csromanspacer{} ภควา เอตทโวจ\csromandash{}}

\csromanheader{2}{4. มหามาลุกฺยสุตฺต (64)}
\proseitem{130}{\csromanfolio{97}อิธานนฺท อสฺสุตวา ปุถุชฺชโน อริยานํ อทสฺสาวี อริยธมฺมสฺส อโกวิโท อริยธมฺเม อวินีโต สปฺปุริสานํ อทสฺสาวี สปฺปุริธมฺมสฺส อโกวิโท สปฺปุริสธมฺเม อวินีโต สกฺกายทิฏฺฐิปริยุฏฺฐิเตน เจตสา วิหรติ สกฺกายทิฏฺฐิปเรเตน, อุปฺปนฺนาย จ สกฺกายทิฏฺฐิยา นิสฺสรณํ ยถาภูตํ นปฺปชานาติ, ตสฺส สา สกฺกายทิฏฺฐิ ถามคตา อปฺปฏิวินีตา โอรมฺภาคิยํ สํโยชนํ.\csromanspacer{} วิจิกิจฺฉาปริยุฏฺฐิเตน เจตสา วิหรติ วิจิกิจฺฉาปเรเตน, อุปฺปนฺนาย จ วิจิกิจฺฉาย นิสฺสรณํ ยถาภูตํ นปฺปชานาติ, ตสฺส สา วิจิกิจฺฉา ถามคตา อปฺปฏิวินีตา โอรมฺภาคิยํ สํโยชนํ.\csromanspacer{} สีลพฺพตปรามาสปริยุฏฺฐิเตน เจตสา วิหรติ สีลพฺพตปรามาสปเรเตน, อุปฺปนฺนสฺส จ สีลพฺพตปรามาสสฺส นิสฺสรณํ ยถาภูตํ นปฺปชานาติ, ตสฺส โส สีลพฺพตปรามาโส ถามคโต อปฺปฏิวินีโต โอรมฺภาคิยํ สํโยชนํ.\csromanspacer{} กามราคปริยุฏฺฐิเตน เจตสา วิหรติ กามราคปเรเตน, อุปฺปนฺนสฺส จ กามราคสฺส นิสฺสรณํ ยถาภูตํ นปฺปชานาติ, ตสฺส โส กามราโค ถามคโต อปฺปฏิวินีโต โอรมฺภาคิยํ สํโยชนํ.\csromanspacer{} พฺยาปาทปริยุฏฺฐิเตน เจตสา วิหรติ พฺยาปาทปเรเตน, อุปฺปนฺนสฺส จ พฺยาปาทสฺส นิสฺสรณํ ยถาภูตํ นปฺปชานาติ.\csromanspacer{} ตสฺส โส พฺยาปาโท ถามคโต อปฺปฏิวินีโต โอรมฺภาคิยํ สํโยชนํ.}

\proseitem{131}{สุตวา จ โข อานนฺท อริยสาวโก อริยานํ ทสฺสาวี อริยธมฺมสฺส โกวิโท สุวินีโต สปฺปุริสานํ ทสฺสาวี สปฺปุริสธมฺมสฺส โกวิโท สปฺปุริสธมฺเม สุวินีโต น สกฺกายทิฏฺฐิปริยุฏฺฐิเตน เจตสา วิหรติ น สกฺกายทิฏฺฐิปเรเตน, อุปฺปนฺนาย จ สกฺกายทิฏฺฐิยา นิสฺสรณํ ยถาภูตํ ปชานาติ, ตสฺส สา สกฺกายทิฏฺฐิ สานุสยา ปหียติ.\csromanspacer{} น วิจิกิจฺฉาปริยุฏฺฐิเตน เจตสา วิหรติ น วิจิกิจฺฉาปเรเตน, อุปฺปนฺนาย จ วิจิกิจฺฉาย นิสฺสรณํ ยถาภูตํ ปชานาติ, ตสฺส สา วิจิกิจฺฉา สานุสยา ปหียติ.\csromanspacer{} น สีลพฺพตปรามาสปริยุฏฺฐิเตน เจตสา วิหรติ น สีลพฺพตปรามาสปเรเตน, อุปฺปนฺนสฺส จ สีลพฺพตปรามาสสฺส นิสฺสรณํ ยถาภูตํ ปชานาติ, ตสฺส โส สีลพฺพตปรามาโส สานุสโย ปหียติ.\csromanspacer{} น กามราคปริยุฏฺฐิเตน เจตสา วิหรติ น กามราคปเรเตน, อุปฺปนฺนสฺส จ กามราคสฺส นิสฺสรณํ ยถาภูตํ ปชานาติ, ตสฺส โส กามราโค สานุสโย \csromanfolio{98}ปหียติ.\csromanspacer{} น พฺยาปาทปริยุฏฺฐิเตน เจตสา วิหรติ น พฺยาปาทปเรเตน อุปฺปนฺนสฺส จ พฺยาปาทสฺส นิสฺสรณํ ยถาภูตํ ปชานาติ, ตสฺส โส พฺยาปาโท สานุสโย ปหียติ.}

\proseitem{132}{โย อานนฺท มคฺโค ยา ปฏิปทา ปญฺจนฺนํ โอรมฺภาคิยานํ สํโยชนานํ ปหานาย, ตํ มคฺคํ ตํ ปฏิปทํ อนาคมฺม ปญฺโจรมฺภาคิยานิ สํโยชนานิ ญสฺสติ วา ทกฺขติ วา ปชหิสฺสติ วาติ เนตํ ฐานํ วิชฺชติ.\csromanspacer{} เสยฺยถาปิ อานนฺท มหโต รุกฺขสฺส ติฏฺฐโต สารวโต ตจํ อจฺเฉตฺวา เผคฺคุํ อจฺเฉตฺวา สารจฺเฉโท ภวิสฺสตีติ เนตํ ฐานํ วิชฺชติ.\csromanspacer{} เอวเมว โข อานนฺท โย มคฺโค ยา ปฏิปทา ปญฺจนฺนํ โอรมฺภาคิยานํ สํโยชนานํ ปหานาย, ตํ มคฺคํ ตํ ปฏิปทํ อนาคมฺม ปญฺโจรมฺภาคิยานิ สํโยชนานิ ญสฺสติ วา ทกฺขติ วา ปชหิสฺสติ วาติ เนตํ ฐานํ วิชฺชติ.}

\prose{โย จ โข อานนฺท มคฺโค ยา ปฏิปทา ปญฺจนฺนํ โอรมฺภาคิยานํ สํโยชนานํ ปหานาย, ตํ มคฺคํ ตํ ปฏิปทํ อาคมฺม ปญฺโจรมฺภาคิยานิ สํโยชนานิ ญสฺสติ วา ทกฺขติ วา ปชหิสฺสติ วาติ ฐานเมตํ วิชฺชติ.\csromanspacer{} เสยฺยถาปิ อานนฺท มหโต รุกฺขสฺส ติฏฺฐโต สารวโต ตจํ เฉตฺวา เผคฺคุํ เฉตฺวา สารจฺเฉโท ภวิสฺสตีติ ฐานเมตํ วิชฺชติ, เอวเมว โข อานนฺท โย มคฺโค ยา ปฏิปทา ปญฺจนฺนํ โอรมฺภาคิยานํ สํโยชนานํ ปหานาย, ตํ มคฺคํ ตํ ปฏิปทํ อาคมฺม ปญฺโจรมฺภาคิยานิ สํโยชนานิ ญสฺสติ วา ทกฺขติ วา ปชหิสฺสติ วาติ ฐานเมตํ วิชฺชติ.\csromanspacer{} เสยฺยถาปิ อานนฺท คงฺคา นที ปูรา อุทกสฺส สมติตฺติกา กากเปยฺยา, อถ ทุพฺพลโก ปุริโส อาคจฺเฉยฺย “อหํ อิมิสฺสา คงฺคาย นทิยา ติริยํ พาหาย โสตํ เฉตฺวา โสตฺถินา ปารํ คจฺฉิสฺสามี”ติ\footnote{คจฺฉามีติ (สี, อิ)}.\csromanspacer{} โส น สกฺกุเณยฺย คงฺคาย นทิยา ติริยํ พาหาย โสตํ เฉตฺวา โสตฺถินา ปารํ คนฺตุํ, เอวเมว โข อานนฺท เยสํ เกสํจิ\footnote{ยสฺส กสฺสจิ (สพฺพตฺถ)} สกฺกายนิโรธาย ธมฺเม เทสิยมาเน จิตฺตํ น ปกฺขนฺทติ นปฺปสีทติ น สนฺติฏฺฐติ น วิมุจฺจติ, เสยฺยถาปิ โส ทุพฺพลโก ปุริโส เอวเมเต ทฏฺฐพฺพา.\csromanspacer{} เสยฺยถาปิ อานนฺท คงฺคา นที ปูรา อุทกสฺส สมติตฺติกา กากเปยฺยา, อถ พลวา ปุริโส อาคจฺเฉยฺย “อหํ อิมิสฺสา คงฺคาย นทิยา ติริยํ}

\vspace{\tipitakaparskip}
\csromancenter{มหามาลุกฺยสุตฺต (64) พาหาย โสตํ เฉตฺวา โสตฺถินา ปารํ คจฺฉิสฺสามี”ติ.\csromanspacer{} โส สกฺกุเณยฺย คงฺคาย นทิยา ติริยํ พาหาย โสตํ เฉตฺวา โสตฺถินา ปารํ คนฺตุํ, เอวเมว โข อานนฺท เยสํ เกสํจิ สกฺกายนิโรธาย ธมฺเม เทสิยมาเน จิตฺตํ ปกฺขนฺทติ ปสีทติ สนฺติฏฺฐติ วิมุจฺจติ, เสยฺยถาปิ โส พลวา ปุริโส เอวเมเต ทฏฺฐพฺพา.}
\proseitem{133}{\csromanfolio{99}กตโม จานนฺท มคฺโค กตมา ปฏิปทา ปญฺจนฺนํ โอรมฺภาคิยานํ สํโยชนานํ ปหานาย.\csromanspacer{} อิธานนฺท ภิกฺขุ อุปธิวิเวกา อกุสลานํ ธมฺมานํ ปหานา สพฺพโส กายทุฏฺฐุลฺลานํ ปฏิปฺปสฺสทฺธิยา วิวิจฺเจว กาเมหิ วิวิจฺจ อกุสเลหิ ธมฺเมหิ สวิตกฺกํ สวิจารํ วิเวกชํ ปีติสุขํ ปฐมํ ฌานํ อุปสมฺปชฺช วิหรติ.\csromanspacer{} โส ยเทว ตตฺถ โหติ รูปคตํ เวทนาคตํ สญฺญาคตํ สงฺขารคตํ วิญฺญาณคตํ, เต ธมฺเม อนิจฺจโต ทุกฺขโต โรคโต คณฺฑโต สลฺลโต อฆโต อาพาธโต ปรโต ปโลกโต สุญฺญโต อนตฺตโต สมนุปสฺสติ, โส เตหิ ธมฺเมหิ จิตฺตํ ปฏิวาเปติ\footnote{ปฏิปาเปติ (สฺยา), ปติฏฺฐาเปติ (ก)}, โส เตหิ ธมฺเมหิ จิตฺตํ ปฏิวาเปตฺวา อมตาย ธาตุยา จิตฺตํ อุปสํหรติ “เอตํ สนฺตํ เอตํ ปณีตํ, ยทิทํ สพฺพสงฺขารสมโถ สพฺพูปธิปฏินิสฺสคฺโค ตณฺหากฺขโย วิราโค นิโรโธ นิพฺพานนฺ”ติ, โส ตตฺถ ฐิโต อาสวานํ ขยํ ปาปุณาติ, โน เจ อาสวานํ ขยํ ปาปุณาติ, เตเนว ธมฺมราเคน ตาย ธมฺมนนฺทิยา ปญฺจนฺนํ โอรมฺภาคิยานํ สํโยชนานํ ปริกฺขยา โอปปาติโก โหติ ตตฺถ ปรินิพฺพายี อนาวตฺติธมฺโม ตสฺมา โลกา.\csromanspacer{} อยมฺปิ โข อานนฺท มคฺโค อยํ ปฏิปทา ปญฺจนฺนํ โอรมฺภาคิยานํ สํโยชนานํ ปหานาย.}

\prose{ปุน จปรํ อานนฺท ภิกฺขุ วิตกฺกวิจารานํ วูปสมา ฯเปฯ ทุติยํ ฌานํ อุปสมฺปชฺช วิหรติ ฯเปฯ ตติยํ ฌานํ ฯเปฯ จตุตฺถํ ฌานํ อุปสมฺปชฺช วิหรติ.\csromanspacer{} โส ยเทว ตตฺถ โหติ รูปคตํ เวทนาคตํ สญฺญาคตํ สงฺขารคตํ วิญฺญาณคตํ ฯเปฯ อนาวตฺติธมฺโม ตสฺมา โลกา.\csromanspacer{} อยมฺปิ โข อานนฺท มคฺโค อยํ ปฏิปทา ปญฺจนฺนํ โอรมฺภาคิยานํ สํโยชนานํ ปหานาย.}

\prose{ปุน จปรํ อานนฺท ภิกฺขุ สพฺพโส รูปสญฺญานํ สมติกฺกมา ปฏิฆสญฺญานํ อตฺถงฺคมา นานตฺตสญฺญานํ อมนสิการา “อนนฺโต อากาโส”ติ อากาสานญฺจายตนํ อุปสมฺปชฺช วิหรติ.\csromanspacer{} โส ยเทว ตตฺถ โหติ \csromanfolio{100}เวทนาคตํ สญฺญาคตํ สงฺขารคตํ วิญฺญาณคตํ ฯเปฯ อนาวตฺติธมฺโม ตสฺมา โลกา.\csromanspacer{} อยมฺปิ โข อานนฺท มคฺโค อยํ ปฏิปทา ปญฺจนฺนํ โอรมฺภาคิยานํ สํโยชนานํ ปหานาย.}

\prose{ปุน จปรํ อานนฺท ภิกฺขุ สพฺพโส อากาสานญฺจายตนํ สมติกฺกมฺม “อนนฺตํ วิญฺญาณนฺ”ติ วิญฺญาณญฺจายตนํ อุปสมฺปชฺช วิหรติ.\csromanspacer{} โส ยเทว ตตฺถ โหติ เวทนาคตํ สญฺญาคตํ สงฺขารคตํ วิญฺญาณคตํ ฯเปฯ อนาวตฺติธมฺโม ตสฺมา โลกา.\csromanspacer{} อยมฺปิ โข อานนฺท มคฺโค อยํ ปฏิปทา ปญฺจนฺนํ โอรมฺภาคิยานํ สํโยชนานํ ปหานาย.}

\prose{ปุน จปรํ อานนฺท ภิกฺขุ สพฺพโส วิญฺญาณญฺจายตนํ สมติกฺกมฺม “นตฺถิ กิญฺจี”ติ อากิญฺจญฺญายตนํ อุปสมฺปชฺช วิหรติ.\csromanspacer{} โส ยเทว ตตฺถ โหติ เวทนาคตํ สญฺญาคตํ สงฺขารคตํ วิญฺญาณคตํ ฯเปฯ อนาวตฺติธมฺโม ตสฺมา โลกา.\csromanspacer{} อยมฺปิ โข อานนฺท มคฺโค อยํ ปฏิปทา ปญฺจนฺนํ โอรมฺภาคิยานํ สํโยชนานํ ปหานายาติ.}

\prose{เอโส เจ ภนฺเต มคฺโค เอสา ปฏิปทา ปญฺจนฺนํ โอรมฺภาคิยานํ สํโยชนานํ ปหานาย, อถ กิญฺจรหิ อิเธกจฺเจ ภิกฺขู เจโต วิมุตฺติโน, เอกจฺเจ ภิกฺขู ปญฺญาวิมุตฺติโนติ.\csromanspacer{} เอตฺถ โข ปเนสาหํ\footnote{เอตฺถ โข เตสาหํ (สี, สฺยา, กํ, อิ)} อานนฺท อินฺทฺริยเวมตฺตตํ วทามีติ.}

\prose{อิทมโวจ ภควา.\csromanspacer{} อตฺตมโน อายสฺมา อานนฺโท ภควโต ภาสิตํ อภินนฺทีติ.}

\nitthitamruled{มหามาลุกฺยสุตฺตํ นิฏฺฐิตํ จตุตฺถํ.}
\csromanmatikaanchor{mtk.18}
\csromantocmarkref{subsection}{5. ภทฺทาลิสุตฺต}{mtk.18}
\bookmark[dest={mtk.18},level=2]{5. ภทฺทาลิสุตฺต}
\csromanheader{2}{5. ภทฺทาลิสุตฺต}
\proseitem{134}{เอวํ เม สุตํ\csromandash{}เอกํ สมยํ ภควา สาวตฺถิยํ วิหรติ เชตวเน อนาถปิณฺฑิกสฺส อาราเม.\csromanspacer{} ตตฺร โข ภควา ภิกฺขู อามนฺเตสิ “ภิกฺขโว”ติ. “ภทนฺเต”ติ เต ภิกฺขู ภควโต ปจฺจสฺโสสุํ.\csromanspacer{} ภควา เอตทโวจ “อหํ โข ภิกฺขเว เอกาสนโภชนํ ภุญฺชามิ, เอกาสนโภชนํ โข อหํ ภิกฺขเว ภุญฺชมาโน อปฺปาพาธตญฺจ}

\vspace{\tipitakaparskip}
\csromancenter{พทฺทาลิสุตฺต (65) สญฺชานามิ อปฺปาตงฺกตญฺจ ลหุฏฺฐานญฺจ พลญฺจ ผาสุวิหารญฺจ.\csromanspacer{} เอถ ตุเมฺหปิ ภิกฺขเว เอกาสนโภชนํ ภุญฺชถ, เอกาสนโภชนํ โข ภิกฺขเว ตุเมฺหปิ ภุญฺชมานา อปฺปาพาธตญฺจ สญฺชานิสฺสถ อปฺปาตงฺกตญฺจ ลหุฏฺฐานญฺจ พลญฺจ ผาสุวิหารญฺจา”ติ.\csromanspacer{} เอวํ วุตฺเต อายสฺมา ภทฺทาลิ ภควนฺตํ เอตทโวจ “อหํ โข ภนฺเต น อุสฺสหามิ เอกาสนโภชนํ ภุญฺชิตุํ, เอกาสนโภชนญฺหิ เม ภนฺเต ภุญฺชโต สิยา กุกฺกุจฺจํ, สิยา วิปฺปฏิสาโร”ติ.\csromanspacer{} เตน หิ ตฺวํ ภทฺทาลิ ยตฺถ นิมนฺติโต อสฺสสิ, ตตฺถ เอกเทสํ ภุญฺชิตฺวา เอกเทสํ นีหริตฺวาปิ ภุญฺเชยฺยาสิ, เอวมฺปิ โข ตฺวํ ภทฺทาลิ ภุญฺชมาโน เอกาสโน ยาเปสฺสสี”ติ\footnote{ภุญฺชมาโน ยาเปสฺสสีติ (สี, สฺยา, กํ, อิ)}.\csromanspacer{} เอวมฺปิ โข อหํ ภนฺเต น อุสฺสหามิ ภุญฺชิตุํ, เอวมฺปิ หิ เม ภนฺเต ภุญฺชโต สิยา กุกฺกุจฺจํ, สิยา วิปฺปฏิสาโรติ.\csromanspacer{} อถ โข อายสฺมา ภทฺทาลิ ภควตา สิกฺขาปเท ปญฺญาปิยมาเน ภิกฺขุสํเฆ สิกฺขํ สมาทิยมาเน อนุสฺสาหํ ปเวเทสิ.\csromanspacer{} อถ โข อายสฺมา ภทฺทาลิ สพฺพํ ตํ เตมาสํ น ภควโต สมฺมุขีภาวํ อทาสิ, ยถา ตํ สตฺถุสาสเน สิกฺขาย อปริปูรการี.}
\proseitem{135}{\csromanfolio{101}เตน โข ปน สมเยน สมฺพหุลา ภิกฺขู ภควโต จีวรกมฺมํ กโรนฺติ “นิฏฺฐิตจีวโร ภควา เตมาสจฺจเยน จาริกํ ปกฺกมิสฺสตี”ติ.\csromanspacer{} อถ โข อายสฺมา ภทฺทาลิ เยน เต ภิกฺขู เตนุปสงฺกมิ, อุปสงฺกมิตฺวา เตหิ ภิกฺขูหิ สทฺธิํ สมฺโมทิ, สมฺโมทนียํ กถํ สารณียํ วีติสาเรตฺวา เอกมนฺตํ นิสีทิ, เอกมนฺตํ นิสินฺนํ โข อายสฺมนฺตํ ภทฺทาลิํ เต ภิกฺขู เอตทโวจุํ “อิทํ โข อาวุโส ภทฺทาลิ ภควโต จีวรกมฺมํ กรียติ\footnote{กรณียํ (ก)}, นิฏฺฐิตจีวโร ภควา เตมาสจฺจเยน จาริกํ ปกฺกมิสฺสติ.\csromanspacer{} อิงฺฆาวุโส ภทฺทาลิ เอตํ โทสกํ สาธุกํ มนสิ กโรหิ, มา เต ปจฺฉา ทุกฺกรตรํ อโหสี”ติ. “เอวมาวุโส”ติ โข อายสฺมา ภทฺทาลิ เตสํ ภิกฺขูนํ ปฏิสฺสุตฺวา เยน ภควา เตนุปสงฺกมิ, อุปสงฺกมิตฺวา ภควนฺตํ อภิวาเทตฺวา เอกมนฺตํ นิสีทิ, เอกมนฺตํ นิสินฺโน โข อายสฺมา ภทฺทาลิ ภควนฺตํ เอตทโวจ “อจฺจโย มํ ภนฺเต อจฺจคมา ยถาพาลํ ยถามูฬฺหํ ยถา-อกุสลํ, โยหํ ภควตา สิกฺขาปเท ปญฺญาปิยมาเน ภิกฺขุสํเฆ สิกฺขํ สมาทิยมาเน อนุสฺสาหํ ปเวเทสิํ, ตสฺส เม ภนฺเต ภควา อจฺจยํ อจฺจยโต ปฏิคฺคณฺหาตุ อายติํ สํวรายา”ติ}

\prose{\csromanfolio{102}ตคฺฆ ตฺวํ ภทฺทาลิ อจฺจโย อจฺจคมา ยถาพาลํ ยถามูฬฺหํ ยถา-อกุสลํ, ยํ ตฺวํ มยา สิกฺขาปเท ปญฺญาปิยมาเน ภิกฺขุสํเฆ สิกฺขํ สมาทิยมาเน อนุสฺสาหํ ปเวเทสิ.\csromanspacer{} สมโยปิ โข เต ภทฺทาลิ อปฺปฏิวิทฺโธ อโหสิ, ภควา โข สาวตฺถิยํ วิหรติ, ภควาปิ มํ ชานิสฺสติ “ภทฺทาลิ นาม ภิกฺขุ สตฺถุสาสเน สิกฺขาย อปริปูรการี”ติ, อยมฺปิ โข เต ภทฺทาลิ สมโย อปฺปฏิวิทฺโธ อโหสิ.\csromanspacer{} สมโยปิ โข เต ภทฺทาลิ อปฺปฏิวิทฺโธ อโหสิ, สมฺพหุลา โข ภิกฺขู สาวตฺถิยํ วสฺสํ อุปคตา, เตปิ มํ ชานิสฺสนฺติ “ภทฺทาลิ นาม ภิกฺขุ สตฺถุสาสเน สิกฺขาย อปริปูรการี”ติ, อยมฺปิ โข เต ภทฺทาลิ สมโย อปฺปฏิวิทฺโธ อโหสิ.\csromanspacer{} สมโยปิ โข เต ภทฺทาลิ อปฺปฏิวิทฺโธ อโหสิ, สมฺพหุลา โข ภิกฺขุนิโย สาวตฺถิยํ วสฺสํ อุปคตา, ตาปิ มํ ชานิสฺสนฺติ “ภทฺทาลิ นาม ภิกฺขุ สตฺถุสาสเน สิกฺขาย อปริปูรการี”ติ, อยมฺปิ โข เต ภทฺทาลิ สมโย อปฺปฏิวิทฺโธ อโหสิ.\csromanspacer{} สมโยปิ โข เต ภทฺทาลิ อปฺปฏิวิทฺโธ อโหสิ, สมฺพหุลา โข อุปาสกา สาวตฺถิยํ ปฏิวสนฺติ, เตปิ มํ ชานิสฺสนฺติ “ภทฺทาลิ นาม ภิกฺขุ สตฺถุสาสเน สิกฺขาย อปริปูรการี”ติ, อยมฺปิ โข เต ภทฺทาลิ สมโย อปฺปฏิวิทฺโธ อโหสิ.\csromanspacer{} สมโยปิ โข เต ภทฺทาลิ อปฺปฏิวิทฺโธ อโหสิ, สมฺพหุลา โข อุปาสิกา สวตฺถิยํ ปฏิวสนฺติ, ตาปิ มํ ชานิสฺสนฺติ “ภทฺทาลิ นาม ภิกฺขุ สตฺถุสาสเน สิกฺขาย อปริปูรการี”ติ, อยมฺปิ โข เต ภทฺทาลิ สมโย อปฺปฏิวิทฺโธ อโหสิ.\csromanspacer{} สมโยปิ โข เต ภทฺทาลิ อปฺปฏิวิทฺโธ อโหสิ, สมฺพหุลา โข นานาติตฺถิยา สมณพฺราหฺมณา สาวตฺถิยํ วสฺสํ อุปคตา, เตปิ มํ ชานิสฺสนฺติ “ภทฺทาลิ นาม ภิกฺขุ สมณสฺส โคตมสฺส สาวโก เถรญฺญตโร ภิกฺขุ สาสเน สิกฺขาย อปริปูรการี”ติ, อยมฺปิ โข เต ภทฺทาลิ สมโย อปฺปฏิวิทฺโธ อโหสีติ.\csromanspacer{} อจฺจโย มํ ภนฺเต อจฺจคมา ยถาพาลํ ยถามูฬฺหํ ยถา-อกุสลํ, โยหํ ภควตา สิกฺขาปเท ปญฺญาปิยมาเน ภิกฺขุสํเฆ สิกฺขํ สมาทิยมาเน อนุสฺสาหํ ปเวเทสิํ, ตสฺส เม ภนฺเต ภควา อจฺจยํ อจฺจยโต ปฏิคฺคณฺหาตุ อายติํ สํวรายาติ.\csromanspacer{} ตคฺฆ ตฺวํ ภทฺทาลิ อจฺจโย อจฺจคมา ยถาพาลํ ยถามูฬฺหํ ยถา-อกุสลํ, ยํ ตฺวํ มยา}

\vspace{\tipitakaparskip}
\csromancenter{พทฺทาลิสุตฺต (65) สิกฺขาปเท ปญฺญาปิยมาเน ภิกฺขุสํเฆ สิกฺขํ สมาทิยมาเน อนุสฺสาหํ ปเวเทสิ.}
\proseitem{136}{\csromanfolio{103}ตํ กิํ มญฺญสิ ภทฺทาลิ, อิธสฺส ภิกฺขุ อุภโตภาควิมุตฺโต, ตมหํ เอวํ วเทยฺยํ “เอหิ เม ตฺวํ ภิกฺขุ ปงฺเก สงฺกโม โหหี”ติ.\csromanspacer{} อปิ นุ โข โส สงฺกเมยฺย วา อญฺเญน วา กายํ สนฺนาเมยฺย โนติ วา วเทยฺยาติ.\csromanspacer{} โน เหตํ ภนฺเต.\csromanspacer{} ตํ กิํ มญฺญสิ ภทฺทาลิ, อิธสฺส ภิกฺขุ ปญฺญาวิมุตฺโต.\csromanspacer{} กายสกฺขิ.\csromanspacer{} ทิฏฺฐิปฺปตฺโต.\csromanspacer{} สทฺธาวิมุตฺโต.\csromanspacer{} ธมฺมานุสารี.\csromanspacer{} สทฺธานุสารี, ตมหํ เอวํ วเทยฺยํ “เอหิ เม ตฺวํ ภิกฺขุ ปงฺเก สงฺกโม โหหี”ติ.\csromanspacer{} อปิ นุ โข โส สงฺกเมยฺย วา อญฺเญน วา กายํ สนฺนาเมยฺย โนติ วา วเทยฺยาติ.\csromanspacer{} โน เหตํ ภนฺเต.\csromanspacer{} ตํ กิํ มญฺญสิ ภทฺทาลิ, อปิ นุ ตฺวํ ภทฺทาลิ ตสฺมิํ สมเย อุภโตภาควิมุตฺโต วา โหสิ, ปญฺญาวิมุตฺโต วา กายสกฺขิ วา ทิฏฺฐิปฺปตฺโต วา สทฺธาวิมุตฺโต วา ธมฺมานุสารี วา สทฺธานุสารี วาติ.\csromanspacer{} โน เหตํ ภนฺเต.\csromanspacer{} นนุ ตฺวํ ภทฺทาลิ ตสฺมิํ สมเย ริตฺโต ตุจฺโฉ อปรทฺโธติ.\csromanspacer{} เอวํ ภนฺเต.\csromanspacer{} อจฺจโย มํ ภนฺเต อจฺจคมา ยถาพาลํ ยถามูฬฺหํ ยถา-อกุสลํ, โยหํ ภควตา สิกฺขาปเท ปญฺญาปิยมาเน ภิกฺขุสํเฆ สิกฺขํ สมาทิยมาเน อนุสฺสาหํ ปเวเทสิํ, ตสฺส เม ภนฺเต ภควา อจฺจยํ อจฺจยโต ปฏิคฺคณฺหาตุ อายติํ สํวรายาติ.\csromanspacer{} ตคฺฆ ตฺวํ ภทฺทาลิ อจฺจโย อจฺจคมา ยถาพาลํ ยถามูฬฺหํ ยถา-อกุสลํ, ยํ ตฺวํ มยา สิกฺขาปเท ปญฺญาปิยมาเน ภิกฺขุสํเฆ สิกฺขํ สมาทิยมาเน อนุสฺสาหํ ปเวเทสิ.\csromanspacer{} ยโต จ โข ตฺวํ ภทฺทาลิ อจฺจยํ อจฺจยโต ทิสฺวา ยถาธมฺมํ ปฏิกโรสิ, ตํ เต มยํ ปฏิคฺคณฺหาม.\csromanspacer{} วุทฺธิเหสา ภทฺทาลิ อริยสฺส วินเย, โย อจฺจยํ อจฺจยโต ทิสฺวา ยถาธมฺมํ ปฏิกโรติ, อายติํ สํวรํ อาปชฺชติ.}

\proseitem{137}{อิธ ภทฺทาลิ เอกจฺโจ ภิกฺขุ สตฺถุสาสเน สิกฺขาย อปริปูรการี โหติ, ตสฺส เอวํ โหติ “ยํนูนาหํ วิวิตฺตํ เสนาสนํ ภเชยฺยํ อรญฺญํ รุกฺขมูลํ ปพฺพตํ กนฺทรํ คิริคุหํ สุสานํ วนปตฺถํ อพฺโภกาสํ ปลาลปุญฺชํ, อปฺเปว นามาหํ อุตฺตริ\footnote{อุตฺตริํ (สี, สฺยา, กํ, อิ)} มนุสฺสธมฺมา อลมริยญาณทสฺสนวิเสสํ สจฺฉิกเรยฺยนฺ”ติ, โส วิวิตฺตํ เสนาสนํ ภชติ อรญฺญํ รุกฺขมูลํ \csromanfolio{104}ปพฺพตํ กนฺทรํ คิริคุหํ สุสานํ วนปตฺถํ อพฺโภกาสํ ปลาลปุญฺชํ, ตสฺส ตถาวูปกฏฺฐสฺส วิหรโต สตฺถาปิ อุปวทติ, อนุวิจฺจปิ วิญฺญู สพฺรหฺมจารี อุปวทนฺติ, เทวตาปิ อุปวทนฺติ, อตฺตาปิ อตฺตานํ อุปวทติ.\csromanspacer{} โส สตฺถาราปิ อุปวทิโต อนุวิจฺจปิ วิญฺญูหิ สพฺรหฺมจารีหิ อุปวทิโต เทวตาหิปิ อุปวทิโต อตฺตนาปิ อตฺตานํ อุปวทิโต น อุตฺตริ มนุสฺสธมฺมา อลมริยญาณทสฺสนวิเสสํ สจฺฉิกโรติ.\csromanspacer{} ตํ กิสฺส เหตุ, เอวญฺหิ ตํ ภทฺทาลิ โหติ, ยถา ตํ สตฺถุสาสเน สิกฺขาย อปริปูรการิสฺส.}

\proseitem{138}{อิธ ปน ภทฺทาลิ เอกจฺโจ ภิกฺขุ สตฺถุสาสเน สิกฺขาย ปริปูรการี โหติ, ตสฺส เอวํ โหติ “ยํนูนาหํ วิวิตฺตํ เสนาสนํ ภเชยฺยํ อรญฺญํ รุกฺขมูลํ ปพฺพตํ กนฺทรํ คิริคุหํ สุสานํ วนปตฺถํ อพฺโภกาสํ ปลาลปุญฺชํ, อปฺเปว นามาหํ อุตฺตริ มนุสฺสธมฺมา อลมริยญาณทสฺสนวิเสสํ สจฺฉิกเรยฺยนฺ”ติ, โส วิวิตฺตํ เสนาสนํ ภชติ อรญฺญํ รุกฺขมูลํ ปพฺพตํ กนฺทรํ คิริคุหํ สุสานํ วนตฺถํ อพฺโภกาสํ ปลาลปุญฺชํ, ตสฺส ตถาวูปกฏฺฐสฺส วิหรโต สตฺถาปิ น อุปวทติ, อนุวิจฺจปิ วิญฺญู สพฺรหฺมจารี น อุปวทนฺติ, เทวตาปิ น อุปวทนฺติ, อตฺตาปิ อตฺตานํ น อุปวทติ.\csromanspacer{} โส สตฺถาราปิ อนุปวทิโต อนุวิจฺจปิ วิญฺญูหิ สพฺรหฺมจารีหิ อนุปวทิโต เทวตาหิปิ อนุปวทิโต อตฺตนาปิ อตฺตานํ อนุปวทิโต อุตฺตริ มนุสฺสธมฺมา อลมริยญาณทสฺสนวิเสสํ สจฺฉิกโรติ, โส วิวิจฺเจว กาเมหิ วิวิจฺจ อกุสเลหิ ธมฺเมหิ สวิตกฺกํ สวิจารํ วิเวกชํ ปีติสุขํ ปฐมํ ฌานํ อุปสมฺปชฺช วิหรติ.\csromanspacer{} ตํ กิสฺส เหตุ, เอวญฺหิ ตํ ภทฺทาลิ โหติ, ยถา ตํ สตฺถุสาสเน สิกฺขาย ปริปูรการิสฺส.}

\proseitem{139}{ปุน จปรํ ภทฺทาลิ ภิกฺขุ วิตกฺกวิจารานํ วูปสมา อชฺฌตฺตํ สมฺปสาทนํ เจตโส เอโกทิภาวํ อวิตกฺกํ อวิจารํ สมาธิชํ ปีติสุขํ ทุติยํ ฌานํ อุปสมฺปชฺช วิหรติ.\csromanspacer{} ตํ กิสฺส เหตุ, เอวญฺหิ ตํ ภทฺทาลิ โหติ, ยถา ตํ สตฺถุสาสเน สิกฺขาย ปริปูรการิสฺส.}

\prose{ปุน จปรํ ภทฺทาลิ ภิกฺขุ ปีติยา จ วิราคา อุเปกฺขโก จ วิหรติ สโต จ สมฺปชาโน, สุขญฺจ กาเยน ปฏิสํเวเทติ, ยํ ตํ อริยา อาจิกฺขนฺติ “อุเปกฺขโก สติมา สุขวิหารี”ติ ตติยํ}

\vspace{\tipitakaparskip}
\csromancenter{พทฺทาลิสุตฺต (65) ฌานํ อุปสมฺปชฺช วิหรติ.\csromanspacer{} ตํ กิสฺส เหตุ, เอวญฺหิ ตํ ภทฺทาลิ โหติ, ยถา ตํ สตฺถุสาสเน สิกฺขาย ปริปูรการิสฺส.}
\prose{\csromanfolio{105}ปุน จปรํ ภทฺทาลิ ภิกฺขุ สุขสฺส จ ปหานา ทุกฺขสฺส จ ปหานา ปุพฺเพว โสมนสฺสโทมนสฺสานํ อตฺถงฺคมา อทุกฺขมสุขํ อุเปกฺขาสติปาริสุทฺธิํ จตุตฺถํ ฌานํ อุปสมฺปชฺช วิหรติ.\csromanspacer{} ตํ กิสฺส เหตุ, เอวญฺหิ ตํ ภทฺทาลิ โหติ, ยถา ตํ สตฺถุสาสเน สิกฺขาย ปริปูรการิสฺส.}

\prose{โส เอวํ สมาหิเต จิตฺเต ปริสุทฺเธ ปริโยทาเต อนงฺคเณ วิคตูปกฺกิเลเส มุทุภูเต กมฺมนิเย ฐิเต อาเนญฺชปฺปตฺเต ปุพฺเพนิวาสานุสฺสติญาณาย จิตฺตํ อภินินฺนาเมติ.\csromanspacer{} โส อเนกวิหิตํ ปุพฺเพนิวาสํ อนุสฺสรติ, เสยฺยถิทํ, เอกมฺปิ ชาติํ เทฺวปิ ชาติโย ฯเปฯ อิติ สาการํ ส-อุทฺเทสํ อเนกวิหิตํ ปุพฺเพนิวาสํ อนุสฺสรติ.\csromanspacer{} ตํ กิสฺส เหตุ, เอวญฺหิ ตํ ภทฺทาลิ โหติ, ยถา ตํ สตฺถุสาสเน สิกฺขาย ปริปูรการิสฺส.}

\prose{โส เอวํ สมาหิเต จิตฺเต ปริสุทฺเธ ปริโยทาเต อนงฺคเณ วิคตูปกฺกิเลเส มุทุภูเต กมฺมนิเย ฐิเต อาเนญฺชปฺปตฺเต สตฺตานํ จุตูปปาตญาณาย จิตฺตํ อภินินฺนาเมติ.\csromanspacer{} โส ทิพฺเพน จกฺขุนา วิสุทฺเธน อติกฺกนฺตมานุสเกน สตฺเต ปสฺสติ จวมาเน อุปปชฺชมาเน หีเน ปณีเต สุวณฺเณ ทุพฺพณฺเณ สุคเต ทุคฺคเต, ยถากมฺมูปเค สตฺเต ปชานาติ “อิเม วต โภนฺโต สตฺตา กายทุจฺจริเตน สมนฺนาคตา ฯเปฯ วินิปาตํ นิรยํ อุปปนฺนา.\csromanspacer{} อิเม วา ปน โภนฺโต สตฺตา กายสุจริเตน สมนฺนาคตา ฯเปฯ สุคติํ สคฺคํ โลกํ อุปปนฺนา”ติ.\csromanspacer{} อิติ ทิพฺเพน จกฺขุนา วิสุทฺเธน อติกฺกนฺตมานุสเกน ฯเปฯ ยถากมฺมูปเค สตฺเต ปชานาติ.\csromanspacer{} ตํ กิสฺส เหตุ, เอวญฺหิ ตํ ภทฺทาลิ โหติ, ยถา ตํ สตฺถุสาสเน สิกฺขาย ปริปูรการิสฺส.}

\prose{โส เอวํ สมาหิเต จิตฺเต ปริสุทฺเธ ปริโยทาเต อนงฺคเณ วิคตูปกฺกิเลเส มุทุภูเต กมฺมนิเย ฐิเต อาเนญฺชปฺปตฺเต อาสวานํ ขยญาณาย จิตฺตํ อภินินฺนาเมติ.\csromanspacer{} โส อิทํ ทุกฺขนฺติ ยถาภูตํ ปชานาติ, อยํ ทุกฺขสมุทโยติ ยถาภูตํ ปชานาติ, อยํ ทุกฺขนิโรโธติ ยถาภูตํ ปชานาติ, อยํ ทุกฺขนิโรธคามินี ปฏิปทาติ ยถาภูตํ ปชานาติ, อิเม อาสวาติ ยถาภูตํ ปชานาติ, \csromanfolio{106}อยํ อาสวสมุทโยติ ยถาภูตํ ปชานาติ, อยํ อาสวนิโรโธติ ยถาภูตํ ปชานาติ, อยํ อาสวนิโรธคามินี ปฏิปทาติ ยถาภูตํ ปชานาติ.\csromanspacer{} ตสฺส เอวํ ชานโต เอวํ ปสฺสโต กามาสวาปิ จิตฺตํ วิมุจฺจติ, ภวาสวาปิ จิตฺตํ วิมุจฺจติ, อวิชฺชาสวาปิ จิตฺตํ วิมุจฺจติ, วิมุตฺตสฺมิํ วิมุตฺตมิติ ญาณํ โหติ, “ขีณา ชาติ, วุสิตํ พฺรหฺมจริยํ, กตํ กรณียํ, นาปรํ อิตฺถตฺตายา”ติ ปชานาติ.\csromanspacer{} ตํ กิสฺส เหตุ, เอวญฺหิ ตํ ภทฺทาลิ โหติ, ยถา ตํ สตฺถุสาสเน สิกฺขาย ปริปูรการิสฺสาติ.}

\proseitem{140}{เอวํ วุตฺเต อายสฺมา ภทฺทาลิ ภควนฺตํ เอตทโวจ “โก นุ โข ภนฺเต เหตุ, โก ปจฺจโย, เยน มิเธกจฺจํ ภิกฺขุํ ปสยฺห ปสยฺห\footnote{ปวยฺห ปวยฺห (สี, สฺยา, กํ, อิ)} การณํ กโรนฺติ, โก ปน ภนฺเต เหตุ, โก ปจฺจโย, เยน มิเธกจฺจํ ภิกฺขุํ โน ตถา ปสยฺห ปสยฺห การณํ กโรนฺตี”ติ.\csromanspacer{} อิธ ภทฺทาลิ เอกจฺโจ ภิกฺขุ อภิณฺหาปตฺติโก โหติ, อาปตฺติพหุโล.\csromanspacer{} โส ภิกฺขูหิ วุจฺจมาโน อญฺเญนญฺญํ ปฏิจรติ, พหิทฺธา กถํ อปนาเมติ, โกปญฺจ โทสญฺจ อปฺปจฺจยญฺจ ปาตุกโรติ, น สมฺมา วตฺตติ, น โลมํ ปาเตติ, น เนตฺถารํ วตฺตติ, เยน สํโฆ อตฺตมโน โหติ, ตํ กโรมีติ นาห.\csromanspacer{} ตตฺร ภทฺทาลิ ภิกฺขูนํ เอวํ โหติ “อยํ โข อาวุโส ภิกฺขุ อภิณฺหาปตฺติโก, อาปตฺติพหุโล.\csromanspacer{} โส ภิกฺขูหิ วุจฺจมาโน อญฺเญนญฺญํ ปฏิจรติ, พหิทฺธา กถํ อปนาเมติ, โกปญฺจ โทสญฺจ อปฺปจฺจยญฺจ ปาตุกโรติ, น สมฺมา วตฺตติ, น โลมํ ปาเตหิ, น เนตฺถารํ วตฺตติ, เยน สํโฆ อตฺตมโน โหติ, ตํ กโรมีติ นาห.\csromanspacer{} สาธุ วตายสฺมนฺโต อิมสฺส ภิกฺขุโน ตถา ตถา อุปปริกฺขถ, ยถาสฺสิทํ\footnote{ยถยิทํ (สฺยา, กํ, ก)} อธิกรณํ น ขิปฺปเมว วูปสเมยฺยา”ติ.\csromanspacer{} ตสฺส โข เอวํ ภทฺทาลิ ภิกฺขุโน ภิกฺขู ตถา ตถา อุปปริกฺขนฺติ, ยถาสฺสิทํ อธิกรณํ น ขิปฺปเมว วูปสมฺมติ.}

\proseitem{141}{อิธ ปน ภทฺทาลิ เอกจฺโจ ภิกฺขุ อภิณฺหาปตฺติโก โหติ, อาปตฺติพหุโล.\csromanspacer{} โส ภิกฺขูหิ วุจฺจมาโน นาญฺเญนญฺญํ ปฏิจรติ, พหิทฺธา กถํ น อปนาเมติ, น โกปญฺจ โทสญฺจ อปฺปจฺจยญฺจ ปาตุ กโรติ, สมฺมา วตฺตติ, โลมํ ปาเตติ, เนตฺถารํ วตฺตติ, เยน สํโฆ}

\vspace{\tipitakaparskip}
\csromancenter{พทฺทาลิสุตฺต (65) อตฺตมโน โหติ, ตํ กโรมีติ อาห.\csromanspacer{} ตตฺร ภทฺทาลิ ภิกฺขูนํ เอวํ โหติ “อยํ โข อาวุโส ภิกฺขุ อภิณฺหาปตฺติโก, อาปตฺติพหุโล.\csromanspacer{} โส ภิกฺขูหิ วุจฺจมาโน นาญฺเญนญฺญํ ปฏิจรติ, พหิทฺธา กถํ น อปนาเมติ, น โกปญฺจ โทสญฺจ อปฺปจฺจยญฺจ ปาตุกโรติ, สมฺมา วตฺตติ, โลมํ ปาเตติ, เนตฺถารํ วตฺตติ, เยน สํโฆ อตฺตมโน โหติ, ตํ กโรมีติ อาห.\csromanspacer{} สาธุ วตายสฺมนฺโต อิมสฺส ภิกฺขุโน ตถา ตถา อุปปริกฺขถ, ยถาสฺสิทํ อธิกรณํ ขิปฺปเมว วูปสเมยฺยา”ติ.\csromanspacer{} ตสฺส โข เอวํ ภทฺทาลิ ภิกฺขุโน ภิกฺขู ตถา ตถา อุปปริกฺขนฺติ, ยถาสฺสิทํ อธิกรณํ ขิปฺปเมว วูปสมฺมติ.}
\proseitem{142}{\csromanfolio{107}อิธ ภทฺทาลิ เอกจฺโจ ภิกฺขุ อธิจฺจาปตฺติโก โหติ, อนาปตฺติพหุโล.\csromanspacer{} โส ภิกฺขูหิ วุจฺจมาโน อญฺเญนญฺญํ ปฏิจรติ, พหิทฺธา กถํ อปนาเมติ, โกปญฺจ โทสญฺจ อปฺปจฺจยญฺจ ปาตุกโรติ, น สมฺมา วตฺตติ, น โลมํ ปาเตติ, น เนตฺถารํ วตฺตติ, เยน สํโฆ อตฺตมโน โหติ, ตํ กโรมีติ นาห.\csromanspacer{} ตตฺร ภทฺทาลิ ภิกฺขูนํ เอวํ โหติ “อยํ โข อาวุโส ภิกฺขุ อธิจฺจาปตฺติโก, อนาปตฺติพหุโล.\csromanspacer{} โส ภิกฺขูหิ วุจฺจมาโน อญฺเญนญฺญํ ปฏิจรติ, พหิทฺธา กถํ อปนาเมติ, โกปญฺจ โทสญฺจ อปฺปจฺจยญฺจ ปาตุกโรติ, น สมฺมา วตฺตติ, น โลมํ ปาเตติ, น เนตฺถารํ วตฺตติ, เยน สํโฆ อตฺตมโน โหติ, ตํ กโรมีติ นาห.\csromanspacer{} สาธุ วตายสฺมนฺโต อิมสฺส ภิกฺขุโน ตถา ตถา อุปปริกฺขถ, ยถาสฺสิทํ อธิกรณํ น ขิปฺปเมว วูปสเมยฺยา”ติ.\csromanspacer{} ตสฺส โข เอวํ ภทฺทาลิ ภิกฺขุโน ภิกฺขู ตถา ตถา อุปปริกฺขนฺติ, ยถาสฺสิทํ อธิกรณํ น ขิปฺปเมว วูปสมฺมติ.}

\proseitem{143}{อิธ ปน ภทฺทาลิ เอกจฺโจ ภิกฺขุ อธิจฺจาปตฺติโก โหติ, อนาปตฺถิพหุโล.\csromanspacer{} โส ภิกฺขูหิ วุจฺจมาโน นาญฺเญนญฺญํ ปฏิจรติ, น พหิทฺธา กถํ อปนาเมติ, น โกปญฺจ โทสญฺจ อปฺปจฺจยญฺจ ปาตุกโรติ, สมฺมา วตฺตติ, โลมํ ปาเตติ, เนตฺถารํ วตฺตติ, เยน สํโฆ อตฺตมโน โหติ, ตํ กโรมีติ อาห.\csromanspacer{} ตตฺร ภทฺทาลิ ภิกฺขูนํ เอวํ โหติ “อยํ โข อาวุโส ภิกฺขุ อธิจฺจาปตฺติโก อนาปตฺติพหุโล.\csromanspacer{} โส ภิกฺขูหิ วุจฺจมาโน นาญฺเญนญฺญํ ปฏิจรติ, น พหิทฺธา กถํ อปนาเมติ, น โกปญฺจ โทสญฺจ อปฺปจฺจยญฺจ ปาตุกโรติ, สมฺมา วตฺตติ, โลมํ ปาเตติ, เนตฺถารํ วตฺตติ, เยน \csromanfolio{108}สํโฆ อตฺตมโน โหติ, ตํ กโรมีติ อาห.\csromanspacer{} สาธุ วตายสฺมนฺโต อิมสฺส ภิกฺขุโน ตถา ตถา อุปปริกฺขถ, ยถาสฺสิทํ อธิกรณํ ขิปฺปเมว วูปสเมยฺยา”ติ.\csromanspacer{} ตสฺส โข เอวํ ภทฺทาลิ ภิกฺขุโน ภิกฺขู ตถา ตถา อุปปริกฺขนฺติ, ยถาสฺสิทํ อธิกรณํ ขิปฺปเมว วูปสมฺมติ.}

\proseitem{144}{อิธ ภทฺทาลิ เอกจฺโจ ภิกฺขุ สทฺธามตฺตเกน วหติ เปมมตฺตเกน.\csromanspacer{} ตตฺร ภทฺทาลิ ภิกฺขูนํ เอวํ โหติ “อยํ โข อาวุโส ภิกฺขุ สทฺธามตฺตเกน วหติ เปมมตฺตเกน, สเจ มยํ อิมํ ภิกฺขุํ ปสยฺห ปสยฺห การณํ กริสฺสาม, มา ยํปิสฺส ตํ สทฺธามตฺตกํ เปมมตฺตกํ, ตมฺหาปิ ปริหายี”ติ.\csromanspacer{} เสยฺยถาปิ ภทฺทาลิ ปุริสสฺส เอกํ จกฺขุํ, ตสฺส มิตฺตามจฺจา ญาติสาโลหิตา ตํ เอกํ จกฺขุํ รกฺเขยฺยุํ “มา ยํปิสฺส ตํ เอกํ จกฺขุํ, ตมฺหาปิ ปริหายี”ติ.\csromanspacer{} เอวเมว โข ภทฺทาลิ อิเธกจฺโจ ภิกฺขุ สทฺธามตฺตเกน วหติ เปมมตฺตเกน.\csromanspacer{} ตตฺร ภทฺทาลิ ภิกฺขูนํ เอวํ โหติ “อยํ โข อาวุโส ภิกฺขุ สทฺธามตฺตเกน วหติ เปมมตฺตเกน.\csromanspacer{} สเจ มยํ อิมํ ภิกฺขุํ ปสยฺห ปสยฺห การณํ กริสฺสาม, มา ยํปิสฺส ตํ สทฺธามตฺตกํ เปมมตฺตกํ, ตมฺหาปิ ปริหายี”ติ.\csromanspacer{} อยํ โข ภทฺทาลิ เหตุ อยํ ปจฺจโย, เยน มิเธกจฺจํ ภิกฺขุํ ปสยฺห ปสยฺห การณํ กโรนฺติ.\csromanspacer{} อยํ ปน ภทฺทาลิ เหตุ อยํ ปจฺจโย, เยน มิเธกจฺจํ ภิกฺขุํ โน ตถา ปสยฺห ปสยฺห การณํ กโรนฺตีติ.}

\proseitem{145}{โก นุ โข ภนฺเต เหตุ, โก ปจฺจโย, เยน ปุพฺเพ อปฺปตรานิ เจว สิกฺขาปทานิ อเหสุํ, พหุตรา จ ภิกฺขู อญฺญาย สณฺฐหิํสุ.\csromanspacer{} โก ปน ภนฺเต เหตุ, โก ปจฺจโย, เยน เอตรหิ พหุตรานิ เจว สิกฺขาปทานิ โหนฺติ, อปฺปตรา จ ภิกฺขู อญฺญาย สณฺฐหนฺตีติ.\csromanspacer{} เอวเมตํ ภทฺทาลิ โหติ, สตฺเตสุ หายมาเนสุ สทฺธมฺเม อนฺตรธายมาเน พหุตรานิ เจว สิกฺขาปทานิ โหนฺติ, อปฺปตรา จ ภิกฺขู อญฺญาย สณฺฐหนฺตีติ.\csromanspacer{} น ตาว ภทฺทาลิ สตฺถา สาวกานํ สิกฺขาปทํ ปญฺญาเปติ, ยาว น อิเธกจฺเจ อาสวฏฺฐานียา ธมฺมา สํเฆ ปาตุภวนฺติ, ยโต จ โข ภทฺทาลิ อิเธกจฺเจ อาสวฏฺฐานียา ธมฺมา สํเฆ ปาตุภวนฺติ, อถ สตฺถา สาวกานํ สิกฺขาปทํ ปญฺญาเปติ เตสํเยว อาสวฏฺฐานียานํ ธมฺมานํ ปฏิฆาตาย.\csromanspacer{} น ตาว ภทฺทาลิ อิเธกจฺเจ อาสวฏฺฐานียา ธมฺมา สํเฆ ปาตุภวนฺติ, ยาว น สํโฆ มหตฺตํ ปตฺโต โหติ.\csromanspacer{} ยโต จ โข ภทฺทาลิ สํโฆ มหตฺตํ ปตฺโต โหติ, อถ}

\vspace{\tipitakaparskip}
\csromancenter{พทฺทาลิสุตฺต (65) อิเธกจฺเจ อาสวฏฺฐานียา ธมฺมา สํเฆ ปาตุภวนฺติ, อถ สตฺถา สาวกานํ สิกฺขาปทํ ปญฺญาเปติ เตสํเยว อาสวฏฺฐานียานํ ธมฺมานํ ปฏิฆาตาย.\csromanspacer{} น ตาว ภทฺทาลิ อิเธกจฺเจ อาสวฏฺฐานียา ธมฺมา สํเฆ ปาตุภวนฺติ, ยาว น สํโฆ ลาภคฺคํ ปตฺโต โหติ.\csromanspacer{} ยสคฺคํ ปตฺโต โหติ.\csromanspacer{} พาหุสจฺจํ ปตฺโต โหติ.\csromanspacer{} รตฺตญฺญุตํ ปตฺโต โหติ.\csromanspacer{} ยโต จ โข ภทฺทาลิ สํโฆ รตฺตญฺญุตํ ปตฺโต โหติ, อถ อิเธกจฺเจ อาสวฏฺฐานียา ธมฺมา สํเฆ ปาตุภวนฺติ, อถ สตฺถา สาวกานํ สิกฺขาปทํ ปญฺญาเปติ เตสํเยว อาสวฏฺฐานียานํ ธมฺมานํ ปฏิฆาตาย.}
\proseitem{146}{\csromanfolio{109}อปฺปกา โข ตุเมฺห ภทฺทาลิ เตน สมเยน อหุวตฺถ, ยทา โว อหํ อาชานียสุสูปมํ ธมฺมปริยายํ เทเสสิํ, ตํ สรสิ\footnote{สรสิ ตฺวํ (สี, อิ), สรสิ ตํ (?)} ภทฺทาลีติ.\csromanspacer{} โน เหตํ ภนฺเต.\csromanspacer{} ตตฺร ภทฺทาลิ กํ เหตุํ ปจฺเจสีติ.\csromanspacer{} โส หิ นูนาหํ ภนฺเต ทีฆรตฺตํ สตฺถุสาสเน สิกฺขาย อปริปูรการี อโหสินฺติ.\csromanspacer{} น โข ภทฺทาลิ เอเสว เหตุ เอส ปจฺจโย, อปิ จ เม ตฺวํ ภทฺทาลิ ทีฆรตฺตํ เจตสา เจโตปริจฺจ วิทิโต “น จายํ โมฆปุริโส มยา ธมฺเม เทสิยมาเน อฏฺฐิํ กตฺวา มนสิ กตฺวา สพฺพเจตโส\footnote{สพฺพํ เจตโส (ก)} สมนฺนาหริตฺวา โอหิตโสโต ธมฺมํ สุณาตี”ติ.\csromanspacer{} อปิ จ เต อหํ ภทฺทาลิ อาชานียสุสูปมํ ธมฺมปริยายํ เทเสสฺสามิ, ตํ สุณาหิ สาธุกํ มนสิ กโรหิ, ภาสิสฺสามีติ.\csromanspacer{} เอวํ ภนฺเตติ โข อายสฺมา ภทฺทาลิ ภควโต ปจฺจสฺโสสิ.\csromanspacer{} ภควา เอตทโวจ\csromandash{}}

\proseitem{147}{เสยฺยถาปิ ภทฺทาลิ ทกฺโข อสฺสทมโก ภทฺรํ อสฺสาชานียํ ลภิตฺวา ปฐเมเนว มุขาธาเน การณํ กาเรติ.\csromanspacer{} ตสฺส มุขาธาเน การณํ การิยมานสฺส โหนฺติเยว วิสูกายิตานิ วิเสวิตานิ วิปฺผนฺทิตานิ กานิจิ กานิจิ, ยถา ตํ อการิตปุพฺพํ การณํ การิยมานสฺส.\csromanspacer{} โส อภิณฺหการณา อนุปุพฺพการณา ตสฺมิํ ฐาเน ปรินิพฺพายติ.\csromanspacer{} ยโต โข ภทฺทาลิ ภโทฺร อสฺสาชานีโย อภิณฺหการณา อนุปุพฺพการณา ตสฺมิํ ฐาเน ปรินิพฺพุโต โหติ, ตเมนํ อสฺสทมโก อุตฺตริ การณํ กาเรติ ยุคาธาเน.\csromanspacer{} ตสฺส ยุคาธาเน การณํ การิยมานสฺส โหนฺติเยว วิสูกายิตานิ วิเสวิตานิ วิปฺผนฺทิตานิ กานิจิ กานิจิ, ยถา ตํ อการิตปุพฺพํ การณํ การิยมานสฺส.\csromanspacer{} โส อภิณฺหการณา อนุปุพฺพการณา ตสฺมิํ ฐาเน \csromanfolio{110}ปรินิพฺพายติ.\csromanspacer{} ยโต โข ภทฺทาลิ ภโทฺร อสฺสาชานีโย อภิณฺหการณา อนุปุพฺพการณา ตสฺมิํ ฐาเน ปรินิพฺพุโต โหติ, ตเมนํ อสฺสทมโก อุตฺตริ การณํ กาเรติ อนุกฺกเม, มณฺฑเล, ขุรกาเส\footnote{ขุรกาเย (สี, อิ)}, ธาเว, ทวตฺเต\footnote{รวตฺเถ (สี, สฺยา, กํ, อิ)}, ราชคุเณ, ราชวํเส, อุตฺตเม ชเว อุตฺตเม หเย อุตฺตเม สาขเลฺย.\csromanspacer{} ตสฺส อุตฺตเม ชเว อุตฺตเม หเย อุตฺตเม สาขเลฺย การณํ การิยมานสฺส โหนฺติเยว วิสูกายิตานิ วิเสวิตานิ วิปฺผนฺทิตานิ กานิจิ กานิจิ, ยถา ตํ อการิตปุพฺพํ การณํ การิยมานสฺส, โส อภิณฺหการณา อนุปุพฺพการณา ตสฺมิํ ฐาเน ปรินิพฺพายติ.\csromanspacer{} ยโต โข ภทฺทาลิ ภโทฺร อสฺสชานีโย อภิณฺห การณา อนุปุพฺพการณา ตสฺมิํ ฐาเน ปรินิพฺพุโต โหติ, ตเมนํ อสฺสทมโก อุตฺตริ วณฺณิยญฺจ ปาณิยญฺจ\footnote{วลิยญฺจ (สี, อิ), พลิยญฺจ (สฺยา, กํ)} อนุปฺปเวจฺฉติ.\csromanspacer{} อิเมหิ โข ภทฺทาลิ ทสหงฺเคหิ สมนฺนาคโต ภโทฺร อสฺสาชานีโย ราชารโห โหติ ราชโภคฺโค, รญฺโญ องฺคนฺเตว สงฺขฺยํ คจฺฉติ.}

\prose{เอวเมว โข ภทฺทาลิ ทสหิ ธมฺเมหิ สมนฺนาคโต ภิกฺขุ อาหุเนยฺโย โหติ ปาหุเนยฺโย ทกฺขิเณยฺโย อญฺชลิกรณีโย อนุตฺตรํ ปุญฺญกฺเขตฺตํ โลกสฺส.\csromanspacer{} กตเมหิ ทสหิ.\csromanspacer{} อิธ ภทฺทาลิ ภิกฺขุ อเสขาย สมฺมาทิฏฺฐิยา สมนฺนาคโต โหติ, อเสเขน สมฺมาสงฺกปฺเปน สมนฺนาคโต โหติ, อเสขาย สมฺมาวาจาย สมนฺนาคโต โหติ, อเสเขน สมฺมากมฺมนฺเตน สมนฺนาคโต โหติ, อเสเขน สมฺมา-อาชีเวน สมนฺนาคโต โหติ, อเสเขน สมฺมาวายาเมน สมนฺนาคโต โหติ, อเสขาย สมฺมาสติยา สมนฺนาคโต โหติ, อเสเขน สมฺมาสมาธินา สมนฺนาคโต โหติ, อเสเขน สมฺมาญาเณน สมนฺนาคโต โหติ, อเสขาย สมฺมาวิมุตฺติยา สมนฺนาคโต โหติ, อิเมหิ โข ภทฺทาลิ ทสหิ ธมฺเมหิ สมนฺนาคโต ภิกฺขุ อาหุเนยฺโย โหติ ปาหุเนยฺโย ทกฺขิเณยฺโย อญฺชลิกรณีโย อนุตฺตรํ ปุญฺญกฺเขตฺตํ โลกสฺสาติ.}

\prose{อิทมโวจ ภควา.\csromanspacer{} อตฺตมโน อายสฺมา ภทฺทาลิ ภควโต ภาสิตํ อภินนฺทีติ.}

\nitthitam{ภทฺทาลิสุตฺตํ นิฏฺฐิตํ ปญฺจมํ.}
\csromanmatikaanchor{mtk.19}
\csromantocmarkref{subsection}{6. ลฏุกิโกปมสุตฺต}{mtk.19}
\bookmark[dest={mtk.19},level=2]{6. ลฏุกิโกปมสุตฺต}
\csromanheader{2}{6. ลฏุกิโกปมสุตฺต (66) \textbf{6. ลฏุกิโกปมสุตฺต}}
\proseitem{148}{\csromanfolio{111}เอวํ เม สุตํ\csromandash{}เอกํ สมยํ ภควา องฺคุตฺตราเปสุ วิหรติ อาปณํ นาม องฺคุตฺตราปานํ นิคโม.\csromanspacer{} อถ โข ภควา ปุพฺพณฺหสมยํ นิวาเสตฺวา ปตฺตจีวรมาทาย อาปณํ ปิณฺฑาย ปาวิสิ, อาปเณ ปิณฺฑาย จริตฺวา ปจฺฉาภตฺตํ ปิณฺฑปาตปฏิกฺกนฺโต เยนญฺญตโร วนสณฺโฑ เตนุปสงฺกมิ ทิวาวิหาราย, ตํ วนสณฺฑํ อชฺโฌคาเหตฺวา อญฺญตรสฺมิํ รุกฺขมูเล ทิวาวิหารํ นิสีทิ.\csromanspacer{} อายสฺมาปิ โข อุทายี ปุพฺพณฺหสมยํ นิวาเสตฺวา ปตฺตจีวรมาทาย อาปณํ ปิณฺฑาย ปาวิสิ, อาปเณ ปิณฺฑาย จริตฺวา ปจฺฉาภตฺตํ ปิณฺฑปาตปฏิกฺกนฺโต เยน โส วนสณฺโฑ เตนุปสงฺกมิ ทิวาวิหาราย, ตํ วนสณฺฑํ อชฺโฌคาเหตฺวา อญฺญตรสฺมิํ รุกฺขมูเล ทิวาวิหารํ นิสีทิ.\csromanspacer{} อถ โข อายสฺมโต อุทายิสฺส รโหคตสฺส ปฏิสลฺลีนสฺส เอวํ เจตโส ปริวิตกฺโก อุทปาทิ “พหูนํ\footnote{พหุนฺนํ (สี, สฺยา, กํ, อิ) เอวมีทิเส อวิญฺญาณกปฺปกรเณ.} วต โน ภควา ทุกฺขธมฺมานํ อปหตฺตา, พหูนํ วต โน ภควา สุขธมฺมานํ อุปหตฺตา, พหูนํ วต โน ภควา อกุสลานํ ธมฺมานํ อปหตฺตา, พหูนํ วต โน ภควา กุสลานํ ธมฺมานํ อุปหตฺตา”ติ.\csromanspacer{} อถ โข อายสฺมา อุทายี สายนฺหสมยํ ปฏิสลฺลานา วุฏฺฐิโต เยน ภควา เตนุปสงฺกมิ, อุปสงฺกมิตฺวา ภควนฺตํ อภิวาเทตฺวา เอกมนฺตํ นิสีทิ.}

\proseitem{149}{เอกมนฺตํ นิสินฺโน โข อายสฺมา อุทายี ภควนฺตํ เอตทโวจ “อิธ มยฺหํ ภนฺเต รโหคกสฺส ปฏิสลฺลีนสฺส เอวํ เจตโส ปริวิตกฺโก อุทปาทิ ‘พหูนํ วต โน ภควา ทุกฺขธมฺมานํ อปหตฺตา, พหูนํ วต โน ภควา สุขธมฺมานํ อุปหตฺตา, พหูนํ วต โน ภควา อกุสลานํ ธมฺมานํ อปหตฺตา, พหูนํ วต โน ภควา กุสลานํ ธมฺมานํ อุปหตฺตา’ติ, มยญฺหิ ภนฺเต ปุพฺเพ สายญฺเจว ภุญฺชาม ปาโต จ ทิวา จ วิกาเล, อหุ โข โส ภนฺเต สมโย ยํ ภควา ภิกฺขู อามนฺเตสิ ‘อิงฺฆ ตุเมฺห ภิกฺขเว เอตํ ทิวาวิกาลโภชนํ ปชหถา’ติ, ตสฺส มยฺหํ ภนฺเต อหุเทว อญฺญถตฺตํ อหุเทว\footnote{อหุ (สี, อิ)} โทมนสฺสํ ‘ยมฺปิ โน สทฺธา คหปติกา ทิวา วิกาเล ปณีตํ ขาทนียํ โภชนียํ เทนฺติ, ตสฺสปิ โน ภควา ปหานมาห, ตสฺสปิ โน สุคโต ปฏินิสฺสคฺคมาหา’ติ, \csromanfolio{112}เต มยํ ภนฺเต ภควติ เปมญฺจ คารวญฺจ หิริญฺจ โอตฺตปฺปญฺจ สมฺปสฺสมานา เอวํ ตํ ทิวาวิกาลโภชนํ ปชหิมฺหา, เต มยํ ภนฺเต สายญฺเจว ภุญฺชาม ปาโต จ.\csromanspacer{} อหุ โข โส ภนฺเต สมโย, ยํ ภควา ภิกฺขู อามนฺเตสิ ‘อิงฺฆ ตุเมฺห ภิกฺขเว เอตํ รตฺติํวิกาลโภชนํ ปชหถา’ติ, ตสฺส มยฺหํ ภนฺเต อหุเทว อญฺญถตฺตํ อหุเทว โทมนสฺสํ ‘ยมฺปิ โน อิเมสํ ทฺวินฺนํ ภตฺตานํ ปณีตสงฺขาตตรํ, ตสฺสปิ โน ภควา ปหานมาห, ตสฺสปิ โน สุคโต ปฏินิสฺสคฺคมาหา’ติ.\csromanspacer{} ภูตปุพฺพํ ภนฺเต อญฺญตโร ปุริโส ทิวา สูเปยฺยํ ลภิตฺวา เอวมาห ‘หนฺท จ อิมํ นิกฺขิปถ, สายํ สพฺเพว สมคฺคา ภุญฺชิสฺสามา’ติ.\csromanspacer{} ยา กาจิ ภนฺเต สงฺขติโย, สพฺพา ตา รตฺติํ, อปฺปา ทิวา.\csromanspacer{} เต มยํ ภนฺเต ภควติ เปมญฺจ คารวญฺจ หิริญฺจ โอตฺตปฺปญฺจ สมฺปสฺสมานา เอวํ ตํ รตฺติํวิกาลโภชนํ ปชหิมฺหา.\csromanspacer{} ภูตปุพฺพํ ภนฺเต ภิกฺขู รตฺตนฺธการติมิสายํ ปิณฺฑาย จรนฺตา จนฺทนิกํปิ ปวิสนฺติ, โอลิคลฺเลปิ ปปตนฺติ, กณฺฏกาวาฏํปิ\footnote{กณฺฏกวตฺตมฺปิ (สี, อิ), กณฺฏกราชิมฺปิ (สฺยา, กํ)} อาโรหนฺติ, สุตฺตํปิ คาวิํ อาโรหนฺติ, มาณเวหิปิ สมาคจฺฉนฺติ กตกมฺเมหิปิ อกตกมฺเมหิปิ, มาตุคาโมปิ เต\footnote{เตน (ก)} อสทฺธมฺเมน นิมนฺเตติ.\csromanspacer{} ภูตปุพฺพาหํ ภนฺเต รตฺตนฺธการติมิสายํ ปิณฺฑาย จรามิ.\csromanspacer{} อทฺทสา โข มํ ภนฺเต อญฺญตรา อิตฺถี วิชฺชนฺตริกาย ภาชนํ โธวนฺตี, ทิสฺวา มํ ภีตา วิสฺสรมกาสิ ‘อภุมฺเม\footnote{อพฺภุมฺเม (สี, อิ)} ปิสาโจ วต มนฺ’ติ.\csromanspacer{} เอวํ วุตฺเต อหํ ภนฺเต ตํ อิตฺถิํ เอตทโวจํ ‘นาหํ ภคินิ ปิสาโจ, ภิกฺขุ ปิณฺฑาย ฐิโต’ติ.\csromanspacer{} ภิกฺขุสฺส อาตุมารี, ภิกฺขุสฺส มาตุมารี\footnote{ฐิโต’ติ. ภิกฺขุสฺส อาตุมาตุมารี (ก)}, วรํ เต ภิกฺขุ ติเณฺหน โควิกนฺตเนน กุจฺฉิ ปริกนฺโต, น เตฺวว วรํ, ยํ\footnote{น เตฺวว ยา (สี, อิ)} รตฺตนฺธการติมิสายํ กุจฺฉิเหตุ ปิณฺฑาย จรสีติ\footnote{จรสาติ (สี, อิ)}.\csromanspacer{} ตสฺส มยฺหํ ภนฺเต ตทนุสฺสรโต เอวํ โหติ ‘พหูนํ วต โน ภควา ทุกฺขธมฺมานํ อปหตฺตา, พหูนํ วต โน ภควา สุขธมฺมานํ อุปหตฺตา, พหูนํ วต โน ภควา อกุสลานํ ธมฺมานํ อปหตฺตา, พหูนํ วต โน ภควา กุสลานํ ธมฺมานํ อุปหตฺตา’ติ”.}

\proseitem{150}{เอวเมว ปนุทายิ, อิเธกจฺเจ โมฆปุริสา “อิทํ ปชถา”ติ มยา วุจฺจมานา เต เอวมาหํสุ “กิํ ปนิมสฺส อปฺปมตฺตกสฺส โอรมตฺตกสฺส, อธิสลฺลิขเตวายํ สมโณ”ติ.\csromanspacer{} เต ตญฺเจว นปฺปชหนฺติ,}

\vspace{\tipitakaparskip}
\csromancenter{ลฏุกิโกปมสุตฺต (66) มยิ จ อปฺปจฺจยํ อุปฏฺฐาเปนฺติ เย จ ภิกฺขู สิกฺขากามา.\csromanspacer{} เตสํ ตํ อุทายิ โหติ พลวํ พนฺธนํ ทฬฺหํ พนฺธนํ ถิรํ พนฺธนํ อปูติกํ พนฺธนํ ถูโล กลิงฺคโร.\csromanspacer{} เสยฺยถาปิ อุทายิ ลฏุกิกา สกุณิกา ปูติลตาย พนฺธเนน พทฺธา ตตฺเถว วธํ วา พนฺธํ วา มรณํ วา อาคเมติ.\csromanspacer{} โย นุ โข อุทายิ เอวํ วเทยฺย “เยน สา ลฏุกิกา สกุณิกา ปูติลตาย พนฺธเนน พทฺธา ตตฺเถว วธํ วา พนฺธํ วา มรณํ วา อาคเมติ.\csromanspacer{} ตญฺหิ ตสฺสา อพลํ พนฺธนํ ทุพฺพลํ พนฺธนํ ปูติกํ พนฺธนํ อสารกํ พนฺธนนฺ”ติ, สมฺมา นุ โข โส อุทายิ วทมาโน วเทยฺยาติ.\csromanspacer{} โน เหตํ ภนฺเต.\csromanspacer{} เยน สา ภนฺเต ลฏุกิกา สกุณิกา ปูติลตาย พนฺธเนน พทฺธา ตตฺเถว วธํ วา พนฺธํ วา มรณํ วา อาคเมติ.\csromanspacer{} ตญฺหิ ตสฺสา พลวํ พนฺธนํ ทฬฺหํ พนฺธนํ ถิรํ พนฺธนํ อปูติกํ พนฺธนํ ถูโล กลิงฺคโรติ.\csromanspacer{} เอวเมว โข อุทายิ อิเธกจฺเจ โมฆปุริสา “อิทํ ปชหถา”ติ มยา วุจฺจมานา เต เอวมาหํสุ “กิํ ปนิมสฺส อปฺปมตฺตกสฺส โอรมตฺตกสฺส, อธิสลฺลิขเตวายํ สมโณ”ติ.\csromanspacer{} เต ตญฺเจว นปฺปชหนฺติ, มยิ จ อปฺปจฺจยํ อุปฏฺฐาเปนฺติ เย จ ภิกฺขู สิกฺขากามา.\csromanspacer{} เตสํ ตํ อุทายิ โหติ พลวํ พนฺธนํ ทฬฺหํ พนฺธนํ ถิรํ พนฺธนํ อปูติกํ พนฺธนํ ถูโล กลิงฺคโร.}
\proseitem{151}{\csromanfolio{113}อิธ ปนุทายิ เอกจฺเจ กุลปุตฺตา “อิทํ ปชหถา”ติ มยา วุจฺจมานา เต เอวมาหํสุ “กิํ ปนิมสฺส อปฺปมตฺตกสฺส โอรมตฺตกสฺส ปหาตพฺพสฺส, ยสฺส โน ภควา ปหานมาห, ยสฺส โน สุคโต ปฏินิสฺสคฺคมาหา”ติ.\csromanspacer{} เต ตญฺเจว ปชหนฺติ, มยิ จ น อปฺปจฺจยํ อุปฏฺฐาเปนฺติ เย จ ภิกฺขู สิกฺขากามา.\csromanspacer{} เต ตํ ปหาย อปฺโปสฺสุกฺกา ปนฺนโลมา ปรทตฺตวุตฺตา\footnote{ปรทวุตฺตา (สี, สฺยา, กํ, อิ)} มิคภูเตน เจตสา วิหรนฺติ, เตสํ ตํ อุทายิ โหติ อพลํ พนฺธนํ ทุพฺพลํ พนฺธนํ ปูติกํ พนฺธนํ อสารกํ พนฺธนํ.\csromanspacer{} เสยฺยถาปิ อุทายิ รญฺโญ นาโค อีสาทนฺโต อุรูฬฺหวา อภิชาโต สงฺคามาวจโร ทเฬฺหหิ วรตฺเตหิ พนฺธเนหิ พทฺโธ อีสกํเยว กายํ สนฺนาเมตฺวา ตานิ พนฺธนานิ สํฉินฺทิตฺวา สํปทาเลตฺวา เยน กามํ ปกฺกมติ.\csromanspacer{} โย นุ โข อุทายิ เอวํ วเทยฺย “เยหิ โส รญฺโญ นาโค อีสาทนฺโต อุรูฬฺหวา อภิชาโต สงฺคามาวจโร ทเฬฺหหิ วรตฺเตหิ พนฺธเนหิ พทฺโธ อีสกํเยว กายํ สนฺนาเมตฺวา ตานิ พนฺธนานิ สํฉินฺทิตฺวา สํปทาเลตฺวา เยน กามํ ปกฺกมติ.\csromanspacer{} ตญฺหิ ตสฺส พลวํ พนฺธนํ ทฬฺหํ พนฺธนํ ถิรํ พนฺธนํ \csromanfolio{114}อปูติกํ พนฺธนํ ถูโล กลิงฺคโร”ติ, สมฺมา นุ โข โส อุทายิ วทมาโน วเทยฺยาติ.\csromanspacer{} โน เหตํ ภนฺเต.\csromanspacer{} เยหิ โส ภนฺเต รญฺโญ นาโค อีสาทนฺโต อุรูฬฺหวา อภิชาโต สงฺคามาวจโร ทเฬฺหหิ วรตฺเตหิ พนฺธเนหิ พทฺโธ อีสกํเยว กายํ สนฺนาเมตฺวา ตานิ พนฺธนานิ สํฉินฺทิตฺวา สํปทาเลตฺวา เยน กามํ ปกฺกมติ.\csromanspacer{} ตญฺหิ ตสฺส อพลํ พนฺธนํ ฯเปฯ อสารกํ พนฺธนนฺติ.\csromanspacer{} เอวเมว โข อุทายิ อิเธกจฺเจ กุลปุตฺตา “อิทํ ปชหถา”ติ มยา วุจฺจมานา เต เอวมาหํสุ “กิํ ปนิมสฺส อปฺปมตฺตกสฺส โอรมตฺตกสฺส ปหาตพฺพสฺส, ยสฺส โน ภควา ปหานมาห, ยสฺส โน สุคโต ปฏินิสฺสคฺคมาหา”ติ.\csromanspacer{} เต ตญฺเจว ปชหนฺติ, มยิ จ น อปฺปจฺจยํ อุปฏฺฐาเปนฺติ เย จ ภิกฺขู สิกฺขากามา.\csromanspacer{} เต ตํ ปหาย อปฺโปสฺสุกฺกา ปนฺนโลมา ปรทตฺตวุตฺตา มิคภูเตน เจตสา วิหรนฺติ, เตสํ ตํ อุทายิ โหติ อพลํ พนฺธนํ ทุพฺพลํ พนฺธนํ ปูติกํ พนฺธนํ อสารกํ พนฺธนํ.}

\proseitem{152}{เสยฺยถาปิ อุทายิ ปุริโส ทลิทฺโท อสฺสโก อนาฬฺหิโย, ตสฺส’สฺส เอกํ อคารกํ โอลุคฺควิลุคฺคํ กากาติทายิํ\footnote{กากาติฑายิํ (?)} นปรมรูปํ, เอกา ขโฏปิกา\footnote{กโฬปิกา (ก)} โอลุคฺควิลุคฺคา นปรมรูปา, เอกิสฺสา กุมฺภิยา ธญฺญสมวาปกํ นปรมรูปํ, เอกา ชายิกา นปรมรูปา.\csromanspacer{} โส อารามคตํ ภิกฺขุํ ปสฺเสยฺย สุโธตหตฺถปาทํ มนุญฺญํ โภชนํ ภุตฺตาวิํ สีตาย ฉายาย นิสินฺนํ อธิจิตฺเต ยุตฺตํ, ตสฺส เอวมสฺส “สุขํ วต โภ สามญฺญํ, อาโรคฺยํ วต โภ สามญฺญํ, โส วตสฺสํ\footnote{โส วตสฺส (ก)}, โยหํ เกสมสฺสุํ โอหาเรตฺวา กาสายานิ วตฺถานิ อจฺฉาเทตฺวา อคารสฺมา อนคาริยํ ปพฺพเชยฺยนฺ”ติ.\csromanspacer{} โส น สกฺกุเณยฺย เอกํ อคารกํ โอลุคฺควิลุคฺคํ กากาติทายิํ นปรมรูปํ ปหาย, เอกํ ขโฏปิกํ โอลุคฺควิลุคฺคํ นปรมรูปํ ปหาย, เอกิสฺสา กุมฺภิยา ธญฺญสมวาปกํ นปรมรูปํ ปหาย, เอกํ ชายิกํ นปรมรูปํ ปหาย เกสมสฺสุํ โอหาเรตฺวา กาสายานิ วตฺถานิ อจฺฉาเทตฺวา อคารสฺมา อนคาริยํ ปพฺพชิตุํ.\csromanspacer{} โย นุ โข อุทายิ เอวํ วเทยฺย “เยหิ โส ปุริโส พนฺธเนหิ พทฺโธ น สกฺโกติ เอกํ อคารกํ โอลุคฺควิลุคฺคํ กากาติทายิํ นปรมรูปํ ปหาย, เอกํ ขโฏปิกํ โอลุคฺควิลุคฺคํ นปรมรูปํ ปหาย, เอกิสฺสา กุมฺภิยา ธญฺญสมวาปกํ นปรมรูปํ ปหาย, เอกํ ชายิกํ นปรมรูปํ ปหาย}

\vspace{\tipitakaparskip}
\csromancenter{ลฏุกิโกปมสุตฺต (66) เกสมสฺสุํ โอหาเรตฺวา กาสายานิ วตฺถานิ อจฺฉาเทตฺวา อคารสฺมา อนคาริยํ ปพฺพชิตุํ.\csromanspacer{} ตญฺหิ ตสฺส อพลํ พนฺธนํ ทุพฺพลํ พนฺธนํ ปูติกํ พนฺธนํ อสารกํ พนฺธนนฺ”ติ, สมฺมา นุ โข โส อุทายิ วทมาโน วเทยฺยาติ.\csromanspacer{} โน เหตํ ภนฺเต.\csromanspacer{} เยหิ โส ภนฺเต ปุริโส พนฺธเนหิ พทฺโธ น สกฺโกติ เอกํ อคารกํ โอลุคฺควิลุคฺคํ กากาติทายิํ นปรมรูปํ ปหาย, เอกํ ขโฏปิกํ โอลุคฺควิลุคฺคํ นปรมรูปํ ปหาย, เอกิสฺสา กุมฺภิยา ธญฺญสมวาปกํ นปรมรูปํ ปหาย, เอกํ ชายิกํ นปรมรูปํ ปหาย เกสมสฺสุํ โอหาเรตฺวา กาสายานิ วตฺถานิ อจฺฉาเทตฺวา อคารสฺมา อนคาริยํ ปพฺพชิตุํ.\csromanspacer{} ตญฺหิ ตสฺส พลวํ พนฺธนํ ทฬฺหํ พนฺธนํ ถิรํ พนฺธนํ อปูติกํ พนฺธนํ ถูโล กลิงฺคโรติ.\csromanspacer{} เอวเมว โข อุทายิ อิเธกจฺเจ โมฆปุริสา “อิทํ ปชหถา”ติ มยา วิจฺจมานา เต เอวมาหํสุ “กิํ ปนิมสฺส อปฺปมตฺตกสฺส โอรมตฺตกสฺส, อธิสลฺลิขเตวายํ สมโณ”ติ.\csromanspacer{} เต ตญฺเจว นปฺปชหนฺติ, มยิ จ อปฺปจฺจยํ อุปฏฺฐาเปนฺติ เย จ ภิกฺขู สิกฺขากามา.\csromanspacer{} เตสํ ตํ อุทายิ โหติ พลวํ พนฺธนํ ทฬฺหํ พนฺธนํ ถิรํ พนฺธนํ อปูติกํ พนฺธนํ ถูโล กลิงฺคโร.}
\proseitem{153}{\csromanfolio{115}เสยฺยถาปิ อุทายิ คหปติ วา คหปติปุตฺโต วา อฑฺโฒ มหทฺธโน มหาโภโค เนกานํ นิกฺขคณานํ จโย เนกานํ ธญฺญคณานํ จโย เนกานํ เขตฺตคณานํ จโย เนกานํ วตฺถุคณานํ จโย เนกานํ ภริยคณานํ จโย เนกานํ ทาสคณานํ จโย เนกานํ ทาสิคณานํ จโย.\csromanspacer{} โส อารามคตํ ภิกฺขุํ ปสฺเสยฺย, สุโธตหตฺถปาทํ มนุญฺญํ โภชนํ ภุตฺตาวิํ สีตาย ฉายาย นิสินฺนํ อธิจิตฺเต ยุตฺตํ, ตสฺส เอวมสฺส “สุขํ วต โภ สามญฺญํ, อาโรคฺยํ วต โภ สามญฺญํ, โส วตสฺสํ, โยหํ เกสมสฺสุํ โอหาเรตฺวา กาสายานิ วตฺถานิ อจฺฉาเทตฺวา อคารสฺมา อนคาริยํ ปพฺพเชยฺยนฺ”ติ.\csromanspacer{} โส สกฺกุเณยฺย เนกานิ นิกฺขคณานิ ปหาย, เนกานิ ธญฺญคณานิ ปหาย, เนกานิ เขตฺตคณานิ ปหาย, เนกานิ วตฺถุคณานิ ปหาย, เนกานิ ภริยคณานิ ปหาย, เนกานิ ทาสคณานิ ปหาย, เนกานิ ทาสิคณานิ ปหาย เกสมสฺสุํ โอหาเรตฺวา กาสายนิ วตฺถานิ อจฺฉาเทตฺวา อคารสฺมา อนคาริยํ ปพฺพชิตุํ.\csromanspacer{} โย นุ โข อุทายิ เอวํ วเทยฺย “เยหิ โส คหปติ วา คหปติปุตฺโต วา พนฺธเนหิ พทฺโธ สกฺโกติ เนกานิ นิกฺขคณานิ ปหาย, เนกานิ ธญฺญคณานิ ปหาย, \csromanfolio{116}เนกานิ เขตฺตคณานิ ปหาย, เนกานิ วตฺถุคณานิ ปหาย, เนกานิ ภริยคณานิ ปหาย, เนกานิ ทาสคณานิ ปหาย, เนกานิ ทาสิคณานิ ปหาย เกสมสฺสุํ โอหาเรตฺวา กสายานิ วตฺถานิ อจฺฉาเทตฺวา อคารสฺมา อนคาริยํ ปพฺพชิตุํ.\csromanspacer{} ตญฺหิ ตสฺส พลวํ พนฺธนํ ทฬฺหํ พนฺธนํ ถิรํ พนฺธนํ อปูติกํ พนฺธนํ ถูโล กลิงฺคโร”ติ, สมฺมา นุ โข โส อุทายิ วทมาโน วเทยฺยาติ.\csromanspacer{} โน เหตํ ภนฺเต.\csromanspacer{} เยหิ โส ภนฺเต คหปติ วา คหปติปุตฺโต วา พนฺธเนหิ พทฺโธ สกฺโกติ เนกานิ นิกฺขคณานิ ปหาย, เนกานิ ธญฺญคณานิ ปหาย, เนกานิ เขตฺตคณานิ ปหาย, เนกานิ วตฺถุคณานิ ปหาย, เนกานิ ภริยคณานิ ปหาย, เนกานิ ทาสคณานิ ปหาย, เนกานิ ทาสิคณานิ ปหาย เกสมสฺสุํ โอหาเรตฺวา กาสายานิ วตฺถานิ อจฺฉาเทตฺวา อคารสฺมา อนคาริยํ ปพฺพชิตุํ.\csromanspacer{} ตญฺหิ ตสฺส อพลํ พนฺธนํ ทุพฺพลํ พนฺธนํ ปูติกํ พนฺธนํ อสารกํ พนฺธนนฺติ.\csromanspacer{} เอวเมว โข อุทายิ อิเธกจฺเจ กุลปุตฺตา “อิทํ ปชหถา”ติ มยา วุจฺจมานา เต เอวมาหํสุ “กิํ ปนิมสฺส อปฺปมตฺตกสฺส โอรมตฺตกสฺส ปหาตพฺพสฺส, ยสฺส โน ภควา ปหานมาห, ยสฺส โน สุคโต ปฏินิสฺสคฺคมาหา”ติ.\csromanspacer{} เต ตญฺเจว ปชหนฺติ, มยิ จ น อปฺปจฺจยํ อุปฏฺฐาเปนฺติ เย จ ภิกฺขู สิกฺขากามา.\csromanspacer{} เต ตํ ปหาย อปฺโปสฺสุกฺกา ปนฺนโลมา ปรทตฺตวุตฺตา มิคภูเตน เจตสา วิหรนฺติ, เตสํ ตํ อุทายิ โหติ อพลํ พนฺธนํ ทุพฺพลํ พนฺธนํ ปูติกํ พนฺธนํ อสารกํ พนฺธนํ.}

\proseitem{154}{จตฺตาโรเม อุทายิ ปุคฺคลา สนฺโต สํวิชฺชมาโน โลกสฺมิํ.\csromanspacer{} กตเม จตฺตาโร.\csromanspacer{} อิธุทายิ เอกจฺโจ ปุคฺคโล อุปธิปหานาย ปฏิปนฺโน โหติ อุปธิปฏินิสฺสคฺคาย, ตเมนํ อุปธิปหานาย ปฏิปนฺนํ อุปธิปฏินิสฺสคฺคาย อุปธิปฏิสํยุตฺตา สรสงฺกปฺปา สมุทาจรนฺติ.\csromanspacer{} โส เต อธิวาเสติ นปฺปชหติ น วิโนเทติ น พฺยนฺตีกโรติ น อนภาวํ คเมติ.\csromanspacer{} อิมํ โข อหํ อุทายิ ปุคฺคลํ สํยุตฺโตติ วทามิ, โน วิสํยุตฺโต.\csromanspacer{} ตํ กิสฺส เหตุ, อินฺทฺริยเวมตฺตตา หิ เม อุทายิ อิมสฺมิํ ปุคฺคเล วิทิตา.}

\prose{อิธ ปนุทายิ เอกจฺโจ ปุคฺคโล อุปธิปหานาย ปฏิปนฺโน โหติ อุปธิปฏินิสฺสคฺคาย, ตเมนํ อุปธิปหานาย ปฏิปนฺนํ อุปธิปฏินิสฺสคฺคาย อุปธิปฏิสํยุตฺตา สรสงฺกปฺปา สมุทาจรนฺติ.\csromanspacer{} โส เต นาธิวาเสติ ปชหติ วิโนเทติ พฺยนฺตีกโรติ อนภาวํ คเมติ.\csromanspacer{} อิมมฺปิ โข อหํ}

\vspace{\tipitakaparskip}
\csromancenter{ลฏุกิโกปมสุตฺต (66) อุทายิ ปุคฺคลํ สํยุตฺโตติ วทามิ, โน วิสํยุตฺโต.\csromanspacer{} ตํ กิสฺส เหตุ, อินฺทฺริยเวมตฺตตา หิ เม อุทายิ อิมสฺมิํ ปุคฺคเล วิทิตา.}
\prose{\csromanfolio{117}อิธ ปนุทายิ เอกจฺโจ ปุคฺคโล อุปธิปหานาย ปฏิปนฺโน โหติ อุปธิปฏินิสฺสคฺคาย, ตเมนํ อุปธิปหานาย ปฏิปนฺนํ อุปธิปฏินิสฺสคฺคาย กทาจิ กรหจิ สติสมฺโมสา อุปธิปฏิสํยุตฺตา สรสงฺกปฺปา สมุทาจรนฺติ.\csromanspacer{} ทนฺโธ อุทายิ สตุปฺปาโท, อถ โข นํ ขิปฺปเมว ปชหติ วิโนเทติ พฺยนฺตีกโรติ อนภาวํ คเมติ.\csromanspacer{} เสยฺยถาปิ อุทายิ ปุริโส ทิวสํสนฺตตฺเต\footnote{ทิวสสนฺตตฺเต (สี, สฺยา, กํ, อิ)} อโยกฏาเห เทฺว วา ตีณิ วา อุทกผุสิตานิ นิปาเตยฺย.\csromanspacer{} ทนฺโธ อุทายิ อุทกผุสิตานํ นิปาโต, อถ โข นํ ขิปฺปเมว ปริกฺขยํ ปริยาทานํ คจฺเฉยฺย.\csromanspacer{} เอวเมว โข อุทายิ อิเธกจฺโจ ปุคฺคโล อุปธิปหานาย ปฏิปนฺโน โหติ อุปธิปฏินิสฺสคฺคาย, ตเมนํ อุปธิปหานาย ปฏิปนฺนํ อุปธิปฏินิสฺสคฺคาย กทาจิ กรหจิ สติสมฺโมสา อุปธิปฏิสํยุตฺตา สรสงฺกปฺปา สมุทาจรนฺติ.\csromanspacer{} ทนฺโธ อุทายิ สตุปฺปาโท, อถ โข นํ ขิปฺปเมว ปชหติ วิโนเทติ พฺยนฺตีกโรติ อนภาวํ คเมติ.\csromanspacer{} อิมมฺปิ โข อหํ อุทายิ ปุคฺคลํ สํยุตฺโตติ วทามิ, โน วิสํยุตฺโต.\csromanspacer{} ตํ กิสฺส เหตุ, อินฺทฺริยเวมตฺตตา หิ เม อุทายิ อิมสฺมิํ ปุคฺคเล วิทิตา.}

\prose{อิธ ปนุทายิ เอกจฺโจ ปุคฺคโล อุปธิ ทุกฺขสฺส มูลนฺติ อิติ วิทิตฺวา นิรุปธิ โหติ อุปธิสงฺขเย วิมุตฺโต.\csromanspacer{} อิมํ โข อหํ อุทายิ ปุคฺคลํ วิสํยุตฺโตติ วทามิ, โน สํยุตฺโตติ.\csromanspacer{} ตํ กิสฺส เหตุ, อินฺทฺริยเวมตฺตตา หิ เม อุทายิ อิมสฺมิํ ปุคฺคเล วิทิตา.\csromanspacer{} อิเม โข อุทายิ จตฺตาโร ปุคฺคลา สนฺโต สํวิชฺชมานา โลกสฺมิํ.}

\proseitem{155}{ปญฺจ โข อิเม อุทายิ กามคุณา.\csromanspacer{} กตเม ปญฺจ, จกฺขุวิญฺเญยฺยา รูปา อิฏฺฐา กนฺตา มนาปา ปิยรูปา กามูปสํหิตา รชนียา.\csromanspacer{} โสตวิญฺเญยฺยา สทฺทา ฯเปฯ.\csromanspacer{} ฆานวิญฺเญยฺยา คนฺธา.\csromanspacer{} ชิวฺหาวิญฺเญยฺยา รสา.\csromanspacer{} กายวิญฺเญยฺยา โผฏฺฐพฺพา อิฏฺฐา กนฺตา มนาปา ปิยรูปา กามูปสํหิตา รชนียา.\csromanspacer{} อิเม โข อุทายิ ปญฺจ กามคุณา.\csromanspacer{} ยํ โข อุทายิ อิเม ปญฺจ กามคุเณ ปฏิจฺจ อุปฺปชฺชติ สุขํ โสมนสฺสํ, อิทํ วุจฺจติ กามสุขํ มิฬฺหสุขํ\footnote{มีฬฺหสุขํ (สี, อิ)} ปุถุชฺชนสุขํ อนริยสุขํ น เสวิตพฺพํ น ภาเวตพฺพํ น พหุลีกาตพฺพํ, ภายิตพฺพํ เอตสฺส สุขสฺสาติ วทามิ.}

\proseitem{156}{\csromanfolio{118}อิธุทายิ ภิกฺขุ วิวิจฺเจว กาเมหิ ฯเปฯ ปฐมํ ฌานํ อุปสมฺปชฺช วิหรติ.\csromanspacer{} วิตกฺกวิจารานํ วูปสมา ฯเปฯ ทุติยํ ฌานํ อุปสมฺปชฺช วิหรติ.\csromanspacer{} ปีติยา จ วิราคา ฯเปฯ ตติยํ ฌานํ อุปสมฺปชฺช วิหรติ.\csromanspacer{} สุขสฺส จ ปหานา ฯเปฯ จตุตฺถํ ฌานํ อุปสมฺปชฺช วิหรติ.\csromanspacer{} อิทํ วุจฺจติ เนกฺขมฺมสุขํ ปวิเวกสุขํ อุปสมสุขํ สมฺโพธสุขํ อาเสวิตพฺพํ ภาเวตพฺพํ พหุลีกาตพฺพํ, น ภายิตพฺพํ เอตสฺส สุขสฺสาติ วทามิ.}

\prose{อิธุทายิ ภิกฺขุ วิวิจฺเจว กาเมหิ ฯเปฯ ปฐมํ ฌานํ อุปสมฺปชฺช วิหรติ.\csromanspacer{} อิทํ โข อหํ อุทายิ อิญฺชิตสฺมิํ วทามิ.\csromanspacer{} กิญฺจ ตตฺถ อิญฺชิตสฺมิํ.\csromanspacer{} ยเทว ตตฺถ วิตกฺกวิจารา อนิรุทฺธา โหนฺติ, อิทํ ตตฺถ อิญฺชิตสฺมิํ.\csromanspacer{} อิธุทายิ ภิกฺขุ วิตกฺกวิจารานํ วูปสมา ฯเปฯ ทุติยํ ฌานํ อุปสมฺปชฺช วิหรติ, อิทมฺปิ โข อหํ อุทายิ อิญฺชิตสฺมิํ วทามิ.\csromanspacer{} กิญฺจ ตตฺถ อิญฺชิตสฺมิํ.\csromanspacer{} ยเทว ตตฺถ ปีติสุขํ อนิรุทฺธํ โหติ, อิทํ ตตฺถ อิญฺชิตสฺมิํ.\csromanspacer{} อิธุทายิ ภิกฺขุ ปีติยา จ วิราคา ฯเปฯ ตติยํ ฌานํ อุปสมฺปชฺช วิหรติ, อิทมฺปิ โข อหํ อุทายิ อิญฺชิตสฺมิํ วทามิ.\csromanspacer{} กิญฺจ ตตฺถ อิญฺชิตสฺมิํ.\csromanspacer{} ยเทว ตตฺถ อุเปกฺขาสุขํ อนิรุทฺธํ โหติ, อิทํ ตตฺถ อิญฺชิตสฺมิํ.\csromanspacer{} อิธุทายิ ภิกฺขุ สุขสฺส จ ปหานา ฯเปฯ จตุตฺถํ ฌานํ อุปสมฺปชฺช วิหรติ.\csromanspacer{} อิทํ โข อหํ อุทายิ อนิญฺชิตสฺมิํ วทามิ.}

\prose{อิธุทายิ ภิกฺขุ วิวิจฺเจว กาเมหิ ฯเปฯ ปฐมํ ฌานํ อุปสมฺปชฺช วิหรติ.\csromanspacer{} อิทํ โข อหํ อุทายิ อนลนฺติ วทามิ, ปชหถาติ วทามิ, สมติกฺกมถาติ วาทามิ.\csromanspacer{} โก จ ตสฺส สมติกฺกโม.\csromanspacer{} อิธุทายิ ภิกฺขุ วิตกฺกวิจารานํ วูปสมา ฯเปฯ ทุติยํ ฌานํ อุปสมฺปชฺช วิหรติ, อยํ ตสฺส สมติกฺกโม.\csromanspacer{} อิทมฺปิ โข อหํ อุทายิ อนลนฺติ วทามิ, ปชหถาติ วทามิ, สมติกฺกมถาติ วทามิ.\csromanspacer{} โก จ ตสฺส สมติกฺกโม.\csromanspacer{} อิธุทายิ ภิกฺขุ ปีติยา จ วิราคา ฯเปฯ ตติยํ ฌานํ อุปสมฺปชฺช วิหรติ, อยํ ตสฺส สมติกฺกโม.\csromanspacer{} อิทมฺปิ โข อหํ อุทายิ อนลนฺติ วทามิ, ปชหถาติ วทามิ, สมติกฺกมถาติ วทามิ.\csromanspacer{} โก จ ตสฺส สมติกฺกโม.\csromanspacer{} อิธุทายิ ภิกฺขุ สุขสฺส จ ปหานา ฯเปฯ จตุตฺถํ ฌานํ อุปสมฺปชฺช วิหรติ, อยํ ตสฺส สมติกฺกโม.\csromanspacer{} อิทมฺปิ โข อหํ อุทายิ อนลนฺติ วทามิ, ปชหถาติ วทามิ, สมติกฺกมถาติ วทามิ.\csromanspacer{} โก จ ตสฺส สมติกฺกโม.\csromanspacer{} อิธุทายิ ภิกฺขุ สพฺพโส รูปสญฺญานํ สมติกฺกมา ปฏิฆสญฺญานํ อตฺถงฺคมา นานตฺตสญฺญานํ อมนสิการา “อนนฺโต อากาโส”ติ อากาสานญฺจายตนํ อุปสมฺปชฺช วิหรติ, อยํ ตสฺส สมติกฺกโม.\csromanspacer{} อิทมฺปิ โข อหํ อุทายิ}

\vspace{\tipitakaparskip}
\csromancenter{จาตุมสุตฺต (67) อนลนฺติ วทามิ, ปชหถาติ วทามิ, สมติกฺกมถาติ วทามิ.\csromanspacer{} โก จ ตสฺส สมติกฺกโม.\csromanspacer{} อิธุทายิ ภิกฺขุ สพฺพโส อากาสานญฺจายตนํ สมติกฺกมฺม “อนนฺตํ วิญฺญาณนฺ”ติ วิญฺญาณญฺจายตนํ อุปสมฺปชฺช วิหรติ, อยํ ตสฺส สมติกฺกโม.\csromanspacer{} อิทมฺปิ โข อหํ อุทายิ อนลนฺติ วทามิ, ปชหถาติ วทามิ, สมติกฺกมถาติ วทามิ.\csromanspacer{} โก จ ตสฺส สมติกฺกโม.\csromanspacer{} อิธุทายิ ภิกฺขุ สพฺพโส วิญฺญาณญฺจายตนํ สมติกฺกมฺม “นตฺถิ กิญฺจี”ติ อากิญฺจญฺญายตนํ อุปสมฺปชฺช วิหรติ, อยํ ตสฺส สมติกฺกโม.\csromanspacer{} อิทมฺปิ โข อหํ อุทายิ อนลนฺติ วทามิ, ปชหถาติ วาทามิ, สมติกฺกมถาติ วทามิ.\csromanspacer{} โก จ ตสฺส สมติกฺกโม.\csromanspacer{} อิธุทายิ ภิกฺขุ สพฺพโส อากิญฺจญฺญายตนํ สมติกฺกมฺม เนวสญฺญานาสญฺญายตนํ อุปสมฺปชฺช วิหรติ, อยํ ตสฺส สมติกฺกโม.\csromanspacer{} อิทมฺปิ โข อหํ อุทายิ อนลนฺติ วทามิ, ปชหถาติ วทามิ, สมติกฺกมถาติ วทามิ.\csromanspacer{} โก จ ตสฺส สมติกฺกโม.\csromanspacer{} อิธุทายิ ภิกฺขุ สพฺพโส เนวสญฺญานาสญฺญายตนํ สมติกฺกมฺม สญฺญาเวทยิตนิโรธํ อุปสมฺปชฺช วิหรติ, อยํ ตสฺส สมติกฺกโม.\csromanspacer{} อิติ โข อหํ อุทายิ เนวสญฺญานาสญฺญายตนสฺสปิ ปหานํ วทามิ.\csromanspacer{} ปสฺสสิ โน ตฺวํ อุทายิ ตํ สํโยชนํ อณุํ วา ถูลํ วา, ยสฺสาหํ โน ปหานํ วทามีติ.\csromanspacer{} โน เหตํ ภนฺเตติ.}
\csromancenter{อิทมโวจ ภควา.\csromanspacer{} อตฺตมโน อายสฺมา อุทายี ภควโต ภาสิตํ อภินนฺทีติ.}
\nitthitamruled{ลฏุกิโกปมสุตฺตํ นิฏฺฐิตํ ฉฏฺฐํ.}
\csromanmatikaanchor{mtk.20}
\csromantocmarkref{subsection}{7. จาตุมสุตฺต}{mtk.20}
\bookmark[dest={mtk.20},level=2]{7. จาตุมสุตฺต}
\csromanheader{2}{7. จาตุมสุตฺต}
\proseitem{157}{\csromanfolio{119}เอวํ เม สุตํ\csromandash{}เอกํ สมยํ ภควา จาตุมายํ วิหรติ อามลกีวเน.\csromanspacer{} เตน โข ปน สมเยน สาริปุตฺตโมคฺคลฺลานปฺปมุขานิ ปญฺจมตฺตานิ ภิกฺขุสตานิ จาตุมํ อนุปฺปตฺตานิ โหนฺติ ภควนฺตํ ทสฺสนาย.\csromanspacer{} เต จ อาคนฺตุกา ภิกฺขู เนวาสิเกหิ ภิกฺขูหิ สทฺธิํ ปฏิสมฺโมทมานา เสนาสนานิ ปญฺญาปยมานา ปตฺตจีวรานิ ปฏิสามยมานา อุจฺจาสทฺทา มหาสทฺทา อเหสุํ.\csromanspacer{} อถ โข ภควา อายสฺมนฺตํ อานนฺทํ อามนฺเตสิ “เก ปเนเต อานนฺท อุจฺจาสทฺทา มหาสทฺทา เกวฏฺฏา มญฺเญ มจฺฉวิโลเป”ติ.\csromanspacer{} เอตานิ ภนฺเต สาริปุตฺตโมคฺคลฺลานปฺปมุขานิ ปญฺจมตฺตานิ ภิกฺขุสตานิ \csromanfolio{120}จาตุมํ อนุปฺปตฺตานิ ภควนฺตํ ทสฺสนาย.\csromanspacer{} เต อาคนฺตุกา ภิกฺขู เนวาสิเกหิ ภิกฺขูหิ สทฺธิํ ปฏิสมฺโมทมานา เสนาสนานิ ปญฺญาปยมานา ปตฺตจีวรานิ ปฏิสามยมานา อุจฺจาสทฺทา มหาสทฺทาติ.\csromanspacer{} เตนหานนฺท มม วจเนน เต ภิกฺขู อามนฺเตหิ “สตฺถา อายสฺมานฺเต อามนฺเตตี”ติ. “เอวํ ภนฺเต”ติ โข อายสฺมา อานนฺโท ภควโต ปฏิสฺสุตฺวา เยน เต ภิกฺขู เตนุปสงฺกมิ, อุปสงฺกมิตฺวา เต ภิกฺขู เอตทโวจ “สตฺถา อายสฺมนฺเต อามนฺเตตี”ติ. “เอวมาวุโส”ติ โข เต ภิกฺขู อายสฺมโต อานนฺทสฺส ปฏิสฺสุตฺวา เยน ภควา เตนุปสงฺกมิํสุ, อุปสงฺกมิตฺวา ภควนฺตํ อภิวาเทตฺวา เอกมนฺตํ นิสีทิํสุ.\csromanspacer{} เอกมนฺตํ นิสินฺเน โข เต ภิกฺขู ภควา เอตทโวจ “กิํ นุ ตุเมฺห ภิกฺขเว อุจฺจาสทฺทา มหาสทฺทา เกวฏฺฏา มญฺเญ มจฺฉวิโลเป”ติ.\csromanspacer{} อิมานิ ภนฺเต สาริปุตฺตโมคฺคลฺลานปฺปมุขานิ ปญฺจมตฺตานิ ภิกฺขุสตานิ จาตุมํ อนุปฺปตฺตานิ ภควนฺตํ ทสฺสนาย.\csromanspacer{} เตเม อาคนฺตุกา ภิกฺขู เนวาสิเกหิ ภิกฺขูหิ สทฺธิํ ปฏิสมฺโมทมานา เสนาสนานิ ปญฺญาปยมานา ปตฺตจีวรานิ ปฏิสามยมานา อุจฺจาสทฺทา มหาสทฺทาติ.\csromanspacer{} คจฺฉถ ภิกฺขเว, ปณาเมมิ โว, น โว มม สนฺติเก วตฺถพฺพนฺติ. “เอวํ ภนฺเต”ติ โข เต ภิกฺขู ภควโต ปฏิสฺสุตฺวา อุฏฺฐายาสนา ภควนฺตํ อภิวาเทตฺวา ปทกฺขิณํ กตฺวา เสนาสนํ สํสาเมตฺวา ปตฺตจีวรมาทาย ปกฺกมิํสุ.}

\proseitem{158}{เตน โข ปน สมเยน จาตุเมยฺยกา สกฺยา สนฺถาคาเร\footnote{สนฺธาคาเร (ก)} สนฺนิปติตา โหนฺติ เกนจิเทว กรณีเยน.\csromanspacer{} อทฺทสํสุ โข จาตุเมยฺยกา สกฺยา เต ภิกฺขู ทูรโตว อาคจฺฉนฺเต, ทิสฺวาน เยน เต ภิกฺขู เตนุปสงฺกมิํสุ, อุปสงฺกมิตฺวา เต ภิกฺขู เอตทโวจุํ “หนฺท กหํ ปน ตุเมฺห อายสฺมนฺโต คจฺฉถา”ติ.\csromanspacer{} ภควตา โข อาวุโส ภิกฺขุสํโฆ ปณามิโตติ.\csromanspacer{} เตนหายสฺมานฺโต มุหุตฺตํ นิสีทถ, อปฺเปว นาม มยํ สกฺกุเณยฺยาม ภควนฺตํ ปสาเทตุนฺติ. “เอวมาวุโส”ติ โข เต ภิกฺขู จาตุเมยฺยกานํ สกฺยานํ ปจฺจสฺโสสุํ.\csromanspacer{} อถ โข จาตุเมยฺยกา สกฺยา เยน ภควา เตนุปสงฺกมิํสุ, อุปสงฺกมิตฺวา ภควนฺตํ อภิวาเทตฺวา เอกมนฺตํ นิสีทิํสุ, เอกมนฺตํ นิสินฺนา โข จาตุเมยฺยกา สกฺยา ภควนฺตํ เอตทโวจุํ “อภินนฺทตุ ภนฺเต ภควา ภิกฺขุสํฆํ, อภิวทตุ พนฺเต ภควา ภิกฺขุสํฆํ.\csromanspacer{} เสยฺยถาปิ}

\vspace{\tipitakaparskip}
\csromancenter{จาตุมสุตฺต (67) ภนฺเต ภควตา ปุพฺเพ ภิกฺขุสํโฆ อนุคฺคหิโต, เอวเมว ภควา เอตรหิ อนุคฺคณฺหาตุ ภิกฺขุสํฆํ, สนฺเตตฺถ ภนฺเต ภิกฺขู นวา อจิรปพฺพชิตา อธุนาคตา อิมํ ธมฺมวินยํ, เตสํ ภควนฺตํ ทสฺสนาย อลภนฺตานํ สิยา อญฺญถตฺตํ สิยา วิปริณาโม.\csromanspacer{} เสยฺยถาปิ ภนฺเต พีชานํ ตรุณานํ อุทกํ อลภนฺตานํ สิยา อญฺญถตฺตํ สิยา วิปริณาโม, เอวเมว โข ภนฺเต สนฺเตตฺถ ภิกฺขู นวา อจิรปพฺพชิตา อธุนาคตา อิมํ ธมฺมวินยํ, เตสํ ภควนฺตํ ทสฺสนาย อลภนฺตานํ สิยา อญฺญถตฺตํ สิยา วิปริณาโม.\csromanspacer{} เสยฺยถาปิ ภนฺเต วจฺฉสฺส ตรุณสฺส มาตรํ อปสฺสนฺตสฺส สิยา อญฺญถตฺตํ สิยา วิปริณาโม, เอวเมว โข ภนฺเต สนฺเตตฺถ ภิกฺขู นวา อจิรปพฺพชิตา อธุนาคตา อิมํ ธมฺมวินยํ, เตสํ ภควนฺตํ อปสฺสนฺตานํ สิยา อญฺญถตฺตํ สิยา วิปริณาโม.\csromanspacer{} อภินนฺทตุ ภนฺเต ภควา ภิกฺขุสํฆํ, อภิวทตุ ภนฺเต ภควา ภิกฺขุสํฆํ.\csromanspacer{} เสยฺยถาปิ ภนฺเต ภควตา ปุพฺเพ ภิกฺขุสํโฆ อนุคฺคหิโต, เอวเมว ภควา เอตรหิ อนุคฺคณฺหาตุ ภิกฺขุสํฆนฺ”ติ.}
\proseitem{159}{\csromanfolio{121}อถ โข พฺรหฺมา สหมฺปติ ภควโต เจตสา เจโตปริวิตกฺกมญฺญาย เสยฺยถาปิ นาม พลวา ปุริโส สมิญฺชิตํ\footnote{สมฺมิญฺชิตํ (สี, สฺยา, กํ, อิ)} วา พาหํ ปสาเรยฺย, ปสาริตํ วา พาหํ สมิญฺเชยฺย, เอวเมว พฺรหฺมโลเก อนฺตรหิโต ภควโต ปุรโต ปาตุรโหสิ.\csromanspacer{} อถ โข พฺรหฺมา สหมฺปติ เอกํสํ อุตฺตราสงฺคํ กริตฺวา เยน ภควา เตนญฺชลิํ ปณาเมตฺวา ภควนฺตํ เอตทโวจ “อภินนฺทตุ ภนฺเต ภควา ภิกฺขุสํฆํ, อภิวทตุ ภนฺเต ภควา ภิกฺขุสํฆํ.\csromanspacer{} เสยฺยถาปิ ภนฺเต ภควตา ปุพฺเพ ภิกฺขุสํโฆ อนุคฺคหิโต, เอวเมว ภควา เอตรหิ อนุคฺคณฺหาตุ ภิกฺขุสํฆํ, สนฺเตตฺถ ภนฺเต ภิกฺขู นวา อจิรปพฺพชิตา อธุนาคตา อิมํ ธมฺมวินยํ, เตสํ ภควนฺตํ ทสฺสนาย อลภนฺตานํ สิยา อญฺญถตฺตํ สิยา วิปริณาโม.\csromanspacer{} เสยฺยถาปิ ภนฺเต พีชานํ ตรุณานํ อุทกํ อลภนฺตานํ สิยา อญฺญถตฺตํ สิยา วิปริณาโม, เอวเมว โข ภนฺเต สนฺเตตฺถ ภิกฺขู นวา อจิรปพฺพชิตา อธุนาคตา อิมํ ธมฺมวินยํ, เตสํ ภควนฺตํ ทสฺสนาย อลภนฺตานํ สิยา อญฺญถตฺตํ สิยา วิปริณาโม.\csromanspacer{} เสยฺยถาปิ ภนฺเต วจฺฉสฺส ตรุณสฺส มาตรํ อปสฺสนฺตสฺส สิยา อญฺญถตฺตํ สิยา วิปริณาโม, เอวเมว โข ภนฺเต สนฺเตตฺถ ภิกฺขู \csromanfolio{122}นวา อจิรปพฺพชิตา อธุนาคตา อิมํ ธมฺมวินยํ, เตสํ ภควนฺตํ อปสฺสนฺตานํ สิยา อญฺญถตฺตํ สิยา วิปริณาโม.\csromanspacer{} อภินนฺทตุ ภนฺเต ภควา ภิกฺขุสํฆํ, อภิวทตุ ภนฺเต ภควา ภิกฺขุสํฆํ.\csromanspacer{} เสยฺยถาปิ ภนฺเต ภควตา ปุพฺเพ ภิกฺขุสํโฆ อนุคฺคหิโต, เอวเมว ภควา เอตรหิ อนุคฺคณฺหาตุ ภิกฺขุสํฆนฺ”ติ.}

\proseitem{160}{อสกฺขิํสุ โข จาตุเมยฺยกา จ สกฺยา พฺรหฺมา จ สหมฺปติ ภควนฺตํ ปสาเทตุํ พีชูปเมน จ ตรุณูปเมน จ.\csromanspacer{} อถ โข อายสฺมา มหาโมคฺคลฺลาโน ภิกฺขู อามนฺเตสิ “อุฏฺเฐถาวุโส, คณฺหถ ปตฺตจีวรํ ปสาทิโต ภควา จาตุเมยฺยเกหิ จ สเกฺยหิ พฺรหฺมุนา จ สหมฺปตินา พีชูปเมน จ ตรุณูปเมน จา”ติ. “เอวมาวุโส”ติ โข เต ภิกฺขู อายสฺมโต มหาโมคฺคลฺลานสฺส ปฏิสฺสุตฺวา อุฏฺฐายาสนา ปตฺตจีวรมาทาย เยน ภควา เตนุปสงฺกมิํสุ, อุปสงฺกมิตฺวา ภควนฺตํ อภิวาเทตฺวา เอกมนฺตํ นิสีทิํสุ.\csromanspacer{} เอกมนฺตํ นิสินฺนํ โข อายสฺมนฺตํ สาริปุตฺตํ ภควา เอตทโวจ “กินฺติ เต สาริปุตฺต อโหสิ มยา ภิกฺขุสํเฆ ปณามิเต”ติ.\csromanspacer{} เอวํ โข เม ภนฺเต อโหสิ “ภควโต ภิกฺขุสํโฆ ปณามิโต, อปฺโปสฺสุกฺโก ทานิ ภควา ทิฏฺฐธมฺมสุขวิหารํ อนุยุตฺโต วิหริสฺสติ, มยมฺปิ ทานิ อปฺโปสฺสุกฺกา ทิฏฺฐธมฺมสุขวิหารมนุยุตฺโต วิหริสฺสามา”ติ.\csromanspacer{} อาคเมหิ ตฺวํ สาริปุตฺต, อาคเมหิ ตฺวํ สาริปุตฺต ทิฏฺฐธมฺมสุขวิหารนฺติ.\csromanspacer{} อถ โข ภควา อายสฺมนฺตํ มหาโมคฺคลฺลานํ อามนฺเตสิ “กินฺติ เต โมคฺคลฺลาน อโหสิ มยา ภิกฺขุสํเฆ ปณามิเต”ติ.\csromanspacer{} เอวํ โข เม ภนฺเต อโหสิ “ภควตา ภิกฺขุสํโฆ ปณามิโต, อปฺโปสฺสุกฺโก ทานิ ภควา ทิฏฺฐธมฺมสุขวิหารํ อนุยุตฺโต วิหริสฺสติ, อหญฺจ ทานิ อายสฺมา จ สาริปุตฺโต ภิกฺขุสํฆํ ปริหริสฺสามา”ติ.\csromanspacer{} สาธุ สาธุ โมคฺคลฺลาน, อหํ วา หิ โมคฺคลฺลาน ภิกฺขุสํฆํ ปริหเรยฺยํ สาริปุตฺตโมคฺคลฺลานา วาติ.}

\proseitem{161}{อถ โข ภควา ภิกฺขู อามนฺเตสิ\csromandash{}จตฺตาริมานิ ภิกฺขเว ภยานิ อุทโกโรหนฺเต ปาฏิกงฺขิตพฺพานิ.\csromanspacer{} กตมานิ จตฺตาริ, อูมิภยํ\footnote{อุมฺมิภยํ (สฺยา, กํ)} กุมฺภีลภยํ อาวฏฺฏภยํ สุสุกาภยํ.\csromanspacer{} อิมานิ ภิกฺขเว จตฺตาริ ภยานิ อุทโกโรหนฺเต ปาฏิกงฺขิตพฺพานิ.\csromanspacer{} เอวเมว โข ภิกฺขเว จตฺตาริมานิ}

\vspace{\tipitakaparskip}
\csromancenter{จาตุมสุตฺต (67) ภยานิ อิเธกจฺเจ ปุคฺคเล อิมสฺมิํ ธมฺมวินเย อคารสฺมา อนคาริยํ ปพฺพชิเต ปาฏิกงฺขิตพฺพานิ.\csromanspacer{} กตมานิ จตฺตาริ, อูมิภยํ กุมฺภีลภยํ อาวฏฺฏภยํ สุสุกาภยํ.}
\proseitem{162}{\csromanfolio{123}กตมญฺจ ภิกฺขเว อูมิภยํ.\csromanspacer{} อิธ ภิกฺขเว เอกจฺโจ กุลปุตฺโต สทฺธา อคารสฺมา อนคาริยํ ปพฺพชิโต โหติ “โอติณฺโณมฺหิ ชาติยา ชราย มรเณน โสเกหิ ปริเทเวหิ ทุกฺเขหิ โทมนสฺเสหิ อุปายาเสหิ, ทุกฺโขติณฺโณ ทุกฺขปเรโต, อปฺเปว นาม อิมสฺส เกวลสฺส ทุกฺขกฺขนฺธสฺส อนฺตกิริยา ปญฺญาเยถา”ติ.\csromanspacer{} ตเมนํ ตถา ปพฺพชิตํ สมานํ สพฺรหฺมจารี โอวทนฺติ อนุสาสนฺติ “เอวํ เต อภิกฺกมิตพฺพํ, เอวํ เต ปฏิกฺกมิตพฺพํ, เอวํ เต อาโลกิตพฺพํ, เอวํ เต วิโลกิตพฺพํ, เอวํ เต สมิญฺชิตพฺพํ, เอวํ เต ปสาริตพฺพํ, เอวํ เต สํฆาฏิปตฺตจีวรํ ธาเรตพฺพนฺ”ติ.\csromanspacer{} ตสฺส เอวํ โหติ “มยํ โข ปุพฺเพ อคาริยภูตา สมานา อญฺเญ โอวทาม อนุสาสาม\footnote{โอวทามปิ อนุสาสามปิ (สี, สฺยา, กํ, อิ)}, อิเม ปนมฺหากํ ปุตฺตมตฺตา มญฺเญ นตฺตมตฺตา มญฺเญ อเมฺห\footnote{เอวํ (ก)} โอวทิตพฺพํ อนุสาสิตพฺพํ มญฺญนฺตี”ติ.\csromanspacer{} โส สิกฺขํ ปจฺจกฺขาย หีนายาวตฺตติ.\csromanspacer{} อยํ วุจฺจติ ภิกฺขเว อูมิภยสฺส ภีโต, สิกฺขํ ปจฺจกฺขาย หีนายาวตฺโต.\csromanspacer{} อูมิภยนฺติ โข ภิกฺขเว โกธุปายาสสฺเสตํ อธิวจนํ.}

\proseitem{163}{กตมญฺจ ภิกฺขเว กุมฺภีลภยํ.\csromanspacer{} อิธ ภิกฺขเว เอกจฺโจ กุลปุตฺโต สทฺธา อคารสฺมา อนคาริยํ ปพฺพชิโต โหติ “โอติณฺโณมฺหิ ชาติยา ชราย มรเณน โสเกหิ ปริเทเวหิ ทุกฺเขหิ โทมนสฺเสหิ อุปายาเสหิ, ทุกฺโขติณฺโณ ทุกฺขปเรโต, อปฺเปว นาม อิมสฺส เกวลสฺส ทุกฺขกฺขนฺธสฺส อนฺตกิริยา ปญฺญาเยถา”ติ.\csromanspacer{} ตเมนํ ตถา ปพฺพชิตํ สมานํ สพฺรหฺมจารี โอวทนฺติ อนุสาสนฺติ “อิทํ เต ขาทิตพฺพํ, อิทํ เต น ขาทิตพฺพํ, อิทํ เต ภุญฺชิตพฺพํ, อิทํ เต น ภุญฺชิตพฺพํ, อิทํ เต สายิตพฺพํ, อิทํ เต น สายิตพฺพํ, อิทํ เต ปาตพฺพํ, อิทํ เต น ปาตพฺพํ, กปฺปิยํ เต ขาทิตพฺพํ, อกปฺปิยํ เต น ขาทิตพฺพํ, กปฺปิยํ เต ภุญฺชิตพฺพํ, อกปฺปิยํ เต น ภุญฺชิตพฺพํ, กปฺปิยํ เต สายิตพฺพํ, อกปฺปิยํ เต น สายิตพฺพํ, กปฺปิยํ เต ปาตพฺพํ, อกปฺปิยํ เต น ปาตพฺพํ, กาเล เต ขาทิตพฺพํ, วิกาเล เต น ขาทิตพฺพํ, กาเล เต ภุญฺชิตพฺพํ, วิกาเล \csromanfolio{124}เต น ภุญฺชิตพฺพํ, กาเล เต สายิตพฺพํ, วิกาเล เต น สายิตพฺพํ, กาเล เต ปาตพฺพํ, วิกาเล เต น ปาตพฺพนฺ”ติ.\csromanspacer{} ตสฺส เอวํ โหติ “มยํ โข ปุพฺเพ อคาริยภูตา สมานา ยํ อิจฺฉาม, ตํ ขาทาม.\csromanspacer{} ยํ น อิจฺฉาม, น ตํ ขาทาม.\csromanspacer{} ยํ อิจฺฉาม, ตํ ภุญฺชาม.\csromanspacer{} ยํ น อิจฺฉาม, น ตํ ภุญฺชาม.\csromanspacer{} ยํ อิจฺฉาม, ตํ สายาม.\csromanspacer{} ยํ น อิจฺฉาม, น ตํ สายาม.\csromanspacer{} ยํ อิจฺฉาม, ตํ ปิวาม\footnote{ปิปาม (สี, อิ)}.\csromanspacer{} ยํ น อิจฺฉาม, น ตํ ปิวาม.\csromanspacer{} กปฺปิยํปิ ขาทาม, อกปฺปิยํปิ ขาทาม.\csromanspacer{} กปฺปิยํปิ ภุญฺชาม, อกปฺปิยํปิ ภุญฺชาม.\csromanspacer{} กปฺปิยํปิ สายาม, อกปฺปิยํปิ สายาม.\csromanspacer{} กปฺปิยํปิ ปิวาม, อกปฺปิยํปิ ปิวาม.\csromanspacer{} กาเลปิ ขาทาม, วิกาเลปิ ขาทาม.\csromanspacer{} กาเลปิ ภุญฺชาม, วิกาเลปิ ภุญฺชาม.\csromanspacer{} กาเลปิ สยาม, วิกาเลปิ สายาม.\csromanspacer{} กาเลปิ ปิวาม, วิกาเลปิ ปิวาม.\csromanspacer{} ยํปิ โน สทฺธา คหปติกา ทิวา วิกาเล ปณีตํ ขาทนียํ โภชนียํ เทนฺติ, ตตฺถปิเม มุขาวรณํ มญฺเญ กโรนฺตี”ติ.\csromanspacer{} โส สิกฺขํ ปจฺจกฺขาย หีนายาวตฺตติ.\csromanspacer{} อยํ วุจฺจติ ภิกฺขเว กุมฺภีลภยสฺส ภีโต สิกฺขํ ปจฺจกฺขาย หีนายาวตฺโต.\csromanspacer{} กุมฺภีลภยนฺติ โข ภิกฺขเว โอทริกตฺตสฺเสตํ อธิวจนํ.}

\proseitem{164}{กตมญฺจ ภิกฺขเว อาวฏฺฏภยํ.\csromanspacer{} อิธ ภิกฺขเว เอกจฺโจ กุลปุตฺโต สทฺธา อคารสฺมา อนคาริยํ ปพฺพชิโต โหติ “โอติณฺโณมฺหิ ชาติยา ชราย มรเณน โสเกหิ ปริเทเวหิ ทุกฺเขหิ โทมนสฺเสหิ อุปายาเสหิ, ทุกฺโขติณฺโณ ทุกฺขปเรโต, อปฺเปว นาม อิมสฺส เกวลสฺส ทุกฺขกฺขนฺธสฺส อนฺตกิริยา ปญฺญาเยถา”ติ.\csromanspacer{} โส เอวํ ปพฺพชิโต สมาโน ปุพฺพณฺหสมยํ นิวาเสตฺวา ปตฺตจีวรมาทาย คามํ วา นิคมํ วา ปิณฺฑาย ปวิสติ อรกฺขิเตเนว กาเยน อรกฺขิตาย วาจาย อนุปฏฺฐิตาย สติยา อสํวุเตหิ อินฺทฺริเยหิ.\csromanspacer{} โส ตตฺถ ปสฺสติ คหปติํ วา คหปติปุตฺตํ วา ปญฺจหิ กามคุเณหิ สมปฺปิตํ สมงฺคีภูตํ ปริจารยมานํ\footnote{ปริจาริยมานํ (สฺยา, กํ, ก)}.\csromanspacer{} ตสฺส เอวํ โหติ “มยํ โข ปุพฺเพ อคาริยภูตา สมานา ปญฺจหิ กามคุเณหิ สมปฺปิตา สมงฺคีภูตา ปริจาริมฺหา, สํวิชฺชนฺติ โข ปน เม กุเล\footnote{สํวิชฺชนฺติ โข กุเล (สี, สฺยา, กํ, อิ)} โภคา, สกฺกา โภเค จ ภุญฺชิตุํ ปุญฺญานิ จ กาตุนฺ”ติ.\csromanspacer{} โส สิกฺขํ ปจฺจกฺขาย หีนายาวตฺตติ.\csromanspacer{} อยํ วุจฺจติ ภิกฺขเว อาวฏฺฏภยสฺส ภีโต สิกฺขํ ปจฺจกฺขาย หีนายาวตฺโต.\csromanspacer{} อาวฏฺฏภยนฺติ โข ภิกฺขเว ปญฺจนฺเนตํ กามคุณานํ อธิวจนํ.}

\csromanheader{2}{8. นฬกปานสุตฺต (68)}
\proseitem{165}{\csromanfolio{125}กตมญฺจ ภิกฺขเว สุสุกาภยํ.\csromanspacer{} อิธ ภิกฺขเว เอกจฺโจ กุลปุตฺโต สทฺธา อคารสฺมา อนคาริยํ ปพฺพชิโต โหติ “โอติณฺโณมฺหิ ชาติยา ชราย มรเณน โสเกหิ ปริเทเวหิ ทุกฺเขหิ โทมนสฺเสหิ อุปายาเสหิ, ทุกฺโขติณฺโณ ทุกฺขปเรโต, อปฺเปว นาม อิมสฺส เกวลสฺส ทุกฺขกฺขนฺธสฺส อนฺตกิริยา ปญฺญาเยถา”ติ.\csromanspacer{} โส เอวํ ปพฺพชิโต สมาโน ปุพฺพณฺหสมยํ นิวาเสตฺวา ปตฺตจีวรมาทาย คามํ วา นิคมํ วา ปิณฺฑาย ปวิสติ อรกฺขิเตเนว กาเยน อรกฺขิตาย วาจาย อนุปฏฺฐิตาย สติยา อสํวุเตหิ อินฺทฺริเยหิ.\csromanspacer{} โส ตตฺถ ปสฺสติ มาตุคามํ ทุนฺนิวตฺถํ วา ทุปฺปารุตํ วา, ตสฺส มาตุคามํ ทิสฺวา ทุนฺนิวตฺถํ วา ทุปฺปารุตํ วา ราโค จิตฺตํ อนุทฺธํเสติ, โส ราคานุทฺธํเสน\footnote{อนุทฺธสฺเตน (สี, อิ)} จิตฺเตน สิกฺขํ ปจฺจกฺขาย หีนายาวตฺตติ.\csromanspacer{} อยํ วุจฺจติ ภิกฺขเว สุสุกาภยสฺส ภีโต สิกฺขํ ปจฺจกฺขาย หีนายาวตฺโต.\csromanspacer{} สุสุกาภยนฺติ โข ภิกฺขเว มาตุคามสฺเสตํ อธิวจนํ.\csromanspacer{} อิมานิ โข ภิกฺขเว จตฺตาริ ภยานิ อิเธกจฺเจ ปุคฺคเล อิมสฺมิํ ธมฺมวินเย อคารสฺมา อนคาริยํ ปพฺพชิเต ปาฏิกงฺขิตพฺพานีติ.}

\prose{อิทมโวจ ภควา.\csromanspacer{} อตฺตมนา เต ภิกฺขู ภควโต ภาสิตํ อภินนฺทุนฺติ.}

\nitthitamruled{จาตุมสุตฺตํ นิฏฺฐิตํ สตฺตมํ.}
\csromanmatikaanchor{mtk.21}
\csromantocmarkref{subsection}{8. นฬกปานสุตฺต}{mtk.21}
\bookmark[dest={mtk.21},level=2]{8. นฬกปานสุตฺต}
\csromanheader{2}{8. นฬกปานสุตฺต}
\proseitem{166}{เอวํ เม สุตํ\csromandash{}เอกํ สมยํ ภควา โกสเลสุ วิหรติ นฬกปาเน ปลาสเวน.\csromanspacer{} เตน โข ปน สมเยน สมฺพหุลา อภิญฺญาตา อภิญฺญาตา กุลปุตฺตา ภควนฺตํ อุทฺทิสฺส สทฺธา อคารสฺมา อนคาริยํ ปพฺพชิตา โหนฺติ.\csromanspacer{} อายสฺมา จ อนุรุทฺโธ อายสฺมา จ ภทฺทิโย\footnote{นนฺทิโย (สี, อิ) วินเย จ ม 1. จูฬโคสิงฺเค จ.} อายสฺมา จ กิมิโล\footnote{กิมฺพิโล (สี, สฺยา, กํ, อิ)} อายสฺมา จ ภคุ อายสฺมา จ โกณฺฑญฺโญ\footnote{กุณฺฑธาโน (สี, อิ)} อายสฺมา จ เรวโต อายสฺมา จ อานนฺโท อญฺเญ จ อภิญฺญาตา อภิญฺญาตา กุลปุตฺตา.\csromanspacer{} เตน โข ปน สมเยน ภควา ภิกฺขุสํฆปุริวุโต อพฺโภกาเส นิสินฺโน โหติ.\csromanspacer{} อถ โข ภควา เต \csromanfolio{126}กุลปุตฺเต อารพฺภ ภิกฺขู อามนฺเตสิ “เย เต ภิกฺขเว กุลปุตฺตา มมํ อุทฺทิสฺส สทฺธา อคารสฺมา อนคาริยํ ปพฺพชิตา, กจฺจิ เต ภิกฺขเว ภิกฺขู อภิรตา พฺรหฺมจริเย”ติ.\csromanspacer{} เอวํ วุตฺเต เต ภิกฺขู ตุณฺหี อเหสุํ.\csromanspacer{} ทุติยมฺปิ โข ภควา เต กุลปุตฺเต อารพฺภ ภิกฺขู อามนฺเตสิ “เย เต ภิกฺขเว กุลปุตฺตา มมํ อุทฺทิสฺส สทฺธา อคารสฺมา อนคาริยํ ปพฺพชิตา, กจฺจิ เต ภิกฺขเว ภิกฺขู อภิรตา พฺรหฺมจริเย”ติ.\csromanspacer{} ทุติยมฺปิ โข เต ภิกฺขู ตุณฺหี อเหสุํ.\csromanspacer{} ตติยมฺปิ โข ภควา เต กุลปุตฺเต อารพฺภ ภิกฺขู อามนฺเตสิ “เย เต ภิกฺขเว กุลปุตฺตา มมํ อุทฺทิสฺส สทฺธา อคารสฺมา อนคาริยํ ปพฺพชิตา, กจฺจิ เต ภิกฺขเว ภิกฺขู อภิรตา พฺรหฺมจริเย”ติ.\csromanspacer{} ตติยมฺปิ โข เต ภิกฺขู ตุณฺหี อเหสุํ.}

\proseitem{167}{อถ โข ภควโต เอตทโหสิ “ยํนูนาหํ เต กุลปุตฺเต ปุจฺเฉยฺยนฺ”ติ.\csromanspacer{} อถ โข ภควา อายสฺมนฺตํ อนุรุทฺธํ อามนฺเตสิ “กจฺจิ ตุเมฺห อนุรุทฺธา อภิรตา พฺรหฺมจริเย”ติ.\csromanspacer{} ตคฺฆ มยํ ภนฺเต อภิรตา พฺรหฺมจริเยติ.\csromanspacer{} สาธุ สาธุ อนุรุทฺธา, เอตํ โข อนุรุทฺธา ตุมฺหากํ ปติรูปํ กุลปุตฺตานํ สทฺธา อคารสฺมา อนคาริยํ ปพฺพชิตานํ.\csromanspacer{} ยํ ตุเมฺห อภิรเมยฺยาถ พฺรหฺมจริเย.\csromanspacer{} เยน ตุเมฺห อนุรุทฺธา ภเทฺรน โยพฺพเนน สมนฺนาคตา ปฐเมน วยสา สุสุกาฬเกสา กาเม ปริภุญฺเชยฺยาถ.\csromanspacer{} เตน ตุเมฺห อนุรุทฺธา ภเทฺรนปิ โยพฺพเนน สมนฺนาคตา ปฐเมน วยสา สุสุกาฬเกสา อคารสฺมา อนคาริยํ ปพฺพชิตา.\csromanspacer{} เต จ โข ปน ตุเมฺห อนุรุทฺธา เนว ราชาภินีตา อคารสฺมา อนคาริยํ ปพฺพชิตา, น โจราภินีตา อคารสฺมา อนคาริยํ ปพฺพชิตา, น อิณฏฺฏา อคารสฺมา อนคาริยํ ปพฺพชิตา, น ภยฏฺฏา อคารสฺมา อนคาริยํ ปพฺพชิตา, นาชีวิกาปกตา อคารสฺมา อนคาริยํ ปพฺพชิตา, อปิ จ โขมฺหิ โอติณฺโณ ชาติยา ชราย มรเณน โสเกหิ ปริเทเวหิ ทุกฺเขหิ โทมนสฺเสหิ อุปายาเสหิ, ทุกฺโขติณฺโณ ทุกฺขปเรโต, อปฺเปว นาม อิมสฺส เกวลสฺส ทุกฺขกฺขนฺธสฺส อนฺตกิริยา ปญฺญาเยถาติ.\csromanspacer{} นนุ ตุเมฺห อนุรุทฺธา เอวํ สทฺธา อคารสฺมา อนคาริยํ ปพฺพชิตาติ.\csromanspacer{} เอวํ ภนฺเต.\csromanspacer{} เอวํ ปพฺพชิเตน จ ปน อนุรุทฺธา กุลปุตฺเตน กิมสฺส กรณียํ.\csromanspacer{} วิเวกํ อนุรุทฺธา กาเมหิ วิเวกํ อกุสเลหิ ธมฺเมหิ ปีติสุขํ นาธิคจฺฉติ อญฺญํ วา\footnote{อญฺญํ จ (ก)} ตโต สนฺตตรํ.\csromanspacer{} ตสฺส อภิชฺฌาปิ จิตฺตํ ปริยาทาย ติฏฺฐติ, พฺยาปาโทปิ}

\vspace{\tipitakaparskip}
\csromancenter{นฬกปานสุตฺต (68) จิตฺตํ ปริยาทาย ติฏฺฐติ, ถินมิทฺธํปิ\footnote{ถีนมิทฺธมฺปิ (สี, สฺยา, กํ, อิ)} จิตฺตํ ปริยาทาย ติฏฺฐติ, อุทฺธจฺจกุกฺกุจฺจํปิ จิตฺตํ ปริยาทาย ติฏฺฐติ, วิจิกิจฺฉาปิ จิตฺตํ ปริยาทาย ติฏฺฐติ, อรตีปิ จิตฺตํ ปริยาทาย ติฏฺฐติ, ตนฺทีปิ จิตฺตํ ปริยาทาย ติฏฺฐติ.\csromanspacer{} วิเวกํ อนุรุทฺธา กาเมหิ วิเวกํ อกุสเลหิ ธมฺเมหิ ปีติสุขํ นาธิคจฺฉติ อญฺญํ วา ตโต สนฺตตรํ.}
\prosecont{\csromanfolio{127}วิเวกํ อนุรุทฺธา กาเมหิ วิเวกํ อกุสเลหิ ธมฺเมหิ ปีติสุขํ อธิคจฺฉติ อญฺญํ วา ตโต สนฺตตรํ.\csromanspacer{} ตสฺส อภิชฺฌาปิ จิตฺตํ น ปริยาทาย ติฏฺฐติ, พฺยาปาโทปิ จิตฺตํ น ปริยาทาย ติฏฺฐติ, ถินมิทฺธํปิ จิตฺตํ น ปริยาทาย ติฏฺฐติ, อุทฺธจฺจกุกฺกุจฺจํปิ จิตฺตํ น ปริยาทาย ติฏฺฐติ, วิจิกิจฺฉาปิ จิตฺตํ น ปริยาทาย ติฏฺฐติ, อรตีปิ จิตฺตํ น ปริยาทาย ติฏฺฐติ, ตนฺทีปิ จิตฺตํ น ปริยาทาย ติฏฺฐติ.\csromanspacer{} วิเวกํ อนุรุทฺธา กาเมหิ วิเวกํ อกุสเลหิ ธมฺเมหิ ปีติสุขํ อธิคจฺฉติ, อญฺญํ วา ตโต สนฺตตรํ.}

\proseitem{168}{กินฺติ โว อนุรุทฺธา มยิ โหติ, “เย อาสวา สํกิเลสิกา โปโนพฺภวิกา\footnote{โปโนภวิกา (สี, อิ)} สทรา ทุกฺขวิปากา อายติํ ชาติชรามรณิยา, อปฺปหีนา เต ตถาคตสฺส.\csromanspacer{} ตสฺมา ตถาคโต สงฺขาเยกํ ปฏิเสวติ, สงฺขาเยกํ อธิวาเสติ, สงฺขาเยกํ ปริวชฺเชติ, สงฺขาเยกํ วิโนเทตี”ติ.\csromanspacer{} น โข โน ภนฺเต ภควติ เอวํ โหติ “เย อาสวา สํกิเลสิกา โปโนพฺภวิกา สทรา ทุกฺขวิปากา อายติํ ชาติชรามรณิยา, อปฺปหีนา เต ตถาคตสฺส.\csromanspacer{} ตสฺมา ตถาคโต สงฺขาเยกํ ปฏิเสวติ, สงฺขาเยกํ อธิวาเสติ, สงฺขาเยกํ ปริวชฺเชติ, สงฺขาเยกํ วิโนเทตี”ติ.\csromanspacer{} เอวํ โข โน ภนฺเต ภควติ โหติ “เย อาสวา สํกิเลสิกา โปโนพฺภวิกา สทรา ทุกฺขวิปากา อายติํ ชาติชรามรณิยา, ปหีนา เต ตถาคตสฺส.\csromanspacer{} ตสฺมา ตถาคโต สงฺขาเยกํ ปฏิเสวติ, สงฺขาเยกํ อธิวาเสติ, สงฺขาเยกํ ปริวชฺเชติ, สงฺขาเยกํ วิโนเทตี”ติ.\csromanspacer{} สาธุ สาธุ อนุรุทฺธา.\csromanspacer{} ตถาคตสฺส อนุรุทฺธา เย อาสวา สํกิเลสิกา โปโนพฺภวิกา สทรา ทุกฺขวิปากา อายติํ ชาติชรามรณิยา, ปหีนา เต อุจฺฉินฺนมูลา ตาลาวตฺถุกตา อนภาวํกตา อายติํ อนุปฺปาทธมฺมา.\csromanspacer{} เสยฺยถาปิ อนุรุทฺธา ตาโล มตฺถกจฺฉินฺโน อภพฺโพ ปุนวิรูฬฺหิยา, เอวเมว โข \csromanfolio{128}อนุรุทฺธา ตถาคตสฺส เย อาสวา สํกิเลสิกา โปโนพฺภวิกา สทรา ทุกฺขวิปากา อายติํ ชาติชรามรณิยา, ปหีนา เต อุจฺฉินฺนมูลา ตาลาวตฺถุกตา อนภาวํกตา อายติํ อนุปฺปาทธมฺมา.\csromanspacer{} ตสฺมา ตถาคโต สงฺขาเยกํ ปฏิเสวติ, สงฺขาเยกํ อธิวาเสติ, สงฺขาเยกํ ปริวชฺเชติ, สงฺขาเยกํ วิโนเทติ.}

\prose{ตํ กิํ มญฺญสิ อนุรุทฺธา.\csromanspacer{} กํ อตฺถวสํ สมฺปสฺสมาโน ตถาคโต สาวเก อพฺภตีเต กาลงฺกเต อุปปตฺตีสุ พฺยากโรติ “อสุ อมุตฺร อุปปนฺโน, อสุ อมุตฺร อุปปนฺโน”ติ.\csromanspacer{} ภควํมูลกา โน ภนฺเต ธมฺมา ภควํเนตฺติกา ภควํปฏิสรณา, สาธุ วต ภนฺเต ภควนฺตํเยว ปฏิภาตุ เอตสฺส ภาสิตสฺส อตฺโถ.\csromanspacer{} ภควโต สุตฺวา ภิกฺขู ธาเรสฺสนฺตีติ.\csromanspacer{} น โข อนุรุทฺธา ตถาคโต ชนกุหนตฺถํ น ชนลปนตฺตํ น ลาภสกฺการสิโลกานิสํสตฺถํ น ‘อิติ มํ ชโน ชานาตู’ติ สาวเก อพฺภตีเต กาลงฺกเต อุปปตฺตีสุ พฺยากโรติ “อสุ อมุตฺร อุปปนฺโน, อสุ อมุตฺร อุปปนฺโน”ติ.\csromanspacer{} สนฺติ จ โข อนุรุทฺธา กุลปุตฺตา สทฺธา อุฬารเวทา อุฬารปาโมชฺชา, เต ตํ สุตฺวา ตทตฺถาย จิตฺตํ อุปสํหรนฺติ.\csromanspacer{} เตสํ ตํ อนุรุทฺธา โหติ ทีฆรตฺตํ หิตาย สุขาย.}

\proseitem{169}{อิธานุรุทฺธา ภิกฺขุ สุณาติ “อิตฺถนฺนาโม ภิกฺขุ กาลงฺกโต\footnote{กาลกโต (สี, สฺยา, กํ, อิ)}, โส ภควตา พฺยากโต ‘อญฺญาย สณฺฐหี’ติ”.\csromanspacer{} โส โข ปนสฺส อายสฺมา สามํ ทิฏฺโฐ วา โหติ อนุสฺสวสฺสุโต วา “เอวํสีโล โส อายสฺมา อโหสิ อิติปิ, เอวํธมฺโม โส อายสฺมา อโหสิ อิติปิ, เอวํปญฺโญ โส อายสฺมา อโหสิ อิติปิ, เอวํวิหารี โส อายสฺมา อโหสิ อิติปิ, เอวํวิมุตฺโต โส อายสฺมา อโหสิ อิติปี”ติ.\csromanspacer{} โส ตสฺส สทฺธญฺจ สีลญฺจ สุตญฺจ จาคญฺจ ปญฺญญฺจ อนุสฺสรนฺโต ตทตฺถาย จิตฺตํ อุปสํหรติ.\csromanspacer{} เอวมฺปิ โข อนุรุทฺธา ภิกฺขุโน ผาสุวิหาโร โหติ.}

\prose{อิธานุรุทฺธา ภิกฺขุ สุณาติ “อิตฺถนฺนาโม ภิกฺขุ กาลงฺกโต, โส ภควตา พฺยากโต ‘ปญฺจนฺนํ โอรมฺภาคิยานํ สํโยชนานํ ปริกฺขยา โอปปาติโก, ตตฺถ ปรินิพฺพายี อนาวตฺติธมฺโม ตสฺมา โลกา’ติ”.\csromanspacer{} โส โข ปนสฺส อายสฺมา สามํ ทิฏฺโฐ วา โหติ อนุสฺสวสฺสุโต}

\vspace{\tipitakaparskip}
\csromancenter{นฬกปานสุตฺต (68) วา “เอวํสีโล โส อายสฺมา อโหสิ อิติปิ, เอวํธมฺโม ฯเปฯ เอวํปญฺโญ, เอวํวิหารี, เอวํวิมุตฺโต โส อายสฺมา อโหสิ อิติปี”ติ.\csromanspacer{} โส ตสฺส สทฺธญฺจ สีลญฺจ สุตญฺจ จาคญฺจ ปญฺญญฺจ อนุสฺสรนฺโต ตทตฺถาย จิตฺตํ อุปสํหรติ.\csromanspacer{} เอวมฺปิ โข อนุรุทฺธา ภิกฺขุโน ผาสุวิหาโร โหติ.}
\prose{\csromanfolio{129}อิธานุรุทฺธา ภิกฺขุ สุณาติ “อิตฺถนฺนาโม ภิกฺขุ กาลงฺกโต, โส ภควตา พฺยากโต ‘ติณฺณํ สํโยชนานํ ปริกฺขยา ราคโทสโมหานํ ตนุตฺตา สกทาคามี สกิเทว อิมํ โลกํ อาคนฺตฺวา ทุกฺขสฺสนฺตํ กริสฺสตี’ติ”.\csromanspacer{} โส โข ปนสฺส อายสฺมา สามํ ทิฏฺโฐ วา โหติ อนุสฺสวสฺสุโต วา “เอวํสีโล โส อายสฺมา อโหสิ อิติปิ, เอวํธมฺโม ฯเปฯ เอวํปญฺโญ, เอวํวิหารี, เอวํวิมุตฺโต โส อายสฺมา อโหสิ อิติปี”ติ.\csromanspacer{} โส ตสฺส สทฺธญฺจ สีลญฺจ สุตญฺจ จาคญฺจ ปญฺญญฺจ อนุสฺสรนฺโต ตทตฺถาย จิตฺตํ อุปสํหรติ.\csromanspacer{} เอวมฺปิ โข อนุรุทฺธา ภิกฺขุโน ผาสุวิหาโร โหติ.}

\prose{อิธานุรุทฺธา ภุกฺขุ สุณาติ “อิตฺถนฺนาโม ภิกฺขุ กาลงฺกโต, โส ภควตา พฺยากโต ‘ติณฺณํ สํโยชนานํ ปริกฺขยา โสตาปนฺโน อวินิปาตธมฺโม นิยโต สมฺโพธิปรายโณ’ติ”.\csromanspacer{} โส โข ปนสฺส อายสฺมา สามํ ทิฏฺโฐ วา โหติ อนุสฺสวสฺสุโต วา “เอวํสีโล โส อายสฺมา อโหสิ อิติปิ, เอวํธมฺโม ฯเปฯ เอวํปญฺโญ, เอวํวิหารี, เอวํวิมุตฺโต โส อายสฺมา อโหสิ อิติปี”ติ.\csromanspacer{} โส ตสฺส สทฺธญฺจ สีลญฺจ สุตญฺจ จาคญฺจ ปญฺญญฺจ อนุสฺสรนฺโต ตทตฺถาย จิตฺตํ อุปสํหรติ.\csromanspacer{} เอวมฺปิ โข อนุรุทฺธา ภิกฺขุโน ผาสุวิหาโร โหติ.}

\proseitem{170}{อิธานุรุทฺธา ภิกฺขุนี สุณาติ “อิตฺถนฺนามา ภิกฺขุนี กาลงฺกตา, สา ภควตา พฺยากตา ‘อญฺญาย สณฺฐหี’ติ”.\csromanspacer{} สา โข ปนสฺสา ภคินี สามํ ทิฏฺฐา วา โหติ อนุสฺสวสฺสุตา วา “เอวํสีลา สา ภคินี อโหสิ อิติปิ, เอวํธมฺมา สา ภคินี อโหสิ อิติปิ, เอวํปญฺญา สา ภคินี อโหสิ อิติปิ, เอวํวิหารินี สา ภคินี อโหสิ อิติปิ, เอวํ วิมุตฺตา สา ภคินี อโหสิ อิติปี”ติ.\csromanspacer{} สา ตสฺสา สทฺธญฺจ สีลญฺจ สุตญฺจ จาคญฺจ ปญฺญญฺจ อนุสฺสรนฺตี ตทตฺถาย จิตฺตํ อุปสํหรติ.\csromanspacer{} เอวมฺปิ โข อนุรุทฺธา ภิกฺขุนิยา ผาสุวิหาโร โหติ. \csromanfolio{130}อิธานุรุทฺธา ภิกฺขุนี สุณาติ “อิตฺถนฺนามา ภิกฺขุนี กาลงฺกตา, สา ภควตา พฺยากตา ‘ปญฺจนฺนํ โอรมฺภาคิยานํ สํโยชนานํ ปริกฺขยา โอปปาติกา, ตตฺถ ปรินิพฺพายินี อนาวตฺติธมฺมา ตสฺมา โลกา’ติ.\csromanspacer{} สา โข ปนสฺสา ภคินี สามํ ทิฏฺฐา วา โหติ อนุสฺสวสฺสุตา วา “เอวํสีลา สา ภคินี อโหสิ อิติปิ, เอวํธมฺมา ฯเปฯ เอวํปญฺญา, เอวํวิหารินี, เอวํวิมุตฺตา สา ภคินี อโหสิ อิติปี”ติ.\csromanspacer{} สา ตสฺสา สทฺธญฺจ สีลญฺจ สุตญฺจ จาคญฺจ ปญฺญญฺจ อนุสฺสรนฺตี ตทตฺถาย จิตฺตํ อุปสํหรติ.\csromanspacer{} เอวมฺปิ โข อนุรุทฺธา ภิกฺขุนิยา ผาสุวิหาโร โหติ. อิธานุรุทฺธา ภิกฺขุนี สุณาติ “อิตฺถนฺนามา ภิกฺขุนี กาลงฺกตา, สา ภควตา พฺยากตา ‘ติณฺณํ สํโยชนานํ ปริกฺขยา ราคโทสโมหานํ ตนุตฺตา สกทาคามินี สกิเทว อิมํ โลกํ อาคนฺตฺวา ทุกฺขสฺสนฺตํ กริสฺสตี’ติ”.\csromanspacer{} สา โข ปนสฺสา ภคินี สามํ ทิฏฺฐา วา โหติ อนุสฺสวสฺสุตา วา “เอวํสีลา สา ภคินี อโหสิ อิติปิ, เอวํธมฺมา ฯเปฯ เอวํปญฺญา, เอวํวิหารินี, เอวํวิมุตฺตา สา ภคินี อโหสิ อิติปี”ติ.\csromanspacer{} สา ตสฺสา สทฺธญฺจ สีลญฺจ สุตญฺจ จาคญฺจ ปญฺญญฺจ อนุสฺสรนฺตี ตทตฺถาย จิตฺตํ อุปสํหรติ.\csromanspacer{} เอวํมฺปิ โข อนุรุทฺธา ภิกฺขุนิยา ผาสุวิหาโร โหติ. อิธานุรุทฺธา ภิกฺขุนี สุณาติ “อิตฺถนฺนามา ภิกฺขุนี กาลงฺกตา, สา ภควตา พฺยากตา ‘ติณฺณํ สํโยชนานํ ปริกฺขยา โสตาปนฺนา อวินิปาตธมฺมา นิยตา สมฺโพธิปรายณา’ติ”.\csromanspacer{} สา โข ปนสฺส ภคินี สามํ ทิฏฺฐา วา โหติ อนุสฺสวสฺสุตา วา “เอวํสีลา สา ภคินี อโหสิ อิติปิ, เอวํธมฺมา, เอวํปญฺญา, เอวํวิหารินี, เอวํวิมุตฺตา สา ภคินี อโหสิ อิติปี”ติ.\csromanspacer{} สา ตสฺสา สทฺธญฺจ สีลญฺจ สุตญฺจ จาคญฺจ ปญฺญญฺจ อนุสฺสรนฺตี ตทตฺถาย จิตฺตํ อุปสํหรติ.\csromanspacer{} เอวมฺปิ โข อนุรุทฺธา ภิกฺขุนิยา ผาสุวิหาโร โหติ.}

\proseitem{171}{อิธานุรุทฺธา อุปาสโก สุณาติ “อิตฺถนฺนาโม อุปาสโก กาลงฺกโต, โส ภควตา พฺยากโต ‘ปญฺจนฺนํ โอรมฺภาคิยานํ สํโยชนานํ ปริกฺขยา โอปปาติโก, ตตฺถ ปรินิพฺพายี อนาวตฺติธมฺโม ตสฺมา โลกา’ติ”.\csromanspacer{} โส โข ปนสฺส อายสฺมา สามํ ทิฏฺโฐ วา โหติ อนุสฺสวสฺสุโต วา “เอวํสีโล โส อายสฺมา อโหสิ อิติปิ,}

\vspace{\tipitakaparskip}
\csromancenter{นฬกปานสุตฺต (68) เอวํธมฺโม โส อายสฺมา อโหสิ อิติปิ, เอวํปญฺโญ โส อายสฺมา อโหสิ อิติปิ, เอวํวิหารี โส อายสฺมา อโหสิ อิติปิ, เอวํวิมุตฺโต โส อายสฺมา อโหสิ อิติปี”ติ.\csromanspacer{} โส ตสฺส สทฺธญฺจ สุตญฺจ จาคญฺจ ปญฺญญฺจ อนุสฺสรนฺโต ตทตฺถาย จิตฺตํ อุปสํหรติ.\csromanspacer{} เอวมฺปิ โข อนุรุทฺธา อุปาสกสฺส ผาสุวิหาโร โหติ.}
\prose{\csromanfolio{131}อิธานุรุทฺธา อุปาสโก สุณาติ “อิตฺถนฺนาโม อุปาสโก กาลงฺกโต, โส ภควตา พฺยากโต ‘ติณฺณํ สํโยชนานํ ปริกฺขยา ราคโทสโมหานํ ตนุตฺตา สกทาคามี สกิเทว อิมํ โลกํ อาคนฺตฺวา ทุกฺขสฺสนฺตํ กริสฺสตี’ติ”.\csromanspacer{} โส โข ปนสฺส อายสฺมา สามํ ทิฏฺโฐ วา โหติ อนุสฺสวสฺสุโต วา “เอวํสีโล โส อายสฺมา อโหสิ อิติปิ, เอวํธมฺโม, เอวํปญฺโญ, เอวํวิหารี, เอวํวิมุตฺโต โส อายสฺมา อโหสิ อิติปี’ติ”.\csromanspacer{} โส ตสฺส สทฺธญฺจ สีลญฺจ สุตญฺจ จาคญฺจ ปญฺญญฺจ อนุสฺสรนฺโต ตทตฺถาย จิตฺตํ อุปสํหรติ.\csromanspacer{} เอวมฺปิ โข อนุรุทฺธา อุปาสกสฺส ผาสุวิหาโร โหติ.}

\prose{อิธานุรุทฺธา อุปาสโก สุณาติ “อิตฺถนฺนาโม อุปาสโก กาลงฺกโต, โส ภควตา พฺยากโต ‘ติณฺณํ สํโยชนานํ ปริกฺขยา โสตาปนฺโน อวินิปาตธมฺโม นิยโต สมฺโพธิปรายโณ’ติ”.\csromanspacer{} โส โข ปนสฺส อายสฺมา สามํ ทิฏฺโฐ วา โหติ อนุสฺสวสฺสุโต วา “เอวํสีโล โส อายสฺมา อโหสิ อิติปิ, เอวํธมฺโม, เอวํปญฺโญ, เอวํวิหารี, เอวํวิมุตฺโต โส อายสฺมา อโหสิ อิติปี”ติ.\csromanspacer{} โส ตสฺส สทฺธญฺจ สีลญฺจ สุตญฺจ จาคญฺจ ปญฺญญฺจ อนุสฺสรนฺโต ตทตฺถาย จิตฺตํ อุปสํหรติ.\csromanspacer{} เอวมฺปิ โข อนุรุทฺธา อุปาสกสฺส ผาสุวิหาโร โหติ.}

\proseitem{172}{อิธานุรุทฺธา อุปาสิกา สุณาติ “อิตฺถนฺนามา อุปาสิกา กาลงฺกตา, สา ภควตา พฺยากตา ‘ปญฺจนฺนํ โอรมฺภาคิยานํ สํโยชนานํ ปริกฺขยา โอปปาติกา, ตตฺถ ปรินิพฺพายินี อนาวตฺติธมฺมา ตสฺมา โลกา’ติ”.\csromanspacer{} สา โข ปนสฺสา ภคินี สามํ ทิฏฺฐา วา โหติ อนุสฺสวสฺสุตา วา “เอวํสีลา สา ภคินี อโหสิ อิติปิ, เอวํธมฺมา, เอวํ ปญฺญา, เอวํวิหารินี, เอวํวิมุตฺตา สา ภคินี อโหสิ อิติปี”ติ.\csromanspacer{} สา ตสฺสา สทฺธญฺจ สีลญฺจ สุตญฺจ จาคญฺจ ปญฺญญฺจ อนุสฺสรนฺตี ตทตฺถาย \csromanfolio{132}จิตฺตํ อุปสํหรติ.\csromanspacer{} เอวมฺปิ โข อนุรุทฺธา อุปาสิกาย ผาสุวิหาโร โหติ.}

\prose{อิธานุรุทฺธา อุปาสิกา สุณาติ “อิตฺถนฺนามา อุปาสิกา กาลงฺกตา, สา ภควตา พฺยากตา ‘ติณฺณํ สํโยชนานํ ปริกฺขยา ราคโทสโมหานํ ตนุตฺตา สกทาคามินี สกิเทว อิมํ โลกํ อาคนฺตฺวา ทุกฺขสฺสนฺตํ กริสฺสตี’ติ”.\csromanspacer{} สา โข ปนสฺสา ภคินี สามํ ทิฏฺฐา วา โหติ อนุสฺสวสฺสุตา วา “เอวํสีลา สา ภคินี อโหสิ อิติปิ, เอวํธมฺมา, เอวํปญฺญา, เอวํวิหารินี, เอวํวิมุตฺตา สา ภคินี อโหสิ อิติปี”ติ.\csromanspacer{} สา ตสฺสา สทฺธญฺจ สีลญฺจ สุตญฺจ จาคญฺจ ปญฺญญฺจ อนุสฺสรนฺตี ตทตฺถาย จิตฺตํ อุปสํหรติ.\csromanspacer{} เอวมฺปิ โข อนุรุทฺธา อุปาสิกาย ผาสุวิหาโร โหติ.}

\prose{อิธานุรุทฺธา อุปาสิกา สุณาติ “อิตฺถนฺนามา อุปาสิกา กาลงฺกตา, สา ภควตา พฺยากตา ‘ติณฺณํ สํโยชนานํ ปริกฺขยา โสตาปนฺนา อวินิปาตธมฺมา นิยตา สมฺโพธิปรายณา’ติ”.\csromanspacer{} สา โข ปนสฺสา ภคินี สามํ ทิฏฺฐา วา โหติ อนุสฺสวสฺสุตา วา “เอวํสีลา สา ภคินี อโหสิ อิติปิ, เอวํธมฺมา สา ภคินี อโหสิ อิติปิ, เอวํปญฺญา สา ภคินี อโหสิ อิติปิ, เอวํวิหารินี สา ภคินี อโหสิ อิติปิ, เอวํวิมุตฺตา สา ภคินี อโหสิ อิติปี”ติ.\csromanspacer{} สา ตสฺสา สทฺธญฺจ สีลญฺจ สุตญฺจ จาคญฺจ ปญฺญญฺจ อนุสฺสรนฺตี ตทตฺถาย จิตฺตํ อุปสํหรติ.\csromanspacer{} เอวมฺปิ โข อนุรุทฺธา อุปาสิกาย ผาสุวิหาโร โหติ.}

\prose{อิติ โข อนุรุทฺธา ตถาคโต น ชนกุหนตฺถํ น ชนลปนตฺถํ น ลาภสกฺการสิโลกานิสํสตฺถํ น ‘อิติ มํ ชโน ชานาตู’ติ สาวเก อพฺภตีเต กาลงฺกเต อุปปตฺตีสุ พฺยากโรติ “อสุ อมุตฺร อุปปนฺโน อสุ อมุตฺร อุปปนฺโน”ติ.\csromanspacer{} สนฺติ จ โข อนุรุทฺธา กุลปุตฺตา สทฺธา อุฬารเวทา อุฬาปารโมชฺชา, เต ตํ สุตฺวา ตทตฺถาย จิตฺตํ อุปสํหรนฺติ.\csromanspacer{} เตสํ ตํ อนุรุทฺธา โหติ ทีฆรตฺตํ หิตาย สุขายาติ.}

\prose{อิทมโวจ ภควา.\csromanspacer{} อตฺตมโน อายสฺมา อนุรุทฺโธ ภควโต ภาสิตํ อภินนฺทีติ.}

\nitthitam{นฬกปานสุตฺตํ นิฏฺฐิตํ อฏฺฐมํ.}
\csromanmatikaanchor{mtk.22}
\csromantocmarkref{subsection}{9. โคลิยานิสุตฺต}{mtk.22}
\bookmark[dest={mtk.22},level=2]{9. โคลิยานิสุตฺต}
\csromanheader{2}{9. โคลิยานิสุตฺต (69) \textbf{9. โคลิยานิสุตฺต}}
\proseitem{173}{\csromanfolio{133}เอวํ เม สุตํ\csromandash{}เอกํ สมยํ ภควา ราชคเห วิหรติ เวฬุวเน กลนฺทกนิวาเป.\csromanspacer{} เตน โข ปน สมเยน โคลิยานิ\footnote{คุลิสฺสานิ (สี, อิ), โคลิสฺสานิ (สฺยา, กํ)} นาม ภิกฺขุ อารญฺญิโก\footnote{อารญฺญโก (สพฺพตฺถ)} ปทสมาจาโร\footnote{ปทรสมาจาโร (สี, สฺยา, กํ, อิ)} สํฆมชฺเฌ โอสโฏ โหติ เกนจิเทว กรณีเยน.\csromanspacer{} ตตฺร โข อายสฺมา สาริปุตฺโต โคลิยานิํ ภิกฺขุํ อารพฺภ ภิกฺขู อามนฺเตสิ\csromandash{}}

\prose{อารญฺญิเกนาวุโส ภิกฺขุนา สํฆคเตน สํเฆ วิหรนฺเตน สพฺรหฺมจารีสุ สคารเวน ภวิตพฺพํ สปฺปติสฺเสน.\csromanspacer{} สเจ อาวุโส อารญฺญิโก ภิกฺขุ สํฆคโต สํเฆ วิหรนฺโต สพฺรหฺมจารีสุ อคารโว โหติ อปฺปติสฺโส, ตสฺส ภวนฺติ วตฺตาโร “กิํ ปนิมสฺสายสฺมโต อารญฺญิกสฺส เอกสฺสารญฺเญ เสริวิหาเรน, โย อยมายสฺมา สพฺรหฺมจารีสุ อคารโว โหติ อปฺปติสฺโส”ติ, ตสฺส\footnote{อปฺปติสฺโสติสฺส (สี, อิ)} ภวนฺติ วตฺตาโร.\csromanspacer{} ตสฺมา อารญฺญิเกน ภิกฺขุนา สํฆคเตน สํเฆ วิหรนฺเตน สพฺรหฺมจารีสุ สคารเวน ภวิตพฺพํ สปฺปติสฺเสน. (1)}

\prose{อารญฺญิเกนาวุโส ภิกฺขุนา สํฆคเตน สํเฆ วิหรนฺเตน อาสนกุสเลน ภวิตพฺพํ “อิติ เถเร จ ภิกฺขู นานุปขชฺช นิสีทิสฺสามิ นเว จ ภิกฺขู น อาสเนน ปฏิพาหิสฺสามี”ติ.\csromanspacer{} สเจ อาวุโส อารญฺญิโก ภิกฺขุ สํฆคโต สํเฆ วิหรนฺโต น อาสนกุสโล โหติ.\csromanspacer{} ตสฺส ภวนฺติ วตฺตาโร “กิํ ปนิมสฺสายสฺมโต อารญฺญิกสฺส เอกสฺสารญฺเญ เสริวิหาเรน, โย อยมายสฺมา อาสนกุสโล น โหตี”ติ\footnote{โย อยมายสฺมา อาภิสมาจาริกมฺปิ ธมฺมํ น ชานาตีติ (สี, สฺยา, กํ, อิ)}.\csromanspacer{} ตสฺส ภวนฺติ วตฺตาโร.\csromanspacer{} ตสฺมา อารญฺญิเกน ภิกฺขุนา สํฆคเตน สํเฆ วิหรนฺเตน อาสนกุสเลน ภวิตพฺพํ. (2)}

\prose{อารญฺญิเกนาวุโส ภิกฺขุนา สํฆคเตน สํเฆ วิหรนฺเตน อาภิสมาจาริโกปิ ธมฺโม ชานิตพฺโพ.\csromanspacer{} สเจ อาวุโส อารญฺญิโก ภิกฺขุ สํฆคโต สํเฆ วิหรนฺโต อาภิสมาจาริกมฺปิ ธมฺมํ น ชานาติ.\csromanspacer{} ตสฺส ภวนฺติ วตฺตาโร “กิํ ปนิมสฺสายสฺมโต อารญฺญิกสฺส เอกสฺสารญฺเญ เสริวิหาเรน, โย อยมายสฺมา อาภิสมาจาริกมฺปิ \csromanfolio{134}ธมฺมํ น ชานาตี”ติ.\csromanspacer{} ตสฺส ภวนฺติ วตฺตาโร.\csromanspacer{} ตสฺมา อารญฺญิเกน ภิกฺขุนา สํฆคเตน สํเฆ วิหรนฺเตน อาภิสมาจาริโกปิ ธมฺโม ชานิตพฺโพ\footnote{อยํ อาภิสมาจาริกตติยวาโร สี-สฺยา-กํ-อิ-โปตฺถเกสุ น ทิสฺสติ.}. (3)}

\prose{อารญฺญิเกนาวุโส ภิกฺขุนา สํฆคเตน สํเฆ วิหรนฺเตน นาติกาเลน คาโม ปวิสิตพฺโพ, นาติทิวา\footnote{น ทิวา (สฺยา, กํ, อิ, ก)} ปฏิกฺกมิตพฺพํ.\csromanspacer{} สเจ อาวุโส อารญฺญิโก ภิกฺขุ สํฆคโต สํเฆ วิหรนฺโต อติกาเลน คามํ ปวิสติ, อติทิวา ปฏิกฺกมติ.\csromanspacer{} ตสฺส ภวนฺติ วตฺตาโร “กิํ ปนิมสฺสายสฺมโต อารญฺญิกสฺส เอกสฺสารญฺเญ เสริวิหาเรน, โย อยมายสฺมา อติกาเลน คามํ ปวิสติ, อติทิวา ปฏิกฺกมตี”ติ.\csromanspacer{} ตสฺส ภวนฺติ วตฺตาโร.\csromanspacer{} ตสฺมา อารญฺญิเกน ภิกฺขุนา สํฆคเตน สํเฆ วิหรนฺเตน นาติกาเลน คาโม ปวิสิตพฺโพ, นาติทิวา ปฏิกฺกมิตพฺพํ. (4)}

\prose{อารญฺญิเกนาวุโส ภิกฺขุนา สํฆคเตน สํเฆ วิหรนฺเตน น ปุเรภตฺตํ ปจฺฉาภตฺตํ กุเลสุ จาริตฺตํ อาปชฺชิตพฺพํ.\csromanspacer{} สเจ อาวุโส อารญฺญิโก ภิกฺขุ สํฆคโต สํเฆ วิหรนฺโต ปุเรภตฺตํ ปจฺฉาภตฺตํ กุเลสุ จาริตฺตํ อาปชฺชติ.\csromanspacer{} ตสฺส ภวนฺติ วตฺตาโร “อยํ นูนิมสฺสายสฺมโต อารญฺญิกสฺส เอกสฺสารญฺเญ เสริวิหาเรน วิหรโต วิกาลจริยา พหุลีกตา, ตเมนํ สํฆคตํปิ สมุทาจรตี”ติ.\csromanspacer{} ตสฺส ภวนฺติ วตฺตาโร.\csromanspacer{} ตสฺมา อารญฺญิเกน ภิกฺขุนา สํฆคเตน สํเฆ วิหรนฺเตน น ปุเรภตฺตํ ปจฺฉาภตฺตํ กุเลสุ จาริตฺตํ อาปชฺชิตพฺพํ. (5)}

\prose{อารญฺญิเกนาวุโส ภิกฺขุนา สํฆคเตน สํเฆ วิหรนฺเตน อนุทฺธเตน ภวิตพฺพํ อจปเลน.\csromanspacer{} สเจ อาวุโส อารญฺญิโก ภิกฺขุ สํฆคโต สํเฆ วิหรนฺโต อุทฺธโต โหติ จปโล.\csromanspacer{} ตสฺส ภวนฺติ วตฺตาโร “อิทํ นูนิมสฺสายสฺมโต อารญฺญิกสฺส เอกสฺสารญฺเญ เสริวิหาเรน วิหรโต อุทฺธจฺจํ จาปลฺยํ พหุลีกตํ, ตเมนํ สํฆคตํปิ สมุทาจรตี”ติ.\csromanspacer{} ตสฺส ภวนฺติ วตฺตาโร.\csromanspacer{} ตสฺมา อารญฺญิเกน ภิกฺขุนา สํฆคเตน สํเฆ วิหรนฺเตน อนุทฺธเตน ภวิตพฺพํ อจปเลน. (6)}

\csromanheader{2}{9. โคลิยานิสุตฺต (69)}
\prose{\csromanfolio{135}อารญฺญิเกนาวุโส ภิกฺขุนา สํฆคเตน สํเฆ วิหรนฺเตน อมุขเรน ภวิตพฺพํ อวิกิณฺณวาเจน.\csromanspacer{} สเจ อาวุโส อารญฺญิโก ภิกฺขุ สํฆคโต สํเฆ วิหรนฺโต มุขโร โหติ วิกิณฺณวาโจ.\csromanspacer{} ตสฺส ภวนฺติ วตฺตาโร “กิํ ปนิมสฺสายสฺมโต อารญฺญิกสฺส เอกสฺสารญฺเญ เสริวิหาเรน, โย อยมายสฺมา มุขโร วิกิณฺณวาโจ”ติ.\csromanspacer{} ตสฺส ภวนฺติ วตฺตาโร.\csromanspacer{} ตสฺมา อารญฺญิเกน ภิกฺขุนา สํฆคเตน สํเฆ วิหรนฺเตน อมุขเรน ภวิตพฺพํ อวิกิณฺณวาเจน. (7)}

\prose{อารญฺญิเกนาวุโส ภิกฺขุนา สํฆคเตน สํเฆ วิหรนฺเตน สุวเจน\footnote{สุพฺพเจน (สี, ก)} ภวิตพฺพํ กลฺยาณมิตฺเตน.\csromanspacer{} สเจ อาวุโส อารญฺญิโก ภิกฺขุ สํฆคโต สํเฆ วิหรนฺโต ทุพฺพโจ โหติ ปาปมิตฺโต.\csromanspacer{} ตสฺส ภวนฺติ วตฺตาโร “กิํ ปนิมสฺสายสฺมโต อารญฺญิกสฺส เอกสฺสารญฺเญ เสริวิหาเรน, โย อยมายสฺมา ทุพฺพโจ ปาปมิตฺโต”ติ.\csromanspacer{} ตสฺส ภวนฺติ วตฺตาโร.\csromanspacer{} ตสฺมา อารญฺญิเกน ภิกฺขุนา สํฆคเตน สํเฆ วิหรนฺเตน สุวเจน ภวิตพฺพํ กลฺยาณมิตฺเตน. (8)}

\prose{อารญฺญิเกนาวุโส ภิกฺขุนา อินฺทฺริเยสุ คุตฺตทฺวาเรน ภวิตพฺพํ.\csromanspacer{} สเจ อาวุโส อารญฺญิโก ภิกฺขุ อินฺทฺริเยสุ อคุตฺตทฺวาโร โหติ.\csromanspacer{} ตสฺส ภวนฺติ วตฺตาโร “กิํ ปนิมสฺสายสฺมโต อารญฺญิกสฺส เอกสฺสารญฺเญ เสริวิหาเรน, โย อยมายสฺมา อินฺทฺริเยสุ อคุตฺตทฺวาโร”ติ.\csromanspacer{} ตสฺส ภวนฺติ วตฺตาโร.\csromanspacer{} ตสฺมา อารญฺญิเกน ภิกฺขุนา อินฺทฺริเยสุ คุตฺตทฺวาเรน ภวิตพฺพํ. (9)}

\prose{อารญฺญิเกนาวุโส ภิกฺขุนา โภชเน มตฺตญฺญุนา ภวิตพฺพํ.\csromanspacer{} สเจ อาวุโส อารญฺญิโก ภิกฺขุ โภชเน อมตฺตญฺญู โหติ.\csromanspacer{} ตสฺส ภวนฺติ วตฺตาโร “กิํ ปนิมสฺสายสฺมโต อารญฺญิกสฺส เอกสฺสารญฺเญ เสริวิหาเรน, โย อยมายสฺมา โภชเน อมตฺตญฺญู”ติ.\csromanspacer{} ตสฺส ภวนฺติ วตฺตาโร.\csromanspacer{} ตสฺมา อารญฺญิเกน ภิกฺขุนา โภชเน มตฺตญฺญุนา ภวิตพฺพํ. (10)}

\prose{อารญฺญิเกนาวุโส ภิกฺขุนา ชาคริยํ อนุยุตฺเตน ภวิตพฺพํ.\csromanspacer{} สเจ อาวุโส อารญฺญิโก ภิกฺขุ ชาคริยํ อนนุยุตฺโต โหติ.}

\prose{\csromanfolio{136}ตสฺส ภวนฺติ วตฺตาโร “กิํ ปนิมสฺสายสฺมโต อารญฺญิกสฺส เอกสฺสารญฺเญ เสริวิหาเรน, โย อยมายสฺมา ชาคริยํ อนนุยุตฺโต”ติ.\csromanspacer{} ตสฺส ภวนฺติ วตฺตาโร.\csromanspacer{} ตสฺมา อารญฺญิเกน ภิกฺขุนา ชาคริยํ อนุยุตฺเตน ภวิตพฺพํ. (11)}

\prose{อารญฺญิเกนาวุโส ภิกฺขุนา อารทฺธวีริเยน ภวิตพฺพํ.\csromanspacer{} สเจ อาวุโส อารญฺญิโก ภิกฺขุ กุสีโต โหติ.\csromanspacer{} ตสฺส ภวนฺติ วตฺตาโร “กิํ ปนิมสฺสายสฺมโต อารญฺญิกสฺส เอกสฺสารญฺเญ เสริวิหาเรน, โย อยมายสฺมา กุสีโต”ติ.\csromanspacer{} ตสฺส ภวนฺติ วตฺตาโร.\csromanspacer{} ตสฺมา อารญฺญิเกน ภิกฺขุนา อารทฺธวีริเยน ภวิตพฺพํ. (12)}

\prose{อารญฺญิเกนาวุโส ภิกฺขุนา อุปฏฺฐิตสฺสตินา ภวิตพฺพํ.\csromanspacer{} สเจ อาวุโส อารญฺญิโก ภิกฺขุ มุฏฺฐสฺสตี โหติ.\csromanspacer{} ตสฺส ภวนฺติ วตฺตาโร “กิํ ปนิมสฺสายสฺมโต อารญฺญิกสฺส เอกสฺสารญฺเญ เสริวิหาเรน, โย อยมายสฺมา มุฏฺฐสฺสตี”ติ.\csromanspacer{} ตสฺส ภวนฺติ วตฺตาโร.\csromanspacer{} ตสฺมา อารญฺญิเกน ภิกฺขุนา อุปฏฺฐิตสฺสตินา ภวิตพฺพํ. (13)}

\prose{อารญฺญิเกนาวุโส ภิกฺขุนา สมาหิเตน ภวิตพฺพํ.\csromanspacer{} สเจ อาวุโส อารญฺญิโก ภิกฺขุ อสมาหิโต โหติ.\csromanspacer{} ตสฺส ภวนฺติ วตฺตาโร “กิํ ปนิมสฺสายสฺมโต อารญฺญิกสฺส เอกสฺสารญฺเญ เสริวิหาเรน, โย อยมายสฺมา อสมาหิโต”ติ.\csromanspacer{} ตสฺส ภวนฺติ วตฺตาโร.\csromanspacer{} ตสฺมา อารญฺญิเกน ภิกฺขุนา สมาหิเตน ภวิตพฺพํ. (14)}

\prose{อารญฺญิเกนาวุโส ภิกฺขุนา ปญฺญวตา ภวิตพฺพํ.\csromanspacer{} สเจ อาวุโส อารญฺญิโก ภิกฺขุ ทุปฺปญฺโญ โหติ.\csromanspacer{} ตสฺส ภวนฺติ วตฺตาโร “กิํ ปนิมสฺสายสฺมโต อารญฺญิกสฺส เอกสฺสารญฺเญ เสริวิหาเรน, โย อยมายสฺมา ทุปฺปญฺโญ”ติ.\csromanspacer{} ตสฺส ภวนฺติ วตฺตาโร.\csromanspacer{} ตสฺมา อารญฺญิเกน ภิกฺขุนา ปญฺญวตา ภวิตพฺพํ. (15)}

\prose{อารญฺญิเกนาวุโส ภิกฺขุนา อภิธมฺเม อภิวินเย โยโค กรณีโย.\csromanspacer{} สนฺตาวุโส อารญฺญิกํ ภิกฺขุํ อภิธมฺเม อภิวินิเย ปญฺหํ ปุจฺฉิตาโร.\csromanspacer{} สเจ อาวุโส อารญฺญิโก ภิกฺขุ อภิธมฺเม อภิวินเย ปญฺหํ ปุฏฺโฐ น สมฺปายติ.\csromanspacer{} ตสฺส ภวนฺติ วตฺตาโร “กิํ ปนิมสฺสายสฺมโต อารญฺญิกสฺส เอกสฺสารญฺเญ เสริวิหาเรน, โย อยมายสฺมา อภิธมฺเม อภิวินเย ปญฺหํ ปุฏฺโฐ น สมฺปายตี”ติ.\csromanspacer{} ตสฺส ภวนฺติ วตฺตาโร.}

\vspace{\tipitakaparskip}
\csromancenter{โคลิยานิสุตฺต (69) ตสฺมา อารญฺญิเกน ภิกฺขุนา อภิธมฺเม อภิวินเย โยโค กรณีโย. (16)}
\prose{\csromanfolio{137}อารญฺญิเกนาวุโส ภิกฺขุนา เย เต สนฺตา วิโมกฺขา อติกฺกมฺม รูเป อารุปฺปา, ตตฺถ โยโค กรณีโย.\csromanspacer{} สนฺตาวุโส อารญฺญิกํ ภิกฺขุํ เย เต สนฺตา วิโมกฺขา อติกฺกมฺม รูเป อารุปฺปา, ตตฺถ ปญฺหํ ปุจฺฉิตาโร.\csromanspacer{} สเจ อาวุโส อารญฺญิโก ภิกฺขุ เย เต สนฺตา วิโมกฺขา อติกฺกมฺม รูเป อารุปฺปา, ตตฺถ ปญฺหํ ปุฏฺโฐ น สมฺปายติ.\csromanspacer{} ตสฺส ภวนฺติ วตฺตาโร “กิํ ปนิมสฺสายสฺมโต อารญฺญิกสฺส เอกสฺสารญฺเญ เสริวิหาเรน, โย อยมายสฺมา เย เต สนฺตา วิโมกฺขา อติกฺกมฺม รูเป อารุปฺปา, ตตฺถ ปญฺหํ ปุฏฺโฐ น สมฺปายตี”ติ.\csromanspacer{} ตสฺส ภวนฺติ วตฺตาโร.\csromanspacer{} ตสฺมา อารญฺญิเกน ภิกฺขุนา เย เต สนฺตา วิโมกฺขา อติกฺกมฺม รูเป อารุปฺปา, ตตฺถ โยโค กรณีโย. (17)}

\prose{อารญฺญิเกนาวุโส ภิกฺขุนา อุตฺตริ มนุสฺสธมฺเม โยโค กรณีโย.\csromanspacer{} สนฺตาวุโส อารญฺญิกํ ภิกฺขุํ อุตฺตริ มนุสฺสธมฺเม ปญฺหํ ปุจฺฉิตาโร.\csromanspacer{} สเจ อาวุโส อารญฺญิโก ภิกฺขุ อุตฺตริ มนุสฺสธมฺเม ปญฺหํ ปุฏฺโฐ น สมฺปายติ.\csromanspacer{} ตสฺส ภวนฺติ วตฺตาโร “กิํ ปนิมสฺสายสฺมโต อารญฺญิกสฺส เอกสฺสารญฺเญ เสริวิหาเรน, โย อยมายสฺมา ยสฺสตฺถาย ปพฺพชิโต, ตมตฺถํ น ชานาตี”ติ.\csromanspacer{} ตสฺส ภวนฺติ วตฺตาโร.\csromanspacer{} ตสฺมา อารญฺญิเกน ภิกฺขุนา อุตฺตริ มนุสฺสธมฺเม โยโค กรณีโยติ. (18)}

\prose{เอวํ วุตฺเต อายสฺมา มหาโมคฺคลฺลาโน\footnote{มหาโมคฺคลาโน (ก)} อายสฺมนฺตํ สาริปุตฺตํ เอตทโวจ “อารญฺญิเกเนว นุ โข อาวุโส สาริปุตฺต ภิกฺขุนา อิเม ธมฺมา สมาทาย วตฺติตพฺพา, อุทาหุ คามนฺตวิหารินาปีติ.\csromanspacer{} อารญฺญิเกนาปิ โข อาวุโส โมคฺคลฺลาน ภิกฺขุนา อิเม ธมฺมา สมาทาย วตฺติตพฺพา, ปเคว คามนฺตวิหารินา”ติ.}

\nitthitam{โคลิยานิสุตฺตํ นิฏฺฐิตํ นวมํ.}
\csromanmatikaanchor{mtk.23}
\csromantocmarkref{subsection}{10. กีฏาคิริสุตฺต}{mtk.23}
\bookmark[dest={mtk.23},level=2]{10. กีฏาคิริสุตฺต}
\csromanheader{2}{10. กีฏาคิริสุตฺต}
\proseitem{174}{\csromanfolio{138}เอวํ เม สุตํ\csromandash{}เอกํ สมยํ ภควา กาสีสุ จาริกํ จรติ มหตา ภิกฺขุสํเฆน สทฺธิํ.\csromanspacer{} ตตฺร โข ภควา ภิกฺขู อามนฺเตสิ “อหํ โข ภิกฺขเว อญฺญเตฺรว รตฺติโภชนา\footnote{รตฺติโภชนํ (ก)} ภุญฺชามิ.\csromanspacer{} อญฺญตฺร โข ปนาหํ ภิกฺขเว รตฺติโภชนา ภุญฺชมาโน อปฺปาพาธตญฺจ สญฺชานามิ อปฺปาตงฺกตญฺจ ลหุฏฺฐานญฺจ พลญฺจ ผาสุวิหารญฺจ.\csromanspacer{} เอถ ตุเมฺหปิ ภิกฺขเว อญฺญเตฺรว รตฺติโภชนา ภุญฺชถ.\csromanspacer{} อญฺญตฺร โข ปน ภิกฺขเว ตุเมฺหปิ รตฺติโภชนา ภุญฺชมานา อปฺปาพาธตญฺจ สญฺชานิสฺสถ อปฺปาตงฺกตญฺจ ลหุฏฺฐานญฺจ พลญฺจ ผาสุวิหารญฺจา”ติ. “เอวํ ภนฺเต”ติ โข เต ภิกฺขู ภควโต ปจฺจสฺโสสุํ.\csromanspacer{} อถ โข ภควา กาสีสุ อนุปุพฺเพน จาริกํ จรมาโน เยน กีฏาคิริ นาม กาสีนํ นิคโม ตทวสริ.\csromanspacer{} ตตฺร สุทํ ภควา กีฏาคิริสฺมิํ วิหรติ กาสีนํ นิคเม.}

\proseitem{175}{เตน โข ปน สมเยน อสฺสชิปุนพฺพสุกา นาม ภิกฺขู กีฏาคิริสฺมิํ อาวาสิกา โหนฺติ.\csromanspacer{} อถ โข สมฺพหุลา ภิกฺขู เยน อสฺสชิปุนพฺพสุกา ภิกฺขู เตนุปสงฺกมิํสุ, อุปสงฺกมิตฺวา อสฺสชิปุนพฺพสุเก ภิกฺขู เอตทโวจุํ “ภควา โข อาวุโส อญฺญเตฺรว รตฺติโภชนา ภุญฺชติ ภิกฺขุสํโฆ จ.\csromanspacer{} อญฺญตฺร โข ปนาวุโส รตฺติโภชนา ภุญฺชมานา อปฺปาพาธตญฺจ สญฺชานนฺติ อปฺปาตงฺกตญฺจ ลหุฏฺฐานญฺจ พลญฺจ ผาสุวิหารญฺจ.\csromanspacer{} เอถ ตุเมฺหปิ อาวุโส อญฺญเตฺรว รตฺติโภชนา ภุญฺชถ.\csromanspacer{} อญฺญตฺร โข ปนาวุโสตุเมฺหปิ รตฺติโภชนา ภุญฺชมานา อปฺปาพาธตญฺจ สญฺชานิสฺสถ อปฺปาตงฺกตญฺจ ลหุฏฺฐานญฺจ พลญฺจ ผาสุวิหารญฺจา”ติ.\csromanspacer{} เอวํ วุตฺเต อสฺสชิปุนพฺพสุกา ภิกฺขู เต ภิกฺขู เอตทโวจุํ “มยํ โข อาวุโส สายญฺเจว ภุญฺชาม ปาโต จ ทิวา จ วิกาเล, เต มยํ สายญฺเจว ภุญฺชมานา ปาโต จ ทิวา จ วิกาเล อปฺปาพาธตญฺจ สญฺชานาม อปฺปาตงฺกตญฺจ ลหุฏฺฐานญฺจ พลญฺจ ผาสุวิหารญฺจ, เต มยํ กิํ สนฺทิฏฺฐิกํ หิตฺวา กาลิกํ อนุธาวิสฺสาม, สายญฺเจว มยํ ภุญฺชิสฺสาม ปาโต จ ทิวา จ วิกาเล”ติ.}

\prose{ยโต โข เต ภิกฺขู นาสกฺขิํสุ อสฺสชิปุนพฺพสุเก ภิกฺขู สญฺญาเปตุํ, อถ เยน ภควา เตนุปสงฺกมิํสุ, อุปสงฺกมิตฺวา ภควนฺตํ}

\vspace{\tipitakaparskip}
\csromancenter{กีฏาคิริสุตฺต (70) อภิวาเทตฺวา เอกมนฺตํ นิสีทิํสุ, เอกมนฺตํ นิสินฺนา โข เต ภิกฺขู ภควนฺตํ เอตทโวจุํ “อิธ มยํ ภนฺเต เยน อสฺสชิปุนพฺพสุกา ภิกฺขู เตนุปสงฺกมิมฺห, อุปสงฺกมิตฺวา อสฺสชิปุนพฺพสุเก ภิกฺขู เอตทโวจุมฺห ‘ภควา โข อาวุโส อญฺญเตฺรว รตฺติโภชนา ภุญฺชติ ภิกฺขุสํโฆ จ, อญฺญตฺร โข ปนาวุโส รตฺติโภชนา ภุญฺชมานา อปฺปาพาธตญฺจ สญฺชานนฺติ อปฺปาตงฺกตญฺจ ลหุฏฺฐานญฺจ พลญฺจ ผาสุวิหารญฺจ.\csromanspacer{} เอถ ตุเมฺหปิ อาวุโส อญฺญเตฺรว รตฺติโภชนา ภุญฺชถ, อญฺญตฺร โข ปนาวุโส ตุเมฺหปิ รตฺติโภชนา ภุญฺชมานา อปฺปาพาธตญฺจ สญฺชานิสฺสถ อปฺปาตงฺกตญฺจ ลหุฏฺฐานญฺจ พลญฺจ ผาสุวิหารญฺจา’ติ.\csromanspacer{} เอวํ วุตฺเต ภนฺเต อสฺสชิปุนพฺพสุกา ภิกฺขู อเมฺห เอตทโวจุํ ‘มยํ โข อาวุโส สายญฺเจว ภุญฺชาม ปาโต จ ทิวา จ วิกาเล, เต มยํ สายญฺเจว ภุญฺชมานา ปาโต จ ทิวา จ วิกาเล อปฺปาพาธตญฺจ สญฺชานาม อปฺปาตงฺกตญฺจ ลหุฏฺฐานญฺจ พลญฺจ ผาสุวิหารญฺจ, เต มยํ กิํ สนฺทิฏฺฐิกํ หิตฺวา กาลิกํ อนุธาวิสฺสาม, สายญฺเจว มยํ ภุญฺชิสฺสาม ปาโต จ ทิวา จ วิกาเล’ติ.\csromanspacer{} ยโต โข มยํ ภนฺเต นาสกฺขิมฺห อสฺสชิปุนพฺพสุเก ภิกฺขู สญฺญาเปตุํ, อถ มยํ เอตมตฺถํ ภควโต อาโรเจมา”ติ.}
\proseitem{176}{\csromanfolio{139}อถ โข ภควา อญฺญตรํ ภิกฺขุํ อามนฺเตสิ “เอหิ ตฺวํ ภิกฺขุ มม วจเนน อสฺสชิปุนพฺพสุเก ภิกฺขู อามนฺเตหิ ‘สตฺถา อายสฺมนฺเต อามนฺเตตี’ติ”. “เอวํ ภนฺเต”ติ โข โส ภิกฺขุ ภควโต ปฏิสฺสุตฺวา เยน อสฺสชิปุนพฺพสุกา ภิกฺขู เตนุปสงฺกมิ, อุปสงฺกมิตฺวา อสฺสชิปุนพฺพสุเก ภิกฺขู เอตทโวจ “สตฺถา อายสฺมนฺเต อามนฺเตตี”ติ. “เอวมาวุโส”ติ โข อสฺสชิปุนพฺพสุกา ภิกฺขู ตสฺส ภิกฺขุโน ปฏิสฺสุตฺวา เยน ภควา เตนุปสงฺกมิํสุ, อุปสงฺกมิตฺวา ภควนฺตํ อภิวาเทตฺวา เอกมนฺตํ นิสีทิํสุ.\csromanspacer{} เอกมนฺตํ นิสินฺเน โข อสฺสชิปุนพฺพสุเก ภิกฺขู ภควา เอตทโวจ “สจฺจํ กิร ภิกฺขเว สมฺพหุลา ภิกฺขู ตุเมฺห อุปสงฺกมิตฺวา เอตทโวจุํ ‘ภควา โข อาวุโส อญฺญเตฺรว รตฺติโภชนา ภุญฺชติ ภิกฺขุสํโฆ จ, อญฺญตฺร โข ปนาวุโส รตฺติโภชนา ภุญฺชมานา อปฺปาพาธตญฺจ สญฺชานนฺติ อปฺปาตงฺกตญฺจ ลหุฏฺฐานญฺจ พลญฺจ ผาสุวิหารญฺจ.\csromanspacer{} เอถ ตุเมฺหปิ อาวุโส อญฺญเตฺรว รตฺติโภชนา ภุญฺชถ, อญฺญตฺร โข ปนาวุโส ตุเมฺหปิ รตฺติโภชนา ภุญฺชมานา อปฺปาพาธตญฺจ สญฺชานิสฺสถ อปฺปาตงฺกตญฺจ ลหุฏฺฐานญฺจ พลญฺจ ผาสุวิหารญฺจา’ติ.\csromanspacer{} เอวํ \csromanfolio{140}วุตฺเต กิร\footnote{กิํ นุ (ก)} ภิกฺขเว ตุเมฺห เต ภิกฺขู เอวํ อวจุตฺถ ‘มยํ โข ปนาวุโส สายญฺเจว ภุญฺชาม ปาโต จ ทิวา จ วิกาเล, เต มยํ สายญฺเจว ภุญฺชมานา ปาโต จ ทิวา จ วิกาเล อปฺปาพาธตญฺจ สญฺชานาม อปฺปาตงฺกตญฺจ ลหุฏฺฐานญฺจ พลญฺจ ผาสุวิหารญฺจ, เต มยํ กิํ สนฺทิฏฺฐิกํ หิตฺวา กาลิกํ อนุธาวิสฺสาม, สายญฺเจว มยํ ภุญฺชิสฺสาม ปาโต จ ทิวา จ วิกาเล’ติ”.\csromanspacer{} เอวํ ภนฺเต.}

\proseitem{177}{กิํ นุ เม ตุเมฺห ภิกฺขเว เอวํ ธมฺมํ เทสิตํ อาชานาถ “ยํ กิญฺจายํ ปุริสปุคฺคโล ปฏิสํเวเทติ สุขํ วา ทุกฺขํ วา อทุกฺขมสุขํ วา, ตสฺส อกุสลา ธมฺมา ปริหายนฺติ, กุสลา ธมฺมา อภิวฑฺฒนฺตี”ติ.\csromanspacer{} โน เหตํ ภนฺเต.\csromanspacer{} นนุ เม ตุเมฺห ภิกฺขเว เอวํ ธมฺมํ เทสิตํ อาชานาถ “อิเธกจฺจสฺส ยํ เอวรูปํ สุขํ เวทนํ เวทยโต อกุสลา ธมฺมา อภิวฑฺฒนฺติ, กุสลา ธมฺมา ปริหายนฺติ.\csromanspacer{} อิธ ปเนกจฺจสฺส เอวรูปํ สุขํ เวทนํ เวทยโต อกุสลา ธมฺมา ปริหายนฺติ, กุสลา ธมฺมา อภิวฑฺฒนฺติ.\csromanspacer{} อิเธกจฺจสฺส เอวรูปํ ทุกฺขํ เวทนํ เวทยโต อกุสลา ธมฺมา อภิวฑฺฒนฺติ, กุสลา ธมฺมา ปริหายนฺติ.\csromanspacer{} อิธ ปเนกจฺจสฺส เอวรูปํ ทุกฺขํ เวทนํ เวทยโต อกุสลา ธมฺมา ปริหายนฺติ, กุสลา ธมฺมา อภิวฑฺฒนฺติ.\csromanspacer{} อิเธกจฺจสฺส เอวรูปํ อทุกฺขมสุขํ เวทนํ เวทยโต อกุสลา ธมฺมา อภิวฑฺฒนฺติ, กุสลา ธมฺมา ปริหายนฺติ.\csromanspacer{} อิธ ปเนกจฺจสฺส เอวรูปํ อทุกฺขมสุขํ เวทนํ เวทยโต อกุสลา ธมฺมา ปริหายนฺติ กุสลา ธมฺมา อภิวฑฺฒนฺตี”ติ.\csromanspacer{} เอวํ ภนฺเต.}

\proseitem{178}{สาธุ ภิกฺขเว, มยา เจตํ ภิกฺขเว อญฺญาตํ อภวิสฺส อทิฏฺฐํ อวิทิตํ อสจฺฉิกตํ อผสฺสิตํ ปญฺญาย “อิเธกจฺจสฺส เอวรูปํ สุขํ เวทนํ เวทยโต อกุสลา ธมฺมา อภิวฑฺฒนฺติ, กุสลา ธมฺมา ปริหายนฺตี”ติ, เอวาหํ อชานนฺโต “เอวรูปํ สุขํ เวทนํ ปชหถา”ติ วเทยฺยํ, อปิ นุ เม เอตํ ภิกฺขเว ปติรูปํ อภวิสฺสาติ.\csromanspacer{} โน เหตํ ภนฺเต.\csromanspacer{} ยสฺมา จ โข เอตํ ภิกฺขเว มยา ญาตํ ทิฏฺฐํ วิทิตํ สจฺฉิกตํ ผสฺสิตํ ปญฺญาย “อิเธกจฺจสฺส เอวรูปํ สุขํ เวทนํ เวทยโต อกุสลา ธมฺมา อภิวฑฺฒนฺติ, กุสลา ธมฺมา ปริหายนฺตี”ติ, ตสฺมาหํ “เอวรูปํ สุขํ เวทนํ ปชหถา”ติ วทามิ.\csromanspacer{} มยา เจตํ ภิกฺขเว อญฺญาตํ อภวิสฺส อทิฏฺฐํ อวิทิตํ}

\vspace{\tipitakaparskip}
\csromancenter{กีฏาคิริสุตฺต (70) อสจฺฉิกตํ อผสฺสิตํ ปญฺญาย “อิเธกจฺจสฺส เอวรูปํ สุขํ เวทนํ เวทยโต อกุสลา ธมฺมา ปริหายนฺติ, กุสลา ธมฺมา อภิวฑฺฒนฺตี”ติ, เอวาหํ อชานนฺโต “เอวรูปํ สุขํ เวทนํ อุปสมฺปชฺช วิหรถา”ติ วเทยฺยํ, อปิ นุ เม เอตํ ภิกฺขเว ปติรูปํ อภวิสฺสาติ.\csromanspacer{} โน เหตํ ภนฺเต, ยสฺมา จ โข เอตํ ภิกฺขเว มยา ญาตํ ทิฏฺฐํ วิทิตํ สจฺฉิกตํ ผสฺสิตํ ปญฺญาย “อิเธกจฺจสฺส เอวรูปํ สุขํ เวทนํ เวทยโต อกุสลา ธมฺมา ปริหายนฺติ, กุสลา ธมฺมา อภิวฑฺฒนฺตี”ติ, ตสฺมาหํ “เอวรูปํ สุขํ เวทนํ อุปสมฺปชฺช วิหรถา”ติ วทามิ.}
\proseitem{179}{\csromanfolio{141}มยา เจตํ ภิกฺขเว อญฺญาตํ อภวิสฺส อทิฏฺฐํ อวิทิตํ อสจฺฉิกตํ อผสฺสิตํ ปญฺญาย “อิเธกจฺจสฺส เอวรูปํ ทุกฺขํ เวทนํ เวทยโต อกุสลา ธมฺมา อภิวฑฺฒนฺติ, กุสลา ธมฺมา ปริหายนฺตี”ติ, เอวาหํ อชานนฺโต “เอวรูปํ ทุกฺขํ เวทนํ ปชหถา”ติ วเทยฺยํ, อปิ นุ เม เอตํ ภิกฺขเว ปติรูปํ อภวิสฺสาติ.\csromanspacer{} โน เหตํ ภนฺเต.\csromanspacer{} ยสฺมา จ โข เอตํ ภิกฺขเว มยา ญาตํ ทิฏฺฐํ วิทิตํ สจฺฉิกตํ ผสฺสิตํ ปญฺญาย “อิเธกจฺจสฺส เอวรูปํ ทุกฺขํ เวทนํ เวทยโต อกุสลา ธมฺมา อภิวฑฺฒนฺติ, กุสลา ธมฺมา ปริหายนฺตี”ติ, ตสฺมาหํ “เอวรูปํ ทุกฺขํ เวทนํ ปชหถา”ติ วทามิ.\csromanspacer{} มยา เจตํ ภิกฺขเว อญฺญาตํ อภวิสฺส อทิฏฺฐํ อวิทิตํ อสจฺฉิกตํ อผสฺสิตํ ปญฺญาย “อิเธกจฺจสฺส เอวรูปํ ทุกฺขํ เวทนํ เวทยโต อกุสลา ธมฺมา ปริหายนฺติ, กุสลา ธมฺมา อภิวฑฺฒนฺตี”ติ, เอวาหํ อชานนฺโต “เอวรูปํ ทุกฺขํ เวทนํ อุปสมฺปชฺช วิหรถา”ติ วเทยฺยํ, อปิ นุ เม เอตํ ภิกฺขเว ปติรูปํ อภวิสฺสาติ.\csromanspacer{} โน เหตํ ภนฺเต.\csromanspacer{} ยสฺมา จ โข เอตํ ภิกฺขเว มยา ญาตํ ทิฏฺฐํ วิทิตํ สจฺฉิกตํ ผสฺสิตํ ปญฺญาย “อิเธกจฺจสฺส เอวรูปํ ทุกฺขํ เวทนํ เวทยโต อกุสลา ธมฺมา ปริหายนฺติ, กุสลา ธมฺมา อภิวฑฺฒนฺตี”ติ, ตสฺมาหํ “เอวรูปํ ทุกฺขํ เวทนํ อุปสมฺปชฺช วิหรถา”ติ วทามิ.}

\proseitem{180}{มยา เจตํ ภิกฺขเว อญฺญาตํ อภวิสฺส อทิฏฺฐํ อวิทิตํ อสจฺฉิกตํ อผสฺสิตํ ปญฺญาย “อิเธกจฺจสฺส เอวรูปํ อทุกฺขมสุขํ เวทนํ เวทยโต อกุสลา ธมฺมา อภิวฑฺฒนฺติ, กุสลา ธมฺมา ปริหายนฺตี”ติ, เอวาหํ อชานนฺโต “เอวรูปํ อทุกฺขมสุขํ เวทนํ ปชหถา”ติ วเทยฺยํ, อปิ นุ เม เอตํ ภิกฺขเว ปติรูปํ อภวิสฺสาติ.\csromanspacer{} โน เหตํ ภนฺเต ยสฺมา จ โข เอตํ ภิกฺขเว มยา ญาตํ ทิฏฺฐํ วิทิตํ สจฺฉิกตํ ผสฺสิตํ \csromanfolio{142}ปญฺญาย “อิเธกจฺจสฺส เอวรูปํ อทุกฺขมสุขํ เวทนํ เวทยโต อกุสลา ธมฺมา อภิวฑฺฒนฺติ, กุสลา ธมฺมา ปริหายนฺตี”ติ, ตสฺมาหํ “เอวรูปํ อทุกฺขมสุขํ เวทนํ ปชหถา”ติ วทามิ.\csromanspacer{} มยา เจตํ ภิกฺขเว อญฺญาตํ อภวิสฺส อทิฏฺฐํ อวิทิตํ อสจฺฉิกตํ อผสฺสิตํ ปญฺญาย “อิเธกจฺจสฺส เอวรูปํ อทุกฺขมสุขํ เวทนํ เวทยโต อกุสลา ธมฺมา ปริหายนฺติ, กุสลา ธมฺมา อภิวฑฺฒนฺตี”ติ, เอวาหํ อชานนฺโต “เอวรูปํ อทุกฺขมสุขํ เวทนํ อุปสมฺปชฺช วิหรถา”ติ วเทยฺยํ, อปิ นุ เม เอตํ ภิกฺขเว ปติรูปํ อภวิสฺสาติ.\csromanspacer{} โน เหตํ ภนฺเต.\csromanspacer{} ยสฺมา จ โข เอตํ ภิกฺขเว มยา ญาตํ ทิฏฺฐํ วิทิตํ สจฺฉิกตํ ผสฺสิตํ ปญฺญาย “อิเธกจฺจสฺส เอวรูปํ อทุกฺขมสุขํ เวทนํ เวทยโต อกุสลา ธมฺมา ปริหายนฺติ, กุสลา ธมฺมา อภิวฑฺฒนฺตี”ติ, ตสฺมาหํ “เอวรูปํ อทุกฺขมสุขํ เวทนํ อุปสมฺปชฺช วิหรถา”ติ วทามิ.}

\proseitem{181}{นาหํ ภิกฺขเว สพฺเพสํเยว ภิกฺขูนํ อปฺปมาเทน กรณียนฺติ วทามิ.\csromanspacer{} น ปนาหํ ภิกฺขเว สพฺเพสํเยว ภิกฺขูนํ น อปฺปมาเทน กรณียนฺติ วทามิ.\csromanspacer{} เย เต ภิกฺขเว ภิกฺขู อรหนฺโต ขีณาสวา วุสิตวนฺโต กตกรณียา โอหิตภารา อนุปฺปตฺตสทตฺถา ปริกฺขีณภวสํโยชนา สมฺมทญฺญา วิมุตฺตา, ตถารูปานาหํ ภิกฺขเว ภิกฺขูนํ น อปฺปมาเทน กรณียนฺติ วทามิ.\csromanspacer{} ตํ กิสฺส เหตุ, กตํ เตสํ อปฺปมาเทน, อภพฺพา เต ปมชฺชิตุํ.\csromanspacer{} เย จ โข เต ภิกฺขเว ภิกฺขู เสกฺขา อปฺปตฺตมานสา อนุตฺตรํ โยคกฺเขมํ ปตฺถยมานา วิหรนฺติ, ตถารูปานาหํ ภิกฺขเว ภิกฺขูนํ อปฺปมาเทน กรณียนฺติ วทามิ.\csromanspacer{} ตํ กิสฺส เหตุ, อปฺเปว นามิเม อายสฺมนฺโต อนุโลมิกานิ เสนาสนานิ ปฏิเสวมานา กลฺยาณมิตฺเต ภชมานา อินฺทฺริยานิ สมนฺนานยมานา, ยสฺสตฺถาย กุลปุตฺตา สมฺมเทว อคารสฺมา อนคาริยํ ปพฺพชนฺติ, ตทนุตฺตรํ พฺรหฺมจริยปริโยสานํ ทิฏฺเฐว ธมฺเม สยํ อภิญฺญา สจฺฉิกตฺวา อุปสมฺปชฺช วิหเรยฺยุนฺติ.\csromanspacer{} อิมํ โข อหํ ภิกฺขเว อิเมสํ ภิกฺขูนํ อปฺปมาทผลํ สมฺปสฺสมาโน อปฺปมาเทน กรณียนฺติ วทามิ.}

\proseitem{182}{สตฺติเม ภิกฺขเว ปุคฺคลา สนฺโต สํวิชฺชมานา โลกสฺมิํ.\csromanspacer{} กตเม สตฺต, อุภโตภาควิมุตฺโต ปญฺญาวิมุตฺโต กายสกฺขิ ทิฏฺฐิปฺปตฺโต สทฺธาวิมุตฺโต ธมฺมานุสารี สทฺธานุสารี.}

\csromanheader{2}{10. กีฏาคิริสุตฺต (70)}
\prose{\csromanfolio{143}กตโม จ ภิกฺขเว ปุคฺคโล อุภโตภาควิมุตฺโต.\csromanspacer{} อิธ ภิกฺขเว เอกจฺโจ ปุคฺคโล เย เต สนฺตา วิโมกฺขา อติกฺกมฺม รูเป อารุปฺปา, เต กาเยน ผุสิตฺวา\footnote{ผสฺสิตฺวา (สี, อิ)} วิหรติ, ปญฺญาย จสฺส ทิสฺวา อาสวา ปริกฺขีณา โหนฺติ.\csromanspacer{} อยํ วุจฺจติ ภิกฺขเว ปุคฺคโล อุภโตภาควิมุตฺโต.\csromanspacer{} อิมสฺส โข อหํ ภิกฺขเว ภิกฺขุโน น อปฺปมาเทน กรณียนฺติ วทามิ.\csromanspacer{} ตํ กิสฺส เหตุ, กตํ ตสฺส อปฺปมาเทน, อภพฺโพ โส ปมชฺชิตุํ. (1)}

\prose{กตโม จ ภิกฺขเว ปุคฺคโล ปญฺญาวิมุตฺโต.\csromanspacer{} อิธ ภิกฺขเว เอกจฺโจ ปุคฺคโล เย เต สนฺตา วิโมกฺขา อติกฺกมฺม รูเป อารุปฺปา, เต น กาเยน ผุสิตฺวา วิหรติ, ปญฺญาย จสฺส ทิสฺวา อาสวา ปริกฺขีณา โหนฺติ.\csromanspacer{} อยํ วุจฺจติ ภิกฺขเว ปุคฺคโล ปญฺญาวิมุตฺโต.\csromanspacer{} อิมสฺสปิ โข อหํ ภิกฺขเว ภิกฺขุโน น อปฺปมาเทน กรณียนฺติ วทามิ.\csromanspacer{} ตํ กิสฺส เหตุ, กตํ ตสฺส อปฺปมาเทน, อภพฺโพ โส ปมชฺชิตุํ. (2)}

\prose{กตโม จ ภิกฺขเว ปุคฺคโล กายสกฺขิ.\csromanspacer{} อิธ ภิกฺขเว เอกจฺโจ ปุคฺคโล เย เต สนฺตา วิโมกฺขา อติกฺกมฺม รูเป อารุปฺปา, เต กาเยน ผุสิตฺวา วิหรติ, ปญฺญาย จสฺส ทิสฺวา เอกจฺเจ อาสวา ปริกฺขีณา โหนฺติ.\csromanspacer{} อยํ วุจฺจติ ภิกฺขเว ปุคฺคโล กายสกฺขิ.\csromanspacer{} อิมสฺส โข อหํ ภิกฺขเว ภิกฺขุโน อปฺปมาเทน กรณียนฺติ วทามิ.\csromanspacer{} ตํ กิสฺส เหตุ, อปฺเปว นาม อยมายสฺมา อนุโลมิกานิ เสนาสนานิ ปฏิเสวมาโน กลฺยาณมิตฺเต ภชมาโน อินฺทฺริยานิ สมนฺนานยมาโน, ยสฺสตฺถาย กุลปุตฺตา สมฺมเทว อคารสฺมา อนคาริยํ ปพฺพชนฺติ, ตทนุตฺตรํ พฺรหฺมจริยปริโยสานํ ทิฏฺเฐว ธมฺเม สยํ อภิญฺญา สจฺฉิกตฺวา อุปสมฺปชฺช วิหเรยฺยาติ.\csromanspacer{} อิมํ โข อหํ ภิกฺขเว อิมสฺส ภิกฺขุโน อปฺปมาทผลํ สมฺปสฺสมาโน อปฺปมาเทน กรณียนฺติ วทามิ. (3)}

\prose{กตโม จ ภิกฺขเว ปุคฺคโล ทิฏฺฐิปฺปตฺโต.\csromanspacer{} อิธ ภิกฺขเว เอกจฺโจ ปุคฺคโล เย เต สนฺตา วิโมกฺขา อติกฺกมฺม รูเป อารุปฺปา, เต น กาเยน ผุสิตฺวา วิหรติ, ปญฺญาย จสฺส ทิสฺวา เอกจฺเจ อาสวา ปริกฺขีณา โหนฺติ, ตถาคตปฺปเวทิตา จสฺส ธมฺมา ปญฺญาย โวทิฏฺฐา โหนฺติ โวจริตา.\csromanspacer{} อยํ วุจฺจติ ภิกฺขเว ปุคฺคโล ทิฏฺฐิปฺปตฺโต.\csromanspacer{} อิมสฺสปิ โข อหํ ภิกฺขเว ภิกฺขุโน อปฺปมาเทน กรณียนฺติ วทามิ.\csromanspacer{} ตํ}

\prose{\csromanfolio{144}กิสฺส เหตุ, อปฺเปว นาม อยมายสฺมา อนุโลมิกานิ เสนาสนานิ ปฏิเสวมาโน กลฺยาณมิตฺเต ภชมาโน อินฺทฺริยานิ สมนฺนานยมาโน, ยสฺสตฺถาย กุลปุตฺตา สมฺมเทว อคารสฺมา อนคาริยํ ปพฺพชนฺติ, ตทนุตฺตรํ พฺรหฺมจริยปริโยสานํ ทิฏฺเฐว ธมฺเม สยํ อภิญฺญา สจฺฉิกตฺวา อุปสมฺปชฺช วิหเรยฺยาติ.\csromanspacer{} อิมํ โข อหํ ภิกฺขเว อิมสฺส ภิกฺขุโน อปฺปมาทผลํ สมฺปสฺสมาโน อปฺปมาเทน กรณียนฺติ วทามิ. (4)}

\prose{กตโม จ ภิกฺขเว ปุคฺคโล สทฺธาวิมุตฺโต.\csromanspacer{} อิธ ภิกฺขเว เอกจฺโจ ปุคฺคโล เย เต สนฺตา วิโมกฺขา อติกฺกมฺม รูเป อารุปฺปา, เต น กาเยน ผุสิตฺวา วิหรติ, ปญฺญาย จสฺส ทิสฺวา เอกจฺเจ อาสวา ปริกฺขีณา โหนฺติ, ตถาคเต จสฺส สทฺธา นิวิฏฺฐา โหติ มูลชาตา ปติฏฺฐิตา.\csromanspacer{} อยํ วุจฺจติ ภิกฺขเว ปุคฺคโล สทฺธาวิมุตฺโต.\csromanspacer{} อิมสฺสปิ โข อหํ ภิกฺขเว ภิกฺขุโน อปฺปมาเทน กรณียนฺติ วทามิ.\csromanspacer{} ตํ กิสฺส เหตุ, อปฺเปว นาม อยมายสฺมา อนุโลมิกานิ เสนาสนานิ ปฏิเสวมาโน กลฺยาณมิตฺเต ภชมาโน อินฺทฺริยานิ สมนฺนานยมาโน, ยสฺสตฺถาย กุลปุตฺตา สมฺมเทว อคารสฺมา อนคาริยํ ปพฺพชนฺติ, ตทนุตฺตรํ พฺรหฺมจริยปริโยสานํ ทิฏฺเฐว ธมฺเม สยํ อภิญฺญา สจฺฉิกตฺวา อุปสมฺปชฺช วิหเรยฺยาติ.\csromanspacer{} อิมํ โข อหํ ภิกฺขเว อิมสฺส ภิกฺขุโน อปฺปมาทผลํ สมฺปสฺสมาโน อปฺปมาเทน กรณียนฺติ วทามิ. (5)}

\prose{กตโม จ ภิกฺขเว ปุคฺคโล ธมฺมานุสารี.\csromanspacer{} อิธ ภิกฺขเว เอกจฺโจ ปุคฺคโล เย เต สนฺตา วิโมกฺขา อติกฺกมฺม รูเป อารุปฺปา, เต น กาเยน ผุสิตฺวา วิหรติ, ปญฺญาย จสฺส ทิสฺวา เอกจฺเจ อาสวา ปริกฺขีณา\footnote{ทิสฺวา อาสวา อปริกฺขีณา (สี, อิ)} โหนฺติ, ตถาคตปฺปเวทิตา จสฺส ธมฺมา ปญฺญาย มตฺตโส นิชฺฌานํ ขมนฺติ, อปิ จสฺส อิเม ธมฺมา โหนฺติ, เสยฺยถิทํ, สทฺธินฺทฺริยํ วีริยินฺทฺริยํ สตินฺทฺริยํ สมาธินฺทฺริยํ ปญฺญินฺทฺริยํ.\csromanspacer{} อยํ วุจฺจติ ภิกฺขเว ปุคฺคโล ธมฺมานุสารี.\csromanspacer{} อิมสฺสปิ โข อหํ ภิกฺขเว ภิกฺขุโน อปฺปมาเทน กรณียนฺติ วทามิ.\csromanspacer{} ตํ กิสฺส เหตุ, อปฺเปว นาม อยมายสฺมา อนุโลมิกานิ เสนาสนานิ ปฏิเสวมาโน กลฺยาณมิตฺเต ภชมาโน อินฺทฺริยานิ สมนฺนานยมาโน, ยสฺสตฺถาย กุลปุตฺตา สมฺมเทว อคารสฺมา อนคาริยํ ปพฺพชนฺติ, ตทนุตฺตรํ พฺรหฺมจริยปริโยสานํ ทิฏฺเฐว ธมฺเม สยํ อภิญฺญา สจฺฉิกตฺวา อุปสมฺปชฺช}

\vspace{\tipitakaparskip}
\csromancenter{กีฏาคิริสุตฺต (70) วิหเรยฺยาติ.\csromanspacer{} อิมํ โข อหํ ภิกฺขเว อิมสฺส ภิกฺขุโน อปฺปมาทผลํ สมฺปสฺสมาโน อปฺปมาเทน กรณียนฺติ วทามิ. (6)}
\prose{\csromanfolio{145}กตโม จ ภิกฺขเว ปุคฺคโล สทฺธานุสารี.\csromanspacer{} อิธ ภิกฺขเว เอกจฺโจ ปุคฺคโล เย เต สนฺตา วิโมกฺขา อติกฺกมฺม รูเป อารุปฺปา, เต น กาเยน ผุสิตฺวา วิหรติ, ปญฺญาย จสฺส ทิสฺวา เอกจฺเจ อาสวา ปริกฺขีณา\footnote{ทิสฺวา อาสวา อปริกฺขีณา (สี, อิ)} โหนฺติ, ตถาคเต จสฺส สทฺธามตฺตํ โหติ เปมมตฺตํ, อปิ จสฺส อิเม ธมฺมา โหนฺติ.\csromanspacer{} เสยฺยถิทํ, สทฺธินฺทฺริยํ วีริยินฺทฺริยํ สตินฺทฺริยํ สมาธินฺทฺริยํ ปญฺญินฺทฺริยํ.\csromanspacer{} อยํ วุจฺจติ ภิกฺขเว ปุคฺคโล สทฺธานุสารี.\csromanspacer{} อิมสฺสปิ โข อหํ ภิกฺขเว ภิกฺขุโน อปฺปมาเทน กรณียนฺติ วทามิ.\csromanspacer{} ตํ กิสฺส เหตุ, อปฺเปว นาม อยมายสฺมา อนุโลมิกานิ เสนาสนานิ ปฏิเสวมาโน กลฺยาณมิตฺเต ภชมาโน อินฺทฺริยานิ สมนฺนานยมาโน, ยสฺสตฺถาย กุลปุตฺตา สมฺมเทว อคารสฺมา อนคาริยํ ปพฺพชนฺติ, ตทนุตฺตรํ พฺรหฺมจริยปริโยสานํ ทิฏฺเฐว ธมฺเม สยํ อภิญฺญา สจฺฉิกตฺวา อุปสมฺปชฺช วิหเรยฺยาติ.\csromanspacer{} อิมํ โข อหํ ภิกฺขเว อิมสฺส ภิกฺขุโน อปฺปมาทผลํ สมฺปสฺสมาโน อปฺปมาเทน กรณียนฺติ วทามิ. (7)}

\proseitem{183}{นาหํ ภิกฺขเว อาทิเกเนว อญฺญาราธนํ วทามิ.\csromanspacer{} อปิ จ ภิกฺขเว อนุปุพฺพสิกฺขา อนุปุพฺพกิริยา อนุปุพฺพปฏิปทา อญฺญาราธนา โหติ.\csromanspacer{} กถญฺจ ภิกฺขเว อนุปุพฺพสิกฺขา อนุปุพฺพกิริยา อนุปุพฺพปฏิปทา อญฺญาราธนา โหติ.\csromanspacer{} อิธ ภิกฺขเว สทฺธาชาโต อุปสงฺกมติ, อุปสงฺกมนฺโต ปยิรุปาสติ, ปยิรุปาสนฺโต โสตํ โอทหติ, โอหิตโสโต ธมฺมํ สุณาติ, สุตฺวา ธมฺมํ ธาเรติ, ธตานํ\footnote{ธาตานํ (ก)} ธมฺมานํ อตฺถํ อุปปริกฺขติ, อตฺถํ อุปปริกฺขโต ธมฺมา นิชฺฌานํ ขมนฺติ, ธมฺมนิชฺฌานกฺขนฺติยา สติ ฉนฺโท ชายติ, ฉนฺทชาโต อุสฺสหติ, อุสฺสาเหตฺวา ตุเลติ, ตุลยิตฺวา ปทหติ, ปหิตตฺโต สมาโน กาเยน เจว ปรมสจฺจํ สจฺฉิกโรติ, ปญฺญาย จ นํ อติวิชฺฌ ปสฺสติ.\csromanspacer{} สาปิ นาม ภิกฺขเว สทฺธา นาโหสิ, ตมฺปิ นาม ภิกฺขเว อุปสงฺกมนํ นาโหสิ, สาปิ นาม ภิกฺขเว ปยิรุปาสนา นาโหสิ, ตมฺปิ นาม ภิกฺขเว โสตาวธานํ \csromanfolio{146}นาโหสิ, ตมฺปิ นาม ภิกฺขเว ธมฺมสฺสวนํ นาโหสิ, สาปิ นาม ภิกฺขเว ธมฺมธารณา นาโหสิ, สาปิ นาม ภิกฺขเว อตฺถูปปริกฺขา นาโหสิ, สาปิ นาม ภิกฺขเว ธมฺมนิชฺฌานกฺขนฺติ นาโหสิ, โสปิ นาม ภิกฺขเว ฉนฺโท นาโหสิ, โสปิ นาม ภิกฺขเว อุสฺสาโห นาโหสิ, สาปิ นาม ภิกฺขเว ตุลนา นาโหสิ, ตมฺปิ นาม ภิกฺขเว ปธานํ นาโหสิ.\csromanspacer{} วิปฺปฏิปนฺนาตฺถ ภิกฺขเว, มิจฺฉาปฏิปนฺนาตฺถ ภิกฺขเว, กีว ทูเรวิเม ภิกฺขเว โมฆปุริสา อปกฺกนฺตา อิมมฺหา ธมฺมวินยา.}

\proseitem{184}{อตฺถิ ภิกฺขเว จตุปฺปทํ เวยฺยากรณํ.\csromanspacer{} ยสฺสุทฺทิฏฺฐสฺส วิญฺญู ปุริโส นจิรสฺเสว ปญฺญายตฺถํ อาชาเนยฺย, อุทฺทิสิสฺสามิ โว\footnote{อุทฺทิฏฺฐสฺสาปิ (ก)} ภิกฺขเว, อาชานิสฺสถ เม ตนฺติ.\csromanspacer{} เก จ มยํ ภนฺเต, เก จ ธมฺมสฺส อญฺญาตาโรติ.\csromanspacer{} โยปิ โส ภิกฺขเว สตฺถา อามิสครุ อามิสทายาโท อามิเสหิ สํสฏฺโฐ วิหรติ, ตสฺส ปายํ เอวรูปี ปโณปณวิยา น อุเปติ “เอวญฺจ โน อสฺส, อถ นํ กเรยฺยาม.\csromanspacer{} น จ โน เอวมสฺส, น นํ กเรยฺยามา”ติ.\csromanspacer{} กิํ ปน ภิกฺขเว ยํ ตถาคโต สพฺพโส อามิเสหิ วิสํสฏฺโฐ วิหรติ.\csromanspacer{} สทฺธสฺส ภิกฺขเว สาวกสฺส สตฺถุสาสเน ปริโยคาหิย\footnote{ปริโยคาย (สี, อิ, ก), ปริโยคยฺห (สฺยา, กํ)} วตฺตโต อยมนุธมฺโม โหติ “สตฺถา ภควา, สาวโกหมสฺมิ, ชานาติ ภควา, นาหํ ชานามี”ติ.\csromanspacer{} สทฺธสฺส ภิกฺขเว สาวกสฺส สตฺถุสาสเน ปริโยคาหิย วตฺตโต รุฬฺหนียํ\footnote{รุมฺหนิยํ (สี, อิ)} สตฺถุสาสนํ โหติ โอชวนฺตํ.\csromanspacer{} สทฺธสฺส ภิกฺขเว สาวกสฺส สตฺถุสาสเน ปริโยคาหิย วตฺตโต อยมนุธมฺโม โหติ “กามํ ตโจ จ นฺหารุ จ อฏฺฐิ จ อวสิสฺสตุ, สรีเร อุปสุสฺสตุ\footnote{อุปสุสฺสตุ สรีเร (สี), สรีเร อวสุสฺสตุ (ก)} มํสโลหิตํ, ยํ ตํ ปุริสถาเมน ปุริสวีริเยน ปุริสปรกฺกเมน ปตฺตพฺพํ, น ตํ อปาปุณิตฺวา วีริยสฺส สณฺฐานํ\footnote{สนฺถานํ (สี, สฺยา, อิ)} ภวิสฺสตี”ติ.\csromanspacer{} สทฺธสฺส ภิกฺขเว สาวกสฺส สตฺถุสาสเน ปริโยคาหิย วตฺตโต ทฺวินฺนํ ผลานํ อญฺญตรํ ผลํ ปาฏิกงฺขํ “ทิฏฺเฐว ธมฺเม อญฺญา, สติ วา อุปาทิเสเส อนาคามิตา”ติ}

\csromanheader{2}{10. กีฏาคิริสุตฺต (70)}
\vspace{\tipitakaparskip}
\csromancenter{อิทมโวจ ภควา.\csromanspacer{} อตฺตมนา เต ภิกฺขู ภควโต ภาสิตํ อภินนฺทุนฺติ.}
\nitthitam{กีฏาคิริสุตฺตํ นิฏฺฐิตํ ทสมํ.}
\nitthitam{\textbf{ภิกฺขุวคฺโค นิฏฺฐิโต ทุติโย.}}
\titlehead{\textbf{ตสฺสุทฺทานํ}}
\csromanmatikaanchor{mtk.24}
\csromantocmarkref{subsection}{อุทฺทานคาถาโย}{mtk.24}
\bookmark[dest={mtk.24},level=2]{อุทฺทานคาถาโย}
\csromangathameasure{กุญฺชร ราหุล สสฺสตโลโก, \\ มาลุกฺยปุตฺโต จ ภทฺทาลิ นาโม. \\ ขุทฺท ทิชาถ สหมฺปติยาจํ. \\ นาฬก รญฺญิกิฏาคิรินาโม.}
\csromangathagroup{\csromangathastanza{\csromanfolio{147}กุญฺชร ราหุล สสฺสตโลโก, \\ มาลุกฺยปุตฺโต จ ภทฺทาลิ นาโม. \\ ขุทฺท ทิชาถ สหมฺปติยาจํ. \\ นาฬก รญฺญิกิฏาคิรินาโม.}}

\csromanmatikaanchor{mtk.25}
\csromantocmarknopage{chapter}{3. ปริพฺพาชกวคฺค}{mtk.25}
\bookmark[dest={mtk.25},level=0]{3. ปริพฺพาชกวคฺค}
\chapterheadpageread{3. ปริพฺพาชกวคฺค}
\csromanmatikaanchor{mtk.26}
\csromantocmarkref{subsection}{1. เตวิชฺชวจฺฉสุตฺต}{mtk.26}
\bookmark[dest={mtk.26},level=2]{1. เตวิชฺชวจฺฉสุตฺต}
\csromanheader{2}{1. เตวิชฺชวจฺฉสุตฺต}
\proseitem{185}{\csromanfolio{148}เอวํ เม สุตํ\csromandash{}เอกํ สมยํ ภควา เวสาลิยํ วิหรติ มหาวเน กูฏาคารสาลายํ.\csromanspacer{} เตน โข ปน สมเยน วจฺฉโคตฺโต ปริพฺพาชโก เอกปุณฺฑรีเก ปริพฺพาชการาเม ปฏิวสติ.\csromanspacer{} อถ โข ภควา ปุพฺพณฺหสมยํ นิวาเสตฺวา ปตฺตจีวรมาทาย เวสาลิํ ปิณฺฑาย ปาวิสิ.\csromanspacer{} อถ โข ภควโต เอตทโหสิ “อติปฺปโค โข ตาว เวสาลิยํ ปิณฺฑาย จริตุํ, ยํนูนาหํ เยน เอกปุณฺฑรีโก ปริพฺพาชการาโม, เยน วจฺฉโคตฺโต ปริพฺพาชโก เตนุปสงฺกเมยฺยนฺ”ติ.\csromanspacer{} อถ โข ภควา เยน เอกปุณฺฑรีโก ปริพฺพาชการาโม, เยน วจฺฉโคตฺโต ปริพฺพาชโก เตนุปสงฺกมิ, อทฺทสา โข วจฺฉโคตฺโต ปริพฺพาชโก ภควนฺตํ ทูรโตว อาคจฺฉนฺตํ, ทิสฺวาน ภควนฺตํ เอตทโวจ “เอตุ โข ภนฺเต ภควา, สฺวาคตํ\footnote{สาคตํ (สี, อิ)} ภนฺเต ภควโต, จิรสฺสํ โข ภนฺเต ภควา อิมํ ปริยายมกาสิ ยทิทํ อิธาคมนาย, นิสีทตุ ภนฺเต ภควา อิทมาสนํ ปญฺญตฺตนฺ”ติ.\csromanspacer{} นิสีทิ ภควา ปญฺญตฺเต อาสเน.\csromanspacer{} วจฺฉโคตฺโตปิ โข ปริพฺพาชโก อญฺญตรํ นีจํ อาสนํ คเหตฺวา เอกมนฺตํ นิสีทิ, เอกมนฺตํ นิสินฺโน โข วจฺฉโคตฺโต ปริพฺพาชโก ภควนฺตํ เอตทโวจ “สุตํ เมตํ ภนฺเต สมโณ โคตโม สพฺพญฺญู สพฺพทสฺสาวี อปริเสสํ ญาณทสฺสนํ ปฏิชานาติ, จรโต จ เม ติฏฺฐโต จ สุตฺตสฺส จ ชาครสฺส จ สตตํ สมิตํ ญาณทสฺสนํ ปจฺจุปฏฺฐิตนฺ”ติ.\csromanspacer{} เย เต ภนฺเต เอวมาหํสุ “สมโณ โคตโม สพฺพญฺญู สพฺพทสฺสาวี อปริเสสํ ญาณทสฺสนํ ปฏิชานาติ, จรโต จ เม ติฏฺฐโต จ สุตฺตสฺส จ ชาครสฺส จ สตตํ สมิตํ ญาณทสฺสนํ ปจฺจุปฏฺฐิตนฺ”ติ.\csromanspacer{} กจฺจิ เต ภนฺเต ภควโต วุตฺตวาทิโน, น จ ภควนฺตํ อภูเตน อพฺภาจิกฺขนฺติ, ธมฺมสฺส จานุธมฺมํ พฺยากโรนฺติ, น จ โกจิ สหธมฺมิโก วาทานุวาโท คารยฺหํ ฐานํ อาคจฺฉตีติ.\csromanspacer{} เย เต วจฺฉ เอวมาหํสุ “สมโณ โคตโม สพฺพญฺญู สพฺพทสฺสาวี อปริเสสํ ญาณทสฺสนํ ปฏิชานาติ, จรโต จ เม ติฏฺฐโต จ สุตฺตสฺส จ ชาครสฺส จ สตตํ สมิตํ ญาณทสฺสนํ ปจฺจุปฏฺฐิตนฺ”ติ.\csromanspacer{} น เม เต วุตฺตวาทิโน, อพฺภาจิกฺขนฺติ จ ปน มํ อสตา อภูเตนาติ.}

\csromanheader{2}{1. เตวิชฺชวจฺฉสุตฺต (71)}
\proseitem{186}{\csromanfolio{149}กถํ พฺยากรมานา ปน มยํ ภนฺเต วุตฺตวาทิโน เจว ภควโต อสฺสาม, น จ ภควนฺตํ อภูเตน อพฺภาจิกฺเขยฺยาม, ธมฺมสฺส จานุธมฺมํ พฺยากเรยฺยาม, น จ โกจิ สหธมฺมิโก วาทานุวาโท คารยฺหํ ฐานํ อาคจฺเฉยฺยาติ.}

\prose{“เตวิชฺโช สมโณ โคตโม”ติ โข วจฺฉ พฺยากรมาโน วุตฺตวาที เจว เม อสฺส, น จ มํ อภูเตน อพฺภาจิกฺเขยฺย, ธมฺมสฺส จานุธมฺมํ พฺยากเรยฺย, น จ โกจิ สหธมฺมิโก วาทานุวาโท คารยฺหํ ฐานํ อาคจฺเฉยฺย.\csromanspacer{} อหํ หิ วจฺฉ ยาวเทว อากงฺขามิ, อเนกวิหิตํ ปุพฺเพนิวาสํ อนุสฺสรามิ.\csromanspacer{} เสยฺยถิทํ, เอกมฺปิ ชาติํ เทฺวปิ ชาติโย ฯเปฯ อิติ สาการํ ส-อุทฺเทสํ อเนกวิหิตํ ปุพฺเพนิวาสํ อนุสฺสรามิ.\csromanspacer{} อหํ หิ วจฺฉ ยาวเทว อากงฺขามิ, ทิพฺเพน จกฺขุนา วิสุทฺเธน อติกฺกนฺตมานุสเกน สตฺเต ปสฺสามิ จวมาเน อุปปชฺชมาเน หีเน ปณีเต สุวณฺเณ ทุพฺพณฺเณ สุคเต ทุคฺคเต ฯเปฯ ยถากมฺมูปเค สตฺเต ปชานามิ.\csromanspacer{} อหํ หิ วจฺฉ อาสวานํ ขยา อนาสวํ เจโตวิมุตฺติํ ปญฺญาวิมุตฺติํ ทิฏฺเฐว ธมฺเม สยํ อภิญฺญา สจฺฉิกตฺวา อุปสมฺปชฺช วิหรามิ.}

\prose{เตวิชฺโช สมโณ โคตโมติ โข วจฺฉ พฺยากรมาโน วุตฺตวาที เจว เม อสฺส, น จ มํ อภูเตน อพฺภาจิกฺเขยฺย, ธมฺมสฺส จานุธมฺมํ พฺยากเรยฺย, น จ โกจิ สหธมฺมิโก วาทานุวาโท คารยฺหํ ฐานํ อาคจฺเฉยฺยาติ.}

\prose{เอวํ วุตฺเต วจฺฉโคตฺโต ปริพฺพาชโก ภควนฺตํ เอตทโวจ “อตฺถิ นุ โข โภ โคตม โกจิ คิหี คิหิสํโยชนํ อปฺปหาย กายสฺส เภทา ทุกฺขสฺสนฺตกโร”ติ.\csromanspacer{} นตฺถิ โข วจฺฉ โกจิ คิหี คิหิสํโยชนํ อปฺปหาย กายสฺส เภทา ทุกฺขสฺสนฺตกโรติ.}

\prose{อตฺถิ ปน โภ โคตม โกจิ คิหี คิหิสํโยชนํ อปฺปหาย กายสฺส เภทา สคฺคูปโคติ. “น โข วจฺฉ เอกํเยว สตํ น เทฺว สตานิ น ตีณิ สตานิ น จตฺตาริ สตานิ น ปญฺจ สตานิ, อถ โข ภิยฺโยว เย คิหี คิหิสํโยชนํ อปฺปหาย กายสฺส เภทา สคฺคูปคาติ”\footnote{“อตฺถิ โข วจฺฉ โกจิ คิหี คิหิสํโยชนํ อปฺปหาย กายสฺส เภทา สคฺคูปโคติ”. (ก)}}

\prose{\csromanfolio{150}อตฺถิ นุ โข โภ โคตม โกจิ อาชีวโก\footnote{อาชีวิโก (ก)} กายสฺส เภทา ทุกฺขสฺสนฺตกโรติ.\csromanspacer{} นตฺถิ โข วจฺฉ โกจิ อาชีวโก กายสฺส เภทา ทุกฺขสฺสนฺตกโรติ.}

\prose{อตฺถิ ปน โภ โคตม โกจิ อาชีวโก กายสฺส เภทา สคฺคูปโคติ.\csromanspacer{} อิโต โข โส วจฺฉ เอกนวุโต กปฺโป\footnote{อิโต โข วจฺฉ เอกนวุเต กปฺเป (ก)} ยมหํ อนุสฺสรามิ, นาภิชานามิ กญฺจิ อาชีวกํ สคฺคูปคํ อญฺญตฺร เอเกน, โสปาสิ กมฺมวาที กิริยวาทีติ.\csromanspacer{} เอวํ สนฺเต โภ โคตม สุญฺญํ อทุํ ติตฺถายตนํ อนฺตมโส สคฺคูปเคนปีติ.\csromanspacer{} เอวํ วจฺฉ สุญฺญํ อทุํ ติตฺถายตนํ อนฺตมโส สคฺคูปเคนปีติ.}

\prose{อิทมโวจ ภควา.\csromanspacer{} อตฺตมโน วจฺฉโคตฺโต ปริพฺพาชโก ภควโต ภาสิตํ อภินนฺทีติ.}

\nitthitamruled{เตวิชฺชวจฺฉสุตฺตํ นิฏฺฐิตํ ปฐมํ.}
\csromanmatikaanchor{mtk.27}
\csromantocmarkref{subsection}{2. อคฺคิวจฺฉสุตฺต}{mtk.27}
\bookmark[dest={mtk.27},level=2]{2. อคฺคิวจฺฉสุตฺต}
\csromanheader{2}{2. อคฺคิวจฺฉสุตฺต}
\proseitem{187}{เอวํ เม สุตํ\csromandash{}เอกํ สมยํ ภควา สาวตฺถิยํ วิหรติ เชตวเน อนาถปิณฺฑิกสฺส อาราเม.\csromanspacer{} อถ โข วจฺฉโคตฺโต ปริพฺพาชโก เยน ภควา เตนุปสงฺกมิ, อุปสงฺกมิตฺวา ภควตา สทฺธิํ สมฺโมทิ, สมฺโมทนียํ กถํ สารณียํ วีติสาเรตฺวา เอกมนฺตํ นิสีทิ, เอกมนฺตํ นิสินฺโน โข วจฺฉโคตฺโต ปริพฺพาชโก ภควนฺตํ เอตทโวจ “กิํ นุ โข โภ โคตม ‘สสฺสโต โลโก อิทเมว สจฺจํ โมฆมญฺญนฺ’ติ เอวํทิฏฺฐิ\footnote{เอวํทิฏฺฐี (สี, สฺยา, กํ, อิ)} ภวํ โคตโม”ติ.\csromanspacer{} น โข อหํ วจฺฉ เอวํทิฏฺฐิ “สสฺสโต โลโก อิทเมว สจฺจํ โมฆมญฺญนฺ”ติ.}

\prose{กิํ ปน โภ โคตม “อสสฺสโต โลโก อิทเมว สจฺจํ โมฆมญฺญนฺ”ติ เอวํทิฏฺฐิ ภวํ โคตโมติ.\csromanspacer{} น โข อหํ วจฺฉ เอวํทิฏฺฐิ “อสสฺสโต โลโก อิทเมว สจฺจํ โมฆมญฺญนฺ”ติ.}

\csromanheader{2}{2. อคฺคิวจฺฉสุตฺต (72)}
\prose{\csromanfolio{151}กิํ นุ โข โภ โคตม “อนฺตวา โลโก อิทเมว สจฺจํ โมฆมญฺญนฺ”ติ เอวํทิฏฺฐิ ภวํ โคตโมติ.\csromanspacer{} น โข อหํ วจฺฉ เอวํทิฏฺฐิ “อนฺตวา โลโก อิทเมว สจฺจํ โมฆมญฺญนฺ”ติ.}

\prose{กิํ ปน โภ โคตม “อนนฺตวา โลโก อิทเมว สจฺจํ โมฆมญฺญนฺ”ติ เอวํทิฏฺฐิ ภวํ โคตโมติ.\csromanspacer{} น โข อหํ วจฺฉ เอวํทิฏฺฐิ “อนนฺตวา โลโก อิทเมว สจฺจํ โมฆมญฺญนฺ”ติ.}

\prose{กิํ นุ โข โภ โคตม “ตํ ชีวํ ตํ สรีรํ อิทเมว สจฺจํ โมฆมญฺญนฺ”ติ เอวํทิฏฺฐิ ภวํ โคตโมติ.\csromanspacer{} น โข อหํ วจฺฉ เอวํทิฏฺฐิ “ตํ ชีวํ ตํ สรีรํ อิทเมว สจฺจํ โมฆมญฺญนฺ”ติ.}

\prose{กิํ ปน โภ โคตม “อญฺญํ ชีวํ อญฺญํ สรีรํ อิทเมว สจฺจํ โมฆมญฺญนฺ”ติ เอวํทิฏฺฐิ ภวํ โคตโมติ.\csromanspacer{} น โข อหํ วจฺฉ เอวํทิฏฺฐิ “อญฺญํ ชีวํ อญฺญํ สรีรํ อิทเมว สจฺจํ โมฆมญฺญนฺ”ติ.}

\prose{กิํ นุ โข โภ โคตม “โหติ ตถาคโต ปรํ มรณา อิทเมว สจฺจํ โมฆมญฺญนฺ”ติ เอวํทิฏฺฐิ ภวํ โคตโมติ.\csromanspacer{} น โข อหํ วจฺฉ เอวํทิฏฺฐิ “โหติ ตถาคโต ปรํ มรณา อิทเมว สจฺจํ โมฆมญฺญนฺ”ติ.}

\prose{กิํ ปน โภ โคตม “น โหติ ตถาคโต ปรํ มรณา อิทเมว สจฺจํ โมฆมญฺญนฺ”ติ เอวํทิฏฺฐิ ภวํ โคตโมติ.\csromanspacer{} น โข อหํ วจฺฉ เอวํทิฏฺฐิ “น โหติ ตถาคโต ปรํ มรณา อิทเมว สจฺจํ โมฆมญฺญนฺ”ติ.}

\prose{กิํ นุ โข โภ โคตม “โหติ จ น จ โหติ ตถาคโต ปรํ มรณา อิทเมว สจฺจํ โมฆมญฺญนฺ”ติ เอวํทิฏฺฐิ ภวํ โคตโมติ.\csromanspacer{} น โข อหํ วจฺฉ เอวํทิฏฺฐิ “โหติ จ น จ โหติ ตถาคโต ปรํ มรณา อิทเมว สจฺจํ โมฆมญฺญนฺ”ติ.}

\prose{กิํ ปน โภ โคตม “เนว โหติ น น โหติ ตถาคโต ปรํ มรณา อิทเมว สจฺจํ โมฆมญฺญนฺ”ติ เอวํทิฏฺฐิ ภวํ โคตโมติ.\csromanspacer{} น โข อหํ วจฺฉ เอวํทิฏฺฐิ “เนว โหติ น น โหติ ตถาคโต ปรํ มรณา อิทเมว สจฺจํ โมฆมญฺญนฺ”ติ.}

\proseitem{188}{กิํ นุ โข โภ โคตม “สสฺสโต โลโก อิทเมว สจฺจํ โมฆมญฺญนฺ”ติ เอวํทิฏฺฐิ ภวํ โคตโมติ อิติ ปุฏฺโฐ สมาโน \csromanfolio{152}“น โข อหํ วจฺฉ เอวํทิฏฺฐิ ‘สสฺสโต โลโก อิทเมว สจฺจํ โมฆมญฺญนฺ’ติ” วเทสิ\footnote{โมฆมญฺญนฺตีติ วเทสิ (สี), โมฆมญฺญนฺติ อิติ วเทสิ (?)}.\csromanspacer{} กิํ ปน โภ โคตม “อสสฺสโต โลโก อิทเมว สจฺจํ โมฆมญฺญนฺ”ติ เอวํทิฏฺฐิ ภวํ โคตโมติ อิติ ปุฏฺโฐ สมาโน “น โข อหํ วจฺฉ เอวํทิฏฺฐิ ‘อสสฺสโต โลโก อิทเมว สจฺจํ โมฆมญฺญนฺ’ติ” วเทสิ.\csromanspacer{} กิํ นุ โข โภ โคตม “อนฺตวา โลโก อิทเมว สจฺจํ โมฆมญฺญนฺ”ติ เอวํทิฏฺฐิ ภวํ โคตโมติ อิติ ปุฏฺโฐ สมาโน “น โข อหํ วจฺฉ เอวํทิฏฺฐิ ‘อนฺตวา โลโก อิทเมว สจฺจํ โมฆมญฺญนฺ’ติ” วเทสิ.\csromanspacer{} กิํ ปน โภ โคตม “อนนฺตวา โลโก อิทเมว สจฺจํ โมฆมญฺญนฺ”ติ เอวํทิฏฺฐิ ภวํ โคตโมติ อิติ ปุฏฺโฐ สมาโน “น โข อหํ วจฺฉ เอวํทิฏฺฐิ ‘อนนฺตวา โลโก อิทเมว สจฺจํ โมฆมญฺญนฺ’ติ” วเทสิ.\csromanspacer{} กิํ นุ โข โภ โคตม “ตํ ชีวํ ตํ สรีรํ อิทเมว สจฺจํ โมฆมญฺญนฺ”ติ เอวํทิฏฺฐิ ภวํ โคตโมติ อิติ ปุฏฺโฐ สมาโน “น โข อหํ วจฺฉ เอวํทิฏฺฐิ ‘ตํ ชีวํ ตํ สรีรํ อิทเมว สจฺจํ โมฆมญฺญนฺ’ติ” วเทสิ.\csromanspacer{} กิํ ปน โภ โคตม “อญฺญํ ชีวํ อญฺญํ สรีรํ อิทเมว สจฺจํ โมฆมญฺญนฺ”ติ เอวํทิฏฺฐิ ภวํ โคตโมติ อิติ ปุฏฺโฐ สมาโน “น โข อหํ วจฺฉ เอวํทิฏฺฐิ ‘อญฺญํ ชีวํ อญฺญํ สรีรํ อิทเมว สจฺจํ โมฆมญฺญนฺ’ติ” วเทสิ.\csromanspacer{} กิํ นุ โข โภ โคตม “โหติ ตถาคโต ปรํ มรณา อิทเมว สจฺจํ โมฆมญฺญนฺ”ติ เอวํทิฏฺฐิ ภวํ โคตโมติ อิติ ปุฏฺโฐ สมาโน “น โข อหํ วจฺฉ เอวํทิฏฺฐิ ‘โหติ ตถาคโต ปรํ มรณา อิทเมว สจฺจํ โมฆมญฺญนฺ’ติ” วเทสิ.}

\prose{กิํ ปน โภ โคตม “น โหติ ตถาคโต ปรํ มรณา อิทเมว สจฺจํ โมฆมญฺญนฺ”ติ เอวํทิฏฺฐิ ภวํ โคตโมติ อิติ ปุฏฺโฐ สมาโน “น โข อหํ วจฺฉ เอวํทิฏฺฐิ ‘น โหติ ตถาคโต ปรํ มรณา อิทเมว สจฺจํ โมฆมญฺญนฺ’ติ” วเทสิ.\csromanspacer{} กิํ นุ โข โภ โคตม “โหติ จ น จ โหติ ตถาคโต ปรํ มรณา อิทเมว สจฺจํ โมฆมญฺญนฺ”ติ เอวํทิฏฺฐิ ภวํ โคตโมติ อิติ ปุฏฺโฐ สมาโน “น โข อหํ วจฺฉ เอวํทิฏฺฐิ ‘โหติ จ น จ โหติ ตถาคโต ปรํ มรณา อิทเมว สจฺจํ โมฆมญฺญนฺ’ติ” วเทสิ.\csromanspacer{} กิํ ปน โภ โคตม “เนว โหติ น น โหติ ตถาคโต ปรํ มรณา อิทเมว สจฺจํ โมฆมญฺญนฺ”ติ เอวํทิฏฺฐิ ภวํ โคตโมติ อิติ ปุฏฺโฐ สมาโน “น}

\vspace{\tipitakaparskip}
\csromancenter{อคฺคิวจฺฉสุตฺต (72) โข อหํ วจฺฉ เอวํทิฏฺฐิ ‘เนว โหติ น น โหติ ตถาคโต ปรํ มรณา อิทเมว สจฺจํ โมฆมญฺญนฺ’ติ” วเทสิ.}
\prose{\csromanfolio{153}กิํ ปน โภ โคตม อาทีนวํ สมฺปสฺสมาโน เอวํ อิมานิ สพฺพโส ทิฏฺฐิคตานิ อนุปคโตติ.}

\proseitem{189}{สสฺสโต โลโกติ โข วจฺฉ ทิฏฺฐิคตเมตํ ทิฏฺฐิคหนํ ทิฏฺฐิกนฺตาโร\footnote{ทิฏฺฐิกนฺตารํ (สี, อิ)} ทิฏฺฐิวิสูกํ ทิฏฺฐิวิปฺผนฺทิตํ ทิฏฺฐิสํโยชนํ สทุกฺขํ สวิฆาตํ ส-อุปายาสํ สปริฬาหํ น นิพฺพิทาย น วิราคาย น นิโรธาย น อุปสมาย น อภิญฺญาย น สมฺโพธาย น นิพฺพานาย สํวตฺตติ.\csromanspacer{} อสสฺสโต โลโกติ โข วจฺฉ ฯเปฯ.\csromanspacer{} อนฺตวา โลโกติ โข วจฺฉ ฯเปฯ.\csromanspacer{} อนนฺตวา โลโกติ โข วจฺฉ ฯเปฯ.\csromanspacer{} ตํ ชีวํ ตํ สรีรนฺติ โข วจฺฉ ฯเปฯ.\csromanspacer{} อญฺญํ ชีวํ อญฺญํ สรีรนฺติ โข วจฺฉ ฯเปฯ.\csromanspacer{} โหติ ตถาคโต ปรํ มรณาติ โข วจฺฉ ฯเปฯ.\csromanspacer{} น โหติ ตถาคโต ปรํ มรณาติ โข วจฺฉ ฯเปฯ.\csromanspacer{} โหติ จ น จ โหติ ตถาคโต ปรํ มรณาติ โข วจฺฉ ฯเปฯ.\csromanspacer{} เนว โหติ น น โหติ ตถาคโต ปรํ มรณาติ โข วจฺฉ ทิฏฺฐิคตเมตํ ทิฏฺฐิคหนํ ทิฏฺฐิกนฺตาโร ทิฏฺฐิวิสูกํ ทิฏฺฐิวิปฺผนฺทิตํ ทิฏฺฐิสํโยชนํ สทุกฺขํ สวิฆาตํ ส-อุปายาสํ สปริฬาหํ น นิพฺพิทาย น วิราคาย น นิโรธาย น อุปสมาย น อภิญฺญาย น สมฺโพธาย น นิพฺพานาย สํวตฺตติ.\csromanspacer{} อิมํ โข อหํ วจฺฉ อาทีนวํ สมฺปสฺสมาโน เอวํ อิมานิ สพฺพโส ทิฏฺฐิคตานิ อนุปคโตติ.}

\prose{อตฺถิ ปน โภโต โคตมสฺส กิญฺจิ ทิฏฺฐิคตนฺติ.\csromanspacer{} ทิฏฺฐิคตนฺติ โข วจฺฉ อปนีตเมตํ ตถาคตสฺส.\csromanspacer{} ทิฏฺฐํ เหตํ วจฺฉ ตถาคเตน “อิติ รูปํ อิติ รูปสฺส สมุทโย อิติ รูปสฺส อตฺถงฺคโม, อิติ เวทนา อิติ เวทนาย สมุทโย อิติ เวทนาย อตฺถงฺคโม, อิติ สญฺญา อิติ สญฺญาย สมุทโย อิติ สญฺญาย อตฺถงฺคโม, อิติ สงฺขารา อิติ สงฺขารานํ สมุทโย อิติ สงฺขารานํ อตฺถงฺคโม, อิติ วิญฺญาณํ อิติ วิญฺญาณสฺส สมุทโย อิติ วิญฺญาณสฺส อตฺถงฺคโม”ติ.\csromanspacer{} ตสฺมา ตถาคโต สพฺพมญฺญิตานํ สพฺพมถิตานํ สพฺพ-อหํการมมํการมานานุสยานํ ขยา วิราคา นิโรธา จาคา ปฏินิสฺสคฺคา อนุปาทา วิมุตฺโตติ วทามีติ.}

\proseitem{190}{\csromanfolio{154}เอวํ วิมุตฺตจิตฺโต ปน โภ โคตม ภิกฺขุ กุหิํ อุปปชฺชตีติ.\csromanspacer{} อุปปชฺชตีติ โข วจฺฉ น อุเปติ.\csromanspacer{} เตน หิ โภ โคตม น อุปปชฺชตีติ.\csromanspacer{} น อุปปชฺชตีติ โข วจฺฉ น อุเปติ.\csromanspacer{} เตน หิ โภ โคตม อุปปชฺชติ จ น จ อุปปชฺชตีติ.\csromanspacer{} อุปปชฺชติ จ น จ อุปปชฺชตีติ โข วจฺฉ น อุเปติ.\csromanspacer{} เตน หิ โภ โคตม เนว อุปปชฺชติ น น อุปปชฺชตีติ.\csromanspacer{} เนว อุปปชฺชติ น น อุปปชฺชตีติ โข วจฺฉ น อุเปติ.}

\prose{“เอวํ วิมุตฺตจิตฺโต ปน โภ โคตม ภิกฺขุ กุหิํ อุปปชฺชตี”ติ อิติ ปุฏฺโฐ สมาโน “อุปปชฺชตีติ โข วจฺฉ น อุเปตี”ติ วเทสิ. “เตน หิ โภ โคตม น อุปปชฺชตี”ติ อิติ ปุฏฺโฐ สมาโน “น อุปปชฺชตีติ โข วจฺฉ น อุเปตี”ติ วเทสิ. “เตน หิ โภ โคตม อุปปชฺชติ จ น จ อุปปชฺชตี”ติ อิติ ปุฏฺโฐ สมาโน “อุปปชฺชติ จ น จ อุปปชฺชตีติ โข วจฺฉ น อุเปตี”ติ วเทสิ. “เตน หิ โภ โคตม เนว อุปปชฺชติ น น อุปปชฺชตี”ติ อิติ ปุฏฺโฐ สมาโน “เนว อุปปชฺชติ น น อุปปชฺชตีติ โข วจฺฉ น อุเปตี”ติ วเทสิ.\csromanspacer{} เอตฺถาหํ โภ โคตม อญฺญาณมาปาทิํ, เอตฺถ สมฺโมหมาปาทิํ.\csromanspacer{} ยาปิ เม เอสา โภโต โคตมสฺส ปุริเมน กถาสลฺลาเปน อหุ ปสาทมตฺตา, สาปิ เม เอตรหิ อนฺตรหิตาติ.\csromanspacer{} อลํ หิ เต วจฺฉ อญฺญาณาย อลํ สมฺโมหาย.\csromanspacer{} คมฺภีโร หายํ วจฺฉ ธมฺโม ทุทฺทโส ทุรนุโพโธ สนฺโต ปณีโต อตกฺกาวจโร นิปุโณ ปณฺฑิตเวทนีโย, โส ตยา ทุชฺชาโน อญฺญทิฏฺฐิเกน อญฺญขนฺติเกน อญฺญรุจิเกน อญฺญตฺรโยเคน\footnote{อญฺญตฺราโยเคน (ที 1. 173)} อญฺญตฺราจริยเกน\footnote{อญฺญตฺถาจริยเกน (สี, สฺยา, กํ, อิ)}.}

\proseitem{191}{เตน หิ วจฺฉ ตญฺเญเวตฺถ ปฏิปุจฺฉิสฺสามิ, ยถา เต ขเมยฺย ตถา นํ พฺยากเรยฺยาสิ.\csromanspacer{} ตํ กิํ มญฺญสิ วจฺฉ สเจ เต ปุรโต อคฺคิ ชเลยฺย, ชาเนยฺยาสิ ตฺวํ “อยํ เม ปุรโต อคฺคิ ชลตี”ติ.\csromanspacer{} สเจ เม โภ โคตม ปุรโต อคฺคิ ชเลยฺย, ชาเนยฺยาหํ “อยํ เม ปุรโต อคฺคิ ชลตี”ติ.}

\prose{สเจ ปน ตํ วจฺฉ เอวํ ปุจฺเฉยฺย “โย เต อยํ ปุรโต อคฺคิ ชลติ, อยํ อคฺคิ กิํ ปฏิจฺจ ชลตี”ติ, เอวํ ปุฏฺโฐ ตฺวํ วจฺฉ กินฺติ พฺยากเรยฺยาสีติ.\csromanspacer{} สเจ มํ โภ โคตม เอวํ ปุจฺเฉยฺย “โย เต อยํ ปุรโต อคฺคิ ชลติ, อยํ อคฺคิ กิํ ปฏิจฺจ ชลตี”ติ, เอวํ ปุฏฺโฐ อหํ โภ โคตม เอวํ}

\vspace{\tipitakaparskip}
\csromancenter{อคฺคิวจฺฉสุตฺต (72) พฺยากเรยฺยํ “โย เม อยํ ปุรโต อคฺคิ ชลติ, อยํ อคฺคิ ติณกฏฺฐุปาทานํ ปฏิจฺจ ชลตี”ติ.}
\prose{\csromanfolio{155}สเจ เต วจฺฉ ปุรโต โส อคฺคิ นิพฺพาเยยฺย, ชาเนยฺยาสิ ตฺวํ “อยํ เม ปุรโต อคฺคิ นิพฺพุโต”ติ.\csromanspacer{} สเจ เม โภ โคตม ปุรโต โส อคฺคิ นิพฺพาเยยฺย, ชาเนยฺยาหํ “อยํ เม ปุรโต อคฺคิ นิพฺพุโต”ติ.}

\prose{สเจ ปน ตํ วจฺฉ เอวํ ปุจฺเฉยฺย “โย เต อยํ ปุรโต อคฺคิ นิพฺพุโต, โส อคฺคิ อิโต กตมํ ทิสํ คโต, ปุรตฺถิมํ วา ทกฺขิณํ วา ปจฺฉิมํ วา อุตฺตรํ วา”ติ, เอวํ ปุฏฺโฐ ตฺวํ วจฺฉ กินฺติ พฺยากเรยฺยาสีติ.\csromanspacer{} น อุเปติ โภ โคตม.\csromanspacer{} ยญฺหิ โส โภ โคตม อคฺคิ ติณกฏฺฐุปาทานํ ปฏิจฺจ อชลิ\footnote{ชลติ (สฺยา, กํ, ก)}, ตสฺส จ ปริยาทานา อญฺญสฺส จ อนุปหารา อนาหาโร นิพฺพุโต เตฺวว สงฺขฺยํ คจฺฉตีติ.}

\proseitem{192}{เอวเมว โข วจฺฉ เยน รูเปน ตถาคตํ ปญฺญาปยมาโน ปญฺญาเปยฺย, ตํ รูปํ ตถาคตสฺส ปหีนํ อุจฺฉินฺนมูลํ ตาลาวตฺถุกตํ อนภาวํกตํ อายติํ อนุปฺปาทธมฺมํ.\csromanspacer{} รูปสงฺขยวิมุตฺโต\footnote{รูปสงฺขาวิมุตฺโต (สี, สฺยา, กํ, อิ) เอวํ เวทนาสงฺขยาทีสุปิ.} โข วจฺฉ ตถาคโต คมฺภีโร อปฺปเมยฺโย ทุปฺปริโยคาโฬฺห, เสยฺยถาปิ มหาสมุทฺโท.\csromanspacer{} อุปปชฺชตีติ น อุเปติ, น อุปปชฺชตีติ น อุเปติ, อุปปชฺชติ จ น จ อุปปชฺชตีติ น อุเปติ, เนว อุปปชฺชติ น น อุปปชฺชตีติ น อุเปติ.}

\prose{ยาย เวทนาย ตถาคตํ ปญฺญาปยมาโน ปญฺญาเปยฺย, สา เวทนา ตถาคตสฺส ปหีนา อุจฺฉินฺนมูลา ตาลาวตฺถุกตา อนภาวํกตา อายติํ อนุปฺปาทธมฺมา.\csromanspacer{} เวทนาสงฺขยวิมุตฺโต โข วจฺฉ ตถาคโต คมฺภีโร อปฺปเมยฺโย ทุปฺปริโยคาโฬฺห, เสยฺยถาปิ มหาสมุทฺโท.\csromanspacer{} อุปปชฺชตีติ น อุเปติ, น อุปปชฺชตีติ น อุเปติ, อุปปชฺชติ จ น จ อุปปชฺชตีติ น อุเปติ, เนว อุปปชฺชติ น น อุปปชฺชตีติ น อุเปติ.}

\prose{ยาย สญฺญาย ตถาคตํ ปญฺญาปยมาโน ปญฺญาเปยฺย, สา สญฺญา ตถาคตสฺส ปหีนา อุจฺฉินฺนมูลา ตาลาวตฺถุกตา อนภาวํกตา อายติํ อนุปฺปาทธมฺมา, สญฺญาสงฺขยวิมุตฺโต โข วจฺฉ ตถาคโต คมฺภีโร อปฺปเมยฺโย ทุปฺปริโยคาโฬฺห, เสยฺยถาปิ มหาสมุทฺโท. \csromanfolio{156}อุปปชฺชตีติ น อุเปติ, น อุปปชฺชตีติ น อุเปติ, อุปปชฺชติ จ น จ อุปปชฺชตีติ น อุเปติ, เนว อุปปชฺชติ น น อุปปชฺชตีติ น อุเปติ.}

\prose{เยหิ สงฺขาเรหิ ตถาคตํ ปญฺญาปยมาโน ปญฺญาเปยฺย, เต สงฺขารา ตถาคตสฺส ปหีนา อุจฺฉินฺนมูลา ตาลาวตฺถุกตา อนภาวํกตา อายติํ อนุปฺปาทธมฺมา.\csromanspacer{} สงฺขารสงฺขยวิมุตฺโต โข วจฺฉ ตถาคโต คมฺภีโร อปฺปเมยฺโย ทุปฺปริโยคาโฬฺห, เสยฺยถาปิ มหาสมุทฺโท.\csromanspacer{} อุปปชฺชตีติ น อุเปติ, น อุปปชฺชตีติ น อุเปติ, อุปปชฺชติ จ น จ อุปปชฺชตีติ น อุเปติ, เนว อุปปชฺชติ น น อุปปชฺชตีติ น อุเปติ.}

\prose{เยน วิญฺญาเณน ตถาคตํ ปญฺญาปยมาโน ปญฺญาเปยฺย, ตํ วิญฺญาณํ ตถาคตสฺส ปหีนํ อุจฺฉินฺนมูลํ ตาลาวตฺถุกตํ อนภาวํกตํ อายติํ อนุปฺปาทธมฺมํ.\csromanspacer{} วิญฺญาณสงฺขยวิมุตฺโต โข วจฺฉ ตถาคโต คมฺภีโร อปฺปเมยฺโย ทุปฺปริโยคาโฬฺห, เสยฺยถาปิ มหาสมุทฺโท.\csromanspacer{} อุปปชฺชตีติ น อุเปติ, น อุปปชฺชตีติ น อุเปติ, อุปปชฺชติ จ น จ อุปปชฺชตีติ น อุเปติ, เนว อุปปชฺชติ น น อุปปชฺชตีติ น อุเปติ.}

\prose{เอวํ วุตฺเต วจฺฉโคตฺโต ปริพฺพาชโก ภควนฺตํ เอตทโวจ “เสยฺยถาปิ โภ โคตม คามสฺส วา นิคมสฺส วา อวิทูเร มหาสาลรุกฺโข, ตสฺส อนิจฺจตา สาขาปลาสา ปลุชฺเชยฺยุํ\footnote{สาขาปลาสํ ปลุชฺเชยฺย}, ตจปปฏิกา ปลุชฺเชยฺยุํ, เผคฺคู ปลุชฺเชยฺยุํ\footnote{เผคฺคุ ปลุชฺเชยฺย (สี, สฺยา, กํ, อิ)}.\csromanspacer{} โส อปเรน สมเยน อปคตสาขาปลาโส อปคตตจปปฏิโก อปคตเผคฺคุโก สุทฺโธ อสฺส สาเร ปติฏฺฐิโต.\csromanspacer{} เอวเมว โภโต โคตมสฺส ปาวจนํ อปคตสาขาปลาสํ อปคตตจปปฏิกํ อปคตเผคฺคุกํ สุทฺธํ สาเร ปติฏฺฐิตํ.\csromanspacer{} อภิกฺกนฺตํ โภ โคตม ฯเปฯ อุปาสกํ มํ ภวํ โคตโม ธาเรตุ อชฺชตคฺเค ปาณุเปตํ สรณํ คตนฺ”ติ.}

\nitthitamruled{อคฺคิวจฺฉ สุตฺตํ นิฏฺฐิตํ ทุติยํ.}
\csromanmatikaanchor{mtk.28}
\csromantocmarkref{subsection}{3. มหาวจฺฉสุตฺต}{mtk.28}
\bookmark[dest={mtk.28},level=2]{3. มหาวจฺฉสุตฺต}
\csromanheader{2}{3. มหาวจฺฉสุตฺต}
\proseitem{193}{เอวํ เม สุตํ\csromandash{}เอกํ สมยํ ภควา ราชคเห วิหรติ เวฬุวเน กลนฺทกนิวาเป.\csromanspacer{} อถ โข วจฺฉโคตฺโต ปริพฺพาชโก เยน}

\vspace{\tipitakaparskip}
\csromancenter{มหาวจฺฉสุตฺต (73) ภควา เตนุปสงฺกมิ, อุปสงฺกมิตฺวา ภควตา สทฺธิํ สมฺโมทิ, สมฺโมทนียํ กถํ สารณียํ วีติสาเรตฺวา เอกมนฺตํ นิสีทิ, เอกมนฺตํ นิสินฺโน โข วจฺฉโคตฺโต ปริพฺพาชโก ภควนฺตํ เอตทโวจ “ทีฆรตฺตาหํ โภตา โคตเมน สหกถี, สาธุ เม ภวํ โคตโม สํขิตฺเตน กุสลากุสลํ เทเสตู”ติ.\csromanspacer{} สํขิตฺเตนปิ โข เต อหํ วจฺฉ กุสลากุสลํ เทเสยฺยํ, วิตฺถาเรนปิ โข เต อหํ วจฺฉ กุสลากุสลํ เทเสยฺยํ.\csromanspacer{} อปิ จ เต อหํ วจฺฉ สํขิตฺเตน กุสลากุสลํ เทเสสฺสามิ, ตํ สุณาหิ สาธุกํ มนสิ กโรหิ ภาสิสฺสามีติ. “เอวํ โภ”ติ โข วจฺฉโคตฺโต ปริพฺพาชโก ภควโต ปจฺจสฺโสสิ.\csromanspacer{} ภควา เอตทโวจ\csromandash{}}
\proseitem{194}{\csromanfolio{157}โลโภ โข วจฺฉ อกุสลํ, อโลโภ กุสลํ.\csromanspacer{} โทโส โข วจฺฉ อกุสลํ, อโทโส กุสลํ.\csromanspacer{} โมโห โข วจฺฉ อกุสลํ, อโมโห กุสลํ.\csromanspacer{} อิติ โข วจฺฉ อิเม ตโย ธมฺมา อกุสลา, ตโย ธมฺมา กุสลา.}

\prose{ปาณาติปาโต โข วจฺฉ อกุสลํ, ปาณาติปาตา เวรมณี กุสลํ.\csromanspacer{} อทินฺนาทานํ โข วจฺฉ อกุสลํ, อทินฺนาทานา เวรมณี กุสลํ.\csromanspacer{} กาเมสุมิจฺฉาจาโร โข วจฺฉ อกุสลํ, กาเมสุมิจฺฉาจารา เวรมณี กุสลํ.\csromanspacer{} มุสาวาโท โข วจฺฉ อกุสลํ, มุสาวาทา เวรมณี กุสลํ.\csromanspacer{} ปิสุณา วาจา โข วจฺฉ อกุสลํ, ปิสุณาย วาจาย เวรมณี กุสลํ.\csromanspacer{} ผรุสา วาจา โข วจฺฉ อกุสลํ, ผรุสาย วาจาย เวรมณี กุสลํ.\csromanspacer{} สมฺผปฺปลาโป โข วจฺฉ อกุสลํ, สมฺผปฺปลาปา เวรมณี กุสลํ.\csromanspacer{} อภิชฺฌา โข วจฺฉ อกุสลํ, อนภิชฺฌา กุสลํ.\csromanspacer{} พฺยาปาโท โข วจฺฉ อกุสลํ, อพฺยาปาโท กุสลํ.\csromanspacer{} มิจฺฉาทิฏฺฐิ โข วจฺฉ อกุสลํ, สมฺมาทิฏฺฐิ กุสลํ.\csromanspacer{} อิติ โข วจฺฉ อิเม ทส ธมฺมา อกุสลา, ทส ธมฺมา กุสลา.}

\prose{ยโต โข วจฺฉ ภิกฺขุโน ตณฺหา ปหีนา โหติ อุจฺฉินฺนมูลา ตาลาวตฺถุกตา อนภาวํกตา อายติํ อนุปฺปาทธมฺมา.\csromanspacer{} โส โหติ ภิกฺขุ อรหํ ขีณาสโว วุสิตวา กตกรณีโย โอหิตภาโร อนุปฺปตฺตสทตฺโถ ปริกฺขีณภวสํโยชโน สมฺมทญฺญา วิมุตฺโตติ.}

\proseitem{195}{ติฏฺฐตุ ภวํ โคตโม.\csromanspacer{} อตฺถิ ปน เต โภโต โคตมสฺส เอกภิกฺขุปิ สาวโก, โย อาสวานํ ขยา\footnote{สาวโก อาสวานํ ขยา (สี, สฺยา, กํ, อิ) เอวมุปริปิ.} อนาสวํ เจโตวิมุตฺติํ \csromanfolio{158}ปญฺญาวิมุตฺติํ ทิฏฺเฐว ธมฺเม สยํ อภิญฺญา สจฺฉิกตฺวา อุปสมฺปชฺชวิหรตีติ.\csromanspacer{} น โข วจฺฉ เอกํเยว สตํ น เทฺว สตานิ น ตีณิ สตานิ น จตฺตาริ สตานิ น ปญฺจ สตานิ, อถ โข ภิยฺโยว เย ภิกฺขู มม สาวกา อาสวานํ ขยา อนาสวํ เจโตวิมุตฺติํ ปญฺญาวิมุตฺติํ ทิฏฺเฐว ธมฺเม สยํ อภิญฺญา สจฺฉิกตฺวา อุปสมฺปชฺช วิหรนฺตีติ.}

\prose{ติฏฺฐตุ ภวํ โคตโม, ติฏฺฐนฺตุ ภิกฺขู.\csromanspacer{} อตฺถิ ปน โภโต โคตมสฺส เอกา ภิกฺขุนีปิ สาวิกา, ยา อาสวานํ ขยา อนาสวํ เจโตมุตฺติํ ปญฺญาวิมุตฺติํ ทิฏฺเฐว ธมฺเม สยํ อภิญฺญา สจฺฉิกตฺวา อุปสมฺปชฺช วิหรตีติ.\csromanspacer{} น โข วจฺฉ เอกํเยว สตํ น เทฺว สตานิ น ตีณิ สตานิ น จตฺตาริ สตานิ น ปญฺจ สตานิ, อถ โข ภิยฺโยว ยา ภิกฺขุนิโย มม สาวิกา อาสวานํ ขยา อนาสวํ เจโตวิมุตฺติํ ปญฺญาวิมุตฺติํ ทิฏฺเฐว ธมฺเม สยํ อภิญฺญา สจฺฉิกตฺวา อุปสมฺปชฺช วิหรนฺตีติ.}

\prose{ติฏฺฐตุ ภวํ โคตโม, ติฏฺฐนฺตุ ภิกฺขู, ติฏฺฐนฺตุ ภิกฺขุนิโย.\csromanspacer{} อตฺถิ ปน โภโต โคตมสฺส เอกุปาสโกปิ สาวโก คิหี โอทาตวสโน พฺรหฺมจารี, โย ปญฺจนฺนํ โอรมฺภาคิยานํ สํโยชนานํ ปริกฺขยา โอปปาติโก ตตฺถ ปรินิพฺพายี อนาวตฺติธมฺโม ตสฺมา โลกาติ.\csromanspacer{} น โข วจฺฉ เอกํเยว สตํ น เทฺว สตานิ น ตีณิ สตานิ น จตฺตาริ สตานิ น ปญฺจ สตานิ, อถ โข ภิยฺโยว เย อุปาสกา มม สาวกา คิหี โอทาตวสนา พฺรหฺมจาริโน ปญฺจนฺนํ โอรมฺภาคิยานํ สํโยชนานํ ปริกฺขยา โอปปาติกา ตตฺถ ปรินิพฺพายิโน อนาวตฺติธมฺมา ตสฺมา โลกาติ.}

\prose{ติฏฺฐตุ ภวํ โคตโม, ติฏฺฐนฺตุ ภิกฺขู, ติฏฺฐนฺตุ ภิกฺขุนิโย, ติฏฺฐนฺตุ อุปาสกา คิหี โอทาตวสนา พฺรหฺมจาริโน.\csromanspacer{} อตฺถิ ปน โภโต โคตมสฺส เอกุปาสโกปิ สาวโก คิหี โอทาตวสโน กามโภคี สาสนกโร โอวาทปฺปฏิกโร, โย ติณฺณวิจิกิจฺโฉ วิคตกถํกโถ เวสารชฺชปฺปตฺโต อปรปฺปจฺจโย สตฺถุสาสเน วิหรตีติ.\csromanspacer{} น โข วจฺฉ เอกํเยว สตํ น เทฺว สตานิ น ตีณิ สตานิ น จตฺตาริ สตานิ น ปญฺจ สตานิ, อถ โข ภิยฺโยว เย อุปาสกา มม สาวกา คิหี โอทาตวสนา กามโภคิโน สาสนกรา โอวาทปฺปฏิกรา ติณฺณวิจิกิจฺฉา วิคตกถํกถา เวสารชฺชปฺปตฺตา อปรปฺปจฺจยา สตฺถุสาสเน วิหรนฺตีติ.}

\csromanheader{2}{3. มหาวจฺฉสุตฺต (73)}
\prose{\csromanfolio{159}ติฏฺฐตุ ภวํ โคตโม, ติฏฺฐนฺตุ ภิกฺขู, ติฏฺฐนฺตุ ภิกฺขุนิโย, ติฏฺฐนฺตุ อุปาสกา คิหี โอทาตวสนา พฺรหฺมจาริโน, ติฏฺฐนฺตุ อุปาสกา คิหี โอทาตวสนา กามโภคิโน.\csromanspacer{} อตฺถิ ปน โภโต โคตมสฺส เอกุปาสิกาปิ สาวิกา คิหินี โอทาตวสนา พฺรหฺมจารินี, ยา ปญฺจนฺนํ โอรมฺภาคิยานํ สํโยชนานํ ปริกฺขยา โอปปาติกา ตตฺถ ปรินิพฺพายินี อนาวตฺติธมฺมา ตสฺมา โลกาติ.\csromanspacer{} น โข วจฺฉ เอกํเยว สตํ น เทฺว สตานิ น ตีณิ สตานิ น จตฺตาริ สตานิ น ปญฺจ สตานิ, อถ โข ภิยฺโยว ยา อุปาสิกา มม สาวิกา คิหินิโย โอทาตวสนา พฺรหฺมจารินิโย ปญฺจนฺนํ โอรมฺภาคิยานํ สํโยชนานํ ปริกฺขยา โอปปาติกา ตตฺถ ปรินิพฺพายินิโย อนาวตฺติธมฺมา ตสฺมา โลกาติ.}

\prose{ติฏฺฐตุ ภวํ โคตโม, ติฏฺฐนฺตุ ภิกฺขู, ติฏฺฐนฺตุ ภิกฺขุนิโย, ติฏฺฐนฺตุ อุปาสกา คิหี โอทาตวสนา พฺรหฺมจาริโน, ติฏฺฐนฺตุ อุปาสกา คิหี โอทาตวสนา กามโภคิโน, ติฏฺฐนฺตุ อุปาสิกา คิหินิโย โอทาตวสนา พฺรหฺมจารินิโย.\csromanspacer{} อตฺถิ ปน โภโต โคตมสฺส เอกุปาสิกาปิ สาวิกา คิหินี โอทาตวสนา กามโภคินี สาสนกรา โอวาทปฺปฏิกรา, ยา ติณฺณวิจิกิจฺฉา วิคตกถํกถา เวสารชฺชปฺปตฺตา อปรปฺปจฺจยา สตฺถุสาสเน วิหรตีติ.\csromanspacer{} น โข วจฺฉ เอกํเยว สตํ น เทฺว สตานิ น ตีณิ สตานิ น จตฺตาริ สตานิ น ปญฺจ สตานิ, อถ โข ภิยฺโยว ยา อุปาสิกา มม สาวิกา คิหินิโย โอทาตวสนา กามโภคินิโย สาสนกรา โอวาทปฺปฏิกรา ติณฺณวิจิกิจฺฉา วิคตกถํกถา เวสารชฺชปฺปตฺตา อปรปฺปจฺจยา สตฺถุสาสเน วิหรนฺตีติ.}

\proseitem{196}{สเจ หิ โภ โคตม อิมํ ธมฺมํ ภวํเยว โคตโม อาราธโก อภวิสฺส, โน จ โข ภิกฺขู อาราธกา อภวิสฺสํสุ, เอวมิทํ พฺรหฺมจริยํ อปริปูรํ อภวิสฺส เตนงฺเคน.\csromanspacer{} ยสฺมา จ โข โภ โคตม อิมํ ธมฺมํ ภวํ เจว โคตโม อาราธโก ภิกฺขู จ อาราธกา, เอวมิทํ พฺรหฺมจริยํ ปริปูรํ เตนงฺเคน.}

\prose{สเจ หิ โภ โคตม อิมํ ธมฺมํ ภวํ เจว โคตโม อาราธโก อภวิสฺส, ภิกฺขู จ อาราธกา อภวิสฺสํสุ, โน จ โข ภิกฺขุนิโย อาราธิกา อภวิสฺสํสุ, เอวมิทํ พฺรหฺมจริยํ อปริปูรํ อภวิสฺส เตนงฺเคน.\csromanspacer{} ยสฺมา จ โข โภ โคตม อิมํ ธมฺมํ ภวํ เจว โคตโม อาราธโก \csromanfolio{160}ภิกฺขู จ อาราธกา ภิกฺขุนิโย จ อาราธิกา, เอวมิทํ พฺรหฺมจริยํ ปริปูรํ เตนงฺเคน.}

\prose{สเจ หิ โภ โคตม อิมํ ธมฺมํ ภวํ เจว โคตโม อาราธโก อภวิสฺส, ภิกฺขู จ อาราธกา อภวิสฺสํสุ, ภิกฺขุนิโย จ อาราธิกา อภวิสฺสํสุ, โน จ โข อุปาสกา คิหี โอทาตวสนา พฺรหฺมจาริโน อาราธกา อภวิสฺสํสุ, เอวมิทํ พฺรหฺมจริยํ อปริปูรํ อภวิสฺส เตนงฺเคน.\csromanspacer{} ยสฺมา จ โข โภ โคตม อิมํ ธมฺมํ ภวํ เจว โคตโม อาราธโก ภิกฺขู จ อาราธกา ภิกฺขุนิโย จ อาราธิกา อุปาสกา จ คิหี โอทาตวสนา พฺรหฺมจาริโน อาราธกา, เอวมิทํ พฺรหฺมจริยํ ปริปูรํ เตนงฺเคน.}

\prose{สเจ หิ โภ โคตม อิมํ ธมฺมํ ภวํ เจว โคตโม อาราธโก อภวิสฺส, ภิกฺขู จ อาราธกา อภวิสฺสํสุ, ภิกฺขุนิโย จ อาราธิกา อภวิสฺสํสุ, อุปาสกา จ คิหี โอทาตวสนา พฺรหฺมจาริโน อาราธกา อภวิสฺสํสุ, โน จ โข อุปาสกา คิหี โอทาตวสนา กามโภคิโน อาราธกา อภวิสฺสํสุ, เอวมิทํ พฺรหฺมจริยํ อปริปูรํ อภวิสฺส เตนงฺเคน.\csromanspacer{} ยสฺมา จ โข โภ โคตม อิมํ ธมฺมํ ภวํ เจว โคตโม อาราธโก ภิกฺขู จ อาราธกา ภิกฺขุนิโย จ อาราธิกา อุปาสกา จ คิหี โอทาตวสนา พฺรหฺมจาริโน อาราธกา อุปาสกา จ คิหี โอทาตวสนา กามโภคิโน อาราธกา, เอวมิทํ พฺรหฺมจริยํ ปริปูรํ เตนงฺเคน.}

\prose{สเจ หิ โภ โคตม อิมํ ธมฺมํ ภวํ เจว โคตโม อาราธโก อภวิสฺส, ภิกฺขู จ อาราธกา อภวิสฺสํสุ, ภิกฺขุนิโย จ อาราธิกา อภวิสฺสํสุ, อุปาสกา จ คิหี โอทาตวสนา พฺรหฺมจาริโน อาราธกา อภวิสฺสํสุ, อุปาสกา จ คิหี โอทาตวสนา กามโภคิโน อาราธกา อภวิสฺสํสุ, โน จ โข อุปาสิกา คิหินิโย โอทาตวสนา พฺรหฺมจารินิโย อาราธิกา อภวิสฺสํสุ, เอวมิทํ พฺรหฺมจริยํ อปริปูรํ อภวิสฺส เตนงฺเคน.\csromanspacer{} ยสฺมา จ โข โภ โคตม อิมํ ธมฺมํ ภวํ เจว โคตโม อาราธโก ภิกฺขู จ อาราธกา ภิกฺขุนิโย จ อาราธิกา อุปาสกา จ คิหี โอทาตวสนา พฺรหฺมจาริโน อาราธกา อุปาสกา จ คิหี โอทาตวสนา กามโภคิโน อาราธกา อุปาสิกา จ คิหินิโย}

\vspace{\tipitakaparskip}
\csromancenter{มหาวจฺฉสุตฺต (73) โอทาตวสนา พฺรหฺมจารินิโย อาราธิกา, เอวมิทํ พฺรหฺมจริยํ ปริปูรํ เตนงฺเคน.}
\prose{\csromanfolio{161}สเจ หิ โภ โคตม อิมํ ธมฺมํ ภวํ เจว โคตโม อาราธโก อภวิสฺส, ภิกฺขู จ อาราธกา อภวิสฺสํสุ, ภิกฺขุนิโย จ อาราธิกา อภวิสฺสํสุ, อุปาสกา จ คิหี โอทาตวสนา พฺรหฺมจาริโน อาราธกา อภวิสฺสํสุ, อุปาสกา จ คิหี โอทาตวสนา กามโภคิโน อาราธกา อภวิสฺสํสุ, อุปาสิกา จ คิหินิโย โอทาตวสนา พฺรหฺมจารินิโย อาราธิกา อภวิสฺสํสุ, โน จ โข อุปาสิกา คิหินิโย โอทาตวสนา กามโภคินิโย อาราธิกา อภวิสฺสํสุ, เอวมิทํ พฺรหฺมจริยํ อปริปูรํ อภวิสฺส เตนงฺเคน.\csromanspacer{} ยสฺมา จ โข โภ โคตม อิมํ ธมฺมํ ภวํ เจว โคตโม อาราธโก ภิกฺขู จ อาราธกา ภิกฺขุนิโย จ อาราธิกา อุปาสกา จ คิหี โอทาตวสนา พฺรหฺมจาริโน อาราธกา อุปาสกา จ คิหี โอทาตวสนา กามโภคิโน อาราธกา อุปาสิกา จ คิหินิโย โอทาตวสนา พฺรหฺมจารินิโย อาราธิกา อุปาสิกา จ คิหินิโย โอทาตวสนา กามโภคินิโย อาราธิกา, เอวมิทํ พฺรหฺมจริยํ ปริปูรํ เตนงฺเคน.}

\proseitem{197}{เสยฺยถาปิ โภ โคตม คงฺคา นที สมุทฺทนินฺนา สมุทฺทโปณา สมุทฺทปพฺภารา สมุทฺทํ อาหจฺจ ติฏฺฐติ.\csromanspacer{} เอวเมวายํ โภโต โคตมสฺส ปริสา สคหฏฺฐปพฺพชิตา นิพฺพานนินฺนา นิพฺพานโปณา นิพฺพานปพฺภารา นิพฺพานํ อาหจฺจ ติฏฺฐติ.\csromanspacer{} อภิกฺกนฺตํ โภ โคตม ฯเปฯ.\csromanspacer{} เอสาหํ ภวนฺตํ โคตมํ สรณํ คจฺฉามิ ธมฺมญฺจ ภิกฺขุสํฆญฺจ.\csromanspacer{} ลเภยฺยาหํ โภโต โคตมสฺส สนฺติเก ปพฺพชฺชํ, ลเภยฺยํ อุปสมฺปทนฺติ.\csromanspacer{} โย โข วจฺฉ อญฺญติตฺถิยปุพฺโพ อิมสฺมิํ ธมฺมวินเย อากงฺขติ ปพฺพชฺชํ อากงฺขติ อุปสมฺปทํ, โส จตฺตาโร มาเส ปริวสติ, จตุนฺนํ มาสานํ อจฺจเยน อารทฺธจิตฺตา ภิกฺขู ปพฺพาเชนฺติ อุปสมฺปาเทนฺติ ภิกฺขุภาวาย.\csromanspacer{} อปิ จ เมตฺถ ปุคฺคลเวมตฺตตา วิทิตาติ.\csromanspacer{} สเจ ภนฺเต อญฺญติตฺถิยปุพฺพา อิมสฺมิํ ธมฺมวินเย อากงฺขนฺตา ปพฺพชฺชํ อากงฺขนฺตา อุปสมฺปทํ จตฺตาโร มาเส ปริวสนฺติ, จตุนฺนํ มาสานํ อจฺจเยน อารทฺธจิตฺตา ภิกฺขู ปพฺพาเชนฺติ อุปสมฺปาเทนฺติ ภิกฺขุภาวาย.\csromanspacer{} อหํ จตฺตาริ วสฺสานิ ปริวสิสฺสามิ, จตุนฺนํ วสฺสานํ อจฺจเยน อารทฺธจิตฺตา ภิกฺขู ปพฺพาเชนฺตุ อุปสมฺปาเทนฺตุ ภิกฺขุภาวายาติ.\csromanspacer{} อลตฺถ โข วจฺฉโคตฺโต ปริพฺพาชโก ภควโต สนฺติเก ปพฺพชฺชํ, อลตฺถ อุปสมฺปทํ.}

\prose{\csromanfolio{162}อจิรูปสมฺปนฺโน โข ปนายสฺมา วจฺฉโคตฺโต อทฺธมาสูปสมฺปนฺโน เยน ภควา เตนุปสงฺกมิ, อุปสงฺกมิตฺวา ภควนฺตํ อภิวาเทตฺวา เอกมนฺตํ นิสีทิ, เอกมนฺตํ นิสินฺโน โข อายสฺมา วจฺฉโคตฺโต ภควนฺตํ เอตทโวจ “ยาวตกํ ภนฺเต เสเขน ญาเณน เสขาย วิชฺชาย ปตฺตพฺพํ, อนุปฺปตฺตํ ตํ มยา, อุตฺตริ จ เม\footnote{อุตฺตริํ เม (สี, สฺยา, กํ, อิ)} ภควา ธมฺมํ เทเสตู”ติ.\csromanspacer{} เตน หิ ตฺวํ วจฺฉ เทฺว ธมฺเม อุตฺตริ ภาเวหิ สมถญฺจ วิปสฺสนญฺจ, อิเม โข เต วจฺฉ เทฺว ธมฺมา อุตฺตริ ภาวิตา สมโถ จ วิปสฺสนา จ อเนกธาตุปฏิเวธาย สํวตฺติสฺสนฺติ.}

\proseitem{198}{โส ตฺวํ วจฺฉ ยาวเทว\footnote{ยาวเท (อิ)} อากงฺขิสฺสสิ “อเนกวิหิตํ อิทฺธิวิธํ ปจฺจนุภเวยฺยํ, เอโกปิ หุตฺวา พหุธา อสฺสํ, พหุธาปิ หุตฺวา เอโก อสฺสํ, อาวิภาวํ ติโรภาวํ ติโรกุฏฺฏํ ติโรปาการํ ติโรปพฺพตํ อสชฺชมาโน คจฺเฉยฺยํ, เสยฺยถาปิ อากาเส.\csromanspacer{} ปถวิยาปิ อุมฺมุชฺชนิมุชฺชํ กเรยฺยํ, เสยฺยถาปิ อุทเก.\csromanspacer{} อุทเกปิ อภิชฺชมาเน คจฺเฉยฺยํ, เสยฺยถาปิ ปถวิยํ.\csromanspacer{} อากาเสปิ ปลฺลงฺเกน กเมยฺยํ, เสยฺยถาปิ ปกฺขี สกุโณ.\csromanspacer{} อิเมปิ จนฺทิมสูริเย เอวํมหิทฺธิเก เอวํมหานุภาเว ปาณินา ปริมเสยฺยํ ปริมชฺเชยฺยํ, ยาวพฺรหฺมโลกาปิ กาเยน วสํ วตฺเตยฺยนฺ”ติ.\csromanspacer{} ตตฺร ตเตฺรว สกฺขิภพฺพตํ ปาปุณิสฺสสิ สติ สติ-อายตเน. (1)}

\prose{โส ตฺวํ วจฺฉ ยาวเทว อากงฺขิสฺสสิ “ทิพฺพาย โสตธาตุยา วิสุทฺธาย อติกฺกนฺตมานิสิกาย อุโภ สทฺเท สุเณยฺยํ ทิพฺเพ จ มานุเส จ เย ทูเร สนฺติเก จา”ติ.\csromanspacer{} ตตฺร ตเตฺรว สกฺขิภพฺพตํ ปาปุณิสฺสสิ สติ สติ- อายตเน. (2)}

\prose{โส ตฺวํ วจฺฉ ยาวเทว อากงฺขิสฺสสิ “ปรสตฺตานํ ปรปุคฺคลานํ เจตสา เจโต ปริจฺจ ปชาเนยฺยํ.\csromanspacer{} สราคํ วา จิตฺตํ สราคํ จิตฺตนฺติ ปชาเนยฺยํ, วีตราคํ วา จิตฺตํ วีตราคํ จิตฺตนฺติ ปชาเนยฺยํ.\csromanspacer{} สโทสํ วา จิตฺตํ สโทสํ จิตฺตนฺติ ปชาเนยฺยํ, วีตโทสํ วา จิตฺตํ วีตโทสํ จิตฺตนฺติ ปชาเนยฺยํ.\csromanspacer{} สโมหํ วา จิตฺตํ สโมหํ จิตฺตนฺติ ปชาเนยฺยํ, วีตโมหํ วา จิตฺตํ วีตโมหํ จิตฺตนฺติ ปชาเนยฺยํ.\csromanspacer{} สํขิตฺตํ วา จิตฺตํ สํขิตฺตํ จิตฺตนฺติ ปชาเนยฺยํ, วิกฺขิตฺตํ วา จิตฺตํ วิกฺขิตฺตํ จิตฺตนฺติ ปชาเนยฺยํ.\csromanspacer{} มหคฺคตํ วา จิตฺตํ}

\vspace{\tipitakaparskip}
\csromancenter{มหาวจฺฉสุตฺต (73) มหคฺคตํ จิตฺตนฺติ ปชาเนยฺยํ, อมหคฺคตํ วา จิตฺตํ อมหคฺคตํ จิตฺตนฺติ ปชาเนยฺยํ.\csromanspacer{} ส-อุตฺตรํ วา จิตฺตํ ส-อุตฺตรํ จิตฺตนฺติ ปชาเนยฺยํ, อนุตฺตรํ วา จิตฺตํ อนุตฺตรํ จิตฺตนฺติ ปชาเนยฺยํ.\csromanspacer{} สมาหิตํ วา จิตฺตํ สมาหิตํ จิตฺตนฺติ ปชาเนยฺยํ, อสมาหิตํ วา จิตฺตํ อสมาหิตํ จิตฺตนฺติ ปชาเนยฺยํ.\csromanspacer{} วิมุตฺตํ วา จิตฺตํ วิมุตฺตํ จิตฺตนฺติ ปชาเนยฺยํ, อวิมุตฺตํ วา จิตฺตํ อวิมุตฺตํ จิตฺตนฺติ ปชาเนยฺยนฺ”ติ.\csromanspacer{} ตตฺร ตเตฺรว สกฺขิภพฺพตํ ปาปุณิสฺสสิ สติ สติ-อายตเน. (3)}
\prose{\csromanfolio{163}โส ตฺวํ วจฺฉ ยาวเทว อากงฺขิสฺสสิ “อเนกวิหิตํ ปุพฺเพนิวาสํ อนุสฺสเรยฺยํ.\csromanspacer{} เสยฺยถิทํ, เอกมฺปิ ชาติํ เทฺวปิ ชาติโย ติสฺโสปิ ชาติโย จตสฺโสปิ ชาติโย ปญฺจปิ ชาติโย ทสปิ ชาติโย วีสมฺปิ ชาติโย ติํสมฺปิ ชาติโย จตฺตาลีสมฺปิ ชาติโย ปญฺญาสมฺปิ ชาติโย ชาติสตมฺปิ ชาติสหสฺสมฺปิ ชาติสตสหสฺสมฺปิ อเนเกปิ สํวฏฺฏกปฺเป อเนเกปิ วิวฏฺฏกปฺเป อเนเกปิ สํวฏฺฏวิวฏฺฏกปฺเป, ‘อมุตฺราสิํ เอวํนาโม เอวํโคตฺโต เอวํวณฺโณ เอวมาหาโร เอวํสุขทุกฺขปฺปฏิสํเวที เอวมายุปริยนฺโต, โส ตโต จุโต อมุตฺร อุทปาทิํ.\csromanspacer{} ตตฺราปาสิํ เอวํนาโม เอวํโคตฺโต เอวํวณฺโณ เอวมาหาโร เอวํสุขทุกฺขปฺปฏิสํเวที เอวมายุปริยนฺโต, โส ตโต จุโต อิธูปปนฺโน’ติ.\csromanspacer{} อิติ สาการํ ส-อุทฺเทสํ อเนกวิหิตํ ปุพฺเพนิวาสํ อนุสฺสเรยฺยนฺ”ติ.\csromanspacer{} ตตฺร ตเตฺรว สกฺขิภพฺพตํ ปาปุณิสฺสสิ สติ สติ- อายตเน. (4)}

\prose{โส ตฺวํ วจฺฉ ยาวเทว อากงฺขิสฺสสิ “ทิพฺเพน จกฺขุนา วิสุทฺเธน อติกฺกนฺตมานุสเกน สตฺเต ปสฺเสยฺยํ จวมาเน อุปปชฺชมาเน หีเน ปณีเต สุวณฺเณ ทุพฺพณฺเณ สุคเต ทุคฺคเต, ยถากมฺมูปเค สตฺเต ปชาเนยฺยํ ‘อิเม วต โภนฺโต สตฺตา กายทุจฺจริเตน สมนฺนาคตา วจีทุจฺจริเตน สมนฺนาคตา มโนทุจฺจริเตน สมนฺนาคตา อริยานํ อุปวาทกา มิจฺฉาทิฏฺฐิกา มิจฺฉาทิฏฺฐิกมฺมสมาทานา, เต กายสฺส เภทา ปรํ มรณา อปายํ ทุคฺคติํ วินิปาตํ นิรยํ อุปปนฺนา.\csromanspacer{} อิเม วา ปน โภนฺโต สตฺตา กายสุจริเตน สมนฺนาคตา วจีสุจริเตน สมนฺนาคตา มโนสุจริเตน สมนฺนาคตา อริยานํ อนุปวาทกา สมฺมาทิฏฺฐิกา สมฺมาทิฏฺฐิกมฺมสมาทานา, เต กายสฺส เภทา ปรํ มรณา สุคติํ สคฺคํ โลกํ อุปปนฺนา’ติ.\csromanspacer{} อิติ ทิพฺเพน จกฺขุนา วิสุทฺเธน อติกฺกนฺตมานุสเกน \csromanfolio{164}สตฺเต ปสฺเสยฺยํ จวมาเน อุปปชฺชมาเน หีเน ปณีเต สุวณฺเณ ทุพฺพณฺเณ สุคเต ทุคฺคเต, ยถากมฺมูปเค สตฺเต ปชาเนยฺยนฺ”ติ.\csromanspacer{} ตตฺร ตเตฺรว สกฺขิภพฺพตํ ปาปุณิสฺสสิ สติ สติ-อายตเน. (5)}

\prose{โส ตฺวํ วจฺฉ ยาวเทว อากงฺขิสฺสสิ “อาสวานํ ขยา อนาสวํ เจโตวิมุตฺติํ ปญฺญาวิมุตฺติํ ทิฏฺเฐว ธมฺเม สยํ อภิญฺญา สจฺฉิกตฺวา อุปสมฺปชฺช วิหเรยฺยนฺ”ติ.\csromanspacer{} ตตฺร ตเตฺรว สกฺขิภพฺพตํ ปาปุณิสฺสสิ สติ สติ-อายตเนติ. (6)}

\proseitem{199}{อถ โข อายสฺมา วจฺฉโคตฺโต ภควโต ภาสิตํ อภินนฺทิตฺวา อนุโมทิตฺวา อุฏฺฐายาสนา ภควนฺตํ อภิวาเทตฺวา ปทกฺขิณํ กตฺวา ปกฺกามิ.\csromanspacer{} อถ โข อายสฺมา วจฺฉโคตฺโต เอโก วูปกฏฺโฐ อปฺปมตฺโต อาตาปี ปหิตตฺโต วิหรนฺโต นจิรสฺเสว, ยสฺสตฺถาย กุลปุตฺตา สมฺมเทว อคารสฺมา อนคาริยํ ปพฺพชนฺติ.\csromanspacer{} ตทนุตฺตรํ พฺรหฺมจริยปริโยสานํ ทิฏฺเฐว ธมฺเม สยํ อภิญฺญา สจฺฉิกตฺวา อุปสมฺปชฺช วิหาสิ, “ขีณา ชาติ, วุสิตํ พฺรหฺมจริยํ, กตํ กรณียํ, นาปรํ อิตฺถตฺตายา”ติ อพฺภญฺญาสิ.\csromanspacer{} อญฺญตโร โข ปนายสฺมา วจฺฉโคตฺโต อรหตํ อโหสิ.}

\proseitem{200}{เตน โข ปน สมเยน สมฺพหุลา ภิกฺขู ภควนฺตํ ทสฺสนาย คจฺฉนฺติ.\csromanspacer{} อทฺทสา โข อายสฺมา วจฺฉโคตฺโต เต ภิกฺขู ทูรโตว อาคจฺฉนฺเต, ทิสฺวาน เยน เต ภิกฺขู เตนุปสงฺกมิ, อุปสงฺกมิตฺวา เต ภิกฺขู เอตทโวจ “หนฺท กหํ ปน ตุเมฺห อายสฺมนฺโต คจฺฉถา”ติ.\csromanspacer{} ภควนฺตํ โข มยํ อาวุโส ทสฺสนาย คจฺฉามาติ.\csromanspacer{} เตนหายสฺมนฺโต มม วจเนน ภควโต ปาเท สิรสา วนฺทถ, เอวญฺจ วเทถ “วจฺฉโคตฺโต ภนฺเต ภิกฺขุ ภควโต ปาเท สิรสา วนฺทติ, เอวญฺจ วเทติ ‘ปริจิณฺโณ เม ภควา, ปริจิณฺโณ เม สุคโต’ติ”. “เอวมาวุโส”ติ โข เต ภิกฺขู อายสฺมโต วจฺฉโคตฺตสฺส ปจฺจสฺโสสุํ.\csromanspacer{} อถ โข เต ภิกฺขู เยน ภควา เตนุปสงฺกมิํสุ, อุปสงฺกมิตฺวา ภควนฺตํ อภิวาเทตฺวา เอกมนฺตํ นิสีทิํสุ, เอกมนฺตํ นิสินฺนา โข เต ภิกฺขู ภควนฺตํ เอตทโวจุํ “อายสฺมา ภนฺเต วจฺฉโคตฺโต ภควโต ปาเท สิรสา วนฺทติ, เอวญฺจ วเทติ ‘ปริจิณฺโณ เม ภควา, ปริจิณฺโณ เม สุคโต’ติ”.\csromanspacer{} ปุพฺเพว เม ภิกฺขเว วจฺฉโคตฺโต ภิกฺขุ เจตสา เจโต ปริจฺจ วิทิโต “เตวิชฺโช}

\vspace{\tipitakaparskip}
\csromancenter{ทีฆนขสุตฺต (74) วจฺฉโคตฺโต ภิกฺขุ มหิทฺธิโก มหานุภาโว”ติ.\csromanspacer{} เทวตาปิ เม เอตมตฺถํ อาโรเจสุํ “เตวิชฺโช ภนฺเต วจฺฉโคตฺโต ภิกฺขุ มหิทฺธิโก มหานุภาโว”ติ.}
\csromancenter{อิทมโวจ ภควา.\csromanspacer{} อตฺตมนา เต ภิกฺขู ภควโต ภาสิตํ อภินนฺทุนฺติ.}
\nitthitamruled{มหาวจฺฉสุตฺตํ นิฏฺฐิตํ ตติยํ.}
\csromanmatikaanchor{mtk.29}
\csromantocmarkref{subsection}{4. ทีฆนขสุตฺต}{mtk.29}
\bookmark[dest={mtk.29},level=2]{4. ทีฆนขสุตฺต}
\csromanheader{2}{4. ทีฆนขสุตฺต}
\proseitem{201}{\csromanfolio{165}เอวํ เม สุตํ\csromandash{}เอกํ สมยํ ภควา ราชคเห วิหรติ คิชฺฌกูเฏ ปพฺพเต สูกรขตายํ.\csromanspacer{} อถ โข ทีฆนโข ปริพฺพาชโก เยน ภควา เตนุปสงฺกมิ, อุปสงฺกมิตฺวา ภควตา สทฺธิํ สมฺโมทิ, สมฺโมทนียํ กถํ สารณียํ วีติสาเรตฺวา เอกมนฺตํ อฏฺฐาสิ, เอกมนฺตํ ฐิโต โข ทีฆนโข ปริพฺพาชโก ภควนฺตํ เอตทโวจ “อหํ หิ โภ โคตม เอวํวาที เอวํทิฏฺฐิ สพฺพํ เม นกฺขมตี”ติ.\csromanspacer{} ยาปิ โข เต เอสา อคฺคิเวสฺสน ทิฏฺฐิ “สพฺพํ เม นกฺขมตี”ติ, เอสาปิ เต ทิฏฺฐิ นกฺขมตีติ.\csromanspacer{} เอสา เจ\footnote{เอสาปิ (ก)} เม โภ โคตม ทิฏฺฐิ ขเมยฺย “ตํปสฺส ตาทิสเมว, ตํปสฺส ตาทิสเมวา”ติ.\csromanspacer{} อโต โข เต อคฺคิเวสฺสน พหู หิ พหุตรา โลกสฺมิํ, เย เอวมาหํสุ “ตํปสฺส ตาทิสเมว, ตํปสฺส ตาทิสเมวา”ติ, เต ตํ เจว ทิฏฺฐิํ นปฺปชหนฺติ, อญฺญญฺจ ทิฏฺฐิํ อุปาทิยนฺติ.\csromanspacer{} อโต โข เต อคฺคิเวสฺสน ตนู หิ ตนุตรา โลกสฺมิํ, เย เอวมาหํสุ “ตํปสฺส ตาทิสเมว, ตํปสฺส ตาทิสเมวา”ติ, เต ตญฺเจว ทิฏฺฐิํ ปชหนฺติ, อญฺญญฺจ ทิฏฺฐิํ น อุปาทิยนฺติ.\csromanspacer{} สนฺตคฺคิเวสฺสน เอเก สมณพฺราหฺมณา เอวํวาทิโน เอวํทิฏฺฐิโน “สพฺพํ เม ขมตี”ติ.\csromanspacer{} สนฺตคฺคิเวสฺสน เอเก สมณพฺราหฺมณา เอวํวาทิโน เอวํทิฏฺฐิโน “สพฺพํ เม นกฺขมตี”ติ.\csromanspacer{} สนฺตคฺคิเวสฺสน เอเก สมณพฺราหฺมณา เอวํวาทิโน เอวํทิฏฺฐิโน “เอกจฺจํ เม ขมติ, เอกจฺจํ เม นกฺขมตี”ติ.\csromanspacer{} ตตฺรคฺคิเวสฺสน เย เต สมณพฺราหฺมณา เอวํวาทิโน เอวํทิฏฺฐิโน “สพฺพํ เม ขมตี”ติ.\csromanspacer{} เตสมยํ ทิฏฺฐิ สาราคาย สนฺติเก สญฺโญคาย สนฺติเก อภินนฺทนาย สนฺติเก อชฺโฌสานาย \csromanfolio{166}สนฺติเก อุปาทานาย สนฺติเก.\csromanspacer{} ตตฺรคฺคิเวสฺสน เย เต สมณพฺราหฺมณา เอวํวาทิโน เอวํทิฏฺฐิโน “สพฺพํ เม นกฺขมตี”ติ.\csromanspacer{} เตสมยํ ทิฏฺฐิ อสาราคาย สนฺติเก อสญฺโญคาย สนฺติเก อนภินนฺทนาย สนฺติเก อนชฺโฌสานาย สนฺติเก อนุปาทานาย สนฺติเกติ.}

\proseitem{202}{เอวํ วุตฺเต ทีฆนโข ปริพฺพาชโก ภควนฺตํ เอตทโวจ “อุกฺกํเสติ\footnote{อุกฺกํสติ (สี, อิ, ก)} เม ภวํ โคตโม ทิฏฺฐิคตํ, สมุกฺกํเสติ\footnote{สมฺปหํสติ (ก)} เม ภวํ โคตโม ทิฏฺฐิคตนฺ”ติ.\csromanspacer{} ตตฺรคฺคิเวสฺสน เย เต สมณพฺราหฺมณา เอวํวาทิโน เอวํทิฏฺฐิโน “เอกจฺจํ เม ขมติ, เอกจฺจํ เม นกฺขมตี”ติ.\csromanspacer{} ยา หิ เตสํ ขมติ, สายํ ทิฏฺฐิ สาราคาย สนฺติเก สญฺโญคาย สนฺติเก อภินนฺทนาย สนฺติเก อชฺโฌสานาย สนฺติเก อุปาทานาย สนฺติเก.\csromanspacer{} ยา หิ เตสํ นกฺขมติ, สายํ ทิฏฺฐิ อสาราคาย สนฺติเก อสญฺโญคาย สนฺติเก อนภินนฺทนาย สนฺติเก อนชฺโฌสานาย สนฺติเก อนุปาทานาย สนฺติเก.\csromanspacer{} ตตฺรคฺคิเวสฺสน เย เต สมณพฺราหฺมณา เอวํวาทิโน เอวํทิฏฺฐิโน “สพฺพํ เม ขมตี”ติ.\csromanspacer{} ตตฺถ วิญฺญู ปุริโส อิติ ปฏิสญฺจิกฺขติ\csromandash{}ยา โข เม อยํ ทิฏฺฐิ “สพฺพํ เม ขมตี”ติ, อิมญฺเจ อหํ ทิฏฺฐิํ ถามสา ปรามาสา อภินิวิสฺส โวหเรยฺยํ “อิทเมว สจฺจํ โมฆมญฺญนฺ”ติ.\csromanspacer{} ทฺวีหิ เม อสฺส วิคฺคโห, โย จายํ สมโณ วา พฺราหฺมโณ วา เอวํวาที เอวํทิฏฺฐิ “สพฺพํ เม นกฺขมตี”ติ, โย จายํ สมโณ วา พฺราหฺมโณ วา เอวํวาที เอวํทิฏฺฐิ “เอกจฺจํ เม ขมติ เอกจฺจํ เม นกฺขมตี”ติ.\csromanspacer{} อิเมหิ อสฺส ทฺวีหิ วิคฺคโห.\csromanspacer{} อิติ วิคฺคเห สติ วิวาโท, วิวาเท สติ วิฆาโต, วิฆาเต สติ วิเหสา.\csromanspacer{} อิติ โส วิคฺคหญฺจ วิวาทญฺจ วิฆาตญฺจ วิเหสญฺจ อตฺตนิ สมฺปสฺสมาโน ตญฺเจว ทิฏฺฐิํ ปชหติ, อญฺญญฺจ ทิฏฺฐิํ น อุปาทิยติ.\csromanspacer{} เอวเมตาสํ ทิฏฺฐีนํ ปหานํ โหติ, เอวเมตาสํ ทิฏฺฐีนํ ปฏินิสฺสคฺโค โหติ.}

\proseitem{203}{ตตฺรคฺคิเวสฺสน เย เต สมณพฺราหฺมณา เอวํวาทิโน เอวํทิฏฺฐิโน “สพฺพํ เม นกฺขมตี”ติ.\csromanspacer{} ตตฺถ วิญฺญู ปุริโส อิติ ปฏิสญฺจิกฺขติ\csromandash{}ยา โข เม อยํ ทิฏฺฐิ “สพฺพํ เม นกฺขมตี”ติ, อิมญฺเจ อหํ ทิฏฺฐิํ ถามสา ปรามาสา อภินิวิสฺส โวหเรยฺยํ “อิทเมว สจฺจํ โมฆมญฺญนฺ”ติ.\csromanspacer{} ทฺวีหิ เม อสฺส วิคฺคโห, โย จายํ สมโณ วา พฺราหฺมโณ วา เอวํวาที}

\vspace{\tipitakaparskip}
\csromancenter{ทีฆนขสุตฺต (74) เอวํทิฏฺฐิ “สพฺพํ เม ขมตี”ติ, โย จายํ สมโณ วา พฺราหฺมโณ วา เอวํวาที เอวํทิฏฺฐิ “เอกจฺจํ เม ขมติ เอกจฺจํ เม นกฺขมตี”ติ.\csromanspacer{} อิเมหิ อสฺส ทฺวีหิ วิคฺคโห.\csromanspacer{} อิติ วิคฺคเห สติ วิวาโท, วิวาเท สติ วิฆาโต, วิฆาเต สติ วิเหสา.\csromanspacer{} อิติ โส วิคฺคหญฺจ วิวาทญฺจ วิฆาตญฺจ วิเหสญฺจ อตฺตนิ สมฺปสฺสมาโน ตญฺเจว ทิฏฺฐิํ ปชหติ, อญฺญญฺจ ทิฏฺฐิํ น อุปาทิยติ.\csromanspacer{} เอวเมตาสํ ทิฏฺฐีนํ ปหานํ โหติ, เอวเมตาสํ ทิฏฺฐีนํ ปฏินิสฺสคฺโค โหติ.}
\proseitem{204}{\csromanfolio{167}ตตฺรคฺคิเวสฺสน เย เต สมณพฺราหฺมณา เอวํวาทิโน เอวํทิฏฺฐิโน “เอกจฺจํ เม ขมติ เอกจฺจํ เม นกฺขมตี”ติ.\csromanspacer{} ตตฺถ วิญฺญู ปุริโส อิติ ปฏิสญฺจิกฺขทิ\csromandash{}ยา โข เม อยํ ทิฏฺฐิ “เอกจฺจํ เม ขมติ เอกจฺจํ เม นกฺขมตี”ติ, อิมญฺจ อหํ ทิฏฺฐิํ ถามสา ปรามาสา อภินิวิสฺส โวหเรยฺยํ “อิทเมว สจฺจํ โมฆมญฺญนฺ”ติ.\csromanspacer{} ทฺวีหิ เม อสฺส วิคฺคโห, โย จายํ สมโณ วา พฺราหฺมโณ วา เอวํวาที เอวํทิฏฺฐิ “สพฺพํ เม ขมตี”ติ, โย จายํ สมโณ วา พฺราหฺมโณ วา เอวํวาที เอวํทิฏฺฐิ “สพฺพํ เม นกฺขมตี”ติ.\csromanspacer{} อิเมหิ อสฺส ทฺวีหิ วิคฺคโห.\csromanspacer{} อิติ วิคฺคเห สติ วิวาโท, วิวาเท สติ วิฆาโต, วิฆาเต สติ วิเหสา.\csromanspacer{} อิติ โส วิคฺคหญฺจ วิวาทญฺจ วิฆาตญฺจ วิเหสญฺจ อตฺตนิ สมฺปสฺสมาโน ตญฺเจว ทิฏฺฐิํ ปชหติ, อญฺญญฺจ ทิฏฺฐิํ น อุปาทิยติ.\csromanspacer{} เอวเมตาสํ ทิฏฺฐีนํ ปหานํ โหติ, เอวเมตาสํ ทิฏฺฐีนํ ปฏินิสฺสคฺโค โหติ.}

\proseitem{205}{อยํ โข ปนคฺคิเวสฺสน กาโย รูปี จาตุมหาภูติโก\footnote{จาตุมฺมหาภูติโก (สี, สฺยา)} มาตาเปตฺติกสมฺภโว โอทนกุมฺมาสุปจโย อนิจฺจุจฺฉาทนปริมทฺทนเภทนวิทฺธํสนธมฺโม อนิจฺจโต ทุกฺขโต โรคโต คณฺฑโต สลฺลโต อฆโต อาพาธโต ปรโต ปโลกโต สุญฺญโต อนตฺตโต สมนุปสฺสิตพฺโพ, ตสฺสิมํ กายํ อนิจฺจโต ทุกฺขโต โรคโต คณฺฑโต สลฺลโต อฆโต อาพาธโต ปรโต ปโลกโต สุญฺญโต อนตฺตโต สมนุปสฺสโต โย กายสฺมิํ กายฉนฺโท กายสฺเนโห กายนฺวยตา, สา ปหียติ.}

\prose{ติสฺโส โข อิมา อคฺคิเวสฺสน เวทนา, สุขา เวทนา ทุกฺขา เวทนา อทุกฺขมสุขา เวทนา.\csromanspacer{} ยสฺมิํ อคฺคิเวสฺสน สมเย สุขํ เวทนํ \csromanfolio{168}เวเทติ, เนว ตสฺมิํ สมเย ทุกฺขํ เวทนํ เวเทติ, น อทุกฺขมสุขํ เวทนํ เวเทติ, สุขํเยว ตสฺมิํ สมเย เวทนํ เวเทติ.\csromanspacer{} ยสฺมิํ อคฺคิเวสฺสน สมเย ทุกฺขํ เวทนํ เวเทติ, เนว ตสฺมิํ สมเย สุขํ เวทนํ เวเทติ, น อทุกฺขมสุขํ เวทนํ เวเทติ, ทุกฺขํเยว ตสฺมิํ สมเย เวทนํ เวเทติ.\csromanspacer{} ยสฺมิํ อคฺคิเวสฺสน สมเย อทุกฺขมสุขํ เวทนํ เวเทติ, เนว ตสฺมิํ สมเย สุขํ เวทนํ เวเทติ, น ทุกฺขํ เวทนํ เวเทติ, อทุกฺขมสุขํเยว ตสฺมิํ สมเย เวทนํ เวเทติ.\csromanspacer{} สุขาปิ โข อคฺคิเวสฺสน เวทนา อนิจฺจา สงฺขตา ปฏิจฺจสมุปฺปนฺนา ขยธมฺมา วยธมฺมา วิราคธมฺมา นิโรธธมฺมา, ทุกฺขาปิ โข อคฺคิเวสฺสน เวทนา อนิจฺจา สงฺขตา ปฏิจฺจสมุปฺปนฺนา ขยธมฺมา วยธมฺมา วิราคธมฺมา นิโรธธมฺมา, อทุกฺขมสุขาปิ โข อคฺคิเวสฺสน เวทนา อนิจฺจา สงฺขตา ปฏิจฺจสมุปฺปนฺนา ขยธมฺมา วยธมฺมา วิราคธมฺมา นิโรธธมฺมา.\csromanspacer{} เอวํ ปสฺสํ อคฺคิเวสฺสน สุตวา อริยสาวโก สุขายปิ เวทนาย นิพฺพินฺทติ, ทุกฺขายปิ เวทนาย นิพฺพินฺทติ, อทุกฺขมสุขายปิ เวทนาย นิพฺพินฺทติ, นิพฺพินฺทํ วิรชฺชติ, วิราคา วิมุจฺจติ, วิมุตฺตสฺมิํ วิมุตฺตมิติ ญาณํ โหติ, “ขีณา ชาติ, วุสิตํ พฺรหฺมจริยํ, กตํ กรณียํ, นาปรํ อิตฺถตฺตายา”ติ ปชานาติ.\csromanspacer{} เอวํ วิมุตฺตจิตฺโต โข อคฺคิเวสฺสน ภิกฺขุ น เกนจิ สํวทติ, น เกนจิ วิวทติ.\csromanspacer{} ยญฺจ โลเก วุตฺตํ, เตน โวหรติ อปรามสนฺติ.}

\proseitem{206}{เตน โข ปน สมเยน อายสฺมา สาริปุตฺโต ภควโต ปิฏฺฐิโต ฐิโต โหติ ภควนฺตํ พีชยมาโน\footnote{วีชยมาโน (สี, อิ)}.\csromanspacer{} อถ โข อายสฺมโต สาริปุตฺตสฺส เอตทโหสิ “เตสํ เตสํ กิร โน ภควา ธมฺมานํ อภิญฺญา ปหานมาห, เตสํ เตสํ กิร โน สุคโต ธมฺมานํ อภิญฺญา ปฏินิสฺสคฺคมาหา”ติ.\csromanspacer{} อิติ หิทํ อายสฺมโต สาริปุตฺตสฺส ปฏิสญฺจิกฺขโต อนุปาทาย อาสเวหิ จิตฺตํ วิมุจฺจิ.\csromanspacer{} ทีฆนขสฺส ปน ปริพฺพาชกสฺส วิรชํ วีตมลํ ธมฺมจกฺขุํ อุทปาทิ “ยํ กิญฺจิ สมุทยธมฺมํ, สพฺพํ ตํ นิโรธธมฺมนฺ”ติ.\csromanspacer{} อถ โข ทีฆนโข ปริพฺพาชโก ทิฏฺฐธมฺโม ปตฺตธมฺโม วิทิตธมฺโม ปริโยคาฬฺหธมฺโม ติณฺณวิจิกิจฺโฉ วิคตกถํกโถ เวสารชฺชปฺปตฺโต อปรปฺปจฺจโย สตฺถุสาสเน ภควนฺตํ เอตทโวจ “อภิกฺกนฺตํ โภ โคตม, อภิกฺกนฺตํ โภ โคตม, เสยฺยถาปิ โภ โคตม นิกฺกุชฺชิตํ วา อุกฺกุชฺเชยฺย, ปฏิจฺฉนฺนํ วา วิวเรยฺย, มูฬฺหสฺส วา}

\vspace{\tipitakaparskip}
\csromancenter{มาคณฺฑิยสุตฺต (75) มคฺคํ อาจิกฺเขยฺย, อนฺธกาเร วา เตลปชฺโชตํ ธาเรยฺย ‘จกฺขุมนฺโต รูปานิ ทกฺขนฺตี’ติ, เอวเมว โข โภตา โคตเมน อเนกปริยาเยน ธมฺโม ปกาสิโต.\csromanspacer{} เอสาหํ ภวนฺตํ โคตมํ สรณํ คจฺฉามิ ธมฺมญฺจ ภิกฺขุสํฆญฺจ, อุปาสกํ มํ ภวํ โคตโม ธาเรตุ อชฺชตคฺเค ปาณุเปตํ สรณํ คตนฺ”ติ.}
\nitthitamruled{ทีฆนขสุตฺตํ นิฏฺฐิตํ จตุตฺถํ.}
\csromanmatikaanchor{mtk.30}
\csromantocmarkref{subsection}{5. มาคณฺฑิยสุตฺต}{mtk.30}
\bookmark[dest={mtk.30},level=2]{5. มาคณฺฑิยสุตฺต}
\csromanheader{2}{5. มาคณฺฑิยสุตฺต}
\proseitem{207}{\csromanfolio{169}เอวํ เม สุตํ\csromandash{}เอกํ สมยํ ภควา กุรูสุ วิหรติ กมฺมาสธมฺมํ นาม กุรูนํ นิคโม ภารทฺวาชโคตฺตสฺส พฺราหฺมณสฺส อคฺยาคาเร ติณสนฺถารเก\footnote{ติณสนฺถรเก (สี, สฺยา, กํ, อิ)}.\csromanspacer{} อถ โข ภควา ปุพฺพณฺหสมยํ นิวาเสตฺวา ปตฺตจีวรมาทาย กมฺมาสธมฺมํ ปิณฺฑาย ปาวิสิ, กมฺมาสธมฺมํ ปิณฺฑาย จริตฺวา ปจฺฉาภตฺตํ ปิณฺฑปาตปฏิกฺกนฺโต เยน อญฺญตโร วนสณฺโฑ เตนุปสงฺกมิ ทิวาวิหาราย, ตํ วนสณฺฑํ อชฺโฌคาเหตฺวา อญฺญตรสฺมิํ รุกฺขมูเล ทิวาวิหารํ นิสีทิ.\csromanspacer{} อถ โข มาคณฺฑิโย\footnote{มาคนฺทิโย (สี, อิ)} ปริพฺพาชโก ชงฺฆาวิหารํ อนุจงฺกมมาโน อนุวิจรมาโน เยน ภารทฺวาชโคตฺตสฺส พฺราหฺมณสฺส อคฺยาคารํ เตนุปสงฺกมิ, อทฺทสา โข มาคณฺฑิโย ปริพฺพาชโก ภารทฺวาชโคตฺตสฺส พฺราหฺมณสฺส อคฺยาคาเร ติณสนฺถารกํ ปญฺญตฺตํ, ทิสฺวาน ภารทฺวาชโคตฺตํ พฺราหฺมณํ เอตทโวจ “กสฺส นฺวยํ โภโต ภารทฺวาชสฺส อคฺยาคาเร ติณสนฺถารโก ปญฺญตฺโต, สมณเสยฺยานุรูปํ\footnote{สมณเสยฺยารูปํ (สี, อิ)} มญฺเญ”ติ.\csromanspacer{} อตฺถิ โภ มาคณฺฑิย สมโณ โคตโม สกฺยปุตฺโต สกฺยกุลา ปพฺพชิโต, ตํ โข ปน ภวนฺตํ โคตมํ เอวํ กลฺยาโณ กิตฺติสทฺโท อพฺภุคฺคโต “อิติปิ โส ภควา อรหํ สมฺมาสมฺพุทฺโธ วิชฺชาจรณสมฺปนฺโน สุคโต โลกวิทู อนุตฺตโร ปุริสทมฺมสารถิ สตฺถา เทวมนุสฺสานํ พุทฺโธ ภควา”ติ, ตสฺเสสา โภโต โคตมสฺส เสยฺยา ปญฺญตฺตาติ.\csromanspacer{} ทุทฺทิฏฺฐํ วต โภ ภารทฺวาช อทฺทสาม, ทุทฺทิฏฺฐํ วต โภ ภารทฺวาช อทฺทสาม, เย มยํ ตสฺส โภโต โคตมสฺส ภูนหุโน\footnote{ภูนหนสฺส (สฺยา, กํ)} เสยฺยํ อทฺทสามาติ.\csromanspacer{} รกฺขสฺเสตํ มาคณฺฑิย วาจํ, \csromanfolio{170}รกฺขสฺเสตํ มาคณฺฑิย วาจํ, พหู หิ ตสฺส โภโต โคตมสฺส ขตฺติยปณฺฑิตาปิ พฺราหฺมณปณฺฑิตาปิ คหปติปณฺฑิตาปิ สมณปณฺฑิตาปิ อภิปฺปสนฺนา วินีตา อริเย ญาเย ธมฺเม กุสเลติ.\csromanspacer{} สมฺมุขา เจปิ มยํ โภ ภารทฺวาช ตํ ภวนฺตํ โคตมํ ปสฺเสยฺยาม.\csromanspacer{} สมฺมุขาปิ นํ วเทยฺยาม “ภูนหุ\footnote{ภูนหโน (สฺยา, กํ)} สมโณ โคตโม”ติ.\csromanspacer{} ตํ กิสฺส เหตุ.\csromanspacer{} เอวํ หิ โน สุตฺเต โอจรตีติ.\csromanspacer{} สเจ ตํ โภโต มาคณฺฑิยสฺส อครุ, อาโรเจยฺยามิ ตํ\footnote{อาโรเจยฺยเมตํ (สี, อิ), อาโรเจสฺสามิ ตสฺส (สฺยา, กํ)} สมณสฺส โคตมสฺสาติ.\csromanspacer{} อปฺโปสฺสุกฺโก ภวํ ภารทฺวาโช วุตฺโตว นํ วเทยฺยาติ.}

\proseitem{208}{อสฺโสสิ โข ภควา ทิพฺพาย โสตธาตุยา วิสุทฺธาย อติกฺกนฺตมานุสิกาย ภารทฺวาชโคตฺตสฺส พฺราหฺมณสฺส มาคณฺฑิเยน ปริพฺพาชเกน สทฺธิํ อิมํ กถาสลฺลาปํ.\csromanspacer{} อถ โข ภควา สายนฺหสมยํ ปฏิสลฺลานา วุฏฺฐิโต เยน ภารทฺวาชโคตฺตสฺส พฺราหฺมณสฺส อคฺยาคารํ เตนุปสงฺกมิ, อุปสงฺกมิตฺวา นิสีทิ ภควา ปญฺญตฺเต ติณสนฺถารเก.\csromanspacer{} อถ โข ภารทฺวาชโคตฺโต พฺราหฺมโณ เยน ภควา เตนุปสงฺกมิ, อุปสงฺกมิตฺวา ภควตา สทฺธิํ สมฺโมทิ, สมฺโมทนียํ กถํ สารณียํ วีติสาเรตฺวา เอกมนฺตํ นิสีทิ, เอกมนฺตํ นิสินฺนํ โข ภารทฺวาชโคตฺตํ พฺราหฺมณํ ภควา เอตทโวจ “อหุ ปน เต ภารทฺวาช มาคณฺฑิเยน ปริพฺพาชเกน สทฺธิํ อิมํเยว ติณสนฺถารกํ อารพฺภ โกจิเทว กถาสลฺลาโป”ติ.\csromanspacer{} เอวํ วุตฺเต ภารทฺวาชโคตฺโต พฺราหฺมโณ สํวิคฺโค โลมหฏฺฐชาโต ภควนฺตํ เอตทโวจ “เอตเทว โข ปน มยํ โภโต โคตมสฺส อาโรเจตุกามา, อถ จ ปน ภวํ โคตโม อนกฺขาตํเยว อกฺขาสี”ติ.\csromanspacer{} อยญฺจ หิ\footnote{อยญฺจ หิทํ (สี, สฺยา, กํ, อิ)} ภควโต ภารทฺวาชโคตฺเตน พฺราหฺมเณน สทฺธิํ อนฺตรากถา วิปฺปกตา โหติ.\csromanspacer{} อถ โข มาคณฺฑิโย ปริพฺพาชโก ชงฺฆาวิหารํ อนุจงฺกมมาโน อนุวิจรมาโน เยน ภารทฺวาชโคตฺตสฺส พฺราหฺมณสฺส อคฺยาคารํ เยน ภควา เตนุปสงฺกมิ, อุปสงฺกมิตฺวา ภควตา สทฺธิํ สมฺโมทิ, สมฺโมทนียํ กถํ สารณียํ วีติสาเรตฺวา เอกมนฺตํ นิสีทิ.\csromanspacer{} เอกมนฺตํ นิสินฺนํ โข มาคณฺฑิยํ ปริพฺพาชกํ ภควา เอตทโวจ\csromandash{}}

\csromanheader{2}{5. มาคณฺฑิยสุตฺต (75)}
\proseitem{209}{\csromanfolio{171}“จกฺขุํ โข มาคณฺฑิย รูปารามํ รูปรตํ รูปสมฺมุทิตํ, ตํ ตถาคตสฺส ทนฺตํ คุตฺตํ รกฺขิตํ สํวุตํ, ตสฺส จ สํวราย ธมฺมํ เทเสติ, อิทํ นุ เต เอตํ มาคณฺฑิย สนฺธาย ภาสิตํ ภูนหุ สมโณ โคตโม”ติ.\csromanspacer{} เอตเทว โข ปน เม โภ โคตม สนฺธาย ภาสิตํ ภูนหุ สมโณ โคตโมติ.\csromanspacer{} ตํ กิสฺส เหตุ.\csromanspacer{} เอวํ หิ โน สุตฺเต โอจรตีติ.\csromanspacer{} โสตํ โข มาคณฺฑิย สทฺทารามํ ฯเปฯ.\csromanspacer{} ฆานํ โข มาคณฺฑิย คนฺธารามํ.\csromanspacer{} ชิวฺหา โข มาคณฺฑิย รสารามา รสรตา รสสมฺมุทิตา, สา ตถาคตสฺส ทนฺตา คุตฺตา รกฺขิตา สํวุตา, ตสฺสา จ สํวราย ธมฺมํ เทเสติ, อิทํ นุเต เอตํ มาคณฺฑิย สนฺธาย ภาสิตํ ภูนหุ สมโณ โคตโมติ.\csromanspacer{} เอตเทว โข ปน เม โภ โคตม สนฺธาย ภาสิตํ ภูนหุ สมโณ โคตโมติ.\csromanspacer{} ตํ กิสฺส เหตุ.\csromanspacer{} เอวํ หิ โน สุตฺเต โอจรตีติ.\csromanspacer{} กาโย โข มาคณฺฑิย โผฏฺฐพฺพาราโม โผฏฺฐพฺพรโต ฯเปฯ.\csromanspacer{} มโน โข มาคณฺฑิย ธมฺมาราโม ธมฺมรโต ธมฺมสมฺมุทิโต, โส ตถาคตสฺส ทนฺโต คุตฺโต รกฺขิโต สํวุโต, ตสฺส จ สํวราย ธมฺมํ เทเสติ, อิทํ นุ เต เอตํ มาคณฺฑิย สนฺธาย ภาสิตํ ภูนหุ สมโณ โคตโมติ.\csromanspacer{} เอตเทว โข ปน เม โภ โคตม สนฺธาย ภาสิตํ ภูนหุ สมโณ โคตโมติ.\csromanspacer{} ตํ กิสฺส เหตุ.\csromanspacer{} เอวํ หิ โน สุตฺเต โอจรตีติ.}

\proseitem{210}{ตํ กิํ มญฺญสิ มาคณฺฑิย, อิเธกจฺโจ จกฺขุวิญฺเญยฺเยหิ รูเปหิ ปริจาริตปุพฺโพ อสฺส อิฏฺเฐหิ กนฺเตหิ มนาเปหิ ปิยรูเปหิ กามูปสํหิเตหิ รชนีเยหิ, โส อปเรน สมเยน รูปานํเยว สมุทยญฺจ อตฺถงฺคมญฺจ อสฺสาทญฺจ อาทีนวญฺจ นิสฺสรณญฺจ ยถาภูตํ วิทิตฺวา รูปตณฺหํ ปหาย รูปปริฬาหํ ปฏิวิโน เทตฺวา วิคตปิปาโส อชฺฌตฺตํ วูปสนฺตจิตฺโต วิหเรยฺย, อิมสฺส ปน เต มาคณฺฑิย กิมสฺส วจนียนฺติ.\csromanspacer{} น กิญฺจิ โภ โคตม.\csromanspacer{} ตํ กิํ มญฺญสิ มาคณฺฑิย, อิเธกจฺโจ โสตวิญฺเญยฺเยหิ สทฺเทหิ ฯเปฯ ฆานวิญฺเญยฺเยหิ คนฺเธหิ.\csromanspacer{} ชิวฺหาวิญฺเญยฺเยหิ รเสหิ.\csromanspacer{} กายวิญฺเญยฺเยหิ โผฏฺฐพฺเพหิ ปริจาริตปุพฺโพ อสฺส อิฏฺเฐหิ กนฺเตหิ มนาเปหิ ปิยรูเปหิ กามูปสํหิเตหิ รชนีเยหิ, โส อปเรน สมเยน โผฏฺฐพฺพานํเยว สมุทยญฺจ อตฺถงฺคมญฺจ อสฺสาทญฺจ-อาทีนวญฺจ นิสฺสรณญฺจ ยถาภูตํ วิทิตฺวา โผฏฺฐพฺพตณฺหํ ปหาย โผฏฺฐพฺพปริฬาหํ ปฏิวิโนเทตฺวา วิคตปิปาโส-อชฺฌตฺตํ วูปสนฺตจิตฺโต วิหเรยฺย, อิมสฺส ปน เต มาคณฺฑิย กิมสฺส วจนียนฺติ.\csromanspacer{} น กิญฺจิ โภ โคตม.}

\proseitem{211}{\csromanfolio{172}อหํ โข ปน มาคณฺฑิย ปุพฺเพ อคาริยภูโต สมาโน ปญฺจหิ กามคุเณหิ สมปฺปิโต สมงฺคีภูโต ปริจาเรสิํ จกฺขุวิญฺเญยฺเยหิ รูเปหิ อิฏฺเฐหิ กนฺเตหิ มนาเปหิ ปิยรูเปหิ กามูปสํหิเตหิ รชนีเยหิ, โสตวิญฺเญยฺเยหิ สทฺเทหิ ฯเปฯ ฆานวิญฺเญยฺเยหิ คนฺเธหิ.\csromanspacer{} ชิวฺหาวิญฺเญยฺเยหิ รเสหิ.\csromanspacer{} กายวิญฺเญยฺเยหิ โผฏฺฐพฺเพหิ อิฏฺฐหิ กนฺเตหิ มนาเปหิ ปิยรูเปหิ กามูปสํหิเตหิ รชนีเยติ, ตสฺส มยฺหํ มาคณฺฑิย ตโย ปาสาทา อเหสุํ, เอโก วสฺสิโก เอโก เหมนฺติโก เอโก คิมฺหิโก, โส โข อหํ มาคณฺฑิย วสฺสิเก ปาสาเท วสฺสิเก จตฺตาโร\footnote{วสฺสิเก ปาสาเท จตฺตาโร (สฺยา, กํ)} มาเส นิปฺปุริเสหิ ตูริเยหิ\footnote{ตุริเยหิ (สี, สฺยา, กํ, อิ)} ปริจารยมาโน\footnote{ปริจาริยมาโน (สพฺพตฺถ)} น เหฏฺฐาปาสาทํ โอโรหามิ, โส อปเรน สมเยน กามานํเยว สมุทยญฺจ อตฺถงฺคมญฺจ อสฺสาทญฺจ อาทีนวญฺจ นิสฺสรณญฺจ ยถาภูตํ วิทิตฺวา กามตณฺหํ ปหาย กามปริฬาหํ ปฏิวิโนเทตฺวา วิคตปิปาโส อชฺฌตฺตํ วูปสนฺตจิตฺโต วิหรามิ, โส อญฺเญ สตฺเต ปสฺสามิ กาเมสุ อวีตราเค กามตณฺหาหิ ขชฺชมาเน กามปริฬาเหน ปริฑยฺหมาเน กาเม ปฏิเสวนฺเต, โส เตสํ น ปิเหมิ, น ตตฺถ อภิรมามิ.\csromanspacer{} ตํ กิสฺส เหตุ.\csromanspacer{} ยาหยํ มาคณฺฑิย รติ อญฺญเตฺรว กาเมหิ อญฺญตฺร อกุสเลหิ ธมฺเมหิ อปิ ทิพฺพํ สุขํ สมธิคยฺห ติฏฺฐติ, ตาย รติยา รมมาโน หีนสฺส น ปิเหมิ, น ตตฺถ อภิรมามิ.}

\proseitem{212}{เสยฺยถาปิ มาคณฺฑิย คหปติ วา คหปติปุตฺโต วา อฑฺโฒ มหทฺธโน มหาโภโค ปญฺจหิ กามคุเณหิ สมปฺปิโต สมงฺคีภูโต ปริจาเรยฺย จกฺขุวิญฺเญยฺเยหิ รูเปหิ ฯเปฯ โผฏฺฐพฺเพหิ อิฏฺเฐหิ กนฺเตหิ มนาเปหิ ปิยรูเปหิ กามูปสํหิเตหิ รชนีเยหิ, โส กาเยน สุจริตํ จริตฺวา วาจาย สุจริตํ จริตฺวา มนสา สุจริตํ จริตฺวา กายสฺส เภทา ปรํ มรณา สุคติํ สคฺคํ โลกํ อุปปชฺเชยฺย เทวานํ ตาวติํสานํ สหพฺยตํ, โส ตตฺถ นนฺทเน วเน อจฺฉราสํฆปริวุโต ทิพฺเพหิ ปญฺจหิ กามคุเณหิ สมปฺปิโต สมงฺคีภูโต ปริจาเรยฺย, โส ปสฺเสยฺย คหปติํ วา คหปติปุตฺตํ วา ปญฺจหิ กามคุเณหิ สมปฺปิตํ สมงฺคีภูตํ ปริจารยมานํ.}

\csromanheader{2}{5. มาคณฺฑิยสุตฺต (75)}
\prose{\csromanfolio{173}ตํ กิํ มญฺญสิ มาคณฺฑิย, อปิ นุ โส เทวปุตฺโต นนฺทเน วเน อจฺฉราสํฆปริวุโต ทิพฺพหิ ปญฺจหิ กามคุเณหิ สมปฺปิโต สมงฺคีภูโต ปริจารยมาโน อมุสฺส คหปติสฺส วา คหปติปุตฺตสฺส วา ปิเหยฺย, มานุสกานํ วา ปญฺจนฺนํ กามคุณานํ, มานุสเกหิ วา กาเมหิ อาวฏฺเฏยฺยาติ.\csromanspacer{} โน หิทํ โภ โคตม.\csromanspacer{} ตํ กิสฺส เหตุ.\csromanspacer{} มานุสเกหิ โภ โคตม กาเมหิ ทิพฺพกามา อภิกฺกนฺตตรา จ ปณีตตรา จาติ.\csromanspacer{} เอวเมว โข อหํ มาคณฺฑิย ปุพฺเพ อคาริยภูโต สมาโน ปญฺจหิ กามคุเณหิ สมปฺปิโต สมงฺคีภูโต ปริจาเรสิํ จกฺขุวิญฺเญยฺเยหิ รูเปหิ อิฏฺเฐหิ กนฺเตหิ มนาเปหิ ปิยรูเปหิ กามูปสํหิเตหิ รชนียหิ, โสตวิญฺเญยฺเยหิ สทฺเทหิ ฯเปฯ ฆานวิญฺเญยฺเยหิ คนฺเธหิ.\csromanspacer{} ชิวฺหาวิญฺเญยฺเยหิ รเสหิ.\csromanspacer{} กายวิญฺเญยฺเยหิ โผฏฺฐพฺเพหิ อิฏฺเฐหิ กนฺเตหิ มนาเปหิ ปิยรูเปหิ กามูปสํหิเตหิ รชนีเยหิ, โส อปเรน สมเยน กามานํเยว สมุทยญฺจ อตฺถงฺคมญฺจ อสฺสาทญฺจ อาทีนวญฺจ นิสฺสรณญฺจ ยถาภูตํ วิทิตฺวา กามตณฺหํ ปหาย กามปริฬาหํ ปฏิวิโนเทตฺวา วิคตปิปาโส อชฺฌตฺตํ วูปสนฺตจิตฺโต วิหรามิ, โส อญฺเญ สตฺเต ปสฺสามิ กาเมสุ อวีตราเค กามตณฺหาหิ ขชฺชมาเน กามปริฬาเหน ปริฑยฺหมาเน กาเม ปฏิเสวนฺเต, โส เตสํ น ปิเหมิ, น ตตฺถ อภิรมามิ.\csromanspacer{} ตํ กิสฺส เหตุ.\csromanspacer{} ยาหยํ มาคณฺฑิย รติ อญฺญเตฺรว กาเมหิ อญฺญตฺร อกุสเลหิ ธมฺเมหิ อปิ ทิพฺพํ สุขํ สมธิคยฺห ติฏฺฐติ.\csromanspacer{} ตาย รติยา รมมาโน หีนสฺส น ปิเหมิ, น ตตฺถ อภิรมามิ.}

\proseitem{213}{เสยฺยถาปิ มาคณฺฑิย กุฏฺฐี ปุริโส อรุคตฺโต ปกฺกคตฺโต กิมีหิ ขชฺชมาโน นเขหิ วณมุขานิ วิปฺปตจฺฉมาโน องฺคารกาสุยา กายํ ปริตาเปยฺย.\csromanspacer{} ตสฺส มิตฺตามจฺจา ญาติสาโลหิตา ภิสกฺกํ สลฺลกตฺตํ อุปฐาเปยฺยุํ.\csromanspacer{} ตสฺส โส ภิสกฺโก สลฺลกตฺโต เภสชฺชํ กเรยฺย, โส ตํ เภสชฺชํ อาคมฺม กุฏฺเฐหิ ปริมุจฺเจยฺย, อโรโค อสฺส สุขี เสรี สยํวสี เยน กามํ คโม, โส อญฺญํ กุฏฺฐิํ ปุริสํ ปสฺเสยฺย อรุคตฺตํ ปกฺกคตฺตํ กิมีหิ ขชฺชมานํ นเขหิ วณมุขานิ วิปฺปตจฺฉมานํ องฺคารกาสุยา กายํ ปริตาเปนฺตํ.}

\prose{ตํ กิํ มญฺญสิ มาคณฺฑิย, อปิ นุ โส ปุริโส อมุสฺส กุฏฺฐิสฺส ปุริสสฺส ปิเหยฺย องฺคารกาสุยา วา เภสชฺชํ ปฏิเสวนาย วาติ. \csromanfolio{174}โน หิทํ โภ โคตม.\csromanspacer{} ตํ กิสฺส เหตุ.\csromanspacer{} โรเค หิ โภ โคตม สติ เภสชฺเชน กรณียํ โหติ, โรเค อสติ น เภสชฺเชน กรณียํ โหตีติ.\csromanspacer{} เอวเมว โข อหํ มาคณฺฑิย ปุพฺเพ อคาริยภูโต สมาโน ปญฺจหิ กามคุเณหิ สมปฺปิโต สมงฺคีภูโต ปริจาเรสิํ จกฺขุวิญฺเญยฺเยหิ รูเปหิ อิฏฺเฐหิ กนฺเตหิ มนาเปหิ ปิยรูเปหิ กามูปสํหิเตหิ รชนีเยหิ, โสตวิญฺเญยฺเยหิ สทฺเทหิ ฯเปฯ ฆานวิญฺเญยฺเยหิ คนฺเธหิ.\csromanspacer{} ชิวฺหาวิญฺเญยฺเยหิ รเสหิ.\csromanspacer{} กายวิญฺเญยฺเยหิ โผฏฺฐพฺเพหิ อิฏฺเฐหิ กนฺเตหิ มนาเปหิ กามูปสํหิเตหิ รชนีเยหิ, โส อปเรน สมเยน กามานํเยว สมุทยญฺจ อตฺถงฺคมญฺจ อสฺสาทญฺจ อาทีนวญฺจ นิสฺสรณญฺจ ยถาภูตํ วิทิตฺวา กามตณฺหํ ปหาย กามปริฬาหํ ปฏิวิโนเทตฺวา วิคตปิปาโส อชฺฌตฺตํ วูปสนฺตจิตฺโต วิหรามิ, โส อญฺเญ สตฺเต ปสฺสามิ กาเมสุ อวีตราเค กามตณฺหาหิ ขชฺชมาเน กามปริฬาเหน ปริฑยฺหมาเน กาเม ปฏิเสวนฺเต, โส เตสํ น ปิเหมิ, น ตตฺถ อภิรมามิ.\csromanspacer{} ตํ กิสฺส เหตุ.\csromanspacer{} ยาหยํ มาคณฺฑิย รติ อญฺญเตฺรว กาเมหิ อญฺญตฺร อกุสเลหิ ธมฺเมหิ อปิ ทิพฺพํ สุขํ สมธิคยฺห ติฏฺฐติ.\csromanspacer{} ตาย รติยา รมมาโน หีนสฺส น ปิเหมิ, น ตตฺถ อภิรมามิ.}

\proseitem{214}{เสยฺยถาปิ มาคณฺฑิย กุฏฺฐี ปุริโส อรุคตฺโต ปกฺกคตฺโต กิมีหิ ขชฺชมาโน นเขหิ วณมุขานิ วิปฺปตจฺฉมาโน องฺคารกาสุยา กายํ ปริตาเปยฺย.\csromanspacer{} ตสฺส มิตฺตามจฺจา ญาติสาโลหิตา ภิสกฺกํ สลฺลกตฺตํ อุปฏฺฐาเปยฺยุํ.\csromanspacer{} ตสฺส โส ภิสกฺโก สลฺลกตฺโต เภสชฺชํ กเรยฺย, โส ตํ เภสชฺชํ อาคมฺม กุฏฺเฐหิ ปริมุจฺเจยฺย, อโรโค อสฺส สุขี เสรี สยํวสี เยน กามํ คโม, ตเมนํ เทฺว พลวนฺโต ปุริสา นานาพาหาสุ คเหตฺวา องฺคารกาสุํ อุปกฑฺเฒยฺยุํ.}

\prose{ตํ กิํ มญฺญสิ มาคณฺฑิย, อปิ นุ โส ปุริโส อิติ จิติเจว กายํ สนฺนาเมยฺยาติ.\csromanspacer{} เอวํ โภ โคตม.\csromanspacer{} ตํ กิสฺส เหตุ.\csromanspacer{} อสุ หิ โภ โคตม อคฺคิ ทุกฺขสมฺผสฺโส เจว มหาภิตาโป จ มหาปริฬาโห จาติ.\csromanspacer{} ตํ กิํ มญฺญสิ มาคณฺฑิย, อิทาเนว นุ โข โส อคฺคิ ทุกฺขสมฺผสฺโส เจว มหาภิตาโป จ มหาปริฬาโห จ, อุทาหุ ปุพฺเพปิ โส อคฺคิ ทุกฺขสมฺผสฺโส เจว มหาภิตาโป จ มาหาปริฬาโห จาติ.\csromanspacer{} อิทานิ เจว โภ โคตม โส อคฺคิ ทุกฺขสมฺผสฺโส เจว}

\vspace{\tipitakaparskip}
\csromancenter{มาคณฺฑิยสุตฺต (75) มหาภิตาโป จ มหา ปริฬาโห จ, ปุพฺเพปิ โส อคฺคิ ทุกฺขสมฺผสฺโส เจว มหาภิตาโป จ มหาปริฬาโห จ, อสุ จ\footnote{อสุ หิ จ (สี, อิ)} โภ โคตม กุฏฺฐี ปุริโส อรุคตฺโต ปกฺกคตฺโต กิมีหิ ขชฺชมาโน นเขหิ วณมุขานิ วิปฺปตจฺฉมาโน อุปหตินฺทฺริโย ทุกฺขสมฺผสฺเสเยว อคฺคิสฺมิํ สุขมิติ วิปรีตสญฺญํ ปจฺจลตฺถาติ.\csromanspacer{} เอวเมว โข มาคณฺฑิย อตีตมฺปิ อทฺธานํ กามา ทุกฺขสมฺผสฺสา เจว มหาภิตาปา จ มหาปริฬาหา จ, อนาคตมฺปิ อทฺธานํ กามา ทุกฺขสมฺผสฺสา เจว มหาภิตาปา จ มหาปริฬาหา จ, เอตรหิปิ ปจฺจุปฺปนฺนํ อทฺธานํ กามา ทุกฺขสมฺผสฺสา เจว มหาภิตาปา จ มหาปริฬาหา จ, อิเม จ มาคณฺฑิย สตฺตา กาเมสุ อวีตราคา กามตณฺหาหิ ขชฺชมานา กามปริฬาเหน ปริฑยฺหมานา อุปหตินฺทฺริยา ทุกฺขสมฺผสฺเสสุเยว กาเมสุ สุขมิติ วิปรีตสญฺญํ ปจฺจลตฺถุํ.}
\proseitem{215}{\csromanfolio{175}เสยฺยถาปิ มาคณฺฑิย กุฏฺฐี ปุริโส อรุคตฺโต ปกฺกคตฺโต กิมีหิ ขชฺชมาโน นเขหิ วณมุขานิ วิปฺปตจฺฉมาโน องฺคารกาสุยา กายํ ปริตาเปติ.\csromanspacer{} ยถา ยถา โข มาคณฺฑิย อสุ กุฏฺฐี ปุริโส อรุคตฺโต ปกฺกคตฺโต กิมีหิ ขชฺชมาโน นเขหิ วณมุขานิ วิปฺปตจฺฉมาโน องฺคารกาสุยา กายํ ปริตาเปติ.\csromanspacer{} ตถา ตถา’สฺส\footnote{ตตฺถา ตตฺถา ตสฺเสว (สฺยา, กํ, ก)} ตานิ วณมุขานิ อสุจิตรานิ เจว โหนฺติ ทุคฺคนฺธตรานิ จ ปูติกตรานิ จ.\csromanspacer{} โหติ เจว กาจิ สาตมตฺตา อสฺสาทมตฺตา, ยทิทํ วณมุขานํ กณฺฑูวนเหตุ.\csromanspacer{} เอวเมว โข มาคณฺฑิย สตฺตา กาเมสุ อวีตราคา กามตณฺหาหิ ขชฺชมานา กามปริฬาเหน จ ปริฑยฺหมานา กาเม ปฏิเสวนฺติ.\csromanspacer{} ยถา ยถา โข มาคณฺฑิย สตฺตา กาเมสุ อวีตราคา กามตณฺหาหิ ขชฺชมานา กามปริฬาเหน จ ปริฑยฺหมานา กาเม ปฏิเสวนฺติ.\csromanspacer{} ตถา ตถา เตสํ เตสํ สตฺตานํ กามตณฺหา เจว ปวฑฺฒติ, กามปริฬาเหน จ ปริฑยฺหนฺติ.\csromanspacer{} โหติ เจว สาตมตฺตา อสฺสาทมตฺตา, ยทิทํ ปญฺจกามคุเณ ปฏิจฺจ.}

\prose{ตํ กิํ มญฺญสิ มาคณฺฑิย, อปิ นุ เต ทิฏฺโฐ วา สุโต วา ราชา วา ราชมหามตฺโต วา ปญฺจหิ กามคุเณหิ สมปฺปิโต สมงฺคีภูโต ปริจารยมาโน กามตณฺหํ อปฺปหาย กามปริฬาหํ อปฺปฏิวิโนเทตฺวา วิคตปิปาโส อชฺฌตฺตํ วูปสนฺตจิตฺโต วิหาสิ วา วิหรติ วา วิหริสฺสติ \csromanfolio{176}วาติ.\csromanspacer{} โน หิทํ โภ โคตม.\csromanspacer{} สาธุ มาคณฺฑิย, มยาปิ โข เอตํ มาคณฺฑิย เนว ทิฏฺฐํ น สุตํ ราชา วา ราชมหามตฺโต วา ปญฺจหิ กามคุเณหิ สมปฺปิโต สมงฺคีภูโต ปริจารยมาโน กามตณฺหํ อปฺปหาย กามปริฬาหํ อปฺปฏิวิโนเทตฺวา วิคตปิปาโส อชฺฌตฺตํ วูปสนฺตจิตฺโต วิหาสิ วา วิหรติ วา วิหริสฺสติ วา.\csromanspacer{} อถ โข มาคณฺฑิย เย หิ เกจิ สมณา วา พฺราหฺมณา วา วิคตปิปาสา อชฺฌตฺตํ วูปสนฺตจิตฺตา วิหาสุํ วา วิหรนฺติ วา วิหริสฺสนฺติ วา, สพฺเพ เต กามานํเยว สมุทยญฺจ อตฺถงฺคมญฺจ อสฺสาทญฺจ อาทีนวญฺจ นิสฺสรณญฺจ ยถาภูตํ วิทิตฺวา กามตณฺหํ ปหาย กามปริฬาหํ ปฏิวิโนเทตฺวา วิคตปิปาสา อชฺฌตฺตํ วูปสนฺตจิตฺตา วิหาสุํ วา วิหรนฺติ วา วิหริสฺสนฺติ วาติ.\csromanspacer{} อถ โข ภควา ตายํ เวลายํ อิมํ อุทานํ อุทาเนสิ.}

\csromangathameasure{อาโรคฺยปรมา ลาภา, นิพฺพานํ ปรมํ สุขํ. \\ อฏฺฐงฺคิโก จ มคฺคานํ, เขมํ อมตคามินนฺติ.}
\csromangathagroup{\csromangathastanza{อาโรคฺยปรมา ลาภา, นิพฺพานํ ปรมํ สุขํ. \\ อฏฺฐงฺคิโก จ มคฺคานํ, เขมํ อมตคามินนฺติ.}}

\proseitem{216}{เอวํ วุตฺเต มาคณฺฑิโย ปริพฺพาชโก ภควนฺตํ เอตทโวจ “อจฺฉริยํ โภ โคตม, อพฺภุตํ โภ โคตม, ยาว สุภาสิตํ จิทํ โภตา โคตเมน ‘อาโรคฺยปรมา ลาภา, นิพฺพานํ ปรมํ สุขนฺ’ติ, มยาปิ โข เอตํ โภ โคตม สุตํ ปุพฺพกานํ ปริพฺพาชกานํ อาจริยปาจริยานํ ภาสมานานํ ‘อาโรคฺยปรมา ลาภา, นิพฺพานํ ปรมํ สุขนฺ’ติ, ตยิทํ โภ โคตม สเมตี”ติ.\csromanspacer{} ยํ ปน เต เอตํ มาคณฺฑิย สุตํ ปุพฺพกานํ ปริพฺพาชกานํ อาจริยปาจริยานํ ภาสมานานํ ‘อาโรคฺยปรมา ลาภา, นิพฺพานํ ปรมํ สุขนฺ’ติ, กตมํ ตํ อาโรคฺยํ กตมํ ตํ นิพฺพานนฺติ.\csromanspacer{} เอวํ วุตฺเต มาคณฺฑิโย ปริพฺพาชโก สกาเนว สุทํ คตฺตานิ ปาณินา อโนมชฺชติ “อิทนฺตํ โภ โคตม อาโรคฺยํ, อิทนฺตํ นิพฺพานํ.\csromanspacer{} อหญฺหิ โภ โคตม เอตรหิ อโรโค สุขี, น มํ กิญฺจิ อาพาธตี”ติ.}

\proseitem{217}{เสยฺยถาปิ มาคณฺฑิย ชจฺจนฺโธ ปุริโส, โส น ปสฺเสยฺย กณฺหสุกฺกานิ รูปานิ, น ปสฺเสยฺย นีลกานิ รูปานิ, น ปสฺเสยฺย ปีตกานิ รูปานิ, น ปสฺเสยฺย โลหิตกานิ รูปานิ, น ปสฺเสยฺย มญฺชิฏฺฐกานิ\footnote{มญฺเจฏฺฐิกานิ (สี, สฺยา, กํ, อิ), มญฺเชฏฺฐกานิ (ก)} รูปานิ, น ปสฺเสยฺย สมวิสมํ, น ปสฺเสยฺย ตารกรูปานิ, น ปสฺเสยฺย สมวิสฺสมํ, น ปสฺเสยฺย ตารกรูปานิ, น ปสฺเสยฺย จนฺทิมสูริเย, โส สุเณยฺย จกฺขุมโต ภาสมานสฺส “เฉกํ วต}

\vspace{\tipitakaparskip}
\csromancenter{มาคณฺฑิยสุตฺต (75) โภ โอทาตํ วตฺถํ อภิรูปํ นิมฺมลํ สุจี”ติ.\csromanspacer{} โส โอทาตปริเยสนํ จเรยฺย, ตเมนํ อญฺญตโร ปุริโส เตลมลิเตน สาหุฬิจีเรน\footnote{เตลมสิกเตน สาหุฬจีวเรน (สี, สฺยา, กํ, อิ)} วญฺเจยฺย “อิทํ เต อมฺโภ ปุริส โอทาตํ วตฺถํ อภิรูปํ นิมฺมลํ สุจี”ติ.\csromanspacer{} โส ตํ ปฏิคฺคเณฺหยฺย, ปฏิคฺคเหตฺวา ปารุเปยฺย, ปารุเปตฺวา อตฺตมโน อตฺตมนวาจํ นิจฺฉาเรยฺย “เฉกํ วต โภ โอทาตํ วตฺถํ อภิรูปํ นิมฺมลํ สุจี”ติ.}
\prose{\csromanfolio{177}ตํ กิํ มญฺญสิ มาคณฺฑิย, อปิ นุ โส ชจฺจนฺโธ ปุริโส ชานนฺโต ปสฺสนฺโต อมุํ เตลมลิกตํ สาหุฬิจีรํ ปฏิคฺคเณฺหยฺย, ปฏิคฺคเหตฺวา ปารุเปยฺย, ปารุเปตฺวา อตฺตมโน อตฺตมนวาจํ นิจฺฉาเรยฺย “เฉกํ วต โภ โอทาตํ วตฺถํ อภิรูปํ นิมฺมลํ สุจี”ติ, อุทาหุ จกฺขุมโต สทฺธายาติ.\csromanspacer{} อชานนฺโต หิ โภ โคตม อปสฺสนฺโต โส ชจฺจนฺโธ ปุริโส อมุํ เตลมลิกตํ สาหุฬิจีรํ ปฏิคฺคเณฺหยฺย, ปฏิคฺคเหตฺวา ปารุเปยฺย, ปารุเปตฺวา อตฺตมโน อตฺตมนวาจํ นิจฺฉาเรยฺย “เฉกํ วต โภ โอทาตํ วตฺถํ อภิรูปํ นิมฺมลํ สุจี”ติ.\csromanspacer{} จกฺขุมโต สทฺธายาติ.\csromanspacer{} เอวเมว โข มาคณฺฑิย อญฺญติตฺถิยา ปริพฺพาชกา อนฺธา อจกฺขุกา อชานนฺตา อาโรคฺยํ อปสฺสนฺตา นิพฺพานํ.\csromanspacer{} อถ จ ปนิมํ คาถํ ภาสนฺติ “อาโรคฺยปรมา ลาภา, นิพฺพานํ ปรมํ สุขนฺ”ติ.\csromanspacer{} ปุพฺพเกเหสา มาคณฺฑิย อรหนฺเตหิ สมฺมาสมฺพุทฺเธหิ คาถา ภาสิตา.}

\csromangathameasure{อาโรคฺยปรมา ลาภา, นิพฺพานํ ปรมํ สุขํ. \\ อฏฺฐงฺคิโก จ มคฺคานํ, เขมํ อมตคามินนฺติ.}
\csromangathagroup{\csromangathastanza{อาโรคฺยปรมา ลาภา, นิพฺพานํ ปรมํ สุขํ. \\ อฏฺฐงฺคิโก จ มคฺคานํ, เขมํ อมตคามินนฺติ.}}

\proseitem{218}{สา เอตรหิ อนุปุพฺเพน ปุถุชฺชนคาถา\footnote{ปุถุชฺชนคตา (สี, อิ)}.\csromanspacer{} อยํ โข ปน มาคณฺฑิย กาโย โรคภูโต คณฺฑภูโต สลฺลภูโต อฆภูโต อาพาธภูโต, โส ตฺวํ อิมํ กายํ โรคภูตํ คณฺฑภูตํ สลฺลภูตํ อฆภูตํ อาพาธภูตํ “อิทนฺตํ โภ โคตม อาโรคฺยํ, อิทนฺตํ นิพฺพานนฺ”ติ วเทสิ.\csromanspacer{} ตํ หิ เต มาคณฺฑิย อริยํ จกฺขุํ นตฺถิ, เยน ตฺวํ อริเยน จกฺขุนา อาโรคฺยํ ชาเนยฺยาสิ, นิพฺพานํ ปสฺเสยฺยาสีติ.\csromanspacer{} เอวํ ปสนฺโน อหํ โภโต โคตมสฺส, ปโหติ เม ภวํ โคตโม ตถา ธมฺมํ เทเสตุํ, ยถาหํ อาโรคฺยํ ชาเนยฺยํ, นิพฺพานํ ปสฺเสยฺยนฺติ.}

\proseitem{219}{\csromanfolio{178}เสยฺยถาปิ มาคณฺฑิย ชจฺจนฺโธ ปุริโส, โส น ปสฺเสยฺย กณฺหสุกฺกานิ รูปานิ, น ปสฺเสยฺย นีลกานิ รูปานิ, น ปสฺเสยฺย ปีตกานิ รูปานิ, น ปสฺเสยฺย โลหิตกานิ รูปานิ, น ปสฺเสยฺย มญฺชิฏฺฐกานิ รูปานิ, น ปสฺเสยฺย สมวิสมํ, น ปสฺเสยฺย ตารกรูปานิ, น ปสฺเสยฺย จนฺทิมสูริเย, ตสฺส มิตฺตามจฺจา ญาติสาโลหิตา ภิสกฺกํ สลฺลกตฺตํ อุปฏฺฐาเปยฺยุํ, ตสฺส โส ภิสกฺโก สลฺลกตฺโต เภสชฺชํ กเรยฺย, โส ตํ เภสชฺชํ อาคมฺม น จกฺขูนิ อุปฺปาเทยฺย, น จกฺขูนิ วิโสเธยฺย.\csromanspacer{} ตํ กิํ มญฺญสิ มาคณฺฑิย, นนุ โส เวชฺโช ยาวเทว กิลมถสฺส วิฆาตสฺส ภาคี อสฺสาติ.\csromanspacer{} เอวํ โภ โคตม.\csromanspacer{} เอวเมว โข มาคณฺฑิย อหํ เจ เต ธมฺมํ เทเสยฺยํ “อิทนฺตํ อาโรคฺยํ, อิทนฺตํ นิพฺพานนฺ”ติ.\csromanspacer{} โส ตฺวํ อาโรคฺยํ น ชาเนยฺยาสิ, นิพฺพานํ น ปสฺเสยฺยาสิ.\csromanspacer{} โส มมสฺส กิลมโถ สา มมสฺส วิเหสาติ.\csromanspacer{} เอวํ ปสนฺโน อหํ โภโต โคตมสฺส, ปโหติ เม ภวํ โคตโม ตถา ธมฺมํ เทเสตุํ, ยถาหํ อาโรคฺยํ ชาเนยฺยํ, นิพฺพานํ ปสฺเสยฺยนฺติ.}

\proseitem{220}{เสยฺยถาปิ มาคณฺฑิย ชจฺจนฺโธ ปุริโส, โส น ปสฺเสยฺย กณฺหสุกฺกานิ รูปานิ, น ปสฺเสยฺย นีลกานิ รูปานิ, น ปสฺเสยฺย ปีตกานิ รูปานิ, น ปสฺเสยฺย โลหิตกานิ รูปานิ, น ปสฺเสยฺย มญฺชิฏฺฐกานิ รูปานิ, น ปสฺเสยฺย สมวิสมํ, น ปสฺเสยฺย ตารกรูปานิ, น ปสฺเสยฺย จนฺทิมสูริเย, โส สุเณยฺย จกฺขุมโต ภาสมานสฺส “เฉกํ วต โภ โอทาตํ วตฺถํ อภิรูปํ นิมฺมลํ สุจี”ติ.\csromanspacer{} โส โอทาตปริเยสนํ จเรยฺย, ตเมนํ อญฺญตโร ปุริโส เตลมลิกเตน สาหุฬิจีเรน วญฺเจยฺย “อิทํ เต อมฺโภ ปุริส โอทาตํ วตฺถํ อภิรูปํ นิมฺมลํ สุจี”ติ.\csromanspacer{} โส ตํ ปฏิคฺคเณฺหยฺย, ปฏิคฺคเหตฺวา ปารุเปยฺย, ตสฺส มิตฺตามจฺจา ญาติสาโลหิตา ภิสกฺกํ สลฺลกตฺตํ อุปฏฺฐาเปยฺยุํ, ตสฺส โส ภิสกฺโก สลฺลกตฺโต เภสชฺชํ กเรยฺย อุทฺธํวิเรจนํ อโธวิเรจนํ อญฺชนํ ปจฺจญฺชนํ นตฺถุกมฺมํ โส ตํ เภสชฺชํ อาคมฺม จกฺขูนิ อุปฺปาเทยฺย จกฺขูนิ วิโสเธยฺย.\csromanspacer{} ตสฺส สห จกฺขุปฺปาทา โย อมุสฺมิํ เตลมลิกเต สาหุฬิจีเร ฉนฺทราโค.\csromanspacer{} โส ปหีเยถ, ตญฺจ นํ ปุริสํ อมิตฺตโตปิ ทเหยฺย, ปจฺจตฺถิกโตปิ ทเหยฺย, อปิ จ ชีวิตา โวโรเปตพฺพํ มญฺเญยฺย.\csromanspacer{} ทีฆรตฺตํ วต โภ อหํ อิมินา ปุริเสน เตลมลิกเตน สาหุฬิจีเรน นิกโต วญฺจิโต ปลุทฺโธ “อิทํ เต อมฺโภ ปุริส โอทาตํ}

\vspace{\tipitakaparskip}
\csromancenter{มาคณฺฑิยสุตฺต (75) วตฺถํ อภิรูปํ นิมฺมลํ สุจี”ติ.\csromanspacer{} เอวเมว โข มาคณฺฑิย อหํ เจ เต ธมฺมํ เทเสยฺยํ “อิทนฺตํ อาโรคฺยํ อิทนฺตํ นิพฺพานนฺ”ติ, โส ตฺวํ อาโรคฺยํ ชาเนยฺยาสิ, นิพฺพานํ ปสฺเสยฺยาสิ.\csromanspacer{} ตสฺส เต สห จกฺขุปฺปาทา โย ปญฺจสุปาทานกฺขนฺเธสุ ฉนฺทราโค, โส ปหีเยถ, อปิ จ เต เอวมสฺส”ทีฆรตฺตํ วต โภ อหํ อิมินา จิตฺเตน นิกโต วญฺจิโต ปลุทฺโธ\footnote{ปลทฺโธ (สี, อิ)}, อหํ หิ รูปํเยว อุปาทิยมาโน อุปาทิยิํ, เวทนํเยว อุปาทิยมาโน อุปาทิยิํ, สญฺญํเยว อุปาทิยมาโน อุปาทิยิํ, สงฺขาเรเยว อุปาทิยมาโน อุปาทิยิํ, วิญฺญาณํเยว อุปาทิยมาโน อุปาทิยิํ, ตสฺส เม อุปาทานปจฺจยา ภโว, ภวปจฺจยา ชาติ, ชาติปจฺจยา ชารามรณํ โสกปริเทวทุกฺขโทมนสฺสุปายาสา สมฺภวนฺติ, เอวเมตสฺส เกวลสฺส ทุกฺขกฺขนฺธสฺส สมุทโย โหตี”ติ.\csromanspacer{} เอวํ ปสนฺโน อหํ โภโต โคตมสฺส, ปโหติ เม ภวํ โคตโม ตถา ธมฺมํ เทเสตุํ, ยถาหํ อิมมฺหา อาสนา อนนฺโธ วิฏฺฐเหยฺยนฺติ.}
\proseitem{221}{\csromanfolio{179}เตน หิ ตฺวํ มาคณฺฑิย สปฺปุริเส ภเชยฺยาสิ, ยโต โข ตฺวํ มาคณฺฑิย สปฺปุริเส ภชิสฺสสิ, ตโต ตฺวํ มาคณฺฑิย สทฺธมฺมํ โสสฺสสิ, ยโต โข ตฺวํ มาคณฺฑิย สทฺธมฺมํ โสสฺสสิ, ตโต ตฺวํ มาคณฺฑิย ธมฺมานุธมฺมํ ปฏิปชฺชิสฺสสิ, ยโต โข ตฺวํ มาคณฺฑิย ธมฺมานุธมฺมํ ปฏิปชฺชิสฺสสิ, ตโต ตฺวํ มาคณฺฑิย สามํเยว ญสฺสสิ, สามํ ทกฺขิสฺสสิ “อิเม โรคา คณฺฑา สลฺลา อิธ โรคา คณฺฑา สลฺลา อปริเสสา นิรุชฺฌนฺติ.\csromanspacer{} ตสฺส เม อุปาทานนิโรธา ภวนิโรโธ, ภวนิโรธา ชาตินิโรโธ, ชาตินิโรธา ชรามรณํ โสกปริเทว ทุกฺขโทมนสฺสุปายาสา นิรุชฺฌนฺติ, เอวเมตสฺส เกวลสฺส ทุกฺขกฺขนฺธสฺส นิโรโธ โหตี”ติ.}

\proseitem{222}{เอวํ วุตฺเต มาคณฺฑิโย ปริพฺพาชโก ภควนฺตํ เอตทโวจ “อภิกฺกนฺตํ โภ โคตม, อภิกฺกนฺตํ โภ โคตม, เสยฺยถาปิ โภ โคตม นิกฺกุชฺชิตํ วา อุกฺกุชฺเชยฺย, ปฏิจฺฉนฺนํ วา วิวเรยฺย, มูฬฺหสฺส วา มคฺคํ อาจิกฺเขยฺย, อนฺธกาเร วา เตลปชฺโชตํ ธาเรยฺย ‘จกฺขุมนฺโต รูปานิ ทกฺขนฺตี’ติ, เอวเมวํ โภตา โคตเมน อเนกปริยาเยน ธมฺโม ปกาสิโต.\csromanspacer{} เอสาตํ ภวนฺตํ โคตมํ สรณํ คจฺฉามิ ธมฺมญฺจ \csromanfolio{180}ภิกฺขุสํฆญฺจ, ลเภยฺยาหํ โภโต โคตมสฺส สนฺติเก ปพฺพชฺชํ, ลเภยฺยํ อุปสมฺปทนฺ”ติ.\csromanspacer{} โย โข มาคณฺฑิย อญฺญติตฺถิยปุพฺโพ อิมสฺมิํ ธมฺมวินเย อากงฺขติ ปพฺพชฺชํ, อากงฺขติ อุปสมฺปทํ, โส จตฺตาโร มาเส ปริวสติ, จตุนฺนํ มาสานํ อจฺจเยน อารทฺธจิตฺตา ภิกฺขู ปพฺพาเชนฺติ อุปสมฺปาเทนฺติ ภิกฺขุภาวาย, อปิ จ เมตฺถ ปุคฺคลเวมตฺตตา วิทิตาติ.\csromanspacer{} สเจ ภนฺเต อญฺญติตฺถิยปุพฺพา อิมสฺมิํ ธมฺมวินเย อากงฺขนฺตา ปพฺพชฺชํ อากงฺขนฺตา อุปสมฺปทํ จตฺตาโร มาเส ปริวสนฺติ, จตุนฺนํ มาสานํ อจฺจเยน อารทฺธจิตฺตา ภิกฺขู ปพฺพาเชนฺติ อุปสมฺปาเทนฺติ ภิกฺขุภาวาย, อหํ จตฺตาริ วสฺสานิ ปริวสิสฺสามิ, จตุนฺนํ วสฺสานํ อจฺจเยน อารทฺธจิตฺตา ภิกฺขู ปพฺพาเชนฺตุ อุปสมฺปาเทนฺตุ ภิกฺขุภาวายาติ.\csromanspacer{} อลตฺถ โข มาคณฺฑิโย ปริพฺพาชโก ภควโต สนฺติเก ปพฺพชฺชํ, อลตฺถ อุปสมฺปทํ.\csromanspacer{} อจิรูปสมฺปนฺโน โข ปนายสฺมา มาคณฺฑิโย เอโก วูปกฏฺโฐ อปฺปมตฺโต อาตาปี ปหิตตฺโต วิหรนฺโต นจิรสฺเสว, ยสฺสตฺถาย กุลปุตฺตา สมฺมเทว อคารสฺมา อนคาริยํ ปพฺพชนฺติ, ตทนุตฺตรํ พฺรหฺมจริยปริโยสานํ ทิฏฺเฐว ธมฺเม สยํ อภิญฺญา สจฺฉิกตฺวา อุปสมฺปชฺช วิหาสิ, “ขีณา ชาติ, วุสิตํ พฺรหฺมจริยํ, กตํ กรณียํ, นาปรํ อิตฺถตฺตายา”ติ อพฺภญฺญาสิ.\csromanspacer{} อญฺญตโร โข ปนายสฺมา มาคณฺฑิโย อรหตํ อโหสีติ.}

\nitthitamruled{มาคณฺฑิยสุตฺตํ นิฏฺฐิตํ ปญฺจมํ.}
\csromanmatikaanchor{mtk.31}
\csromantocmarkref{subsection}{6. สนฺทกสุตฺต}{mtk.31}
\bookmark[dest={mtk.31},level=2]{6. สนฺทกสุตฺต}
\csromanheader{2}{6. สนฺทกสุตฺต}
\proseitem{223}{เอวํ เม สุตํ\csromandash{}เอกํ สมยํ ภควา โกสมฺพิยํ วิหรติ โฆสิตาราเม.\csromanspacer{} เตน โข ปน สมเยน สนฺทโก ปริพฺพาชโก ปิลกฺขคุหายํ ปฏิวสติ มหติยา ปริพฺพาชกปริสาย สทฺธิํ ปญฺจมตฺเตหิ ปริพฺพาชกสเตหิ.\csromanspacer{} อถ โข อายสฺมา อานนฺโท สายนฺหสมยํ ปฏิสลฺลานา วุฏฺฐิโต ภิกฺขู อามนฺเตสิ “อายามาวุโส เยน เทวกตโสพฺโภ เตนุปสงฺกมิสฺสาม คุหาทสฺสนายา”ติ. “เอวมาวุโส”ติ โข เต ภิกฺขู อายสฺมโต อานนฺทสฺส ปจฺจสฺโสสุํ.\csromanspacer{} อถ โข อายสฺมา อานนฺโท สมฺพหุเลหิ ภิกฺขูหิ สทฺธิํ เยน เทวกตโสพฺโภ เตนุปสงฺกมิ.\csromanspacer{} เตน โข ปน สมเยน สนฺทโก ปริพฺพาชโก}

\vspace{\tipitakaparskip}
\csromancenter{สนฺทกสุตฺต (76) มหติยา ปริพฺพาชกปริสาย สทฺธิํ นิสินฺโน โหติ อุนฺนาทินิยา อุจฺจาสทฺทมหาสทฺทาย อเนกวิหิตํ ติรจฺฉานกถํ กเถนฺติยา, เสยฺยถิทํ, ราชกถํ โจรกถํ มหามตฺตกถํ เสนากถํ ภยกถํ ยุทฺธกถํ อนฺนกถํ ปานกถํ วตฺถกถํ สยนกถํ มาลากถํ คนฺธกถํ ญาติกถํ ยานกถํ คามกถํ นิคมกถํ นครกถํ ชนปทกถํ อิตฺถิกถํ สูรกถํ วิสิขากถํ กุมฺภฏฺฐานกถํ ปุพฺพเปตกถํ นานตฺตกถํ โลกกฺขายิกํ สมุทฺทกฺขายิกํ อิติภวาภวกถํ อิติ วา.\csromanspacer{} อทฺทสา โข สนฺทโก ปริพฺพาชโก อายสฺมนฺตํ อานนฺทํ ทูรโตว อาคจฺฉนฺตํ, ทิสฺวาน สกํ ปริสํ สณฺฐาเปสิ “อปฺปสทฺทา โภนฺโต โหนฺตุ, มา โภนฺโต สทฺทมกตฺถ, อยํ สมณสฺส โคตมสฺส สาวโก อาคจฺฉติ สมโณ อานนฺโท, ยาวตา โข ปน สมณสฺส โคตมสฺส สาวกา โกสมฺพิยํ ปฏิวสนฺติ, อยํ เตสํ อญฺญตโร สมโณ อานนฺโท, อปฺปสทฺทกามา โข ปน เต อายสฺมนฺโต อปฺปสทฺทวินีตา อปฺปสทฺทสฺส วณฺณวาทิโน, อปฺเปว นาม อปฺปสทฺทํ ปริสํ วิทิตฺวา อุปสงฺกมิตพฺพํ มญฺเญยฺยา”ติ.\csromanspacer{} อถ โข เต ปริพฺพาชกา ตุณฺหี อเหสุํ.}
\proseitem{224}{\csromanfolio{181}อถ โข อายสฺมา อานนฺโท เยน สนฺทโก ปริพฺพาชโก เตนุปสงฺกมิ.\csromanspacer{} อถ โข สนฺทโก ปริพฺพาชโก อายสฺมนฺตํ อานนฺทํ เอตทโวจ “เอตุ โข ภวํ อานนฺโท, สฺวาคตํ โภโต อานนฺทสฺส, จิรสฺสํ โข ภวํ อานนฺโท อิมํ ปริยายมกาสิ, ยทิทํ อิธาคมนาย, นิสีทตุ ภวํ อานนฺโท, อิทมาสนํ ปญฺญตฺตนฺ”ติ.\csromanspacer{} นิสีทิ โข อายสฺมา อานนฺโท ปญฺญตฺเต อาสเน, สนฺทโกปิ โข ปริพฺพาชโก อญฺญตรํ นีจํ อาสนํ คเหตฺวา เอกมนฺตํ นิสีทิ.\csromanspacer{} เอกมนฺตํ นิสินฺนํ โข สนฺทกํ ปริพฺพาชกํ อายสฺมา อานนฺโท เอตทโวจ “กายนุตฺถ สนฺทก เอตรหิ กถาย สนฺนิสินฺนา, กา จ ปน โว อนฺตรากถา วิปฺปกตา”ติ.\csromanspacer{} ติฏฺฐเตสา โภ อานนฺท กถา, ยาย มยํ เอตรหิ กถาย สนฺนิสินฺนา, เนสา โภโต อานนฺทสฺส กถา ทุลฺลภา ภวิสฺสติ ปจฺฉาปิ สวนาย, สาธุ วต ภวนฺตํเยว อานนฺทํ ปฏิภาตุ สเก อาจริยเก ธมฺมีกถาติ.\csromanspacer{} เตน หิ สนฺทก สุณาหิ สาธุกํ มนสิ กโรหิ ภาสิสฺสามีติ. “เอวํ โภ”ติ โข สนฺทโก ปริพฺพาชโก อายสฺมโต อานนฺทสฺส ปจฺจสฺโสสิ.\csromanspacer{} อายสฺมา อานนฺโท เอตทโวจ \csromanfolio{182}“จตฺตาโรเม สนฺทก เตน ภควตา ชานตา ปสฺสตา อรหตา สมฺมาสมฺพุทฺเธน อพฺรหฺมจริยวาสา อกฺขาตา, จตฺตาริ จ อนสฺสาสิกานิ พฺรหฺมจริยานิ อกฺขาตานิ, ยตฺถ วิญฺญู ปุริโส สสกฺกํ พฺรหฺมจริยํ น วเสยฺย, วสนฺโต จ\footnote{วสนฺโต วา (สี, อิ) เอวมุปริปิ อนาราธนปกฺเข.} นาราเธยฺย ญายํ ธมฺมํ กุสลนฺ”ติ.\csromanspacer{} กตเม ปน เต โภ อานนฺท เตน ภควตา ชานตา ปสฺสตา อรหตา สมฺมาสมฺพุทฺเธน จตฺตาโร อพฺรหฺมจริยวาสา อกฺขาตา, ยตฺถ วิญฺญู ปุริโส สสกฺกํ พฺรหฺมจริยํ น วเสยฺย, วสนฺโต จ นาราเธยฺย ญายํ ธมฺมํ กุสลนฺติ.}

\proseitem{225}{อิธ สนฺทก เอกจฺโจ สตฺถา เอวํวาที โหติ เอวํทิฏฺฐิ “นตฺถิ ทินฺนํ, นตฺถิ ยิฏฺฐํ, นตฺถิ หุตํ, นตฺถิ สุกตทุกฺกฏานํ กมฺมานํ ผลํ วิปาโก, นตฺถิ อยํ โลโก, นตฺถิ ปโรโลโก, นตฺถิ มาตา, นตฺถิ ปิตา, นตฺถิ สตฺตา โอปปาติกา, นตฺถิ โลเก สมณพฺราหฺมณา สมฺมคฺคตา สมฺมาปฏิปนฺนา เย อิมญฺจ โลกํ ปรญฺจ โลกํ สยํ อภิญฺญา สจฺฉิกตฺวา ปเวเทนฺติ.\csromanspacer{} จาตุมหาภูติโก อยํ ปุริโส, ยทา กาลํ กโรติ, ปถวี ปถวีกายํ อนุเปติ อนุปคจฺฉติ, อาโป อาโปกายํ อนุเปติ อนุปคจฺฉติ, เตโช เตโชกายํ อนุเปติ อนุปคจฺฉติ, วาโย วาโยกายํ อนุเปติ อนุปคจฺฉติ, อากาสํ อินฺทฺริยานิ สงฺกมนฺติ, อาสนฺทิปญฺจมา ปุริสา มตํ อาทาย คจฺฉนฺติ ยาวาฬาหนา, ปทานิ ปญฺญายนฺติ, กาโปตกานิ อฏฺฐีนิ ภวนฺติ, ภสฺสนฺตา อาหุติโย, ทตฺตุปญฺญตฺตํ ยทิทํ ทานํ.\csromanspacer{} เตสํ ตุจฺฉา มุสา วิลาโป, เย เกจิ อตฺถิกวาทํ วทนฺติ.\csromanspacer{} พาเล จ ปณฺฑิเต จ กายสฺส เภทา อุจฺฉิชฺชนฺติ วินสฺสนฺติ, น โหนฺติ ปรํ มรณา”ติ.}

\prose{ตตฺร สนฺทก วิญฺญู ปุริโส อิติ ปฏิสญฺจิกฺขติ\csromandash{}อยํ โข ภวํ สตฺถา เอวํวาที เอวํทิฏฺฐิ “นตฺถิ ทินฺนํ, นตฺถิ ยิฏฺฐํ, นตฺถิ หุตํ, นตฺถิ สุกตทุกฺกฏานํ กมฺมานํ ผลํ วิปาโก, นตฺถิ อยํ โลโก, นตฺถิ ปโรโลโก, นตฺถิ มาตา, นตฺถิ ปิตา, นตฺถิ สตฺตา โอปปาติกา, นตฺถิ โลเก สมณพฺราหฺมณา สมฺมคฺคตา สมฺมาปฏิปนฺนา เย อิมญฺจ โลกํ ปรญฺจ โลกํ สยํ อภิญฺญา สจฺฉิกตฺวา ปเวเทนฺติ.\csromanspacer{} จาตุมหาภูติโก อยํ ปุริโส, ยทา กาลํกโรติ, ปถวี ปถวีกายํ อนุเปติ อนุปคจฺฉติ, อาโป}

\vspace{\tipitakaparskip}
\csromancenter{สนฺทกสุตฺต (76) อาโปกายํ อนุเปติ อนุปคจฺฉติ, เตโช เตโชกายํ อนุเปติ อนุปคจฺฉติ, วาโย วาโยกายํ อนุเปติ อนุปคจฺฉติ, อากาสํ อินฺทฺริยานิ สงฺกมนฺติ, อาสนฺทิปญฺจมา ปุริสา มตํ อาทาย คจฺฉนฺติ ยาวาฬาหนา, ปทานิ ปญฺญายนฺติ, กาโปตกานิ อฏฺฐีนิ ภวนฺติ, ภสฺสนฺตา อาหุติโย, ทตฺตุปญฺญตฺตํ ยทิทํ ทานํ.\csromanspacer{} เตสํ ตุจฺฉา มุสา วิลาโป, เยเกจิ อตฺถิกวาทํ วทนฺติ.\csromanspacer{} พาเล จ ปณฺฑิเต จ กายสฺส เภทา อุจฺฉิชฺชนฺติ วินสฺสนฺติ, น โหนฺติ ปรํ มรณา”ติ.\csromanspacer{} สเจ อิมสฺส โภโต สตฺถุโน สจฺจํ วจนํ.\csromanspacer{} อกเตน เม เอตฺถ กตํ, อวุสิเตน เม เอตฺถ วุสิตํ.\csromanspacer{} อุโภปิ มยํ เอตฺถ สมสมา สามญฺญํ ปตฺตา, โย จาหํ น วทามิ “อุโภ กายสฺส เภทา อุจฺฉิชฺชิสฺสาม วินสฺสิสฺสาม น ภวิสฺสาม ปรํ มรณา”ติ.\csromanspacer{} อติเรกํ โข ปนิมสฺส โภโต สตฺถุโน นคฺคิยํ มุณฺฑิยํ อุกฺกุฏิกปฺปธานํ เกสมสฺสุโลจนํ.\csromanspacer{} โยหํ ปุตฺตสมฺพาธสยนํ\footnote{ปุตฺตสมฺมาธวสนํ (สี)} อชฺฌาวสนฺโต กาสิกจนฺทนํ ปจฺจนุโภนฺโต มาลาคนฺธวิเลปนํ ธาเรนฺโต ชาตรูปรชตํ สาทิยนฺโต อิมินา โภตา สตฺถารา สมยมคติโก ภวิสฺสามิ อภิสมฺปรายํ.\csromanspacer{} โสหํ กิํ ชานนฺโต กิํ ปสฺสนฺโต อิมสฺมิํ สตฺถริ พฺรหฺมจริยํ จริสฺสามิ.\csromanspacer{} โส “อพฺรหฺมจริยวาโส อยนฺ”ติ อิติ วิทิตฺวา ตสฺมา พฺรหฺมจริยา นิพฺพิชฺช ปกฺกมติ\footnote{นิพฺพิชฺชาปกฺกมติ (สี)}.\csromanspacer{} อยํ โข สนฺทก เตน ภควตา ชานตา ปสฺสตา อรหตา สมฺมาสมฺพุทฺเธน ปฐโม อพฺรหฺมจริยวาโส อกฺขาโต, ยตฺถ วิญฺญู ปุริโส สสกฺกํ พฺรหฺมจริยํ น วเสยฺย, วสนฺโต จ นาราเธยฺย ญายํ ธมฺมํ กุสลํ.}
\proseitem{226}{\csromanfolio{183}ปุน จปรํ สนฺทก อิเธกจฺโจ สตฺถา เอวํวาที โหติ เอวํทิฏฺฐิ “กโรโต การยโต, ฉินฺทโต เฉทาปยโต, ปจโต ปาจาปยโต, โสจยโต โสจาปยโต, กิลมโต กิลมาปยโต, ผนฺทโต ผนฺทาปยโต, ปาณมติปาตยโต, อทินฺนํ อาทิยโต, สนฺธิํ ฉินฺทโต, นิลฺโลปํ หรโต, เอกาคาริกํ กโรโต, ปริปนฺเถ ติฏฺฐโต, ปรทารํ คจฺฉโต, มุสา ภณโต, กโรโต น กรียติ ปาปํ, ขุรปริยนฺเตน เจปิ จกฺเกน โย อิมิสฺสา ปถวิยา ปาเณ เอกํ มํสขลํ เอกํ มํสปุญฺชํ กเรยฺย, นตฺถิ ตโตนิทานํ ปาปํ, นตฺถิ ปาปสฺส อาคโม.\csromanspacer{} ทกฺขิณญฺเจปิ คงฺคาย ตีรํ คจฺเฉยฺย หนนฺโต ฆาเตนฺโต ฉินฺทนฺโต เฉทาเปนฺโต ปจนฺโต ปจาเปนฺโต, นตฺถิ ตโตนิทานํ ปาปํ, \csromanfolio{184}นิตฺถิ ปาปสฺส อาคโม.\csromanspacer{} อุตฺตรญฺเจปิ คงฺคาย ตีรํ คจฺเฉยฺย ททนฺโต ทาเปนฺโต ยชนฺโต ยชาเปนฺโต, นตฺถิ ตโตนิทานํ ปุญฺญํ, นตฺถิ ปุญฺญสฺส อาคโม.\csromanspacer{} ทาเนน ทเมน สํยเมน สจฺจวชฺเชน นตฺถิ ปุญฺญํ, นตฺถิ ปุญฺญสฺส อาคโม”ติ.}

\prose{ตตฺร สนฺทก วิญฺญู ปุริโส อิติ ปฏิสญฺจิกฺขติ\csromandash{}อยํ โข ภวํ สตฺถา เอวํวาที เอวํทิฏฺฐิ “กโรโต การยโต, ฉินฺทโต เฉทาปยโต, ปจโต ปาจาปยโต, โสจโต โสจาปยโต, กิลมโต, กิลมาปยโต, ผนฺทโต ผนฺทาปยโต, ปาณมติปาตยโต, อทินฺนํ อาทิยโต, สนฺธิํ ฉินฺทโต, นิลฺโลปํ หรโต, เอกาคาริกํ กโรโต, ปริปนฺเถ ติฏฺฐโต, ปรทารํ คจฺฉโต, มุสา ภณโต, กโรโต น กรียติ ปาปํ, ขุรปริยนฺเตน เจปิ จกฺเกน โย อิมิสฺสา ปถวิยา ปาเณ เอกํ มํสขลํ เอกํ มํสปุญฺชํ กเรยฺย, นตฺถิ ตโตนิทานํ ปาปํ, นตฺถิ ปาปสฺส อาคโม.\csromanspacer{} ทกฺขิณญฺเจปิ คงฺคาย ตีรํ คจฺเฉยฺย หนนฺโต ฆาเตนฺโต ฉินฺทนฺโต เฉทาเปนฺโต ปจนฺโต ปจาเปนฺโต, นตฺถิ ตโตนิทานํ ปาปํ, นตฺถิ ปาปสฺส อาคโม.\csromanspacer{} อุตฺตรญฺเจปิ คงฺคาย ตีรํ คจฺเฉยฺย ททนฺโต ทาเปนฺโต ยชนฺโต ยชาเปนฺโต, นตฺถิ ตโตนิทานํ ปุญฺญํ, นตฺถิ ปุญฺญสฺส อาคโม.\csromanspacer{} ทาเนน ทเมน สํยเมน สจฺจวชฺเชน นตฺถิ ปุญฺญํ, นตฺถิ ปุญฺญสฺส อาคโม”ติ.\csromanspacer{} สเจ อิมสฺส โภโต สตฺถุโน สจฺจํ วจนํ.\csromanspacer{} อกเตน เม เอตฺถ กตํ, อวุสิเตน เม เอตฺถ วุสิตํ, อุโภปิ มยํ เอตฺถ สมสมา สามญฺญํ ปตฺตา.\csromanspacer{} โย จาหํ น วทามิ “อุภินฺนํ กุรุตํ น กรียติ ปาปนฺ”ติ.\csromanspacer{} อติเรกํ โข ปนิมสฺส โภโต สตฺถุโน นคฺคิยํ มุณฺฑิยํ อุกฺกุฏิกปฺปธานํ เกสมสฺสุโลจนํ.\csromanspacer{} โยหํ ปุตฺตสมฺพาธสยนํ อชฺฌาวสนฺโต กาสิกจนฺทนํ ปจฺจนุโภนฺโต มาลาคนฺธวิเลปนํ ธาเรนฺโต ชาตรูปรชตํ สาทิยนฺโต อิมินา โภตา สตฺถารา สมสมคติโก ภวิสฺสามิ อภิสมฺปรายํ.\csromanspacer{} โสหํ กิํ ชานนฺโต กิํ ปสฺสนฺโต อิมสฺมิํ สตฺถริ พฺรหฺมจริยํ จริสฺสามิ.\csromanspacer{} โส “อพฺรหฺมจริยวาโส อยนฺ”ติ อิติ วิทิตฺวา ตสฺมา พฺรหฺมจริยา นิพฺพิชฺช ปกฺกมติ.\csromanspacer{} อยํ โข สนฺทก เตน ภควตา ชานตา ปสฺสตา อรหตา สมฺมาสมฺพุทฺเธน ทุติโย อพฺรหฺมจริยวาโส อกฺขาโต, ยตฺถ วิญฺญู ปุริโส สสกฺกํ พฺรหฺมจริยํ น วเสยฺย, วสนฺโต จ นาราเธยฺย ญายํ ธมฺมํ กุสลํ.}

\proseitem{227}{ปุน จปรํ สนฺทก อิเธกจฺโจ สตฺถา เอวํวาที โหติ เอวํทิฏฺฐิ “นตฺถิ เหตุ นตฺถิ ปจฺจโย สตฺตานํ สํกิเลสาย, อเหตู อปฺปจฺจยา}

\vspace{\tipitakaparskip}
\csromancenter{สนฺทกสุตฺต (76) สตฺตา สํกิลิสฺสนฺติ.\csromanspacer{} นตฺถิ เหตุ นตฺถิ ปจฺจโย สตฺตานํ วิสุทฺธิยา, อเหตู อปฺปจฺจยา สตฺตา วิสุชฺฌนฺติ.\csromanspacer{} นตฺถิ พลํ, นตฺถิ วีริยํ, นตฺถิ ปุริสถาโม, นตฺถิ ปุริสปรกฺกโม, สพฺเพ สตฺตา สพฺเพ ปาณา สพฺเพ ภูตา สพฺเพ ชีวา อวสา อพลา อวีริยา นิยติ สงฺคติภาวปริณตา ฉเสฺววาภิชาตีสุ สุขทุกฺขํ ปฏิสํเวเทนฺตี”ติ.}
\prose{\csromanfolio{185}ตตฺร สนฺทก วิญฺญู ปุริโส อิติ ปฏิสญฺจิกฺขติ\csromandash{}อยํ โข ภวํ สตฺถา เอวํวาที เอวํทิฏฺฐิ “นตฺถิ เหตุ นตฺถิ ปจฺจโย สตฺตานํ สํกิเลสาย, อเหตู อปฺปจฺจยา สตฺตา สํกิลิสฺสนฺติ.\csromanspacer{} นตฺถิ เหตุ นตฺถิ ปจฺจโย สตฺตานํ วิสุทฺธิยา, อเหตู อปฺปจฺจยา สตฺตา วิสุชฺฌนฺติ.\csromanspacer{} นตฺถิ พลํ, นตฺถิ วีริยํ, นตฺถิ ปุริสถาโม, นตฺถิ ปุริสปรกฺกโม, สพฺเพ สตฺตา สพฺเพ ปาณา สพฺเพ ภูตา สพฺเพ ชีวา อวสา อพลา อวีริยา นิยติสงฺคติภาวปริณตา ฉเสฺววาภิชาตีสุ สุขทุกฺขํ ปฏิสํเวเทนฺตี”ติ.\csromanspacer{} สเจ อิมสฺส โภโต สตฺถุโน สจฺจํ วจนํ.\csromanspacer{} อกเตน เม เอตฺถ กตํ, อวุสิเตน เม เอตฺถ วุสิตํ, อุโภปิ มยํ เอตฺถ สมสมา สามญฺญํ ปตฺตา.\csromanspacer{} โย จาหํ น วทามิ “อุโภ อเหตู อปฺปจฺจยา วิสุชฺฌิสฺสามา”ติ.\csromanspacer{} อติเรกํ โข ปนิมสฺส โภโต สตฺถุโน นคฺคิยํ มุณฺฑิยํ อุกฺกุฏิกปฺปธานํ เกสมสฺสุโลจนํ.\csromanspacer{} โยหํ ปุตฺตสมฺพาธสยนํ อชฺฌาวสนฺโต กาสิกจนฺทนํ ปจฺจนุโภนฺโต มาลาคนฺธวิเลปนํ ธาเรนฺโต ชาตรูปรชตํ สาทิยนฺโต อิมินา โภตา สตฺถารา สมสมคติโก ภวิสฺสามิ อภิสมฺปรายํ.\csromanspacer{} โสหํ กิํ ชานนฺโต กิํ ปสฺสนฺโต อิมสฺมิํ สตฺถริ พฺรหฺมจริยํ จริสฺสามิ.\csromanspacer{} โส “อพฺรหฺมจริยวาโส อยนฺ”ติ อิติ วิทิตฺวา ตสฺมา พฺรหฺมจริยา นิพฺพิชฺช ปกฺกมติ.\csromanspacer{} อยํ โข สนฺทก เตน ภควตา ชานตา ปสฺสตา อรหตา สมฺมาสมฺพุทฺเธน ตติโย อพฺรหฺมจริยวาโส อกฺขาโต, ยตฺถ วิญฺญู ปุริโส สสกฺกํ พฺรหฺมจริยํ น วเสยฺย, วสนฺโต จ นาราเธยฺย ญายํ ธมฺมํ กุสลํ.}

\proseitem{228}{ปุน จปรํ สนฺทก อิเธกจฺโจ สตฺถา เอวํวาที โหติ เอวํทิฏฺฐิ “สตฺติเม กายา อกฏา อกฏวิธา อนิมฺมิตา อนิมฺมาตา วญฺฌา กูฏฏฺฐา เอสิกฏฺฐายิฏฺฐิตา.\csromanspacer{} เต น อิญฺชนฺติ, น วิปริณมนฺติ, น อญฺญมญฺญํ พฺยาพาเธนฺติ.\csromanspacer{} นาลํ อญฺญมญฺญสฺส สุขาย วา ทุกฺขาย วา สุขทุกฺขาย วา.\csromanspacer{} กตเม สตฺต, ปถวีกาโย อาโปกาโย เตโชกาโย วาโยกาโย สุเข ทุกฺเข ชีเว สตฺตเม.\csromanspacer{} อิเม สตฺตกายา อกฏา อกฏวิธา \csromanfolio{186}อนิมฺมิตา อนิมฺมาตา วญฺฌา กูฏฏฺฐา เอสิกฏฺฐายิฏฺฐิตา.\csromanspacer{} เต น อิญฺชนฺติ, น วิปริณมนฺติ, น อญฺญมญฺญํ พฺยาพาเธนฺติ.\csromanspacer{} นาลํ อญฺญมญฺญสฺส สุขาย วา ทุกฺขาย วา สุขทุกฺขาย วา.\csromanspacer{} ตตฺถ นตฺถิ หนฺตา วา ฆาเตตา วา โสตา วา สาเวตา วา วิญฺญาตา วา วิญฺญาเปตา วา.\csromanspacer{} โยปิ ติเณฺหน สตฺเถน สีสํ ฉินฺทติ, น โกจิ กญฺจิ\footnote{กิญฺจิ (ก)} ชีวิตา โวโรเปติ.\csromanspacer{} สตฺตนฺนํเตฺวว กายานมนฺตเรน สตฺถํ วิวรมนุปตติ.\csromanspacer{} จุทฺทส โข ปนิมานิ โยนิปมุขสตสหสฺสานิ สฏฺฐิ จ สตานิ ฉ จ สตานิ ปญฺจ จ กมฺมุโน สตานิ ปญฺจ จ กมฺมานิ ตีณิ จ กมฺมานิ กมฺเม จ อฑฺฒกมฺเม จ ทฺวฏฺฐิปฏิปทา ทฺวฏฺฐนฺตรกปฺปา ฉฬาภิชาติโย อฏฺฐ ปุริสภูมิโย เอกูนปญฺญาส อาชีวกสเต เอกูนปญฺญาส ปริพฺพาชกสเต เอกูนปญฺญาย นาคาวาสสเต วีเส อินฺทฺริยสเต ติํเส นิรยสเต ฉตฺติํส รโชธาตุโย สตฺต สญฺญีคพฺภา สตฺต อสญฺญีคพฺภา สตฺต นิคณฺฐิคพฺภา สตฺต เทวา สตฺต มานุสา สตฺต เปสาจา สตฺต สรา สตฺต ปวุฏา สตฺต ปปาตา สตฺต ปปาตสตานิ สตฺต สุปินา สตฺต สุปินสตานิ จุลฺลาสีติ\footnote{จูฬาสีติ (สี, สฺยา, กํ, อิ)} มหากปฺปิโน\footnote{มหากปฺปุโน (สี, อิ)} สตสหสฺสานิ, ยานิ พาเล จ ปณฺฑิเต จ สนฺธาวิตฺวา สํสริตฺวา ทุกฺขสฺสนฺตํ กริสฺสนฺติ.\csromanspacer{} ตตฺถ นตฺถิ ‘อิมินาหํ สีเลน วา วเตน วา ตเปน วา พฺรหฺมจริเยน วา อปริปกฺกํ วา กมฺมํ ปริปาเจสฺสามิ, ปริปกฺกํ วา กมฺมํ ผุสฺส ผุสฺส พฺยนฺติํ กริสฺสามี’ติ, เหวํ นตฺถิ, โทณมิเต สุขทุกฺเข, ปริยนฺตกเต สํสาเร, นตฺถิ หายนวฑฺฒเน, นตฺถิ อุกฺกํสาวกํเส.\csromanspacer{} เสยฺยถาปิ นาม สุตฺตคุเฬ ขิตฺเต นิพฺเพฐิยมานเมว ปเลติ, เอวเมว พาเล จ ปณฺฑิเต จ สนฺธาวิตฺวา สํสริตฺวา ทุกฺขสฺสนฺตํ กริสฺสนฺตี”ติ.}

\prose{ตตฺร สนฺทก วิญฺญู ปุริโส อิติ ปฏิสญฺจิกฺขติ\csromandash{}อยํ โข ภวํ สตฺถา เอวํ วาที เอวํทิฏฺฐิ “สตฺติเม กายา อกฏา อกฏวิธา อนิมฺมิตา อนิมฺมาตา วญฺฌา กูฏฏฺฐา เอสิกฏฺฐายิฏฺฐิตา.\csromanspacer{} เต น อิญฺชนฺติ, น วิปริณมนฺติ, น อญฺญมญฺญํ พฺยาพาเธนฺติ.\csromanspacer{} นาลํ อญฺญมญฺญสฺส สุขาย วา ทุกฺขาย วา สุขทุกฺขาย วา.\csromanspacer{} กตเม สตฺต, ปถวีกาโย อาโปกาโย เตโชกาโย วาโยกาโย สุเข ทุกฺเข ชีเว สตฺตเม.\csromanspacer{} อิเม สตฺต กายา อกฏา อกฏวิธา อนิมฺมิตา อนิมฺมาตา วญฺฌา กูฏฏฺฐา เอสิกฏฺฐายิฏฺฐิตา.\csromanspacer{} เต น อิญฺชนฺติ น วิปริณมนฺติ น อญฺญมญฺญํ พฺยาพาเธนฺติ.\csromanspacer{} นาลํ อญฺญมญฺญสฺส สุขาย วา ทุกฺขาย วา สุขทุกฺขาย วา.\csromanspacer{} ตตฺถ นตฺถิ หนฺตา}

\vspace{\tipitakaparskip}
\csromancenter{สนฺทกสุตฺต (76) วา ฆาเตตา วา โสตา วา สาเวตา วา วิญฺญาตา วา วิญฺญาเปตา วา.\csromanspacer{} โยปิ ติเณฺหน สตฺเถน สีสํ ฉินฺทติ, น โกจิ กญฺจิ ชีวิตา โวโรเปติ.\csromanspacer{} สตฺตนฺนํเตฺวว กายานมนฺตเรน สตฺถํ วิวรมนุปตติ.\csromanspacer{} จุทฺทส โข ปนิมานิ โยนิปมุขสตสหสฺสานิ สฏฺฐิ จ สตานิ ฉ จ สตานิ ปญฺจ จ กมฺมุโน สตานิ ปญฺจ จ กมฺมานิ ตีณิ จ กมฺมานิ, กมฺเม จ อฑฺฒกมฺเม จ ทฺวฏฺฐิปฏิปทา ทฺวฏฺฐนฺตรกปฺปา ฉฬาภิชาติโย อฏฺฐ ปุริสภูมิโย เอกูนปญฺญาส อาชีวกสเต เอกูนปญฺญาส ปริพฺพาชกสเต เอกูนปญฺญาส นาคาวาสสเต วีเส อินฺทฺริยสเต ติํเส นิรยสเต ฉตฺติํส รโชธาตุโย สตฺต สญฺญีคพฺภา สตฺต อสญฺญีคพฺภา สตฺต นิคณฺฐิคพฺภา สตฺต เทวา สตฺต มานุสา สตฺต เปสาจา สตฺต สรา สตฺต ปวุฏา สตฺต ปปาตา สตฺต ปปาตสตานิ สตฺต สุปินา สตฺต สุปินสตานิ จุลฺลาสีติ มหากปฺปิโน สตสหสฺสานิ, ยานิ พาเล จ ปณฺฑิเต จ สนฺธาวิตฺวา สํสริตฺวา ทุกฺขสฺสนฺตํ กริสฺสนฺติ.\csromanspacer{} ตตฺถ นตฺถิ ‘อิมินาหํ สีเลน วา วเตน วา ตเปน วา พฺรหฺมจริเยน วา อปริปกฺกํ วา กมฺมํ ปริปาเจสฺสามิ, ปริปกฺกํ วา กมฺมํ ผุสฺส ผุสฺส พฺยนฺติํ กริสฺสามี’ติ, เหวํ นตฺถิ, โทณมิเต สุขทุกฺเข, ปริยนฺตกเต สํสาเร, นตฺถิ หายนวฑฺฒเน, นตฺถิ อุกฺกํสาวกํเส.\csromanspacer{} เสยฺยถาปิ นาม สุตฺตคุเฬ ขิตฺเต นิพฺเพฐิยมานเมว ปเลติ, เอวเมว พาเล จ ปณฺฑิเต จ สนฺธาวิตฺวา สํสริตฺวา ทุกฺขสฺสนฺตํ กริสฺสนฺตี”ติ.\csromanspacer{} สเจ ปน อิมสฺส โภโต สตฺถุโน สจฺจํ วจนํ.\csromanspacer{} อกเตน เม เอตฺถ กตํ, อวุสิเตน เม เอตฺถ วุสิตํ, อุโภปิ มยํ เอตฺถ สมสมา สามญฺญํ ปตฺตา.\csromanspacer{} โย จาหํ น วทามิ “อุโภ สนฺธาวิตฺวา สํสริตฺวา ทุกฺขสฺสนฺตํ กริสฺสามา”ติ.\csromanspacer{} อติเรกํ โข ปนิมสฺส โภโต สตฺถุโน นคฺคิยํ มุณฺฑิยํ อุกฺกุฏิกปฺปธานํ เกสมสฺสุโลจนํ.\csromanspacer{} โยหํ ปุตฺตสมฺพาธสยนํ อชฺฌาวสนฺโต กาสิกจนฺทนํ ปจฺจนุโภนฺโต มาลาคนฺธวิเลปนํ ธาเรนฺโต ชาตรูปรชตํ สาทิยนฺโต อิมินา โภตา สตฺถารา สมสมคติโก ภวิสฺสามิ อภิสมฺปรายํ.\csromanspacer{} โสหํ กิํ ชานนฺโต กิํ ปสฺสนฺโต อิมสฺมิํ สตฺถริ พฺรหฺมจริยํ จริสฺสามิ.\csromanspacer{} โส “อพฺรหฺมจริยวาโส อยนฺ”ติ อิติ วิทิตฺวา ตสฺมา พฺรหฺมจริยา นิพฺพิชฺช ปกฺกมติ.\csromanspacer{} อยํ โข สนฺทก เตน ภควตา ชานตา ปสฺสตา อรหตา สมฺมาสมฺพุทฺเธน จตุตฺโถ อพฺรหฺมจริยวาโส อกฺขาโต, ยตฺถ วิญฺญู ปุริโส สสกฺกํ พฺรหฺมจริยํ น วเสยฺย, วสนฺโต จ นาราเธยฺย ญายํ ธมฺมํ กุสลํ.}
\prose{\csromanfolio{188}อิเม โข เต สนฺทก เตน ภควตา ชานตา ปสฺสตา อรหตา สมฺมาสมฺพุทฺเธน จตฺตาโร อพฺรหฺมจริยวาสา อกฺขาตา, ยตฺถ วิญฺญู ปุริโส สสกฺกํ พฺรหฺมจริยํ น วเสยฺย, วสนฺโต จ นาราเธยฺย ญายํ ธมฺมํ กุสลนฺติ.}

\prose{อจฺฉริยํ โภ อานนฺท, อพฺภุตํ โภ อานนฺท, ยาวญฺจิทํ เตน ภควตา ชานตา ปสฺสตา อรหตา สมฺมาสมฺพุทฺเธน จตฺตาโร อพฺรหฺมจริยวาสาว สมานา อพฺรหฺมจริยวาสาติ อกฺขาตา “ยตฺถ วิญฺญู ปุริโส สสกฺกํ พฺรหฺมจริยํ น วเสยฺย, วสนฺโต จ นาราเธยฺย ญายํ ธมฺมํ กุสลนฺ”ติ.\csromanspacer{} กตมานิ ปน ตานิ โภ อานนฺท เตน ภควตา ชานตา ปสฺสตา อรหตา สมฺมาสมฺพุทฺเธน จตฺตาริ อนสฺสาสิกานิ พฺรหฺมจริยานิ อกฺขาตานิ “ยตฺถ วิญฺญู ปุริโส สสกฺกํ พฺรหฺมจริยํ น วเสยฺย, วสนฺโต จ นาราเธยฺย ญายํ ธมฺมํ กุสลนฺ”ติ.}

\proseitem{229}{อิธ สนฺทก เอกจฺโจ สตฺถา สพฺพญฺญู สพฺพทสฺสาวี อปริเสสํ ญาณทสฺสนํ ปฏิชานาติ “จรโต จ เม ติฏฺฐโต จ สุตฺตสฺส จ ชาครสฺส จ สตตํ สมิตํ ญาณทสฺสนํ ปจฺจุปฏฺฐิตนฺ”ติ.\csromanspacer{} โส สุญฺญํปิ อคารํ ปวิสติ, ปิณฺฑํปิ น ลภติ, กุกฺกุโรปิ ฑํสติ, จณฺเฑนปิ หตฺถินา สมาคจฺฉติ, จณฺเฑนปิ อสฺเสน สมาคจฺฉติ, จณฺเฑนปิ โคเณน สมาคจฺฉติ, อิตฺถิยาปิ ปุริสสฺสปิ นามํปิ โคตฺตํปิ ปุจฺฉติ, คามสฺสปิ นิคมสฺสปิ นามํปิ มคฺคํปิ ปุจฺฉติ.\csromanspacer{} โส “กิมิทนฺ”ติ ปุฏฺโฐ สมาโน สุญฺญํ เม อคารํ ปวิสิตพฺพํ อโหสิ, เตน ปาวิสิํ.\csromanspacer{} ปิณฺฑํปิ อลทฺธพฺพํ อโหสิ, เตน นาลตฺถํ.\csromanspacer{} กุกฺกุเรน ฑํสิตพฺพํ อโหสิ, เตนมฺหิ\footnote{เตน (ก), เตนาสิํ (?)} ทฏฺโฐ.\csromanspacer{} จณฺเฑน หตฺถินา สมาคนฺตพฺพํ อโหสิ, เตน สมาคมิํ.\csromanspacer{} จณฺเฑน อสฺเสน สมาคนฺตพฺพํ อโหสิ, เตน สมาคมิํ.\csromanspacer{} จณฺเฑน โคเณน สมาคนฺตพฺพํ อโหสิ, เตน สมาคมิํ.\csromanspacer{} อิตฺถิยาปิ ปุริสสฺสปิ นามํปิ โคตฺตํปิ ปุจฺฉิตพฺพํ อโหสิ, เตน ปุจฺฉิํ.\csromanspacer{} คามสฺสปิ นิคมสฺสปิ นามํปิ มคฺคํปิ ปุจฺฉิตพฺพํ อโหสิ, เตน ปุจฺฉินฺติ.\csromanspacer{} ตตฺร สนฺทก วิญฺญู ปุริโส อิติ ปฏิสญฺจิกฺขติ\csromandash{}อยํ โข ภวํ สตฺถา สพฺพญฺญู สพฺพทสฺสาวี อปริเสสํ ญาณทสฺสนํ ปฏิชานาติ ฯเปฯ.\csromanspacer{} คามสฺสปิ นิคมสฺสปิ นามํปิ มคฺคํปิ ปุจฺฉิตพฺพํ อโหสิ, เตน}

\vspace{\tipitakaparskip}
\csromancenter{สนฺทกสุตฺต (76) ปุจฺฉินฺติ.\csromanspacer{} โส อนสฺสาสิกํ “อิทํ พฺรหฺมจริยนฺ”ติ อิติ วิทิตฺวา ตสฺมา พฺรหฺมจริยา นิพฺพิชฺช ปกฺกมติ.\csromanspacer{} อิทํ โข สนฺทก เตน ภควตา ชานตา ปสฺสตา อรหตา สมฺมาสมฺพุทฺเธน ปฐมํ อนสฺสาสิกํ พฺรหฺมจริยํ อกฺขาตํ, ยตฺถ วิญฺญู ปุริโส สสกฺกํ พฺรหฺมจริยํ น วเสยฺย, วสนฺโต จ นาราเธยฺย ญายํ ธมฺมํ กุสลํ.}
\proseitem{230}{\csromanfolio{189}ปุน จปรํ สนฺทก อิเธกจฺโจ สตฺถา อนุสฺสวิโก โหติ อนุสฺสวสจฺโจ.\csromanspacer{} โส อนุสฺสเวน อิติหิติหปรมฺปราย ปิฏกสมฺปทาย ธมฺมํ เทเสติ.\csromanspacer{} อนุสฺสวิกสฺส โข ปน สนฺทก สตฺถุโน อนุสฺสวสจฺจสฺส สุสฺสุตมฺปิ โหติ, ทุสฺสุตมฺปิ โหติ.\csromanspacer{} ตถาปิ โหติ, อญฺญถาปิ โหติ.\csromanspacer{} ตตฺร สนฺทก วิญฺญู ปุริโส อิติ ปฏิสญฺจิกฺขติ\csromandash{}อยํ โข ภวํ สตฺถา อนุสฺสวิโก อนุสฺสวสจฺโจ, โส อนุสฺสเวน อิติหิติหปรมฺปราย ปิฏกสมฺปทาย ธมฺมํ เทเสติ.\csromanspacer{} อนุสฺสวิกสฺส โข ปน สตฺถุโน อนุสฺสวสจฺจสฺส สุสฺสุตมฺปิ โหติ, ทุสฺสุตมฺปิ โหติ.\csromanspacer{} ตถาปิ โหติ, อญฺญถาปิ โหติ.\csromanspacer{} โส อนสฺสาสิกํ “อิทํ พฺรหฺมจริยนฺ”ติ อิติ วิทิตฺวา ตสฺมา พฺรหฺมจริยา นิพฺพิชฺช ปกฺกมติ.\csromanspacer{} อิทํ โข สนฺทก เตน ภควตา ชานตา ปสฺสตา อรหตา สมฺมาสมฺพุทฺเธน ทุติยํ อนสฺสาสิกํ พฺรหฺมจริยํ อกฺขาตํ.\csromanspacer{} ยตฺถ วิญฺญู ปุริโส สสกฺกํ พฺรหฺมจริยํ น วเสยฺย, วสนฺโต จ นาราเธยฺย ญายํ ธมฺมํ กุสลํ.}

\proseitem{231}{ปุน จปรํ สนฺทก อิเธกจฺโจ สตฺถา ตกฺกี โหติ วีมํสี, โส ตกฺกปริยาหตํ วีมํสานุจริตํ สยํปฏิภานํ ธมฺมํ เทเสติ.\csromanspacer{} ตกฺกิสฺส โข ปน สนฺทก สตฺถุโน วีมํสิสฺส สุตกฺกิตมฺปิ โหติ, ทุตฺตกฺกิตมฺปิ โหติ.\csromanspacer{} ตถาปิ โหติ, อญฺญถาปิ โหติ.\csromanspacer{} ตตฺร สนฺทก วิญฺญู ปุริโส อิติ ปฏิสญฺจิกฺขติ\csromandash{}อยํ โข ภวํ สตฺถา ตกฺกี วีมํสี, โส ตกฺกปริยาหตํ วีมํสานุจริตํ สยํปฏิภานํ ธมฺมํ เทเสติ.\csromanspacer{} ตกฺกิสฺส โข ปน สตฺถุโน วีมํสิสฺส สุตกฺกิตมฺปิ โหติ, ทุตฺตกฺกิตมฺปิ โหติ.\csromanspacer{} ตถาปิ โหติ, อญฺญถาปิ โหติ.\csromanspacer{} โส อนสฺสาสิกํ “อิทํ พฺรหฺมจริยนฺ”ติ อิติ วิทิตฺวา ตสฺมา พฺรหฺมจริยา นิพฺพิชฺช ปกฺกมติ.\csromanspacer{} อิทํ โข สนฺทก เตน ภควตา ชานตา ปสฺสตา อรหตา สมฺมาสมฺพุทฺเธน ตติยํ อนสฺสาสิกํ พฺรหฺมจริยํ อกฺขาตํ, ยตฺถ วิญฺญู ปุริโส สสกฺกํ พฺรหฺมจริยํ น วเสยฺย, วสนฺโต จ นาราเธยฺย ญายํ ธมฺมํ กุสลํ.}

\proseitem{232}{\csromanfolio{190}ปุน จปรํ สนฺทก อิเธกจฺโจ สตฺถา มนฺโท โหติ โมมูโห, โส มนฺทตฺตา โมมูหตฺตา ตตฺถ ตตฺถ\footnote{ตถา ตถา (สี, สฺยา, กํ, อิ)} ปญฺหํ ปุฏฺโฐ สมาโน วาจาวิกฺเขปํ อาปชฺชติ อมราวิกฺเขปํ “เอวนฺติปิ\footnote{เอวมฺปิ (สี, อิ)} เม โน, ตถาติปิ\footnote{ตถาปิ (สี, อิ)} เม โน, อญฺญถาติปิ\footnote{อญฺญถาปิ (สี, อิ) ( ) สพฺพตฺถ นตฺถิ.} เม โน, โนติปิ เม โน, โน โนติปิ เม โน”ติ.\csromanspacer{} ตตฺร สนฺทก วิญฺญู ปุริโส อิติ ปฏิสญฺจิกฺขติ\csromandash{}อยํ โข ภวํ สตฺถา มนฺโท โมมูโห, โส มนฺทตฺตา โมมูหตฺตา ตตฺถ ตตฺถ ปญฺหํ ปุฏฺโฐ สมาโน วาจาวิกฺเขปํ อาปชฺชติ อมราวิกฺเขปํ “เอวนฺติปิ เม โน, ตถาติปิ เม โน, อญฺญถาติปิ เม โน, โนติปิ เม โน, โน โนติปิ เม โน”ติ.\csromanspacer{} โส อนสฺสาสิกํ “อิทํ พฺรหฺมจริยนฺ”ติ อิติ วิทิตฺวา ตสฺมา พฺรหฺมจริยา นิพฺพิชฺช ปกฺกมติ.\csromanspacer{} อิทํ โข สนฺทก เตน ภควตา ชานตา ปสฺสตา อรหตา สมฺมาสมฺพุทฺเธน จตุตฺถํ อนสฺสาสิกํ พฺรหฺมจริยํ อกฺขาตํ, ยตฺถ วิญฺญู ปุริโส สสกฺกํ พฺรหฺมจริยํ น วเสยฺย, วสนฺโต จ นาราเธยฺย ญายํ ธมฺมํ กุสลํ.}

\prose{อิมานิ โข (ตานิ) สนฺทก เตน ภควตา ชานตา ปสฺสตา อรหตา สมฺมาสมฺพุทฺเธน จตฺตาริ อนสฺสาสิกานิ พฺรหฺมจริยานิ อกฺขาตานิ, ยตฺถ วิญฺญู ปุริโส สสกฺกํ พฺรหฺมจริยํ น วเสยฺย, วสนฺโต จ นาราเธยฺย ญายํ ธมฺมํ กุสลนฺติ.}

\prose{อจฺฉริยํ โภ อานนฺท, อพฺภุตํ โภ อานนฺท, ยาวญฺจิทํ เตน ภควตา ชานตา ปสฺสตา อรหตา สมฺมาสมฺพุทฺเธน จตฺตาริ อนสฺสาสิกาเนว พฺรหฺมจริยานิ อนสฺสาสิกานิ พฺรหฺมจริยานีติ อกฺขาตานิ, ยตฺถ วิญฺญู ปุริโส สสกฺกํ พฺรหฺมจริยํ น วเสยฺย, วสนฺโต จ นาราเธยฺย ญายํ ธมฺมํ กุสลํ.\csromanspacer{} โส ปน โภ อานนฺท สตฺถา กิํ วาที กิํ อกฺขายี, ยตฺถ วิญฺญู ปุริโส สสกฺกํ พฺรหฺมจริยํ วเสยฺย, วสนฺโต จ อาราเธยฺย ญายํ ธมฺมํ กุสลนฺติ.}

\proseitem{233}{อิธ สนฺทก ตถาคโต โลเก อุปฺปชฺชติ อรหํ สมฺมาสมฺพุทฺโธ วิชฺชาจรณสมฺปนฺโน สุคโต โลกวิทู อนุตฺตโร ปุริสทมฺมสารถิ สตฺถา เทวมนุสฺสานํ พุทฺโธ ภควา ฯเปฯ.\footnote{วิตฺถาโร กนฺทรกสุตฺเต (7) ปิฏฺเฐ.} โส อิเม ปญฺจ นีวรเณ ปหาย เจตโส อุปกฺกิเลเส ปญฺญาย ทุพฺพลีกรเณ วิวิจฺเจว กาเมหิ วิวิจฺจ}

\vspace{\tipitakaparskip}
\csromancenter{สนฺทกสุตฺต (76) อกุสเลหิ ธมฺเมหิ สวิตกฺกํ สวิจารํ วิเวกชํ ปีติสุขํ ปฐมํ ฌานํ อุปสมฺปชฺช วิหรติ.\csromanspacer{} ยสฺมิํ โข\footnote{ยสฺมิํ โข ปน (สฺยา, กํ, ก)} สนฺทก สตฺถริ สาวโก เอวรูปํ อุฬารวิเสสํ อธิคจฺฉติ, ตตฺถ วิญฺญู ปุริโส สสกฺกํ พฺรหฺมจริยํ วเสยฺย, วสนฺโต จ อาราเธยฺย ญายํ ธมฺมํ กุสลํ.}
\prose{\csromanfolio{191}ปุน จปรํ สนฺทก ภิกฺขุ วิตกฺกวิจารานํ วูปสมา ฯเปฯ ทุติยํ ฌานํ อุปสมฺปชฺช วิหรติ.\csromanspacer{} ยสฺมิํ โข สนฺทก สตฺถริ สาวโก เอวรูปํ อุฬารวิเสสํ อธิคจฺฉติ, ตตฺถ วิญฺญู ปุริโส สสกฺกํ พฺรหฺมจริยํ วเสยฺย, วสนฺโต จ อาราเธยฺย ญายํ ธมฺมํ กุสลํ.}

\prose{ปุน จปรํ สนฺทก ภิกฺขุ ปีติยา จ วิราคา อุเปกฺขโก จ วิหรติ ฯเปฯ ตติยํ ฌานํ อุปสมฺปชฺช วิหรติ.\csromanspacer{} ยสฺมิํ โข สนฺทก สตฺถริ สาวโก เอวรูปํ อุฬารวิเสสํ อธิคจฺฉติ, ตตฺถ วิญฺญู ปุริโส สสกฺกํ พฺรหฺมจริยํ วเสยฺย, วสนฺโต จ อาราเธยฺย ญายํ ธมฺมํ กุสลํ.}

\prose{ปุน จปรํ สนฺทก ภิกฺขุ สุขสฺส จ ปหานา ฯเปฯ จตุตฺถํ ฌานํ อุปสมฺปชฺช วิหรติ.\csromanspacer{} ยสฺมิํ โข สนฺทก สตฺถริ สาวโก เอวรูปํ อุฬารวิเสสํ อธิคจฺฉติ, ตตฺถ วิญฺญู ปุริโส สสกฺกํ พฺรหฺมจริยํ วเสยฺย, วสนฺโต จ อาราเธยฺย ญายํ ธมฺมํ กุสลํ.}

\prose{โส เอวํ สมาหิเต จิตฺเต ปริสุทฺเธ ปริโยทาเต อนงฺคเณ วิคตูปกฺกิเลเส มุทุภูเต กมฺมนิเย ฐิเต อาเนญฺชปฺปตฺเต ปุพฺเพนิวาสานุสฺสติญาณาย จิตฺตํ อภินินฺนาเมติ, โส อเนกวิหิตํ ปุพฺเพนิวาสํ อนุสฺสรติ.\csromanspacer{} เสยฺยถิทํ, เอกมฺปิ ชาติํ เทฺวปิ ชาติโย ฯเปฯ อิติ สาการํ ส-อุทฺเทสํ อเนกวิหิตํ ปุพฺเพนิวาสํ อนุสฺสรติ.\csromanspacer{} ยสฺมิํ โข สนฺทก สตฺถริ สาวโก เอวรูปํ อุฬารวิเสสํ อธิคจฺฉติ, ตตฺถ วิญฺญู ปุริโส สสกฺกํ พฺรหฺมจริยํ วเสยฺย, วสนฺโต จ อาราเธยฺย ญายํ ธมฺมํ กุสลํ.}

\prose{โส เอวํ สมาหิเต จิตฺเต ปริสุทฺเธ ปริโยทาเต อนงฺคเณ วิคตูปกฺกิเลเส มุทุภูเต กมฺมนิเย ฐิเต อาเนญฺชปฺปตฺเต สตฺตานํ จุตูปปาตญาณาย จิตฺตํ อภินินฺนาเมติ, โส ทิพฺเพน จกฺขุนา วิสุทฺเธน อติกฺกนฺตมานุสเกน สตฺเต ปสฺสติ จวมาเน อุปปชฺชมาเน หีเน ปณีเต สุวณฺเณ ทุพฺพณฺเณ สุคเต ทุคฺคเต ฯเปฯ ยถากมฺมูปเค สตฺเต ปชานาติ.\csromanspacer{} ยสฺมิํ โข สนฺทก สตฺถริ สาวโก เอวรูปํ อุฬารวิเสสํ \csromanfolio{192}อธิคจฺฉติ, ตตฺถ วิญฺญู ปุริโส สสกฺกํ พฺรหฺมจริยํ วเสยฺย, วสนฺโต จ อาราเธยฺย ญายํ ธมฺมํ กุสลํ.}

\prose{โส เอวํ สมาหิเต จิตฺเต ปริสุทฺเธ ปริโยทาเต อนงฺคเณ วิคตูปกฺกิเลเส มุทุภูเต กมฺมนิเย ฐิเต อาเนญฺชปฺปตฺเต อาสวานํ ขยญาณาย จิตฺตํ อภินินฺนาเมติ.\csromanspacer{} โส อิทํ ทุกฺขนฺติ ยถาภูตํ ปชานาติ, อยํ ทุกฺขสมุทโยติ ยถาภูตํ ปชานาติ, อยํ ทุกฺขนิโรโธติ ยถาภูตํ ปชานาติ, อยํ ทุกฺขนิโรธคามินี ปฏิปทาติ ยถาภูตํ ปชานาติ.\csromanspacer{} อิเม อาสวาติ ยถาภูตํ ปชานาติ, อยํ อาสวสมุทโยติ ยถาภูตํ ปชานาติ, อยํ อาสวนิโรโธติ ยถาภูตํ ปชานาติ, อยํ อาสวนิโรธคามินี ปฏิปทาติ ยถาภูตํ ปชานาติ.\csromanspacer{} ตสฺส เอวํ ชานโต เอวํ ปสฺสโต กามาสวาปิ จิตฺตํ วิมุจฺจติ, ภวาสวาปิ จิตฺตํ วิมุจฺจติ, อวิชฺชาสวาปิ จิตฺตํ วิมุจฺจติ วิมุตฺตสฺมิํ วิมุตฺตมิติ ญาณํ โหติ, “ขีณา ชาติ, วุสิตํ พฺรหฺมจริยํ กตํ กรณียํ, นาปรํ อิตฺถตฺตายา”ติ ปชานาติ.\csromanspacer{} ยสฺมิํ โข สนฺทก สตฺถริ สาวโก เอวรูปํ อุฬารวิเสสํ อธิคจฺฉติ, ตตฺถ วิญฺญู ปุริโส สสกฺกํ พฺรหฺมจริยํ วเสยฺย, วสนฺโต จ อาราเธยฺย ญายํ ธมฺมํ กุสลนฺติ.}

\proseitem{234}{โย ปน โส โภ อานนฺท ภิกฺขุ อรหํ ขีณาสโว วุสิตวา กตกรณีโย โอหิตภาโร อนุปฺปตฺตสทตฺโถ ปริกฺขีณภวสํโยชโน สมฺมทญฺญา วิมุตฺโต, ปริภุญฺเชยฺย โส กาเมติ? โย โส สนฺทก ภิกฺขุ อรหํ ขีณาสโว วุสิตวา กตกรณีโย โอหิตภาโร อนุปฺปตฺตสทตฺโถ ปริกฺขีณภวสํโยชโน สมฺมทญฺญา วิมุตฺโต, อภพฺโพ โส ปญฺจฏฺฐานานิ อชฺฌาจริตุํ, อภพฺโพ ขีณาสโว ภิกฺขุ สญฺจิจฺจ ปาณํ ชีวิตา โวโรเปตุํ, อภพฺโพ ขีณาสโว ภิกฺขุ อทินฺนํ เถยฺยสงฺขาตํ อาทาตุํ, อภพฺโพ ขีณาสโว ภิกฺขุ เมถุนํ ธมฺมํ ปฏิเสเวตุํ, อภพฺโพ ขีณาสโว ภิกฺขุ สมฺปชานมุสา ภาสิตุํ อภพฺโพ ขีณาสโว ภิกฺขุ สนฺนิธิการกํ กาเม ปริภุญฺชิตุํ เสยฺยถาปิ ปุพฺเพ อคาริยภูโต.\csromanspacer{} โย โส สนฺทก ภิกฺขุ อรหํ ขีณาสโว วุสิตวา กตกรณีโย โอหิตภาโร อนุปฺปตฺตสทตฺโถ ปริกฺขีณภวสํโยชโน สมฺมทญฺญา วิมุตฺโต, อภพฺโพ อิมานิ ปญฺจฏฺฐานานิ อชฺฌา จริตุนฺติ.}

\csromanheader{2}{6. สนฺทกสุตฺต (76)}
\proseitem{235}{\csromanfolio{193}โย ปน โส โภ อานนฺท ภิกฺขุ อรหํ ขีณาสโว วุสิตวา กตกรณีโย โอหิตภาโร อนุปฺปตฺตสทตฺโถ ปริกฺขีณ ภวสํโยชโน สมฺมทญฺญา วิมุตฺโต, ตสฺส จรโต เจว ติฏฺฐโต จ สุตฺตสฺส จ ชาครสฺส จ สตตํ สมิตํ ญาณทสฺสนํ ปจฺจุปฏฺฐิตํ “ขีณา เม อาสวา”ติ? เตน หิ สนฺทก อุปมํ เต กริสฺสามิ.\csromanspacer{} อุปมายปิเธกจฺเจ วิญฺญู ปุริสา ภาสิตสฺส อตฺถํ อาชานนฺติ.\csromanspacer{} เสยฺยถาปิ สนฺทก ปุริสสฺส หตฺถปาทา ฉินฺนา, ตสฺส จรโต เจว ติฏฺฐิตา จ สุตฺตสฺส จ ชาครสฺส จ สตตํ สมิตํ (ชานาติ “ฉินฺนา เม หตฺถปาทา”ติ.\csromanspacer{} อุทาหุ ปจฺจเวกฺขมาโน ชานาติ “ฉินฺนา เม หตฺถปาทา”ติ? น โข โภ อานนฺท โส ปุริโส สตตํ สมิตํ ชานาติ “ฉินฺนา เม หตฺถปาทา”ติ.)\footnote{(ฉินฺนาว หตฺถปาทา,) (สี, สฺยา, กํ, อิ)} อปิ จ โข ปน นํ ปจฺจเวกฺขมาโน ชานาติ “ฉินฺนา เม หตฺถปาทา”ติ.\csromanspacer{} เอวเมว โข สนฺทก โย โส ภิกฺขุ อรหํ ขีณาสโว วุสิตวา กตกรณีโย โอหิตภาโร อนุปฺปตฺตสทตฺโถ ปริกฺขีณภวสํโยชโน สมฺมทญฺญา วิมุตฺโต, ตสฺส จรโต เจว ติฏฺฐิโต จ สุตฺตสฺส จ สุตฺตสฺส จ ชาครสฺส จ สตตํ สมิตํ (ญาณทสฺสนํ น ปจฺจุปฏฺฐิตํ “ขีณา เม อาสวา”ติ.)\footnote{(ขีณาว อาสวา,) (สี, สฺยา, กํ, อิ)} อปิ จ โข ปน นํ ปจฺจเวกฺขมาโน ชานาติ “ขีณา เม อาสวา”ติ.}

\proseitem{236}{กีวพหุกา ปน โภ อานนฺท อิมสฺมิํ ธมฺมวินเย นิยฺยาตาโรติ.\csromanspacer{} น โข สนฺทก เอกํเยว สตํ น เทฺว สตานิ น ตีณิ สตานิ น จตฺตาริ สตานิ น ปญฺจ สตานิ, อถ โข ภิยฺโยว เย อิมสฺมิํ ธมฺมวินเย นิยฺยาตาโรติ.\csromanspacer{} อจฺฉริยํ โภ อานนฺท, อพฺภุตํ โภ อานนฺท, น จ นาม สธมฺโมกฺกํสนา ภวิสฺสติ, น ปรธมฺมวมฺภนา.\csromanspacer{} อายตเน จ ธมฺมเทสนา, ตาว พหุกา จ นิยฺยาตาโร ปญฺญายิสฺสนฺติ.\csromanspacer{} อิเม ปนาชีวกา ปุตฺตมตาย ปุตฺตา อตฺตานญฺเจว อุกฺกํเสนฺติ ปเร จ วมฺเภนฺติ.\csromanspacer{} ตโย เจว นิยฺยาตาโร ปญฺญเปนฺติ.\csromanspacer{} เสยฺยถิทํ, นนฺทํ วจฺฉํ, กิสํ สํกิจฺจํ, มกฺขลิํ โคสาลนฺติ.\csromanspacer{} อถ โข สนฺทโก ปริพฺพาชโก สกํ ปริสํ อามนฺเตสิ “จรนฺตุ โภนฺโต, สมเณ โคตเม พฺรหฺมจริยวาโส, น ทานิ สุกรํ อเมฺหหิ ลาภสกฺการสิโลเก ปริจฺจชิตุนฺ”ติ.\csromanspacer{} อิติ หิทํ สนฺทโก ปริพฺพาชโก สกํ ปริสํ อุยฺโยเชสิ ภควติ พฺรหฺมจริเยติ.}

\nitthitam{สนฺทกสุตฺตํ นิฏฺฐิตํ ฉฏฺฐํ.}
\csromanmatikaanchor{mtk.32}
\csromantocmarkref{subsection}{7. มหาสกุลุทายิสุตฺต}{mtk.32}
\bookmark[dest={mtk.32},level=2]{7. มหาสกุลุทายิสุตฺต}
\csromanheader{2}{7. มหาสกุลุทายิสุตฺต}
\proseitem{237}{\csromanfolio{194}เอวํ เม สุตํ\csromandash{}เอกํ สมยํ ภควา ราชคเห วิหรติ เวฬุวเน กลนฺทกนิวาเป.\csromanspacer{} เตน โข ปน สมเยน สมฺพหุลา อภิญฺญาตา อภิญฺญาตา ปริพฺพาชกา โมรนิวาเป ปริพฺพาชการาเม ปฏิวสนฺติ.\csromanspacer{} เสยฺยถิทํ, อนฺนภาโร วรธโร สกุลุทายี จ ปริพฺพาชโก อญฺเญ จ อภิญฺญาตา อภิญฺญาตา ปริพฺพาชกา.\csromanspacer{} อถ โข ภควา ปุพฺพณฺหสมยํ นิวาเสตฺวา ปตฺตจีวรมาทาย ราชคหํ ปิณฺฑาย ปาวิสิ.\csromanspacer{} อถ โข ภควโต เอตทโหสิ “อติปฺปโค โข ตาว ราชคเห ปิณฺฑาย จริตุํ, ยํนูนาหํ เยน โมรนิวาโป ปริพฺพาชการาโม, เยน สกุลุทายี ปริพฺพาชโก เตนุปสงฺกเมยฺยนฺ”ติ.\csromanspacer{} อถ โข ภควา เยน โมรนิวาโป ปริพฺพาชการาโม เตนุปสงฺกมิ.\csromanspacer{} เตน โข ปน สมเยน สกุลุทายี ปริพฺพาชโก มหติยา ปริพฺพาชกปริสาย สทฺธิํ นิสินฺโน โหติ อุนฺนาทินิยา อุจฺจาสทฺทมหาสทฺทาย อเนกวิหิตํ ติรจฺฉานกถํ กเถนฺติยา.\csromanspacer{} เสยฺยถิทํ, ราชกถํ โจรกถํ มหามตฺตกถํ เสนากถํ ภยกถํ ยุทฺธกถํ อนฺนกถํ ปานกถํ วตฺถกถํ สยนกถํ มาลากถํ คนฺธกถํ ญาติกถํ ยานกถํ คามกถํ นิคมกถํ นครกถํ ชนปทกถํ อิตฺถิกถํ สูรกถํ วิสิขากถํ กุมฺภฏฺฐานกถํ ปุพฺพเปตกถํ นานตฺตกถํ โลกกฺขายิกํ สมุทฺทกฺขายิกํ อิติภวาภวกถํ อิติ วา.\csromanspacer{} อทฺทสา โข สกุลุทายี ปริพฺพาชโก ภควนฺตํ ทูรโตว อาคจฺฉนฺตํ, ทิสฺวาน สกํ ปริสํ สณฺฐาเปติ “อปฺปสทฺทา โภนฺโต โหนฺตุ, มา โภนฺโต สทฺทมกตฺถ, อยํ สมโณ โคตโม อาคจฺฉติ, อปฺปสทฺทกาโม โข ปน โส อายสฺมา อปฺปสทฺทสฺส วณฺณวาที, อปฺเปว นาม อปฺปสทฺทํ ปริสํ วิทิตฺวา อุปสงฺกมิตพฺพํ มญฺเญยฺยา”ติ.\csromanspacer{} อถ โข เต ปริพฺพาชกา ตุณฺหี อเหสุํ.\csromanspacer{} อถ โข ภควา เยน สกุลุทายี ปริพฺพาชโก เตนุปสงฺกมิ.\csromanspacer{} อถ โข สกุลุทายี ปริพฺพาชโก ภควนฺตํ เอตทโวจ “เอตุ โข ภนฺเต ภควา, สฺวาคตํ ภนฺเต ภควโต, จิรสฺสํ โข ภนฺเต ภควา อิมํ ปริยายมกาสิ ยทิทํ อิธาคมนาย, นิสีทตุ ภนฺเต ภควา อิทมาสนํ ปญฺญตฺตนฺ”ติ.\csromanspacer{} นิสีทิ ภควา ปญฺญตฺเต อาสเน.\csromanspacer{} สกุลุทายีปิ โข ปริพฺพาชโก อญฺญตรํ นีจํ อาสนํ คเหตฺวา เอกมนฺตํ นิสีทิ.\csromanspacer{} เอกมนฺตํ นิสินฺนํ โข สกุลุทายิํ ปริพฺพาชกํ ภควา เอตทโวจ\csromandash{}}

\csromanheader{2}{7. มหาสกุลุทายีสุตฺต (77)}
\proseitem{238}{\csromanfolio{195}กายนุตฺถ อุทายิ เอตฺรหิ กถาย สนฺนิสินฺนา, กา จ ปน โว อนฺตรากถา วิปฺปกตาติ.\csromanspacer{} ติฏฺฐเตสา ภนฺเต กถา, ยาย มยํ เอตรหิ กถาย สนฺนิสินฺนา.\csromanspacer{} เนสา ภนฺเต กถา ภควโต ทุลฺลภา ภวิสฺสติ ปจฺฉาปิ สวนาย.\csromanspacer{} ปุริมานิ ภนฺเต ทิวสานิ ปุริมตรานิ นานาติตฺถิยานํ สมณพฺราหฺมณานํ กุตูหลสาลายํ สนฺนิสินฺนานํ สนฺนิปติตานํ อยมนฺตรากถา อุทปาทิ “ลาภา วต โภ องฺคมคธานํ, สุลทฺธลาภา วต โภ องฺคมคธานํ.\csromanspacer{} ตตฺริเม\footnote{ยตฺถิเม (สี)} สมณพฺราหฺมณา สงฺฆิโน คณิโน คณาจริยา ญาตา ยสสฺสิโน ติตฺถกรา สาธุสมฺมตา พหุชนสฺส, ราชคหํ วสฺสาวาสํ โอสฏา.\csromanspacer{} อยมฺปิ โข ปูรโณ กสฺสโป สงฺฆี เจว คณี จ คณาจริโย จ ญาโต ยสสฺสี ติตฺถกโร สาธุสมฺมโต พหุชนสฺส, โสปิ ราชคหํ วสฺสาวาสํ โอสโฏ.\csromanspacer{} อยมฺปิ โข มกฺขลิ โคสาโล ฯเปฯ อชิโต เกสกมฺพโล.\csromanspacer{} ปกุโธ กจฺจายโน.\csromanspacer{} สญฺชโย เพลฏฺฐปุตฺโต.\csromanspacer{} นิคณฺโฐ นาฏปุตฺโต สงฺฆี เจว คณี จ คณาจริโย จ ญาโต ยสสฺสี ติตฺถกโร สาธุสมฺมโต พหุชนสฺส, โสปิ ราชคหํ วสฺสาวาสํ โอสโฏ.\csromanspacer{} อยมฺปิ โข สมโณ โคตโม สงฺฆี เจว คณี จ คณาจริโย จ ญาโต ยสสฺสี ติตฺถกโร สาธุสมฺมโต พหุชนสฺส, โสปิ ราชคหํ วสฺสาวาสํ โอสโฏ.\csromanspacer{} โก นุ โข อิเมสํ ภวตํ สมณพฺราหฺมณานํ สงฺฆีนํ คณีนํ คณาจริยานํ ญาตานํ ยสสฺสีนํ ติตฺถกรานํ สาธุสมฺมตานํ พหุชนสฺส สาวกานํ สกฺกโต ครุกโต มานิโต ปูชิโต.\csromanspacer{} กญฺจ ปน สาวกา สกฺกตฺวา ครุํ กตฺวา\footnote{ครุกตฺวา (สี, สฺยา, กํ, อิ)} อุปนิสฺสาย วิหรนฺตี”ติ.}

\proseitem{239}{ตเตฺรกจฺเจ เอวมาหํสุ “อยํ โข ปูรโณ กสฺสโป สงฺฆี เจว คณี จ คณาจริโย จ ญาโต ยสสฺสี ติตฺถกโร สาธุสมฺมโต พหุชนสฺส.\csromanspacer{} โส จ โข สาวกานํ น สกฺกโต น ครุกโต น มานิโต น ปูชิโต.\csromanspacer{} น จ ปน ปูรณํ กสฺสปํ สาวกา สกฺกตฺวา ครุํ กตฺวา อุปนิสฺสาย วิหรนฺติ.\csromanspacer{} ภูตปุพฺพํ ปูรโณ กสฺสโป อเนกสตาย ปริสาย ธมฺมํ เทเสติ.\csromanspacer{} ตตฺรญฺญตโร ปูรณสฺส กสฺสปสฺส สาวโก สทฺทมกาสิ ‘มา โภนฺโต ปูรณํ กสฺสปํ เอตมตฺถํ ปุจฺฉิตฺถ, เนโส เอตํ ชานาติ, มยเมตํ ชานาม, อเมฺห เอตมตฺถํ ปุจฺฉถ, มยเมตํ ภวนฺตานํ พฺยากริสฺสามา’ติ.\csromanspacer{} ภูตปุพฺพํ ปูรโณ \csromanfolio{196}กสฺสโป พาหา ปคฺคยฺห กนฺทนฺโต น ลภติ ‘อปฺปสทฺทา โภนฺโต โหนฺตุ, มา โภนฺโต สทฺทมกตฺถ, เนเต ภวนฺเต ปุจฺฉนฺติ.\csromanspacer{} อเมฺห เอเต ปุจฺฉนฺติ, มยเมเตสํ พฺยากริสฺสามา’ติ.\csromanspacer{} พหู โข ปน ปูรณสฺส กสฺสปสฺส สาวกา วาทํ อาโรเปตฺวา อปกฺกนฺตา ‘น ตฺวํ อิมํ ธมฺมวินยํ อาชานาสิ, อหํ อิมํ ธมฺมวินยํ อาชานามิ.\csromanspacer{} กิํ ตฺวํ อิมํ ธมฺมวินยํ อาชานิสฺสสิ, มิจฺฉาปฏิปนฺโน ตฺวมสิ, อหมสฺมิ สมฺมาปฏิปนฺโน, สหิตํ เม, อสหิตํ เต, ปุเรวจนียํ ปจฺฉา อวจ, ปจฺฉาวจนียํ ปุเร อวจ, อธิจิณฺณํ เต วิปราวตฺตํ, อาโรปิโต เต วาโท, นิคฺคหิโตสิ.\csromanspacer{} จร วาทปฺปโมกฺขาย, นิพฺเพเฐหิ วา สเจ ปโหสี’ติ.\csromanspacer{} อิติ ปูรโณ กสฺสโป สาวกานํ น สกฺกโต น ครุกโต น มานิโต น ปูชิโต.\csromanspacer{} น จ ปน ปูรณํ กสฺสปํ สาวกา สกฺกตฺวา ครุํ กตฺวา อุปนิสฺสาย วิหรนฺติ.\csromanspacer{} อกฺกุฏฺโฐ จ ปน ปูรโณ กสฺสโป ธมฺมกฺโกเสนา”ติ.}

\prose{เอกจฺเจ เอวมาหํสุ “อยมฺปิ โข มกฺขลิ โคสาโล ฯเปฯ อชิโต เกสกมฺพโล.\csromanspacer{} ปกุโธ กจฺจายโน.\csromanspacer{} สญฺชโย เพลฏฺฐปุตฺโต นิคณฺโฐ นาฏปุตฺโต สงฺฆี เจว คณี จ คณาจริโย จ ญาโต ยสสฺสี ติตฺถกโร สาธุสมฺมโต พหุชนสฺส.\csromanspacer{} โส จ โข สาวกานํ น สกฺกโต น ครุกโต น มานิโต น ปูชิโต.\csromanspacer{} น จ ปน นิคณฺฐํ นาฏปุตฺตํ สาวกา สกฺกตฺวา ครุํ กตฺวา อุปนิสฺสาย วิหรนฺติ.\csromanspacer{} ภูตปุพฺพํ นิคณฺโฐ นาฏปุตฺโต อเนกสตาย ปริสาย ธมฺมํ เทเสติ.\csromanspacer{} ตตฺรญฺญตโร นิคณฺฐสฺส นาฏปุตฺตสฺส สาวโก สทฺทมกาสิ ‘มา โภนฺโต นิคณฺฐํ นาฏปุตฺตํ เอตมตฺถํ ปุจฺฉิตฺถ, เนโส เอตํ ชานาติ, มยเมตํ ชานาม, อเมฺห เอตมตฺถํ ปุจฺฉถ, มยเมตํ ภวนฺตานํ พฺยากริสฺสามา’ติ.\csromanspacer{} ภูตปุพฺพํ นิคณฺโฐ นาฏปุตฺโต พาหา ปคฺคยฺห กนฺทนฺโต น ลภติ “อปฺปสทฺทา โภนฺโต โหนฺตุ, มา โภนฺโต สทฺทมกตฺถ, เนเต ภวนฺเต ปุจฺฉนฺติ, อเมฺห เอเต ปุจฺฉนฺติ, มยเมเตสํ พฺยากริสฺสามา”ติ.\csromanspacer{} พหู โข ปน นิคณฺฐสฺส นาฏปุตฺตสฺส สาวกา วาทํ อาโรเปตฺวา อปกฺกนฺตา ‘น ตฺวํ อิมํ ธมฺมวินยํ อาชานาสิ, อหํ อิมํ ธมฺมวินยํ อาชานามิ.\csromanspacer{} กิํ ตฺวํ อิมํ ธมฺมวินยํ อาชานิสฺสสิ, มิจฺฉาปฏิปนฺโน ตฺวมสิ.\csromanspacer{} อหมสฺมิ สมฺมาปฏิปนฺโน.\csromanspacer{} สหิตํ เม อสหิตํ เต, ปุเรวจนียํ ปจฺฉา อวจ, ปจฺฉาวจนียํ ปุเร อวจ, อธิจิณฺณํ เต วิปราวตฺตํ, อาโรปิโต เต วาโท, นิคฺคหิโตสิ, จร วาทปฺปโมกฺขาย,}

\vspace{\tipitakaparskip}
\csromancenter{มหาสกุลุทายีสุตฺต (77) นิพฺเพเฐหิ วา สเจ ปโหสี’ติ.\csromanspacer{} อิติ นิคณฺโฐ นาฏปุตฺโต สาวกานํ น สกฺกโต น ครุกโต น มานิโต น ปูชิโต.\csromanspacer{} น จ ปน นิคณฺฐํ นาฏปุตฺตํ สาวกา สกฺกตฺวา ครุํ กตฺวา อุปนิสฺสาย วิหรนฺติ.\csromanspacer{} อกฺกุฏฺโฐ จ ปน นิคณฺโฐ นาฏปุตฺโต ธมฺมกฺโกเสนา”ติ.}
\proseitem{240}{\csromanfolio{197}เอกจฺเจ เอวมาหํสุ “อยมฺปิ โข สมโณ โคตโม สงฺฆี เจว คณี จ คณาจริโย จ ญาโต ยสสฺสี ติตฺถกโร สาธุสมฺมโต พหุชนสฺส.\csromanspacer{} โส จ โข สาวกานํ สกฺกโต ครุกโต มานิโต ปูชิโต.\csromanspacer{} สมณญฺจ ปน โคตมํ สาวกา สกฺกตฺวา ครุํ กตฺวา อุปนิสฺสาย วิหรนฺติ.\csromanspacer{} ภูตปุพฺพํ สมโณ โคตโม อเนกสตาย ปริสาย ธมฺมํ เทเสสิ.\csromanspacer{} ตตฺรญฺญตโร สมณสฺส โคตมสฺส สาวโก อุกฺกาสิ.\csromanspacer{} ตเมนาญฺญตโร สพฺรหฺมจารี ชณฺณุเกน\footnote{ชณฺณุเก (สี)} ฆฏฺเฏสิ ‘อปฺปสทฺโท อายสฺมา โหตุ, มายสฺมา สทฺทมกาสิ, สตฺถา โน ภควา ธมฺมํ เทเสสี’ติ.\csromanspacer{} ยสฺมิํ สมเย สมโณ โคตโม อเนกสตาย ปริสาย ธมฺมํ เทเสติ.\csromanspacer{} เนว ตสฺมิํ สมเย สมณสฺส โคตมสฺส สาวกานํ ขิปิตสทฺโท วา โหติ อุกฺกาสิตสทฺโท วา, ตเมนํ มหาชนกาโย ปจฺจาสีสมานรูโป\footnote{ปจฺจาสิํ สมานรูโป (สี, สฺยา, กํ, อิ) 3. ขุทฺทํ มธุํ (สี, สฺยา, กํ, อิ)} ปจฺจุปฏฺฐิโต โหติ ‘ยํ โน ภควา ธมฺมํ ภาสิสฺสติ, ตํ โน โสสฺสามา’ติ.\csromanspacer{} เสยฺยถาปิ นาม ปุริโส จาตุมฺมหาปเถ ขุทฺทมธุํ3 อเนลกํ ปีเฬยฺย\footnote{อุปฺปีเฬยฺย (สี)}, ตเมนํ มหาชนกาโย ปจฺจาสีสมานรูโป ปจฺจุปฏฺฐิโต อสฺส.\csromanspacer{} เอวเมว ยสฺมิํ สมเย สมโณ โคตโม อเนกสตาย ปริสาย ธมฺมํ เทเสติ.\csromanspacer{} เนว ตสฺมิํ สมเย สมณสฺส โคตมสฺส สาวกานํ ขิปิตสทฺโท วา โหติ อุกฺกาสิตสทฺโท วา, ตเมนํ มหาชนกาโย ปจฺจาสีสมานรูโป ปจฺจุปฏฺฐิโต โหติ ‘ยํ โน ภควา ธมฺมํ ภาสิสฺสติ, ตํ โน โสสฺสามา’ติ.\csromanspacer{} เยปิ สมณสฺส โคตมสฺส สาวกา สพฺรหฺมจารีหิ สมฺปโยเชตฺวา สิกฺขํ ปจฺจกฺขาย หีนายาวตฺตนฺติ, เตปิ สตฺถุ เจว วณฺณวาทิโน โหนฺติ, ธมฺมสฺส จ วณฺณวาทิโน โหนฺติ, สํฆสฺส จ วณฺณวาทิโน โหนฺติ.\csromanspacer{} อตฺตครหิโนเยว โหนฺติ อนญฺญครหิโน ‘มยเมวมฺหา อลกฺขิกา, มยํ อปฺปปุญฺญา, เต มยํ เอวํ สฺวากฺขาเต ธมฺมวินเย ปพฺพชิตฺวา นาสกฺขิมฺหา \csromanfolio{198}ยาวชีวํ ปริปุณฺณํ ปริสุทฺธํ พฺรหฺมจริยํ จริตุนฺ’ติ.\csromanspacer{} เต อารามิกภูตา วา อุปาสกภูตา วา ปญฺจสิกฺขาปเท สมาทาย วตฺตนฺติ.\csromanspacer{} อิติ สมโณ โคตโม สาวกานํ สกฺกโต ครุกโต มานิโต ปูชิโต.\csromanspacer{} สมณญฺจ ปน โคตมํ สาวกา สกฺกตฺวา ครุํ กตฺวา อุปนิสฺสาย วิหรนฺตี”ติ.}

\proseitem{241}{กติ ปน ตฺวํ อุทายิ มยิ ธมฺเม สมนุปสฺสสิ, เยหิ มมํ\footnote{มม (สพฺพตฺถ)} สาวกา สกฺกโรนฺติ ครุํ กโรนฺติ\footnote{ครุกโรนฺติ (สี, สฺยา, กํ, อิ)} มาเนนฺติ ปูเชนฺติ, สกฺกตฺวา ครุํ กตฺวา อุปนิสฺสาย วิหรนฺตีติ.\csromanspacer{} ปญฺจ โข อหํ ภนฺเต ภควติ ธมฺเม สมนุปสฺสามิ, เยหิ ภควนฺตํ สาวกา สกฺกโรนฺติ ครุํ กโรนฺติ มาเนนฺติ ปูเชนฺติ, สกฺกตฺวา ครุํ กตฺวา อุปนิสฺสาย วิหรนฺติ.\csromanspacer{} กตเม ปญฺจ, ภควา หิ ภนฺเต อปฺปาหาโร, อปฺปาหารตาย จ วณฺณวาที.\csromanspacer{} ยมฺปิ ภนฺเต ภควา อปฺปาหาโร, อปฺปาหารตาย จ วณฺณวาที, อิมํ โข อหํ ภนฺเต ภควติ ปฐมํ ธมฺมํ สมนุปสฺสามิ, เยน ภควนฺตํ สาวกา สกฺกโรนฺติ ครุํ กโรนฺติ มาเนนฺติ ปูเชนฺติ, สกฺกตฺวา ครุํ กตฺวา อุปนิสฺสาย วิหรนฺติ. (1)}

\prose{ปุน จปรํ ภนฺเต ภควา สนฺตุฏฺโฐ อิตรีตเรน จีวเรน, อิตรีตรจีวรสนฺตุฏฺฐิยา จ วณฺณวาที.\csromanspacer{} ยมฺปิ ภนฺเต ภควา สนฺตุฏฺโฐ อิตรีตเรน จีวเรน, อิตรีตรจีวรสนฺตุฏฺฐิยา จ วณฺณวาที, อิมํ โข อหํ ภนฺเต ภควติ ทุติยํ ธมฺมํ สมนุปสฺสามิ, เยน ภควนฺตํ สาวกา สกฺกโรนฺติ ครุํ กโรนฺติ มาเนนฺติ ปูเชนฺติ, สกฺกตฺวา ครุํ กตฺวา อุปนิสฺสาย วิหรนฺติ. (2)}

\prose{ปุน จปรํ ภนฺเต ภควา สนฺตุฏฺโฐ อิตรีตเรน ปิณฺฑปาเตน, อิตรีตรปิณฺฑปาตสนฺตุฏฺฐิยา จ วณฺณวาที.\csromanspacer{} ยมฺปิ ภนฺเต ภควา สนฺตุฏฺโฐ อิตรีตเรน ปิณฺฑปาเตน, อิตรีตรปิณฺฑปาตสนฺตุฏฺฐิยา จ วณฺณวาที, อิมํ โข อหํ ภนฺเต ภควติ ตติยํ ธมฺมํ สมนุปสฺสามิ, เยน ภควนฺตํ สาวกา สกฺกโรนฺติ ครุํ กโรนฺติ มาเนนฺติ ปูเชนฺติ, สกฺกตฺวา ครุํ กตฺวา อุปนิสฺสาย วิหรนฺติ. (3)}

\prose{ปุน จปรํ ภนฺเต ภควา สนฺตุฏฺโฐ อิตรีตเรน เสนาสเนน, อิตรีตรเสนาสนสนฺตุฏฺฐิยา จ วณฺณวาที.\csromanspacer{} ยมฺปิ ภนฺเต ภควา สนฺตุฏฺโฐ อิตรีตเรน เสนาสเนน, อิตรีตรเสนาสนสนฺตุฏฺฐิยา จ วณฺณวาที,}

\vspace{\tipitakaparskip}
\csromancenter{มหาสกุลุทายีสุตฺต (77) อิมํ โข อหํ ภนฺเต ภควติ จตุตฺถํ ธมฺมํ สมนุปสฺสามิ, เยน ภควนฺตํ สาวกา สกฺกโรนฺติ ครุํ กโรนฺติ มาเนนฺติ ปูเชนฺติ, สกฺกตฺวา ครุํ กตฺวา อุปนิสฺสาย วิหรนฺติ. (4)}
\prose{\csromanfolio{199}ปุน จปรํ ภนฺเต ภควา ปวิวิตฺโต, ปวิเวกสฺส จ วณฺณวาที.\csromanspacer{} ยมฺปิ ภนฺเต ภควา ปวิวิตฺโต, ปวิเวกสฺส จ วณฺณวาที, อิมํ โข อหํ ภนฺเต ภควติ ปญฺจมํ ธมฺมํ สมนุปสฺสามิ, เยน ภควนฺตํ สาวกา สกฺกโรนฺติ ครุํ กโรนฺติ มาเนนฺติ ปูเชนฺติ, สกฺกตฺวา ครุํ กตฺวา อุปนิสฺสาย วิหรนฺติ. (5)}

\prose{อิเม โข อหํ ภนฺเต ภควติ ปญฺจ ธมฺเม สมนุปสฺสามิ, เยหิ ภควนฺตํ สาวกา สกฺกโรนฺติ ครุํ กโรนฺติ มาเนนฺติ ปูเชนฺติ, สกฺกตฺวา ครุํ กตฺวา อุปนิสฺสาย วิหรนฺตีติ.}

\proseitem{242}{“อปฺปาหาโร สมโณ โคตโม, อปฺปาหารตาย จ วณฺณวาที”ติ อิติ เจ มํ อุทายิ สาวกา สกฺกเรยฺยุํ ครุํ กเรยฺยุํ มาเนยฺยุํ ปูเชยฺยุํ, สกฺกตฺวา ครุํ กตฺวา อุปนิสฺสาย วิหเรยฺยุํ.\csromanspacer{} สนฺติ โข ปน เม อุทายิ สาวกา โกสกาหาราปิ อฑฺฒโกสกาหาราปิ เพลุวาหาราปิ อฑฺฒเพลุวาหาราปิ.\csromanspacer{} อหํ โข ปน อุทายิ อปฺเปกทา อิมินา ปตฺเตน สมติตฺติกํปิ ภุญฺชามิ, ภิยฺโยปิ ภุญฺชามิ. “อปฺปาหาโร สมโณ โคตโม อปฺปาหารตาย จ วณฺณวาที”ติ อิติ เจ มํ อุทายิ สาวกา สกฺกเรยฺยุํ ครุํ กเรยฺยุํ มาเนยฺยุํ ปูเชยฺยุํ, สกฺกตฺวา ครุํ กตฺวา อุปนิสฺสาย วิหเรยฺยุํ.\csromanspacer{} เย เต อุทายิ มม สาวกา โกสกาหาราปิ อฑฺฒโกสกาหาราปิ เพลุวาหาราปิ อฑฺฒเพลุวาหาราปิ.\csromanspacer{} น มํ เต อิมินา ธมฺเมน สกฺกเรยฺยุํ ครุํ กเรยฺยุํ มาเนยฺยุํ ปูเชยฺยุํ, สกฺกตฺวา ครุํ กตฺวา อุปนิสฺสาย วิหเรยฺยุํ. (1)}

\prose{“สนฺตุฏฺโฐ สมโณ โคตโม อิตรีตเรน จีวเรน, อิตรีตรจีวรสนฺตุฏฺฐิยา จ วณฺณวาที”ติ อิติ เจ มํ อุทายิ สาวกา สกฺกเรยฺยุํ ครุํ กเรยฺยุํ มาเนยฺยุํ ปูเชยฺยุํ, สกฺกตฺวา ครุํ กตฺวา อุปนิสฺสาย วิหเรยฺยุํ.\csromanspacer{} สนฺติ โข ปน เม อุทายิ สาวกา ปํสุกูลิกา ลูขจีวรธรา, เต สุสานา วา สงฺการกูฏา วา ปาปณิกา วา นนฺตกานิ\footnote{ปาปณิกานิ วา นนฺตกานิ วา (สี)} อุจฺจินิตฺวา\footnote{อุจฺฉินฺทิตฺวา (ก)} สงฺฆาฏิํ กริตฺวา ธาเรนฺติ.\csromanspacer{} อหํ โข ปนุทายิ อปฺเปกทา คหปติจีวรานิ \csromanfolio{200}ธาเรมิ ทฬฺหานิ สตฺถลูขานิ อลาพุโลมสานิ. “สนฺตุฏฺโฐ สมโณ โคตโม อิตรีตเรน จีวเรน, อิตรีตรจีวรสนฺตุฏฺฐิยา จ วณฺณวาที”ติ อิติ เจ มํ อุทายิ สาวกา สกฺกเรยฺยุํ ครุํ กเรยฺยุํ มาเนยฺยุํ ปูเชยฺยุํ, สกฺกตฺวา ครุํ กตฺวา อุปนิสฺสาย วิหเรยฺยุํ.\csromanspacer{} เย เต อุทายิ มม สาวกา ปํสุกูลิกา ลูขจีวรธรา, เต สุสานา วา สงฺการกูฏา วา ปาปณิกา วา นนฺตกานิ อุจฺจินิตฺวา สงฺฆาฏิํ กริตฺวา ธาเรนฺติ.\csromanspacer{} น มํ เต อิมินา ธมฺเมน สกฺกเรยฺยุํ ครุํ กเรยฺยุํ มาเนยฺยุํ ปูเชยฺยุํ, สกฺกตฺวา ครุํ กตฺวา อุปนิสฺสาย วิหเรยฺยุํ. (2)}

\prose{“สนฺตุฏฺโฐ สมโณ โคตโม อิตรีตเรน ปิณฺฑปาเตน, อิตรีตรปิณฺฑปาตสนฺตุฏฺฐิยา จ วณฺณวาที”ติ อิติ เจ มํ อุทายิ สาวกา สกฺกเรยฺยุํ ครุํ กเรยฺยุํ มาเนยฺยุํ ปูเชยฺยุํ, สกฺกตฺวา ครุํ กตฺวา อุปนิสฺสาย วิหเรยฺยุํ.\csromanspacer{} สนฺติ โข ปน เม อุทายิ สาวกา ปิณฺฑปาติกา สปทานจาริโน อุญฺฉาสเก วเต รตา, เต อนฺตรฆรํ ปวิฏฺฐา สมานา อาสเนนปิ นิมนฺติยมานา น สาทิยนฺติ.\csromanspacer{} อหํ โข ปนุทายิ อปฺเปกทา นิมนฺตเนปิ\footnote{นิมนฺตนสฺสาปิ (ก)} ภุญฺชามิ สาลีนํ โอทนํ วิจิตกาฬกํ อเนกสูปํ อเนกพฺยญฺชนํ. “สนฺตุฏฺโฐ สมโณ โคตโม อิตรีตเรน ปิณฺฑปาเตน, อิตรีตรปิณฺฑปาตสนฺตุฏฺฐิยา จ วณฺณวาที”ติ อิติ เจ มํ อุทายิ สาวกา สกฺกเรยฺยุํ ครุํ กเรยฺยุํ มาเนยฺยุํ ปูเชยฺยุํ, สกฺกตฺวา ครุํ กตฺวา อุปนิสฺสาย วิหเรยฺยุํ.\csromanspacer{} เย เต อุทายิ มม สาวกา ปิณฺฑปาติกา สปทานจาริโน อุญฺฉาสเก วเต รตา, เต อนฺตรฆรํ ปวิฏฺฐา สมานา อาสเนนปิ นิมนฺติยมานา น สาทิยนฺติ.\csromanspacer{} น มํ เต อิมินา ธมฺเมน สกฺกเรยฺยุํ ครุํ กเรยฺยุํ มาเนยฺยุํ ปูเชยฺยุํ, สกฺกตฺวา ครุํ กตฺวา อุปนิสฺสาย วิหเรยฺยุํ. (3)}

\prose{“สนฺตุฏฺโฐ สมโณ โคตโม อิตรีตเรน เสนาสเนน, อิตรีตรเสนาสนสนฺตุฏฺฐิยา จ วณฺณวาที”ติ อิติ เจ มํ อุทายิ สาวกา สกฺกเรยฺยุํ ครุํ กเรยฺยุํ มาเนยฺยุํ ปูเชยฺยุํ, สกฺกตฺวา ครุํ กตฺวา อุปนิสฺสาย วิหเรยฺยุํ.\csromanspacer{} สนฺติ โข ปน เม อุทายิ สาวกา รุกฺขมูลิกา อพฺโภกาสิกา, เต อฏฺฐมาเส ฉนฺนํ น อุเปนฺติ.\csromanspacer{} อหํ โข ปนุทายิ อปฺเปกทา กูฏาคาเรสุปิ วิหรามิ อุลฺลิตฺตาวลิตฺเตสุ นิวาเตสุ ผุสิตคฺคเฬสุ\footnote{ผุสฺสิตคฺคเฬสุ (สี, อิ)} ปิหิตวาตปาเนสุ. “สนฺตุฏฺโฐ สมโณ โคตโม}

\vspace{\tipitakaparskip}
\csromancenter{มหาสกุลุทายีสุตฺต (77) อิตรีตเรน เสนาสเนน, อิตรีตรเสนาสนสนฺตุฏฺฐิยา จ วณฺณวาที”ติ อิติ เจ มํ อุทายิ สาวกา สกฺกเรยฺยุํ ครุํ กเรยฺยุํ มาเนยฺยุํ ปูเชยฺยุํ, สกฺกตฺวา ครุํ กตฺวา อุปนิสฺสาย วิหเรยฺยุํ.\csromanspacer{} เย เต อุทายิ มม สาวกา รุกฺขมูลิกา อพฺโภกาสิกา, เต อฏฺฐมาเส ฉนฺนํ น อุเปนฺติ.\csromanspacer{} น มํ เต อิมินา ธมฺเมน สกฺกเรยฺยุํ ครุํ กเรยฺยุํ มาเนยฺยุํ ปูเชยฺยุํ, สกฺกตฺวา ครุํ กตฺวา อุปนิสฺสาย วิหเรยฺยุํ. (4)}
\prose{\csromanfolio{201}“ปวิวิตฺโต สมโณ โคตโม, ปวิเวกสฺส จ วณฺณวาที”ติ อิติ เจ มํ อุทายิ สาวกา สกฺกเรยฺยุํ ครุํ กเรยฺยุํ มาเนยฺยุํ ปูเชยฺยุํ, สกฺกตฺวา ครุํ กตฺวา อุปนิสฺสาย วิหเรยฺยุํ.\csromanspacer{} สนฺติ โข ปน เม อุทายิ สาวกา อารญฺญิกา ปนฺตเสนาสนา อรญฺญวนปตฺถานิ ปนฺตานิ เสนาสนานิ อชฺโฌคาเหตฺวา วิหรนฺติ, เต อนฺวทฺธมาสํ สํฆมชฺเฌ โอสรนฺติ ปาติโมกฺขุทฺเทสาย.\csromanspacer{} อหํ โข ปนุทายิ อปฺเปกทา อากิณฺโณ วิหรามิ ภิกฺขูหิ ภิกฺขุนีหิ อุปาสเกหิ อุปาสิกาหิ รญฺญา ราชมหามตฺเตหิ ติตฺถิเยหิ ติตฺถิยสาวเกหิ. “ปวิวิตฺโต สมโณ โคตโม, ปวิเวกสฺส จ วณฺณวาที”ติ อิติ เจ มํ อุทายิ สาวกา สกฺกเรยฺยุํ ครุํ กเรยฺยุํ มาเนยฺยุํ ปูเชยฺยุํ, สกฺกตฺวา ครุํ กตฺวา อุปนิสฺสาย วิหเรยฺยุํ.\csromanspacer{} เย เต อุทายิ มม สาวกา อารญฺญกา ปนฺตเสนาสนา, อรญฺญวนปตฺถานิ ปนฺตานิ เสนาสนานิ อชฺโฌคาเหตฺวา วิหรนฺติ, เต อนฺวทฺธมาสํ สํฆมชฺเฌ โอสรนฺติ ปาติโมกฺขุทฺเทสาย.\csromanspacer{} น มํ เต อิมินา ธมฺเมน สกฺกเรยฺยุํ ครุํ กเรยฺยุํ มาเนยฺยุํ ปูเชยฺยุํ, สกฺกตฺวา ครุํ กตฺวา อุปนิสฺสาย วิหเรยฺยุํ. (5)}

\prose{อิติ โข อุทายิ น มมํ สาวกา อิเมหิ ปญฺจหิ ธมฺเมหิ สกฺกโรนฺติ ครุํ กโรนฺติ มาเนนฺติ ปูเชนฺติ, สกฺกตฺวา ครุํ กตฺวา อุปนิสฺสาย วิหรนฺติ.}

\proseitem{243}{อตฺถิ โข อุทายิ อญฺเญ จ ปญฺจ ธมฺมา, เยหิ ปญฺจหิ ธมฺเมหิ มมํ สาวกา สกฺกโรนฺติ ครุํ กโรนฺติ มาเนนฺติ ปูเชนฺติ, สกฺกตฺวา ครุํ กตฺวา อุปนิสฺสาย วิหรนฺติ.\csromanspacer{} กตเม ปญฺจ, อิธุทายิ มมํ สาวกา อธิสีเล สมฺภาเวนฺติ “สีลวา สมโณ โคตโม, ปรเมน สีลกฺขนฺเธน สมนฺนาคโต”ติ.\csromanspacer{} ยมฺปุทายิ\footnote{ยมุทายิ (สฺยา, ก)} มมํ สาวกา อธิสีเล สมฺภาเวนฺติ “สีลวา สมโณ โคตโม, ปรเมน สีลกฺขนฺเธน สมนฺนาคโต”ติ.\csromanspacer{} อยํ โข \csromanfolio{202}อุทายิ ปฐโม ธมฺโม, เยน มมํ สาวกา สกฺกโรนฺติ ครุํ กโรนฺติ มาเนนฺติ ปูเชนฺติ, สกฺกตฺวา ครุํ กตฺวา อุปนิสฺสาย วิหรนฺติ.}

\proseitem{244}{ปุน จปรํ อุทายิ มมํ สาวกา อภิกฺกนฺเต ญาณทสฺสเน สมฺภาเวนฺติ “ชานํเยวาห สมโณ โคตโม ‘ชานามี’ติ, ปสฺสํเยวาห สมโณ โคตโม ‘ปสฺสามี’ติ, อภิญฺญาย สมโณ โคตโม ธมฺมํ เทเสติ, โน อนภิญฺญาย.\csromanspacer{} สนิทานํ สมโณ โคตโม ธมฺมํ เทเสติ, โน อนิทานํ.\csromanspacer{} สปฺปาฏิหาริยํ สมโณ โคตโม ธมฺมํ เทเสติ, โน อปฺปาฏิหาริยนฺ”ติ.\csromanspacer{} ยมฺปุทายิ มมํ สาวกา อภิกฺกนฺเต ญาณทสฺสเน สมฺภาเวนฺติ “ชานํเยวาห สมโณ โคตโม ‘ชานามี’ติ, ปสฺสํเยวาห สมโณ โคตโม ‘ปสฺสามี’ติ อภิญฺญาย สมโณ โคตโม ธมฺมํ เทเสติ, โน อนภิญฺญาย.\csromanspacer{} สนิทานํ สมโณ โคตโม ธมฺมํ เทเสติ, โน อนิทานํ.\csromanspacer{} สปฺปาฏิหาริยํ สมโณ โคตโม ธมฺมํ เทเสติ, โน อปฺปาฏิหาริยนฺ”ติ.\csromanspacer{} อยํ โข อุทายิ ทุติโย ธมฺโม, เยน มมํ สาวกา สกฺกโรนฺติ ครุํ กโรนฺติ มาเนนฺติ ปูเชนฺติ, สกฺกตฺวา ครุํ กตฺวา อุปนิสฺสาย วิหรนฺติ.}

\proseitem{245}{ปุน จปรํ อุทายิ มมํ สาวกา อธิปญฺญาย สมฺภาเวนฺติ “ปญฺญวา สมโณ โคตโม, ปรเมน ปญฺญากฺขนฺเธน สมนฺนาคโต.\csromanspacer{} ตํ วต อนาคตํ วาทปถํ น ทกฺขติ, อุปฺปนฺนํ วา ปรปฺปวาทํ น สหธมฺเมน สุนิคฺคหิตํ นิคฺคณฺหิสฺสตีติ เนตํ ฐานํ วิชฺชติ”.\csromanspacer{} ตํ กิํ มญฺญสิ อุทายิ, อปิ นุ เม สาวกา เอวํ ชานนฺตา เอวํ ปสฺสนฺตา อนฺตรนฺตรา กถํ โอปาเตยฺยุนฺติ.\csromanspacer{} โน เหตํ ภนฺเต.\csromanspacer{} น โข ปนาหํ อุทายิ สาวเกสุ อนุสาสนิํ ปจฺจาสีสามิ\footnote{ปจฺจาสิํสามิ (สี, สฺยา, กํ, อิ)}, อญฺญทตฺถุ มมเยว สาวกา อนุสาสนิํ ปจฺจาสีสนฺติ, ยมฺปุทายิ มมํ สาวกา อธิปญฺญาย สมฺภาเวนฺติ “ปญฺญวา สมโณ โคตโม, ปรเมน ปญฺญากฺขนฺเทน สมนฺนาคโต.\csromanspacer{} ตํ วต อนาคตํ วาทปถํ น ทกฺขติ, อุปฺปนฺนํ วา ปรปฺปวาทํ น สหธมฺเมน นิคฺคหิตํ นิคฺคณฺหิสฺสตีติ เนตํ ฐานํ วิชฺชติ”.\csromanspacer{} อยํ โข อุทายิ ตติโย ธมฺโม, เยน มมํ สาวกา สกฺกโรนฺติ ครุํ กโรนฺติ มาเนนฺติ ปูเชนฺติ, สกฺกตฺวา ครุํ กตฺวา อุปนิสฺสาย วิหรนฺติ.}

\csromanheader{2}{7. มหาสกุลุทายีสุตฺต (77)}
\proseitem{246}{\csromanfolio{203}ปุน จปรํ อุทายิ มม สาวกา เยน ทุกฺเขน ทุกฺโขติณฺณา ทุกฺขปเรตา.\csromanspacer{} เต มํ อุปสงฺกมิตฺวา ทุกฺขํ อริยสจฺจํ ปุจฺฉนฺติ.\csromanspacer{} เตสาหํ ทุกฺขํ อริยสจฺจํ ปุฏฺโฐ พฺยากโรมิ, เตสาหํ จิตฺตํ อาราเธมิ ปญฺหสฺส เวยฺยากรเณน.\csromanspacer{} เต มํ ทุกฺขสมุทยํ.\csromanspacer{} ทุกฺขนิโรธํ.\csromanspacer{} ทุกฺขนิโรธคามินิํ ปฏิปทํ อริยสจฺจํ ปุจฺฉนฺติ.\csromanspacer{} เตสาหํ ทุกฺขนิโรธคามินิํ ปฏิปทํ อริยสจฺจํ ปุฏฺโฐ พฺยากโรมิ, เตสาหํ จิตฺตํ อาราเธมิ ปญฺหสฺส เวยฺยากรเณน.\csromanspacer{} ยมฺปุทายิ มม สาวกา เยน ทุกฺเขน ทุกฺโขติณฺณา ทุกฺขปเรตา.\csromanspacer{} เต มํ อุปสงฺกมิตฺวา ทุกฺขํ อริยสจฺจํ ปุจฺฉนฺติ.\csromanspacer{} เตสาหํ ทุกฺขํ อริยสจฺจํ ปุฏฺโฐ พฺยากโรมิ, เตสาหํ จิตฺตํ อาราเธมิ ปญฺหสฺส เวยฺยากรเณน.\csromanspacer{} เต มํ ทุกฺขสมุทยํ.\csromanspacer{} ทุกฺขนิโรธํ.\csromanspacer{} ทุกฺขนิโรธคามินิํ ปฏิปทํ อริยสจฺจํ ปุจฺฉนฺติ.\csromanspacer{} เตสาหํ ทุกฺขนิโรธคามินิํ ปฏิปทํ อริยสจฺจํ ปุฏฺโฐ พฺยากโรมิ, เตสาหํ จิตฺตํ อาราเธมิ ปญฺหสฺส เวยฺยากรเณน.\csromanspacer{} อยํ โข อุทายิ จตุตฺโถ ธมฺโม, เยน มมํ สาวกา สกฺกโรนฺติ ครุํ กโรนฺติ มาเนนฺติ ปูเชนฺติ, สกฺกตฺวา ครุํ กตฺวา อุปนิสฺสาย วิหรนฺติ.}

\proseitem{247}{ปุน จปรํ อุทายิ อกฺขาตา มยา สาวกานํ ปฏิปทา, ยถาปฏิปนฺนา เม สาวกา จตฺตาโร สติปฏฺฐาเน ภาเวนฺติ.\csromanspacer{} อิธุทายิ ภิกฺขุ กาเย กายานุปสฺสี วิหรติ อาตาปี สมฺปชาโน สติมา วิเนยฺย โลเก อภิชฺฌาโทมนสฺสํ.\csromanspacer{} เวทนาสุ เวทนานุปสฺสี วิหรติ.\csromanspacer{} จิตฺเต จิตฺตานุปสฺสี วิหรติ.\csromanspacer{} ธมฺเมสุ ธมฺมานุปสฺสี วิหรติ อาตาปี สมฺปชาโน สติมา วิเนยฺย โลเก อภิชฺฌาโทมนสฺสํ.\csromanspacer{} ตตฺร จ ปน เม สาวกา พหู อภิญฺญาโวสานปารมิปฺปตฺตา วิหรนฺติ.}

\prose{ปุน จปรํ อุทายิ อกฺขาตา มยา สาวกานํ ปฏิปทา, ยถาปฏิปนฺนา เม สาวกา จตฺตาโร สมฺมปฺปธาเน ภาเวนฺติ.\csromanspacer{} อิธุทายิ ภิกฺขุ อนุปฺปนฺนานํ ปาปกานํ อกุสลานํ ธมฺมานํ อนุปฺปาทาย ฉนฺทํ ชเนติ วายมติ วีริยํ อารภติ จิตฺตํ ปคฺคณฺหาติ ปทหติ.\csromanspacer{} อุปฺปนฺนานํ ปาปกานํ อกุสลานํ ธมฺมานํ ปหานาย ฉนฺทํ ชเนติ วายมติ วีริยํ อารภติ จิตฺตํ ปคฺคณฺหาติ ปทหติ.\csromanspacer{} อนุปฺปนฺนานํ กุสลานํ ธมฺมานํ อุปฺปาทาย ฉนฺทํ ชเนติ วายมติ วีริยํ อารภติ จิตฺตํ ปคฺคณฺหาติ ปทหติ.\csromanspacer{} อุปฺปนฺนานํ กุสลานํ ธมฺมานํ ฐิติยา อสมฺโมสาย ภิยฺโยภาวาย เวปุลฺลาย ภาวนาย ปาริปูริยา ฉนฺทํ ชเนติ วายมติ วีริยํ อารภติ จิตฺตํ ปคฺคณฺหาติ ปทหติ.\csromanspacer{} ตตฺร จ ปน เม สาวกา พหู อภิญฺญาโวสานปารมิปฺปตฺตา วิหรนฺติ.}

\prose{\csromanfolio{204}ปุน จปรํ อุทายิ อกฺขาตา มยา สาวกานํ ปฏิปทา, ยถาปฏิปนฺนา เม สาวกา จตฺตาโร อิทฺธิปาเท ภาเวนฺติ.\csromanspacer{} อิธุทายิ ภิกฺขุ ฉนฺทสมาธิปธานสงฺขารสมนฺนาคตํ อิทฺธิปาทํ ภาเวติ, วีริยสมาธิปธานสงฺขารสมนฺนาคตํ อิทฺธิปาทํ ภาเวติ, จิตฺตสมาธิปธานสงฺขารสมนฺนาคตํ อิทฺธิปาทํ ภาเวติ, วีมํสสมาธิปธานสงฺขารสมนฺนาคตํ อิทฺธิปาทํ ภาเวติ.\csromanspacer{} ตตฺร จ ปน เม สาวกา พหู อภิญฺญาโวสานปารมิตฺตา วิหรนฺติ.}

\prose{ปุน จปรํ อุทายิ อกฺขาตา มยา สาวกานํ ปฏิปทา, ยถาปฏิปนฺนา เม สาวกา ปญฺจินฺทฺริยานิ ภาเวนฺติ.\csromanspacer{} อิธุทายิ ภิกฺขุ สทฺธินฺทฺริยํ ภาเวติ อุปสมคามิํ สมฺโพธคามิํ.\csromanspacer{} วีริยินฺทฺริยํ ภาเวติ ฯเปฯ.\csromanspacer{} สตินฺทฺริยํ ภาเวติ.\csromanspacer{} สมาธินฺทฺริยํ ภาเวติ.\csromanspacer{} ปญฺญินฺทฺริยํ ภาเวติ อุปสมคามิํ สมฺโพธคามิํ.\csromanspacer{} ตตฺร จ ปน เม สาวกา พหู อภิญฺญาโวสานปารมิปฺปตฺตา วิหรนฺติ.}

\prose{ปุน จปรํ อุทายิ อกฺขาตา มยา สาวกานํ ปฏิปทา, ยถาปฏิปนฺนา เม สาวกา ปญฺจ พลานิ ภาเวนฺติ.\csromanspacer{} อิธุทายิ ภิกฺขุ สทฺธาพลํ ภาเวติ อุปสมคามิํ สมฺโพธคามิํ.\csromanspacer{} วีริยพลํ ภาเวติ ฯเปฯ.\csromanspacer{} สติพลํ ภาเวติ.\csromanspacer{} สมาธิพลํ ภาเวติ.\csromanspacer{} ปญฺญาพลํ ภาเวติ อุปสมคามิํ สมฺโพธคามิํ.\csromanspacer{} ตตฺร จ ปน เม สาวกา พหู อภิญฺญาโวสานปารมิปฺปตฺตา วิหรนฺติ.}

\prose{ปุน จปรํ อุทายิ อกฺขาตา มยา สาวกานํ ปฏิปทา, ยถาปฏิปนฺนา เม สาวกา สตฺตโพชฺฌงฺเค ภาเวนฺติ.\csromanspacer{} อิธุทายิ ภิกฺขุ สติสมฺโพชฺฌงฺคํ ภาเวติ วิเวกนิสฺสิตํ วิราคนิสฺสิตํ นิโรธนิสฺสิตํ โวสฺสคฺคปริณามิํ.\csromanspacer{} ธมฺมวิจยสมฺโพชฺฌงฺคํ ภาเวติ ฯเปฯ.\csromanspacer{} วีริยสมฺโพชฺฌงฺคํ ภาเวติ.\csromanspacer{} ปีติสมฺโพชฺฌงฺคํ ภาเวติ.\csromanspacer{} ปสฺสทฺธิสมฺโพชฺฌงฺคํ ภาเวติ.\csromanspacer{} สมาธิสมฺโพชฺฌงฺคํ ภาเวติ.\csromanspacer{} อุเปกฺขาสมฺโพชฺฌงฺคํ ภาเวติ วิเวกนิสฺสิตํ วิราคนิสฺสิตํ นิโรธนิสฺสิตํ โวสฺสคฺคปริณามิํ.\csromanspacer{} ตตฺร จ ปน เม สาวกา พหู อภิญฺญาโวสานปารมิปฺปตฺตา วิหรนฺติ.}

\prose{ปุน จปรํ อุทายิ อกฺขาตา มยา สาวกานํ ปฏิปทา, ยถาปฏิปนฺนา เม สาวกา อริยํ อฏฺฐงฺคิกํ มคฺคํ ภาเวนฺติ.\csromanspacer{} อิธุทายิ ภิกฺขุ สมฺมาทิฏฺฐิํ ภาเวติ, สมฺมาสงฺกปฺปํ ภาเวติ, สมฺมาวาจํ ภาเวติ, สมฺมากมฺมนฺตํ ภาเวติ, สมฺมา-อาชีวํ ภาเวติ, สมฺมาวายามํ ภาเวติ,}

\vspace{\tipitakaparskip}
\csromancenter{มหาสกุลุทายีสุตฺต (77) สมฺมาสติํ ภาเวติ, สมฺมาสมาธิํ ภาเวติ.\csromanspacer{} ตตฺร จ ปน เม สาวกา พหู อภิญฺญาโวสานปารมิปฺปตฺตา วิหรนฺติ.}
\proseitem{248}{\csromanfolio{205}ปุน จปรํ อุทายิ อกฺขาตา มยา สาวกานํ ปฏิปทา, ยถาปฏิปนฺนา เม สาวกา อฏฺฐ วิโมกฺเข ภาเวนฺติ.\csromanspacer{} รูปี รูปานิ ปสฺสติ, อยํ ปฐโม วิโมกฺโข.\csromanspacer{} อชฺฌตฺตํ อรูปสญฺญี พหิทฺธา รูปานิ ปสฺสติ, อยํ ทุติโย วิโมกฺโข.\csromanspacer{} สุภนฺเตว อธิมุตฺโต โหติ, อยํ ตติโย วิโมกฺโข.\csromanspacer{} สพฺพโส รูปสญฺญานํ สมติกฺกมา ปฏิฆสญฺญานํ อตฺถงฺคมา นานตฺตสญฺญานํ อมนสิการา “อนนฺโต อากาโส”ติ อากาสานญฺจายตนํ อุปสมฺปชฺช วิหรติ, อยํ จตุตฺโถ วิโมกฺโข.\csromanspacer{} สพฺพโส อากาสานญฺจายตนํ สมติกฺกมฺม “อนนฺตํ วิญฺญาณนฺ”ติ วิญฺญาณญฺจายตนํ อุปสมฺปชฺช วิหรติ, อยํ ปญฺจโม วิโมกฺโข.\csromanspacer{} สพฺพโส วิญฺญาณญฺจายตนํ สมติกฺกมฺม “นตฺถิ กิญฺจี”ติ อากิญฺจญฺญายตนํ อุปสมฺปชฺช วิหรติ, อยํ ฉฏฺโฐ วิโมกฺโข.\csromanspacer{} สพฺพโส อากิญฺจญฺญายตนํ สมติกฺกมฺม เนวสญฺญานาสญฺญายตนํ อุปสมฺปชฺช วิหรติ, อยํ สตฺตโม วิโมกฺโข.\csromanspacer{} สพฺพโส เนวสญฺญานาสญฺญายตนํ สมติกฺกมฺม สญฺญาเวทยิตนิโรธํ อุปสมฺปชฺช วิหรติ, อยํ อฏฺฐโม วิโมกฺโข.\csromanspacer{} ตตฺร จ ปน เม สาวกา พหู อภิญฺญาโวสานปารมิปฺปตฺตา วิหรนฺติ.}

\proseitem{249}{ปุน จปรํ อุทายิ อกฺขาตา มยา สาวกานํ ปฏิปทา, ยถาปฏิปนฺนา เม สาวกา อฏฺฐ อภิภายตนานิ ภาเวนฺติ.\csromanspacer{} อชฺฌตฺตํ รูปสญฺญี เอโก พหิทฺธา รูปานิ ปสฺสติ ปริตฺตานิ สุวณฺณทุพฺพณฺณานิ, ตานิ อภิภุยฺย ชานามิ ปสฺสามีติ เอวํ สญฺญี โหติ, อิทํ ปฐมํ อภิภายตนํ.}

\prose{อชฺฌตฺตํ รูปสญฺญี เอโก พหิทฺธา รูปานิ ปสฺสติ อปฺปมาณานิ สุวณฺณทุพฺพณฺณานิ, ตานิ อภิภุยฺย ชานามิ ปสฺสามีติ เอวํ สญฺญี โหติ, อิทํ ทุติยํ อภิภายตนํ.}

\prose{อชฺฌตฺตํ อรูปสญฺญี เอโก พหิทฺธา รูปานิ ปสฺสติ ปริตฺตานิ สุวณฺณทุพฺพณฺณานิ, ตานิ อภิภุยฺย ชานามิ ปสฺสามีติ เอวํ สญฺญี โหติ, อิทํ ตติยํ อภิภายตนํ.}

\prose{อชฺฌตฺตํ อรูปสญฺญี เอโก พหิทฺธา รูปานิ ปสฺสติ อปฺปมาณานิ สุวณฺณทุพฺพณฺณานิ, ตานิ อภิภุยฺย ชานามิ ปสฺสามีติ เอวํ สญฺญี โหติ, อิทํ จตุตฺถํ อภิภายตนํ.}

\prose{\csromanfolio{206}อชฺฌตฺตํ อรูปสญฺญี เอโก พหิทฺธา รูปานิ ปสฺสติ นีลานิ นีลวณฺณานิ นีลนิทสฺสนานิ นีลนิภาสานิ.\csromanspacer{} เสยฺยถาปิ นาม อุมาปุปฺผํ นีลํ นีลวณฺณํ นีลนิทสฺสนํ นีลนิภาสํ.\csromanspacer{} เสยฺยถา วา ปน ตํ วตฺถํ พาราณเสยฺยกํ อุภโตภาควิมฏฺฐํ นีลํ นีลวณฺณํ นีลนิทสฺสนํ นีลนิภาสํ.\csromanspacer{} เอวเมว อชฺฌตฺตํ อรูปสญฺญี เอโก พหิทฺธา รูปานิ ปสฺสติ นีลานิ นีลวณฺณานิ นีลนิทสฺสนานิ นีลนิภาสานิ, ตานิ อภิภุยฺย ชานามิ ปสฺสามีติ เอวํ สญฺญี โหติ, อิทํ ปญฺจมํ อภิภายตนํ.}

\prose{อชฺฌตฺตํ อรูปสญฺญี เอโก พหิทฺธา รูปานิ ปสฺสติ ปีตานิ ปีตวณฺณานิ ปีตนิทสฺสนานิ ปีตนิภาสานิ.\csromanspacer{} เสยฺยถาปิ นาม กณิการปุปฺผํ ปีตํ ปีตวณฺณํ ปีตนิทสฺสนํ ปีตนิภาสํ.\csromanspacer{} เสยฺยถา วา ปน ตํ วตฺถํ พาราณเสยฺยกํ อุภโตภาควิมฏฺฐํ ปีตํ ปีตวณฺณํ ปีตนิทสฺสนํ ปีตนิภาสํ.\csromanspacer{} เอวเมว อชฺฌตฺตํ อรูปสญฺญี เอโก พหิทฺธา รูปานิ ปสฺสติ ปีตานิ ปีตวณฺณานิ ปีตนิทสฺสนานิ ปีตนิภาสานิ, ตานิ อภิภุยฺย ชานามิ ปสฺสามีติ เอวํ สญฺญี โหติ, อิทํ ฉฏฺฐํ อภิภายตนํ.}

\prose{อชฺฌตฺตํ อรูปสญฺญี เอโก พหิทฺธา รูปานิ ปสฺสติ โลหิตกานิ โลหิตกวณฺณานิ โลหิตกนิทสฺสนานิ โลหิตกนิภาสานิ.\csromanspacer{} เสยฺยถาปิ นาม พนฺธุชีวกปุปฺผํ โลหิตกํ โลหิตกวณฺณํ โลหิตกนิทสฺสนํ โลหิตกนิภาสํ.\csromanspacer{} เสยฺยถา วา ปน ตํ วตฺตํ พาราณเสยฺยกํ อุภโตภาควิมฏฺฐํ โลหิตกํ โลหิตกวณฺณํ โลหิตกนิทสฺสนํ โลหิตกนิภาสํ.\csromanspacer{} เอวเมว อชฺฌตฺตํ อรูปสญฺญี เอโก พหิทฺธา รูปานิ ปสฺสติ โลหิตกานิ โลหิตกวณฺณานิ โลหิตกนิทสฺสนานิ โลหิตกนิภาสานิ, ตานิ อภิภุยฺย ชานามิ ปสฺสามีติ เอวํ สญฺญี โหติ, อิทํ สตฺตมํ อภิภายตนํ.}

\prose{อชฺฌตฺตํ อรูปสญฺญี เอโก พหิทฺธา รูปานิ ปสฺสติ โอทาตานิ โอทาตวณฺณานิ โอทาตนิทสฺสนานิ โอทาตนิภาสานิ.\csromanspacer{} เสยฺยถาปิ นาม โอสธิตารกา โอทาตา โอทาตวณฺณา โอทาตนิทสฺสนา โอทาตนิภาสา.\csromanspacer{} เสยฺยถา วา ปน ตํ วตฺถํ พาราณเสยฺยกํ อุภโตภาควิมฏฺฐํ โอทาตํ โอทาตวณฺณํ โอทาตนิทสฺสนํ โอทาตนิภาสํ.\csromanspacer{} เอวเมว อชฺฌตฺตํ อรูปสญฺญี เอโก พหิทฺธา รูปานิ ปสฺสติ โอทาตานิ โอทาตวณฺณานิ โอทาตนิทสฺสนานิ โอทาตนิภาสานิ, ตานิ อภิภุยฺย}

\vspace{\tipitakaparskip}
\csromancenter{มหาสกุลุทายีสุตฺต (77) ชานามิ ปสฺสามีติ เอวํสญฺญี โหติ, อิทํ อฏฺฐมํ อภิภายตนํ.\csromanspacer{} ตตฺร จ ปน เม สาวกา พหู อภิญฺญาโวสานปารมิปฺปตฺตา วิหรนฺติ.}
\proseitem{250}{\csromanfolio{207}ปุน จปรํ อุทายิ อกฺขาตา มยา สาวกานํ ปฏิปทา, ยถาปฏิปนฺนา เม สาวกา ทส กสิณายตนานิ ภาเวนฺติ, ปถวีกสิณเมโก สญฺชานาติ อุทฺธมโธ ติริยํ อทฺวยํ อปฺปมาณํ.\csromanspacer{} อาโปกสิณ เมโก สญฺชานาติ ฯเปฯ.\csromanspacer{} เตโชกสิณ เมโก สญฺชานาติ.\csromanspacer{} วาโยกสิณเมโก สญฺชานาติ.\csromanspacer{} นีลกสิณเมโก สญฺชานาติ.\csromanspacer{} ปีตกสิณเมโก สญฺชานาติ.\csromanspacer{} โลหิตกสิณเมโก สญฺชนาติ.\csromanspacer{} โอทาตกสิณเมโก สญฺชานาติ.\csromanspacer{} อากาสกสิณเมโก สญฺชานาติ.\csromanspacer{} วิญฺญาณกสิณเมโก สญฺชานาติ อุทฺธมโธ ติริยํ อทฺวยํ อปฺปมาณํ.\csromanspacer{} ตตฺร จ ปน เม สาวกา พหู อภิญฺญาโวสานปารมิปฺปตฺตา วิหรนฺติ.}

\proseitem{251}{ปุน จปรํ อุทายิ อกฺขาตา มยา สาวกานํ ปฏิปทา, ยถาปฏิปนฺนา เม สาวกา จตฺตาริ ฌานานิ ภาเวนฺติ.\csromanspacer{} อิธุทายิ ภิกฺขุ วิวิจฺเจว กาเมหิ วิวิจฺจ อกุสเลหิ ธมฺเมหิ สวิตกฺกํ สวิจารํ วิเวกชํ ปีติสุขํ ปฐมํ ฌานํ อุปสมฺปชฺช วิหรติ, โส อิมเมว กายํ วิเวกเชน ปีติสุเขน อภิสนฺเทติ ปริสนฺเทติ ปริปูเรติ ปริปฺผรติ, นาสฺส กิญฺจิ สพฺพาวโต กายสฺส วิเวกเชน ปีติสุเขน อปฺผุฏํ โหติ.\csromanspacer{} เสยฺยถาปิ อุทายิ ทกฺโข นฺหาปโก\footnote{นหาปโก (สี, อิ)} วา นฺหาปกนฺเตวาสี วา กํสถาเล นฺหานียจุณฺณานิ\footnote{นหานียจุณฺณานิ (สี, อิ)} อากิริตฺวา อุทเกน ปรฺùปฺโผสกํ ปริปฺโผสกํ สนฺเนยฺย, สายํ นฺหานียปิณฺฑิ\footnote{สาสฺส นหานียปิณฺฑี (สี, สฺยา, กํ)} สฺเนหานุคตา สฺเนหปเรตา สนฺตรพาหิรา ผุฏา สฺเนเหน, น จ ปคฺฆริณี.\csromanspacer{} เอวเมว โข อุทายิ ภิกฺขุ อิมเมว กายํ วิเวกเชน ปีติสุเขน อภิสนฺเทติ ปริสนฺเทติ ปริปูเรติ ปริปฺผรติ, นาสฺส กิญฺจิ สพฺพาวโต กายสฺส วิเวกเชน ปีติสุเขน อปฺผุฏํ โหติ.}

\prose{ปุน จปรํ อุทายิ ภิกฺขุ วิตกฺกวิจารานํ วูปสมา อชฺฌตฺตํ สมฺปสาทนํ ฯเปฯ ทุติยํ ฌานํ อุปสมฺปชฺช วิหรติ, โส อิมเมว กายํ สมาธิเชน ปีติสุเขน อภิสนฺเทติ ปริสนฺเทติ ปริปูเรติ ปริปฺผรติ, นาสฺส กิญฺจิ สพฺพาวโต กายสฺส สมาธิเชน ปีติสุเขน อปฺผุฏํ \csromanfolio{208}โหติ.\csromanspacer{} เสยฺยถาปิ อุทายิ อุทกรหโท คมฺภีโร อุพฺภิโททโก\footnote{อุพฺภิโตทโก (สฺยา, กํ, ก)}, ตสฺส เนวสฺส ปุรตฺถิมาย ทิสาย อุทกสฺส อายมุขํ, น ปจฺฉิมาย ทิสาย อุทกสฺส อายมุขํ, น อุตฺตราย ทิสาย อุทกสฺส อายมุขํ, น ทกฺขิณาย ทิสาย อุทกสฺส อายมุขํ, เทโว จ น กาเลน กาลํ สมฺมา ธารํ อนุปฺปเวจฺเฉยฺย.\csromanspacer{} อถ โข ตมฺหาว อุทกรหทา สีตา วาริธารา อุพฺภิชฺชิตฺวา ตเมว อุทกรหทํ สีเตน วารินา อภิสนฺเทยฺย ปริสนฺเทยฺย ปริปูเรยฺย ปริปฺผเรยฺย, นาสฺส กิญฺจิ สพฺพาวโต อุทกรหทสฺส สีเตน วารินา อปฺผุฏํ อสฺส.\csromanspacer{} เอวเมว โข อุทายิ ภิกฺขุ อิมเมว กายํ สมาธิเชน ปีติสุเขน อภิสนฺเทติ ปริสนฺเทติ ปริปูเรติ ปริปฺผรติ, นาสฺส กิญฺจิ สพฺพาวโต กายสฺส สมาธิเชน ปีติสุเขน อปฺผุฏํ โหติ.}

\prose{ปุน จปรํ อุทายิ ภิกฺขุ ปีติยา จ วิราคา ฯเปฯ ตติยํ ฌานํ อุปสมฺปชฺช วิหรติ.\csromanspacer{} โส อิมเมว กายํ นิปฺปีติเกน สุเขน อภิสนฺเทติ ปริสนฺเทติ ปริปูเรติ ปริปฺผรติ, นาสฺส กิญฺจิ สพฺพาวโต กายสฺส นิปฺปีติเกน สุเขน อปฺผุฏํ โหติ.\csromanspacer{} เสยฺยถาปิ อุทายิ อุปฺปลินิยํ วา ปทุมินิยํ วา ปุณฺฑรีกินิยํ วา อปฺเปกจฺจานิ อุปฺปลานิ วา ปทุมานิ วา ปุณฺฑรีกานิ วา อุทเก ชาตานิ อุทเก สํวฑฺฒานิ อุทกานุคฺคตานิ อนฺโต นิมุคฺคโปสีนิ, ตานิ ยาว จคฺคา ยาว จ มูลา สีเตน วารินา อภิสนฺนานิ ปริสนฺนานิ ปริปูรานิ ปริปฺผุฏานิ, นาสฺส\footnote{น เนสํ (สี)} กิญฺจิ สพฺพาวตํ อุปฺปลานํ วา ปทุมานํ วา ปุณฺฑรีกานํ วา สีเตน วารินา อปฺผุฏํ อสฺส.\csromanspacer{} เอวเมว โข อุทายิ ภิกฺขุ อิมเมว กายํ นิปฺปีติเกน สุเขน อภิสนฺเทติ ปริสนฺเทติ ปริปูเรติ ปริปฺผรติ, นาสฺส กิญฺจิ สพฺพาวโต กายสฺส นิปฺปีติเกน สุเขน อปฺผุฏํ โหติ.}

\prose{ปุน จปรํ อุทายิ ภิกฺขุ สุขสฺส จ ปหานา ทุกฺขสฺส จ ปหานา ปุพฺเพว โสมนสฺสโทมนสฺสานํ อตฺถงฺคมา อทุกฺขมสุขํ อุเปกฺขาสติปาริสุทฺธิํ จตุตฺถํ ฌานํ อุปสมฺปชฺช วิหรติ, โส อิมเมว กายํ ปริสุทฺเธน เจตสา ปริโยทาเตน ผริตฺวา นิสินฺโน โหติ, นาสฺส กิญฺจิ สพฺพาวโต กายสฺส ปริสุทฺเธน เจตสา ปริโยทาเตน อปฺผุฏํ โหติ.\csromanspacer{} เสยฺยถาปิ อุทายิ ปุริโส โอทาเตน วตฺเถน สสีตํ ปารุปิตฺวา นิสินฺโน อสฺส, นาสฺส กิญฺจิ สพฺพาวโต กายสฺส โอทาเตน}

\vspace{\tipitakaparskip}
\csromancenter{มหาสกุลุทายีสุตฺต (77) วตฺเถน อปฺผุฏํ อสฺส.\csromanspacer{} เอวเมว โข อุทายิ ภิกฺขุ อิมเมว กายํ ปริสุทฺเธน เจตสา ปริโยทาเตน ผริตฺวา นิสินฺโน โหติ, นาสฺส กิญฺจิ สพฺพาวโต กายสฺส ปริสุทฺเธน เจตสา ปริโยทาเตน อปฺผุฏํ โหติ.\csromanspacer{} ตตฺร จ ปน เม สาวกา พหู อภิญฺญาโวสานปารมิปฺปตฺตา วิหรนฺติ.}
\proseitem{252}{\csromanfolio{209}ปุน จปรํ อุทายิ อกฺขาตา มยา สาวกานํ ปฏิปทา, ยถา ปฏิปนฺนา เม สาวกา เอวํ ปชานนฺติ “อยํ โข เม กาโย รูปี จาตุมหาภูติโก มาตาเปตฺติกสมฺภโว โอทนกุมฺมาสูปจโย อนิจฺจุจฺฉาทนปริมทฺทนเภทนวิทฺธํสนธมฺโม, อิทญฺจ ปน เม วิญฺญาณํ เอตฺถ สิตํ เอตฺถ ปฏิพทฺธํ.\csromanspacer{} เสยฺยถาปิ อุทายิ มณิ เวฬุริโย สุโภ ชาติมา อฏฺฐํโส สุปริกมฺมกโต อจฺโฉ วิปฺปสนฺโน สพฺพาการสมฺปนฺโน, ตตฺริทํ สุตฺตํ อาวุตํ นีลํ วา ปีตํ วา โลหิตํ วา โอทาตํ วา ปณฺฑุสุตฺตํ วา, ตเมนํ จกฺขุมา ปุริโส หตฺเถ กริตฺวา ปจฺจเวกฺเขยฺย ‘อยํ โข มณิ เวฬุริโย สุโภ ชาติมา อฏฺฐํโส สุปริกมฺมกโต อจฺโฉ วิปฺปสนฺโน สพฺพาการสมฺปนฺโน, ตตฺริทํ สุตฺตํ อาวุตํ นีลํ วา ปีตํ วา โลหิตํ วา โอทาตํ วา ปณฺฑุสุตฺตํ วา’ติ”.\csromanspacer{} เอวเมว โข อุทายิ อกฺขาตา มยา สาวกานํ ปฏิปทา, ยถาปฏิปนฺนา เม สาวกา เอวํ ปชานนฺติ “อยํ โข เม กาโย รูปี จาตุมหาภูติโก มาตาเปตฺติกสมฺภโว โอทนกุมฺมาสูปจโย อนิจฺจุจฺฉาทนปริมทฺทนเภทนวิทฺธํสนธมฺโม, อิทญฺจ ปน เม วิญฺญาณํ เอตฺถ สิตํ เอตฺถ ปฏิพทฺธนฺ”ติ.\csromanspacer{} ตตฺร จ ปน เมสาวกา พหู อภิญฺญาโวสานปารมิปฺปตฺตา วิหรนฺติ.}

\proseitem{253}{ปุน จปรํ อุทายิ อกฺขาตา มยา สาวกานํ ปฏิปทา, ยถาปฏิปนฺนา เม สาวกา อิมมฺหา กายา อญฺญํ กายํ อภินิมฺมินนฺติ รูปิํ มโนมยํ สพฺพงฺคปจฺจงฺคิํ อหีนินฺทฺริยํ.\csromanspacer{} เสยฺยถาปิ อุทายิ ปุริโส มุญฺชมฺหา อีสิกํ ปพฺพาเหยฺย, ตสฺส เอวมสฺส “อยํ มุญฺโช, อยํ อีสิกา, อญฺโญ มุญฺโช, อญฺญา อีสิกา, มุญฺชมฺหาเตฺวว อีสิกา ปพฺพาฬฺหา”ติ.\csromanspacer{} เสยฺยถา วา ปนุทายิ ปุริโส อสิํ โกสิยา ปพฺพาเหยฺย, ตสฺส เอวมสฺส “อยํ อสิ, อยํ โกสิ, อญฺโญ อสิ, อญฺญา โกสิ, โกสิยาเตฺวว อสิ ปพฺพาโฬฺห”ติ.\csromanspacer{} เสยฺยถา วา ปนุทายิ ปุริโส อหิํ กรณฺฑา อุทฺธเรยฺย, ตสฺส เอวมสฺส “อยํ อหิ, อยํ กรณฺโฑ, \csromanfolio{210}อญฺโญ อหิ, อญฺโญ กรณฺโฑ, กรณฺฑาเตฺวว อหิ อุพฺภโต”ติ.\csromanspacer{} เอวเมว โข อุทายิ อกฺขาตา มยา สาวกานํ ปฏิปทา, ยถาปฏิปนฺนา เม สาวกา อิมมฺหา กายา อญฺญํ กายํ อภินิมฺมินนฺติ รูปิํ มโนมยํ สพฺพงฺคปจฺจงฺคิํ อหีนินฺทฺริยํ.\csromanspacer{} ตตฺร จ ปน เม สาวกา พหู อภิญฺญาโวสนปารมิปฺปตฺตา วิหรนฺติ.}

\proseitem{254}{ปุน จปรํ อุทายิ อกฺขาตา มยา สาวกานํ ปฏิปทา, ยถาปฏิปนฺนา เม สาวกา อเนกวิหิตํ อิทฺธิวิธํ ปจฺจนุโภนฺติ, เอโกปิ หุตฺวา พหุธา โหนฺติ, พหุธาปิ หุตฺวา เอโก โหติ, อาวิภาวํ ติโรภาวํ ติโรกุฏฺฏํ ติโรปาการํ ติโรปพฺพตํ อสชฺชมานา คจฺฉนฺติ เสยฺยถาปิ อากาเส, ปถวิยาปิ อุมฺมุชฺชนิมุชฺชํ กโรนฺติ เสยฺยถาปิ อุทเก, อุทเกปิ อภิชฺชมาเน\footnote{อภิชฺชมานา (ก)} คจฺฉนฺติ เสยฺยถาปิ ปถวิยํ, อากาเสปิ ปลฺลงฺเกน กมนฺติ เสยฺยถาปิ ปกฺขี สกุโณ, อิเมปิ จนฺทิมสูริเย เอวํมหิทฺธิเก เอวํมหานุภาเว ปาณินา ปริมสนฺติ ปริมชฺชนฺติ, ยาว พฺรหฺมโลกาปิ กาเยน วสํ วตฺเตนฺติ.\csromanspacer{} เสยฺยถาปิ อุทายิ ทกฺโข กุมฺภกาโร วา กุมฺภการนฺเตวาสี วา สุปริกมฺมกตาย มตฺติกาย ยํ ยเทว ภาชนวิกติํ อากงฺเขยฺย, ตํ ตเทว กเรยฺย อภินิปฺผาเทยฺย.\csromanspacer{} เสยฺยถา วา ปนุทายิ ทกฺโข ทนฺตกาโร วา ทนฺตการนฺเตวาสี วา สุปริกมฺมกตสฺมิํ ทนฺตสฺมิํ ยํ ยเทว ทนฺตวิกติํ อากงฺเขยฺย, ตํ ตเทว กเรยฺย อภินิปฺผาเทยฺย.\csromanspacer{} สยฺยถา วา ปนุทายิ ทกฺโข สุวณฺณกาโร วา สุวณฺณการนฺเตวาสี วา สุปริกมฺมกตสฺมิํ สุวณฺณสฺมิํ ยํ ยเทว สุวณฺณวิกติํ อากงฺเขยฺย, ตํ ตเทว กเรยฺย อภินิปฺผาเทยฺย.\csromanspacer{} เอวเมว โข อุทายิ อกฺขาตา มยา สาวกานํ ปฏิปทา, ยถาปฏิปนฺนา เม สาวกา อเนกวิหิตํ อิทฺธิวิธํ ปจฺจนุโภนฺติ, เอโกปิ หุตฺวา พหุธา โหนฺติ, พหุธาปิ หุตฺวา เอโก โหติ, อาวิภาวํ ติโรภาวํ ติโรกุฏฺฏํ ติโรปาการํ ติโรปพฺพตํ อสชฺชมานา คจฺฉนฺติ เสยฺยถาปิ อากาเส, ปถวิยาปิ อุมฺมุชฺชนิมุชฺชํ กโรนฺติ เสยฺยถาปิ อุทเก, อุทเกปิ อภิชฺชมาเน คจฺฉนฺติ เสยฺยถาปิ ปถวิยํ, อากาเสปิ ปลฺลงฺเกน กมนฺติ เสยฺยถาปิ ปกฺขี สกุโณ, อิเมปิ จนฺทิมสูริเย เอวํมหิทฺธิเก เอวํมหานุภาเว ปาณินา ปริมชฺชนฺติ, ยาว พฺรหฺมโลกปิ}

\vspace{\tipitakaparskip}
\csromancenter{มหาสกุลุทายีสุตฺต (77) กาเยน วสํ วตฺเตนฺติ.\csromanspacer{} ตตฺร จ ปน เม สาวกา พหู อภิญฺญาโวสานปารมิปฺปตฺตา วิหรนฺติ.}
\proseitem{255}{\csromanfolio{211}ปุน จปรํ อุทายิ อกฺขาตา มยา สาวกานํ ปฏิปทา, ยถาปฏิปนฺนา เม สาวกา ทิพฺพาย โสตธาตุยา วิสุทฺธาย อติกฺกนฺตมานุสิกาย อุโภ สทฺเท สุณนฺติ ทิพฺเพ จ มานุเส จ เย ทูเร สนฺติเก จ.\csromanspacer{} เสยฺยถาปิ อุทายิ พลวา สงฺขธโม อปฺปกสิเรเนว จาตุทฺทิสา วิญฺญาเปยฺย.\csromanspacer{} เอวเมว โข อุทายิ อกฺขาตา มยา สาวกานํ ปฏิปทา, ยถาปฏิปนฺนา เม สาวกา ทิพฺพาย โสตธาตุยา วิสุทฺธาย อติกฺกนฺตมานุสิกาย อุโภ สทฺเท สุณนฺติ ทิพฺเพ จ มานุเส จ เย ทูเร สนฺติเก จ.\csromanspacer{} ตตฺร จ ปน เม สาวกา พหู อภิญฺญาโวสานปารมิปฺปตฺตา วิหรนฺติ.}

\proseitem{256}{ปุน จปรํ อุทายิ อกฺขาตา มยา สาวกานํ ปฏิปทา, ยถาปฏิปนฺนา เม สาวกา ปรสตฺตานํ ปรปุคฺคลานํ เจตสา เจโต ปริจฺจ ปชานนฺติ, สราคํ วา จิตฺตํ สราคํ จิตฺตนฺติ ปชานนฺติ, วีตราคํ วา จิตฺตํ วีตราคํ จิตฺตนฺติ ปชานนฺติ, สโทสํ วา จิตฺตํ สโทสํ จิตฺตนฺติ ปชานนฺติ, วีตโทสํ วา จิตฺตํ วีตโทสํ จิตฺตนฺติ ปชานนฺติ, สโมหํ วา จิตฺตํ สโมหํ จิตฺตนฺติ ปชานนฺติ, วีตโมหํ วา จิตฺตํ วีตโมหํ จิตฺตนฺติ ปชานนฺติ, สํขิตฺตํ วา จิตฺตํ สํขิตฺตํ จิตฺตนฺติ ปชานนฺติ, วิกฺขิตฺตํ วา จิตฺตํ วิกฺขิตฺตํ จิตฺตนฺติ ปชานนฺติ, มหคฺคตํ วา จิตฺตํ มหคฺคตํ จิตฺตนฺติ ปชานนฺติ, อมหคฺคตํ วา จิตฺตํ อมหคฺคตํ จิตฺตนฺติ ปชานนฺติ, ส-อุตฺตรํ วา จิตฺตํ ส-อุตฺตรํ จิตฺตนฺติ ปชานนฺติ, อนุตฺตรํ วา จิตฺตํ อนุตฺตรํ จิตฺตนฺติ ปชานนฺติ, สมาหิตํ วา จิตฺตํ สมาหิตํ จิตฺตนฺติ ปชานนฺติ, อสมาหิตํ วา จิตฺตํ อสมาหิตํ จิตฺตนฺติ ปชานนฺติ, วิมุตฺตํ วา จิตฺตํ วิมุตฺตํ จิตฺตนฺติ ปชานนฺติ, อวิมุตฺตํ วา จิตฺตํ อวิมุตฺตํ จิตฺตนฺติ ปชานนฺติ.\csromanspacer{} เสยฺยถาปิ อุทายิ อิตฺถี วา ปุริโส วา ทหโร ยุวา มณฺฑนกชาติโก อาทาเส วา ปริสุทฺเธ ปริโยทาเต อจฺเฉ วา อุทกปตฺเต สกํ มุขนิมิตฺตํ ปจฺจเวกฺขมาโน สกณิกํ วา สกณิกนฺติ\footnote{สกณิกงฺคํ วา สกณิกงฺคนฺติ (สี)} ชาเนยฺย, อกณิกํ วา อกณิกนฺติ\footnote{อกณิกงฺคํ วา อกณิกงฺคนฺติ (สี)} ชาเนยฺย.\csromanspacer{} เอวเมว โข อุทายิ อกฺขาตา มยา สาวกานํ ปฏิปทา, ยถาปฏิปนฺนา เม สาวกา ปรสตฺตานํ ปรปุคฺคลานํ เจตสา เจโต ปริจฺจ ปชานนฺติ, \csromanfolio{212}สราคํ วา จิตฺตํ สราคํ จิตฺตนฺติ ปชานนฺติ, วีตราคํ วา จิตฺตํ ฯเปฯ สโทสํ วา จิตฺตํ.\csromanspacer{} วีตโทสํ วา จิตฺตํ.\csromanspacer{} สโมหํ วา จิตฺตํ.\csromanspacer{} วีตโมหํ วาจิตฺตํ.\csromanspacer{} สํขิตฺตํ วา จิตฺตํ.\csromanspacer{} วิกฺขิตฺตํ วา จิตฺตํ.\csromanspacer{} มหคฺคตํ วา จิตฺตํ.\csromanspacer{} อมหคฺคตํ วา จิตฺตํ.\csromanspacer{} ส-อุตฺตรํ วา จิตฺตํ.\csromanspacer{} อนุตฺตรํ วา จิตฺตํ.\csromanspacer{} สมาหิตํ วา จิตฺตํ.\csromanspacer{} อสมาหิตํ วา จิตฺตํ.\csromanspacer{} วิมุตฺตํ วา จิตฺตํ.\csromanspacer{} อวิมุตฺตํ วา จิตฺตํ อวิมุตฺตํ จิตฺตนฺติ ปชานนฺติ.\csromanspacer{} ตตฺร จ ปน เม สาวกา พหู อภิญฺญาโวสานปารมิปฺปตฺตา วิหรนฺติ.}

\proseitem{257}{ปุน จปรํ อุทายิ อกฺขาตา มยา สาวกานํ ปฏิปทา, ยถาปฏิปนฺนา เม สาวกา อเนกวิหิตํ ปุพฺเพนิวาสํ อนุสฺสรนฺติ.\csromanspacer{} เสยฺยถิทํ, เอกมฺปิ ชาติํ เทฺวปิ ชาติโย ติสฺโสปิ ชาติโย จตสฺโสปิ ชาติโย ปญฺจปิ ชาติโย ทสปิ ชาติโย วีสมฺปิ ชาติโย ติํสมฺปิ ชาติโย จตฺตาลีสมฺปิ ชาติโย ปญฺญาสมฺปิ ชาติโย ชาติสตมฺปิ ชาติสหสฺสมฺปิ ชาติสตสหสฺสมฺปิ อเนเกปิ สํวฏฺฏกปฺเป อเนเกปิ วิวฏฺฏกปฺเป อเนเกปิ สํวฏฺฏวิวฏฺฏกปฺเป, “อมุตฺราสิํ เอวํนาโม เอวํโคตฺโต เอวํวณฺโณ เอวมาหาโร เอวํสุขทุกฺขปฺปฏิสํเวที เอวมายุปริยนฺโต, โส ตโต จุโต อมุตฺร อุทปาทิํ, ตตฺราปาสิํ เอวํนาโม เอวํโคตฺโต เอวํวณฺโณ เอวมาหาโร เอวํสุขทุกฺขปฺปฏิสํเวที เอวมายุปริยนฺโต, โส ตโต จุโต อิธูปปนฺโน”ติ.\csromanspacer{} อิติ สาการํ ส-อุทฺเทสํ อเนกวิหิตํ ปุพฺเพนิวาสํ อนุสฺสรติ.\csromanspacer{} เสยฺยถาปิ อุทายิ ปุริโส สกมฺหา คามา อญฺญํ คามํ คจฺเฉยฺย, ตมฺหาปิ คามา อญฺญํ คามํ คจฺเฉยฺย, โส ตมฺหา คามา สกํเยว คามํ ปจฺจาคจฺเฉยฺย.\csromanspacer{} ตสฺส เอวมสฺส “อหํ โข สกมฺหา คามา อญฺญํ คามํ อคจฺฉิํ, ตตฺร เอวํ อฏฺฐาสิํ เอวํ นิสีทิํ เอวํ อภาสิํ เอวํ ตุณฺหี อโหสิํ, ตมฺหาปิ คามา อมุํ คามํ อคจฺฉิํ, ตตฺราปิ เอวํ อฏฺฐาสิํ เอวํ นิสีทิํ เอวํ อภาสิํ เอวํ ตุณฺหี อโหสิํ, โสมฺหิ ตมฺหา คามา สกํเยว คามํ ปจฺจาคโต”ติ.\csromanspacer{} เอวเมว โข อุทายิ อกฺขาตา มยา สาวกานํ ปฏิปทา, ยถาปฏิปนฺนา เม สาวกา อเนกวิหิตํ ปุพฺเพนิวาสํ อนุสฺสรนฺติ.\csromanspacer{} เสยฺยถิทํ, เอกมฺปิ ชาติํ ฯเปฯ อิติ สาการํ ส-อุทฺเทสํ อเนกวิหิตํ ปุพฺเพนิวาสํ อนุสฺสรนฺติ.\csromanspacer{} ตตฺร จ ปน เม สาวกา พหู อภิญฺญาโวสานปารมิปฺปตฺตา วิหรนฺติ.}

\proseitem{258}{ปุน จปรํ อุทายิ อกฺขาตา มยา สาวกานํ ปฏิปทา, ยถาปฏิปนฺนา เม สาวกา ทิพฺเพน จกฺขุนา วิสุทฺเธน อติกฺกนฺตมานุสเกน}

\vspace{\tipitakaparskip}
\csromancenter{มหาสกุลุทายีสุตฺต (77) สตฺเต ปสฺสนฺติ จวมาเน อุปปชฺชมาเน หีเน ปณีเต สุวณฺเณ ทุพฺพณฺเณ สุคเต ทุคฺคเต, ยถากมฺมูปเค สตฺเต ปชานนฺติ “อิเม วต โภนฺโต สตฺตา กายทุจฺจริเตน สมนฺนาคตา วจีทุจฺจริเตน สมนฺนาคตา มโนทุจฺจริเตน สมนฺนาคตา อริยานํ อุปวาทกา มิจฺฉาทิฏฺฐิกา มิจฺฉาทิฏฺฐิกมฺมสมาทานา, เต กายสฺส เภทา ปรํ มรณา อปายํ ทุคฺคติํ วินิปาตํ นิรยํ อุปปนฺนา.\csromanspacer{} อิเม วา ปน โภนฺโต สตฺตา กายสุจริเตน สมนฺนาคตา วจีสุจริเตน สมนฺนาคตา มโนสุจริเตน สมนฺนาคตา อริยานํ อนุปวาทกา สมฺมาทิฏฺฐิกา สมฺมาทิฏฺฐิกมฺมสมาทานา, เต กายสฺส เภทา ปรํ มรณา สุคติํ สคฺคํ โลกํ อุปปนฺนา”ติ.\csromanspacer{} อิติ ทิพฺเพน จกฺขุนา วิสุทฺเธน อติกฺกนฺตมานุสเกน สตฺเต ปสฺสนฺติ จวมาเน อุปปชฺชมาเน หีเน ปณีเต สุวณฺเณ ทุพฺพณฺเณ สุคเต ทุคฺคเต, ยถากมฺมูปเค สตฺเต ปชานนฺติ.\csromanspacer{} เสยฺยถาปิ อุทายิ เทฺว อคารา สทฺวารา\footnote{สนฺนทฺวารา (ก)}, ตตฺร จกฺขุมา ปุริโส มชฺเฌ ฐิโต ปสฺเสยฺย มนุสฺเส เคหํ ปวิสนฺเตปิ นิกฺขมนฺเตปิ อนุจงฺกมนฺเตปิ อนุวิจรนฺเตปิ.\csromanspacer{} เอวเมว โข อุทายิ อกฺขาตา มยา สาวกานํ ปฏิปทา, ยถาปฏิปนฺนา เม สาวกา ทิพฺเพน จกฺขุนา วิสุทฺเธน อติกฺกนฺตมานุสเกน สตฺเต ปสฺสนฺติ จวมาเน อุปปชฺชมาเน หีเน ปณีเต สุวณฺเณ ทุพฺพณฺเณ สุคเต ทุคฺคเต, ยถากมฺมูปเค สตฺเต ปชานนฺติ ฯเปฯ.\csromanspacer{} ตตฺร จ ปน เม สาวกา พหู อภิญฺญาโวสานปารมิปฺปตฺตา วิหรนฺติ.}
\proseitem{259}{\csromanfolio{213}ปุน จปรํ อุทายิ อกฺขาตา มยา สาวกานํ ปฏิปทา, ยถาปฏิปนฺนา เม สาวกา อาสวานํ ขยา อนาสวํ เจโตวิมุตฺติํ ปญฺญาวิมุตฺติํ ทิฏฺเฐว ธมฺเม สยํ อภิญฺญา สจฺฉิกตฺวา อุปสมฺปชฺช วิหรนฺติ.\csromanspacer{} เสยฺยถาปิ อุทายิ ปพฺพตสงฺเขเป อุทกรหโท อจฺโฉ วิปฺปสนฺโน อนาวิโล, ตตฺถ จกฺขุมา ปุริโส ตีเร ฐิโต ปสฺเสยฺย สิปฺปิสมฺพุกมฺปิ\footnote{สิปฺปิกสมฺพุกมฺปิ (สฺยา, กํ, ก)} สกฺขรกฐลมฺปิ มจฺฉคุมฺพมฺปิ จรนฺตมฺปิ ติฏฺฐนฺตมฺปิ.\csromanspacer{} ตสฺส เอวมสฺส “อยํ โข อุทกรหโท อจฺโฉ วิปฺปสนฺโน อนาวิโล, ตตฺริเม สิปฺปิสมฺพุกาปิ สกฺขรกฐลาปิ มจฺฉคุมฺพาปิ จรนฺติปิ ติฏฺฐนฺติปี”ติ.\csromanspacer{} เอวเมว โข อุทายิ อกฺขาตา มยา สาวกานํ ปฏิปทา, ยถาปฏิปนฺนา เม สาวกา อาสวานํ ขยา อนาสวํ เจโตวิมุตฺติํ ปญฺญามุตฺติํ ทิฏฺเฐว ธมฺเม สยํ อภิญฺญา สจฺฉิกตฺวา อุปสมฺปชฺช วิหรนฺติ.\csromanspacer{} ตตฺร จ ปน \csromanfolio{214}เม สาวกา พหู อภิญฺญาโวสานปารมิปฺปตฺตา วิหรนฺติ.\csromanspacer{} อยํ โข อุทายิ ปญฺจโม ธมฺโม.\csromanspacer{} เยน มม สาวกา สกฺกโรนฺติ ครุํ กโรนฺติ มาเนนฺติ ปูเชนฺติ, สกฺกตฺวา ครุํ กตฺวา อุปนิสฺสาย วิหรนฺติ.}

\prose{อิเม โข อุทายิ ปญฺจ ธมฺมา, เยหิ มมํ สาวกา สกฺกโรนฺติ ครุํ กโรนฺติ มาเนนฺติ ปูเชนฺติ, สกฺกตฺวา ครุํ กตฺวา อุปนิสฺสาย วิหรนฺตีติ.}

\prose{อิทมโวจ ภควา.\csromanspacer{} อตฺตมโน สกุลุทายี ปริพฺพาชโก ภควโต ภาสิตํ อภินนฺทีติ.}

\nitthitamruled{มหาสกุลุทายิสุตฺตํ นิฏฺฐิตํ สตฺตมํ.}
\csromanmatikaanchor{mtk.33}
\csromantocmarkref{subsection}{8. สมณมุณฺฑิกสุตฺต}{mtk.33}
\bookmark[dest={mtk.33},level=2]{8. สมณมุณฺฑิกสุตฺต}
\csromanheader{2}{8. สมณมุณฺฑิกสุตฺต}
\proseitem{260}{เอวํ เม สุตํ\csromandash{}เอกํ สมยํ ภควา สาวตฺถิยํ วิหรติ เชตวเน อนาถปิณฺฑิกสฺส อาราเม.\csromanspacer{} เตน โข ปน สมเยน อุคฺคาหมาโน ปริพฺพาชโก สมณมุณฺฑิกาปุตฺโต\footnote{สมณมณฺฑิกาปุตฺโต (สี, อิ)} สมยปฺปวาทเก ตินฺทุกาจีเร เอกสาลเก มลฺลิกาย อาราเม ปฏิวสติ มหติยา ปริพฺพาชกปริสาย สทฺธิํ ปญฺจมตฺเตหิ ปริพฺพาชกสเตหิ.\csromanspacer{} อถ โข ปญฺจกงฺโค ถปติ สาวตฺถิยา นิกฺขมิ ทิวา ทิวสฺส ภควนฺตํ ทสฺสนาย.\csromanspacer{} อถ โข ปญฺจกงฺคสฺส ถปติสฺส เอตทโหสิ “อกาโล โข ตาว ภควนฺตํ ทสฺสนาย, ปฏิสลฺลีโน ภควา, มโนภาวนิยานมฺปิ ภิกฺขูนํ อสมโย ทสฺสนาย, ปฏิสลฺลีนา มโนภาวนิยา ภิกฺขู, ยํนูนาหํ เยน สมยปฺปวาทโก ตินฺทุกาจีโร เอกสาลโก มลฺลิกาย อาราโม, เยน อุคฺคาหมาโน ปริพฺพาชโก สมณมุณฺฑิกาปุตฺโต เตนุปสงฺกเมยฺยนฺ”ติ.\csromanspacer{} อถ โข ปญฺจกงฺโค ถปติ เยน สมยปฺปวาทโก ตินฺทุกาจีโร เอกสาลโก มลฺลิกาย อาราโม, เยน อุคฺคาหมาโน ปริพฺพาชโก สมณมุณฺฑิกาปุตฺโต เตนุปสงฺกมิ.}

\prose{เตน โข ปน สมเยน อุคฺคาหมาโน ปริพฺพาชโก สมณมุณฺฑิกาปุตฺโต มหติยา ปริพฺพาชกปริสาย สทฺธิํ นิสินฺโน โหติ อุนฺนาทินิยา อุจฺจาสทฺทมหาสทฺทาย อเนกวิหิตํ ติรจฺฉานกถํ กเถนฺติยา.\csromanspacer{} เสยฺยถิทํ, ราชกถํ โจรกถํ มหามตฺตกถํ เสนากถํ ภยกถํ}

\vspace{\tipitakaparskip}
\csromancenter{สมณมุณฺฑิกสุตฺต (78) ยุทฺธกถํ อนฺนกถํ ปานกถํ วตฺถกถํ สยนกถํ มาลากถํ คนฺธกถํ ญาติกถํ ยานกถํ คามกถํ นิคมกถํ นครกถํ ชนปทกถํ อิตฺถิกถํ สูรกถํ วิสิขากถํ กุมฺภฏฺฐานกถํ ปุพฺพเปตกถํ นานตฺตกถํ โลกกฺขายิกํ สมุทฺทกฺขายิกํ อิติภวาภวกถํ อิติ วา.}
\prose{\csromanfolio{215}อทฺทสา โข อุคฺคาหมาโน ปริพฺพาชโก สมณมุณฺฑิกาปุตฺโต ปญฺจกงฺคํ ถปติํ ทูรโตว อาคจฺฉนฺตํ, ทิสฺวาน สกํ ปริสํ สณฺฐาเปสิ “อปฺปสทฺทา โภนฺโต โหนฺตุ, มา โภนฺโต สทฺทมกตฺต.\csromanspacer{} อยํ สมณสฺส โคตมสฺส สาวโก อาคจฺฉติ ปญฺจกงฺโค ถปติ.\csromanspacer{} ยาวตา โข ปน สมณสฺส โคตมสฺส สาวกา คิหี โอทาตวสนา สาวตฺถิยํ ปฏิวสนฺติ, อยํ เตสํ อญฺญตโร ปญฺจกงฺโค ถปติ.\csromanspacer{} อปฺปสทฺทกามา โข ปน เต อายสฺมนฺโต อปฺปสทฺทวินีตา อปฺปสทฺทสฺส วณฺณวาทิโน, อปฺเปว นาม อปฺปสทฺทํ ปริสํ วิทิตฺวา อุปสงฺกมิตพฺพํ มญฺเญยฺยา”ติ.\csromanspacer{} อถ โข เต ปริพฺพาชกา ตุณฺหี อเหสุํ.}

\proseitem{261}{อถ โข ปญฺจกงฺโค ถปติ เยน อุคฺคาหมาโน ปริพฺพาชโก สมณมุณฺฑิกาปุตฺโต เตนุปสงฺกมิ, อุปสงฺกมิตฺวา อุคฺคาหมาเนน ปริพฺพาชเกน สมณมุณฺฑิกาปุตฺเตน สทฺธิํ สมฺโมทิ, สมฺโมทนียํ กถํ สารณียํ วีติสาเรตฺวา เอกมนฺตํ นิสีทิ, เอกมนฺตํ นิสินฺนํ โข ปญฺจกงฺคํ ถปติํ อุคฺคาหมาโน ปริพฺพาชโก สมณมุณฺฑิกาปุตฺโต เอตทโวจ “จตูหิ โข อหํ คหปติ ธมฺเมหิ สมนฺนาคตํ ปุริสปุคฺคลํ ปญฺญเปมิ สมฺปนฺนกุสลํ ปรมกุสลํ อุตฺตมปตฺติปตฺตํ สมณํ อโยชฺฌํ.\csromanspacer{} กตเมหิ จตูหิ.\csromanspacer{} อิธ คหปติ น กาเยน ปาปกมฺมํ กโรติ, น ปาปกํ วาจํ ภาสติ, น ปาปกํ สงฺกปฺปํ สงฺกปฺเปติ, น ปาปกํ อาชีวํ อาชีวติ.\csromanspacer{} อิเมหิ โข อหํ คหปติ จตูหิ ธมฺเมหิ สมนฺนาคตํ ปุริสปุคฺคลํ ปญฺญเปมิ สมฺปนฺนกุสลํ ปรมกุสลํ อุตฺตมปตฺติปตฺตํ สมณํ อโยชฺฌนฺ”ติ.}

\prose{อถ โข ปญฺจกงฺโค ถปติ อุคฺคาหมานสฺส ปริพฺพาชกสฺส สมณมุณฺฑิกาปุตฺตสฺส ภาสิตํ เนว อภินนฺทิ นปฺปฏิกฺโกสิ, อนภินนฺทิตฺวา อปฺปฏิกฺโกสิตฺวา อุฏฺฐายาสนา ปกฺกามิ “ภควโต สนฺติเก เอตสฺส ภาสิตสฺส อตฺถํ อาชานิสฺสามี”ติ.\csromanspacer{} อถ โข ปญฺจกงฺโค ถปติ เยน ภควา เตนุปสงฺกมิ, อุปสงฺกมิตฺวา ภควนฺตํ อภิวาเทตฺวา เอกมนฺตํ นิสีทิ, เอกมนฺตํ นิสินฺโน โข ปญฺจกงฺโค ถปติ ยาวตโก อโหสิ \csromanfolio{216}อุคฺคาหมาเนน ปริพฺพาชเกน สมณมุณฺฑิกาปุตฺเตน สทฺธิํ กถาสลฺลาโป, ตํ สพฺพํ ภควโต อาโรเจสิ.}

\proseitem{262}{เอวํ วุตฺเต ภควา ปญฺจกงฺคํ ถปติํ เอตทโวจ\csromandash{}เอวํ สนฺเต โข ถปติ ทหโร กุมาโร มนฺโท อุตฺตานเสยฺยโก สมฺปนฺนกุสโล ภวิสฺสติ ปรมกุสโล อุตฺตมปตฺติปตฺโต สมโณ อโยชฺโฌ, ยถา อุคฺคาหมานสฺส ปริพฺพาชกสฺส สมณมุณฺฑิกาปุตฺตสฺส วจนํ.\csromanspacer{} ทหรสฺส หิ ถปติ กุมารสฺส มนฺทสฺส อุตฺตานเสยฺยกสฺส กาโยติปิ น โหติ, กุโต ปน กาเยน ปาปกมฺมํ กริสฺสติ อญฺญตฺร ผนฺทิตมตฺตา.\csromanspacer{} ทหรสฺส หิ ถปติ กุมารสฺส มนฺทสฺส อุตฺตานเสยฺยกสฺส วาจาติปิ น โหติ, กุโต ปน ปาปกํ วาจํ ภาสิสฺสติ อญฺญตฺร โรทิตมตฺตา.\csromanspacer{} ทหรสฺส หิ ถปติ กุมารสฺส มนฺทสฺส อุตฺตานเสยฺยกสฺส สงฺกปฺโปติปิ น โหติ, กุโต ปน ปาปกํ สงฺกปฺปํ สงฺกปฺปิสฺสติ อญฺญตฺร วิกูชิตมตฺตา\footnote{วิกุชฺชิตมตฺตา (สี, สฺยา, กํ, อิ)}.\csromanspacer{} ทหรสฺส หิ ถปติ กุมารสฺส มนฺทสฺส อุตฺตานเสยฺยกสฺส อาชีโวติปิ น โหติ, กุโต ปน ปาปกํ อาชีวํ อาชีวิสฺสติ อญฺญตฺร มาตุถญฺญา.\csromanspacer{} เอวํ สนฺเต โข ถปติ ทหโร กุมาโร มนฺโท อุตฺตานเสยฺยโก สมฺปนฺนกุสโล ภวิสฺสติ ปรมกุสโล อุตฺตมปตฺติปตฺโต สมโณ อโยชฺโฌ, ยถา อุคฺคาหมานสฺส ปริพฺพาชกสฺส สมณมุณฺฑิกาปุตฺตสฺส วจนํ.}

\proseitem{263}{จตูหิ โข อหํ ถปติ ธมฺเมหิ สมนฺนาคตํ ปุริสปุคฺคลํ ปญฺญเปมิ น เจว สมฺปนฺนกุสลํ น ปรมกุสลํ น อุตฺตมปตฺติปตฺตํ สมณํ อโยชฺฌํ, อปิ จิมํ ทหรํ กุมารํ มนฺทํ อุตฺตานเสยฺยกํ สมธิคยฺห ติฏฺฐติ.\csromanspacer{} กตเมหิ จตูหิ.\csromanspacer{} อิธ ถปติ น กาเยน ปาปกมฺมํ กโรติ, น ปาปกํ วาจํ ภาสติ, น ปาปกํ สงฺกปฺปํ สงฺกปฺเปติ, น ปาปกํ อาชีวํ อาชีวติ.\csromanspacer{} อิเมหิ โข อหํ ถปติ จตูหิ ธมฺเมหิ สมนฺนาคตํ ปุริสปุคฺคลํ ปญฺญเปมิ น เจว สมฺปนฺนกุสลํ น ปรมกุสลํ น อุตฺตมปตฺติปตฺตํ สมณํ อโยชฺฌํ, อปิ จิมํ ทหรํ กุมารํ มนฺทํ อุตฺตานเสยฺยกํ สมธิคยฺห ติฏฺฐติ.}

\prose{ทสหิ โข อหํ ถปติ ธมฺเมหิ สมนฺนาคตํ ปุริสปุคฺคลํ ปญฺญเปมิ สมฺปนฺนกุสลํ ปรมกุสลํ อุตฺตมปตฺติปตฺตํ สมณํ อโยชฺฌํ.\csromanspacer{} อิเม อกุสลา สีลา, ตมหํ\footnote{กหํ (สี), ตหํ (อิ)} ถปติ เวทิตพฺพนฺติ วทามิ.\csromanspacer{} อิโตสมุฏฺฐานา}

\vspace{\tipitakaparskip}
\csromancenter{สมณมุณฺฑิกสุตฺต (78) อกุสลา สีลา, ตมหํ ถปติ เวทิตพฺพนฺติ วทามิ.\csromanspacer{} อิธ อกุสลา สีลา อปริเสสา นิรุชฺฌนฺติ, ตมหํ ถปติ เวทิตพฺพนฺติ วทามิ.\csromanspacer{} เอวํ ปฏิปนฺโน อกุสลานํ สีลานํ นิโรธาย ปฏิปนฺโน โหติ, ตมหํ ถปติ เวทิตพฺพนฺติ วทามิ.}
\prose{\csromanfolio{217}อิเม กุสลา สีลา, ตมหํ ถปติ เวทิตพฺพนฺติ วทามิ.\csromanspacer{} อิโตสมุฏฺฐานา กุสลา สีลา, ตมหํ ถปติ เวทิตพฺพนฺติ วทามิ.\csromanspacer{} อิธ กุสลา สีลา อปริเสสา นิรุชฺฌนฺติ, ตมหํ ถปติ เวทิตพฺพนฺติ วทามิ.\csromanspacer{} เอวํ ปฏิปนฺโน กุสลานํ สีลานํ นิโรธาย ปฏิปนฺโน โหติ, ตมหํ ถปติ เวทิตพฺพนฺติ วทามิ.}

\prose{อิเม อกุสลา สงฺกปฺปา, ตมหํ ถปติ เวทิตพฺพนฺติ วทามิ.\csromanspacer{} อิโตสมุฏฺฐานา อกุสลา สงฺกปฺปา, ตมหํ ถปติ เวทิตพฺพนฺติ วทามิ.\csromanspacer{} อิธ อกุสลา สงฺกปฺปา อปริเสสา นิรุชฺฌนฺติ, ตมหํ ถปติ เวทิตพฺพนฺติ วทามิ.\csromanspacer{} เอวํ ปฏิปนฺโน อกุสลานํ สงฺกปฺปานํ นิโรธาย ปฏิปนฺโน โหติ, ตมหํ ถปติ เวทิตพฺพนฺติ วทามิ.}

\prose{อิเม กุสลา สงฺกปฺปา, ตมหํ ถปติ เวทิตพฺพนฺติ วทามิ.\csromanspacer{} อิโตสมุฏฺฐานา กุสลา สงฺกปฺปา, ตมหํ ถปติ เวทิตพฺพนฺติ วทามิ.\csromanspacer{} อิธ กุสลา สงฺกปฺปา อปริเสสา นิรุชฺฌนฺติ, ตมหํ ถปติ เวทิตพฺพนฺติ วทามิ.\csromanspacer{} เอวํ ปฏิปนฺโน กุสลานํ สงฺกปฺปานํ นิโรธาย ปฏิปนฺโน โหติ, ตมหํ ถปติ เวทิตพฺพนฺติ วทามิ.}

\proseitem{264}{กตเม จ ถปติ อกุสลา สีลา.\csromanspacer{} อกุสลํ กายกมฺมํ อกุสลํ วจีกมฺมํ ปาปโก อาชีโว.\csromanspacer{} อิเม วุจฺจนฺติ ถปติ อกุสลา สีลา.}

\prose{อิเม จ ถปติ อกุสลา สีลา กิํสมุฏฺฐานา.\csromanspacer{} สมุฏฺฐานมฺปิ เนสํ วุตฺตํ, จิตฺตสมุฏฺฐานาติสฺส วจนียํ.\csromanspacer{} กตมํ จิตฺตํ.\csromanspacer{} จิตฺตมฺปิ หิ พหุํ อเนกวิธํ นานปฺปการกํ, ยํ จิตฺตํ สราคํ สโทสํ สโมหํ.\csromanspacer{} อิโตสมุฏฺฐานา อกุสลา สีลา.}

\prose{อิเม จ ถปติ อกุสลา สีลา กุหิํ อปริเสสา นิรุชฺฌนฺติ.\csromanspacer{} นิโรโธปิ เนสํ วุตฺโต, อิธ ถปติ ภิกฺขุ กายทุจฺจริตํ ปหาย \csromanfolio{218}กายสุจริตํ ภาเวติ, วจีทุจฺจริตํ ปหาย วจีสุจริตํ ภาเวติ, มโนทุจฺจริตํ ปหาย มโนสุจริตํ ภาเวติ, มิจฺฉาชีวํ ปหาย สมฺมาชีเวน ชีวิตํ กปฺเปติ.\csromanspacer{} เอตฺเถเต อกุสลา สีลา อปริเสสา นิรุชฺฌนฺติ.}

\prose{กถํ ปฏิปนฺโน ถปติ อกุสลานํ สีลานํ นิโรธาย ปฏิปนฺโน โหติ.\csromanspacer{} อิธ ถปติ ภิกฺขุ อนุปฺปนฺนานํ ปาปกานํ อกุสลานํ ธมฺมานํ อนุปฺปาทาย ฉนฺทํ ชเนติ วายมติ วีริยํ อารภติ จิตฺตํ ปคฺคณฺหาติ ปทหติ, อุปฺปนฺนานํ ปาปกานํ อกุสลานํ ธมฺมานํ ปหานาย ฉนฺทํ ชเนติ วายมติ วีริยํ อารภติ จิตฺตํ ปคฺคณฺหาติ ปทหติ, อนุปฺปนฺนานํ กุสลานํ ธมฺมานํ อุปฺปาทาย ฉนฺทํ ชเนติ วายมติ วีริยํ อารภติ จิตฺตํ ปคฺคณฺหาติ ปทหติ, อุปฺปนฺนานํ กุสลานํ ธมฺมานํ ฐิติยา อสมฺโมสาย ภิยฺโยภาวาย เวปุลฺลาย ภาวนาย ปาริปูริยา ฉนฺทํ ชเนติ วายมติ วีริยํ อารภติ จิตฺตํ ปคฺคณฺหาติ ปทหติ.\csromanspacer{} เอวํ ปฏิปนฺโน โข ถปติ อกุสลานํ สีลานํ นิโรธาย ปฏิปนฺโน โหติ.}

\proseitem{265}{กตเม จ ถปติ กุสลา สีลา.\csromanspacer{} กุสลํ กายกมฺมํ กุสลํ วจีกมฺมํ อาชีวปริสุทฺธํปิ โข อหํ ถปติ สีลสฺมิํ วทามิ.\csromanspacer{} อิเม วุจฺจนฺติ ถปติ กุสลา สีลา.}

\prose{อิเม จ ถปติ กุสลา สีลา กิํสมุฏฺฐานา.\csromanspacer{} สมุฏฺฐานมฺปิ เนสํ วุตฺตํ, จิตฺตสมุฏฺฐานาติสฺส วจนียํ.\csromanspacer{} กตมํ จิตฺตํ.\csromanspacer{} จิตฺตมฺปิ หิ พหุํ อเนกวิธํ นานปฺปการกํ, ยํ จิตฺตํ วีตราคํ วีตโทสํ วีตโมหํ.\csromanspacer{} อิโตสมุฏฺฐานา กุสลา สีลา.}

\prose{อิเม จ ถปติ กุสลา สีลา กุหิํ อปริเสสา นิรุชฺฌนฺติ.\csromanspacer{} นิโรโธปิ เนสํ วุตฺโต, อิธ ถปติ ภิกฺขุ สีลวา โหติ โน จ สีลมโย, ตญฺจ เจโตวิมุตฺติํ ปญฺญาวิมุตฺติํ ยถาภูตํ ปชานาติ.\csromanspacer{} ยตฺถสฺส เต กุสลา สีลา อปริเสสา นิรุชฺฌนฺติ.}

\prose{กถํ ปฏิปนฺโน จ ถปติ กุสลานํ สีลานํ นิโรธาย ปฏิปนฺโน โหติ.\csromanspacer{} อิธ ถปติ ภิกฺขุ อนุปฺปนฺนานํ ปาปกานํ อกุสลานํ ธมฺมานํ อนุปฺปาทาย ฉนฺทํ ชเนติ วายมติ วีริยํ อารภติ จิตฺตํ ปคฺคณฺหาติ}

\vspace{\tipitakaparskip}
\csromancenter{สมณมุณฺฑิกสุตฺต (78) ปทหติ.\csromanspacer{} อุปฺปนฺนานํ ปาปกานํ อกุสลานํ ธมฺมานํ ปหานาย ฯเปฯ.\csromanspacer{} อนุปฺปนฺนานํ กุสลานํ ธมฺมานํ อุปฺปาทาย ฯเปฯ.\csromanspacer{} อุปฺปนฺนานํ กุสลานํ ธมฺมานํ ฐิติยา อสมฺโมสาย ภิยฺโยภาวาย เวปุลฺลาย ภาวนาย ปาริปูริยา ฉนฺทํ ชเนติ วายมติ วีริยํ อารภติ จิตฺตํ ปคฺคณฺหาติ ปทหติ.\csromanspacer{} เอวํ ปฏิปนฺโน โข ถปติ กุสลานํ สีลานํ นิโรธาย ปฏิปนฺโน โหติ.}
\proseitem{266}{\csromanfolio{219}กตเม จ ถปติ อกุสลา สงฺกปฺปา.\csromanspacer{} กามสงฺกปฺโป พฺยาปาทสงฺกปฺโป วิหิํสาสงฺกปฺโป.\csromanspacer{} อิเม วุจฺจนฺติ ถปติ อกุสลา สงฺกปฺปา.}

\prose{อิเม จ ถปติ อกุสลา สงฺกปฺปา กิํสมุฏฺฐานา.\csromanspacer{} สมุฏฺฐานมฺปิ เนสํ วุตฺตํ, สญฺญาสมุฏฺฐานาติสฺส วจนียํ.\csromanspacer{} กตมา สญฺญา.\csromanspacer{} สญฺญาปิ หิ พหู อเนกวิธา นานปฺปการกา, กามสญฺญา พฺยาปาทสญฺญา วิหิํสาสญฺญา.\csromanspacer{} อิโตสมุฏฺฐานา อกุสลา สงฺกปฺปา.}

\prose{อิเม จ ถปติ อกุสลา สงฺกปฺปา กุหิํ อปริเสสา นิรุชฺฌนฺติ.\csromanspacer{} นิโรโธปิ เนสํ วุตฺโต, อิธ ถปติ ภิกฺขุ วิวิจฺเจว กาเมหิ ฯเปฯ ปฐมํ ฌานํ อุปสมฺปชฺช วิหรติ.\csromanspacer{} เอตฺเถเต อกุสลา สงฺกปฺปา อปริเสสา นิรุชฺฌนฺติ.}

\prose{กถํ ปฏิปนฺโน จ ถปติ อกุสลานํ สงฺกปฺปานํ นิโรธาย ปฏิปนฺโน โหติ.\csromanspacer{} อิธ ถปติ ภิกฺขุ อนุปฺปนฺนานํ ปาปกานํ อกุสลานํ ธมฺมานํ อนุปฺปาทาย ฉนฺทํ ชเนติ วายมติ วีริยํ อารภติ จิตฺตํ ปคฺคณฺหาติ ปทหติ, อุปฺปนฺนานํ ปาปกานํ อกุสลานํ ธมฺมานํ ปหานาย ฯเปฯ อนุปฺปนฺนานํ กุสลานํ ธมฺมานํ อุปฺปาทาย ฯเปฯ อุปฺปนฺนานํ กุสลานํ ธมฺมานํ ฐิติยา อสมฺโมสาย ภิยฺโยภาวาย เวปุลฺลาย ภาวนาย ปาริปูริยา ฉนฺทํ ชเนติ วายมติ วีริยํ อารภติ จิตฺตํ ปคฺคณฺหาติ ปทหติ.\csromanspacer{} เอวํ ปฏิปนฺโน โข ถปติ อกุสลานํ สงฺกปฺปานํ นิโรธาย ปฏิปนฺโน โหติ.}

\proseitem{267}{กตเม จ ถปติ กุสลา สงฺกปฺปา.\csromanspacer{} เนกฺขมฺมสงฺกปฺโป อพฺยาปาทสงฺกปฺโป อวิหิํสาสงฺกปฺโป.\csromanspacer{} อิเม วุจฺจนฺติ ถปติ กุสลา สงฺกปฺปา.}

\prose{อิเม จ ถปติ กุสลา สงฺกปฺปา กิํสมุฏฺฐานา.\csromanspacer{} สมุฏฺฐานมฺปิ เนสํ วุตฺตํ, สญฺญาสมุฏฺฐานาติสฺส วจนียํ.\csromanspacer{} กตมา สญฺญา.\csromanspacer{} สญฺญาปิ หิ พหู \csromanfolio{220}อเนกวิธา นานปฺปการกา, เนกฺขมฺมสญฺญา อพฺยาปาทสญฺญา อวิหิํสาสญฺญา.\csromanspacer{} อิโตสมุฏฺฐานา กุสลา สงฺกปฺปา.}

\prose{อิเม จ ถปติ กุสลา สงฺกปฺปา กุหิํ อปริเสสา นิรุชฺฌนฺติ.\csromanspacer{} นิโรโธปิ เนสํ วุตฺโต, อิธ ถปติ ภิกฺขุ วิตกฺกวิจารานํ วูปสมา ฯเปฯ ทุติยํ ฌานํ อุปสมฺปชฺช วิหรติ.\csromanspacer{} เอตฺเถเต กุสลา สงฺกปฺปา อปริเสสา นิรุชฺฌนฺติ.}

\prose{กถํ ปฏิปนฺโน จ ถปติ กุสลานํ สงฺกปฺปานํ นิโรธาย ปฏิปนฺโน โหติ.\csromanspacer{} อิธ ถปติ ภิกฺขุ อนุปฺปนฺนานํ ปาปกานํ อกุสลานํ ธมฺมานํ อนุปฺปาทาย ฉนฺทํ ชเนติ วายมติ วีริยํ อารภติ จิตฺตํ ปคฺคณฺหาติ ปทหติ, อุปฺปนฺนานํ ปาปกานํ อกุสลานํ ธมฺมานํ ปหานาย ฯเปฯ อนุปฺปนฺนานํ กุสลานํ ธมฺมานํ อุปฺปาทาย ฯเปฯ อุปฺปนฺนานํ กุสลานํ ธมฺมานํ ฐิติยา อสมฺโมสาย ภิยฺโยภาวาย เวปุลฺลาย ภาวนาย ปาริปูริยา ฉนฺทํ ชเนติ วายมติ วีริยํ อารภติ จิตฺตํ ปคฺคณฺหาติ ปทหติ.\csromanspacer{} เอวํ ปฏิปนฺโน โข ถปติ กุสลานํ สงฺกปฺปานํ นิโรธาย ปฏิปนฺโน โหติ.}

\proseitem{268}{กตเมหิ จาหํ ถปติ ทสหิ ธมฺเมหิ สมนฺนาคตํ ปุริสปุคฺคลํ ปญฺญเปมิ สมฺปนฺนกุสลํ ปรมกุสลํ อุตฺตมปตฺติปตฺตํ สมณํ อโยชฺฌํ.\csromanspacer{} อิธ ถปติ ภิกฺขุ อเสขาย สมฺมาทิฏฺฐิยา สมนฺนาคโต โหติ, อเสเขน สมฺมาสงฺกปฺเปน สมนฺนาคโต โหติ, อเสขาย สมฺมาวาจาย สมนฺนาคโต โหติ, อเสเขน สมฺมากมฺมนฺเตน สมนฺนาคโต โหติ, อเสเขน สมฺมา-อาชีเวน สมนฺนาคโต โหติ, อเสเขน สมฺมาวายาเมน สมนฺนาคโต โหติ, อเสขาย สมฺมาสติยา สมนฺนาคโต โหติ, อเสเขน สมฺมาสมาธินา สมนฺนาคโต โหติ, อเสเขน สมฺมาญาเณน สมนฺนาคโต โหติ, อเสขาย สมฺมาวิมุตฺติยา สมนฺนาคโต โหติ.\csromanspacer{} อิเมหิ โข อหํ ถปติ ทสหิ ธมฺเมหิ สมนฺนาคตํ ปุริสปุคฺคลํ ปญฺญเปมิ สมฺปนฺนกุสลํ ปรมกุสลํ อุตฺตมปตฺติปตฺตํ สมณํ อโยชฺฌนฺติ.}

\prose{อิทมโวจ ภควา.\csromanspacer{} อตฺตมโน ปญฺจกงฺโค ถปติ ภควโต ภาสิตํ อภินนฺทีติ.}

\nitthitam{สมณมุณฺฑิกสุตฺตํ นิฏฺฐิตํ อฏฺฐมํ.}
\csromanmatikaanchor{mtk.34}
\csromantocmarkref{subsection}{9. จูฬสกุลุทายิสุตฺต}{mtk.34}
\bookmark[dest={mtk.34},level=2]{9. จูฬสกุลุทายิสุตฺต}
\csromanheader{2}{9. จูฬสกุลุทายีสุตฺต (79) \textbf{9. จูฬสกุลุทายิสุตฺต}}
\proseitem{269}{\csromanfolio{221}เอวํ เม สุตํ\csromandash{}เอกํ สมยํ ภควา ราชคเห วิหรติ เวฬุวเน กลนฺทกนิวาเป.\csromanspacer{} เตน โข ปน สมเยน สกุลุทายี ปริพฺพาชโก โมรนิวาเป ปริพฺพาชการาเม ปฏิวสติ มหติยา ปริพฺพาชกปริสาย สทฺธิํ.\csromanspacer{} อถ โข ภควา ปุพฺพณฺหสมยํ นิวาเสตฺวา ปตฺตจีวรมาทาย ราชคหํ ปิณฺฑาย ปาวิสิ.\csromanspacer{} อถ โข ภควโต เอตทโหสิ “อติปฺปโค โข ตาว ราชคเห ปิณฺฑาย จริตุํ.\csromanspacer{} ยํนูนาหํ เยน โมรนิวาโป ปริพฺพาชการาโม เยน สกุลุทายี ปริพฺพาชโก เตนุปสงฺกเมยฺยนฺ”ติ.\csromanspacer{} อถ โข ภควา เยน โมรนิวาโป ปริพฺพาชการาโม เตนุปสงฺกมิ.}

\prose{เตน โข ปน สมเยน สกุลุทายี ปริพฺพาชโก มหติยา ปริพฺพาชกปริสาย สทฺธิํ นิสินฺโน โหติ อุนฺนาทินิยา อุจฺจาสทฺทมหาสทฺทาย อเนกวิหิตํ ติรจฺฉานกถํ กเถนฺติยา.\csromanspacer{} เสยฺยถิทํ, ราชกถํ โจรกถํ มหามตฺตกถํ เสนากถํ ภยกถํ ยุทฺธกถํ อนฺนกถํ ปานกถํ วตฺถกถํ สยนกถํ มาลากถํ คนฺธกถํ ญาติกถํ ยานกถํ คามกถํ นิคมกถํ นครกถํ ชนปทกถํ อิตฺถิกถํ สูรกถํ วิสิขากถํ กุมฺภฏฺฐานกถํ ปุพฺพเปตกถํ นานตฺตกถํ โลกกฺขายิกํ สมุทฺทกฺขายิกํ อิติภวาภวกถํ อิติ วา.\csromanspacer{} อทฺทสา โข สกุลุทายี ปริพฺพาชโก ภควนฺตํ ทูรโตว อาคจฺฉนฺตํ, ทิสฺวาน สกํ ปริสํ สณฺฐาเปสิ “อปฺปสทฺทา โภนฺโต โหนฺตุ, มา โภนฺโต สทฺทมกตฺถ, อยํ สมโณ โคตโม อาคจฺฉติ.\csromanspacer{} อปฺปสทฺทกาโม โข ปน โส อายสฺมา อปฺปสทฺทสฺส วณฺณวาที, อปฺเปว นาม อปฺปสทฺทํ ปริสํ วิทิตฺวา อุปสงฺกมิตพฺพํ มญฺเญยฺยา”ติ.\csromanspacer{} อถ โข เต ปริพฺพาชกา ตุณฺหี อเหสุํ.}

\proseitem{270}{อถ โข ภควา เยน สกุลุทายี ปริพฺพาชโก เตนุปสงฺกมิ.\csromanspacer{} อถ โข สกุลุทายี ปริพฺพาชโก ภควนฺตํ เอตทโวจ “เอตุ โข ภนฺเต ภควา, สฺวาคตํ ภนฺเต ภควโต.\csromanspacer{} จิรสฺสํ โข ภนฺเต ภควา อิมํ ปริยายมกาสิ, ยทิทํ อิธาคมนาย.\csromanspacer{} นิสีทตุ ภนฺเต ภควา, อิทมาสนํ ปญฺญตฺตนฺ”ติ.\csromanspacer{} นสีทิ ภควา ปญฺญตฺเต อาสเน.\csromanspacer{} สกุลุทายีปิ โข ปริพฺพาชโก อญฺญตรํ นีจํ อาสนํ คเหตฺวา เอกมนฺตํ นิสีทิ, เอกมนฺตํ นิสินฺนํ โข สกุลุทายิํ ปริพฺพาชกํ ภควา เอตทโวจ “กาย นุตฺถ อุทายิ \csromanfolio{222}เอตรหิ กถาย สนฺนิสินฺนา, กา จ ปน โว อนฺตรากถา วิปฺปกตา”ติ.\csromanspacer{} ติฏฺฐเตสา ภนฺเต กถา, ยาย มยํ เอตรหิ กถาย สนฺนิสินฺนา, เนสา ภนฺเต กถา ภควโต ทุลฺลภา ภวิสฺสติ ปจฺฉาปิ สวนาย.\csromanspacer{} ยทาหํ ภนฺเต อิมํ ปริสํ อนุปสงฺกนฺโต โหมิ, อถายํ ปริสา อเนกวิหิตํ ติรจฺฉานกถํ กเถนฺตี นิสินฺนา โหติ.\csromanspacer{} ยทา จ โข อหํ ภนฺเต อิมํ ปริสํ อุปสงฺกนฺโต โหมิ, อถายํ ปริสา มมญฺเญว มุขํ อุลฺโลเกนฺตี นิสินฺนา โหติ “ยํ โน สมโณ อุทายี ธมฺมํ ภาสิสฺสติ, ตํ\footnote{ตํ โน (สี, สฺยา, กํ, อิ)} โสสฺสามา”ติ.\csromanspacer{} ยทา ปน ภนฺเต ภควา อิมํ ปริสํ อุปสงฺกนฺโต โหติ, อถาหญฺเจว อยญฺจ ปริสา ภควโต มุขํ อุลฺโลเกนฺตา\footnote{โอโลเกนฺตี (สฺยา, กํ, ก)} นิสินฺนา โหม “ยํ โน ภควา ธมฺมํ ภาสิสฺสติ, ตํ1 โสสฺสามา”ติ.}

\proseitem{271}{เตนหุทายิ ตํเยเวตฺถ ปฏิภาตุ, ยถา มํ ปฏิภาเสยฺยาติ.\csromanspacer{} ปุริมานิ ภนฺเต ทิวสานิ ปุริมตรานิ สพฺพญฺญู สพฺพทสฺสาวี อปริเสสํ ญาณทสฺสนํ ปฏิชานมาโน “จรโต จ เม ติฏฺฐโต จ สุตฺตสฺส จ ชาครสฺส จ สตตํ สมิตํ ญาณทสฺสนํ ปจฺจุปฏฺฐิตนฺ”ติ, โส มยา\footnote{ปจฺจุปฏฺฐิตนฺ”ติ มยา (?)} ปุพฺพนฺตํ อารพฺภ ปญฺหํ ปุฏฺโฐ สมาโน อญฺเญนญฺญํ ปฏิจริ, พหิทฺธา กถํ อปนาเมสิ, โกปญฺจ โทสญฺจ อปฺปจฺจยญฺจ ปาตฺวากาสิ.\csromanspacer{} ตสฺส มยฺหํ ภนฺเต ภควนฺตํเยว อารพฺภ สติ อุทปาทิ “อโห นูน ภควา, อโห นูน สุคโต, โย อิเมสํ ธมฺมานํ สุกุสโล”ติ.\csromanspacer{} โก ปน โส อุทายิ สพฺพญฺญู สพฺพทสฺสาวี อปริเสสํ ญาณทสฺสนํ ปฏิชานมาโน “จรโต จ เม ติฏฺฐโต จ สุตฺตสฺส จ ชาครสฺส จ สตตํ สมิตํ ญาณทสฺสนํ ปจฺจุปฏฺฐิตนฺ”ติ, โย ตยา ปุพฺพนฺตํ อารพฺภ ปญฺหํ ปุฏฺโฐ สมาโน อญฺเญนญฺญํ ปฏิจริ, พหิทฺธา กถํ อปนาเมสิ, โกปญฺจ โทสญฺจ อปฺปจฺจยญฺจ ปาตฺวากาสีติ.\csromanspacer{} นิคณฺโฐ ภนฺเต นาฏปุตฺโตติ.}

\prose{โย โข อุทายิ อเนกวิหิตํ ปุพฺเพนิวาสํ อนุสฺสเรยฺย.\csromanspacer{} เสยฺยถิทํ, เอกมฺปิ ชาติํ เทฺวปิ ชาติโย ฯเปฯ อิติ สาการํ ส-อุทฺเทสํ อเนกวิหิตํ ปุพฺเพนิวาสํ อนุสฺสเรยฺย.\csromanspacer{} โส วา มํ ปุพฺพนฺตํ อารพฺภ ปญฺหํ ปุจฺเฉยฺย, ตํ วาหํ ปุพฺพนฺตํ อารพฺภ ปญฺหํ ปุจฺเฉยฺยํ, โส วา เม ปุพฺพนฺตํ อารพฺภ}

\vspace{\tipitakaparskip}
\csromancenter{จูฬสกุลุทายีสุตฺต (79) ปญฺหสฺส เวยฺยากรเณน จิตฺตํ อาราเธยฺย, ตสฺส วาหํ ปุพฺพนฺตํ อารพฺภ ปญฺหสฺส เวยฺยากรเณน จิตฺตํ อาราเธยฺยํ.}
\prose{\csromanfolio{223}โย\footnote{โส (สี, อิ)} โข อุทายิ ทิพฺเพน จกฺขุนา วิสุทฺเธน อติกฺกนฺตมานุสเกน สตฺเต ปสฺเสยฺย จวมาเน อุปปชฺชมาเน หีเน ปณีเต สุวณฺเณ ทุพฺพณฺเณ สุคเต ทุคฺคเต, ยถากมฺมูปเค สตฺเต ปชาเนยฺย.\csromanspacer{} โส วา มํ อปรนฺตํ อารพฺภ ปญฺหํ ปุจฺเฉยฺย, ตํ วาหํ อปรนฺตํ อารพฺภ ปญฺหํ ปุจฺเฉยฺยํ.\csromanspacer{} โส วา เม อปรนฺตํ อารพฺภ ปญฺหสฺส เวยฺยากรเณน จิตฺตํ อาราเธยฺย, ตสฺส วาหํ อปรนฺตํ อารพฺภ ปญฺหสฺส เวยฺยากรเณน จิตฺตํ อาราเธยฺยํ.}

\prose{อปิ จ อุทายิ ติฏฺฐตุ ปุพฺพนฺโต, ติฏฺฐตุ อปรนฺโต ธมฺมํ เต เทเสสฺสามิ “อิมสฺมิํ สติ อิทํ โหติ, อิมสฺสุปฺปาทา อิทํ อุปฺปชฺชติ, อิมสฺมิํ อสติ อิทํ น โหติ, อิมสฺส นิโรธา อิทํ นิรุชฺฌตี”ติ.}

\prose{อหํ หิ ภนฺเต ยาวตกํปิ เม อิมินา อตฺตภาเวน ปจฺจนุภูตํ, ตมฺปิ นปฺปโหมิ สาการํ ส-อุทฺเทสํ อนุสฺสริตุํ, กุโต ปนาหํ อเนกวิหิตํ ปุพฺเพนิวาสํ อนุสฺสริสฺสามิ.\csromanspacer{} เสยฺยถิทํ, เอกมฺปิ ชาติํ เทฺวปิ ชาติโย ฯเปฯ อิติ สาการํ ส-อุทฺเทสํ อเนกวิหิตํ ปุพฺเพนิวาสํ อนุสฺสริสฺสามิ, เสยฺยถาปิ ภควา.\csromanspacer{} อหํ หิ ภนฺเต เอตรหิ ปํสุปิสาจกมฺปิ น ปสฺสามิ, กุโต ปนาหํ ทิพฺเพน จกฺขุนา วิสุทฺเธน อติกฺกนฺตมานุสเกน สตฺเต ปสฺสิสฺสามิ จวมาเน อุปปชฺชมาเน หีเน ปณีเต สุวณฺเณ ทุพฺพณฺเณ สุคเต ทุคฺคเต, ยถากมฺมูปเค สตฺเต ปชานิสฺสามิ, เสยฺยถาปิ ภควา.\csromanspacer{} ยํ ปน มํ ภนฺเต ภควา เอวมาห “อปิ จ อุทายิ ติฏฺฐตุ ปุพฺพนฺโต, ติฏฺฐตุ อปรนฺโต, ธมฺมํ เต เทเสสฺสามิ ‘อิมสฺมิํ สติ อิทํ โหติ, อิมสฺสุปฺปาทา อิทํ อุปฺปชฺชติ, อิมสฺมิํ อสติ อิทํ น โหติ, อิมสฺส นิโรธา อิทํ นิรุชฺฌตี’ติ”.\csromanspacer{} ตญฺจ ปน เม ภิยฺโยโส มตฺตาย น ปกฺขายติ, อปฺเปว นามาหํ ภนฺเต สเก อาจริยเก ภควโต จิตฺตํ อาราเธยฺยํ ปญฺหสฺส เวยฺยากรเณนาติ.}

\proseitem{272}{กินฺติ ปน เต อุทายิ สเก อาจริยเก โหตีติ.\csromanspacer{} อมฺหากํ ภนฺเต สเก อาจริยเก เอวํ โหติ “อยํ ปรโม วณฺโณ อยํ ปรโม วณฺโณ”ติ.}

\prose{\csromanfolio{224}ยํ ปน เต เอตํ อุทายิ สเก อาจริยเก เอวํ โหติ “อยํ ปรโม วณฺโณ อยํ ปรโม วณฺโณ”ติ, กตโม โส ปรโม วณฺโณติ.\csromanspacer{} ยสฺมา ภนฺเต วณฺณา อญฺโญ วณฺโณ อุตฺตริตโร วา ปณีตตโร วา นตฺถิ, โส ปรโม วณฺโณติ.}

\prose{กตโม ปน โส ปรโม วณฺโณ, ยสฺมา วณฺณา อญฺโญ วณฺโณ อุตฺตริตโร วา ปณีตตโร วา นตฺถีติ.\csromanspacer{} ยสฺมา ภนฺเต วณฺณา อญฺโญ วณฺโณ อุตฺตริตโร วา ปณีตตโร วา นตฺถิ, โส ปรโม วณฺโณติ.}

\prose{ทีฆาปิ โข เต เอสา อุทายิ ผเรยฺย. “ยสฺมา ภนฺเต วณฺณา อญฺโญ วณฺโณ อุตฺตริตโร วา ปณีตตโร วา นตฺถิ, โส ปรโม วณฺโณ”ติ วเทสิ, ตญฺจ วณฺณํ น ปญฺญเปสิ.\csromanspacer{} เสยฺยถาปิ อุทายิ ปุริโส เอวํ วเทยฺย “อหํ ยา อิมสฺมิํ ชนปเท ชนปทกลฺยาณี, ตํ อิจฺฉามิ ตํ กาเมมี”ติ.\csromanspacer{} ตเมนํ เอวํ วเทยฺยุํ “อมฺโภ ปุริส ยํ ตฺวํ ชนปทกลฺยาณิํ อิจฺฉสิ กาเมสิ, ชานาสิ ตํ ชนปทกลฺยาณิํ ‘ขตฺติยี วา พฺราหฺมณี วา เวสฺสี วา สุทฺที วา’ติ”, อิติ ปุฏฺโฐ โนติ วเทยฺย.\csromanspacer{} ตเมนํ เอวํ วเทยฺยุํ “อมฺโภ ปุริส ยํ ตฺวํ ชนปทกลฺยาณิํ อิจฺฉสิ กาเมสิ, ชานาสิ ตํ ชนปทกลฺยาณิํ ‘เอวํนามา เอวํโคตฺตาติ วา’ติ ฯเปฯ ‘ทีฆา วา รสฺสา วา มชฺฌิมา วา กาฬี วา สามา วา มงฺคุรจฺฉวี วา’ติ. ‘อมุกสฺมิํ วาเม วา นิคเม วา นคเร วา’ติ”, อิติ ปุฏฺโฐ โนติ วเทยฺย.\csromanspacer{} ตเมนํ เอวํ วเทยฺยุํ “อมฺโภ ปุริส ยํ ตฺวํ น ชานาสิ น ปสฺสสิ, ตํ ตฺวํ อิจฺฉสิ กาเมสี”ติ, อิติ ปุฏฺโฐ อามาติ วเทยฺย.}

\prose{ตํ กิํ มญฺญสิ อุทายิ, นนุ เอวํ สนฺเต ตสฺส ปุริสสฺส อปฺปาฏิหีรกตํ ภาสิตํ สมฺปชฺชตีติ.\csromanspacer{} อทฺธา โข ภนฺเต เอวํ สนฺเต ตสฺส ปุริสสฺส อปฺปาฏิหีรกตํ ภาสิตํ สมฺปชฺชตีติ.}

\prose{เอวเมว โข ตฺวํ อุทายิ “ยสฺมา ภนฺเต วณฺณา อญฺโญ วณฺโณ อุตฺตริตโร วา ปณีตตโร วา นตฺถิ, โส ปรโม วณฺโณ”ติ วเทสิ, ตญฺจ วณฺณํ น ปญฺญเปสีติ.}

\prose{เสยฺยถาปิ ภนฺเต มณิ เวฬุริโย สุโภ ชาติมา อฏฺฐํโส สุปริกมฺมกโต ปณฺฑุกมฺพเล นิกฺขิตฺโต ภาสเต จ ตปเต จ วิโรจติ จ, เอวํ วณฺโณ อตฺตา โหติ อโรโค ปรํ มรณาติ.}

\csromanheader{2}{9. จูฬสกุลุทายีสุตฺต (79)}
\proseitem{273}{\csromanfolio{225}ตํ กิํ มญฺญสิ อุทายิ, โย วา มณิ เวฬุริโย สุโภ ชาติมา อฏฺฐํโส สุปริกมฺมกโต ปณฺฑุกมฺพเล นิกฺขิตฺโต ภาสเต จ ตปเต จ วิโรจติ จ, โย วา รตฺตนฺธการติมิสาย กิมิ ขชฺโชปนโก, อิเมสํ อุภินฺนํ วณฺณานํ กตโม วณฺโณ อภิกฺกนฺตตโร จ ปณีตตโร จาติ.\csromanspacer{} ยฺวายํ ภนฺเต รตฺตนฺธการติมิสาย กิมิ ขชฺโชปนโก, อยํ อิเมสํ อุภินฺนํ วณฺณานํ อภิกฺกนฺตตโร จ ปณีตตโร จาติ.}

\prose{ตํ กิํ มญฺญสิ อุทายิ, โย วา รตฺตนฺธการติมิสาย กิมิ ขชฺโชปนโก, โย วา รตฺตนฺธการติมิสาย เตลปฺปทีโป, อิเมสํ อุภินฺนํ วณฺณานํ กตโม วณฺโณ อภิกฺกนฺตตโร จ ปณีตตโร จาติ.\csromanspacer{} ยฺวายํ ภนฺเต รตฺตนฺธการติมิสาย เตลปฺปทีโป, อยํ อิเมสํ อุภินฺนํ วณฺณานํ อภิกฺกนฺตตโร จ ปณีตตโร จาติ.}

\prose{ตํ กิํ มญฺญสิ อุทายิ, โย วา รตฺตนฺธการติมิสาย เตลปฺปทีโป, โย วา รตฺตนฺธการติมิสาย มหา-อคฺคิกฺขนฺโธ, อิเมสํ อุภินฺนํ วณฺณานํ กตโม วณฺโณ อภิกฺกนฺตตโร จ ปณีตตโร จาติ.\csromanspacer{} ยฺวายํ ภนฺเต รตฺตนฺธการติมิสาย มหา-อคฺคิกฺขนฺโธ, อยํ อิเมสํ อุภินฺนํ วณฺณานํ อภิกฺกนฺตตโร ว ปณีตตโร จาติ.}

\prose{ตํ กิํ มญฺญสิ อุทายิ, โย วา รตฺตนฺธการติมิสาย มหา-อคฺคิกฺขนฺโธ, ยา วา รตฺติยา ปจฺจูสสมยํ วิทฺเธ วิคตวลาหเก เทเว โอสธิตารกา, อิเมสํ อุภินฺนํ วณฺณานํ กตโม วณฺโณ อภิกฺกนฺตตโร จ ปณีตตโร จาติ.\csromanspacer{} ยฺวายํ ภนฺเต รตฺติยา ปจฺจูสสมยํ วิทฺเธ วิคตวลาหเก เทเว โอสธิตารกา, อยํ อิเมสํ อุภินฺนํ วณฺณานํ อภิกฺกนฺตตโร จ ปณีตตโร จาติ.}

\prose{ตํ กิํ มญฺญสิ อุทายิ, ยา วา รตฺติยา ปจฺจูสสมยํ วิทฺเธ วิคตวลาหเก เทเว โอสธิตารกา, โย วา ตทหุโปสเถ ปนฺนรเส วิทฺเธ วิคตวลาหเก เทเว อภิโท\footnote{อภิเท (ก-สี), อภิโทสํ (ก) อภิโทติ อภิสทฺเทน สมานตฺถนิปาตปทํ (ฉกฺกงฺคุตฺตรฏีกา มหาวคฺค-อฏฺฐมสุตฺตวณฺณนา)} อฑฺฒรตฺตสมยํ จนฺโท, อิเมสํ อุภินฺนํ วณฺณานํ กตโม วณฺโณ อภิกฺกนฺตตโร จ ปณีตตโร จาติ.\csromanspacer{} ยฺวายํ ภนฺเต ตทหุโปสเถ ปนฺนรเส วิทฺเธ วิคตวลาหเก เทเว \csromanfolio{226}อภิโท อฑฺฒรตฺตสมยํ จนฺโท, อยํ อิเมสํ อุภินฺนํ วณฺณานํ อภิกฺกนฺตตโร จ ปณีตตโร จาติ.}

\prose{ตํ กิํ มญฺญสิ อุทายิ, โย วา ตทหุโปสเถ ปนฺนรเส วิทฺเธ วิคตวลาหเก เทเว อภิโท อฑฺฒรตฺตสมยํ จนฺโท, โย วา วสฺสานํ ปจฺฉิเม มาเส สรทสมเย วิทฺเธ วิคตวลาหเก เทเว อภิโท มชฺฌนฺหิกสมยํ สูริโย, อิเมสํ อุภินฺนํ วณฺณานํ กตโม วณฺโณ อภิกฺกนฺตตโร จ ปณีตตโร จาติ.\csromanspacer{} ยฺวายํ ภนฺเต วสฺสานํ ปจฺฉิเม มาเส สรทสมเย วิทฺเธ วิคตวลาหเก เทเว อภิโท มชฺฌนฺหิกสมยํ สูริโย, อยํ อิเมสํ อุภินฺนํ วณฺณานํ อภิกฺกนฺตตโร จ ปณีตตโร จาติ.}

\prose{อโต โข เต อุทายิ พหู หิ พหุตรา เทวา.\csromanspacer{} เย อิเมสํ จนฺทิมสูริยานํ อาภา นานุโภนฺติ, ตฺยาหํ ปชานามิ.\csromanspacer{} อถ จ ปนาหํ น วทามิ “ยสฺมา วณฺณา อญฺโญ วณฺโณ อุตฺตริตโร วา ปณีตตโร วา นตฺถี”ติ.\csromanspacer{} อถ จ ปน ตฺวํ อุทายิ “ยฺวายํ วณฺโณ กิมินา ขชฺโชปนเกน นีหีนตโร\footnote{หีนตโร (สี, อิ)} จ ปติกิฏฺฐตโร จ, โส ปรโม วณฺโณ”ติ วเทสิ, ตญฺจ วณฺณํ น ปญฺญเปสีติ.\csromanspacer{} อจฺฉิทํ\footnote{อจฺฉิร (ก), อจฺฉิท (?)} ภควา กถํ, อจฺฉิทํ สุคโต กถํติ.}

\prose{กิํ ปน ตฺวํ อุทายิ เอวํ วเทสิ “อจฺฉิทํ ภควา กถํ, อจฺฉิทํ สุคโต กถํ”ติ.\csromanspacer{} อมฺหากํ ภนฺเต สเก อาจริยเก เอวํ โหติ “อยํ ปรโม วณฺโณ อยํ ปรโม วณฺโณ”ติ.\csromanspacer{} เต มยํ ภนฺเต ภควตา สเก อาจริยเก สมนุยุญฺชิยมานา สมนุคฺคาหิยมานา สมนุภาสิยมานา ริตฺตา ตุจฺฉา อปรทฺธาติ.}

\proseitem{274}{กิํ ปนุทายิ อตฺถิ เอกนฺตสุโข โลโก, อตฺถิ อาการวตี ปฏิปทา เอกนฺตสุขสฺส โลกสฺส สจฺฉิกิริยายาติ.\csromanspacer{} อมฺหากํ ภนฺเต สเก อาจริยเก เอวํ โหติ “อตฺถิ เอกนฺตสุโข โลโก, อตฺถิ อาการวตี ปฏิปทา เอกนฺตสุขสฺส โลกสฺส สจฺฉิกิริยายา”ติ.}

\prose{กตมา ปน สา อุทายิ อาการวตี ปฏิปทา เอกนฺตสุขสฺส โลกสฺส สจฺฉิกิริยายาติ.\csromanspacer{} อิธ ภนฺเต เอกจฺโจ ปาณาติปาตํ ปหาย ปาณาติปาตา ปฏิวิรโต โหติ, อทินฺนาทานํ ปหาย อทินฺนาทานา}

\vspace{\tipitakaparskip}
\csromancenter{จูฬสกุลุทายีสุตฺต (79) ปฏิวิรโต โหติ, กาเมสุมิจฺฉาจารํ ปหาย กาเมสุมิจฺฉาจารา ปฏิวิรโต โหติ, มุสาวาทํ ปหาย มุสาวาทา ปฏิวิรโต โหติ, อญฺญตรํ วา ปน ตโปคุณํ สมาทาย วตฺตติ.\csromanspacer{} อยํ โข สา ภนฺเต อาการวตี ปฏิปทา เอกนฺตสุขสฺส โลกสฺส สจฺฉิกิริยายาติ.}
\prose{\csromanfolio{227}ตํ กิํ มญฺญสิ อุทายิ, ยสฺมิํ สมเย ปาณาติปาตํ ปหาย ปาณาติปาตา ปฏิวิรโต โหติ, เอกนฺตสุขี วา ตสฺมิํ สมเย อตฺตา โหติ สุขทุกฺขี วาติ.\csromanspacer{} สุขทุกฺขี ภนฺเต.}

\prose{ตํ กิํ มญฺญสิ อุทายิ, ยสฺมิํ สมเย อทินฺนาทานํ ปหาย อทินฺนาทานา ปฏิวิรโต โหติ, เอกนฺตสุขี วา ตสฺมิํ สมเย อตฺตา โหติ สุขทุกฺขี วาติ.\csromanspacer{} สุขทุกฺขี ภนฺเต.}

\prose{ตํ กิํ มญฺญสิ อุทายิ, ยสฺมิํ สมเย กาเมสุมิจฺฉาจารํ ปหาย กาเมสุมิจฺฉาจารา ปฏิวิรโต โหติ, เอกนฺตสุขี วา ตสฺมิํ สมเย อตฺตา โหติ สุขทุกฺขี วาติ.\csromanspacer{} สุขทุกฺขี ภนฺเต.}

\prose{ตํ กิํ มญฺญสิ อุทายิ, ยสฺมิํ สมเย มุสาวาทํ ปหาย มุสาวาทา ปฏิวิรโต โหติ, เอกนฺตสุขี วา ตสฺมิํ สมเย อตฺตาโหติ สุขทุกฺขี วาติ.\csromanspacer{} สุขทุกฺขี ภนฺเต.}

\prose{ตํ กิํ มญฺญสิ อุทายิ, ยสฺมิํ สมเย อญฺญตรํ ตโปคุณํ สมาทาย วตฺตติ, เอกนฺตสุขี วา ตสฺมิํ สมเย อตฺตา โหติ สุขทุกฺขี วาติ.\csromanspacer{} สุขทุกฺขี ภนฺเต.}

\prose{ตํ กิํ มญฺญสิ อุทายิ, อปิ นุ โข โวกิณฺณสุขทุกฺกํ ปฏิปทํ อาคมฺม เอกนฺตสุขสฺส โลกสฺส สจฺฉิกิริยา โหตีติ\footnote{สจฺฉิกิริยายาติ (ก)}.\csromanspacer{} อจฺฉิทํ ภควา กถํ, อจฺฉิทํ สุคโต กถํติ.}

\prose{กิํ ปน ตฺวํ อุทายิ วเทสิ “อจฺฉิทํ ภควา กถํ, อจฺฉิทํ สุคโต กถํ”ติ.\csromanspacer{} อมฺหากํ ภนฺเต สเก อาจริยเก เอวํ โหติ “อตฺถิ เอกนฺตสุโข โลโก, อตฺถิ อาการวตี ปฏิปทา เอกนฺตสุขสฺส โลกสฺส สจฺฉิกิริยายา”ติ.\csromanspacer{} เต มยํ ภนฺเต ภควตา สเก อาจริยเก สมนุยุญฺชิยมานา สมนุคฺคาหิยมานา สมนุภาสิยมานา ริตฺตา ตุจฺฉา อปรทฺธาติ\footnote{อปรทฺธา (สี), อปรทฺธาปิ (สฺยา, กํ, อิ)}.}

\proseitem{275}{\csromanfolio{228}กิํ ปน ภนฺเต อตฺถิ เอกนฺตสุโข โลโก, อตฺถิ อาการวตี ปฏิปทา เอกนฺตสุขสฺส โลกสฺส สจฺฉิกิริยายาติ.\csromanspacer{} อตฺถิ โข อุทายิ เอกนฺตสุโข โลโก, อตฺถิ อาการวตี ปฏิปทา เอกนฺตสุขสฺส โลกสฺส สจฺฉิกิริยายาติ.}

\prose{กตมา ปน สา ภนฺเต อาการวตี ปฏิปทา เอกนฺตสุขสฺส โลกสฺส สจฺฉิกิริยายาติ.\csromanspacer{} อิธุทายิ ภิกฺขุ วิวิจฺเจว กาเมหิ ฯเปฯ ปฐมํ ฌานํ อุปสมฺปชฺช วิหรติ.\csromanspacer{} วิตกฺกวิจารานํ วูปสมา ฯเปฯ ทุติยํ ฌานํ อุปสมฺปชฺช วิหรติ.\csromanspacer{} ปีติยา จ วิราคา ฯเปฯ ตติย ฌานํ อุปสมฺปชฺช วิหรติ.\csromanspacer{} อยํ โข สา อุทายิ อาการวตี ปฏิปทา เอกนฺตสุขสฺส โลกสฺส สจฺฉิกิริยายาติ.}

\prose{น\footnote{กิํ นุ (สฺยา, กํ, ก)} โข สา ภนฺเต อาการวตี ปฏิปทา เอกนฺตสุขสฺส โลกสฺส สจฺฉิกิริยาย.\csromanspacer{} สจฺฉิกโต หิสฺส ภนฺเต เอตฺตาวตา เอกนฺตสุโข โลโก โหตีติ.\csromanspacer{} น ขฺวาสฺส อุทายิ เอตฺตาวตา เอกนฺตสุโข โลโก สจฺฉิกโต โหติ, อาการวตีเตฺวว สา ปฏิปทา เอกนฺตสุขสฺส โลกสฺส สจฺฉิกิริยายาติ.}

\prose{เอวํ วุตฺเต สกุลุทายิสฺส ปริพฺพาชกสฺส ปริสา อุนฺนาทินี อุจฺจาสทฺทมหาสทฺทา อโหสิ “เอตฺถ มยํ อนสฺสาม สาจริยกา, เอตฺถ มยํ อนสฺสาม\footnote{ปนสฺสาม (สี)} สาจริยกา.\csromanspacer{} น มยํ อิโต ภิยฺโย อุตฺตริตรํ ปชานามา”ติ.}

\prose{อถ โข สกุลุทายี ปริพฺพาชโก เต ปริพฺพาชเก อปฺปสทฺเท กตฺวา ภควนฺตํ เอตทโวจ “กิตฺตาวตา ปนาสฺส ภนฺเต เอกนฺตสุโข โลโก สจฺฉิกโต โหตี”ติ.\csromanspacer{} อิธุทายิ ภิกฺขุ สุขสฺส จ ปหานา ฯเปฯ จตุตฺถํ ฌานํ อุปสมฺปชฺช วิหรติ.\csromanspacer{} ยา ตา เทวตา เอกนฺตสุขํ โลกํ อุปปนฺนา, ตาหิ เทวตาหิ สทฺธิํ สนฺติฏฺฐติ สลฺลปติ สากจฺฉํ สมาปชฺชติ.\csromanspacer{} เอตฺตาวตา ขฺวาสฺส อุทายิ เอกนฺตสุโข โลโก สจฺฉิกโต โหตีติ.}

\proseitem{276}{เอตสฺส นูน ภนฺเต เอกนฺตสุขสฺส โลกสฺส สจฺฉิกิริยาเหตุ ภิกฺขู ภควติ พฺรหฺมจริยํ จรนฺตีติ.\csromanspacer{} น โข อุทายิ เอกนฺตสุขสฺส โลกสฺส สจฺฉิกิริยาเหตุ ภิกฺขู มยิ พฺรหฺมจริยํ จรนฺติ.\csromanspacer{} อตฺถิ โข}

\vspace{\tipitakaparskip}
\csromancenter{จูฬสกุลุทายีสุตฺต (79) อุทายิ อญฺเญว ธมฺมา อุตฺตริตรา จ ปณีตตรา จ, เยสํ สจฺฉิกิริยาเหตุ ภิกฺขู มยิ พฺรหฺมจริยํ จรนฺตีติ.}
\prose{\csromanfolio{229}กตเม ปน เต ภนฺเต ธมฺมา อุตฺตริตรา จ ปณีตตรา จ, เยสํ สจฺฉิกิริยาเหตุ ภิกฺขู ภควติ พฺรหฺมจริยํ จรนฺตีติ.\csromanspacer{} อิธุทายิ ตถาคโต โลเก อุปฺปชฺชติ อรหํ สมฺมาสมฺพุทฺโธ วิชฺชาจรณสมฺปนฺโน สุคโต โลกวิทู อนุตฺตโร ปุริสทมฺมสารถิ สตฺถา เทวมนุสฺสานํ พุทฺโธ ภควา ฯเปฯ.\csromanspacer{} โส อิเม ปญฺจ นีวรเณ ปหาย เจตโส อุปกฺกิเลเส ปญฺญาย ทุพฺพลีกรเณ, วิวิจฺเจว กาเมหิ ฯเปฯ ปฐมํ ฌานํ อุปสมฺปชฺช วิหรติ.\csromanspacer{} อยมฺปิ โข อุทายิ ธมฺโม อุตฺตริตโร จ ปณีตตโร จ, ยสฺส สจฺฉิกิริยาเหตุ ภิกฺขู มยิ พฺรหฺมจริยํ จรนฺติ.}

\prose{ปุน จปรํ อุทายิ ภิกฺขุ วิตกฺกวิจารานํ วูปสมา ทุติยํ ฌานํ.\csromanspacer{} ตติยํ ฌานํ.\csromanspacer{} จตุตฺถํ ฌานํ อุปสมฺปชฺช วิหรติ.\csromanspacer{} อยมฺปิ โข อุทายิ ธมฺโม อุตฺตริตโร จ ปณีตตโร จ, ยสฺส สจฺฉิกิริยาเหตุ ภิกฺขู มยิ พฺรหฺมจริยํ จรนฺติ.}

\prose{โส เอวํ สมาหิเต จิตฺเต ปริสุทฺเธ ปริโยทาเต อนงฺคเณ วิคตูปกฺกิเลเส มุทุภูเต กมฺมนิเย ฐิเต อาเนญฺชปฺปตฺเต ปุพฺเพนิวาสานุสฺสติญาณาย จิตฺตํ อภินินฺนาเมติ, โส อเนกวิหิตํ ปุพฺเพนิวาสํ อนุสฺสรติ.\csromanspacer{} เสยฺยถิทํ, เอกมฺปิ ชาติํ เทฺวปิ ชาติโย ฯเปฯ อิติ สาการํ ส-อุทฺเทสํ อเนกวิหิตํ ปุพฺเพนิวาสํ อนุสฺสรติ.\csromanspacer{} อยมฺปิ โข อุทายิ ธมฺโม อุตฺตริตโร จ ปณีตตโร จ, ยสฺส สจฺฉิกิริยาเหตุ ภิกฺขู มยิ พฺรหฺมจริยํ จรนฺติ.}

\prose{โส เอวํ สมาหิเต จิตฺเต ปริสุทฺเธ ปริโยทาเต อนงฺคเณ วิคตูปกฺกิเลเส มุทุภูเต กมฺมนิเย ฐิเต อาเนญฺชปฺปตฺเต สตฺตานํ จุตูปปาตญาณาย จิตฺตํ อภินินฺนาเมติ.\csromanspacer{} โส ทิพฺเพน จกฺขุนา วิสุทฺเธน อติกฺกนฺตมานุสเกน สตฺเต ปสฺสติ จวมาเน อุปปชฺชมาเน หีเน ปณีเต สุวณฺเณ ทุพฺพณฺเณ สุคเต ทุคฺคเต ฯเปฯ ยถากมฺมูปเค สตฺเต ปชานาติ.\csromanspacer{} อยมฺปิ โข อุทายิ ธมฺโม อุตฺตริตโร จ ปณีตตโร จ, ยสฺส สจฺฉิกิริยาเหตุ ภิกฺขู มยิ พฺรหฺมจริยํ จรนฺติ.}

\prose{โส เอวํ สมาหิเต จิตฺเต ปริสุทฺเธ ปริโยทาเต อนงฺคเณ วิคตูปกฺกิเลเส มุทุภูเต กมฺมนิเย ฐิเต อาเนญฺชปฺปตฺเต อาสวานํ \csromanfolio{230}ขยญาณาย จิตฺตํ อภินินฺนาเมติ.\csromanspacer{} โส อิทํ ทุกฺขนฺติ ยถาภูตํ ปชานาติ, อยํ ทุกฺขสมุทโยติ ฯเปฯ.\csromanspacer{} อยํ ทุกฺขนิโรโธติ.\csromanspacer{} อยํ ทุกฺขนิโรธคามินี ปฏิปทาติ ยถาภูตํ ปชานาติ.\csromanspacer{} อิเม อาสวาติ ยถาภูตํ ปชานาติ, อยํ อาสวสมุทโยติ.\csromanspacer{} อยํ อาสวนิโรโธติ.\csromanspacer{} อยํ อาสวนิโรธคามินี ปฏิปทาติ ยถาภูตํ ปชานาติ.\csromanspacer{} ตสฺส เอวํ ชานโต เอวํ ปสฺสโต กามาสวาปิ จิตฺตํ วิมุจฺจติ, ภวาสวาปิ จิตฺตํ วิมุจฺจติ, อวิชฺชาสวาปิ จิตฺตํ วิมุจฺจติ.\csromanspacer{} วิมุตฺตสฺมิํ วิมุตฺตมิติ ญาณํ โหติ, “ขีณา ชาติ, วุสิตํ พฺรหฺมจริยํ, กตํ กรณียํ, นาปรํ อิตฺถตฺตายา”ติ ปชานาติ.\csromanspacer{} อยมฺปิ โข อุทายิ ธมฺโม อุตฺตริตโร จ ปณีตตโร จ, ยสฺส สจฺฉิกิริยาเหตุ ภิกฺขู มยิ พฺรหฺมจริยํ จรนฺติ.\csromanspacer{} อิเม โข อุทายิ ธมฺมา อุตฺตริตรา จ ปณีตตรา จ, เยสํ สจฺฉิกิริยาเหตุ ภิกฺขู มยิ พฺรหฺมจริยํ จรนฺตีติ.}

\proseitem{277}{เอวํ วุตฺเต สกุลุทายี ปริพฺพาชโก ภควนฺตํ เอตทโวจ “อภิกฺกนฺตํ ภนฺเต อภิกฺกนฺตํ ภนฺเต, เสยฺยถาปิ ภนฺเต นิกฺกุชฺชิตํ วา อุกฺกุชฺเชยฺย, ปฏิจฺฉนฺนํ วา วิวเรยฺย, มูฬฺหสฺส วา มคฺคํ อาจิกฺเขยฺย, อนฺธกาเร วา เตลปชฺโชตํ ธาเรยฺย ‘จกฺขุมนฺโต รูปานิ ทกฺขนฺตี’ติ.\csromanspacer{} เอวเมวํ ภควตา อเนกปริยาเยน ธมฺโม ปกาสิโต, เอสาหํ ภนฺเต ภควนฺตํ สรณํ คจฺฉามิ ธมฺมญฺจ ภิกฺขุสํฆญฺจ, ลเภยฺยาหํ ภนฺเต ภควโต สนฺติเก ปพฺพชฺชํ ลเภยฺยํ อุปสมฺปทนฺ”ติ.}

\prose{เอวํ วุตฺเต สกุลุทายิสฺส ปริพฺพาชกสฺส ปริสา สกุลุทายิํ ปริพฺพาชกํ เอตทโวจุํ “มา ภวํ อุทายิ สมเณ โคตเม พฺรหฺมจริยํ จริ, มา ภวํ อุทายิ อาจริโย หุตฺวา อนฺเตวาสีวาสํ วสิ.\csromanspacer{} เสยฺยถาปิ นาม อุทกมณิโก\footnote{มณิโก (สี, อิ, ก)} หุตฺวา อุทญฺจนิโก\footnote{อุทฺเทกนิโก (สี, สฺยา, กํ, อิ)} อสฺส, เอวํ สมฺปทมิทํ\footnote{เอวํ สมฺปทเมตํ (สี, อิ)} โภโต อุทายิสฺส ภวิสฺสติ.\csromanspacer{} มา ภวํ อุทายิ สมเณ โคตเม พฺรหฺมจริยํ จริ, มา ภวํ อุทายิ อาจริโย หุตฺวา อนฺเตวาสีวาสํ วสี”ติ.\csromanspacer{} อิติ หิทํ สกุลุทายิสฺส ปริพฺพาชกสฺส ปริสา สกุลุทายิํ ปริพฺพาชกํ อนฺตรายมกาสิ ภควติ พฺรหฺมจริเยติ.}

\nitthitam{จูฬสกุลุทายิสุตฺตํ นิฏฺฐิตํ นวมํ.}
\csromanmatikaanchor{mtk.35}
\csromantocmarkref{subsection}{10. เวขนสสุตฺต}{mtk.35}
\bookmark[dest={mtk.35},level=2]{10. เวขนสสุตฺต}
\csromanheader{2}{10. เวขนสสุตฺต (80) \textbf{10. เวขนสสุตฺต}}
\proseitem{278}{\csromanfolio{231}เอวํ เม สุตํ\csromandash{}เอกํ สมยํ ภควา สาวตฺถิยํ วิหรติ เชตวเน อนาถปิณฺฑิกสฺส อาราเม.\csromanspacer{} อถ โข เวขนโส\footnote{เวขนสฺโส (สี, อิ)} ปริพฺพาชโก เยน ภควา เตนุปสงฺกมิ, อุปสงฺกมิตฺวา ภควตา สทฺธิํ สมฺโมทิ, สมฺโมทนียํ กถํ สารณียํ วีติสาเรตฺวา เอกมนฺตํ อฏฺฐาสิ, เอกมนฺตํ ฐิโต โข เวขนโส ปริพฺพาชโก ภควโต สนฺติเก อุทานํ อุทาเนสิ “อยํ ปรโม วณฺโณ อยํ ปรโม วณฺโณ”ติ.}

\prose{กิํ ปน ตฺวํ กจฺจาน เอวํ วเทสิ “อยํ ปรโม วณฺโณ อยํ ปรโม วณฺโณ”ติ.\csromanspacer{} กตโม กจฺจาน โส ปรโม วณฺโณติ.\csromanspacer{} ยสฺมา โภ โคตม วณฺณา อญฺโญ วณฺโณ อุตฺตริตโร วา ปณีตตโร วา นตฺถิ, โส ปรโม วณฺโณติ.\csromanspacer{} กตโม ปน โส กจฺจาน วณฺโณ, ยสฺมา วณฺณา อญฺโญ วณฺโณ อุตฺตริตโร วา ปณีตตโร วา นตฺถีติ.\csromanspacer{} ยสฺมา โภ โคตม วณฺณา อญฺโญ วณฺโณ อุตฺตริตโร วา ปณีตตโร วา นตฺถิ, โส ปรโม วณฺโณติ.\csromanspacer{} ทีฆาปิ โข เต เอสา กจฺจาน ผเรยฺย “ยสฺมา โภ โคตม วณฺณา อญฺโญ วณฺโณ อุตฺตริตโร วา ปณีตตโร วา นตฺถิ, โส ปรโม วณฺโณ”ติ วเทสิ, ตญฺจ วณฺณํ น ปญฺญเปสิ, เสยฺยถาปิ กจฺจาน ปุริโส เอวํ วเทยฺย “อหํ ยา อิมสฺมิํ ชนปเท ชนปทกลฺยาณี, ตํ อิจฺฉามิ ตํ กาเมมี”ติ.\csromanspacer{} ตเมนํ เอวํ วเทยฺยุํ “อมฺโภ ปุริส ยํ ตฺวํ ชนปทกลฺยาณิํ อิจฺฉสิ กาเมสิ, ชานาสิ ตํ ชนปทกลฺยาณิํ ‘ขตฺติยี วา พฺราหฺมณี วา เวสฺสี วา สุทฺที วา’ติ”, อิติ ปุฏฺโฐ โนติ วเทยฺย.\csromanspacer{} ตเมนํ เอวํ วเทยฺยุํ “อมฺโภ ปุริส ยํ ตฺวํ ชนปทกลฺยาณิํ อิจฺฉสิ กาเมสิ, ชานาสิ ตํ ชนปทกลฺยาณิํ ‘เอวํนามา เอวํโคตฺตาติ วา’ติ ฯเปฯ ‘ทีฆา วา รสฺสา วา มชฺฌิมา วา กาฬี วา สามา วา มงฺคุรจฺฉวี วา’ติ. ‘อมุกสฺมิํ คาเม วา นิคเม วา นคเร วา’ติ”, อิติ ปุฏฺโฐ โนติ วเทยฺย.\csromanspacer{} ตเมนํ เอวํ วเทยฺยุํ “อมฺโภ ปุริส ยํ ตฺวํ น ชานาสิ น ปสฺสสิ, ตํ ตฺวํ อิจฺฉสิ กาเมสี”ติ, อิติ ปุฏฺโฐ อามาติ วเทยฺย.}

\prose{ตํ กิํ มญฺญสิ กจฺจาน, นนุ เอวํ สนฺเต ภสฺส ปุริสสฺส อปฺปาฏีหีรกตํ ภาสิตํ สมฺปชฺชตีติ.\csromanspacer{} อทฺธา โข โภ โคตม เอวํ สนฺเต ตสฺส ปุริสสฺส \csromanfolio{232}อปฺปาฏิหีรกตํ ภาสิตํ สมฺปชฺชตีติ.\csromanspacer{} เอวเมว โข ตฺวํ กจฺจาน “ยสฺมา โภ โคตม วณฺณา อญฺโญ วณฺโณ อุตฺตริตโร วา ปณีตตโร วา นตฺถิ, โส ปรโม วณฺโณ”ติ วเทสิ, ตญฺจ วณฺณํ น ปญฺญเปสีติ.\csromanspacer{} เสยฺยถาปิ โภ โคตม มณิ เวฬุริโย สุโภ ชาติมา อฏฺฐํโส สุปริกมฺมกโต ปณฺฑุกมฺพเล นิกฺขิตฺโต ภาสเต จ ตปเต จ วิโรจติ จ, เอวํ วณฺโณ อตฺตา โหติ อโรโค ปรํ มรณาติ.}

\proseitem{279}{ตํ กิํ มญฺญสิ กจฺจาน, โย วา มณิ เวฬุริโย สุโภ ชาติมา อฏฺฐํโส สุปริกมฺมกโต ปณฺฑุกมฺพเล นิกฺขิตฺโต ภาสเต จ ตปเต จ วิโรจติ จ, โย วา รตฺตนฺธการติมิสาย กิมิ ขชฺโชปนโก, อิเมสํ อุภินฺนํ วณฺณานํ กตโม วณฺโณ อภิกฺกนฺตตโร จ ปณีตตโร จาติ.\csromanspacer{} ยฺวายํ โภ โคตม รตฺตนฺธการติมิสาย กิมิ ขชฺโชปนโก, อยํ อิเมสํ อุภินฺนํ วณฺณานํ อภิกฺกนฺตตโร จ ปณีตตโร จาติ.}

\prose{ตํ กิํ มญฺญสิ กจฺจาน, โย วา รตฺตนฺธการติมิสาย กิมิ ขชฺโชปนโก, โย วา รตฺตนฺธการติมิสาย เตลปฺปทีโป, อิเมสํ อุภินฺนํ วณฺณานํ กตโม วณฺโณ อภิกฺกนฺตตโร จ ปณีตตโร จาติ.\csromanspacer{} ยฺวายํ โภ โคตม รตฺตนฺธการติมิสาย เตลปฺปทีโป, อยํ อิเมสํ อุภินฺนํ วณฺณานํ อภิกฺกนฺตตโร จ ปณีตตโร จาติ.}

\prose{ตํ กิํ มญฺญสิ กจฺจาน, โย วา รตฺตนฺธการติมิสาย เตลปฺปทีโป, โย วา รตฺตนฺธการติมิสาย มหา-อคฺคิกฺขนฺโธ, อิเมสํ อุภินฺนํ วณฺณานํ กตโม วณฺโณ อภิกฺกนฺตตโร จ ปณีตตโร จาติ.\csromanspacer{} ยฺวายํ โภ โคตม รตฺตนฺธการติมิสาย มหา-อคฺคิกฺขนฺโธ, อยํ อิเมสํ อุภินฺนํ วณฺณานํ อภิกฺกนฺตตโร จ ปณีตตโร จาติ.}

\prose{ตํ กิํ มญฺญสิ กจฺจาน, โย วา รตฺตนฺธการติมิสาย มหา-อคฺคิกฺขนฺโธ, ยา วา รตฺติยา ปจฺจูสสมยํ วิทฺเธ วิคตวลาหเก เทเว โอสธิตารกา, อิเมสํ อุภินฺนํ วณฺณานํ กตโม วณฺโณ อภิกฺกนฺตตโร จ ปณีตตโร จาติ.\csromanspacer{} ยายํ โภ โคตม รตฺติยา ปจฺจูสสมยํ วิทฺเธ วิคตวลาหเก เทเว โอสธิตารกา, อยํ อิเมสํ อุภินฺนํ วณฺณานํ อภิกฺกนฺตตโร จ ปณีตตโร จาติ.\csromanspacer{} ตํ กิํ มญฺญสิ กจฺจาน, ยา วา รตฺติยา ปจฺจูสสมยํ วิทฺเธ วิคตวลาหเก เทเว โอสธิตารกา, โย วา ตทหุโปสเถ ปนฺนรเส วิทฺเธ}

\vspace{\tipitakaparskip}
\csromancenter{เวขนสสุตฺต (80) วิคตวลาหเก เทเว อภิโท อฑฺธุรตฺตสมยํ จนฺโท, อิเมสํ อุภินฺนํ วณฺณานํ กตโม วณฺโณ อภิกฺกนฺตตโร จ ปณีตตโร จาติ.\csromanspacer{} ยฺวายํ โภ โคตม ตทหุโปสเถ ปนฺนรเส วิทฺเธ วิคตวลาหเก เทเว อภิโท อฑฺฒุรตฺตสมยํ จนฺโท, อยํ อิเมยํ อุภินฺนํ วณฺณานํ อภิกฺกนฺตตโร จ ปณีตตโร จาติ.\csromanspacer{} ตํ กิํ มญฺญสิ กจฺจาน, โย วา ตทหุโปสเถ ปนฺนรเส วิทฺเธ วิคตวลาหเก เทเว อภิโท อฑฺฒรตฺตสมยํ จนฺโท, โย วา วสฺสานํ ปจฺฉิเม มาเส สรทสมเย วิทฺเธ วิคตวลาหเก เทเว อภิโท มชฺฌนฺหิกสมยํ สูริโย, อิเมสํ อุภินฺนํ วณฺณานํ กตโม วณฺโณ อภิกฺกนฺตตโร จ ปณีตตโร จาติ.\csromanspacer{} ยฺวายํ โภ โคตม วสฺสานํ ปจฺฉิเม มาเส สรทสมเย วิทฺเธ วิคตวลาหเก เทเว อภิโท มชฺฌนฺหิกสมยํ สูริโย, อยํ อิเมสํ อุภินฺนํ วณฺณานํ อภิกฺกนฺตตโร จ ปณีตตโร จาติ.\csromanspacer{} อโต โข เต กจฺจาน พหู หิ พหุตรา เทวา, เย อิเมสํ จนฺทิมสูริยานํ อาภา นานุโภนฺติ, ตฺยาหํ ปชานามิ, อถ จ ปนาหํ น วทามิ, ยสฺมา วณฺณา อญฺโญ วณฺโณ อุตฺตริตโร จ ปณีตตโร จ นตฺถีติ, อถ จ ปน ตฺวํ กจฺจาน ยฺวายํ วณฺโณ กิมินา ขชฺโชปนเกน นิหีนตโร จ ปติกิฏฺฐตโร จ โส ปรโม วณฺโณติ วเทติ, ตญฺจ วณฺณํ น ปญฺญเปสิ.}
\proseitem{280}{\csromanfolio{233}ปญฺจ โข อิเม กจฺจาน กามคุณา.\csromanspacer{} กตเม ปญฺจ, จกฺขุวิญฺเญยฺยา รูปา อิฏฺฐา กนฺตา มนาปา ปิยรูปา กามูปสํหิตา รชนียา, โสตวิญฺเญยฺยา สทฺทา ฯเปฯ ฆานวิญฺเญยฺยา คนฺธา.\csromanspacer{} ชิวฺหาวิญฺเญยฺยา รสา.\csromanspacer{} กายวิญฺเญยฺยา โผฏฺฐพฺพา อิฏฺฐา กนฺตา มนาปา ปิยรูปา กามูปสํหิตา รชนียา, อิเม โข กจฺจาน ปญฺจ กามคุณา.\csromanspacer{} ยํ โข กจฺจาน อิเม ปญฺจ กามคุเณ ปฏิจฺจ อุปฺปชฺชติ สุขํ โสมนสฺสํ, อิทํ วุจฺจติ กามสุขํ.\csromanspacer{} อิติ กาเมหิ กามสุขํ, กามสุขา กามคฺคสุขํ ตตฺถ อคฺคมกฺขายตีติ.}

\prose{เอวํ วุตฺเต เวขนโส ปริพฺพาชโก ภควนฺตํ เอตทโวจ “อจฺฉริยํ โภ โคตม, อพฺภุตํ โภ โคตม, ยาว สุภาสิตํ จิทํ โภตา โคตเมน ‘กาเมหิ กามสุขํ, กามสุขา กามคฺคสุขํ ตตฺถ อคฺคมกฺขายตี’ติ. (กาเมหิ โภ โคตม กามสุขํ กามสุขา กามคฺคสุขํ, ตตฺถ อคฺคมกฺขายตี”ติ.) ทุชฺชานํ โข เอตํ กจฺจาน \csromanfolio{234}ตยา อญฺญทิฏฺฐิเกน อญฺญขนฺติเกน อญฺญรุจิเกน อญฺญตฺรโยเคน อญฺญตฺราจริยเกน กามา\footnote{กามํ (สี, สฺยา, กํ, อิ)} วา กามสุขํ วา กามคฺคสุขํ วา, เย โข เต กจฺจาน ภิกฺขู อรหนฺโต ขีณาสวา วุสิตวนฺโต กตกรณียา โอหิตภารา อนุปฺปตฺตสทตฺถา ปริกฺขีณภวสํโยชนา สมฺมทญฺญา วิมุตฺตา, เต โข เอตํ ชาเนยฺยุํ, กามา วา กามสุขํ วา กามคฺคสุขํ วาติ.}

\proseitem{281}{เอวํ วุตฺเต เวขนโส ปริพฺพาชโก กุปิโต อนตฺตมโน ภควนฺตํเยว ขุํเสนฺโต ภควนฺตํเยว วมฺเภนฺโต ภควนฺตํเยว วทมาโน “สมโณ\footnote{สมโณ จ (สี, อิ) 3. ปนิเธเก (สี, อิ), ปนิเมเก (อุปริสุภสุตฺเต)} โคตโม ปาปิโต ภวิสฺสตี”ติ ภควนฺตํ เอตทโวจ “เอวเมว ปนิเธกจฺเจ3 สมณพฺราหฺมณา อชานนฺตา ปุพฺพนฺตํ อปสฺสนฺตา อปรนฺตํ, อถ จ ปน ‘ขีณา ชาติ, วุสิตํ พฺรหฺมจริยํ, กตํ กรณียํ, นาปรํ อิตฺถตฺตายาติ ปชานามา’ติ ปฏิชานนฺติ\footnote{อิตฺถตฺตายาติ ปฏิชานนฺติ (อิ) 5. เตสํ เตสายํ (สี), เตสํเยว โส (?)}, เตสมิทํ ภาสิตํ หสฺสกํเยว สมฺปชฺชติ, นามกํเยว สมฺปชฺชติ, ริตฺตกํเยว สมฺปชฺชติ, ตุจฺฉกํเยว สมฺปชฺชตี”ติ.\csromanspacer{} เย โข เต กจฺจาน สมณพฺราหฺมณา อชานนฺตา ปุพฺพนฺตํ อปสฺสนฺตา อปรนฺตํ, “ขีณา ชาติ, วุสิตํ พฺรหฺมจริยํ, กตํ กรณียํ, นาปรํ อิตฺถตฺตายาติ ปชานามา”ติ ปฏิชานนฺติ, เตสํ โสเยว5 สหธมฺมิโก นิคฺคโห โหติ.\csromanspacer{} อปิ จ กจฺจาน ติฏฺฐตุ ปุพฺพนฺโต, ติฏฺฐตุ อปรนฺโต, เอตุ วิญฺญู ปุริโส อสโฐ อมายาวี อุชุชาติโก, อหมนุสาสามิ, อหํ ธมฺมํ เทเสมิ, ยถานุสิฏฺฐํ ตถา ปฏิปชฺชามาโน\footnote{ยถานุสิฏฺฐํ ปฏิปชฺชมาโน (?) 7. เอวํ กิรายสฺมา (สฺยา, ก)} นจิรสฺเสว สามญฺเญว ญสฺสติ สามํ ทกฺขิติ, เอวํ กิร สมฺมา7 พนฺธนา วิปฺปโมกฺโข โหติ, ยทิทํ อวิชฺชา พนฺธนา, เสยฺยถาปิ กจฺจาน ทหโร กุมาโร มนฺโท อุตฺตานเสยฺยโก กณฺฐปญฺจเมหิ พนฺธเนหิ พทฺโธ อสฺส สุตฺตพนฺธเนหิ, ตสฺส วุทฺธิมนฺวาย อินฺทฺริยานํปริปากมนฺวาย ตานิ พนฺธนานิ มุจฺเจยฺยุํ, โส โมกฺโขมฺหีติ โข ชาเนยฺย โน จ พนฺธนํ, เอวเมว โข กจฺจาน เอตุ วิญฺญู ปุริโส อสโฐ อมายาวี อุชุชาติโก, อหมนุสาสามิ, อหํ ธมฺมํ เทเสมิ, ยถานุสิฏฺฐํ ตถา ปฏิปชฺชมาโน นจิรสฺเสว สามญฺเญว}

\vspace{\tipitakaparskip}
\csromancenter{เวขนสสุตฺต (80) ญสฺสติ สามํ ทกฺขิติ, เอวํ กิร สมฺมา พนฺธนา วิปฺปโมกฺโข โหติ, ยทิทํ อวิชฺชา พนฺธนาติ.}
\csromanmatikaanchor{mtk.36}
\csromantocmarkref{subsection}{อุทฺทานคาถาโย}{mtk.36}
\bookmark[dest={mtk.36},level=2]{อุทฺทานคาถาโย}
\prose{\csromanfolio{235}เอวํ วุตฺเต เวขนโส ปริพฺพาชโก ภควนฺตํ เอตทโวจ “อภิกฺกนฺตํ โภ โคตม ฯเปฯ อุปาสกํ มํ ภวํ โคตโม ธาเรตุ อชฺชตคฺเค ปาณุเปตํ สรณํ คตนฺ”ติ.}

\nitthitam{เวขนสสุตฺตํ นิฏฺฐิตํ ทสมํ.}
\nitthitam{\textbf{ปริพฺพาชกวคฺโค นิฏฺฐิโต ตติโย.}}
\titlehead{\textbf{ตสฺสุทฺทานํ}}
\csromangathameasure{ปุณฺฑรี อคฺคิสห กถินาโม, \\ ทีฆนโข ปุน ภารทฺวาชโคตฺโต. \\ สนฺทก-อุทายิมุณฺฑิกปุตฺโต, \\ มณิโก ตถากจฺจาโน วรวคฺโค.}
\csromangathagroup{\csromangathastanza{ปุณฺฑรี อคฺคิสห กถินาโม, \\ ทีฆนโข ปุน ภารทฺวาชโคตฺโต. \\ สนฺทก-อุทายิมุณฺฑิกปุตฺโต, \\ มณิโก ตถากจฺจาโน วรวคฺโค.}}

\csromanmatikaanchor{mtk.37}
\csromantocmarknopage{chapter}{4. ราชวคฺค}{mtk.37}
\bookmark[dest={mtk.37},level=0]{4. ราชวคฺค}
\chapterheadpageread{4. ราชวคฺค}
\csromanmatikaanchor{mtk.38}
\csromantocmarkref{subsection}{1. ฆฏิการสุตฺต}{mtk.38}
\bookmark[dest={mtk.38},level=2]{1. ฆฏิการสุตฺต}
\csromanheader{2}{1. ฆฏิการสุตฺต}
\proseitem{282}{\csromanfolio{236}เอวํ เม สุตํ\csromandash{}เอกํ สมยํ ภควา โกสเลสุ จาริกํ จรติ มหตา ภิกฺขุสํเฆน สทฺธิํ.\csromanspacer{} อถ โข ภควา มคฺคา โอกฺกมฺม อญฺญตรสฺมิํ ปเทเส สิตํ ปาตฺวากาสิ.\csromanspacer{} อถ โข อายสฺมโต อานนฺทสฺส เอตทโหสิ “โก นุ โข เหตุ โก ปจฺจโย ภควโต สิตสฺส ปาตุกมฺมาย, น อการเณน\footnote{น อการเณ (สี)} ตถาคตา สิตํ ปาตุกโรนฺตี”ติ.\csromanspacer{} อถ โข อายสฺมา อานนฺโท เอกํสํ จีวรํ\footnote{อุตฺตราสงฺคํ (สฺยา, กํ)} กตฺวา เยน ภควา เตนญฺชลิํ ปณาเมตฺวา ภควนฺตํ เอตทโวจ “โก นุ โข ภนฺเต เหตุ โก ปจฺจโย ภควโต สิตสฺส ปาตุกมฺมาย, น อการเณน ตถาคตา สิตํ ปาตุกโรนฺตี”ติ.\csromanspacer{} ภูตปุพฺพํ อานนฺท อิมสฺมิํ ปเทเส เวคฬิงฺคํ\footnote{เวหลิงฺคํ (สี), เวภลิคํ (สฺยา, กํ), เวภลิงฺคํ (อิ)} นาม คามนิคโม อโหสิ, อิทฺโธ เจว ผีโต จ พหุชโน อากิณฺณมนุสฺโส.\csromanspacer{} เวคฬิงฺคํ โข อานนฺท คามนิคมํ กสฺสโป ภควา อรหํ สมฺมาสมฺพุทฺโธ อุปนิสฺสาย วิหาสิ, อิธ สุทํ อานนฺท กสฺสปสฺส ภควโต อรหโต สมฺมาสมฺพุทฺธสฺส อาราโม อโหสิ, อิธ สุทํ อานนฺท กสฺสโป ภควา อรหํ สมฺมาสมฺพุทฺโธ นิสินฺนโก ภิกฺขุสํฆํ โอวทตีติ.\csromanspacer{} อถ โข อายสฺมา อานนฺโท จตุคฺคุณํ สงฺฆาฏิํ ปญฺญเปตฺวา ภควนฺตํ เอตทโวจ “เตน หิ ภนฺเต ภควา นิสีทตุ เอตฺถ, อยํ ภูมิปเทโส ทฺวีหิ อรหนฺเตหิ สมฺมาสมฺพุทฺเธหิ ปริภุตฺโต ภวิสฺสตี”ติ.\csromanspacer{} นิสีทิ ภควา ปญฺญตฺเต อาสเน, นิสชฺช โข ภควา อายสฺมนฺตํ อานนฺทํ อามนฺเตสิ.}

\prose{ภูตปุพฺพํ อานนฺท อิมสฺมิํ ปเทเส เวคฬิงฺคํ นาม คามนิคโม อโหสิ, อิทฺโธ เจว ผีโต จ พหุชโน อากิณฺณมนุสฺโส.\csromanspacer{} เวคฬิงฺคํ โข อานนฺท คามนิคมํ กสฺสโป ภควา อรหํ สมฺมาสมฺพุทฺโธ อุปนิสฺสาย วิหาสิ, อิธ สุทํ อานนฺท กสฺสปสฺส ภควโต อรหโต สมฺมาสมฺพุทฺธสฺส อาราโม อโหสิ, อิธ สุทํ อานนฺท กสฺสโป ภควา อรหํ สมฺมาสมฺพุทฺโธ นิสินฺนโก ภิกฺขุสํฆํ โอวทติ.}

\proseitem{283}{เวคฬิงฺเค โข อานนฺท คามนิคเม ฆฏิกาโร\footnote{ฆฏีกาโร (สี, อิ)} นาม กุมฺภกาโร กสฺสปสฺส ภควโต อรหโต สมฺมาสมฺพุทฺธสฺส อุปฏฺฐาโก}

\vspace{\tipitakaparskip}
\csromancenter{ฆฏิการสุตฺต (81) อโหสิ อคฺคุปฏฺฐาโก.\csromanspacer{} ฆาฏิการสฺส โข อานนฺท กุมฺภการสฺส โชติปาโล นาม มาณโว สหาโย อโหสิ ปิยสหาโย.\csromanspacer{} อถ โข อานนฺท ฆฏิกาโร กุมฺภกาโร โชติปาลํ มาณวํ อามนฺเตสิ “อายาม สมฺม โชติปาล กสฺสปํ ภควนฺตํ อรหนฺตํ สมฺมาสมฺพุทฺธํ ทสฺสนาย อุปสงฺกมิสฺสาม, สาธุสมฺมตํ หิ เม ตสฺส ภควโต ทสฺสนํ อรหโต สมฺมาสมฺพุทฺธสฺสา”ติ.\csromanspacer{} เอวํ วุตฺเต อานนฺท โชติปาโล มาณโว ฆฏิการํ กุมฺภการํ เอตทโวจ “อลํ สมฺม ฆฏิการ, กิํ ปน เตน มุณฺฑเกน สมณเกน ทิฏฺเฐนา”ติ.\csromanspacer{} ทุติยมฺปิ โข อานนฺท ฯเปฯ.\csromanspacer{} ตติยมฺปิ โข อานนฺท ฆฏิกาโร กุมฺภกาโร โชติปาลํ มาณวํ เอตทโวจ “อายาม สมฺม โชติปาล กสฺสปํ ภควนฺตํ อรหนฺตํ สมฺมาสมฺพุทฺธํ ทสฺสนาย อุปสงฺกมิสฺสาม, สาธุสมฺมตํ หิ เม ตสฺส ภควโต ทสฺสนํ อรหโต สมฺมาสมฺพุทฺธสฺสา”ติ.\csromanspacer{} ตติยมฺปิ โข อานนฺท โชติปาโล มาณโว ฆฏิการํ กุมฺภการํ เอตทโวจ “อลํ สมฺม ฆฏิการ, กิํ ปน มุณฺฑเกน สมณเกน ทิฏฺเฐนา”ติ.\csromanspacer{} เตน หิ สมฺม โชติปาล โสตฺติสินานิํ\footnote{โสตฺติํ สินานิํ (สี, อิ), โสตฺติสินานํ (สฺยา, กํ, ก)} อาทาย\footnote{อาหร (ก)} นทิํ คมิสฺสาม สินายิตุนฺติ. “เอวํ สมฺมา”ติ โข อานนฺท โชติปาโล มาณโว ฆฏิการสฺส กุมฺภการสฺส ปจฺจสฺโสสิ.\csromanspacer{} อถ โข อานนฺท ฆฏิกาโร จ กุมฺภกาโร โชติปาโล จ มาณโว โสตฺติสินานิํ อาทาย นทิํ อคมํสุ สินายิตุํ.}
\proseitem{284}{\csromanfolio{237}อถ โข อานนฺท ฆฏิกาโร กุมฺภกาโร โชติปาลํ มาณวํ อามนฺเตสิ “อยํ สมฺม โชติปาล กสฺสปสฺส ภควโต อรหโต สมฺมาสมฺพุทฺธสฺส อวิทูเร อาราโม, อายาม สมฺม โชติปาล กสฺสปํ ภควนฺตํ อรหนฺตํ สมฺมาสมฺพุทฺธํ ทสฺสนาย อุปสงฺกมิสฺสาม, สาธุสมฺมตํ หิ เม ตสฺส ภควโต ทสฺสนํ อรหโต สมฺมาสมฺพุทฺธสฺสา”ติ.\csromanspacer{} เอวํ วุตฺเต อานนฺท โชติปาโล มาณโว ฆฏิการํ กุมฺภการํ เอตทโวจ “อลํ สมฺม ฆฏิการ, กิํ ปน เตน มุณฺฑเกน สมณเกน ทิฏฺเฐนา”ติ.\csromanspacer{} ทุติยมฺปิ โข อานนฺท ฯเปฯ.\csromanspacer{} ตติยมฺปิ โข อานนฺท ฆฏิกาโร กุมฺภกาโร โชติปาลํ มาณวํ เอตทโวจ “อยํ สมฺม โชติปาล กสฺสปสฺส ภควโต อรหโต สมฺมาสมฺพุทฺธสฺส อวิทูเร อาราโม, อายาม สมฺม โชติปาล กสฺสปํ ภควนฺตํ อรหนฺตํ สมฺมาสมฺพุทฺธํ \csromanfolio{238}ทสฺสนาย อุปสงฺกมิสฺสาม, สาธุสมฺมตํ หิ เม ตสฺส ภควโต ทสฺสนํ อรหโต สมฺมาสมฺพุทฺธสฺสา”ติ.\csromanspacer{} ตติยมฺปิ โข อานนฺท โชติปาโล มาณโว ฆฏิการํ กุมฺภการํ เอตทโวจ “อลํ สมฺม ฆฏิการ, กิํ ปน เตน มุณฺฑเกน สมณเกน ทิฏฺเฐนา”ติ.\csromanspacer{} อถ โข อานนฺท ฆฏิกาโร กุมฺพกาโร โชติปาลํ มาณวํ โอวฏฺฏิกายํ ปรามสิตฺวา เอตทโวจ “อยํ สมฺม โชติปาล กสฺสปสฺส ภควโต อรหโต สมฺมาสมฺพุทฺธสฺส อวิทูเร อาราโม, อายาม สมฺม โชติปาล กสฺสปํ ภควนฺตํ อรหนฺตํ สมฺมาสมฺพุทฺธํ ทสฺสนาย อุปสงฺกมิสฺสาม, สาธุสมฺมตํ หิ เม ตสฺส ภควโต ทสฺสนํ อรหโต สมฺมาสมฺพุทฺธสฺสา”ติ.\csromanspacer{} อถ โข อานนฺท โชติปาโล มาณโว โอวฏฺฏิกํ วินิวฏฺเฏตฺวา\footnote{วินิเวเฐตฺวา (สี, สฺยา, กํ, อิ)} ฆฏิการํ กุมฺภการํ เอตทโวจ “อลํ สมฺม ฆฏิการ, กิํ ปน เตน มุณฺฑเกน สมณเกน ทิฏฺเฐนา”ติ.\csromanspacer{} อถ โข อานนฺท ฆฏิกาโร กุมฺภกาโร โชติปาลํ มาณวํ สีสํนฺหาตํ\footnote{สสีสํ นหาตํ (สี), สีสนฺหาตํ (สฺยา, กํ) 3. ยาเวตโทหิปิ (สี, สฺยา, กํ, อิ)} เกเสสุ ปรามสิตฺวา เอตทโวจ “อยํ สมฺม โชติปาล กสฺสปสฺส ภควโต อรหโต สมฺมาสมฺพุทฺธสฺส อวิทูเร อาราโม, อายาม สมฺม โชติปาล กสฺสปํ ภควนฺตํ อรหนฺตํ สมฺมาสมฺพุทฺธํ ทสฺสนาย อุปสงฺกมิสฺสาม, สาธุสมฺมตํ หิ เม ตสฺส ภควโต ทสฺสนํ อรหโต สมฺมาสมฺพุทฺธสฺสา”ติ.\csromanspacer{} อถ โข อานนฺท โชติปาลสฺส มาณวสฺส เอตทโหสิ “อจฺฉริยํ วต โภ, อพฺภุตํ วต โภ, ยตฺร หิ นามายํ ฆฏิกาโร กุมฺภกาโร อิตฺตรชจฺโจ สมาโน อมฺหากํ สีสํนฺหาตานํ เกเสสุ ปรามสิตพฺพํ มญฺญิสฺสติ, น วติทํ กิร โอรกํ มญฺเญ ภวิสฺสตี”ติ ฆฏิการํ กุมฺภการํ เอตทโวจ “ยาวตาโทหิปิ3 สมฺม ฆฏิการา”ติ.\csromanspacer{} ยาวตาโทหิปิ สมฺม โชติปาล, ตถา หิ ปน เม สาธุสมฺมตํ ตสฺส ภควโต ทสฺสนํ อรหโต สมฺมาสมฺพุทฺธสฺสาติ.\csromanspacer{} เตน หิ สมฺม ฆฏิการ มุญฺจ คมิสฺสามาติ.}

\proseitem{285}{อถ โข อานนฺท ฆฏิกาโร จ กุมฺภกาโร โชติปาโล จ มาณโว เยน กสฺสโป ภควา อรหํ สมฺมาสมฺพุทฺโธ เตนุปสงฺกมิํสุ, อุปสงฺกมิตฺวา ฆฏิกาโร กุมฺภกาโร กสฺสปํ ภควนฺตํ อรหนฺตํ สมฺมาสมฺพุทฺธํ อภิวาเทตฺวา เอกมนฺตํ นิสีทิ.\csromanspacer{} โชติปาโล ปน มาณโว}

\vspace{\tipitakaparskip}
\csromancenter{ฆฏิการสุตฺต (81) กสฺสเปน ภควตา อรหตา สมฺมาสมฺพุทฺเธน สทฺธิํ สมฺโมทิ, สมฺโมทนียํ กถํ สารณียํ วีติสาเรตฺวา เอกมนฺตํ นิสีทิ.\csromanspacer{} เอกมนฺตํ นิสินฺโน โข อานนฺท ฆฏิกาโร กุมฺภกาโร กสฺสปํ ภควนฺตํ อรหนฺตํ สมฺมาสมฺพุทฺธํ เอตทโวจ “อยํ เม ภนฺเต โชติปาโล มาณโว สหาโย ปิยสหาโย, อิมสฺส ภควา ธมฺมํ เทเสตู”ติ.\csromanspacer{} อถ โข อานนฺท กสฺสโป ภควา อรหํ สมฺมาสมฺพุทฺโธ ฆฏิการญฺจ กุมฺภการํ โชติปาลญฺจ มาณวํ ธมฺมิยา กถาย สนฺทสฺเสสิ สมาทเปสิ สมุตฺเตเชสิ สมฺปหํเสสิ.\csromanspacer{} อถ โข อานนฺท ฆฏิกาโร จ กุมฺภกาโร โชติปาโล จ มาณโว กสฺสเปน ภควตา อรหตา สมฺมาสมฺพุทฺเธน ธมฺมิยา กถาย สนฺทสฺสิตา สมาทปิตา สมุตฺเตชิตา สมฺปหํสิตา กสฺสปสฺส ภควโต อรหโต สมฺมาสมฺพุทฺธสฺส ภาสิตํ อภินนฺทิตฺวา อนุโมทิตฺวา อุฏฺฐายาสนา กสฺสปํ ภควนฺตํ อรหนฺตํ สมฺมาสมฺพุทฺธํ อภิวาเทตฺวา ปทกฺขิณํ กตฺวา ปกฺกมิํสุ.}
\proseitem{286}{\csromanfolio{239}อถ โข อานนฺท โชติปาโล มาณโว ฆฏิการํ กุมฺภการํ เอตทโวจ “อิมํ นุ ตฺวํ สมฺม ฆฏิการ ธมฺมํ สุณนฺโต, อถ จ ปน อคารสฺมา อนคาริยํ น ปพฺพชิสฺสสี”ติ.\csromanspacer{} นนุ มํ สมฺม โชติปาล ชานาสิ อนฺเธ ชิณฺเณ มาตาปิตโร โปเสมีติ.\csromanspacer{} เตน หิ สมฺม ฆฏิการ อหํ อคารสฺมา อนคาริยํ ปพฺพชิสฺสามีติ.\csromanspacer{} อถ โข อานนฺท ฆฏิกาโร จ กุมฺภกาโร โชติปาโล จ มาณโว เยน กสฺสโป ภควา อรหํ สมฺมาสมฺพุทฺโธ เตนุปสงฺกมิํสุ, อุปสงฺกมิตฺวา กสฺสปํ ภควนฺตํ อรหนฺตํ สมฺมาสมฺพุทฺธํ อภิวาเทตฺวา เอกมนฺตํ นิสีทิํสุ, เอกมนฺตํ นิสินฺโน โข อานนฺท ฆฏิกาโร กุมฺภกาโร กสฺสปํ ภควนฺตํ อรหนฺตํ สมฺมาสมฺพุทฺธํ เอตทโวจ “อยํ เม ภนฺเต โชติปาโล มาณโว สหาโย ปิยสหาโย, อิมํ ภควา ปพฺพาเชตู”ติ.\csromanspacer{} อลตฺถ โข อานนฺท โชติปาโล มาณโว กสฺสปสฺส ภควโต อรหโต สมฺมาสมฺพุทฺธสฺส สนฺติเก ปพฺพชฺชํ, อลตฺต อุปสมฺปทํ.}

\proseitem{287}{อถ โข อานนฺท กสฺสโป ภควา อรหํ สมฺมาสมฺพุทฺโธ อจิรูปสมฺปนฺเน โชติปาเล มาณเว อฑฺฒมาสุปสมฺปนฺเน เวคฬิงฺเค ยถาภิรนฺตํ วิหริตฺวา เยน พาราณสี เตน จาริกํ ปกฺกามิ, อนุปุพฺเพน จาริกํ จรมาโน เยน พาราณสี ตทวสริ, ตตฺร สุทํ อานนฺท กสฺสโป ภควา อรหํ สมฺมาสมฺพุทฺโธ พาราณสิยํ วิหรติ อิสิปตเน \csromanfolio{240}มิคทาเย.\csromanspacer{} อสฺโสสิ โข อานนฺท กิกี กาสิราชา กสฺสโป กิร ภควา อรหํ สมฺมาสมฺพุทฺโธ พาราณสิํ อนุปฺปตฺโต พาราณสิยํ วิหรติ อิสิปตเน มิคทาเยติ.\csromanspacer{} อถ โข อานนฺท กิกี กาสิราชา ภทฺรานิ ภทฺรานิ ยานานิ โยชาเปตฺวา ภทฺรํ\footnote{ภทฺรํ ภทฺรํ (ก)} ยานํ อภิรุหิตฺวา ภเทฺรหิ ภเทฺรหิ ยาเนหิ พาราณสิยา นิยฺยาสิ มหจฺจราชานุภาเวน\footnote{มหจฺจา ราชนุภาเวน (สี), มหตา ราชานุภาเวน (อิ)} กสฺสปํ ภควนฺตํ อรหนฺตํ สมฺมาสมฺพุทฺธํ ทสฺสนาย, ยาวติกา ยานสฺส ภูมิ, ยาเนน คนฺตฺวา ยานา ปจฺโจโรหิตฺวา ปตฺติโกว เยน กสฺสโป ภควา อรหํ สมฺมาสมฺพุทฺโธ เตนุปสงฺกมิ, อุปสงฺกมิตฺวา กสฺสปํ ภควนฺตํ อรหนฺตํ สมฺมาสมฺพุทฺธํ อภิวาเทตฺวา เอกมนฺตํ นิสีทิ, เอกมนฺตํ นิสินฺนํ โข อานนฺท กิกิํ กาสิราชานํ กสฺสโป ภควา อรหํ สมฺมาสมฺพุทฺโธ ธมฺมิยา กถาย สนฺทสฺเสสิ สมาทเปสิ สมุตฺเตเชสิ สมฺปหํเสสิ.\csromanspacer{} อถ โข อานนฺท กิกี กาสิราชา กสฺสเปน ภควตา อรหตา สมฺมาสมฺพุทฺเธน ธมฺมิยา กถาย สนฺทสฺสิโต สมาทปิโต สมุตฺเตชิโต สมฺปหํสิโต กสฺสปํ ภควนฺตํ อรหนฺตํ สมฺมาสมฺพุทฺธํ เอตทโวจ “อธิวาเสตุ เม ภนฺเต ภควา สฺวาตนาย ภตฺตํ สทฺธิํ ภิกฺขุสํเฆนา”ติ.\csromanspacer{} อธิ วาเสสิ โข อานนฺท กสฺสโป ภควา อรหํ สมฺมาสมฺพุทฺโธ ตุณฺหีภาเวน.\csromanspacer{} อถ โข อานนฺท กิกี กาสิราชา กสฺสปสฺส ภควโต สมฺมาสมฺพุทฺธสฺส อธิวาสนํ วิทิตฺวา อุฏฺฐายาสนา กสฺสปํ ภควนฺตํ อรหนฺตํ สมฺมาสมฺพุทฺธํ อภิวาเทตฺวา ปทกฺขิณํ กตฺวา ปกฺกามิ.\csromanspacer{} อถ โข อานนฺท กิกี กาสิราชา ตสฺสา รตฺติยา อจฺจเยน สเก นิเวสเน ปณีตํ ขาทนียํ โภชนียํ ปฏิยาทาเปตฺวา ปณฺฑุปุฏกสฺส\footnote{ปณฺฑุมุฏิกสฺส (สี, อิ), ปณฺฑุมุทิกสฺส (สฺยา, กํ)} สาลิโน วิคตกาฬกํ อเนกสูปํ อเนกพฺยญฺชนํ กสฺสปสฺส ภควโต อรหโต สมฺมาสมฺพุทฺธสฺส กาลํ อาโรจาเปสิ “กาโล ภนฺเต นิฏฺฐิตํ ภตฺตนฺ”ติ.}

\proseitem{288}{อถ โข อานนฺท กสฺสโป ภควา อรหํ สมฺมาสมฺพุทฺโธ ปุพฺพณฺหสมยํ นิวาเสตฺวา ปตฺตจีวรมาทาย เยน กิกิสฺส กาสิรญฺโญ นิเวสนํ เตนุปสงฺกมิ, อุปสงฺกมิตฺวา ปญฺญตฺเต อาสเน นิสีทิ สทฺธิํ ภิกฺขุสํเฆน.\csromanspacer{} อถ โข อานนฺท กิกี กาสิราชา พุทฺธปฺปุมุขํ ภิกฺขุสํฆํ ปณีเตน ขาทนีเยน โภชนีเยน สหตฺถา สนฺตปฺเปสิ สมฺปวาเรสิ.}

\vspace{\tipitakaparskip}
\csromancenter{ฆฏิการสุตฺต (81) อถ โข อานนฺท กิกี กาสิราชา กสฺสปํ ภควนฺตํ อรหนฺตํ สมฺมาสมฺพุทฺธํ ภุตฺตาวิํ โอนีตปตฺตปาณิํ อญฺญตรํ นีจํ อาสนํ คเหตฺวา เอกมนฺตํ นิสีทิ, เอกมนฺตํ นิสินฺโน โข อานนฺท กิกี กาสิราชา กสฺสปํ ภควนฺตํ อรหนฺตํ สมฺมาสมฺพุทฺธํ เอตทโวจ “อธิวาเสตุ เม ภนฺเต ภควา พาราณสิยํ วสฺสาวาสํ, เอวรูปํ สํฆสฺส อุปฏฺฐานํ ภวิสฺสตี”ติ.\csromanspacer{} อลํ มหาราช, อธิวุตฺโถ เม วสฺสาวาโสติ.\csromanspacer{} ทุติยมฺปิ โข อานนฺท.\csromanspacer{} ตติยมฺปิ โข อานนฺท กิกี กาสิราชา กสฺสปํ ภควนฺตํ อรหนฺตํ สมฺมาสมฺพุทฺธํ เอตทโวจ “อธิวาเสตุ เม ภนฺเต ภควา พาราณสิยํ วสฺสาวาสํ, เอวรูปํ สํฆสฺส อุปฏฺฐานํ ภวิสฺสตี”ติ.\csromanspacer{} อลํ มหาราช, อธิวุตฺโถ เม วสฺสาวาโสติ.\csromanspacer{} อถ โข อานนฺท กิกิสฺส กาสิรญฺโญ น เม กสฺสโป ภควา อรหํ สมฺมาสมฺพุทฺโธ อธิวาเสติ พาราณสิยํ วสฺสาวาสนฺติ อหุเทว อญฺญถตฺตํ, อหุ โทมนสฺสํ.\csromanspacer{} อถ โข อานนฺท กิกี กาสิราชา กสฺสปํ ภควนฺตํ อรหนฺตํ สมฺมาสมฺพุทฺธํ เอตทโวจ “อตฺถิ นุ โข ภนฺเต อญฺโญ โกจิ มยา อุปฏฺฐากตโร”ติ.}
\prose{\csromanfolio{241}อตฺถิ มหาราช เวคฬิงฺคํ นาม คามนิคโม, ตตฺถ ฆฏิกาโร นาม กุมฺภกาโร, โส เม อุปฏฺฐาโก อคฺคุปฏฺฐาโก.\csromanspacer{} ตุยฺหํ โข ปน มหาราช น เม กสฺสโป ภควา อรหํ สมฺมาสมฺพุทฺโธ อธิวาเสติ พาราณสิยํ วสฺสาวาสนฺติ อตฺเถว\footnote{อตฺถิ (สี, อิ)} อญฺญถตฺตํ, อตฺถิ โทมนสฺสํ.\csromanspacer{} ตยิทํ ฆฏิการสฺส กุมฺภการสฺส\footnote{ฆฏิกาเร กุมฺภกาเร (สี, สฺยา, กํ, อิ)} นตฺถิ จ น จ ภวิสฺสติ.\csromanspacer{} ฆฏิกาโร โข มหาราช กุมฺภกาโร พุทฺธํ สรณํ คโต ธมฺมํ สรณํ คโต สํฆํ สรณํ คโต.\csromanspacer{} ฆฏิกาโร โข มหาราช กุมฺภกาโร ปาณาติปาตา ปฏิวิรโต, อทินฺนาทานา ปฏิวิรโต, กาเมสุมิจฺฉาจารา ปฏิวิรโต, มุสาวาทา ปฏิวิรโต, สุราเมรยมชฺชปมาทฏฺฐานา ปฏิวิรโต.\csromanspacer{} ฆฏิกาโร โข มหาราช กุมฺภกาโร พุทฺเธ อเวจฺจปฺปสาเทน สมนฺนาคโต, ธมฺเม อเวจฺจปฺปสาเทน สมนฺนาคโต, สํเฆ อเวจฺจปฺปสาเทน สมนฺนาคโต, อริยกนฺเตหิ สีเลหิ สมนฺนาคโต.\csromanspacer{} ฆฏิกาโร โข มหาราช กุมฺภกาโร ทุกฺเข นิกฺกงฺโข, ทุกฺขสมุทเย นิกฺกงฺโข, ทุกฺขนิโรเธ นิกฺกงฺโข, ทุกฺขนิโรธคามินิยา ปฏิปทาย นิกฺกงฺโข.\csromanspacer{} ฆฏิกาโร โข มหาราช กุมฺภกาโร เอกภตฺติโก พฺรหฺมจารี สีลวา กลฺยาณธมฺโม.\csromanspacer{} ฆฏิกาโร โข มหาราช กุมฺภกาโร นิกฺขิตฺตมณิสุวณฺโณ \csromanfolio{242}อเปตชาตรูปรชโต.\csromanspacer{} ฆฏิกาโร โข มหาราช กุมฺภกาโร ปนฺนมุสโล, น สหตฺถา ปถวิํ ขณติ\footnote{กุมฺภกาโร น มุสเลน น สหตฺถา ปฐวิํ ขณติ (สฺยา, กํ, อิ), กุมฺภกาโร น มุสเลน สหตฺถา ปถวิญฺจ ขณติ (ก)}, ยํ โหติ กูลปลุคฺคํ วา มูสิกุกฺกโร\footnote{มูสิกุกฺกุโร (สี, สฺยา, กํ, อิ)} วา, ตํ กาเชน อาหริตฺวา ภาชนํ กริตฺวา เอวมาห “เอตฺถ โย อิจฺฉติ, ตณฺฑุลปฏิภสฺตานิ\footnote{ตณฺฑุล ปภิวตฺตานิ (สี, อิ)} วา มุคฺคปฏิภสฺตานิ วา กฬายปฏิภสฺตานิ วา นิกฺขิปิตฺวา ยํ อิจฺฉติ ตํ หรตู”ติ.\csromanspacer{} ฆฏิกาโร โข มหาราช กุมฺภกาโร อนฺเธ ชิณฺเณ มาตาปิตโร โปเสติ.\csromanspacer{} ฆฏิกาโร โข มหาราช กุมฺภกาโร ปญฺจนฺนํ โอรมฺภาคิยานํ สํโยชนานํ ปริกฺขยา โอปปาติโก ตตฺถ ปรินิพฺพายี อนาวตฺติธมฺโม ตสฺมา โลกา.}

\proseitem{289}{เอกมิทาหํ มหาราช สมยํ เวคฬิงฺเค นาม คามนิคเม วิหรามิ.\csromanspacer{} อถ ขฺวาหํ มหาราช ปุพฺพณฺหสมยํ นิวาเสตฺวา ปตฺตจีวรมาทาย เยน ฆฏิการสฺส กุมฺภการสฺส มาตาปิตโร เตนุปสงฺกมิํ, อุปสงฺกมิตฺวา ฆฏิการสฺส กุมฺภการสฺส มาตาปิตโร เอตทโวจํ “หนฺท โก นุ โข อยํ ภคฺคโว คโต”ติ.\csromanspacer{} นิกฺขนฺโต โข เต ภนฺเต อุปฏฺฐาโก, อนฺโตกุมฺภิยา โอทนํ คเหตฺวา ปริโยคา สูปํ คเหตฺวา ปริภุญฺชาติ.\csromanspacer{} อถ ขฺวาหํ มหาราช กุมฺภิยา โอทนํ คเหตฺวา ปริโยคา สูปํ คเหตฺวา ปริภุญฺชิตฺวา อุฏฺฐายาสนา ปกฺกมิํ\footnote{ปกฺกามิํ (สฺยา, กํ, อิ)}.\csromanspacer{} อถ โข มหาราช ฆฏิกาโร กุมฺภกาโร เยน มาตาปิตโร เตนุปสงฺกมิ, อุปสงฺกมิตฺวา มาตาปิตโร เอตทโวจ “โก กุมฺภิยา โอทนํ คเหตฺวา ปริโยคา สูปํ คเหตฺวา ปริภุญฺชิตฺวา อุฏฺฐายาสนา ปกฺกนฺโต”ติ.\csromanspacer{} กสฺสโป ตาต ภควา อรหํ สมฺมาสมฺภุทฺโธ กุมฺภิยา โอทนํ คเหตฺวา ปริโยคา สูปํ คเหตฺวา ปริภุญฺชิตฺวา อุฏฺฐายาสนา ปกฺกนฺโตติ.\csromanspacer{} อถ โข มหาราช ฆฏิการสฺส กุมฺภการสฺส เอตทโหสิ “ลาภา วต เม, สุลทฺธํ วต เม, ยสฺส เม กสฺสโป ภควา อรหํ สมฺมาสมฺพุทฺโธ เอวํ อภิวิสฺสตฺโถ”ติ.\csromanspacer{} อถ โข มหาราช ฆฏิการํ กุมฺภการํ อฑฺฒมาสํ ปีติสุขํ น วิชหติ\footnote{น วิชหิ (สี, สฺยา, กํ, อิ)}, สตฺตาหํ มาตาปิตูนํ.}

\proseitem{290}{เอกมิทาหํ มหาราช สมยํ ตตฺเถว เวคฬิงฺเค นาม คามนิคเม วิหรามิ.\csromanspacer{} อถ ขฺวาหํ มหาราช ปุพฺพณฺหสมยํ นิวาเสตฺวา ปตฺตจีวรมาทาย}

\vspace{\tipitakaparskip}
\csromancenter{ฆฏิการสุตฺต (81) เยน ฆฏิการสฺส กุมฺภการสฺส มาตาปิตโร เตนุปสงฺกมิํ, อุปสงฺกมิตฺวา ฆฏิการสฺส กุมฺภการสฺส มาตาปิตโร เอตทโวจํ “หนฺท โก นุ โข อยํ ภคฺคโว คโต”ติ.\csromanspacer{} นิกฺขนฺโต โข เต ภนฺเต อุปฏฺฐาโก, อนฺโต กโฬปิยา กุมฺมาสํ คเหตฺวา ปริโยคา สูปํ คเหตฺวา ปริภุญฺชาติ.\csromanspacer{} อถ ขฺวาหํ มหาราช กโฬปิยา กุมฺมาสํ คเหตฺวา ปริโยคา สูปํ คเหตฺวา ปริภุญฺชิตฺวา อุฏฺฐายาสนา ปกฺกมิํ.\csromanspacer{} อถ โข มหาราช ฆฏิกาโร กุมฺภกาโร เยน มาตาปิตโร เตนุปสงฺกมิ, อุปสงฺกมิตฺวา มาตาปิตโร เอตทโวจ “โก กโฬปิยา กุมฺมาสํ คเหตฺวา ปริโยคา สูปํ คเหตฺวา ปริภุญฺชิตฺวา อุฏฺฐายาสนา ปกฺกนฺโต”ติ.\csromanspacer{} กสฺสโป ตาต ภควา อรหํ สมฺมาสมฺพุทฺโธ กโฬปิยา กุมฺมาสํ คเหตฺวา ปริโยคา สูปํ คเหตฺวา ปริภุญฺชิตฺวา อุฏฺฐายาสนา ปกฺกนฺโตติ.\csromanspacer{} อถ โข มหาราช ฆฏิการสฺส กุมฺภการสฺส เอตทโหสิ “ลาภา วต เม, สุลทฺธํ วต เม, ยสฺส เม กสฺสโป ภควา อรหํ สมฺมาสมฺพุทฺโธ เอวํ อภิวิสฺสตฺโถ”ติ.\csromanspacer{} อถ โข มหาราช ฆฏิการํ กุมฺภการํ อฑฺฒุมาสํ ปีติสุขํ น วิชหติ, สตฺตาหํ มาตาปิตูนํ.}
\proseitem{291}{\csromanfolio{243}เอกมิทาหํ มหาราช สมยํ ตตฺเถว เวคฬิงฺเค นาม คามนิคเม วิหรามิ.\csromanspacer{} เตน โข ปน สมเยน กุฏิ\footnote{คนฺธกุฏิ (สี)} โอวสฺสติ.\csromanspacer{} อถ ขฺวาหํ มหาราช ภิกฺขู อามนฺเตสิํ “คจฺฉถ ภิกฺขเว ฆฏิการสฺส กุมฺภการสฺส นิเวสเน ติณํ ชานาถา”ติ.\csromanspacer{} เอวํ วุตฺเต มหาราช เต ภิกฺขู มํ เอตทโวจุํ “นตฺถิ โข ภนฺเต ฆฏิการสฺส กุมฺภการสฺส นิเวสเน ติณํ, อตฺถิ จ ขฺวาสฺส อาเวสเน\footnote{อาเวสนํ (สี, สฺยา, กํ, อิ)} ติณจฺฉทน3นฺ”ติ.\csromanspacer{} คจฺฉถ ภิกฺขเว ฆฏิการสฺส กุมฺภการสฺส อาเวสนํ อุตฺติณํ กโรถาติ.\csromanspacer{} อถ โข เต มหาราช ภิกฺขู ฆฏิการสฺส กุมฺภการสฺส อาเวสนํ อุตฺติณมกํสุ.\csromanspacer{} อถ โข มหาราช ฆฏิการสฺส กุมฺภการสฺส มาตาปิตโร เต ภิกฺขู เอตทโวจุํ “เก อาสเวสนํ อุตฺติณํ กโรนฺตี”ติ.\csromanspacer{} ภิกฺขู ภคินิ, กสฺสปสฺส ภควโต อรหโต สมฺมาสมฺพุทฺธสฺส กุฏิ โอวสฺสตีติ.\csromanspacer{} หรถ ภนฺเต หรถ ภทฺรมุขาติ.\csromanspacer{} อถ โข มหาราช ฆฏิกาโร กุมฺภกาโร เยน มาตาปิตโร เตนุปสงฺกมิ, อุปสงฺกมิตฺวา มาตาปิตโร เอตทโวจ “เก อาเวสนํ อุตฺติณมกํสู”ติ.\csromanspacer{} ภิกฺขู ตาต กสฺสปสฺส กิร \csromanfolio{244}ภควโต อรหโต สมฺมาสมฺพุทฺธสฺส กุฏิ โอวสฺสตีติ.\csromanspacer{} อถ โข มหาราช ฆฏิการสฺส กุมฺภการสฺส เอตทโหสิ “ลาภา วต เม, สุลทฺธํ วต เม, ยสฺส เม กสฺสโป ภควา อรหํ สมฺมาสมฺพุทฺโธ เอวํ อภิวิสฺสตฺโถ”ติ.\csromanspacer{} อถ โข มหาราช ฆฏิการํ กุมฺภการํ อฑฺฒมาสํ ปีติสุขํ น วิชหติ, สตฺตาหํ มาตาปิตูนํ.\csromanspacer{} อถ โข มหาราช อาเวสนํ สพฺพนฺตํ เตมาสํ อากาสจฺฉทนํ อฏฺฐาสิ, น เทโวติวสฺสิ\footnote{น จาติวสฺสิ (สี, สฺยา, กํ, อิ)}, เอวรูโป จ มหาราช ฆฏิกาโร กุมฺภกาโรติ.\csromanspacer{} ลาภา ภนฺเต ฆฏิการสฺส กุมฺภการสฺส, สุลทฺธา ภนฺเต ฆฏิการสฺส กุมฺภการสฺส, ยสฺส ภควา เอวํ อภิวิสฺสตฺโถติ.}

\proseitem{292}{อถ โข อานนฺท กิกี กาสิราชา ฆฏิการสฺส กุมฺภการสฺส ปญฺจมตฺตานิ ตณฺฑุลวาหสตานิ ปาเหสิ ปณฺฑุปุฏกสฺส สาลิโน ตทุปิยญฺจ สูเปยฺยํ.\csromanspacer{} อถ โข เต อานนฺท ราชปุริสา ฆฏิการํ กุมฺภการํ อุปสงฺกมิตฺวา เอตทโวจุํ “อิมานิ โข ภนฺเต ปญฺจมตฺตานิ ตณฺฑุลวาหสตานิ กิกินา กาสิราเชน ปหิตานิ ปณฺฑุปุฏกสฺส สาลิโน ตทุปิยญฺจ สูเปยฺยํ, ตานิ ภนฺเต ปฏิคฺคณฺหถา”ติ\footnote{ปติคฺคณฺหาตูติ (สี, อิ), ปฏิคฺคณฺหาตูติ (สฺยา, กํ)}.\csromanspacer{} ราชา โข พหุกิจฺโจ พหุกรณีโย, อลํ เม รญฺโญว โหตูติ.\csromanspacer{} สิยา โข ปน เต อานนฺท เอวมสฺส “อญฺโญ นูน เตน สมเยน โชติปาโล มาณโว อโหสี”ติ.\csromanspacer{} น โข ปเนตํ อานนฺท เอวํ ทฏฺฐพฺพํ, อหํ เตน สมเยน โชติปาโล มาณโว อโหสินฺติ.}

\prose{อิทมโวจ ภควา.\csromanspacer{} อตฺตมโน อายสฺมา อานนฺโท ภควโต ภาสิตํ อภินนฺทีติ.}

\csromanheader{2}{ฆฏิการสุตฺตํ นิฏฺฐํตํ ปฐมํ.}
\csromansectionrule
\csromanmatikaanchor{mtk.39}
\csromantocmarkref{subsection}{2. รฏฺฐปาลสุตฺต}{mtk.39}
\bookmark[dest={mtk.39},level=2]{2. รฏฺฐปาลสุตฺต}
\csromanheader{2}{2. รฏฺฐปาลสุตฺต}
\proseitem{293}{เอวํ เม สุตํ\csromandash{}เอกํ สมยํ ภควา กุรูสุ จาริกํ จรมาโน มหตา ภิกฺขุสํเฆน สทฺธิํ เยน ถุลฺลโกฏฺฐิกํ\footnote{ถูลโกฏฺฐิตํ (สี, สฺยา, กํ, อิ)} นาม กุรูนํ นิคโม ตทวสริ.\csromanspacer{} อสฺโสสุํ โข ถุลฺลโกฏฺฐิกา\footnote{ถูลโกฏฺฐิตกา (สี, สฺยา, กํ, อิ)} พฺราหฺมณคหปติกา}

\vspace{\tipitakaparskip}
\csromancenter{รฏฺฐปาลสุตฺต (82) “สมโณ ขลุ โภ โคตโม สกฺยปุตฺโต สกฺยกุลา ปพฺพชิโต กุรูสุ จาริกํ จรมาโน มหตา ภิกฺขุสํเฆน สทฺธิํ ถุลฺลโกฏฺฐิกํ อนุปฺปตฺโต, ตํ โข ปน ภวนฺตํ โคตมํ เอวํ กลฺยาโณ กิตฺติสทฺโท อพฺภุคฺคโต ‘อิติปิ โส ภควา อรหํ สมฺมาสมฺพุทฺโธ วิชฺชาจรณสมฺปนฺโน สุคโต โลกวิทู อนุตฺตโร ปุริสทมฺมสารถิ สตฺถา เทวมนุสฺสานํ พุทฺโธ ภควา’ติ, โส อิมํ โลกํ สเทวกํ สมารกํ สพฺรหฺมกํ สสฺสมณพฺราหฺมณิํ ปชํ สเทวมนุสฺสํ สยํ อภิญฺญา สจฺฉิกตฺวา ปเวเทติ, โส ธมฺมํ เทเสติ อาทิกลฺยาณํ มชฺเฌกลฺยาณํ ปริโยสานกลฺยาณํ สาตฺถํ สพฺยญฺชนํ เกวลปริปุณฺณํ ปริสุทฺธํ พฺรหฺมจริยํ ปกาเสติ, สาธุ โข ปน ตถารูปานํ อรหตํ ทสฺสนํ โหตี”ติ.\csromanspacer{} อถ โข ถุลฺลโกฏฺฐิกา พฺราหฺมณคหปติกา เยน ภควา เตนุปสงฺกมิํสุ, อุปสงฺกมิตฺวา อปฺเปกจฺเจ ภควนฺตํ อภิวาเทตฺวา เอกมนฺตํ นิสีทิํสุ, อปฺเปกจฺเจ ภควตา สทฺธิํ สมฺโมทิํสุ, สมฺโมทนียํ กถํ สารณียํ วีติสาเรตฺวา เอกมนฺตํ นิสีทิํสุ, อปฺเปกจฺเจ เยน ภควา เตนญฺชลิํ ปณาเมตฺวา เอกมนฺตํ นิสีทิํสุ, อปฺเปกจฺเจ ภควโต สนฺติเก นามโคตฺตํ สาเวตฺวา เอกมนฺตํ นิสีทิํสุ, อปฺเปกจฺเจ ตุณฺหีภูตา เอกมนฺตํ นิสีทิํสุ.\csromanspacer{} เอกมนฺตํ นิสินฺเน โข ถุลฺลโกฏฺฐิเก พฺราหฺมณคหปติเก ภควา ธมฺมิยา กถาย สนฺทสฺเสสิ สมาทเปสิ สมุตฺเตเชสิ สมฺปหํเสสิ.}
\proseitem{294}{\csromanfolio{245}เตน โข ปน สมเยน รฏฺฐปาโล นาม กุลปุตฺโต ตสฺมิํเยว ถุลฺลโกฏฺฐิเก อคฺคกุลสฺส\footnote{อคฺคกุลิกสฺส (สี, สฺยา, กํ, อิ)} ปุตฺโต ติสฺสํ ปริสายํ นิสินฺโน โหติ.\csromanspacer{} อถ โข รฏฺฐปาลสฺส กุลปุตฺตสฺส เอตทโหสิ “ยถา ยถา ขฺวาหํ ภควตา ธมฺมํ เทสิตํ อาชานามิ\footnote{ยถา ยถา โข ภควา ธมฺมํ เทเสติ (สี)}, นยิทํ สุกรํ อคารํ อชฺฌาวสตา เอกนฺตปริปุณฺณํ เอกนฺตปริสุทฺธํ สงฺขลิขิตํ พฺรหฺมจริยํ จริตุํ, ยํนูนาหํ เกสมสฺสุํ โอหาเรตฺวา กาสายานิ วตฺถานิ อจฺฉาเทตฺวา อคารสฺมา อนคาริยํ ปพฺพเชยฺยนฺ”ติ.\csromanspacer{} อถ โข ถุลฺลโกฏฺฐิกา พฺราหฺมณคหปติกา ภควตา ธมฺมิยา กถาย สนฺทสฺสิตา สมาทปิตา สมุตฺเตชิตา สมฺปหํสิตา ภควโต ภาสิตํ อภินนฺทิตฺวา อนุโมทิตฺวา อุฏฺฐายาสนา ภควนฺตํ อภิวาเทตฺวา ปทกฺขิณํ กตฺวา ปกฺกมิํสุ.\csromanspacer{} อถ โข รฏฺฐปาโล กุลปุตฺโต อจิรปกฺกนฺเตสุ \csromanfolio{246}ถุลฺลโกฏฺฐิเกสุ พฺราหฺมณคหปติเกสุ เยน ภควา เตนุปสงฺกมิ, อุปสงฺกมิตฺวา ภควนฺตํ อภิวาเทตฺวา เอกมนฺตํ นิสีทิ, เอกมนฺตํ นิสินฺโน โข รฏฺฐปาโล กุลปุตฺโต ภควนฺตํ เอตทโวจ “ยถา ยถาหํ ภนฺเต ภควตา ธมฺมํ เทสิตํ อาชานามิ, นยิทํ สุกรํ อคารํ อชฺฌาวสตา เอกนฺตปริปุณฺณํ เอกนฺตปริสุทฺธํ สงฺขลิขิตํ พฺรหฺมจริยํ จริตุํ, อิจฺฉามหํ ภนฺเต เกสมสฺสุํ โอหาเรตฺวา กาสายานิ วตฺถานิ อจฺฉาเทตฺวา อคารสฺมา อนคาริยํ ปพฺพชิตุํ, ลเภยฺยาหํ ภนฺเต ภควโต สนฺติเก ปพฺพชฺชํ, ลเภยฺยํ อุปสมฺปทํ, ปพฺพาเชตุ มํ ภควา”ติ\footnote{เอตฺถ “ลเภยฺยาหํ ฯเปฯ อุปสมฺปทํ”ติ วากฺยทฺวยํ สพฺเพสุปิ มูลโปตฺถเกสุ ทิสฺสติ, ปาราชิกปาฬิยํ ปน สุทินฺนภาณวาเร เอตํ นตฺถิ. “ปพฺพาเชตุ มํ ภควา”ติ อิทํ ปน วากฺยํ มรมฺมโปตฺเถเก เยว ทิสฺสติ, ปาราชิกปาฬิยญฺจ ตเทว อตฺถิ.}.\csromanspacer{} อนุญฺญาโตสิ ปน ตฺวํ รฏฺฐปาล มาตาปิตูหิ อคารสฺมา อนคาริยํ ปพฺพชฺชายาติ.\csromanspacer{} น โขหํ ภนฺเต อนุญฺญาโต มาตาปิตูหิ อคารสฺมา อนคาริยํ ปพฺพชฺชายาติ.\csromanspacer{} น โข รฏฺฐปาล ตถาคตา อนนุญฺญาตํ มาตาปิตูหิ ปุตฺตํ ปพฺพาเชนฺตีติ.\csromanspacer{} สฺวาหํ ภนฺเต ตถา กริสฺสามิ, ยถา มํ มาตาปิตโร อนุชานิสฺสนฺติ อคารสฺมา อนคาริยํ ปพฺพชฺชายาติ.}

\proseitem{295}{อถ โข รฏฺฐปาโล กุลปุตฺโต อุฏฺฐายาสนา ภควนฺตํ อภิวาเทตฺวา ปทกฺขิณํ กตฺวา เยน มาตาปิตโร เตนุปสงฺกมิ, อุปสงฺกมิตฺวา มาตาปิตโร เอตทโวจ “อมฺมตาตา ยถา ยถาหํ ภควตา ธมฺมํ เทสิตํ อาชานามิ, นยิทํ สุกรํ อคารํ อชฺฌาวสตา เอกนฺตปริปุณฺณํ เอกนฺตปริสุทฺธํ สงฺขลิขิตํ พฺรหฺมจริยํ จริตุํ, อิจฺฉามหํ เกสมสฺสุํ โอหาเรตฺวา กาสายานิ วตฺถานิ อจฺฉาเทตฺวา อคารสฺมา อนคาริยํ ปพฺพชิตุํ, อนุชานาถ มํ อคารสฺมา อนคาริยํ ปพฺพชฺชายา”ติ.\csromanspacer{} เอวํ วุตฺเต รฏฺฐปาลสฺส กุลปุตฺตสฺส มาตาปิตโร รฏฺฐปาลํ กุลปุตฺตํ เอตทโวจุํ “ตฺวํ โขสิ ตาต รฏฺฐปาล อมฺหากํ เอกปุตฺตโก ปิโย มนาโป สุเขธิโต สุขปริภโต\footnote{สุขปริหโต (สฺยา, กํ, ก) (เอหิ ตฺวํ ตาต รฏฺฐปาล ภุญฺช จ ปิว จ ปริจาเร หิ จ, ภุญฺชนฺโต ปิวนฺโต ปริจาเรนฺโต กาเม ปริภุญฺชนฺโต ปุญฺญานิ กโรนฺโต อภิรมสฺสุ, น ตํ มยํ อนุชานาม อคารสฺมา อนคาริยํ ปพฺพชฺชาย,) สพฺพตฺถ ทิสฺสติ, สุทินฺนกณฺเฑ ปน นตฺถิ, อฏฺฐกถาสุปิ น ทสฺสิตํ.}, น ตฺวํ ตาต รฏฺฐปาล กสฺสจิ ทุกฺขสฺส ชานาสิ, ( ) มรเณนปิ เต มยํ อกามกา วินา ภวิสฺสาม, กิํ ปน}

\vspace{\tipitakaparskip}
\csromancenter{รฏฺฐปาลสุตฺต (82) มยํ ตํ ชีวนฺตํ อนุชานิสฺสาม อคารสฺมา อนคาริยํ ปพฺพชฺชายา”ติ.\csromanspacer{} ทุติยมฺปิ โข รฏฺฐปาโล กุลปุตฺโต ฯเปฯ.\csromanspacer{} ตติยมฺปิ โข รฏฺฐปาโล กุลปุตฺโต มาตาปิตโร เอตทโวจ “อมฺมตาตา ยถา ยถาหํ ภควตา ธมฺมํ เทสิตํ อาชานามิ, นยิทํ สุกรํ อคารํ อชฺฌาวสตา เอกนฺตปริปุณฺณํ เอกนฺตปริสุทฺธํ สงฺขลิขิตํ พฺรหฺมจริยํ จริตุํ, อิจฺฉามหํ เกสมสฺสุํ โอหาเรตฺวา กาสายานิ วตฺถานิ อจฺฉาเทตฺวา อคารสฺมา อนคาริยํ ปพฺพชิตุํ, อนุชานาถ มํ อคารสฺมา อนคาริยํ ปพฺพชฺชายา”ติ.\csromanspacer{} ตติยมฺปิ โข รฏฺฐปาลสฺส กุลปุตฺตสฺส มาตาปิตโร รฏฺฐปาลํ กุลปุตฺตํ เอตทโวจุํ “ตฺวํ โขสิ ตาต รฏฺฐปาล อมฺหากํ เอกปุตฺตโก ปิโย มนาโป สุเขธิโต สุขปริภโต, น ตฺวํ ตาต รฏฺฐปาล กสฺสจิ ทุกฺขสฺส ชานาสิ, มรเณนปิ เต มยํ อกามกา วินา ภวิสฺสาม, กิํ ปน มยํ ตํ ชีวนฺตํ อนุชานิสฺสาม อคารสฺมา อนคาริยํ ปพฺพชฺชายา”ติ.}
\proseitem{296}{\csromanfolio{247}อถ โข รฏฺฐปาโล กุลปุตฺโต “น มํ มาตาปิตโร อนุชานนฺติ อคารสฺมา อนคาริยํ ปพฺพชฺชายา”ติ ตตฺเถว อนนฺตรหิตาย ภูมิยา นิปชฺชิ “อิเธว เม มรณํ ภวิสฺสติ ปพฺพชฺชา วา”ติ.\csromanspacer{} อถ โข รฏฺฐปาโล กุลปุตฺโต เอกํปิ ภตฺตํ น ภุญฺชิ, เทฺวปิ ภตฺตานิ น ภุญฺชิ, ตีณิปิ ภตฺตานิ น ภุญฺชิ, จตฺตาริปิ ภตฺตานิ น ภุญฺชิ, ปญฺจปิ ภตฺตานิ น ภุญฺชิ, ฉปิ ภตฺตานิ น ภุญฺชิ, สตฺตปิ ภุตฺตานิ น ภุญฺชิ.\csromanspacer{} อถ โข รฏฺฐปาลสฺส กุลปุตฺตสฺส มาตาปิตโร รฏฺฐปาลํ กุลปุตฺตํ เอตทโวจุํ “ตฺวํ โขสิ ตาต รฏฺฐปาล อมฺหากํ เอกปุตฺตโก ปิโย มนาโป สุเขธิโต สุขปริภโต, น ตฺวํ ตาต รฏฺฐปาล กสฺสจิ ทุกฺขสฺส ชานาสิ, x มรเณนปิ เต มยํ อกามกา วินา ภวิสฺสาม, กิํ ปน มยํ ตํ ชีวนฺตํ อนุชานิสฺสาม อคารสฺมา อนคาริยํ ปพฺพชฺชาย, อุฏฺเฐหิ ตาต รฏฺฐปาล ภุญฺช จ ปิว จ ปริจาเรหิ จ, ภุญฺชนฺโต ปิวนฺโต ปริจาเรนฺโต กาเม ปริภุญฺชนฺโต ปุญฺญานิ กโรนฺโต อภิรมสฺสุ, น ตํ มยํ อนุชานาม อคารสฺมา อนคาริยํ ปพฺพชฺชาย, + มรเณนปิ เต มยํ อกามยา วินา ภวิสฺสาม, กิํ ปน มยํ ตํ ชีวนฺตํ อนุชานิสฺสาม อคารสฺมา อนคาริยํ ปพฺพชฺชายา”ติ.\csromanspacer{} เอวํ วุตฺเต รฏฺฐปาโล กุลปุตฺโต ตุณฺหี อโหสิ.\csromanspacer{} ทุติยมฺปิ โข รฏฺฐปาลสฺส กุลปุตฺตสฺส \csromanfolio{248}มาตาปิตโร รฏฺฐปาลํ กุลปุตฺตํ เอตทโวจุํ ฯเปฯ.\csromanspacer{} ทุติยมฺปิ โข รฏฺฐปาโล กุลปุตฺโต ตุณฺหี อโหสิ.\csromanspacer{} ตติยมฺปิ โข รฏฺฐปาลสฺส กุลปุตฺตสฺส มาตาปิตโร รฏฺฐปาลํ กุลปุตฺตํ เอตทโวจุํ “ตฺวํ โขสิ ตาต รฏฺฐปาล อมฺหากํ เอกปุตฺตโก ปิโย มนาโป สุเขธิโต สุขปริภโต, น ตฺวํ ตาต รฏฺฐปาล กสฺสจิ ทุกฺขสฺส ชานาสิ, มรเณนปิ เต มยํ อกามกา วินา ภวิสฺสาม, กิํ ปน มยํ ตํ ชีวนฺตํ อนุชานิสฺสาม อคารสฺมา อนคาริยํ ปพฺพชฺชาย, อุฏฺเฐหิ ตาต รฏฺฐปาล ภุญฺช จ ปิว จ ปริจาเรหิ จ, ภุญฺชนฺโต ปิวนฺโต ปริจาเรนฺโต กาเม ปริภุญฺชนฺโต ปุญฺญานิ กโรนฺโต อภิรมสฺสุ, น ตํ มยํ อนุชานาม อคารสฺมา อนคาริยํ ปพฺพชฺชาย, มรเณนปิ เต มยํ อกามกา วินา ภวิสฺสาม, กิํ ปน มยํ ตํ ชีวนฺตํ อนุชานิสฺสาม อคารสฺมา อนคาริยํ ปพฺพชฺชายา”ติ.\csromanspacer{} ตติยมฺปิ โข รฏฺฐปาโล กุลปุตฺโต ตุณฺหี อโหสิ.}

\proseitem{297}{อถ โข รฏฺฐปาลสฺส กุลปุตฺตสฺส สหายกา เยน รฏฺฐปาโล กุลปุตฺโต เตนุปสงฺกมิํสุ, อุปสงฺกมิตฺวา รฏฺฐปาลํ กุลปุตฺตํ เอตทโวจุํ “ตฺวํ โขสิ\footnote{ตฺวํ โข (สี, อิ)} สมฺม รฏฺฐปาล มาตาปิตูนํ เอกปุตฺตโก ปิโย มนาโป สุเขธิโต สุขปริภโต, น ตฺวํ สมฺม รฏฺฐปาล กสฺสจิ ทุกฺขสฺส ชานาสิ, มรเณนปิ เต มาตาปิตโร อกามกา วินา ภวิสฺสนฺติ, กิํ ปน เต ตํ ชีวนฺตํ อนุชานิสฺสนฺติ อคารสฺมา อนคาริยํ ปพฺพชฺชาย, อุฏฺเฐหิ สมฺม รฏฺฐปาล ภุญฺช จ ปิว จ ปริจาเรหิ จ, ภุญฺชนฺโต ปิวนฺโต ปริจาเรนฺโต กาเม ปริภุญฺชนฺโต ปุญฺญานิ กโรนฺโต อภิรมสฺสุ, น ตํ มาตาปิตโร อนุชานิสฺสนฺติ\footnote{อนุชานนฺติ (สี, สฺยา, กํ, อิ)} อคารสฺมา อนคาริยํ ปพฺพชฺชาย, มรเณนปิ เต มาตาปิตโร อกามกา วินา ภวิสฺสนฺติ, กิํ ปน เต ตํ ชีวนฺตํ อนุชานิสฺสนฺติ อคารสฺมา อนคาริยํ ปพฺพชฺชายา”ติ.\csromanspacer{} เอวํ วุตฺเต รฏฺฐปาโล กุลปุตฺโต ตุณฺหี อโหสิ.\csromanspacer{} ทุติยมฺปิ โข.\csromanspacer{} ตติยมฺปิ โข รฏฺฐปาลสฺส กุลปุตฺตสฺส สหายกา รฏฺฐปาลํ กุลปุตฺตํ เอตทโวจุํ “ตฺวํ โขสิ สมฺม รฏฺฐปาล มาตาปิตูนํ เอกปุตฺตโก ปิโย มนาโป สุเขธิโต สุขปริภโต, น ตฺวํ สมฺม รฏฺฐปาล กสฺสจิ ทุกฺขสฺส ชานาสิ, มรเณนปิ เต มาตาปิตโร อกามกา วินา ภวิสฺสนฺติ, กิํ ปน เต ตํ ชีวนฺตํ อนุชานิสฺสนฺติ อคารสฺมา อนคาริยํ ปพฺพชฺชาย, อุฏฺเฐหิ สมฺม รฏฺฐปาล ภุญฺช จ ปิว จ ปริจาเรหิ จ, ภุญฺชนฺโต ปิวนฺโต ปริจาเรนฺโต กาเม ปริภุญฺชนฺโต}

\vspace{\tipitakaparskip}
\csromancenter{รฏฺฐปาลสุตฺต (82) ปุญฺญานิ กโรนฺโต อภิรมสฺสุ, น ตํ มาตาปิตโร อนุชานิสฺสนฺติ อคารสฺมา อนคาริยํ ปพฺพชฺชาย, มรเณนปิ เต มาตาปิตโร อกามกา วินา ภวิสฺสนฺติ, กิํ ปน เต ตํ ชีวนฺตํ อนุชานิสฺสนฺติ อคารสฺมา อนคาริยํ ปพฺพชฺชายา”ติ.\csromanspacer{} ตติยมฺปิ โข รฏฺฐปาโล กุลปุตฺโต ตุณฺหี อโหสิ.}
\proseitem{298}{\csromanfolio{249}อถ โข รฏฺฐปาลสฺส กุลปุตฺตสฺส สหายกา เยน รฏฺฐปาลสฺส กุลปุตฺตสฺส มาตาปิตโร เตนุปสงฺกมิํสุ, อุปสงฺกมิตฺวา รฏฺฐปาลสฺส กุลปุตฺตสฺส มาตาปิตโร เอตทโวจุํ “อมฺมตาตา เอโส รฏฺฐปาโล กุลปุตฺโต ตตฺเถว อนนฺตรหิตาย ภูมิยา นิปนฺโน ‘อิเธว เม มรณํ ภวิสฺสติ ปพฺพชฺชา วา’ติ.\csromanspacer{} สเจ ตุเมฺห รฏฺฐปาลํ กุลปุตฺตํ นานุชานิสฺสถ อคารสฺมา อนคาริยํ ปพฺพชฺชาย, ตตฺเถว\footnote{ตตฺเถวสฺส (สี)} มรณํ อาคมิสฺสติ, สเจ ปน ตุเมฺห รฏฺฐปาลํ กุลปุตฺตํ อนุชานิสฺสถ อคารสฺมา อนคาริยํ ปพฺพชฺชาย, ปพฺพชิตมฺปิ นํ ทกฺขิสฺสถ.\csromanspacer{} สเจ รฏฺฐปาโล กุลปุตฺโต นาภิรมิสฺสติ อคารสฺมา อนคาริยํ ปพฺพชฺชาย, กา ตสฺส\footnote{กา จสฺส (สี)} อญฺญา คติ ภวิสฺสติ, อิเธว ปจฺจาคมิสฺสติ, อนุชานาถ รฏฺฐปาลํ กุลปุตฺตํ อคารสฺมา อนคาริยํ ปพฺพชฺชายา”ติ.\csromanspacer{} อนุชานาม ตาตา รฏฺฐปาลํ กุลปุตฺตํ อคารสฺมา อนคาริยํ ปพฺพชฺชาย, ปพฺพชิเตน จ ปน\footnote{ปน เต (สฺยา, กํ, ก)} มาตาปิตโร อุทฺทสฺเสตพฺพาติ.\csromanspacer{} อถ โข รฏฺฐปาลสฺส กุลปุตฺตสฺส สหายกา เยน รฏฺฐปาโล กุลปุตฺโต เตนุปสงฺกมิํสุ, อุปสงฺกมิตฺวา รฏฺฐปาลํ กุลปุตฺตํ เอตทโวจุํ “อุฏฺเฐหิ สมฺม รฏฺฐปาล\footnote{“ตฺวํ โขสิ สมฺม รฏฺฐปาล มาตาปิตูนํ เอกปุตฺตโก ปิโย มนาโป สุเขธิโต สุขปริหโต, น ตฺวํ สมฺม รฏฺฐปาล กสฺสจิ ทุกฺขสฺส ชานาสิ, อุฏฺเฐหิ สมฺม รฏฺฐปาล ภุญฺช จ ปิว จ ปริจาเรหิ จ, ภุญฺชนฺโต ปิวนฺโต ปริจาเรนฺโต กาเม ปริภุญฺชนฺโต ปุญฺญานิ กโรนฺโต อภิรมสฺสุ (สี, อิ, ก)}, อนุญฺญาโตสิ มาตาปิตูหิ อคารสฺมา อนคาริยํ ปพฺพชฺชาย, ปพฺพชิเตน จ ปน เต มาตาปิตโร อุทฺทสฺเสตพฺพา”ติ.}

\proseitem{299}{อถ โข รฏฺฐปาโล กุลปุตฺโต อุฏฺฐหิตฺวา พลํ คาเหตฺวา เยน ภควา เตนุปสงฺกมิ, อุปสงฺกมิตฺวา ภควนฺตํ อภิวาเทตฺวา เอกมนฺตํ นิสีทิ, เอกมนฺตํ นิสินฺโน โข รฏฺฐปาโล กุลปุตฺโต ภควนฺตํ เอตทโวจ “อนุญฺญาโต อหํ ภนฺเต มาตาปิตูหิ อคารสฺมา อนคาริยํ ปพฺพชฺชาย, ปพฺพาเชตุ มํ ภควา”ติ.\csromanspacer{} อลตฺถ โข รฏฺฐปาโล กุลปุตฺโต \csromanfolio{250}ภควโต สนฺติเก ปพฺพชฺชํ, อลตฺถ อุปสมฺปทํ.\csromanspacer{} อถ โข ภควา อจิรูปสมฺปนฺเน อายสฺมนฺเต รฏฺฐปาเล อฑฺฒมาสูปสมฺปนฺเน ถุลฺลโกฏฺฐิเก ยถาภิรนฺตํ วิหรตฺวา เยน สาวตฺถิ เตน จาริกํ ปกฺกามิ, อนุปุพฺเพน จาริกํ จรมาโน เยน สาวตฺถิ ตทวสริ, ตตฺร สุทํ ภควา สาวตฺถิยํ วิหรติ เชตวเน อนาถปิณฺฑิกสฺส อาราเม.\csromanspacer{} อถ โข อายสฺมา รฏฺฐปาโล เอโก วูปกฏฺโฐ อปฺปมตฺโต อาตาปี ปหิตตฺโต วิหรนฺโต นจิรสฺเสว, ยสฺสตฺถาย กุลปุตฺตา สมฺมเทว อคารสฺมา อนคาริยํ ปพฺพชนฺติ.\csromanspacer{} ตทนุตฺตรํ พฺรหฺมจริยปริโยสานํ ทิฏฺเฐว ธมฺเม สยํ อภิญฺญา สจฺฉิกตฺวา อุปสมฺปชฺช วิหาสิ, “ขีณา ชาติ, วุสิตํ พฺรหฺมจริยํ, กตํ กรณียํ, นาปรํ อิตฺถตฺตายา”ติ อพฺภญฺญาสิ, อญฺญตโร โข ปนายสฺมา รฏฺฐปาโล อรหตํ อโหสิ.}

\prose{อถ โข อายสฺมา รฏฺฐปาโล เยน ภควา เตนุปสงฺกมิ, อุปสงฺกมิตฺวา ภควนฺตํ อภิวาเทตฺวา เอกมนฺตํ นิสีทิ, เอกมนฺตํ นิสินฺโน โข อายสฺมา รฏฺฐปาโล ภควนฺตํ เอตทโวจ “อิจฺฉามหํ ภนฺเต มาตาปิตโร อุทฺทสฺเสตุํ, สเจ มํ ภควา อนุชานาตี”ติ.\csromanspacer{} อถ โข ภควา อายสฺมโต รฏฺฐปาลสฺส เจตสฺสา เจโต ปริจฺจ\footnote{เจโตปริวิตกฺกํ (สี, อิ)} มนสากาสิ, ยถา\footnote{ยทา (สี, อิ)} ภควา อญฺญาสิ “อภพฺโพ โข รฏฺฐปาโล กุลปุตฺโต สิกฺขํ ปจฺจกฺขาย หีนายาวตฺติตุนฺ”ติ.\csromanspacer{} อถ โข ภควา อายสฺมนฺตํ รฏฺฐปาลํ เอตทโวจ “ยสฺสทานิ ตฺวํ รฏฺฐปาล กาลํ มญฺญสี”ติ.\csromanspacer{} อถ โข อายสฺมา รฏฺฐปาโล อุฏฺฐายาสนา ภควนฺตํ อภิวาเทตฺวา ปทกฺขิณํ กตฺวา เสนาสนํ สํสาเมตฺวา ปตฺตจีวรมาทาย เยน ถุลฺลโกฏฺฐิกํ เตน จาริกํ ปกฺกามิ, อนุปุพฺเพน จาริกํ จรมาโน เยน ถุลฺลโกฏฺฐิโก ตทวสริ, ตตฺร สุทํ อายสฺมา รฏฺฐปาโล ถุลฺลโกฏฺฐิเก วิหรติ รญฺโญ โกรพฺยสฺส มิคจีเร.\csromanspacer{} อถ โข อายสฺมา รฏฺฐปาโล ปุพฺพณฺหสมยํ นิวาเสตฺวา ปตฺตจีวรมาทาย ถุลฺลโกฏฺฐิกํ ปิณฺฑาย ปาวิสิ, ถุลฺลโกฏฺฐิเก สปทานํ ปิณฺฑาย จรมาโน เยน สกปิตุ นิเวสนํ เตนุปสงฺกมิ.\csromanspacer{} เตน โข ปน สมเยน อายสฺมโต รฏฺฐปาลสฺส ปิตา มชฺฌิมาย ทฺวารสาลาย อุลฺลิขาเปติ.\csromanspacer{} อทฺทสา โข อายสฺมโต รฏฺฐปาลสฺส ปิตา อายสฺมนฺตํ รฏฺฐปาลํ ทูรโตว อาคจฺฉนฺตํ, ทิสฺวาน เอตทโวจ “อิเมหิ มุณฺฑเกหิ มสณเกหิ อมฺหากํ เอกปุตฺตโก ปิโย มนาโป ปพฺพาชิโต”ติ.\csromanspacer{} อถ โข อายสฺมา}

\vspace{\tipitakaparskip}
\csromancenter{รฏฺฐปาลสุตฺต (82) รฏฺฐปาโล สกปิตุ นิเวสเน เนว ทานํ อลตฺถ, น ปจฺจกฺขานํ, อญฺญทตฺถุ อกฺโกสเมว อลตฺถ.\csromanspacer{} เตน โข ปน สมเยน อายสฺมโต รฏฺฐปาลสฺส ญาติทาสี อาภิโทสิกํ กุมฺมาสํ ฉฑฺเฑตุกามา โหติ.\csromanspacer{} อถ โข อายสฺมา รฏฺฐปาโล ตํ ญาติทาสิํ เอตทโวจ “สเจตํ ภคินิ ฉฑฺฑนียธมฺมํ, อิธ เม ปตฺเต อากิรา”ติ.\csromanspacer{} อถ โข อายสฺมโต รฏฺฐปาลสฺส ญาติทาสี ตํ อาภิโทสิกํ กุมฺมาสํ อายสฺมโต รฏฺฐปาลสฺส ปตฺเต อากิรนฺตี หตฺถานญฺจ ปาทานญฺจ สรสฺส จ นิมิตฺตํ อคฺคเหสิ.}
\proseitem{300}{\csromanfolio{251}อถ โข อายสฺมโต รฏฺฐปาลสฺส ญาติทาสี เยนายสฺมโต รฏฺฐปาลสฺส มาตา เตนุปสงฺกมิ, อุปสงฺกมิตฺวา อายสฺมโต รฏฺฐปาลสฺส มาตรํ เอตทโวจ “ยคฺเฆยฺเย ชาเนยฺยาสิ อยฺยปุตฺโต รฏฺฐปาโล อนุปฺปตฺโต”ติ.\csromanspacer{} สเจ เช สจฺจํ ภณสิ, อทาสิํ ตํ กโรมีติ\footnote{สจฺจํ วทสิ, อทาสี ภวสีติ (สี, อิ), สจฺจํ วทสิ, อทาสี ภวิสฺสสิ (ก)}.\csromanspacer{} อถ โข อายสฺมโต รฏฺฐปาลสฺส มาตา เยนายสฺมโต รฏฺฐปาลสฺส ปิตา เตนุปสงฺกมิ, อุปสงฺกมิตฺวา อายสฺมโต รฏฺฐปาลสฺส ปิตรํ เอตทโวจ “ยคฺเฆ คหปติ ชาเนยฺยาสิ รฏฺฐปาโล กิร กุลปุตฺโต อนุปฺปตฺโต”ติ.\csromanspacer{} เตน โข ปน สมเยน อายสฺมา รฏฺฐปาโล ตํ อาภิโทสิกํ กุมฺมาสํ อญฺญตรํ กุฏฺฏมูลํ\footnote{กุฑฺฑํ (สี, สฺยา, กํ, อิ)} นิสฺสาย ปริภุญฺชติ.\csromanspacer{} อถ โข อายสฺมโต รฏฺฐปาลสฺส ปิตา เยนายสฺมา รฏฺฐปาโล เตนุปสงฺกมิ, อุปสงฺกมิตฺวา อายสฺมนฺตํ รฏฺฐปาลํ เอตทโวจ “อตฺถิ นาม ตาต รฏฺฐปาล อาภิโทสิกํ กุมฺมาสํ ปริภุญฺชิสฺสสิ, นนุ ตาต รฏฺฐปาล สกํ เคหํ คนฺตพฺพนฺ”ติ.\csromanspacer{} กุโต โน คหปติ, อมฺหากํ เคหํ อคารสฺมา อนคาริยํ ปพฺพชิตานํ.\csromanspacer{} อนคารา มยํ คหปติ, อคมมฺห โข เต คหปติ เคหํ, ตตฺถ เนว ทานํ อลตฺถมฺห, น ปจฺจกฺขานํ, อญฺญทตฺถุ อกฺโกสเมว อลตฺถมฺหาติ.\csromanspacer{} เอหิ ตาต รฏฺฐปาล ฆรํ คมิสฺสามาติ.\csromanspacer{} อลํ คหปติ, กตํ เม อชฺช ภตฺตกิจฺจํ.\csromanspacer{} เตน หิ ตาต รฏฺฐปาล อธิวาเสหิ สฺวาตนาย ภตฺตนฺติ.\csromanspacer{} อธิวาเสสิ โข อายสฺมา รฏฺฐปาโล ตุณฺหีภาเวน.\csromanspacer{} อถ โข อายสฺมโต รฏฺฐปาลสฺส ปิตา อายสฺมโต รฏฺฐปาลสฺส อธิวาสนํ วิทิตฺวา เยน สกํ นิเวสนํ เตนุปสงฺกมิ, อุปสงฺกมิตฺวา มหนฺตํ หิรญฺญสุวณฺณสฺส ปุญฺชํ การาเปตฺวา \csromanfolio{252}กิลญฺเชหิ ปฏิจฺฉาเทตฺวา อายสฺมโต รฏฺฐปาลสฺส ปุราณทุติยิกา อามนฺเตสิ “เอถ ตุเมฺห วธุโย เยน อลงฺกาเรน อลงฺกตา ปุพฺเพ รฏฺฐปาลสฺส กุลปุตฺตสฺส ปิยา โหถ มนาปา, เตน อลงฺกาเรน อลงฺกโรถา”ติ.}

\proseitem{301}{อถ โข อายสฺมโต รฏฺฐปาลสฺส ปิตา ตสฺสา รตฺติยา อจฺจเยน สเก นิเวสเน ปณีตํ ขาทนียํ โภชนียํ ปฏิยาทาเปตฺวา อายสฺมโต รฏฺฐปาลสฺส กาลํ อาโรเจสิ “กาโล ตาต รฏฺฐปาล นิฏฺฐิตํ ภตฺตนฺ”ติ.\csromanspacer{} อถ โข อายสฺมา รฏฺฐปาโล ปุพฺพณฺหสมยํ นิวาเสตฺวา ปตฺตจีวรมาทาย เยน สกปิตุ นิเวสนํ เตนุปสงฺกมิ, อุปสงฺกมิตฺวา ปญฺญตฺเต อาสเน นิสีทิ.\csromanspacer{} อถ โข อายสฺมโต รฏฺฐปาลสฺส ปิตา ตํ หิรญฺญสุวณฺณสฺส ปุญฺชํ วิวราเปตฺวา อายสฺมนฺตํ รฏฺฐปาลํ เอตทโวจ “อิทํ เต ตาต รฏฺฐปาล มาตุ มตฺติกํ ธนํ, อญฺญํ เปตฺติกํ, อญฺญํ ปิตามหํ.\csromanspacer{} สกฺกา ตาต รฏฺฐปาล โภเค จ ภุญฺชิตุํ ปุญฺญานิ จ กาตุํ, เอหิ ตฺวํ ตาต รฏฺฐปาล\footnote{รฏฺฐปาล สิกฺขํ ปจฺจกฺขาย (สพฺพตฺถ)} หีนายาวตฺติตฺวา โภเค จ ภุญฺชสฺสุ ปุญฺญานิ จ กโรหี”ติ.\csromanspacer{} สเจ เม ตฺวํ คหปติ วจนํ กเรยฺยาสิ, อิมํ หิรญฺญสุวณฺณสฺส ปุญฺชํ สกเฏ อาโรเปตฺวา นิพฺพาหาเปตฺวา มชฺเฌคงฺคาย นทิยา โสเต โอปิลาเปยฺยาสิ.\csromanspacer{} ตํ กิสฺส เหตุ, เย อุปฺปชฺชิสฺสนฺติ หิ เต คหปติ ตโตนิทานํ โสกปริเทวทุกฺขโทมนสฺสุปายาสาติ.\csromanspacer{} อถ โข อายสฺมโต รฏฺฐปาลสฺส ปุราณทุติยิกา ปจฺเจกํ ปาเทสุ คเหตฺวา อายสฺมนฺตํ รฏฺฐปาลํ เอตทโวจุํ “กีทิสา นาม ตา อยฺยปุตฺต อจฺฉราโย, ยาสํ ตฺวํ เหตุ พฺรหฺมจริยํ จรสี”ติ.\csromanspacer{} น โข มยํ ภคินี อจฺฉรานํ เหตุ พฺรหฺมจริยํ จรามาติ.\csromanspacer{} ภคินิวาเทน โน อยฺยปุตฺโต รฏฺฐปาโล สมุทาจรตีติ ตา ตตฺเถว มุจฺฉิตา ปปติํสุ.\csromanspacer{} อถ โข อายสฺมา รฏฺฐปาโล ปิตรํ เอตทโวจ “สเจ คหปติ โภชนํ ทาตพฺพํ, เทถ มา โน วิเหเฐถา”ติ.\csromanspacer{} ภุญฺช ตาต รฏฺฐปาล นิฏฺฐิตํ ภตฺตนฺติ.\csromanspacer{} อถ โข อายสฺมโต รฏฺฐปาลสฺส ปิตา อายสฺมนฺตํ รฏฺฐปาลํ ปณีเตน ขาทนีเยน โภชนีเยน สหตฺถา สนฺตปฺเปสิ สมฺปวาเรสิ.}

\csromanheader{2}{302. อถ โข อายสฺมา รฏฺฐปาโล ภุตฺตาวี โอนีตปตฺตปาณี ฐิตโกว อิมา คาถา อภาสิ\csromandash{}}
\csromanheader{2}{2. รฏฺฐปาลสุตฺต (82)}
\csromangathameasure{“ปสฺส จิตฺตีกตํ พิมฺพํ, อรุกายํ สมุสฺสิตํ. \\ อาตุรํ พหุสงฺกปฺปํ, ยสฺส นตฺถิ ธุวํ ฐิติ. \\ ปสฺส จิตฺตีกตํ รูปํ, มณินา กุณฺฑเลน จ. \\ อฏฺฐิ ตเจน โอนทฺธํ, สห วตฺเถภิ โสภติ. \\ อลตฺตกกตา ปาทา, มุขํ จุณฺณกมกฺขิตํ. \\ อลํ พาลสฺส โมหาย, โน จ ปารคเวสิโน. \\ อฏฺฐาปทกตา เกสา, เนตฺตา อญฺชนมกฺขิตา. \\ อลํ พาลสฺส โมหาย, โน จ ปารคเวสิโน. \\ อญฺชนีว นวา จิตฺตา, ปูติกาโย อลงฺกโต. \\ อลํ พาลสฺส โมหาย, โน จ ปารคเวสิโน. \\ โอทหิ มิคโว ปาสํ, นาสทา วากรํ มิโค. \\ ภุตฺวา นิวาปํ คจฺฉาม, กนฺทนฺเต มิคพนฺธเก”ติ.}
\csromangathagroup{\csromangathastanza{\csromanfolio{253}“ปสฺส จิตฺตีกตํ พิมฺพํ, อรุกายํ สมุสฺสิตํ. \\ อาตุรํ พหุสงฺกปฺปํ, ยสฺส นตฺถิ ธุวํ ฐิติ.}\csromangathastanza{ปสฺส จิตฺตีกตํ รูปํ, มณินา กุณฺฑเลน จ. \\ อฏฺฐิ ตเจน โอนทฺธํ, สห วตฺเถภิ โสภติ.}\csromangathastanza{อลตฺตกกตา ปาทา, มุขํ จุณฺณกมกฺขิตํ. \\ อลํ พาลสฺส โมหาย, โน จ ปารคเวสิโน.}\csromangathastanza{อฏฺฐาปทกตา เกสา, เนตฺตา อญฺชนมกฺขิตา. \\ อลํ พาลสฺส โมหาย, โน จ ปารคเวสิโน.}\csromangathastanza{อญฺชนีว นวา\footnote{อญฺชนีวณฺณวา (ก)} จิตฺตา, ปูติกาโย อลงฺกโต. \\ อลํ พาลสฺส โมหาย, โน จ ปารคเวสิโน.}\csromangathastanza{โอทหิ มิคโว ปาสํ, นาสทา วากรํ มิโค. \\ ภุตฺวา นิวาปํ คจฺฉาม\footnote{คจฺฉามิ (สฺยา, ก)}, กนฺทนฺเต มิคพนฺธเก”ติ.}}

\prose{อถ โข อายสฺมา รฏฺฐปาโล ฐิตโกว อิมา คาถา ภาสิตฺวา เยน รญฺโญ โกรพฺยสฺส มิคจีรํ เตนุปสงฺกมิ, อุปสงฺกมิตฺวา อญฺญตรสฺมิํ รุกฺขมูเล ทิวาวิหารํ นิสีทิ.}

\proseitem{303}{อถ โข ราชา โกรโพฺย มิควํ อามนฺเตสิ “โสเธหิ สมฺม มิคว มิคจีรํ อุยฺยานภูมิํ, คจฺฉาม สุภูมิํ ทสฺสนายา”ติ. “เอวํ เทวา”ติ โข มิคโว รญฺโญ โกรพฺยสฺส ปฏิสฺสุตฺวา มิคจีรํ โสเธนฺโต อทฺทส อายสฺมนฺตํ รฏฺฐปาลํ อญฺญตรสฺมิํ รุกฺขมูเล ทิวาวิหารํ นิสินฺนํ, ทิสฺวาน เยน ราชา โกรโพฺย เตนุปสงฺกมิ, อุปสงฺกมิตฺวา ราชานํ โกรพฺยํ เอตทโวจ “สุทฺธํ โข เต เทว มิคจีรํ, อตฺถิ เจตฺถ รฏฺฐปาโล นาม กุลปุตฺโต อิมสฺมิํเยว ถุลฺลโกฏฺฐิเก อคฺคกุลสฺส ปุตฺโต, ยสฺส ตฺวํ อภิณฺหํ กิตฺตยมาโน อโหสิ, โส อญฺญตรสฺมิํ รุกฺขมูเล ทิวาวิหารํ นิสินฺโน”ติ.\csromanspacer{} เตน หิ สมฺม มิคว อลํ ทานชฺช อุยฺยานภูมิยา ตเมว ทานิ มยํ ภวนฺตํ รฏฺฐปาลํ ปยิรุปาสิสฺสามาติ.\csromanspacer{} อถ โข ราชา โกรโพฺย ยํ ตตฺถ ขาทนียํ โภชนียํ ปฏิยตฺตํ, ตํ สพฺพํ วิสฺสชฺเชถาติ วตฺวา ภทฺรานิ ภทฺรานิ ยานานิ โยชาเปตฺวา ภทฺรํ ยานํ อภิรุหิตฺวา ภเทฺรหิ ภเทฺรหิ ยาเนหิ ถุลฺลโกฏฺฐิกมฺหา \csromanfolio{254}นิยฺยาสิ มหจฺจราชานุภาเวน\footnote{มหจฺจา ราชานุภาเวน (สี)} อายสฺมนฺตํ รฏฺฐปาลํ ทสฺสนาย.\csromanspacer{} ยาวติกา ยานสฺส ภูมิ, ยาเนน คนฺตฺวา ยานา ปจฺโจโรหิตฺวา ปตฺติโกว อุสฺสฏาย อุสฺสฏาย ปริสาย เยนายสฺมา รฏฺฐปาโล เตนุปสงฺกมิ, อุปสงฺกมิตฺวา อายสฺมตา รฏฺฐปาเลน สทฺธิํ สมฺโมทิ, สมฺโมทนียํ กถํ สารณียํ วีติสาเรตฺวา เอกมนฺตํ อฏฺฐาสิ, เอกมนฺตํ ฐิโต โข ราชา โกรโพฺย อายสฺมนฺตํ รฏฺฐปาลํ เอตทโวจ “อิธ ภวํ รฏฺฐปาล หตฺถตฺถเร\footnote{กฏฺฐตฺถเร (สฺยา, กํ)} นิสีทตู”ติ.\csromanspacer{} อลํ มหาราช, นิสีท ตฺวํ, นิสินฺโน อหํ สเก อาสเนติ.\csromanspacer{} นิสีทิ ราชา โกรโพฺย ปญฺญตฺเต อาสเน.\csromanspacer{} นิสชฺช โข ราชา โกรโพฺย อายสฺมนฺตํ รฏฺฐปาลํ เอตทโวจ\csromandash{}}

\proseitem{304}{จตฺตาริมานิ โภ รฏฺฐปาล ปาริชุญฺญานิ, เยหิ ปาริชุญฺเญหิ สมนฺนาคตา อิเธกจฺเจ เกสมสฺสุํ โอหาเรตฺวา กาสายานิ วตฺถานิ อจฺฉาเทตฺวา อคารสฺมา อนคาริยํ ปพฺพชนฺติ.\csromanspacer{} กตมานิ จตฺตาริ, ชราปาริชุญฺญํ พฺยาธิปาริชุญฺญํ โภคปาริชุญฺญํ ญาติปาริชุญฺญํ.\csromanspacer{} กตมญฺจ โภ รฏฺฐปาล ชราปาริชุญฺญํ.\csromanspacer{} อิธ โภ รฏฺฐปาล เอกจฺโจ ชิณฺโณ โหติ วุฑฺโฒ มหลฺลโก อทฺธคโต วโย-อนุปฺปตฺโต.\csromanspacer{} โส อิติ ปฏิสญฺจิกฺขติ “อหํ โขมฺหิ เอตรหิ ชิณฺโณ วุฑฺโฒ มหลฺลโก อทฺธคโต วโย-อนุปฺปตฺโต, น โข ปน มยา สุกรํ อนธิคตํ วา โภคํ อธิคนฺตุํ อธิคตํ วา โภคํ ผาติํ กาตุํ\footnote{ผาติกตฺตุํ (สี)}, ยํนูนาหํ เกสมสฺสุํ โอหาเรตฺวา กาสายานิ วตฺถานิ อจฺฉาเทตฺวา อคารสฺมา อนคาริยํ ปพฺพเชยฺยนฺ”ติ.\csromanspacer{} โส เตน ชราปาริชุญฺเญน สมนฺนาคโต เกสมสฺสุํ โอหาเรตฺวา กาสายานิ วตฺถานิ อจฺฉาเทตฺวา อคารสฺมา อนคาริยํ ปพฺพชติ.\csromanspacer{} อิทํ วุจฺจติ โภ รฏฺฐปาล ชราปาริชุญฺญํ.\csromanspacer{} ภวํ โข ปน รฏฺฐปาโล เอตรหิ ทหโร ยุวา สุสุกาฬเกโส ภเทฺรน โยพฺพเนน สมนฺนาคโต ปฐเมน วยสา, ตํ โภโต รฏฺฐปาลสฺส ชราปาริชุญฺญํ นตฺถิ.\csromanspacer{} กิํ ภวํ รฏฺฐปาโล ญตฺวา วา ทิสฺวา วา สุตฺวา วา อคารสฺมา อนคาริยํ ปพฺพชิโต. (1)}

\prose{กตมญฺจ โภ รฏฺฐปาล พฺยาธิปาริชุญฺญํ.\csromanspacer{} อิธ โภ รฏฺฐปาล เอกจฺโจ อาพาธิโก โหติ ทุกฺขิโต พาฬฺหคิลาโน.\csromanspacer{} โส อิติ}

\vspace{\tipitakaparskip}
\csromancenter{รฏฺฐปาลสุตฺต (82) ปฏิสญฺจิกฺขติ “อหํ โขมฺหิ เอตรหิ อาพาธิโก ทุกฺขิโต พาฬฺหคิลาโน, น โข ปน มยา สุกรํ อนธิคตํ วา โภคํ อธิคนฺตุํ อธิคตํ วา โภคํ ผาติํ กาตุํ, ยํนูนาหํ เกสมสฺสุํ โอหาเรตฺวา กาสายานิ วตฺถานิ อจฺฉาเทตฺวา อคารสฺมา อนคาริยํ ปพฺพเชยฺยนฺ”ติ.\csromanspacer{} โส เตน พฺยาธิปาริชุญฺเญน สมนฺนาคโต เกสมสฺสุํ โอหาเรตฺวา กาสายานิ วตฺถานิ อจฺฉาเทตฺวา อคารสฺมา อนคาริยํ ปพฺพชติ.\csromanspacer{} อิทํ วุจฺจติ โภ รฏฺฐปาล พฺยาธิปาริชุญฺญํ.\csromanspacer{} ภวํ โข ปน รฏฺฐปาโล เอตรหิ อปฺปาพาโธ อปฺปาตงฺโก สมเวปากินิยา คหณิยา สมนฺนาคโต นาติสีตาย นาจฺจุณฺหาย, ตํ โภโต รฏฺฐปาลสฺส พฺยาธิปาริชุญฺญํ นตฺถิ.\csromanspacer{} กิํ ภวํ รฏฺฐปาโล ญตฺวา วา ทิสฺวา วา สุตฺวา วา อคารสฺมา อนคาริยํ ปพฺพชิโต. (2)}
\prose{\csromanfolio{255}กตมญฺจ โภ รฏฺฐปาล โภคปาริชุญฺญํ.\csromanspacer{} อิธ โภ รฏฺฐปาล เอกจฺโจ อฑฺโฒ โหติ มหทฺธโน มหาโภโค.\csromanspacer{} ตสฺส เต โภคา อนุปุพฺเพน ปริกฺขยํ คจฺฉนฺติ.\csromanspacer{} โส อิติ ปฏิสญฺจิกฺขติ “อหํ โข ปุพฺเพ อฑฺโฒ อโหสิํ มหทฺธโน มหาโภโค, ตสฺส เม เต โภคา อนุปุพฺเพน ปริกฺขยํ คตา.\csromanspacer{} น โข ปน มายา สุกรํ อนธิคตํ วา โภคํ อธิคนฺตุํ อธิคตํ วา โภคํ ผาติํ กาตุํ, ยํนูนาหํ เกสมสฺสุํ โอหาเรตฺวา กาสายานิ วตฺถานิ อจฺฉาเทตฺวา อคารสฺมา อนคาริยํ ปพฺพเชยฺยนฺ”ติ.\csromanspacer{} โส เตน โภคปาริชุญฺเญน สมนฺนาคโต เกสมสฺสุํ โอหาเรตฺวา กาสายานิ วตฺถานิ อจฺฉาเทตฺวา อคารสฺมา อนคาริยํ ปพฺพชติ.\csromanspacer{} อิทํ วุจฺจติ โภ รฏฺฐปาล โภคปาริชุญฺญํ.\csromanspacer{} ภวํ โข ปน รฏฺฐปาโล อิมสฺมิํเยว ถุลฺลโกฏฺฐิเก อคฺคกุลสฺส ปุตฺโต, ตํ โภโต รฏฺฐปาลสฺส โภคปาริชุญฺญํ นตฺถิ.\csromanspacer{} กิํ ภวํ รฏฺฐปาโล ญตฺวา วา ทิสฺวา วา สุตฺวา วา อคารสฺมา อนคาริยํ ปพฺพชิโต. (3)}

\prose{กตมญฺจ โภ รฏฺฐปาล ญาติปาริชุญฺญํ.\csromanspacer{} อิธ โภ รฏฺฐปาล เอกจฺจสฺส พหู โหนฺติ มิตฺตามจฺจา ญาติสาโลหิตา, ตสฺส เต ญาตกา อนุปุพฺเพน ปริกฺขยํ คจฺฉนฺติ.\csromanspacer{} โส อิติ ปฏิสญฺจิกฺขติ “มมํ โข ปุพฺเพ พหู อเหสุํ มิตฺตามจฺจา ญาติสาโลหิตา, ตสฺส เม เต อนุปุพฺเพน ปริกฺขยํ คตา.\csromanspacer{} น โข ปน มยา สุกรํ อนธิคตํ วา โภคํ อธิคนฺตุํ อธิคตํ วา โภคํ ผาติํ กาตุํ, ยํนูนาหํ เกสมสฺสุํ โอหาเรตฺวา กาสายานิ วตฺถานิ อจฺฉาเทตฺวา อคารสฺมา อนคาริยํ ปพฺพเชยฺยนฺ”ติ. \csromanfolio{256}โส เตน ญาติปาริชุญฺเญน สมนฺนาคโต เกสมสฺสุํ โอหาเรตฺวา กาสายานิ วตฺถานิ อจฺฉาเทตฺวา อคารสฺมา อนคาริยํ ปพฺพชติ.\csromanspacer{} อิทํ วุจฺจติ โภ รฏฺฐปาล ญาติปาริชุญฺญํ.\csromanspacer{} โภโต โข ปน รฏฺฐปาลสฺส อิมสฺมิํ เยว ถุลฺลโกฏฺฐิเก พหู มิตฺตามจฺจา ญาติสาโลหิตา, ตํ โภโต รฏฺฐปาลสฺส ญาติปาริชุญฺญํ นตฺถิ.\csromanspacer{} กิํ ภวํ รฏฺฐปาโล ญตฺวา วา ทิสฺวา วา สุตฺวา วา อคารสฺมา อนคาริยํ ปพฺพชิโต. (4)}

\prose{อิมานิ โข โภ รฏฺฐปาล จตฺตาริ ปาริชุญฺญานิ.\csromanspacer{} เยหิ ปาริชุญฺเญหิ สมนฺนาคตา อิเธกจฺเจ เกสมสฺสุํ โอหาเรตฺวา กาสายานิ วตฺถานิ อจฺฉาเทตฺวา อคารสฺมา อนคาริยํ ปพฺพชนฺติ, ตานิ โภโต รฏฺฐปาลสฺส นตฺถิ.\csromanspacer{} กิํ ภวํ รฏฺฐปาโล ญตฺวา วา ทิสฺวา วา สุตฺวา วา อคารสฺมา อนคาริยํ ปพฺพชิโตติ.}

\proseitem{305}{อตฺถิ โข มหาราช เตน ภควตา ชานตา ปสฺสตา อรหตา สมฺมาสมฺพุทฺเธน จตฺตาโร ธมฺมุทฺเทสา อุทฺทิฏฺฐา.\csromanspacer{} เย อหํ\footnote{ยมหํ (สฺยา, กํ, ก)} ญตฺวา จ ทิสฺวา จ สุตฺวา จ อคารสฺมา อนคาริยํ ปพฺพชิโต.\csromanspacer{} กตเม จตฺตาโร, “อุปนิยฺยติ โลโก อทฺธุโว”ติ โข มหาราช เตน ภควตา ชานตา ปสฺสตา อรหตา สมฺมาสมฺพุทฺเธน ปฐโม ธมฺมุทฺเทโส อุทฺทิฏฺโฐ, “ยมหํ ญตฺวา จทิสฺวา จ สุตฺวา จ อคารสฺมา อนคาริยํ ปพฺพชิโต.\csromanspacer{} อตาโณ โลโก อนภิสฺสโร”ติ โข มหาราช เตน ภควตา ชานตา ปสฺสตา อรหตา สมฺมาสมฺพุทฺเธน ทุติโย ธมฺมุทฺเทโส อุทฺทิฏฺโฐ, ยมหํ ญตฺวา จ ทิสฺวา จ สุตฺวา จ อคารสฺมา อนคาริยํ ปพฺพชิโต. “อสฺสโก โลโก สพฺพํ ปหาย คมนียนฺ”ติ โข มหาราช เตน ภควตา ชานตา ปสฺสตา อรหตา สมฺมาสมฺพุทฺเธน ตติโย ธมฺมุทฺเทโส อุทฺทิฏฺโฐ, ยมหํ ญตฺวา จ ทิตฺวา จ สุตฺวา จ อคารสฺมา อนคาริยํ ปพฺพชิโต. “อูโน โลโก อติตฺโต ตณฺหาทาโส”ติ โข มหาราช เตน ภควตา ชานตา ปสฺสตา อรหตา สมฺมาสมฺพุทฺเธน จตุตฺโถ ธมฺมุทฺเทโส อุทฺทิฏฺโฐ, ยมหํ ญตฺวา จ ทิสฺวา จ สุตฺวา จ อคารสฺมา อนคาริยํ ปพฺพชิโต.\csromanspacer{} อิเม โข มหาราช เตน ภควตา ชานตา ปสฺสตา อรหตา สมฺมาสมฺพุทฺเธน จตฺตาโร ธมฺมุทฺเทสา อุทฺทิฏฺฐา, เย อหํ ญตฺวา จ ทิสฺวา จ สุตฺวา จ อคารสฺมา อนคาริยํ ปพฺพชิโตติ.}

\csromanheader{2}{2. รฏฺฐปาลสุตฺต (82)}
\proseitem{306}{\csromanfolio{257}“อุปนิยฺยติ โลโก อทฺธุโว”ติ ภวํ รฏฺฐปาโล อาห, อิมสฺส โภ รฏฺฐปาล ภาสิตสฺส กถํ อตฺโถ ทฏฺฐพฺโพติ.\csromanspacer{} ตํ กิํ มญฺญสิ มหาราช ตฺวํ วีสติวสฺสุทฺเทสิโกปิ ปณฺณวีสติวสฺสุทฺเทสิโกปิ หตฺถิสฺมิํปิ กตาวี อสฺสสฺมิํปิ กตาวี รถสฺมิํปิ กตาวี ธนุสฺมิํปิ กตาวี ถรุสฺมิํปิ กตาวี อูรุพลี พาหุพลี อลมตฺโต สงฺคามาวจโรติ? อโหสิํ อหํ โภ รฏฺฐปาล วีสติวสฺสุทฺเทสิโกปิ ปณฺณวีสติวสฺสุทฺเทสิโกปิ หตฺถิสฺมิํปิ กตาวี อสฺสสฺมิํปิ กตาวี รถสฺมิํปิ กตาวี ธนุสฺมิํปิ กตาวี ถรุสฺมิํปิ กตาวี อูรุพลี พาหุพลี อลมตฺโต สงฺคามาวจโร.\csromanspacer{} อปฺเปกทาหํ โภ รฏฺฐปาล อิทฺธิมาว มญฺเญ, น\footnote{อิทฺธิมา มญฺเญ น (สฺยา, กํ), อิทฺธิมา จ มญฺเญ (สี), น วิย มญฺเญ (ก)} อตฺตโน พเลน สมสมํ สมนุปสฺสามีติ.\csromanspacer{} ตํ กิํ มญฺญสิ มหาราช, เอวเมว ตฺวํ เอตรหิ อูรุพลี พาหุพลี อลมตฺโต สงฺคามาวจโรติ? โน หิทํ โภ รฏฺฐปาล, เอตรหิ ชิณฺโณ วุฑฺโฒ มหลฺลโก อทฺธคโต วโย-อนุปฺปตฺโต, อาสีติโก เม วโย วตฺตติ, อปฺเปกทาหํ โภ รฏฺฐปาล อิธ ปาทํ กริสฺสามีติ อญฺเญเนว ปาทํ กโรมีติ.\csromanspacer{} อิทํ โข ตํ มหาราช เตน ภควตา ชานตา ปสฺสตา อรหตา สมฺมาสมฺพุทฺเธน สนฺธาย ภาสิตํ “อุปนิยฺยติ โลโก อทฺธุโว”ติ.\csromanspacer{} ยมหํ ญตฺวา จ ทิสฺวา จ สุตฺวา จ อคารสฺมา อนคาริยํ ปพฺพชิโตติ, อจฺฉริยํ โภ รฏฺฐปาล, อพฺภุตํ โภ รฏฺฐปาล, ยาว สุภาสิตํ จิทํ เตน ภควตา ชานตา ปสฺสตา อรหตา สมฺมาสมฺพุทฺเธน “อุปนิยฺยติ โลโก อทฺธุโว”ติ, อุปนิยฺยติ หิ โภ รฏฺฐปาล โลโก อทฺธุโว. (1)}

\prose{สํวิชฺชนฺเต โข โภ รฏฺฐปาล อิมสฺมิํ ราชกุเล หตฺถิกายาปิ อสฺสกายาปิ รถกายาปิ ปตฺติกายาปิ, อมฺหากํ อาปทาสุ ปริโยธาย วตฺติสฺสนฺติ. “อตาโณ โลโก อนภิสฺสโร”ติ ภวํ รฏฺฐปาโล อาห, อิมสฺส ปน โภ รฏฺฐปาล ภาสิตสฺส กถํ อตฺโถ ทฏฺฐพฺโพติ.\csromanspacer{} ตํ กิํ มญฺญสิ มหาราช, อตฺถิ เต โกจิ อนุสายิโก อาพาโธติ.\csromanspacer{} อตฺถิ เม โภ รฏฺฐปาล อนุสายิโก อาพาโธ, อปฺเปกทา มํ โภ รฏฺฐปาล มิตฺตามจฺจา ญาติสาโลหิตา ปริวาเรตฺวา ฐิตา โหนฺติ “อิทานิ ราชา โกรโพฺย กาลํ กริสฺสติ, อิทานิ ราชา โกรโพฺย กาลํ กริสฺสตี”ติ.\csromanspacer{} ตํ กิํ มญฺญสิ มหาราช, ลภสิ \csromanfolio{258}ตฺวํ เต มิตฺตามจฺเจ ญาติสาโลหิเต “อายนฺตุ เม โภนฺโต มิตฺตามจฺจา ญาติสาโลหิตา, สพฺเพว สนฺตา อิมํ เวทนํ สํวิภชถ, ยถาหํ ลหุกตริกํ เวทนํ เวทิเยยฺยนฺ”ติ, อุทาหุ ตฺวํเยว ตํ เวทนํ เวทิยสีติ? นาหํ โภ รฏฺฐปาล ลภามิ เต มิตฺตามจฺเจ ญาติสาโลหิเต “อายนฺตุ เม โภนฺโต มิตฺตามจฺจา ญาติสาโลหิตา, สพฺเพว สนฺตา อิมํ เวทนํ สํวิภชถ, ยถาหํ ลหุกตริกํ เวทนํ เวทิเยยฺยนฺ”ติ, อถ โข อหเมว ตํ เวทนํ เวทิยามีติ.\csromanspacer{} อิทํ โข ตํ มหาราช เตน ภควตา ชานตา ปสฺสตา อรหตา สมฺมาสมฺพุทฺเธน สนฺธาย ภาสิตํ “อตาโณ โลโก อนภิสฺสโร”ติ.\csromanspacer{} ยมหํ ญตฺวา จ ทิสฺวา จ สุตฺวา จ อคารสฺมา อนคาริยํ ปพฺพชิโตติ, อจฺฉริยํ โภ รฏฺฐปาล, อพฺภุตํ โภ รฏฺฐปาล, ยาว สุภาสิตํ จิทํ เตน ภควตา ชานตา ปสฺสตา อรหตา สมฺมาสมฺพุทฺเธน “อตาโณ โลโก อนภิสฺสโร”ติ, อตาโณ หิ โภ รฏฺฐปาล โลโก อนภิสฺสโร. (2)}

\prose{สํวิชฺชติ โข โภ รฏฺฐปาล อิมสฺมิํ ราชกุเล ปหูตํ หิรญฺญสุวณฺณํ ภูมิคตญฺจ เวหาสคตญฺจ. “อสฺสโก โลโก สพฺพํ ปหาย คมนียนฺ”ติ ภวํ รฏฺฐปาโล อาห, อิมสฺส ปน โภ รฏฺฐปาล ภาสิตสฺส กถํ อตฺโถ ทฏฺฐพฺโพติ.\csromanspacer{} ตํ กิํ มญฺญสิ มหาราช, ยถา ตฺวํ เอตรหิ ปญฺจหิ กามคุเณหิ สมปฺปิโต สมงฺคีภูโต ปริจาเรสิ, ลจฺฉสิ ตฺวํ ปรตฺถาปิ “เอวเมวาหํ อิเมเหว ปญฺจหิ กามคุเณหิ สมปฺปิโต สมงฺคีภูโต ปริจาเรมี”ติ, อุทาหุ อญฺเญ อิมํ โภคํ ปฏิปชฺชิสฺสนฺติ, ตฺวํ ปน ยถากมฺมํ คมิสฺสสีติ? ยถาหํ โภ รฏฺฐปาล เอตรหิ ปญฺจหิ กามคุเณหิ สมปฺปิโต สมงฺคีภูโต ปริจาเรมิ, นาหํ ลจฺฉามิ ปรตฺถาปิ “เอวเมว อิเมเหว ปญฺจหิ กามคุเณหิ สมปฺปิโต สมงฺคีภูโต ปริจาเรมี”ติ, อถ โข อญฺเญ อิมํ โภคํ ปฏิปชฺชิสฺสนฺติ, อหํ ปน ยถากมฺมํ คมิสฺสามีติ.\csromanspacer{} อิทํ โข ตํ มหาราช เตน ภควตา ชานตา ปสฺสตา อรหตา สมฺมาสมฺพุทฺเธน สนฺธาย ภาสิตํ “อสฺสโก โลโก สพฺพํ ปหาย คมนียนฺ”ติ.\csromanspacer{} ยมหํ ญตฺวา จ ทิสฺวา จ สุตฺวา จ อคารสฺมา อนคาริยํ ปพฺพชิโตติ.\csromanspacer{} อจฺฉริยํ โภ รฏฺฐปาล, อพฺภุตํ โภ รฏฺฐปาล, ยาว สุภาสิตํ จิทํ เตน ภควตา ชานตา ปสฺสตา อรหตา สมฺมาสมฺพุทฺเธน “อสฺสโก โลโก สพฺพํ ปหาย}

\vspace{\tipitakaparskip}
\csromancenter{รฏฺฐปาลสุตฺต (82) คมนียนฺ”ติ, อสฺสโก หิ โภ รฏฺฐปาล โลโก สพฺพํ ปหาย คมนียํ. (3)}
\prose{\csromanfolio{259}“อูโน โลโก อติตฺโต ตณฺหาทาโส”ติ ภวํ รฏฺฐปาโล อาห, อิมสฺส โภ รฏฺฐปาล ภาสิตสฺส กถํ อตฺโถ ทฏฺฐพฺโพติ.\csromanspacer{} ตํ กิํ มญฺญสิ มหาราช, ผีตํ กุรุํ อชฺฌาวสสีติ? เอวํ โภ รฏฺฐปาล ผีตํ กุรุํ อชฺฌาวสามีติ.\csromanspacer{} ตํ กิํ มญฺญสิ มหาราช, อิธ ปุริโส อาคจฺเฉยฺย ปุรตฺถิมาย ทิสาย สทฺธายิโก ปจฺจยิโก, โส ตํ อุปสงฺกมิตฺวา เอวํ วเทยฺย “ยคฺเฆ มหาราช ชาเนยฺยาสิ, อหํ อาคจฺฉามิ ปุรตฺถิมาย ทิสาย, ตตฺถทฺทสํ มหนฺตํ ชนปทํ อิทฺธญฺเจว ผีตญฺจ พหุชนํ อากิณฺณมนุสฺสํ, พหู ตตฺถ หตฺถิกายา อสฺสกายา รถกายา ปตฺติกายา, พหุ ตตฺถ ธนธญฺญํ\footnote{ทนฺตาชินํ (สี, สฺยา, กํ, อิ)}, พหุ ตตฺถ หิรญฺญสุวณฺณํ อกตญฺเจว กตญฺจ, พหุ ตตฺถ อิตฺถิปริคฺคโห, สกฺกา จ ตาวตเกเนว พลมตฺเตน\footnote{พลตฺเถน (สี, สฺยา, กํ, อิ), พหลตฺเถน (ก)} อภิวิชินิตุํ, อภิวิชิน มหาราชา”ติ, กินฺติ นํ กเรยฺยาสีติ.\csromanspacer{} ตมฺปิ มยํ โภ รฏฺฐปาล อภิวิชิย อชฺฌาวเสยฺยามาติ.\csromanspacer{} ตํ กิํ มญฺญสิ มหาราช, อิธ ปุริโส อาคจฺเฉยฺย ปจฺฉิมาย ทิสาย.\csromanspacer{} อุตฺตราย ทิสาย.\csromanspacer{} ทกฺขิณาย ทิสาย.\csromanspacer{} ปรสมุทฺทโต สทฺธายิโก ปจฺจยิโก, โส ตํ อุปสงฺกมิตฺวา เอวํ วเทยฺย “ยคฺเฆ มหาราช ชาเนยฺยาสิ, อหํ อาคจฺฉามิ ปรสมุทฺทโต, ตตฺถทฺทสํ มหนฺตํ ชนปทํ อิทฺธญฺเจว ผีตญฺจ พหุชนํ อากิณฺณมนุสฺสํ, พหู ตตฺถ หตฺถิกายา อสฺสกายา รถกายา ปตฺถิกายา, พหุ ตตฺถ ธนธญฺญํ, พหุ ตตฺถ หิรญฺญสุวณฺณํ อกตญฺเจว กตญฺจ, พหุ ตตฺถ อิตฺถิปริคฺคโห, สกฺกา จ ตาวตเกเนว พลมตฺเตน อภิวิชินิตุํ, อภิวิชิน มหาราชา”ติ, กินฺติ นํ กเรยฺยาสีติ.\csromanspacer{} ตมฺปิ มยํ โภ รฏฺฐปาล อภิวิชิย อชฺฌาวเสยฺยามาติ.\csromanspacer{} อิทํ โข ตํ มหาราช เตน ภควตา ชานตา ปสฺสตา อรหตา สมฺมาสมฺพุทฺเธน สนฺธาย ภาสิตํ “อูโน โลโก อติตฺโต ตณฺหาทาโส”ติ.\csromanspacer{} ยมหํ ญตฺวา จ ทิสฺวา จ สุตฺวา จ อคารสฺมา อนคาริยํ ปพฺพชิโตติ, อจฺฉริยํ โภ รฏฺฐปาล, อพฺภุตํ โภ รฏฺฐปาล, ยาว สุภาสิตํ จิทํ เตน ภควตา ชานตา ปสฺสตา อรหตา สมฺมาสมฺพุทฺเธน “อูโน โลโก อติตฺโต ตณฺหาทาโส”ติ.\csromanspacer{} อูโน หิ โภ รฏฺฐปาล โลโก อติตฺโต ตณฺหาทาโสติ. (4)}

\prose{\csromanfolio{260}อิทมโวจ อายสฺมา รฏฺฐปาโล, อิทํ วตฺวา อถาปรํ เอตทโวจ\csromandash{}}

\csromanhangingbegin
\hangingheaditem{307}{ปสฺสามิ โลเก สธเน มนุสฺเส,}
\hangingbody{ลทฺธาน วิตฺตํ น ททนฺติ โมหา.}
\hangingbody{ลุทฺธา ธนํ1 สนฺนิจยํ กโรนฺติ,}
\hangingbody{ภิยฺโยว กาเม อภิปตฺถยนฺติ.}
\hangingbody{ราชา ปสยฺหา ปถวิํ วิชิตฺวา,}
\hangingbody{สสาครนฺตํ มหิมาวสนฺโต2.}
\hangingbody{โอรํ สมุทฺทสฺส อติตฺตรูโป,}
\hangingbody{ปารํ สมุทฺทสฺสปิ ปตฺถเยถ.}
\hangingbody{ราชา จ อญฺเญ จ พหู มนุสฺสา,}
\hangingbody{อวีตตณฺหา3 มรณํ อุเปนฺติ.}
\hangingbody{อูนาว หุตฺวาน ชหนฺติ เทหํ,}
\hangingbody{กาเมหิ โลกมฺหิ น หตฺถิ ติตฺติ.}
\hangingbody{กนฺทนฺติ นํ ญาตี ปกิริย เกเส,}
\hangingbody{อโหวตา โน อมราติ จาหุ.}
\hangingbody{วตฺเถน นํ ปารุตํ นีหริตฺวา,}
\hangingbody{จิตํ สมาทาย4 ตโตฑหนฺติ.}
\hangingbody{โส ฑยฺหติ สูเลหิ ตุชฺชมาโน,}
\hangingbody{เอเกน วตฺเถน ปหาย โภเค.}
\hangingbody{น มียมานสฺส ภวนฺติ ตาณา,}
\hangingbody{ญาตีธ มิตฺตา อถ วา สหายา.}
\hangingbody{ทายาทกา ตสฺส ธนํ หรนฺติ,}
\hangingbody{สตฺโต ปน คจฺฉติ เยน กมฺมํ.}
\hangingbody{น มียมานํ ธนมเนฺวติ กิญฺจิ,}
\hangingbody{ปุตฺตา จ ทารา จ ธนํ จ รฏฺฐํ.}
\csromanhangingend

\csromanheader{2}{2. รฏฺฐปาลสุตฺต (82)}
\csromangathameasure{น ทีฆมายุํ ลภเต ธเนน, น จาปิ วิตฺเตน ชรํ วิหนฺติ. \\ อปฺปํ หิทํ ชีวิตมาหุ ธีรา, อสสฺสตํ วิปฺปริณามธมฺมํ. \\ อฑฺฒา ทลิทฺทา จ ผุสนฺติ ผสฺสํ, \\ พาโล จ ธีโร จ ตเถว ผุฏฺโฐ. \\ พาโล จ พาลฺยา วธิโตว เสติ, \\ ธีโร จ น เวธติ ผสฺสผุฏฺโฐ. \\ ตสฺมา หิ ปญฺญาว ธเนน เสยฺโย, \\ ยาย โวสานมิธาธิคจฺฉติ. \\ อโพฺยสิตตฺตา หิ ภวาภเวสุ, \\ ปาปานิ กมฺมานิ กโรนฺติ โมหา. \\ อุเปติ คพฺภญฺจ ปรญฺจ โลกํ, \\ สํสารมาปชฺช ปรมฺปราย. \\ ตสฺสปฺปปญฺโญ อภิสทฺทหนฺโต, \\ อุเปติ คพฺภญฺจ ปรญฺจ โลกํ. \\ โจโร ยถา สนฺธิมุเข คหิโต, \\ สกมฺมุนา หญฺญติ ปาปธมฺโม. \\ เอวํ ปชา เปจฺจ ปรมฺหิ โลเก, \\ สกมฺมุนา หญฺญติ ปาปธมฺโม. \\ กามาหิ จิตฺรา มธุรา มโนรมา, \\ วิรูปรูเปน มเถนฺติ จิตฺตํ. \\ อาทีนวํ กามคุเณสุ ทิสฺวา, \\ ตสฺมา อหํ ปพฺพชิโตมฺหิ ราช. \\ ทุมปฺผลาเนว ปตนฺติ มาณวา, \\ ทหรา จ วุฑฺฒา จ สรีรเภทา. \\ เอตํปิ ทิสฺวา ปพฺพชิโตมฺหิ ราช, \\ อปณฺณกํ สามญฺญเมว เสยฺโยติ.}
\csromangathagroup{\csromangathastanza{\csromanfolio{261}น ทีฆมายุํ ลภเต ธเนน, น จาปิ วิตฺเตน ชรํ วิหนฺติ. \\ อปฺปํ หิทํ ชีวิตมาหุ ธีรา, อสสฺสตํ วิปฺปริณามธมฺมํ.}\csromangathastanza{อฑฺฒา ทลิทฺทา จ ผุสนฺติ ผสฺสํ, \\ พาโล จ ธีโร จ ตเถว ผุฏฺโฐ. \\ พาโล จ พาลฺยา วธิโตว เสติ, \\ ธีโร จ\footnote{ธีโรว (ก)} น เวธติ ผสฺสผุฏฺโฐ.}\csromangathastanza{ตสฺมา หิ ปญฺญาว ธเนน เสยฺโย, \\ ยาย โวสานมิธาธิคจฺฉติ. \\ อโพฺยสิตตฺตา\footnote{อโสสิตตฺตา (สี, อิ)} หิ ภวาภเวสุ, \\ ปาปานิ กมฺมานิ กโรนฺติ โมหา.}\csromangathastanza{อุเปติ คพฺภญฺจ ปรญฺจ โลกํ, \\ สํสารมาปชฺช ปรมฺปราย. \\ ตสฺสปฺปปญฺโญ อภิสทฺทหนฺโต, \\ อุเปติ คพฺภญฺจ ปรญฺจ โลกํ.}\csromangathastanza{โจโร ยถา สนฺธิมุเข คหิโต, \\ สกมฺมุนา หญฺญติ ปาปธมฺโม. \\ เอวํ ปชา เปจฺจ ปรมฺหิ โลเก, \\ สกมฺมุนา หญฺญติ ปาปธมฺโม.}\csromangathastanza{กามาหิ จิตฺรา มธุรา มโนรมา, \\ วิรูปรูเปน มเถนฺติ จิตฺตํ. \\ อาทีนวํ กามคุเณสุ ทิสฺวา, \\ ตสฺมา อหํ ปพฺพชิโตมฺหิ ราช.}\csromangathastanza{ทุมปฺผลาเนว ปตนฺติ มาณวา, \\ ทหรา จ วุฑฺฒา จ สรีรเภทา. \\ เอตํปิ ทิสฺวา\footnote{เอวมฺปิ ทิสฺวา (สี), เอตํ วิทิตฺวา (สฺยา, กํ)} ปพฺพชิโตมฺหิ ราช, \\ อปณฺณกํ สามญฺญเมว เสยฺโยติ.}}

\nitthitam{รฏฺฐปาลสุตฺตํ นิฏฺฐิตํ ทุติยํ.}
\csromanheader{2}{3. มฆเทวสุตฺต}
\proseitem{308}{\csromanfolio{262}เอวํ เม สุตํ\csromandash{}เอกํ สมยํ ภควา มิถิลายํ วิหรติ มฆเทว-อมฺพวเน\footnote{มขาเทว-อมฺพวเน (สี, อิ), มคฺฆเทว-อมฺพวเน (ก)}.\csromanspacer{} อถ โข ภควา อญฺญตรสฺมิํ ปเทเส สิตํ ปาตฺวากาสิ.\csromanspacer{} อถ โข อายสฺมโต อานนฺทสฺส เอตทโหสิ “โก นุ โข เหตุ โก ปจฺจโย ภควโต สิตสฺส ปาตุกมฺมาย, น อการเณน ตถาคตา สิตํ ปาตุกโรนฺตี”ติ.\csromanspacer{} อถ โข อายสฺมา อานนฺโท เอกํสํ จีวรํ กตฺวา เยน ภควา เตนญฺชลิํ ปณาเมตฺวา ภควนฺตํ เอตทโวจ “โก นุ โข ภนฺเต เหตุ โก ปจฺจโย ภควโต สิตสฺส ปาตุกมฺมาย, น อการเณน ตถาคตา สิตํ ปาตุกโรนฺตี”ติ.\csromanspacer{} ภูตปุพฺพํ อานนฺท อิมิสฺสาเยว มิถิลายํ ราชา อโหสิ มฆเทโว นาม ธมฺมิโก ธมฺมราชา ธมฺเม ฐิโต มหาราชา, ธมฺมํ จรติ พฺราหฺมณคหปติเกสุ เนคเมสุ เจว ชานปเทสุ จ, อุโปสถญฺจ อุปวสติ จาตุทฺทสิํ ปญฺจทสิํ อฏฺฐมิํ จ ปกฺขสฺส.\csromanspacer{} อถ โข อานนฺท ราชา มฆเทโว พหูนํ วสฺสานํ พหูนํ วสฺสสตานํ พหูนํ วสฺสสหสฺสานํ อจฺจเยน กปฺปกํ อามนฺเตสิ “ยทา เม สมฺม กปฺปก ปสฺเสยฺยาสิ สิรสฺมิํ ปลิตานิ ชาตานิ, อถ เม อาโรเจยฺยาสี”ติ. “เอวํ เทวา”ติ โข อานนฺท กปฺปโก รญฺโญ มฆเทวสฺส ปจฺจสฺโสสิ.\csromanspacer{} อทฺทสา โข อานนฺท กปฺปโก พหูนํ วสฺสานํ พหูนํ วสฺสสตานํ พหูนํ วสฺสสหสฺสานํ อจฺจเยน รญฺโญ มฆเทวสฺส สิรสฺมิํ ปลิตานิ ชาตานิ, ทิสฺวาน ราชานํ มฆเทวํ เอตทโวจ “ปาตุภูตา โข เทวสฺส เทวทูตา, ทิสฺสนฺติ สิรสฺมิํ ปลิตานิ ชาตานี”ติ.\csromanspacer{} เตน หิ สมฺม กปฺปก ตานิ ปลิตานิ สาธุกํ สณฺฑาเสน อุทฺธริตฺวา มม อญฺชลิสฺมิํ ปติฏฺฐาเปหีติ. “เอวํ เทวา”ติ โข อานนฺท กปฺปโก รญฺโญ มฆเทวสฺส ปฏิสฺสุตฺวา ตานิ ปลิตานิ สาธุกํ สณฺฑาเสน อุทฺธริตฺวา รญฺโญ มฆเทวสฺส อญฺชลิสฺมิํ ปติฏฺฐาเปสิ.}

\proseitem{309}{อถ โข อานนฺท ราชา มฆเทโว กปฺปกสฺส คามวรํ ทตฺวา เชฏฺฐปุตฺตํ กุมารํ อามนฺตาเปตฺวา เอตทโวจ “ปาตุภูตา โข เม ตาต กุมาร เทวทูตา, ทิสฺสนฺติ สิรสฺมิํ ปลิตานิ ชาตานิ, ภุตฺตา โข ปน เม มานุสกา กามา, สมโย ทิพฺเพ กาเม ปริเยสิตุํ.\csromanspacer{} เอหิ}

\vspace{\tipitakaparskip}
\csromancenter{มฆเทวสุตฺต (83) ตฺวํ ตาต กุมาร อิมํ รชฺชํ ปฏิปชฺช, อหํ ปน เกสมสฺสุํ โอหาเรตฺวา กาสายานิ วตฺถานิ อจฺฉาเทตฺวา อคารสฺมา อนคาริยํ ปพฺพชิสฺสามิ.\csromanspacer{} เตน หิ ตาต กุมาร ยทา ตฺวํปิ ปสฺเสยฺยาสิ สิรสฺมิํ ปลิตานิ ชาตานิ, อถ กปฺปกสฺส คามวรํ ทตฺวา เชฏฺฐปุตฺตํ กุมารํ สาธุกํ รชฺเช สมนุสาสิตฺวา เกสมสฺสุํ โอหาเรตฺวา กาสายานิ วตฺถานิ อจฺฉาเทตฺวา อคารสฺมา อนคาริยํ ปพฺพเชยฺยาสิ, เยน เม อิทํ กลฺยาณํ วตฺตํ นิหิตํ อนุปฺปวตฺเตยฺยาสิ, มา โข เม ตฺวํ อนฺติมปุริโส อโหสิ.\csromanspacer{} ยสฺมิํ โข ตาต กุมาร ปุริสยุเค วตฺตมาเน เอวรูปสฺส กลฺยาณสฺส วตฺตสฺส สมุจฺเฉโท โหติ, โส เตสํ อนฺติมปุริโส โหติ.\csromanspacer{} ตํ ตาหํ ตาต กุมาร เอวํ วทามิ ‘เยน เม อิทํ กลฺยาณํ วตฺตํ นิหิตํ อนุปฺปวตฺเตยฺยาสิ, มา โข เม ตฺวํ อนฺติมปุริโส อโหสี’ติ”.\csromanspacer{} อถ โข อานนฺท ราชา มฆเทโว กปฺปกสฺส คามวรํ ทตฺวา เชฏฺฐปุตฺตํ กุมารํ สาธุกํ รชฺเช สมนุสาสิตฺวา อิมสฺมิํ เยว มฆเทว-อมฺพวเน เกสมสฺสุํ โอหาเรตฺวา กาสายานิ วตฺถานิ อจฺฉาเทตฺวา อคารสฺมา อนคาริยํ ปพฺพชิ.\csromanspacer{} โส เมตฺตาสหคเตน เจตสา เอกํ ทิสํ ผริตฺวา วิหาสิ.\csromanspacer{} ตถา ทุติยํ.\csromanspacer{} ตถา ตติยํ.\csromanspacer{} ตถา จตุตฺถํ.\csromanspacer{} อิติ อุทฺธมโธ ติริยํ สพฺพธิ สพฺพตฺตตาย สพฺพาวนฺตํ โลกํ เมตฺตาสหคเตน เจตสา วิปุเลน มหคฺคเตน อปฺปมาเณน อเวเรน อพฺยาพชฺเฌน\footnote{อพฺยาปชฺเฌน (สี, สฺยา, กํ, อิ), อพฺยาปชฺเชน (ก)} ผริตฺวา วิหาสิ.\csromanspacer{} กรุณา สหคเตน เจตสา.\csromanspacer{} มุทิตาสหคเตน เจตสา.\csromanspacer{} อุเปกฺขาสหคเตน เจตสา เอกํ ทิสํ ผริตฺวา วิหาสิ.\csromanspacer{} ตถา ทุติยํ.\csromanspacer{} ตถา ตติยํ.\csromanspacer{} ตถา จตุตฺถํ.\csromanspacer{} อิติ อุทฺธมโธ ติริยํ สพฺพธิ สพฺพตฺตตาย สพฺพาวนฺตํ โลกํ อุเปกฺขาสหคเตน เจตสา วิปุเลน มหคฺคเตน อปฺปมาเณน อเวเรน อพฺยาพชฺเฌน ผริตฺวา วิหาสิ.}
\prose{\csromanfolio{263}ราชา โข ปนานนฺท มฆเทโว จตุราสีติวสฺสสหสฺสานิ กุมารกีฬิตํ กีฬิ, จตุราสีติวสฺสสหสฺสานิ โอปรชฺชํ กาเรสิ, จตุราสีติวสฺสสหสฺสานิ รชฺชํ กาเรสิ, จตุราสีติวสฺสสหสฺสานิ อิมสฺมิํเยว มฆเทว-อมฺพวเน อคารสฺมา อนคาริยํ ปพฺพชิโต พฺรหฺมจริยมจริ.\csromanspacer{} โส จตฺตาโร พฺรหฺมวิหาเร ภาเวตฺวา กายสฺส เภทา ปรํ มรณา พฺรหฺมโลกูปโค อโหสิ.}

\proseitem{310}{\csromanfolio{264}อถ โข รญฺโญ อานนฺท มฆเทวสฺส ปุตฺโต พหูนํ วสฺสานํ พหูนํ วสฺสสตานํ พหูนํ วสฺสสหสฺสานํ อจฺจเยน กปฺปกํ อามนฺเตสิ “ยทา เม สมฺม กปฺปก ปสฺเสยฺยาสิ สิรสฺมิํ ปลิตานิ ชาตานิ, อถ โข อาโรเจยฺยาสี”ติ. “เอวํ เทวา”ติ โข อานนฺท กปฺปโก รญฺโญ มฆเทวสฺส ปุตฺตสฺส ปจฺจสฺโสสิ.\csromanspacer{} อทฺทสา โข อานนฺท กปฺปโก พหูนํ วสฺสานํ พหูนํ วสฺสสตานํ พหูนํ วสฺสสหสฺสานํ อจฺจเยน รญฺโญ มฆเทวสฺส ปุตฺตสฺส สิรสฺมิํ ปลิตานิ ชาตานิ, ทิสฺวาน รญฺโญ มฆเทวสฺส ปุตฺตํ เอตทโวจ “ปาตุภูตา โข เทวสฺส เทวทูตา, ทิสฺสนฺติ สิรสฺมิํ ปลิตานิ ชาตานี”ติ.\csromanspacer{} เตน หิ สมฺม กปฺปก ตานิ ปลิตานิ สาธุกํ สณฺฑาเสน อุทฺธริตฺวา มม อญฺชลิสฺมิํ ปติฏฺฐาเปหีติ. “เอวํ เทวา”ติ โข อานนฺท กปฺปโก รญฺโญ มฆเทวสฺส ปุตฺตสฺส ปฏิสฺสุตฺวา ตานิ ปลิตานิ สาธุกํ สณฺฑาเสน อุทฺธริตฺวา รญฺโญ มฆเทวสฺส ปุตฺตสฺส อญฺชลิสฺมิํ ปติฏฺฐาเปสิ.}

\prose{อถ โข อานนฺท รญฺโญ มฆเทวสฺส ปุตฺโต กปฺปกสฺส คามวรํ ทตฺวา เชฏฺฐปุตฺตํ กุมารํ อามนฺตาเปตฺวา เอตทโวจ “ปาตุภูตา โข เม ตาต กุมาร เทวทูตา, ทิสฺสนฺติ สิรสฺมิํ ปลิตานิ ชาตานิ, ภุตฺตา โข ปน เม มานุสกฺกา กามา, สมโย ทิพฺเพ กาเม ปริเยสิตุํ.\csromanspacer{} เอหิ ตฺวํ ตาต กุมาร อิมํ รชฺชํ ปฏิปชฺช, อหํ ปน เกสมสฺสุํ โอหาเรตฺวา กาสายานิ วตฺถานิ อจฺฉาเทตฺวา อคารสฺมา อนคาริยํ ปพฺพชิสฺสามิ.\csromanspacer{} เตน หิ ตาต กุมาร ยทา ตฺวํปิ ปสฺเสยฺยาสิ สิรสฺมิํ ปลิตานิ ชาตานิ, อถ กปฺปกสฺส คามวรํ ทตฺวา เชฏฺฐปุตฺตํ กุมารํ สาธุกํ รชฺเช สมนุสาสิตฺวา เกสมสฺสุํ โอหาเรตฺวา กาสายานิ วตฺถานิ อจฺฉาเทตฺวา อคารสฺมา อนคาริยํ ปพฺพเชยฺยาสิ, เยน เม อิทํ กลฺยาณํ วตฺตํ นิหิตํ อนุปฺปวตฺเตยฺยาสิ, มา โข เม ตฺวํ อนฺติมปุริโส อโหสิ.\csromanspacer{} ยสฺมิํ โข ตาต กุมาร ปุริสยุเค วตฺตมาเน เอวรูปสฺส กลฺยาณสฺส วตฺตสฺส สมุจฺเฉโท โหติ, โส เตสํ อนฺติมปุริโส โหติ, ตํ ตาหํ ตาต กุมาร เอวํ วทามิ ‘เยน เม อิทํ กลฺยาณํ วตฺตํ นิหิตํ อนุปฺปวตฺเตยฺยาสิ, มา โข เม ตฺวํ อนฺติมปุริโส อโหสี’ติ.\csromanspacer{} อถ โข อานนฺท รญฺโญ มฆเทวสฺส ปุตฺโต กปฺปกสฺส คามวรํ ทตฺวา เชฏฺฐปุตฺตํ กุมารํ สาธุกํ รชฺเช สมนุสาสิตฺวา อิมสฺมิํเยว มฆเทว-อมฺพวเน เกสมสฺสุํ โอหาเรตฺวา}

\vspace{\tipitakaparskip}
\csromancenter{มฆเทวสุตฺต (83) กาสายานิ วตฺถานิ อจฺฉาเทตฺวา อคารสฺมา อนคาริยํ ปพฺพชิ.\csromanspacer{} โส เมตฺตาสหคเตน เจตสา เอกํ ทิสํ ผริตฺวา วิหาสิ.\csromanspacer{} ตถา ทุติยํ.\csromanspacer{} ตถา ตติยํ.\csromanspacer{} ตถา จตุตฺถํ.\csromanspacer{} อิติ อุทฺธมโธ ติริยํ สพฺพธิ สพฺพตฺตตาย สพฺพาวนฺตํ โลกํ เมตฺตาสหคเตน เจตสา วิปุเลน มหคฺคเตน อปฺปมาเณน อเวเรน อพฺยาพชฺเฌน ผริตฺวา วิหาสิ.\csromanspacer{} กรุณาสหคเตน เจตสา.\csromanspacer{} มุทิตาสหคเตน เจตสา.\csromanspacer{} อุเปกฺขาสหคเตน เจตสา เอกํ ทิสํ ผริตฺวา วิหาสิ.\csromanspacer{} ตถา ทุติยํ.\csromanspacer{} ตถา ตติยํ.\csromanspacer{} ตถา จตุตฺถํ.\csromanspacer{} อิติ อุทฺธมโธ ติริยํ สพฺพธิ สพฺพตฺตตาย สพฺพาวนฺตํ โลกํ อุเปกฺขาสหคเตน เจตสา วิปุเลน มหคฺคเตน อปฺปมาเณน อเวเรน อพฺยาพชฺเฌน ผริตฺวา วิหาสิ.\csromanspacer{} รญฺโญ โข ปนานนฺท มฆเทวสฺส ปุตฺโต จตุราสีติวสฺสสหสฺสานิ กุมารกีฬิตํ กีฬิ, จตุราสีติวสฺสสหสฺสานิ โอปรชฺชํ กาเรสิ, จตุราสีติวสฺสสหสฺสานิ รชฺชํ กาเรสิ, จตุราสีติวสฺสสหสฺสานิ อิมสฺมิํเยว มฆเทว-อมฺพวเน อคารสฺมา อนคาริยํ ปพฺพชิโต พฺรหฺมจริยมจริ.\csromanspacer{} โส จตฺตาโร พฺรหฺมวิหาเร ภาเวตฺวา กายสฺส เภทา ปรํ มรณา พฺรหฺมโลกูปโค อโหสิ.}
\proseitem{311}{\csromanfolio{265}รญฺโญ โข ปนานนฺท มฆเทวสฺส ปุตฺตปปุตฺตกา ตสฺส ปรมฺปรา จตุราสีติราชสหสฺสานิ\footnote{จตุราสีติขตฺติยสหสฺสานิ (สี, อิ), จตุราสีติสหสฺสานิ (สฺยา, กํ)} อิมสฺมิํเยว มฆเทว-อมฺพเวน เกสมสฺสุํ โอหาเรตฺวา กาสายานิ วตฺถานิ อจฺฉาเทตฺวา อคารสฺมา อนคาริยํ ปพฺพชิํสุ.\csromanspacer{} เต เมตฺตาสหคเตน เจตสา เอกํ ทิสํ ผริตฺวา วิหริํสุ.\csromanspacer{} ตถา ทุติยํ.\csromanspacer{} ตถา ตติยํ.\csromanspacer{} ตถา จตุตฺตํ.\csromanspacer{} อิติ อุทฺธมโธ ติริยํ สพฺพธิ สพฺพตฺตตาย สพฺพาวนฺตํ โลกํ เมตฺตาสหคเตน เจตสา วิปุเลน มหคฺคเตน อปฺปมาเณน อเวเรน อพฺยาพชฺเฌน ผริตฺวา วิหริํสุ.\csromanspacer{} กรุณาสหคเตน เจตสา.\csromanspacer{} มุทิตาสหคเตน เจตสา.\csromanspacer{} อุเปกฺขาสหคเตน เจตสา เอกํ ทิสํ ผริตฺวา วิหริํสุ.\csromanspacer{} ตถา ทุติยํ.\csromanspacer{} ตถา ตติยํ.\csromanspacer{} ตถา จตุตฺถํ.\csromanspacer{} อิติ อุทฺธมโธ ติริยํ สพฺพธิ สพฺพตฺตตาย สพฺพาวนฺตํ โลกํ อุเปกฺขาสหคเตน เจตสา วิปุเลน มหคฺคเตน อปฺปมาเณน อเวเรน อพฺยาพชฺเฌน ผริตฺวา วิหริํสุ.\csromanspacer{} จตุราสีติวสฺสสหสฺสานิ กุมารกีฬิตํ กีฬิํสุ, จตุราสีติวสฺสสหสฺสานิ โอปรชฺชํ กาเรสุํ, จตุราสีติวสฺสสหสฺสานิ รชฺชํ กาเรสุํ, จตุราสีติวสฺสสหสฺสานิ อิมสฺมิํเยว มฆเทว-อมฺพวเน \csromanfolio{266}อคารสฺมา อนคาริยํ ปพฺพชิตา พฺรหฺมจริยมจริํสุ.\csromanspacer{} เต จตฺตาโร พฺรหฺมวิหาเร ภาเวตฺวา กายสฺส เภทา ปรํ มรณา พฺรหฺมโลกูปคา อเหสุํ.\csromanspacer{} นิมิ เตสํ ราชา\footnote{ราชานํ (สี, อิ)} ปจฺฉิมโก อโหสิ ธมฺมิโก ธมฺมราชา ธมฺเม ฐิโต มหาราชา, ธมฺมํ จรติ พฺราหฺมณคหปติเกสุ เนคเมสุ เจว ชานปเทสุ จ, อุโปสถญฺจ อุปวสติ จาตุทฺทสิํ ปญฺจทสิํ อฏฺฐมิํ จ ปกฺขสฺส.}

\proseitem{312}{ภูตปุพฺพํ อานนฺท เทวานํ ตาวติํสานํ สุธมฺมายํ สภายํ สนฺนิสินฺนานํ สนฺนิปติตานํ อยมนฺตรากถา อุทปาทิ “ลาภา วต โภ วิเทหานํ, สุลทฺธํ วต โภ วิเทหานํ, เยสํ นิมิ ราชา ธมฺมิโก ธมฺมราชา ธมฺเม ฐิโต มหาราชา, ธมฺมํ จรติ พฺราหฺมณคหปติเกสุ เนคเมสุ เจว ชานปเทสุ จ, อุโปสถญฺจ อุปวสติ จาตุทฺทสิํ ปญฺจทสิํ อฏฺฐมิํ จ ปกฺขสฺสา”ติ.\csromanspacer{} อถ โข อานนฺท สกฺโก เทวานมินฺโท เทเว ตาวติํเส อามนฺเตสิ “อิจฺเฉยฺยาถ โน ตุเมฺห มาริสา นิมิํ ราชานํ ทฏฺฐุนฺ”ติ.\csromanspacer{} อิจฺฉาม มยํ มาริส นิมิํ ราชานํ ทฏฺฐุนฺติ.\csromanspacer{} เตน โข ปน อานนฺท สมเยน นิมิ ราชา ตทหุโปสเถ ปนฺนรเส สีสํนฺหาโต\footnote{สสีสํ นหาโต (สี), สีสนฺหาโต (สฺยา, กํ)} อุโปสถิโก อุปริปาสาทวรคโต นิสินฺโน โหติ.\csromanspacer{} อถ โข อานนฺท สกฺโก เทวานมินฺโท เสยฺยถาปิ นาม พลวา ปุริโส สมิญฺชิตํ วา พาหํ ปสาเรยฺย, ปสาริตํ วา พาหํ สมิญฺเชยฺย.\csromanspacer{} เอวเมว เทเวสุ ตาวติํเสสุ อนฺตรหิโต นิมิสฺส รญฺโญ ปมุเข ปาตุรโหสิ.\csromanspacer{} อถ โข อานนฺท สกฺโก เทวานมินฺโท นิมิํ ราชานํ เอตทโวจ “ลาภา เต มหาราช, สุลทฺธํ เต มหาราช,เทวา มหาราช ตาวติํสา สุธมฺมายํ สภายํ กิตฺตยมานรูปา สนฺนิสินฺนา ‘ลาภา วต โภ วิเทหานํ, สุลทฺธํ วต โภ วิเทหานํ, เยสํ นิมิ ราชา ธมฺมิโก ธมฺมราชา ธมฺเม ฐิโต มหาราชา, ธมฺมํ จรติ พฺราหฺมณคหปติเกสุ เนคเมสุ เจว ชานปเทสุ จ, อุโปสถญฺจ อุปวสติ จาตุทฺทสิํ ปญฺจทสิํ อฏฺฐมิํ จ ปกฺขสฺสา’ติ.\csromanspacer{} เทวา เต มหาราช ตาวติํสา ทสฺสนกามา, ตสฺส เต อหํ มหาราช สหสฺสยุตฺตํ อาชญฺญรถํ ปหิณิสฺสามิ, อภิรุเหยฺยาสิ มหาราช ทิพฺพํ ยานํ อวิกมฺปมาโน”ติ.\csromanspacer{} อธิวาเสสิ โข อานนฺท นิมิ ราชา ตุณฺหีภาเวน.}

\csromanheader{2}{3. มฆเทวสุตฺต (83)}
\proseitem{313}{\csromanfolio{267}อถ โข อานนฺท สกฺโก เทวานมินฺโท นิมิสฺส รญฺโญ อธิวาสนํ วิทิตฺวา เสยฺยถาปิ นาม พลวา ปุริโส สมิญฺชิตํ วา พาหํ ปสาเรยฺย, ปสาริตํ วา พาหํ สมิญฺเชยฺย.\csromanspacer{} เอวเมว นิมิสฺส รญฺโญ ปมุเข อนฺตรหิโต เทเวสุ ตาวติํเสสุ ปาตุรโหสิ.\csromanspacer{} อถ โข อานนฺท สกฺโก เทวานมินฺโท มาตลิํ สงฺคาหกํ อามนฺเตสิ “เอหิ ตฺวํ สมฺม มาตลิ สหสฺสยุตฺตํ อาชญฺญรถํ โยเชตฺวา นิมิํ ราชานํ อุปสงฺกมิตฺวา เอวํ วเทหิ ‘อยํ เต มหาราช สหสฺสยุตฺโต อาชญฺญรโถ สกฺเกน เทวานมินฺเทน เปสิโต, อภิรุเหยฺยาสิ มหาราช ทิพฺพํ ยานํ อวิกมฺปมาโน’ติ”. “เอวํ ภทฺทนฺตวา”ติ โข อานนฺท มาตลิ สงฺคาหโก สกฺกสฺส เทวานมินฺทสฺส ปฏิสฺสุตฺวา สหสฺสยุตฺตํ อาชญฺญรถํ โยเชตฺวา นิมิํ ราชานํ อุปสงฺกมิตฺวา เอตทโวจ “อยํ เต มหาราช สหสฺสยุตฺโต อาชญฺญรโถ สกฺเกน เทวานมินฺเทน เปสิโต, อภิรุห มหาราช ทิพฺพํ ยานํ อวิกมฺปมาโน, อปิ จ มหาราช กตเมน ตํ เนมิ, เยน วา ปาปกมฺมา ปาปกานํ กมฺมานํ วิปากํ ปฏิสํเวเทนฺติ, เยน วา กลฺยาณกมฺมา กลฺยาณกมฺมานํ วิปากํ ปฏิสํเวเทนฺตี”ติ.\csromanspacer{} อุภเยเนว มํ มาตลิ เนหีติ.\csromanspacer{} สมฺปเวเสสิ\footnote{สมฺปาเปสิ (สี, อิ)} โข อานนฺท มาตลิ สงฺคหโก นิมิํ ราชานํ สุธมฺมํ สภํ.\csromanspacer{} อทฺทสา โข อานนฺท สกฺโก เทวานมินฺโท นิมิํ ราชานํ ทูรโตว อาคจฺฉนฺตํ, ทิสฺวาน นิมิํ ราชา นํ เอตทโวจ “เอหิ โข มหาราช, สฺวาคตํ มหาราช, เทวา เต ทสฺสนกามา มหาราช, ตาวติํสา สุธมฺมายํ สภายํ กิตฺตยมานรูปา สนฺนิสินฺนา ‘ลาภา วต โภ วิเทหานํ, สุลทฺธํ วต โภ วิเทหานํ, เยสํ นิมิ ราชา ธมฺมิโก ธมฺมราชา ธมฺเม ฐิโต มหาราชา, ธมฺมํ จรติ พฺราหฺมณคหปติเกสุ เนคเมสุ เจว ชานปเทสุ จ, อุโปสถญฺจ อุปวสติ จาตุทฺทสิํ ปญฺจทสิํ อฏฺฐมิํ จ ปกฺขสฺสา’ติ, เทวา เต มหาราช ตาวติํสา ทสฺสนกามา, อภิรม มหาราช เทวสุ เทวานุภาเวนา”ติ.\csromanspacer{} อลํ มาริส, ตตฺเถว มํ มิถิลํ ปฏิเนตุ, ตถาหํ ธมฺมํ จริสฺสามิ พฺราหฺมณคหปติเกสุ เนคเมสุ เจว ชานปเทสุ จ, อุโปสถญฺจ อุปวสามิ จาตุทฺทสิํ ปญฺจทสิํ อฏฺฐมิํ จ ปกฺขสฺสาติ.}

\proseitem{314}{อถ โข อานนฺท สกฺโก เทวานมินฺโท มาตลิํ สงฺคาหกํ อามนฺเตสิ “เอหิ ตฺวํ สมฺม มาตลิ สหสฺสยุตฺตํ อาชญฺญรถํ โยเชตฺวา \csromanfolio{268}นิมิํ ราชานํ ตตฺเถว มิถิลํ ปฏิเนหี”ติ. “เอวํ ภทฺทนฺตวา”ติ โข อานนฺท มาตลิ สงฺคาหโก สกฺกสฺส เทวานมินฺทสฺส ปฏิสฺสุตฺวา สหสฺสยุตฺตํ อาชญฺญรถํ โยเชตฺวา นิมิํ ราชานํ ตตฺเถว มิถิลํ ปฏิเนสิ.\csromanspacer{} ตตฺร สุทํ อานนฺท นิมิ ราชา ธมฺมํ จรติ พฺราหฺมณคหปติเกสุ เนคเมสุ เจว ชานปเทสุ จ, อุโปสถญฺจ อุปวสติ จาตุทฺทสิํ ปญฺจทสิํ อฏฺฐมิํ จ ปกฺขสฺสาติ.\csromanspacer{} อถ โข อานนฺท นิมิ ราชา พหูนํ วสฺสานํ พหูนํ วสฺสสตานํ พหูนํ วสฺสสหสฺสานํ อจฺจเยน กปฺปกํ อามนฺเตสิ “ยทา เม สมฺม กปฺปก ปสฺเสยฺยาสิ สิรสฺมิํ ปลิตานิ ชาตานิ, อถ เม อาโรเจยฺยาสี”ติ. “เอวํ เทวา”ติ โข อานนฺท กปฺปโก นิมิสฺส รญฺโญ ปจฺจสฺโสสิ.\csromanspacer{} อทฺทสา โข อานนฺท กปฺปโก พหูนํ วสฺสานํ พหูนํ วสฺสสตานํ พหูนํ วสฺสสหสฺสานํ อจฺจเยน นิมิสฺส รญฺโญ สิรสฺมิํ ปลิตานิ ชาตานิ, ทิสฺวาน นิมิํ ราชานํ เอตทโวจ “ปาตุภูตา โข เทวสฺส เทวทูตา, ทิสฺสนฺติ สิรสฺมิํ ปลิตานิ ชาตานี”ติ.\csromanspacer{} เตน หิ สมฺม กปฺปก ตานิ ปลิตานิ สาธุกํ สณฺฑาเสน อุทฺธริตฺวา มม อญฺชลิสฺมิํ ปติฏฺฐาเปหีติ. “เอวํ เทวา”ติ โข อานนฺท กปฺปโก นิมิสฺส รญฺโญ ปฏิสฺสุตฺวา ตานิ ปลิตานิ สาธุกํ สณฺฑาเสน อุทฺธริตฺวา นิมิสฺส รญฺโญ อญฺชลิสฺมิํ ปติฏฺฐาเปสิ.\csromanspacer{} อถ โข อานนฺท นิมิ ราชา กปฺปกสฺส คามวรํ ทตฺวา เชฏฺฐปุตฺตํ กุมารํ อามนฺตาเปตฺวา เอตทโวจ “ปาตุภูตา โข เม ตาต กุมาร เทวทูตา, ทิสฺสนฺติ สิรสฺมิํ ปลิตานิ ชาตานิ, ภุตฺตา โข ปน เม มานุสกา กามา, สมโย ทิพฺเพ กาเม ปริเยสิตุํ.\csromanspacer{} เอหิ ตฺวํ ตาต กุมาร อิมํ รชฺชํ ปฏิปชฺช.\csromanspacer{} อหํ ปน เกสมสฺสุํ โอหาเรตฺวา กาสายานิ วตฺถานิ อจฺฉาเทตฺวา อคารสฺมา อนคาริยํ ปพฺพชิสฺสามิ.\csromanspacer{} เตน หิ ตาต กุมาร ยทา ตฺวมฺปิ ปสฺเสยฺยาสิ สิรสฺมิํ ปลิตานิ ชาตานิ, อถ กปฺปกสฺส คามวรํ ทตฺวา เชฏฺฐปุตฺตํ กุมารํ สาธุกํ รชฺเช สมนุสาสิตฺวา เกสมสฺสุํ โอหาเรตฺวา กาสายานิ วตฺถานิ อจฺฉาเทตฺวา อคารสฺมา อนคาริยํ ปพฺพเชยฺยาสิ.\csromanspacer{} เยน เม อิทํ กลฺยาณํ วตฺตํ นิหิตํ อนุปฺปวตฺเตยฺยาสิ, มา โข เม ตฺวํ อนฺติมปุริโส อโหสิ.\csromanspacer{} ยสฺมิํ โข ตาต กุมาร ปุริสยุเค วตฺตมาเน เอวรูปสฺส กลฺยาณสฺส วตฺตสฺส สมุจฺเฉโท โหติ, โส เตสํ อนฺติมปุริโส โหติ.\csromanspacer{} ตํ ตาหํ ตาต กุมาร เอวํ วทามิ ‘เยน เม อิทํ กลฺยาณํ วตฺตํ นิหิตํ อนุปฺปวตฺเตยฺยาสิ, มา โข เม ตฺวํ อนฺติมปุริโส อโหสี’ติ”.}

\csromanheader{2}{3. มฆเทวสุตฺต (83)}
\proseitem{315}{\csromanfolio{269}อถ โข อานนฺท นิมิ ราชา กปฺปกสฺส คามวรํ ทตฺวา เชฏฺฐปุตฺตํ กุมารํ สาธุกํ รชฺเช สมนุสาสิตฺวา อิมสฺมิํเยว มฆเทว-อมฺพวเน เกสมสฺสุํ โอหาเรตฺวา กาสายานิ วตฺถานิ อจฺฉาเทตฺวา อคารสฺมา อนคาริยํ ปพฺพชิ.\csromanspacer{} โส เมตฺตาสหคเตน เจตสา เอกํ ทิสํ ผริตฺวา วิหาสิ.\csromanspacer{} ตถา ทุติยํ.\csromanspacer{} ตถา ตติยํ.\csromanspacer{} ตถา จตุตฺถํ.\csromanspacer{} อิติ อุทฺธมโธ ติริยํ สพฺพธิ สพฺพตฺตตาย สพฺพาวนฺตํ โลกํ เมตฺตาสหคเตน เจตสา วิปุเลน มหคฺคเตน อปฺปมาเณน อเวเรน อพฺยาพชฺเฌน ผริตฺวา วิหาสิ.\csromanspacer{} กรุณาสหคเตน เจตสา.\csromanspacer{} มุทิตาสหคเตน เจตสา.\csromanspacer{} อุเปกฺขาสหคเตน เจตสา เอกํ ทิสํ ผริตฺวา วิหาสิ.\csromanspacer{} ตถา ทุติยํ.\csromanspacer{} ตถา ตติยํ.\csromanspacer{} ตถา จตุตฺถํ.\csromanspacer{} อิติ อุทฺธมโธ ติริยํ สพฺพธิ สพฺพตฺตตาย สพฺพาวนฺตํ โลกํ อุเปกฺขาสหคเตน เจตสา วิปุเลน มหคฺคเตน อปฺปมาเณน อเวเรน อพฺยาพชฺเฌน ผริตฺวา วิหาสิ.\csromanspacer{} นิมิ โข ปนานนฺท ราชา จตุราสีติวสฺสสหสฺสานิ กุมารกีฬิตํ กีฬิ, จตุราสีติวสฺสสหสฺสานิ โอปรชฺชํ กาเรสิ, จตุราสีติวสฺสสหสฺสานิ รชฺชํ กาเรสิ, จตุราสีติวสฺสสหสฺสานิ อิมสฺมิํเยว มฆเทว-อมฺพวเน อคารสฺมา อนคาริยํ ปพฺพชิโต พฺรหฺมจริยมจริ.\csromanspacer{} โส จตฺตาโร พฺรหฺมวิหาเร ภาเวตฺวา กายสฺส เภทา ปรํ มรณา พฺรหฺมโลกูปโค อโหสิ.\csromanspacer{} นิมิสฺส โข ปนานนฺท รญฺโญ กฬารชนโก นาม ปุตฺโต อโหสิ, น โส อคารสฺมา อนคาริยํ ปพฺพชิ.\csromanspacer{} โส ตํ กลฺยาณํ วตฺตํ สมุจฺฉินฺทิ, โส เตสํ อนฺติมปุริโส อโหสิ.}

\csromanmatikaanchor{mtk.40}
\csromantocmarkref{subsection}{3. มฆเทวสุตฺต}{mtk.40}
\bookmark[dest={mtk.40},level=2]{3. มฆเทวสุตฺต}
\proseitem{316}{สิยา โข ปน เต อานนฺท เอวมสฺส “อญฺโญ นูน เตน สมเยน ราชา มฆเทโว อโหสิ, เยน ตํ กลฺยาณํ วตฺตํ นิหิตนฺ”ติ\footnote{โย ตํ กลฺยาณํ วตฺตํ นิหินีติ (สี) ( ) นตฺถิ (ก)}.\csromanspacer{} น โข ปเนตํ อานนฺท เอวํ ทฏฺฐพฺพํ.\csromanspacer{} อหํ เตน สมเยน ราชา มฆเทโว อโหสิํ, (อหํ ตํ กลฺยาณํ วตฺตํ นิหินิํ,) มยา ตํ กลฺยาณํ วตฺตํ นิหิตํ ปจฺฉิมา ชนตา อนุปฺปวตฺเตสิ.\csromanspacer{} ตํ โข ปนานนฺท กลฺยาณํ วตฺตํ น นิพฺพิทาย น วิราคาย น นิโรธาย น อุปสมาย น อภิญฺญาย น สมฺโพธาย น นิพฺพานาย สํวตฺตติ ยาวเทว พฺรหฺมโลกูปปตฺติยา.\csromanspacer{} อิทํ โข ปนานนฺท เอตรหิ มยา กลฺยาณํ วตฺตํ นิหิตํ เอกนฺตนิพฺพิทาย วิราคาย นิโรธาย อุปสมาย อภิญฺญาย สมฺโพธาย นิพฺพานาย สํวตฺตติ.\csromanspacer{} กตมญฺจานนฺท เอตรติ มยา กลฺยาณํ \csromanfolio{270}วตฺตํ นิหิตํ เอกนฺตนิพฺพิทาย วิราคาย นิโรธาย อุปสมาย อภิญฺญาย สมฺโพธาย นิพฺพานาย สํวตฺตติ.\csromanspacer{} อยเมว อริโย อฏฺฐงฺคิโก มคฺโค.\csromanspacer{} เสยฺยถิทํ, สมฺมาทิฏฺฐิ สมฺมาสงฺกปฺโป สมฺมาวาจา สมฺมากมฺมนฺโต สมฺมา อาชีโว สมฺมาวายาโม สมฺมาสติ สมฺมาสมาธิ.\csromanspacer{} อิทํ โข อานนฺท เอตรหิ มยา กลฺยาณํ วตฺตํ นิหิตํ เอกนฺตนิพฺพิทาย วิราคาย นิโรธาย อุปสมาย อภิญฺญาย สมฺโพธาย นิพฺพานาย สํวตฺตติ.\csromanspacer{} ตํ โว อหํ อานนฺท เอวํ วทามิ, เยน เม อิทํ กลฺยาณํ วตฺตํ นิหิตํ อนุปฺปวตฺเตยฺยาถ, มา โข เม ตุเมฺห อนฺติมปุริสา อหุวตฺถ.\csromanspacer{} ยสฺมิํ โข อานนฺท ปุริสยุเค วตฺตมาเน เอวรูปสฺส กลฺยาณสฺส วตฺตสฺส สมุจฺเฉโท โหติ, โส เตสํ อนฺติมปุริโส โหติ.\csromanspacer{} ตํ โว อหํ อานนฺท เอวํ วทามิ “เยน เม อิทํ กลฺยาณํ วตฺตํ นิหิตํ อนุปฺปวตฺเตยฺยาถ, มา โข เม ตุเมฺห อนฺติมปุริสา อหุวตฺถา”ติ.}

\prose{อิทมโวจ ภควา.\csromanspacer{} อตฺตมโน อายสฺมา อานนฺโท ภควโต ภาสิตํ อภินนฺทีติ.}

\nitthitamruled{มฆเทวสุตฺตํ นิฏฺฐิตํ ตติยํ.}
\csromanheader{2}{4. มธุรสุตฺต}
\proseitem{317}{เอวํ เม สุตํ\csromandash{}เอกํ สมยํ อายสฺมา มหากจฺจาโน มธุรายํ วิหรติ คุนฺทาวเน.\csromanspacer{} อสฺโสสิ โข ราชา มาธุโร อวนฺติปุตฺโต “สมโณ ขลุ โภ กจฺจาโน มธุรายํ\footnote{มถุรายํ (ฏีกา)} วิหรติ คุนฺทาวเน, ตํ โข ปน ภวนฺตํ กจฺจานํ เอวํ กลฺยาโณ กิตฺติสทฺโท อพฺภุคฺคโต ‘ปณฺฑิโต วิยตฺโต เมธาวี พหุสฺสุโต จิตฺตกถี กลฺยาณปฏิภาโน วุทฺโธ เจว อรหา จ, สาธุ โข ปน ตถารูปานํ อรหตํ ทสฺสนํ โหตี’ติ”.\csromanspacer{} อถ โข ราชา มาธุโร อวนฺติปุตฺโต ภทฺรานิ ภทฺรานิ ยานานิ โยชาเปตฺวา ภทฺรํ ยานํ อภิรุหิตฺวา ภเทฺรหิ ภเทฺรหิ ยาเนหิ มธุราย นิยฺยาสิ มหจฺจราชานุภาเวน อายสฺมนฺตํ มหากจฺจานํ ทสฺสนาย, ยาวติกา ยานสฺส ภูมิ, ยาเนน คนฺตฺวา ยานา ปจฺโจโรหิตฺวา ปตฺติโกว เยนายสฺมา มหากจฺจาโน เตนุปสงฺกมิ,}

\vspace{\tipitakaparskip}
\csromancenter{มธุรสุตฺต (84) อุปสงฺกมิตฺวา อายสฺมตา มหากจฺจาเนน สทฺธิํ สมฺโมทิ, สมฺโมทนียํ กถํ สารณียํ วีติสาเรตฺวา เอกมนฺตํ นิสีทิ, เอกมนฺตํ นิสินฺโน โข ราชา มาธุโร อวนฺติปุตฺโต อายสฺมนฺตํ มหากจฺจานํ เอตทโวจ “พฺราหฺมณา โภ กจฺจาน เอวมาหํสุ ‘พฺราหฺมโณว เสฏฺโฐ วณฺโณ, หีโน อญฺโญ วณฺโณ.\csromanspacer{} พฺราหฺมโณว สุกฺโก วณฺโณ, กโณฺห อญฺโญ วณฺโณ.\csromanspacer{} พฺราหฺมณาว สุชฺฌนฺติ โน อพฺราหฺมณา.\csromanspacer{} พฺราหฺมณาว พฺรหฺมุโน ปุตฺตา โอรสา มุขโต ชาตา พฺรหฺมชา พฺรหฺมนิมฺมิตา พฺรหฺมทายาทา’ติ.\csromanspacer{} อิธ ภวํ กจฺจาโน กิมกฺขายี”ติ.\csromanspacer{} โฆโสเยว โข เอโส มหาราช โลกสฺมิํ “พฺราหฺมโณว เสฏฺโฐ วณฺโณ, หีโน อญฺโญ วณฺโณ.\csromanspacer{} พฺราหฺมโณว สุกฺโก วณฺโณ, กโณฺห อญฺโญ วณฺโณ.\csromanspacer{} พฺราหฺมณาว สุชฺฌนฺติ โน อพฺราหฺมณา.\csromanspacer{} พฺราหฺมณาว พฺรหฺมุโน ปุตฺตา โอรสา มุขโต ชาตา พฺรหฺมชา พฺรหฺมนิมฺมิตา พฺรหฺมทายาทา”ติ.\csromanspacer{} ตทมินาเปตํ มหาราช ปริยาเยน เวทิตพฺพํ, ยถา “โฆโสเยเวโส โลกสฺมิํ, พฺราหฺมโณว เสฏฺโฐ วณฺโณ, หีโน อญฺโญ วณฺโณ ฯเปฯ พฺรหฺมทายาทา”ติ.}
\proseitem{318}{\csromanfolio{271}ตํ กิํ มญฺญสิ มหาราช, ขตฺติยสฺส เจปิ อิชฺเฌยฺย ธเนน วา ธญฺเญน วา รชเตน วา ชาตรูเปน วา, ขตฺติโยปิสฺสาสฺส ปุพฺพุฏฺฐายี ปจฺฉานิปาตี กิํการปฏิสฺสาวี มนาปจารี วิยวาที.\csromanspacer{} พฺราหฺมโณปิสฺสาสฺส.\csromanspacer{} เวสฺโสปิสฺสาสฺส.\csromanspacer{} สุทฺโทปิสฺสาสฺส ปุพฺพุฏฺฐายี ปจฺฉานิปาตี กิํการปฏิสฺสาวี มนาปจารี ปิยวาทีติ? ขตฺติยสฺส เจปิ โภ กจฺจาน อิชฺเฌยฺย ธเนน วา ธญฺเญน วา รชเตน วา ชาตรูเปน วา.\csromanspacer{} ขตฺถิโยปิสฺสาสฺส ปุพฺพุฏฺฐายี ปจฺฉานิปาตี กิํการปฏิสฺสาวี มนาปจารี ปิยวาที.\csromanspacer{} พฺราหฺมโณปิสฺสาสฺส.\csromanspacer{} เวสฺโสปิสฺสาสฺส.\csromanspacer{} สุทฺโทปิสฺสาสฺส ปุพฺพุฏฺฐายี ปจฺฉานิปาตี กิํการปฏิสฺสาวี มนาปจารี ปิยวาทีติ.}

\prose{ตํ กิํ มญฺญสิ มหาราช, พฺราหฺมณสฺส เจปิ อิชฺเฌยฺย ธเนน วา ธญฺเญน วา รชเตน วา ชาตรูเปน วา, พฺราหฺมโณปิสฺสาสฺส ปุพฺพุฏฺฐายี ปจฺฉานิปาตี กิํการปฏิสฺสาวี มนาปจารี ปิยวาที.\csromanspacer{} เวสฺโสปิสฺสาสฺส.\csromanspacer{} สุทฺโท ปิสฺสาสฺส.\csromanspacer{} ขตฺติโยปิสฺสาสฺส ปุพฺพุฏฺฐายี ปจฺฉานิปาตี กิํการปฏิสฺสาวี มนาปสารี ปิยวาทีติ? พฺราหฺมณสฺส เจปิ โภ กจฺจาน อิชฺเฌยฺย ธเนน วา ธญฺเญน วา รชเตน วา ชาตรูเปน วา, พฺราหฺมโณปิสฺสาสฺส ปุพฺพุฏฺฐายี ปจฺฉานิปาตี กิํการปฏิสฺสาวี มนาปจารี ปิยวาที.\csromanspacer{} เวสฺโสปิสฺสาสฺส. \csromanfolio{272}สุทฺโทปิสฺสาสฺส.\csromanspacer{} ขตฺติโยปิสฺสาสฺส ปุพฺพุฏฺฐายี ปจฺฉานิปาตี กิํการปฏิสฺสาวี มนาปจารี ปิยวาทีติ.}

\prose{ตํ กิํ มญฺญสิ มหาราช, เวสฺสสฺส เจปิ อิชฺเฌยฺย ธเนน วา ธญฺเญน วา รชเตน วา ชาตรูเปน วา, เวสฺโสปิสฺสาสฺส ปุพฺพุฏฺฐายี ปจฺฉานิปาตี กิํการปฏิสฺสาวี มนาปจารี ปิยวาที.\csromanspacer{} สุทฺโทปิสฺสาสฺส.\csromanspacer{} ขตฺถิโยปิสฺสาสฺส.\csromanspacer{} พฺราหฺมโณปิสฺสาสฺส ปุพฺพุฏฺฐายี ปจฺฉานิปาตี กิํการปฏิสฺสาวี มนาปจารี ปิยวาทีติ? เวสฺสสฺส เจปิ โภ กจฺจาน อิชฺเฌยฺย ธเนน วา ธญฺเญน วา รชเตน วา ชาตรูเปน วา, เวสฺโสปิสฺสาสฺส ปุพฺพุฏฺฐายี ปจฺฉานิปาตี กิํการปฏิสฺสาวี มนาปจารี ปิยวาที.\csromanspacer{} สุทฺโทปิสฺสาสฺส.\csromanspacer{} ขตฺติโยปิสฺสาสฺส.\csromanspacer{} พฺราหฺมโณปิสฺสาสฺส ปุพฺพุฏฺฐายี ปจฺฉานิปาตี กิํการปฏิสฺสาวี มนาปจารี ปิยวาทีติ.}

\prose{ตํ กิํ มญฺญสิ มหาราช, สุทฺทสฺส เจปิ อิชฺเฌยฺย ธเนน วา ธญฺเญน วา รชเตน วา ชาตรูเปน วา.\csromanspacer{} สุทฺโธปิสฺสาสฺส ปุพฺพุฏฺฐายี ปจฺฉานิปาตี กิํการปฏิสฺสาวี มนาปจารี ปิยวาที.\csromanspacer{} ขตฺติโยปิสฺสาสฺส.\csromanspacer{} พฺราหฺมโณปิสฺสาสฺส.\csromanspacer{} เวสฺโสปิสฺสาสฺส ปุพฺพุฏฺฐายี ปจฺฉานิปาตี กิํการปฏิสฺสาวี มนาปจารี ปิยวาทีติ? สุทฺทสฺส เจติ โภ กจฺจาน อิชฺเฌยฺย ธเนน วา ธญฺเญน วา รชเตน วา ชาตรูเปน วา, สุทฺโทปิสฺสาสฺส ปุพฺพุฏฺฐายี ปจฺฉานิปาตี กิํการปฏิสฺสาวี มนาปจารี ปิยวาทีติ.\csromanspacer{} ขตฺติโยปิสฺสาสฺส.\csromanspacer{} พฺราหฺมโณปิสฺสาสฺส.\csromanspacer{} เวสฺโสปิสฺสาสฺส ปุพฺพุฏฺฐายี ปจฺฉานิปาตี กิํการปฏิสฺสาวี มนาปจารี ปิยวาทีติ.}

\prose{ตํ กิํ มญฺญสิ มหาราช, ยทิ เอวํ สนฺเต อิเม จตฺตาโร วณฺณา สมสมา โหนฺติ, โน วา, กถํ วา เต เอตฺถ โหตีติ.\csromanspacer{} อทฺธา โข โภ กจฺจาน เอวํ สนฺเต อิเม จตฺตาโร วณฺณา สมสมา โหนฺติ, เนสํ\footnote{นา สํ (สี), นาหํ (สฺยา, กํ)} เอตฺถ กิญฺจิ นานากรณํ สมนุปสฺสามีติ.\csromanspacer{} อิมินาปิ โข เอตํ มหาราช ปริยาเยน เวทิตพฺพํ, ยถา “โฆโสเยเวโส โลกสฺมิํ, พฺราหฺมโณว เสฏฺโฐ วณฺโณ, หีโน อญฺโญ วณฺโณ ฯเปฯ พฺรหฺมทายาทา”ติ.}

\csromanheader{2}{4. มธุรสุตฺต (84)}
\proseitem{319}{\csromanfolio{273}ตํ กิํ มญฺญสิ มหาราช, อิธสฺส ขตฺติโย ปาณาติปาตี อทินฺนาทายี กาเมสุมิจฺฉาจารี มุสาวาที ปิสุณวาโจ ผรุสวาโจ สมฺผปฺปลาปี อภิชฺฌาลุ พฺยาปนฺนจิตฺโต มิจฺฉาทิฏฺถิ\footnote{มิจฺฉาทิฏฺฐี (สพฺพตฺถ)}, กายสฺส เภทา ปรํ มรณา อปายํ ทุคฺคติํ วินิปาตํ นิรยํ อุปปชฺเชยฺย, โน วา, กถํ วา เต เอตฺถ โหตีติ.\csromanspacer{} ขตฺติโยปิ หิ โภ กจฺจาน ปาณาติปาตี อทินฺนาทายี กาเมสุมิจฺฉาจารี มุสาวาที ปิสุณาวาโจ ผรุสวาโจ สมฺผปฺปลาปี อภิชฺฌาลุ พฺยาปนฺนจิตฺโต มิจฺฉาทิฏฺฐิ, กายสฺส เภทา ปรํ มรณา อปายํ ทุคฺคติํ วินิปาตํ นิรยํ อุปปชฺเชยฺย, เอวํ เม เอตฺถ โหติ, เอวญฺจ ปน เม เอตํ อรหตํ สุตนฺติ.}

\prose{สาธุ สาธุ มหาราช, สาธุ โข เต เอตํ มหาราช เอวํ โหติ, สาธุ จ ปน เต เอตํ อรหตํ สุตํ.\csromanspacer{} ตํ กิํ มญฺญสิ มหาราช, อิธสฺส พฺราหฺมโณ ฯเปฯ อิธสฺส เวสฺโส ฯเปฯ อิธสฺส สุทฺโท ปาณาติปาตี อทินฺนาทายี ฯเปฯ มิจฺฉาทิฏฺฐิ, กายสฺส เภทา ปรํ มรณา อปายํ ทุคฺคติํ วินิปาตํ นิรยํ อุปปชฺเชยฺย, โน วา, กถํ วา เต เอตฺถ โหตีติ.\csromanspacer{} สุทฺโทปิ หิ โภ กจฺจาน ปาณาติปาตี อทินฺนาทายี ฯเปฯ มิจฺฉาทิฏฺฐิ, กายสฺส เภทา ปรํ มรณา อปายํ ทุคฺคติํ วินิปาตํ นิรยํ อุปปชฺเชยฺย, เอวํ เม เอตฺถ โหติ, เอวญฺจ ปน เม เอตํ อรหตํ สุตนฺติ.}

\prose{สาธุ สาธุ มหาราช, สาธุ โข เต เอตํ มหาราช เอวํ โหติ, สาธุ จ ปน เต เอตํ อรหตํ สุตํ.\csromanspacer{} ตํ กิํ มญฺญสิ มหาราช, ยทิ เอวํ สนฺเต อิเม จตฺตาโร วณฺณา สมสมา โหนฺติ, โน วา, กถํ วา เต เอตฺถ โหตีติ.\csromanspacer{} อทฺธา โข โภ กจฺจาน เอวํ สนฺเต อิเม จตฺตาโร วณฺณา สมสมา โหนฺติ, เนสํ เอตฺถ กิญฺจิ นานากรณํ สมนุปสฺสามีติ.\csromanspacer{} อิมินาปิ โข เอตํ มหาราช ปริยาเยน เวทิตพฺพํ, ยถา “โฆโสเยเวโส โลกสฺมิํ, พฺราหฺมโณว เสฏฺโฐ วณฺโณ, หีโน อญฺโญ วณฺโณ ฯเปฯ พฺรหฺมทายาทา”ติ.}

\proseitem{320}{ตํ กิํ มญฺญสิ มหาราช, อิธสฺส ขตฺติโย ปาณาติปาตา ปฏิวิรโต, อทินฺนาทานา ปฏิวิรโต, กาเมสุมิจฺฉาจารา ปฏิวิรโต, มุสาวาทา ปฏิวิรโต, ปิสุณาย วาจาย ปฏิวิรโต, ผรุสาย วาจาย ปฏิวิรโต, สมฺผปฺปลาปา ปฏิวิรโต, อนภิชฺฌาลุ อพฺยาปนฺนจิตฺโต \csromanfolio{274}สมฺมาทิฏฺฐิ\footnote{สมฺมาทิฏฺฐี (สฺยา, กํ, อิ, ก)}, กายสฺส เภทา ปรํ มรณา สุคติํ สคฺคํ โลกํ อุปปชฺเชยฺย โน วา, กถํ วา เต เอตฺถ โหตีติ.\csromanspacer{} ขตฺถิโยปิ หิ โภ กจฺจาน ปาณาติปาตา ปฏิวิรโต, อทินฺนาทานา ปฏิวิรโต, กาเมสุมิจฺฉาจารา ปฏิวิรโต, มุสาวาทา ปฏิวิรโต, ปิสุณาย วาจาย ปฏิวิรโต, ผรุสาย วาจาย ปฏิวิรโต, สมฺผปฺปลาปา ปฏิวิรโต, อนภิชฺฌาลุ อพฺยาปนฺนจิตฺโต สมฺมาทิฏฺฐิ, กายสฺส เภทา ปรํ มรณา สุคติํ สคฺคํ โลกํ อุปปชฺเชยฺย, เอวํ เม เอตฺถ โหติ, เอวญฺจ ปน เม เอตํ อรหตํ สุตนฺติ.}

\prose{สาธุ สาธุ มหาราช, สาธุ โข เต เอตํ มหาราช เอวํ โหติ, สาธุ จ ปน เต เอตํ อรหตํ สุตํ.\csromanspacer{} ตํ กิํ มญฺญสิ มหาราช, อิธสฺส พฺราหฺมโณ.\csromanspacer{} อิธสฺส เวสฺโส.\csromanspacer{} อิธสฺส สุทฺโท ปาณาติปาตา ปฏิวิรโต, อทินฺนาทานา ปฏิวิรโต ฯเปฯ สมฺมาทิฏฺฐิ, กายสฺส เภทา ปรํ มรณา สุคติํ สคฺคํ โลกํ อุปปชฺเชยฺย, โน วา, กถํ วา เต เอตฺถ โหตีติ.\csromanspacer{} สุทฺโทปิ หิ โภ กจฺจาน ปาณาติปาตา ปฏิวิรโต, อทินฺนาทานา ปฏิวิรโต ฯเปฯ สมฺมาทิฏฺฐิ, กายสฺส เภทา ปรํ มรณา สุคติํ สคฺคํ โลกํ อุปปชฺเชยฺย, เอวํ เม เอตฺถ โหติ, เอวญฺจ ปน เม เอตํ อรหตํ สุตนฺติ.}

\prose{สาธุ สาธุ มหาราช, สาธุ โข เต เอตํ มหาราช เอวํ โหติ, สาธุ จ ปน เต เอตํ อรหตํ สุตํ.\csromanspacer{} ตํ กิํ มญฺญสิ มหาราชา, ยทิ เอวํ สนฺเต อิเม จตฺตาโร วณฺณา สมสมา โหนฺติ, โน วา, กถํ วา เต เอตฺถ โหตีติ.\csromanspacer{} อทฺธา โข โภ กจฺจาน เอวํ สนฺเต อิเม จตฺตาโร วณฺณา สมสมา โหนฺติ, เนสํ เอตฺถ กิญฺจิ นานากรณํ สมนุปสฺสามีติ.\csromanspacer{} อิมินาปิ โข เอตํ มหาราช ปริยาเยน เวทิตพฺพํ, ยถา “โฆโสเยเวโส โลกสฺมิํ, พฺราหฺมโณว เสฏฺโฐ วณฺโณ, หีโน อญฺโญ วณฺโณ ฯเปฯ พฺรหฺมทายาทา”ติ.}

\proseitem{321}{ตํ กิํ มญฺญสิ มหาราช, อิธ ขตฺติโย สนฺธิํ วา ฉินฺเทยฺย, นิลฺโลปํ วา หเรยฺย, เอกาคาริกํ วา กเรยฺย, ปริปนฺเถ วา ติฏฺเฐยฺย, ปรทารํ วา คจฺเฉยฺย.\csromanspacer{} ตญฺเจ เต ปุริสา คเหตฺวา ทสฺเสยฺยุํ “อยํ เต เทว โจโร อาคุจารี, อิมสฺส ยํ อิจฺฉสิ, ตํ ทณฺฑํ ปเณหี”ติ,}

\csromanmatikaanchor{mtk.41}
\csromantocmarkref{subsection}{4. มธุรสุตฺต}{mtk.41}
\bookmark[dest={mtk.41},level=2]{4. มธุรสุตฺต}
\vspace{\tipitakaparskip}
\csromancenter{มธุรสุตฺต (84) กินฺติ นํ กเรยฺยาสีติ.\csromanspacer{} ฆาเตยฺยาม วา โภ กจฺจาน ชาเปยฺยาม วา ปพฺพาเชยฺยาม วา ยถาปจฺจยํ วา กเรยฺยาม, ตํ กิสฺส เหตุ, ยา หิสฺส โภ กจฺจาน ปุพฺเพ ขตฺติโยติ สมญฺญา, สาสฺส อนฺตรหิตา, โจโรเตฺวว สงฺขฺยํ\footnote{สงฺขํ (สี, สฺยา, กํ, อิ)} คจฺฉติ.}
\prose{\csromanfolio{275}ตํ กิํ มญฺญสิ มหาราช, อิธ พฺราหฺมโณ.\csromanspacer{} อิธ เวสฺโส.\csromanspacer{} อิธ สุทฺโท สนฺธิํ วา ฉินฺเทยฺย, นิลฺโลปํ วา หเรยฺย, เอกาคาริกํ วา กเรยฺย, ปริปนฺเถ วา ติฏฺเฐยฺย, ปรทารํ วา คจฺเฉยฺย.\csromanspacer{} ตญฺเจ เต ปุริสา คเหตฺวา ทสฺเสยฺยุํ “อยํ เต เทว โจโร อาคุจารี, อิมสฺส ยํ อิจฺฉสิ, ตํ ทณฺฑํ ปเณหี”ติ, กินฺติ นํ กเรยฺยาสีติ.\csromanspacer{} ฆาเตยฺยาม วา โภ กจฺจาน ชาเปยฺยาม วา ปพฺพาเชยฺยาม วา ยถาปจฺจยํ วา กเรยฺยาม, ตํ กิสฺส เหตุ, ยา หิสฺส โภ กจฺจาน ปุพฺเพ สุทฺโทติ สมญฺญา, สาสฺส อนฺตรหิตา, โจโรเตฺวว สงฺขฺยํ คจฺฉตีติ.}

\prose{ตํ กิํ มญฺญสิ มหาราช, ยทิ เอวํ สนฺเต อิเม จตฺตาโร วณฺณา สมสมา โหนฺติ, โน วา, กถํ วา เต เอตฺถ โหตีติ.\csromanspacer{} อทฺธา โข โภ กจฺจาน เอวํ สนฺเต อิเม จตฺตาโร วณฺณา สมสมา โหนฺติ, เนสํ เอตฺถ กิญฺจิ นานากรณํ สมนุปสฺสามีติ.\csromanspacer{} อิมินาปิ โข เอตํ มหาราช ปริยาเยน เวทิตพฺพํ, ยถา “โฆโสเยเวโส โลกสฺมิํ, พฺราหฺมโณว เสฏฺโฐ วณฺโณ, หีโน อญฺโญ วณฺโณ ฯเปฯ พฺรหฺมทายาทา”ติ.}

\proseitem{322}{ตํ กิํ มญฺญสิ มหาราช, อิธ ขตฺติโย เกสมสฺสุํ โอหาเรตฺวา กาสายานิ วตฺตานิ อจฺฉาเทตฺวา อคารสฺมา อนคาริยํ ปพฺพชิโต อสฺส.\csromanspacer{} วิรโต ปาณาติปาตา, วิรโต อทินฺนาทานา, วิรโต มุสาวาทา, รตฺตูปรโต เอกภตฺติโก, พฺรหฺมจารี สีลวา กลฺยาณธมฺโม, กินฺติ นํ กเรยฺยาสีติ.\csromanspacer{} อภิวาเทยฺยาม วา\footnote{ปิ (ที 1 สามญฺญผเล 57-58)} โภ กจฺจาน ปจฺจุฏฺเฐยฺยาม วา อาสเนน วา นิมนฺเตยฺยาม, อภินิมนฺเตยฺยาม วา นํ จีวรปิณฺฑปาตเสนาสนคิลานปฺปจฺจยเภสชฺชปริกฺขาเรหิ, ธมฺมิกํ วา อสฺส รกฺขาวรณคุตฺติํ สํวิทเหยฺยาม, ตํ กิสฺส เหตุ, ยา หิสฺส โภ กจฺจาน ปุพฺเพ ขตฺติโยติ สมญฺญา, สาสฺส อนฺตรหิตา, สมโณเตฺวว สงฺขฺยํ คจฺฉตีติ.}

\prose{\csromanfolio{276}ตํ กิํ มญฺญสิ มหาราช, อิธ พฺราหฺมโณ.\csromanspacer{} อิธ เวสฺโส.\csromanspacer{} อิธ สุทฺโท เกสมสฺสุํ โอหาเรตฺวา กาสายานิ วตฺถานิ อจฺฉาเทตฺวา อคารสฺมา อนคาริยํ ปพฺพชิโต อสฺส.\csromanspacer{} วิรโต ปาณาติปาตา, วิรโต อทินฺนาทานา, วิรโต มุสาวาทา, รตฺถูปรโต เอกภตฺติโก, พฺรหฺมจารี สีลวา กลฺยาณธมฺโม, กินฺติ นํ กเรยฺยาสีติ.\csromanspacer{} อภิวาเทยฺยาม วา โภ กจฺจาน ปจฺจุฏฺเฐยฺยาม วา อาสเนน วา นิมนฺเตยฺยาม, อภินิมนฺเตยฺยาม วา นํ จีวรปิณฺฑปาตเสนาสนคิลานปฺปจฺจยเภสชฺชปริกฺขาเรหิ, ธมฺมิกํ วา อสฺส รกฺขาวรณคุตฺติํ สํวิทเหยฺยาม, ตํ กิสฺส เหตุ, ยา หิสฺส โภ กจฺจาน ปุพฺเพ สุทฺโทติ สมญฺญา, สาสฺส อนฺตรหิตา, สมโณเตฺวว สงฺขฺยํ คจฺฉตีติ.}

\prose{ตํ กิํ มญฺญสิ มหาราช, ยทิ เอวํ สนฺเต อิเม จตฺตาโร วณฺณา สมสมา โหนฺติ, โน วา, กถํ วา เต เอตฺถ โหตีติ.\csromanspacer{} อทฺธา โข โภ กจฺจาน เอวํ สนฺเต อิเม จตฺตาโร วณฺณา สมสมา โหนฺติ, เนสํ เอตฺถ กิญฺจิ นานากรณํ สมนุปสฺสามีติ.\csromanspacer{} อิมินาปิ โข เอตํ มหาราช ปริยาเยน เวทิตพฺพํ, ยถา “โฆโสเยเวโส โลกสฺมิํ, พฺราหฺมโณว เสฏฺโฐ วณฺโณ, หีโน อญฺโญ วณฺโณ.\csromanspacer{} พฺราหฺมโณว สุกฺโก วณฺโณ, กโณฺห อญฺโญ วณฺโณ.\csromanspacer{} พฺราหฺมณาว สุชฺฌนฺติ โน อพฺราหฺมณา.\csromanspacer{} พฺราหฺมณาว พฺรหฺมุโน ปุตฺตา โอรสา มุขโต ชาตา พฺรหฺมชา พฺรหฺมนิมฺมิตา พฺรหฺมทายาทา”ติ.}

\proseitem{323}{เอวํ วุตฺเต ราชา มาธุโร อวนฺติปุตฺโต อายสฺมนฺตํ มหากจฺจานํ เอตทโวจ “อภิกฺกนฺตํ โภ กจฺจาน, อภิกฺกนฺตํ โภ กจฺจาน, เสยฺยถาปิ โภ กจฺจาน นิกฺกุชฺชิตํ วา อุกฺกุชฺเชยฺย, ปฏิจฺฉนฺนํ วา วิวเรยฺย, มูฬฺหสฺส วา มคฺคํ อาจิกฺเขยฺย, อนฺธกาเร วา เตลปชฺโชตํ ธาเรยฺย ‘จกฺขุมนฺโต รูปานิ ทกฺขนฺตี’ติ.\csromanspacer{} เอวเมวํ โภตา กจฺจาเนน อเนกปริยาเยน ธมฺโม ปกาสิโต, เอสาหํ ภวนฺตํ กจฺจานํ สรณํ คจฺฉามิ ธมฺมญฺจ ภิกฺขุสํฆญฺจ, อุปาสกํ มํ ภวํ กจฺจาโน ธาเรตุ อชฺชตคฺเค ปาณุเปตํ สรณํ คตนฺ”ติ.\csromanspacer{} มา โข มํ ตฺวํ มหาราช สรณํ อคมาสิ, ตเมว ตฺวํ\footnote{ตเมตํ ตฺวํ (สฺยา, กํ), ตเมตํ (ก)} ภควนฺตํ สรณํ คจฺฉ, ยมหํ สรณํ คโตติ.\csromanspacer{} กหํ ปน โภ กจฺจาน เอตรหิ โส ภควา วิหรติ อรหํ สมฺมาสมฺพุทฺโธติ.\csromanspacer{} ปรินิพฺพุโต โข มหาราช เอตรหิ}

\vspace{\tipitakaparskip}
\csromancenter{โพธิราชกุมารสุตฺต (85) โส ภควา อรหํ สมฺมาสมฺพุทฺโธติ.\csromanspacer{} สเจปิ มยํ โภ กจฺจาน สุเณยฺยาม ตํ ภควนฺตํ ทสสุ โยชเนสุ, ทสปิ มยํ โยชนานิ คจฺเฉยฺยาม ตํ ภควนฺตํ ทสฺสนาย อรหนฺตํ สมฺมาสมฺพุทฺธํ.\csromanspacer{} สเจปิ มยํ โภ กจฺจาน สุเณยฺยาม ตํ ภควนฺตํ วีสติยา โยชเนสุ.\csromanspacer{} ติํสาย โยชเนสุ.\csromanspacer{} จตฺตารีสาย โยชเนสุ.\csromanspacer{} ปญฺญาสาย โยชเนสุ, ปญฺญาสมฺปิ มยํ โยชนานิ คจฺเฉยฺยาม ตํ ภควนฺตํ ทสฺสนาย อรหนฺตํ สมฺมาสมฺพุทฺธํ.\csromanspacer{} โยชนสเต เจปิ มยํ โภ กจฺจาน สุเณยฺยาม ตํ ภควนฺตํ, โยชนสตมฺปิ มยํ คจฺเฉยฺยาม ตํ ภควนฺตํ ทสฺสนาย อรหนฺตํ สมฺมาสมฺพุทฺธํ.\csromanspacer{} ยโต จ โภ กจฺจาน ปรินิพฺพุโต โส ภควา, ปรินิพฺพุตมฺปิ มยํ ภควนฺตํ สรณํ คจฺฉาม ธมฺมญฺจ หิกฺขุสํฆญฺจ, อุปาสกํ มํ ภวํ กจฺจาโน ธาเรตุ อชฺชตคฺเค ปาณุเปตํ สรณํ คตนฺติ.}
\nitthitamruled{มธุรสุตฺตํ นิฏฺฐิตํ จตุตฺถํ.}
\csromanheader{2}{5. โพธิราชกุมารสุตฺต}
\proseitem{324}{\csromanfolio{277}เอวํ เม สุตํ\csromandash{}เอกํ สมยํ ภควา ภคฺเคสุ วิหรติ สุสุมารคิเร เภสกฬาวเน มิคทาเย.\csromanspacer{} เตน โข ปน สมเยน โพธิสฺส ราชกุมารสฺส โกกนโท\footnote{โกกนุโท (สฺยา, กํ, ก)} นาม ปาสาโท อจิรการิโต โหติ อนชฺฌาวุฏฺโฐ สมเณน วา พฺราหฺมเณน วา เกนจิ วา มนุสฺสภูเตน.\csromanspacer{} อถ โข โพธิ ราชกุมาโร สญฺชิกาปุตฺตํ มาณวํ อามนฺเตสิ “เอหิ ตฺวํ สมฺม สญฺชิกาปุตฺต เยน ภควา เตนุปสงฺกม, อุปสงฺกมิตฺวา มม วจเนน ภควโต ปาเท สิรสา วนฺท, อปฺปาพาธํ อปฺปาตงฺกํ ลหุฏฺฐานํ พลํ ผาสุวิหารํ ปุจฺฉ ‘โพธิ ภนฺเต ราชกุมาโร ภควโต ปาเท สิรสา วนฺทติ, อปฺปาพาธํ อปฺปาตงฺกํ ลหุฏฺฐานํ พลํ ผาสุวิหารํ ปุจฺฉตี’ติ, เอวญฺจ วเทหิ ‘อธิวาเสตุ กิร ภนฺเต ภควา โพธิสฺส ราชกุมารสฺส สฺวาตนาย ภตฺตํ สทฺธิํ ภิกฺขุสํเฆนา’ติ”. “เอวํ โภ”ติ โข สญฺชิกาปุตฺโต มาณโว โพธิสฺส ราชกุมารสฺส ปฏิสฺสุตฺวา เยน ภควา เตนุปสงฺกมิ, อุปสงฺกมิตฺวา ภควตา สทฺธิํ สมฺโมทิ.\csromanspacer{} สมฺโมทนียํ กถํ สายณียํ วีติสาเรตฺวา เอกมนฺตํ นิสีทิ. \csromanfolio{278}เอกมนฺตํ นิสินฺโน โข สญฺชิกาปุตฺโต มาณโว ภควนฺตํ เอตทโวจ “โพธิ โข\footnote{โพธิ โภ โคตม (สี, สฺยา, กํ, อิ)} ราชกุมาโร โภโต โคตมสฺส ปาเท สิรสา วนฺทติ, อปฺปาพาธํ อปฺปาตงฺกํ ลหุฏฺฐานํ พลํ ผาสุวิหารํ ปุจฺฉติ, เอวญฺจ วเทติ ‘อธิวาเสตุ กิร ภวํ โคตโม โพธิสฺส ราชกุมารสฺส สฺวาตนาย ภตฺตํ สทฺธิํ ภิกฺขุสํเฆนา’ติ”.\csromanspacer{} อธิวาเสสิ ภควา ตุณฺหีภาเวน.\csromanspacer{} อถ โข สญฺชิกาปุตฺโต มาณโว ภควโต อธิวาสนํ วิทิตฺวา อุฏฺฐายาสนา เยน โพธิ ราชกุมาโร เตนุปสงฺกมิ, อุปสงฺกมิตฺวา โพธิํ ราชกุมารํ เอตทโวจ “อโวจุมฺหา โภโต วจเนน ตํ ภวนฺตํ โคตมํ, โพธิ โข ราชกุมาโร โภโต โคตมสฺส ปาเท สิรสา วนฺทติ, อปฺปาพาธํ อปฺปาตงฺกํ ลหุฏฺฐานํ พลํ ผาสุวิหารํ ปุจฺฉติ, เอวญฺจ วเทติ ‘อธิวาเสตุ กิร ภวํ โคตโม โพธิสฺส ราชกุมารสฺส สฺวาตนาย ภตฺตํ สทฺธิํ ภิกฺขุสํเฆนา’ติ, อธิวุฏฺฐญฺจ ปน สมเณน โคตเมนา”ติ.}

\proseitem{325}{อถ โข โพธิ ราชกุมาโร ตสฺสา รตฺติยา อจฺจเยน สเก นิเวสเน ปณีตํ ขาทนียํ โภชนียํ ปฏิยาทาเปตฺวา โกกนทญฺจ ปาสาทํ โอทาเตหิ ทุสฺเสหิ สนฺถราเปตฺวา ยาว ปจฺฉิมโสปานกเฬวรา\footnote{กเฬพรา (สี)} สญฺชิกาปุตฺตํ มาณวํ อามนฺเตสิ “เอหิ ตฺวํ สมฺม สญฺชิกาปุตฺต เยน ภควา เตนุปสงฺกม, อุปสงฺกมิตฺวา ภควโต กาลํ อาโรเจหิ ‘กาโล ภนฺเต นิฏฺฐิตํ ภตฺตนฺ’ติ”. “เอวํ โภ”ติ โข สญฺชิกาปุตฺโต มาณโว โพธิสฺส ราชกุมารสฺส ปฏิสฺสุตฺวา เยน ภควา เตนุปสงฺกมิ, อุปสงฺกมิตฺวา ภควโต กาลํ อาโรเจสิ “กาโล โภ โคตม นิฏฺฐิตํ ภตฺตนฺ”ติ.\csromanspacer{} อถ โข ภควา ปุพฺพณฺหสมยํ นิวาเสตฺวาปตฺตจีวรมาทาย เยน โพธิสฺส ราชกุมารสฺส นิเวสนํ เตนุปสงฺกมิ.\csromanspacer{} เตน โข ปน สมเยน โพธิ ราชกุมาโร พหิทฺวารโกฏฺฐเก ฐิโต โหติ ภควนฺตํ อาคมยมาโน.\csromanspacer{} อทฺทสา โข โพธิ ราชกุมาโร ภควนฺตํ ทูรโตว อาคจฺฉนฺตํ, ทิสฺวาน ปจฺจุคฺคนฺตฺวา ภควนฺตํ อภิวาเทตฺวา ปุรกฺขตฺวา เยน โกกนโท ปาสาโท เตนุปสงฺกมิ.\csromanspacer{} อถ โข ภควา ปจฺฉิมํ โสปานกเฬวรํ นิสฺสาย อฏฺฐาสิ.\csromanspacer{} อถ โข โพธิ ราชกุมาโร ภควนฺตํ เอตทโวจ “อภิรุหตุ\footnote{อภิรูหตุ (สฺยา, กํ, อิ) อกฺกมตุ (วิ 4. 267 ปิฏฺเฐ)} ภนฺเต ภควา ทุสฺสานิ, อภิรุหตุ สุคโต ทุสฺสานิ, ยํ มม อสฺส ทีฆรตฺตํ}

\vspace{\tipitakaparskip}
\csromancenter{โพธิราชกุมารสุตฺต (85) หิตาย สุขายา”ติ.\csromanspacer{} เอวํ วุตฺเต ภควา ตุณฺหี อโหสิ.\csromanspacer{} ทุติยมฺปิ โข ฯเปฯ.\csromanspacer{} ตติยมฺปิ โข โพธิ ราชกุมาโร ภควนฺตํ เอตทโวจ “อภิรุหตุ ภนฺเต ภควา ทุสฺสานิ, อภิรุหตุ สุคโต ทุสฺสานิ, ยํ มม อสฺส ทีฆรตฺตํ หิตาย สุขายา”ติ.}
\proseitem{326}{\csromanfolio{279}อถ โข ภควา อายสฺมนฺตํ อานนฺทํ อปโลเกสิ.\csromanspacer{} อถ โข อายสฺมา อานนฺโท โพธิํ ราชกุมารํ เอตทโวจ “สํหรตุ ราชกุมาร ทุสฺสานิ, น ภควา เจลปฏิกํ\footnote{เจลปตฺติกํ (สี, อิ)} อกฺกมิสฺสติ, ปจฺฉิมํ ชนตํ ตถาคโต อนุกมฺปตี”ติ\footnote{อปโลเกตีติ (สพฺพตฺถ)}.\csromanspacer{} อถ โข โพธิ ราชกุมาโร ทุสฺสานิ สํหราเปตฺวา อุปริโกกนทปาสาเท\footnote{อุปริโกกนเท ปาสาเท (สี, อิ, วินเย จ), อุปริโกกนเท (สฺยา, กํ)} อาสนานิ ปญฺญเปสิ.\csromanspacer{} อถ โข ภควา โกกนทํ ปาสาทํ อภิรุหิตฺวา ปญฺญตฺเต อาสเน นิสีทิ สทฺธิํ ภิกฺขุสํเฆน.\csromanspacer{} อถ โข โพธิ ราชกุมาโร พุทฺธปฺปมุขํ ภิกฺขุสํฆํ ปณีเตน ขาทนีเยน โภชนีเยน สหตฺถา สนฺตปฺเปสิ สมฺปวาเรสิ.\csromanspacer{} อถ โข โพธิ ราชกุมาโร ภควนฺตํ ภุตฺตาวิํ โอนีตปตฺตปาณิํ อญฺญตรํ นีจํ อาสนํ คเหตฺวา เอกมนฺตํ นิสีทิ, เอกมนฺตํ นิสินฺโน โข โพธิ ราชกุมาโร ภควนฺตํ เอตทโวจ “มยฺหํ โข ภนฺเต เอวํ โหติ ‘น โข สุเขน สุขํ อธิคนฺตพฺพํ, ทุกฺเขน โข สุขํ อธิคนฺตพฺพนฺ’ติ”.}

\proseitem{327}{มยฺหมฺปิ โข ราชกุมาร ปุพฺเพว สมฺโพธา อนภิสมฺพุทฺธสฺส โพธิสตฺตสฺเสว สโต เอตทโหสิ “น โข สุเขน สุขํ อธิคนฺตพฺพํ, ทุกฺเขน โข สุขํ อธิคนฺตพฺพนฺ”ติ.\csromanspacer{} โส โข อหํ ราชกุมาร อปเรน สมเยน ทหโรว สมาโน สุสุกาฬเกโส ภเทฺรน โยพฺพเนน สมนฺนาคโต, ปฐเมน วยสา อกามกานํ มาตาปิตูนํ อสฺสุมุขานํ รุทนฺตานํ เกสมสฺสุํ โอหาเรตฺวา กาสายานิ วตฺถานิ อจฺฉาเทตฺวา อคารสฺมา อนคาริยํ ปพฺพชิํ, โส เอวํ ปพฺพชิโต สมาโน กิํกุสลคเวสี\footnote{กิํกุสลํคเวสี (ก)} อนุตฺตรํ สนฺติวรปทํ ปริเยสมาโน เยน อาฬาโร กาลาโม เตนุปสงฺกมิํ, อุปสงฺกมิตฺวา อาฬารํ กาลามํ เอตทโวจํ “อิจฺฉามหํ อาวุโส กาลาม อิมสฺมิํ ธมฺมวินเย พฺรหฺมจริยํ จริตุนฺ”ติ.\csromanspacer{} เอวํ วุตฺเต ราชกุมาร อาฬาโร กาลาโม มํ เอตทโวจ “วิหรตายสฺมา, \csromanfolio{280}ตาทิโส อยํ ธมฺโม, ยตฺถ วิญฺญู ปุริโส นจิรสฺเสว สกํ อาจริยกํ สยํ อภิญฺญา สจฺฉิกตฺวา อุปสมฺปชฺช วิหเรยฺยา”ติ.\csromanspacer{} โส โข อหํ ราชกุมาร นจิรสฺเสว ขิปฺปเมว ตํ ธมฺมํ ปริยาปุณิํ, โส โข อหํ ราชกุมาร ตาวตเกเนว โอฏฺฐปหตมตฺเตน ลปิตลาปนมตฺเตน ญาณวาทญฺจ วทามิ เถรวาทญฺจ, “ชานามิ ปสฺสามี”ติ จ ปฏิชานามิ อหญฺเจว อญฺเญ จ.\csromanspacer{} ตสฺส มยฺหํ ราชกุมาร เอตทโหสิ “น โข อาฬาโร กาลาโม อิมํ ธมฺมํ เกวลํ สทฺธามตฺตเกน ‘สยํ อภิญฺญา สจฺฉิกตฺวา อุปสมฺปชฺช วิหรามี’ติ ปเวเทติ, อทฺธา อาฬาโร กาลาโม อิมํ ธมฺมํ ชานํ ปสฺสํ วิหรตี”ติ.}

\prose{อถ ขฺวาหํ ราชกุมาร เยน อาฬาโร กาลาโม เตนุปสงฺกมิํ, อุปสงฺกมิตฺวา อาฬารํ กาลามํ เอตทโวจํ “กิตฺตาวตา โน อาวุโส กาลาม อิมํ ธมฺมํ สยํ อภิญฺญา สจฺฉิกตฺวา อุปสมฺปชฺช วิหรามีติ ปเวเทสี”ติ\footnote{อุปสมฺปชฺช ปเวเทสีติ (สี, สฺยา, กํ, อิ)}.\csromanspacer{} เอวํ วุตฺเต ราชกุมาร อาฬาโร กาลาโม อากิญฺจญฺญายตนํ ปเวเทสิ.\csromanspacer{} ตสฺส มยฺหํ ราชกุมาร เอตทโหสิ “น โข อาฬารสฺเสว กาลามสฺส อตฺถิ สทฺธา, มยฺหํปตฺถิ สทฺธา.\csromanspacer{} น โข อาฬารสฺเสว กาลามสฺส อตฺถิ วีริยํ ฯเปฯ สติ.\csromanspacer{} สมาธิ.\csromanspacer{} ปญฺญา, มยฺหํปตฺถิ ปญฺญา.\csromanspacer{} ยํ นูนาหํ ยํ ธมฺมํ อาฬาโร กาลาโม สยํ อภิญฺญา สจฺฉิกตฺวา อุปสมฺปชฺช วิหรามีติ ปเวเทติ, ตสฺส ธมฺมสฺส สจฺฉิกิริยาย ปทเหยฺยนฺ”ติ.\csromanspacer{} โส โข อหํ ราชกุมาร นจิรสฺเสว ขิปฺปเมว ตํ ธมฺมํ สยํ อภิญฺญา สจฺฉิกตฺวา อุปสมฺปชฺช วิหาสิํ.\csromanspacer{} อถ ขฺวาหํ ราชกุมาร เยน อาฬาโร กาลาโม เตนุปสงฺกมิํ, อุปสงฺกมิตฺวา อาฬารํ กาลามํ เอตทโวจํ “เอตฺตาวตา โน อาวุโส กาลาม อิมํ ธมฺมํ สยํ อภิญฺญา สจฺฉิกตฺวา อุปสมฺปชฺช ปเวเทสี”ติ.\csromanspacer{} เอตฺตาวตา โข อหํ อาวุโส อิมํ ธมฺมํ สยํ อภิญฺญา สจฺฉิกตฺวา อุปสมฺปชฺช ปเวเทมีติ.\csromanspacer{} อหมฺปิ โข อาวุโส เอตฺตาวตา อิมํ ธมฺมํ สยํ อภิญฺญา สจฺฉิกตฺวา อุปสมฺปชฺช วิหรามีติ.\csromanspacer{} ลาภา โน อาวุโส, สุลทฺธํ โน อาวุโส, เย มยํ อายสฺมนฺตํ ตาทิสํ สพฺรหฺมจาริํ ปสฺสาม, อิติ ยาหํ ธมฺมํ สยํ อภิญฺญา สจฺฉิกตฺวา อุปสมฺปชฺช ปเวเทมิ, ตํ ตฺวํ ธมฺมํ สยํ อภิญฺญา สจฺฉิกตฺวา อุปสมฺปชฺช วิหรสิ.\csromanspacer{} ยํ ตฺวํ ธมฺมํ สยํ อภิญฺญา สจฺฉิกตฺวา อุปสมฺปชฺช วิหรสิ, ตมหํ ธมฺมํ สยํ อภิญฺญา}

\vspace{\tipitakaparskip}
\csromancenter{โพธิราชกุมารสุตฺต (85) สจฺฉิกตฺวา อุปสมฺปชฺช ปเวเทมิ.\csromanspacer{} อิติ ยาหํ ธมฺมํ ชานามิ, ตํ ตฺวํ ธมฺมํ ชานาสิ.\csromanspacer{} ยํ ตฺวํ ธมฺมํ ชานาสิ, ตมหํ ธมฺมํ ชานามิ.\csromanspacer{} อิติ ยาทิโส อหํ ตาทิโส ตุวํ.\csromanspacer{} ยาทิโส ตุวํ, ตาทิโส อหํ.\csromanspacer{} เอหิ ทานิ อาวุโส อุโภว สนฺตา อิมํ คณํ ปริหรามาติ.\csromanspacer{} อิติ โข ราชกุมาร อาฬาโร กาลาโม อาจริโย เม สมาโน (อตฺตโน)\footnote{( ) นตฺถิ (สี, สฺยา, กํ, อิ)} อนฺเตวาสิํ มํ สมานํ อตฺตนา\footnote{อตฺตโน (สี, อิ)} สมสมํ ฐเปสิ, อุฬาราย จ มํ ปูชาย ปูเชสิ.\csromanspacer{} ตสฺส มยฺหํ ราชกุมาร เอตทโหสิ “นายํ ธมฺโม นิพฺพิทาย น วิราคาย น นิโรธาย น อุปสมาย น อภิญฺญาย น สมฺโพธาย น นิพฺพานาย สํวตฺตติ, ยาวเทว อากิญฺจญฺญายตนูปปตฺติยา”ติ.\csromanspacer{} โส โข อหํ ราชกุมาร ตํ ธมฺมํ อนลงฺกริตฺวา ตสฺมา ธมฺมา นิพฺพิชฺช อปกฺกมิํ.}
\proseitem{328}{\csromanfolio{281}โส โข อหํ ราชกุมาร กิํกุสลคเวสี อนุตฺตรํ สนฺติวรปทํ ปริเยสมาโน เยน อุทโก\footnote{อุทฺทโก (สี, สฺยา, กํ, อิ)} รามปุตฺโต เตนุปสงฺกมิํ, อุปสงฺกมิตฺวา อุทกํ รามปุตฺตํ เอตทโวจํ “อิจฺฉามหํ อาวุโส\footnote{อาวุโส ราม (สี, สฺยา, กํ, ก) ปสฺส ม 1. ปาสราสิสุตฺเต (221) ปิฏฺเฐ.} อิมสฺมิํ ธมฺมวินเย พฺรหฺมจริยํ จริตุนฺ”ติ.\csromanspacer{} เอวํ วุตฺเต ราชกุมาร อุทโก รามปุตฺโต มํ เอตทโวจ “วิหรตายสฺมา, ตาทิโส อยํ ธมฺโม, ยตฺถ วิญฺญู ปุริโส นจิรสฺเสว สกํ อาจริยกํ สยํ อภิญฺญา สจฺฉิกตฺวา อุปสมฺปชฺช วิหเรยฺยา”ติ.\csromanspacer{} โส โข อหํ ราชกุมาร นจิรสฺเสว ขิปฺปเมว ตํ ธมฺมํ ปริยาปุณิํ.\csromanspacer{} โส โข อหํ ราชกุมาร ตาวตเกเนว โอฏฺฐปหตมตฺเตน ลปิตลาปนมตฺเตน ญาณวาทญฺจ วทามิ เถรวาทญฺจ, “ชานามิ ปสฺสามี”ติ จ ปฏิชานามิ อหญฺเจว อญฺเญ จ.\csromanspacer{} ตสฺส มยฺหํ ราชกุมาร เอตทโหสิ “น โข ราโม อิมํ ธมฺมํ เกวลํ สทฺธามตฺตเกน ‘สยํ อภิญฺญา สจฺฉิกตฺวา อุปสมฺปชฺช วิหรามี’ติ ปเวเทสิ, อทฺธา ราโม อิมํ ธมฺมํ ชานํ ปสฺสํ วิหาสี”ติ.\csromanspacer{} อถ ขฺวาหํ ราชกุมาร เยน อุทโก รามปุตฺโต เตนุปสงฺกมิํ, อุปสงฺกมิตฺวา อุทกํ รามปุตฺตํ เอตทโวจํ “กิตฺตาวตา โน อาวุโส ราโม อิมํ ธมฺมํ สยํ อภิญฺญา สจฺฉิกตฺวา อุปสมฺปชฺช วิหรามีติ ปเวเทสี”ติ.\csromanspacer{} เอวํ วุตฺเต ราชกุมาร อุทโก รามปุตฺโต เนวสญฺญานาสญฺญายตนํ ปเวเทสิ.\csromanspacer{} ตสฺส มยฺหํ ราชกุมาร เอตทโหสิ “น โข รามสฺเสว อโหสิ สทฺธา, มยฺหํปตฺถิ สทฺธา.\csromanspacer{} น โข รามสฺเสว อโหสิ วีริยํ ฯเปฯ สติ.\csromanspacer{} สมาธิ.\csromanspacer{} ปญฺญา, \csromanfolio{282}มยฺหํปตฺถิ ปญฺญา, ยํนูนาหํ ยํธมฺมํ ราโม สยํ อภิญฺญา สจฺฉิกตฺวา อุปสมฺปชฺช วิหรามีติ ปเวเทสิ, ตสฺส ธมฺมสฺส สจฺฉิกิริยาย ปทเหยฺยนฺ”ติ.\csromanspacer{} โส โข อหํ ราชกุมาร นจิรสฺเสว ขิปฺปเมว ตํ ธมฺมํ สยํ อภิญฺญา สจฺฉิกตฺวา อุปสมฺปชฺช วิหาสิํ.}

\prose{อถ ขฺวาหํ ราชกุมาร เยน อุทโก รามปุตฺโต เตนุปสงฺกมิํ, อุปสงฺกมิตฺวา อุทกํ รามปุตฺตํ เอตทโวจํ “เอตฺตาวตา โน อาวุโส ราโม อิมํ ธมฺมํ สยํ อภิญฺญา สจฺฉิกตฺวา อุปสมฺปชฺช ปเวเทสี”ติ.\csromanspacer{} เอตฺตาวตา โข อาวุโส ราโม อิมํ ธมฺมํ สยํ อภิญฺญา สจฺฉิกตฺวา อุปสมฺปชฺช ปเวเทสีติ.\csromanspacer{} อหมฺปิ โข อาวุโส เอตฺตาวตา อิมํ ธมฺมํ สยํ อภิญฺญา สจฺฉิกตฺวา อุปสมฺปชฺช วิหรามีติ.\csromanspacer{} ลาภา โน อาวุโส, สุลทฺธํ โน อาวุโส, เย มยํ อายสฺมนฺตํ ตาทิสํ สพฺรหฺมจาริํ ปสฺสาม.\csromanspacer{} อิติ ยํ ธมฺมํ ราโม สยํ อภิญฺญา สจฺฉิกตฺวา อุปสมฺปชฺช ปเวเทสิ, ตํ ตฺวํ ธมฺมํ สยํ อภิญฺญา สจฺฉิกตฺวา อุปสมฺปชฺช วิหรสิ.\csromanspacer{} ยํ ตฺวํ ธมฺมํ สยํ อภิญฺญา สจฺฉิกตฺวา อุปสมฺปชฺช วิหรสิ, ตํ ธมฺมํ ราโม สยํ อภิญฺญา สจฺฉิกตฺวา อุปสมฺปชฺช ปเวเทสิ.\csromanspacer{} อิติ ยํ ธมฺมํ ราโม อภิญฺญาสิ, ตํ ตฺวํ ธมฺมํ ชานาสิ.\csromanspacer{} ยํ ตฺวํ ชานาสิ, ตํ ธมฺมํ ราโม อภิญฺญาสิ.\csromanspacer{} อิติ ยาทิโส ราโม อโหสิ, ตาทิโส ตุวํ.\csromanspacer{} ยาทิโส ตุวํ, ตาทิโส ราโม อโหสิ.\csromanspacer{} เอหิ ทานิ อาวุโส ตุวํ อิมํ คณํ ปริหราติ.\csromanspacer{} อิติ โข ราชกุมาร อุทโก รามปุตฺโต สพฺรหฺมจารี เม สมาโน อาจริยฏฺฐาเน มํ ฐเปสิ, อุฬาราย จ มํ ปูชาย ปูเชสิ.\csromanspacer{} ตสฺส มยฺหํ ราชกุมาร เอตทโหสิ “นายํ ธมฺโม นิพฺพิทาย น วิราคาย น นิโรธาย น อุปสมาย น อภิญฺญาย น สมฺโพธาย น นิพฺพานาย สํวตฺตติ, ยาวเทว เนวสญฺญานาสญฺญายตนูปปตฺติยา”ติ.\csromanspacer{} โส โข อหํ ราชกุมาร ตํ ธมฺมํ อนลงฺกริตฺวา ตสฺมา ธมฺมา นิพฺพิชฺช อปกฺกมิํ.}

\proseitem{329}{โส โข อหํ ราชกุมาร กิํกุสลคเวสี อนุตฺตรํ สนฺติวรปทํ ปริเยสมาโน มคเธสุ อนุปุพฺเพน จาริกํ จรมาโน เยน อุรุเวลา เสนานิคโม ตทวสริํ.\csromanspacer{} ตตฺถทฺทสํ รมณียํ ภูมิภาคํ ปาสาทิกญฺจ วนสณฺฑํ, นทิญฺจ สนฺทนฺติํ เสตกํ สุปติตฺถํ รมณียํ สมนฺตา จ โคจรคามํ.\csromanspacer{} ตสฺส มยฺหํ ราชกุมาร เอตทโหสิ “รมณีโย วต โภ ภูมิภาโค, ปาสาทิโก จ วนสณฺโฑ, นที จ สนฺทติ เสตกา}

\vspace{\tipitakaparskip}
\csromancenter{โพธิราชกุมารสุตฺต (85) สุปติตฺถา รมณียา, สมนฺตา\footnote{สามนฺตา (?) ปุริมปิฏฺเฐปิ.} จ โคจรคาโม, ‘อลํ วติทํ กุลปุตฺตสฺส ปธานตฺถิกสฺส ปธานายา’ติ”.\csromanspacer{} โส โข อหํ ราชกุมาร ตตฺเถว นิสีทิํ “อลมิทํ ปธานายา”ติ.\csromanspacer{} อปิสฺสุ มํ ราชกุมาร ติสฺโส อุปมา ปฏิภํสุ อนจฺฉริยา ปุพฺเพ อสฺสุตปุพฺพา\csromandash{}}
\prose{\csromanfolio{283}เสยฺยถาปิ ราชกุมาร อลฺลํ กฏฺฐํ สสฺเนหํ อุทเก นิกฺขิตฺตํ.\csromanspacer{} อถ ปุริโส อาคจฺเฉยฺย อุตฺตรารณิํ อาทาย “อคฺคิํ อภินิพฺพตฺเตสฺสามิ เตโช ปาตุกริสฺสามี”ติ.\csromanspacer{} ตํ กิํ มญฺญสิ ราชกุมาร, อปิ นุ โส ปุริโส อมุํ อลฺลํ กฏฺฐํ สสฺเนหํ อุทเก นิกฺขิตฺตํ อุตฺตรารณิํ อาทาย อภิมนฺเถนฺโต\footnote{อภิมตฺถนฺโต (สฺยา, กํ, ก)} อคฺคิํ อภินิพฺพตฺเตยฺย เตโช ปาตุกเรยฺยาติ.\csromanspacer{} โน หิทํ ภนฺเต.\csromanspacer{} ตํ กิสฺส เหตุ, อทุํ หิ ภนฺเต อลฺลํ กฏฺฐํ สสฺเนหํ, ตญฺจ ปน อุทเก นิกฺขิตฺตํ, ยาวเทว จ ปน โส ปุริโส กิลมถสฺส วิฆาตสฺส ภาคี อสฺสาติ.\csromanspacer{} เอวเมว โข ราชกุมาร เย หิ เกจิ สมณา วา พฺราหฺมณา วา กาเยน เจว จิตฺเตน จ กาเมหิ อวูปกฏฺฐา วิหรนฺติ.\csromanspacer{} โย จ เนสํ กาเมสุ กามจฺฉนฺโท กามสฺเนโห กามมุจฺฉา กามปิปาสา กามปริฬาโห, โส จ อชฺฌตฺตํ น สุปฺปหีโน โหติ น สุปฺปฏิปฺปสฺสทฺโธ.\csromanspacer{} โอปกฺกมิกา เจปิ เต โภนฺโต สมณพฺราหฺมณา ทุกฺขา ติพฺพา ขรา กฏุกา เวทนา เวทยนฺติ, อภพฺพาว เต ญาณาย ทสฺสนาย อนุตฺตราย สมฺโพธาย.\csromanspacer{} โน เจปิ เต โภนฺโต สมณพฺราหฺมณา โอปกฺกมิกา ทุกฺขา ติพฺพา ขรา กฏุกา เวทนา เวทยนฺติ, อภพฺพาว เต ญาณาย ทสฺสนาย อนุตฺตราย สมฺโพธาย.\csromanspacer{} อยํ โข มํ ราชกุมาร ปฐมา อุปมา ปฏิภาสิ อนจฺฉริยา ปุพฺเพ อสฺสุตปุพฺพา.}

\proseitem{330}{อปราปิ โข มํ ราชกุมาร ทุติยา อุปมา ปฏิภาสิ อนจฺฉริยา ปุพฺเพ อสฺสุตปุพฺพา\csromandash{}}

\csromanmatikaanchor{mtk.42}
\csromantocmarkref{subsection}{5. โพธิราชกุมารสุตฺต}{mtk.42}
\bookmark[dest={mtk.42},level=2]{5. โพธิราชกุมารสุตฺต}
\prose{เสยฺยถาปิ ราชกุมาร อลฺลํ กฏฺฐํ สสฺเนหํ อารกา อุทกา ถเล นิกฺขิตฺตํ.\csromanspacer{} อถ ปุริโส อาคจฺเฉยฺย อุตฺตรารณิํ อาทาย “อคฺคิํ อภินิพฺพตฺเตสฺสามิ เตโช ปาตุกริสฺสามี”ติ.\csromanspacer{} ตํ กิํ มญฺญสิ ราชกุมาร, อปิ นุ โส ปุริโส อมุํ อลฺลํ กฏฺฐํ สสฺเนหํ อารกา อุทกา ถเล นิกฺขิตฺตํ อุตฺตรารณิํ อาทาย อภิมนฺเถนฺโต อคฺคิํ \csromanfolio{284}อภินิพฺพตฺเตยฺย เตโช ปาตุกเรยฺยาติ.\csromanspacer{} โน หิทํ ภนฺเต.\csromanspacer{} ตํ กิสฺส เหตุ, อทุํ หิ ภนฺเต อลฺลํ กฏฺฐํ สสฺเนหํ กิญฺจาปิ อารกา อุทกา ถเล นิกฺขิตฺตํ, ยาวเทว จ ปน โส ปุริโส กิลมถสฺส วิฆาตสฺส ภาคี อสฺสาติ.\csromanspacer{} เอวเมว โข ราชกุมาร เย หิ เกจิ สมณา วา พฺราหฺมณา วา กาเยน เจว จิตฺเตน จ กาเมหิ วูปกฏฺฐา วิหรนฺติ.\csromanspacer{} โย จ เนสํ กาเมสุ กามจฺฉนฺโท กามสฺเนโห กามมุจฺฉา กามปิปาสา กามปริฬาโห, โส จ อชฺฌตฺตํ น สุปฺปหีโน โหติ น สุปฺปฏิปฺปสฺสทฺโธ.\csromanspacer{} โอปกฺกมิกา เจปิ เต โภนฺโต สมณพฺราหฺมณา ทุกฺขา ติพฺพา ขรา กฏุกา เวทนา เวทยนฺติ, อภพฺพาว เต ญาณาย ทสฺสนาย อนุตฺตราย สมฺโพธาย, โน เจปิ เต โภนฺโต สมณพฺราหฺมณา โอปกฺกมิกา ทุกฺขา ติพฺพา ขรา กฏุกา เวทนา เวทยนฺติ, อภพฺพาว เต ญาณาย ทสฺสนาย อนุตฺตราย สมฺโพธาย.\csromanspacer{} อยํ โข มํ ราชกุมาร ทุติยา อุปมา ปฏิภาสิ อนจฺฉริยา ปุพฺเพ อสฺสุตปุพฺพา.}

\proseitem{331}{อปราปิ โข มํ ราชกุมาร ตติยา อุปมา ปฏิภาสิ อนจฺฉริยา ปุพฺเพ อสฺสุตปุพฺพา\csromandash{}}

\prose{เสยฺยถาปิ ราชกุมาร สุกฺขํ กฏฺฐํ โกฬาปํ อารกา อุทกา ถเล นิกฺขิตฺตํ.\csromanspacer{} อถ ปุริโส อาคจฺเฉยฺย อุตฺตรารณิํ อาทาย “อคฺคิํ อภินิพฺพตฺเตสฺสามิ เตโช ปาตุกริสฺสามี”ติ.\csromanspacer{} ตํ กิํ มญฺญสิ ราชกุมาร, อปิ นุ โส ปุริโส อมุํ สุกฺขํ กฏฺฐํ โกฬาปํ อารกา อุทกา ถเล นิกฺขิตฺตํ อุตฺตรารณิํ อาทาย อภิมนฺเถนฺโต อคฺคิํ อภินิพฺพตฺเตยฺย เตโช ปาตุกเรยฺยาติ.\csromanspacer{} เอวํ ภนฺเต.\csromanspacer{} ตํ กิสฺส เหตุ, อทุํ หิ ภนฺเต สุกฺขํ กฏฺฐํ โกฬาปํ, ตญฺจ ปน อารกา อุทกา ถเล นิกฺขิตฺตนฺติ.\csromanspacer{} เอวเมว โข ราชกุมาร เย หิ เกจิ สมณา วา พฺราหฺมณา วา กาเยน เจว จิตฺเตน จ กาเมหิ วูปกฏฺฐา วิหรนฺติ.\csromanspacer{} โย จ เนสํ กาเมสุ กามจฺฉนฺโท กามสฺเนโห กามมุจฺฉา กามปิปาสา กามปริฬาโห, โส จ อชฺฌตฺตํ สุปฺปหีโน โหติ สุปฺปฏิปฺปสฺสทฺโธ.\csromanspacer{} โอปกฺกมิกา เจปิ เต โภนฺโต สมณพฺราหฺมณา ทุกฺขา ติพฺพา ขรา กฏุกา เวทนา เวทยนฺติ, ภพฺพาว เต ญาณาย ทสฺสนาย อนุตฺตราย สมฺโพธาย.\csromanspacer{} โน เจปิ เต โภนฺโต สมณพฺราหฺมณา โอปกฺกมิกา ทุกฺขา ติพฺพา ขรา กฏุกา}

\vspace{\tipitakaparskip}
\csromancenter{โพธิราชกุมารสุตฺต (85) เวทนา เวทยนฺติ, ภพฺพาว เต ญาณาย ทสฺสนาย อนุตฺตราย สมฺโพธาย.\csromanspacer{} อยํ โข มํ ราชกุมาร ตติยา อุปมา ปฏิภาสิ อนจฺฉริยา ปุพฺเพ อสฺสุตปุพฺพา.\csromanspacer{} อิมา โข มํ ราชกุมาร ติสฺโส อุปมา ปฏิภํสุ อนจฺฉริยา ปุพฺเพ อสฺสุตปุพฺพา.}
\proseitem{332}{\csromanfolio{285}ตสฺส มยฺหํ ราชกุมาร เอตทโหสิ “ยํนูนาหํ ทนฺเตภิทนฺตมาธาย\footnote{ปสฺส ม 1 วิตกฺกสณฺฐานสุตฺเต (170) ปิฏฺเฐ.} ชิวฺหาย ตาลุํ อาหจฺจ เจตสา จิตฺตํ อภินิคฺคเณฺหยฺยํ อภินิปฺปีเฬยฺยํ อภิสนฺตาเปยฺยนฺ”ติ, โส โข อหํ ราชกุมาร ทนฺเตภิทนฺตมาธาย ชิวฺหาย ตาลุํ อาหจฺจ เจตสา จิตฺตํ อภินิคฺคณฺหามิ อภินิปฺปีเฬมิ อภิสนฺตาเปมิ.\csromanspacer{} ตสฺส มยฺหํ ราชกุมาร ทนฺเตภิทนฺตมาธาย ชิวฺหาย ตาลุํ อาหจฺจ เจตสา จิตฺตํ อภินิคฺคณฺหาโต อภินิปฺปีฬยโต อภิสนฺตาปยโต กจฺเฉหิ เสทา มุจฺจนฺติ.\csromanspacer{} เสยฺยถาปิ ราชกุมาร พลวา ปุริโส ทุพฺพลตรํ ปุริสํ สีเส วา คเหตฺวา ขนฺเธ วา คเหตฺวา อภินิคฺคเณฺหยฺย อภินิปฺปีเฬยฺย อภิสนฺตาเปยฺย.\csromanspacer{} เอวเมว โข เม ราชกุมาร ทนฺเตภิทนฺตมาธาย ชิวฺหาย ตาลุํ อาหจฺจ เจตสา จิตฺตํ อภินิคฺคณฺหโต อภินิปฺปีฬยโต อภิสนฺตาปยโต กจฺเฉหิ เสทา มุจฺจนฺติ.\csromanspacer{} อารทฺธํ โข ปน เม ราชกุมาร วีริยํ โหติ อสลฺลีนํ, อุปฏฺฐิตา สติ อสมฺมุฏฺฐา, สารทฺโธ จ ปน เม กาโย โหติ อปฺปฏิปฺปสฺสทฺโธ เตนว ทุกฺขปฺปธาเนน ปธานาภิตุนฺนสฺส สโต.}

\proseitem{333}{ตสฺส มยฺหํ ราชกุมาร เอตทโหสิ “ยํนูนาหํ อปฺปาณกํเยว ฌานํ ฌาเยยฺยนฺ”ติ, โส โข อหํ ราชกุมาร มุขโต จ นาสโต จ อสฺสาสปสฺสาเส อุปรุนฺธิํ.\csromanspacer{} ตสฺส มยฺหํ ราชกุมาร มุขโต จ นาสโต จ อสฺสาสปสฺสาเสสุ อุปรุทฺเธสุ กณฺณโสเตหิ วาตานํ นิกฺขมนฺตานํ อธิมตฺโต สทฺโท โหติ.\csromanspacer{} เสยฺยถาปิ นาม กมฺมารคคฺคริยา ธมมานาย อธิมตฺโต สทฺโท โหติ.\csromanspacer{} เอวเมว โข เม ราชกุมาร มุขโต จ นาสโต จ อสฺสาสปสฺสาเสสุ อุปรุทฺเธสุ กณฺณโสเตหิ วาตานํ นิกฺขมนฺตานํ อธิมตฺโต สทฺโท โหติ.\csromanspacer{} อารทฺธํ โข ปน เม ราชกุมาร วีริยํ โหติ อสลฺลีนํ, \csromanfolio{286}อุปฏฺฐิตา สติ อสมฺมุฏฺฐา, สารทฺโธ จ ปน เม กาโย โหติ อปฺปฏิปฺปสฺสทฺโธ เตเนว ทุกฺขปฺปธาเนน ปธานาภิตุนฺนสฺส สโต.}

\prose{ตสฺส มยฺหํ ราชกุมาร เอตทโหสิ “ยํนูนาหํ อปฺปาณกํเยว ฌานํ ฌาเยยฺยนฺ”ติ.\csromanspacer{} โส โข อหํ ราชกุมาร มุขโต จ นาสโต จ กณฺณโต จ อสฺสาสปสฺสาเส อุปรุนฺธิํ.\csromanspacer{} ตสฺส มยฺหํ ราชกุมาร มุขโต จ นาสโต จ กณฺณโต จ อสฺสาสปสฺสาเสสุ อุปรุทฺเธสุ อธิมตฺตา วาตา มุทฺธนิ อูหนนฺติ\footnote{อูหนฺติ (สี), โอหนนฺติ (สฺยา, กํ), อุหนนฺติ (ก)}.\csromanspacer{} เสยฺยถาปิ ราชกุมาร พลวา ปุริโส ติเณฺหน สิขเรน มุทฺธนิ อภิมตฺเถยฺย\footnote{มุทฺธานํ อภิมนฺเถยฺย (สี, อิ), มุทฺธานํ อภิมตฺเถยฺย (สฺยา, กํ)}.\csromanspacer{} เอวเมว โข เม ราชกุมาร มุขโต จ นาสโต จ กณฺณโต จ อสฺสาสปสฺสาเสสุ อุปรุทฺเธสุ อธิมตฺตา วาตา มุทฺธนิ อูหนนฺติ.\csromanspacer{} อารทฺธํ โข ปน เม ราชกุมาร วีริยํ โหติ อสลฺลีนํ, อุปฏฺฐิตา สติ อสมฺมุฏฺฐา, สารทฺโธ จ ปน เม กาโย โหติ อปฺปฏิปฺปสฺสทฺโธ เตเนว ทุกฺขปฺปธาเนน ปธานาภิตุนฺนสฺส สโต.}

\prose{ตสฺส มยฺหํ ราชกุมาร เอตทโหสิ “ยํนูนาหํ อปฺปาณกํเยว ฌานํ ฌาเยยฺยนฺ”ติ.\csromanspacer{} โส โข อหํ ราชกุมาร มุขโต จ นาสโต จ กณฺณโต จ อสฺสาสปสฺสาเส อุปรุนฺธิํ.\csromanspacer{} ตสฺส มยฺหํ ราชกุมาร มุขโต จ นาสโต จ กณฺณโต จ อสฺสาสปสฺสาเสสุ อุปรุทฺเธสุ อธิมตฺตา สีเส สีสเวทนา โหนฺติ.\csromanspacer{} เสยฺยถาปิ ราชกุมาร พลวา ปุริโส ทเฬฺหน วรตฺตกฺขณฺเฑน\footnote{วรตฺตกพนฺธเนน (สี)} สีเส สีสเวฐํ ทเทยฺย.\csromanspacer{} เอวเมว โข เม ราชกุมาร มุขโต จ นาสโต จ กณฺณโต จ อสฺสาสปสฺสาเสสุ อุปรุทฺเธสุ อธิมตฺตา สีเส สีสเวทนา โหนฺติ.\csromanspacer{} อารทฺธํ โข ปน เม ราชกุมาร วีริยํ โหติ อสลฺลีนํ, อุปฏฺฐิตา สติ อสมฺมุฏฺฐา, สารทฺโธ จ ปน เม กาโย โหติ อปฺปฏิปฺปสฺสทฺโธ เตเนว ทุกฺขปฺปธาเนน ปธานาภิตุนฺนสฺส สโต.}

\prose{ตสฺส มยฺหํ ราชกุมาร เอตทโหสิ “ยํนูนาหํ อปฺปาณกํเยว ฌานํ ฌาเยยฺยนฺ”ติ.\csromanspacer{} โส โข อหํ ราชกุมาร มุขโต จ นาสโต จ กณฺณโต จ อสฺสาสปสฺสาเส อุปรุนฺธิํ.\csromanspacer{} ตสฺส มยฺหํ ราชกุมาร มุขโต จ นาสโต จ กณฺณโต จ อสฺสาสปสฺสาเสสุ}

\vspace{\tipitakaparskip}
\csromancenter{โพธิราชกุมารสุตฺต (85) อุปรุทฺเธสุ อธิมตฺตา วาตา กุจฺฉิํ ปริกนฺตนฺติ.\csromanspacer{} เสยฺยถาปิ ราชกุมาร ทกฺโข โคฆาตโก วา โคฆาตกนฺเตวาสี วา ติเณฺหน โควิกนฺตเนน กุจฺฉิํ ปริกนฺเตยฺย.\csromanspacer{} เอวเมว โข เม ราชกุมาร มุขโต จ นาสโต จ กณฺณโต จ อสฺสาสปสฺสาเสสุ อุปรุทฺเธสุ อธิมตฺตา วาตา กุจฺฉิํ ปริกนฺตนฺติ.\csromanspacer{} อารทฺธํ โข ปน เม ราชกุมาร วีริยํ โหติ อสลฺลีนํ, อุปฏฺฐิตา สติ อสมฺมุฏฺฐา, สารทฺโธ จ ปน เม กาโย โหติ อปฺปฏิปฺปสฺสทฺโธ เตเนว ทุกฺขปฺปธาเนน ปธานาภิตุนฺนสฺส สโต.}
\prose{\csromanfolio{287}ตสฺส มยฺหํ ราชกุมาร เอตทโหสิ “ยํนูนาหํ อปฺปาณกํเยว ฌานํ ฌาเยยฺยนฺ”ติ.\csromanspacer{} โส โข อหํ ราชกุมาร มุขโต จ นาสโต จ กณฺณโต จ อสฺสาสปสฺสาเส อุปรุนฺธิํ.\csromanspacer{} ตสฺส มยฺหํ ราชกุมาร มุขโต จ นาสโต จ กณฺณโต จ อสฺสาสปสฺสาเสสุ อุปรุทฺเธสุ อธิมตฺโต กายสฺมิํ ฑาโห โหติ.\csromanspacer{} เสยฺยถาปิ ราชกุมาร เทฺว พลวนฺโต ปุริสา ทุพฺพลตรํ ปุริสํ นานาพาหาสุ คเหตฺวา องฺคารกาสุยา สนฺตาเปยฺยุํ สมฺปริตาเปยฺยุํ.\csromanspacer{} เอวเมว โข เม ราชกุมาร มุขโต จ นาสโต จ กณฺณโต จ อสฺสาสปสฺสาเสสุ อุปรุทฺเธสุ อธิมตฺโต กายสฺมิํ ฑาโห โหติ.\csromanspacer{} อารทฺธํ โข ปน เม ราชกุมาร วีริยํ โหติ อสลฺลีนํ, อุปฏฺฐิตา สติ อสมฺมุฏฺฐา, สารทฺโธ จ ปน เม กาโย โหติ อปฺปฏิปฺปสฺสทฺโธ เตเนว ทุกฺขปฺปธาเนน ปธานาภิตุนฺนสฺส สโต.}

\prose{อปิสฺสุ มํ ราชกุมาร เทวตา ทิสฺวา เอวมาหํสุ “กาลงฺกโต สมโณ โคตโม”ติ.\csromanspacer{} เอกจฺจา เทวตา เอวมาหํสุ “น กาลงฺกโต สมโณ โคตโม, อปิ จ กาลงฺกโรตี”ติ.\csromanspacer{} เอกจฺจา เทวตา เอวมาหํสุ “น กาลงฺกโต สมโณ โคตโม, นาปิ กาลงฺกโรติ, อรหํ สมโณ โคตโม, วิหาโรเตฺวว โส\footnote{วิหาโรเตฺวเวโส (สี)} อรหโต เอวรูโป โหตี”ติ\footnote{วิหาโรเตฺวเวโส อรหโต”ติ (?)}.}

\proseitem{334}{ตสฺส มยฺหํ ราชกุมาร เอตทโหสิ “ยํนูนาหํ สพฺพโส อาหารุปจฺเฉทาย ปฏิปชฺเชยฺยนฺ”ติ.\csromanspacer{} อถ โข มํ ราชกุมาร เทวตา \csromanfolio{288}อุปสงฺกมิตฺวา เอตทโวจุํ “มา โข ตฺวํ มาริส สพฺพโส อาหารุปจฺเฉทาย ปฏิปชฺชิ.\csromanspacer{} สเจ โข ตฺวํ มาริส สพฺพโส อาหารุปจฺเฉทาย ปฏิปชฺชิสฺสสิ, ตสฺส เต มยํ ทิพฺพํ โอชํ โลมกูเปหิ อชฺโฌหาเรสฺสาม\footnote{อชฺโฌหริสฺสาม (สฺยา, กํ, อิ, ก)}, ตาย ตฺวํ ยาเปสฺสสี”ติ.\csromanspacer{} ตสฺส มยฺหํ ราชกุมาร เอตทโหสิ “อหญฺเจว โข ปน สพฺพโส อชชฺชิตํ\footnote{อชทฺธุกํ (สี, อิ), ชทฺธุกํ (สฺยา, กํ) 3. อชฺโฌหเรยฺยุํ (สฺยา, กํ, อิ, ก) 4. สมฺผุสิโต (สฺยา, กํ), สํปุฏิโต (ก) สํผุฏิโตติ เอตฺถ สงฺกุจิโตติ อตฺโถ.} ปฏิชาเนยฺยํ, อิมา จ เม เทวตา ทิพฺพํ โอชํ โลมกูเปหิ อชฺโฌหาเรยฺยุํ3, ตาย จาหํ ยาเปยฺยํ, ตํ มมสฺส มุสา”ติ.\csromanspacer{} โส โข อหํ ราชกุมาร ตา เทวตา ปจฺจาจิกฺขามิ หลนฺติ วทามิ. ตสฺส มยฺหํ ราชกุมาร เอตทโหสิ “ยํนูนาหํ โถกํ โถกํ อาหารํ อาหาเรยฺยํ ปสตํ ปสตํ, ยทิ วา มุคฺคยูสํ ยทิ วา กุลตฺถยูสํ ยทิ วา กฬายยูสํ ยทิ วา หเรณุกยูสนฺ”ติ.\csromanspacer{} โส โข อหํ ราชกุมาร โถกํ โถกํ อาหารํ อาหาเรสิํ ปสตํ ปสตํ, ยทิ วา มุคฺคยูสํ ยทิ วา กุลตฺถยูสํ ยทิ วา กฬายยูสํ ยทิ วา หเรณุกยูสํ.\csromanspacer{} ตสฺส มยฺหํ ราชกุมาร โถกํ โถกํ อาหารํ อาหารยโต ปสตํ ปสตํ, ยทิ วา มุคฺคยูสํ ยทิ วา กุลตฺถยูสํ ยทิ วา กฬายยูสํ ยทิ วา หเรณุกยูสํ, อธิมตฺตกสิมานํ ปตฺโต กาโย โหติ เสยฺยถาปิ นาม อาสีติกปพฺพานิ วา กาฬปพฺพานิ วา.\csromanspacer{} เอวเมวสฺสุ เม องฺคปจฺจงฺคานิ ภวนฺติ ตาเยวปฺปาหารตาย.\csromanspacer{} เสยฺยถาปิ นาม โอฏฺฐปทํ.\csromanspacer{} เอวเมวสฺสุ เม อานิสทํ โหติ ตาเยวปฺปาหารตาย.\csromanspacer{} เสยฺยถาปิ นาม วฏฺฏนาวฬี.\csromanspacer{} เอวเมวสฺสุ เม ปิฏฺฐิกณฺฏโก อุณฺณตาวนโต โหติ ตาเยวปฺปาหารตาย.\csromanspacer{} เสยฺยถาปิ นาม ชรสาลาย โคปานสิโย โอลุคฺควิลุคฺคา ภวนฺติ.\csromanspacer{} เอวเมวสฺสุ เม ผาสุฬิโย โอลุคฺควิลุคฺคา ภวนฺติ ตาเยวปฺปาหา- รตาย.\csromanspacer{} เสยฺยถาปิ นาม คมฺภีเร อุทปาเน อุทกตารกา คมฺภีรคตา โอกฺขายิกา ทิสฺสนฺติ.\csromanspacer{} เอวเมวสฺสุ เม อกฺขิกูเปสุ อกฺขิตารกา คมฺภีรคตา โอกฺขายิกา ทิสฺสนฺติ ตาเยวปฺปาหารตาย.\csromanspacer{} เสยฺยถาปิ นาม ติตฺตกาลาพุ อามกจฺฉินฺโน วาตาตเปน สํผุฏิโต4 โหติ สมฺมิลาโต.\csromanspacer{} เอวเมวสฺสุ เม สีสจฺฉวิ สํผุฏิตา โหติ สมฺมิลาตา ตาเยวปฺปาหารตาย.\csromanspacer{} โส}

\vspace{\tipitakaparskip}
\csromancenter{โพธิราชกุมารสุตฺต (85) โข อหํ ราชกุมาร “อุทรจฺฉวิํ ปริมสิสฺสามี”ติ ปิฏฺฐิกณฺฏกํเยว ปริคฺคณฺหามิ, “ปิฏฺฐกณฺฏกํ ปริมสิสฺสามี”ติ อุทรจฺฉวิํ เยว ปริคฺคณฺหามิ, ยาวสฺสุ เม ราชกุมาร อุทรจฺฉวิ ปิฏฺฐิกณฺฏกํ อลฺลีนา โหติ ตาเยวปฺปาหารตาย.\csromanspacer{} โส โข อหํ ราชกุมาร “วจฺจํ วา มุตฺตํ วา กริสฺสามี”ติ ตตฺเถว อวกุชฺโช ปปตามิ ตาเยวปฺปาหารตาย.\csromanspacer{} โส โข อหํ ราชกุมาร อิมเมว กายํ อสฺสาเสนฺโต ปาณินา คตฺตานิ อนุมชฺชามิ.\csromanspacer{} ตสฺส มยฺหํ ราชกุมาร ปาณินา คตฺตานิ อนุมชฺชโต ปูติมูลานิ โลมานิ กายสฺมา ปปตนฺติ ตาเยวปฺปาหารตาย.\csromanspacer{} อปิสฺสุ มํ ราชกุมาร มนุสฺสา ทิสฺวา เอวมาหํสุ “กาโฬ สมโณ โคตโม”ติ.\csromanspacer{} เอกจฺเจ มนุสฺสา เอวมาหํสุ “น กาโฬ สมโณ โคตโม, สาโม สมโณ โคตโม”ติ.\csromanspacer{} เอกจฺเจ มนุสฺสา เอวมาหํสุ “น กาโฬ สมโณ นปิ สาโม, มงฺคุรจฺฉวิ สมโณ โคตโม”ติ.\csromanspacer{} ยาวสฺสุ เม ราชกุมาร ตาว ปริสุทฺโธ ฉวิวณฺโณ ปริโยทาโต อุปหโต โหติ ตาเยวปฺปาหารตาย.}
\proseitem{335}{\csromanfolio{289}ตสฺส มยฺหํ ราชกุมาร เอตทโหสิ “เย โข เกจิ อตีตมทฺธานํ สมณา วา พฺราหฺมณา วา โอปกฺกมิกา ทุกฺขา ติพฺพา\footnote{ติปฺปา (สี, อิ)} ขรา กฏุกา เวทนา เวทยิํสุ, เอตาวปรมํ นยิโต ภิยฺโย.\csromanspacer{} เยปิ หิ เกจิ อนาคตมทฺธานํ สมณา วา พฺราหฺมณา วา โอปกฺกมิกา ทุกฺขา ติพฺพา ขรา กฏุกา เวทนา เวทยิสฺสนฺติ, เอตาวปรมํ นยิโต ภิยฺโย.\csromanspacer{} เยปิ หิ เกจิ เอตรหิ สมณา วา พฺราหฺมณา วา โอปกฺกมิกา ทุกฺขา ติพฺพา ขรา กฏุกา เวทนา เวทยนฺติ, เอตาวปรมํ นยิโต ภิยฺโย.\csromanspacer{} น โข ปนาหํ อิมาย กฏุกาย ทุกฺกรการิกาย อธิคจฺฉามิ อุตฺตริ มนุสฺสธมฺมา อลมริยญาณทสฺสนวิเสสํ.\csromanspacer{} สิยา นุ โข อญฺโญ มคฺโค โพธายา”ติ.\csromanspacer{} ตสฺส มยฺหํ ราชกุมาร เอตทโหสิ “อภิชานามิ โข ปนาหํ ปิตุ อ กมฺมนฺเต สีตาย ชมฺพุจฺฉายาย นิสินฺโน วิวิจฺเจว กาเมหิ วิวิจฺจ อกุสเลหิ ธมฺเมหิ สวิตกฺกํ สวิจารํ วิเวกชํ ปีติสุขํ ปฐมํ ฌานํ อุปสมฺปชฺช วิหริตา.\csromanspacer{} สิยา นุ โข เอโส มคฺโค โพธายา”ติ.\csromanspacer{} ตสฺส มยฺหํ ราชกุมาร สตานุสาริ วิญฺญาณํ อโหสิ “เอเสว มคฺโค โพธายา”ติ.\csromanspacer{} ตสฺส มยฺหํ ราชกุมาร เอตทโหสิ “กิํ นุ โข อหํ ตสฺส สุขสฺส ภายามิ, ยํ ตํ สุขํ อญฺญเตฺรว กาเมหิ อญฺญตฺร \csromanfolio{290}อกุสเลหิ ธมฺเมหี”ติ.\csromanspacer{} ตสฺส มยฺหํ ราชกุมาร เอตทโหสิ “น โข อหํ ตสฺส สุขสฺส ภายามิ, ยํ ตํ สุขํ อญฺญเตฺรว กาเมหิ อญฺญตฺร อกุสเลหิ ธมฺเมหี”ติ. ตสฺส มยฺหํ ราชกุมาร เอตทโหสิ “น โข ตํ สุกรํ สุขํ อธิคนฺตุํ เอวํ อธิมตฺตกสิมานํ ปตฺตกาเยน.\csromanspacer{} ยํนูนาหํ โอฬาริกํ อาหารํ อาหาเรยฺยํ โอทนกุมฺมาสนฺ”ติ.\csromanspacer{} โส โข อหํ ราชกุมาร โอฬาริกํ อาหารํ อาหาเรสิํ โอทนกุมฺมาสํ.\csromanspacer{} เตน โข ปน มํ ราชกุมาร สมเยน ปญฺจวคฺคิยา ภิกฺขู ปจฺจุปฏฺฐิตา โหนฺติ “ยํ โข สมโณ โคตโม ธมฺมํ อธิคมิสฺสติ, ตํ โน อาโรเจสฺสตี”ติ.\csromanspacer{} ยโต โข อหํ ราชกุมาร โอฬาริกํ อาหารํ อาหาเรสิํ โอทนกุมฺมาสํ, อถ เม เต ปญฺจวคฺคิยา ภิกฺขู นิพฺพิชฺช ปกฺกมิํสุ “พาหุลฺลิโก\footnote{พาหุลิโก (สี, อิ) สารตฺถฏีกาย สํฆเภทสิกฺขาปทวณฺณนาย สเมติ.} สมโณ โคตโม ปธานวิพฺภนฺโต อาวตฺโต พาหุลฺลายา”ติ.}

\proseitem{336}{โส โข อหํ ราชกุมาร โอฬาริกํ อาหารํ อาหาเรตฺวา พลํ คเหตฺวา วิวิจฺเจว กาเมหิ ฯเปฯ ปฐมํ ฌานํ อุปสมฺปชฺช วิหาสิํ.\csromanspacer{} วิตกฺกวิจารานํ วูปสมา ทุติยํ ฌานํ.\csromanspacer{} ตติยํ ฌานํ.\csromanspacer{} จตุตฺถํ ฌานํ อุปสมฺปชฺช วิหาสิํ.\csromanspacer{} โส เอวํ สมาหิเต จิตฺเต ปริสุทฺเธ ปริโยทาเต อนงฺคเณ วิคตูปกฺกิเลเส มุทุภูเต กมฺมนิเย ฐิเต อาเนญฺชปฺปตฺเต ปุพฺเพนิวาสานุสฺสติญาณาย จิตฺตํ อภินินฺนาเมสิํ.\csromanspacer{} โส อเนกวิหิตํ ปุพฺเพนิวาสํ อนุสฺสรามิ.\csromanspacer{} เสยฺยถิทํ, เอกมฺปิ ชาติํ เทฺวปิ ชาติโย ฯเปฯ อิติ สาการํ ส-อุทฺเทสํ อเนกวิหิตํ ปุพฺเพนิวาสํ อนุสฺสรามิ.\csromanspacer{} อยํ โข เม ราชกุมาร รตฺติยา ปฐเม ยาเม ปฐมา วิชฺชา อธิคตา, อวิชฺชา วิหตา, วิชฺชา อุปฺปนฺนา, ตโม วิหโต, อาโลโก อุปฺปนฺโน, ยถา ตํ อปฺปมตฺตสฺส อาตาปิโน ปหิตตฺตสฺส วิหรโต.}

\prose{โส เอวํ สมาหิเต จิตฺเต ปริสุทฺเธ ปริโยทาเต อนงฺคเณ วิคตูกฺกิเลเส มุทุภูเต กมฺมนิเย ฐิเต อาเนญฺชปฺปตฺเต สตฺตานํ จุตูปปาตญาณาย จิตฺตํ อภินินฺนาเมสิํ.\csromanspacer{} โส ทิพฺเพน จกฺขุนา วิสุทฺเธน อติกฺกนฺตมานุสเกน สตฺเต ปสฺสามิ จวมาเน อุปปชฺชมาเน หีเน ปณีเต สุวณฺเณ ทุพฺพณฺเณ สุคเต ทุคฺคเต, ยถากมฺมูปเค สตฺเต}

\vspace{\tipitakaparskip}
\csromancenter{โพธิราชกุมารสุตฺต (85) ปชานามิ ฯเปฯ.\csromanspacer{} อหํ โข เม ราชกุมาร รตฺติยา มชฺฌิเม ยาเม ทุติยา วิชฺชา อธิคตา, อวิชฺชา วิหตา, วิชฺชา อุปฺปนฺนา, ตโม วิหโต, อาโลโก อุปฺปนฺโน, ยถา ตํ อปฺปมตฺตสฺส อาตาปิโน ปหิตตฺตสฺส วิหรโต.}
\prose{\csromanfolio{291}โส เอวํ สมาหิเต จิตฺเต ปริสุทฺเธ ปริโยทาเต อนงฺคเณ วิคตูปกฺกิเลเส มุทุภูเต กมฺมนิเย ฐิเต อาเนญฺชปฺปตฺเต อาสวานํ ขยญาณาย จิตฺตํ อภินินฺนาเมสิํ.\csromanspacer{} โส อิทํ ทุกฺขนฺติ ยถาภูตํ อพฺภญฺญาสิํ ฯเปฯ อยํ ทุกฺขนิโรธคามินี ปฏิปทาติ ยถาภูตํ อพฺภญฺญาสิํ, อิเม อาสวาติ ยถาภูตํ อพฺภญฺญาสิํ ฯเปฯ อยํ อาสวนิโรธคามินี ปฏิปทาติ ยถาภูตํ อพฺภญฺญาสิํ.\csromanspacer{} ตสฺส เม เอวํ ชานโต เอวํ ปสฺสโต กามาสวาปิ จิตฺตํ วิมุจฺจิตฺถ, ภวาสวาปิ จิตฺตํ วิมุจฺจิตฺถ, อวิชฺชาสวาปิ จิตฺตํ วิมุจฺจิตฺถ, วิมุตฺตสฺมิํ “วิมุตฺตมฺ”อิติ ญาณํ อโหสิ, “ขีณา ชาติ, วุสิตํ พฺรหฺมจริยํ, กตํ กรณียํ, นาปรํ อิตฺถตฺตายา”ติ อพฺภญฺญาสิํ.\csromanspacer{} อยํ โข เม ราชกุมาร รตฺติยา ปจฺฉิเม ยาเม ตติยา วิชฺชา อธิคตา, อวิชฺชา วิหตา, วิชฺชา อุปฺปนฺนา, ตโม วิหโต, อาโลโก อุปฺปนฺโน, ยถา ตํ อปฺปมตฺตสฺส อาตาปิโน ปหิตตฺตสฺส วิหรโต.}

\proseitem{337}{ตสฺส มยฺหํ ราชกุมาร เอตทโหสิ “อธิคโต โข มฺยายํ ธมฺโม คมฺภีโร ทุทฺทโส ทุรนุโพโธ สนฺโต ปณีโต อตกฺกาวจโร นิปุโณ ปณฺฑิตเวทนีโย, อาลยรามา โข ปนายํ ปชา อาลยรตา อาลยสมฺมุทิตา, อาลยรามาย โข ปน ปชาย อาลยรตาย อาลยสมฺมุทิตาย ทุทฺทสํ อิทํ ฐานํ ยทิทํ อิทปฺปจฺจยตาปฏิจฺจสมุปฺปาโท, อิทมฺปิ โข ฐานํ ทุทฺทสํ ยทิทํ สพฺพสงฺขารสมโถ สพฺพูปธิปฏินิสฺสคฺโค ตณฺหากฺขโย วิราโค นิโรโธ นิพฺพานํ.\csromanspacer{} อหญฺเจว โข ปน ธมฺมํ เทเสยฺยํ ปเร จ เม น อาชาเนยฺยุํ, โส มมสฺส กิลมโถ, สา มมสฺส วิเหสา”ติ, อปิสฺสุ มํ ราชกุมาร อิมา อนจฺฉริยา คาถาโย ปฏิภํสุ ปุพฺเพ อสฺสุตปุพฺพา\csromandash{}}

\csromangathameasure{“กิจฺเฉน เม อธิคตํ, หลํ ทานิ ปกาสิตุํ. \\ ราคโทสปเรเตหิ, นายํ ธมฺโม สุสมฺพุโธ. \\ ปฏิโสตคามิํ นิปุณํ, คมฺภีรํ ทุทฺทสํ อณุํ. \\ ราครตฺตา น ทกฺขนฺติ, ตโมขนฺเธน อาวุฏา”ติ.}
\csromangathagroup{\csromangathastanza{“กิจฺเฉน เม อธิคตํ, หลํ ทานิ ปกาสิตุํ. \\ ราคโทสปเรเตหิ, นายํ ธมฺโม สุสมฺพุโธ.}\csromangathastanza{ปฏิโสตคามิํ นิปุณํ, คมฺภีรํ ทุทฺทสํ อณุํ. \\ ราครตฺตา น ทกฺขนฺติ, ตโมขนฺเธน อาวุฏา\footnote{อาวฏา (สี), อาวุตา (สฺยา, กํ)}”ติ.}}

\prose{\csromanfolio{292}อิติห เม ราชกุมาร ปฏิสญฺจิกฺขโต อปฺโปสฺสุกฺกตาย จิตฺตํ นมติ, โน ธมฺมเทสนาย.}

\proseitem{338}{อถ โข ราชกุมาร พฺรหฺมุโน สหมฺปติสฺส มม เจตสา เจโตปริวิตกฺกมญฺญาย เอตทโหสิ “นสฺสติ วต โภ โลโก, วินสฺสติ วต โภ โลโก, ยตฺร หิ นาม ตถาคตสฺส อรหโต สมฺมาสมฺพุทฺธสฺส อปฺโปสฺสุกฺกตาย จิตฺตํ นมติ\footnote{นมิสฺสติ (?)}, โน ธมฺมเทสนายา”ติ.\csromanspacer{} อถ โข ราชกุมาร พฺรหฺมา สหมฺปติ เสยฺยถาปิ นาม พลวา ปุริโส สมิญฺชิตํ วา พาหํ ปสาเรยฺย, ปสาริตํ วา พาหํ สมิญฺเชยฺย, เอวเมว พฺรหฺมโลเก อนฺตรหิโต มม ปุรโต ปาตุรโหสิ.\csromanspacer{} อถ โข ราชกุมาร พฺรหฺมา สหมฺปติ เอกํสํ อุตฺตราสงฺคํ กริตฺวา เยนาหํ เตนญฺชลิํ ปณาเมตฺวา มํ เอตทโวจ “เทเสตุ ภนฺเต ภควา ธมฺมํ, เทเสตุ สุคโต ธมฺมํ, สนฺติ สตฺตา อปฺปรชกฺขชาติกา, อสฺสวนตาย ธมฺมสฺส ปริหายนฺติ, ภวิสฺสนฺติ ธมฺมสฺส อญฺญาตาโร”ติ.\csromanspacer{} อิทมโวจ ราชกุมาร พฺรหฺมา สหมฺปติ, อิทํ วตฺวา อถาปรํ เอตทโวจ\csromandash{}}

\csromangathameasure{“ปาตุรโหสิ มคเธสุ ปุพฺเพ, \\ ธมฺโม อสุทฺโธ สมเลหิ จินฺติโต. \\ อปาปุเรตํ อมตสฺส ทฺวารํ, \\ สุณนฺตุ ธมฺมํ วิมเลนานุพุทฺธํ.}
\csromangathagroup{\csromangathastanza{“ปาตุรโหสิ มคเธสุ ปุพฺเพ, \\ ธมฺโม อสุทฺโธ สมเลหิ จินฺติโต. \\ อปาปุเรตํ\footnote{อวาปุเรตํ (สี)} อมตสฺส ทฺวารํ, \\ สุณนฺตุ ธมฺมํ วิมเลนานุพุทฺธํ.}}

\nitthitam{เสเล ยถา ปพฺพตมุทฺธนิฏฺฐิโต,}
\vspace{\tipitakaparskip}
\csromancenter{ยถาปิ ปสฺเส ชนตํ สมนฺตโต.}
\vspace{0.5\baselineskip}
\csromangathameasure{ตถูปมํ ธมฺมมยํ สุเมธ, \\ ปาสาทมารุยฺห สมนฺตจกฺขุ. \\ โสกาวติณฺณํ ชนตมเปตโสโก, \\ อเวกฺขสฺสุ ชาติชราภิภูตํ. \\ อุฏฺเฐหิ วีร วิชิตสงฺคาม, \\ สตฺถวาห อณณ วิจร โลเก. \\ เทสสฺสุ ภควา ธมฺมํ, \\ อญฺญาตาโร ภวิสฺสนฺตี”ติ.}
\csromangathagroup{\csromangathastanza{ตถูปมํ ธมฺมมยํ สุเมธ, \\ ปาสาทมารุยฺห สมนฺตจกฺขุ. \\ โสกาวติณฺณํ\footnote{โสกาวกิณฺณํ (สฺยา)} ชนตมเปตโสโก, \\ อเวกฺขสฺสุ ชาติชราภิภูตํ.}\csromangathastanza{อุฏฺเฐหิ วีร วิชิตสงฺคาม, \\ สตฺถวาห อณณ\footnote{อนณ (สี, สฺยา, กํ, อิ, ก)} วิจร โลเก. \\ เทสสฺสุ\footnote{อนณ (สี, สฺยา, กํ, อิ, ก)} ภควา ธมฺมํ, \\ อญฺญาตาโร ภวิสฺสนฺตี”ติ.}}

\csromanheader{2}{5. โพธิราชกุมารสุตฺต (85)}
\proseitem{339}{\csromanfolio{293}อถ ขฺวาหํ ราชกุมาร พฺรหฺมุโน จ อชฺเฌสนํ วิทิตฺวา สตฺเตสุ จ การุญฺญตํ ปฏิจฺจ พุทฺธจกฺขุนา โลกํ โวโลเกสิํ, อทฺทสํ โข อหํ ราชกุมาร พุทฺธจกฺขุนา โลกํ โวโลเกนฺโต สตฺเต อปฺปรชกฺเข มหารชกฺเข ติกฺขินฺทฺริเย มุทินฺทฺริเย สฺวากาเร ทฺวากาเร สุวิญฺญาปเย ทุวิญฺญาปเย อปฺเปกจฺเจ ปรโลกวชฺชภยทสฺสาวิเน\footnote{ทสฺสาวิโน (สฺยา, กํ, ก)} วิหรนฺเต, อปฺเปกจฺเจ น ปรโลกวชฺชภยทสฺสาวิเน วิหรนฺเต, เสยฺยถาปิ นาม อุปฺปลินิยํ วา ปทุมินิยํ วา ปุณฺฑรีกินิยํ วา อปฺเปกจฺจานิ อุปฺปลานิ วา ปทุมานิ วา ปุณฺฑรีกานิ วา อุทเก ชาตานิ อุทเก สํวฑฺฒานิ อุทกานุคฺคตานิ อนฺโตนิมุคฺคโปสีนิ, อปฺเปกจฺจานิ อุปฺปลานิ วา ปทุมานิ วา ปุณฺฑรีกานิ วา อุทเก ชาตานิ อุทเก สํวฑฺฒานิ อุทกานุคฺคตานิ สโมทกํ ฐิตานิ, อปฺเปกจฺจานิ อุปฺปลานิ วา ปทุมานิ วา ปุณฺฑรีกานิ วา อุทเก ชาตานิ อุทเก สํวฑฺฒานิ อุทกา อจฺจุคฺคมฺม ฐิตานิ\footnote{ติฏฺฐนฺติ (สี, สฺยา, กํ, อิ)} อนุปลิตฺตานิ อุทเกน, เอวเมว โข อหํ ราชกุมาร พุทฺธจกฺขุนา โลกํ โวโลเกนฺโต อทฺทสํ สตฺเต อปฺปรชกฺเข มหารชกฺเข ติกฺขินฺทฺริเย มุทินฺทฺริเย สฺวากาเร ทฺวากาเร สุวิญฺญาปเย ทุวิญฺญาปเย อปฺเปกจฺเจ ปรโลกวชฺชภยทสฺสาวิเน วิหรนฺเต อปฺเปกจฺเจ น ปรโลกวชฺชภยทสฺสาวิเน วิหรนฺเต.\csromanspacer{} อถ ขฺวาหํ ราชกุมาร พฺรหฺมานํ สหมฺปติํ คาถาย ปจฺจภาสิํ\csromandash{}}

\csromangathameasure{“อปารุตา เตสํ อมตสฺส ทฺวารา, \\ เย โสตวนฺโต ปมุญฺจนฺตุ สทฺธํ. \\ วิหิํส สญฺญี ปคุณํ น ภาสิํ, \\ ธมฺมํ ปณีตํ มนุเชสุ พฺรหฺเม”ติ.}
\csromangathagroup{\csromangathastanza{“อปารุตา เตสํ อมตสฺส ทฺวารา, \\ เย โสตวนฺโต ปมุญฺจนฺตุ สทฺธํ. \\ วิหิํส สญฺญี ปคุณํ น ภาสิํ, \\ ธมฺมํ ปณีตํ มนุเชสุ พฺรหฺเม”ติ.}}

\proseitem{340}{อถ โข ราชกุมาร พฺรหฺมา สหมฺปติ “กตาวกาโส โขมฺหิ ภควตา ธมฺมเทสนายา”ติ มํ อภิวาเทตฺวา ปทกฺขิณํ กตฺวา ตตฺเถวนฺตรธายิ.}

\prose{ตสฺส มยฺหํ ราชกุมาร เอตทโหสิ “กสฺส นุ โข อหํ ปฐมํ ธมฺมํ เทเสยฺยํ, โก อิมํ ธมฺมํ ขิปฺปเมว อาชานิสฺสตี”ติ.\csromanspacer{} ตสฺส มยฺหํ ราชกุมาร เอตทโหสิ “อยํ โข อาฬาโร กาลาโม ปณฺฑิโต วิยตฺโต เมธาวี ทีฆรตฺตํ อปฺปรชกฺขชาติโก, ยํนูนาหํ อาฬารสฺส \csromanfolio{294}กาลามสฺส ปฐมํ ธมฺมํ เทเสยฺยํ, โส อิมํ ธมฺมํ ขิปฺปเมว อาชานิสฺสตี”ติ.\csromanspacer{} อถ โข มํ ราชกุมาร เทวตา อุปสงฺกมิตฺวา เอตทโวจ “สตฺตาหกาลงฺกโต ภนฺเต อาฬาโร กาลาโม”ติ, ญาณญฺจ ปน เม ทสฺสนํ อุทปาทิ “สตฺตาหกาลงฺกโต อาฬาโร กาลาโม”ติ.\csromanspacer{} ตสฺส มยฺหํ ราชกุมาร เอตทโหสิ “มหาชานิโย โข อาฬาโร กาลาโม, สเจ หิ โส อิมํ ธมฺมํ สุเณยฺย, ขิปฺปเมว อาชาเนยฺยา”ติ.\csromanspacer{} ตสฺส มยฺหํ ราชกุมาร เอตทโหสิ “กสฺส นุ โข อหํ ปฐมํ ธมฺมํ เทเสยฺยํ, โก อิมํ ธมฺมํ ขิปฺปเมว อาชานิสฺสตี”ติ.\csromanspacer{} ตสฺส มยฺหํ ราชกุมาร เอตทโหสิ “อยํ โข อุทโก รามปุตฺโต ปณฺฑิโต วิยตฺโต เมธาวี ทีฆรตฺตํ อปฺปรชกฺขชาติโก, ยํนูนาหํ อุทกสฺส รามปุตฺตสฺส ปฐมํ ธมฺมํ เทเสยฺยํ, โส อิมํ ธมฺมํ ขิปฺปเมว อาชานิสฺสตี”ติ.\csromanspacer{} อถ โข มํ ราชกุมาร เทวตา อุปสงฺกมิตฺวา เอตทโวจ “อภิโทสกาลงฺกโต ภนฺเต อุทโก รามปุตฺโต”ติ.\csromanspacer{} ญาณญฺจ ปน เม ทสฺสนํ อุทปาทิ “อภิโทสกาลงฺกโต อุทโก รามปุตฺโต”ติ.\csromanspacer{} ตสฺส มยฺหํ ราชกุมาร เอตทโหสิ “มหาชานิโย โข อุทโก รามปุตฺโต, สเจ หิ โส อิมํ ธมฺมํ สุเณยฺย, ขิปฺปเมว อาชาเนยฺยา”ติ.}

\proseitem{341}{ตสฺส มยฺหํ ราชกุมาร เอตทโหสิ “กสฺส นุ โข อหํ ปฐมํ ธมฺมํ เทเสยฺยํ, โก อิมํ ธมฺมํ ขิปฺปเมว อาชานิสฺสตี”ติ.\csromanspacer{} ตสฺส มยฺหํ ราชกุมาร เอตทโหสิ “พหุการา โข เม ปญฺจวคฺคิยา ภิกฺขู, เย มํ ปธานปหิตตฺตํ อุปฏฺฐหิํสุ, ยํนูนาหํ ปญฺจวคฺคิยานํ ภิกฺขูนํ ปฐมํ ธมฺมํ เทเสยฺยนฺ”ติ.\csromanspacer{} ตสฺส มยฺหํ ราชกุมาร เอตทโหสิ “กหํ นุ โข เอตรหิ ปญฺจวคฺคิยา ภิกฺขู วิหรนฺตี”ติ.\csromanspacer{} อทฺทสํ ขฺวาหํ ราชกุมาร ทิพฺเพน จกฺขุนา วิสุทฺเธน อติกฺกนฺตมานุสเกน ปญฺจวคฺคิเย ภิกฺขู พาราณสิยํ วิหรนฺเต อิสิปตเน มิคทาเย.\csromanspacer{} อถ ขฺวาหํ ราชกุมาร อุรุเวลายํ ยถาภิรนฺตํ วิหริตฺวา เยน พาราณสี เตน จาริกํ ปกฺกมิํ.}

\prose{อทฺทสา โข มํ ราชกุมาร อุปโก อาชีวโก อนฺตรา จ คยํ อนฺตรา จ โพธิํ อทฺธานมคฺคปฺปฏิปนฺนํ, ทิสฺวาน มํ เอตทโวจ “วิปฺปสนฺนานิ โข เต อาวุโส อินฺทฺริยานิ, ปริสุทฺโธ ฉวิวณฺโณ ปริโยทาโต, กํสิ ตฺวํ อาวุโส อุทฺทิสฺส ปพฺพชิโต, โก วา เต สตฺถา, กสฺส วา ตฺวํ ธมฺมํ โรเจสี”ติ.\csromanspacer{} เอวํ วุตฺเต อหํ ราชกุมาร อุปกํ อาชีวกํ คาถาติ อชฺฌภาสิํ\csromandash{}}

\csromanheader{2}{5. โพธิราชกุมารสุตฺต (85)}
\csromangathameasure{“สพฺพาภิภู สพฺพวิทูหมสฺมิ, \\ สพฺเพสุ ธมฺเมสุ อนูปลิตฺโต. \\ สพฺพญฺชโห ตณฺหากฺขเย วิมุตฺโต, \\ สยํ อภิญฺญาย กมุทฺทิเสยฺยํ. \\ น เม อาจริโย อตฺถิ, สทิโส เม น วิชฺชติ. \\ สเทวกสฺมิํ โลกสฺมิํ, นตฺถิ เม ปฏิปุคฺคโล. \\ อหํ หิ อรหา โลเก, อหํ สตฺถา อนุตฺตโร. \\ เอโกมฺหิ สมฺมาสมฺพุทฺโธ, สีติภูโตสฺมิ นิพฺพุโต. \\ ธมฺมจกฺกํ ปวตฺเตตุํ, คจฺฉามิ กาสินํ ปุรํ. \\ อนฺธีภูตสฺมิํ โลกสฺมิํ, อาหญฺฉํ อมตทุนฺทุภินฺ”ติ.}
\csromangathagroup{\csromangathastanza{\csromanfolio{295}“สพฺพาภิภู สพฺพวิทูหมสฺมิ, \\ สพฺเพสุ ธมฺเมสุ อนูปลิตฺโต. \\ สพฺพญฺชโห ตณฺหากฺขเย วิมุตฺโต, \\ สยํ อภิญฺญาย กมุทฺทิเสยฺยํ. \\ น เม อาจริโย อตฺถิ, สทิโส เม น วิชฺชติ. \\ สเทวกสฺมิํ โลกสฺมิํ, นตฺถิ เม ปฏิปุคฺคโล.}\csromangathastanza{อหํ หิ อรหา โลเก, อหํ สตฺถา อนุตฺตโร. \\ เอโกมฺหิ สมฺมาสมฺพุทฺโธ, สีติภูโตสฺมิ นิพฺพุโต.}\csromangathastanza{ธมฺมจกฺกํ ปวตฺเตตุํ, คจฺฉามิ กาสินํ ปุรํ. \\ อนฺธีภูตสฺมิํ\footnote{อนฺธภูตสฺมิํ (สี, สฺยา, อิ)} โลกสฺมิํ, อาหญฺฉํ\footnote{อาหญฺญิํ (สฺยา, กํ, ก)} อมตทุนฺทุภินฺ”ติ.}}

\prose{ยถา โข ตฺวํ อาวุโส ปฏิชานาสิ “อรหสิ อนนฺตชิโน”ติ.}

\csromangathameasure{“มาทิสา เว ชินา โหนฺติ, เย ปตฺตา อาสวกฺขยํ. \\ ชิตา เม ปาปกา ธมฺมา, ตสฺมาหมุปก ชิโน”ติ.}
\csromangathagroup{\csromangathastanza{“มาทิสา เว ชินา โหนฺติ, เย ปตฺตา อาสวกฺขยํ. \\ ชิตา เม ปาปกา ธมฺมา, ตสฺมาหมุปก\footnote{ตสฺมาหํ อุปกา (สี, สฺยา, กํ, อิ)} ชิโน”ติ.}}

\prose{เอวํ วุตฺเต ราชกุมาร อุปโก อาชีวโก “หุเปยฺยปาวุโส”ติ\footnote{หุเวยฺยปาวุโส (สี, อิ), หุเวยฺยาวุโส (สฺยา, กํ)} วตฺวา สีสํ โอกมฺเปตฺวา อุมฺมคฺคํ คเหตฺวา ปกฺกามิ.}

\proseitem{342}{อถ ขฺวาหํ ราชกุมาร อนุปุพฺเพน จาริกํ จรมาโน เยน พาราณสี อิสิปตนํ มิคทาโย, เยน ปญฺจคคฺคิยา ภิกฺขู เตนุปสงฺกมิํ, อทฺทสํสุ โข มํ ราชกุมาร ปญฺจวคฺคิยา ภิกฺขู ทูรโตว อาคจฺฉนฺตํ, ทิสฺวาน อญฺญมญฺญํ สณฺฐเปสุํ “อยํ โข อาวุโส สมโณ โคตโม อาคจฺฉติ, พาหุลฺลิโก ปธานวิพฺภนฺโต อาวตฺโต พาหุลฺลาย, โส เนว อภิวาเทตพฺโพ, น ปจฺจุฏฺฐาตพฺโพ, นาสฺส ปตฺตจีวรํ ปฏิคฺคเหตพฺพํ, อปิ จ โข อาสนํ ฐเปตพฺพํ, สเจ โส อากงฺขิสฺสติ, นิสีทิสฺสตี”ติ.\csromanspacer{} ยถา ยถา โข อหํ ราชกุมาร ปญฺจวคฺคิเย ภิกฺขู อุปสงฺกมิํ\footnote{อุปสงฺกมามิ (สี, อิ)}, ตถา ตถา ปญฺจวคฺคิยา ภิกฺขู นาสกฺขิํสุ สกาย กติกาย สณฺฐาตุํ.\csromanspacer{} อปฺเปกจฺเจ มํ ปจฺจุคฺคนฺตฺวา ปตฺตจีวรํ ปฏิคฺคเหสุํ, อปฺเปกจฺเจ อาสนํ ปญฺญเปสุํ, อปฺเปกจฺเจ ปาโททกํ อุปฏฺฐเปสุํ, อปิ จ โข มํ นาเมน จ \csromanfolio{296}อาวุโสวาเทน จ สมุทาจรนฺติ.\csromanspacer{} เอวํ วุตฺเต อหํ ราชกุมาร ปญฺจวคฺคิเย ภิกฺขู เอตทโวจํ “มา ภิกฺขเว ตถาคตํ นาเมน จ อาวุโสวาเทน จ สมุทาจรถ\footnote{สมุทาจริตฺถ (สี, สฺยา, กํ, อิ)}.\csromanspacer{} อรหํ ภิกฺขเว ตถาคโต สมฺมาสมฺพุทฺโธ, โอทหถ ภิกฺขเว โสตํ, อมตมธิคตํ อหมนุสาสามิ, อหํ ธมฺมํ เทเสมิ, ยถานุสิฏฺฐํ ตถา ปฏิปชฺชมานา นจิรสฺเสว, ยสฺสตฺถาย กุลปุตฺตา สมฺมเทว อคารสฺมา อนคาริยํ ปพฺพชนฺติ, ตทนุตฺตรํ พฺรหฺมจริยปริโยสานํ ทิฏฺเฐว ธมฺเม สยํ อภิญฺญา สจฺฉิกตฺวา อุปสมฺปชฺช วิหริสฺสถา”ติ.\csromanspacer{} เอวํ วุตฺเต ราชกุมาร ปญฺจวคฺคิยา ภิกฺขู มํ เอตทโวจุํ “ตายปิ โข ตฺวํ อาวุโส โคตม อิริยาย\footnote{จริยาย (สฺยา, กํ)} ตาย ปฏิปทาย ตาย ทุกฺกรการิกาย นาชฺฌคมา อุตฺตริ มนุสฺสธมฺมา อลมริยญาณทสฺสนวิเสสํ, กิํ ปน ตฺวํ เอตรหิ พาหุลฺลิโก ปธานวิพฺภนฺโต อาวตฺโต พหุลฺลาย อธิคมิสฺสสิ อุตฺตริ มนุสฺสธมฺมา อลมริยญาณทสฺสนวิเสสนฺ”ติ.\csromanspacer{} เอวํ วุตฺเต อหํ ราชกุมาร ปญฺจวคฺคิเย ภิกฺขู เอตทโวจํ “น ภิกฺขเว ตถาคโต พาหุลฺลิโก, น ปธานวิพฺภนฺโต, น อาวตฺโต พาหุลฺลาย, อรหํ ภิกฺขเว ตถาคโต สมฺมาสมฺพุทฺโธ, โอทหถ ภิกฺขเว โสตํ, อมตมธิคตํ อหมนุสาสามิ, อหํ ธมฺมํ เทเสมิ, ยถานุสิฏฺฐํ ตถา ปฏิปชฺชมานา นจิรสฺเสว, ยสฺสตฺถาย กุลปุตฺตา สมฺมเทว อคารสฺมา อนคาริยํ ปพฺพชนฺติ, ตทนุตฺตรํ พฺรหฺมจริยปริโยสานํ ทิฏฺเฐว ธมฺเม สยํ อภิญฺญา สจฺฉิกตฺวา อุปสมฺปชฺช วิหริสฺสถา”ติ.\csromanspacer{} ทุติยมฺปิ โข ราชกุมาร ปญฺจวคฺคิยา ภิกฺขู มํ เอตทโวจุํ “ตายปิ โข ตฺวํ อาวุโส โคตม อิริยาย ตาย ปฏิปทาย ตาย ทุกฺกรการิกาย นาชฺฌคมา อุตฺตริ มนุสฺสธมฺมา อลมริยญาณทสฺสนวิเสสํ, กิํ ปน ตฺวํ เอตรหิ พาหุลฺลิโก ปธานวิพฺภนฺโต อาวตฺโต พาหุลฺลาย อธิคมิสฺสสิ อุตฺตริ มนุสฺสธมฺมา อลมริยญาณทสฺสนวิเสสนฺ”ติ.\csromanspacer{} ทุติยมฺปิ โข อหํ ราชกุมาร ปญฺจวคฺคิเย ภิกฺขู เอเตทโวจํ “น ภิกฺขเว ตถาคโต พาหุลฺลิโก, น ปธานวิพฺภนฺโต, น อาวตฺโต พาหุลฺลาย, อรหํ ภิกฺขเว ตถาคโต สมฺมาสมฺพุทฺโธ, โอทหถ ภิกฺขเว โสตํ, อมตมธิคตํ อหมนุสาสามิ, อหํ ธมฺมํ เทเสมิ, ยถานุสิฏฺฐํ ตถา ปฏิปชฺชมานา นจิรสฺเสว, ยสฺสตฺถาย กุลปุตฺตา สมฺมเทว อคารสฺมา อนคาริยํ ปพฺพชนฺติ, ตทนุตฺตรํ พฺรหฺมจริยปริโยสานํ ทิฏฺเฐว ธมฺเม สยํ อภิญฺญา สจฺฉิกตฺวา อุปสมฺปชฺช}

\vspace{\tipitakaparskip}
\csromancenter{โพธิราชกุมารสุตฺต (85) วิหริสฺสถา”ติ.\csromanspacer{} ตติยมฺปิ โข ราชกุมาร ปญฺจวคฺคิยา ภิกฺขู มํ เอตทโวจุํ “ตายปิ โข ตฺวํ อาวุโส โคตม อิริยาย ตาย ปฏิปทาย ตาย ทุกฺกรการิกาย นาชฺฌคมา อุตฺตริ มนุสฺสธมฺมา อลมริยญาณทสฺสนวิเสสํ, กิํ ปน ตฺวํ เอตรหิ พาหุลฺลิโก ปธานวิพฺภนฺโต อาวตฺโต พาหุลฺลาย อธิคมิสฺสสิ อุตฺตริ มนุสฺสธมฺมา อลมริยญาณทสฺสนวิเสสนฺ”ติ.\csromanspacer{} เอวํ วุตฺเต อหํ ราชกุมาร ปญฺจวคฺคิเย ภิกฺขู เอตทโวจํ “อภิชานาถ เม โน ตุเมฺห ภิกฺขเว อิโต ปุพฺเพ เอวรูปํ ปภาวิตเมตนฺ”ติ\footnote{ภาสิตเมตนฺติ (สี, สฺยา, วินเยปิ)}.\csromanspacer{} โน เหตํ ภนฺเต.\csromanspacer{} อรหํ ภิกฺขเว ตถาคโต สมฺมาสมฺพุทฺโธ, โอทหถ ภิกฺขเว โสตํ, อมตมธิคตํ อหมนุสาสามิ, อหํ ธมฺมํ เทเสมิ, ยถานุสิฏฺฐํ ตถา ปฏิปชฺชมานา นจิรสฺเสว, ยสฺสตฺถาย กุลปุตฺตา สมฺมเทว อคารสฺมา อนคาริยํ ปพฺพชนฺติ, ตทนุตฺตรํ พฺรหฺมจริยปริโยสานํ ทิฏฺเฐว ธมฺเม สยํ อภิญฺญา สจฺฉิกตฺวา อุปสมฺปชฺช วิหริสฺสถาติ.}
\prose{\csromanfolio{297}อสกฺขิํ โข อหํ ราชกุมาร ปญฺจวคฺคิเย ภิกฺขู สญฺญาเปตุํ.\csromanspacer{} เทฺวปิ สุทํ ราชกุมาร ภิกฺขู โอวทามิ, ตโย ภิกฺขู ปิณฺฑาย จรนฺติ, ยํ ตโย ภิกฺขู ปิณฺฑาย จริตฺวา อาหรนฺติ, เตน ฉพฺพคฺคิยา\footnote{ฉพฺพคฺคา (สี, สฺยา, กํ), ฉพฺพคฺโค (อิ) 3. นายกํ (?)} ยาเปม.\csromanspacer{} ตโยปิ สุทํ ราชกุมาร ภิกฺขู โอวทามิ, เทฺว ภิกฺขู ปิณฺฑาย จรนฺติ, ยํ เทฺว ภิกฺขู ปิณฺฑาย จริตฺวา อาหรนฺติ, เตน ฉพฺพคฺคิยา ยาเปม.}

\proseitem{343}{อถ โข ราชกุมาร ปญฺจวคฺคิยา ภิกฺขู มยา เอวํ โอวทิยมานา เอวํ อนุสาสิยมานา นจิรสฺเสว, ยสฺสตฺถาย กุลปุตฺตา สมฺมเทว อคารสฺมา อนคาริยํ ปพฺพชนฺติ.\csromanspacer{} ตทนุตฺตรํ พฺรหฺมจริยปริโยสานํ ทิฏฺเฐว ธมฺเม สยํ อภิญฺญา สจฺฉิกตฺวา อุปสมฺปชฺช วิหริํสูติ.\csromanspacer{} เอวํ วุตฺเต โพธิ ราชกุมาโร ภควนฺตํ เอตทโวจ “กีว จิเรน นุ โข ภนฺเต ภิกฺขุ ตถาคตํ วินายกํ3 ลภมาโน, ยสฺสตฺถาย กุลปุตฺตา สมฺมเทว อคารสฺมา อนคาริยํ ปพฺพชนฺติ, ตทนุตฺตรํ พฺรหฺมจริยปริโยสานํ ทิฏฺเฐว ธมฺเม สยํ อภิญฺญา สจฺฉิกตฺวา อุปสมฺปชฺช วิหเรยฺยา”ติ.\csromanspacer{} เตน หิ ราชกุมาร ตํเยเวตฺถ ปฏิปุจฺฉิสฺสามิ, ยถา เต ขเมยฺย, ตถา นํ พฺยากเรยฺยาสิ.\csromanspacer{} ตํ กิํ มญฺญสิ ราชกุมาร, กุสโล ตฺวํ หตฺถารูเฬฺห\footnote{หตฺถารูเยฺห (สี, อิ) 5. องฺกุสคเณฺห (สฺยา, กํ)} องฺกุสคเยฺห5สิปฺเปติ.\csromanspacer{} เอวํ ภนฺเต กุสโล อหํ หตฺถารูเฬฺห องฺกุสคเยฺห \csromanfolio{298}สิปฺเปติ.\csromanspacer{} ตํ กิํ มญฺญสิ ราชกุมาร, อิธ ปุริโส อาคจฺเฉยฺย “โพธิ ราชกุมาโร หตฺถารูฬฺหํ องฺกุสคยฺหํ สิปฺปํ ชานาติ, ตสฺสาหํ สนฺติเก หตฺถารูฬฺหํ องฺกุสคยฺหํ สิปฺปํ สิกฺขิสฺสามี”ติ.\csromanspacer{} โส จสฺส อสฺสทฺโธ, ยาวตกํ สทฺเธน ปตฺตพฺพํ, ตํ น สมฺปาปุเณยฺย.\csromanspacer{} โส จสฺส พหฺวาพาโธ, ยาวตกํ อปฺปาพาเธน ปตฺตพฺพํ, ตํ น สมฺปาปุเณยฺย.\csromanspacer{} โส จสฺส สโฐ มายาวี, ยาวตกํ อสเฐน อมายาวินา ปตฺตพฺพํ, ตํ น สมฺปาปุเณยฺย.\csromanspacer{} โส จสฺส กุสีโต, ยาวตกํ อารทฺธวีริเยน ปตฺตพฺพํ, ตํ น สมฺปาปุเณยฺย.\csromanspacer{} โส จสฺส ทุปฺปญฺโญ, ยาวตกํ ปญฺญวตา ปตฺตพฺพํ, ตํ น สมฺปาปุเณยฺย.\csromanspacer{} ตํ กิํ มญฺญสิ ราชกุมาร, อปิ นุ โส ปุริโส ตว สนฺติเก หตฺถารูฬฺหํ องฺกุสคยฺหํ สิปฺปํ สิกฺเขยฺยาติ.\csromanspacer{} เอกเมเกนาปิ ภนฺเต องฺเคน สมนฺนาคโต โส ปุริโส น มม สนฺติเก หตฺถารูฬฺหํ องฺกุสคยฺหํ สิปฺปํ สิกฺเขยฺย, โก ปน วาโท ปญฺจหงฺเคหีติ.}

\proseitem{344}{ตํ กิํ มญฺญสิ ราชกุมาร, อิธ ปุริโส อาคจฺเฉยฺย “โพธิ ราชกุมาโร หตฺถารูฬฺหํ องฺกุสคยฺหํ สิปฺปํ ชานาติ, ตสฺสาหํ สนฺติเก หตฺถารูฬฺหํ องฺกุสคยฺหํ สิปฺปํ สิกฺขิสฺสามี”ติ.\csromanspacer{} โส จสฺส สทฺโธ, ยาวตกํ สทฺเธน ปตฺตพฺพํ, ตํ สมฺปาปุเณยฺย.\csromanspacer{} โส จสฺส อปฺปาพาโธ, ยาวตกํ อปฺปาพาเธน ปตฺตพฺพํ, ตํ สมฺปาปุเณยฺย.\csromanspacer{} โส จสฺส อสโฐ อมายาวี, ยาวตกํ อสเฐน อมายาวินา ปตฺตพฺพํ, ตํ สมฺปาปุเณยฺย.\csromanspacer{} โส จสฺส อารทฺธวีริโย, ยาวตกํ อารทฺธวีริเยน ปตฺตพฺพํ, ตํ สมฺปาปุเณยฺย.\csromanspacer{} โส จสฺส ปญฺญวา, ยาวตกํ ปญฺญวตา ปตฺตพฺพํ, ตํ สมฺปาปุเณยฺย.\csromanspacer{} ตํ กิํ มญฺญสิ ราชกุมาร, อปิ นุ โส ปุริโส ตว สนฺติเก หตฺถารูฬฺหํ องฺกุสคยฺหํ สิปฺปํ สิกฺเขยฺยาติ.\csromanspacer{} เอกเมเกนาปิ ภนฺเต องฺเคน สมนฺนาคโต โส ปุริโส มม สนฺติเก หตฺถารูฬฺหํ องฺกุสคยฺหํ สิปฺปํ สิกฺเขยฺย, โก ปน วาโท ปญฺจหงฺเคหีติ.\csromanspacer{} เอวเมว โข ราชกุมาร ปญฺจิมานิ ปธานิยงฺคานิ.\csromanspacer{} กตมานิ ปญฺจ.\csromanspacer{} อิธ ราชกุมาร ภิกฺขุ สทฺโธ โหติ, สทฺทหติ ตถาคตสฺส โพธิํ “อิติปิ โส ภควา อรหํ สมฺมาสมฺพุทฺโธ วิชฺชาจรณสมฺปนฺโน สุคโต โลกวิทู อนุตฺตโร ปุริสทมฺมสารถิ สตฺถา เทวมนุสฺสานํ พุทฺโธ ภควา”ติ.\csromanspacer{} อปฺปาพาโธ โหติ อปฺปาตงฺโก สมเวปากินิยา คหณิยา สมนฺนาคโต นาติสีตาย นาจฺจุณฺหาย มชฺฌิมาย ปธานกฺขมาย.\csromanspacer{} อสโฐ โหติ อมายาวี ยถาภูตํ อตฺตานํ อาวิกตฺตา}

\vspace{\tipitakaparskip}
\csromancenter{โพธิราชกุมารสุตฺต (85) สตฺถริ วา วิญฺญูสุ วา สพฺรหฺมจารีสุ.\csromanspacer{} อารทฺธวีริโย วิหรติ อกุสลานํ ธมฺมานํ ปหานาย กุสลานํ ธมฺมานํ อุปสมฺปทาย, ถามวา ทฬฺหปรกฺกโม อนิกฺขิตฺตธุโร กุสเลสุ ธมฺเมสุ.\csromanspacer{} ปญฺญวา โหติ อุทยตฺถคามินิยา ปญฺญาย สมนฺนาคโต อริยาย นิพฺเพธิกาย สมฺมาทุกฺขกฺขยคามินิยา.\csromanspacer{} อิมานิ โข ราชกุมาร ปญฺจ ปธานิยงฺคานิ.}
\proseitem{345}{\csromanfolio{299}อิเมหิ ราชกุมาร ปญฺจหิ ปธานิยงฺเคหิ สมนฺนาคโต ภิกฺขุ ตถาคตํ วินายกํ ลภมาโน, ยสฺสตฺถาย กุลปุตฺตา สมฺมเทว อคารสฺมา อนคาริยํ ปพฺพชนฺติ, ตทนุตฺตรํ พฺรหฺมจริยปริโยสานํ ทิฏฺเฐว ธมฺเม สยํ อภิญฺญา สจฺฉิกตฺวา อุปสมฺปชฺช วิหเรยฺย สตฺต วสฺสานิ.\csromanspacer{} ติฏฺฐนฺตุ ราชกุมาร สตฺต วสฺสานิ.\csromanspacer{} อิเมหิ ปญฺจหิ ปธานิยงฺเคหิ สมนฺนาคโต ภิกฺขุ ตถาคตํ วินายกํ ลภมาโน, ยสฺสตฺถาย กุลปุตฺตา สมฺมเทว อคารสฺมา อนคาริยํ ปพฺพชนฺติ, ตทนุตฺตรํ พฺรหฺมจริยปริโยสานํ ทิฏฺเฐว ธมฺเม สยํ อภิญฺญา สจฺฉิกตฺวา อุปสมฺปชฺช วิหเรยฺย ฉพฺพสฺสานิ.\csromanspacer{} ปญฺจ วสฺสานิ.\csromanspacer{} จตฺตาริ วสฺสานิ.\csromanspacer{} ตีณิ วสฺสานิ.\csromanspacer{} เทฺว วสฺสานิ.\csromanspacer{} เอกํ วสฺสํ.\csromanspacer{} ติฏฺฐตุ ราชกุมาร เอกํ วสฺสํ.\csromanspacer{} อิเมหิ ปญฺจหิ ปธานิยงฺเคหิ สมนฺนาคโต ภิกฺขุ ตถาคตํ วินายกํ ลภมาโน, ยสฺสตฺถาย กุลปุตฺตา สมฺมเทว อคารสฺมา อนคาริยํ ปพฺพชนฺติ, ตทนุตฺตรํ พฺรหฺมจริยปริโยสานํ ทิฏฺเฐว ธมฺเม สยํ อภิญฺญา สจฺฉิกตฺวา อุปสมฺปชฺช วิหรยฺย สตฺต มาสานิ.\csromanspacer{} ติฏฺฐนฺตุ ราชกุมาร สตฺต มาสานิ.\csromanspacer{} อิเมหิ ปญฺจหิ ปธานิยงฺเคหิ สมนฺนาคโต ภิกฺขุ ตถาคตํ วินายกํ ลภมาโน, ยสฺสตฺถาย กุลปุตฺตา สมฺมเทว อคารสฺมา อนคาริยํ ปพฺพชนฺติ, ตทนุตฺตรํ พฺรหฺมจริยปริโยสานํ ทิฏฺเฐว ธมฺเม สยํ อภิญฺญา สจฺฉิกตฺวา อุปสมฺปชฺช วิหเรยฺย ฉ มาสานิ.\csromanspacer{} ปญฺจ มาสานิ.\csromanspacer{} จตฺตาริ มาสานิ.\csromanspacer{} ตีณิ มาสานิ.\csromanspacer{} เทฺว มาสานิ.\csromanspacer{} เอกํ มาสํ.\csromanspacer{} อฑฺฒมาสํ.\csromanspacer{} ติฏฺฐตุ ราชกุมาร อฑฺฒมาโส.\csromanspacer{} อิเมหิ ปญฺจหิ ปธานิยงฺเคหิ สมนฺนาคโต ภิกฺขุ ตถาคตํ วินายกํ ลภามาโน, ยสฺสตฺถาย กุลปุตฺโต สมฺมเทว อคารสฺมา อนคาริยํ ปพฺพชนฺติ, ตทนุตฺตรํ พฺรหฺมจริยปริโยสานํ ทิฏฺเฐว ธมฺเม สยํ อภิญฺญา สจฺฉิกตฺวา อุปสมฺปชฺช วิหเรยฺย สตฺต รตฺตินฺทิวานิ.\csromanspacer{} ติฏฺฐนฺตุ ราชกุมาร สตฺต รตฺตินฺทิวานิ.\csromanspacer{} อิเมหิ ปญฺจหิ ปธานิยงฺเคหิ สมนฺนาคโต ภิกฺขุ ตถาคตํ วินายกํ ลภมาโน, ยสฺสตฺถาย กุลปุตฺตา สมฺมเทว อคารสฺมา อนคาริยํ ปพฺพชนฺติ, ตทนุตฺตรํ พฺรหฺมจริยปริโยสานํ ทิฏฺเฐว ธมฺเม สยํ อภิญฺญา \csromanfolio{300}สจฺฉิกตฺวา อุปสมฺปชฺช วิหเรยฺย ฉ รตฺตินฺทิวานิ.\csromanspacer{} ปญฺจ รตฺตินฺทิวานิ.\csromanspacer{} จตฺตาริ รตฺตินฺทิวานิ.\csromanspacer{} ตีณิ รตฺตินฺทิวานิ.\csromanspacer{} เทฺว รตฺถินฺทิวานิ.\csromanspacer{} เอกํ รตฺตินฺทิวํ.\csromanspacer{} ติฏฺฐตุ ราชกุมาร เอโก รตฺตินฺทิโว.\csromanspacer{} อิเมหิ ปญฺจหิ ปธานิยงฺเคหิ สมนฺนาคโต ภิกฺขุ ตถาคตํ วินายกํ ลภมาโน สายมนุสิฏฺโฐ ปาโต วิเสสํ อธิคมิสฺสติ, ปาตมนุสิฏฺโฐ สายํ วิเสสํ อธิคมิสฺสตีติ.\csromanspacer{} เอวํ วุตฺเต โพธิ ราชกุมาโร ภควนฺตํ เอตทโวจ “อโห พุทฺโธ อโห ธมฺโม อโห ธมฺมสฺส สฺวากฺขาตตา, ยตฺร หิ นาม สายมนุสิฏฺโฐ ปาโต วิเสสํ อธิคมิสฺสติ, ปาตมนุสิฏฺโฐ สายํ วิเสสํ อธิคมิสฺสตี”ติ.}

\proseitem{346}{เอวํ วุตฺเต สญฺชิกาปุตฺโต มาณโว โพธิํ ราชกุมารํ เอตทโวจ “เอวเมว ปนายํ ภวํ โพธิ ‘อโห พุทฺโธ อโห ธมฺโม อโห ธมฺมสฺส สฺวากฺขาตตา’ติ จ วเทติ\footnote{วเทสิ (สี), ปเวเทติ (สฺยา, กํ)}, อถ จ ปน น ตํ ภวนฺตํ โคตมํ สรณํ คจฺฉติ, ธมฺมญฺจ ภิกฺขุสํฆญฺจา”ติ.\csromanspacer{} มา เหวํ สมฺม สญฺชิกาปุตฺต อวจ, มา เหวํ สมฺม สญฺชิกาปุตฺต อวจ, สมฺมุขา เมตํ สมฺม สญฺชิกาปุตฺต อยฺยาย สุตํ, สมฺมุขา ปฏิคฺคหิตํ, เอกมิทํ สมฺม สญฺชิกาปุตฺต สมยํ ภควา โกสมฺพิยํ วิหรติ โฆสิตาราเม.\csromanspacer{} อถ โข เม อยฺยา กุจฺฉิมตี เยน ภควา เตนุปสงฺกมิ, อุปสงฺกมิตฺวา ภควนฺตํ อภิวาเทตฺวา เอกมนฺตํ นิสีทิ, เอกมนฺตํ นิสินฺนา โข เม อยฺยา ภควนฺตํ เอตทโวจ “โย เม อยํ ภนฺเต กุจฺฉิคโต กุมารโก วา กุมาริกา วา, โส ภควนฺตํ สรณํ คจฺฉติ ธมฺมญฺจ ภิกฺขุสํฆญฺจ, อุปาสกํ ตํ ภควา ธาเรตุ อชฺชตคฺเค ปาณุเปตํ สรณํ คตนฺ”ติ.\csromanspacer{} เอกมิทํ สมฺม สญฺชิกาปุตฺต สมยํ ภควา อิเธว ภคฺเคสุ วิหรติ สุสุมารคิเร เภสกฬาวเน มิคทาเย.\csromanspacer{} อถ โข มํ ธาติ องฺเกน หริตฺวา เยน ภควา เตนุปสงฺกมิ, อุปสงฺกมิตฺวา ภควนฺตํ อภิวาเทตฺวา เอกมนฺตํ อฏฺฐาสิ, เอกมนฺตํ ฐิตา โข มํ ธาติ ภควนฺตํ เอตทโวจ “อยํ ภนฺเต โพธิ ราชกุมาโร ภควนฺตํ สรณํ คจฺฉติ ธมฺมญฺจ ภิกฺขุสํฆญฺจ, อุปาสกํ ตํ ภควา ธาเรตุ อชฺชตคฺเค ปาณุเปตํ สรณํ คตนฺ”ติ.\csromanspacer{} เอสาหํ สมฺม สญฺชิกาปุตฺต ตติยกมฺปิ ภควนฺตํ สรณํ คจฺฉามิ ธมฺมญฺจ ภิกฺขุสํฆญฺจ, อุปาสกํ มํ ภควา ธาเรตุ อชฺชตคฺเค ปาณุเปตํ สรณํ คตนฺติ.}

\nitthitam{โพธิราชกุมารสุตฺตํ นิฏฺฐิตํ ปญฺจมํ.}
\csromanheader{2}{6. องฺคุลิมาลสุตฺต (86) \textbf{6. องฺคุลิมาลสุตฺต}}
\proseitem{347}{\csromanfolio{301}เอวํ เม สุตํ\csromandash{}เอกํ สมยํ ภควา สาวตฺถิยํ วิหรติ เชตวเน อนาถปิณฺฑิกสฺส อาราเม.\csromanspacer{} เตน โข ปน สมเยน รญฺโญ ปเสนทิสฺส โกสลสฺส วิชิเต โจโร องฺคุลิมาโล นาม โหติ ลุทฺโท โลหิตปาณิ หตปหเต นิวิฏฺโฐ อทยาปนฺโน ปาณภูเตสุ, เตน คามาปิ อคามา กตา, นิคมาปิ อนิคมา กตา, นชปทาปิ อชนปทา กตา, โส มนุสฺเส วธิตฺวา วธิตฺวา องฺคุลีนํ มาลํ ธาเรติ.\csromanspacer{} อถ โข ภควา ปุพฺพณฺหสมยํ นิวาเสตฺวา ปตฺตจีวรมาทาย สาวตฺถิํ ปิณฺฑาย ปาวิสิ สาวตฺถิยํ ปิณฺฑาย จริตฺวา ปจฺฉาภตฺตํ ปิณฺฑปาตปฏิกฺกนฺโต เสนาสนํ สํสาเมตฺวา ปตฺตจีวรมาทาย เยน โจโร องฺคุลิมาโล เตนทฺธานมคฺคํ ปฏิปชฺชิ.\csromanspacer{} อทฺทสาสุํ โข โคปาลกา ปสุปาลกา กสฺสกา ปถาวิโน ภควนฺตํ เยน โจโร องฺคุลิมาโล เตนทฺธานมคฺคปฏิปนฺนํ, ทิสฺวาน ภควนฺตํ เอตทโวจุํ “มา สมณ เอตํ มคฺคํ ปฏิปชฺชิ, เอตสฺมิํ สมณ มคฺเค โจโร องฺคุลิมาโล นาม ลุทฺโท โลหิตปาณิ หตปหเต นิวิฏฺโฐ อทยาปนฺโน ปาณภูเตสุ, เตน คามาปิ อคามา กตา, นิคมาปิ อนิคมา กตา, ชนปทาปิ อชนปทา กตา, โส มนุสฺเส วธิตฺวา วธิตฺวา องฺคุลีนํ มาลํ ธาเรติ, เอตญฺหิ สมณ มคฺคํ ทสปิ ปุริสา วีสมฺปิ ปุริสา ติํสมฺปิ ปุริสา จตฺตารีสมฺปิ ปุริสา ปญฺญาสมฺปิ ปุริสา สงฺกริตฺวา สงฺกริตฺวา\footnote{สํหริตฺวา สํหริตฺวา (สี, อิ), สงฺคริตฺวา (สฺยา, กํ)} ปฏิปชฺชนฺติ, เตปิ โจรสฺส องฺคุลิมาลสฺส หตฺถตฺถํ คจฺฉนฺตี”ติ.\csromanspacer{} เอวํ วุตฺเต ภควา ตุณฺหีภูโต อคมาสิ.\csromanspacer{} ทุติยมฺปิ โข โคปาลกา ฯเปฯ.\csromanspacer{} ตติยมฺปิ โข โคปาลกา ปสุปาลกา กสฺสกา ปถาวิโน ภควนฺตํ เอตทโวจุํ “มา สมณ เอตํ มคฺคํ ปฏิปชฺชิ, เอตสฺมิํ สมณ มคฺเค โจโร องฺคุลิมาโล นาม ลุทฺโท โลหิตปาณิ หตปหเต นิวิฏฺโฐ อทยาปนฺโน ปาณภูเตสุ, เตน คามาปิ อคามา กตา, นิคมาปิ อนิคมา กตา, ชนปทาปิ อชนปทา กตา, โสมนุสฺเส วธิตฺวา วธิตฺวา องฺคุลีนํ มาลํ ธาเรติ, เอตญฺหิ สมณ มคฺคํ ทสปิ ปุริสา วีสมฺปิ ปุริสา ติํสมฺปิ ปุริสา จตฺตารีสมฺปิ ปุริสา ปญฺญาสมฺปิ ปุริสา สงฺกริตฺวา สงฺกริตฺวา ปฏิปชฺชนฺติ, เตปิ โจรสฺส องฺคุลิมาลสฺส หตฺถตฺถํ คจฺฉนฺตี”ติ.}

\proseitem{348}{\csromanfolio{302}อถ โข ภควา ตุณฺหีภูโต อคมาสิ.\csromanspacer{} อทฺทสา โข โจโร องฺคุลิมาโล ภควนฺตํ ทูรโตว อาคจฺฉนฺตํ, ทิสฺวานสฺส เอตทโหสิ “อจฺฉริยํ วต โภ, อพฺภุตํ วต โภ, อิมํ หิ มคฺคํ ทสปิ ปุริสา วีสมฺปิ ปุริสา ติํสมฺปิ ปุริสา จตฺตารีสมฺปิ ปุริสา ปญฺญาสมฺปิ ปุริสา สงฺกริตฺวา สงฺกริตฺวา ปฏิปชฺชนฺติ, เตปิ มม หตฺถตฺถํ คจฺฉนฺติ.\csromanspacer{} อถ จ ปนายํ สมโณ เอโก อทุติโย ปสยฺห มญฺเญ อาคจฺฉติ, ยํนูนาหํ อิมํ สมณํ ชีวิตา โวโรเปยฺยนฺ”ติ.\csromanspacer{} อถ โข โจโร องฺคุลิมาโล อสิจมฺมํ คเหตฺวา ธนุกลาปํ สนฺนยฺหิตฺวา ภควนฺตํ ปิฏฺฐิโต ปิฏฺฐิโต อนุพนฺธิ.\csromanspacer{} อถ โข ภควา ตถารูปํ อิทฺธาภิสงฺขารํ อภิสงฺขาสิ\footnote{อภิสงฺขาเรสิ (สฺยา, กํ, ก)}, ยถา โจโร องฺคุลิมาโล ภควนฺตํ ปกติยา คจฺฉนฺตํ สพฺพถาเมน คจฺฉนฺโต น สกฺโกติ สมฺปาปุณิตุํ.\csromanspacer{} อถ โข โจรสฺส องฺคุลิมาลสฺส เอตทโหสิ “อจฺฉริยํ วต โภ, อพฺภุตํ วต โภ, อหญฺหิ ปุพฺเพ หตฺถิมฺปิ ธาวนฺตํ อนุปติตฺวา คณฺหามิ, อสฺสมฺปิ ธาวนฺตํ อนุปติตฺวา คณฺหามิ, รถมฺปิ ธาวนฺตํ อนุปติตฺวา คณฺหามิ, มิคมฺปิ ธาวนฺตํ อนุปติตฺวา คณฺหามิ, อถ จ ปนาหํ อิมํ สมณํ ปกติยา คจฺฉนฺตํ สพฺพถาเมน คจฺฉนฺโต น สกฺโกมิ สมฺปาปุณิตุนฺ”ติ.\csromanspacer{} ฐิโตว ภควนฺตํ เอตทโวจ “ติฏฺฐ ติฏฺฐ สมณา”ติ.\csromanspacer{} ฐิโต อหํ องฺคุลิมาล, ตฺวํ จ ติฏฺฐาติ.\csromanspacer{} อถ โข โจรสฺส องฺคุลิมาลสฺส เอตทโหสิ “อิเม โข สมณา สกฺยปุตฺติยา สจฺจวาทิโน สจฺจปฏิญฺญา, อถ ปนายํ สมโณ คจฺฉํเยวาห ‘ฐิโต อหํ องฺคุลิมาล, ตฺวํ จ ติฏฺฐา’ติ, ยํ นูนาหํ อิมํ สมณํ ปุจฺเฉยฺยนฺ”ติ.}

\csromanheader{2}{349. อถ โข โจโร องฺคุลิมาโล ภควนฺตํ คาถาย อชฺฌภาสิ\csromandash{}}
\csromangathameasure{“คจฺฉํ วเทสิ สมณ ฐิโตมฺหิ, \\ มมญฺจ พฺรูสิ ฐิตมฏฺฐิโตติ. \\ ปุจฺฉามิ ตํ สมณ เอตมตฺถํ, \\ กถํ ฐิโต ตฺวํ อหมฏฺฐิโตมฺหี”ติ.}
\csromangathagroup{\csromangathastanza{“คจฺฉํ วเทสิ สมณ ฐิโตมฺหิ, \\ มมญฺจ พฺรูสิ ฐิตมฏฺฐิโตติ. \\ ปุจฺฉามิ ตํ สมณ เอตมตฺถํ, \\ กถํ ฐิโต ตฺวํ อหมฏฺฐิโตมฺหี”ติ.}}

\csromanheader{2}{6. องฺคุลิมาลสุตฺต (86)}
\csromangathameasure{ฐิโต อหํ องฺคุลิมาล สพฺพทา, \\ สพฺเพสุ ภูเตสุ นิธาย ทณฺฑํ. \\ ตุวํ จ ปาเณสุ อสญฺญโตสิ, \\ ตสฺมา ฐิโตหํ ตุวมฏฺฐิโตสีติ. \\ จิรสฺสํ วต เม มหิโต มเหสี, \\ มหาวนํ ปาปุณิ สจฺจวาที. \\ โสหํ จริสฺสามิ ปหาย ปาปํ2, \\ สุตฺวาน คาถํ ตว ธมฺมปุตฺตํ. \\ อิเตฺวว โจโร อสิมาวุธญฺจ, \\ โสพฺเภ ปปาเต นรเก อกิริ. \\ อวนฺทิ โจโร สุคตสฺส ปาเท, \\ ตตฺเถว นํ ปพฺพชฺชํ อยาจิ. \\ พุทฺโธ จ โข การุณิโก มเหสิ, \\ โย สตฺถา โลกสฺส สเทวกสฺส. \\ ตเมหิ ภิกฺขูติ ตทา อโวจ, \\ เอเสว ตสฺส อหุ ภิกฺขุภาโวติ.}
\csromangathagroup{\csromangathastanza{\csromanfolio{303}ฐิโต อหํ องฺคุลิมาล สพฺพทา, \\ สพฺเพสุ ภูเตสุ นิธาย ทณฺฑํ. \\ ตุวํ จ ปาเณสุ อสญฺญโตสิ, \\ ตสฺมา ฐิโตหํ ตุวมฏฺฐิโตสีติ.}\csromangathastanza{จิรสฺสํ วต เม มหิโต มเหสี, \\ มหาวนํ ปาปุณิ สจฺจวาที\footnote{มหาวนํ สมโณยํ ปจฺจุปาทิ (สี), มหาวนํ สมณ ปจฺจุปาทิ (สฺยา, กํ) 2. โส หํ จิรสฺสาปิ ปหาสฺสํ ปาปํ (สี), โสหํ จริสฺสามิ ปชาหิสฺสํ ปาปํ(สฺยา,กํ)}. \\ โสหํ จริสฺสามิ ปหาย ปาปํ2, \\ สุตฺวาน คาถํ ตว ธมฺมปุตฺตํ.}\csromangathastanza{อิเตฺวว โจโร อสิมาวุธญฺจ, \\ โสพฺเภ ปปาเต นรเก อกิริ. \\ อวนฺทิ โจโร สุคตสฺส ปาเท, \\ ตตฺเถว นํ ปพฺพชฺชํ อยาจิ.}\csromangathastanza{พุทฺโธ จ โข การุณิโก มเหสิ, \\ โย สตฺถา โลกสฺส สเทวกสฺส. \\ ตเมหิ ภิกฺขูติ ตทา อโวจ, \\ เอเสว ตสฺส อหุ ภิกฺขุภาโวติ.}}

\proseitem{350}{อถ โข ภควา อายสฺมตา องฺคุลิมาเลน ปจฺฉาสมเณน เยน สาวตฺถิ เตน จาริกํ ปกฺกามิ, อนุปุพฺเพน จาริกํ จรมาโน เยน สาวตฺถิ ตทวสริ, ตตฺร สุทํ ภควา สาวตฺถิยํ วิหรติ เชตวเน อนาถปิณฺฑิกสฺส อาราเม.\csromanspacer{} เตน โข ปน สมเยน รญฺโญ ปเสนทิสฺส โกสลสฺส อนฺเตปุรทฺวาเร มหาชนกาโย สนฺนิปติตฺวา อุจฺจาสทฺโท มหาสทฺโท โหติ “โจโร เต เทว วิชิเต องฺคุลิมาโล นาม ลุทฺโท โลหิตปาณิ หตปหเต นิวิฏฺโฐ อทยาปนฺโน ปาณภูเตสุ, เตน คามาปิ อคามา กตา, นิคมาปิ อนิคมา กตา, ชนปทาปิ อชนปทา กตา, โส มนุสฺเส วธิตฺวา วธิตฺวา องฺคุลีนํ มาลํ ธาเรติ, ตํ เทโว ปฏิเสเธตู”ติ.}

\prose{\csromanfolio{304}อถ โข ราชา ปเสนทิ โกสโล ปญฺจมตฺเตหิ อสฺสสเตหิ สาวตฺถิยา นิกฺขมิ ทิวา ทิวสฺส เยน อาราโม เตน ปาวิสิ, ยาวติกา ยานสฺส ภูมิ, ยาเนน คนฺตฺวา ยานา ปจฺโจโรหิตฺวา ปตฺติโกว เยน ภควา เตนุปสงฺกมิ, อุปสงฺกมิตฺวา ภควนฺตํ อภิวาเทตฺวา เอกมนฺตํ นิสีทิ, เอกมนฺตํ นิสินฺนํ โข ราชานํ ปเสนทิํ โกสลํ ภควา เอตทโวจ “กิํ นุ เต มหาราช ราชา วา มาคโธ เสนิโย พิมฺพิสาโร กุปิโต, เวสาลิกา วา ลิจฺฉวี อญฺเญ วา ปฏิราชาโน”ติ.\csromanspacer{} น โข เม ภนฺเต ราชา มาคโธ เสนิโย พิมฺพิสาโร กุปิโต, นาปิ เวสาลิกา ลิจฺฉวี, นาปิ อญฺเญ ปฏิราชาโน, โจโร เม ภนฺเต วิชิเต องฺคุลิมาโล นาม ลุทฺโท โลหิตปาณิ หตปหเต นิวิฏฺโฐ อทยาปนฺโน ปาณภูเตสุ, เตน คามาปิ อคามา กตา, นิคมาปิ อนิคมา กตา, ชนปทาปิ อชนปทา กตา, โส มนุสฺเส วธิตฺวา วธิตฺวา องฺคุลีนํ มาลํ ธาเรติ, ตาหํ ภนฺเต ปฏิเสธิสฺสามีติ.\csromanspacer{} สเจ ปน ตฺวํ มหาราช องฺคุลิมาลํ ปสฺเสยฺยาสิ เกสมสฺสุํ โอหาเรตฺวา กาสายานิ วตฺถานิ อจฺฉาเทตฺวา อคารสฺมา อนคาริยํ ปพฺพชิตํ วิรตํ ปาณาติปาตา วิรตํ อทินฺนาทานา วิรตํ มุสาวาทา เอกภตฺติกํ พฺรหฺมจาริํ สีลวนฺตํ กลฺยาณธมฺมํ, กินฺติ นํ กเรยฺยาสีติ.\csromanspacer{} อภิวาเทยฺยาม วา ภนฺเต ปจฺจุฏฺเฐยฺยาม วา อาสเนน วา นิมนฺเตยฺยาม, อภินิมนฺเตยฺยาม วา นํ จีวรปิณฺฑปาตเสนาสนคิลานปฺปจฺจยเภสชฺชปริกฺขาเรหิ, ธมฺมิกํ วา อสฺส รกฺขาวรณคุตฺติํ สํวิทเหยฺยาม, กุโต ปนสฺส ภนฺเต ทุสฺสีลสฺส ปาปธมฺมสฺส เอวรูโป สีลสํยโม ภวิสฺสตีติ.\csromanspacer{} เตน โข ปน สมเยน อายสฺมา องฺคุลิมาโล ภควโต อวิทูเร นิสินฺโน โหติ.\csromanspacer{} อถ โข ภควา ทกฺขิณํ พาหุํ ปคฺคเหตฺวา ราชานํ ปเสนทิํ โกสลํ เอตทโวจ “เอโส มหาราช องฺคุลิมาโล”ติ.\csromanspacer{} อถ โข รญฺโญ ปเสนทิสฺส โกสลสฺส อหุเทว ภยํ, อหุ ฉมฺภิตตฺตํ, อหุ โลมหํโส.\csromanspacer{} อถ โข ภควา ราชานํ ปเสนทิํ โกสลํ ภีตํ สํวิคฺคํ โลมหฏฺฐชาตํ วิทิตฺวา ราชานํ ปเสนทิํ โกสลํ เอตทโวจ “มา ภายิ มหาราช, นตฺถิ เต อิโต ภยนฺ”ติ.\csromanspacer{} อถ โข รญฺโญ ปเสนทิสฺส โกสลสฺส ยํ อโหสิ ภยํ วา ฉมฺภิตตฺตํ วา โลมหํโส วา, โส ปฏิปฺปสฺสมฺภิ.\csromanspacer{} อถ โข ราชา ปเสนทิ โกสโล เยนายสฺมา}

\vspace{\tipitakaparskip}
\csromancenter{องฺคุลิมาลสุตฺต (86) องฺคุลิมาโล เตนุปสงฺกมิ, อุปสงฺกมิตฺวา อายสฺมนฺตํ องฺคุลิมาลํ เอตทโวจ “อยฺโย โน ภนฺเต องฺคุลิมาโล”ติ. “เอวํ มหาราชา”ติ.\csromanspacer{} กถํโคตฺโต อยฺยสฺส ปิตา, กถํโคตฺตา มาตาติ.\csromanspacer{} คคฺโค โข มหาราช ปิตา, มนฺตาณี มาตาติ.\csromanspacer{} อภิรมตุ ภนฺเต อยฺโย คคฺโค มนฺตาณิปุตฺโต, อหมยฺยสฺส คคฺคสฺส มนฺตาณิปุตฺตสฺส อุสฺสุกฺกํ กริสฺสามิ จีวรปิณฺฑปาตเสนาสนคิลานปฺปจฺจยเภสชฺชปริกฺขารานนฺติ.}
\proseitem{351}{\csromanfolio{305}เตน โข ปน สมเยน อายสฺมา องฺคุลิมาโล อารญฺญิโก โหติ ปิณฺฑปาติโก ปํสุกูลิโก เตจีวริโก.\csromanspacer{} อถ โข อายสฺมา องฺคุลิมาโล ราชานํ ปเสนทิํ โกสลํ เอตทโวจ “อลํ มหาราช, ปริปุณฺณํ เม จีวรนฺ”ติ.\csromanspacer{} อถ โข ราชา ปเสนทิ โกสโล เยน ภควา เตนุปสงฺกมิ, อุปสงฺกมิตฺวา ภควนฺตํ อภิวาเทตฺวา เอกมนฺตํ นิสีทิ, เอกมนฺตํ นิสินฺโน โข ราชา ปเสนทิ โกสโล ภควนฺตํ เอตทโวจ “อจฺฉริยํ ภนฺเต, อพฺภุตํ ภนฺเต, ยาวญฺจิทํ ภนฺเต ภควา อทนฺตานํ ทเมตา, อสนฺตานํ สเมตา, อปรินิพฺพุตานํ ปรินิพฺพาเปตา, ยํ หิ มยํ ภนฺเต นาสกฺขิมฺหา ทณฺเฑนปิ สตฺเถนปิ ทเมตุํ, โส ภควตา อทณฺเฑน อสตฺเถเนว\footnote{อสตฺเถน (สฺยา, กํ)} ทนฺโต.\csromanspacer{} หนฺท จ ทานิ\footnote{หนฺท ทานิ (สฺยา, กํ, อิ)} มยํ ภนฺเต คจฺฉาม, พหุกิจฺจา มยํ พหุกรณียา”ติ.\csromanspacer{} ยสฺสทานิ มหาราช กาลํ มญฺญสีติ.\csromanspacer{} อถ โข ราชา ปเสนาทิ โกสโล อุฏฺฐายาสนา ภควนฺตํ อภิวาเทตฺวา ปทกฺขิณํ กตฺวา ปกฺกามิ.}

\prose{อถ โข อายสฺมา องฺคุลิมาโล ปุพฺพณฺหสมยํ นิวาเสตฺวา ปตฺตจีวรมาทาย สาวตฺถิยํ ปิณฺฑาย ปาวิสิ.\csromanspacer{} อทฺทสา โข อายสฺมา องฺคุลิมาโล สาวตฺถิยํ สปทานํ ปิณฺฑาย จรมาโน อญฺญตรํ อิตฺถิํ มูฬฺหคพฺภํ วิฆาตคพฺภํ\footnote{วิสาหคพฺพํ (สฺยา, กํ, อิ, ก)}, ทิสฺวานสฺส เอตทโหสิ “กิลิสฺสนฺติ วต โภ สตฺตา, กิลิสฺสนฺติ วต โภ สตฺตา”ติ.\csromanspacer{} อถ โข อายสฺมา องฺคุลิมาโล สาวตฺถิยํ ปิณฺฑาย จริตฺวา ปจฺฉาภตฺตํ ปิณฺฑปาตปฏิกฺกนฺโต เยน ภควา เตนุปสงฺกมิ, อุปสงฺกมิตฺวา ภควนฺตํ อภิวาเทตฺวา เอกมนฺตํ นิสีทิ, เอกมนฺตํ นิสินฺโน โข อายสฺมา องฺคุลิมาโล ภควนฺตํ เอตทโวจ “อิธาหํ ภนฺเต ปุพฺพณฺหสมยํ นิวาเสตฺวา ปตฺตจีวรมาทาย \csromanfolio{306}สาวตฺถิํ ปิณฺฑาย ปาวิสิํ, อทฺทสํ โข อหํ ภนฺเต สาวตฺถิยํ สปทานํ ปิณฺฑาย จรมาโน อญฺญตรํ อิตฺถิํ มูฬฺหคพฺภํ วิฆาตคพฺภํ, ทิสฺวาน มยฺหํ เอตทโหสิ “กิลิสฺสนฺติ วต โภ สตฺตา, กิลิสฺสนฺติ วต โภ สตฺตา”ติ.}

\prose{เตน หิ ตฺวํ องฺคุลิมาล เยน สา อิตฺถี เตนุปสงฺกม, อุปสงฺกมิตฺวา ตํ อิตฺถิํ เอวํ วเทหิ “ยโตหํ ภคินิ ชาโต\footnote{ภคินิ ชาติยา ชาโต (สี)} นาภิชานามิ สญฺจิจฺจ ปาณํ ชีวิตา โวโรเปตา, เตน สจฺเจน โสตฺถิ เต โหตุ โสตฺถิ คพฺภสฺสา”ติ.}

\prose{โส หิ นูน เม ภนฺเต สมฺปชานมุสาวาโท ภวิสฺสติ, มยา หิ ภนฺเต พหู สญฺจิจฺจ ปาณา ชีวิตา โวโรปิตาติ.\csromanspacer{} เตน หิ ตฺวํ องฺคุลิมาล เยน สา อิตฺถี เตนุปสงฺกม, อุปสงฺกมิตฺวา ตํ อิตฺถิํ เอวํ วเทหิ “ยโตหํ ภคินิ อริยาย ชาติยา ชาโต นาภิชานามิ สญฺจิจฺจ ปาณํ ชีวิตา โวโรเปตา, เตน สจฺเจน โสตฺถิ เต โหตุ โสตฺถิ คพฺภสฺสา”ติ.}

\csromanmatikaanchor{mtk.43}
\csromantocmarkref{subsection}{6. องฺคุลิมาลสุตฺต}{mtk.43}
\bookmark[dest={mtk.43},level=2]{6. องฺคุลิมาลสุตฺต}
\prose{“เอวํ ภนฺเต”ติ โข อายสฺมา องฺคุลิมาโล ภควโต ปฏิสฺสุตฺวา เยน สา อิตฺถี เตนุปสงฺกมิ, อุปสงฺกมิตฺวา ตํ อิตฺถิํ เอตทโวจ “ยโตหํ ภคินิ อริยาย ชาติยา ชาโต นาภิชานามิ สญฺจิจฺจ ปาณํ ชีวิตา โวโรเปตา, เตน สจฺเจน โสตฺถิ เต โหตุ โสตฺถิ คพฺภสฺสา”ติ.\csromanspacer{} อถ ขฺวาสฺสา อิตฺถิยา โสตฺถิ อโหสิ โสตฺถิ คพฺภสฺส.}

\prose{อถ โข อายสฺมา องฺคุลิมาโล เอโก วูปกฏฺโฐ อปฺปมตฺโต อาตาปี ปหิตตฺโต วิหรนฺโต นจิรสฺเสว, ยสฺสตฺถาย กุลปุตฺตา สมฺมเทว อคารสฺมา อนคาริยํ ปพฺพชนฺติ, ตทนุตฺตรํ พฺรหฺมจริยปริโยสานํ ทิฏฺเฐว ธมฺเม สยํ อภิญฺญา สจฺฉิกตฺวา อุปสมฺปชฺช วิหาสิ, “ขีณา ชาติ, วุสิตํ พฺรหฺมจริยํ, กตํ กรณียํ, นาปรํ อิตฺถตฺตายา”ติ อพฺภญฺญาสิ.\csromanspacer{} อญฺญตโร โข ปนายสฺมา องฺคุลิมาโล อรหตํ อโหสิ.}

\proseitem{352}{อถ โข อายสฺมา องฺคุลิมาโล ปุพฺพณฺหสมยํ นิวาเสตฺวา ปตฺตจีวรมาทาย สาวตฺถิํ ปิณฺฑาย ปาวิสิ.\csromanspacer{} เตน โข ปน สมเยน อญฺเญนปิ เลฑฺฑุ ขิตฺโต อายสฺมโต องฺคุลิมาลสฺส กาเย นิปตติ,}

\vspace{\tipitakaparskip}
\csromancenter{องฺคุลิมาลสุตฺต (86) อญฺเญนปิ ทณฺโฑ ขิตฺโต อายสฺมโต องฺคุลิมาลสฺส กาเย นิปตติ, อญฺเญนปิ สกฺขรา ขิตฺตา อายสฺมโต องฺคุลิมาลสฺส กาเย นิปตติ.\csromanspacer{} อถ โข อายสฺมา องฺคุลิมาโล ภินฺเนน สีเสน โลหิเตน คฬนฺเตน ภินฺเนน ปตฺเตน วิปฺผาลิตาย สงฺฆาฏิยา เยน ภควา เตนุปสงฺกมิ.\csromanspacer{} อทฺทสา โข ภควา อายสฺมนฺตํ องฺคุลิมาลํ ทูรโตว อาคจฺฉนฺตํ, ทิสฺวาน อายสฺมนฺตํ องฺคุลิมาลํ เอตทโวจ “อธิวาเสหิ ตฺวํ พฺราหฺมณ, อธิวาเสหิ ตฺวํ พฺราหฺมณ, ยสฺส โข ตฺวํ พฺราหฺมณ กมฺมสฺส วิปาเกน พหูนิ วสฺสานิ พหูนิ วสฺสสตานิ พหูนิ วสฺสสหสฺสานิ นิรเย ปจฺเจยฺยาสิ, ตสฺส ตฺวํ พฺราหฺมณ กมฺมสฺส วิปากํ ทิฏฺเฐว ธมฺเม ปฏิสํเวเทสี”ติ.\csromanspacer{} อถ โข อายสฺมา องฺคุลิมาโล รโหคโต ปฏิสลฺลีโน วิมุตฺติสุขํ ปฏิสํเวทิ.\csromanspacer{} ตายํ เวลายํ อิมํ อุทานํ อุทาเนสิ\csromandash{}}
\vspace{0.5\baselineskip}
\csromangathameasure{“โย ปุพฺเพว ปมชฺชิตฺวา, ปจฺฉา โส นปฺปมชฺชติ. \\ โสมํ โลกํ ปภาเสติ, อพฺภา มุตฺโตว จนฺทิมา. \\ ยสฺส ปาปํ กตํ กมฺมํ, กุสเลน ปิธียติ. \\ โสมํ โลกํ ปภาเสติ, อพฺภา มุตฺโตว จนฺทิมา. \\ โย หเว ทหโร ภิกฺขุ, ยุญฺชติ พุทฺธสาสเน. \\ โสมํ โลกํ ปภาเสติ, อพฺโภ มุตฺโตว จนฺทิมา. \\ ทิสา หิ เม ธมฺมกถํ สุณนฺตุ, \\ ทิสา หิ เม ยุญฺชนฺตุ พุทฺธสาสเน. \\ ทิสา หิ เม เต มนุชา ภชนฺตุ, \\ เย ธมฺมเมวาทปยนฺติ สนฺโต. \\ ทิสา หิ เม ขนฺติวาทานํ, อวิโรธปฺปสํสีนํ. \\ สุณนฺตุ ธมฺมํ กาเลน, ตญฺจ อนุวิธียนฺตุ. \\ น หิ ชาตุ โส มมํ หิํเส, อญฺญํ วา ปน กิญฺจิ นํ. \\ ปปฺปุยฺย ปรมํ สนฺติํ, รกฺเขยฺย ตสถาวเร. \\ อุทกญฺหิ นยนฺติ เนตฺติกา, อุสุการา นมยนฺติ เตชนํ. \\ ทารุํ นมยนฺติ1 ตจฺฉกา, อตฺตานํ ทมยนฺติ ปณฺฑิตา. \\ ทณฺเฑเนเก ทมยนฺติ, องฺกุเสหิ กสาหิ จ. \\ อทณฺเฑน อสตฺเถน, อหํ ทนฺโตมฺหิ ตาทินา. \\ อหิํสโกติ เม นามํ, หิํสกสฺส ปุเร สโต. \\ อชฺชาหํ สจฺจนาโมมฺหิ, น นํ หิํสามิ กิญฺจิ นํ. \\ โจโร อหํ ปุเร อาสิํ, องฺคุลิมาโลหิ วิสฺสุโต. \\ วุยฺหมาโน มโหเฆน, พุทฺธํ สรณมาคมํ. \\ โลหิตปาณิ ปุเร อาสิํ, องฺคุลิมาโลติ วิสฺสุโต. \\ สรณคมนํ ปสฺส, ภวเนตฺติ สมูหตา. \\ ตาทิสํ กมฺมํ กตฺวาน, พหุํ ทุคฺคติคามินํ. \\ ผุฏฺโฐ กมฺมวิปาเกน, อณโณ ภุญฺชามิ โภชนํ. \\ ปมาทมนุยุญฺชนฺติ, พาลา ทุมฺเมธิโน ชนา. \\ อปฺปมาทญฺจ เมธาวี, ธนํ เสฏฺฐํว รกฺขติ. \\ มา ปมาทมนุยุญฺเชถ, มา กามรติ สนฺถวํ. \\ อปฺปมตฺโต หิ ฌายนฺโต, ปปฺโปติ วิปุลํ สุขํ. \\ สฺวาคตํ นาปคตํ, นยิทํ ทุมฺมนฺติตํ มม. \\ สํวิภตฺเตสุ ธมฺเมสุ, ยํ เสฏฺฐํ ตทุปาคมํ. \\ สฺวาคตํ นาปคตํ, นยิทํ ทุมฺมนฺติตํ มม. \\ ติสฺโส วิชฺชา อนุปฺปตฺตา, กตํ พุทฺธสฺส สาสนนฺ”ติ.}
\csromangathagroup{\csromangathastanza{\csromanfolio{307}“โย ปุพฺเพว\footnote{โย จ ปุพฺเพ (สี, สฺยา, กํ, อิ)} ปมชฺชิตฺวา, ปจฺฉา โส นปฺปมชฺชติ. \\ โสมํ\footnote{โย จ ปุพฺเพ (สี, สฺยา, กํ, อิ)} โลกํ ปภาเสติ, อพฺภา มุตฺโตว จนฺทิมา.}\csromangathastanza{ยสฺส ปาปํ กตํ กมฺมํ, กุสเลน ปิธียติ\footnote{ปิถียติ (สี, สฺยา, กํ, อิ)}. \\ โสมํ โลกํ ปภาเสติ, อพฺภา มุตฺโตว จนฺทิมา.}\csromangathastanza{โย หเว ทหโร ภิกฺขุ, ยุญฺชติ พุทฺธสาสเน. \\ โสมํ โลกํ ปภาเสติ, อพฺโภ มุตฺโตว จนฺทิมา.}\csromangathastanza{ทิสา หิ เม ธมฺมกถํ สุณนฺตุ, \\ ทิสา หิ เม ยุญฺชนฺตุ พุทฺธสาสเน. \\ ทิสา หิ เม เต มนุชา ภชนฺตุ, \\ เย ธมฺมเมวาทปยนฺติ สนฺโต. \\ ทิสา หิ เม ขนฺติวาทานํ, อวิโรธปฺปสํสีนํ. \\ สุณนฺตุ ธมฺมํ กาเลน, ตญฺจ อนุวิธียนฺตุ.}\csromangathastanza{น หิ ชาตุ โส มมํ หิํเส, อญฺญํ วา ปน กิญฺจิ นํ\footnote{กญฺจิ นํ (สี, สฺยา, กํ, อิ), กญฺจนํ (?)}. \\ ปปฺปุยฺย ปรมํ สนฺติํ, รกฺเขยฺย ตสถาวเร.}\csromangathastanza{\csromanfolio{308}อุทกญฺหิ นยนฺติ เนตฺติกา, อุสุการา นมยนฺติ\footnote{ทมยนฺติ (ก)} เตชนํ. \\ ทารุํ นมยนฺติ1 ตจฺฉกา, อตฺตานํ ทมยนฺติ ปณฺฑิตา.}\csromangathastanza{ทณฺเฑเนเก ทมยนฺติ, องฺกุเสหิ กสาหิ จ. \\ อทณฺเฑน อสตฺเถน, อหํ ทนฺโตมฺหิ ตาทินา.}\csromangathastanza{อหิํสโกติ เม นามํ, หิํสกสฺส ปุเร สโต. \\ อชฺชาหํ สจฺจนาโมมฺหิ, น นํ หิํสามิ กิญฺจิ นํ\footnote{กญฺจิ นํ (สี, สฺยา, กํ, อิ), กญฺจนํ (?)}.}\csromangathastanza{โจโร อหํ ปุเร อาสิํ, องฺคุลิมาโลหิ วิสฺสุโต. \\ วุยฺหมาโน มโหเฆน, พุทฺธํ สรณมาคมํ.}\csromangathastanza{โลหิตปาณิ ปุเร อาสิํ, องฺคุลิมาโลติ วิสฺสุโต. \\ สรณคมนํ ปสฺส, ภวเนตฺติ สมูหตา.}\csromangathastanza{ตาทิสํ กมฺมํ กตฺวาน, พหุํ ทุคฺคติคามินํ. \\ ผุฏฺโฐ กมฺมวิปาเกน, อณโณ ภุญฺชามิ โภชนํ.}\csromangathastanza{ปมาทมนุยุญฺชนฺติ, พาลา ทุมฺเมธิโน ชนา. \\ อปฺปมาทญฺจ เมธาวี, ธนํ เสฏฺฐํว รกฺขติ.}\csromangathastanza{มา ปมาทมนุยุญฺเชถ, มา กามรติ สนฺถวํ. \\ อปฺปมตฺโต หิ ฌายนฺโต, ปปฺโปติ วิปุลํ\footnote{ปรมํ (ก)} สุขํ.}\csromangathastanza{สฺวาคตํ\footnote{สาคตํ (สี, อิ)} นาปคตํ\footnote{นาม สคตํ (ก)}, นยิทํ ทุมฺมนฺติตํ มม. \\ สํวิภตฺเตสุ\footnote{สาคตํ (สี, อิ)} ธมฺเมสุ, ยํ เสฏฺฐํ ตทุปาคมํ.}\csromangathastanza{สฺวาคตํ นาปคตํ, นยิทํ ทุมฺมนฺติตํ มม. \\ ติสฺโส วิชฺชา อนุปฺปตฺตา, กตํ พุทฺธสฺส สาสนนฺ”ติ.}}

\nitthitam{องฺคุลิมาลสุตฺตํ นิฏฺฐิตํ ฉฏฺฐํ.}
\csromanheader{2}{7. ปิยชาติกสุตฺต (87) \textbf{7. ปิยชาติกสุตฺต}}
\proseitem{353}{\csromanfolio{309}เอวํ เม สุตํ\csromandash{}เอกํ สมยํ ภควา สาวตฺถิยํ วิหรติ เชตวเน อนาถปิณฺฑิกสฺส อาราเม.\csromanspacer{} เตน โข ปน สมเยน อญฺญตรสฺส คหปติสฺส เอกปุตฺตโก ปิโย มนาโป กาลงฺกโต โหติ.\csromanspacer{} ตสฺส กาลํกิริยาย เนว กมฺมนฺตา ปฏิภนฺติ, น ภตฺตํ ปฏิภาติ.\csromanspacer{} โส อาฬาหนํ คนฺตฺวา กนฺทติ “กหํ เอกปุตฺตก กหํ เอกปุตฺตกา”ติ.\csromanspacer{} อถ โข โส คหปติ เยน ภควา เตนุปสงฺกมิ, อุปสงฺกมิตฺวา ภควนฺตํ อภิวาเทตฺวา เอกมนฺตํ นิสีทิ, เอกมนฺตํ นิสินฺนํ โข ตํ คหปติํ ภควา เอตทโวจ “น โข เต คหปติ สเก จิตฺเต ฐิตสฺส อินฺทฺริยานิ, อตฺถิ เต อินฺทฺริยานํ อญฺญถตฺตนฺ”ติ.\csromanspacer{} กิํ หิ เม ภนฺเต อินฺทฺริยานํ นาญฺญถตฺตํ ภวิสฺสติ, มยฺหญฺหิ ภนฺเต เอกปุตฺโต ปิโย มนาโป กาลงฺกโต, ตสฺส กาลํกิริยาย เนว กมฺมนฺตา ปฏิภนฺติ, น ภตฺตํ ปฏิภาติ.\csromanspacer{} โสหํ อาฬาหนํ คนฺตฺวา กนฺทามิ “กหํ เอกปุตฺตก กหํ เอกปุตฺตกา”ติ.\csromanspacer{} เอวเมตํ คหปติ, เอวเมตํ คหปติ\footnote{เอวเมตํ คหปติ (อิ, สกิเทว), เอวเมว (สี, สกิเทว)}, ปิยชาติกา หิ คหปติ โสกปริเทวทุกฺขโทมนสฺสุปายาสา ปิยปฺปภวิกาติ.\csromanspacer{} กสฺส โข\footnote{กิสฺส นุ โข (สี)} นาเมตํ ภนฺเต เอวํ ภวิสฺสติ ปิยชาติกา โสกปริเทวทุกฺขโทมนสฺสุปายาสา ปิยปฺปภวิกา, ปิยชาติกา หิ โข ภนฺเต อานนฺทโสมนสฺสา ปิยปฺปภวิกาติ.\csromanspacer{} อถ โข โส คหปติ ภควโต ภาสิตํ อนภินนฺทิตฺวา ปฏิกฺโกสิตฺวา อุฏฺฐายาสนา ปกฺกามิ.}

\proseitem{354}{เตน โข ปน สมเยน สมฺพหุลา อกฺขธุตฺตา ภควโต อวิทูเร อกฺเขหิ ทิพฺพนฺติ.\csromanspacer{} อถ โข โส คหปติ เยน เต อกฺขธุตฺตา เตนุปสงฺกมิ, อุปสงฺกมิตฺวา อกฺขธุตฺเต เอตทโวจ “อิธาหํ โภนฺโต เยน สมโณ โคตโม เตนุปสงฺกมิํ, อุปสงฺกมิตฺวา สมณํ โคตมํ อภิวาเทตฺวา เอกมนฺตํ นิสีทิํ, เอกมนฺตํ นิสินฺนํ โข มํ โภนฺโต สมโณ โคตโม เอตทโวจ ‘น โข เต คหปติ สเก จิตฺเต ฐิตสฺส อินฺทฺริยานิ, อตฺถิ เต อินฺทฺริยานํ อญฺญถตฺตนฺ’ติ.\csromanspacer{} เอวํ วุตฺเต อหํ โภนฺโต สมณํ โคตมํ เอตทโวจํ ‘กิํ หิ เม ภนฺเต อินฺทฺริยานํ นาญฺญถตฺตํ ภวิสฺสติ, มยฺหญฺหิ ภนฺเต เอกปุตฺตโก ปิโย มนาโป กาลงฺกโต, ตสฺส กาลํกิริยาย เนว กมฺมนฺตา ปฏิภนฺติ, น ภตฺตํ \csromanfolio{310}ปฏิภาติ, โสหํ อาฬาหนํ คนฺตฺวา กนฺทามิ กหํ เอกปุตฺตก กหํ เอกปุตฺตกา’ติ.\csromanspacer{} เอวเมตํ คหปติ, เอวเมตํ คหปติ, ปิยชาติกา หิ คหปติ โสกปริเทวทุกฺขโทมนสฺสุปายาสา ปิยปฺปภวิกาติ.\csromanspacer{} กสฺส โข นาเมตํ ภนฺเต เอวํ ภวิสฺสติ ปิยชาติกา โสกปริเทวทุกฺขโทมนสฺสุปายาสา ปิยปฺปภวิกา, ปิยชาติกา หิ โข ภนฺเต อานนฺทโสมนสฺสา ปิยปฺปภวิกาติ.\csromanspacer{} อถ ขฺวาหํ โภนฺโต สมณสฺส โคตมสฺส ภาสิตํ อนภินนฺทิตฺวา ปฏิกฺโกสิตฺวา อุฏฺฐายาสนา ปกฺกมินฺ”ติ.\csromanspacer{} เอวเมตํ คหปติ, เอวเมตํ คหปติ, ปิยชาติกา หิ คหปติ อานนฺทโสมนสฺสา ปิยปฺปภวิกาติ.\csromanspacer{} อถ โข โส คหปติ “สเมติ เม อกฺขธุตฺเตหี”ติ ปกฺกามิ.\csromanspacer{} อถ โข อิทํ กถาวตฺถุ อนุปุพฺเพน ราชนฺเตปุรํ ปาวิสิ.}

\proseitem{355}{อถ โข ราชา ปเสนทิ โกสโล มลฺลิกํ เทวิํ อามนฺเตสิ “อิทํ เต มลฺลิเก สมเณน โคตเมน ภาสิตํ ปิยชาติกา โสกปริเทวทุกฺขโทมนสฺสุปายาสา ปิยปฺปภวิกา”ติ.\csromanspacer{} สเจตํ มหาราช ภควตา ภาสิตํ เอวเมตนฺติ.\csromanspacer{} เอวเมว ปนายํ มลฺลิกา ยญฺญเทว สมโณ โคตโม ภาสติ, ตํ ตเทวสฺส อพฺภนุโมทติ.\csromanspacer{} สเจตํ มหาราช ภควตา ภาสิตํ เอวเมตนฺติ.\csromanspacer{} เสยฺยถาปิ นาม ยญฺญเทว อาจริโย อนฺเตวาสิสฺส ภาสติ, ตํ ตเทวสฺส อนฺเตวาสี อพฺภนุโมทติ “เอวเมตํ อาจริย เอวเมตํ อาจริยา”ติ.\csromanspacer{} เอวเมว โข ตฺวํ มลฺลิเก ยญฺญเทว สมโณ โคตโม ภาสติ, ตํ ตเทวสฺส อพฺภนุโมทสิ “สเจตํ มหาราช ภควตา ภาสิตํ เอวเมตนฺ”ติ, จร ปิเร มลฺลิเก วินสฺสาติ.\csromanspacer{} อถ โข มลฺลิกา เทวี นาฬิชงฺฆํ พฺราหฺมณํ อามนฺเตสิ “เอหิ ตฺวํ พฺราหฺมณ เยน ภควา เตนุปสงฺกม, อุปสงฺกมิตฺวา มม วจเนน ภควโต ปาเท สิรสา วนฺทาหิ, อปฺปาพาธํ อปฺปาตงฺกํ ลหุฐานํ พลํ ผาสุวิหารํ ปุจฺฉ ‘มลฺลิกา ภนฺเต เทวี ภควโต ปาเท สิรสา วนฺทติ, อปฺปาพาธํ อปฺปาตงฺกํ ลหุฏฺฐานํ พลํ ผาสุวิหารํ ปุจฺฉตี’ติ, เอวญฺจ วเทหิ ‘ภาสิตา นุ โข ภนฺเต ภควตา เอสา วาจา, ปิยชาติกา โสกปริเทวทุกฺขโทมนสฺสุปายาสา ปิยปฺปภวิกา’ติ, ยถา เต ภควา พฺยากโรติ, ตํ สาธุกํ อุคฺคเหตฺวา มม อาโรเจยฺยาสิ, น หิ ตถาคตา วิตถํ ภณนฺตี”ติ. “เอวํ โภตี”ติ โข นาฬิชงฺโฆ พฺราหฺมโณ มลฺลิกาย เทวิยา ปฏิสฺสุตฺวา เยน}

\vspace{\tipitakaparskip}
\csromancenter{ปิยชาติกสุตฺต (87) ภควา เตนุปสงฺกมิ, อุปสงฺกมิตฺวา ภควตา สทฺธิํ สมฺโมทิ, สมฺโมทนียํ กถํ สารณียํ วีติสาเรตฺวา เอกมนฺตํ นิสีทิ, เอกมนฺตํ นิสินฺโน โข นาฬิชงฺโฆ พฺราหฺมโณ ภควนฺตํ เอตทโวจ “มลฺลิกา โภ โคตม เทวี โภโต โคตมสฺส ปาเท สิรสา วนฺทติ, อปฺปาพาธํ อปฺปาตงฺกํ ลหุฏฺฐานํ พลํ ผาสุวิหารํ ปุจฺฉติ, เอวญฺจ วเทติ ‘ภาสิตา นุ โข ภนฺเต ภควตา เอสา วาจา, ปิยชาติกา โสกปริเทวทุกฺขโทมนสฺสุปายาสา ปิยปฺปภวิกา’ติ”.}
\proseitem{356}{\csromanfolio{311}เอวเมตํ พฺราหฺมณ เอวเมตํ พฺราหฺมณ, ปิยชาติกา หิ พฺราหฺมณ โสกปริเทวทุกฺขโทมนสฺสุปายาสา ปิยปฺปภวิกาติ.\csromanspacer{} ตทมินาเปตํ พฺราหฺมณ ปริยาเยน เวทิตพฺพํ.\csromanspacer{} ยถา ปิยชาติกา โสกปริเทวทุกฺขโทมนสฺสุปายาสา ปิยปฺปภวิกา.\csromanspacer{} ภูตปุพฺพํ พฺราหฺมณ อิมิสฺสาเยว สาวตฺถิยา อญฺญตริสฺสา อิตฺถิยา มาตา กาลมกาสิ, สา ตสฺสา กาลกิริยาย อุมฺมตฺติกา ขิตฺตจิตฺตา รถิกาย รถิกํ\footnote{รถิยาย รถิยํ (สี, สฺยา, กํ, อิ)} สิงฺฆาฏเกน สิงฺฆาฏกํ อุปสงฺกมิตฺวา เอวมาห “อปิ เม มาตรํ อทฺทสฺสถ\footnote{อทฺทสถ (สี, อิ)}, อปิ เม มาตรํ อทฺทสฺสถา”ติ.\csromanspacer{} อิมินาปิ โข เอตํ พฺราหฺมณ ปริยาเยน เวทิตพฺพํ “ยถา ปิยชาติกา โสกปริเทวทุกฺขโทมนสฺสุปายาสา ปิยปฺปภวิกา”ติ.}

\prose{ภูตปุพฺพํ พฺราหฺมณ อิมิสฺสาเยว สาวตฺถิยา อญฺญตริสฺสา อิตฺถิยา ปิตา กาลมกาสิ.\csromanspacer{} ภาตา กาลมกาสิ.\csromanspacer{} ภคินี กาลมกาสิ.\csromanspacer{} ปุตฺโต กาลมกาสิ.\csromanspacer{} ธีตา กาลมกาสิ.\csromanspacer{} สามิโก กาลมกาสิ, สา ตสฺส กาลกิริยาย อุมฺมตฺติกา ขิตฺตจิตฺตา รถิกาย รถิกํ สิงฺฆาฏเกน สิงฺฆาฏกํ อุปสงฺกมิตฺวา เอวมาห “อปิ เม สามิกํ อทฺทสฺสถ, อปิ เม สามิกํ อทฺทสฺสถา”ติ.\csromanspacer{} อิมินาปิ โข เอตํ พฺราหฺมณ ปริยาเยน เวทิตพฺพํ “ยถา ปิยชาติกา โสกปริเทวทุกฺขโทมนสฺสุปายาสา ปิยปฺปภวิกา”ติ.}

\prose{ภูตปุพฺพํ พฺราหฺมณ อิมิสฺสาเยว สาวตฺถิยา อญฺญตรสฺส ปุริสสฺส มาตา กาลมกาสิ, โส ตสฺสา กาลกิริยาย อุมฺมตฺตโก ขิตฺตจิตฺโต รถิกาย รถิกํ สิงฺฆาฏเกน สิงฺฆาฏกํ อุปสงฺกมิตฺวา เอวมาห “อปิ เม มาตรํ อทฺทสฺสถ, อปิ เม มาตรํ \csromanfolio{312}อทฺทสฺสถา”ติ.\csromanspacer{} อิมินาปิ โข เอตํ พฺราหฺมณ ปริยาเยน เวทิตพฺพํ “ยถา ปิยชาติกา โสกปริเทวทุกฺขโทมนสฺสุปายาสา ปิยปฺปภวิกา”ติ.}

\prose{ภูตปุพฺพํ พฺราหฺมณ อิมิสฺสาเยว สาวตฺถิยา อญฺญตรสฺส ปุริสสฺส ปิตา กาลมกาสิ.\csromanspacer{} ภาตา กาลมกาสิ.\csromanspacer{} ภคินี กาลมกาสิ.\csromanspacer{} ปุตฺโต กาลมกาสิ.\csromanspacer{} ธีตา กาลมกาสิ.\csromanspacer{} ปชาปติ กาลมกาสิ, โส ตสฺสา กาลกิริยาย อุมฺมตฺตโก ขิตฺตจิตฺโต รถิกาย รถิกํ สิงฺฆาฏเกน สิงฺฆาฏกํ อุปสงฺกมิตฺวา เอวมาห “อปิ เม ปชาปติํ อทฺทสฺสถ, อปิ เม ปชาปติํ อทฺทสฺสถา”ติ.\csromanspacer{} อิมินาปิ โข เอตํ พฺราหฺมณ ปริยาเยน เวทิตพฺพํ “ยถา ปิยชาติกา โสกปริเทวทุกฺขโทมนสฺสุปายาสา ปิยปฺปภวิกา”ติ.}

\prose{ภูตปุพฺพํ พฺราหฺมณ อิมิสฺสาเยว สาวตฺถิยา อญฺญตรา อิตฺถี ญาติกุลํ อคมาสิ, ตสฺสา เต ญาตกา สามิกํ\footnote{สามิกา (สี)} อจฺฉินฺทิตฺวา อญฺญสฺส ทาตุกามา, สา จ ตํ น อิจฺฉติ.\csromanspacer{} อถ โข สา อิตฺถี สามิกํ เอตทโวจ “อิเม มํ\footnote{มม (สฺยา, กํ, อิ)} อยฺยปุตฺต ญตกา ตฺวํ\footnote{ตยา (สี), ตํ (สฺยา, กํ, อิ)} อจฺฉินฺทิตฺวา อญฺญสฺส ทาตุกามา, อหญฺจ ตํ น อิจฺฉามี”ติ.\csromanspacer{} อถ โข โส ปุริโส ตํ อิตฺถิํ ทฺวิธา เฉตฺวา อตฺตานํ อุปฺผาเลสิ\footnote{อุปฺปาเฏสิ (สี, อิ), โอผาเรสิ (ก)} อุโภ เปจฺจ ภวิสฺสามาติ.\csromanspacer{} อิมินาปิ โข เอตํ พฺราหฺมณ ปริยาเยน เวทิตพฺพํ “ยถา ปิยชาติกา โสกปริเทวทุกฺขโทมนสฺสุปายาสา ปิยปฺปภวิกา”ติ.}

\proseitem{357}{อถ โข นาฬิชงฺโฆ พฺราหฺมโณ ภควโต ภาสิตํ อภินนฺทิตฺวา อนุโมทิตฺวา อุฏฺฐายาสนา เยน มลฺลิกา เทวี เตนุปสงฺกมิ, อุปสงฺกมิตฺวา ยาวตโก อโหสิ ภควตา สทฺธิํ กถาสลฺลาโป, ตํ สพฺพํ มลฺลิกาย เทวิยา อาโรเจสิ.\csromanspacer{} อถ โข มลฺลิกา เทวี เยน ราชา ปเสนทิ โกสโล เตนุปสงฺกมิ, อุปสงฺกมิตฺวา ราชานํ ปเสนทิํ โกสลํ เอตทโวจ “ตํ กิํ มญฺญสิ มหาราช, ปิยา เต วชิรี กุมารี”ติ.\csromanspacer{} เอวํ มลฺลิเก ปิยา เม วชิรี กุมารีติ.\csromanspacer{} ตํ กิํ มญฺญสิ มหาราช, วชิริยา เต กุมาริยา วิปริณามญฺญถาภาวา อุปฺปชฺเชยฺยุํ โสกปริเทวทุกฺขโทมนสฺสุปายาสาติ.\csromanspacer{} วชิริยา เม มลฺลิเก กุมาริยา วิปริณามญฺญถาภาวา ชีวิตสฺสปิ}

\vspace{\tipitakaparskip}
\csromancenter{ปิยชาติกสุตฺต (87) สิยา อญฺญถตฺตํ, กิํ ปน เม น อุปฺปชฺชิสฺสนฺติ โสกปริเทวทุกฺขโทมนสฺสุปายาสาติ.\csromanspacer{} อิทํ โข ตํ มหาราช เตน ภควตา ชานตา ปสฺสตา อรหตา สมฺมาสมฺพุทฺเธน สนฺธาย ภาสิตํ “ปิยชาติกา โสกปริเทวทุกฺขโทมนสฺสุปายาสา ปิยปฺปภวิกา”ติ.}
\prose{\csromanfolio{313}ตํ กิํ มญฺญสิ มหาราช, ปิยา เต วาสภา ขตฺติยาติ.\csromanspacer{} เอวํ มลฺลิเก ปิยา เม วาสภา ขตฺติยาติ.\csromanspacer{} ตํ กิํ มญฺญสิ มหาราช, วาสภาย เต ขตฺติยาย วิปริณามญฺญถาภาวา อุปฺปชฺเชยฺยุํ โสกปริเทวทุกฺขโทมนสฺสุปายาสาติ.\csromanspacer{} วาสภาย เม มลฺลิเก ขตฺติยาย วิปริณามญฺญถาภาวา ชีวิตสฺสปิ สิยา อญฺญถตฺตํ, กิํ ปน เม น อุปฺปชฺชิสฺสนฺติ โสกปริเทวทุกฺขโทมนสฺสุปายาสาติ.\csromanspacer{} อิทํ โข ตํ มหาราช เตน ภควตา ชานตา ปสฺสตา อรหตา สมฺมาสมฺพุทฺเธน สนฺธาย ภาสิตํ “ปิยชาติกา โสกปริเทวทุกฺขโทมนสฺสุปายาสา ปิยปฺปภวิกา”ติ.}

\prose{ตํ กิํ มญฺญสิ มหาราช, ปิโย เต วิฏฏูโภ\footnote{วิฑูฑโภ (สี, สฺยา, กํ, อิ)} เสนาปตีติ.\csromanspacer{} เอวํ มลฺลิเก ปิโย เม วิฏฏูโภ เสนาปตีติ.\csromanspacer{} ตํ กิํ มญฺญสิ มหาราช, วิฏฏูภสฺส เต เสนาปติสฺส วิปริณามญฺญถาภาวา อุปฺปชฺเชยฺยุํ โสกปริเทวทุกฺขโทมนสฺสุปายาสาติ.\csromanspacer{} วิฏฏูภสฺส เม มลฺลิเก เสนาปติสฺส วิปริณามญฺญถาภาวา ชีวิตสฺสปิ สิยา อญฺญถตฺตํ, กิํ ปน เม น อุปฺปชฺชิสฺสนฺติ โสกปริเทวทุกฺขโทมนสฺสุปายาสาติ.\csromanspacer{} อิทํ โข ตํ มหาราช เตน ภควตา ชานตา ปสฺสตา อรหตา สมฺมาสมฺพุทฺเธน สนฺธาย ภาสิตํ “ปิยชาติกา โสกปริเทวทุกฺขโทมนุสฺสุปายาสา วิยปฺปภวิกา”ติ.}

\prose{ตํ กิํ มญฺญสิ มหาราช, ปิยา เต อหนฺติ.\csromanspacer{} เอวํ มลฺลิเก ปิยา เมสิ ตฺวนฺติ.\csromanspacer{} ตํ กิํ มญฺญสิ มหาราช, มยฺหํ เต วิปริณามญฺญถาภาวา อุปฺปชฺเชยฺยุํ โสกปริเทวทุกฺขโทมนสฺสุปายาสาติ.\csromanspacer{} ตุยฺหญฺหิ เม มลฺลิเก วิปริณามญฺญถาภาวา ชีวิตสฺสปิ สิยา อญฺญถตฺตํ, กิํ ปน เม น อุปฺปชฺชิสฺสนฺติ โสกปริเทวทุกฺขโทมนสฺสุปายาสาติ.\csromanspacer{} อิทํ โข ตํ มหาราช เตน ภควตา ชานตา ปสฺสตา อรหตา สมฺมาสมฺพุทฺเธน สนฺธาย ภาสิตํ “ปิยชาติกา โสกปริเทวทุกฺขโทมนสฺสุปายาสา ปิยปฺปภวิกา”ติ.}

\prose{\csromanfolio{314}ตํ กิํ มญฺญสิ มหาราช, ปิยา เต กาสิโกสลาติ.\csromanspacer{} เอวํ มลฺลิเก ปิยา เม กาสิโกสลา, กาสิโกสลานํ มลฺลิเก อานุภาเวน กาสิกจนฺทนํ ปจฺจนุโภม, มาลาคนฺธวิเลปนํ ธาเรมาติ.\csromanspacer{} ตํ กิํ มญฺญสิ มหาราช, กาสิโกสลานํ เต วิปริณามญฺญถาภาวา อุปฺปชฺเชยฺยุํ โสกปริเทวทุกฺขโทมนสฺสุปายาสาติ.\csromanspacer{} กาสิโกสลานญฺหิ มลฺลิเก วิปริณามญฺญถาภาวา ชีวิตสฺสปิ สิยา อญฺญถตฺตํ, กิํ ปน เม น อุปฺปชฺชิสฺสนฺติ โสกปริเทวทุกฺขโทมนสฺสุปายาสาติ.\csromanspacer{} อิทํ โข ตํ มหาราช เตน ภควตา ชานตา ปสฺสตา อรหตา สมฺมาสมฺพุทฺเธน สนฺธาย ภาสิตํ “ปิยชาติกา โสกปริเทวทุกฺขโทมนสฺสุปายาสา ปิยปฺปภวิกา”ติ.}

\csromanmatikaanchor{mtk.44}
\csromantocmarkref{subsection}{7. ปิยชาติกสุตฺต}{mtk.44}
\bookmark[dest={mtk.44},level=2]{7. ปิยชาติกสุตฺต}
\prose{อจฺฉริยํ มลฺลิเก, อพฺภุตํ มลฺลิเก, ยาวญฺจ โส ภควา ปญฺญาย อติวิชฺฌ มญฺเญ\footnote{ปฏิวิชฺฌ ปญฺญาย (ก)} ปสฺสติ, เอหิ มลฺลิเก อาจเมหีติ\footnote{อาจาเมหีติ (สี, อิ)}.\csromanspacer{} อถ โข ราชา ปเสนทิ โกสโล อุฏฺฐายาสนา เอกํสํ อุตฺตราสงฺคํ กริตฺวา เยน ภควา เตนญฺชลิํ ปณาเมตฺวา ติกฺขตฺตุํ อุทานํ อุทาเนสิ “นโม ตสฺส ภควโต อรหโต สมฺมาสมฺพุทฺธสฺส, นโม ตสฺส ภควโต อรหโต สมฺมาสมฺพุทฺธสฺส, นโม ตสฺส ภควโต อรหโต สมฺมาสมฺพุทฺธสฺสา”ติ.}

\nitthitamruled{ปิยชาติกสุตฺตํ นิฏฺฐิตํ สตฺตมํ.}
\csromanheader{2}{8. พาหิติกสุตฺต}
\proseitem{358}{เอวํ เม สุตํ\csromandash{}เอกํ สมยํ ภควา สาวตฺถิยํ วิหรติ เชตวเน อนาถปิณฺฑิกสฺส อาราเม.\csromanspacer{} อถ โข อายสฺมา อานนฺโท ปุพฺพณฺหสมยํ นิวาเสตฺวา ปตฺตจีวรมาทาย สาวตฺถิยํ ปิณฺฑาย ปาวิสิ, สาวตฺถิยํ ปิณฺฑาย จริตฺวา ปจฺฉาภตฺตํ ปิณฺฑปาตปฏิกฺกนฺโต เยน ปุพฺพาราโม มิคารมาตุปาสาโท เตนุปสงฺกมิ ทิวาวิหาราย.\csromanspacer{} เตน โข ปน สมเยน ราชา ปเสนทิ โกสโล เอกปุณฺฑรีกํ นาคํ อภิรุหิตฺวา สาวตฺถิยา นิยฺยาติ ทิวา ทิวสฺส.\csromanspacer{} อทฺทสา โข ราชา ปเสนทิ โกสโล อายสฺมนฺตํ อานนฺทํ ทูรโตว อาคจฺฉนฺตํ, ทิสฺวาน สิริวฑฺฒํ มหามตฺตํ อามนฺเตสิ “อายสฺมา โน เอโส สมฺม สิริวฑฺฒ}

\vspace{\tipitakaparskip}
\csromancenter{พาหิติกสุตฺต (88) อานนฺโท”ติ.\csromanspacer{} เอวํ มหาราช อายสฺมา เอโส อานนฺโทติ.\csromanspacer{} อถ โข ราชา ปเสนทิ โกสโล อญฺญตรํ ปุริสํ อามนฺเตสิ “เอหิ ตฺวํ อมฺโภ ปุริส เยนายสฺมา อานนฺโท เตนุปสงฺกม, อุปสงฺกมิตฺวา มม วจเนน อายสฺมโต อานนฺทสฺส ปาเท สิรสา วนฺทาหิ ‘ราชา ภนฺเต ปเสนทิ โกสโล อายสฺมโต อานนฺทสฺส ปาเท สิรสา วนฺทตี’ติ, เอวญฺจ วเทหิ ‘สเจ กิร ภนฺเต อายสฺมโต อานนฺทสฺส น กิญฺจิ อจฺจายิกํ กรณียํ, อาคเมตุ กิร ภนฺเต อายสฺมา อานนฺโท มุหุตฺตํ อนุกมฺปํ อุปาทายา’ติ”. “เอวํ เทวา”ติ โข โส ปุริโส รญฺโญ ปเสนทิสฺส โกสลสฺส ปฏิสฺสุตฺวา เยนายสฺมา อานนฺโท เตนุปสงฺกมิ, อุปสงฺกมิตฺวา อายสฺมนฺตํ อานนฺทํ อภิวาเทตฺวา เอกมนฺตํ อฏฺฐาสิ, เอกมนฺตํ ฐิโต โข โส ปุริโส อายสฺมนฺตํ อานนฺทํ เอตทโวจ “ราชา ภนฺเต ปเสนทิ โกสโล อายสฺมโต อานนฺทสฺส ปาเท สิรสา วนฺทติ, เอวญฺจ วเทติ ‘สเจ กิร ภนฺเต อายสฺมโต อานนฺทสฺส น กิญฺจิ อจฺจายิกํ กรณียํ, อาคเมตุ กิร ภนฺเต อายสฺมา อานนฺโท มุหุตฺตํ อนุกมฺปํ อุปาทายา’ติ”.\csromanspacer{} อธิวาเสสิ โข อายสฺมา อานนฺโท ตุณฺหีภาเวน.\csromanspacer{} อถ โข ราชา ปเสนทิ โกสโล ยาวติกา นาคสฺส ภูมิ, นาเคน คนฺตฺวา นาคา ปจฺโจโรหิตฺวา ปตฺติโกว เยนายสฺมา อานนฺโท เตนุปสงฺกมิ, อุปสงฺกมิตฺวา อายสฺมนฺตํ อานนฺทํ อภิวาเทตฺวา เอกมนฺตํ อฏฺฐาสิ, เอกมนฺตํ ฐิโต โข ราชา ปเสนทิ โกสโล อายสฺมนฺตํ อานนฺทํ เอตทโวจ “สเจ ภนฺเต อายสฺมโต อานนฺทสฺส น กิญฺจิ อจฺจายิกํ กรณียํ, สาธุ ภนฺเต, อายสฺมา อานนฺโท เยน อจิรวติยา นทิยา ตีรํ เตนุปสงฺกมตุ อนุกมฺปํ อุปาทายา”ติ.\csromanspacer{} อธิวาเสสิ โข อายสฺมา อานนฺโท ตุณฺหีภาเวน.}
\proseitem{359}{\csromanfolio{315}อถ โข อายสฺมา อานนฺโท เยน อจิรวติยา นทิยา ตีรํ เตนุปสงฺกมิ, อุปสงฺกมิตฺวา อญฺญตรสฺมิํ รุกฺขมูเล ปญฺญตฺเต อาสเน นิสีทิ.\csromanspacer{} อถ โข ราชา ปเสนทิ โกสโล ยาวติกา นาคสฺส ภูมิ, นาเคน คนฺตฺวา นาคา ปจฺโจ โรหิตฺวา ปตฺติโกว เยนายสฺมา อานนฺโท เตนุปสงฺกมิ, อุปสงฺกมิตฺวา อายสฺมนฺตํ อานนฺทํ อภิวาเทตฺวา เอกมนฺตํ อฏฺฐาสิ, เอกมนฺตํ ฐิโต โข ราชา ปเสนทิ โกสโล อายสฺมนฺตํ อานนฺทํ เอตทโวจ “อิธ ภนฺเต \csromanfolio{316}อายสฺมา อานนฺโท หตฺถตฺถเร นิสีทตู”ติ.\csromanspacer{} อลํ มหาราช, นิสีท ตฺวํ, นิสินฺโน อหํ สเก อาสเนติ.\csromanspacer{} นิสีทิ โข ราชา ปเสนทิ โกสโล ปญฺญตฺเต อาสเน, นิสชฺช โข ราชา ปเสนทิ โกสโล อายสฺมนฺตํ อานนฺทํ เอตทโวจ “กิํ นุ โข ภนฺเต อานนฺท โส ภควา ตถารูปํ กายสมาจารํ สมาจเรยฺย, ยฺวาสฺส กายสมาจาโร โอปารมฺโภ สมเณหิ พฺราหฺมเณหี”ติ\footnote{พฺราหฺมเณหิ วิญฺญูหีติ (สพฺพตฺถ) อฏฺฐกถา ฏีกา โอโลเกตพฺพา.}.\csromanspacer{} น โข มหาราช โส ภควา ตถารูปํ กายสมาจารํ สมาจเรยฺย, ยฺวาสฺส กายสมาจาโร โอปารมฺโภ สมเณหิ พฺราหฺมเณหิ วิญฺญูหีติ.}

\prose{กิํ ปน ภนฺเต อานนฺท โส ภควา ตถารูปํ วจีสมาจารํ ฯเปฯ มโนสมาจารํ สมาจเรยฺย, ยฺวาสฺส มโนสมาจาโร โอปารมฺโภ สมเณหิ พฺราหฺมเณหีติ1.\csromanspacer{} น โข มหาราช โส ภควา ตถารูปํ มโนสมาจารํ สมาจเรยฺย, ยฺวาสฺส มโนสมาจาโร โอปารมฺโภ สมเณหิ พฺราหฺมเณหิ วิญฺญูหีติ.}

\prose{อจฺฉริยํ ภนฺเต, อพฺภุตํ ภนฺเต, ยญฺหิ มยํ ภนฺเต นาสกฺขิมฺหา ปเญฺหน ปริปูเรตุํ.\csromanspacer{} ตํ ภนฺเต อายสฺมตา อานนฺเทน ปญฺหสฺส เวยฺยากรเณน ปริปูริตํ.\csromanspacer{} เย เต ภนฺเต พาลา อพฺยตฺตา อนนุวิจฺจ อปริโยคาเหตฺวา ปเรสํ วณฺณํ วา อวณฺณํ วา ภาสนฺติ, น มยํ ตํ สารโต ปจฺจาคจฺฉาม.\csromanspacer{} เย ปน\footnote{เย จ โข (สี, สฺยา, กํ, อิ)} เต ภนฺเต ปณฺฑิตา วิยตฺตา\footnote{พฺยตฺตา (สี, สฺยา, กํ, อิ)} เมธาวิโน อนุวิจฺจ ปริโยคาเหตฺวา ปเรสํ วณฺณํ วา อวณฺณํ วา ภาสนฺติ, มยํ ตํ สารโต ปจฺจาคจฺฉาม.}

\proseitem{360}{กตโม ปน ภนฺเต อานนฺท กายสมาจาโร โอปารมฺโภ สมเณหิ พฺราหฺมเณหิ วิญฺญูหีติ.\csromanspacer{} โย โข มหาราช กายาสมาจาโร อกุสโล.}

\prose{กตโม ปน ภนฺเต กายสมาจาโร อกุสโล.\csromanspacer{} โย โข มหาราช กายสมาจาโร สาวชฺโช.}

\prose{กตโม ปน ภนฺเต กายสมาจาโร สาวชฺโช.\csromanspacer{} โย โข มหาราช กายสมาจาโร สพฺยาพชฺโฌ\footnote{สพฺยาปชฺโฌ (สี, สฺยา, กํ, อิ), สพฺยาปชฺโช (ก)}.}

\csromanheader{2}{8. พาหิติกสุตฺต (88)}
\prose{\csromanfolio{317}กตโม ปน ภนฺเต กายสมาจาโร สพฺยาพชฺโฌ.\csromanspacer{} โย โข มหาราช กายสมาจาโร ทุกฺขวิปาโก.}

\prose{กตโม ปน ภนฺเต กายสมาจาโร ทุกฺขวิปาโก.\csromanspacer{} โย โข มหาราช กายสมาจาโร อตฺตพฺยาพาธายปิ สํวตฺตติ, ปรพฺยาพาธายปิ สํวตฺตติ, อุภยพฺยาพาธายปิ สํวตฺตติ.\csromanspacer{} ตสฺส อกุสลา ธมฺมา อภิวฑฺฒนฺติ, กุสลา ธมฺมา ปริหายนฺติ.\csromanspacer{} เอวรูโป โข มหาราช กายสมาจาโร โอปารมฺโภ สมเณหิ พฺราหฺมเณหิ วิญฺญูหีติ.}

\prose{กตโม ปน ภนฺเต อานนฺท วจีสมาจาโร ฯเปฯ มโนสมาจาโร โอปารมฺโภ สมเณหิ พฺราหฺมเณหิ วิญฺญูหีติ.\csromanspacer{} โย โข มหาราช มโนสมาจาโร อกุสโล.}

\prose{กตโม ปน ภนฺเต มโนสมาจาโร อกุสโล.\csromanspacer{} โย โข มหาราช มโนสมาจาโร สาวชฺโช.}

\prose{กตโม ปน ภนฺเต มโนสมาจาโร สาวชฺโช.\csromanspacer{} โย โข มหาราช มโนสมาจาโร สพฺยาพชฺโฌ.}

\prose{กตโม ปน ภนฺเต มโนสมาจาโร สพฺยาพชฺโฌ.\csromanspacer{} โย โข มหาราช มโนสมาจาโร ทุกฺขวิปาโก.}

\prose{กตโม ปน ภนฺเต มโนสมาจาโร ทุกฺขวิปาโก.\csromanspacer{} โย โข มหาราช มโนสมาจาโร อตฺตพฺยาพาธายปิ สํวตฺตติ, ปรพฺยาพาธายปิ สํวตฺตติ, อุภยพฺยาพาธายปิ สํวตฺตติ.\csromanspacer{} ตสฺส อกุสลา ธมฺมา อภิวฑฺฒนฺติ, กุสลา ธมฺมา ปริหายนฺติ.\csromanspacer{} เอวรูโป โข มหาราช มโนสมาจาโร โอปารมฺโภ สมเณหิ พฺราหฺมเณหิ วิญฺญูหีติ.}

\csromanmatikaanchor{mtk.45}
\csromantocmarkref{subsection}{8. พาหิติกสุตฺต}{mtk.45}
\bookmark[dest={mtk.45},level=2]{8. พาหิติกสุตฺต}
\prose{กิํ นุ โข ภนฺเต อานนฺท โส ภควา สพฺเพสํเยว อกุสลานํ ธมฺมานํ ปหานํ วณฺเณตีติ.\csromanspacer{} สพฺพากุสลธมฺมปหีโน โข มหาราช ตถาคโต กุสลธมฺมสมนฺนาคโตติ.}

\proseitem{361}{กตโม ปน ภนฺเต อานนฺท กายสมาจาโร อโนปารมฺโภ สมเณหิ พฺราหฺมเณหิ วิญฺญูหีติ.\csromanspacer{} โย โข มหาราช กายสมาจาโร กุสโล.}

\prose{\csromanfolio{318}กตโม ปน ภนฺเต กายสมาจาโร กุสโล.\csromanspacer{} โย โข มหาราช กายสมาจาโร อนวชฺโช.}

\prose{กตโม ปน ภนฺเต กายสมาจาโร อนวชฺโช.\csromanspacer{} โย โข มหาราช กายสมาจาโร อพฺยาพชฺโฌ.}

\prose{กตโม ปน ภนฺเต กายสมาจาโร อพฺยาพชฺโฌ.\csromanspacer{} โย โข มหาราช กายสมาจาโร สุขวิปาโก.}

\prose{กตโม ปน ภนฺเต กายสมาจาโร สุขวิปาโก.\csromanspacer{} โย โข มหาราช กายสมาจาโร เนวตฺตพฺยาพาธายปิ สํวตฺตติ, น ปรพฺยาพาธายปิ สํวตฺตติ, น อุภยพฺยาพาธายปิ สํวตฺตติ.\csromanspacer{} ตสฺส อกุสลา ธมฺมา ปริหายนฺติ, กุสลา ธมฺมา อภิวฑฺฒนฺติ.\csromanspacer{} เอวรูโป โข มหาราช กายสมาจาโร อโนปารมฺโภ สมเณหิ พฺราหฺมเณหิ วิญฺญูหีติ.}

\prose{กตโม ปน ภนฺเต อานนฺท วจีสมาจาโร ฯเปฯ มโนสมาโร อโนปารมฺโภ สมเณหิ พฺราหฺมเณหิ วิญฺญูหีติ.\csromanspacer{} โย โข มหาราช มโนสมาจาโร กุสโล.}

\prose{กตโม ปน ภนฺเต มโนสมาจาโร กุสโล.\csromanspacer{} โย โข มหาราช มโนสมาจาโร อนวชฺโช.}

\prose{กตโม ปน ภนฺเต มโนสมาจาโร อนวชฺโช.\csromanspacer{} โย โข มหาราช มโนสมาจาโร อพฺยาพชฺโฌ.}

\prose{กตโม ปน ภนฺเต มโนสมาจาโร อพฺยาพชฺโฌ.\csromanspacer{} โย โข มหาราช มโนสมาจาโร สุขวิปาโก.}

\prose{กตโม ปน ภนฺเต มโนสมาจาโร สุขวิปาโก.\csromanspacer{} โย โข มหาราช มโนสมาจาโร เนวตฺตพฺยาพาธายปิ สํวตฺตติ, น ปรพฺยาพาธายปิ สํวตฺตติ, น อุภยพฺยาพาธายปิ สํวตฺตติ.\csromanspacer{} ตสฺส อกุสลา ธมฺมา ปริหายนฺติ, กุสลา ธมฺมา อภิวฑฺฒนฺติ.\csromanspacer{} เอวรูโป โข มหาราช มโนสมาจาโร อโนปารมฺโภ สมเณหิ พฺราหฺมเณหิ วิญฺญูหีติ.}

\prose{กิํ ปน ภนฺเต อานนฺท โส ภควา สพฺเพสํเยว กุสลานํ ธมฺมานํ อุปสมฺปทํ วณฺเณตีติ.\csromanspacer{} สพฺพากุสลธมฺมปหีโน โข มหาราช ตถาคโต กุสลธมฺมสมนฺนาคโตติ.}

\csromanheader{2}{8. พาหิติกสุตฺต (88)}
\proseitem{362}{\csromanfolio{319}อจฺฉริยํ ภนฺเต, อพฺภุตํ ภนฺเต, ยาว สุภาสิตํ จิทํ\footnote{สุภาสิตมิทํ (สี)} ภนฺเต อายสฺมตา อานนฺเทน, อิมินา จ มยํ ภนฺเต อายสฺมโต อานนฺทสฺส สุภาสิเตน อตฺตมนาภิรทฺธา, เอวํ อตฺตมนาภิรทฺธา จ มยํ ภนฺเต อายสฺมโต อานนฺทสฺส สุภาสิเตน.\csromanspacer{} สเจ ภนฺเต อายสฺมโต อานนฺทสฺส หตฺถิรตนํ กปฺเปยฺย, หตฺถิรตนมฺปิ มยํ อายสฺมโต อานนฺทสฺส ทเทยฺยาม.\csromanspacer{} สเจ ภนฺเต อายสฺมโต อานนฺทสฺส อสฺสรตนํ กปฺเปยฺย, อสฺสรตนมฺปิ มยํ อายสฺมโต อานนฺทสฺส ทเทยฺยาม.\csromanspacer{} สเจ ภนฺเต อายสฺมโต อานนฺทสฺส คามวรํ กปฺเปยฺย, คามวรมฺปิ มยํ อายสฺมโต อานนฺทสฺส ทเทยฺยาม.\csromanspacer{} อปิ จ ภนฺเต มยมฺเปตํ\footnote{มยเมว ตํ (สี), มยมฺปเนตํ (สฺยา, กํ)} ชานาม “เนตํ อายสฺมโต อานนฺทสฺส กปฺปตี”ติ.\csromanspacer{} อยํ เม ภนฺเต พาหิติกา รญฺญา มาคเธน อชาตสตฺตุนา เวเทหิปุตฺเตน วตฺถนาฬิยา\footnote{ฉตฺตนาฬิยา (สฺยา, กํ, อิ)} ปกฺขิปิตฺวา ปหิตา โสฬสสมา อายาเมน อฏฺฐสมา วิตฺตาเรน.\csromanspacer{} ตํ ภนฺเต อายสฺมา อานนฺโท ปฏิคฺคณฺหาตุ อนุกมฺปํ อุปาทายาติ.\csromanspacer{} อลํ มหาราช, ปริปุณฺณํ เม ติจีวรนฺติ.}

\prose{อยํ ภนฺเต อจิรวตี นที ทิฏฺฐา อายสฺมตา เจว อานนฺเทน อเมฺหหิ จ, ยทา อุปริปพฺพเต มหาเมโฆ อภิปฺปวุฏฺโฐ โหติ.\csromanspacer{} อถายํ อจิรวตี นที อุภโต กูลานิ สํวิสฺสนฺทนฺตี คจฺฉติ.\csromanspacer{} เอวเมว โข ภนฺเต อายสฺมา อานนฺโท อิมาย พาหิติกาย อตฺตโน ติจีวรํ กริสฺสติ.\csromanspacer{} ยํ ปนายสฺมโต อานนฺทสฺส ปุราณํ ติจีวรํ, ตํ สพฺรหฺมจารีหิ สํวิภชิสฺสติ.\csromanspacer{} เอวายํ อมฺหากํ ทกฺขิณา สํวิสฺสนฺทนฺตี มญฺเญ คมิสฺสติ.\csromanspacer{} ปฏิคฺคณฺหาตุ ภนฺเต อายสฺมา อานนฺโท พาหิติกนฺติ.\csromanspacer{} ปฏิคฺคเหสิ โข อายสฺมา อานนฺโท พาหิติกํ.}

\prose{อถ โข ราชา ปเสนทิ โกสโล อายสฺมนฺตํ อานนฺทํ เอตโวจ “หนฺท จ ทานิ มยํ ภนฺเต อานนฺท คจฺฉาม, พหุกิจฺจา มยํ พหุกรณียา”ติ.\csromanspacer{} ยสฺสทานิ ตฺวํ มหาราช กาลํ มญฺญสีติ.\csromanspacer{} อถ โข ราชา ปเสนทิ โกสโล อายสฺมโต อานนฺทสฺส ภาสิตํ อภินนฺทิตฺวา อนุโมทิตฺวา อุฏฺฐายาสนา อายสฺมนฺตํ อานนฺทํ อภิวาเทตฺวา ปทกฺขิณํ กตฺวา ปกฺกามิ.}

\proseitem{363}{\csromanfolio{320}อถ โข อายสฺมา อานนฺโท อจิรปกฺกนฺตสฺส รญฺโญ ปเสนทิสฺส โกสลสฺส เยน ภควา เตนุปสงฺกมิ, อุปสงฺกมิตฺวา ภควนฺตํ อภิวาเทตฺวา เอกมนฺตํ นิสีทิ, เอกมนฺตํ นิสินฺโน โข อายสฺมา อานนฺโท ยาวตโก อโหสิ รญฺญา ปเสนทินา โกสเลน สิทฺธิํ กถาสลฺลาโป, ตํ สพฺพํ ภควโต อาโรเจสิ, ตญฺจ พาหิติกํ ภควโต ปาทาสิ.\csromanspacer{} อถ โข ภควา ภิกฺขู อามนฺเตสิ “ลาภา ภิกฺขเว รญฺโญ ปเสนทิสฺส โกสลสฺส, สุลทฺธลาภา ภิกฺขเว รญฺโญ ปเสนทิสฺส โกสลสฺส, ยํ ราชา ปเสนทิ โกสโล ลภติ อานนฺทํ ทสฺสนาย, ลภติ ปยิรุปาสนายา”ติ.}

\prose{อิทมโวจ ภควา.\csromanspacer{} อตฺตมนา เต ภิกฺขู ภควโต ภาสิตํ อภินนฺทุนฺติ.}

\nitthitamruled{พาหิติกสุตฺตํ นิฏฺฐิตํ อฏฺฐมํ.}
\csromanheader{2}{9. ธมฺมเจติยสุตฺต}
\proseitem{364}{เอวํ เม สุตํ\csromandash{}เอกํ สมยํ ภควา สกฺเกสุ วิหรติ เมทาฬุปํ\footnote{เมตฬูปํ (สี), เมทฬุมฺปํ (อิ)} นาม สกฺยานํ นิคโม.\csromanspacer{} เตน โข ปน สมเยน ราชา ปเสนทิ โกสโล นครกํ อนุปฺปตฺโต โหติ เกนจิเทว กรณีเยน.\csromanspacer{} อถ โข ราชา ปเสนทิ โกสโล ทีฆํ การายนํ อามนฺเตสิ “โยเชหิ สมฺม การายน ภทฺรานิ ภทฺรานิ ยานานิ, อุยฺยานภูมิํ คจฺฉาม สุภูมิํ ทสฺสนายา”ติ\footnote{สุภูมิทสฺสนายาติ (ที 2. 18 ปิฏฺเฐ)}. “เอวํ เทวา”ติ โข ทีโฆ การายโน รญฺโญ ปเสนทิสฺส โกสลสฺส ปฏิสฺสุตฺวา ภทฺรานิ ภทฺรานิ ยานานิ โยชาเปตฺวา รญฺโญ ปเสนทิสฺส โกสลสฺส ปฏิเวเทสิ “ยุตฺตานิ โข เต เทว ภทฺรานิ ภทฺรานิ ยานานิ, ยสฺสทานิ กาลํ มญฺญสี”ติ.\csromanspacer{} อถ โข ราชา ปเสนทิ โกสโล ภทฺรํ ยานํ อภิรุหิตฺวา ภเทฺรหิ ภเทฺรหิ ยาเนหิ นครกมฺหา นิยฺยาสิ มหจฺจา ราชานุภาเวน.\csromanspacer{} เยน อาราโม เตน ปายาสิ.\csromanspacer{} ยาวติกา ยานสฺส ภูมิ, ยาเนน คนฺตฺวา ยานา ปจฺโจโรหิตฺวา ปตฺติโกว อารามํ ปาวิสิ.\csromanspacer{} อทฺทสา โข ราชา ปเสนทิ โกสโล อาราเม ชงฺฆาวิหารํ อนุจงฺกมมาโน}

\vspace{\tipitakaparskip}
\csromancenter{ธมฺมเจติยสุตฺต (89) อนุวิจรมาโน รุกฺขมูลานิ ปาสาทิกานิ ปสาทนียานิ อปฺปสทฺทานิ อปฺปนิคฺโฆสานิ วิชนวาตานิ มนุสฺสราหสฺเสยฺยกานิ\footnote{มนุสฺสราหเสยฺยกานิ (สี, อิ)} ปฏิสลฺลานสารุปฺปานิ.\csromanspacer{} ทิสฺวาน ภควนฺตํเยว อารพฺภ สติ อุทปาทิ “อิมานิ โข ตานิ รุกฺขมูลานิ ปาสาทิกานิ ปสาทนียานิ อปฺปสทฺทานิ อปฺปนิคฺโฆสานิ วิชนวาตานิ มนุสฺสราหสฺเสยฺยกานิ ปฏิสลฺลานสารุปฺปานิ.\csromanspacer{} ยตฺถ สุทํ มยํ ตํ ภควนฺตํ ปยิรุปาสาม อรหนฺตํ สมฺมาสมฺพุทฺธนฺ”ติ.}
\proseitem{365}{\csromanfolio{321}อถ โข ราชา ปเสนทิ โกสโล ทีฆํ การายนํ อามนฺเตสิ “อิมานิ โข สมฺม การายน ตานิ รุกฺขมูลานิ ปาสาทิกานิ ปสาทนียานิ อปฺปสทฺทานิ อปฺปนิคฺโฆสานิ วิชนวาตานิ มนุสฺสราหสฺเสยฺยกานิ ปฏิสลฺลานสารุปฺปานิ.\csromanspacer{} ยตฺถ สุทํ มยํ ตํ ภควนฺตํ ปยิรุปาสาม อรหนฺตํ สมฺมาสมฺพุทฺธํ.\csromanspacer{} กหํ นุ โข สมฺม การายน เอตรหิ โส ภควา วิหรติ อรหํ สมฺมาสมฺพุทฺโธ”ติ.\csromanspacer{} อตฺถิ มหาราช เมทาฬุปํ นาม สกฺยานํ นิคโม, ตตฺถ โส ภควา เอตรหิ วิหรติ อรหํ สมฺมาสมฺพุทฺโธติ.\csromanspacer{} กีวทูเร\footnote{กีวทูโร (สี, สฺยา, กํ, อิ)} ปน สมฺม การายน นครกมฺหา เมทาฬุปํ นาม สกฺยานํ นิคโม โหตีติ.\csromanspacer{} น ทูเร มหาราช ตีณิ โยชนานิ สกฺกา ทิวสาวเสเสน คนฺตุนฺติ.\csromanspacer{} เตน หิ สมฺม การายน โยเชหิ ภทฺรานิ ภทฺรานิ ยานานิ, คมิสฺสาม มยํ ตํ ภควนฺตํ ทสฺสนาย อรหนฺตํ สมฺมาสมฺพุทฺธนฺติ. “เอวํ เทวา”ติ โข ทีโฆ การายโน รญฺโญ ปเสนทิสฺส โกสลสฺส ปฏิสฺสุตฺวา ภทฺรานิ ภทฺรานิ ยานานิ โยชาเปตฺวา รญฺโญ ปเสนทิสฺส โกสลสฺส ปฏิเวเทสิ “ยุตฺตานิ โข เต เทว ภทฺรานิ ภทฺรานิ ยานานิ, ยสฺสทานิ กาลํ มญฺญสี”ติ.\csromanspacer{} อถ โข ราชา ปเสนทิ โกสโล ภทฺรํ ยานํ อภิรุหิตฺวา ภเทฺรหิ ภเทฺรหิ ยาเนหิ นครกมฺหา เยน เมทาฬุปํ นาม สกฺยานํ นิคโม เตน ปายาสิ.\csromanspacer{} เตนว ทิวสาเสเสน เมทาฬุปํ นาม สกฺยานํ นิคมํ สมฺปาปุณิ.\csromanspacer{} เยน อาราโม เตน ปายาสิ, ยาวติกา ยานสฺส ภูมิ, ยาเนน คนฺตฺวา ยานา ปจฺโจโรหิตฺวา ปตฺติโกว อารามํ ปาวิสิ.}

\proseitem{366}{เตน โข ปน สมเยน สมฺพหุลา ภิกฺขู อพฺโภกาเส จงฺกมนฺติ.\csromanspacer{} อถ โข ราชา ปเสนทิ โกสโล เยน เต ภิกฺขู เตนุปสงฺกมิ, อุปสงฺกมิตฺวา เต ภิกฺขู เอตทโวจ “กหํ นุ โข ภนฺเต \csromanfolio{322}เอตรหิ โส ภควา วิหรติ อรหํ สมฺมาสมฺพุทฺโธ, ทสฺสนกามา หิ มยํ ตํ ภควนฺตํ อรหนฺตํ สมฺมาสมฺพุทฺธนฺ”ติ.\csromanspacer{} เอโส มหาราช วิหาโร สํวุตทฺวาโร, เตน อปฺปสทฺโท อุปสงฺกมิตฺวา อตรมาโน อาฬินฺทํ ปวิสิตฺวา อุกฺกาสิตฺวา อคฺคฬํ อาโกเฏหิ, วิวริสฺสติ ภควา เต ทฺวารนฺติ.\csromanspacer{} อถ โข ราชา ปเสนทิ โกสโล ตตฺเถว ขคฺคญฺจ อุณฺหีสญฺจ ทีฆสฺส การายนสฺส ปาทาสิ.\csromanspacer{} อถ โข ทีฆสฺส การายนสฺส เอตทโหสิ “รหายติ โข ทานิ ราชา\footnote{มหาราชา (สี, สฺยา, กํ, อิ)} อิเธว\footnote{เตนิเธว (สี)} ทานิ มยา ฐาตพฺพนฺ”ติ.\csromanspacer{} อถ โข ราชา ปเสนทิ โกสโล เยน โส วิหาโร สํวุตทฺวาโร, เตน อปฺปสทฺโท อุปสงฺกมิตฺวา อตรมาโน อาฬินฺทํ ปวิสิตฺวา อุกฺกาสิตฺวา อคฺคฬํ อาโกเฏสิ.\csromanspacer{} วิวริ ภควา ทฺวารํ.\csromanspacer{} อถ โข ราชา ปเสนทิ โกสโล วิหารํ ปวิสิตฺวา ภควโต ปาเทสุ สิรสา นิปติตฺวา ภควโต ปาทานิ มุเขน จ ปริจุมฺพติ, ปาณีหิ จ ปริสมฺพาหติ, นามญฺจ สาเวติ “ราชาหํ ภนฺเต ปเสนทิ โกสโล, ราชาหํ ภนฺเต ปเสนทิ โกสโล”ติ.}

\proseitem{367}{กิํ ปน ตฺวํ มหาราช อตฺถวสํ สมฺปสฺสมาโน อิมสฺมิํ สรีเร เอวรูปํ ปรมนิปจฺจการํ กโรสิ, มิตฺตูปหารํ\footnote{จิตฺตูปหารํ (สี)} อุปทํเสสีติ.\csromanspacer{} อตฺถิ โข เม ภนฺเต ภควติ ธมฺมนฺวโย โหติ “สมฺมาสมฺพุทฺโธ ภควา, สฺวากฺขาโต ภควตา ธมฺโม, สุปฺปฏิปนฺโน ภควโต สาวกสํโฆ”ติ.\csromanspacer{} อิธาหํ ภนฺเต ปสฺสามิ เอเก สมณพฺราหฺมเณ ปริยนฺตกตํ พฺรหฺมจริยํ จรนฺเต ทสปิ วสฺสานิ วีสมฺปิ วสฺสานิ ติํสมฺปิ วสฺสานิ จตฺตารีสมฺปิ วสฺสานิ.\csromanspacer{} เต อปเรน สมเยน สุนฺหาตา สุวิลิตฺตา กปฺปิตเกสมสฺสู ปญฺจหิ กามคุเณหิ สมปฺปิตา สมงฺคีภูตา ปริจาเรนฺติ.\csromanspacer{} อิธ ปนาหํ ภนฺเต ภิกฺขู ปสฺสามิ ยาวชีวํ อาปาณโกฏิกํ ปริปุณฺณํ ปริสุทฺธํ พฺรหฺมจริยํ จรนฺเต.\csromanspacer{} น โข ปนาหํ ภนฺเต อิโต พหิทฺธา อญฺญํ เอวํ ปริปุณฺณํ ปริสุทฺธํ พฺรหฺมจริยํ สมนุปสฺสามิ.\csromanspacer{} อยมฺปิ โข เม ภนฺเต ภควติ ธมฺมนฺวโย โหติ “สมฺมาสมฺพุทฺโธ ภควา, สฺวากฺขาโต ภควตา ธมฺโม, สุปฺปฏิปนฺโน ภควโต สาวกสํโฆ”ติ.}

\proseitem{368}{ปุน จปรํ ภนฺเต ราชาโนปิ ราชูหิ วิวทนฺติ, ขตฺติยาปิ ขตฺติเยหิ วิวทนฺติ, พฺราหฺมณาปิ พฺราหฺมเณหิ วิวทนฺติ, คหปตโยปิ}

\vspace{\tipitakaparskip}
\csromancenter{ธมฺมเจติยสุตฺต (89) คหปตีหิ วิวทนฺติ, มาตาปิ ปุตฺเตน วิวทติ, ปุตฺโตปิ มาตรา วิวทติ, ปิตาปิ ปุตฺเตน วิวทติ, ปุตฺโตปิ ปิตรา วิวทติ, ภาตาปิ ภคินิยา วิวทติ, ภคินีปิ ภาตรา วิวทติ, สหาโยปิ สหาเยน วิวทติ.\csromanspacer{} อิธ ปนาหํ ภนฺเต ภิกฺขู ปสฺสามิ สมคฺเค สมฺโมทมาเน อวิวทมาเน ขีโรทกีภูเต อญฺญมญฺญํ ปิยจกฺขูหิ สมฺปสฺสนฺเต วิหรนฺเต.\csromanspacer{} น โข ปนาหํ ภนฺเต อิโต พหิทฺธา อญฺญํ เอวํ สมคฺคํ ปริสํ สมนุปสฺสามิ.\csromanspacer{} อยมฺปิ โข เม ภนฺเต ภควติ ธมฺมนฺวโย โหติ “สมฺมาสมฺพุทฺโธ ภควา, สฺวากฺขาโต ภควตา ธมฺโม, สุปฺปฏิปนฺโน ภควโต สาวกสํโฆ”ติ.}
\proseitem{369}{\csromanfolio{323}ปุน จปราหํ ภนฺเต อาราเมน อารามํ อุยฺยาเนน อุยฺยานํ อนุจงฺกมามิ อนุวิจรามิ, โสหํ ตตฺถ ปสฺสามิ เอเก สมณพฺราหฺมเณ กิเส ลูเข ทุพฺพณฺเณ อุปฺปณฺฑุปฺปณฺฑุกชาเต ธมนิสนฺถตคตฺเต น วิย มญฺเญ จกฺขุํ พนฺธนฺเต ชนสฺส ทสฺสนาย.\csromanspacer{} ตสฺส มยฺหํ ภนฺเต เอตทโหสิ “อทฺธา อิเม อายสฺมนฺโต อนภิรตา วา พฺรหฺมจริยํ จรนฺติ.\csromanspacer{} อตฺถิ วา เตสํ กิญฺจิ ปาปํ กมฺมํ กตํ ปฏิจฺฉนฺนํ, ตถา หิ อิเม อายสฺมนฺโต กิสา ลูขา ทุพฺพณฺณา อุปฺปณฺฑุปฺปณฺฑุกชาตา ธมนิสนฺถตคตฺตา น วิย มญฺเญ จกฺขุํ พนฺธนฺติ ชนสฺส ทสฺสนายา”ติ.\csromanspacer{} ตฺยาหํ อุปสงฺกมิตฺวา เอวํ วทามิ “กิํ นุ โข ตุเมฺห อายสฺมนฺโต กิสา ลูขา ทุพฺพณฺณา อุปฺปณฺฑุปฺปณฺฑุกชาตา ธมญิสนฺถตคตฺตา น วิย มญฺเญ จกฺขุํ พนฺธถ ชนสฺส ทสฺสนายา”ติ.\csromanspacer{} เต เอวมาหํสุ “พนฺธุกโรโค โน\footnote{ปณฺฑุกโรคิโน (ก)} มหาราชา”ติ.\csromanspacer{} อิธ ปนาหํ ภนฺเต ภิกฺขู ปสฺสามิ หฏฺฐปหฏฺเฐ อุทคฺคุทคฺเค อภิรตรูเป ปีณินฺทฺริเย\footnote{ปีณิตินฺทฺริเย (สี, อิ)} อปฺโปสฺสุกฺเก ปนฺนโลเม ปรทตฺตวุตฺเต มิคภูเตน เจตสา วิหรนฺเต.\csromanspacer{} ตสฺส มยฺหํ ภนฺเต เอตทโหสิ “อทฺโธ อิเม อายสฺมนฺโต ตสฺส ภควโต สาสเน อุฬารํ ปุพฺเพนาปรํ วิเสสํ ชานนฺติ, ตถา หิ อิเม อายสฺมนฺโต หฏฺฐปหฏฺฐา อุทคฺคุทคฺคา อภิรตรูปา ปีณินฺทฺริยา อปฺโปสฺสุกฺกา ปนฺนโลมา ปรทตฺตวุตฺตา มิคภูเตน เจตสา วิหรนฺตี”ติ.\csromanspacer{} อยมฺปิ โข เม ภนฺเต ภควติ ธมฺมนฺวโย โหติ “สมฺมาสมฺพุทฺโธ ภควา, สฺวากฺขาโต ภควตา ธมฺโม, สุปฺปฏิปนฺโน ภควโต สาวกสํโฆ”ติ.}

\proseitem{370}{ปุน จปราหํ ภนฺเต ราชา ขตฺติโย มุทฺธาวสิตฺโต ปโหมิ ฆาเตตายํ วา ฆาเตตุํ ชาเปตายํ วา ชาเปตุํ ปพฺพาเชตายํ วา \csromanfolio{324}ปพฺพาเชตุํ, ตสฺส มยฺหํ ภนฺเต อฑฺฑกรเณ นิสินฺนสฺส อนฺตรนฺตรา กถํ โอปาเตนฺติ.\csromanspacer{} โสหํ น ลภามิ, มา เม โภนฺโต อฑฺฑกรเณ นิสินฺนสฺส อนฺตรนฺตรา กถํ โอปาเตถ\footnote{โอปาเตนฺตุ (สี) อุปริเสลสุตฺเต ปน “โอปาเตถา”ติเยว ทิสฺสติ.}.\csromanspacer{} กถาปริโยสานํ เม โภนฺโต อาคเมนฺตูติ.\csromanspacer{} ตสฺส มยฺหํ ภนฺเต อนฺตรนฺตรา กถํ โอปาเตนฺติ.\csromanspacer{} อิธ ปนาหํ ภนฺเต ภิกฺขู ปสฺสามิ, ยสฺมิํ สมเย ภควา อเนกสตาย ปริสาย ธมฺมํ เทเสติ, เนว ตสฺมิํ สมเย ภควโต สาวกานํ ขิปิตสทฺโท วา โหติ อุกฺกาสิตสทฺโท วา.\csromanspacer{} ภูตปุพฺพํ ภนฺเต ภควา อเนกสตาย ปริสาย ธมฺมํ เทเสติ, ตตฺรญฺญตโร ภควโต สาวโก อุกฺกาสิ, ตเมนํ อญฺญตโร สพฺรหฺมจารี ชณฺณุเกน ฆฏฺเฏสิ “อปฺปสทฺโท อายสฺมา โหตุ, มายสฺมา สทฺทมกาสิ, สตฺถา โน ภควา ธมฺมํ เทเสตี”ติ.\csromanspacer{} ตสฺส มยฺหํ ภนฺเต เอตทโหสิ “อจฺฉริยํ วต โภ, อพฺภุตํ วต โภ, อทณฺเฑน วต กิร โภ อสตฺเตน เอวํ สุวินีตา ปริสา ภวิสฺสตี”ติ, น โข ปนาหํ ภนฺเต อิโต พหิทฺธา อญฺญํ เอวํ สุวินีตํ ปริสํ สมนุปสฺสามิ, อยมฺปิ โข เม ภนฺเต ภควติ ธมฺมนฺวโย โหติ “สมฺมาสมฺพุทฺโธ ภควา, สฺวากฺขาโต ภควตา ธมฺโม, สุปฺปฏิปนฺโน ภควโต สาวกสํโฆ”ติ.}

\proseitem{371}{ปุน จปราหํ ภนฺเต ปสฺสามิ อิเมกจฺเจ ขตฺติยปณฺฑิเต นิปุเณ กตปรปฺปวาเท วาลเวธิรูเป, เต ภินฺทนฺตา\footnote{โวภินฺทนฺตา (สี)} มญฺเญ จรนฺติ ปญฺญาคเตน ทิฏฺฐิคตานิ.\csromanspacer{} เต สุณนฺติ “สมโณ ขลุ โภ โคตโม อมุกํ นาม คามํ วา นิคมํ วา โอสริสฺสตี”ติ.\csromanspacer{} เต ปญฺหํ อภิสงฺขโรนฺติ “อิมํ มยํ ปญฺหํ สมณํ โคตมํ อุปสงฺกมิตฺวา ปุจฺฉิสฺสาม, เอวํ เจ โน ปุฏฺโฐ เอวํ พฺยากริสฺสติ, เอวมสฺส มยํ วาทํ อาโรเปสฺสาม.\csromanspacer{} เอวํ เจปิ โน ปุฏฺโฐ เอวํ พฺยากริสฺสติ, เอวมฺปิสฺส มยํ วาทํ อาโรเปสฺสามา”ติ.\csromanspacer{} เต สุณนฺติ “สมโณ ขลุ โภ โคตโม อมุกํ นาม คามํ วา นิคมํ วา โอสโฏ”ติ, เต เยน ภควา เตนุปสงฺกมนฺติ.\csromanspacer{} เต ภควา ธมฺมิยา กถาย สนฺทสฺเสติ สมาทเปติ สมุตฺเตเชติ สมฺปหํเสติ.\csromanspacer{} เต ภควตา ธมฺมิยา กถาย สนฺทสฺสิตา สมาทปิตา สมุตฺเตชิตา สมฺปหํสิตา น เจว ภควนฺตํ ปญฺหํ ปุจฺฉนฺติ, กุโต วาทํ อาโรเปสฺสนฺติ, อญฺญทตฺถุ ภควโต สาวกา สมฺปชฺชนฺติ.\csromanspacer{} อยมฺปิ โข เม ภนฺเต ภควโต ธมฺมนฺวโย โหติ “สมฺมา-}

\vspace{\tipitakaparskip}
\csromancenter{ธมฺมเจติยสุตฺต (89) สมฺพุทฺโธ ภควา, สฺวากฺขาโต ภควตา ธมฺโม, สุปฺปฏิปนฺโน ภควโต สาวกสํโฆ”ติ.}
\proseitem{372}{\csromanfolio{325}ปุน จปราหํ ภนฺเต ปสฺสามิ อิเธกจฺเจ พฺราหฺมณปณฺฑิเต ฯเปฯ คหปติปณฺฑิเต.\csromanspacer{} สมณปณฺฑิเต นิปุเณ กตปรปฺปวาเท วาลเวธิรูเป.\csromanspacer{} เต ภินฺทนฺตา มญฺเญ จรนฺติ ปญฺญาคเตน ทิฏฺฐิคตานิ.\csromanspacer{} เต สุณนฺติ “สมโณ ขลุ โภ โคตโม อมุกํ นาม คามํ วา นิคมํ วา โอสริสฺสตี”ติ.\csromanspacer{} เต ปญฺหํ อภิสงฺขโรนฺติ “อิมํ มยํ ปญฺหํ สมณํ โคตมํ อุปสงฺกมิตฺวา ปุจฺฉิสฺสาม, เอวํ เจ โน ปุฏฺโฐ เอวํ พฺยากริสฺสติ, เอวมสฺส มยํ วาทํ อาโรเปสฺสาม.\csromanspacer{} เอวํ เจปิ โน ปุฏฺโฐ เอวํ พฺยากริสฺสติ, เอวมฺปิสฺส มยํ วาทํ อาโรเปสฺสามา”ติ.\csromanspacer{} เต สุณนฺติ “สมโณ ขลุ โภ โคตโม อมุกํ นาม คามํ วานิคมํ วา โอสโฏ”ติ, เต เยน ภควา เตนุปสงฺกมนฺติ.\csromanspacer{} เต ภควา ธมฺมิยา กถาย สนฺทสฺเสติ สมาทเปติ สมุตฺเตเชติ สมฺปหํเสติ.\csromanspacer{} เต ภควตา ธมฺมิยา กถาย สนฺทสฺสิตา สมาทปิตา สมุตฺเตชิตา สมฺปหํสิตา น เจว ภควนฺตํ ปญฺหํ ปุจฺฉนฺติ, กุโต วาทํ อาโรเปสฺสนฺติ, อญฺญทตฺถุ ภควนฺตํเยว โอกาสํ ยาจนฺติ อคารสฺมา อนคาริยํ ปพฺพชฺชาย.\csromanspacer{} เต ภควา ปพฺพาเชติ, เต ตถาปพฺพชิตา สมานา เอกา วูปกฏฺฐา อปฺปมตฺตา อาตาปิโน ปหิตตฺตา วิหรนฺตา นจิรสฺเสว, ยสฺสตฺถาย กุลปุตฺตา สมฺมเทว อคารสฺมา อนคาริยํ ปพฺพชนฺติ.\csromanspacer{} ตทนุตฺตรํ พฺรหฺมจริยปริโยสานํ ทิฏฺเฐว ธมฺเม สยํ อภิญฺญา สจฺฉิกตฺวา อุปสมฺปชฺช วิหรนฺติ.\csromanspacer{} เต เอวมาหํสุ “มนํ วต โภ อนสฺสาม, มนํ วต โภ ปนสฺสาม.\csromanspacer{} มยํ หิ ปุพฺเพ อสฺสมณาว สมานา สมณามฺหาติ ปฏิชานิมฺหา, อพฺราหฺมณาว สมานา พฺราหฺมณมฺหาติ ปฏิชานิมฺหา, อนรหนฺโตว สมานา อรหนฺตามฺหาติ ปฏิชานิมฺหา, อิทานิ โขมฺห สมณา, อิทานิ โขมฺห พฺราหฺมณา, อิทานิ โขมฺห อรหนฺโต”ติ.\csromanspacer{} อยมฺปิ โข เม ภนฺเต ภควติ ธมฺมนฺวโย โหติ “สมฺมาสมฺพุทฺโธ ภควา, สฺวากฺขาโต ภควตา ธมฺโม, สุปฺปฏิปนฺโน ภควโต สาวกสํโฆ”ติ.}

\proseitem{373}{ปุน จปราหํ ภนฺเต อิเม อิสิทตฺตปุราณา ถปตโย มมภตฺตา มมยานา, อหํ เนสํ ชีวิกาย\footnote{ชีวิตสฺส(สี), ชีวิกํ (สี-ฏฺฐ), ชีวิตํ (สฺยา, กํ, อิ, ก)} ทาตา, ยสสฺส อาหตฺตา. \csromanfolio{326}อถ จ ปน โน ตถา มยิ นิปจฺจการํ กโรนฺติ, ยถา ภควติ.\csromanspacer{} ภูตปุพฺพาหํ ภนฺเต เสนํ อพฺภุยฺยาโต สมาโน อิเม จ อิสิทตฺตปุราณา ถปตโย วีมํสมาโน อญฺญตรสฺมิํ สมฺพาเธ อาวสเถ วาสํ อุปคจฺฉิํ.\csromanspacer{} อถ โข ภนฺเต อิเม อิสิทตฺตปุราณา ถปตโย พหุเทว รตฺติํ ธมฺมิยา กถาย วีตินาเมตฺวา ยโต อโหสิ ภควา\footnote{อสฺโสสุํ โข ภควนฺตํ (สี, สฺยา, กํ, อิ)}, ตโต สีสํ กตฺวา มํ ปาทโต กริตฺวา นิปชฺชิํสุ.\csromanspacer{} ตสฺส มยฺหํ ภนฺเต เอตทโหสิ “อจฺฉริยํ วต โภ, อพฺภุตํ วต โภ, อิเม อิสิทตฺตปุราณา ถปตโย มมภตฺตา มมยานา, อหํ เนสํ ชีวิกาย ทาตา, ยสสฺส อาหตฺตา.\csromanspacer{} อถ จ ปน โน ตถา มยิ นิปจฺจการํ กโรนฺติ, ยถา ภควติ.\csromanspacer{} อทฺธา อิเม อายสฺมนฺโต ตสฺส ภควโต สาสเน อุฬารํ ปุพฺเพนาปรํ วิเสสํ ชานนฺตี”ติ.\csromanspacer{} อยมฺปิ โข เม ภนฺเต ภควติ ธมฺมนฺวโย โหติ “สมฺมาสมฺพุทฺโธ ภควา, สฺวากฺขาโต ภควตา ธมฺโม, สุปฺปฏิปนฺโน ภควโต สาวกสํโฆ”ติ.}

\proseitem{374}{ปุน จปรํ ภนฺเต ภควาปิ ขตฺติโย, อหมฺปิ ขตฺติโย, ภควาปิ โกสโล, อหมฺปิ โกสโล, ภควาปิ อาสีติโก, อหมฺปิ อาสีติโก.\csromanspacer{} ยมฺปิ ภนฺเต ภควาปิ ขตฺติโย, อหมฺปิ ขตฺติโย, ภควาปิ โกสโล, อหมฺปิ โกสโล, ภควาปิ อาสีติโก, อหมฺปิ อาสีติโก, อิมินาวารหาเมวาหํ\footnote{อิมินาปาหํ (ก)} ภนฺเต ภควติ ปรมนิปจฺจการํ กาตุํ มิตฺตูปหารํ อุปทํเสตุํ.\csromanspacer{} หนฺท จ ทานิ มยํ ภนฺเต คจฺฉาม พหุกิจฺจา มยํ พหุกรณียาติ.\csromanspacer{} ยสฺสทานิ ตฺวํ มหาราช กาลํ มญฺญสีติ.\csromanspacer{} อถ โข ราชา ปเสนทิ โกสโล อุฏฺฐายาสนา ภควนฺตํ อภิวาเทตฺวา ปทกฺขิณํ กตฺวา ปกฺกามิ.\csromanspacer{} อถ โข ภควา อจิรปกฺกนฺตสฺส รญฺโญ ปเสนทิสฺส โกสลสฺส ภิกฺขู อามนฺเตสิ “เอโส ภิกฺขเว ราชา ปเสนทิ โกสโล ธมฺมเจติยานิ ภาสิตฺวา อุฏฺฐายาสนา ปกฺกนฺโต, อุคฺคณฺหถ ภิกฺขเว ธมฺมเจติยานิ, ปริยาปุณาถ ภิกฺขเว ธมฺมเจติยานิ, ธาเรถ ภิกฺขเว ธมฺมเจติยานิ, อตฺถสํหิตานิ ภิกฺขเว ธมฺมเจติยานิ อาทิพฺรหฺมจริยกานี”ติ.}

\csromanmatikaanchor{mtk.46}
\csromantocmarkref{subsection}{9. ธมฺมเจติยสุตฺต}{mtk.46}
\bookmark[dest={mtk.46},level=2]{9. ธมฺมเจติยสุตฺต}
\prose{อิทมโวจ ภควา.\csromanspacer{} อตฺตมนา เต ภิกฺขู ภควโต ภาสิตํ อภินนฺทุนฺติ.}

\nitthitam{ธมฺมเจติยสุตฺตํ นิฏฺฐิตํ นวมํ.}
\csromanheader{2}{10. กณฺณกตฺถลสุตฺต (90) \textbf{10. กณฺณกตฺถลสุตฺต}}
\proseitem{375}{\csromanfolio{327}เอวํ เม สุตํ\csromandash{}เอกํ สมยํ ภควา อุรุญฺญายํ\footnote{อุชุญฺญายํ (สี, อิ), อุทญฺญายํ (สฺยา, กํ)} วิหรติ กณฺณกตฺถเล มิคทาเย.\csromanspacer{} เตน โข ปน สมเยน ราชา ปเสนทิ โกสโล อุรุญฺญํ อนุปฺปตฺโต โหติ เกนจิเทว กรณีเยน.\csromanspacer{} อถ โข ราชา ปเสนทิ โกสโล อญฺญตรํ ปุริสํ อามนฺเตสิ “เอหิ ตฺวํ อมฺโภ ปุริส เยน ภควา เตนุปสงฺกม, อุปสงฺกมิตฺวา มม วจเนน ภควโต ปาเท สิรสา วนฺทาหิ, อปฺปาพาธํ อปฺปาตงฺกํ ลหุฏฺฐานํ พลํ ผาสุวิหารํ ปุจฺฉ ‘ราชา ภนฺเต ปเสนทิ โกสโล ภควโต ปาเท สิรสา วนฺทติ, อปฺปาพาธํ อปฺปาตงฺกํ ลหุฏฺฐานํ พลํ ผาสุวิหารํ ปุจฺฉตี’ติ.\csromanspacer{} เอวญฺจ วเทหิ ‘อชฺช กิร ภนฺเต ราชา ปเสนทิ โกสโล ปจฺฉาภตฺตํ ภุตฺตปาตราโส ภควนฺตํ ทสฺสนาย อุปสงฺกมิสฺสตี’ติ”. “เอวํ เทวา”ติ โข โส ปุริโส รญฺโญ ปเสนทิสฺส โกสลสฺส ปฏิสฺสุตฺวา เยน ภควา เตนุปสงฺกมิ, อุปสงฺกมิตฺวา ภควนฺตํ อภิวาเทตฺวา เอกมนฺตํ นิสีทิ, เอกมนฺตํ นิสินฺโน โข โส ปุริโส ภควนฺตํ เอตทโวจ “ราชา ภนฺเต ปเสนทิ โกสโล ภควโต ปาเท สิรสา วนฺทติ, อปฺปาพาธํ อปฺปาตงฺกํ ลหุฏฺฐานํ พลํ ผสุวิหารํ ปุจฺฉติ, เอวญฺจ วเทติ ‘อชฺช กิร ภนฺเต ราชา ปเสนทิ โกสโล ปจฺฉาภตฺตํ ภุตฺตปาตราโส ภควนฺตํ ทสฺสนาย อุปสงฺกมิสฺสตี’ติ”.\csromanspacer{} อสฺโสสุํ โข โสมา จ ภคินี สกุลา จ ภคินี “อชฺช กิร ราชา ปเสนทิ โกสโล ปจฺฉาภตฺตํ ภุตฺตปาตราโส ภควนฺตํ ทสฺสนาย อุปสงฺกมิสฺสตี”ติ.\csromanspacer{} อถ โข โสมา จ ภคินี สกุลา จ ภคินี ราชานํ ปเสนทิํ โกสลํ ภตฺตาภิหาเร อุปสงฺกมิตฺวา เอตทโวจุํ “เตน หิ มหาราช อมฺหากมฺปิ วจเนน ภควโต ปาเท สิรสา วนฺทาหิ, อปฺปาพาธํ อปฺปาตงฺกํ ลหุฏฺฐานํ พลํ ผาสุวิหารํ ปุจฺฉ, โสมา จ ภนฺเต ภคินี สกุลา จ ภคินี ภควโต ปาเท สิรสา วนฺทติ, อปฺปาพาธํ อปฺปาตงฺกํ ลหุฏฺฐานํ พลํ ผาสุวิหารํ ปุจฺฉตี”ติ.}

\proseitem{376}{อถ โข ราชา ปเสนทิ โกสโล ปจฺฉาภตฺตํ ภุตฺตปาตราโส เยน ภควา เตนุปสงฺกมิ, อุปสงฺกมิตฺวา ภควนฺตํ อภิวาเทตฺวา เอกมนฺตํ นิสีทิ, เอกมนฺตํ นิสินฺโน โข ราชา ปเสนทิ \csromanfolio{328}โกสโล ภควนฺตํ เอตทโวจ “โสมา จ ภนฺเต ภคินี สกุลา จ ภนฺเต ภคินี ภควโต ปาเท สิรสา วนฺทติ\footnote{วนฺทนฺติ (สี, สฺยา, กํ, อิ)}, อปฺปาพาธํ อปฺปาตงฺกํ ลหุฏฺฐานํ พลํ ผาสุวิหารํ ปุจฺฉตี”ติ\footnote{ปุจฺฉนฺตีติ (สี, สฺยา, กํ, อิ)}.\csromanspacer{} กิํ ปน มหาราช โสมา จ ภคินี สกุลา จ ภิคินี อญฺญํ ทูตํ นาลตฺถุนฺติ.\csromanspacer{} อสฺโสสุํ โข ภนฺเต โสมา จ ภคินี สกุลา จ ภคินี “อชฺช กิร ราชา ปเสนทิ โกสโล ปจฺฉาภตฺตํ ภุตฺตปาตราโส ภควนฺตํ ทสฺสนาย อุปสงฺกมิสฺสตี”ติ.\csromanspacer{} อถ โข ภนฺเต โสมา จ ภคินี สกุลา จ ภคินี มํ ภตฺตาภิหาเร อุปสงฺกมิตฺวา เอตทโวจุํ “เตน หิ มหาราช อมฺหากมฺปิ วจเนน ภควโต ปาเท สิรสา วนฺทาหิ, อปฺปาพาธํ อปฺปาตงฺกํ ลหุฏฺฐานํ พลํ ผาสุวิหารํ ปุจฺฉ, โสมา จ ภคินี สกุลา จ ภคินี ภควโต ปาเท สิรสา วนฺทติ, อปฺปาพาธํ อปฺปาตงฺกํ ลหุฏฺฐานํ พลํ ผาสุวิหารํ ปุจฺฉตี”ติ.\csromanspacer{} สุขินิโย โหนฺตุ ตา มหาราช โสมา จ ภคินี สกุลา จ ภคินีติ.}

\proseitem{377}{อถ โข ราชา ปเสนทิ โกสโล ภควนฺตํ เอตทโวจ “สุตํ เมตํ ภนฺเต, สมโณ โคตโม เอวมาห ‘นตฺถิ โส สมโณ วา พฺราหฺมโณ วา, โย สพฺพญฺญู สพฺพทสฺสาวี อปริเสสํ ญาณทสฺสนํ ปฏิชานิสฺสติ, เนตํ ฐานํ วิชฺชตี’ติ, เย เต ภนฺเต เอวมาหํสุ ‘สมโณ โคตโม เอวมาห, นตฺถิ โส สมโณ วา พฺราหฺมโณ วา โย สพฺพญฺญู สพฺพทสฺสาวี อปริเสสํ ญาณทสฺสนํ ปฏิชานิสฺสติ, เนตํ ฐานํ วิชฺชตี’ติ, กจฺจิ เต ภนฺเต ภควโต วุตฺตวาทิโน, น จ ภควนฺตํ อภูเตน อพฺภาจิกฺขนฺติ, ธมฺมสฺส จานุธมฺมํ พฺยากโรนฺติ, น จ โกจิ สหธมฺมิโก วาทานุวาโท คารยฺหํ ฐานํ อาคจฺฉตี”ติ.\csromanspacer{} เย เต มหาราช เอวมาหํสุ “สมโณ โคตโม เอวมาห ‘นตฺถิ โส สมโณ วา พฺราหฺมโณ วา, โย สพฺพญฺญู สพฺพทสฺสาวี อปริเสสํ ญาณทสฺสนํ ปฏิชานิสฺสติ, เนตํ ฐานํ วิชฺชตี’ติ, น เม เต วุตฺตวาทิโน อพฺภาจิกฺขนฺติ จ ปน มํ เต อสตา อภูเตนา”ติ.}

\proseitem{378}{อถ โข ราชา ปเสนทิ โกสโล วิฏฏูภํ เสนาปติํ อามนฺเตสิ “โก นุ โข เสนาปติ อิมํ กถาวตฺถุํ ราชนฺเตปุเร อพฺภุทาหาสี”ติ.\csromanspacer{} สญฺชโย มหาราช พฺราหฺมโณ อากาสโคตฺโตติ.\csromanspacer{} อถ โข ราชา ปเสนทิ โกสโล อญฺญตรํ ปุริสํ อามนฺเตสิ “เอหิ}

\vspace{\tipitakaparskip}
\csromancenter{กณฺณกตฺถลสุตฺต (90) ตฺวํ อมฺโภ ปุริส มม วจเนน สญฺชยํ พฺราหฺมณํ อากาสโคตฺตํ อามนฺเตหิ ‘ราชา ตํ ภนฺเต ปเสนทิ โกสโล อามนฺเตตี’ติ”. “เอวํ เทวา”ติ โข โส ปุริโส รญฺโญ ปเสนทิสฺส โกสลสฺส ปฏิสฺสุตฺวา เยน สญฺชโย พฺราหฺมโณ อากาสโคตฺโต เตนุปสงฺกมิ, อุปสํกมิตฺวา สญฺชยํ พฺราหฺมณํ อากาสโคตฺตํ เอตทโวจ “ราชา ตํ ภนฺเต ปเสนทิ โกสโล อามนฺเตตี”ติ.\csromanspacer{} อถ โข ราชา ปเสนทิ โกสโล ภควนฺตํ เอตทโวจ “สิยา นุ โข ภนฺเต ภควตา อญฺญเทว กิญฺจิ สนฺธาย ภาสิตํ, ตญฺจ ชโน อญฺญถาปิ ปจฺจาคจฺเฉยฺย\footnote{ปจฺจาคจฺเฉยฺยาติ, อภิชานามิ มหาราช วาจํ ภาสิตาติ (สี)}, ยถา กถํ ปน ภนฺเต ภควา อภิชานาติ วาจํ ภาสิตา”ติ.\csromanspacer{} เอวํ โข อหํ มหาราช อภิชานามิ วาจํ ภาสิตา, นตฺถิ โส สมโณ วา พฺราหฺมโณ วา, โย สกิเทว สพฺพํ ญสฺสติ สพฺพํ ทกฺขิติ, เนตํ ฐานํ วิชฺชตีติ.\csromanspacer{} เหตุรูปํ ภนฺเต ภควา อาห, สเหตุรูปํ ภนฺเต ภควา อาห, นตฺถิ โส สมโณ วา พฺราหมโณ วา, โย สกิเทว สพฺพํ ญสฺสติ สพฺพํ ทกฺขิติ, เนตํ ฐานํ วิชฺชตีติ.\csromanspacer{} จตฺตาโรเม ภนฺเต วณฺณา ขตฺติยา พฺราหฺมณา เวสฺสา สุทฺทา, อิเมสํ นุ โข ภนฺเต จตุนฺนํ วณฺณานํ สิยา วิเสโส, สิยา นานากรณนฺติ.\csromanspacer{} จตฺตาโรเม มหาราช วณฺณา ขตฺติยา พฺราหฺมณา เวสฺสา สุทฺทา, อิเมสํ โข มหาราช จตุนฺนํ วณฺณานํ เทฺว วณฺณา อคฺคมกฺขายนฺติ ขตฺติยา จ พฺราหฺมณา จ, ยทิทํ อภิวาทนปจฺจุฏฺฐาน-อญฺชลิกมฺมสามีจิกมฺมานีติ\footnote{สามิจิกมฺมานนฺติ (สี)}.\csromanspacer{} นาหํ ภนฺเต ภควนฺตํ ทิฏฺฐธมฺมิกํ ปุจฺฉามิ, สมฺปรายิกาหํ ภนฺเต ภควนฺตํ ปุจฺฉามิ, จตฺตาโรเม ภนฺเต วณฺณา ขตฺติยา พฺราหฺมณา เวสฺสา สุทฺทา, อิเมสํ นุ โข ภนฺเต จตุนฺนํ วณฺณานํ สิยา วิเสโส, สิยา นานากรณนฺติ.}
\proseitem{379}{\csromanfolio{329}ปญฺจิมานิ มหาราช ปธานิยงฺคานิ.\csromanspacer{} กตมานิ ปญฺจ, อิธ มหาราช ภิกฺขุ สทฺโธ โหติ, สทฺทหติ ตถาคตสฺส โพธิํ “อิติปิ โส ภควา อรหํ สมฺมาสมฺพุทฺโธ วิชฺชาจรณสมฺปนฺโน สุคโต โลกวิทู อนุตฺตโร ปุริสทมฺมสารถิ สตฺถา เทวมนุสฺสานํ พุทฺโธ ภควา”ติ.\csromanspacer{} อปฺปาพาโธ โหติ อปฺปาตงฺโก, สมเวปากินิยา คหณิยา สมนฺนาคโต นาติสีตาย นาจฺจุณฺหาย มชฺฌิมาย ปธานกฺขมาย.\csromanspacer{} อสโฐ โหติ อมายาวี ยถาภูตํ อตฺตานํ อาวิกตฺตา สตฺถริ วา วิญฺญูสุ วา สพฺรหฺมจารีสุ.\csromanspacer{} อารทฺธวีริโย วิหรติ อกุสลานํ \csromanfolio{330}ธมฺมานํ ปหานาย กุสลานํ ธมฺมานํ อุปสมฺปทาย ถามวา ทฬฺหปรกฺกโม อนิกฺขิตฺตธุโร กุสเลสุ ธมฺเมสุ.\csromanspacer{} ปญฺญวา โหติ อุทยตฺถคามินิยา ปญฺญาย สมนฺนาคโต อริยาย นิพฺเพธิกาย สมฺมาทุกฺขกฺขยคามินิยา.\csromanspacer{} อิมานิ โข มหาราช ปญฺจ ปธานิยงฺคานิ.\csromanspacer{} จตฺตาโรเม มหาราช วณฺณา ขตฺติยา พฺราหฺมณา เวสฺสา สุทฺทา, เต จสฺสุ อิเมหิ ปญฺจหิ ปธานิยงฺเคหิ สมนฺนาคตา.\csromanspacer{} เอตฺถ ปน เนสํ อสฺส ทีฆรตฺตํ หิตาย สุขายาติ.\csromanspacer{} จตฺตาโรเม ภนฺเต วณฺณา ขตฺติยา พฺราหฺมณา เวสฺสา สุทฺทา, เต จสฺสุ อิเมหิ ปญฺจหิ ปธานิยงฺเคหิ สมนฺนาคตา.\csromanspacer{} เอตฺถ ปน เนสํ ภนฺเต สิยา วิเสโส, สิยา นานากรณนฺติ.\csromanspacer{} เอตฺถ โข เนสาหํ มหาราช ปธานเวมตฺตตํ วทามิ, เสยฺยถาปิสฺสุ มหาราช เทฺว หตฺถิทมฺมา วา อสฺสทมฺมา วา โคทมฺมา วา สุทนฺตา สุวินีตา, เทฺว หตฺถิทมฺมา วา อสฺสทมฺมา วา โคทมฺมา วา อทนฺตา อวินีตา.\csromanspacer{} ตํ กิํ มญฺญสิ มหาราช, เย เต เทฺว หตฺถิทมฺมา วา อสฺสทมฺมา วา โคทมฺมา วา สุทนฺตา สุวินีตา, อปิ นุ เต ทนฺตาว ทนฺตการณํ คจฺเฉยฺยุํ, ทนฺตาว ทนฺตภูมิํ สมฺปาปุเณยฺยุนฺติ.\csromanspacer{} เอวํ ภนฺเต.\csromanspacer{} เย ปน เต เทฺว หตฺถิทมฺมา วา อสฺสทมฺมา วา โคทมฺมา วา อทนฺตา อวินีตา, อปิ นุ เต อทนฺตาว ทนฺตการณํ คจฺเฉยฺยุํ, อทนฺตาว ทนฺตภูมิํ สมฺปาปุเณยฺยุํ.\csromanspacer{} เสยฺยถาปิ เต เทฺว หตฺถิทมฺมา วา อสฺสทมฺมา วา โคทมฺมา วา สุทนฺตา สุวินีตาติ.\csromanspacer{} โน เหตํ ภนฺเต.\csromanspacer{} เอวเมว โข มหาราช ยํ ตํ สทฺเธน ปตฺตพฺพํ อปฺปาพาเธน อสเฐน อมายาวินา อารทฺธวีริเยน ปญฺญวตา.\csromanspacer{} ตํ วต\footnote{ตํ ตถา โส (ก)} อสฺสทฺโธ พหฺวาพาโธ สโฐ มายาวี กุสีโต ทุปฺปญฺโญ ปาปุณิสฺสตีติ, เนตํ ฐานํ วิชฺชตีติ.}

\proseitem{380}{เหตุรูปํ ภนฺเต ภควา อาห, สเหตุรูปํ ภนฺเต ภควา อาห.\csromanspacer{} จตฺตาโรเม ภนฺเต วณฺณา ขตฺติยา พฺราหฺมณา เวสฺสา สุทฺทา, เต จสฺสุ อิเมหิ ปญฺจหิ ปธานิยงฺเคหิ สมนฺนาคตา, เต จสฺสุ สมฺมปฺปธานา.\csromanspacer{} เอตฺถ ปน เนสํ ภนฺเต สิยา วิเสโส, สิยา นานากรณนฺติ.\csromanspacer{} เอตฺถ โข\footnote{เอตฺถ โข ปน (สี)} เนสาหํ มหาราช น กิญฺจิ นานากรณํ วทามิ, ยทิทํ วิมุตฺติยา วิมุตฺติํ.\csromanspacer{} เสยฺยถาปิ มหาราช ปุริโส สุกฺขํ สากกฏฺฐํ อาทาย อคฺคิํ อภินิพฺพตฺเตยฺย, เตโช ปาตุกเรยฺย.\csromanspacer{} อถาปโร ปุริโส สุกฺขํ สาลกฏฺฐํ อาทาย อคฺคิํ อภินิพฺพตฺเตยฺย, เตโช}

\vspace{\tipitakaparskip}
\csromancenter{กณฺณกตฺถลสุตฺต (90) ปาตุกเรยฺย.\csromanspacer{} อถาปโร ปุริโส สุกฺขํ อมฺพกฏฺฐํ อาทาย อคฺคิํ อภินิพฺพตฺเตยฺย, เตโช ปาตุกเรยฺย.\csromanspacer{} อถาปโร ปุริโส สุกฺขํ อุทุมฺพรกฏฺฐํ อาทาย อคฺคิํ อภินิพฺพตฺเตยฺย, เตโช ปาตุกเรยฺย.\csromanspacer{} ตํ กิํ มญฺญสิ มหาราช, สิยา นุ โข เตสํ อคฺคีนํ นานาทารุโต อภินิพฺพตฺตานํ กิญฺจิ นานากรณํ, อจฺจิยา วา อจฺจิํ วณฺเณน วา วณฺณํ อาภาย วา อาภนฺติ.\csromanspacer{} โน เหตํ ภนฺเต.\csromanspacer{} เอวเมว โข มหาราช ยํ ตํ เตชํ วีริยา นิมฺมถิตํ ปธานาภินิพฺพตฺตํ\footnote{วิริยํ นิปฺผรติ, ตํ ปจฺฉาภินิพฺพตฺตํ (สี)}.\csromanspacer{} นาหํ ตตฺถ กิญฺจิ นานากรณํ วทามิ.\csromanspacer{} ยทิทํ วิมุตฺติยา วิมุตฺตินฺติ.\csromanspacer{} เหตุรูปํ ภนฺเต ภควา อาห, สเหตุรูปํ ภนฺเต ภควา อาห.\csromanspacer{} กิํ ปน ภนฺเต อตฺถิ เทวาติ.\csromanspacer{} กิํ ปน ตฺวํ มหาราช เอวํ วเทสิ, กิํ ปน ภนฺเต อตฺถิ เทวาติ.\csromanspacer{} ยทิ วา เต ภนฺเต เทวา อาคนฺตาโร อิตฺถตฺตํ, ยทิ วา อนาคนฺตาโร อิตฺถตฺตํ.\csromanspacer{} เย เต มหาราช เทวา สพฺยาพชฺฌา, เต เทวา อาคนฺตาโร อิตฺถตฺตํ.\csromanspacer{} เย เต เทวา อพฺยาพชฺฌา, เต เทวา อนาคนฺตาโร อิตฺถตฺตนฺติ.}
\proseitem{381}{\csromanfolio{331}เอวํ วุตฺเต วิฏฏูโภ เสนาปติ ภควนฺตํ เอตทโวจ “เย เต ภนฺเต เทวา สพฺยาพชฺฌา อาคนฺตาโร อิตฺถตฺตํ, เต เทวา.\csromanspacer{} เย เต เทวา อพฺยาพชฺฌา อนาคนฺตาโร อิตฺถตฺตํ, เต เทเว ตมฺหา ฐานา จาเวสฺสนฺติ วา ปพฺพาเชสฺสนฺติ วา”ติ.}

\prose{อถ โข อายสฺมโต อานนฺทสฺส เอตทโหสิ “อยํ โข วิฏฏูโภ เสนาปติ รญฺโญ ปเสนทิสฺส โกสลสฺส ปุตฺโต, อหํ ภควโต ปุตฺโต.\csromanspacer{} อยํ โข กาโล ยํ ปุตฺโต ปุตฺเตน มนฺเตยฺยา”ติ.\csromanspacer{} อถ โข อายสฺมา อานนฺโท วิฏฏูภํ เสนาปติํ อามนฺเตสิ “เตน หิ เสนาปติ ตํเยเวตฺถ ปฏิปุจฺฉิสฺสามิ, ยถา เต ขเมยฺย, ตถา นํ พฺยากเรยฺยาสิ.\csromanspacer{} ตํ กิํ มญฺญสิ เสนาปติ, ยาวตา รญฺโญ ปเสนทิสฺส โกสลสฺส วิชิตํ, ยตฺถ จ ราชา ปเสนทิ โกสโล อิสฺสริยาธิปจฺจํ รชฺชํ กาเรติ.\csromanspacer{} ปโหติ ตตฺถ ราชา ปเสนทิ โกสโล สมณํ วา พฺราหฺมณํ วา ปุญฺญวนฺตํ วา อปุญฺญวนฺตํ วา พฺรหฺมจริยวนฺตํ วา อพฺรหฺมจริยวนฺตํ วา ตมฺหา ฐานา จาเวตุํ วา ปพฺพาเชตุํ วา”ติ.\csromanspacer{} ยาวตา โภ รญฺโญ ปเสนทิสฺส โกสลสฺส วิชิตํ, ยตฺถ จ ราชา ปเสนทิ โกสโล อิสฺสริยาธิปจฺจํ รชฺชํ กาเรติ, ปโหติ ตตฺถ ราชา ปเสนทิ \csromanfolio{332}โกสโล สมณํ วา พฺราหฺมณํ วา ปุญฺญวนฺตํ วา อปุญฺญวนฺตํ วา พฺรหฺมจริยวนฺตํ วา อพฺรหฺมจริยวนฺตํ วา ตมฺหา ฐานา จาเวตุํ วา ปพฺพาเชตุํ วาติ.}

\prose{ตํ กิํ มญฺญสิ เสนาปติ, ยาวตา รญฺโญ ปเสนทิสฺส โกสลสฺส อวิชิตํ, ยตฺถ จ ราชา ปเสนทิ โกสโล น อิสฺสริยาธิปจฺจํ รชฺชํ กาเรติ, ตตฺถ ปโหติ ราชา ปเสนทิ โกสโล สมณํ วา พฺราหฺมณํ วา ปุญฺญวนฺตํ วา อปุญฺญวนฺตํ วา พฺรหฺมจริยวนฺตํ วา อพฺรหฺมจริยวนฺตํ วา ตมฺหา ฐานา จาเวตุํ วา ปพฺพาเชตุํ วาติ.\csromanspacer{} ยาวตา โภ รญฺโญ ปเสนทิสฺส โกสลสฺส อวิชิตํ, ยตฺถ จ ราชา ปเสนทิ โกสโล น อิสฺสริยาธิปจฺจํ รชฺชํ กาเรติ, น ตตฺถ ปโหติ ราชา ปเสนทิ โกสโล สมณํ วา พฺราหฺมณํ วา ปุญฺญวนฺตํ วา อปุญฺญวนฺตํ วา พฺรหฺมจริยวนฺตํ วา อพฺรหฺมจริยวนฺตํ วา ตมฺหา ฐานา จาเวตุํ วา ปพฺพาเชตุํ วาติ.}

\prose{ตํ กิํ มญฺญสิ เสนาปติ, สุตา เต เทวา ตาวติํสาติ.\csromanspacer{} เอวํ โภ สุตา เม เทวา ตาวติํสา.\csromanspacer{} อิธาปิ โภตา รญฺญา ปเสนทินา โกสเลน สุตา เทวา ตาวติํสาติ.\csromanspacer{} ตํ กิํ มญฺญสิ เสนาปติ, ปโหติ ราชา ปเสนทิ โกสโล เทเว ตาวติํเส ตมฺหา ฐานา จาเวตุํ วา ปพฺพาเชตุํ วาติ.\csromanspacer{} ทสฺสนมฺปิ โภ ราชา ปเสนทิ โกสโล เทเว ตาวติํเส นปฺปโหติ, กุโต ปน ตมฺหา ฐานา จาเวสฺสติ วา ปพฺพาเชสฺสติ วาติ.\csromanspacer{} เอวเมว โข เสนาปติ เย เต เทวา สพฺยาพชฺฌา อาคนฺตาโร อิตฺถตฺตํ, เต เทวา, เย เต เทวา อพฺยาพชฺฌา อนาคนฺตาโร อิตฺถตฺตํ, เต เทเว ทสฺสนายปิ นปฺปโหนฺติ, กุโต ปน ตมฺหา ฐานา จาเวสฺสนฺติ วา ปพฺพาเชสฺสนฺติ วาติ.}

\proseitem{382}{อถ โข ราชา ปเสนทิ โกสโล ภควนฺตํ เอตทโวจ “โกนาโม อยํ ภนฺเต ภิกฺขู”ติ.\csromanspacer{} อานนฺโท นาม มหาราชาติ.\csromanspacer{} อานนฺโท วต โภ, อานนฺทรูโป วต โภ, เหตุรูปํ ภนฺเต อายสฺมา อานนฺโท อาห, สเหตุรูปํ ภนฺเต อายสฺมา อานนฺโท อาห.\csromanspacer{} กิํ ปน ภนฺเต อตฺถิ พฺรหฺมาติ.\csromanspacer{} กิํ ปน ตฺวํ มหาราช เอวํ วเทสิ “กิํ ปน ภนฺเต อตฺถิ พฺรหฺมา”ติ.\csromanspacer{} ยทิ วา โส ภนฺเต พฺรหฺมา อาคนฺตา อิตฺถตฺตํ, ยทิ วา อนาคนฺตา อิตฺถตฺตนฺติ.\csromanspacer{} โย โส มหาราช พฺรหฺมา สพฺยาพชฺโฌ, โส พฺรหฺมา อาคนฺตา อิตฺถตฺตํ.\csromanspacer{} โย โส พฺรหฺม อพฺยาพชฺโฌ, โส พฺรหฺมา อนาคนฺตา อิตฺถตฺตนฺติ.\csromanspacer{} อถ โข อญฺญตโร ปุริโส ราชานํ}

\vspace{\tipitakaparskip}
\csromancenter{กณฺณกตฺถลสุตฺต (90) ปเสนทิํ โกสลํ เอตทโวจ “สญฺชโย มหาราช พฺราหฺมโณ อากาสโคตฺโต อาคโต”ติ.\csromanspacer{} อถ โข ราชา ปเสนทิ โกสโล สญฺชยํ พฺราหฺมณํ อากาสโคตฺตํ เอตทโวจ “โก นุ โข พฺราหฺมณ อิมํ กถาวตฺถุํ ราชนฺเตปุเร อพฺภุทาหาสี”ติ.\csromanspacer{} วิฏฏูโภ มหาราช เสนาปตีติ.\csromanspacer{} วิฏฏูโภ เสนาปติ เอวมาห “สญฺชโย มหาราช พฺราหฺมโณ อากาสโคตฺโต”ติ.\csromanspacer{} อถ โข อญฺญตโร ปุริโส ราชานํ ปเสนทิํ โกสลํ เอตทโวจ “ยานกาโล มหาราชา”ติ.}
\csromanmatikaanchor{mtk.47}
\csromanmatikaanchor{mtk.48}
\csromantocmarkref{subsection}{10. กณฺณกตฺถลสุตฺต}{mtk.47}
\bookmark[dest={mtk.47},level=2]{10. กณฺณกตฺถลสุตฺต}
\csromantocmarkref{subsection}{อุทฺทานคาถาโย}{mtk.48}
\bookmark[dest={mtk.48},level=2]{อุทฺทานคาถาโย}
\prosecont{\csromanfolio{333}อถ โข ราชา ปเสนทิ โกสโล ภควนฺตํ เอตทโวจ “สพฺพญฺญุตํ มยํ ภนฺเต ภควนฺตํ อปุจฺฉิมฺหา, สพฺพญฺญุตํ ภควา พฺยากาสิ, ตญฺจ ปนมฺหากํ รุจฺจติ เจว ขมติ จ, เตน จมฺหา อตฺตมนา.\csromanspacer{} จาตุวณฺณิสุทฺธิํ มยํ ภนฺเต ภควนฺตํ อปุจฺฉิมฺหา, จาตุวณฺณิสุทฺธิํ ภควา พฺยากาสิ, ตญฺจ ปนมฺหากํ รุจฺจติ เจว ขมติ จ, เตน จมฺหา อตฺตมนา.\csromanspacer{} อธิเทเว มยํ ภนฺเต ภควนฺตํ อปุจฺฉิมฺหา, อธิเทเว ภควา พฺยากาสิ, ตญฺจ ปนมฺหากํ รุจฺจติ เจว ขมติ จ, เตน จมฺหา อตฺตมนา.\csromanspacer{} อธิพฺรหฺมานํ มยํ ภนฺเต ภควนฺตํ อปุจฺฉิมฺหา, อธิพฺรหฺมานํ ภควา พฺยากาสิ, ตญฺจ ปนมฺหากํ รุจฺจติ เจว ขมติ จ, เตน จมฺหา อตฺตมนา.\csromanspacer{} ยํ ยเทว จ มยํ ภควนฺตํ อปุจฺฉิมฺหา, ตํ ตเทว ภควา พฺยากาสิ, ตญฺจ ปนมฺหากํ รุจฺจติ เจว ขมติ จ, เตน จมฺหา อตฺตมนา.\csromanspacer{} หนฺท จ ทานิ มยํ ภนฺเต คจฺฉาม พหุกิจฺจา มยํ พหุกรณียา”ติ.\csromanspacer{} ยสฺสทานิ ตฺวํ มหาราช กาลํ มญฺญสีติ.\csromanspacer{} อถ โข ราชา ปเสนทิ โกสโล ภควโต ภาสิตํ อภินนฺทิตฺวา อนุโมทิตฺวา อุฏฺฐายาสนา ภควนฺตํ อภิวาเทตฺวา ปทกฺขิณํ กตฺวา ปกฺกามีติ.}

\nitthitam{กณฺณกตฺถลสุตฺตํ นิฏฺฐิตํ ทสมํ.}
\nitthitam{\textbf{ราชวคฺโค นิฏฺฐิโต จตุตฺโถ.}}
\titlehead{\textbf{ตสฺสุทฺทานํ}}
\csromangathameasure{ฆฏิกาโร รฏฺฐปาโล, มฆเทโว มธุริยํ. \\ โพธิ องฺคุลิมาโล จ, ปิยชาตํ พาหิติกํ. \\ ธมฺมเจติยสุตฺตญฺจ, ทสมํ กณฺณกตฺถลํ.}
\csromangathagroup{\csromangathastanza{ฆฏิกาโร รฏฺฐปาโล, มฆเทโว มธุริยํ. \\ โพธิ องฺคุลิมาโล จ, ปิยชาตํ พาหิติกํ. \\ ธมฺมเจติยสุตฺตญฺจ, ทสมํ กณฺณกตฺถลํ.}}

\chapterheadpageread{5. พฺราหฺมณวคฺค}
\csromanheader{2}{1. พฺรหฺมายุสุตฺต}
\proseitem{383}{\csromanfolio{334}เอวํ เม สุตํ\csromandash{}เอกํ สมยํ ภควา วิเทเหสุ จาริกํ จรติ มหตา ภิกฺขุสํเฆน สทฺธิํ ปญฺจมตฺเตหิ ภิกฺขุสเตหิ.\csromanspacer{} เตน โข ปน สมเยน พฺรหฺมายุ พฺราหฺมโณ มิถิลายํ ปฏิวสติ ชิณฺโณ วุฑฺโฒ มหลฺลโก อทฺธคโต วโย-อนุปฺปตฺโต วีสวสฺสสติโก ชาติยา, ติณฺณํ เวทานํ\footnote{เพทานํ (ก)} ปารคู สนิฆณฺฑุเกฏุภานํ สากฺขรปฺปเภทานํ อิติหาสปญฺจมานํ ปทโก เวยฺยากรโณ โลกายตมหาปุริสลกฺขเณสุ อนวโย.\csromanspacer{} อสฺโสสิ โข พฺรหฺมายุ พฺราหฺมโณ “สมโณ ขลุ โภ โคตโม สกฺยปุตฺโต สกฺยกุลา ปพฺพชิโต วิเทเหสุ จาริกํ จรติ มหตา ภิกฺขุสํเฆน สทฺธิํ ปญฺจมตฺเตหิ ภิกฺขุสเตหิ, ตํ โข ปน ภวนฺตํ โคตมํ เอวํ กลฺยาโณ กิตฺติสทฺโท อพฺภุคฺคโต ‘อิติปิ โส ภควา อรหํ สมฺมาสมฺพุทฺโธ วิชฺชาจรณสมฺปนฺโน สุคโต โลกวิทู อนุตฺตโร ปุริสทมฺมสารถิ สตฺถา เทวมนุสฺสานํ พุทฺโธ ภควา’ติ, โส อิมํ โลกํ สเทวกํ สมารกํ สพฺรหฺมกํ สสฺสมณพฺราหฺมณิํ ปชํ สเทวมนุสฺสํ สยํ อภิญฺญา สจฺฉิกตฺวา ปเวเทติ, โส ธมฺมํ เทเสติ อาทิกลฺยาณํ มชฺเฌกลฺยาณํ ปริโยสานกลฺยาณํ สาตฺถํ สพฺยญฺชนํ เกวลปริปุณฺณํ ปริสุทฺธํ พฺรหฺมจริยํ ปกาเสติ, สาธุ โข ปน ตาถารูปานํ อรหตํ ทสฺสนํ โหตี”ติ.}

\proseitem{384}{เตน โข ปน สมเยน พฺรหฺมายุสฺส พฺราหฺมณสฺส อุตฺตโร นาม มาณโว อนฺเตวาสี โหติ ติณฺณํ เวทานํ ปารคู สนิฆณฺฑุเกฏุภานํ สากฺขรปฺปเภทานํ อิติหาสปญฺจมานํ ปทโก เวยฺยากรโณ โลกายตมหาปุริสลกฺขเณสุ อนวโย.\csromanspacer{} อถ โข พฺรหฺมายุ พฺราหฺมโณ อุตฺตรํ มาณวํ อามนฺเตสิ “อยํ ตาต อุตฺตร สมโณ โคตโม สกฺยปุตฺโต สกฺยกุลา ปพฺพชิโต วิเทเหสุ จาริกํ จรติ มหตา ภิกฺขุสํเฆน สทฺธิํ ปญฺจมตฺเตหิ ภิกฺขุสเตหิ.\csromanspacer{} ตํ โข ปน ภวนฺตํ โคตมํ เอวํ กลฺยาโณ กิตฺติสทฺโท อพฺภุคฺคโต อิติปิ โส ภควา อรหํ สมฺมาสมฺพุทฺโธ ฯเปฯ สาธุ โข ปน ตถารูปานํ อรหตํ}

\vspace{\tipitakaparskip}
\csromancenter{พฺรหฺมายุสุตฺต (91) ทสฺสนํ โหตี”ติ.\csromanspacer{} เอหิ ตฺวํ ตาต อุตฺตร เยน สมโณ โคตโม เตนุปสงฺกม, อุปสงฺกมิตฺวา สมณํ โคตมํ ชานาหิ, ยทิ วา ตํ ภวนฺตํ โคตมํ ตถา สนฺตํเยว สทฺโท อพฺภุคฺคโต, ยทิ วา โน ตถา.\csromanspacer{} ยทิ วา โส ภวํ โคตโม ตาทิโส, ยทิ วา น ตาทิโส.\csromanspacer{} ตถา มยํ ตํ ภวนฺตํ โคตมํ เวทิสฺสามา”ติ.\csromanspacer{} ยถา กถํ ปนาหํ โภ ตํ ภวนฺตํ โคตมํ ชานิสฺสามิ “ยทิ วา ตํ ภวนฺตํ โคตมํ ตถา สนฺตํเยว สทฺโท อพฺภุคฺคโต, ยทิ วา โน ตถา.\csromanspacer{} ยทิ วา โส ภวํ โคตโม ตาทิโส, ยทิ วา น ตาทิโส”ติ.\csromanspacer{} อาคตานิ โข ตาต อุตฺตร อมฺหากํ มนฺเตสุ ทฺวตฺติํสมหาปุริสลกฺขณานิ.\csromanspacer{} เยหิ สมนฺนาคตสฺส มหาปุริสสฺส เทฺวเยว คติโย ภวนฺติ อนญฺญา.\csromanspacer{} สเจ อคารํ อชฺฌาวสติ, ราชา โหติ จกฺกวตฺตี ธมฺมิโก ธมฺมราชา จาตุรนฺโต วิชิตาวี ชนปทตฺถาวริยปฺปตฺโต สตฺตรตนสมนฺนาคโต.\csromanspacer{} ตสฺสิมานิ สตฺต รตนานิ ภวนฺติ.\csromanspacer{} เสยฺยถิทํ, จกฺกรตนํ หตฺถิรตนํ อสฺสรตนํ มณิรตนํ อิตฺถิรตนํ คหปติรตนํ ปริณายกรตนเมว สตฺตมํ.\csromanspacer{} ปโรสหสฺสํ โข ปนสฺส ปุตฺตา ภวนฺติ สูรา วีรงฺครูปา ปรเสนปฺปมทฺทนา, โส อิมํ ปถวิํ สาครปริยนฺตํ อทณฺเฑน อสตฺเถน ธมฺเมน\footnote{ธมฺเมน สเมน (ก)} อภิวิชิย อชฺฌาวสติ.\csromanspacer{} สเจ โข ปน อคารสฺมา อนคาริยํ ปพฺพชติ, อรหํ โหติ สมฺมาสมฺพุทฺโธ โลเก วิวฏฺฏจฺฉโท.\csromanspacer{} อหํ โข ปน ตาต อุตฺตร มนฺตานํ ทาตา, ตฺวํ มนฺตานํ ปฏิคฺคเหตาติ.}
\proseitem{385}{\csromanfolio{335}“เอวํ โภ”ติ โข อุตฺตโร มาณโว พฺรหฺมายุสฺส พฺราหฺมณสฺส ปฏิสฺสุตฺวา อุฏฺฐายาสนา พฺรหฺมายุํ พฺราหฺมณํ อภิวาเทตฺวา ปทกฺขิณํ กตฺวา วิเทเหสุ เยน ภควา เตน จาริกํ ปกฺกามิ, อนุปุพฺเพน จาริกํ จรมาโน เยน ภควา เตนุปสงฺกมิ, อุปสงฺกมิตฺวา ภควตา สทฺธิํ สมฺโมทิ, สมฺโมทนียํ กถํ สารณียํ วีติสาเรตฺวา เอกมนฺตํ นิสีทิ, เอกมนฺตํ นิสินฺโน โข อุตฺตโร มาณโว ภควโต กาเย ทฺวตฺติํสมหาปุริสลกฺขณานิ สมนฺเนสิ.\csromanspacer{} อทฺทสา โข อุตฺตโร มาณโว ภควโต กาเย ทฺวตฺติํสมหาปุริสลกฺขณานิ เยภุยฺเยน ถเปตฺวา เทฺว, ทฺวีสุ มหาปุริสลกฺขเณสุ กงฺขติ วิจิกิจฺฉติ นาธิมุจฺจติ น สมฺปสีทติ โกโสหิเต จ วตฺถคุเยฺห ปหูตชิวฺหตาย จ.\csromanspacer{} อถ โข ภควโต เอตทโหสิ “ปสฺสติ โข เม อยํ อุตฺตโร มาณโว \csromanfolio{336}ทฺวตฺติํสมหาปุริสลกฺขณานิ เยภุยฺเยน ถเปตฺวา เทฺว, ทฺวีสุ มหาปุริสลกฺขเณสุ กงฺขติ วิจิกิจฺฉติ นาธิมุจฺจติ น สมฺปสีทติ โกโสหิเต จ วตฺถคุเยฺห ปหูตชิวฺหตาย จา”ติ.\csromanspacer{} อถ โข ภควา ตถารูปํ อิทฺธาภิสงฺขารํ อภิสงฺขาสิ, ยถา อทฺทส อุตฺตโร มาณโว ภควโต โกโสหิตํ วตฺถคุยฺหํ.\csromanspacer{} อถ โข ภควา ชิวฺหํ นินฺนาเมตฺวา อุโภปิ กณฺณโสตานิ อนุมสิ ปฏิมสิ\footnote{ปริมสิ (สี, ก)}, อุโภปิ นาสิกโสตานิ\footnote{นาสิกาโสตานิ (สี)} อนุมสิ ปฏิมสิ, เกวลมฺปิ นลาฏมณฺฑลํ ชิวฺหาย ฉาเทสิ.\csromanspacer{} อถ โข อุตฺตรสฺส มาณวสฺส เอตทโหสิ “สมนฺนาคโต โข สมโณ โคตโม ทฺวตฺติํสมหาปุริสลกฺขเณหิ, ยํนูนาหํ สมณํ โคตมํ อนุพฺพนฺเธยฺยํ, อิริยาปถมสฺส ปสฺเสยฺยนฺ”ติ.\csromanspacer{} อถ โข อุตฺตโร มาณโว สตฺตมาสานิ ภควนฺตํ อนุพนฺธิ ฉายาว อนปายินี\footnote{อนุปายินี (สฺยา, กํ, ก)}.}

\proseitem{386}{อถ โข อุตฺตโร มาณโว สตฺตนฺนํ มาสานํ อจฺจเยน วิเทเหสุ เยน มิถิลา เตน จาริกํ ปกฺกามิ, อนุปุพฺเพน จาริกํ จรมาโน เยน มิถิลา, เยน พฺรหฺมายุ พฺราหฺมโณ เตนุปสงฺกมิ, อุปสงฺกมิตฺวา พฺรหฺมายุํ พฺราหฺมณํ อภิวาเทตฺวา นิสีทิ, เอกมนฺตํ นิสินฺนํ โข อุตฺตรํ มาณวํ พฺรหฺมายุ พฺราหฺมโณ เอตทโวจ “กจฺจิ ตาต อุตฺตร ตํ ภวนฺตํ โคตมํ ตถา สนฺตํเยว สทฺโท อพฺภุคฺคโต, โน อญฺญถา.\csromanspacer{} กจฺจิ ปน โส ภวํ โคตโม ตาทิโส, โน อญฺญาทิโส”ติ.\csromanspacer{} ตถา สนฺตํเยว โภ ตํ ภวนฺตํ โคตมํ สทฺโท อพฺภุคฺคโต, โน อญฺญถา.\csromanspacer{} ตาทิโสว\footnote{ตาทิโสว โภ (สี, อิ), ตาทิโส จ โข (สฺยา, กํ, ก)} โส ภวํ โคตโม, โน อญฺญาทิโส.\csromanspacer{} สมนฺนาคโต จ\footnote{สมนฺนาคโต จ โภ (สพฺพตฺถ)} โส ภวํ โคตโม ทฺวตฺติํสมหาปุริสลกฺขเณหิ.}

\prose{สุปฺปติฏฺฐิตปาโท โข ปน โส ภวํ โคตโม.\csromanspacer{} อิทมฺปิ ตสฺส โภโต โคตมสฺส มหาปุริสสฺส มหาปุริสลกฺขณํ ภวติ. (1)}

\prose{เหฏฺฐา โข ปน ตสฺส โภโต โคตมสฺส ปาทตเลสุ จกฺกานิ ชาตานิ สหสฺสารานิ สเนมิกานิ สนาภิกานิ สพฺพาการปริปูรานิ. (2)}

\csromanheader{2}{1. พฺรหฺมายุสุตฺต (91)}
\csromanheader{2}{อายตปณฺหิ โข ปน โส ภวํ โคตโม. (3)}
\csromanheader{2}{ทีฆงฺคุลิ โข ปน โส ภวํ โคตโม. (4)}
\csromanheader{2}{มุทุตลุนหตฺถปาโท โข ปน โส ภวํ โคตโม. (5)}
\csromanheader{2}{ชาลหตฺถปาโท โข ปน โส ภวํ โคตโม. (6)}
\csromanheader{2}{อุสฺสงฺขปาโท โข ปน โส ภวํ โคตโม. (7)}
\csromanheader{2}{เอณิชงฺโฆ โข ปน โส ภวํ โคตโม. (8)}
\prose{\csromanfolio{337}ฐิตโก โข ปน โส ภวํ โคตโม อโนนมนฺโต อุโภหิ ปาณิตเลหิ ชณฺณุกานิ ปริมสติ ปริมชฺชติ. (9)}

\csromanmatikaanchor{mtk.49}
\csromantocmarknopage{chapter}{5. พฺราหฺมณวคฺค}{mtk.49}
\bookmark[dest={mtk.49},level=0]{5. พฺราหฺมณวคฺค}
\csromanheader{2}{โกโสหิตวตฺถคุโยฺห โข ปน โส ภวํ โคตโม. (10)}
\csromanmatikaanchor{mtk.50}
\csromantocmarkref{subsection}{1. พฺรหฺมายุสุตฺต}{mtk.50}
\bookmark[dest={mtk.50},level=2]{1. พฺรหฺมายุสุตฺต}
\prose{สุวณฺณวณฺโณ โข ปน โส ภวํ โคตโม กญฺจนสนฺนิภตฺตโจ. (11)}

\prose{สุขุมจฺฉวิ โข ปน โส ภวํ โคตโม, สุขุมตฺตา ฉวิยา รโชชลฺลํ กาเย น อุปลิมฺปติ. (12)}

\prose{เอเกกโลโม โข ปน โส ภวํ โคตโม, เอเกกานิ โลมานิ โลมกูเปสุ ชาตานิ. (13)}

\prose{อุทฺธคฺคโลโม โข ปน โส ภวํ โคตโม, อุทฺธคฺคานิ โลมานิ ชาตานิ นีลานิ อญฺชนวณฺณานิ กุณฺฑลาวฏฺฏานิ ทกฺขิณาวฏฺฏกชาตานิ. (14)}

\csromanheader{2}{พฺรหฺมุชุคตฺโต โข ปน โส ภวํ โคตโม. (15)}
\csromanheader{2}{สตฺถุสฺสโท โข ปน โส ภวํ โคตโม. (16)}
\csromanheader{2}{สีหปุพฺพทฺธกาโย โข ปน โส ภวํ โคตโม. (17)}
\csromanheader{2}{จิตนฺตรํโส โข ปน โส ภวํ โคตโม. (18)}
\prose{นิโคฺรธปริมณฺฑโล โข ปน โส ภวํ โคตโม, ยาวตกฺวสฺส กาโย, ตาวตกฺวสฺส พฺยาโม, ยาวตกฺวสฺส พฺยาโม, ตาวตกฺวสฺส กาโย. (19)}

\csromanheader{2}{สมวฏฺฏกฺขนฺโธ โข ปน โส ภวํ โคตโม. (20)}
\csromanheader{2}{รสคฺคสคฺคี โข ปน โส ภวํ โคตโม. (21)}
\csromanheader{2}{สีหหนุ โข ปน โส ภวํ โคตโม. (22)}
\csromanheader{2}{จตฺตาลีสทนฺโต โข ปน โส ภวํ โคตโม. (23)}
\csromanheader{2}{สมทนฺโต โข ปน โส ภวํ โคตโม. (24)}
\csromanheader{2}{อวิรฬทนฺโต โข ปน โส ภวํ โคตโม. (25)}
\csromanheader{2}{สุสุกฺกทาโฐ โข ปน โส ภวํ โคตโม. (26)}
\csromanheader{2}{ปหูตชิโวฺห โข ปน โส ภวํ โคตโม. (27)}
\csromanheader{2}{พฺรหฺมสฺสโร โข ปน โส ภวํ โคตโม กรวิกภาณี. (28)}
\csromanheader{2}{อภินีลเนตฺโต โข ปน โส ภวํ โคตโม. (29)}
\csromanheader{2}{โคปขุโม โข ปน โส ภวํ โคตโม. (30)}
\prose{\csromanfolio{338}อุณฺณา โข ปนสฺส โภโต โคตมสฺส ภมุกนฺตเร ชาตา โอทาตา มุทุตูลสนฺนิภา. (31)}

\prose{อุณฺหีสสีโส โข ปน โส ภวํ โคตโม, อิทมฺปิ ตสฺส โภโต โคตมสฺส มหาปุริสสฺส มหาปุริสลกฺขณํ ภวติ. (32)}

\prose{อิเมหิ โข โภ โส ภวํ โคตโม ทฺวตฺติํสมหาปุริสลกฺขเณหิ สมนฺนาคโต.}

\proseitem{387}{คจฺฉนฺโต โข ปน โส ภวํ โคตโม ทกฺขิเณเนว ปาเทน ปฐมํ ปกฺกมติ, โส นาติทูเร ปาทํ อุทฺธรติ, นาจฺจาสนฺเน ปาทํ นิกฺขิปติ.\csromanspacer{} โส นาติสีฆํ คจฺฉติ, นาติสณิกํ คจฺฉติ, น จ อทฺทุเวน อทฺทุวํ สงฺฆฏฺเฏนฺโต คจฺฉติ, น จ โคปฺผเกน โคปฺผกํ สงฺฆฏฺเฏนฺโต คจฺฉติ.\csromanspacer{} โส คจฺฉนฺโต น สตฺถิํ อุนฺนาเมติ, น สตฺถิํ โอนาเมติ, น สตฺถิํ สนฺนาเมติ, น สตฺถิํ วินาเมติ.\csromanspacer{} คจฺฉโต โข ปน ตสฺส โภโต โคตมสฺส อธรกาโยว\footnote{อฑฺฒกาโยว (ก) อารทฺธกาโยว (สฺยา, กํ)} อิญฺชติ, น จ กายพเลน คจฺฉติ.\csromanspacer{} อปโลเกนฺโต โข ปน โส ภวํ โคตโม สพฺพกาเยเนว อปโลเกติ, โส น อุทฺธํ อุลฺโลเกติ, น อโธ โอโลเกติ, น จ วิเปกฺขมาโน คจฺฉติ, ยุคมตฺตญฺจ เปกฺขติ, ตโต จสฺส อุตฺตริ อนาวฏํ ญาณทสฺสนํ ภวติ.\csromanspacer{} โส อนฺตรฆรํ ปวิสนฺโต น กายํ}

\vspace{\tipitakaparskip}
\csromancenter{พฺรหฺมายุสุตฺต (91) อุนฺนาเมติ, น กายํ โอนาเมติ, น กายํ สนฺนาเมติ, น กายํ วินาเมติ.\csromanspacer{} โส นาติทูเร นาจฺจาสนฺเน อาสนสฺส ปริวตฺตติ, น จ ปาณินา อาลมฺพิตฺวา อาสเน นิสีทติ, น จ อาสนสฺมิํ กายํ ปกฺขิปติ.\csromanspacer{} โส อนฺตรฆเร นิสินฺโน สมาโน น หตฺถกุกฺกุจฺจํ อาปชฺชติ, น ปาทกุกฺกุจฺจํ อาปชฺชติ, น อทฺทุเวน อทฺทุวํ อาโรเปตฺวา นิสีทติ, น จ โคปฺผเกน โคปฺผกํ อาโรเปตฺวา นิสีทติ, น จ ปาณินา หนุกํ อุปทหิตฺวา\footnote{อุปาทิยิตฺวา (สี, อิ)} นิสีทติ.\csromanspacer{} โส อนฺตรฆเร นิสินฺโน สมาโน น ฉมฺภติ น กมฺปติ น เวธติ น ปริตสฺสติ, โส อฉมฺภี อกมฺปี อเวธี อปริตสฺสี วิคตโลมหํโส วิเวกวตฺโต จ โส ภวํ โคตโม อนฺตรฆเร นิสินฺโน โหติ.\csromanspacer{} โส ปตฺโตทกํ ปฏิคฺคณฺหนฺโต น ปตฺตํ อุนฺนาเมติ, น ปตฺตํ โอนาเมติ, น ปตฺถํ สนฺนาเมติ, น ปตฺตํ วินาเมติ.\csromanspacer{} โส ปตฺโตทกํ ปฏิคฺคณฺหาติ นาติโถกํ นาติพหุํ.\csromanspacer{} โส น ขุลุขุลุการกํ\footnote{พุลุพุลุการกํ (สี)} ปตฺตํ โธวติ, น สมฺปริวตฺตกํ ปตฺตํ โธวติ.\csromanspacer{} น ปตฺตํ ภูมิยํ นิกฺขิปิตฺวา หตฺเถ โธวติ, หตฺเถสุ โธเตสุ ปตฺโต โธโต โหติ, ปตฺเต โธเต หตฺถา โธตา โหนฺติ.\csromanspacer{} โส ปตฺโตทกํ ฉฑฺเฑติ นาติทูเร นาจฺจาสนฺเน, น จ วิจฺฉฑฺฑยมาโน.\csromanspacer{} โส โอทนํ ปฏิคฺคณฺหนฺโต น ปตฺตํ อุนฺนาเมติ, น ปตฺตํ โอนาเมติ, น ปตฺตํ สนฺนาเมติ, น ปตฺตํ วินาเมติ, โส โอทนํ ปฏิคฺคณฺหาติ นาติโถกํ นาติพหุํ.\csromanspacer{} พฺยญฺชนํ โข ปน ภวํ โคตโม พฺยญฺชนมตฺตาย อาหาเรติ, น จ พฺยญฺชเนน อาโลปํ อตินาเมติ.\csromanspacer{} ทฺวตฺติกฺขตฺตุํ โข ภวํ โคตโม มุเข อาโลปํ สมฺปริวตฺเตตฺวา อชฺโฌหรติ.\csromanspacer{} น จสฺส กาจิ โอทนมิญฺชา อสมฺภินฺนา กายํ ปวิสติ, น จสฺส กาจิ โอทนมิญฺชา มุเข อวสิฏฺฐา โหติ, อถาปรํ อาโลปํ อุปนาเมติ.\csromanspacer{} รสปฏิสํเวที โข ปน โส ภวํ โคตโม อาหารํ อาหาเรติ, โน จ รสราคปฏิสํเวที.}
\prose{\csromanfolio{339}อฏฺฐงฺคสมนฺนาคตํ\footnote{อฏฺฐงฺคสมนฺนาคโต (ก)} โข ปน โส ภวํ โคตโม อาหารํ อาหาเรติ, เนว ทวาย น มทาย น มณฺฑนาย น วิภูสนาย, ยาวเทว อิมสฺส กายสฺส ฐิติยา ยาปนาย วิหิํสูปรติยา พฺรหฺมจริยานุคฺคหาย, อิติ ปุราณญฺจ เวทนํ ปฏิหงฺขามิ, นวญฺจ เวทนํ น อุปฺปาเทสฺสามิ, ยาตฺรา จ เม ภวิสฺสติ อนวชฺชตา จ ผาสุวิหาโร \csromanfolio{340}จาติ.\csromanspacer{} โส ภุตฺตาวี ปตฺโตทกํ ปฏิคฺคณฺหนฺโต น ปตฺตํ อุนฺนาเมติ, น ปตฺตํ โอนาเมติ, น ปตฺตํ สนฺนาเมติ, น ปตฺตํ วินาเมติ.\csromanspacer{} โส ปตฺโตทกํ ปฏิคฺคณฺหาติ นาติโถกํ นาติพหุํ.\csromanspacer{} โส น ขุลุขุลุการกํ ปตฺตํ โธวติ, น สมฺปริวตฺตกํ ปตฺตํ โธวติ.\csromanspacer{} น ปตฺตํ ภูมิยํ นิกฺขิปิตฺวา หตฺเถ โธวติ, หตฺเถสุ โธเตสุ ปตฺโต โธโต โหติ, ปตฺเต โธเต หตฺถา โธตา โหนฺติ.\csromanspacer{} โส ปตฺโตทกํ ฉฑฺเฑติ นาติทูเร นาจฺจาสนฺเน, น จ วิจฺฉฑฺฑยมาโน.\csromanspacer{} โส ภุตฺตาวี น ปตฺตํ ภูมิยํ นิกฺขิปติ นาติทูเร นาจฺจาสนฺเน, น จ อนตฺถิโก ปตฺเตน โหติ, น จ อติเวลานุรกฺขี ปตฺตสฺมิํ.\csromanspacer{} โส ภุตฺตาวี มุหุตฺตํ ตุณฺหี นิสีทติ, น จ อนุโมทนสฺส กาลมตินาเมติ, โส ภุตฺตาวี อนุโมทติ.\csromanspacer{} น ตํ ภตฺตํ ครหติ, น อญฺญํ ภตฺตํ ปฏิกงฺขติ, อญฺญทตฺถุ ธมฺมิยา กถาย ตํ ปริสํ สนฺทสฺเสติ สมาทเปติ สมุตฺเตเชติ สมฺปหํเสติ, โส ตํ ปริสํ ธมฺมิยา กถาย สนฺทสฺเสตฺวา สมาทเปตฺวา สมุตฺเตเชตฺวา สมฺปหํเสตฺวา อุฏฺฐายาสนา ปกฺกมติ, โส นาติสีฆํ คจฺฉติ, นาติสณิกํ คจฺฉติ, น จ มุจฺจิตุกาโม คจฺฉติ.\csromanspacer{} น จ ตสฺส โภโต โคตมสฺส กาเย จีวรํ อจฺจุกฺกฏฺฐํ โหติ, น จ อจฺโจกฺกฏฺฐํ, น จ กายสฺมิํ อลฺลีนํ, น จ กายสฺมา อปกฏฺฐํ, น จ ตสฺส โภโต โคตมสฺส กายมฺหา วาโต จีวรํ อปวหติ.\csromanspacer{} น จ ตสฺส โภโต โคตมสฺส กาเย รโชชลฺลํ อุปลิมฺปติ.\csromanspacer{} โส อารามคโต นิสีทติ ปญฺญตฺเต อาสเน, นิสชฺช ปาเท ปกฺขาเลติ.\csromanspacer{} น จ โส ภวํ โคตโม ปาทมณฺฑนานุโยคมนุยุตฺโต วิหรติ.\csromanspacer{} โส ปาเท ปกฺขาเลตฺวา นิสีทติ ปลฺลงฺกํ อาภุชิตฺวา อุชุํ กายํ ปณิธาย ปริมุขํ สติํ อุปฏฺฐเปตฺวา.\csromanspacer{} โส เนว อตฺตพฺยาพาธาย เจเตติ, น ปรพฺยาพาธาย เจเตติ, น อุภยพฺยาพาธาย เจเตติ, อตฺตหิตปรหิต-อุภยหิตสพฺพโลกหิตเมว โส ภวํ โคตโม จินฺเตนฺโต นิสินฺโน โหติ.\csromanspacer{} โส อารามคโต ปริสติ ธมฺมํ เทเสติ, น ตํ ปริสํ อุสฺสาเทติ, น ตํ ปริสํ อปสาเทติ, อญฺญทตฺถุ ธมฺมิยา กถาย ตํ ปริสํ สนฺทสฺเสติ สมาทเปติ สมุตฺเตเชติ สมฺปหํเสติ.}

\prose{อฏฺฐงฺคสมนฺนาคโต โข ปนสฺส โภโต โคตมสฺส มุขโต โฆโส นิจฺฉรติ วิสฺสฏฺโฐ จ วิญฺเญยฺโย จ มญฺชุ จ สวนีโย จ พินฺทุ จ อวิสารี จ คมฺภีโร จ นินฺนาที จ.\csromanspacer{} ยถาปริสํ โข ปน โส}

\vspace{\tipitakaparskip}
\csromancenter{พฺรหฺมายุสุตฺต (91) ภวํ โคตโม สเรน วิญฺญาเปติ, น จสฺส พหิทฺธา ปริสาย โฆโส นิจฺฉรติ, เต เตน โภตา โคตเมน ธมฺมิยา กถาย สนฺทสฺสิตา สมาทปิตา สมุตฺเตชิตา สมฺปหํสิตา อุฏฺฐายาสนา ปกฺกมนฺติ, อวโลกยมานาเยว\footnote{อปโลกยมานาเยว (สี, ก)} อวิชหิตตฺตา\footnote{อวิชหนฺตาภาเวน (สี, สฺยา, กํ, อิ)}.\csromanspacer{} อทฺทสาม โข มยํ โภ ตํ ภวนฺตํ โคตมํ คจฺฉนฺตํ, อทฺทสาม ฐิตํ, อทฺทสาม อนฺตรฆรํ ปวิสนฺตํ, อทฺทสาม อนฺตรฆเร นิสินฺนํ ตุณฺหีภูตํ, อทฺทสาม อนฺตรฆเร ภุญฺชนฺตํ, อทฺทสาม ภุตฺตาวิํ นิสินฺนํ ตุณฺหีภูตํ, อทฺทสาม ภุตฺตาวิํ อนุโมทนฺตํ, อทฺทสาม อารามํ คจฺฉนฺตํ, อทฺทสาม อารามคตํ นิสินฺนํ ตุณฺหีภูตํ, อทฺทสาม อารามคตํ ปริสติ ธมฺมํ เทเสนฺตํ, เอทิโส จ เอทิโส จ โส ภวํ โคตโม, ตโต จ ภิยฺโยติ.}
\proseitem{388}{\csromanfolio{341}เอวํ วุตฺเต พฺรหฺมายุ พฺราหฺมโณ อุฏฺฐายาสนา เอกํสํ อุตฺตราสงฺคํ กริตฺวา เยน ภควา เตนญฺชลิํ ปณาเมตฺวา ติกฺขตฺตุํ อุทานํ อุทาเนติ\csromandash{}}

\namakkaram{นโม ตสฺส ภควโต อรหโต สมฺมาสมฺพุทฺธสฺส,}
\namakkaram{นโม ตสฺส ภควโต อรหโต สมฺมาสมฺพุทฺธสฺส,}
\namakkaram{นโม ตสฺส ภควโต อรหโต สมฺมาสมฺพุทฺธสฺสาติ.}
\prose{อปฺเปว นาม มยํ กทาจิ กรหจิ เตน โภตา โคตเมน สมาคจฺเฉยฺยาม, อปฺเปว นาม สิยา โกจิเทว กาถาสลฺลาโปติ.}

\proseitem{389}{อถ โข ภควา วิเทเหสุ อนุปุพฺเพน จาริกํ จรมาโน เยน มิถิลา ตทวสริ, ตตฺร สุทํ ภควา มิถิลายํ วิหรติ มฆเทวมฺพวเน.\csromanspacer{} อสฺโสสุํ โข มิถิเลยฺยกา\footnote{เมถิเลยฺยกา (สี, อิ)} พฺราหฺมณคหปติกา “สมโณ ขลุ โภ โคตโม สกฺยปุตฺโต สกฺยกุลา ปพฺพชิโต วิเทเหสุ จาริกํ จรมาโน มหตา ภิกฺขุสํเฆน สทฺธิํ ปญฺจมตฺเตหิ ภิกฺขุสเตหิ มิถิลํ อนุปฺปตฺโต มิถิลายํ วิหรติ มฆเทวมฺพวเน.\csromanspacer{} ตํ โข ปน ภวนฺตํ โคตมํ เอวํ กลฺยาโณ กิตฺติสทฺโท อพฺภุคฺคโต ‘อิติปิ โส ภควา อรหํ สมฺมาสมฺพุทฺโธ วิชฺชาจรณสมฺปนฺโน สุคโต โลกวิทู อนุตฺตโร ปุริสทมฺมสารถิ สตฺถา เทวมนุสฺสานํ พุทฺโธ ภควา’ติ, โส อิมํ โลกํ สเทวกํ สมารกํ สพฺรหฺมกํ สสฺสมณพฺราหฺมณิํ \csromanfolio{342}ปชํ สเทวมนุสฺสํ สยํ อภิญฺญา สจฺฉิกตฺวา ปเวเทติ, โส ธมฺมํ เทเสติ อาทิกลฺยาณํ มชฺเฌกลฺยาณํ ปริโยสานกลฺยาณํ สาตฺถํ สพฺยาญฺชนํ เกวลปริปุณฺณํ ปริสุทฺธํ พฺรหฺมจริยํ ปกาเสติ, สาธุ โข ปน ตถารูปานํ อรหตํ ทสฺสนํ โหตี”ติ.}

\prose{อถ โข มิถิเลยฺยกา พฺราหฺมณคหปติกา เยน ภควา เตนุปสงฺกมิํสุ, อุปสงฺกมิตฺวา อปฺเปกจฺเจ ภควนฺตํ อภิวาเทตฺวา เอกมนฺตํ นิสีทิํสุ, อปฺเปกจฺเจ ภควตา สทฺธิํ สมฺโมทิํสุ, สมฺโมทนียํ กถํ สารณียํ วีติสาเรตฺวา เอกมนฺตํ นิสีทิํสุ, อปฺเปกจฺเจ เยน ภควา เตนญฺชลิํ ปณาเมตฺวา เอกมนฺตํ นิสีทิํสุ, อปฺเปกจฺเจ ภควโต สนฺติเก นามโคตฺตํ สาเวตฺวา เอกมนฺตํ นิสีทิํสุ, อปฺเปกจฺเจ ตุณีภูตา เอกมนฺตํ นิสีทิํสุ.}

\proseitem{390}{อสฺโสสิ โข พฺรหฺมายุ พฺราหฺมโณ “สมโณ ขลุ โภ โคตโม สกฺยปุตฺโต สกฺยกุลา ปพฺพชิโต มิถิลํ อนุปฺปตฺโต มิถิลายํ วิหรติ มฆเทวมฺพวเน”ติ.\csromanspacer{} อถ โข พฺรหฺมายุ พฺราหฺมโณ สมฺพหุเลหิ สาวเกหิ สทฺธิํ เยน มฆเทวมฺพวนํ เตนุปสงฺกมิ, อถ โข พฺรหฺมายุโน พฺราหฺมณสฺส อวิทูเร อมฺพวนสฺส เอตทโหสิ “น โข เมตํ ปติรูปํ, โยหํ ปุพฺเพ อปฺปฏิสํวิทิโต สมณํ โคตมํ ทสฺสนาย อุปสงฺกเมยฺยนฺ”ติ.\csromanspacer{} อถ โข พฺรหฺมายุ พฺราหฺมโณ อญฺญตรํ มาณวกํ อามนฺเตสิ “เอหิ ตฺวํ มาณวก เยน สมโณ โคตโม เตนุปสงฺกม, อุปสงฺกมิตฺวา มม วจเนน สมณํ โคตมํ อปฺปาพาธํ อปฺปาตงฺกํ ลหุฏฺฐานํ พลํ ผสุวิหารํ ปุจฺฉ ‘พฺรหฺมายุ โภ โคตม พฺราหฺมโณ ภวนฺตํ โคตมํ อปฺปาพาธํ อปฺปาตงฺกํ ลหุฏฺฐานํ พลํ ผาสุวิหารํ ปุจฺฉตี’ติ, เอวญฺจ วเทหิ ‘พฺรหฺมายุ โภ โคตม พฺราหฺมโณ ชิณฺโณ วุฑฺโฒ มหลฺลโก อทฺธคโต วโย-อนุปฺปตฺโต วีสวสฺสสติโก ชาติยา, ติณฺณํ เวทานํ ปารคู สนิฆณฺฑุเกฏุภานํ สากฺขรปฺปเภทานํ อิติหาสปญฺจมานํ ปทโก เวยฺยากรโณ โลกายตมหาปุริสลกฺขเณสุ อนวโย, ยาวตา โภ พฺราหฺมณคหปติกา มิถิลายํ ปฏิวสนฺติ, พฺรหฺมายุ เตสํ พฺราหฺมโณ อคฺคมกฺขายติ ยทิทํ โภเคหิ, พฺรหฺมายุ เตสํ พฺราหฺมโณ อคฺคมกฺขายติ ยทิทํ มนฺเตหิ, พฺรหฺมายุ เตสํ พฺราหฺมโณ อคฺคมกฺขายติ ยทิทํ อายุนา เจว ยสสา จ.\csromanspacer{} โส โภโต โคตมสฺส ทสฺสนกาโม’ติ”.}

\csromanheader{2}{1. พฺรหฺมายุสุตฺต (91)}
\prose{\csromanfolio{343}“เอวํ โภ”ติ โข โส มาณวโก พฺรหฺมายุสฺส พฺราหฺมณสฺส ปฏิสฺสุตฺวา เยน ภควา เตนุปสงฺกมิ, อุปสงฺกมิตฺวา ภควตา สทฺธิํ สมฺโมทิ, สมฺโมทนียํ กถํ สารณียํ วีติสาเรตฺวา เอกมนฺตํ อฏฺฐาสิ, เอกมนฺตํ ฐิโต โข โส มาณวโก ภควนฺตํ เอตทโวจ “พฺรหฺมายุ โภ โคตม พฺราหฺมโณ ภวนฺตํ โคตมํ อปฺปาพาธํ อปฺปาตงฺกํ ลหุฏฺฐานํ พลํ ผาสุวิหารํ ปุจฺฉติ, เอวญฺจ วเทติ ‘พฺรหฺมายุ โภ โคตม พฺราหฺมโณ ชิณฺโณ วุฑฺโฒ มหลฺลโก อทฺธคโต วโย-อนุปฺปตฺโต วีสวสฺสสติโก ชาติยา, ติณฺณํ เวทานํ ปารคู สนิฆณฺฑุเกฏุภานํ สากฺขรปฺปเภทานํ อิติหาสปญฺจมานํ ปทโก เวยฺยากรโณ โลกายตมหาปุริสลกฺขเณสุ อนวโย, ยาวตา โภ พฺราหฺมณคหปติกา มิถิลายํ ปฏิวสนฺติ พฺรหฺมายุ เตสํ พฺราหฺมโณ อคฺคมกฺขายติ ยทิทํ โภเคหิ, พฺรหฺมายุ เตสํ พฺราหฺมโณ อคฺคมกฺขายติ ยทิทํ มนฺเตหิ, พฺรหฺมายุ เตสํ พฺราหฺมโณ อคฺคมกฺขายติ ยทิทํ อายุนา เจว ยสสา จ.\csromanspacer{} โส โภโต โคตมสฺส ทสฺสนกาโม’ติ”.\csromanspacer{} ยสฺสทานิ มาณว พฺรหฺมายุ พฺราหฺมโณ กาลํ มญฺญตีติ.\csromanspacer{} อถ โข โส มาณวโก เยน พฺรหฺมายุ พฺราหฺมโณ เตนุปสงฺกมิ, อุปสงฺกมิตฺวา พฺรหฺมายุํ พฺราหฺมณํ เอตทโวจ “กตาวกาโส โขมฺหิ ภวตา สมเณน โคตเมน, ยสฺสทานิ ภวํ กาลํ มญฺญตี”ติ.}

\proseitem{391}{อถ โข พฺรหฺมายุ พฺราหฺมาโณ เยน ภควา เตนุปสงฺกมิ.\csromanspacer{} อทฺทสา โข สา ปริสา พฺรหฺมายุํ พฺราหฺมณํ ทูรโตว อาคจฺฉนฺตํ.\csromanspacer{} ทิสฺวาน โอรมิย\footnote{โอรมตฺถ (สฺยา, กํ, อิ), โอรมถ, โอรมติ(ก), อถ นํ (สี), โอรมิยาติ ปน ตฺวาปจฺจยนฺตตฺถสํวณฺณนานุรูปํ วิโสธิตปทํ.} โอกาสมกาสิ, ยถา ตํ ญาตสฺส ยสสฺสิโน.\csromanspacer{} อถ โข พฺรหฺมายุ พฺราหฺมโณ ตํ ปริสํ เอตทโวจ “อลํ โภ นิสีทถ ตุเมฺห สเก อาสเน, อิธาหํ สมณสฺส โคตมสฺส สนฺติเก นิสีทิสฺสามี”ติ.}

\prose{อถ โข พฺรหฺมายุ พฺราหฺมโณ เยน ภควา เตนุปสงฺกมิ, อุปสงฺกมิตฺวา ภควตา สทฺธิํ สมฺโมทิ, สมฺโมทนียํ กถํ สารณียํ วีติสาเรตฺวา เอกมนฺตํ นิสีทิ, เอกมนฺตํ นิสินฺโน โข พฺรหฺมายุ พฺราหฺมโณ \csromanfolio{344}ภควโต กาเย ทฺวตฺติํสมหาปุริสลกฺขณานิ สมนฺเนสิ.\csromanspacer{} อทฺทสา โข พฺรหฺมายุ พฺราหฺมโณ ภควโต กาเย ทฺวตฺติํสมหาปุริสลกฺขณานิ เยภุยฺเยน ฐเปตฺวา เทฺว, ทฺวีสุ มหาปุริสลกฺขเณสุ กงฺขติ วิจิกิจฺฉติ นาธิมุจฺจติ น สมฺปสีทติ โกโสหิเต จ วตฺถคุเยฺห ปหูตชิวฺหตาย จ.\csromanspacer{} อถ โข พฺรหฺมายุ พฺราหฺมโณ ภควนฺตํ คาถาหิ อชฺฌภาสิ\csromandash{}}

\csromangathameasure{“เย เม ทฺวตฺติํสาติ สุตา, มหาปุริสลกฺขณา. \\ ทุเว เตสํ น ปสฺสามิ, โภโต กายสฺมิํ โคตม. \\ กจฺจิ โกโสหิตํ โภโต, วตฺถุคุยฺหํ นรุตฺตม. \\ นารีสมานสวฺหยา, กจฺจิ ชิวฺหา น ทสฺสกา. \\ กจฺจิ ปหูตชิโวฺหสิ, ยถา ตํ ชานิยามเส. \\ นินฺนามเยตํ ปหูตํ, กงฺขํ วินย โน อิเส. \\ ทิฏฺฐธมฺมหิตตฺถาย, สมฺปรายสุขาย จ. \\ กตาวกาสา ปุจฺฉาม, ยํ กิญฺจิ อภิปตฺถิตนฺ”ติ.}
\csromangathagroup{\csromangathastanza{“เย เม ทฺวตฺติํสาติ สุตา, มหาปุริสลกฺขณา. \\ ทุเว เตสํ น ปสฺสามิ, โภโต กายสฺมิํ โคตม.}\csromangathastanza{กจฺจิ โกโสหิตํ โภโต, วตฺถุคุยฺหํ นรุตฺตม. \\ นารีสมานสวฺหยา, กจฺจิ ชิวฺหา น ทสฺสกา\footnote{นารีสหนาม สวฺหยา, กจฺจิ ชิวฺหา นรสฺสิกา. (สี, สฺยา, กํ, อิ)}.}\csromangathastanza{กจฺจิ ปหูตชิโวฺหสิ, ยถา ตํ ชานิยามเส. \\ นินฺนามเยตํ ปหูตํ, กงฺขํ วินย โน อิเส.}\csromangathastanza{ทิฏฺฐธมฺมหิตตฺถาย, สมฺปรายสุขาย จ. \\ กตาวกาสา ปุจฺฉาม, ยํ กิญฺจิ อภิปตฺถิตนฺ”ติ.}}

\proseitem{392}{อถ โข ภควโต เอตทโหสิ “ปสฺสติ โข เม อยํ พฺรหฺมายุ พฺราหฺมโณ ทฺวตฺติํสมหาปุริสลกฺขณานิ เยภุยฺเยน ฐเปตฺวา เทฺว, ทฺวีสุ มหาปุริสลกฺขเณสุ กงฺขติ วิจิกิจฺฉติ นาธิมุจฺจติ น สมฺปสีทติ โกโสหิเต จ วตฺถคุเยฺห ปหูตชิวฺหตาย จา”ติ.\csromanspacer{} อถ โข ภควา ตถารูปํ อิทฺธาภิสงฺขารํ อภิสงฺขาสิ, ยถา อทฺทส พฺรหฺมายุ พฺราหฺมโณ ภควโต โกโสหิตํ วตฺถคุยฺหํ.\csromanspacer{} อถ โข ภควา ชิวฺหํ นินฺนาเมตฺวา อุโภปิ กณฺณโสตานิ อนุมสิ ปฏิมสิ, อุโภปิ นาสิกโสตานิ อนุมสิ ปฏิมสิ, เกวลมฺปิ นลาฏมณฺฑลํ ชิวฺหาย ฉาเทสิ.\csromanspacer{} อถ โข ภควา พฺรหฺมายุํ พฺราหฺมณํ คาถาหิ ปจฺจภาสิ\csromandash{}}

\csromangathameasure{“เย เต ทฺวตฺติํสาติ สุตา, มหาปุริสลกฺขณา. \\ สพฺเพ เต มม กายสฺมิํ, มา เต กงฺขาหุ พฺราหฺมณ. \\ อภิญฺเญยฺยํ อภิญฺญาตํ, ภาเวตพฺพญฺจ ภาวิตํ. \\ ปหาตพฺพํ ปหีนํ เม, ตสฺมา พุทฺโธสฺมิ พฺราหฺมณ.}
\csromangathagroup{\csromangathastanza{“เย เต ทฺวตฺติํสาติ สุตา, มหาปุริสลกฺขณา. \\ สพฺเพ เต มม กายสฺมิํ, มา เต\footnote{มา โว (ก)} กงฺขาหุ พฺราหฺมณ.}\csromangathastanza{อภิญฺเญยฺยํ อภิญฺญาตํ, ภาเวตพฺพญฺจ ภาวิตํ. \\ ปหาตพฺพํ ปหีนํ เม, ตสฺมา พุทฺโธสฺมิ พฺราหฺมณ.}}

\csromanheader{2}{1. พฺรหฺมายุสุตฺต (91)}
\csromangathameasure{ทิฏฺฐธมฺมหิตตฺถาย, สมฺปรายสุขาย จ. \\ กตาวกาโส ปุจฺฉสฺสุ, ยํ กิญฺจิ อภิปตฺถิตนฺ”ติ.}
\csromangathagroup{\csromangathastanza{\csromanfolio{345}ทิฏฺฐธมฺมหิตตฺถาย, สมฺปรายสุขาย จ. \\ กตาวกาโส ปุจฺฉสฺสุ, ยํ กิญฺจิ อภิปตฺถิตนฺ”ติ.}}

\proseitem{393}{อถ โข พฺรหฺมายุสฺส พฺราหฺมณสฺส เอตทโหสิ “กตาวกาโส โขมฺหิ สมเณน โคตเมน.\csromanspacer{} กิํ นุ โข อหํ สมณํ โคตมํ ปุจฺเฉยฺยํ ทิฏฺฐธมฺมิกํ วา อตฺถํ สมฺปรายิกํ วา”ติ.\csromanspacer{} อถ โข พฺรหฺมายุสฺส พฺราหฺมณสฺส เอตทโหสิ “กุสโล โข อหํ ทิฏฺฐธมฺมิกานํ อตฺถานํ, อญฺเญปิ มํ ทิฏฺฐธมฺมิกํ อตฺถํ ปุจฺฉนฺติ.\csromanspacer{} ยํนูนาหํ สมณํ โคตมํ สมฺปรายิกํเยว อตฺถํ ปุจฺเฉยฺยนฺ”ติ.\csromanspacer{} อถ โข พฺรหฺมายุ พฺราหฺมโณ ภควนฺตํ คาถาหิ อชฺฌภาสิ\csromandash{}}

\csromangathameasure{“กถํ โข พฺราหฺมโณ โหติ, กถํ ภวติ เวทคู. \\ เตวิชฺโช โภ กถํ โหติ, โสตฺติโย กินฺติ วุจฺจติ. \\ อรหํ โภ กถํ โหติ, กถํ ภวติ เกวลี. \\ มุนิ จ โภ กถํ โหติ, พุทฺโธ กินฺติ ปวุจฺจตี”ติ.}
\csromangathagroup{\csromangathastanza{“กถํ โข พฺราหฺมโณ โหติ, กถํ ภวติ เวทคู. \\ เตวิชฺโช โภ กถํ โหติ, โสตฺติโย กินฺติ วุจฺจติ.}\csromangathastanza{อรหํ โภ กถํ โหติ, กถํ ภวติ เกวลี. \\ มุนิ จ โภ กถํ โหติ, พุทฺโธ กินฺติ ปวุจฺจตี”ติ.}}

\csromanheader{2}{394. อถ โข ภควา พฺรหฺมายุํ พฺราหฺมณํ คาถาหิ ปจฺจภาสิ\csromandash{}}
\csromangathameasure{“ปุพฺเพนิวาสํ โย เวทิ, สคฺคาปายญฺจ ปสฺสติ. \\ อโถ ชาติกฺขยํ ปตฺโต, อภิญฺญา โวสิโต มุนิ. \\ จิตฺตํ วิสุทฺธํ ชานาติ, มุตฺตํ ราเคหิ สพฺพโส. \\ ปหีนชาติมรโณ, พฺรหฺมจริยสฺส เกวลี. \\ ปารคู สพฺพธมฺมานํ, พุทฺโธ ตาที ปวุจฺจตี”ติ.}
\csromangathagroup{\csromangathastanza{“ปุพฺเพนิวาสํ โย เวทิ, สคฺคาปายญฺจ ปสฺสติ. \\ อโถ ชาติกฺขยํ ปตฺโต, อภิญฺญา โวสิโต มุนิ.}\csromangathastanza{จิตฺตํ วิสุทฺธํ ชานาติ, มุตฺตํ ราเคหิ สพฺพโส. \\ ปหีนชาติมรโณ, พฺรหฺมจริยสฺส เกวลี. \\ ปารคู สพฺพธมฺมานํ, พุทฺโธ ตาที ปวุจฺจตี”ติ.}}

\prose{เอวํ วุตฺเต พฺรหฺมายุ พฺราหฺมโณ อุฏฺฐายาสนา เอกํสํ อุตฺตราสงฺคํ กริตฺวา ภควโต ปาเทสุ สิรสา นิปติตฺวา ภควโต ปาทานิ มุเขน จ ปริจุมฺพติ, ปาณีหิ จ ปริสมฺพาหติ, นามญฺจ สาเวติ “พฺรหฺมายุ อหํ โภ โคตม พฺราหฺมโณ, พฺรหฺมายุ อหํ โภ โคตม พฺราหฺมโณ”ติ.\csromanspacer{} อถ โข สา ปริสา อจฺฉริยพฺภุตจิตฺตชาตา อโหสิ “อจฺฉริยํ วต โภ, อพฺภุตํ วต โภ, ยตฺร หิ นามายํ พฺรหฺมายุ พฺราหฺมโณ ญาโต ยสสฺสี เอวรูปํ ปรมนิปจฺจการํ กริสฺสตี”ติ.\csromanspacer{} อถ โข ภควา พฺรหฺมายุํ พฺราหฺมณํ เอตทโวจ “อลํ พฺราหฺมณ อุฏฺฐห, นิสีท ตฺวํ สเก อาสเน, ยโต เต มยิ จิตฺตํ ปสนฺนนฺ”ติ.\csromanspacer{} อถ โข พฺรหฺมายุ พฺราหฺมโณ อุฏฺฐหิตฺวา สเก อาสเน นิสีทิ.}

\proseitem{395}{\csromanfolio{346}อถ โข ภควา พฺรหฺมายุสฺส พฺราหฺมณสฺส อนุปุพฺพิํกถํ กเถสิ.\csromanspacer{} เสยฺยถิทํ, ทานกถํ สีลกถํ สคฺคกถํ กามานํ อาทีนวํ โอการํ สํกิเลสํ เนกฺขมฺเม อานิสํสํ ปกาเสสิ.\csromanspacer{} ยทา ภควา อญฺญาสิ พฺรหฺมายุํ พฺราหฺมณํ กลฺลจิตฺตํ มุทุจิตฺตํ วินีวรณจิตฺตํ อุทคฺคจิตฺตํ ปสนฺนจิตฺตํ.\csromanspacer{} อถ ยา พุทฺธานํ สามุกฺกํสิกา ธมฺมเทสนา, ตํ ปกาเสสิ ทุกฺขํ สมุทยํ นิโรธํ มคฺคํ.\csromanspacer{} เสยฺยถาปิ นาม สุทฺธํ วตฺถํ อปคตกาฬกํ สมฺมเทว รชนํ ปฏิคฺคเณฺหยฺย, เอวเมว พฺราหฺมายุสฺส พฺราหฺมณสฺส ตสฺมิํ เยว อาสเน วิรชํ วีตมลํ ธมฺมจกฺขุํ อุทปาทิ “ยํ กิญฺจิ สมุทยธมฺมํ, สพฺพํ ตํ นิโรธธมฺมนฺ”ติ.\csromanspacer{} อถ โข พฺรหฺมายุ พฺราหฺมโณ ทิฏฺฐธมฺโม ปตฺตธมฺโม วิทิตธมฺโม ปริโยคาฬฺหธมฺโม ติณฺณวิจิกิจฺโฉ วิคตกถํกโถ เวสารชฺชปฺปตฺโต อปรปฺปจฺจโย สตฺถุสาสเน ภควนฺตํ เอตทโวจ “อภิกฺกนฺตํ โภ โคตม, อภิกฺกนฺตํ โภ โคตม, เสยฺยถาปิ โภ โคตม นิกฺกุชฺชิตํ วา อุกฺกุชฺเชยฺย, ปฏิจฺฉนฺนํ วา วิวเรยฺย, มูฬฺหสฺส วา มคฺคํ อาจิกฺเขยฺย, อนฺธกาเร วา เตลปชฺโชตํ ธาเรยฺย ‘จกฺขุมนฺโต รูปานิ ทกฺขนฺตี’ติ, เอวเมวํ โภตา โคตเมน อเนกปริยาเยน ธมฺโม ปกาสิโต, เอสาหํ ภวนฺตํ โคตมํ สรณํ คจฺฉามิ ธมฺมญฺจ ภิกฺขุสํฆญฺจ, อุปาสกํ มํ ภวํ โคตโม ธาเรตุ อชฺชตคฺเค ปาณุเปตํ สรณํ คตํ, อธิวาเสตุ จ เม ภวํ โคตโม สฺวาตนาย ภตฺตํ สทฺธิํ ภิกฺขุสํเฆนา”ติ.\csromanspacer{} อธิวาเสสิ ภควา ตุณฺหีภาเวน.\csromanspacer{} อถ โข พฺรหฺมายุ พฺราหฺมโณ ภควโต อธิวาสนํ วิทิตฺวา อุฏฺฐายาสนา ภควนฺตํ อภิวาเทตฺวา ปทกฺขิณํ กตฺวา ปกฺกามิ.\csromanspacer{} อถ โข พฺรหฺมายุ พฺราหฺมโณ ตสฺสา รตฺติยา อจฺจเยน สเก นิเวสเน ปณีตํ ขาทนียํ โภชนียํ ปฏิยาทาเปตฺวา ภควโต กาลํ อาโรจาเปสิ “กาโล โภ โคตม นิฏฺฐิตํ ภตฺตนฺ”ติ.}

\prose{อถ โข ภควา ปุพฺพณฺหสมยํ นิวาเสตฺวา ปตฺตจีวรมาทาย เยน พฺรหฺมายุสฺส พฺราหฺมณสฺส นิเวสนํ เตนุปสงฺกมิ, อุปสงฺกมิตฺวา ปญฺญตฺเต อาสเน นิสีทิ สทฺธิํ ภิกฺขุสํเฆน.\csromanspacer{} อถ โข พฺรหฺมายุ พฺราหฺมโณ สตฺตาหํ พุทฺธปฺปมุขํ ภิกฺขุสํฆํ ปณีเตน ขาทนีเยน โภชนีเยน สหตฺถา สนฺตปฺเปสิ สมฺปวาเรสิ.\csromanspacer{} อถ โข ภควา ตสฺส สตฺตาหสฺส อจฺจเยน วิเทเหสุ จาริกํ ปกฺกามิ.\csromanspacer{} อถ โข พฺรหฺมายุ พฺราหฺมโณ อจิรปกฺกนฺตสฺส ภควโต กาลมกาสิ.\csromanspacer{} อถ โข}

\vspace{\tipitakaparskip}
\csromancenter{เสลสุตฺต (92) สมฺพหุลา ภิกฺขู เยน ภควา เตนุปสงฺกมิํสุ, อุปสงฺกมิตฺวา ภควนฺตํ อภิวาเทตฺวา เอกมนฺตํ นิสีทิํสุ, เอกมนฺตํ นิสินฺนา โข เต ภิกฺขู ภควนฺตํ เอตทโวจุํ “พฺรหฺมายุ ภนฺเต พฺราหฺมโณ กาลงฺกโต, ตสฺส กา คติ, โก อภิสมฺปราโย”ติ.\csromanspacer{} ปณฺฑิโต ภิกฺขเว พฺรหฺมายุ พฺราหฺมโณ ปจฺจปาทิ ธมฺมสฺสานุธมฺมํ, น จ มํ ธมฺมาธิกรณํ วิเหเสสิ, พฺรหฺมายุ ภิกฺขเว พฺราหฺมโณ ปญฺจนฺนํ โอรมฺภาคิยานํ สํโยชนานํ ปริกฺขยา โอปปาติโก โหติ ตตฺถ ปรินิพฺพายี, อนาวตฺติธมฺโม ตสฺมา โลกาติ.}
\csromancenter{อิทมโวจ ภควา.\csromanspacer{} อตฺตมนา เต ภิกฺขู ภควโต ภาสิตํ อภินนฺทุนฺติ.}
\nitthitamruled{พฺรหฺมายุสุตฺตํ นิฏฺฐิตํ ปฐมํ.}
\csromanheader{2}{2. เสลสุตฺต}
\proseitem{396}{\csromanfolio{347}เอวํ เม สุตํ\csromandash{}เอกํ สมยํ ภควา องฺคุตฺตราเปสุ จาริกํ จรมาโน มหตา ภิกฺขุสํเฆน สทฺธิํ อฑฺฒเตฬเสหิ ภิกฺขุสเตหิ.\csromanspacer{} เยน อาปณํ นาม องฺคุตฺตราปานํ นิคโม ตทวสริ.\csromanspacer{} อสฺโสสิ โข เกณิโย ชฏิโล “สมโณ ขลุ โภ โคตโม สกฺยปุตฺโต สกฺยกุลา ปพฺพชิโต องฺคุตฺตราเปสุ จาริกํ จรมาโน มหตา ภิกฺขุสํเฆน สทฺธิํ อฑฺฒเตฬเสหิ ภิกฺขุสเตหิ อาปณํ อนุปฺปตฺโต.\csromanspacer{} ตํ โข ปน ภวนฺตํ โคตมํ เอวํ กลฺยาโณ กิตฺติสทฺโท อพฺภุคฺคโต ‘อิติปิ โส ภควา อรหํ สมฺมาสมฺพุทฺโธ วิชฺชาจรณสมฺปนฺโน สุคโต โลกวิทู อนุตฺตโร ปุริสทมฺมสารถิ สตฺถา เทวมนุสฺสานํ พุทฺโธ ภควา’ติ, โส อิมํ โลกํ สเทวกํ สมารกํ สพฺรหฺมกํ สสฺสมณพฺราหฺมณิํ ปชํ สเทวมนุสฺสํ สยํ อภิญฺญา สจฺฉิกตฺวา ปเวเทติ, โส ธมฺมํ เทเสติ อาทิกลฺยาณํ มชฺเฌกลฺยาณํ ปริโยสานกลฺยาณํ สาตฺถํ สพฺยญฺชนํ เกวลปริปุณฺณํ ปริสุทฺธํ พฺรหฺมจริยํ ปกาเสติ, สาธุ โข ปน ตถารูปานํ อรหตํ ทสฺสนํ โหตี”ติ.}

\prose{อถ โข เกณิโย ชฏิโล เยน ภควา เตนุปสงฺกมิ, อุปสงฺกมิตฺวา ภควตา สทฺธิํ สมฺโมทิ, สมฺโมทนียํ กถํ สารณียํ \csromanfolio{348}วีติสาเรตฺวา เอกมนฺตํ นิสีทิ, เอกมนฺตํ นิสินฺนํ โข เกณิยํ ชฏิลํ ภควา ธมฺมิยา กถาย สนฺทสฺเสสิ สมาทเปสิ สมุตฺเตเชสิ สมฺปหํเสสิ.\csromanspacer{} อถ โข เกณิโย ชฏิโล ภควตา ธมฺมิยา กถาย สนฺทสฺสิโต สมาทปิโต สมุตฺเตชิโต สมฺปหํสิโต ภควนฺตํ เอตทโวจ “อธิวาเสตุ เม ภวํ โคตโม สฺวาตนาย ภตฺตํ สทฺธิํ ภิกฺขุสํเฆนา”ติ.\csromanspacer{} เอวํ วุตฺเต ภควา เกณิยํ ชฏิกํ เอตทโวจ “มหา โข เกณิย ภิกฺขุสํโฆ อฑฺฒเตฬสานิ ภิกฺขุสตานิ, ตฺวญฺจ พฺราหฺมเณสุ อภิปฺปสนฺโน”ติ.\csromanspacer{} ทุติยมฺปิ โข เกณิโย ชฏิโล ภควนฺตํ เอตทโวจ “กิญฺจาปิ โข โภ โคตม มหา ภิกฺขุสํโฆ อฑฺฒเตฬสานิ ภิกฺขุสตานิ, อหญฺจ พฺราหฺมเณสุ อภิปฺปสนฺโน, อธิวาเสตุ เม ภวํ โคตโม สฺวาตนาย ภตฺตํ สทฺธิํ ภิกฺขุสํเฆนา”ติ.\csromanspacer{} ทุติยมฺปิ โข ภควา เกณิยํ ชฏิลํ เอตทโวจ “มหา โข เกณิย ภิกฺขุสํโฆ อฑฺฒเตฬสานิ ภิกฺขุสตานิ, ตฺวญฺจ พฺราหฺมเณสุ อภิปฺปสนฺโน”ติ.\csromanspacer{} ตติยมฺปิ โข เกณิโย ชฏิโล ภควนฺตํ เอตทโวจ “กิญฺจาปิ โข โภ โคตม มหา ภิกฺขุสํโฆ อฑฺฒเตฬสานิ ภิกฺขุสตานิ, อหญฺจ พฺราหฺมเณสุ อภิปฺปสนฺโน, อธิวาเสตุ เม ภวํ โคตโม สฺวาตนาย ภตฺตํ สทฺธิํ ภิกฺขุสํเฆนา”ติ.\csromanspacer{} อธิวาเสสิ ภควา ตุณฺหีภาเวน.\csromanspacer{} อถ โข เกณิโย ชฏิโล ภควโต อธิวาสนํ วิทิตฺวา อุฏฺฐายาสนา เยน สโก อสฺสโม เตนุปสงฺกมิ, อุปสงฺกมิตฺวา มิตฺตามจฺเจ ญาติสาโลหิเต อามนฺเตสิ “สุณนฺตุ เม โภนฺโต มิตฺตามจฺจา ญาติสาโลหิตา, สมโณ เม โคตโม นิมนฺติโต สฺวาตนาย ภตฺตํ สทฺธิํ ภิกฺขุสํเฆน, เยน เม กายเวยฺยาวฏิกํ\footnote{กายเวยาวฏฺฏิกํ (สี, สฺยา, กํ), กายเวยฺยาวติกํ (ก)} กเรยฺยาถา”ติ. “เอวํ โภ”ติ โข เกณิยสฺส ชฏิลสฺส มิตฺตามจฺจา ญาติสาโลหิตา เกณิยสฺส ชฏิลสฺส ปฏิสฺสุตฺวา อปฺเปกจฺเจ อุทฺธนานิ ขณนฺติ, อปฺเปกจฺเจ กฏฺฐานิ ผาเลนฺติ, อปฺเปกจฺเจ ภาชนานิ โธวนฺติ, อปฺเปกจฺเจ อุทกมณิกํ ปติฏฺฐาเปนฺติ, อปฺเปกจฺเจ อาสนานิ ปญฺญเปนฺติ.\csromanspacer{} เกณิโย ปน ชฏิโล สามํเยว มณฺฑลมาลํ ปฏิยาเทติ.}

\proseitem{397}{เตน โข ปน สมเยน เสโล พฺราหฺมโณ อาปเณ ปฏิวสติ, ติณฺณํ เวทานํ ปารคู สนิฆณฺฑุเกฏุภานํ สากฺขรปฺปเภทานํ}

\vspace{\tipitakaparskip}
\csromancenter{เสลสุตฺต (92) อิติหาสปญฺจมานํ ปทโก เวยฺยากรโณ โลกายตมหาปุริสลกฺขเณสุ อนวโย, ตีณิ จ มาณวกสตานิ มนฺเต วาเจติ.\csromanspacer{} เตน โข ปน สมเยน เกณิโย ชฏิโล เสเล พฺราหฺมเณ อภิปฺปสนฺโน โหติ.\csromanspacer{} อถ โข เสโล พฺราหฺมโณ ตีหิ มาณวกสเตหิ ปริวุโต ชงฺฆาวิหารํ อนุจงฺกมมาโน อนุวิจรมาโน เยน เกณิยสฺส ชฏิลสฺส อสฺสโม เตนุปสงฺกมิ.\csromanspacer{} อทฺทสา โข เสโล พฺราหฺมโณ เกณิยสฺส ชฏิลสฺส อสฺสเม อปฺเปกจฺเจ อุทฺธนานิ ขณนฺเต อปฺเปกจฺเจ กฏฺฐานิ ผาเลนฺเต อปฺเปกจฺเจ ภาชนานิ โธวนฺเต อปฺเปกจฺเจ อุทกมณิกํ ปติฏฺฐาเปนฺเต อปฺเปกจฺเจ อาสนานิ ปญฺญเปนฺเต.\csromanspacer{} เกณิยํ ปน ชฏิลํ สามํเยว มณฺฑลมาลํ ปฏิยาเทนฺตํ.\csromanspacer{} ทิสฺวาน เกณิยํ ชฏิลํ เอตทโวจ “กิํ นุ โภโต เกณิยสฺส อาวาโห วา ภวิสฺสติ, วิวาโห วา ภวิสฺสติ, มหายญฺโญ วา ปจฺจุปฏฺฐิโต, ราชา วา มาคโธ เสนิโย พิมฺพิสาโร นิมนฺติโต สฺวาตนาย สทฺธิํ พลกาเยนา”ติ.\csromanspacer{} น เม โภ เสล อาวาโห ภวิสฺสติ, นปิ วิวาโห ภวิสฺสติ, นปิ ราชา มาคโธ เสนิโย พิมฺพิสาโร นิมนฺติโต สฺวาตนาย สทฺธิํ พลกาเยน.\csromanspacer{} อปิ จ โข เม มหายญฺโญ ปจฺจุปฏฺฐิโต, อตฺถิ โภ สมโณ โคตโม สกฺยปุตฺโต สกฺยกุลา ปพฺพชิโต, องฺคุตฺตราเปสุ จาริกํ จรมาโน มหตา ภิกฺขุสํเฆน สทฺธิํ อฑฺฒเตฬเสหิ ภิกฺขุสเตหิ อาปณํ อนุปฺปตฺโต.\csromanspacer{} ตํ โข ปน ภวนฺตํ โคตมํ เอวํ กลฺยาโณ กิตฺติสทฺโท อพฺภุคฺคโต “อิติปิ โส ภควา อรหํ สมฺมาสมฺพุทฺโธ วิชฺชาจรณสมฺปนฺโน สุคโต โลกวิทู อนุตฺตโร ปุริสทมฺมสารถิ สตฺถา เทวมนุสฺสานํ พุทฺโธ ภควา”ติ.\csromanspacer{} โส เม นิมนฺติโต สฺวาตนาย ภตฺตํ สทฺธิํ ภิกฺขุสํเฆนาติ. “พุทฺโธ”ติ โภ เกณิเย วเทสิ. “พุทฺโธ”ติ โภ เสล วทามิ. “พุทฺโธ”ติ โภ เกณิย วเทสิ. “พุทฺโธ”ติ โภ เสล วทามีติ.}
\proseitem{398}{\csromanfolio{349}อถ โข เสลสฺส พฺราหฺมณสฺส เอตทโหสิ “โฆโสปิ โข เอโส ทุลฺลโภ โลกสฺมิํ ยทิทํ พุทฺโธ”ติ\footnote{ยทิทํ พุทฺโธ พุทฺโธติ (ก)}.\csromanspacer{} อาคตานิ โข ปนมฺหากํ มนฺเตสุ ทฺวตฺติํสมหาปุริสลกฺขณานิ, เยหิ สมนฺนาคตสฺส \csromanfolio{350}มหาปุริสสฺส เทฺวเยว คติโย ภวนฺติ อนญฺญา.\csromanspacer{} สเจ อคารํ อชฺฌาวสติ, ราชา โหติ จกฺกวตฺตี ธมฺมิโก ธมฺมราชา จาตุรนฺโต วิชิตาวี ชปนทตฺถาวริยปฺปตฺโต สตฺตรตนสมนฺนาคโต.\csromanspacer{} ตสฺสิมานิ สตฺตรตนานิ ภวนฺติ.\csromanspacer{} เสยฺยถิทํ, จกฺกรตนํ หตฺถิรตนํ อสฺสรตนํ มณิรตนํ อิตฺถิรตนํ คหปติรตนํ ปริณายกรตนเมว สตฺตมํ.\csromanspacer{} ปโรสหสฺสํ โข ปนสฺส ปุตฺตา ภวนฺติ สุรา วีรงฺครูปา ปรเสนปฺปมทฺทนา, โส อิมํ ปถวิํ สาครปริยนฺตํ อทณฺเฑน อสตฺเถน ธมฺเมน อภิวิชิย อชฺฌาวสติ.\csromanspacer{} สเจ ปน อคารสฺมา อนคาริยํ ปพฺพชติ, อรหํ โหติ สมฺมาสมฺพุทฺโธ โลเก วิวฏฺฏจฺฉโท.}

\prose{กหํ ปน โภ เกณิย เอตรหิ โส ภวํ โคตโม วิหรติ อรหํ สมฺมาสมฺพุทฺโธติ.\csromanspacer{} เอวํ วุตฺเต เกณิโย ชฏิโล ทกฺขิณํ พาหุํ ปคฺคเหตฺวา เสลํ พฺราหฺมณํ เอตทโวจ “เยเนสา โภ เสล นีลวนราชี”ติ.\csromanspacer{} อถ โข เสโล พฺราหฺมโณ ตีหิ มาณวกสเตหิ สทฺธิํ เยน ภควา เตนุปสงฺกมิ.\csromanspacer{} อถ โข เสโล พฺราหฺมโณ เต มาณวเก อามนฺเตสิ “อปฺปสทฺทา โภนฺโต อาคจฺฉนฺตุ ปเท ปทํ\footnote{ปาเท ปาทํ (สี)} นิกฺขิปนฺตา, ทุราสทา\footnote{ทูรสทฺทา (ก)} หิ เต ภควนฺโต, สีหาว เอกจรา, ยทา จาหํ โภ สมเณน โคตเมน สทฺธิํ มนฺเตยฺยํ ‘มา เม โภนฺโต อนฺตรนฺตรา กถํ โอปาเตถ, กถาปริโยสานํ เม ภวนฺโต อาคเมนฺตู’ติ”.\csromanspacer{} อถ โข เสโล พฺราหฺมโณ เยน ภควา เตนุปสงฺกมิ, อุปสงฺกมิตฺวา ภควตา สทฺธิํ สมฺโมทิ, สมฺโมทนียํ กถํ สารณียํ วีติสาเรตฺวา เอกมนฺตํ นิสีทิ, เอกมนฺตํ นิสินฺโน โข เสโล พฺราหฺมโณ ภควโต กาเย ทฺวตฺติํสมหาปุริสลกฺขณานิ สมนฺเนสิ.}

\prose{อทฺทสา โข เสโล พฺราหฺมโณ ภควโต กาเย ทฺวตฺติํสมหาปุริสลกฺขณานิ เยภุยฺเยน ฐเปตฺวา เทฺว, ทฺวีสุ มหาปุริสลกฺขเณสุ กงฺขติ วิจิกิจฺฉติ นาธิมุจฺจติ น สมฺปสีทติ โกโสหิเต จ วตฺถคุเยฺห ปหูตชิวฺหตาย จ.\csromanspacer{} อถ โข ภควโต เอตทโหสิ “ปสฺสติ โข เม อยํ เสโล พฺราหฺมโณ ทฺวตฺติํสมหาปุริสลกฺขณานิ เยภุยฺเยน ฐเปตฺวา เทฺว, ทฺวีสุ มหาปุริสลกฺขเณสุ กงฺขติ วิจิกิจฺฉติ นาธิมุจฺจติ น สมฺปสีทติ}

\vspace{\tipitakaparskip}
\csromancenter{เสลสุตฺต (92) โกโสหิเต จ วตฺถคุเยฺห ปหูตชิวฺหตาย จา”ติ.\csromanspacer{} อถ โข ภควา ตถารูปํ อิทฺธาภิสงฺขารํ อภิสงฺขาสิ, ยถา อทฺทส เสโล พฺราหฺมโณ ภควโต โกโสหิตํ วตฺถคุยฺหํ.\csromanspacer{} อถ โข ภควา ชิวฺหํ นินฺนาเมตฺวา อุโภปิ กณฺณโสตานิ อนุมสิ ปฏิมสิ, อุโภปิ นาสิกโสตานิ อนุมสิ ปฏิมสิ, เกวลมฺปิ นลาฏมณฺฑลํ ชิวฺหาย ฉาเทสิ.\csromanspacer{} อถ โข เสลสฺส พฺราหฺมณสฺส เอตทโหสิ “สมนฺนาคโต โข สมโณ โคตโม ทฺวตฺติํสมหาปุริสลกฺขเณหิ ปริปุณฺเณหิ โน อปริปุณฺเณหิ, โน จ โข นํ ชานามิ ‘พุทฺโธ วา โน วา’, สุตํ โข ปน เมตํ พฺราหฺมณานํ วุทฺธานํ มหลฺลกานํ อาจริยปาจริยานํ ภาสมานานํ ‘เย เต ภวนฺติ อรหนฺโต สมฺมาสมฺพุทฺธา, เต สเก วณฺเณ ภญฺญามาเน อตฺตานํ ปาตุกโรนฺตี”ติ, ยํนูนาหํ สมณํ โคตมํ สมฺมุขา สารุปฺปาหิ คาถาหิ อภิตฺถเวยฺยนฺ”ติ.}
\proseitem{399}{\csromanfolio{351}อถ โข เสลา พฺราหฺมโณ ภควนฺตํ สมฺมุขา สารุปฺปาหิ คาถาหิ อภิตฺถวิ\csromandash{}}

\csromanmatikaanchor{mtk.51}
\csromantocmarkref{subsection}{2. เสลสุตฺต}{mtk.51}
\bookmark[dest={mtk.51},level=2]{2. เสลสุตฺต}
\csromangathameasure{ปริปุณฺณกาโย สุรุจิ, สุชาโต จารุทสฺสโน. \\ สุวณฺณวณฺโณสิ ภควา, สุสุกฺกทาโฐสิ วีริยวา. \\ นรสฺส หิ สุชาตสฺส, เย ภวนฺติ วิยญฺชนา. \\ สพฺเพ เต ตว กายสฺมิํ, มหาปุริสลกฺขณา. \\ ปสนฺนเนตฺโต สุมุโข, พฺรหา อุชุ ปตาปวา. \\ มชฺเฌ สมณสํฆสฺส, อาทิจฺโจว วิโรจสิ. \\ กลฺยาณทสฺสโน ภิกฺขุ, กญฺจนสนฺนิภตฺตโจ. \\ กิํ เต สมณภาเวน, เอวํ อุตฺตมวณฺณิโน. \\ ราชา อรหสิ ภวิตุํ, จกฺกวตฺตี รเถสโภ. \\ จาตุรนฺโต วิชิตาวี, ชมฺพุสณฺฑสฺส อิสฺสโร. \\ ขตฺติยา โภคิราชาโน, อนุยนฺตา ภวนฺตุ เต. \\ ราชาภิราชา มนุชินฺโท, รชฺชํ กาเรหิ โคตม. \\ ราชาหมสฺมิ เสลาติ, ธมฺมราชา อนุตฺตโร. \\ ธมฺเมน จกฺกํ วตฺเตมิ, จกฺกํ อปฺปฏิวตฺติยํ. \\ สมฺพุทฺโธ ปฏิชานาสิ, ธมฺมราชา อนุตฺตโร. \\ ธมฺเมน จกฺกํ วตฺเตมิ, อิติ ภาสสิ โคตม. \\ โก นุ เสนาปติ โภโต, สาวโก สตฺถุรนฺวโย. \\ โก เตตมนุวตฺเตติ, ธมฺมจกฺกํ ปวตฺติตํ. \\ มยา ปวตฺติตํ จกฺกํ, (เสลาติ ภควา) ธมฺมจกฺกํ อนุตฺตรํ. \\ สาริปุตฺโต อนุวตฺเตติ, อนุชาโต ตถาคตํ. \\ อภิญฺเญยฺยํ อภิญฺญาตํ, ภาเวตพฺพญฺจ ภาวิตํ. \\ ปหาตพฺพํ ปหีนํ เม, ตสฺมา พุทฺโธสฺมิ พฺราหฺมณ. \\ วินยสฺสุ มยิ กงฺขํ, อธิมุจฺจสฺสุ พฺราหฺมณ. \\ ทุลฺลภํ ทสฺสนํ โหติ, สมฺพุทฺธานํ อภิณฺหโส. \\ เยสํ เว ทุลฺลโภ โลเก, ปาตุภาโว อภิณฺหโส. \\ โสหํ พฺราหฺมณ สมฺพุทฺโธ, สลฺลกตฺโต อนุตฺตโร. \\ พฺรหฺมภูโต อติตุโล, มารเสนปฺปมทฺทโน. \\ สพฺพามิตฺเต วสี กตฺวา, โมทามิ อกุโตภโย. \\ อิมํ โภนฺโต นิสาเมถ, ยถา ภาสติ จกฺขุมา. \\ สลฺลกตฺโต มหาวีโร, สีโหว นทตี วเน. \\ พฺรหฺมภูตํ อติตุลํ, มารเสนปฺปมทฺทนํ. \\ โก ทิสฺวา นปฺปสีเทยฺย, อปิ กณฺหาภิชาติโก. \\ โย มํ อิจฺฉติ อเนฺวตุ, โย วา นิจฺฉติ คจฺฉตุ. \\ อิธาหํ ปพฺพชิสฺสามิ, วรปญฺญสฺส สนฺติเก. \\ เอตญฺเจ รุจฺจติ โภโต, สมฺมาสมฺพุทฺธสาสนํ. \\ มยมฺปิ ปพฺพชิสฺสาม, วรปญฺญสฺส สนฺติเก. \\ พฺราหฺมณา ติสตา อิเม, ยาจนฺติ ปญฺชลีกตา. \\ พฺรหฺมจริยํ จริสฺสาม, ภควา ตว สนฺติเก.}
\csromangathagroup{\csromangathastanza{ปริปุณฺณกาโย สุรุจิ, สุชาโต จารุทสฺสโน. \\ สุวณฺณวณฺโณสิ ภควา, สุสุกฺกทาโฐสิ วีริยวา\footnote{วิริยวา (สี, สฺยา, กํ, อิ)}.}\csromangathastanza{นรสฺส หิ สุชาตสฺส, เย ภวนฺติ วิยญฺชนา. \\ สพฺเพ เต ตว กายสฺมิํ, มหาปุริสลกฺขณา.}\csromangathastanza{ปสนฺนเนตฺโต สุมุโข, พฺรหา\footnote{พฺรหฺมา (สฺยา, กํ, ก)} อุชุ ปตาปวา. \\ มชฺเฌ สมณสํฆสฺส, อาทิจฺโจว วิโรจสิ.}\csromangathastanza{กลฺยาณทสฺสโน ภิกฺขุ, กญฺจนสนฺนิภตฺตโจ. \\ กิํ เต สมณภาเวน, เอวํ อุตฺตมวณฺณิโน.}\csromangathastanza{ราชา อรหสิ ภวิตุํ, จกฺกวตฺตี รเถสโภ. \\ จาตุรนฺโต วิชิตาวี, ชมฺพุสณฺฑสฺส\footnote{ชมฺพุมณฺฑสฺส (ก)} อิสฺสโร.}\csromangathastanza{ขตฺติยา โภคิราชาโน, อนุยนฺตา\footnote{อนุยุตฺตา (สี, สฺยา, กํ, อิ)} ภวนฺตุ เต. \\ ราชาภิราชา มนุชินฺโท, รชฺชํ กาเรหิ โคตม.}\csromangathastanza{\csromanfolio{352}ราชาหมสฺมิ เสลาติ, ธมฺมราชา อนุตฺตโร. \\ ธมฺเมน จกฺกํ วตฺเตมิ, จกฺกํ อปฺปฏิวตฺติยํ.}\csromangathastanza{สมฺพุทฺโธ ปฏิชานาสิ, ธมฺมราชา อนุตฺตโร. \\ ธมฺเมน จกฺกํ วตฺเตมิ, อิติ ภาสสิ โคตม.}\csromangathastanza{โก นุ เสนาปติ โภโต, สาวโก สตฺถุรนฺวโย. \\ โก เตตมนุวตฺเตติ, ธมฺมจกฺกํ ปวตฺติตํ.}\csromangathastanza{มยา ปวตฺติตํ จกฺกํ, (เสลาติ ภควา) ธมฺมจกฺกํ อนุตฺตรํ. \\ สาริปุตฺโต อนุวตฺเตติ, อนุชาโต ตถาคตํ.}\csromangathastanza{อภิญฺเญยฺยํ อภิญฺญาตํ, ภาเวตพฺพญฺจ ภาวิตํ. \\ ปหาตพฺพํ ปหีนํ เม, ตสฺมา พุทฺโธสฺมิ พฺราหฺมณ.}\csromangathastanza{วินยสฺสุ มยิ กงฺขํ, อธิมุจฺจสฺสุ พฺราหฺมณ. \\ ทุลฺลภํ ทสฺสนํ โหติ, สมฺพุทฺธานํ อภิณฺหโส.}\csromangathastanza{เยสํ เว ทุลฺลโภ โลเก, ปาตุภาโว อภิณฺหโส. \\ โสหํ พฺราหฺมณ สมฺพุทฺโธ, สลฺลกตฺโต อนุตฺตโร.}\csromangathastanza{พฺรหฺมภูโต อติตุโล, มารเสนปฺปมทฺทโน. \\ สพฺพามิตฺเต วสี กตฺวา, โมทามิ อกุโตภโย.}\csromangathastanza{อิมํ โภนฺโต นิสาเมถ, ยถา ภาสติ จกฺขุมา. \\ สลฺลกตฺโต มหาวีโร, สีโหว นทตี วเน.}\csromangathastanza{พฺรหฺมภูตํ อติตุลํ, มารเสนปฺปมทฺทนํ. \\ โก ทิสฺวา นปฺปสีเทยฺย, อปิ กณฺหาภิชาติโก.}\csromangathastanza{โย มํ อิจฺฉติ อเนฺวตุ, โย วา นิจฺฉติ คจฺฉตุ. \\ อิธาหํ ปพฺพชิสฺสามิ, วรปญฺญสฺส สนฺติเก.}\csromangathastanza{เอตญฺเจ\footnote{เอวญฺเจ (สฺยา, กํ)} รุจฺจติ โภโต, สมฺมาสมฺพุทฺธสาสนํ\footnote{สมฺมาสมฺพุทฺธสาสเน (กตฺถจิ สุตฺตนิปาเต)}. \\ มยมฺปิ ปพฺพชิสฺสาม, วรปญฺญสฺส สนฺติเก.}\csromangathastanza{พฺราหฺมณา ติสตา อิเม, ยาจนฺติ ปญฺชลีกตา. \\ พฺรหฺมจริยํ จริสฺสาม, ภควา ตว สนฺติเก.}}

\csromanheader{2}{2. เสลสุตฺต (92)}
\csromangathameasure{สฺวากฺขาตํ พฺรหฺมจริยํ, (เสลาติ ภควา) สนฺติฏฺฐิกมกาลิกํ. \\ ยตฺถ อโมฆา ปพฺพชฺชา, อปฺปมตฺตสฺส สิกฺขโตติ.}
\csromangathagroup{\csromangathastanza{\csromanfolio{353}สฺวากฺขาตํ พฺรหฺมจริยํ, (เสลาติ ภควา) สนฺติฏฺฐิกมกาลิกํ. \\ ยตฺถ อโมฆา ปพฺพชฺชา, อปฺปมตฺตสฺส สิกฺขโตติ.}}

\prose{อลตฺถ โข เสโล พฺราหฺมโณ สปริโส ภควโต สนฺติเก ปพฺพชฺชํ, อลตฺถ อุปสมฺปทํ.}

\proseitem{400}{อถ โข เกณิโย ชฏิโล ตสฺสา รตฺติยา อจฺจเยน สเก อสฺสเม ปณีตํ ขาทนียํ โภชนียํ ปฏิยาทาเปตฺวา ภควโต กาลํ อาโรจาเปสิ “กาโล โภ โคตม นิฏฺฐิตํ ภตฺตนฺ”ติ.\csromanspacer{} อถ โข ภควา ปุพฺพณฺหสมยํ นิวาเสตฺวา ปตฺตจีวรมาทาย เยน เกณิยสฺส ชฏิลสฺส อสฺสโม เตนุปสงฺกมิ, อุปสงฺกมิตฺวา ปญฺญตฺเต อาสเน นิสีทิ สทฺธิํ ภิกฺขุสํเฆน.\csromanspacer{} อถ โข เกณิโย ชฏิโล พุทฺธปฺปมุขํ ภิกฺขุสํฆํ ปณีเตน ขาทนีเยน โภชนีเยน สหตฺถา สนฺตปฺเปสิ สมฺปวาเรสิ.\csromanspacer{} อถ โข เกณิโย ชฏิโล ภควนฺตํ ภุตฺตาวิํ โอนีตปตฺตปาณิํ อญฺญตรํ นีจํ อาสนํ คเหตฺวา เอกมนฺตํ นิสีทิ, เอกมนฺตํ นิสินฺนํ โข เกณิยํ ชฏิลํ ภควา อิมาหิ คาถาหิ อนุโมทิ\csromandash{}}

\csromangathameasure{“อคฺคิหุตฺตมุขา ยญฺญา, สาวิตฺตี ฉนฺทโส มุขํ. \\ ราชา มุขํ มนุสฺสานํ, นทีนํ สาคโร มุขํ. \\ นกฺขตฺตานํ มุขํ จนฺโท, อาทิจฺโจ ตปตํ มุขํ. \\ ปุญฺญํ อากงฺขมานานํ, สํโฆ เว ยชตํ มุขนฺ”ติ.}
\csromangathagroup{\csromangathastanza{“อคฺคิหุตฺตมุขา ยญฺญา, สาวิตฺตี ฉนฺทโส มุขํ. \\ ราชา มุขํ มนุสฺสานํ, นทีนํ สาคโร มุขํ.}\csromangathastanza{นกฺขตฺตานํ มุขํ จนฺโท, อาทิจฺโจ ตปตํ มุขํ. \\ ปุญฺญํ อากงฺขมานานํ, สํโฆ เว ยชตํ มุขนฺ”ติ.}}

\prose{อถ โข ภควา เกณิยํ ชฏิลํ อิมาหิ คาถาหิ อนุโมทิตฺวา อุฏฺฐายาสนา ปกฺกามิ.}

\prose{อถ โข อายสฺมา เสโล สปริโส เอโก วูปกฏฺโฐ อปฺปมตฺโต อาตาปี ปหิตตฺโต วิหรนฺโต นจิรสฺเสว, ยสฺสตฺถาย กุลปุตฺตา สมฺมเทว อคารสฺมา อนคาริยํ ปพฺพชนฺติ, ตทนุตฺตรํ พฺรหฺมจริยปริโยสานํ ทิฏฺเฐว ธมฺเม สยํ อภิญฺญา สจฺฉิกตฺวา อุปสมฺปชฺช วิหาสิ, “ขีณา ชาติ, วุสิตํ พฺรหฺมจริยํ, กตํ กรณียํ, นาปรํ อิตฺถตฺตายา”ติ อพฺภญฺญาสิ.\csromanspacer{} อญฺญตโร โข ปนายสฺมา เสโล สปริโส อรหตํ อโหสิ.\csromanspacer{} อถ โข อายสฺมา เสโล สปริโส เยน ภควา เตนุปสงฺกมิ, อุปสงฺกมิตฺวา เอกํสํ จีวรํ กตฺวา เยน ภควา เตนญฺชลิํ ปณาเมตฺวา ภควนฺตํ คาถาหิ อชฺฌภาสิ\csromandash{}}

\csromangathameasure{“ยํ ตํ สรณมาคมฺม, อิโต อฏฺฐมิ จกฺขุมา. \\ สตฺตรตฺเตน ภควา, ทนฺตมฺห ตว สาสเน. \\ ตุวํ พุทฺโธ ตุวํ สตฺถา, ตุวํ มาราภิภู มุนิ. \\ ตุวํ อนุสเย เฉตฺวา, ติณฺโณ ตาเรสิมํ ปชํ. \\ อุปธี เต สมติกฺกนฺตา, อาสวา เต ปทาลิตา. \\ สีโหว อนุปาทาโน, ปหีนภยเภรโว. \\ ภิกฺขโว ติสตา อิเม, ติฏฺฐนฺติ ปญฺชลีกตา. \\ ปาเท วีร ปสาเรหิ, นาคา วนฺทนฺตุ สตฺถุโน”ติ.}
\csromangathagroup{\csromangathastanza{\csromanfolio{354}“ยํ ตํ สรณมาคมฺม, อิโต อฏฺฐมิ จกฺขุมา. \\ สตฺตรตฺเตน\footnote{อนุตฺตเรน (ก)} ภควา, ทนฺตมฺห ตว สาสเน.}\csromangathastanza{ตุวํ พุทฺโธ ตุวํ สตฺถา, ตุวํ มาราภิภู มุนิ. \\ ตุวํ อนุสเย เฉตฺวา, ติณฺโณ ตาเรสิมํ ปชํ.}\csromangathastanza{อุปธี เต สมติกฺกนฺตา, อาสวา เต ปทาลิตา. \\ สีโหว อนุปาทาโน, ปหีนภยเภรโว.}\csromangathastanza{ภิกฺขโว ติสตา อิเม, ติฏฺฐนฺติ ปญฺชลีกตา. \\ ปาเท วีร ปสาเรหิ, นาคา วนฺทนฺตุ สตฺถุโน”ติ.}}

\nitthitamruled{เสลสุตฺตํ นิฏฺฐิตํ ทุติยํ.}
\csromanheader{2}{3. อสฺสลายนสุตฺต}
\proseitem{401}{เอวํ เม สุตํ\csromandash{}เอกํ สมยํ ภควา สาวตฺถิยํ วิหรติ เชตวเน อนาถปิณฺฑิกสฺส อาราเม.\csromanspacer{} เตน โข ปน สมเยน นานาเวรชฺชกานํ พฺราหฺมณานํ ปญฺจมตฺตานิ พฺราหฺมณสตานิ สาวตฺถิยํ ปฏิวสนฺติ เกนจิเทว กรณีเยน.\csromanspacer{} อถ โข เตสํ พฺราหฺมณานํ เอตทโหสิ “อยํ โข สมโณ โคตโม จาตุวณฺณิํ สุทฺธิํ ปญฺญเปติ, โก นุ โข ปโหติ สมเณน โคตเมน สทฺธิํ อสฺมิํ วจเน ปฏิมนฺเตตุนฺ”ติ.\csromanspacer{} เตน โข ปน สมเยน อสฺสลายโน นาม มาณโว สาวตฺถิยํ ปฏิวสติ ทหโร วุตฺตสิโร โสฬสวสฺสุทฺเทสิโก ชาติยา, ติณฺณํ เวทานํ ปารคู สนิฆณฺฑุเกฏุภานํ สากฺขรปฺปเภทานํ อิติหาสปญฺจมานํ ปทโก เวยฺยากรโณ โลกายตมหาปุริสลกฺขเณสุ อนวโย.\csromanspacer{} อถ โข เตสํ พฺราหฺมณานํ เอตทโหสิ “อยํ โข อสฺสลายโน มาณโว สาวตฺถิยํ ปฏิวสติ ทหโร วุตฺตสิโร โสฬสวสฺสุทฺเทสิโก ชาติยา, ติณฺณํ เวทานํ ปารคู ฯเปฯ อนวโย, โส โข ปโหติ สมเณน โคตเมน สทฺธิํ อสฺมิํ วจเน ปฏิมนฺเตตุนฺ”ติ.}

\prose{อถ โข เต พฺราหฺมณา เยน อสฺสลายโน มาณโว เตนุปสงฺกมิํสุ, อุปสงฺกมิตฺวา อสฺสลายนํ มาณวํ เอตทโวจุํ “อยํ โภ}

\vspace{\tipitakaparskip}
\csromancenter{อสฺสลายนสุตฺต (93) อสฺสลายน สมโณ โคตโม จาตุวณฺณิํ สุทฺธิํ ปญฺญเปติ, เอตุ ภวํ อสฺสลายโน สมเณน โคตเมน สทฺธิํ อสฺมิํ วจเน ปฏิมนฺเตตู”ติ\footnote{ปฏิมนฺเตตุนฺติ (อิ, ก)}.}
\prose{\csromanfolio{355}เอวํ วุตฺเต อสฺสลายโน มาณโว เต พฺราหฺมเณ เอตทโวจ “สมโณ ขลุ โภ โคตโม ธมฺมวาที, ธมฺมวาทิโน จ ปน ทุปฺปฏิมนฺติยา ภวนฺติ, นาหํ สกฺโกมิ สมเณน โคตเมน สทฺธิํ อสฺมิํ วจเน ปฏิมนฺเตตุนฺ”ติ.\csromanspacer{} ทุติยมฺปิ โข เต พฺราหฺมณา อสฺสลายนํ มาณวํ เอตทโวจุํ “อยํ โภ อสฺสลายน สมโณ โคตโม จาตุวณฺณิํ สุทฺธิํ ปญฺญเปติ, เอตุ ภวํ อสฺสลายโน สมเณน โคตเมน สทฺธิํ อสฺมิํ วจเน ปฏิมนฺเตตุ\footnote{ปฏิมนฺเตตุํ (สี, อิ, ก)}.\csromanspacer{} จริตํ โข ปน โภตา อสฺสลายเนน ปริพฺพาชกนฺ”ติ.\csromanspacer{} ทุติยมฺปิ โข อสฺสลายโน มาณโว เต พฺราหฺมเณ เอตทโวจ “สมโณ ขลุ โภ โคตโม ธมฺมวาที, ธมฺมวาทิโน จ ปน ทุปฺปฏิมนฺติยา ภวนฺติ, นาหํ สกฺโกมิ สมเณน โคตเมน สทฺธิํ อสฺมิํ วจเน ปฏิมนฺเตตุนฺ”ติ.\csromanspacer{} ตติยมฺปิ โข เต พฺราหฺมณา อสฺสลายนํ มาณวํ เอตทโวจุํ “อยํ โภ อสฺสลายน สมโณ โคตโม จาตุวณฺณิํ สุทฺธิํ ปญฺญเปติ, เอตุ ภวํ อสฺสลายโน สมเณน โคตเมน สทฺธิํ อสฺมิํ วจเน ปฏิมนฺเตตุ2.\csromanspacer{} จริตํ โข ปน โภตา อสฺสลายเนน ปริพฺพาชกํ, มา ภวํ อสฺสลายโน อยุทฺธปราชิตํ ปราชยี”ติ.}

\prose{เอวํ วุตฺเต อสฺสลายโน มาณโว เต พฺราหฺมเณ เอตทโวจ “อทฺธา โข อหํ ภวนฺโต น ลภามิ ‘สมโณ ขลุ โภ โคตโม ธมฺมวาที, ธมฺมวาทิโน จ ปน ทุปฺปฏิมนฺติยา ภวนฺติ, นาหํ สกฺโกมิ สมเณน โคตเมน สทฺธิํ อสฺมิํ วจเน ปฏิมนฺเตตุนฺ’ติ, อปิ จาหํ ภวนฺตานํ วจเนน คมิสฺสามี”ติ.}

\proseitem{402}{อถ โข อสฺสลายโน มาณโว มหตา พฺราหฺมณคเณน สทฺธิํ เยน ภควา เตนุปสงฺกมิ, อุปสงฺกมิตฺวา ภควตา สทฺธิํ สมฺโมทิ, สมฺโมทนียํ กถํ สารณียํ วีติสาเรตฺวา เอกมนฺตํ นิสีทิ, เอกมนฺตํ นิสินฺโน โข อสฺสาลายโน มาณโว ภควนฺตํ เอตทโวจ “พฺราหฺมณา โภ โคตม เอวมาหํสุ ‘พฺราหฺมโณว เสฏฺโฐ วณฺโณ, หีโน อญฺโญ วณฺโณ.\csromanspacer{} พฺราหฺมโณว สุกฺโก วณฺโณ, กโณฺห อญฺโญ วณฺโณ.\csromanspacer{} พฺราหฺมณาว สุชฺฌนฺติ โน อพฺราหฺมณา, พฺราหฺมณาว พฺรหฺมุโน \csromanfolio{356}ปุตฺตา โอรสา มุขโต ชาตา พฺรหฺมชา พฺรหฺมนิมฺมิตา พฺรหฺมทายาทา’ติ.\csromanspacer{} อิธ ภวํ โคตโม กิมาหา”ติ.\csromanspacer{} ทิสฺสนฺติ\footnote{ทิสฺสนฺเต (สี, สฺยา, กํ, อิ)} โข ปน อสฺสลายน พฺราหฺมณานํ พฺราหฺมณิโย อุตุนิโยปิ คพฺภินิโยปิ วิชายมานาปิ ปายมานาปิ, เต จ พฺราหฺมณิโยนิชาว สมานา เอวมาหํสุ “พฺราหฺมโณว เสฏฺโฐ วณฺโณ, หีโน อญฺโญ วณฺโณ.\csromanspacer{} พฺราหฺมโณว สุกฺโก วณฺโณ, กโณฺห อญฺโญ วณฺโณ.\csromanspacer{} พฺราหฺมณาว สุชฺฌนฺติ โน อพฺราหฺมณา, พฺราหฺมณาว พฺรหฺมุโน ปุตฺตา โอรสา มุขโต ชาตา พฺรหฺมชา พฺรหฺมนิมฺมิตา พฺรหฺมทายาทา”ติ.\csromanspacer{} กิญฺจาปิ ภวํ โคตโม เอวมาห, อถ เขฺวตฺถ พฺราหฺมณา เอวเมตํ มญฺญนฺติ “พฺราหฺมโณว เสฏฺโฐ วณฺโณ, หีโน อญฺโญ วณฺโณ ฯเปฯ พฺรหฺมทายาทา”ติ.}

\proseitem{403}{ตํ กิํ มญฺญสิ อสฺสลายน, สุตํ เต “โยนกมฺโพเชสุ อญฺเญสุ จ ปจฺจนฺติเมสุ ชนปเทสุ เทฺวว วณฺณา อยฺโย เจว ทาโส จ.\csromanspacer{} อยฺโย หุตฺวา ทาโส โหติ, ทาโส หุตฺวา อยฺโย โหตี”ติ.\csromanspacer{} เอวํ โภ สุตํ ตํ เม “โยนกมฺโพเชสุ อญฺเญสุ จ ปจฺจนฺติเมสุ ชนปเทสุ เทฺวว วณฺณา อยฺโย เจว ทาโส จ.\csromanspacer{} อยฺโย หุตฺวา ทาโส โหติ, ทาโส หุตฺวา อยฺโย โหตี”ติ.\csromanspacer{} เอตฺถ อสฺสลายน พฺราหฺมณานํ กิํ พลํ, โก อสฺสาโส, ยเทตฺถ พฺราหฺมณา เอวมาหํสุ “พฺราหฺมโณว เสฏฺโฐ วณฺโณ, หีโน อญฺโญ วณฺโณ ฯเปฯ พฺรหฺมทายาทา”ติ.\csromanspacer{} กิญฺจาปิ ภวํ โคตโม เอวมาห, อถ เขฺวตฺถ พฺราหฺมณา เอวเมตํ มญฺญนฺติ “พฺราหฺมโณว เสฏฺโฐ วณฺโณ, หีโน อญฺโญ วณฺโณ ฯเปฯ พฺรหฺมทายาทา”ติ.}

\proseitem{404}{ตํ กิํ มญฺญสิ อสฺสลายน, ขตฺติโยว นุ โข ปาณาติปาตี อทินฺนาทายี กาเมสุมิจฺฉาจารี มุสาวาที ปิสุณวาโจ ผรุสวาโจ สมฺผปฺปลาปี อภิชฺฌาลุ พฺยาปนฺนจิตฺโต มิจฺฉาทิฏฺฐิ กายสฺส เภทา ปรํ มรณา อปายํ ทุคฺคติํ วินิปาตํ นิรยํ อุปปชฺเชยฺย, โน พฺราหฺมโณ.\csromanspacer{} เวสฺโสว นุ โข ฯเปฯ.\csromanspacer{} สุทฺโทว นุ โข ปาณาติปาตี อทินฺนาทายี กาเมสุมิจฺฉาจารี มุสาวาที ปิสุณวาโจ ผรุสวาโจ สมฺผปฺปลาปี อภิชฺฌาลุ พฺยาปนฺนจิตฺโต มิจฺฉาทิฏฺฐิ กายสฺส เภทา ปรํ มรณา อปายํ ทุคฺคติํ วินิปาตํ นิรยํ อุปปชฺเชยฺย, โน พฺราหฺมโณติ.\csromanspacer{} โน หิทํ โภ โคตม, ขตฺติโยปิ หิ}

\vspace{\tipitakaparskip}
\csromancenter{อสฺสลายนสุตฺต (93) โภ โคตม ปาณาติปาตี อทินฺนาทายี กาเมสุมิจฺฉาจารี มุสาวาที ปิสุณวาโจ ผรุสวาโจ สมฺผปฺปลาปี อภิชฺฌาลุ พฺยาปนฺนจิตฺโต มิจฺฉาทิฏฺฐิ กายสฺส เภทา ปรํ มรณา อปายํ ทุคฺคติํ วินิปาตํ นิรยํ อุปปชฺเชยฺย.\csromanspacer{} พฺราหฺมโณปิ หิ โภ โคตม ฯเปฯ.\csromanspacer{} เวสฺโสปิ หิ โภ โคตม ฯเปฯ.\csromanspacer{} สุทฺโทปิ หิ โภ โคตม ฯเปฯ.\csromanspacer{} สพฺเพปิ หิ โภ โคตม จตฺตาโร วณฺณา ปาณาติปาติโน อทินฺนาทายิโน กาเมสุมิจฺฉาจาริโน มุสาวาทิโน ปิสุณวาจา ผรุสวาจา สมฺผปฺปลาปิโน อภิชฺฌาลู พฺยาปนฺนจิตฺตา มิจฺฉาทิฏฺฐี กายสฺส เภทา ปรํ มรณา อปายํ ทุคฺคติํ วินิปาตํ นิรยํ อุปปชฺเชยฺยุนฺติ.\csromanspacer{} เอตฺถ อสฺสลายน พฺราหฺมณานํ กิํ พลํ, โก อสฺสาโส, ยเทตฺถ พฺราหฺมณา เอวมาหํสุ “พฺราหฺมโณว เสฏฺโฐ วณฺโณ, หีโน อญฺโญ วณฺโณ ฯเปฯ พฺรหฺมทายาทา”ติ.\csromanspacer{} กิญฺจาปิ ภวํ โคตโม เอวมาห, อถ เขฺวตฺถ พฺราหฺมณา เอวเมตํ มญฺญนฺติ “พฺราหฺมโณว เสฏฺโฐ วณฺโณ, หีโน อญฺโญ วณฺโณ ฯเปฯ พฺรหฺมทายาทา”ติ.}
\proseitem{405}{\csromanfolio{357}ตํ กิํ มญฺญสิ อสฺสาลายน, พฺราหฺมโณว นุ โข ปาณาติปาตา ปฏิวิรโต อทินฺนาทานา ปฏิวิรโต กาเมสุมิจฺฉาจารา ปฏิวิรโต มุสาวาทา ปฏิวิรโต ปิสุณาย วาจาย ปฏิวิรโต ผรุสาย วาจาย ปฏิวิรโต สมฺผปฺปลาปา ปฏิวิรโต อนภิชฺฌาลุ อพฺยาปนฺนจิตฺโต สมฺมาทิฏฺฐิ กายสฺส เภทา ปรํ มรณา สุคติํ สคฺคํ โลกํ อุปปชฺเชยฺย, โน\footnote{โน จ (ก)} ขตฺติโย, โน1 เวสฺโส, โน1 สุทฺโทติ.\csromanspacer{} โน หิทํ โภ โคตม, ขตฺติโยปิ หิ โภ โคตม ปาณาติปาตา ปฏิวิรโต อทินฺนาทานา ปฏิวิรโต กาเมสุมิจฺฉาจารา ปฏิวิรโต มุสาวาทา ปฏิวิรโต ปิสุณาย วาจาย ปฏิวิรโต ผรุสาย วาจาย ปฏิวิรโต สมฺผปฺปลาปา ปฏิวิรโต อนภิชฺฌาลุ อพฺยาปนฺนจิตฺโต สมฺมาทิฏฺฐิ กายสฺส เภทา ปรํ มรณา สุคติํ สคฺคํ โลกํ อุปปชฺเชยฺย.\csromanspacer{} พฺราหฺมโณปิ หิ โภ โคตม ฯเปฯ.\csromanspacer{} เวสฺโสปิ หิ โภ โคตม ฯเปฯ.\csromanspacer{} สุทฺโทปิ หิ โภ โคตม ฯเปฯ.\csromanspacer{} สพฺเพปิ หิ โภ โคตม จตฺตาโร วณฺณา ปาณาติปาตา ปฏิวิรตา อทินฺนาทานา ปฏิวิรตา กาเมสุมิจฺฉาจารา ปฏิวิรตา มุสาวาทา ปฏิวิรตา ปิสุณาย วาจาย ปฏิวิรตา ผรุสาย วาจาย ปฏิวิรตา สมฺผปฺปลาปา ปฏิวิรตา อนภิชฺฌาลู อพฺยาปนฺนจิตฺตา สมฺมาทิฏฺฐี กายสฺส เภทา ปรํ มรณา สุคติํ สคฺคํ โลกํ อุปปชฺเชยฺยุนฺติ.\csromanspacer{} เอตฺถ \csromanfolio{358}อสฺสลายน พฺราหฺมณานํ กิํ พลํ, โก อสฺสาโส, ยเทตฺถ พฺราหฺมณา เอวมาหํสุ “พฺราหฺมโณว เสฏฺโฐ วณฺโณ, หีโน อญฺโญ วณฺโณ ฯเปฯ พฺรหฺมทายาทา”ติ.\csromanspacer{} กิญฺจาปิ ภวํ โคตโม เอวมาห, อถ เขฺวตฺถ พฺราหฺมณา เอวเมตํ มญฺญนฺติ “พฺราหฺมโณว เสฏฺโฐ วณฺโณ, หีโน อญฺโญ วณฺโณ ฯเปฯ พฺรหฺมทายาทา”ติ.}

\proseitem{406}{ตํ กิํ มญฺญสิ อสฺสลายน, พฺราหฺมโณว นุ โข ปโหติ อสฺมิํ ปเทเส อเวรํ อพฺยาพชฺฌํ เมตฺตจิตฺตํ ภาเวตุํ, โน ขตฺติโย, โน เวสฺโส, โน สุทฺโทติ.\csromanspacer{} โน หิทํ โภ โคตม, ขตฺติโยปิ หิ โภ โคตม ปโหติ อสฺมิํ ปเทเส อเวรํ อพฺยาพชฺฌํ เมตฺตจิตฺตํ ภาเวตุํ.\csromanspacer{} พฺราหฺมโณปิ หิ โภ โคตม.\csromanspacer{} เวสฺโสปิ หิ โภ โคตม.\csromanspacer{} สุทฺโทปิ หิ โภ โคตม.\csromanspacer{} สพฺเพปิ หิ โภ โคตม จตฺตาโร วณฺณา ปโหนฺติ อสฺมิํ ปเทเส อเวรํ อพฺยาพชฺฌํ เมตฺตจิตฺตํ ภาเวตุนฺติ.\csromanspacer{} เอตฺถ อสฺสลายน พฺราหฺมณานํ กิํ พลํ, โก อสฺสาโส, ยเทตฺถ พฺราหฺมณา เอวมาหํสุ “พฺราหฺมโณว เสฏฺโฐ วณฺโณ, หีโน อญฺโญ วณฺโณ ฯเปฯ พฺรหฺมทายาทา”ติ.\csromanspacer{} กิญฺจาปิ ภวํ โคตโม เอวมาห, อถ เขฺวตฺถ พฺราหฺมณา เอวเมตํ มญฺญนฺติ “พฺราหฺมโณว เสฏฺโฐ วณฺโณ, หีโน อญฺโญ วณฺโณ ฯเปฯ พฺรหฺมทายาทา”ติ.}

\proseitem{407}{ตํ กิํ มญฺญสิ อสฺสลายน, พฺราหฺมโณว นุ โข ปโหติ โสตฺติสินานิํ อาทาย นทิํ คนฺตฺวา รโชชลฺลํ ปวาเหตุํ, โน ขตฺติโย, โน เวสฺโส, โน สุทฺโทติ.\csromanspacer{} โน หิทํ โภ โคตม, ขตฺติโยปิ หิ โภ โคตม ปโหติ โสตฺติสินานิํ อาทาย นทิํ คนฺตฺวา รโชชลฺลํ ปวาเหตุํ.\csromanspacer{} พฺราหฺมโณปิ หิ โภ โคตม.\csromanspacer{} เวสฺโสปิ หิ โภ โคตม.\csromanspacer{} สุทฺโทปิ หิ โภ โคตม.\csromanspacer{} สพฺเพปิ หิ โภ โคตม จตฺตาโร วณฺณา ปโหนฺติ โสตฺติสินานิํ อาทาย นทิํ คนฺตฺวา รโชชลฺลํ ปวาเหตุนฺติ.\csromanspacer{} เอตฺถ อสฺสลายน พฺราหฺมณานํ กิํ พลํ, โก อสฺสาโส, ยเทตฺถ พฺราหฺมณา เอวมาหํสุ “พฺราหฺมโณว เสฏฺโฐ วณฺโณ, หีโน อญฺโญ วณฺโณ ฯเปฯ พฺรหฺมทายาทา”ติ.\csromanspacer{} กิญฺจาปิ ภวํ โคตโม เอวมาห, อถ เขฺวตฺถ พฺราหฺมณา เอวเมตํ มญฺญนฺติ “พฺราหฺมโณว เสฏฺโฐ วณฺโณ, หีโน อญฺโญ วณฺโณ ฯเปฯ พฺรหฺมทายาทา”ติ.}

\proseitem{408}{ตํ กิํ มญฺญสิ อสฺสลายน, อิธ ราชา ขตฺติโย มุทฺธาวสิตฺโต นานาชจฺจานํ ปุริสานํ ปุริสสตํ สนฺนิปาเตยฺย “อายนฺตุ โภนฺโต,}

\vspace{\tipitakaparskip}
\csromancenter{อสฺสลายนสุตฺต (93) เย ตตฺถ ขตฺติยกุลา พฺราหฺมณกุลา ราชญฺญกุลา อุปฺปนฺนา สากสฺส วา สาลสฺส วา\footnote{อุปฺปนฺนา สาลสฺส วา (สี, อิ)} สลฬสฺส วา จนฺทนสฺส วา ปทุมกสฺส วา อุตฺตรารณิํ อาทาย อคฺคิํ อภินิพฺพตฺเตนฺตุ, เตโช ปาตุกโรนฺตุ.\csromanspacer{} อายนฺตุ ปน โภนฺโต, เย ตตฺถ จณฺฑาลกุลา เนสาทกุลา เวนกุลา\footnote{เวณกุลา (สี, อิ), เวณุกุลา(สฺยา, กํ)} รถการกุลา ปุกฺกุสลุลา อุปฺปนฺนา สาปานโทณิยา วา สูกรโทณิยา วา รชกโทณิยา วา เอรณฺฑกฏฺฐสฺส วา อุตฺตรารณิํ อาทาย อคฺคิํ อภินิพฺพตฺเตนฺตุ, เตโช ปาตุกโรนฺตู”ติ.}
\csromanmatikaanchor{mtk.52}
\csromantocmarkref{subsection}{3. อสฺสลายนสุตฺต}{mtk.52}
\bookmark[dest={mtk.52},level=2]{3. อสฺสลายนสุตฺต}
\prose{\csromanfolio{359}ตํ กิํ มญฺญสิ อสฺสลายน, โย เอวํ นุ โข โส\footnote{โย จ นุ โข (สฺยา, กํ, ก)} ขตฺติยกุลา พฺราหฺมณกุลา ราชญฺญกุลา อุปฺปนฺเนหิ สากสฺส วา สาลสฺส วา สลฬสฺส วา จนฺทนสฺส วา ปทุมกสฺส วา อุตฺตรารณิํ อาทาย อคฺคิ อภินิพฺพตฺโต เตโช ปาตุกโต, โส เอว นุ ขฺวาสฺส อคฺคิ อจฺจิมา เจว\footnote{จ (สี, อิ)} วณฺณวา\footnote{วณฺณิมา (สฺยา, กํ, อิ, ก)} จ ปภสฺสโร จ, เตน จ สกฺกา อคฺคินา อคฺคิกรณียํ กาตุํ.\csromanspacer{} โย ปน โส จณฺฑาลกุลา เนสาทกุลา เวนกุลา รถการกุลา ปุกฺกุสกุลา อุปฺปนฺเนหิ สาปานโทณิยา วา สูกรโทณิยา วา รชกโทณิยา วา เอรณฺฑกฏฺฐสฺส วา อุตฺตรารณิํ อาทาย อคฺคิ อภินิพฺพตฺโต เตโช ปาตุกโต, สฺวาสฺส อคฺคิ น เจว อจฺจิมา น จ วณฺณวา น จ ปภสฺสโร, น จ เตน สกฺกา อคฺคินา อคฺคิกรณียํ กาตุนฺติ.\csromanspacer{} โน หิทํ โภ โคตม, โยปิ หิ โส\footnote{โย โส (สี, อิ)} โภ โคตม ขตฺติยกุลา พฺราหฺมณกุลา ราชญฺญกุลา อุปฺปนฺเนหิ สากสฺส วา สาลสฺส วา สลฬสฺส วา จนฺทนสฺส วา ปทุมกสฺส วา อุตฺตรารณิํ อาทาย อคฺคิ อภินิพฺพตฺโต เตโช ปาตุกโต, สฺวาสฺส อคฺคิ อจฺจิมา เจว วณฺณวา จ ปภสฺสโร จ, เตน จ สกฺกา อคฺคินา อคฺคิกรณียํ กาตุํ.\csromanspacer{} โยปิ โส จณฺฑาลกุลา เนสาทกุลา เวนกุลา รถการกุลา ปุกฺกุสกุลา อุปฺปนฺเนหิ สาปานโทณิยา วา สูกรโทณิยา วา รชกโทณิยา วา เอรณฺฑกฏฺฐสฺส วา อุตฺตรารณิํ อาทาย อคฺคิ อภินิพฺพตฺโต เตโช ปาตุกโต, สฺวาสฺส\footnote{โส จสฺส (สี, อิ), โสปิสฺส (สฺยา, กํ)} อคฺคิ อจฺจิมา เจว วณฺณวา จ ปภสฺสโร จ, เตน จ สกฺกา อคฺคินา อคฺคิกรณียํ กาตุํ.\csromanspacer{} สพฺโพปิ หิ โภ โคตม อคฺคิ อจฺจิมา เจว วณฺณวา จ ปภสฺสโร จ, สพฺเพนปิ \csromanfolio{360}สกฺกา อคฺคินา อคฺคิกรณียํ กาตุนฺติ.\csromanspacer{} เอตฺถ อสฺสลายน พฺราหฺมณานํ กิํ พลํ, โก อสฺสาโส, ยเทตฺถ พฺราหฺมณา เอวมาหํสุ “พฺราหฺมโณว เสฏฺโฐ วณฺโณ, หีโน อญฺโญ วณฺโณ.\csromanspacer{} พฺราหฺมโณว สุกฺโก วณฺโณ, กโณฺห อญฺโญ วณฺโณ.\csromanspacer{} พฺราหฺมณาว สุชฺฌนฺติ โน อพฺราหฺมณา, พฺราหฺมณาว พฺรหฺมุโน ปุตฺตา โอรสา มุขโต ชาตา พฺรหฺมชา พฺรหฺมนิมฺมิตา พฺรหฺมทายาทา”ติ.\csromanspacer{} กิญฺจาปิ ภวํ โคตโม เอวมาห, อถ เขฺวตฺถ พฺราหฺมณา เอวเมตํ มญฺญนฺติ “พฺราหฺมโณว เสฏฺโฐ วณฺโณ, หีโน อญฺโญ วณฺโณ ฯเปฯ พฺรหฺมทายาทา”ติ.}

\proseitem{409}{ตํ กิํ มญฺญสิ อสฺสลายน, อิธ ขตฺติยกุมาโร พฺราหฺมณกญฺญาย สทฺธิํ สํวาสํ กปฺเปยฺย, เตสํ สํวาสมนฺวาย ปุตฺโต ชาเยถ.\csromanspacer{} โย โส ขตฺติยกุมาเรน พฺราหฺมณกญฺญาย ปุตฺโต อุปฺปนฺโน สิยา, โส มาตุปิ สทิโส, ปิตุปิ สทิโส.\csromanspacer{} ขตฺติโยติปิ วตฺตพฺโพ, พฺราหฺมโณติปิ วตฺตพฺโพติ.\csromanspacer{} โย โส โภ โคตม ขตฺติยกุมาเรน พฺราหฺมณกญฺญาย ปุตฺโต อุปฺปนฺโน สิยา, โส มาตุปิ สทิโส, ปิตุปิ สทิโส.\csromanspacer{} ขตฺติโยติปิ วตฺตพฺโพ, พฺราหฺมโณติปิ วตฺตพฺโพติ.}

\prose{ตํ กิํ มญฺญสิ อสฺสลายน, อิธ พฺราหฺมณกุมาโร ขตฺติยกญฺญาย สทฺธิํ สํวาสํ กปฺเปยฺย, เตสํ สํวาสมนฺวาย ปุตฺโต ชาเยถ.\csromanspacer{} โย โส พฺราหฺมณกุมาเรน ขตฺติยกญฺญาย ปุตฺโต อุปฺปนฺโน สิยา, โส มาตุปิ สทิโส, ปิตุปิ สทิโส.\csromanspacer{} ขตฺติโยติปิ วตฺตพฺโพ, พฺราหฺมโณติปิ วตฺตพฺโพติ.\csromanspacer{} โย โส โภ โคตม พฺราหฺมณกุมาเรน ขตฺติยกญฺญาย ปุตฺโต อุปฺปนฺโน สิยา, โส มาตุปิ สทิโส, ปิตุปิ สทิโส.\csromanspacer{} ขตฺติโยติปิ วตฺตพฺโพ, พฺราหฺมโณติปิ วตฺตพฺโพติ.}

\prose{ตํ กิํ มญฺญสิ อสฺสลายน, อิธ วฬวํ คทฺรเภน สมฺปโยเชยฺยุํ\footnote{สํโยเชยฺย (ก)}, เตสํ สมฺปโยคมนฺวาย กิโสโร ชาเยถ.\csromanspacer{} โย โส วฬวาย คทฺรเภน กิโสโร อุปฺปนฺโน สิยา, โส มาตุปิ สทิโส, ปิตุปิ สทิโส.\csromanspacer{} อสฺโสติปิ วตฺตพฺโพ, คทฺรโภติปิ วตฺตพฺโพติ.\csromanspacer{} กุณฺฑํ หิ โส\footnote{เวกุรญฺชาย หิ โส (สี, อิ), โส กุมารณฺฑุปิ โส (สฺยา, กํ), เวกุลโช หิ โส (?)} โภ โคตม อสฺสตโร โหติ, อิทํ หิสฺส โภ โคตม}

\vspace{\tipitakaparskip}
\csromancenter{อสฺสลายนสุตฺต (93) นานากรณํ ปสฺสามิ, อมุตฺร จ ปเนสานํ น กิญฺจิ นานากรณํ ปสฺสามีติ.}
\prose{\csromanfolio{361}ตํ กิํ มญฺญสิ อสฺสลายน, อิธาสฺสุ เทฺว มาณวกา ภาตโร ส-อุทริยา.\csromanspacer{} เอโก อชฺฌายโก อุปนีโต, เอโก อนชฺฌายโก อนุปนีโต.\csromanspacer{} กเมตฺถ พฺราหฺมณา ปฐมํ โภเชยฺยุํ สทฺเธ วา ถาลิปาเก วา ยญฺเญ วา ปาหุเน วาติ.\csromanspacer{} โย โส โภ โคตม มาณวโก อชฺฌายโก อุปนีโต, ตเมตฺถ พฺราหฺมณา ปฐมํ โภเชยฺยุํ สทฺเธ วา ถาลิปาเก วา ยญฺเญ วา ปาหุเน วา.\csromanspacer{} กิํ หิ โภ โคตม อนชฺฌายเก อนุปนีเต ทินฺนํ มหปฺผลํ ภวิสฺสตีติ.}

\prose{ตํ กิํ มญฺญสิ อสฺสลายน, อิธาสฺสุ เทฺว มาณวกา ภาตโร โสทริยา.\csromanspacer{} เอโก อชฺฌายโก อุปนีโต ทุสฺสีโล ปาปธมฺโม, เอโก อนชฺฌายโก อนุปนีโต สีลวา กลฺยาณธมฺโม.\csromanspacer{} กเมตฺถ พฺราหฺมณา ปฐมํ โภเชยฺยุํ สทฺเธ วา ถาลิปาเก วา ยญฺเญ วา ปาหุเน วาติ.\csromanspacer{} โย โส โภ โคตม มาณวโก อนชฺฌายโก อนุปนีโต สีลวา กลฺยาณธมฺโม, ตเมตฺถ พฺราหฺมณา ปฐมํ โภเชยฺยุํ สทฺเธ วา ถาลิปาเก วา ยญฺเญ วา ปาหุเน วา.\csromanspacer{} กิํ หิ โภ โคตม ทุสฺสีเล ปาปธมฺเม ทินฺนํ มหปฺผลํ ภวิสฺสตีติ.}

\prose{ปุพฺเพ โข ตฺวํ อสฺสลายน ชาติํ อคมาสิ, ชาติํ คนฺตฺวา มนฺเต อคมาสิ, มนฺเต คนฺตฺวา ตเป อคมาสิ, ตเป คนฺตฺวา\footnote{มนฺเต คนฺตฺวา ตเมตํ ตฺวํ (สี, อิ), มนฺเต คนฺตฺวา ตเมว ฐเปตฺวา (สฺยา, กํ)} จาตุวณฺณิํ สุทฺธิํ ปจฺจาคโต, ยมหํ ปญฺญเปมีติ.\csromanspacer{} เอวํ วุตฺเต อสฺสลายโน มาณโว ตุณฺหีภูโต มงฺกุภูโต ปตฺตกฺขนฺโธ อโธมุโข ปชฺฌายนฺโต อปฺปฏิภาโน นิสีทิ.}

\proseitem{410}{อถ โข ภควา อสฺสลายนํ มาณวํ ตุณฺหีภูตํ มงฺกุภูตํ ปตฺตกฺขนฺธํ อโธมุขํ ปชฺฌายนฺตํ อปฺปฏิภานํ วิทิตฺวา อสฺสลายนํ มาณวํ เอตทโวจ\csromandash{}ภูตปุพฺพํ อสฺสลายน สตฺตนฺนํ พฺราหฺมณิสีนํ อรญฺญายตเน ปณฺณกุฏีสุ สมฺมนฺตานํ\footnote{วสนฺตานํ (สี)} เอวรูปํ ปาปกํ ทิฏฺฐิคตํ อุปฺปนฺนํ โหติ “พฺราหฺมโณว เสฏฺโฐ วณฺโณ, หีโน อญฺโญ วณฺโณ ฯเปฯ พฺรหฺมทายาทา”ติ.\csromanspacer{} อสฺโสสิ \csromanfolio{362}โข อสฺสลายน อสิโต เทวโล อิสิ “สตฺตนฺนํ กิร พฺราหฺมณิสีนํ อรญฺญายตเน ปณฺณกุฏีสุ สมฺมนฺตานํ เอวรูปํ ปาปกํ ทิฏฺฐิคตํ อุปฺปนฺนํ ‘พฺราหฺมโณว เสฏฺโฐ วณฺโณ ฯเปฯ พฺรหฺมทายาทา’ติ”.\csromanspacer{} อถ โข อสฺสลายน อสิโต เทวโล อิสิ เกสมสฺสุํ กปฺเปตฺวา มญฺชิฏฺฐวณฺณานิ ทุสฺสานิ นิวาเสตฺวา ปฏลิโย\footnote{อฏลิโย (สี, อิ), อคลิโย (สฺยา, กํ)} อุปาหนา อารุหิตฺวา ชาตรูปมยํ ทณฺฑํ คเหตฺวา สตฺตนฺนํ พฺราหฺมณิสีนํ ปตฺถณฺฑิเล ปาตุรโหสิ.\csromanspacer{} อถ โข อสฺสลายน อสิโต เทวโล อิสิ สตฺตนฺนํ พฺราหฺมณิสีนํ ปตฺถณฺฑิเล จงฺกมมาโน เอวมาห “หนฺท โก นุ โข อิเม ภวนฺโต พฺราหฺมณิสโย คตา\footnote{คนฺตา (สฺยา, กํ, ก)}, หนฺท โก นุ โข อิเม ภวนฺโต พฺราหฺมณิสโย คตา”ติ.\csromanspacer{} อถ โข อสฺสลายน สตฺตนฺนํ พฺราหฺมณิสีนํ เอตทโหสิ “โก นายํ คามณฺฑลรูโป วิย สตฺตนฺนํ พฺราหฺมณิสีนํ ปตฺถณฺฑิเล จงฺกมมาโน เอวมาห ‘หนฺท โก นุ โข อิเม ภวนฺโต พฺราหฺมณิสโย คตา, หนฺท โก นุ โข อิเม ภวนฺโต พฺราหฺมณิสโย คตา’ติ.\csromanspacer{} หนฺท นํ อภิสปามา”ติ.\csromanspacer{} อถ โข อสฺสลายน สตฺต พฺราหฺมณิสโย อสิตํ เทวลํ อิสิํ อภิสปิํสุ “ภสฺมา วสล\footnote{วสลี (อิ), วสลิ (ก), จปลี (สฺยา, กํ)} โหติ, ภสฺมา วสล โหตี”ติ\footnote{ภสฺมา วสล โหหีติ อภิสปวจนํ สี-อิ-โปตฺถเกสุ สกิเทว อาคตํ.}.\csromanspacer{} ยถา ยถา โข อสฺสลายน สตฺต พฺราหฺมณิสโย อสิตํ เทวลํ อิสิํ อภิสปิํสุ, ตถา ตถา อสิโต เทวโล อิสิ อภิรูปตโร เจว โหติ ทสฺสนียตโร จ ปาสาทิกตโร จ.\csromanspacer{} อถ โข อสฺสลายน สตฺตนฺนํ พฺราหฺมณิสีนํ เอตทโหสิ “โมฆํ วต โน ตโป, อผลํ พฺรหฺมจริยํ, มยํ หิ ปุพฺเพ ยํ อภิสปาม ‘ภสฺมา วสล โหติ, ภสฺมา วสล โหหี’ติ, ภสฺมาว ภวติ เอกจฺโจ.\csromanspacer{} อิมํ ปน มยํ ยถา ยถา อภิสปาม, ตถา ตถา อภิรูปตโร เจว โหติ ทสฺสนียตโร จ ปาสาทิกตโร จา”ติ.\csromanspacer{} น ภวนฺตานํ โมฆํ ตโป, นาผลํ พฺรหฺมจริยํ.\csromanspacer{} อิงฺฆ ภวนฺโต โย มยิ มโนปโทโส, ตํ ปชหถาติ.\csromanspacer{} โย ภวติ มโนปโทโส, ตํ ปชหาม.\csromanspacer{} โก นุ ภวํ โหตีติ.\csromanspacer{} สุโต นุ ภวตํ “อสิโต เทวโล อิสี”ติ.\csromanspacer{} เอวํ โภ.\csromanspacer{} โส ขฺวาหํ โภ โหมีติ.\csromanspacer{} อถ โข อสฺสลายน สตฺต พฺราหฺมณิสโย อสิตํ เทวลํ อิสิํ อภิวาเทตุํ อุปกฺกมิํสุ.}

\csromanheader{2}{3. อสฺสลายนสุตฺต (93)}
\proseitem{411}{\csromanfolio{363}อถ โข อสฺสลายน อสิโต เทวโล อิสิ สตฺต พฺราหฺมณิสโย เอตทโวจ “สุตํ เมตํ โภ ‘สตฺตนฺนํ กิร พฺราหฺมณิสีนํ อรญฺญายตเน ปณฺณกุฏีสุ สมฺมนฺตานํ เอวรูปํ ปาปกํ ทิฏฺฐิคตํ อุปฺปนฺนํ ‘พฺราหฺมโณว เสฏฺโฐ วณฺโณ, หีโน อญฺโญ วณฺโณ.\csromanspacer{} พฺราหฺมโณว สุกฺโก วณฺโณ, กโณฺห อญฺโญ วณฺโณ.\csromanspacer{} พฺราหฺมณาว สุชฺฌนฺติ โน อพฺราหฺมณา, พฺราหฺมณาว พฺรหฺมุโน ปุตฺตา โอรสา มุขโต ชาตา พฺรหฺมชา พฺรหฺมนิมฺมิตา พฺรหฺมทายาทา’ติ”.\csromanspacer{} เอวํ โภ.}

\prose{ชานนฺติ ปน โภนฺโต “ยา ชนิกา มาตา\footnote{ชนิมาตา (สี, สฺยา, กํ, อิ)} พฺราหฺมณํเยว อคมาสิ, โน อพฺราหฺมณนฺ”ติ.\csromanspacer{} โน หิทํ โภ.}

\prose{ชานนฺติ ปน โภนฺโต “ยา ชนิกามาตุ\footnote{ชนิมาตุ (สี, สฺยา,กํ, อิ)} มาตา ยาว สตฺตมา มาตุมาตามหยุคา พฺราหฺมณํเยว อคมาสิ, โน อพฺราหฺมณนฺ”ติ.\csromanspacer{} โน หิทํ โภ.}

\prose{ชานนฺติ ปน โภนฺโต “โย ชนโก ปิตา\footnote{ชนิปิตา (สี, สฺยา, กํ, อิ)} พฺราหฺมณิํเยว อคมาสิ, โน อพฺราหฺมณินฺ”ติ.\csromanspacer{} โน หิทํ โภ.}

\prose{ชานนฺติ ปน โภนฺโต “โย ชนกปิตุ\footnote{ชนิปิตุ (สี, สฺยา, กํ, อิ) 5-5. น มยํ ชานาม โภ ยถา คพฺภสฺส อวกฺกนฺติ โหตีติ. ยถา กถํ ปน โภ คพฺภสฺส อวกฺกนฺติ โหตีติ. (ก)} ปิตา ยาว สตฺตมา ปิตุปิตามหยุคา พฺราหฺมณิํเยว อคมาสิ, โน อพฺราหฺมณินฺ”ติ.\csromanspacer{} โน หิทํ โภ.}

\prose{ชานนฺติ ปน โภนฺโต ยถา คพฺภสฺส อวกฺกนฺติ โหตีติ.5 ชานาม มยํ โภ ยถา คพฺภสฺส อวกฺกนฺติ โหติ5.\csromanspacer{} อิธ มาตาปิตโร จ สนฺนิปติตา โหนฺติ, มาตา จ อุตุนี โหติ, คนฺธพฺโพ จ ปจฺจุปฏฺฐิโต โหติ.\csromanspacer{} เอวํ ติณฺณํ สนฺนิปาตา คพฺภสฺส อวกฺกนฺติ โหตีติ.}

\prose{ชานนฺติ ปน โภนฺโต “ตคฺฆ\footnote{ยคฺเฆ (สี, สฺยา, กํ, อิ)} โส คนฺธพฺโพ ขตฺติโย วา พฺราหฺมโณ วา เวสฺโส วา สุทฺโท วา”ติ.\csromanspacer{} น มยํ โภ ชานาม “ตคฺฆ โส คนฺธพฺโพ ขตฺติโย วา พฺราหฺมโณ วา เวสฺโส วา สุทฺโท วา”ติ.\csromanspacer{} เอวํ สนฺเต โภ ชานาถ “เก ตุเมฺห โหถา”ติ.\csromanspacer{} เอวํ สนฺเต โภ น มยํ \csromanfolio{364}ชานาม “เก มยํ โหมา”ติ.\csromanspacer{} เต หิ นาม อสฺสลายน สตฺต พฺราหฺมณิสโย อสิเตน เทวเลน อิสินา สเก ชาติวาเท สมนุยุญฺชียมานา สมนุคฺคาหียมานา สมนุภาสียมานา น สมฺปายิสฺสนฺติ.\csromanspacer{} กิํ ปน ตฺวํ เอตรหิ มยา สกสฺมิํ ชาติวาเท สมนุยุญฺชียมาโน สมนุคฺคาหียมาโน สมนุภาสียมาโน สมฺปายิสฺสสิ.\csromanspacer{} เยสํ ตฺวํ สาจริยโก น ปุณฺโณ ทพฺพิคาโหติ.}

\prose{เอวํ วุตฺเต อสฺสลายโน มาณโว ภควนฺตํ เอตทโวจ “อภิกฺกนฺตํ โภ โคตม ฯเปฯ อุปาสกํ มํ ภวํ โคตโม ธาเรตุ อชฺชตคฺเค ปาณุเปตํ สรณํ คตนฺ”ติ.}

\nitthitamruled{อสฺสลายนสุตฺตํ นิฏฺฐิตํ ตติยํ.}
\csromanheader{2}{4. โฆฏมุขสุตฺต}
\proseitem{412}{เอวํ เม สุตํ\csromandash{}เอกํ สมยํ อายสฺมา อุเทโน พาราณสิยํ วิหรติ เขมิยมฺพวเน.\csromanspacer{} เตน โข ปน สมเยน โฆฏมุโข พฺราหฺมโณ พาราณสิํ อนุปฺปตฺโต โหติ เกนจิเทว กรณีเยน.\csromanspacer{} อถ โข โฆฏมุโข พฺราหฺมโณ ชงฺฆาวิหารํ อนุจงฺกมมาโน อนุวิจรมาโน เยน เขมิยมฺพวนํ เตนุปสงฺกมิ.\csromanspacer{} เตน โข ปน สมเยน อายสฺมา อุเทโน อพฺโภกาเส จงฺกมติ.\csromanspacer{} อถ โข โฆฏมุโข พฺราหฺมโณ เยนายสฺมา อุเทโน เตนุปสงฺกมิ, อุปสงฺกมิตฺวา อายสฺมตา อุเทเนน สทฺธิํ สมฺโมทิ, สมฺโมทนียํ กถํ สารณียํ วีติสาเรตฺวา อายสฺมนฺตํ อุเทนํ จงฺกมนฺตํ อนุจงฺกมมาโน เอวมาห “อมฺโภ สมณ นตฺถิ ธมฺมิโก ปริพฺพโช\footnote{ปริพฺพาโช (สี, อิ)}, เอวํ เม เอตฺถ โหติ, ตญฺจ โข ภวนฺตรูปานํ วา อทสฺสนา, โย วา ปเนตฺถ ธมฺโมติ.}

\prose{เอวํ วุตฺเต อายสฺมา อุเทโน จงฺกมา โอโรหิตฺวา วิหารํ ปวิสิตฺวา ปญฺญตฺเต อาสเน นิสีทิ.\csromanspacer{} โฆฏมุโขปิ โข พฺราหฺมโณ จงฺกมา โอโรหิตฺวา วิหารํ ปวิสิตฺวา เอกมนฺต อฏฺฐาสิ, เอกมนฺตํ ฐิตํ โข โฆฏมุขํ พฺราหฺมณํ อายสฺมา อุเทโน เอตทโวจ “สํวิชฺชนฺติ\footnote{สํวิชฺชนฺเต (พหูสุ)} โข พฺราหฺมณ อาสนานิ, สเจ อากงฺขสิ นิสีทา”ติ.\csromanspacer{} เอตเทว โข ปน}

\vspace{\tipitakaparskip}
\csromancenter{โฆฏมุขสุตฺต (94) มยํ โภโต อุเทนสฺส อาคมยมานา นิสีทาม, กถํ หิ นาม มาทิโส ปุพฺเพ อนิมนฺติโต อาสเน นิสีทิตพฺพํ มญฺเญยฺยาติ.\csromanspacer{} อถ โข โฆฏมุโข พฺราหฺมโณ อญฺญตรํ นีจํ อาสนํ คเหตฺวา เอกมนฺตํ นิสีทิ, เอกมนฺตํ นิสินฺโน โข โฆฏมุโข พฺราหฺมโณ อายสฺมนฺตํ อุเทนํ เอตทโวจ “อพฺโภ สมณ นตฺถิ ธมฺมิโก ปริพฺพโช, เอวํ เม เอตฺถ โหติ, ตญฺจ โข ภวนฺตรูปานํ วา อทสฺสนา, โย วา ปเนตฺถ ธมฺโม”ติ.\csromanspacer{} สเจ โข ปน เม ตฺวํ พฺราหฺมณ อนุญฺเญยฺยํ อนุชาเนยฺยาสิ, ปฏิกฺโกสิตพฺพญฺจ ปฏิกฺโกเสยฺยาสิ, ยสฺส จ ปน เม ภาสิตสฺส อตฺถํ น ชาเนยฺยาสิ, มมํเยว ตตฺถ อุตฺตริ ปฏิปุจฺเฉยฺยาสิ “อิทํ โภ อุเทน กถํ, อิมสฺส กฺวตฺโถ”ติ, เอวํ กตฺวา สิยา โน เอตฺถ กถาสลฺลาโปติ.\csromanspacer{} อนุญฺเญยฺยํ ขฺวาหํ โภโต อุเทนสฺส อนุชานิสฺสามิ, ปฏิกฺโกสิตพฺพญฺจ ปฏิกฺโกสิสฺสามิ, ยสฺส จ ปนาหํ โภโต อุเทนสฺส ภาสิตสฺส อตฺถํ น ชานิสฺสามิ, ภวนฺตํเยว ตตฺถ อุเทนํ อุตฺตริ ปฏิปุจฺฉิสฺสามิ “อิทํ โภ อุเทน กถํ, อิมสฺส กฺวตฺโถ”ติ, เอวํ กตฺวา โหตุ โน เอตฺถ กถาสลฺลาโปติ.}
\proseitem{413}{\csromanfolio{365}จตฺตาโรเม พฺราหฺมณ ปุคฺคลา สนฺโต สํวิชฺชมานา โลกสฺมิํ.\csromanspacer{} กตเม จตฺตาโร.\csromanspacer{} อิธ พฺราหฺมณ เอกจฺโจ ปุคฺคโล อตฺตนฺตโป โหติ อตฺตปริตาปนานุโยคมนุยุตฺโต.\csromanspacer{} อิธ ปน พฺราหฺมณ เอกจฺโจ ปุคฺคโล ปรนฺตโป โหติ ปรปริตาปนานุโยคมนุยุตฺโต.\csromanspacer{} อิธ ปน พฺราหฺมณ เอกจฺโจ ปุคฺคโล อตฺตนฺตโป จ โหติ อตฺตปริตาปนานุโยคมนุยุตฺโต ปรนฺตโป จ ปรปริตาปนานุโยคมนุยุตฺโต.\csromanspacer{} อิธ ปน พฺราหฺมณ เอกจฺโจ ปุคฺคโล เนวตฺตนฺตโป โหติ นาตฺตปริตาปนานุโยคมนุยุตฺโต น ปรนฺตโป น ปรปริตาปนานุโยคมนุยุตฺโต, โส อนตฺตนฺตโป อปรนฺตโป ทิฏฺเฐว ธมฺเม นิจฺฉาโต นิพฺพุโต สีตีภูโต สุขปฺปฏิสํเวที พฺรหฺมภูเตน อตฺตนา วิหรติ.\csromanspacer{} อิเมสํ พฺราหฺมณ จตุนฺนํ ปุคฺคลานํ กตโม เต ปุคฺคโล จิตฺตํ อาราเธตีติ.}

\prose{ยฺวายํ โภ อุเทน ปุคฺคโล อตฺตนฺตโป อตฺตปริตาปนานุโยคมนุยุตฺโต, อยํ เม ปุคฺคโล จิตฺตํ นาราเธติ.\csromanspacer{} โยปายํ โภ อุเทน ปุคฺคโล ปรนฺตโป ปรปริตาปนานุโยคมนุยุตฺโต, อยมฺปิ เม ปุคฺคโล จิตฺตํ นาราเธติ.\csromanspacer{} โยปายํ โภ อุเทน ปุคฺคโล อตฺตนฺตโป จ \csromanfolio{366}อตฺตปริตาปนานุโยคมนุยุตฺโต, ปรนฺตโป จ ปรปริตาปนานุโยคมนุยุตฺโต, อยมฺปิ เม ปุคฺคโล จิตฺตํ นาราเธติ.\csromanspacer{} โย จ โข อยํ โภ อุเทน ปุคฺคโล เนวตฺตนฺตโป นาตฺตปริตาปนานุโยคมนุยุตฺโต น ปรนฺตโป น ปรปริตาปนานุโยคมนุยุตฺโต, โส อนตฺตนฺตโป อปรนฺตโป ทิฏฺเฐว ธมฺเม นิจฺฉาโต นิพฺพุโต สีตีภูโต สุขปฺปฏิสํเวที พฺรหฺมภูเตน อตฺตนา วิหรติ, อยเมว เม ปุคฺคโล จิตฺตํ อาราเธตีติ.}

\prose{กสฺมา ปน เต พฺราหฺมณ อิเม ตโย ปุคฺคลา จิตฺตํ นาราเธนฺตีติ.\csromanspacer{} ยฺวายํ โภ อุเทน ปุคฺคโล อตฺตนฺตโป อตฺตปริตาปนานุโยคมนุยุตฺโต, โส อตฺตานํ สุขกามํ ทุกฺขปฏิกฺกูลํ อาตาเปติ ปริตาเปติ, อิมินา เม อยํ ปุคฺคโล จิตฺตํ นาราเธติ.\csromanspacer{} โยปายํ โภ อุเทน ปุคฺคโล ปรนฺตโป ปรปริตาปนานุโยคมนุยุตฺโต, โส ปรํ สุขกามํ ทุกฺขปฏิกฺกูลํ อาตาเปติ ปริตาเปติ, อิมินา เม อยํ ปุคฺคโล จิตฺตํ นาราเธติ.\csromanspacer{} โยปายํ โภ อุเทน ปุคฺคโล อตฺตนฺตโป จ อตฺตปริตาปนานุยุคมนุยุตฺโต ปรนฺตโป จ ปรปริตาปนานุโยคมนุยุตฺโต, โส อตฺตานญฺจ ปรญฺจ สุขกามํ ทุกฺขปฏิกฺกูลํ อาตาเปติ ปริตาเปติ, อิมินา เม อยํ ปุคฺคโล จิตฺตํ นาราเธติ.\csromanspacer{} โย จ โข อยํ โภ อุเทน ปุคฺคโล เนวตฺตนฺตโป นาตฺตปริตาปนานุโยคมนุยุตฺโต น ปรนฺตโป น ปรปริตาปนานุโยคมนุยุตฺโต, โส อนตฺตนฺตโป อปรนฺตโป ทิฏฺเฐว ธมฺเม นิจฺฉาโต นิพฺพุโต สีตีภูโต สุขปฺปฏิสํเวที พฺรหฺมภูเตน อตฺตนา วิหรติ, โส อตฺตานญฺจ ปรญฺจ สุขกามํ ทุกฺขปฏิกฺกูลํ เนว อาตาเปติ น ปริตาเปติ, อิมินา เม อยํ ปุคฺคโล จิตฺตํ อาราเธตีติ.}

\proseitem{414}{เทฺวมา พฺราหฺมณ ปริสา.\csromanspacer{} กตมา เทฺว, อิธ พฺราหฺมณ เอกจฺจา ปริสา สารตฺตรตฺตา มณิกุณฺฑเลสุ ปุตฺตภริยํ ปริเยสติ, ทาสิทาสํ ปริเยสติ, เขตฺตวตฺถุํ ปริเยสติ, ชาตรูปรชตํ ปริเยสติ.}

\prose{อิธ ปน พฺราหฺมณ เอกจฺจา ปริสา อสารตฺตรตฺตา มณิกุณฺฑเลสุ ปุตฺตภริยํ ปหาย ทาสิทาสํ ปหาย เขตฺตวตฺถุํ ปหาย ชาตรูปรชตํ ปหาย อคารสฺมา อนคาริยํ ปพฺพชิตา, สฺวายํ พฺราหฺมณ ปุคฺคโล เนวตฺตนฺตโป นาตฺตปริตาปนานุโยคมนุยุตฺโต น ปรนฺตโป น ปรปริตาปนานุโยคมนุยุตฺโต, โส อนตฺตนฺตโป อปรนฺตโป ทิฏฺเฐว ธมฺเม}

\vspace{\tipitakaparskip}
\csromancenter{โฆฏมุขสุตฺต (94) นิจฺฉาโต นิพฺพุโต สีติภูโต สุขปฺปฏิสํเวที พฺรหฺมภูเตน อตฺตนา วิหรติ.\csromanspacer{} อิธ กตมํ ตฺวํ พฺราหฺมณ ปุคฺคลํ กตมาย ปริสาย พหุลํ สมนุปสฺสสิ, ยา จายํ ปริสา สารตฺตรตฺตา มณิกุณฺฑเลสุ ปุตฺตภริยํ ปริเยสติ, ทาสิทาสํ ปริเยสติ, เขตฺตวตฺถุํ ปริเยสติ, ชาตรูปรชตํ ปริเยสติ, ยา จายํ ปริสา อสารตฺตรตฺตา มณิกุณฺฑเลสุ ปุตฺตภริยํ ปหาย ทาสิทาสํ ปหาย เขตฺตวตฺถุํ ปหาย ชาตรูปรชตํ ปหาย อคารสฺมา อนคาริยํ ปพฺพชิตาติ.}
\prose{\csromanfolio{367}ยฺวายํ โภ อุเทน ปุคฺคโล เนวตฺตนฺตโป นาตฺตปริตาปนานุโยคมนุยุตฺโต น ปรนฺตโป น ปรปริตาปนานุโยคมนุยุตฺโต, โส อนตฺตนฺตโป อปรนฺตโป ทิฏฺเฐว ธมฺเม นิจฺฉาโต นิพฺพุโต สีตีภูโต สุขปฺปฏิสํเวที พฺรหฺมภูเตน อตฺตนา วิหรติ.\csromanspacer{} อิมาหํ ปุคฺคลํ ยายํ ปริสา อสารตฺตรตฺตา มณิกุณฺฑเลสุ ปุตฺตภริยํ ปหาย ทาสิทาสํ ปหาย เขตฺตวตฺถุํ ปหาย ชาตรูปรชตํ ปหาย อคารสฺมา อนคาริยํ ปพฺพชิตา, อิมิสฺสํ ปริสายํ พหุลํ สมนุปสฺสามีติ.}

\prose{อิทาเนว โข ปน เต พฺราหฺมณ ภาสิตํ มยํ เอวํ อาชานาม “อมฺโภ สมณ นตฺถิ ธมฺมิโก ปริพฺพโช, เอวํ เม เอตฺถ โหติ, ตญฺจ โข ภวนฺตรูปานํ วา อทสฺสนา, โย วา ปเนตฺถ ธมฺโม”ติ.\csromanspacer{} อทฺธา เมสา โภ อุเทน สานุคฺคหา วาจา ภาสิตา, อตฺถิ ธมฺมิโก ปริพฺพโช, เอวํ เม เอตฺถ โหติ, เอวญฺจ ปน มํ ภวํ อุเทโน ธาเรตุ, เย จ เม โภตา อุเทเนน จตฺตาโร ปุคฺคลา สํขิตฺเตน วุตฺตา วิตฺถาเรน อวิภตฺตา, สาธุ เม ภวํ อุเทโน อิเม จตฺตาโร ปุคฺคเล วิตฺถาเรน วิภชตุ อนุกมฺปํ อุปาทายาติ.\csromanspacer{} เตน หิ พฺราหฺมณ สุณาหิ สาธุกํ มนสิ กโรหิ ภาสิสฺสามีติ.\csromanspacer{} เอวํ โภติ โข โฆฏมุโข พฺราหฺมโณ อายสฺมโต อุเทนสฺส ปจฺจสฺโสสิ.\csromanspacer{} อายสฺมา อุเทโน เอตทโวจ\csromandash{}}

\proseitem{415}{กตโม จ พฺราหฺมณ ปุคฺคโล อตฺตนฺตโป อตฺตปริตาปนานุโยคมนุยุตฺโต.\csromanspacer{} อิธ พฺราหฺมณ เอกจฺโจ ปุคฺคโล อเจลโก โหติ มุตฺตาจาโร, หตฺถาปเลขโน, น-เอหิภทฺทนฺติโก, นติฏฺฐภทฺทนฺติโก, นาภิหฏํ, น อุทฺทิสฺสกตํ, น นิมนฺตนํ สาทิยติ, โส น กุมฺภิมุขา ปฏิคฺคณฺหาติ, น กโฬปิมุขา ปฏิคฺคณฺหาติ, น เอฬกมนฺตรํ, น ทณฺฑมนฺตรํ, น มุสลมนฺตรํ, น ทฺวินฺนํ ภุญฺชมานานํ, น คพฺภินิยา, น \csromanfolio{368}ปายมานาย, น ปุริสนฺตรคตาย, น สงฺกิตฺตีสุ, น ยตฺถ สา อุปฏฺฐิโต โหติ, น ยตฺถ มกฺขิกา สณฺฑสณฺฑจารินี, น มจฺฉํ, น มํสํ, น สุรํ, น เมรยํ, น ถุโสทกํ ปิวติ.\csromanspacer{} โส เอกาคาริโก วา โหติ เอกาโลปิโก, ทฺวาคาริโก วา โหติ ทฺวาโลปิโก ฯเปฯ สตฺตาคาริโก วา โหติ สตฺตาโลปิโก.\csromanspacer{} เอกิสฺสาปิ ทตฺติยา ยาเปติ, ทฺวีหิปิ ทตฺตีหิ ยาเปติ ฯเปฯ สตฺตหิปิ ทตฺตีหิ ยาเปติ.\csromanspacer{} เอกาหิกมฺปิ อาหารํ อาหาเรติ.\csromanspacer{} ทฺวีหิกมฺปิ อาหารํ อาหาเรติ ฯเปฯ.\csromanspacer{} สตฺตาหิกมฺปิ อาหารํ อาหาเรติ.\csromanspacer{} อิติ เอวรูปํ อทฺธมาสิกํ ปริยายภตฺตโภชนานุโยคมนุยุตฺโต วิหรติ.\csromanspacer{} โส สากภกฺโข วา โหติ, สามากภกฺโข วา โหติ, นีวารภกฺโข วา โหติ, ททฺทุลภกฺโข วา โหติ, หฏภกฺโข วา โหติ, กณภกฺโข วา โหติ, อาจามภกฺโข วา โหติ, ปิญฺญากภกฺโข วา โหติ, ติณภกฺโข วา โหติ, โคมยภกฺโข วา โหติ, วนมูลผลาหาโร ยาเปติ ปวตฺตผลโภชี.\csromanspacer{} โส สาณานิปิ ธาเรติ, มสาณานิปิ ธาเรติ, ฉวทุสฺสานิปิ ธาเรติ, ปํสุกูลานิปิ ธาเรติ, ติรีฏานิปิ ธาเรติ, อชินมฺปิ ธาเรติ, อชินกฺขิปมฺปิ ธาเรติ, กุสจีรมฺปิ ธาเรติ, วากจีรมฺปิ ธาเรติ, ผลกจีรมฺปิ ธาเรติ, เกสกมฺพลมฺปิ ธาเรติ, วาฬกมฺพลมฺปิ ธาเรติ, อุลูกปกฺขมฺปิ ธาเรติ, เกสมสฺสุโลจโกปิ โหติ เกสมสฺสุโลจนานุโยคมนุยุตฺโต, อุพฺภฏฺฐโกปิ โหติ อาสนปฏิกฺขิตฺโต, อุกฺกุฏิโกปิ โหติ อุกฺกุฏิกปฺปธานมนุยุตฺโต, กณฺฏกาปสฺสยิโกปิ โหติ กณฺฏกาปสฺสเย เสยฺยํ กปฺเปติ, สายตติยกมฺปิ อุทโกโรหนานุโยคมนุยุตฺโต วิหรติ.\csromanspacer{} อิติ เอวรูปํ อเนกวิหิตํ กายสฺส อาตาปนปริตาปนานุโยคมนุยุตฺโต วิหรติ.\csromanspacer{} อยํ วุจฺจติ พฺราหฺมณ ปุคฺคโล อตฺตนฺตโป อตฺตปริตาปนานุโยคมนุยุตฺโต.}

\proseitem{416}{กตโม จ พฺราหฺมณ ปุคฺคโล ปรนฺตโป ปรปริตาปนานุโยคมนุยุตฺโต.\csromanspacer{} อิธ พฺราหฺมณ เอกจฺโจ ปุคฺคโล โอรพฺภิโก โหติ สูกริโก สากุณิโก มาควิโก ลุทฺโท มจฺฉฆาตโก โจโร โจรฆาตโก โคฆาตโก พนฺธนาคาริโก, เย วา ปนญฺเญปิ เกจิ กุรูรกมฺมนฺตา.\csromanspacer{} อยํ วุจฺจติ พฺราหฺมณ ปุคฺคโล ปรนฺตโป ปรปริตาปนานุโยคมนุยุตฺโต.}

\csromanheader{2}{4. โฆฏมุขสุตฺต (94)}
\proseitem{417}{\csromanfolio{369}กตโม จ พฺราหฺมณ ปุคฺคโล อตฺตนฺตโป จ อตฺตปริตาปนานุโยคมนุยุตฺโต ปรนฺตโป จ ปรปริตาปนานุโยคมนุยุตฺโต.\csromanspacer{} อิธ พฺราหฺมณ เอกจฺโจ ปุคฺคโล ราชา วา โหติ ขตฺติโย มุทฺธาวสิตฺโต, พฺราหฺมโณ วา มหาสาโล, โส ปุรตฺถิเมน นครสฺส นวํ สนฺถาคารํ การาเปตฺวา เกสมสฺสุํ โอหาเรตฺวา ขราชินํ นิวาเสตฺวา สปฺปิเตเลน กายํ อพฺภญฺชิตฺวา มควิสาเณน ปิฏฺฐิํ กณฺฑุวมาโน นวํ สนฺถาคารํ ปวิสติ สทฺธิํ มเหสิยา พฺราหฺมเณน จ ปุโรหิเตน, โส ตตฺถ อนนฺตรหิตาย ภูมิยา หริตุปลิตฺตาย เสยฺยํ กปฺเปติ, เอกิสฺสาย คาวิยา สรูปวจฺฉาย ยํ เอกสฺมิํ ถเน ขีรํ โหติ, เตน ราชา ยาเปติ, ยํ ทุติยสฺมิํ ถเน ขีรํ โหติ, เตน มเหสี ยาเปติ, ยํ ตติยสฺมิํ ถเน ขีรํ โหติ, เตน พฺราหฺมโณ ปุโรหิโต ยาเปติ, ยํ จตุตฺถสฺมิํ ถเน ขีรํ โหติ, เตน อคฺคิํ ชุหติ, อวเสเสน วจฺฉโก ยาเปติ.\csromanspacer{} โส เอวมาห “เอตฺตกา อุสภา หญฺญนฺตุ ยญฺญตฺถาย, เอตฺตกา วจฺฉตรา หญฺญนฺตุ ยญฺญตฺถาย, เอตฺตกา วจฺฉตริโย หญฺญนฺตุ ยญฺญตฺถาย, เอตฺตกา อชา หญฺญนฺตุ ยญฺญตฺถาย, เอตฺตกา อุรพฺภา หญฺญนฺตุ ยญฺญตฺถาย, เอตฺตกา อสฺสา หญฺญนฺตุ ยญฺญตฺถาย, เอตฺตกา รุกฺขา ฉิชฺชนฺตุ ยูปตฺถาย, เอตฺตกา ทพฺภา ลูยนฺตุ พริหิสตฺถายา”ติ.\csromanspacer{} เยปิสฺส เต โหนฺติ ทาสาติ วา เปสฺสาติ วา กมฺมกราติ วา, เตปิ ทณฺฑตชฺชิตา ภยตชฺชิตา อสฺสุมุขา รุทมานา ปริกมฺมานิ กโรนฺติ.\csromanspacer{} อยํ วุจฺจติ พฺราหฺมณ ปุคฺคโล อตฺตนฺตโป จ อตฺตปริตาปนานุโยคมนุยุตฺโต ปรนฺตโป จ ปรปริตาปนานุโยคมนุยุตฺโต.}

\csromanmatikaanchor{mtk.53}
\csromantocmarkref{subsection}{4. โฆฏมุขสุตฺต}{mtk.53}
\bookmark[dest={mtk.53},level=2]{4. โฆฏมุขสุตฺต}
\proseitem{418}{กตโม จ พฺราหฺมณ ปุคฺคโล เนวตฺตนฺตโป นาตฺตปริตาปนานุโยคมนุยุตฺโต น ปรนฺตโป น ปรปริตาปนานุยุคมนุยุตฺโต, โส อนตฺตนฺตโป อปรนฺตโป ทิฏฺเฐว ธมฺเม นิจฺฉาโต นิพฺพุโต สีตีภูโต สุขปฺปฏิสํเวที พฺรหฺมภูเตน อตฺตนา วิหรติ.\csromanspacer{} อิธ พฺราหฺมณ ตถาคโต โลเก อุปฺปชฺชติ อรหํ สมฺมาสมฺพุทฺโธ วิชฺชาจรณสมฺปนฺโน สุคโต โลกวิทู อนุตฺตโร ปุริสทมฺมสารถิ สตฺถา เทวมนุสฺสานํ พุทฺโธ ภควา, โส อิมํ โลกํ สเทวกํ สมารกํ สพฺรหฺมกํ สสฺสมณพฺราหฺมณิํ ปชํ สเทวมนุสฺสํ สยํ อภิญฺญา สจฺฉิกตฺวา ปเวเทติ, โส ธมฺมํ เทเสติ อาทิกลฺยาณํ \csromanfolio{370}มชฺเฌกลฺยาณํ ปริโยสานกลฺยาณํ สาตฺถํ สพฺยญฺชนํ เกวลปริปุณฺณํ ปริสุทฺธํ พฺรหฺมจริยํ ปกาเสติ, ตํ ธมฺมํ สุณาติ คหปติ วา คหปติปุตฺโต วา อญฺญตรสฺมิํ วา กุเล ปจฺจาชาโต, โส ตํ ธมฺมํ สุตฺวา ตถาคเต สทฺธํ ปฏิลภติ, โส เตน สทฺธาปฏิลาเภน สมนฺนาคโต อิติ ปฏิสญฺจิกฺขติ “สมฺพาโธ ฆราวาโส รโชปโถ, อพฺโภกาโส ปพฺพชฺชา, นยิทํ สุกรํ อคารํ อชฺฌาวสตา เอกนฺตปริปุณฺณํ เอกนฺตปริสุทฺธํ สงฺขลิขิตํ พฺรหฺมจริยํ จริตุํ, ยํนูนาหํ เกสมสฺสุํ โอหาเรตฺวา กาสายานิ วตฺถานิ อจฺฉาเทตฺวา อคารสฺมา อนคาริยํ ปพฺพเชยฺยนฺ”ติ, โส อปเรน สมเยน อปฺปํ วา โภคกฺขนฺธํ ปหาย มหนฺตํ วา โภคกฺขนฺธํ ปหาย อปฺปํ วา ญาติปริวฏฺฏํ ปหาย มหนฺตํ วา ญาติปริวฏฺฏํ ปหาย เกสมสฺสุํ โอหาเรตฺวา กาสายานิ วตฺถานิ อจฺฉาเทตฺวา อคารสฺมา อนคาริยํ ปพฺพชติ, โส เอวํ ปพฺพชิโต สมาโน ภิกฺขูนํ สิกฺขาสาชีวสมาปนฺโน ปาณาติปาตํ ปหาย ปาณาติปาตา ปฏิวิรโต โหติ, นิหิตทณฺโฑ นิหิตสตฺโถ ลชฺชี ทยาปนฺโน, สพฺพปาณภูตหิตานุกมฺปี วิหรติ.}

\prose{อทินฺนาทานํ ปหาย อทินฺนาทานา ปฏิวิรโต โหติ ทินฺนาทายี ทินฺนปาฏิกงฺขี, อเถเนน สุจิภูเตน อตฺตนา วิหรติ.}

\prose{อพฺรหฺมจริยํ ปหาย พฺรหฺมจารี โหติ อาราจารี วิรโต เมถุนา คามธมฺมา.}

\prose{มุสาวาทํ ปหาย มุสาวาทา ปฏิวิรโต โหติ สจฺจวาที สจฺจสนฺโธ เถโต ปจฺจยิโก อวิสํวาทโก โลกสฺส.}

\prose{ปิสุณํ วาจํ ปหาย ปิสุณาย วาจาย ปฏิวิรโต โหติ, อิโต สุตฺวา น อมุตฺร อกฺขาตา อิเมสํ เภทาย, อมุตฺร วา สุตฺวา น อิเมสํ อกฺขาตา อมูสํ เภทาย, อิติ ภินฺนานํ วา สนฺธาตา, สหิตานํ วา อนุปฺปทาตา, สมคฺคาราโม สมคฺครโต สมคฺคนนฺที สมคฺคกรณิํ วาจํ ภาสิตา โหติ.}

\prose{ผรุสํ วาจํ ปหาย ผรุสาย วาจาย ปฏิวิรโต โหติ, ยา สา วาจา เนลา กณฺณสุขา เปมนียา หทยงฺคมา โปรี พหุชนกนฺตา พหุชนมนาปา, ตถารูปิํ วาจํ ภาสิตา โหติ.}

\csromanheader{2}{4. โฆฏมุขสุตฺต (94)}
\prose{\csromanfolio{371}สมฺผปฺปลาปํ ปหาย สมฺผปฺปลาปา ปฏิวิรโต โหติ กาลวาที ภูตวาที อตฺถวาที ธมฺมวาที วินยวาที, นิธานวติํ วาจํ ภาสิตา กาเลน สาปเทสํ ปริยนฺตวติํ อตฺถสํหิตํ.}

\prose{โส พีชคามภูตคามสมารมฺภา ปฏิวิรโต โหติ.\csromanspacer{} เอกภตฺติโก โหติ รตฺตูปรโต วิรโต วิกาลโภชนา.\csromanspacer{} นจฺจคีตวาทิตวิสูกทสฺสนา ปฏิวิรโต โหติ.\csromanspacer{} มาลาคนฺธวิเลปนธารณมณฺฑนวิภูสนฏฺฐานา ปฏิวิรโต โหติ.\csromanspacer{} อุจฺจาสยนมหาสยนา ปฏิวิรโต โหติ.\csromanspacer{} ชาตรูปรชตปฏิคฺคหณา ปฏิวิรโต โหติ.\csromanspacer{} อามกธญฺญปฏิคฺคหณา ปฏิวิรโต โหติ.\csromanspacer{} อามกมํสปฏิคฺคหณา ปฏิวิรโต โหติ.\csromanspacer{} อิตฺถิกุมาริกปฏิคฺคหณา ปฏิวิรโต โหติ.\csromanspacer{} ทาสิทาสปฏิคฺคหณา ปฏิวิรโต โหติ.\csromanspacer{} อเชฬกปฏิคฺคหณา ปฏิวิรโต โหติ.\csromanspacer{} กุกฺกุฏสูกรปฏิคฺคหณา ปฏิวิรโต โหติ.\csromanspacer{} หตฺถิควสฺสวฬวปฏิคฺคหณา ปฏิวิรโต โหติ.\csromanspacer{} เขตฺตวตฺถุปฏิคฺคหณา ปฏิวิรโต โหติ.\csromanspacer{} ทูเตยฺยปหิณคมนานุโยคา ปฏิวิรโต โหติ.\csromanspacer{} กยวิกฺกยา ปฏิวิรโต โหติ.\csromanspacer{} ตุลากูฏกํสกูฏมานกูฏา ปฏิวิรโต โหติ.\csromanspacer{} อุกฺโกฏนวญฺจนนิกติสาจิโยคา ปฏิวิรโต โหติ.\csromanspacer{} เฉทนวธพนฺธนวิปราโมส-อาโลปสหสาการา ปฏิวิรโต โหติ.}

\prose{โส สนฺตุฏฺโฐ โหติ กายปริหาริเกน จีวเรน กุจฺฉิปริหาริเกน ปิณฺฑปาเตน, โส เยน เยเนว ปกฺกมติ, สมาทาเยว ปกฺกมติ.\csromanspacer{} เสยฺยถาปิ นาม ปกฺขี สกุโณ เยน เยเนว เฑติ, สปตฺตภาโรว เฑติ.\csromanspacer{} เอวเมว ภิกฺขุ สนฺตุฏฺโฐ โหติ กายปริหาริเกน จีวเรน กุจฺฉิปริหาริเกน ปิณฺฑปาเตน, โส เยน เยเนว ปกฺกมติ, สมาทาเยว ปกฺกมติ.\csromanspacer{} โส อิมินา อริเยน สีลกฺขนฺเธน สมนฺนาคโต อชฺฌตฺตํ อนวชฺชสุขํ ปฏิสํเวเทติ.}

\proseitem{419}{โส จกฺขุนา รูปํ ทิสฺวา น นิมิตฺตคฺคาหี โหติ นานุพฺยญฺชนคฺคาหี, ยตฺวาธิกรณเมนํ จกฺขุนฺทฺริยํ อสํวุตํ วิหรนฺตํ อภิชฺฌา โทมนสฺสา ปาปกา อกุสลา ธมฺมา อนฺวาสฺสเวยฺยุํ, ตสฺส สํวราย ปฏิปชฺชติ รกฺขติ จกฺขุนฺทฺริยํ, จกฺขุนฺทฺริเย สํวรํ อาปชฺชติ.\csromanspacer{} โสเตน สทฺทํ สุตฺวา ฯเปฯ.\csromanspacer{} ฆาเนน คนฺธํ ฆายิตฺวา.\csromanspacer{} ชิวฺหาย รสํ สายิตฺวา.\csromanspacer{} กาเยน โผฏฺฐพฺพํ ผุสิตฺวา.\csromanspacer{} มนสา ธมฺมํ วิญฺญาย น นิมิตฺตคฺคาหี โหติ นานุพฺยญฺชนคฺคาหี, \csromanfolio{372}ยตฺวาธิกรณเมนํ มนินฺทฺริยํ อสํวุตํ วิหรนฺตํ อภิชฺฌาโทมนสฺสา ปาปกา อกุสลา ธมฺมา อนฺวาสฺสเวยฺยุํ, ตสฺส สํวราย ปฏิปชฺชติ, รกฺขติ มนินฺทฺริยํ, มนินฺทฺริเย สํวรํ อาปชฺชติ, โส อิมินา อริเยน อินฺทฺริยสํวเรน สมนฺนาคโต อชฺฌตฺตํ อพฺยาเสกสุขํ ปฏิสํเวเทติ.}

\prose{โส อภิกฺกนฺเต ปฏิกฺกนฺเต สมฺปชานการี โหติ, อาโลกิเต วิโลกิเต สมฺปชานการี โหติ, สมิญฺชิเต ปสาริเต สมฺปชานการี โหติ, สํฆาฏิปตฺตจีวรธารเณ สมฺปชานการี โหติ, อสิเต ปีเต ขายิเต สายิเต สมฺปชานการี โหติ, อุจฺจารปสฺสาวกมฺเม สมฺปชานการี โหติ, คเต ฐิเต นิสินฺเน สุตฺเต ชาคริเต ภาสิเต ตุณฺหีภาเว สมฺปชานการี โหติ.}

\prose{โส อิมินา จ อริเยน สีลกฺขนฺเธน สมนฺนาคโต, (อิมาย จ อริยาย สนฺตุฏฺฐิยา สมนฺนาคโต,)\footnote{ปสฺส ม 1. 239 ปิฏฺเฐ.} อิมินา จ อริเยน อินฺทฺริยสํวเรน สมนฺนาคโต, อิมินา จ อริเยน สติสมฺปชญฺเญน สมนฺนาคโต วิวิตฺตํ เสนาสนํ ภชติ อรญฺญํ รุกฺขมูลํ ปพฺพตํ กนฺทรํ คิริคุหํ สุสานํ วนปตฺถํ อพฺโภกาสํ ปลาลปุญฺชํ.\csromanspacer{} โส ปจฺฉาภตฺตํ ปิณฺฑปาตปฏิกฺกนฺโต นิสีทติ ปลฺลงฺกํ อาภุชิตฺวา อุชุํ กายํ ปณิธาย ปริมุขํ สติํ อุปฏฺฐเปตฺวา.\csromanspacer{} โส อภิชฺฌํ โลเก ปหาย วิคตาภิชฺเฌน เจตสา วิหรติ, อภิชฺฌาย จิตฺตํ ปริโสเธติ.\csromanspacer{} พฺยาปาทปโทสํ ปหาย อพฺยาปนฺนจิตฺโต วิหรติ สพฺพปาณภูตหิตานุกมฺปี, พฺยาปาทปโทสา จิตฺตํ ปริโสเธติ.\csromanspacer{} ถินมิทฺธํ ปหาย วิคตถินมิทฺโธ วิหรติ อาโลกสญฺญี สโต สมฺปชาโน, ถินมิทฺธา จิตฺตํ ปริโสเธติ.\csromanspacer{} อุทฺธจฺจกุกฺกุจฺจํ ปหาย อนุทฺธโต วิหรติ อชฺฌตฺตํ วูปสนฺตจิตฺโต, อุทฺธจฺจกุกฺกุจฺจา จิตฺตํ ปริโสเธติ.\csromanspacer{} วิจิกิจฺฉํ ปหาย ติณฺณวิจิกิจฺโฉ วิหรติ อกถํกถี กุสเลสุ ธมฺเมสุ, วิจิกิจฺฉาย จิตฺตํ ปริโสเธติ.}

\prose{โส อิเม ปญฺจ นีวรเณ ปหาย เจตโส อุปกฺกิเลเส ปญฺญาย ทุพฺพลีกรเณ วิวิจฺเจว กาเมหิ วิวิจฺจ อกุสเลหิ ธมฺเมหิ สวิตกฺกํ สวิจารํ วิเวกชํ ปีติสุขํ ปฐมํ ฌานํ อุปสมฺปชฺช วิหรติ.\csromanspacer{} วิตกฺกวิจารานํ วูปสมา อชฺฌตฺตํ สมฺปสาทนํ เจตโส เอโกทิภาวํ อวิตกฺกํ อวิจารํ สมาธิชํ ปีติสุขํ ทุติยํ ฌานํ อุปสมฺปชฺช วิหรติ.\csromanspacer{} ปีติยา จ วิราคา อุเปกฺขโก จ วิหรติ สโต จ สมฺปชาโน, สุขญฺจ กาเยน}

\vspace{\tipitakaparskip}
\csromancenter{โฆฏมุขสุตฺต (94) ปฏิสํเวเทติ, ยํ ตํ อริยา อาจิกฺขนฺติ “อุเปกฺขโก สติมา สุขวิหารี”ติ, ตติยํ ฌานํ อุปสมฺปชฺช วิหรติ.\csromanspacer{} สุขสฺส จ ปหานา ทุกฺขสฺส จ ปหานา ปุพฺเพว โสมนสฺสโทมนสฺสานํ อตฺถงฺคมา อทุกฺขมสุขํ อุเปกฺขาสติปาริสุทฺธิํ จตุตฺถํ ฌานํ อุปสมฺปชฺช วิหรติ.}
\proseitem{420}{\csromanfolio{373}โส เอวํ สมาหิเต จิตฺเต ปริสุทฺเธ ปริโยทาเต อนงฺคเณ วิคตูปกฺกิเลเส มุทุภูเต กมฺมนิเย ฐิเต อาเนญฺชปฺปตฺเต ปุพฺเพนิวาสานุสฺสติญาณาย จิตฺตํ อภินินฺนาเมติ, โส อเนกวิหิตํ ปุพฺเพนิวาสํ อนุสฺสรติ, เสยฺยถิทํ, เอกมฺปิ ชาติํ เทฺวปิ ชาติโย ติสฺโสปิ ชาติโย จตสฺโสปิ ชาติโย ปญฺจปิ ชาติโย ทสปิ ชาติโย วีสมฺปิ ชาติโย ติํสมฺปิ ชาติโย จตฺตาลีสมฺปิ ชาติโย ปญฺญาสมฺปิ ชาติโย ชาติสตมฺปิ ชาติสหสฺสมฺปิ ชาติสตสหสฺสมฺปิ อเนเกปิ สํวฏฺฏกปฺเป อเนเกปิ วิวฏฺฏกปฺเป อเนเกปิ สํวฏฺฏวิวฏฺฏกปฺเป “อมุตฺราสิํ เอวํนาโม เอวํโคตฺโต เอวํวณฺโณ เอวมาหาโร เอวํสุขทุกฺขปฺปฏิสํเวที เอวมายุปริยนฺโต, โส ตโต จุโต อมุตฺร อุทปาทิํ, ตตฺราปาสิํ เอวํนาโม เอวํโคตฺโต เอวํวณฺโณ เอวมาหาโร เอวํสุขทุกฺขปฺปฏิสํเวที เอวมายุปริยนฺโต, โส ตโต จุโต อิธูปปนฺโน”ติ, อิติ สาการํ ส-อุทฺเทสํ อเนกวิหิตํ ปุพฺเพนิวาสํ อนุสฺสรติ.}

\prose{โส เอวํ สมาหิเต จิตฺเต ปริสุทฺเธ ปริโยทาเต อนงฺคเณ วิคตูปกฺกิเลเส มุทุภูเต กมฺมนิเย ฐิเต อาเนญฺชปฺปตฺเต สตฺตานํ จุตูปปาตญาณาย จิตฺตํ อภินินฺนาเมติ, โส ทิพฺเพน จกฺขุนา วิสุทฺเธน อติกฺกนฺตมานุสเกน สตฺเต ปสฺสติ จวมาเน อุปปชฺชมาเน หีเน ปณีเต สุวณฺเณ ทุพฺพณฺเณ สุคเต ทุคฺคเต, ยถากมฺมูปเค สตฺเต ปชานาติ “อิเม วต โภนฺโต สตฺตา กายทุจฺจริเตน สมนฺนาคตา ฯเปฯ อริยานํ อุปวาทกา มิจฺฉาทิฏฺฐิกา มิจฺฉาทิฏฺฐิกมฺมสมาทานา, เต กายสฺส เภทา ปรํ มรณา อปายํ ทุคฺคติํ วินิปาตํ นิรยํ อุปปนฺนา.\csromanspacer{} อิเม วา ปน โภนฺโต สตฺตา กายสุจริเตน สมนฺนาคตา ฯเปฯ อริยานํ อนุปวาทกา สมฺมาทิฏฺฐิกา สมฺมาทิฏฺฐิกมฺมสมาทานา, เต กายสฺส เภทา ปรํ มรณา สุคติํ สคฺคํ โลกํ อุปปนฺนา”ติ, อิติ ทิพฺเพน จกฺขุนา วิสุทฺเธน อติกฺกนฺตมานุสเกน สตฺเต ปสฺสติ จวมาเน อุปปชฺชมาเน หีเน ปณีเต สุวณฺเณ ทุพฺพณฺเณ สุคเต ทุคฺคเต, ยถากมฺมูปเค สตฺเต ปชานาติ.}

\prose{\csromanfolio{374}โส เอวํ สมาหิเต จิตฺเต ปริสุทฺเธ ปริโยทาเต อนงฺคเณ วิคตูปกฺกิเลเส มุทุภูเต กมฺมนิเย ฐิเต อาเนญฺชปฺปตฺเต อาสวานํ ขยญาณาย จิตฺตํ อภินินฺนาเมติ, โส อิทํ ทุกฺขนฺติ ยถาภูตํ ปชานาติ, อยํ ทุกฺขสมุทโยติ ยถาภูตํ ปชานาติ, อยํ ทุกฺขนิโรโธติ ยถาภูตํ ปชานาติ, อยํ ทุกฺขนิโรธคามินี ปฏิปทาติ ยถาภูตํ ปชานาติ.\csromanspacer{} อิเม อาสวาติ ยถาภูตํ ปชานาติ, อยํ อาสวสมุทโยติ ยถาภูตํ ปชานาติ, อยํ อาสวนิโรโธติ ยถาภูตํ ปชานาติ, อยํ อาสวนิโรธคามินี ปฏิปทาติ ยถาภูตํ ปชานาติ.\csromanspacer{} ตสฺส เอวํ ชานโต เอวํ ปสฺสโต กามาสวาปิ จิตฺตํ วิมุจฺจติ, ภวาสวาปิ จิตฺตํ วิมุจฺจติ, อวิชฺชาสวาปิ จิตฺตํ วิมุจฺจติ, วิมุตฺตสฺมิํ วิมุตฺตมิติ ญาณํ โหติ, “ขีณา ชาติ, วุสิตํ พฺรหฺมจริยํ, กตํ กรณียํ, นาปรํ อิตฺถตฺตายา”ติ ปชานาติ.}

\prose{อยํ วุจฺจติ พฺราหฺมณ ปุคฺคโล เนวตฺตนฺตโป นาตฺตปริตาปนานุโยคมนุยุตฺโต น ปรนฺตโป น ปรปริตาปนานุโยคมนุยุตฺโต, โส อนตฺตนฺตโป อปรนฺตโป ทิฏฺเฐว ธมฺเม นิจฺฉาโต นิพฺพุโต สีตีภูโต สุขปฺปฏิสํเวที พฺรหฺมภูเตน อตฺตนา วิหรตีติ.}

\proseitem{421}{เอวํ วุตฺเต โฆฏมุโข พฺราหฺมโณ อายสฺมนฺตํ อุเทนํ เอตทโวจ “อภิกฺกนฺตํ โภ อุเทน, อภิกฺกนฺตํ โภ อุเทน, เสยฺยถาปิ โภ อุเทน นิกฺกุชฺชิตํ วา อุกฺกุชฺเชยฺย, ปฏิจฺฉนฺนํ วา วิวเรยฺย, มูฬฺหสฺส วา มคฺคํ อาจิกฺเขยฺย, อนฺธกาเร วา เตลปชฺโชตํ ธาเรยฺย ‘จกฺขุมนฺโต รูปานิ ทกฺขนฺตี’ติ.\csromanspacer{} เอวเมวํ โภตา อุเทเนน อเนกปริยาเยน ธมฺโม ปกาสิโต, เอสาหํ ภวนฺตํ อุเทนํ สรณํ คจฺฉามิ ธมฺมญฺจ ภิกฺขุสํฆญฺจ, อุปาสกํ มํ ภวํ อุเทโน ธาเรตุ อชฺชตคฺเค ปาณุเปตํ สรณํ คตนฺ”ติ.\csromanspacer{} มา โข มํ ตฺวํ พฺราหฺมณ สรณํ อคมาสิ, ตเมว ภควนฺตํ สรณํ คจฺฉาหิ, ยมหํ สรณํ คโตติ.\csromanspacer{} กหํ ปน โภ อุเทน เอตรหิ โส ภวํ โคตโม วิหรติ อรหํ สมฺมาสมฺพุทฺโธติ.\csromanspacer{} ปรินิพฺพุโต โข พฺราหฺมณ เอตรหิ โส ภควา อรหํ สมฺมาสมฺพุทฺโธติ.}

\prose{สเจปิ\footnote{สเจ หิ (สี, สฺยา, กํ, อิ)} มยํ โภ อุเทน สุเณยฺยาม ตํ ภวนฺตํ โคตมํ ทสสุ โยชเนสุ, ทสปิ มยํ โยชนานิ คจฺเฉยฺยาม ตํ ภวนฺตํ โคตมํ}

\vspace{\tipitakaparskip}
\csromancenter{จงฺกีสุตฺต (95) ทสฺสนาย อรหนฺตํ สมฺมาสมฺพุทฺธํ.\csromanspacer{} สเจปิ\footnote{สเจ (สี, อิ), สเจ หิ (สฺยา, กํ)} มยํ โภ อุเทน สุเณยฺยาม ตํ ภวนฺตํ โคตมํ วีสติยา โยชเนสุ.\csromanspacer{} ติํสาย โยชเนสุ.\csromanspacer{} จตฺตารีสาย โยชเนสุ.\csromanspacer{} ปญฺญาสาย โยชเนสุ, ปญฺญาสมฺปิ มยํ โยชนานิ คจฺเฉยฺยาม ตํ ภวนฺตํ โคตมํ ทสฺสนาย อรหนฺตํ สมฺมาสมฺพุทฺธํ.\csromanspacer{} โยชนสเต เจปิ\footnote{โยชนสเตปิ (สี, สฺยา, กํ, อิ)} มยํ โภ อุเทน สุเณยฺยาม ตํ ภวนฺตํ โคตมํ โยชนสตมฺปิ มยํ คจฺเฉยฺยาม ตํ ภวนฺตํ โคตมํ ทสฺสนาย อรหนฺตํ สมฺมาสมฺพุทฺธํ.}
\prose{\csromanfolio{375}ยโต จ โข โภ อุเทน ปรินิพฺพุโต โส ภวํ โคตโม, ปรินิพฺพุตมฺปิ มยํ ตํ ภวนฺตํ โคตมํ สรณํ คจฺฉาม ธมฺมญฺจ ภิกฺขุสํฆญฺจ, อุปาสกํ มํ ภวํ อุเทโน ธาเรตุ อชฺชตคฺเค ปาณุเปตํ สรณํ คตํ.\csromanspacer{} อตฺถิ จ เม โภ อุเทน องฺคราชา เทวสิกํ นิจฺจภิกฺขํ ททาติ, ตโต อหํ โภโต อุเทนสฺส เอกํ นิจฺจภิกฺขํ ททามีติ.\csromanspacer{} กิํ ปน เต พฺราหฺมณ องฺคราชา เทวสิกํ นิจฺจภิกฺขํ ททาตีติ.\csromanspacer{} ปญฺจ โภ อุเทน กหาปณสตานีติ.\csromanspacer{} น โข โน พฺราหฺมณ กปฺปติ ชาตรูปรชตํ ปฏิคฺคเหตุนฺติ.\csromanspacer{} สเจ ตํ โภโต อุเทนสฺส น กปฺปติ, วิหารํ โภโต อุเทนสฺส การาเปสฺสามีติ.\csromanspacer{} สเจ โข เม ตฺวํ พฺราหฺมณ วิหารํ การาเปตุ กาโม, ปาฏลิปุตฺเต สํฆสฺส อุปฏฺฐานสาลํ การาเปหีติ.\csromanspacer{} อิมินาปาหํ โภโต อุเทนสฺส ภิยฺโยโส มตฺตาย อตฺตมโน อภิรทฺโธ, ยํ มํ ภวํ อุเทโน สํเฆ ทาเน สมาทเปติ, เอสาหํ โภ อุเทน เอติสฺสา จ นิจฺจภิกฺขาย อปราย จ นิจฺจภิกฺขาย ปาฏลิปุตฺเต สํฆสฺส อุปฏฺฐานสาลํ การาเปสฺสามีติ.\csromanspacer{} อถ โข โฆฏมุโข พฺราหฺมโณ เอติสฺสา จ นิจฺจาภิกฺขาย อปราย จ นิจฺจภิกฺขาย ปาฏลิปุตฺเต สํฆสฺส อุปฏฺฐานสาลํ การาเปสิ.\csromanspacer{} สา เอตรหิ โฆฏมุขีติ วุจฺจตีติ.}

\nitthitamruled{โฆฏมุขสุตฺตํ นิฏฺฐิตํ จตุตฺถํ.}
\csromanheader{2}{5. จงฺกีสุตฺต}
\proseitem{422}{เอวํ เม สุตํ\csromandash{}เอกํ สมยํ ภควา โกสเลสุ จาริกํ จรมาโน มหตา ภิกฺขุสํเฆน สทฺธิํ เยน โอปาสาทํ นาม \csromanfolio{376}โกสลานํ พฺราหฺมณคาโม ตทวสริ, ตตฺร สุทํ ภควา โอปาสาเท วิหรติ อุตฺตเรน โอปาสาทํ เทววเน สาลวเน.\csromanspacer{} เตน โข ปน สมเยน จงฺกี พฺราหฺมโณ โอปาสาทํ อชฺฌาวสติ สตฺตุสฺสทํ สติณกฏฺโฐทกํ สธญฺญํ ราชโภคฺคํ รญฺญา ปเสนทินา โกสเลน ทินฺนํ ราชทายํ พฺรหฺมเทยฺยํ.\csromanspacer{} อสฺโสสุํ โข โอปาสาทกา พฺราหฺมณคหปติกา “สมโณ ขลุ โภ โคตโม สกฺยปุตฺโต สกฺยกุลา ปพฺพชิโต โกสเลสุ จาริกํ จรมาโน มหตา ภิกฺขุสํเฆน สทฺธิํ โอปาสาทํ อนุปฺปตฺโต โอปาสาเท วิหรติ อุตฺตเรน โอปาสาทํ เทววเน สาลวเน.\csromanspacer{} ตํ โข ปน ภวนฺตํ โคตมํ เอวํ กลฺยาโณ กิตฺติสทฺโท อพฺภุคฺคโต ‘อิติปิ โส ภควา อรหํ สมฺมาสมฺพุทฺโธ วิชฺชาจรณสมฺปนฺโน สุคโต โลกวิทู อนุตฺตโร ปุริสทมฺมสารถิ สตฺถา เทวมนุสฺสานํ พุทฺโธ ภควา’ติ, โส อิมํ โลกํ สเทวกํ สมารกํ สพฺรหฺมกํ สสฺสมณพฺราหฺมณิํ ปชํ สเทวมนุสฺสํ สยํ อภิญฺญา สจฺฉิกตฺวา ปเวเทติ, โส ธมฺมํ เทเสติ อาทิกลฺยาณํ มชฺเฌกลฺยาณํ ปริโยสานกลฺยาณํ สาตฺถํ สพฺยญฺชนํ เกวลปริปุณฺณํ ปริสุทฺธํ พฺรหฺมจริยํ ปกาเสติ, สาธุ โข ปน ตถารูปานํ อรหตํ ทสฺสนํ โหตี”ติ.}

\proseitem{423}{อถ โข โอปาสาทกา พฺราหฺมณคหปติกา โอปาสาทา นิกฺขมิตฺวา สงฺฆสงฺฆี คณีภูตา อุตฺตเรนมุขา คจฺฉนฺติ เยน เทววนํ สาลวนํ.\csromanspacer{} เตน โข ปน สมเยน จงฺกี พฺราหฺมโณ อุปริปาสาเท ทิวาเสยฺยํ อุปคโต.\csromanspacer{} อทฺทสา โข จงฺกี พฺราหฺมโณ โอปาสาทเก พฺราหฺมณคหปติเก โอปาสาทา นิกฺขมิตฺวา สงฺฆสงฺฆี คณีภูเต อุตฺตเรนมุขํ เยน เทววนํ สาลวนํ เตนุปสงฺกมนฺเต, ทิสฺวา ขตฺตํ อามนฺเตสิ “กิํ นุ โข โภ ขตฺเต โอปาสาทกา พฺราหฺมณคหปติกา โอปาสาทา นิกฺขมิตฺวา สงฺฆสงฺฆี คณีภูตา อุตฺตเรนมุขา คจฺฉนฺติ เยน เทววนํ สาลวนนฺ”ติ.\csromanspacer{} อตฺถิ โภ จงฺกี สมโณ โคตโม สกฺยปุตฺโต สกฺยกุลา ปพฺพชิโต โกสเลสุ จาริกํ จรมาโน มหตา ภิกฺขุสํเฆน สทฺธิํ โอปาสาทํ อนุปฺปตฺโต โอปาสาเท วิหรติ อุตฺตเรน โอปาสาทํ เทววเน สาลวเน.\csromanspacer{} ตํ โข ปน ภวนฺตํ โคตมํ เอวํ กลฺยาโณ กิตฺติสทฺโท อพฺภุคฺคโต “อิติปิ โส ภควา อรหํ สมฺมาสมฺพุทฺโธ วิชฺชาวรณสมฺปนฺโน สุคโต โลกวิทู อนุตฺตโร ปุริสทมฺมสารถิ สตฺถา เทวมนุสฺสานํ พุทฺโธ ภควา”ติ.\csromanspacer{} ตเมเต}

\vspace{\tipitakaparskip}
\csromancenter{จงฺกีสุตฺต (95) ภวนฺตํ โคตมํ ทสฺสนาย คจฺฉนฺตีติ.\csromanspacer{} เตน หิ โภ ขตฺเต เยน โอปาสาทกา พฺราหฺมณคหปติกา เตนุปสงฺกม, อุปสงฺกมิตฺวา โอปาสาทเก พฺราหฺมณคหปติเก เอวํ วเทหิ “จงฺกี โภ พฺราหฺมโณ เอวมาห อาคเมนฺตุ กิร โภนฺโต, จงฺกีปิ พฺราหฺมโณ สมณํ โคตมํ ทสฺสนาย อุปสงฺกมิสฺสตี”ติ. “เอวํ โภ”ติ โข โส ขตฺโต จงฺกิสฺส พฺราหฺมณสฺส ปฏิสฺสุตฺวา เยน โอปาสาทกา พฺราหฺมณคหปติกา เตนุปสงฺกมิ, อุปสงฺกมิตฺวา โอปาสาทเก พฺราหฺมณคหปติเก เอตทโวจ “จงฺกี โภ พฺราหฺมโณ เอวมาห อาคเมนฺตุ กิร โภนฺโต, จงฺกีปิ พฺราหฺมโณ สมณํ โคตมํ ทสฺสนาย อุปสงฺกมิสฺสตี”ติ.}
\proseitem{424}{\csromanfolio{377}เตน โข ปน สมเยน นานาเวรชฺชกานํ พฺราหฺมณานํ ปญฺจมตฺตานิ พฺราหฺมสตานิ โอปาสาเท ปฏิวสนฺติ เกนจิเทว กรณีเยน, อสฺโสสุํ โข เต พฺราหฺมณา “จงฺกี กิร พฺราหฺมโณ สมณํ โคตมํ ทสฺสนาย อุปสงฺกมิสฺสตี”ติ.\csromanspacer{} อถ โข เต พฺราหฺมณา เยน จงฺกี พฺราหฺมโณ เตนุปสงฺกมิํสุ, อุปสงฺกมิตฺวา จงฺกิํ พฺราหฺมณํ เอตทโวจุํ “สจฺจํ กิร ภวํ จงฺกี สมณํ โคตมํ ทสฺสนาย อุปสงฺกมิสฺสตี”ติ.\csromanspacer{} เอวํ โข เม โภ โหติ “อหํ สมณํ โคตมํ ทสฺสนาย อุปสงฺกมิสฺสามี”ติ.\csromanspacer{} มา ภวํ จงฺกี สมณํ โคตมํ ทสฺสนาย อุปสงฺกมิ, น อรหติ ภวํ จงฺกี สมณํ โคตมํ ทสฺสนาย อุปสงฺกมิตุํ, สมโณเตฺวว โคตโม อรหติ ภวนฺตํ จงฺกิํ ทสฺสนาย อุปสงฺกมิตุํ.\csromanspacer{} ภวํ หิ จงฺกี อุภโต สุชาโต มาติโต จ ปิติโต จ สํสุทฺธคหณิโก ยาว สตฺตมา ปิตามหยุคา อกฺขิตฺโต อนุปกฺกุฏฺโฐ ชาติวาเทน.\csromanspacer{} ยมฺปิ ภวํ จงฺกี อุภโต สุชาโต มาติโต จ ปิติโต จ สํสุทฺธคหณิโก ยาว สตฺตมา ปิตามหยุคา อกฺขิตฺโต อนุปกฺกุฏฺโฐ ชาติวาเทน, อิมินาปงฺเคน น อรหติ ภวํ จงฺกี สมณํ โคตมํ ทสฺสนาย อุปสงฺกมิตุํ, สมโณ เตฺวว โคตโม อรหติ ภวนฺตํ จงฺกิํ ทสฺสนาย อุปสงฺกมิตุํ.\csromanspacer{} ภวํ หิ จงฺกี อฑฺโฒ มหทฺธโน มหาโภโค ฯเปฯ.\csromanspacer{} ภวํ หิ จงฺกี ติณฺณํ เวทานํ ปารคู สนิฆณฺฑุเกฏุภานํ สากฺขรปฺปเภทานํ อิติหาสปญฺจมานํ ปทโก เวยฺยากรโณ โลกายตมหาปุริสลกฺขเณสุ อนวโย ฯเปฯ.\csromanspacer{} ภวํ หิ จงฺกี อภิรูโป ทสฺสนีโย ปาสาทิโก ปรมาย วณฺณโปกฺขรตาย \csromanfolio{378}สมนฺนาคโต พฺรหฺมวณฺณี พฺรหฺมวจฺฉสี\footnote{พฺรหฺมวจฺจสี (สี, อิ)} อขุทฺทาวกาโส ทสฺสนาย ฯเปฯ.\csromanspacer{} ภวํ หิ จงฺกี สีลวา วุทฺธสีลี วุทฺธสีเลน สมนฺนาคโต ฯเปฯ.\csromanspacer{} ภวํ หิ จงฺกี กลฺยาณวาโจ กลฺยาณวากฺกรโณ โปริยา วาจาย สมนฺนาคโต วิสฺสฏฺฐาย อเนลคลาย อตฺถสฺส วิญฺญาปนิยา ฯเปฯ.\csromanspacer{} ภวํ หิ จงฺกี พหูนํ อาจริยปาจริโย ตีณิ มาณวกสตานิ มนฺเต วาเจติ ฯเปฯ.\csromanspacer{} ภวํ หิ จงฺกี รญฺโญ ปเสนทิสฺส โกสลสฺส สกฺกโต ครุกโต มานิโต ปูชิโต อปจิโต ฯเปฯ.\csromanspacer{} ภวํ หิ จงฺกี พฺราหฺมณสฺส โปกฺขรสาติสฺส สกฺกโต ครุกโต มานิโต ปูชิโต อปจิโต ฯเปฯ.\csromanspacer{} ภวํ หิ จงฺกี โอปาสาทํ อชฺฌาวสติ สตฺตุสฺสทํ สติณกฏฺโฐทกํ สธญฺญํ ราชโภคฺคํ รญฺญา ปเสนทินา โกสเลน ทินฺนํ ราชทายํ พฺรหฺมเทยฺยํ.\csromanspacer{} ยมฺปิ ภวํ จงฺกี โอปาสาทํ อชฺฌาวสติ สตฺตุสฺสทํ สติณกฏฺโฐทกํ สธญฺญํ ราชโภคฺคํ รญฺญา ปเสนทินา โกสเลน ทินฺนํ ราชทายํ พฺรหฺมเทยฺยํ, อิมินาปงฺเคน น อรหติ ภวํ จงฺกี สมณํ โคตมํ ทสฺสนาย อุปสงฺกมิตุํ, สมโณเตฺวว โคตโม อรหติ ภวนฺตํ จงฺกิํ ทสฺสนาย อุปสงฺกมิตุนฺติ.}

\proseitem{425}{เอวํ วุตฺเต จงฺกี พฺราหฺมโณ เต พฺราหฺมเณ เอตทโวจ “เตน หิ โภ มมปิ สุณาถ, ยถา มยเมว อรหาม ตํ สมณํ โคตมํ ทสฺสนาย อุปสงฺกมิตุํ, นเตฺวว อรหติ โส ภวํ โคตโม อมฺหากํ ทสฺสนาย อุปสงฺกมิตุํ.\csromanspacer{} สมโณ ขลุ โภ โคตโม อุภโต สุชาโต มาติโต จ ปิติโต จ สํสุทฺธคหณิโก ยาว สตฺตมา ปิตามหยุคา อกฺขิตฺโต อนุปกฺกุฏฺโฐ ชาติวาเทน.\csromanspacer{} ยมฺปิ โภ สมโณ โคตโม อุภโต สุชาโต มาติโต จ ปิติโต จ สํสุทฺธคหณิโก ยาว สตฺตมา ปิตามหยุคา อกฺขิตฺโต อนุปกฺกุฏฺโฐ ชาติวาเทน, อิมินาปงฺเคน น อรหติ โส ภวํ โคตโม อมฺหากํ ทสฺสนาย อุปสงฺกมิตุํ, อถ โข มยเมว อรหาม ตํ ภวนฺตํ โคตมํ ทสฺสานาย อุปสงฺกมิตุํ\footnote{เอตฺถ ที 1. 108 ปิฏฺเฐ อญฺญมฺปิ คุณปทํ ทิสฺสติ.}.\csromanspacer{} สมโณ ขลุ โภ โคตโม ปหูตํ หิรญฺญสุวณฺณํ โอหาย ปพฺพชิโต ภูมิคตญฺจ เวหาสฏฺฐญฺจ ฯเปฯ.\csromanspacer{} สมโณ ขลุ โภ โคตโม ทหโรว สมาโน ยุวา สุสุกาฬเกโส ภเทฺรน โยพฺพเนน สมนฺนาคโต ปฐเมน วยสา อคารสฺมา อนคาริยํ ปพฺพชิโต ฯเปฯ.\csromanspacer{} สมโณ ขลุ โภ โคตโม}

\vspace{\tipitakaparskip}
\csromancenter{จงฺกีสุตฺต (95) อกามกานํ มาตาปิตูนํ อสฺสุมุขานํ รุทนฺตานํ เกสมสฺสุํ โอหาเรตฺวา กาสายานิ วตฺถานิ อจฺฉาเทตฺวา อคารสฺมา อนคาริยํ ปพฺพชิโต ฯเปฯ.\csromanspacer{} สมโณ ขลุ โภ โคตโม อภิรูโป ทสฺสนีโย ปาสาทิโก ปรมาย วณฺณโปกฺขรตาย สมนฺนาคโต พฺรหฺมวณฺณี พฺรหฺมวจฺฉสี อขุทฺทาวกาโส ทสฺสนาย ฯเปฯ.\csromanspacer{} สมโณ ขลุ โภ โคตโม สีลวา อริยสีลี กุสลสีลี กุสเลน สีเลน สมนฺนาคโต ฯเปฯ.\csromanspacer{} สมโณ ขลุ โภ โคตโม กลฺยาณวาโจ กลฺยาณวากฺกรโณ โปริยา วาจาย สมนฺนาคโต วิสฺสฏฺฐาย อเนลคลาย อตฺถสฺส วิญฺญาปนิยา ฯเปฯ.\csromanspacer{} สมโณ ขลุ โภ โคตโม พหูนํ อาจริยปาจริโย ฯเปฯ.\csromanspacer{} สมโณ ขลุ โภ โคตโม ขีณกามราโค วิคตจาปลฺโล ฯเปฯ.\csromanspacer{} สมโณ ขลุ โภ โคตโม กมฺมวาที กิริยวาที อปาปปุเรกฺขาโร พฺรหฺมญฺญาย ปชาย ฯเปฯ.\csromanspacer{} สมโณ ขลุ โภ โคตโม อุจฺจา กุลา ปพฺพชิโต อสมฺภินฺนา ขตฺติยกุลา ฯเปฯ.\csromanspacer{} สมโณ ขลุ โภ โคตโม อฑฺฒา กุลา ปพฺพชิโต มหทฺธนา มหาโภคา ฯเปฯ.\csromanspacer{} สมณํ ขลุ โภ โคตมํ ติโรรฏฺฐา ติโรชนปทา สํปุจฺฉิตุํ อาคจฺฉนฺติ ฯเปฯ.\csromanspacer{} สมณํ ขลุ โภ โคตมํ อเนกานิ เทวตาสหสฺสานิ ปาเณหิ สรณํ คตานิ ฯเปฯ.\csromanspacer{} สมณํ ขลุ โภ โคตมํ เอวํ กลฺยาโณ กิตฺติสทฺโท อพฺภุคฺคโต “อิติปิ โส ภควา อรหํ สมฺมาสมฺพุทฺโธ วิชฺชาจรณสมฺปนฺโน สุคโต โลกวิทู อนุตฺตโร ปุริสทมฺมสารถิ สตฺถา เทวมนุสฺสานํ พุทฺโธ ภควา”ติ ฯเปฯ.\csromanspacer{} สมโณ ขลุ โภ โคตโม ทฺวตฺติํสมหาปุริสลกฺขเณหิ สมนฺนาคโต ฯเปฯ.\footnote{เอตฺถาปิ ที 1. 109 ปิฏฺเฐ อญฺญานิปิ คุณปทานิ ทิสฺสนฺติ.} สมณํ ขลุ โภ โคตมํ ราชา มาคโธ เสนิโย พิมฺพิสาโร สปุตฺตทาโร ปาเณหิ สรณํ คโต ฯเปฯ.\csromanspacer{} สมณํ ขลุ โภ โคตมํ ราชา ปเสนทิ โกสโล สปุตฺตทาโร ปาเณหิ สรณํ คโต ฯเปฯ.\csromanspacer{} สมณํ ขลุ โภ โคตมํ พฺราหฺมโณ โปกฺขรสาติ สปุตฺตทาโร ปาเณหิ สรณํ คโต ฯเปฯ.\csromanspacer{} สมโณ ขลุ โภ โคตโม โอปาสาทํ อนุปฺปตฺโต โอปาสาเท วิหรติ อุตฺตเรน โอปาสาทํ เทววเน สาลวเน, เย โข เต สมณา วา พฺราหฺมณา วา อมฺหากํ คามกฺเขตฺตํ อาคจฺฉนฺติ, อติถี โน เต โหนฺติ, อติถี โข ปนเมฺหหิ สกฺกาตพฺพา ครุกาตพฺพา มาเนตพฺพา ปูเชตพฺพา.\csromanspacer{} ยมฺปิ สมโณ โคตโม โอปาสาทํ}
\prosecont{\csromanfolio{380}อนุปฺปตฺโต โอปาสาเท วิหรติ อุตฺตเรน โอปาสาทํ เทววเน สาลวเน, อติถิมฺหากํ สมโณ โคตโม, อติถิ โข ปนเมฺหหิ สกฺกาตพฺโพ ครุกาตพฺโพ มาเนตพฺโพ ปูเชตพฺโพ, อิมินาปงฺเคน น อรหติ โส ภวํ โคตโม อมฺหากํ ทสฺสนาย อุปสงฺกมิตุํ, อถ โข มยเมว อรหาม ตํ ภวนฺตํ โคตมํ ทสฺสนาย อุปสงฺกมิตุํ.\csromanspacer{} เอตฺตเก โข อหํ โภ ตสฺส โภโต โคตมสฺส วณฺเณ ปริยาปุณามิ, โน จ โข โส ภวํ โคตโม เอตฺตกวณฺโณ, อปริมาณวณฺโณ หิ โส ภวํ โคตโม, เอกเมเกนปิ เตน\footnote{เอกเมเกนปิ โภ (สี, สฺยา, กํ, อิ)} องฺเคน สมนฺนาคโต น อรหติ โส ภวํ โคตโม อมฺหากํ ทสฺสนาย อุปสงฺกมิตุํ, อถ โข มยเมว อรหาม ตํ ภวนฺตํ โคตมํ ทสฺสนาย อุปสงฺกมิตุนฺติ, เตน หิ โภ สพฺเพว มยํ สมณํ โคตมํ ทสฺสนาย อุปสงฺกมิสฺสามาติ.}

\proseitem{426}{อถ โข จงฺกี พฺราหฺมโณ มหตา พฺราหฺมณคเณน สทฺธิํ เยน ภควา เตนุปสงฺกมิ, อุปสงฺกมิตฺวา ภควตา สทฺธิํ สมฺโมทิ, สมฺโมทนียํ กถํ สารณียํ วีติสาเรตฺวา เอกมนฺตํ นิสีทิ.\csromanspacer{} เตน โข ปน สมเยน ภควา วุทฺเธหิ วุทฺเธหิ พฺราหฺมเณหิ สทฺธิํ กิญฺจิ กิญฺจิ กถํ สารณียํ วีติสาเรตฺวา นิสินฺโน โหติ.\csromanspacer{} เตน โข ปน สมเยน กาปฏิโก\footnote{กาปฐิโก (สี, อิ), กาปทิโก (สฺยา, กํ)} นาม มาณโว ทหโร วุตฺตสิโร โสฬสวสฺสุทฺเทสิโก ชาติยา, ติณฺณํ เวทานํ ปารคู สนิฆณฺฑุเกฏุภานํ สากฺขรปฺปเภทานํ อิติหาสปญฺจมานํ ปทโก เวยฺยากรโณ โลกายตมหาปุริสลกฺขเณสุ อนวโย, ตสฺสํ ปริสายํ นิสินฺโน โหติ, โส วุทฺธานํ วุทฺธานํ พฺราหฺมณานํ ภควตา สทฺธิํ มนฺตยมานานํ อนฺตรนฺตรา กถํ โอปาเตติ.\csromanspacer{} อถ โข ภควา กาปฏิกํ มาณวํ อปสาเทติ “มายสฺมา ภารทฺวาโช วุทฺธานํ วุทฺธานํ พฺราหฺมณานํ มนฺตยมานานํ อนฺตรนฺตรา กถํ โอปาเตตุ, กถาปริโยสานํ อายสฺมา ภารทฺวาโช อาคเมตู”ติ.\csromanspacer{} เอวํ วุตฺเต จงฺกี พฺราหฺมโณ ภควนฺตํ เอตทโวจ “มา ภวํ โคตโม กาปฏิกํ มาณวํ อปสาเทสิ, กุลปุตฺโต จ กาปฏิโก มาณโว, พหุสฺสุโต จ กาปฏิโก มาณโว, ปณฺฑิโต จ กาปฏิโก มาณโว, กลฺยาณวากฺกรโณ จ กาปฏิโก มาณโว, ปโหติ จ กาปฏิโก มาณโว โภตา โคตเมน สทฺธิํ อสฺมิํ วจเน ปฏิมนฺเตตุนฺ”ติ.\csromanspacer{} อถ โข}

\vspace{\tipitakaparskip}
\csromancenter{จงฺกีสุตฺต (95) ภควโต เอตทโหสิ “อทฺธา โข กาปฏิกสฺส\footnote{เอตทโหสิ “กาปฏิกสฺส (ก)} มาณวสฺส เตวิชฺชเก ปาวจเน กถา\footnote{กถํ (สี, ก), กตํ (สฺยา, กํ, อิ)} ภวิสฺสติ, ตถา หิ นํ พฺราหฺมณา สํปุเรกฺขโรนฺตี”ติ.\csromanspacer{} อถ โข กาปฏิกสฺส มาณวสฺส เอตทโหสิ “ยทา เม สมโณ โคตโม จกฺขุํ อุปสํหริสฺสติ, อถาหํ สมณํ โคตมํ ปญฺหํ ปุจฺฉิสฺสามี”ติ.\csromanspacer{} อถ โข ภควา กาปฏิกสฺส มาณวสฺส เจตสา เจโตปริวิตกฺกมญฺญาย เยน กาปฏิโก มาณโว เตน จกฺขูนิ อุปสํหาสิ.}
\proseitem{427}{\csromanfolio{381}อถ โข กาปฏิกสฺส มาณวสฺส เอตทโหสิ “สมนฺนาหรติ โข มํ สมโณ โคตโม, ยํนูนาหํ สมณํ โคตมํ ปญฺหํ ปุจฺเฉยฺยนฺ”ติ.\csromanspacer{} อถ โข กาปฏิโก มาณโว ภควนฺตํ เอตทโวจ “ยทิทํ โภ โคตม พฺราหฺมณานํ โปราณํ มนฺตปทํ อิติหิติหปรมฺปราย ปิฏกสมฺปทาย, ตตฺถ จ พฺราหฺมณา เอกํเสน นิฏฺฐํ คจฺฉนฺติ ‘อิทเมว สจฺจํ โมฆมญฺญนฺ’ติ, อิธ ภวํ โคตโม กิมาหา”ติ.\csromanspacer{} กิํ ปน ภารทฺวาช อตฺถิ โกจิ พฺราหฺมณานํ เอกพฺราหฺมโณปิ, โย เอวมาห “อหเมตํ ชานามิ, อหเมตํ ปสฺสามิ, อิทเมว สจฺจํ โมฆมญฺญนฺ”ติ.\csromanspacer{} โน หิทํ โภ โคตม.\csromanspacer{} กิํ ปน ภารทฺวาช อตฺถิ โกจิ พฺราหฺมณานํ เอกาจริโยปิ เอกาจริยปาจริโยปิ ยาว สตฺตมา อาจริยมหยุคาปิ, โย เอวมาห “อหเมตํ ชานามิ, อหเมตํ ปสฺสามิ, อิทเมว สจฺจํ โมฆมญฺญนฺ”ติ.\csromanspacer{} โน หิทํ โภ โคตม.\csromanspacer{} กิํ ปน ภารทฺวาช เยปิ เต พฺราหฺมณานํ ปุพฺพกา อิสโย มนฺตานํ กตฺตาโร มนฺตานํ ปวตฺตาโร, เยสมิทํ เอตรหิ พฺราหฺมณา โปราณํ มนฺตปทํ คีตํ ปวุตฺตํ สมิหิตํ, ตทนุคายนฺติ ตทนุภาสนฺติ ภาสิตมนุภาสนฺติ วาจิตมนุวาเจนฺติ, เสยฺยถิทํ, อฏฺฐโก วามโก วามเทโว เวสฺสามิตฺโต ยมตคฺคิ องฺคีรโส ภารทฺวาโช วาเสฏฺโฐ กสฺสโป ภคุ, เตปิ เอวมาหํสุ “มยเมตํ ชานาม, มยเมตํ ปสฺสาม, อิทเมว สจฺจํ โมฆมญฺญนฺ”ติ.\csromanspacer{} โน หิทํ โภ โคตม.}

\prose{อิติ กิร ภารทฺวาช นตฺถิ โกจิ พฺราหฺมณานํ เอกพฺราหฺมโณปิ, โย เอวมาห “อหเมตํ ชานามิ, อหเมตํ ปสฺสามิ, อิทเมว สจฺจํ โมฆมญฺญนฺ”ติ.\csromanspacer{} นตฺถิ โกจิ พฺราหฺมณานํ เอกาจริโยปิ เอกาจริยปาจริโยปิ ยาว สตฺตมา อาจริยมหยุคาปิ, โย เอวมาห \csromanfolio{382}“อหเมตํ ชานามิ, อหเมตํ ปสฺสามิ, อิทเมว สจฺจํ โมฆมญฺญนฺ”ติ.\csromanspacer{} เยปิ เต พฺราหฺมณานํ ปุพฺพกา อิสโย มนฺตานํ กตฺตาโร มนฺตานํ ปวตฺตาโร, เยสมิทํ เอตรหิ พฺราหฺมณา โปราณํ มนฺตปทํ คีตํ ปวุตฺตํ สมิหิตํ, ตทนุคายนฺติ ตทนุภาสนฺติ ภาสิตมนุภาสนฺติ วาจิตมนุวาเจนฺติ, เสยฺยถิทํ, อฏฺฐโก วามโก วามเทโว เวสฺสามิตฺโต ยมตคฺคิ องฺคีรโส ภารทฺวาโช วาเสฏฺโฐ กสฺสโป ภคุ, เตปิ น เอวมาหํสุ “มยเมตํ ชานาม, มยเมตํ ปสฺสาม, อิทเมว สจฺจํ โมฆมญฺญนฺ”ติ.}

\proseitem{428}{เสยฺยถาปิ ภารทฺวาช อนฺธเวณิ ปรมฺปราสํสตฺตา ปุริโมปิ น ปสฺสติ, มชฺฌิโมปิ น ปสฺสติ, ปจฺฉิโมปิ น ปสฺสติ, เอวเมว โข ภารทฺวาช อนฺธเวณูปมํ มญฺเญ พฺราหฺมณานํ ภาสิตํ สมฺปชฺชติ.\csromanspacer{} ปุริโมปิ น ปสฺสติ, มชฺฌิโมปิ น ปสฺสติ, ปจฺฉิโมปิ น ปสฺสติ.\csromanspacer{} ตํ กิํ มญฺญสิ ภารทฺวาช, นนุ เอวํ สนฺเต พฺราหฺมณานํ อมูลิกา สทฺธา สมฺปชฺชตีติ.\csromanspacer{} น เขฺวตฺถ โภ โคตม พฺราหฺมณา สทฺธาเยว ปยิรุปาสนฺติ.\csromanspacer{} อนุสฺสวาเปตฺถ พฺราหฺมณา ปยิรุปาสนฺตีติ.\csromanspacer{} ปุพฺเพว โข ตฺวํ ภารทฺวาช สทฺธํ อคมาสิ.\csromanspacer{} อนุสฺสวํ อิทานิ วเทสิ.\csromanspacer{} ปญฺจ โข อิเม ภารทฺวาช ธมฺมา ทิฏฺเฐว ธมฺเม เทฺวธา วิปากา.\csromanspacer{} กตเม ปญฺจ.\csromanspacer{} สทฺธา, รุจิ, อนุสฺสโว, อาการปริวิตกฺโก, ทิฏฺฐินิชฺฌานกฺขนฺติ.\csromanspacer{} อิเม โข ภารทฺวาช ปญฺจ ธมฺมา ทิฏฺเฐว ธมฺเม เทฺวธา วิปากา.\csromanspacer{} อปิ จ ภารทฺวาช สุสทฺทหิตํเยว โหติ, ตญฺจ โหติ ริตฺตํ ตุจฺฉํ มุสา, โน เจปิ สุสทฺทหิตํ โหติ.\csromanspacer{} ตญฺจ โหติ ภูตํ ตจฺฉํ อนญฺญถา, อปิ จ ภารทฺวาช สุรุจิตํเยว โหติ ฯเปฯ สฺวานุสฺสุตํเยว โหติ ฯเปฯ สุปริวิตกฺกิตํเยว โหติ ฯเปฯ สุนิชฺฌายิตํเยว โหติ, ตญฺจ โหติ ริตฺตํ ตุจฺฉํ มุสา, โน เจปิ สุนิชฺฌายิตํ โหติ, ตญฺจ โหติ, ภูตํ ตจฺฉํ อนญฺญถา.\csromanspacer{} สจฺจมนุรกฺขตา ภารทฺวาช วิญฺญุนา ปุริเสน นาลเมตฺถ เอกํเสน นิฏฺฐํ คนฺตุํ “อิทเมว สจฺจํ โมฆมญฺญนฺ”ติ.}

\proseitem{429}{กิตฺตาวตา ปน โภ โคตม สจฺจานุรกฺขณา โหติ, กิตฺตาวตา สจฺจมนุรกฺขติ, สจฺจานุรกฺขณํ มยํ ภวนฺตํ โคตมํ ปุจฺฉามาติ.\csromanspacer{} สทฺธา เจปิ ภารทฺวาช ปุริสสฺส โหติ, เอวํ เม สทฺธาติ อิติ วทํ สจฺจมนุรกฺขติ\footnote{เอวเมว สิชฺฌตีติ อิติ วา, ตํ สจฺจมนุรกฺขติ (ก)}, นเตฺวว ตาว เอกํเสน นิฏฺฐํ คจฺฉติ “อิทเมว สจฺจํ}

\vspace{\tipitakaparskip}
\csromancenter{จงฺกีสุตฺต (95) โมฆมญฺญนฺ”ติ. ( )\footnote{(เอตฺตาวตา โข ภารทฺวาช สจฺจานุรกฺขณา โหติ, เอตฺตาวตา สจฺจมนุรกฺขติ, เอตฺตาวตา จ มยํ สจฺจานุรกฺขณํ ปญฺญาเปม, น เตฺวว ตาว สจฺจานุโพโธ โหติ) (สี, สฺยา, กํ, อิ) 2. อิธ กิร (สฺยา, กํ, ก)} รุจิเจปิ ภารทฺวาช ปุริสสฺส โหติ ฯเปฯ.\csromanspacer{} อนุสฺสโว เจปิ ภารทฺวาช ปุริสสฺส โหติ ฯเปฯ.\csromanspacer{} อาการปริวิตกฺโก เจปิ ภารทฺวาช ปุริสสฺส โหติ ฯเปฯ.\csromanspacer{} ทิฏฺฐินิชฺฌานกฺขนฺติ เจปิ ภารทฺวาช ปุริสสฺส โหติ, เอวํ เม ทิฏฺฐินิชฺฌานกฺขนฺตีติ อิติ วทํ สจฺจมนุรกฺขติ, นเตฺวว ตาว เอกํเสน นิฏฺฐํ คจฺฉติ “อิทเมว สจฺจํ โมฆมญฺญนฺ”ติ.\csromanspacer{} เอตฺตาวตา โข ภารทฺวาช สจฺจานุรกฺขณา โหติ, เอตฺตาวตา สจฺจมนุรกฺขติ, เอตฺตาวตา จ มยํ สจฺจานุรกฺขณํ ปญฺญเปม, น เตฺวว ตาว สจฺจานุโพโธ โหตีติ.}
\csromanmatikaanchor{mtk.54}
\csromantocmarkref{subsection}{5. จงฺกีสุตฺต}{mtk.54}
\bookmark[dest={mtk.54},level=2]{5. จงฺกีสุตฺต}
\proseitem{430}{\csromanfolio{383}เอตฺตาวตา โภ โคตม สจฺจานุรกฺขณา โหติ, เอตฺตาวตา สจฺจมนุรกฺขติ, เอตฺตาวตา จ มยํ สจฺจานุรกฺขณํ เปกฺขาม.\csromanspacer{} กิตฺตาวตา ปน โภ โคตม สจฺจานุโพโธ โหติ, กิตฺตาวตา สจฺจมนุพุชฺฌติ, สจฺจานุโพธํ มยํ ภวนฺตํ โคตมํ ปุจฺฉามาติ.\csromanspacer{} อิธ2 ภารทฺวาช ภิกฺขุ อญฺญตรํ คามํ วา นิคมํ วา อุปนิสฺสาย วิหรติ, ตเมนํ คหปติ วา คหปติปุตฺโต วา อุปสงฺกมิตฺวา ตีสุ ธมฺเมสุ สมนฺเนสติ โลภนีเยสุ ธมฺเมสุ โทสนีเยสุ ธมฺเมสุ โมหนีเยสุ ธมฺเมสุ, “อตฺถิ นุ โข อิมสฺสายสฺมโต ตถารูปา โลภนียา ธมฺมา, ยถารูเปหิ โลภนีเยหิ ธมฺเมหิ ปริยาทินฺนจิตฺโต อชานํ วา วเทยฺย ‘ชานามี’ติ, อปสฺสํ วา วเทยฺย ‘ปสฺสามี’ติ ปรํ วา ตทตฺถาย สมาทเปยฺย, ยํ ปเรสํ อสฺส ทีฆรตฺตํ อหิตาย ทุกฺขายา”ติ.\csromanspacer{} ตเมนํ สมนฺเนสมาโน เอวํ ชานาติ “นตฺถิ โข อิมสฺสายสฺมโต ตถารูปา โลภนียา ธมฺมา, ยถารูเปหิ โลภนีเยหิ ธมฺเมหิ ปริยาทินฺนจิตฺโต อชานํ วา วเทยฺย’ชนามี’ติ, อปสฺสํ วา วเทยฺย ‘ปสฺสามี’ติ, ปรํ วา ตทตฺถาย สมาทเปยฺย, ยํ ปเรสํ อสฺส ทีฆรตฺตํ อหิตาย ทุกฺขาย\footnote{ทุกฺขายาติ (สพฺพตฺถ)}.\csromanspacer{} ตถารูโป\footnote{ตถา (สี, สฺยา, กํ, อิ)} โข ปนิมสฺสายสฺมโต กายสมาจาโร ตถารูโป4 วจีสมาจาโร, ยถา ตํ อลุทฺธสฺส.\csromanspacer{} ยํ โข ปน อยมายสฺมา ธมฺมํ เทเสติ, คมฺภีโร โส ธมฺโม ทุทฺทโส \csromanfolio{384}ทุรนุโพโธ สนฺโต ปณีโต อตกฺกาวจโร นิปุโณ ปณฺฑิตเวทนีโย, น โส ธมฺโม สุเทสิโย ลุทฺเธนา”ติ.}

\proseitem{431}{ยโต นํ สมนฺเนสมาโน วิสุทฺธํ โลภนีเยหิ ธมฺเมหิ สมนุปสฺสติ, ตโต นํ อุตฺตริ สมนฺเนสติ โทสนีเยสุ ธมฺเมสุ “อตฺถิ นุ โข อิมสฺสายสฺมโต ตถารูปา โทสนียา ธมฺมา, ยถารูเปหิ โทสนีเยหิ ธมฺเมหิ ปริยาทินฺนจิตฺโต อชานํ วา วเทยฺย ‘ชานามี’ติ, อปสฺสํ วา วเทยฺย ‘ปสฺสามี’ติ, ปรํ วา ตทตฺถาย สมาทเปยฺย, ยํ ปเรสํ อสฺส ทีฆรตฺตํ อหิตาย ทุกฺขายา”ติ.\csromanspacer{} ตเมนํ สมนฺเนสมาโน เอวํ ชานาติ “นตฺถิ โข อิมสฺสายสฺมโต ตถารูปา โทสนียา ธมฺมา, ยถารูเปหิ โทสนีเยหิ ธมฺเมหิ ปริยาทินฺนจิตฺโต อชานํ วา วเทยฺย ‘ชานามี’ติ, อปสฺสํ วา วเทยฺย ‘ปสฺสามี’ติ, ปรํ วา ตทตฺถาย สมาทเปยฺย, ยํ ปเรสํ อสฺส ทีฆรตฺตํ อหิตาย ทุกฺขาย.\csromanspacer{} ตถารูโป โข ปนิมสฺสายสฺมโต กายสมาจาโร ตถารูโป วจีสมาจาโร, ยถา ตํ อทุฏฺฐสฺส.\csromanspacer{} ยํ โข ปน อยมายสฺมา ธมฺมํ เทเสติ, คมฺภีโร โส ธมฺโม ทุทฺทโส ทุรนุโพโธ สนฺโต ปณีโต อตกฺกาวจโร นิปุโณ ปณฺฑิตเวทนีโย, น โส ธมฺโม สุเทสิโย ทุฏฺเฐนา”ติ.}

\proseitem{432}{ยโต นํ สมนฺเนสมาโน วิสุทฺธํ โทสนีเยหิ ธมฺเมหิ สมนุปสฺสติ, ตโต นํ อุตฺตริ สมนฺเนสติ โมหนีเยสุ ธมฺเมสุ “อตฺถิ นุ โข อิมสฺสายสฺมโต ตถารูปา โมหนียา ธมฺมา, ยถารูเปหิ โมหนีเยหิ ธมฺเมหิ ปริยาทินฺนจิตฺโต อชานํ วา วเทยฺย ‘ชานามี’ติ, อปสฺสํ วา วเทยฺย ‘ปสฺสามี’ติ, ปรํ วา ตทตฺถาย สมาทเปยฺย, ยํ ปเรสํ อสฺส ทีฆรตฺตํ อหิตาย ทุกฺขายา”ติ, ตเมนํ สมนฺเนสมาโน เอวํ ชานาติ “นตฺถิ โข อิมสฺสายสฺมโต ตถารูปา โมหนียา ธมฺมา, ยถารูเปหิ โมหนีเยหิ ธมฺเมหิ ปริยาทินฺนจิตฺโต อชานํ วา วเทยฺย ‘ชานามี’ติ, อปสฺสํ วา วเทยฺย ‘ปสฺสามี’ติ, ปรํ วา ตทตฺถาย สมาทเปยฺย, ยํ ปเรสํ อสฺส ทีฆรตฺตํ อหิตาย ทุกฺขาย.\csromanspacer{} ตถารูโป โข ปนิมสฺสายสฺมโต กายสมาจาโร ตถารูโป วจีสมาจาโร, ยถา ตํ อมูฬฺหสฺส.\csromanspacer{} ยํ โข ปน อยมายสฺมา ธมฺมํ เทเสติ คมฺภีโร โส ธมฺโม ทุทฺทโส ทุรนุโพโธ สนฺโต}

\vspace{\tipitakaparskip}
\csromancenter{จงฺกีสุตฺต (95) ปณีโต อตกฺกาวจโร นิปุโณ ปณฺฑิตเวทนีโย, น โส ธมฺโม สุเทสิโย มูเฬฺหนา”ติ.}
\prose{\csromanfolio{385}ยโต นํ สมนฺเนสมาโน วิสุทฺธํ โมหนีเยหิ ธมฺเมหิ สมนุปสฺสติ, อถ ตมฺหิ สทฺธํ นิเวเสติ, สทฺธาชาโต อุปสงฺกมติ, อุปสงฺกมนฺโต ปยิรุปาสติ, ปยิรุปาสนฺโต โสตํ โอทหติ, โอหิตโสโต ธมฺมํ สุณาติ, สุตฺวา ธมฺมํ ธาเรติ, ธตานํ\footnote{ธาริตานํ (ก)} ธมฺมานํ อตฺถํ อุปปริกฺขติ, อตฺถํ อุปปริกฺขโต ธมฺมา นิชฺฌานํ ขมนฺติ, ธมฺมนิชฺฌานกฺขนฺติยา สติ ฉนฺโท ชายติ, ฉนฺทชาโต อุสฺสหติ, อุสฺสหิตฺวา ตุเลติ, ตุลยิตฺวา ปทหติ, ปหิตตฺโต สมาโน กาเยน เจว ปรมสจฺจํ สจฺฉิกโรติ, ปญฺญาย จ นํ อติวิชฺฌ ปสฺสติ.\csromanspacer{} เอตฺตาวตา โข ภารทฺวาช สจฺจานุโพโธ โหติ.\csromanspacer{} เอตฺตาวตา สจฺจมนุพุชฺฌติ.\csromanspacer{} เอตฺตาวตา จ มยํ สจฺจานุโพธํ ปญฺญเปม, น เตฺวว ตาว สจฺจานุปฺปตฺติ โหตีติ.}

\proseitem{433}{เอตฺตาวตา โภ โคตม สจฺจานุโพโธ โหติ, เอตฺตาวตา สจฺจมนุพุชฺฌติ.\csromanspacer{} เอตฺตาวตา จ มยํ สจฺจานุโพธํ เปกฺขาม.\csromanspacer{} กิตฺตาวตา ปน โภ โคตม สจฺจานุปฺปตฺติ โหติ, กิตฺตาวตา สจฺจมนุปาปุณาติ, สจฺจานุปฺปตฺติํ มยํ ภวนฺตํ โคตมํ ปุจฺฉามาติ.\csromanspacer{} เตสํเยว ภารทฺวาช ธมฺมานํ อาเสวนา ภาวนา พหุลีกมฺมํ สจฺจานุปฺปตฺติ โหติ.\csromanspacer{} เอตฺตาวตา โข ภารทฺวาช สจฺจานุปฺปตฺติ โหติ, เอตฺตาวตา สจฺจมนุปาปุณาติ, เอตฺตาวตา จ มยํ สจฺจานุปฺปตฺติํ ปญฺญเปมาติ.}

\proseitem{434}{เอตฺตาวตา โภ โคตม สจฺจานุปฺปตฺติ โหติ, เอตฺตาวตา สจฺจมนุปาปุณาติ, เอตฺตาวตา จ มยํ สจฺจานุปฺปตฺติํ เปกฺขาม.\csromanspacer{} สจฺจานุปฺปตฺติยา ปน โภ โคตม กตโม ธมฺโม พหุกาโร, สจฺจานุปฺปตฺติยา พหุการํ ธมฺมํ มยํ ภวนฺตํ โคตมํ ปุจฺฉามาติ.\csromanspacer{} สจฺจานุปฺปตฺติยา โข ภารทฺวาช ปธานํ พหุการํ, โน เจตํ ปทเหยฺย, นยิทํ สจฺจมนุปาปุเณยฺย, ยสฺมา จ โข ปทหติ, ตสฺมา สจฺจมนุปาปุณาติ, ตสฺมา สจฺจานุปฺปตฺติยา ปธานํ พหุการนฺติ.}

\prose{ปธานสฺส ปน โภ โคตม กตโม ธมฺโม พหุกาโร, ปธานสฺส พหุการํ ธมฺมํ มยํ ภวนฺตํ โคตมํ ปุจฺฉามาติ.\csromanspacer{} ปธานสฺส โข ภารทฺวาช \csromanfolio{386}ตุลนา พหุการา, โน เจตํ ตุเลยฺย, นยิทํ ปทเหยฺย, ยสฺมา จ โข ตุเลติ, ตสฺมา ปทหติ, ตสฺมา ปธานสฺส ตุลนา พหุการาติ.}

\prose{ตุลนาย ปน โภ โคตม กตโม ธมฺโม พหุกาโร, ตุลนาย พหุการํ ธมฺมํ มยํ ภวนฺตํ โคตมํ ปุจฺฉามาติ.\csromanspacer{} ตุลนาย โข ภารทฺวาช อุสฺสาโห พหุกาโร, โน เจตํ อุสฺสเหยฺย, นยิทํ ตุเลยฺย, ยสฺมา จ โข อุสฺสหติ, ตสฺมา ตุเลติ, ตสฺมา ตุลนาย อุสฺสาโห พหุกาโรติ.}

\prose{อุสฺสาหสฺส ปน โภ โคตม กตโม ธมฺโม พหุกาโร, อุสฺสาหสฺส พหุการํ ธมฺมํ มยํ ภวนฺตํ โคตมํ ปุจฺฉามาติ.\csromanspacer{} อุสฺสาหสฺส โข ภารทฺวาช ฉนฺโท พหุกาโร, โน เจตํ ฉนฺโท ชาเยถ, นยิทํ อุสฺสเหยฺย, ยสฺมา จ โข ฉนฺโท ชายติ, ตสฺมา อุสฺสหติ, ตสฺมา อุสฺสาหสฺส ฉนฺโท พหุกาโรติ.}

\prose{ฉนฺทสฺส ปน โภ โคตม กตโม ธมฺโม พหุกาโร, ฉนฺทสฺส พหุการํ ธมฺมํ มยํ ภวนฺตํ โคตมํ ปุจฺฉามาติ.\csromanspacer{} ฉนฺทสฺส โข ภารทฺวาช ธมฺมนิชฺฌานกฺขนฺติ พหุการา, โน เจเต ธมฺมา นิชฺฌานํ ขเมยฺยุํ, นยิทํ ฉนฺโท ชาเยถ, ยสฺมา จ โข ธมฺมา นิชฺฌานํ ขมนฺติ, ตสฺมา ฉนฺโท ชายติ, ตสฺมา ฉนฺทสฺส ธมฺมนิชฺฌานกฺขนฺติ พหุการาติ.}

\prose{ธมฺมนิชฺฌานกฺขนฺติยา ปน โภ โคตม กตโม ธมฺโม พหุกาโร, ธมฺมนิชฺฌานกฺขนฺติยา พหุการํ ธมฺมํ มยํ ภวนฺตํ โคตมํ ปุจฺฉามาติ.\csromanspacer{} ธมฺมนิชฺฌานกฺขนฺติยา โข ภารทฺวาช อตฺถูปปริกฺขา พหุการา โน เจตํ อตฺถํ อุปปริกฺเขยฺย, นยิทํ ธมฺมา นิชฺฌานํ ขเมยฺยุํ, ยสฺมา จ โข อตฺถํ อุปปริกฺขติ, ตสฺมา ธมฺมา นิชฺฌานํ ขมนฺติ, ตสฺมา ธมฺมนิชฺฌานกฺขนฺติยา อตฺถูปปริกฺขา พหุการาติ.}

\prose{อตฺถูปปริกฺขาย ปน โภ โคตม กตโม ธมฺโม พหุกาโร, อตฺถูปปริกฺขาย พหุการํ ธมฺมํ มยํ ภวนฺตํ โคตมํ ปุจฺฉามาติ.\csromanspacer{} อตฺถูปปริกฺขาย โข ภารทฺวาช ธมฺมธารณา พหุการา, โน เจตํ ธมฺมํ ธาเรยฺย, นยิทํ อตฺถํ อุปปริกฺเขยฺย, ยสฺมา จ โข ธมฺมํ ธาเรติ, ตสฺมา อตฺถํ อุปปริกฺขติ, ตสฺมา อตฺถูปปริกฺขาย ธมฺมธารณา พหุการาติ.}

\prose{ธมฺมธารณาย ปน โภ โคตม กตโม ธมฺโม พหุกาโร, ธมฺมธารณาย พหุการํ ธมฺมํ มยํ ภวนฺตํ โคตมํ ปุจฺฉามาติ.}

\vspace{\tipitakaparskip}
\csromancenter{จงฺกีสุตฺต (95) ธมฺมธารณาย โข ภารทฺวาช ธมฺมสฺสวนํ พหุการํ, โน เจตํ ธมฺมํ สุเณยฺย นยิทํ ธมฺมํ ธาเรยฺย, ยสฺมา จ โข ธมฺมํ สุณาติ, ตสฺมา ธมฺมํ ธาเรติ, ตสฺมา ธมฺมธารณาย ธมฺมสฺสวนํ พหุการนฺติ.}
\prose{\csromanfolio{387}ธมฺมสฺสวนสฺส ปน โภ โคตม กตโม ธมฺโม พหุกาโร, ธมฺมสฺสวนสฺส พหุการํ ธมฺมํ มยํ ภวนฺตํ โคตมํ ปุจฺฉามาติ.\csromanspacer{} ธมฺมสฺสวนสฺส โข ภารทฺวาช โสตาวธานํ พหุการํ, โน เจตํ สิตํ โอทเหยฺย, นยิทํ ธมฺมํ สุเณยฺย, ยสฺมา จ โข โสตํ โอทหติ, ตสฺมา ธมฺมา สุณาติ, ตสฺมา ธมฺมสฺสวนสฺส โสตาวธานํ พหุการนฺติ.}

\prose{โสตาวธานสฺส ปน โภ โคตม กตโม ธมฺโม พหุกาโร, โสตาวธานสฺส พหุการํ ธมฺมํ มยํ ภวนฺตํ โคตมํ ปุจฺฉามาติ.\csromanspacer{} โสตาวธานสฺส โข ภารทฺวาช ปยิรุปาสนา พหุการา, โน เจตํ ปยิรุปาเสยฺย, นยิทํ โสตํ โอทเหยฺย, ยสฺมา จ โข ปยิรุปาสติ, ตสฺมา โสตํ โอทหติ, ตสฺมา โสตาวธานสฺส ปยิรุปาสนา พหุการาติ.}

\prose{ปยิรุปาสนาย ปน โภ โคตม กตโม ธมฺโม พหุกาโร, ปยิรุปาสนาย พหุการํ ธมฺมํ มยํ ภวนฺตํ โคตมํ ปุจฺฉามาติ.\csromanspacer{} ปยิรุปาสนาย โข ภารทฺวาช อุปสงฺกมนํ พหุการํ, โน เจตํ อุปสงฺกเมยฺย, นยิทํ ปยิรุปาเสยฺย, ยสฺมา จ โข อุปสงฺกมติ, ตสฺมา ปยิรุปาสติ, ตสฺมา ปยิรุปาสนาย อุปสงฺกมนํ พหุการนฺติ.}

\prose{อุปสงฺกมนสฺส ปน โภ โคตม กตโม ธมฺโม พหุกาโร, อุปสงฺกมนสฺส พหุการํ ธมฺมํ มยํ ภวนฺตํ โคตมํ ปุจฺฉามาติ.\csromanspacer{} อุปสงฺกมนสฺส โข ภารทฺวาช สทฺธา พหุการา, โน เจตํ สทฺธา ชาเยถ, นยิทํ อุปสงฺกเมยฺย, ยสฺมา จ โข สทฺธา ชายติ, ตสฺมา อุปสงฺกมติ, ตสฺมา อุปสงฺกมนสฺส สทฺธา พหุการาติ.}

\proseitem{435}{สจฺจานุรกฺขณํ มยํ ภวนฺตํ โคตมํ อปุจฺฉิมฺห, สจฺจานุรกฺขณํ ภวํ โคตโม พฺยากาสิ, ตญฺจ ปนมฺหากํ รุจฺจติ เจว ขมติ จ, เตน จมฺห อตฺตมนา.\csromanspacer{} สจฺจานุโพธํ มยํ ภวนฺตํ โคตมํ อปุจฺฉิมฺห, สจฺจานุโพธํ ภวํ โคตโม พฺยากาสิ, ตญฺจ ปนมฺหากํ รุจฺจติ เจว ขมติ จ, เตน จมฺห อตฺตมนา.\csromanspacer{} สจฺจานุปฺปตฺติํ มยํ ภวนฺตํ โคตมํ อปุจฺฉิมฺห, สจฺจานุปฺปตฺติํ ภวํ โคตโม พฺยากาสิ, ตญฺจ ปนมฺหากํ รุจฺจติ เจว ขมติ จ, เตน จมฺห \csromanfolio{388}อตฺตมนา.\csromanspacer{} สจฺจานุปฺปตฺติยา พหุการํ ธมฺมํ มยํ ภวนฺตํ โคตมํ อปุจฺฉิมฺห, สจฺจานุปฺปตฺติยา พหุการํ ธมฺมํ ภวํ โคตโม พฺยากาสิ, ตญฺจ ปนมฺหากํ รุจฺจติ เจว ขมติ จ, เตน จมฺห อตฺตมนา.\csromanspacer{} ยํยเทว จ มยํ ภวนฺตํ โคตมํ อปุจฺฉิมฺห, ตํตเทว ภวํ โคตโม พฺยากาสิ, ตญฺจ ปนมฺหากํ รุจฺจติ เจว ขมติ จ, เตน จมฺห อตฺตมนา.\csromanspacer{} มยญฺหิ โภ โคตม ปุพฺเพ เอวํ ชานาม, เก จ มุณฺฑกา สมณกา อิพฺภา กณฺหา พนฺธุปาทาปจฺจา, เก จ ธมฺมสฺส อญฺญาตาโรติ.\csromanspacer{} อชเนสิ วต เม ภวํ โคตโม สมเณสุ สมณเปมํ, สมเณสุ สมณปสาทํ, สมเณสุ สมณคารวํ.\csromanspacer{} อภิกฺกนฺตํ โภ โคตม ฯเปฯ อุปาสกํ มํ ภวํ โคตโม ธาเรตุ อชฺชตคฺเค ปาณุเปตํ สรณํ คตนฺติ.}

\nitthitamruled{จงฺกีสุตฺตํ นิฏฺฐิตํ ปญฺจมํ.}
\csromanheader{2}{6. เอสุการีสุตฺต}
\proseitem{436}{เอวํ เม สุตํ\csromandash{}เอกํ สมยํ ภควา สาวตฺถิยํ วิหรติ เชตวเน อนาถปิณฺฑิกสฺส อาราเม.\csromanspacer{} อถ โข เอสุการี พฺราหฺมโณ เยน ภควา เตนุปสงฺกมิ, อุปสงฺกมิตฺวา ภควตา สทฺธิํ สมฺโมทิ, สมฺโมทนียํ กถํ สารณียํ วีติสาเรตฺวา เอกมนฺตํ นิสีทิ, เอกมนฺตํ นิสินฺโน โข เอสุการี พฺราหฺมโณ ภควนฺตํ เอตทโวจ “พฺราหฺมณา โภ โคตม จตสฺโส ปาริจริยา ปญฺญเปนฺติ.\csromanspacer{} พฺราหฺมณสฺส ปาริจริยํ ปญฺญเปนฺติ, ขตฺติยสฺส ปาริจริยํ ปญฺญเปนฺติ, เวสฺสสฺส ปาริจริยํ ปญฺญเปนฺติ, สุทฺทสฺส ปาริจริยํ ปญฺญเปนฺติ.\csromanspacer{} ตตฺริทํ โภ โคตม พฺราหฺมณา พฺราหฺมณสฺส ปาริจริยํ ปญฺญเปนฺติ, พฺราหฺมโณ วา พฺราหฺมณํ ปริจเรยฺย, ขตฺติโย วา พฺราหฺมณํ ปริจเรยฺย, เวสฺโส วา พฺราหฺมณํ ปริจเรยฺย, สุทฺโท วา พฺราหฺมณํ ปริจเรยฺยาติ.\csromanspacer{} อิทํ โข โภ โคตม พฺราหฺมณา พฺราหฺมณสฺส ปาริจริยํ ปญฺญเปนฺติ.\csromanspacer{} ตตฺริทํ โภ โคตม พฺราหฺมณา ขตฺติยสฺส ปาริจริยํ ปญฺญเปนฺติ, ขตฺติโย วา ขตฺติยํ ปริจเรยฺย, เวสฺโส วา ขตฺติยํ ปริจเรยฺย, สุทฺโท วา ขตฺติยํ ปริจเรยฺยาติ.\csromanspacer{} อิทํ โข โภ โคตม พฺราหฺมณา ขตฺติยสฺส ปาริจริยํ ปญฺญเปนฺติ.\csromanspacer{} ตตฺริทํ โภ โคตม พฺราหฺมณา เวสฺสสฺส ปาริจริยํ ปญฺญเปนฺติ, เวสฺโส วา เวสฺสํ ปริจเรยฺย, สุทฺโท วา เวสฺสํ ปริจเรยฺยาติ.\csromanspacer{} อิทํ โข โภ โคตม พฺราหฺมณา เวสฺสสฺส ปาริจริยํ}

\vspace{\tipitakaparskip}
\csromancenter{เอสุการีสุตฺต (96) ปญฺญเปนฺติ.\csromanspacer{} ตตฺริทํ โภ โคตม พฺราหฺมณา สุทฺทสฺส ปาริจริยํ ปญฺญเปนฺติ, สุทฺโทว สุทฺทํ ปริจเรยฺย, โก ปนญฺโญ สุทฺทํ ปริจริสฺสตีติ.\csromanspacer{} อิทํ โข โภ โคตม พฺราหฺมณา สุทฺทสฺส ปาริจริยํ ปญฺญเปนฺติ.\csromanspacer{} พฺราหฺมณา โภ โคตม อิมา จตสฺโส ปาริจริยา ปญฺญเปนฺติ, อิธ ภวํ โคตโม กิมาหา”ติ.}
\proseitem{437}{\csromanfolio{389}กิํ ปน พฺราหฺมณ สพฺโพ โลโก พฺราหฺมณานํ เอตทพฺภนุชานาติ “อิมา จตสฺโส ปาริจริยา ปญฺญเปนฺตู”ติ\footnote{ปญฺญเปนฺตีติ (สี, ก)}.\csromanspacer{} โน หิทํ โภ โคตม.\csromanspacer{} เสยฺยถาปิ พฺราหฺมณ ปุริโส ทลิทฺโท\footnote{ทฬิทฺโท (สี, สฺยา, กํ, อิ)} อสฺสโก อนาฬฺหิโย, ตสฺส อกามสฺส พิลํ โอลคฺเคยฺยุํ “อิทํ เต อมฺโภ ปุริส มํสํ ขาทิตพฺพํ, มูลญฺจ อนุปฺปทาตพฺพนฺ”ติ.\csromanspacer{} เอวเมว โข พฺราหฺมณ พฺราหฺมณา อปฺปฏิญฺญาย เตสํ สมณพฺราหฺมณานํ อถ จ ปนิมา จตสฺโส ปาริจริยา ปญฺญเปนฺติ.\csromanspacer{} นาหํ พฺราหฺมณ สพฺพํ ปริจริตพฺพนฺติ วทามิ, นาหํ พฺราหฺมณ สพฺพํ น ปริจริตพฺพนฺติ วทามิ.\csromanspacer{} ยํ หิสฺส พฺราหฺมณ ปริจรโต ปาริจริยาเหตุ ปาปิโย อสฺส น เสยฺโย, นาหํ ตํ ปริจริตพฺพนฺติ วทามิ.\csromanspacer{} ยญฺจ ขฺวาสฺส พฺราหฺมณ ปริจรโต ปาริจริยาเหตุ เสยฺโย อสฺส น ปาปิโย, ตมหํ ปริจริตพฺพนฺติ วทามิ.\csromanspacer{} ขตฺติยํ เจปิ พฺราหฺมณ เอวํ ปุจฺเฉยฺยุํ “ยํ วา เต ปริจรโต ปาริจริยาเหตุ ปาปิโย อสฺส น เสยฺโย, ยํ วา เต ปริจรโต ปาริจริยาเหตุ เสยฺโย อสฺส น ปาปิโย, กเมตฺถ ปริจเรยฺยาสี”ติ.\csromanspacer{} ขตฺติโยปิ หิ พฺราหฺมณ สมฺมา พฺยากรมาโน เอวํ พฺยากเรยฺย “ยํ หิ เม ปริจรโต ปาริจริยาเหตุ ปาปิโย อสฺส น เสยฺโย, นาหํ ตํ ปริจเรยฺยํ.\csromanspacer{} ยญฺจ โข เม ปริจรโต ปาริจริยาเหตุ เสยฺโย อสฺส น ปาปิโย, ตมหํ ปริจเรยฺยนฺ”ติ.\csromanspacer{} พฺราหฺมณํ เจปิ พฺราหฺมณ ฯเปฯ.\csromanspacer{} เวสฺสํ เจปิ พฺราหฺมณ ฯเปฯ.\csromanspacer{} สุทฺทํ เจปิ พฺราหฺมณ เอวํ ปุจฺเฉยฺยุํ “ยํ วา เต ปริจรโต ปาริจริยาเหตุ ปาปิโย อสฺส น เสยฺโย, ยํ วา เต ปริจรโต ปาริจริยาเหตุ เสยฺโย อสฺส น ปาปิโย, กเมตฺถ ปริจเรยฺยาสี”ติ.\csromanspacer{} สุทฺโทปิ หิ พฺราหฺมณ สมฺมา พฺยากรมาโน เอวํ พฺยากเรยฺย “ยํ หิ เม ปริจรโต ปาริจริยาเหตุ ปาปิโย อสฺส น เสยฺโย, นาหํ ตํ ปริจเรยฺยํ.\csromanspacer{} ยญฺจ โข เม ปริจรโต ปาริจริยาเหตุ เสยฺโย อสฺส น ปาปิโย, ตมหํ ปริจเรยฺยนฺ”ติ.\csromanspacer{} นาหํ พฺราหฺมณ อุจฺจากุลีนตา เสยฺยํโสติ วทามิ, น ปนาหํ พฺราหฺมณ อุจฺจากุลีนตา \csromanfolio{390}ปาปิยํโสติ วทามิ.\csromanspacer{} นาหํ พฺราหฺมณ อุฬารวณฺณตา เสยฺยํโสติ วทามิ, น ปานาหํ พฺราหฺมณา อุฬารวณฺณตา ปาปิยํโสติ วทามิ.\csromanspacer{} นาหํ พฺราหฺมณ อุฬารโภคตา เสยฺยํโสติ วทามิ, น ปนาหํ พฺราหฺมณ อุฬารโภคตา ปาปิยํโสติ วทามิ.}

\proseitem{438}{อุจฺจากุลีโนปิ หิ พฺราหฺมณ อิเธกจฺโจ ปาณาติปาตี โหติ, อทินฺนาทายี โหติ, กาเมสุมิจฺฉาจารี โหติ, มุสาวาที โหติ, ปิสุณาวาโจ โหติ, ผรุสาวาโจ โหติ, สมฺผปฺปลาปี โหติ, อภิชฺฌาลุ โหติ, พฺยาปนฺนจิตฺโต โหติ, มิจฺฉาทิฏฺฐิ โหติ.\csromanspacer{} ตสฺมา น อุจฺจากุลีนตา เสยฺยํโสติ วทามิ.\csromanspacer{} อุจฺจากุลีโนปิ หิ พฺราหฺมณ อิเธกจฺโจ ปาณาติปาตา ปฏิวิรโต โหติ, อทินฺนาทานา ปฏิวิรโต โหติ, กาเมสุมิจฺฉาจารา ปฏิวิรโต โหติ, มุสาวาทา ปฏิวิรโต โหติ, ปิสุณาย วาจาย ปฏิวิรโต โหติ, ผรุสาย วาจาย ปฏิวิรโต โหติ, สมฺผปฺปลาปา ปฏิวิรโต โหติ, อนภิชฺฌาลุ โหติ, อพฺยาปนฺนจิตฺโต โหติ, สมฺมาทิฏฺฐิ โหติ.\csromanspacer{} ตสฺมา น อุจฺจากุลีนตา ปาปิยํโสติ วทามิ.}

\proseitem{439}{อุฬารวณฺโณปิ หิ พฺราหฺมณ ฯเปฯ.\csromanspacer{} อุฬารโภโคปิ หิ พฺราหฺมณ อิเธกจฺโจ ปาณาติปาตี โหติ ฯเปฯ มิจฺฉาทิฏฺฐิ โหติ.\csromanspacer{} ตสฺมา น อุฬารโภคตา เสยฺยํโสติ วทามิ.\csromanspacer{} อุฬารโภโคปิ หิ พฺราหฺมณ อิเธกจฺโจ ปาณาติปาตา ปฏิวิรโต โหติ ฯเปฯ สมฺมาทิฏฺฐิ โหติ.\csromanspacer{} ตสฺมา น อุฬารโภคตา ปาปิยํโสติ วทามิ.\csromanspacer{} นาหํ พฺราหฺมณ สพฺพํ ปริจริตพฺพนฺติ วทามิ, น ปนาหํ พฺราหฺมณ สพฺพํ น ปริจริตพฺพนฺติ วทามิ.\csromanspacer{} ยํ หิสฺส พฺราหฺมณ ปริจรโต ปาริจริยาเหตุ สทฺธา วฑฺฒติ สีลํ วฑฺฒติ สุตํ วฑฺฒติ จาโค วฑฺฒติ ปญฺญา วฑฺฒติ, ตมหํ ปริจริตพฺพนฺติ (วทามิ.\csromanspacer{} ยํ หิสฺส พฺราหฺมณ ปริจรโต ปาริจริยาเหตุ น สทฺธา วฑฺฒติ น สีลํ วฑฺฒติ น สุตํ วฑฺฒติ น จาโค วฑฺฒติ น ปญฺญา วฑฺฒติ, นาหํ ตํ ปริจริตพฺพนฺติ) วทามีติ.}

\proseitem{440}{เอวํ วุตฺเต เอสุการี พฺราหฺมโณ ภควนฺตํ เอตทโวจ “พฺราหฺมณา โภ โคตม จตฺตาริ ธนานิ ปญฺญเปนฺติ.\csromanspacer{} พฺราหฺมณสฺส สนฺธนํ ปญฺญเปนฺติ, ขตฺติยสฺส สนฺธนํ ปญฺญเปนฺติ, เวสฺสสฺส สนฺธนํ ปญฺญเปนฺติ,}

\vspace{\tipitakaparskip}
\csromancenter{เอสุการีสุตฺต (96) สุทฺทสฺส สนฺธนํ ปญฺญเปนฺติ.\csromanspacer{} ตตฺริทํ โภ โคตม พฺราหฺมณา พฺราหฺมณสฺส สนฺธนํ ปญฺญเปนฺติ ภิกฺขาจริยํ.\csromanspacer{} ภิกฺขาจริยญฺจ ปน พฺราหฺมโณ สนฺธนํ อติมญฺญมาโน อกิจฺจการี โหติ, โคโปว อทินฺนํ อาทิยมาโนติ.\csromanspacer{} อิทํ โข โภ โคตม พฺราหฺมณา พฺราหฺมณสฺส สนฺธนํ ปญฺญเปนฺติ.\csromanspacer{} ตตฺริทํ โภ โคตม พฺราหฺมณา ขตฺติยสฺส สนฺธนํ ปญฺญเปนฺติ ธนุกลาปํ.\csromanspacer{} ธนุกลาปญฺจ ปน ขตฺติโย สนฺธนํ อติมญฺญมาโน อกิจฺจการี โหติ, โคโปว อทินฺนํ อาทิยมาโนติ.\csromanspacer{} อิทํ โข โภ โคตม พฺราหฺมณา ขตฺติยสฺส สนฺธนํ ปญฺญเปนฺติ.\csromanspacer{} ตตฺริทํ โภ โคตม พฺราหฺมณา เวสฺสสฺส สนฺธนํ ปญฺญเปนฺติ กสิโครกฺขํ.\csromanspacer{} กสิโครกฺขญฺจ ปน เวสฺโส สนฺธนํ อติมญฺญมาโน อกิจฺจการี โหติ, โคโปว อทินฺนํ อาทิยมาโนติ.\csromanspacer{} อิทํ โข โภ โคตม พฺราหฺมณา เวสฺสสฺส สนฺธนํ ปญฺญเปนฺติ.\csromanspacer{} ตตฺริทํ โภ โคตม พฺราหฺมณา สุทฺทสฺส สนฺธนํ ปญฺญเปนฺติ อสิตพฺยาภงฺคิํ.\csromanspacer{} อสิตพฺยาภงฺคิญฺจ ปน สุทฺโท สนฺธนํ อติมญฺญมาโน อกิจฺจการี โหติ, โคโปว อทินฺนํ อาทิยมาโนติ.\csromanspacer{} อิทํ โข โภ โคตม พฺราหฺมณา สุทฺทสฺส สนฺธนํ ปญฺญเปนฺติ.\csromanspacer{} พฺราหฺมณา โภ โคตม อิมานิ จตฺตาริ ธนานิ ปญฺญเปนฺติ.\csromanspacer{} อิธ ภวํ โคตโม กิมาหา”ติ.}
\proseitem{441}{\csromanfolio{391}กิํ ปน พฺราหฺมณ สพฺโพ โลโก พฺราหฺมณานํ เอตทพฺภนุชานาติ “อิมานิ จตฺตาริ ธนานิ ปญฺญเปนฺตู”ติ.\csromanspacer{} โน หิทํ โภ โคตม.\csromanspacer{} เสยฺยถาปิ พฺราหฺมณ ปุริโส ทลิทฺโท อสฺสโก อนาฬฺหิโย, ตสฺส อกามสฺส พิลํ โอลคฺเคยฺยุํ “อิทํ เต อมฺโภ ปุริส มํสํ ขาทิตพฺพํ, มูลญฺจ อนุปฺปทาตพฺพนฺ”ติ.\csromanspacer{} เอวเมว โข พฺราหฺมณ พฺราหฺมณา อปฺปฏิญฺญาย เตสํ สมณพฺราหฺมณานํ อถ จ ปนิมานิ จตฺตาริ ธนานิ ปญฺญเปนฺติ.\csromanspacer{} อริยํ โข อหํ พฺราหฺมณ โลกุตฺตรํ ธมฺมํ ปุริสสฺส สนฺธนํ ปญฺญเปมิ.\csromanspacer{} โปราณํ โข ปนสฺส มาตาเปตฺติกํ กุลวํสํ อนุสฺสรโต ยตฺถ ยตฺเถว อตฺตภาวสฺส อภินิพฺพตฺติ โหติ, เตน เตนว สงฺขฺยํ คจฺฉติ.\csromanspacer{} ขตฺติยกุเล เจ อตฺตภาวสฺส อภินิพฺพตฺติ โหติ, ขตฺติโยเตฺวว สงฺขฺยํ คจฺฉติ.\csromanspacer{} พฺราหฺมณกุเล เจ อตฺตภาวสฺส อภินิพฺพตฺติ โหติ, พฺราหฺมโณเตฺวว สงฺขฺยํ คจฺฉติ.\csromanspacer{} เวสฺสกุเล เจ อตฺตภาวสฺส อภินิพฺพตฺติ โหติ, เวสฺโสเตฺวว สงฺขฺยํ คจฺฉติ.\csromanspacer{} สุทฺทกุเล เจ อตฺตภาวสฺส อภินิพฺพตฺติ โหติ, \csromanfolio{392}สุทฺโทเตฺวว สงฺขฺยํ คจฺฉติ.\csromanspacer{} เสยฺยถาปิ พฺราหฺมณ ยํยเทว ปจฺจยํ ปฏิจฺจ อคฺคิ ชลติ, เตน เตเนว สงฺขฺยํ คจฺฉติ.\csromanspacer{} กฏฺฐญฺเจ ปฏิจฺจ อคฺคิ ชลติ, กฏฺฐคฺคิเตฺวว สงฺขฺยํ คจฺฉติ.\csromanspacer{} สกลิกญฺเจ ปฏิจฺจ อคฺคิ ชลติ, สกลิกคฺคิเตฺวว สงฺขฺยํ คจฺฉติ.\csromanspacer{} ติณญฺเจ ปฏิจฺจ อคฺคิ ชลติ, ติณคฺคิเตฺวว สงฺขฺยํ คจฺฉติ.\csromanspacer{} โคมยญฺเจ ปฏิจฺจ อคฺคิ ชลติ, โคมยคฺคิเตฺวว สงฺขฺยํ คจฺฉติ.\csromanspacer{} เอวเมว โข อหํ พฺราหฺมณ อริยํ โลกุตฺตรํ ธมฺมํ ปุริสสฺส สนฺธนํ ปญฺญเปมิ.\csromanspacer{} โปราณํ โข ปนสฺส มาตาเปตฺติกํ กุลวํสํ อนุสฺสรโต ยตฺถ ยตฺเถว อตฺตภาวสฺส อภินิพฺพตฺติ โหติ, เตน เตเนว สงฺขฺยํ คจฺฉติ.}

\prose{ขตฺติยกุเล เจ อตฺตภาวสฺส อภินิพฺพตฺติ โหติ, ขตฺติโยเตฺวว สงฺขฺยํ คจฺฉติ.\csromanspacer{} พฺราหฺมณกุเล เจ อตฺตภาวสฺส อภินิพฺพตฺติ โหติ, พฺราหฺมโณเตฺวว สงฺขฺยํ คจฺฉติ.\csromanspacer{} เวสฺสกุเล เจ อตฺตภาวสฺส อภินิพฺพตฺติ โหติ, เวสฺโสเตฺวว สงฺขฺยํ คจฺฉติ.\csromanspacer{} สุทฺทกุเล เจ อตฺตภาวสฺส อภินิพฺพตฺติ โหติ, สุทฺโทเตฺวว สงฺขฺยํ คจฺฉติ.}

\prose{ขตฺติยกุลา เจปิ พฺราหฺมณ อคารสฺมา อนคาริยํ ปพฺพชิโต โหติ, โส จ ตถาคตปฺปเวทิตํ ธมฺมวินยํ อาคมฺม ปาณาติปาตา ปฏิวิรโต โหติ, อทินฺนาทานา ปฏิวิรโต โหติ, อพฺรหฺมจริยา ปฏิวิรโต โหติ, มุสาวาทา ปฏิวิรโต โหติ, ปิสุณาย วาจาย ปฏิวิรโต โหติ, ผรุสาย วาจาย ปฏิวิรโต โหติ, สมฺผปฺปลาปา ปฏิวิรโต โหติ, อนภิชฺฌาลุ โหติ, อพฺยาปนฺนจิตฺโต โหติ, สมฺมาทิฏฺฐิ โหติ.\csromanspacer{} อาราธโก โหติ ญายํ ธมฺมํ กุสลํ.}

\prose{พฺราหฺมณกุลา เจปิ พฺราหฺมณ อคารสฺมา อนคาริยํ ปพฺพชิโต โหติ, โส จ ตถาคตปฺปเวทิตํ ธมฺมวินยํ อาคมฺม ปาณาติปาตา ปฏิวิรโต โหติ ฯเปฯ สมฺมาทิฏฺฐิ โหติ.\csromanspacer{} อาราธโก โหติ ญายํ ธมฺมํ กุสลํ.}

\prose{เวสฺสกุลา เจปิ พฺราหฺมณ อคารสฺมา อนคาริยํ ปพฺพชิโต โหติ, โส จ ตถาคตปฺปเวทิตํ ธมฺมวินยํ อาคมฺม ปาณาติปาตา ปฏิวิรโต โหติ ฯเปฯ สมฺมาทิฏฺฐิ โหติ.\csromanspacer{} อาราธโก โหติ ญายํ ธมฺมํ กุสลํ.}

\prose{สุทฺทกุลา เจปิ พฺราหฺมณ อคารสฺมา อนคาริยํ ปพฺพชิโต โหติ, โส จ ตถาคตปฺปเวทิตํ ธมฺมวินยํ อาคมฺม ปาณาติปาตา ปฏิวิรโต โหติ ฯเปฯ สมฺมาทิฏฺฐิ โหติ.\csromanspacer{} อาราธโก โหติ ญายํ ธมฺมํ กุสลํ.}

\csromanheader{2}{6. เอสุการีสุตฺต (96)}
\proseitem{442}{\csromanfolio{393}ตํ กิํ มญฺญสิ พฺราหฺมณ, พฺราหฺมโณว นุ โข ปโหติ อสฺมิํ ปเทเส อเวรํ อพฺยาพชฺฌํ เมตฺตจิตฺตํ ภาเวตุํ, โน ขตฺติโย, โน เวสฺโส, โน สุทฺโทติ.\csromanspacer{} โน หิทํ โภ โคตม.\csromanspacer{} ขตฺติโยปิ หิ โภ โคตม ปโหติ อสฺมิํ ปเทเส อเวรํ อพฺยาพชฺฌํ เมตฺตจิตฺตํ ภาเวตุํ.\csromanspacer{} พฺราหฺมโณปิ หิ โภ โคตม.\csromanspacer{} เวสฺโสปิ หิ โภ โคตม.\csromanspacer{} สุทฺโทปิ หิ โภ โคตม.\csromanspacer{} สพฺเพปิ หิ โภ โคตม จตฺตาโร วณฺณา ปโหนฺติ อสฺมิํ ปเทเส อเวรํ อพฺยาพชฺฌํ เมตฺตจิตฺตํ ภาเวตุนฺติ.\csromanspacer{} เอวเมว โข พฺราหฺมณ ขตฺติยกุลา เจปิ อคารสฺมา อนคาริยํ ปพฺพชิโต โหติ, โส จ ตถาคตปฺปเวทิตํ ธมฺมวินยํ อาคมฺม ปาณาติปาตา ปฏิวิรโต โหติ ฯเปฯ สมฺมาทิฏฺฐิ โหติ.\csromanspacer{} อาราธโก โหติ ญายํ ธมฺมํ กุสลํ.}

\csromanmatikaanchor{mtk.55}
\csromantocmarkref{subsection}{6. เอสุการีสุตฺต}{mtk.55}
\bookmark[dest={mtk.55},level=2]{6. เอสุการีสุตฺต}
\prose{พฺราหฺมณกุลา เจปิ พฺราหฺมณ.\csromanspacer{} เวสฺสกุลา เจปิ พฺราหฺมณ.\csromanspacer{} สุทฺทกุลา เจปิ พฺราหฺมณ อคารสฺมา อนคาริยํ ปพฺพชิโต โหติ, โส จ ตถาคตปฺปเวทิตํ ธมฺมวินยํ อาคมฺม ปาณาติปาตา ปฏิวิรโต โหติ ฯเปฯ สมฺมาทิฏฺฐิ โหติ.\csromanspacer{} อาราธโก โหติ ญายํ ธมฺมํ กุสลํ.}

\proseitem{443}{ตํ กิํ มญฺญสิ พฺราหฺมณ, พฺราหฺมโณว นุ โข ปโหติ โสตฺติสินานิํ อาทาย นทิํ คนฺตฺวา รโชชลฺลํ ปวาเหตุํ, โน ขตฺติโย, โน เวสฺโส, โน สุทฺโทติ.\csromanspacer{} โน หิทํ โภ โคตม.\csromanspacer{} ขตฺติโยปิ หิ โภ โคตม ปโหติ โสตฺติสินานิํ อาทาย นทิํ คนฺตฺวา รโชชลฺลํ ปวาเหตุํ.\csromanspacer{} พฺราหฺมโณปิ หิ โภ โคตม.\csromanspacer{} เวสฺโสปิ หิ โภ โคตม.\csromanspacer{} สุทฺโทปิ หิ โภ โคตม.\csromanspacer{} สพฺเพปิ หิ โภ โคตม จตฺตาโร วณฺณา ปโหนฺติ โสตฺติสินานิํ อาทาย นทิํ คนฺตฺวา รโชชลฺลํ ปวาเหตุนฺติ.\csromanspacer{} เอวเมว โข พฺราหฺมณ ขตฺติยกุลา เจปิ อคารสฺมา อนคาริยํ ปพฺพชิโต โหติ, โส จ ตถาคตปฺปเวทิตํ ธมฺมวินยํ อาคมฺม ปาณาติปาตา ปฏิวิรโต โหติ ฯเปฯ สมฺมาทิฏฺฐิ โหติ.\csromanspacer{} อาราธโก โหติ ญายํ ธมฺมํ กุสลํ.}

\prose{พฺราหฺมณกุลา เจปิ พฺราหฺมณ.\csromanspacer{} เวสฺสกุลา เจปิ พฺราหฺมณ.\csromanspacer{} สุทฺทกุลา เจปิ พฺราหฺมณ อคารสฺมา อนคาริยํ ปพฺพชิโต โหติ, โส จ ตถาคตปฺปเวทิตํ ธมฺมวินยํ อาคมฺม ปาณาติปาตา ปฏิวิรโต โหติ ฯเปฯ สมฺมาทิฏฺฐิ โหติ.\csromanspacer{} อาราธโก โหติ ญายํ ธมฺมํ กุสลํ.}

\proseitem{444}{\csromanfolio{394}ตํ กิํ มญฺญสิ พฺราหฺมณ, อิธ ราชา ขตฺติโย มุทฺธาวสิตฺโต นานาชจฺจานํ ปุริสานํ ปุริสสตํ สนฺนิปาเตยฺย.\csromanspacer{} อายนฺตุ โภนฺโต, เย ตตฺถ ขตฺติยกุลา พฺราหฺมณกุลา ราชญฺญกุลา อุปฺปนฺนา สากสฺส วา สาลสฺส วา สลฬสฺส วา จนฺทนสฺส วา ปทุมกสฺส วา อุตฺตรารณิํ อาทาย อคฺคิํ อภินิพฺพตฺเตนฺตุ, เตโช ปาตุกโรนฺตุ.\csromanspacer{} อายนฺตุ ปน โภนฺโต, เย ตตฺถ จณฺฑาลกุลา เนสาทกุลา เวนกุลา รถการกุลา ปุกฺกุสลุลา อุปฺปนฺนา สาปานโทณิยา วา สูกรโทณิยา วา รชกโทณิยา วา เอรณฺฑกฏฺฐสฺส วา อุตฺตรารณิํ อาทาย อคฺคิํ อภินิพฺพตฺเตนฺตุ, เตโช ปาตุกโรนฺตูติ.}

\prose{ตํ กิํ มญฺญสิ พฺราหฺมณ, โย เอวํ นุ โข โส ขตฺติยกุลา พฺราหฺมณกุลา ราชญฺญกุลา อุปฺปนฺเนหิ สากสฺส วา สาลสฺส วา สลฬสฺส วา จนฺทนสฺส วา ปทุมกสฺส วา อุตฺตรารณิํ อาทาย อคฺคิ อภินิพฺพตฺโต เตโช ปาตุกโต, โส เอว นุ ขฺวาสฺส อคฺคิ อจฺจิมา เจว วณฺณวา จ ปภสฺสโร จ, เตน จ สกฺกา อคฺคินา อคฺคิกรณียํ กาตุํ.\csromanspacer{} โย ปน โส จณฺฑาลกุลา เนสาทกุลา เวนกุลา รถการกุลา ปุกฺกุสกุลา อุปฺปนฺเนหิ สาปานโทณิยา วา สูกรโทณิยา วา รชกโทณิยา วา เอรณฺฑกฏฺฐสฺส วา อุตฺตรารณิํ อาทาย อคฺคิ อภินิพฺพตฺโต เตโช ปาตุกโต, สฺวาสฺส อคฺคิ น เจว อจฺจิมา น จ วณฺณวา น จ ปภสฺสโร, น จ เตน สกฺกา อคฺคินา อคฺคิกรณียํ กาตุนฺติ.\csromanspacer{} โน หิทํ โภ โคตม, โยปิ หิ โส โภ โคตม ขตฺติยกุลา พฺราหฺมณกุลา ราชาญฺญกุลา อุปฺปนฺเนหิ สากสฺส วา สาลสฺส วา สลฬสฺส วา จนฺทนสฺส วา ปทุมกสฺส วา อุตฺตรารณิํ อาทาย อคฺคิ อภินิพฺพตฺโต เตโช ปาตุกโต, สฺวาสฺส อคฺคิ อจฺจิมา เจว วณฺณวา จ ปภสฺสโร จ, เตน จ สกฺกา อคฺคินา อคฺคิกรณียํ กาตุํ.\csromanspacer{} โยปิ โส จณฺฑาลกุลา เนสาทกุลา เวนกุลา รถการกุลา ปุกฺกุสกุลา อุปฺปนฺเนหิ สาปานโทณิยา วา สูกรโทณิยา วา รชกโทณิยา วา เอรณฺฑกฏฺฐสฺส วา อุตฺตรารณิํ อาทาย อคฺคิ อภินิพฺพตฺโต เตโช ปาตุกโต, สฺวาสฺส อคฺคิ อจฺจิมา เจว วณฺณวา จ ปภสฺสโร จ, เตน จ สกฺกา อคฺคินา อคฺคิกรณียํ กาตุํ.\csromanspacer{} สพฺโพปิ หิ โภ โคตม อคฺคิ อจฺจิมา เจว วณฺณวา จ ปภสฺสโร จ, สพฺเพนปิ สกฺกา อคฺคินา อคฺคิกรณียํ กาตุนฺติ.}

\csromanheader{2}{7. ธนญฺชานิสุตฺต (97)}
\prose{\csromanfolio{395}เอวเมว โข พฺราหฺมณ ขตฺติยกุลา เจปิ อคารสฺมา อนคาริยํ ปพฺพชิโต โหติ, โส จ ตถาคตปฺปเวทิตํ ธมฺมวินยํ อาคมฺม ปาณาติปาตา ปฏิวิรโต โหติ ฯเปฯ สมฺมาทิฏฺฐิ โหติ.\csromanspacer{} อาราธโก โหติ, ญายํ ธมฺมํ กุสลํ.\csromanspacer{} พฺราหฺมณกุลา เจปิ พฺราหฺมณ.\csromanspacer{} เวสฺสกุลา เจปิ พฺราหฺมณ.\csromanspacer{} สุทฺทกุลา เจปิ พฺราหฺมณ อคารสฺมา อนคาริยํ ปพฺพชิโต โหติ, โส จ ตถาคตปฺปเวทิตํ ธมฺมวินยํ อาคมฺม ปาณาติปาตา ปฏิวิรโต โหติ, อทินฺนาทานา ปฏิวิรโต โหติ, อพฺรหฺมจริยา ปฏิวิรโต โหติ, มุสาวาทา ปฏิวิรโต โหติ, ปิสุณาย วาจาย ปฏิวิรโต โหติ, ผรุสาย วาจาย ปฏิวิรโต โหติ, สมฺผปฺปลาปา ปฏิวิรโต โหติ, อนภิชฺฌาลุ โหติ, อพฺยาปนฺนจิตฺโต โหติ, สมฺมาทิฏฺฐิ โหติ.\csromanspacer{} อาราธโก โหติ, ญายํ ธมฺมํ กุสลนฺติ.}

\prose{เอวํ วุตฺเต เอสุการี พฺราหฺมโณ ภควนฺตํ เอตทโวจ “อภิกฺกนฺตํ โภ โคตม, อภิกฺกนฺตํ โภ โคตม ฯเปฯ อุปาสกํ มํ ภวํ โคตโม ธาเรตุ อชฺชตคฺเค ปาณุเปตํ สรณํ คตนฺ”ติ.}

\nitthitamruled{เอสุการีสุตฺตํ นิฏฺฐิตํ ฉฏฺฐํ.}
\csromanheader{2}{7. ธนญฺชานิสุตฺต}
\proseitem{445}{เอวํ เม สุตํ\csromandash{}เอกํ สมยํ ภควา ราชคเห วิหรติ เวฬุวเน กลนฺทกนิวาเป.\csromanspacer{} เตน โข ปน สมเยน อายสฺมา สาริปุตฺโต ทกฺขิณาคิริสฺมิํ จาริกํ จรติ มหตา ภิกฺขุสํเฆน สทฺธิํ.\csromanspacer{} อถ โข อญฺญตโร ภิกฺขุ ราชคเห วสฺสํวุฏฺโฐ\footnote{วสฺสํวุตฺโถ (สี, สฺยา, กํ, อิ)} เยน ทกฺขิณาคิริ, เยนายสฺมา สาริปุตฺโต เตนุปสงฺกมิ, อุปสงฺกมิตฺวา อายสฺมตา สาริปุตฺเตน สทฺธิํ สมฺโมทิ, สมฺโมทนียํ กถํ สารณียํ วีติสาเรตฺวา เอกมนฺตํ นิสีทิ, เอกมนฺตํ นิสินฺนํ โข ตํ ภิกฺขุํ อายสฺมา สาริปุตฺโต เอตทโวจ “กจฺจาวุโส ภควา อโรโค จ พลวา จา”ติ.\csromanspacer{} อโรโค จาวุโส ภควา พลวา จาติ.\csromanspacer{} กจฺจิ ปนาวุโส ภิกฺขุสํโฆ อโรโค จ พลวา จาติ.\csromanspacer{} ภิกฺขุสํโฆปิ โข อาวุโส อโรโค จ พลวา จาติ.\csromanspacer{} เอตฺถ อาวุโส ตณฺฑุลปาลิทฺวาราย ธนญฺชานิ\footnote{ธานญฺชานิ (สี, อิ)} นาม พฺราหฺมโณ อตฺถิ, \csromanfolio{396}กจฺจาวุโส ธนญฺชานิ พฺราหฺมโณ อโรโค จ พลวา จาติ.\csromanspacer{} ธนญฺชานิปิ โข อาวุโส พฺราหฺมโณ อโรโค จ พลวา จาติ.\csromanspacer{} กจฺจิ ปนาวุโส ธนญฺชานิ พฺราหฺมโณ อปฺปมตฺโตติ.\csromanspacer{} กุโต ปนาวุโส ธนญฺชานิสฺส พฺราหฺมณสฺส อปฺปมาโท, ธนญฺชานิ อาวุโส พฺราหฺมโณ ราชานํ นิสฺสาย พฺราหฺมณคหปติเก วิลุมฺปติ, พฺราหฺมณคหปติเก นิสฺสาย ราชานํ วิลุมฺปติ.\csromanspacer{} ยาปิสฺส ภริยา สทฺธา สทฺธกุลา อานีตา, สาปิ กาลงฺกตา.\csromanspacer{} อญฺญาสฺส ภริยา อสฺสทฺธา อสฺสทฺธกุลา อานีตา.\csromanspacer{} ทุสฺสุตํ วตาวุโส อสฺสุมฺห, ทุสฺสุตํ วตาวุโส อสฺสุมฺห, เย มยํ ธนญฺชานิํ พฺราหฺมณํ ปมตฺตํ อสฺสุมฺห, อปฺเปว จ นาม มยํ กทาจิ กรหจิ ธนญฺชานินา พฺราหฺมเณน สทฺธิํ สมาคจฺเฉยฺยาม, อปฺเปว นาม สิยา โกจิเทว กถาสลฺลาโปติ.}

\proseitem{446}{อถ โข อายสฺมา สาริปุตฺโต ทกฺขิณาคิริสฺมิํ ยถาภิรนฺตํ วิหริตฺวา เยน ราชคหํ เตน จาริกํ ปกฺกามิ.\csromanspacer{} อนุปุพฺเพน จาริกํ จรมาโน เยน ราชคหํ ตทวสริ, ตตฺร สุทํ อายสฺมา สาริปุตฺโต ราชคเห วิหรติ เวฬุวเน กลนฺทกนิวาเป.\csromanspacer{} อถ โข อายสฺมา สาริปุตฺโต ปุพฺพณฺหสมยํ นิวาเสตฺวา ปตฺตจีวรมาทาย ราชคหํ ปิณฺฑาย ปาวิสิ.\csromanspacer{} เตน โข ปน สมเยน ธนญฺชานิ พฺราหฺมโณ พหินคเร คาโว โคฏฺเฐ ทุหาเปติ.\csromanspacer{} อถ โข อายสฺมา สาริปุตฺโต ราชคเห ปิณฺฑาย จริตฺวา ปจฺฉาภตฺตํ ปิณฺฑปาตปฏิกฺกนฺโต เยน ธนญฺชานิ พฺราหฺมโณ เตนุปสงฺกมิ, อทฺทสา โข ธนญฺชานิ พฺราหฺมโณ อายมนฺตํ สาริปุตฺตํ ทูรโตว อาคจฺฉนฺตํ, ทิสฺวาน เยนายสฺมา สาริปุตฺโต เตนุปสงฺกมิ, อุปสงฺกมิตฺวา อายสฺมนฺตํ สาริปุตฺตํ เอตทโวจ “อิโต โภ สาริปุตฺต ปโย ปียตํ, ตาว ภตฺตสฺส กาโล ภวิสฺสตี”ติ.\csromanspacer{} อลํ พฺราหฺมณ, กตํ เม อชฺช ภตฺตกิจฺจํ, อมุกสฺมิํ เม รุกฺขมูเล ทิวาวิหาโร ภวิสฺสติ, ตตฺถ อาคจฺเฉยฺยาสีติ. “เอวํ โภ”ติ โข ธนญฺชานิ พฺราหฺมโณ อายสฺมโต สาริปุตฺตสฺส ปจฺจสฺโสสิ.\csromanspacer{} อถ โข ธนญฺชานิ พฺราหฺมโณ ปจฺฉาภตฺตํ ภุตฺตปาตราโส เยนายสฺมา สาริปุตฺโต เตนุปสงฺกมิ, อุปสงฺกมิตฺวา อายสฺมตา สาริปุตฺเตน สทฺธิํ สมฺโมทิ, สมฺโมทนียํ กถํ สารณียํ วีติสาเรตฺวา เอกมนฺตํ นิสีทิ, เอกมนฺตํ นิสินฺนํ โข ธนญฺชานิํ พฺราหฺมณํ อายสฺมา สาริปุตฺโต เอตทโวจ “กจฺจาสิ ธนญฺชานิ อปฺปมตฺโต”ติ.\csromanspacer{} กุโต โภ สาริปุตฺต อมฺหากํ อปฺปมาโท, เยสํ โน}

\vspace{\tipitakaparskip}
\csromancenter{ธนญฺชานิสุตฺต (97) มาตาปิตโร โปเสตพฺพา, ปุตฺตทาโร โปเสตพฺโพ, ทาสกมฺมกรา โปเสตพฺพา, มิตฺตามจฺจานํ มิตฺตามจฺจกรณียํ กาตพฺพํ, ญาติสาโลหิตานํ ญาติสาโลหิตกรณียํ กาตพฺพํ, อติถีนํ อติถิกรณียํ กาตพฺพํ, ปุพฺพเปตานํ ปุพฺพเปตกรณียํ กาตพฺพํ, เทวตานํ เทวตากรณียํ กาตพฺพํ, รญฺโญ ราชกรณียํ กาตพฺพํ, อยมฺปิ กาโย ปีเณตพฺโพ พฺรูเหตพฺโพติ.}
\proseitem{447}{\csromanfolio{397}ตํ กิํ มญฺญสิ ธนญฺชานิ, อิเธกจฺโจ มาตาปิตูนํ เหตุ อธมฺมจารี วิสมจารี อสฺส, ตเมนํ อธมฺมจริยาวิสมจริยาเหตุ นิรยํ นิรยปาลา อุปกฑฺเฒยฺยุํ.\csromanspacer{} ลเภยฺย นุ โข โส “อหํ โข มาตาปิตูนํ เหตุ อธมฺมจารี วิสมจารี อโหสิํ, มา มํ นิรยํ นิรยปาลา”ติ.\csromanspacer{} มาตาปิตโร วา ปนสฺส ลเภยฺยุํ “เอโส โข อมฺหากํ เหตุ อธมฺมจารี วิสมจารี อโหสิ, มา นํ นิรยํ นิรยปาลา”ติ.\csromanspacer{} โน หิทํ โภ สาริปุตฺต, อถ โข นํ วิกฺกนฺทนฺตํเยว นิรเย นิรยปาลา ปกฺขิเปยฺยุํ.}

\prose{ตํ กิํ มญฺญสิ ธนญฺชานิ, อิเธกจฺโจ ปุตฺตทารสฺส เหตุ อธมฺมจารี วิสมจารี อสฺส, ตเมนํ อธมฺมจริยาวิสมจริยาเหตุ นิรยํ นิรยปาลา อุปกฑฺเฒยฺยุํ.\csromanspacer{} ลเภยฺย นุ โข โส “อหํ โข ปุตฺตทารสฺส เหตุ อธมฺมจารี วิสมจารี อโหสิํ, มา มํ นิรยํ นิรยปาลา”ติ.\csromanspacer{} ปุตฺตทาโร วา ปนสฺส ลเภยฺย “เอโส โข อมฺหากํ เหตุ อธมฺมจารี วิสมจารี อโหสิ, มา นํ นิรยํ นิรยปาลา”ติ, โน หิทํ โภ สาริปุตฺต, อถ โข นํ วิกฺกนฺทนฺตํเยว นิรเย นิรยปาลา ปกฺขิเปยฺยุํ.}

\prose{ตํ กิํ มญฺญสิ ธนญฺชานิ, อิเธกจฺโจ ทาสกมฺมกรโปริสสฺส เหตุ อธมฺมจารี วิสมจารี อสฺส, ตเมนํ อธมฺมจริยาวิสมจริยาเหตุ นิรยํ นิรยปาลา อุปกฑฺเฒยฺยุํ.\csromanspacer{} ลเภยฺย นุ โข โส “อหํ โข ทาสกมฺมกรโปริสสฺส เหตุ อธมฺมจารี วิสมจารี อโหสิํ, มา มํ นิรยํ นิรยปาลา”ติ.\csromanspacer{} ทาสกมฺมกรโปริสา วา ปนสฺส ลเภยฺยุํ “เอโส โข อมฺหากํ เหตุ อธมฺมจารี วิสมจารี อโหสิ, มา นํ นิรยํ นิรยปาลา”ติ.\csromanspacer{} โน หิทํ โภ สาริปุตฺต, อถ โข นํ วิกฺกนฺทนฺตํเยว นิรเย นิรยปาลา ปกฺขิเปยฺยุํ.}

\prose{ตํ กิํ มญฺญสิ ธนญฺชานิ, อิเธกจฺโจ มิตฺตามจฺจานํ เหตุ อธมฺมจารี วิสมจารี อสฺส, ตเมนํ อธมฺมจริยาวิสมจริยาเหตุ \csromanfolio{398}นิรยํ นิรยปาลา อุปกฑฺเฒยฺยุํ.\csromanspacer{} ลเภยฺย นุ โข โส “อหํ โข มิตฺตามจฺจานํ เหตุ อธมฺมจารี วิสมจารี อโหสิํ, มา มํ นิรยํ นิรยปาลา”ติ.\csromanspacer{} มิตฺตามจฺจา วา ปนสฺส ลเภยฺยุํ “เอโส โข อมฺหากํ เหตุ อธมฺมจารี วิสมจารี อโหสิ, มา นํ นิรยํ นิรยปาลา”ติ.\csromanspacer{} โน หิทํ โภ สาริปุตฺต.\csromanspacer{} อถ โข นํ วิกฺกนฺทนฺตํเยว นิรเย นิรยปาลา ปกฺขิเปยฺยุํ.}

\prose{ตํ กิํ มญฺญสิ ธนญฺชานิ, อิเธกจฺโจ ญาติสาโลหิตานํ เหตุ อธมฺมจารี วิสมจารี อสฺส, ตเมนํ อธมฺมจริยาวิสมจริยาเหตุ นิรยํ นิรยปาลา อุปกฑฺเฒยฺยุํ.\csromanspacer{} ลเภยฺย นุ โข โส “อหํ โข ญาติสาโลหิตานํ เหตุ อธมฺมจารี วิสมจารี อโหสิํ, มา มํ นิรยํ นิรยปาลา”ติ.\csromanspacer{} ญาติสาโลหิตา วา ปนสฺส ลเภยฺยุํ “เอโส โข อมฺหากํ เหตุ อธมฺมจารี วิสมจารี อโหสิ, มา นํ นิรยํ นิรยปาลา”ติ.\csromanspacer{} โน หิทํ โภ สาริปุตฺต, อถ โข นํ วิกฺกนฺทนฺตํเยว นิรเย นิรยปาลา ปกฺขิเปยฺยุํ.}

\prose{ตํ กิํ มญฺญสิ ธนญฺชานิ, อิเธกจฺโจ อติถีนํ เหตุ อธมฺมจารี วิสมจารี อสฺส, ตเมนํ อธมฺมจริยาวิสมจริยาเหตุ นิรยํ นิรยปาลา อุปกฑฺเฒยฺยุํ.\csromanspacer{} ลเภยฺย นุ โข โส “อหํ โข อติถีนํ เหตุ อธมฺมจารี วิสมจารี อโหสิํ, มา มํ นิรยํ นิรยปาลา”ติ.\csromanspacer{} อติถี วา ปนสฺส ลเภยฺยุํ “เอโส โข อมฺหากํ เหตุ อธมฺมจารี วิสมจารี อโหสิ, มา นํ นิรยํ นิรยปาลา”ติ.\csromanspacer{} โน หิทํ โภ สาริปุตฺต, อถ โข นํ วิกฺกนฺทนฺตํเยว นิรเย นิรยปาลา ปกฺขิเปยฺยุํ.}

\prose{ตํ กิํ มญฺญสิ ธญฺชานิ, อิเธกจฺโจ ปุพฺพเปตานํ เหตุ อธมฺมจารี วิสมจารี อสฺส, ตเมนํ อธมฺมจริยาวิสมจริยาเหตุ นิรยํ นิรยปาลา อุปกฑฺเฒยฺยุํ.\csromanspacer{} ลเภยฺย นุ โข โส “อหํ โข ปุพฺพเปตานํ เหตุ อธมฺมจารี วิสมจารี อโหสิํ, มา มํ นิรยํ นิรยปาลา”ติ.\csromanspacer{} ปุพฺพเปตา วา ปนสฺส ลเภยฺยุํ “เอโส โข อมฺหากํ เหตุ อธมฺมจารี วิสมจารี อโหสิ, มา นํ นิรยํ นิรยปาลา”ติ.\csromanspacer{} โน หิทํ โภ สาริปุตฺต, อถ โข นํ วิกฺกนฺทนฺตํเยว นิรเย นิรยปาลา ปกฺขิเปยฺยุํ.}

\csromanheader{2}{7. ธนญฺชานิสุตฺต (97)}
\prose{\csromanfolio{399}ตํ กิํ มญฺญสิ ธนญฺชานิ อิเธกจฺโจ เทวตานํ เหตุ อธมฺมจารี วิสมจารี อสฺส, ตเมนํ อธมฺมจริยาวิสมจริยาเหตุ นิรยํ นิรยปาลา อุปกฑฺเฒยฺยุํ.\csromanspacer{} ลเภยฺย นุ โข โส “อหํ โข เทวตานํ เหตุ อธมฺมจารี วิสมจารี อโหสิํ, มา มํ นิรยํ นิรยปาลา”ติ.\csromanspacer{} เทวตา วา ปนสฺส ลเภยฺยุํ “เอโส โข อมฺหากํ เหตุ อธมฺมจารี วิสมจารี อโหสิ, มา นํ นิรยํ นิรยปาลา”ติ.\csromanspacer{} โน หิทํ โภ สาริปุตฺต, อถ โข นํ วิกฺกนฺทนฺตํเยว นิรเย นิรยปาลา ปกฺขิเปยฺยุํ.}

\prose{ตํ กิํ มญฺญสิ ธนญฺชานิ, อิเธกจฺโจ รญฺโญ เหตุ อธมฺมจารี วิสมจารี อสฺส, ตเมนํ อธมฺมจริยาวิสมจริยาเหตุ นิรยํ นิรยปาลา อุปกฑฺเฒยฺยุํ.\csromanspacer{} ลเภยฺย นุ โข โส “อหํ โข รญฺโญ เหตุ อธมฺมจารี วิสมจารี อโหสิํ, มา มํ นิรยํ นิรยปาลา”ติ.\csromanspacer{} ราชา วา ปนสฺส ลเภยฺย “เอโส โข อมฺหากํ เหตุ อธมฺมจารี วิสมจารี อโหสิ, มา นํ นิรยํ นิรยปาลา”ติ.\csromanspacer{} โน หิทํ โภ สาริปุตฺต, อถ โข นํ วิกฺกนฺทนฺตํเยว นิรเย นิรยปาลา ปกฺขิเปยฺยุํ.}

\prose{ตํ กิํ มญฺญสิ ธนญฺชานิ, อิเธกจฺโจ กายสฺส ปีณนาเหตุ พฺรูหนาเหตุ อธมฺมจารี วิสมจารี อสฺส, ตเมนํ อธมฺมจริยาวิสมจริยาเหตุ นิรยํ นิรยปาลา อุปกฑฺเฒยฺยุํ.\csromanspacer{} ลเภยฺย นุ โข โส “อหํ โข กายสฺส ปีณนาเหตุ พฺรูหนาเหตุ อธมฺมจารี วิสมจารี อโหสิํ, มา มํ นิรยํ นิรยปาลา”ติ.\csromanspacer{} ปเร วา ปนสฺส ลเภยฺยุํ “เอโส โข กายสฺส ปีณนาเหตุ พฺรูหนาเหตุ อธมฺมจารี วิสมจารี อโหสิ, มา นํ นิรยํ นิรยปาลา”ติ.\csromanspacer{} โน หิทํ โภ สาริปุตฺต, อถ โข นํ วิกฺกนฺทนฺตํเยว นิรเย นิรยปาลา ปกฺขิเปยฺยุํ.}

\csromanmatikaanchor{mtk.56}
\csromantocmarkref{subsection}{7. ธนญฺชานิสุตฺต}{mtk.56}
\bookmark[dest={mtk.56},level=2]{7. ธนญฺชานิสุตฺต}
\proseitem{448}{ตํ กิํ มญฺญสิ ธนญฺชานิ, โย วา มาตาปิตูนํ เหตุ อธมฺมจารี วิสมจารี อสฺส, โย วา มาตาปิตูนํ เหตุ ธมฺมจารี สมจารี อสฺส, กตมํ เสยฺโยติ.\csromanspacer{} โย หิ โภ สาริปุตฺต มาตาปิตูนํ เหตุ อธมฺมจารี วิสมจารี อสฺส, น ตํ เสยฺโย.\csromanspacer{} โย จ โข โภ สาริปุตฺต มาตาปิตูนํ เหตุ ธมฺมจารี สมจารี อสฺส, ตเทเวตฺถ เสยฺโย.\csromanspacer{} อธมฺมจริยาวิสมจริยาหิ โภ สาริปุตฺต ธมฺมจริยาสมจริยา เสยฺโยติ.\csromanspacer{} อตฺถิ โข ธนญฺชานิ อญฺเญสํ เหตุกา \csromanfolio{400}ธมฺมิกา กมฺมนฺตา, เยหิ สกฺกา มาตาปิตโร เจว โปเสตุํ, น จ ปาปกมฺมํ กาตุํ, ปุญฺญญฺจ ปฏิปทํ ปฏิปชฺชิตุํ.}

\prose{ตํ กิํ มญฺญสิ ธนญฺชานิ, โย วา ปุตฺตทารสฺส เหตุ อธมฺมจารี วิสมจารี อสฺส, โย วา ปุตฺตทารสฺส เหตุ ธมฺมจารี สมจารี อสฺส, กตมํ เสยฺโยติ.\csromanspacer{} โย หิ โภ สาริปุตฺต ปุตฺตทารสฺส เหตุ อธมฺมจารี วิสมจารี อสฺส, น ตํ เสยฺโย.\csromanspacer{} โย จ โข โภ สาริปุตฺต ปุตฺตทารสฺส เหตุ ธมฺมจารี สมจารี อสฺส, ตเทเวตฺถ เสยฺโย.\csromanspacer{} อธมฺมจริยาวิสมจริยาหิ โภ สาริปุตฺต ธมฺมจริยาสมจริยา เสยฺโยติ.\csromanspacer{} อตฺถิ โข ธนญฺชานิ อญฺเญสํ เหตุกา ธมฺมิกา กมฺมนฺตา, เยหิ สกฺกา ปุตฺตทารญฺเจว โปเสตุํ, น จ ปาปกมฺมํ กาตุํ, ปุญฺญญฺจ ปฏิปทํ ปฏิปชฺชิตุํ.}

\prose{ตํ กิํ มญฺญสิ ธนญฺชานิ, โย วา ทาสกมฺมกรโปริสสฺส เหตุ อธมฺมจารี วิสมจารี อสฺส, โย วา ทาสกมฺมกรโปริสสฺส เหตุ ธมฺมจารี สมจารี อสฺส, กตมํ เสยฺโยติ.\csromanspacer{} โย หิ โภ สาริปุตฺต ทาสกมฺมกรโปริสสฺส เหตุ อธมฺมจารี วิสมจารี อสฺส, น ตํ เสยฺโย.\csromanspacer{} โย จ โข โภ สาริปุตฺต ทาสกมฺมกรโปริสสฺส เหตุ ธมฺมจารี สมจารี อสฺส, ตเทเวตฺถ เสยฺโย.\csromanspacer{} อธมฺมจริยาวิสมจริยาหิ โภ สาริปุตฺต ธมฺมจริยาสมจริยา เสยฺโยติ.\csromanspacer{} อตฺถิ โข ธนญฺชานิ อญฺเญสํ เหตุกา ธมฺมิกา กมฺมนฺตา, เยหิ สกฺกา ทาสกมฺมกรโปริเส เจว โปเสตุํ, น จ ปาปกมฺมํ กาตุํ, ปุญฺญญฺจ ปฏิปทํ ปฏิปชฺชิตุํ.}

\prose{ตํ กิํ มญฺญสิ ธนญฺชานิ, โย วา มิตฺตามจฺจานํ เหตุ อธมฺมจารี วิสมจารี อสฺส, โย วา มิตฺตามจฺจานํ เหตุ ธมฺมจารี สมจารี อสฺส, กตมํ เสยฺโยติ.\csromanspacer{} โย หิ โภ สาริปุตฺต มิตฺตามจฺจานํ เหตุ อธมฺมจารี วิสมจารี อสฺส, น ตํ เสยฺโย.\csromanspacer{} โย จ โข โภ สาริปุตฺต มิตฺตามจฺจานํ เหตุ ธมฺมจารี สมจารี อสฺส, ตเทเวตฺถ เสยฺโย.\csromanspacer{} อธมฺมจริยาวิสมจริยาหิ โภ สาริปุตฺต ธมฺมจริยาสมจริยา เสยฺโยติ.\csromanspacer{} อตฺถิ โข ธนญฺชานิ อญฺเญสํ เหตุกา ธมฺมิกา กมฺมนฺตา, เยหิ สกฺกา มิตฺตามจฺจานญฺเจว มิตฺตามจฺจกรณียํ กาตุํ, น จ ปาปกมฺมํ กาตุํ, ปุญฺญญฺจ ปฏิปทํ ปฏิปชฺชิตุํ.}

\csromanheader{2}{7. ธนญฺชานิสุตฺต (97)}
\prose{\csromanfolio{401}ตํ กิํ มญฺญสิ ธนญฺชานิ, โย วา ญาติสาโลหิตานํ เหตุ อธมฺมจารี วิสมจารี อสฺส, โย วา ญาติสาโลหิตานํ เหตุ ธมฺมจารี สมจารี อสฺส, กตมํ เสยฺโยติ.\csromanspacer{} โย หิ โภ สาริปุตฺต ญาติสาโลหิตานํ เหตุ อธมฺมจารี วิสมจารี อสฺส, น ตํ เสยฺโย.\csromanspacer{} โย จ โข โภ สาริปุตฺต ญาติสาโลหิตานํ เหตุ ธมฺมจารี สมจารี อสฺส, ตเทเวตฺถ เสยฺโย.\csromanspacer{} อธมฺมจริยาวิสมจริยาหิ โภ สาริปุตฺต ธมฺมจริยาสมจริยา เสยฺโยติ.\csromanspacer{} อตฺถิ โข ธนญฺชานิ อญฺเญสํ เหตุกา ธมฺมิกา กมฺมนฺตา, เยหิ สกฺกา ญาติสาโลหิตานญฺเจว ญาติสาโลหิตกรณียํ กาตุํ, น จ ปาปกมฺมํ กาตุํ, ปุญฺญญฺจ ปฏิปทํ ปฏิปชฺชิตุํ.}

\prose{ตํ กิํ มญฺญสิ ธนญฺชานิ, โย วา อติถีนํ เหตุ อธมฺมจารี วิสมจารี อสฺส, โย วา อติถีนํ เหตุ ธมฺมจารี สมจารี อสฺส, กตมํ เสยฺโยติ.\csromanspacer{} โย หิ โภ สาริปุตฺต อติถีนํ เหตุ อธมฺมจารี วิสมจารี อสฺส, น ตํ เสยฺโย.\csromanspacer{} โย จ โข โภ สาริปุตฺต อติถีนํ เหตุ ธมฺมจารี สมจารี อสฺส, ตเทเวตฺถ เสยฺโย.\csromanspacer{} อธมฺมจริยาวิสมจริยาหิ โภ สาริปุตฺต ธมฺมจริยาสมจริยา เสยฺโยติ.\csromanspacer{} อตฺถิ โข ธนญฺชานิ อญฺเญสํ เหตุกา ธมฺมิกา กมฺมนฺตา, เยหิ สกฺกา อติถีนญฺเจว อติถิกรณียํ กาตุํ, น จ ปาปกมฺมํ กาตุํ, ปุญฺญญฺจ ปฏิปทํ ปฏิปชฺชิตุํ.}

\prose{ตํ กิํ มญฺญสิ ธนญฺชานิ, โย วา ปุพฺพเปตานํ เหตุ อธมฺมจารี วิสมจารี อสฺส, โย วา ปุพฺพเปตานํ เหตุ ธมฺมจารี สมจารี อสฺส, กตมํ เสยฺโยติ.\csromanspacer{} โย หิ โภ สาริปุตฺต ปุพฺพเปตานํ เหตุ อธมฺมจารี วิสมจารี อสฺส, น ตํ เสยฺโย.\csromanspacer{} โย จ โข โภ สาริปุตฺต ปุพฺพเปตานํ เหตุ ธมฺมจารี สมจารี อสฺส, ตเทเวตฺถ เสยฺโย.\csromanspacer{} อธมฺมจริยาวิสมจริยาหิ โภ สาริปุตฺต ธมฺมจริยาสมจริยาเสยฺโยติ.\csromanspacer{} อตฺถิ โข ธนญฺชานิ อญฺเญสํ เหตุกา ธมฺมิกา กมฺมนฺตา, เยหิ สกฺกา ปุพฺพเปตานญฺเจว ปุพฺพเปตกรณียํ กาตุํ, น จ ปาปกมฺมํ กาตุํ, ปุญฺญญฺจ ปฏิปทํ ปฏิปชฺชิตุํ.}

\prose{ตํ กิํ มญฺญสิ ธนญฺชานิ, โย วา เทวตานํ เหตุ อธมฺมจารี วิสมจารี อสฺส, โย วา เทวตานํ เหตุ ธมฺมจารี สมจารี อสฺส,}

\prose{\csromanfolio{402}กตมํ เสยฺโยติ.\csromanspacer{} โย หิ โภ สาริปุตฺต เทวตานํ เหตุ อธมฺมจารี วิสมจารี อสฺส, น ตํ เสยฺโย.\csromanspacer{} โย จ โข โภ สาริปุตฺต เทวตานํ เหตุ ธมฺมจารี สมจารี อสฺส, ตเทเวตฺถ เสยฺโย.\csromanspacer{} อธมฺมจริยาวิสมจริยาหิ โภ สาริปุตฺต ธมฺมจริยาสมจริยา เสยฺโยติ.\csromanspacer{} อตฺถิ โข ธนญฺชานิ อญฺเญสํ เหตุกา ธมฺมิกา กมฺมนฺตา, เยหิ สกฺกา เทวตานญฺเจว เทวตากรณียํ กาตุํ, น จ ปาปกมฺมํ กาตุํ, ปุญฺญญฺจ ปฏิปทํ ปฏิปชฺชิตุํ.}

\prose{ตํ กิํ มญฺญสิ ธนญฺชานิ, โย วา รญฺโญ เหตุ อธมฺมจารี วิสมจารี อสฺส, โย วา รญฺโญ เหตุ ธมฺมจารี สมจารี อสฺส, กตมํ เสยฺโยติ.\csromanspacer{} โย หิ โภ สาริปุตฺต รญฺโญ เหตุ อธมฺมจารี วิสมจารี อสฺส, น ตํ เสยฺโย.\csromanspacer{} โย จ โข โภ สาริปุตฺต รญฺโญ เหตุ ธมฺมจารี สมจารี อสฺส, ตเทเวตฺถ เสยฺโย.\csromanspacer{} อธมฺมจริยาวิสมจริยาหิ โภ สาริปุตฺต ธมฺมจริยาสมจริยา เสยฺโยติ.\csromanspacer{} อตฺถิ โข ธนญฺชานิ อญฺเญสํ เหตุกา ธมฺมิกา กมฺมนฺตา, เยหิ สกฺกา รญฺโญ เจว ราชกรณียํ กาตุํ, น จ ปาปกมฺมํ กาตุํ, ปุญฺญญฺจ ปฏิปทํ ปฏิปชฺชิตุํ.}

\prose{ตํ กิํ มญฺญสิ ธนญฺชานิ, โย วา กายสฺส ปีณนาเหตุ พฺรูหนาเหตุ อธมฺมจารี วิสมจารี อสฺส, โย วา กายสฺส ปีณนาเหตุ พฺรูหนาเหตุ ธมฺมจารี สมจารี อสฺส, กตมํ เสยฺโยติ.\csromanspacer{} โย หิ โภ สาริปุตฺต กายสฺส ปีณนาเหตุ พฺรูหนาเหตุ อธมฺมจารี วิสมจารี อสฺส, น ตํ เสยฺโย.\csromanspacer{} โย จ โข โภ สาริปุตฺต กายสฺส ปีณนาเหตุ พฺรูหนาเหตุ ธมฺมจารี สมจารี อสฺส, ตเทเวตฺถ เสยฺโย.\csromanspacer{} อธมฺมจริยาวิสมจริยาหิ โภ สาริปุตฺต ธมฺมจริยาสมจริยา เสยฺโยติ.\csromanspacer{} อตฺถิ โข ธนญฺชานิ อญฺเญสํ เหตุกา ธมฺมิกา กมฺมนฺตา, เยหิ สกฺกา กายญฺเจว ปีเณตุํ พฺรูเหตุํ, น จ ปาปกมฺมํ กาตุํ, ปุญฺญญฺจ ปฏิปทํ ปฏิปชฺชิตุนฺติ.}

\proseitem{449}{อถ โข ธนญฺชานิ พฺราหฺมโณ อายสฺมโต สาริปุตฺตสฺส ภาสิตํ อภินนฺทิตฺวา อนุโมทิตฺวา อุฏฺฐายาสนา ปกฺกามิ.\csromanspacer{} อถ โข ธนญฺชานิ พฺราหฺมโณ อปเรน สมเยน อาพาธิโก อโหสิ ทุกฺขิโต พาฬฺหคิลาโน.\csromanspacer{} อถ โข ธนญฺชานิ พฺราหฺมโณ อญฺญตรํ ปุริสํ}

\vspace{\tipitakaparskip}
\csromancenter{ธนญฺชานิสุตฺต (97) อามนฺเตสิ “เอหิ ตฺวํ อมฺโภ ปุริส เยน ภควา เตนุปสงฺกม, อุปสงฺกมิตฺวา มม วจเนน ภควโต ปาเท สิรสา วนฺทาหิ ‘ธนญฺชานิ ภนฺเต พฺราหฺมโณ อาพาธิโก ทุกฺขิโต พาฬฺหคิลาโน, โส ภควโต ปาเท สิรสา วนฺทตี’ติ, เยน จายสฺมา สาริปุตฺโต เตนุปสงฺกม, อุปสงฺกมิตฺวา มม วจเนน อายสฺมโต สาริปุตฺตสฺส ปาเท สิรสา วนฺทาหิ ‘ธนญฺชานิ ภนฺเต พฺราหฺมโณ อาพาธิโก ทุกฺขิโต พาฬฺหคิลาโน, โส อายสฺมโต สาริปุตฺตสฺส ปาเท สิรสา วนฺทตี’ติ, เอวญฺจ วเทหิ ‘สาธุ กิร ภนฺเต อายสฺมา สาริปุตฺโต เยน ธนญฺชานิสฺส พฺราหฺมณสฺส นิเวสนํ เตนุปสงฺกมตุ อนุกมฺปํ อุปาทายา’ติ”. “เอวํ ภนฺเต”ติ โข โส ปุริโส ธนญฺชานิสฺส พฺราหฺมณสฺส ปฏิสฺสุตฺวา เยน ภควา เตนุปสงฺกมิ, อุปสงฺกมิตฺวา ภควนฺตํ อภิวาเทตฺวา เอกมนฺตํ นิสีทิ, เอกมนฺตํ นิสินฺโน โข โส ปุริโส ภควนฺตํ เอตทโวจ “ธนญฺชานิ ภนฺเต พฺราหฺมโณ อาพาธิโก ทุกฺขิโต พาฬฺหคิลาโน, โส ภควโต ปาเท สิรสา วนฺทตี”ติ.\csromanspacer{} เยน จายสฺมา สาริปุตฺโต เตนุปสงฺกมิ, อุปสงฺกมิตฺวา อายสฺมนฺตํ สาริปุตฺตํ อภิวาเทตฺวา เอกมนฺตํ นิสีทิ, เอกมนฺตํ นิสินฺโน โข โส ปุริโส อายสฺมนฺตํ สาริปุตฺตํ เอตทโวจ “ธนญฺชานิ ภนฺเต พฺราหฺมโณ อาพาธิโก ทุกฺขิโต พาฬฺหคิลาโน, โส อายสฺมโต สาริปุตฺตสฺส ปาเท สิรสา วนฺทติ, เอวญฺจ วเทติ ‘สาธุ กิร ภนฺเต อายสฺมา สาริปุตฺโต เยน ธนญฺชานิสฺส พฺราหฺมณสฺส นิเวสนํ เตนุปสงฺกมตุ อนุกมฺปํ อุปาทายา’ติ”.\csromanspacer{} อธิวาเสสิ โข อายสฺมา สาริปุตฺโต ตุณฺหีภาเวน.}
\proseitem{450}{\csromanfolio{403}อถ โข อายสฺมา สาริปุตฺโต นิวาเสตฺวา ปตฺตจีวรมาทาย เยน ธนญฺชานิสฺส พฺราหฺมณสฺส นิเวสนํ เตนุปสงฺกมิ, อุปสงฺกมิตฺวา ปญฺญตฺเต อาสเน นิสีทิ, นิสชฺช โข อายสฺมา สาริปุตฺโต ธนญฺชานิํ พฺราหฺมณํ เอตทโวจ “กจฺจิ เต ธนญฺชานิ ขมนียํ, กจฺจิ ยาปนียํ, กจฺจิ ทุกฺขา เวทนา ปฏิกฺกมนฺติ, โน อภิกฺกมนฺติ, ปฏิกฺกโมสานํ ปญฺญายติ, โน อภิกฺกโม”ติ.\csromanspacer{} น เม โภ สาริปุตฺต ขมนียํ, น ยาปนียํ, พาฬฺหา เม ทุกฺขา เวทนา อภิกฺกมนฺติ, โน ปฏิกฺกมนฺติ, อภิกฺกโมสานํ ปญฺญายติ, โน ปฏิกฺกโม.\csromanspacer{} เสยฺยถาปิ โภ สาริปุตฺต พลวา ปุริโส ติเณฺหน สิขเรน มุทฺธนิ\footnote{มุทฺธานํ (สี, สฺยา, กํ, อิ)} อภิมตฺเถยฺย.\csromanspacer{} เอวเมว \csromanfolio{404}โข โภ สาริปุตฺต อธิมตฺตา วาตา มุทฺธนิ จ อูหนนฺติ, น เม โภ สาริปุตฺต ขมนียํ, น ยาปนียํ, พาฬฺหา เม ทุกฺขา เวทนา อภิกฺกมนฺติ, โน ปฏิกฺกมนฺติ, อภิกฺกโมสานํ ปญฺญายติ, โน ปฏิกฺกโม.\csromanspacer{} เสยฺยถาปิ โภ สาริปุตฺต พลวา ปุริโส ทเฬฺหน วรตฺตกฺขณฺเฑน\footnote{วรตฺตพนฺธเนน (สี, อิ)} สีเส สีสเวฐํ ทเทยฺย.\csromanspacer{} เอวเมว โข โภ สาริปุตฺต อธิมตฺตา สีเส สีสเวทนา, น เม โภ สาริปุตฺต ขมนียํ, น ยาปนียํ, พาฬฺหา เม ทุกฺขา เวทนา อภิกฺกมนฺติ, โน ปฏิกฺกมนฺติ, อภิกฺกโมสานํ ปญฺญายติ, โน ปฏิกฺกโม.\csromanspacer{} เสยฺยถาปิ โภ สาริปุตฺต ทกฺโข โคฆาตโก วา โคฆาตกนฺเตวาสี วา ติเณฺหน โควิกนฺตเนน กุจฺฉิํ ปริกนฺเตยฺย.\csromanspacer{} เอวเมว โข โภ สาริปุตฺต อธิมตฺตา วาตา กุจฺฉิํ ปริกนฺตนฺติ.\csromanspacer{} น เม โภ สาริปุตฺต ขมนียํ, น ยาปนียํ, พาฬฺหา เม ทุกฺขา เวทนา อภิกฺกมนฺติ, โน ปฏิกฺกมนฺติ, อภิกฺกโมสานํ ปญฺญายติ, โน ปฏิกฺกโม.\csromanspacer{} เสยฺยถาปิ โภ สาริปุตฺต เทฺว พลวนฺโต ปุริสา ทุพฺพลตรํ ปุริสํ นานาพาหาสุ คเหตฺวา องฺคารกาสุยา สนฺตาเปยฺยุํ สมฺปริตาเปยฺยุํ.\csromanspacer{} เอวเมว โข โภ สาริปุตฺต อธิมตฺโต กายสฺมิํ ฑาโห, น เม โภ สาริปุตฺต ขมนียํ, น ยาปนียํ, พาฬฺหา เม ทุกฺขา เวทนา อภิกฺกมนฺติ, โน ปฏิกฺกมนฺติ, อภิกฺกโมสานํ มญฺญายติ, โน ปฏิกฺกโมติ.}

\proseitem{451}{ตํ กิํ มญฺญสิ ธนญฺชานิ, กตมํ เสยฺโย นิรโย วา ติรจฺฉานโยนิ วาติ.\csromanspacer{} นิรยา โภ สาริปุตฺต ติรจฺฉานโยนิ เสยฺโยติ.\csromanspacer{} ตํ กิํ มญฺญสิ ธนญฺชานิ, กตมํ เสยฺโย ติรจฺฉานโยนิ วา เปตฺถิวิสโย วาติ.\csromanspacer{} ติรจฺฉานโยนิยา โภ สาริปุตฺต เปตฺติวิสโย เสยฺโยติ.\csromanspacer{} ตํ กิํ มญฺญสิ ธนญฺชานิ, กตมํ เสยฺโย เปตฺติวิสโย วา มนุสฺสา วาติ.\csromanspacer{} เปตฺติวิสยา โภ สาริปุตฺต มนุสฺสา เสยฺโยติ.\csromanspacer{} ตํ กิํ มญฺญสิ ธนญฺชานิ, กตมํ เสยฺโย มนุสฺสา วา จาตุมหาราชิกา\footnote{จาตุมฺมหาราชิกา (สี, สฺยา, กํ, อิ)} วา เทวาติ.\csromanspacer{} มนุสฺเสหิ โภ สาริปุตฺต จาตุมหาราชิกา เทวา เสยฺโยติ.\csromanspacer{} ตํ กิํ มญฺญสิ ธนญฺชานิ, กตมํ เสยฺโย จาตุมหาราชิกา วา เทวา ตาวติํสา วา เทวาติ.\csromanspacer{} จาตุมหาราชิเกหิ โภ สาริปุตฺต เทเวหิ ตาวติํสา เทวา เสยฺโยติ.\csromanspacer{} ตํ กิํ มญฺญสิ ธนญฺชานิ, กตมํ เสยฺโย ตาวติํสา วา เทวา ยามา วา เทวาติ.}

\vspace{\tipitakaparskip}
\csromancenter{ธนญฺชานิสุตฺต (97) ตาวติํเสหิ โภ สาริปุตฺต เทเวหิ ยามา เทวา เสยฺโยติ.\csromanspacer{} ตํ กิํ มญฺญสิ ธนญฺชานิ, กตมํ เสยฺโย ยามา วา เทวา ตุสิตา วา เทวาติ.\csromanspacer{} ยาเมหิ โภ สาริปุตฺต เทเวหิ ตุสิตา เทวา เสยฺโยติ.\csromanspacer{} ตํ กิํ มญฺญสิ ธนญฺชานิ, กตมํ เสยฺโย ตุสิตา วา เทวา นิมฺมานรตี วา เทวาติ.\csromanspacer{} ตุสิเตหิ โภ สาริปุตฺต เทเวหิ นิมฺมานรตี เทวา เสยฺโยติ.\csromanspacer{} ตํ กิํ มญฺญสิ ธนญฺชานิ, กตมํ เสยฺโย นิมฺมานรตี วา เทวา ปรนิมฺมิตวสวตฺตี วา เทวาติ.\csromanspacer{} นิมฺมานรตีหิ โภ สาริปุตฺต เทเวหิ ปรนิมฺมิตวสวตฺตี เทวา เสยฺโยติ.\csromanspacer{} ตํ กิํ มญฺญสิ ธนญฺชานิ, กตมํ เสยฺโย ปรนิมฺมิตวสวตฺตี วา เทวา พฺรหฺมโลโก วาติ.\csromanspacer{} พฺรหฺมโลโกติ 1 ภวํ สาริปุตฺโต อาห, พฺรหฺมโลโกติ ภวํ สาริปุตฺโต อาหาติ1.}
\prose{\csromanfolio{405}อถ โข อายสฺมโต สาริปุตฺตสฺส เอตทโหสิ “อิเม โข พฺราหฺมณา พฺรหฺมโลกาธิมุตฺตา, ยํนูนาหํ ธนญฺชานิสฺส พฺราหฺมณสฺส พฺรหฺมานํ สหพฺยตาย มคฺคํ เทเสยฺยนฺ”ติ.\csromanspacer{} พฺรหฺมานํ เต ธนญฺชานิ สหพฺยตาย มคฺคํ เทเสสฺสามิ, ตํ สุณาหิ สาธุกํ มนสิ กโรหิ ภาสิสฺสามี”ติ. “เอวํ โภ”ติ โข ธนญฺชานิ พฺราหฺมโณ อายสฺมโต สาริปุตฺตสฺส ปจฺจสฺโสสิ.\csromanspacer{} อายสฺมา สาริปุตฺโต เอตทโวจ “กตโม จ ธนญฺชานิ พฺรหฺมานํ สหพฺยตาย มคฺโค.\csromanspacer{} อิธ ธนญฺชานิ ภิกฺขุ เมตฺตาสหคเตน เจตสา เอกํ ทิสํ ผริตฺวา วิหรติ.\csromanspacer{} ตถา ทุติยํ.\csromanspacer{} ตถา ตติยํ.\csromanspacer{} ตถา จตุตฺถํ.\csromanspacer{} อิติ อุทฺธมโธ ติริยํ สพฺพธิ สพฺพตฺตตาย สพฺพาวนฺตํ โลกํ เมตฺตาสหคเตน เจตสา วิปุเลน มหคฺคเตน อปฺปมาเณน อเวเรน อพฺยาพชฺเฌน ผริตฺวา วิหรติ.\csromanspacer{} อยํ โข ธนญฺชานิ พฺรหฺมานํ สหพฺยตาย มคฺโค.}

\proseitem{452}{ปุน จปรํ ธนญฺชานิ ภิกฺขุ กรุณาสหคเตน เจตสา ฯเปฯ มุทิตาสหคเตน เจตสา.\csromanspacer{} อุเปกฺขาสหคเตน เจตสา เอกํ ทิสํ ผริตฺวา วิหรติ.\csromanspacer{} ตถา ทุติยํ.\csromanspacer{} ตถา ตติยํ.\csromanspacer{} ตถา จตุตฺถํ.\csromanspacer{} อิติ อุทฺธมโธ ติริยํ สพฺพธิ สพฺพตฺตตาย สพฺพาวนฺตํ โลกํ อุเปกฺขาสหคเตน เจตสา วิปุเลน มหคฺคเตน อปฺปมาเณน อเวเรน อพฺยาพชฺเฌน ผริตฺวา วิหรติ.\csromanspacer{} อยํ โข ธนญฺชานิ พฺรหฺมานํ สหพฺยตาย มคฺโค”ติ.\csromanspacer{} เตน หิ โภ สาริปุตฺต มม วจเนน ภควโต ปาเท สิรสา วนฺทาหิ \csromanfolio{406}“ธนญฺชานิ ภนฺเต พฺราหฺมโณ อาพาธิโก ทุกฺขิโต พาฬฺหคิลาโน, โส ภควโต ปาเท สิรสา วนฺทตี”ติ.\csromanspacer{} อถ โข อายสฺมา สาริปุตฺโต ธนญฺชานิํ พฺราหฺมณํ สติ อุตฺตริกรณีเย หีเน พฺรหฺมโลเก ปติฏฺฐาเปตฺวา อุฏฺฐายาสนา ปกฺกามิ.\csromanspacer{} อถ โข ธนญฺชานิ พฺราหฺมโณ อจิรปกฺกนฺเต อายสฺมนฺเต สาริปุตฺเต กาลมกาสิ, พฺรหฺมโลกญฺจ อุปปชฺชิ.}

\proseitem{453}{อถ โข ภควา ภิกฺขู อามนฺเตสิ “เอโส ภิกฺขเว สาริปุตฺโต ธนญฺชานิํ พฺราหฺมณํ สติ อุตฺตริกรณีเย หีเน พฺรหฺมโลเก ปติฏฺฐาเปตฺวา อุฏฺฐายาสนา ปกฺกนฺโต”ติ.\csromanspacer{} อถ โข อายสฺมา สาริปุตฺโต เยน ภควา เตนุปสงฺกมิ, อุปสงฺกมิตฺวา ภควนฺตํ อภิวาเทตฺวา เอกมนฺตํ นิสีทิ, เอกมนฺตํ นิสินฺโน โข อายสฺมา สาริปุตฺโต ภควนฺตํ เอตทโวจ “ธนญฺชานิ ภนฺเต พฺราหฺมโณ อาพาธิโก ทุกฺขิโต พาฬฺหาคิลาโน, โส ภควโต ปาเท สิรสา วนฺทตี”ติ.\csromanspacer{} กิํ ปน ตฺวํ สาริปุตฺต ธนญฺชานิํ พฺราหฺมณํ สติ อุตฺตริกรณีเย หีเน พฺรหฺมโลเก ปติฏฺฐาเปตฺวา อุฏฺฐายาสนา ปกฺกนฺโตติ.\csromanspacer{} มยฺหํ โข ภนฺเต เอวํ อโหสิ “อิเม โข พฺราหฺมณา พฺรหฺมโลกาธิมุตฺตา, ยํนูนาหํ ธนญฺชานิสฺส พฺราหฺมณสฺส พฺรหฺมานํ สหพฺยตาย มคฺคํ เทเสยฺยนฺ”ติ.\csromanspacer{} กาลงฺกโตจ\footnote{กาลงฺกโตว (สฺยา, กํ, ก)} สาริปุตฺต ธนญฺชานิ พฺราหฺมโณ, พฺรหฺมโลกญฺจ อุปปนฺโนติ.}

\nitthitamruled{ธนญฺชานิสุตฺตํ นิฏฺฐิตํ สตฺตมํ.}
\csromanheader{2}{8. วาเสฏฺฐสุตฺต}
\proseitem{454}{เอวํ เม สุตํ\csromandash{}เอกํ สมยํ ภควา อิจฺฉานงฺคเล\footnote{อิจฺฉานงฺกเล (สี, อิ) 3. ชาณุสฺโสณี (อิ), ชาณุโสณี (ก)} วิหรติ อิจฺฉานงฺคลวนสณฺเฑ.\csromanspacer{} เตน โข ปน สมเยน สมฺพหุลา อภิญฺญาตา อภิญฺญาตา พฺราหฺมณมหาสาลา อิจฺฉานงฺคเล ปฏิวสนฺติ.\csromanspacer{} เสยฺยถิทํ, จงฺกี พฺราหฺมโณ ตารุกฺโข พฺราหฺมโณ โปกฺขรสาติ พฺราหฺมโณ ชาณุสฺโสณิ3 พฺราหฺมโณ โตเทยฺโย พฺราหฺมโณ อญฺเญ จ อภิญฺญาตา อภิญฺญาตา พฺราหฺมณมหาสาลา.\csromanspacer{} อถ โข วาเสฏฺฐภารทฺวาชานํ มาณวานํ ชงฺฆาวิหารํ อนุจงฺกมนฺตานํ อนุวิจรนฺตานํ\footnote{อนุจงฺกมมานานํ อนุวิจรมานานํ (สี, อิ)}}

\vspace{\tipitakaparskip}
\csromancenter{วาเสฏฺฐสุตฺต (98) อยมนฺตรากถา อุทปาทิ “กถํ โภ พฺราหฺมโณ โหตี”ติ.\csromanspacer{} ภารทฺวาโช มาณโว เอวมาห “ยโต โข โภ อุภโต สุชาโต มาติโต จ ปิติโต จ สํสุทฺธคหณิโก ยาว สตฺตมา ปิตามหยุคา อกฺขิตฺโต อนุปกฺกุฏฺโฐ ชาติวาเทน, เอตฺตาวตา โข โภ พฺราหฺมโณ โหตี”ติ.\csromanspacer{} วาเสฏฺโฐ มาณโว เอวมาห “ยโต โข โภ สีลวา จ โหติ วตฺตสมฺปนฺโน\footnote{วตสมฺปนฺโน(อิ)} จ, เอตฺตาวตา โข โภ พฺราหฺมโณ โหตี”ติ.\csromanspacer{} เนว โข อสกฺขิ ภารทฺวาโช มาณโว วาเสฏฺฐํ มาณวํ สญฺญาเปตุํ, น ปน อสกฺขิ วาเสฏฺโฐ มาณโว ภารทฺวาชํ มาณวํ สญฺญาเปตุํ.\csromanspacer{} อถ โข วาเสฏฺโฐ มาณโว ภารทฺวาชํ มาณวํ อามนฺเตสิ “อยํ โข โภ ภารทฺวาช สมโณ โคตโม สกฺยปุตฺโต สกฺยกุลา ปพฺพชิโต อิจฺฉานงฺคเล วิหรติ อิจฺฉานงฺคลวนสณฺเฑ, ตํ โข ปน ภวนฺตํ โคตมํ เอวํ กลฺยาโณ กิตฺติสทฺโท อพฺภุคฺคโต ‘อิติปิ โส ภควา อรหํ สมฺมาสมฺพุทฺโธ วิชฺชาจรณสมฺปนฺโน สุคโต โลกวิทู อนุตฺตโร ปุริสทมฺมสารถิ สตฺถา เทวมนุสฺสานํ พุทฺโธ ภควา’ติ, อายาม โภ ภารทฺวาช, เยน สมโณ โคตโม เตนุปสงฺกมิสฺสาม, อุปสงฺกมิตฺวา สมณํ โคตมํ เอตมตฺถํ ปุจฺฉิสฺสาม, ยถา โน สมโณ โคตโม พฺยากริสฺสติ, ตถา นํ ธาเรสฺสามา”ติ. “เอวํ โภ”ติ โข ภารทฺวาโช มาณโว วาเสฏฺฐสฺส มาณวสฺส ปจฺจสฺโสสิ.}
\proseitem{455}{\csromanfolio{407}อถ โข วาเสฏฺฐภารทฺวาชา มาณวา เยน ภควา เตนุปสงฺกมิํสุ, อุปสงฺกมิตฺวา ภควตา สทฺธิํ สมฺโมทิํสุ, สมฺโมทนียํ กถํ สารณียํ วีติสาเรตฺวา เอกมนฺตํ นิสีทิํสุ, เอกมนฺตํ นิสินฺโน โข วาเสฏฺโฐ มาณโว ภควนฺตํ คาถาหิ อชฺฌภาสิ\csromandash{}}

\csromangathameasure{“อนุญฺญาตปฏิญฺญาตา, เตวิชฺชา มยมสฺมุโต. \\ อหํ โปกฺขรสาติสฺส, ตารุกฺขสฺสายํ มาณโว. \\ เตวิชฺชานํ ยทกฺขาตํ, ตตฺร เกวลิโนสฺมเส. \\ ปทกสฺมา เวยฺยากรณา, ชปฺเป อาจริยสาทิสา. \\ เตสํ โน ชาติวาทสฺมิํ, วิวาโท อตฺถิ โคตม. \\ ชาติยา พฺราหฺมโณ โหติ, ภารทฺวาโช อิติ ภาสติ. \\ อหญฺจ กมฺมุนา พฺรูมิ, เอวํ ชานาหิ จกฺขุม. \\ เต น สกฺโกม ญาเปตุํ, อญฺญมญฺญํ มยํ อุโภ. \\ ภวนฺตํ ปุฏฺฐุมาคมา, สมฺพุทฺธํ อิติ วิสฺสุตํ. \\ จนฺทํ ยถา ขยาตีตํ, เปจฺจ ปญฺชลิกา ชนา. \\ วนฺทมานา นมสฺสนฺติ, เอวํ โลกสฺมิํ โคตมํ. \\ จกฺขุํ โลเก สมุปฺปนฺนํ, มยํ ปุจฺฉาม โคตมํ. \\ ชาติยา พฺราหฺมโณ โหติ, อุทาหุ ภวติ กมฺมุนา1. \\ อชานตํ โน ปพฺรูหิ, ยถา ชาเนมุ พฺราหฺมณนฺติ.}
\csromangathagroup{\csromangathastanza{“อนุญฺญาตปฏิญฺญาตา, เตวิชฺชา มยมสฺมุโต. \\ อหํ โปกฺขรสาติสฺส, ตารุกฺขสฺสายํ มาณโว.}\csromangathastanza{เตวิชฺชานํ ยทกฺขาตํ, ตตฺร เกวลิโนสฺมเส. \\ ปทกสฺมา เวยฺยากรณา\footnote{โน พฺยากรณา (สฺยา, กํ, ก)}, ชปฺเป อาจริยสาทิสา.}\csromangathastanza{เตสํ โน ชาติวาทสฺมิํ, วิวาโท อตฺถิ โคตม. \\ ชาติยา พฺราหฺมโณ โหติ, ภารทฺวาโช อิติ ภาสติ.}\csromangathastanza{\csromanfolio{408}อหญฺจ กมฺมุนา\footnote{กมฺมนา (สี, อิ)} พฺรูมิ, เอวํ ชานาหิ จกฺขุม. \\ เต น สกฺโกม ญาเปตุํ\footnote{กมฺมนา (สี, อิ)}, อญฺญมญฺญํ มยํ อุโภ.}\csromangathastanza{ภวนฺตํ ปุฏฺฐุมาคมา, สมฺพุทฺธํ อิติ วิสฺสุตํ. \\ จนฺทํ ยถา ขยาตีตํ, เปจฺจ ปญฺชลิกา ชนา.}\csromangathastanza{วนฺทมานา นมสฺสนฺติ, เอวํ โลกสฺมิํ โคตมํ. \\ จกฺขุํ โลเก สมุปฺปนฺนํ, มยํ ปุจฺฉาม โคตมํ.}\csromangathastanza{ชาติยา พฺราหฺมโณ โหติ, อุทาหุ ภวติ กมฺมุนา1. \\ อชานตํ โน ปพฺรูหิ, ยถา ชาเนมุ พฺราหฺมณนฺติ.}}

\csromanhangingbegin
\hangingheaditem{456}{เตสํ โว อหํ พฺยกฺขิสฺสํ, (วาเสฏฺฐาติ ภควา)}
\hangingbody{อนุปุพฺพํ ยถาตถํ.}
\hangingbody{ชาติวิภงฺคํ ปาณานํ, อญฺญมญฺญาหิ ชาติโย.}
\hangingbody{ติณรุกฺเขปิ ชานาถ, น จาปิ ปฏิชานเร.}
\hangingbody{ลิงฺคํ ชาติมยํ เตสํ, อญฺญมญฺญา หิ ชาติโย.}
\hangingbody{ตโต กีเฏ ปฏงฺเค จ, ยาว กุนฺถกิปิลฺลิเก.}
\hangingbody{ลิงฺคํ ชาติมยํ เตสํ, อญฺญมญฺญา หิ ชาติโย.}
\hangingbody{จตุปฺปเทปิ ชานาถ, ขุทฺทเก จ มหลฺลเก.}
\hangingbody{ลิงฺคํ ชาติมยํ เตสํ, อญฺญมญฺญา หิ ชาติโย.}
\hangingbody{ปาทุทเรปิ ชานาถ, อุรเค ทีฆปิฏฺฐิเก.}
\hangingbody{ลิงฺคํ ชาติมยํ เตสํ, อญฺญมญฺญา หิ ชาติโย.}
\hangingbody{ตโต มจฺเฉปิ ชานาถ, อุทเก วาริโคจเร.}
\hangingbody{ลิงฺคํ ชาติมยํ เตสํ, อญฺญมญฺญา หิ ชาติโย.}
\hangingbody{ตโต ปกฺขีปิ ชานาถ, ปตฺตยาเน วิหงฺคเม.}
\hangingbody{ลิงฺคํ ชาติมยํ เตสํ, อญฺญมญฺญา หิ ชาติโย.}
\hangingbody{ยถา เอตาสุ ชาตีสุ, ลิงฺคํ ชาติมยํ ปุถุ.}
\hangingbody{เอวํ นตฺถิ มนุสฺเสสุ, ลิงฺคํ ชาติมยํ ปุถุ.}
\csromanhangingend

\csromanheader{2}{8. วาเสฏฺฐสุตฺต (98)}
\csromangathameasure{น เกเสหิ น สีเสหิ, น กณฺเณหิ น อกฺขีหิ. \\ น มุเขน น นาสาย, น โอฏฺเฐหิ ภมูหิ วา. \\ น คีวาย น อํเสหิ, น อุทเรน น ปิฏฺฐิยา. \\ น โสณิยา น อุรสา, น สมฺพาเธ น เมถุเน. \\ น หตฺเถหิ น ปาเทหิ, นงฺคุลีหิ นเขหิ วา. \\ น ชงฺฆาหิ น อูรูหิ, น วณฺเณน สเรน วา. \\ ลิงฺคํ ชาติมยํ เนว, ยถา อญฺญาสุ ชาติสุ.}
\csromangathagroup{\csromangathastanza{\csromanfolio{409}น เกเสหิ น สีเสหิ, น กณฺเณหิ น อกฺขีหิ. \\ น มุเขน น นาสาย, น โอฏฺเฐหิ ภมูหิ วา.}\csromangathastanza{น คีวาย น อํเสหิ, น อุทเรน น ปิฏฺฐิยา. \\ น โสณิยา น อุรสา, น สมฺพาเธ น เมถุเน\footnote{น สมฺพาธา น เมถุนา (ก)}.}\csromangathastanza{น หตฺเถหิ น ปาเทหิ, นงฺคุลีหิ นเขหิ วา. \\ น ชงฺฆาหิ น อูรูหิ, น วณฺเณน สเรน วา. \\ ลิงฺคํ ชาติมยํ เนว, ยถา อญฺญาสุ ชาติสุ.}}

\csromanhangingbegin
\hangingheaditem{457}{ปจฺจตฺตญฺจ สรีเรสุ\footnote{ปจฺจตฺตํ สสรีเรสุ (สี, อิ)}, มนุสฺเสเสฺวตํ น วิชฺชติ.}
\hangingbody{โวการญฺจ มนุสฺเสสุ, สมญฺญาย ปวุจฺจติ.}
\hangingbody{โย หิ โกจิ มนุสฺเสสุ, โครกฺขํ อุปชีวติ.}
\hangingbody{เอวํ วาเสฏฺฐ ชานาหิ, กสฺสโก โส น พฺราหฺมโณ.}
\hangingbody{โย หิ โกจิ มนุสฺเสสุ, ปุถุสิปฺเปน ชีวติ.}
\hangingbody{เอวํ วาเสฏฺฐ ชานาหิ, สิปฺปิโก โส น พฺราหฺมโณ.}
\hangingbody{โย หิ โกจิ มนุสฺเสสุ, โวหารํ อุปชีวติ.}
\hangingbody{เอวํ วาเสฏฺฐ ชานาหิ, วาณิโช โส น พฺราหฺมโณ.}
\hangingbody{โย หิ โกจิ มนุสฺเสสุ, ปรเปสฺเสน ชีวติ.}
\hangingbody{เอวํ วาเสฏฺฐ ชานาหิ, เปสฺสโก3 โส น พฺราหฺมโณ.}
\hangingbody{โย หิ โกจิ มนุสฺเสสุ, อทินฺนํ อุปชีวติ.}
\hangingbody{เอวํ วาเสฏฺฐ ชานาหิ, โจโร เอโส น พฺราหฺมโณ.}
\hangingbody{โย หิ โกจิ มนุสฺเสสุ, อิสฺสตฺถํ อุปชีวติ.}
\hangingbody{เอวํ วาเสฏฺฐ ชานาหิ, โยธาชีโว น พฺราหฺมโณ.}
\hangingbody{โย หิ โกจิ มนุสฺเสสุ, โปโรหิจฺเจน ชีวติ.}
\hangingbody{เอวํ วาเสฏฺฐ ชานาหิ, ยาชโก โส น พฺราหฺมโณ.}
\hangingbody{โย หิ โกจิ มนุสฺเสสุ, คามํ รฏฺฐญฺจ ภุญฺชติ.}
\hangingbody{เอวํ วาเสฏฺฐ ชานาหิ, ราชา เอโส น พฺราหฺมโณ.}
\csromanhangingend

\csromangathameasure{น จาหํ พฺราหฺมณํ พฺรูมิ, โยนิชํ มตฺติสมฺภวํ. \\ โภวาทิ นาม โส โหติ, สเจ โหติ สกิญฺจโน. \\ อกิญฺจนํ อนาทานํ, ตมหํ พฺรูมิ พฺราหฺมณํ.}
\csromangathagroup{\csromangathastanza{\csromanfolio{410}น จาหํ พฺราหฺมณํ พฺรูมิ, โยนิชํ มตฺติสมฺภวํ. \\ โภวาทิ\footnote{โภวาที (สฺยา, กํ)} นาม โส โหติ, สเจ โหติ สกิญฺจโน. \\ อกิญฺจนํ อนาทานํ, ตมหํ พฺรูมิ พฺราหฺมณํ.}}

\csromanmatikaanchor{mtk.57}
\csromantocmarkref{subsection}{8. วาเสฏฺฐสุตฺต}{mtk.57}
\bookmark[dest={mtk.57},level=2]{8. วาเสฏฺฐสุตฺต}
\csromanhangingbegin
\hangingheaditem{458}{สพฺพสํโยชนํ เฉตฺวา, โย เว น ปริตสฺสติ.}
\hangingbody{สงฺคาติคํ วิสํยุตฺตํ2, ตมหํ พฺรูมิ พฺราหฺมณํ.}
\hangingbody{เฉตฺวา นทฺธิํ3 วรตฺตญฺจ, สนฺทานํ สหนุกฺกมํ.}
\hangingbody{อุกฺขิตฺตปลิฆํ พุทฺธํ, ตมหํ พฺรูมิ พฺราหฺมณํ.}
\hangingbody{อกฺโกสํ วธพนฺธญฺจ, อทุฏฺโฐ โย ติติกฺขติ.}
\hangingbody{ขนฺตีพลํ พลานีกํ, ตมหํ พฺรูมิ พฺราหฺมณํ.}
\hangingbody{อกฺโกธนํ วตวนฺตํ, สีลวนฺตํ อนุสฺสทํ.}
\hangingbody{ทนฺตํ อนฺติมสารีรํ, ตมหํ พฺรูมิ พฺราหฺมณํ.}
\hangingbody{วาริโปกฺขรปตฺเตว, อารคฺเคริว สาสโป.}
\hangingbody{โย น ลิมฺปติ กาเมสุ, ตมหํ พฺรูมิ พฺราหฺมณํ.}
\hangingbody{โย ทุกฺขสฺส ปชานาติ, อิเธว ขยมตฺตโน.}
\hangingbody{ปนฺนภารํ วิสํยุตฺตํ, ตมหํ พฺรูมิ พฺราหฺมณํ.}
\hangingbody{คมฺภีรปญฺญํ เมธาวิํ, มคฺคามคฺคสฺส โกวิทํ.}
\hangingbody{อุตฺตมตฺถมนุปฺปตฺตํ, ตมหํ พฺรูมิ พฺราหฺมณํ.}
\hangingbody{อสํสฏฺฐํ คหฏฺเฐหิ, อนาคาเรหิ จูภยํ.}
\hangingbody{อโนกสาริมปฺปิจฺฉํ, ตมหํ พฺรูมิ พฺราหฺมณํ.}
\hangingbody{นิธาย ทณฺฑํ ภูเตสุ, ตเสสุ ถาวเรสุ จ.}
\hangingbody{โย น หนฺติ น ฆาเตติ, ตมหํ พฺรูมิ พฺราหฺมณํ.}
\hangingbody{อวิรุทฺธํ วิรุทฺเธสุ, อตฺถทณฺเฑสุ นิพฺพุตํ.}
\hangingbody{สาทาเนสุ อนาทานํ, ตมหํ พฺรูมิ พฺราหฺมณํ.}
\hangingbody{ยสฺส ราโค จ โทโส จ, มาโน มกฺโข จ โอหิโต.}
\hangingbody{สาสโปริว อารคฺคา, ตมหํ พฺรูมิ พฺราหฺมณํ.}
\csromanhangingend

\csromanheader{2}{8. วาเสฏฺฐสุตฺต (98)}
\csromanhangingbegin
\hangingheaditem{459}{\csromanfolio{411}อกกฺกสํ วิญฺญาปนิํ, คิรํ สจฺจํ อุทีรเย.}
\hangingbody{ยาย นาภิสชฺเช กิญฺจิ, ตมหํ พฺรูมิ พฺราหฺมณํ.}
\hangingbody{โย จ ทีฆํ ว รสฺสํ วา, อณุํ ถูลํ สุภาสุภํ.}
\hangingbody{โลเก อทินฺนํ นาเทติ1, ตมหํ พฺรูมิ พฺราหฺมณํ.}
\hangingbody{อาสา ยสฺส น วิชฺชนฺติ, อสฺมิํ โลเก ปรมฺหิ จ.}
\hangingbody{นิราสาสํ2 วิสํยุตฺตํ, ตมหํ พฺรูมิ พฺราหฺมณํ.}
\hangingbody{ยสฺสาลยา น วิชฺชนฺติ, อญฺญาย อกถํกถิํ.}
\hangingbody{อมโตคธํ อนุปฺปตฺตํ, ตมหํ พฺรูมิ พฺราหฺมณํ.}
\hangingbody{โยธปุญฺญญฺจ ปาปญฺจ, อุโภ สงฺคํ อุปจฺจคา.}
\hangingbody{อโสกํ วิรชํ สุทฺธํ, ตมหํ พฺรูมิ พฺราหฺมณํ.}
\hangingbody{จนฺทํว วิมลํ สุทฺธํ, วิปฺปสนฺนํ อนาวิลํ.}
\hangingbody{นนฺทีภวปริกฺขีณํ, ตมหํ พฺรูมิ พฺราหฺมณํ.}
\hangingbody{โย อิมํ ปลิปถํ ทุคฺคํ, สํสารํ โมหมจฺจคา.}
\hangingbody{ติณฺโณ ปารงฺคโต ฌายี, อเนโช อกถํกถี.}
\hangingbody{อนุปาทาย นิพฺพุโต, ตมหํ พฺรูมิ พฺราหฺมณํ.}
\hangingbody{โยธกาเม ปหนฺตฺวาน3, อนาคาโร ปริพฺพเช.}
\hangingbody{กามภวปริกฺขีณํ, ตมหํ พฺรูมิ พฺราหฺมณํ.}
\hangingbody{โยธตณฺหํ ปหนฺตฺวาน, อนาคาโร ปริพฺพเช.}
\hangingbody{ตณฺหาภวปริกฺขีณํ, ตมหํ พฺรูมิ พฺราหฺมณํ.}
\hangingbody{หิตฺวา มานุสกํ โยคํ, ทิพฺพํ โยคํ อุปจฺจคา.}
\hangingbody{สพฺพโยควิสํยุตฺตํ, ตมหํ พฺรูมิ พฺราหฺมณํ.}
\hangingbody{หิตฺวา รติํ จ อรติํ, สีตีภูตํ นิรูปธิํ.}
\hangingbody{สพฺพโลกาภิภุํ วีรํ, ตมหํ พฺรูมิ พฺราหฺมณํ.}
\hangingbody{จุติํ โย เวทิ สตฺตานํ, อุปปตฺติํ จ สพฺพโส.}
\hangingbody{อสตฺตํ สุคตํ พุทฺธํ, ตมหํ พฺรูมิ พฺราหฺมณํ.}
\hangingbody{ยสฺส คติํ น ชานนฺติ, เทวา คนฺธพฺพมานุสา.}
\hangingbody{ขีณาสวํ อรหนฺตํ, ตมหํ พฺรูมิ พฺราหฺมณํ.}
\csromanhangingend

\csromangathameasure{ยสฺส ปุเร จ ปจฺฉา จ, มชฺเฌ จ นตฺถิ กิญฺจนํ. \\ อกิญฺจนํ อนาทานํ, ตมหํ พฺรูมิ พฺราหฺมณํ. \\ อุสภํ ปวรํ วีรํ, มเหสิํ วิชิตาวินํ. \\ อเนชํ นฺหาตกํ พุทฺธํ, ตมหํ พฺรูมิ พฺราหฺมณํ. \\ ปุพฺเพนิวาสํ โย เวทิ, สคฺคาปายญฺจ ปสฺสติ. \\ อโถ ชาติกฺขยํ ปตฺโต, ตมหํ พฺรูมิ พฺราหฺมณํ.}
\csromangathagroup{\csromangathastanza{\csromanfolio{412}ยสฺส ปุเร จ ปจฺฉา จ, มชฺเฌ จ นตฺถิ กิญฺจนํ. \\ อกิญฺจนํ อนาทานํ, ตมหํ พฺรูมิ พฺราหฺมณํ.}\csromangathastanza{อุสภํ ปวรํ วีรํ, มเหสิํ วิชิตาวินํ. \\ อเนชํ นฺหาตกํ\footnote{นหาตกํ (สี, อิ)} พุทฺธํ, ตมหํ พฺรูมิ พฺราหฺมณํ.}\csromangathastanza{ปุพฺเพนิวาสํ โย เวทิ, สคฺคาปายญฺจ ปสฺสติ. \\ อโถ ชาติกฺขยํ ปตฺโต, ตมหํ พฺรูมิ พฺราหฺมณํ.}}

\csromanhangingbegin
\hangingheaditem{460}{สมญฺญา เหสา โลกสฺมิํ, นามโคตฺตํ ปกปฺปิตํ.}
\hangingbody{สมฺมุจฺจา สมุทาคตํ, ตตฺถ ตตฺถ ปกปฺปิตํ.}
\hangingbody{ทีฆรตฺตานุสยิตํ, ทิฏฺฐิคตมชานตํ.}
\hangingbody{อชานนฺตา โน2 ปพฺรุนฺติ3, ชาติยา โหติ พฺราหฺมโณ.}
\hangingbody{น ชจฺจา พฺราหฺมโณ4 โหติ, น ชจฺจา โหติ อพฺราหฺมโณ5.}
\hangingbody{กมฺมุนา พฺราหฺมโณ4 โหติ, กมฺมุนา โหติ อพฺราหฺมโณ5.}
\hangingbody{กสฺสโก กมฺมุนา โหติ, สิปฺปิโก โหติ กมฺมุนา.}
\hangingbody{วาณิโช กมฺมุนา โหติ, เปสฺสโก โหติ กมฺมุนา.}
\hangingbody{โจโรปิ กมฺมุนา โหติ, โยธาชีโวปิ กมฺมุนา.}
\hangingbody{ยาชโก กมฺมุนา โหติ, ราชาปิ โหติ กมฺมุนา.}
\hangingbody{เอวเมตํ ยถาภูตํ, กมฺมํ ปสฺสนฺติ ปณฺฑิตา.}
\hangingbody{ปฏิจฺจสมุปฺปาททสฺสา, กมฺมวิปากโกวิทา.}
\hangingbody{กมฺมุนา วตฺตติ โลโก, กมฺมุนา วตฺตติ ปชา.}
\hangingbody{กมฺมนิพนฺธนา สตฺตา, รถสฺสาณีว ยายโต.}
\hangingbody{ตเปน พฺรหฺมจริเยน, สํยเมน ทเมน จ.}
\hangingbody{เอเตน พฺราหฺมโณ โหติ, เอตํ พฺราหฺมณมุตฺตมํ.}
\hangingbody{ตีหิ วิชฺชาหิ สมฺปนฺโน, สนฺโต ขีณปุนพฺภโว.}
\hangingbody{เอวํ วาเสฏฺฐ ชานาหิ, พฺรหฺมา สกฺโก วิชานตนฺ”ติ.}
\csromanhangingend

\csromanheader{2}{9. สุภสุตฺต (99)}
\proseitem{461}{\csromanfolio{413}เอวํ วุตฺเต วาเสฏฺฐภารทฺวาชา มาณวา ภควนฺตํ เอตทโวจุํ “อภิกฺกนฺตํ โภ โคตม, อภิกฺกนฺตํ โภ โคตม, เสยฺยถาปิ โภ โคตม นิกฺกุชฺชิตํ วา อุกฺกุชฺเชยฺย, ปฏิจฺฉนฺนํ วา วิวเรยฺย, มูฬฺหสฺส วา มคฺคํ อาจิกฺเขยฺย, อนฺธกาเร วา เตลปชฺโชตํ ธาเรยฺย ‘จกฺขุมนฺโต รูปานิ ทกฺขนฺตี’ติ, เอวเมวํ โภตา โคตเมน อเนกปริยาเยน ธมฺโม ปกาสิโต, เอเต มยํ ภวนฺตํ โคตมํ สรณํ คจฺฉาม, ธมฺมญฺจ ภิกฺขุสํฆญฺจ, อุปาสเก โน ภควํ โคตโม ธาเรตุ อชฺชตคฺเค ปาณุเปตํ สรณํ คเต”ติ.}

\nitthitamruled{วาเสฏฺฐสุตฺตํ นิฏฺฐิตํ อฏฺฐมํ.}
\csromanheader{2}{9. สุภสุตฺต}
\proseitem{462}{เอวํ เม สุตํ\csromandash{}เอกํ สมยํ ภควา สาวตฺถิยํ วิหรติ เชตวเน อนาถปิณฺฑิกสฺส อาราเม.\csromanspacer{} เตน โข ปน สมเยน สุโภ มาณโว โตเทยฺยปุตฺโต สาวตฺถิยํ ปฏิวสติ อญฺญตรสฺส คหปติสฺส นิเวสเน เกนจิเทว กรณีเยน.\csromanspacer{} อถ โข สุโภ มาณโว โตเทยฺยปุตฺโต ยสฺส คหปติสฺส นิเวสเน ปฏิวสติ, ตํ คหปติํ เอตทโวจ “สุตํ เมตํ คหปติ ‘อวิวิตฺตา สาวตฺถี อรหนฺเตหี’ติ.\csromanspacer{} กํ นุ ขฺวชฺช สมณํ วา พฺราหฺมณํ วา ปยิรุปาเสยฺยามา”ติ.\csromanspacer{} อยํ ภนฺเต ภควา สาวตฺถิยํ วิหรติ เชตวเน อนาถปิณฺฑิกสฺส อาราเม, ตํ ภนฺเต ภควนฺตํ ปยิรุปาสสฺสูติ.\csromanspacer{} อถ โข สุโภ มาณโว โตเทยฺยปุตฺโต ตสฺส คหปติสฺส ปฏิสฺสุตฺวา เยน ภควา เตนุปสงฺกมิ, อุปสงฺกมิตฺวา ภควตา สทฺธิํ สมฺโมทิ, สมฺโมทนียํ กถํ สารณียํ วีติสาเรตฺวา เอกมนฺตํ นิสีทิ, เอกมนฺตํ นิสินฺโน โข สุโภ มาณโว โตเทยฺยปุตฺโต ภควนฺตํ เอตทโวจ “พฺราหฺมณา โภ โคตม เอวมาหํสุ ‘คหฏฺโฐ อาราธโก โหติ ญายํ ธมฺมํ กุสลํ, น ปพฺพชิโต อาราธโก โหติ ญายํ ธมฺมํ กุสลนฺ’ติ.\csromanspacer{} อิธ ภวํ โคตโม กิมาหา”ติ.}

\proseitem{463}{วิภชฺชวาโท โข อหเมตฺถ มาณว, นาหเมตฺถ เอกํสวาโท.\csromanspacer{} คิหิสฺส วาหํ มาณว ปพฺพชิตสฺส วา มิจฺฉาปฏิปตฺติํ น วณฺเณมิ.\csromanspacer{} คิหี \csromanfolio{414}วา หิ มาณว ปพฺพชิโต วา มิจฺฉาปฏิปนฺโน มิจฺฉาปฏิปตฺตาธิกรณเหตุ น อาราธโก โหติ ญายํ ธมฺมํ กุสลํ.\csromanspacer{} คิหิสฺส วาหํ มาณว ปพฺพชิตสฺส วา สมฺมาปฏิปตฺติํ วณฺเณมิ.\csromanspacer{} คิหี วา หิ มาณว ปพฺพชิโต วา สมฺมาปฏิปนฺโน สมฺมาปฏิปตฺตาธิกรณเหตุ อาราธโก โหติ ญายํ ธมฺมํ กุสลนฺติ.}

\prose{พฺราหฺมณา โภ โคตม เอวมาหํสุ “มหฏฺฐมิทํ มหากิจฺจํ มหาธิกรณํ มหาสมารมฺภํ, ฆราวาสกมฺมฏฺฐานํ มหปฺผลํ โหติ.\csromanspacer{} อปฺปฏฺฐมิทํ อปฺปกิจฺจํ อปฺปาธิกรณํ อปฺปสมารมฺภํ, ปพฺพชฺชา กมฺมฏฺฐานํ อปฺปผลํ โหตี”ติ.\csromanspacer{} อิธ ภวํ โคตโม กิมาหาติ.}

\prose{เอตฺถาปิ โข อหํ มาณว วิภชฺชวาโท, นาหเมตฺถ เอกํสวาโท.\csromanspacer{} อตฺถิ มาณว กมฺมฏฺฐานํ มหฏฺฐํ มหากิจฺจํ มหาธิกรณํ มหาสมารมฺภํ, วิปชฺชมานํ อปฺปผลํ โหติ.\csromanspacer{} อตฺถิ มาณว กมฺมฏฺฐานํ มหฏฺฐํ มหากิจฺจํ มหาธิกรณํ มหาสมารมฺภํ, สมฺปชฺชมานํ มหปฺผลํ โหติ.\csromanspacer{} อตฺถิ มาณว กมฺมฏฺฐานํ อปฺปฏฺฐํ อปฺปกิจฺจํ อปฺปาธิกรณํ อปฺปสมารมฺภํ, วิปชฺชมานํ อปฺปผลํ โหติ.\csromanspacer{} อตฺถิ มาณว กมฺมฏฺฐานํ อปฺปฏฺฐํ อปฺปกิจฺจํ อปฺปาธิกรณํ อปฺปสมารมฺภํ, สมฺปชฺชมานํ มหปฺผลํ โหติ.\csromanspacer{} กตมญฺจ มาณว กมฺมฏฺฐานํ มหฏฺฐํ มหากิจฺจํ มหาธิกรณํ มหาสมารมฺภํ, วิปชฺชมานํ อปฺปผลํ โหติ.\csromanspacer{} กสิ โข มาณว กมฺมฏฺฐานํ มหฏฺฐํ มหากิจฺจํ มหาธิกรณํ มหาสมารมฺภํ, วิปชฺชมานํ อปฺปผลํ โหติ.\csromanspacer{} กตมญฺจ มาณว กมฺมฏฺฐานํ มหฏฺฐํ มหากิจฺจํ มหาธิกรณํ มหาสมารมฺภํ, สมฺปชฺชมานํ มหปฺผลํ โหติ.\csromanspacer{} กสิเยว โข มาณว กมฺมฏฺฐานํ มหฏฺฐํ มหากิจฺจํ มหาธิกรณํ มหาสมารมฺภํ, สมฺปชฺชมานํ มหปฺผลํ โหติ.\csromanspacer{} กตมญฺจ มาณว กมฺมฏฺฐานํ อปฺปฏฺฐํ อปฺปกิจฺจํ อปฺปาธิกรณํ อปฺปสมารมฺภํ, วิปชฺชมานํ อปฺปผลํ โหติ.\csromanspacer{} วณิชฺชา โข มาณว กมฺมฏฺฐานํ อปฺปฏฺฐํ อปฺปกิจฺจํ อปฺปาธิกรณํ อปฺปสมารมฺภํ, วิปชฺชมานํ อปฺปผลํ โหติ.\csromanspacer{} กตมญฺจ มาณว กมฺมฏฺฐานํ อปฺปฏฺฐํ อปฺปกิจฺจํ อปฺปาธิกรณํ อปฺปสมารมฺภํ, สมฺปชฺชมานํ มหปฺผลํ โหติ.\csromanspacer{} วณิชฺชาเยว โข มาณว กมฺมฏฺฐานํ อปฺปฏฺฐํ อปฺปกิจฺจํ อปฺปาธิกรณํ อปฺปสมารมฺภํ, สมฺปชฺชามานํ มหปฺผลํ โหติ.}

\proseitem{464}{เสยฺยถาปิ มาณว กสิ กมฺมฏฺฐานํ มหฏฺฐํ มหากิจฺจํ มหาธิกรณํ มหาสมารมฺภํ, วิปชฺชมานํ อปฺปผลํ โหติ.\csromanspacer{} เอวเมว โข มาณว ฆราวาสกมฺมฏฺฐานํ มหฏฺฐํ มหากิจฺจํ มหาธิกรณํ มหาสมารมฺภํ,}

\vspace{\tipitakaparskip}
\csromancenter{สุภสุตฺต (99) วิปชฺชมานํ อปฺปผลํ โหติ.\csromanspacer{} เสยฺยถาปิ มาณว กสิเยว กมฺมฏฺฐานํ มหฏฺฐํ มหากิจฺจํ มหาธิกรณํ มหาสมารมฺภํ, สมฺปชฺชมานํ มหปฺผลํ โหติ.\csromanspacer{} เอวเมว โข มาณว ฆราวาสกมฺมฏฺฐานํ มหฏฺฐํ มหากิจฺจํ มหาธิกรณํ มหาสมารมฺภํ, สมฺปชฺชมานํ มหปฺผลํ โหติ.\csromanspacer{} เสยฺยถาปิ มาณว วณิชฺชา กมฺมฏฺฐานํ อปฺปฏฺฐํ อปฺปกิจฺจํ อปฺปาธิกรณํ อปฺปสมารมฺภํ, วิปชฺชมานํ อปฺปผลํ โหติ.\csromanspacer{} เอวเมว โข มาณว ปพฺพชฺชา กมฺมฏฺฐานํ อปฺปฏฺฐํ อปฺปกิจฺจํ อปฺปาธิกรณํ อปฺปสมารมฺภํ, วิปชฺชมานํ อปฺปผลํ โหติ.\csromanspacer{} เสยฺยถาปิ มาณว วณิชฺชา เยว กมฺมฏฺฐานํ อปฺปฏฺฐํ อปฺปกิจฺจํ อปฺปาธิกรณํ อปฺปสมารมฺภํ, สมฺปชฺชมานํ มหปฺผลํ โหติ.\csromanspacer{} เอวเมว โข มาณว ปพฺพชฺชา กมฺมฏฺฐานํ อปฺปฏฺฐํ อปฺปกิจฺจํ อปฺปาธิกรณํ อปฺปสมารมฺภํ, สมฺปชฺชมานํ มหปฺผลํ โหตีติ.}
\prose{\csromanfolio{415}พฺราหฺมณา โภ โคตม ปญฺจ ธมฺเม ปญฺญเปนฺติ ปุญฺญสฺส กิริยาย กุสลสฺส อาราธนายาติ.\csromanspacer{} เย เต มาณว พฺราหฺมณา ปญฺจ ธมฺเม ปญฺญเปนฺติ ปุญฺญสฺส กิริยาย กุสลสฺส อาราธนาย.\csromanspacer{} สเจ เต อครุ, สาธุ เต ปญฺจ ธมฺเม อิมสฺมิํ ปริสติ ภาสสฺสูติ.\csromanspacer{} น โข เม โภ โคตม ครุ, ยตฺถสฺสุ ภวนฺโต วา นิสินฺโน ภวนฺตรูโป วาติ\footnote{นิสินฺนา ภวนฺตรูปา วาติ (สี, สฺยา, กํ, อิ)}.\csromanspacer{} เตน หิ มาณว ภาสสฺสูติ.\csromanspacer{} สจฺจํ โข โภ โคตม พฺราหฺมณา ปฐมํ ธมฺมํ ปญฺญเปนฺติ ปุญฺญสฺส กิริยาย กุสลสฺส อาราธนาย.\csromanspacer{} ตปํ โข โภ โคตม พฺราหฺมณา ทุติยํ ธมฺมํ ปญฺญเปนฺติ ปุญฺญสฺส กิริยาย กุสลสฺส อาราธนาย.\csromanspacer{} พฺรหฺมจริยํ โข โภ โคตม พฺราหฺมณา ตติยํ ธมฺมํ ปญฺญเปนฺติ ปุญฺญสฺส กิริยาย กุสลสฺส อาราธนาย.\csromanspacer{} อชฺเฌนํ โข โภ โคตม พฺราหฺมณา จตุตฺถํ ธมฺมํ ปญฺญเปนฺติ ปุญฺญสฺส กิริยาย กุสลสฺส อาราธนาย.\csromanspacer{} จาคํ โข โภ โคตม พฺราหฺมณา ปญฺจมํ ธมฺมํ ปญฺญเปนฺติ ปุญฺญสฺส กิริยาย กุสลสฺส อาราธนาย.\csromanspacer{} พฺราหฺมณา โภ โคตม อิเม ปญฺจ ธมฺเม ปญฺญเปนฺติ ปุญฺญสฺส กิริยาย กุสลสฺส อาราธนายาติ.\csromanspacer{} อิธ ภวํ โคตโม กิมาหาติ.}

\proseitem{465}{กิํ ปน มาณว อตฺถิ โกจิ พฺราหฺมณานํ เอกพฺราหฺมโณปิ, โย เอวมาห “อหํ อิเมสํ ปญฺจนฺนํ ธมฺมานํ สยํ อภิญฺญา สจฺฉิกตฺวา วิปากํ ปเวเทมี”ติ.\csromanspacer{} โน หิทํ โภ โคตม.\csromanspacer{} กิํ ปน มาณว อตฺถิ โกจิ พฺราหฺมณานํ เอกาจริโยปิ เอกาจริยปาจริโยปิ ยาว สตฺตมา อาจริยมหยุคาปิ โย เอวมาห “อหํ อิเมสํ ปญฺจนฺนํ ธมฺมานํ สยํ อภิญฺญา \csromanfolio{416}สจฺฉิกตฺวา วิปากํ ปเวเทมี”ติ.\csromanspacer{} โน หิทํ โภ โคตม.\csromanspacer{} กิํ ปน มาณว เยปิ เต พฺราหฺมณานํ ปุพฺพกา อิสโย มนฺตานํ กตฺตาโร มนฺตานํ ปวตฺตาโร, เยสมิทํ เอตรหิ พฺราหฺมณา โปราณํ มนฺตปทํ คีตํ ปวุตฺตํ สมิหิตํ, ตทนุคายนฺติ ตทนุภาสนฺติ ภาสิตมนุภาสนฺติ วาจิตมนุวาเจนฺติ.\csromanspacer{} เสยฺยถิทํ, อฏฺฐโก วามโก วามเทโว เวสฺสามิตฺโต ยมตคฺคิ องฺคีรโส ภารทฺวาโช วาเสฏฺโฐ กสฺสโป ภคุ.\csromanspacer{} เตปิ เอวมาหํสุ “มยํ อิเมสํ ปญฺจนฺนํ ธมฺมานํ สยํ อภิญฺญา สจฺฉิกตฺวา วิปากํ ปเวเทมา”ติ.\csromanspacer{} โน หิทํ โภ โคตม.}

\prose{อิติ กิร มาณว นตฺถิ โกจิ พฺราหฺมณานํ เอกพฺราหฺมโณปิ, โย เอวมาห “อหํ อิเมสํ ปญฺจนฺนํ ธมฺมานํ สยํ อภิญฺญา สจฺฉิกตฺวา วิปากํ ปเวเทมี”ติ.\csromanspacer{} นตฺถิ โกจิ พฺราหฺมณานํ เอกาจริโยปิ เอกาจริยปาจริโยปิ ยาว สตฺตมา อาจริยมหยุคาปิ โย เอวมาห “อหํ อิเมสํ ปญฺจนฺนํ ธมฺมานํ สยํ อภิญฺญา สจฺฉิกตฺวา วิปากํ ปเวเทมี”ติ.\csromanspacer{} เยปิ เต พฺราหฺมณานํ ปุพฺพกา อิสโย มนฺตานํ กตฺตาโร มนฺตานํ ปวตฺถาโร, เยสมิทํ เอตรหิ พฺราหฺมณา โปราณํ มนฺตปทํ คีตํ ปวุตฺตํ สมิหิตํ, ตทนุคายนฺติ ตทนุภาสนฺติ ภาสิตมนุภาสนฺติ วาจิตมนุวาเจนฺติ.\csromanspacer{} เสยฺยถิทํ, อฏฺฐโก วามโก วามเทโว เวสฺสามิตฺโต ยมตคฺคิ องฺคีรโส ภารทฺวาโช วาเสฏฺโฐ กสฺสโป ภคุ.\csromanspacer{} เตปิ น เอวมาหํสุ “มยํ อิเมสํ ปญฺจนฺนํ ธมฺมานํ สยํ อภิญฺญา สจฺฉิกตฺวา วิปากํ ปเวเทมา”ติ.}

\prose{เสยฺยถาปิ มาณว อนฺธเวณิ ปรมฺปราสํสตฺตา ปุริโมปิ น ปสฺสติ มชฺฌิโมปิ น ปสฺสติ, ปจฺฉิโมปิ น ปสฺสติ.\csromanspacer{} เอวเมว โข มาณว อนฺธเวณูปมํ มญฺเญ พฺราหฺมณานํ ภาสิตํ สมฺปชฺชติ, ปุริโมปิ น ปสฺสติ, มชฺฌิโมปิ น ปสฺสติ, ปจฺฉิโมปิ น ปสฺสตีติ.}

\proseitem{466}{เอวํ วุตฺเต สุโภ มาณโว โตเทยฺยปุตฺโต ภควตา อนฺธเวณูปเมน วุจฺจมาโน กุปิโต อนตฺตมโน ภควนฺตํเยว ขุํเสนฺโต ภควนฺตํเยว วมฺเภนฺโต ภควนฺตํเยว วทมาโน “สมโณ โคตโม ปาปิโต ภวิสฺสตี”ติ ภควนฺตํ เอตทโวจ “พฺราหฺมโณ โภ โคตม โปกฺขรสาติ โอปมญฺโญ สุภควนิโก เอวมาห}

\vspace{\tipitakaparskip}
\csromancenter{สุภสุตฺต (99) ‘เอวเมว ปนิเธกจฺเจ\footnote{ปนิเมเก (สพฺพตฺถ)} สมณพฺราหฺมณา อุตฺตริ มนุสฺสธมฺมา อลมริยญาณทสฺสนวิเสสํ ปฏิชานนฺติ, เตสมิทํ ภาสิตํ หสฺสกํเยว สมฺปชฺชติ, นามกํเยว สมฺปชฺชติ, ริตฺตกํเยว สมฺปชฺชติ, ตุจฺฉกํเยว สมฺปชฺชติ, กถํ หิ นาม มนุสฺสภูโต อุตฺตริ มนุสฺสธมฺมา อลมริยญาณทสฺสนวิเสสํ ญสฺสติ วา ทกฺขติ วา สจฺฉิ วา กริสฺสตี’ติ.\csromanspacer{} เนตํ ฐานํ วิชฺชตี”ติ.}
\prose{\csromanfolio{417}กิํ ปน มาณว พฺราหฺมโณ โปกฺขรสาติ โอปมญฺโญ สุภควนิโก สพฺเพสํเยว สมณพฺราหฺมณานํ เจตสา เจโต ปริจฺจ ปชานาตีติ.\csromanspacer{} สกายปิ หิ โภ โคตม ปุณฺณิกาย ทาสิยา พฺราหฺมโณ โปกฺขรสาติ โอปมญฺโญ สุภควนิโก เจตสา เจโต ปริจฺจ น ปชานาติ, กุโต ปน สพฺเพสํเยว สมณพฺราหฺมณานํ เจตสา เจโต ปริจฺจ ปชานิสฺสตีติ.}

\prose{เสยฺยถาปิ มาณว ชจฺจนฺโธ ปุริโส น ปสฺเสยฺย กณฺหสุกฺกานิ รูปานิ, น ปสฺเสยฺย นีลกานิ รูปานิ, น ปสฺเสยฺย ปีตกานิ รูปานิ, น ปสฺเสยฺย โลหิตกานิ รูปานิ, น ปสฺสยฺย มญฺชิฏฺฐกานิ รูปานิ, น ปสฺเสยฺย สมวิสมํ, น ปสฺเสยฺย ตารกรูปานิ, น ปสฺเสยฺย จนฺทิมสูริเย.\csromanspacer{} โส เอวํ วเทยฺย “นตฺถิ กณฺหสุกฺกานิ รูปานิ, นตฺถิ กณฺหสุกฺกานํ รูปานํ ทสฺสาวี.\csromanspacer{} นตฺถิ นีลกานิ รูปานิ, นตฺถิ นีลกานํ รูปานํ ทสฺสาวี.\csromanspacer{} นตฺถิ ปีตกานิ รูปานิ, นตฺถิ ปีตกานํ รูปานํ ทสฺสาวี.\csromanspacer{} นตฺถิ โลหิตกานิ รูปานิ, นตฺถิ โลหิตกานํ รูปานํ ทสฺสาวี.\csromanspacer{} นตฺถิ มญฺชิฏฺฐกานิ รูปานิ, นตฺถิ มญฺชิฏฺฐกานํ รูปานํ ทสฺสาวี.\csromanspacer{} นตฺถิ สมวิสมํ, นตฺถิ สมวิสมสฺส ทสฺสาวี.\csromanspacer{} นตฺถิ ตารกรูปานิ, นตฺถิ ตารกรูปานํ ทสฺสาวี.\csromanspacer{} นตฺถิ จนฺทิมสูริยา, นตฺถิ จนฺทิมสูริยานํ ทสฺสาวี.\csromanspacer{} อหเมตํ น ชานามิ, อหเมตํ น ปสฺสามิ, ตสฺมา ตํ นตฺถี”ติ.\csromanspacer{} สมฺมา นุ โข โส มาณว วทมาโน วเทยฺยาติ.}

\csromanmatikaanchor{mtk.58}
\csromantocmarkref{subsection}{9. สุภสุตฺต}{mtk.58}
\bookmark[dest={mtk.58},level=2]{9. สุภสุตฺต}
\prose{โน หิทํ โภ โคตม.\csromanspacer{} อตฺถิ กณฺหสุกฺกานิ รูปานิ, อตฺถิ กณฺหสุกฺกานํ รูปานํ ทสฺสาวี.\csromanspacer{} อตฺถิ นีลกานิ รูปานิ, อตฺถิ นีลกานํ รูปานํ ทสฺสาวี.\csromanspacer{} อตฺถิ ปีตกานิ รูปานิ, อตฺถิ ปีตกานํ รูปานํ ทสฺสาวี.\csromanspacer{} อตฺถิ โลหิตกานิ รูปานิ, อตฺถิ โลหิตกานํ รูปานํ ทสฺสาวี.\csromanspacer{} อตฺถิ มญฺชิฏฺฐกานิ รูปานิ, อตฺถิ มญฺชิฏฺฐกานํ รูปานํ ทสฺสาวี.\csromanspacer{} อตฺถิ สมวิสมํ, อตฺถิ สมวิสมสฺส ทสฺสาวี.\csromanspacer{} อตฺถิ ตารกรูปานิ, อตฺถิ ตารกรูปานํ \csromanfolio{418}ทสฺสาวี.\csromanspacer{} อตฺถิ จนฺทิมสูริยา, อตฺถิ จนฺทิมสูริยานํ ทสฺสาวี. “อหเมตํ น ชานามิ, อหเมตํ น ปสฺสามิ, ตสฺมา ตํ นตฺถี”ติ น หิ โส โภ โคตม สมฺมา วทมาโน วเทยฺยาติ.}

\prose{เอวเมว โข มาณว พฺราหฺมโณ โปกฺขรสาติ โอปมญฺโญ สุภควนิโก อนฺโธ อจกฺขุโก.\csromanspacer{} โส วต อุตฺตริ มนุสฺสธมฺมา อลมริยญาณทสฺสนวิเสสํ ญสฺสติ วา ทกฺขติ วา สจฺฉิ วา กริสฺสตีติ เนตํ ฐานํ วิชฺชติ.}

\proseitem{467}{ตํ กิํ มญฺญสิ มาณว, เย เต โกสลกา พฺราหฺมณมหาสาลา.\csromanspacer{} เสยฺยถิทํ, จงฺกี พฺราหฺมโณ ตารุกฺโข พฺราหฺมโณ โปกฺขรสาติ พฺราหฺมโณ ชาณุสฺโสณิ พฺราหฺมโณ ปิตา จ\footnote{วา (สี, สฺยา, กํ, อิ)} เต โตเทยฺโย.\csromanspacer{} กตมา เนสํ เสยฺโย\footnote{เสยฺยา (สฺยา, กํ)}, ยํ วา เต สมฺมุจฺจา\footnote{สมฺมุสา (สี, อิ)} วาจํ ภาเสยฺยุํ ยํ วา อสมฺมุจฺจาติ.\csromanspacer{} สมฺมุจฺจา โภ โคตม.}

\prose{กตมา เนสํ เสยฺโย, ยํ วา เต มนฺตา วาจํ ภาเสยฺยุํ ยํ วา อมนฺตาติ.\csromanspacer{} มนฺตา โภ โคตม.}

\prose{กตมา เนสํ เสยฺโย, ยํ วา เต ปฏิสงฺขาย วาจํ ภาเสยฺยุํ ยํ วา อปฺปฏิสงฺขายาติ.\csromanspacer{} ปฏิสงฺขาย โภ โคตม.}

\prose{กตมา เนสํ เสยฺโย, ยํ วา เต อตฺถสํหิตํ วาจํ ภาเสยฺยุํ ยํ วา อนตฺถ สํหิตนฺติ.\csromanspacer{} อตฺถสํหิตํ โภ โคตม.}

\prose{ตํ กิํ มญฺญสิ มาณว, ยทิ เอวํ สนฺเต พฺราหฺมเณน โปกฺขรสาตินา โอปมญฺเญน สุภควนิเกน สมฺมุจฺจา วา จา ภาสิตา อสมฺมุจฺจาติ\footnote{อสมฺมุสา วาติ (อิ) เอวมิตรปญฺหตฺตเยปิ วาสทฺเทน สห ทิสฺสติ.}? อสมฺมุจฺจา โภ โคตม.}

\prose{มนฺตา วาจา ภาสิตา อมนฺตาติ? อมนฺตา โภ โคตม.}

\prose{ปฏิสงฺขาย วาจา ภาสิตา อปฺปฏิสงฺขายาติ? อปฺปฏิสงฺขาย โภ โคตม.}

\prose{อตฺถสํหิตา วาจา ภาสิตา อนตฺถสํหิตาติ? อนตฺถสํหิตา โภ โคตม.}

\csromanheader{2}{9. สุภสุตฺต (99)}
\prose{\csromanfolio{419}ปญฺจ โข อิเม มาณว นีวรณา.\csromanspacer{} กตเม ปญฺจ, กามจฺฉนฺทนีวรณํ พฺยาปาทนีวรณํ ถินมิทฺธนีวรณํ อุทฺธจฺจกุกฺกุจฺจนีวรณํ วิจิกิจฺฉานีวรณํ.\csromanspacer{} อิเม โข มาณว ปญฺจ นีวรณา.\csromanspacer{} อิเมหิ โข มาณว ปญฺจหิ นีวรเณหิ พฺราหฺมโณ โปกฺขรสาติ โอปมญฺโญ สุภควนิโก อาวุโต นิวุโต โอผุโฏ\footnote{โอวุโต (สี), โอผุโต (สฺยา, กํ, อิ)} ปริโยนทฺโธ.\csromanspacer{} โส วต อุตฺตริ มนุสฺสธมฺมา อลมริยญาณทสฺสนวิเสสํ ญสฺสติ วา ทกฺขติ วา สจฺฉิ วา กริสฺสตีติ เนตํ ฐานํ วิชฺชติ.}

\proseitem{468}{ปญฺจ โข อิเม มาณว กามคุณา.\csromanspacer{} กตเม ปญฺจ, จกฺขุวิญฺเญยฺยา รูปา อิฏฺฐา กนฺตา มนาปา ปิยรูปา กามูปสํหิตา รชนียา.\csromanspacer{} โสตวิญฺเญยฺยา สทฺทา ฯเปฯ.\csromanspacer{} ฆานวิญฺเญยฺยา คนฺธา.\csromanspacer{} ชิวฺหา วิญฺเญยฺยา รสา.\csromanspacer{} กายวิญฺเญยฺยา โผฏฺฐพฺพา อิฏฺฐา กนฺตา มนาปา ปิยรูปา กามูปสํหิตา รชนียา.\csromanspacer{} อิเม โข มาณว ปญฺจ กามคุณา.\csromanspacer{} อิเมหิ โข มาณว ปญฺจหิ กามคุเณหิ พฺราหฺมโณ โปกฺขรสาติ โอปมญฺโญ สุภควนิโก คถิโต มุจฺฉิโต อชฺโฌปนฺโน อนาทีนวทสฺสาวี อนิสฺสรณปญฺโญ ปริภุญฺชติ.\csromanspacer{} โส วต อุตฺตริ มนุสฺสธมฺมา อลมริยญาณทสฺสนวิเสสํ ญสฺสติ วา ทกฺขติ วา สจฺฉิ วา กริสฺสตีติ เนตํ ฐานํ วิชฺชติ.}

\prose{ตํ กิํ มญฺญสิ มาณว, ยํ วา ติณกฏฺฐุปาทานํ ปฏิจฺจ อคฺคิํ ชาเลยฺย, ยํ วา นิสฺสฏฺฐติณกฏฺฐุปาทานํ อคฺคิํ ชาเลยฺย, กตโม นุ ขฺวาสฺส อคฺคิ อจฺจิมา เจว วณฺณวา จ ปภสฺสโร จาติ.\csromanspacer{} สเจ ตํ โภ โคตม ฐานํ นิสฺสฏฺฐติณกฏฺฐุปาทานํ อคฺคิํ ชาเลตุํ, สฺวาสฺส อคฺคิ อจฺจิมา เจว วณฺณวา จ ปภสฺสโร จาติ.\csromanspacer{} อฏฺฐานํ โข เอตํ มาณว อนวกาโส, ยํ นิสฺสฏฺฐติณกฏฺฐุปาทานํ อคฺคิํ ชาเลยฺย อญฺญตฺร อิทฺธิมตา.\csromanspacer{} เสยฺยถาปิ มาณว ติณกฏฺฐุปาทานํ ปฏิจฺจ อคฺคิ ชลติ, ตถูปมาหํ มาณว อิมํ ปีติํ วทามิ, ยายํ ปีติ ปญฺจ กามคุเณ ปฏิจฺจ.\csromanspacer{} เสยฺยถาปิ มาณว นิสฺสฏฺฐติณกฏฺฐุปาทาโน\footnote{นิสฺสฏฺฐติณกฏฺฐุปาทานํ ปฏิจฺจ (สี, อิ, ก)} อคฺคิ ชลติ, ตถูปมาหํ มาณว อิมํ ปีติํ วทามิ, ยายํ ปีติ อญฺญเตฺรว กาเมหิ อญฺญตฺร อกุสเลหิ ธมฺเมหิ.}

\prose{กตมา จ มาณว ปีติ อญฺญเตฺรว กาเมหิ อญฺญตฺร อกุสเลหิ ธมฺเมหิ.\csromanspacer{} อิธ มาณว ภิกฺขุ วิวิจฺเจว กาเมหิ ฯเปฯ ปฐมํ ฌานํ อุปสมฺปชฺช วิหรติ.\csromanspacer{} อยมฺปิ โข มาณว ปีติ อญฺญเตฺรว กาเมหิ อญฺญตฺร \csromanfolio{420}อกุสเลหิ ธมฺเมหิ.\csromanspacer{} ปุน จปรํ มาณว ภิกฺขุ วิตกฺกวิจารานํ วูปสมา ทุติยํ ฌานํ อุปสมฺปชฺช วิหรติ.\csromanspacer{} อยมฺปิ โข มาณว ปีติ อญฺญเตฺรว กาเมหิ อญฺญตฺร อกุสเลหิ ธมฺเมหิ.}

\proseitem{469}{เย เต มาณว พฺราหฺมณา ปญฺจ ธมฺเม ปญฺญเปนฺติ ปุญฺญสฺส กิริยาย กุสลสฺส อาราธนาย, กตเมตฺถ\footnote{กเมตฺถ (ก-สี, สฺยา, กํ, อิ)} พฺราหฺมณา ธมฺมํ มหปฺผลตรํ ปญฺญเปนฺติ ปุญฺญสฺส กิริยาย กุสลสฺส อาราธนายาติ.\csromanspacer{} เยเม โภ โคตม พฺราหฺมณา ปญฺจ ธมฺเม ปญฺญเปนฺติ ปุญฺญสฺส กิริยายกุสลสฺส อาราธนาย, จาคเมตฺถ พฺราหฺมณา ธมฺมํ มหปฺผลตรํ ปญฺญเปนฺติ ปุญฺญสฺส กิริยาย กุสลสฺส อาราธนายาติ.}

\prose{ตํ กิํ มญฺญสิ มาณว, อิธ อญฺญตรสฺส พฺราหฺมณสฺส มหายญฺโญ ปจฺจุปฏฺฐิโต อสฺส, อถ เทฺว พฺราหฺมณา อาคจฺเฉยฺยุํ “อิตฺถนฺนามสฺส พฺราหฺมณสฺส มหายญฺญํ อนุภวิสฺสามา”ติ.\csromanspacer{} ตเตฺรกสฺส\footnote{ตตฺเถกสฺส (อิ)} พฺราหฺมณสฺส เอวมสฺส “อโหวต อหเมว ลเภยฺยํ ภตฺตคฺเค อคฺคาสนํ อคฺโคทกํ อคฺคปิณฺฑํ, น อญฺโญ พฺราหฺมโณ ลเภยฺย ภตฺตคฺเค อคฺคาสนํ อคฺโคทกํ อคฺคปิณฺฑนฺ”ติ.\csromanspacer{} ฐานํ โข ปเนตํ มาณว วิชฺชติ, ยํ อญฺโญ พฺราหฺมโณ ลเภยฺย ภตฺตคฺเค อคฺคาสนํ อคฺโคทกํ อคฺคปิณฺฑํ.\csromanspacer{} น โส พฺราหฺมโณ ลเภยฺย ภตฺตคฺเค อคฺคาสนํ อคฺโคทกํ อคฺคปิณฺฑํ. “อญฺโญ พฺราหฺมโณ ลภติ ภตฺตคฺเค อคฺคาสนํ อคฺโคทกํ อคฺคปิณฺฑํ, นาหํ ลภามิ ภตฺตคฺเค อคฺคาสนํ อคฺโคทกํ อคฺคปิณฺฑนฺ”ติ อิติ โส กุปิโต โหติ อนตฺตมโน.\csromanspacer{} อิมสฺส ปน มาณว พฺราหฺมณา กิํ วิปากํ ปญฺญเปนฺตีติ.\csromanspacer{} น เขฺวตฺถ โภ โคตม พฺราหฺมณา เอวํ ทานํ เทนฺติ “อิมินา ปโร กุปิโต โหตุ อนตฺตมโน”ติ.\csromanspacer{} อถ เขฺวตฺถ พฺราหฺมณา อนุกมฺปาชาติกํเยว\footnote{อนุกมฺปชาติกํเยว (สฺยา, กํ, ก)} ทานํ เทนฺตีติ.\csromanspacer{} เอวํ สนฺเต โข มาณว พฺราหฺมณานํ อิทํ ฉฏฺฐํ ปุญฺญกิริยวตฺถุ โหติ, ยทิทํ อนุกมฺปาชาติกนฺติ.\csromanspacer{} เอวํ สนฺเต โภ โคตม พฺราหฺมณานํ อิทํ ฉฏฺฐํ ปุญฺญกิริยวตฺถุ โหติ, ยทิทํ อนุกมฺปาชาติกนฺติ.}

\prose{เย เต มาณว พฺราหฺมณา ปญฺจ ธมฺเม ปญฺญเปนฺติ ปุญฺญสฺส กิริยาย กุสลสฺส อาราธนาย, อิเม ตฺวํ ปญฺจ ธมฺเม กตฺถ พหุลํ สมนุปสฺสสิ คหฏฺเฐสุ วา ปพฺพชิเตสุ วาติ.\csromanspacer{} เยเม โภ โคตม พฺราหฺมณา ปญฺจ ธมฺเม ปญฺญเปนฺติ ปุญฺญสฺส กิริยาย กุสลสฺส อาราธนาย, อิมาหํ ปญฺจ ธมฺเม}

\vspace{\tipitakaparskip}
\csromancenter{สุภสุตฺต (99) ปพฺพชิเตสุ พหุลํ สมนุปสฺสามิ อปฺปํ คหฏฺเฐสุ.\csromanspacer{} คหฏฺโฐ หิ โภ โคตม มหฏฺโฐ มหากิจฺโจ มหาธิกรโณ มหาสมารมฺโภ, น สตตํ สมิตํ สจฺจวาที โหติ.\csromanspacer{} ปพฺพชิโต โข ปน โภ โคตม อปฺปฏฺโฐ อปฺปกิจฺโจ อปฺปาธิกรโณ อปฺปสมารมฺโภ, สตตํ สมิตํ สจฺจวาที โหติ.\csromanspacer{} คหฏฺโฐ หิ โภ โคตม มหฏฺโฐ มหากิจฺโจ มหาธิกรโณ มหาสมารมฺโภ, น สตตํ สมิตํ ตปสฺสี โหติ.\csromanspacer{} พฺรหฺมจารี โหติ.\csromanspacer{} สชฺฌายพหุโล โหติ.\csromanspacer{} จาคพหุโล โหติ.\csromanspacer{} ปพฺพชิโต โข ปน โภ โคตม อปฺปฏฺโฐ อปฺปกิจฺโจ อปฺปาธิกรโณ อปฺปสมารมฺโภ, สตตํ สมิตํ ตปสฺสี โหติ.\csromanspacer{} พฺรหฺมจารี โหติ.\csromanspacer{} สชฺฌายพหุโล โหติ.\csromanspacer{} จาคพหุโล โหติ.\csromanspacer{} เยเม โภ โคตม พฺราหฺมณา ปญฺจ ธมฺเม ปญฺญเปนฺติ ปุญฺญสฺส กิริยาย กุสลสฺส อาราธนาย, อิมาหํ ปญฺจ ธมฺเม ปพฺพชิเตสุ พหุลํ สมนุปสฺสามิ, อปํ คหฏฺเฐสูติ.}
\prose{\csromanfolio{421}เย เต มาณว พฺราหฺมณา ปญฺจ ธมฺเม ปญฺญเปนฺติ ปุญฺญสฺส กิริยาย กุสลสฺส อาราธนาย, จิตฺตสฺสาหํ เอเต ปริกฺขาเร วทามิ, ยทิทํ จิตฺตํ อเวรํ อพฺยาพชฺฌํ, ตสฺส ภาวนาย.\csromanspacer{} อิธ มาณว ภิกฺขุ สจฺจวาที โหติ, โส “สจฺจวาทีมฺหี”ติ ลภติ อตฺถเวทํ, ลภติ ธมฺมเวทํ, ลภติ ธมฺมูปสํหิตํ ปาโมชฺชํ.\csromanspacer{} ยํ ตํ กุสลูปสํหิตํ ปาโมชฺชํ, จิตฺตสฺสาหํ เอตํ ปริกฺขารํ วทามิ, ยทิทํ จิตฺตํ อเวรํ อพฺยาพชฺฌํ, ตสฺส ภาวนาย.\csromanspacer{} อิธ มาณว ภิกฺขุ ตปสฺสี โหติ ฯเปฯ พฺรหฺมจารี โหติ ฯเปฯ สชฺฌายพหุโล โหติ ฯเปฯ จาคพหุโล โหติ, โส “จาคพหุโลมฺหี”ติ ลภติ อตฺถเวทํ, ลภติ ธมฺมเวทํ, ลภติ ธมฺมูปสํหิตํ ปาโมชฺชํ.\csromanspacer{} ยํ ตํ กุสลูปสํหิตํ ปาโมชฺชํ, จิตฺตสฺสาหํ เอตํ ปริกฺขารํ วทามิ, ยทิทํ จิตฺตํ อเวรํ อพฺยาพชฺฌํ, ตสฺส ภวนาย.\csromanspacer{} เย เต มาณว พฺราหฺมณา ปญฺจ ธมฺเม ปญฺญเปนฺติ ปุญฺญสฺส กิริยาย กุสลสฺส อาราธนาย, จิตฺตสฺสาหํ เอเต ปริกฺขาเร วทามิ, ยทิทํ จิตฺตํ อเวรํ อพฺยาพชฺฌํ, ตสฺส ภาวนายาติ.}

\proseitem{470}{เอวํ วุตฺเต สุโภ มาณโว โตเทยฺยปุตฺโต ภควนฺตํ เอตทโวจ “สุตํ เมตํ โภ โคตม ‘สมโณ โคตโม พฺรหฺมานํ สหพฺยตาย มคฺคํ ชานาตี’ติ”.}

\prose{ตํ กิํ มญฺญสิ มาณว, อาสนฺเน อิโต นฬการคาโม, น ยิโต ทูเร นฬการคาโมติ? เอวํ โภ อาสนฺเน อิโต \csromanfolio{422}นฬการคาโม, น ยิโต ทูเร นฬการาคาโมติ.\csromanspacer{} ตํ กิํ มญฺญสิ มาณว, อิธสฺส ปุริโส นฬการคาเม ชาตวทฺโธ\footnote{ชาตวฑฺโฒ (สฺยา, กํ, ก)}, ตเมนํ นฬการคามโต ตาวเทว อวสฏํ\footnote{อปสกฺกํ (สฺยา, กํ, ก) 3-3. เอวเมว โข มาณว เอวํ ภาวิตาย เมตฺตาย (สี, สฺยา, กํ, อิ, ที 1. 235 ปิฏฺเฐปิ) ตถาปิ อิธ ปาโฐเยว อุปมาย สํสนฺทิยมาโน ปุริปุณฺโณ วิย ทิสฺสติ.} นฬการคามสฺส มคฺคํ ปุจฺเฉยฺยุํ.\csromanspacer{} สิยา นุ โข มาณว ตสฺส ปุริสสฺส นฬการคาเม ชาตวทฺธสฺส นฬการคามสฺส มคฺคํ ปุฏฺฐสฺส ทนฺธายิตฺตํ วา วิตฺถายิตตฺตํ วาติ.\csromanspacer{} โน หิทํ โภ โคตม.\csromanspacer{} ตํ กิสฺส เหตุ, อมุ หิ โภ โคตม ปุริโส นฬการคาเม ชาตวทฺโธ, ตสฺส สพฺพาเนว นฬการคามสฺส มคฺคานิ สุวิทิตานีติ.\csromanspacer{} สิยา นุ โข มาณว ตสฺส ปุริสสฺส นฬการคาเม ชาตวทฺธสฺส นฬการคามสฺส มคฺคํ ปุฏฺฐสฺส ทนฺธายิตตฺตํ วา วิตฺถายิตตฺตํ วาติ.\csromanspacer{} น เตฺวว ตถาคตสฺส พฺราหฺมโลกํ วา พฺราหฺมโลกคามินิํ วา ปฏิปทํ ปุฏฺฐสฺส ทนฺธายิตตฺตํ วา วิตฺถายิตตฺตํ วา.\csromanspacer{} พฺรหฺมานญฺจาหํ มาณว ปชานามิ พฺรหฺมโลกญฺจ พฺรหฺมโลกคามินิญฺจ ปฏิปทํ, ยถาปฏิปนฺโน จ พฺรหฺมโลกํ อุปปนฺโน, ตญฺจ ปชานามีติ.\csromanspacer{} สุตํ เมตํ โภ โคตม “สมโณ โคตโม พฺรหฺมานํ สหพฺยตาย มคฺคํ เทเสตี”ติ, สาธุ เม ภวํ โคตโม พฺรหฺมานํ สหพฺยตาย มคฺคํ เทเสตูติ.\csromanspacer{} เตน หิ มาณว สุณาหิ สาธุกํ มนสิ กโรหิ ภาสิสฺสามีติ. “เอวํ โภ”ติ โข สุโภ มาณโว โตเทยฺยปุตฺโต ภควโต ปจฺจสฺโสสิ.\csromanspacer{} ภควา เอตทโวจ\csromandash{}}

\proseitem{471}{กตโม จ มาณว พฺรหฺมานํ สหพฺยตาย มคฺโค.\csromanspacer{} อิธ มาณว ภิกฺขุ เมตฺตาสหคเตน เจตสา เอกํ ทิสํ ผริตฺวา วิหรติ.\csromanspacer{} ตถา ทุติยํ.\csromanspacer{} ตถา ตติยํ.\csromanspacer{} ตถา จตุตฺถํ.\csromanspacer{} อิติ อุทฺธมโธ ติริยํ สพฺพธิ สพฺพตฺตตาย สพฺพาวนฺตํ โลกํ เมตฺตาสหคเตน เจตสา วิปุเลน มหคฺคเตน อปฺปมาเณน อเวเรน อพฺยาพชฺเฌน ผริตฺวา วิหรติ.\csromanspacer{} เอวํ ภาวิตาย โข มาณว เมตฺตาย เจโตวิมุตฺติยา ยํ ปมาณกตํ กมฺมํ, น ตํ ตตฺราวสิสฺสติ, น ตํ ตตฺราวติฏฺฐติ.\csromanspacer{} เสยฺยถาปิ มาณว พลวา สงฺขธโม อปฺปกสิเรเนว จาตุทฺทิสา วิญฺญาเปยฺย.3 เอวเมว โข มาณว ฯเปฯ.\csromanspacer{} เอวํ ภาวิตาย โข มาณว เมตฺตาย3 เจโตวิมุตฺติยา ยํ ปมาณกตํ กมฺมํ, น ตํ ตตฺราวสิสฺสติ, น ตํ ตตฺราวติฏฺฐติ.\csromanspacer{} อยมฺปิ โข มาณว พฺรหฺมานํ สหพฺยตาย มคฺโค.}

\csromanheader{2}{9. สุภสุตฺต (99)}
\prose{\csromanfolio{423}ปุน จปรํ มาณว ภิกฺขุ กรุณาสหคเตน เจตสา ฯเปฯ มุทิตาสหคเตน เจตสา ฯเปฯ อุเปกฺขาสหคเตน เจตสา เอกํ ทิสํ ผริตฺวา วิหรติ.\csromanspacer{} ตถา ทุติยํ.\csromanspacer{} ตถา ตติยํ.\csromanspacer{} ตถา จตุตฺถํ.\csromanspacer{} อิติ อุทฺธมโธ ติริยํ สพฺพธิ สพฺพตฺตตาย สพฺพาวนฺตํ โลกํ อุเปกฺขาสหคเตน เจตสา วิปุเลน มหคฺคเตน อปฺปมาเณน อเวเรน อพฺยาพชฺเฌน ผริตฺวา วิหรติ.\csromanspacer{} เอวํ ภาวิตาย โข มาณว อุเปกฺขาย เจโตวิมุตฺติยา ยํ ปมาณกตํ กมฺมํ, น ตํ ตตฺราวสิสฺสติ, น ตํ ตตฺราวติฏฺฐติ.\csromanspacer{} เสยฺยถาปิ มาณว พลวา สงฺขธโม อปฺปกสิเรเนว จาตุทฺทิสา วิญฺญาเปยฺย.\csromanspacer{} เอวเมว โข มาณว ฯเปฯ.\csromanspacer{} เอวํ ภาวิตาย โข มาณว อุเปกฺขาย เจโตวิมุตฺติยา ยํ ปมาณกตํ กมฺมํ, น ตํ ตตฺราวสิสฺสติ, น ตํ ตตฺราวติฏฺฐติ.\csromanspacer{} อยมฺปิ โข มาณว พฺรหฺมานํ สหพฺยตาย มคฺโคติ.}

\proseitem{472}{เอวํ วุตฺเต สุโภ มาณโว โตเทยฺยปุตฺโต ภควนฺตํ เอตทโวจ “อภิกฺกนฺตํ โภ โคตม, อภิกฺกนฺตํ โภ โคตม, เสยฺยถาปิ โภ โคตม นิกฺกุชฺชิตํ วา อุกฺกุชฺเชยฺย, ปฏิจฺฉนฺนํ วา วิวรยฺย, มูฬฺหสฺส วา มคฺคํ อาจิกฺเขยฺย, อนฺธกาเร วา เตลปชฺโชตํ ธาเรยฺย ‘จกฺขุมนฺโต รูปานิ ทกฺขนฺตี’ติ.\csromanspacer{} เอวเมวํ โภตา โคตเมน อเนกปริยาเยน ธมฺโม ปกาสิโต, เอสาหํ ภวนฺตํ โคตมํ สรณํ คจฺฉามิ ธมฺมญฺจ ภิกฺขุสํฆญฺจ, อุปาสกํ มํ ภวํ โคตโม ธาเรตุ อชฺชตคฺเค ปาณุเปตํ สรณํ คตํ.\csromanspacer{} หนฺท จ ทานิ มยํ โภ โคตม คจฺฉาม พหุกิจฺจา มยํ พหุกรณียา”ติ.\csromanspacer{} ยสฺสทานิ ตฺวํ มาณว กาลํ มญฺญสีติ.\csromanspacer{} อถ โข สุโภ มาณโว โตเทยฺยปุตฺโต ภควโต ภาสิตํ อภินนฺทิตฺวา อนุโมทิตฺวา อุฏฺฐายาสนา ภควนฺตํ อภิวาเทตฺวา ปทกฺขิณํ กตฺวา ปกฺกามิ.}

\prose{เตน โข ปน สมเยน ชาณุสฺโสณิ พฺราหฺมโณ สพฺพเสเตน วฬวาภิรเถน\footnote{วฬภีรเถน (สี)} สาวตฺถิยา นิยฺยาติ ทิวา ทิวสฺส.\csromanspacer{} อทฺทสา โข ชาณุสฺโสณิ พฺราหฺมโณ สุภํ มาณวํ โตเทยฺยปุตฺตํ ทูรโตว อาคจฺฉนฺตํ, ทิสฺวาน สุภํ มาณวํ โตเทยฺยปุตฺตํ เอตทโวจ “หนฺท กุโต นุ ภวํ ภารทฺวาโช อาคจฺฉติ ทิวา ทิวสฺสา”ติ.\csromanspacer{} อิโต หิ โข อหํ \csromanfolio{424}โภ อาคจฺฉามิ สมณสฺส โคตมสฺส สนฺติกาติ.\csromanspacer{} ตํ กิํ มญฺญสิ ภวํ ภารทฺวาโช สมณสฺส โคตมสฺส ปญฺญาเวยฺยตฺติยํ ปณฺฑิโต มญฺเญติ.\csromanspacer{} โก จาหํ โภ, โก จ สมณสฺส โคตมสฺส ปญฺญาเวยฺยตฺติยํ ชานิสฺสามิ.\csromanspacer{} โสปิ นูนสฺส ตาทิโสว, โย สมณสฺส โคตมสฺส ปญฺญาเวยฺยตฺติยํ ชาเนยฺยาติ.\csromanspacer{} อุฬาราย ขลุ ภวํ ภารทฺวาโช สมณํ โคตมํ ปสํสาย ปสํสตีติ.\csromanspacer{} โก จาหํ โภ, โก จ สมณํ โคตมํ ปสํสิสฺสามิ.\csromanspacer{} ปสตฺถปสตฺโถว โส ภวํ โคตโม เสฏฺโฐ เทวมนุสฺสานํ.\csromanspacer{} เย จิเม โภ พฺราหฺมณา ปญฺจ ธมฺเม ปญฺญเปนฺติ ปุญฺญสฺส กิริยาย กุสลสฺส อาราธนาย, จิตฺตสฺเสเต สมโณ โคตโม ปริกฺขาเร วเทติ, ยทิทํ จิตฺตํ อเวรํ อพฺยาพชฺฌํ, ตสฺส ภาวนายาติ.}

\prose{เอวํ วุตฺเต ชาณุสฺโสณิ พฺราหฺมโณ สพฺพเสตา วฬวาภิรถา โอโรหิตฺวา เอกํสํ อุตฺตราสงฺคํ กริตฺวา เยน ภควา เตนญฺชลิํ ปณาเมตฺวา อุทานํ อุทาเนสิ “ลาภา รญฺโญ ปเสนทิสฺส โกสลสฺส, สุลทฺธลาภา รญฺโญ ปเสนทิสฺส โกสลสฺส, ยสฺส วิชิเต ตถาคโต วิหรติ อรหํ สมฺมาสมฺพุทฺโธ”ติ.}

\nitthitamruled{สุภสุตฺตํ นิฏฺฐิตํ นวมํ.}
\csromanheader{2}{10. สงฺคารวสุตฺต}
\proseitem{473}{เอวํ เม สุตํ\csromandash{}เอกํ สมยํ ภควา โกสเลสุ จาริกํ จรติ มหตา ภิกฺขุสํเฆน สทฺธิํ.\csromanspacer{} เตน โข ปน สมเยน ธนญฺชานี\footnote{ธานญฺชานี (สี, อิ)} นาม พฺราหฺมณี จญฺจลิกปฺเป\footnote{มณฺฑลกปฺเป (สี), ปจฺจลกปฺเป (สฺยา, กํ), จณฺฑลกปฺเป (อิ)} ปฏิวสติ, อภิปฺปสนฺนา พุทฺเธ จ ธมฺเม จ สํเฆ จ.\csromanspacer{} อถ โข ธนญฺชานี พฺราหฺมณี อุปกฺขลิตฺวา ติกฺขตฺตุํ อุทานํ อุทาเนสิ “นโม ตสฺส ภควโต อรหโต สมฺมาสมฺพุทฺธสฺส, นโม ตสฺส ภควโต อรหโต สมฺมาสมฺพุทฺธสฺส, นโม ตสฺส ภควโต อรหโต สมฺมาสมฺพุทฺธสฺสา”ติ.}

\prose{เตน โข ปน สมเยน สงฺคารโว นาม มาณโว จญฺจลิกปฺเป ปฏิวสติ ติณฺณํ เวทานํ ปารคู สนิฆณฺฑุเกฏุภานํ สากฺขรปฺปเภทานํ}

\csromanheader{2}{10. สงฺคารวสุตฺต (100)}
\prose{\csromanfolio{425}อิติหาสปญฺจมานํ ปทโก เวยฺยากรโณ โลกายตมหาปุริสลกฺขเณสุ อนวโย.\csromanspacer{} อสฺโสสิ โข สงฺคารโว มาณโว ธนญฺชานิยา พฺราหฺมณิยา เอวํ วาจํ ภาสมานาย, สุตฺวา ธนญฺชานิํ พฺราหฺมณิํ เอตทโวจ “อวภูตาว อยํ\footnote{อวภูตา จยํ (สี, สฺยา, กํ, อิ)} ธนญฺชานี พฺราหฺมณี, ปรภูตาว อยํ\footnote{ปราภูตา จยํ (สี, สฺยา, กํ, อิ)} ธนญฺชานี พฺราหฺมณี วิชฺชมานานํ (เตวิชฺชานํ)\footnote{( ) สี-สฺยา-กํ-อิ-โปตฺถเกสุ นตฺถิ.} พฺราหฺมณานํ อถ จ ปน ตสฺส มุณฺฑกสฺส สมณกสฺส วณฺณํ ภาสิสฺสตี”ติ\footnote{ภาสตีติ (สี, สฺยา, กํ, อิ)}.\csromanspacer{} น หิ ปน ตฺวํ ตาต ภทฺรมุข ตสฺส ภควโต สีลปญฺญาณํ ชานาสิ, สเจ ตฺวํ ตาต ภทฺรมุข ตสฺส ภควโต สีลปญฺญาณํ ชาเนยฺยาสิ, น ตฺวํ ตาต ภทฺรมุข ตํ ภควนฺตํ อกฺโกสิตพฺพํ ปริภาสิตพฺพํ มญฺเญยฺยาสีติ.\csromanspacer{} เตน หิ โภติ ยทา สมโณ โคตโม จญฺจลิกปฺปํ อนุปฺปตฺโต โหติ, อถ เม อาโรเจยฺยาสีติ. “เอวํ ภทฺรมุขา”ติ โข ธนญฺชานี พฺราหฺมณี สงฺคารวสฺส มาณวสฺส ปจฺจสฺโสสิ.}

\prose{อถ โข ภควา โกสเลสุ อนุปุพฺเพน จาริกํ จรมาโน เยน จญฺจลิกปฺปํ ตทวสริ.\csromanspacer{} ตตฺร สุทํ ภควา จญฺจลิกปฺเป วิหรติ โตเทยฺยานํ พฺราหฺมณานํ อมฺพวเน.\csromanspacer{} อสฺโสสิ โข ธนญฺชานี พฺราหฺมณี “ภควา กิร จญฺจลิกปฺปํ อนุปฺปตฺโต, จญฺจลิกปฺเป วิหรติ โตเทยฺยานํ พฺราหฺมณานํ อมฺพวเน”ติ.\csromanspacer{} อถ โข ธนญฺชานี พฺราหฺมณี เยน สงฺคารโว มาณโว เตนุปสงฺกมิ, อุปสงฺกมิตฺวา สงฺคารวํ มาณวํ เอตทโวจ “อยํ ตาต ภทฺรมุข โส ภควา จญฺจลิกปฺปํ อนุปฺปตฺโต, จญฺจลิกปฺเป วิหรติ โตเทยฺยานํ พฺราหฺมณานํ อมฺพวเน, ยสฺสทานิ ตาต ภทฺรมุข กาลํ มญฺญสี”ติ.}

\proseitem{474}{“เอวํ โภ”ติ โข สงฺคารโว มาณโว ธนญฺชานิยา พฺราหฺมณิยา ปฏิสฺสุตฺวา เยน ภควา เตนุปสงฺกมิ, อุปสงฺกมิตฺวา ภควตา สทฺธิํ สมฺโมทิ, สมฺโมทนียํ กถํ สารณียํ วีติสาเรตฺวา เอกมนฺตํ นิสีทิ, เอกมนฺตํ นิสินฺโน โข สงฺคารโว มาณโว ภควนฺตํ เอตทโวจ “สนฺติ โข โภ โคตม เอเก สมณพฺราหฺมณา ทิฏฺฐธมฺมาภิญฺญาโวสานปารมิปฺปตฺตา อาทิพฺรหฺมจริยํ ปฏิชานนฺติ.\csromanspacer{} ตตฺร โภ โคตม \csromanfolio{426}เย เต สมณพฺราหฺมณา ทิฏฺฐธมฺมาภิญฺญาโวสานปารมิปฺปตฺตา อาทิพฺรหฺมจริยํ ปฏิชานนฺติ, เตสํ ภวํ โคตโม กตโมติ? ทิฏฺฐธมฺมาภิญฺญาโวสานปารมิปฺปตฺตานํ อาทิพฺรหฺมจริยํ ปฏิชานนฺตานมฺปิ โข อหํ ภารทฺวาช เวมตฺตํ วทามิ.\csromanspacer{} สนฺติ ภารทฺวาช เอเก สมณพฺราหฺมณา อนุสฺสวิกา, เต อนุสฺสเวน ทิฏฺฐธมฺมาภิญฺญาโวสานปารมิปฺปตฺตา อาทิพฺรหฺมจริยํ ปฏิชานนฺติ, เสยฺยถาปิ พฺราหฺมณา เตวิชฺชา.\csromanspacer{} สนฺติ ปน ภารทฺวาน เอเก สมณพฺราหฺมณา เกวลํ สทฺธามตฺตเกน ทิฏฺฐธมฺมาภิญฺญาโวสานปารมิปฺปตฺตา อาทิพฺรหฺมจริยํ ปฏิชานนฺติ, เสยฺยถาปิ ตกฺกี วีมํสี.\csromanspacer{} สนฺติ ภารทฺวาช เอเก สมณพฺราหฺมณา ปุพฺเพ อนนุสฺสุเตสุ ธมฺเมสุ สามํเยว ธมฺมํ อภิญฺญาย ทิฏฺฐธมฺมาภิญฺญาโวสานปารมิปฺปตฺตา อาทิพฺรหฺมจริยํ ปฏิชานนฺติ.\csromanspacer{} ตตฺร ภารทฺวาช เย เต สมณพฺราหฺมณา ปุพฺเพ อนนุสฺสุเตสุ ธมฺเมสุ สามํเยว ธมฺมํ อภิญฺญาย ทิฏฺฐธมฺมาภิญฺญาโวสานปารมิปฺปตฺตา อาทิพฺรหฺมจริยํ ปฏิชานนฺติ, เตสาหมสฺมิ.\csromanspacer{} ตทมินาเปตํ ภารทฺวาน ปริยาเยน เวทิตพฺพํ, ยถา เย เต สมณพฺราหฺมณา ปุพฺเพ อนนุสฺสุเตสุ ธมฺเมสุ สามํเยว ธมฺมํ อภิญฺญาย ทิฏฺฐธมฺมาภิญฺญาโวสานปารมิปฺปตฺตา อาทิพฺรหฺมจริยํ ปฏิชานนฺติ, เตสาหมสฺมิ.}

\proseitem{475}{อิธ เม ภารทฺวาช ปุพฺเพว สมฺโพธา อนภิสมฺพุทฺธสฺส โพธิสตฺตสฺเสว สโต เอตทโหสิ “สมฺพาโธ ฆราวาโส รชาปโถ, อพฺโภกาโส ปพฺพชฺชา, น ยิทํ สุกรํ อคารํ อชฺฌาวสตา เอกนฺตปริปุณฺณํ เอกนฺตปริสุทฺธํ สงฺขลิขิตํ พฺรหฺมจริยํ จริตุํ, ยํนูนาหํ เกสมสฺสุํ โอหาเรตฺวา กาสายานิ วตฺถานิ อจฺฉาเทตฺวา อคารสฺมา อนคาริยํ ปพฺพเชยฺยนฺ”ติ.\csromanspacer{} โส โข อหํ ภารทฺวาช อปเรน สมเยน ทหโรว สมาโน สุสุกาฬเกโส ภเทฺรน โยพฺพเนน สมนฺนาคโต, ปฐเมน วยสา อกามกานํ มาตาปิตูนํ อสฺสุมุขานํ รุทนฺตานํ เกสมสฺสุํ โอหาเรตฺวา กาสายานิ วตฺถานิ อจฺฉาเทตฺวา อคารสฺมา อนคาริยํ ปพฺพชิํ, โส เอวํ ปพฺพชิโต สมาโน กิํ กุสลคเวสี อนุตฺตรํ สนฺติวรปทํ ปริเยสมาโน เยน อาฬาโร กาลาโม เตนุปสงฺกมิํ, อุปสงฺกมิตฺวา อาฬารํ กาลามํ เอตทโวจํ “อิจฺฉามหํ อาวุโส กาลาม อิมสฺมิํ ธมฺมวินเย พฺรหฺมจริยํ จริตุนฺ”ติ.\csromanspacer{} เอวํ วุตฺเต ภารทฺวาช อาฬาโร กาลาโม มํ เอตทโวจ “วิหรตายสฺมา, ตาทิโส อยํ ธมฺโม, ยตฺถ วิญฺญู ปุริโส นจิรสฺเสว สกํ อาจริยกํ}

\csromanheader{2}{10. สงฺคารวสุตฺต (100)}
\prose{\csromanfolio{427}สยํ อภิญฺญา สจฺฉิกตฺวา อุปสมฺปชฺช วิหเรยฺยา”ติ.\csromanspacer{} โส โข อหํ ภารทฺวาช นจิรสฺเสว ขิปฺปเมว ตํ ธมฺมํ ปริยาปุณิํ, โส โข อหํ ภารทฺวาช ตาวตเกเนว โอฏฺฐปหตมตฺเตน ลปิตลาปนมตฺเตน ญาณวาทญฺจ วทามิ เถรวาทญฺจ, “ชานามิ ปสฺสามี”ติ จ ปฏิชานามิ อหญฺเจว อญฺเญ จ.\csromanspacer{} ตสฺส มยฺหํ ภารทฺวาช เอตทโหสิ “น โข อาฬาโร กาลาโม อิมํ ธมฺมํ เกวลํ สทฺธามตฺตเกน ‘สยํ อภิญฺญา สจฺฉิกตฺวา อุปสมฺปชฺช วิหรามี’ติ ปเวเทติ, อทฺธา อาฬาโร กาลาโม อิมํ ธมฺมํ ชานํ ปสฺสํ วิหรตี”ติ.}

\prose{อถ ขฺวาหํ ภารทฺวาช เยน อาฬาโร กาลาโม เตนุปสงฺกมิํ, อุปสงฺกมิตฺวา อาฬารํ กาลามํ เอตทโวจํ “กิตฺตาวตา โน อาวุโส กาลาม อิมํ ธมฺมํ สยํ อภิญฺญา สจฺฉิกตฺวา อุปสมฺปชฺช วิหรามีติ ปเวเทสี”ติ.\csromanspacer{} เอวํ วุตฺเต ภารทฺวาช อาฬาโร กาลาโม อากิญฺจญฺญายตนํ ปเวเทสิ.\csromanspacer{} ตสฺส มยฺหํ ภารทฺวาช เอตทโหสิ “น โข อาฬารสฺเสว กาลามสฺส อตฺถิ สทฺธา, มยฺหํปตฺถิ สทฺธา.\csromanspacer{} น โข อาฬารสฺเสว กาลามสฺส อตฺถิ วีริยํ ฯเปฯ สติ.\csromanspacer{} สมาธิ.\csromanspacer{} ปญฺญา, มยฺหํปตฺถิ ปญฺญา.\csromanspacer{} ยํนูนาหํ ยํ ธมฺมํ อาฬาโร กาลาโม ‘สยํ อภิญฺญา สจฺฉิกตฺวา อุปสมฺปชฺช วิหรามี’ติ ปเวเทติ, ตสฺส ธมฺมสฺส สจฺฉิกิริยาย ปทเหยฺยนฺ”ติ.\csromanspacer{} โส โข อหํ ภารทฺวาช นจิรสฺเสว ขิปฺปเมว ตํ ธมฺมํ สยํ อภิญฺญา สจฺฉิกตฺวา อุปสมฺปชฺช วิหาสิํ.\csromanspacer{} อถ ขฺวาหํ ภารทฺวาช เยน อาฬาโร กาลาโม เตนุปสงฺกมิํ, อุปสงฺกมิตฺวา อาฬารํ กาลามํ เอตทโวจํ “เอตฺตาวตา โน อาวุโส กาลาม อิมํ ธมฺมํ สยํ อภิญฺญา สจฺฉิกตฺวา อุปสมฺปชฺช ปเวเทสี”ติ.\csromanspacer{} เอตฺตาวตา โข อหํ อาวุโส อิมํ ธมฺมํ สยํ อภิญฺญา สจฺฉิกตฺวา อุปสมฺปชฺช ปเวเทมีติ.\csromanspacer{} อหมฺปิ โข อาวุโส เอตฺตาวตา อิมํ ธมฺมํ สยํ อภิญฺญา สจฺฉิกตฺวา อุปสมฺปชฺช วิหรามีติ.\csromanspacer{} ลาภา โน อาวุโส, สุลทฺธํ โน อาวุโส, เย มยํ อายสฺมนฺตํ ตาทิสํ สพฺรหฺมจาริํ ปสฺสาม.\csromanspacer{} อิติ ยาหํ ธมฺมํ สยํ อภิญฺญา สจฺฉิกตฺวา อุปสมฺปชฺช ปเวเทมิ, ตํ ตฺวํ ธมฺมํ สยํ อภิญฺญา สจฺฉิกตฺวา อุปสมฺปชฺช วิหรสิ.\csromanspacer{} ยํ ตฺวํ ธมฺมํ สยํ อภิญฺญา สจฺฉิกตฺวา อุปสมฺปชฺช วิหรสิ, ตมหํ ธมฺมํ สยํ อภิญฺญา สจฺฉิกตฺวา อุปสมฺปชฺช ปเวเทมิ.\csromanspacer{} อิติ ยาหํ ธมฺมํ ชานามิ, ตํ ตฺวํ ธมฺมํ ชานาสิ.\csromanspacer{} ยํ ตฺวํ ธมฺมํ ชานาสิ, ตมหํ ธมฺมํ}

\prose{\csromanfolio{428}ชานามิ.\csromanspacer{} อิติ ยาทิโส อหํ, ตาทิโส ตุวํ.\csromanspacer{} ยาทิโส ตุวํ, ตาทิโส อหํ.\csromanspacer{} เอหิ ทานิ อาวุโส อุโภว สนฺตา อิมํ คณํ ปริหรามาติ.\csromanspacer{} อิติ โข ภารทฺวาช อาฬาโร กาลาโม อาจริโย เม สมาโน อตฺตโน อนฺเตวาสิํ มํ สมานํ อตฺตนา สมสมํ ฐเปสิ, อุฬาราย จ มํ ปูชาย ปูเชสิ.\csromanspacer{} ตสฺส มยฺหํ ภารทฺวาช เอตทโหสิ “นายํ ธมฺโม นิพฺพิทาย น วิราคาย น นิโรธาย น อุปสมาย น อภิญฺญาย น สมฺโพธาย น นิพฺพานาย สํวตฺตติ, ยาวเทว อากิญฺจญฺญายตนูปปตฺติยา”ติ.\csromanspacer{} โส โข อหํ ภารทฺวาช ตํ ธมฺมํ อนลงฺกริตฺวา ตสฺมา ธมฺมา นิพฺพิชฺช อปกฺกมิํ.}

\proseitem{476}{โส โข อหํ ภารทฺวาช กิํกุสลคเวสี อนุตฺตรํ สนฺติวรปทํ ปริเยสมาโน เยน อุทโก รามปุตฺโต เตนุปสงฺกมิํ, อุปสงฺกมิตฺวา อุทกํ รามปุตฺตํ เอตทโวจํ “อิจฺฉามหํ อาวุโส\footnote{ปสฺส ม 1 ปาสราสิสุตฺเต (221) ปิฏฺเฐ.} อิมสฺมิํ ธมฺมวินเย พฺรหฺมจริยํ จริตุนฺ”ติ.\csromanspacer{} เอวํ วุตฺเต ภารทฺวาช อุทโก รามปุตฺโต มํ เอตทโวจ “วิหรตายสฺมา, ตาทิโส อยํ ธมฺโม, ยตฺถ วิญฺญู ปุริโส นจิรสฺเสว สกํ อาจริยกํ สยํ อภิญฺญา สจฺฉิกตฺวา อุปสมฺปชฺช วิหเรยฺยา”ติ.\csromanspacer{} โส โข อหํ ภารทฺวาช นจิรสฺเสว ขิปฺปเมว ตํ ธมฺมํ ปริยาปุณิํ.\csromanspacer{} โส โข อหํ ภารทฺวาช ตาวตเกเนว โอฏฺฐปหตมตฺเตน ลปิตลาปนมตฺเตน ญาณวาทญฺจ วทามิ เถรวาทญฺจ, “ชานามิ ปสฺสามี”ติ จ ปฏิชานามิ อหญฺเจว อญฺเญ จ.\csromanspacer{} ตสฺส มยฺหํ ภารทฺวาช เอตทโหสิ “น โข ราโม อิมํ ธมฺมํ เกวลํ สทฺธามตฺตเกน ‘สยํ อภิญฺญา สจฺฉิกตฺวา อุปสมฺปชฺช วิหรามี’ติ ปเวเทสิ, อทฺธา ราโม อิมํ ธมฺมํ ชานํ ปสฺสํ วิหาสี”ติ.\csromanspacer{} อถ ขฺวาหํ ภารทฺวาช เยน อุทโก รามปุตฺโต เตนุปสงฺกมิํ, อุปสงฺกมิตฺวา อุทกํ รามปุตฺตํ เอตทโวจํ “กิตฺตาวตา โน อาวุโส ราโม อิมํ ธมฺมํ สยํ อภิญฺญา สจฺฉิกตฺวา อุปสมฺปชฺช วิหรามีติ ปเวเทสี”ติ.\csromanspacer{} เอวํ วุตฺเต ภารทฺวาช อุทโก รามปุตฺโต เนวสญฺญานาสญฺญายตนํ ปเวเทสิ.\csromanspacer{} ตสฺส มยฺหํ ภารทฺวาช เอตทโหสิ “น โข รามสฺเสว อโหสิ สทฺธา, มยฺหํปตฺถิ สทฺธา.\csromanspacer{} น โข รามสฺเสว อโหสิ วีริยํ ฯเปฯ สติ.\csromanspacer{} สมาธิ.\csromanspacer{} ปญฺญา, มยฺหํปตฺถิ ปญฺญา.\csromanspacer{} ยํนูนาหํ ยํ ธมฺมํ ราโม สยํ อภิญฺญา สจฺฉิกตฺวา}

\csromanheader{2}{10. สงฺคารวสุตฺต (100)}
\prose{\csromanfolio{429}อุปสมฺปชฺช วิหรามีติ ปเวเทสิ, ตสฺส ธมฺมสฺส สจฺฉิกิริยาย ปทเหยฺยนฺ”ติ.\csromanspacer{} โส โข อหํ ภารทฺวาช นจิรสฺเสว ขิปฺปเมว ตํ ธมฺมํ สยํ อภิญฺญา สจฺฉิกตฺวา อุปสมฺปชฺช วิหาสิํ.}

\prose{อถ ขฺวาหํ ภารทฺวาช เยน อุทโก รามปุตฺโต เตนุปสงฺคมิํ, อุปสงฺกมิตฺวา อุทกํ รามปุตฺตํ เอตทโวจํ “เอตฺตาวตา โน อาวุโส ราโม อิมํ ธมฺมํ สยํ อภิญฺญา สจฺฉิกตฺวา อุปสมฺปชฺช ปเวเทสี”ติ.\csromanspacer{} เอตฺตาวตา โข อาวุโส ราโม อิมํ ธมฺมํ สยํ อภิญฺญา สจฺฉิกตฺวา อุปสมฺปชฺช ปเวเทสีติ.\csromanspacer{} อหมฺปิ โข อาวุโส เอตฺตาวตา อิมํ ธมฺมํ สยํ อภิญฺญา สจฺฉิกตฺวา อุปสมฺปชฺช วิหรามีติ.\csromanspacer{} ลาภา โน อาวุโส, สุลทฺธํ โน อาวุโส, เย มยํ อายสฺมนฺตํ ตาทิสํ สพฺรหฺมจาริํ ปสฺสาม.\csromanspacer{} อิติ ยํ ธมฺมํ ราโม สยํ อภิญฺญา สจฺฉิกตฺวา อุปสมฺปชฺช ปเวเทสิ, ตํ ตฺวํ ธมฺมํ สยํ อภิญฺญา สจฺฉิกตฺวา อุปสมฺปชฺช วิหรสิ.\csromanspacer{} ยํ ตฺวํ ธมฺมํ สยํ อภิญฺญา สจฺฉิกตฺวา อุปสมฺปชฺช วิหรสิ, ตํ ธมฺมํ ราโม สยํ อภิญฺญา สจฺฉิกตฺวา อุปสมฺปชฺช ปเวเทสิ.\csromanspacer{} อิติ ยํ ธมฺมํ ราโม อภิญฺญาสิ, ตํ ตฺวํ ธมฺมํ ชานาสิ.\csromanspacer{} ยํ ตฺวํ ธมฺมํ ชานาสิ, ตํ ธมฺมํ ราโม อภิญฺญาสิ.\csromanspacer{} อิติ ยาทิโส ราโม อโหสิ, ตาทิโส ตุวํ.\csromanspacer{} ยาทิโส ตุวํ, ตาทิโส ราโม อโหสิ.\csromanspacer{} เอหิ ทานิ อาวุโส ตุวํ อิมํ คณํ ปริหราติ.\csromanspacer{} อิติ โข ภารทฺวาช อุทโก รามปุตฺโต สพฺรหฺมจารี เม สมาโน อาจริยฏฺฐาเน มํ ฐเปสิ, อุฬาราย จ มํ ปูชาย ปูเชสิ.\csromanspacer{} ตสฺส มยฺหํ ภารทฺวาช เอตทโหสิ “นายํ ธมฺโม นิพฺพิทาย น วิราคาย น นิโรธาย น อุปสมาย น อภิญฺญาย น สมฺโพธาย น นิพฺพานาย สํวตฺตติ, ยาวเทว เนวสญฺญานาสญฺญายตนูปปตฺติยา”ติ.\csromanspacer{} โส โข อหํ ภารทฺวาช ตํ ธมฺมํ อนลงฺกริตฺวา ตสฺมา ธมฺมา นิพฺพิชฺช อปกฺกมิํ.}

\csromanmatikaanchor{mtk.59}
\csromantocmarkref{subsection}{10. สงฺคารวสุตฺต}{mtk.59}
\bookmark[dest={mtk.59},level=2]{10. สงฺคารวสุตฺต}
\proseitem{477}{โส โข อหํ ภารทฺวาช กิํกุสลคเวสี อนุตฺตรํ สนฺติวรปทํ ปริเยสมาโน มคเธสุ อนุปุพฺเพน จาริกํ จรมาโน เยน อุรุเวฬา เสนานิคโม ตทวสริํ.\csromanspacer{} ตตฺถทฺทสํ รมณียํ ภูมิภาคํ ปาสาทิกญฺจ วนสณฺฑํ นทิญฺจ สนฺทนฺติํ เสตกํ สุปติตฺถํ รมณียํ สมนฺตา จ โคจรคามํ.\csromanspacer{} ตสฺส มยฺหํ ภารทฺวาช เอตทโหสิ “รมณีโย วต โภ ภูมิภาโค, ปาสาทิโก จ วนสณฺโฑ, นที จ สนฺทติ เสตกา สุปติตฺถา รมณียา, สมนฺตา จ โคจรคาโม, ‘อลํ วติทํ กุลปุตฺตสฺส ปธานตฺถิกสฺส \csromanfolio{430}ปธานายา’ติ”.\csromanspacer{} โส โข อหํ ภารทฺวาช ตตฺเถว นิสีทิํ “อลมิทํ ปธานายา”ติ.\csromanspacer{} อปิสฺสุ มํ ภารทฺวาช ติสฺโส อุปมา ปฏิภํสุ อนจฺฉริยา ปุพฺเพ อสฺสุตปุพฺพา\csromandash{}}

\prose{เสยฺยถาปิ ภารทฺวาช อลฺลํ กฏฺฐํ สสฺเนหํ อุทเก นิกฺขิตฺตํ.\csromanspacer{} อถ ปุริโส อาคจฺเฉยฺย อุตฺตรารณิํ อาทาย “อคฺคิํ อภินิพฺพตฺเตสฺสามิ เตโช ปาตุกริสฺสามี”ติ.\csromanspacer{} ตํ กิํ มญฺญสิ ภารทฺวาช, อปิ นุ โส ปุริโส อมุํ อลฺลํ กฏฺฐํ สสฺเนหํ อุทเก นิกฺขิตฺตํ อุตฺตรารณิํ อาทาย อภิมนฺเถนฺโต อคฺคิํ อภินิพฺพตฺเตยฺย เตโช ปาตุกเรยฺยาติ.\csromanspacer{} โน หิทํ โภ โคตม.\csromanspacer{} ตํ กิสฺส เหตุ, อทุํ หิ โภ โคตม อลฺลํ กฏฺฐํ สสฺเนหํ, ตญฺจ ปน อุทเก นิกฺขิตฺตํ, ยาวเทว จ ปน โส ปริโส กิลมถสฺส วิฆาตสฺส ภาคี อสฺสาติ.\csromanspacer{} เอวเมว โข ภารทฺวาช เย หิ เกจิ สมณา วา พฺราหฺมณา วา กาเยน เจว จิตฺเตน จ กาเมหิ อวูปกฏฺฐา วิหรนฺติ.\csromanspacer{} โย จ เนสํ กาเมสุ กามจฺฉนฺโท กามสฺเนโห กามมุจฺฉา กามปิปาสา กามปริฬาโห, โส จ อชฺฌตฺตํ น สุปฺปหีโน โหติ น สุปฺปฏิปฺปสฺสทฺโธ.\csromanspacer{} โอปกฺกมิกา เจปิ เต โภนฺโต สมณพฺราหฺมณา ทุกฺขา ติพฺพา ขรา กฏุกา เวทนา เวทยนฺติ, อภพฺพาว เต ญาณาย ทสฺสนาย อนุตฺตราย สมฺโพธาย.\csromanspacer{} โน เจปิ เต โภนฺโต สมณพฺราหฺมณา โอปกฺกมิกา ทุกฺขา ติพฺพา ขรา กฏุกา เวทนา เวทยนฺติ, อภพฺพาว เต ญาณาย ทสฺสนาย อนุตฺตราย สมฺโพธาย.\csromanspacer{} อยํ โข มํ ภารทฺวาช ปฐมา อุปมา ปฏิภาสิ อนจฺฉริยา ปุพฺเพ อสฺสุตปุพฺพา.}

\proseitem{478}{อปราปิ โข มํ ภารทฺวาช ทุติยา อุปมา ปฏิภาสิ อนจฺฉริยา ปุพฺเพ อสฺสุตปุพฺพา\csromandash{}}

\prose{เสยฺยถาปิ ภารทฺวาช อลฺลํ กฏฺฐํ สสฺเนหํ อารกา อุทกา ถเล นิกฺขิตฺตํ.\csromanspacer{} อถ ปุริโส อาคจฺเฉยฺย อุตฺตรารณิํ อาทาย “อคฺคิํ อภินิพฺพตฺเตสฺสามิ เตโช ปาตุกริสฺสามี”ติ.\csromanspacer{} ตํ กิํ มญฺญสิ ภารทฺวาช, อปิ นุ โส ปุริโส อมุํ อลฺลํ กฏฺฐํ สสฺเนหํ อารกา อุทกา ถเล นิกฺขิตฺตํ อุตฺตรารณิํ อาทาย อภิมนฺเถนฺโต อคฺคิํ อภินิพฺพตฺเตยฺย เตโช ปาตุกเรยฺยาติ.\csromanspacer{} โน หิทํ โภ โคตม.\csromanspacer{} ตํ กิสฺส เหตุ, อทุํ หิ โภ โคตม อลฺลํ กฏฺฐํ สสฺเนหํ กิญฺจาปิ อารกา อุทกา ถเล นิกฺขิตฺตํ, ยาวเทว จ ปน โส ปุริโส กิลมถสฺส}

\csromanheader{2}{10. สงฺคารวสุตฺต (100)}
\prose{\csromanfolio{431}วิฆาตสฺส ภาคี อสฺสาติ.\csromanspacer{} เอวเมว โข ภารทฺวาช เย หิ เกจิ สมณา วา พฺราหฺมณา วา กาเยน เจว จิตฺเตน จ กาเมหิ วูปกฏฺฐา วิหรนฺติ.\csromanspacer{} โย จ เนสํ กาเมสุ กามจฺฉนฺโท กามสฺเนโห กามมุจฺฉา กามปิปาสา กามปริฬาโห, โส จ อชฺฌตฺตํ น สุปฺปหีโน โหติ น สุปฺปฏิปฺปสฺสทฺโธ.\csromanspacer{} โอปกฺกมิกา เจปิ เต โภนฺโต สมณพฺราหฺมณา ทุกฺขา ติพฺพา ขรา กฏุกา เวทนา เวทยนฺติ, อภพฺพาว เต ญาณาย ทสฺสนาย อนุตฺตราย สมฺโพธาย.\csromanspacer{} โน เจปิ เต โภนฺโต สมณพฺราหฺมณา โอปกฺกมิกา ทุกฺขา ติพฺพา ขรา กฏุกา เวทนา เวทยนฺติ, อภพฺพาว เต ญาณาย ทสฺสนาย อนุตฺตราย สมฺโพธาย.\csromanspacer{} อยํ โข มํ ภารทฺวาช ทุติยา อุปมา ปฏิภาสิ อนจฺฉริยา ปุพฺเพ อสฺสุตปุพฺพา.}

\proseitem{479}{อปราปิ โข มํ ภารทฺวาช ตติยา อุปมา ปฏิภาสิ อนจฺฉริยา ปุพฺเพ อสฺสุตปุพฺพา\csromandash{}}

\prose{เสยฺยถาปิ ภารทฺวาช สุกฺขํ กฏฺฐํ โกฬาปํ อารกา อุทกา ถเล นิกฺขิตฺตํ.\csromanspacer{} อถ ปุริโส อาคจฺเฉยฺย อุตฺตรารณิํ อาทาย “อคฺคิํ อภินิพฺพตฺเตสฺสามิ เตโช ปาตุกริสฺสามี”ติ.\csromanspacer{} ตํ กิํ มญฺญสิ ภารทฺวาช, อปิ นุ โส ปุริโส อมุํ สุกฺขํ กฏฺฐํ โกฬาปํ อารกา อุทกา ถเล นิกฺขิตฺตํ อุตฺตรารณิํ อาทาย อภิมนฺเถนฺโต อคฺคิํ อภินิพฺพตฺเตยฺย เตโช ปาตุกเรยฺยาติ.\csromanspacer{} เอวํ โภ โคตม.\csromanspacer{} ตํ กิสฺส เหตุ, อทุํ หิ โภ โคตม สุกฺขํ กฏฺฐํ โกฬาปํ, ตญฺจ ปน อารกา อุทกา ถเล นิกฺขิตฺตนฺติ.\csromanspacer{} เอวเมว โข ภารทฺวาช เย หิ เกจิ สมณา วา พฺราหฺมณา วา กาเยน เจว จิตฺเตน จ กาเมหิ วูปกฏฺฐา วิหรนฺติ.\csromanspacer{} โย จ เนสํ กาเมสุ กามจฺฉนฺโท กามสฺเนโห กามมุจฺฉา กามปิปาสา กามปริฬาโห, โส จ อชฺฌตฺตํ สุปฺปหีโน โหติ สุปฺปฏิปฺปสฺสทฺโธ.\csromanspacer{} โอปกฺกมิกา เจปิ เต โภนฺโต สมณพฺราหฺมณา ทุกฺขา ติพฺพา ขรา กฏุกา เวทนา เวทยนฺติ, ภพฺพาว เต ญาณาย ทสฺสนาย อนุตฺตราย สมฺโพธาย.\csromanspacer{} โน เจปิ เต โภนฺโต สมณพฺราหฺมณา โอปกฺกมิกา ทุกฺขา ติพฺพา ขรา กฏุกา เวทนา เวทยนฺติ, ภพฺพาว เต ญาณาย ทสฺสนาย อนุตฺตราย สมฺโพธาย.\csromanspacer{} อยํ โข มํ ภารทฺวาช ตติยา อุปมา ปฏิภาสิ อนจฺฉริยา ปุพฺเพ อสฺสุตปุพฺพา.\csromanspacer{} อิมา โข มํ ภารทฺวาช ติสฺโส อุปมา ปฏิภํสุ อนจฺฉริยา ปุพฺเพ อสฺสุตปุพฺพา.}

\proseitem{480}{\csromanfolio{432}ตสฺส มยฺหํ ภารทฺวาช เอตทโหสิ “ยํนูนาหํ ทนฺเตภิทนฺตมาธาย ชิวฺหาย ตาลุํ อาหจฺจ เจตสา จิตฺตํ อภินิคฺคเณฺหยฺยํ อภินิปฺปีเฬยฺยํ อภิสนฺตาเปยฺยนฺ”ติ, โส โข อหํ ภารทฺวาช ทนฺเตภิทนฺตมาธาย ชิวฺหาย ตาลุํ อาหจฺจ เจตสา จิตฺตํ อภินิคฺคณฺหามิ อภินิปฺปีเฬมิ อภิสนฺตาเปมิ.\csromanspacer{} ตสฺส มยฺหํ ภารทฺวาช ทนฺเตภิทนฺตมาธาย ชิวฺหาย ตาลุํ อาหจฺจ เจตสา จิตฺตํ อภินิคฺคณฺหโต อภินิปฺปีฬยโต อภิสนฺตาปยโต กจฺเฉหิ เสทา มุจฺจนฺติ.\csromanspacer{} เสยฺยถาปิ ภารทฺวาช พลวา ปุริโส ทุพฺพลตรํ ปุริสํ สีเส วา คเหตฺวา ขนฺเธ วา คเหตฺวา อภินิคฺคเณฺหยฺย อภินิปฺปีเฬยฺย อภิสนฺตาเปยฺย.\csromanspacer{} เอวเมว โข เม ภารทฺวาช ทนฺเตภิทนฺตมาธาย ชิวฺหาย ตาลุํ อาหจฺจ เจตสา จิตฺตํ อภินิคฺคณฺหโต อภินิปฺปีฬยโต อภิสนฺตาปยโต กจฺเฉหิ เสทา มุจฺจนฺติ.\csromanspacer{} อารทฺธํ โข ปน เม ภารทฺวาช วีริยํ โหติ อสลฺลีนํ, อุปฏฺฐิตา สติ อสมฺมุฏฺฐา, สารทฺโธ จ ปน เม กาโย โหติ อปฺปฏิปฺปสฺสทฺโธ เตเนว ทุกฺขปฺปธาเนน ปธานาภิตุนฺนสฺส สโต.}

\proseitem{481}{ตสฺส มยฺหํ ภารทฺวาช เอตทโหสิ “ยํนูนาหํ อปฺปาณกํเยว ฌานํ ฌาเยยฺยนฺ”ติ, โส โข อหํ ภารทฺวาช มุขโต จ นาสโต จ อสฺสาสปสฺสาเส อุปรุนฺธิํ.\csromanspacer{} ตสฺส มยฺหํ ภารทฺวาช มุขโต จ นาสโต จ อสฺสาสปสฺสาเสสุ อุปรุทฺเธสุ กณฺณโสเตหิ วาตานํ นิกฺขมนฺตานํ อธิมตฺโต สทฺโท โหติ.\csromanspacer{} เสยฺยถาปิ นาม กมฺมารคคฺคริยา ธมมานาย อธิมตฺโต สทฺโท โหติ.\csromanspacer{} เอวเมว โข เม ภารทฺวาช มุขโต จ นาสโต จ อสฺสาสปสฺสาเสสุ อุปรุทฺเธสุ กณฺณโสเตหิ วาตานํ นิกฺขมนฺตานํ อธิมตฺโต สทฺโท โหติ.\csromanspacer{} อารทฺธํ โข ปน เม ภารทฺวาช วีริยํ โหติ อสลฺลีนํ, อุปฏฺฐิตา สติ อสมฺมุฏฺฐา, สารทฺโธ จ ปน เม กาโย โหติ อปฺปฏิปฺปสฺสทฺโธ เตเนว ทุกฺขปฺปธาเนน ปธานาภิตุนฺนสฺส สโต.}

\prose{ตสฺส มยฺหํ ภารทฺวาช เอตทโหสิ “ยํนูนาหํ อปฺปาณกํเยว ฌานํ ฌาเยยฺยนฺ”ติ.\csromanspacer{} โส โข อหํ ภารทฺวาช มุขโต จ นาสโต จ กณฺณโต จ อสฺสาสปสฺสาเส อุปรุนฺธิํ.\csromanspacer{} ตสฺส มยฺหํ ภารทฺวาช มุขโต จ นาสโต จ กณฺณโต จ อสฺสาสปสฺสาเสสุ อุปรุทฺเธสุ อธิมตฺตา วาตา มุทฺธนิ อูหนนฺติ.\csromanspacer{} เสยฺยถาปิ ภารทฺวาช พลวา ปุริโส ติเณฺหน สิขเรน มุทฺธนิ อภิมตฺเถยฺย.\csromanspacer{} เอวเมว โข เม ภารทฺวาช มุขโต จ}

\csromanheader{2}{10. สงฺคารวสุตฺต (100)}
\prose{\csromanfolio{433}นาสโต จ กณฺณโต จ อสฺสาสปสฺสาเสสุ อุปรุทฺเธสุ อธิมตฺตา วาตา มุทฺธนิ อูหนนฺติ.\csromanspacer{} อารทฺธํ โข ปน เม ภารทฺวาช วีริยํ โหติ อสลฺลีนํ, อุปฏฺฐิตา สติ อสมฺมุฏฺฐา, สารทฺโธ จ ปน เม กาโย โหติ อปฺปฏิปฺปสฺสทฺโธ เตเนว ทุกฺขปฺปธาเนน ปธานาภิตุนฺนสฺส สโต.}

\prose{ตสฺส มยฺหํ ภารทฺวาช เอตทโหสิ “ยํนูนาหํ อปฺปาณกํเยว ฌานํ ฌาเยยฺยนฺ”ติ.\csromanspacer{} โส โข อหํ ภารทฺวาช มุขโต จ นาสโต จ กณฺณโต จ อสฺสาสปสฺสาเส อุปรุนฺธิํ, ตสฺส มยฺหํ ภารทฺวาช มุขโต จ นาสโต จ กณฺณโต จ อสฺสาสปสฺสาเสสุ อุปรุทฺเธสุ อธิมตฺตา สีเส สีสเวทนา โหนฺติ.\csromanspacer{} เสยฺยถาปิ ภารทฺวาช พลวา ปุริโส ทเฬฺหน วรตฺตกฺขณฺเฑน สีเส สีสเวฐํ ทเทยฺย.\csromanspacer{} เอวเมว โข ภารทฺวาช มุขโต จ นาสโต จ กณฺณโต จ อสฺสาสปสฺสาเสสุ อุปรุทฺเธสุ อธิมตฺตา สีเส สีสเวทนา โหนฺติ.\csromanspacer{} อารทฺธํ โข ปน เม ภารทฺวาช วีริยํ โหติ อสลฺลีนํ, อุปฏฺฐิตา สติ อสมฺมุฏฺฐา, สารทฺโธ จ ปน เม กาโย โหติ อปฺปฏิปฺปสฺสทฺโธ เตเนว ทุกฺขปฺปธาเนน ปธานาภิตุนฺนสฺส สโต.}

\prose{ตสฺส มยฺหํ ภารทฺวาช เอตทโหสิ “ยํนูนาหํ อปฺปาณกํเยว ฌานํ ฌาเยยฺยนฺ”ติ.\csromanspacer{} โส โข อหํ ภารทฺวาช มุขโต จ นาสโต จ กณฺณโต จ อสฺสาสปสฺสาเส อุปรุนฺธิํ, ตสฺส มยฺหํ ภารทฺวาช มุขโต จ นาสโต จ กณฺณโต จ อสฺสาสปสฺสาเสสุ อุปรุทฺเธสุ อธิมตฺตา วาตา กุจฺฉิํ ปริกนฺตนฺติ.\csromanspacer{} เสยฺยถาปิ ภารทฺวาช ทกฺโข โคฆาตโก วา โคฆาตกนฺเตวาสี วา ติเณฺหน โควิกนฺตเนน กุจฺฉิํ ปริกนฺเตยฺย.\csromanspacer{} เอวเมว โข เม ภารทฺวาช มุขโต จ นาสโต จ กณฺณโต จ อสฺสาสปสฺสาเสสุ อุปรุทฺเธสุ อธิมตฺตา วาตา กุจฺฉิํ ปริกนฺตนฺติ.\csromanspacer{} อารทฺธิํ โข ปน เม ภารทฺวาช วีริยํ โหติ อสลฺลีนํ, อุปฏฺฐิตา สติ อสมฺมุฏฺฐา, สารทฺโธ จ ปน เม กาโย โหติ อปฺปฏิปฺปสฺสทฺโธ เตเนว ทุกฺขปฺปธาเนน ปธานาภิตุนฺนสฺส สโต.}

\prose{ตสฺส มยฺหํ ภารทฺวาช เอตทโหสิ “ยํนูนาหํ อปฺปาณกํเยว ฌานํ ฌาเยยฺยนฺ”ติ.\csromanspacer{} โส โข อหํ ภารทฺวาช มุขโต จ นาสโต จ กณฺณโต จ อสฺสาสปสฺสาเส อุปรุนฺธิํ, ตสฺส มยฺหํ ภารทฺวาช มุขโต}

\prose{\csromanfolio{434}จ นาสโต จ กณฺณโต จ อสฺสาสปสฺสาเสสุ อุปรุทฺเธสุ อธิมตฺโต กายสฺมิํ ฑาโห โหติ.\csromanspacer{} เสยฺยถาปิ ภารทฺวาช เทฺว พลวนฺโต ปุริสา ทุพฺพลตรํ ปุริสํ นานาพาหาสุ คเหตฺวา องฺคารกาสุยา สนฺตาเปยฺยุํ สมฺปริตาเปยฺยุํ.\csromanspacer{} เอวเมว โข เม ภารทฺวาช มุขโต จ นาสโต จ กณฺณโต จ อสฺสาสปสฺสาเสสุ อุปรุทฺเธสุ อธิมตฺโต กายสฺมิํ ฑาโห โหติ.\csromanspacer{} อารทฺธํ โข ปน เม ภารทฺวาช วีริยํ โหติ อสลฺลีนํ, อุปฏฺฐิตา สติ อสมฺมุฏฺฐา, สารทฺโธ จ ปน เม กาโย โหติ อปฺปฏิปฺปสฺสทฺโธ เตเนว ทุกฺขปฺปธาเนน ปธานาภิตุนฺนสฺส สโต.\csromanspacer{} อปิสฺสุ มํ ภารทฺวาช เทวตา ทิสฺวา เอวมาหํสุ “กาลงฺกโต สมโณ โคตโม”ติ.\csromanspacer{} เอกจฺจา เทวตา เอวมาหํสุ “น กาลงฺกโต สมโณ โคตโม, อปิ จ กาลงฺกโรตี”ติ.\csromanspacer{} เอกจฺจา เทวตา เอวมาหํสุ “น กาลงฺกโต สมโณ โคตโม, นาปิ กาลงฺกโรติ, อรหํ สมโณ โคตโม, วิหาโรเตฺวว โส อรหโต เอวรูโป โหตี”ติ.}

\prose{ตสฺส มยฺหํ ภารทฺวาช เอตทโหสิ “ยํนูนาหํ สพฺพโส อาหารุปจฺเฉทาย ปฏิปชฺเชยฺยนฺ”ติ.\csromanspacer{} อถ โข มํ ภารทฺวาช เทวตา อุปสงฺกมิตฺวา เอตทโวจุํ “มา โข ตฺวํ มาริส สพฺพโส อาหารุปจฺเฉทาย ปฏิปชฺชิ.\csromanspacer{} สเจ โข ตฺวํ มาริส สพฺพโส อาหารุปจฺเฉทาย ปฏิปชฺชิสฺสสิ, ตสฺส เต มยํ ทิพฺพํ โอชํ โลมกูเปหิ อชฺโฌหาเรสฺสาม, ตาย ตฺวํ ยาเปสฺสสี”ติ.\csromanspacer{} ตสฺส มยฺหํ ภารทฺวาช เอตทโหสิ “อหญฺเจว โข ปน สพฺพโส อชชฺชิตํ ปฏิชาเนยฺยํ, อิมา จ เทวตา ทิพฺพํ โอชํ โลมกูเปหิ อชฺโฌหาเรยฺยุํ, ตาย จาหํ ยาเปยฺยํ, ตํ มมสฺส มุสา”ติ.\csromanspacer{} โส โข อหํ ภารทฺวาช ตา เทวตา ปจฺจาจิกฺขามิ หลนฺติ วทามิ.}

\prose{ตสฺส มยฺหํ ภารทฺวาช เอตทโหสิ “ยํนูนาหํ โถกํ โถกํ อาหารํ อาหาเรยฺยํ ปสตํ ปสตํ, ยทิ วา มุคฺคยูสํ ยทิ วา กุลตฺถยูสํ ยทิ วา กฬายยูสํ ยทิ วา หเรณุกยูสนฺ”ติ.\csromanspacer{} โส โข อหํ ภารทฺวาช โถกํ โถกํ อาหารํ อาหาเรสิํ ปสตํ ปสตํ, ยทิ วา มุคฺคยูสํ ยทิ วา กุลตฺถยูสํ ยทิ วา กฬายยูสํ ยทิ วา หเรณุกยูสํ.\csromanspacer{} ตสฺส มยฺหํ ภารทฺวาช โถกํ โถกํ}

\csromanheader{2}{10. สงฺคารวสุตฺต (100)}
\prose{\csromanfolio{435}อาหารํ อาหารยโต ปสตํ ปสตํ, ยทิ วา มุคฺคยูสํ ยทิ วา กุลตฺถยูสํ ยทิ วา กฬายยูสํ ยทิ วา หเรณุกยูสํ อธิมตฺตกสิมานํ ปตฺโต กาโย โหติ, เสยฺยถาปิ นาม อาสีติกปพฺพานิ วา กาฬปพฺพานิ วา.\csromanspacer{} เอวเมวสฺสุ เม องฺคปจฺจงฺคานิ ภวนฺติ ตาเยวปฺปาหารตาย.\csromanspacer{} เสยฺยถาปิ นาม โอฏฺฐปทํ.\csromanspacer{} เอวเมวสฺสุ เม อานิสทํ โหติ ตาเยวปฺปาหารตาย.\csromanspacer{} เสยฺยถาปิ นาม วฏฺฏนาวฬี.\csromanspacer{} เอวเมวสฺสุ เม ปิฏฺฐิกณฺฏโก อุณฺณตาวนโต โหติ ตาเยวปฺปาหารตาย.\csromanspacer{} เสยฺยถาปิ นาม ชรสาลาย โคปานสิโย โอลุคฺควิลุคฺคา ภวนฺติ.\csromanspacer{} เอวเมวสฺสุ เม ผาสุฬิโย โอลุคฺควิลุคฺคา ภวนฺติ ตาเยวปฺปาหารตาย.\csromanspacer{} เสยฺยถาปิ นาม คมฺภีเร อุทปาเน อุทกตารกา คมฺภีรคตา โอกฺขายิกา ทิสฺสนฺติ.\csromanspacer{} เอวมวสฺสุ เม อกฺขิกูเปสุ อกฺขิตารกา คมฺภีรคตา โอกฺขายิกา ทิสฺสนฺติ ตาเยวปฺปาหารตาย.\csromanspacer{} เสยฺยถาปิ นาม ติตฺตกาลาพุ อามกจฺฉินฺโน วาตาตเปน สํผุฏิโต โหติ สมฺมิลาโต.\csromanspacer{} เอวเมวสฺสุ เม สีสจฺฉวิ สํผุฏิตา โหติ สมฺมิลาตา ตาเยวปฺปาหารตาย.\csromanspacer{} โส โข อหํ ภารทฺวาช “อุทรจฺฉวิํ ปริมสิสฺสามี”ติ ปิฏฺฐิกณฺฏกํเยว ปริคฺคณฺหามิ, “ปิฏฺฐิกณฺฏกํ ปริมสิสฺสามี”ติ อุทรจฺฉวิํเยว ปริคฺคณฺหามิ.\csromanspacer{} ยาวสฺสุ เม ภารทฺวาช อุทรจฺฉวิ ปิฏฺฐิกณฺฏกํ อลฺลีนา โหติ ตาเยวปฺปาหารตาย.\csromanspacer{} โส โข อหํ ภารทฺวาช “วจฺจํ วา มุตฺตํ วา กริสฺสามี”ติ ตตฺเถว อวกุชฺโช ปปตามิ ตาเยวปฺปาหารตาย.\csromanspacer{} โส โข อหํ ภารทฺวาช อิมเมว กายํ อสฺสาเสนฺโต ปาณินา คตฺตานิ อนุมชฺชามิ, ตสฺส มยฺหํ ภารทฺวาช ปาณินา คตฺตานิ อนุมชฺชโต ปูติมูลานิ โลมานิ กายสฺมา ปปตนฺติ ตาเยวปฺปาหารตาย.\csromanspacer{} อปิสฺสุ มํ ภารทฺวาช มนุสฺสา ทิสฺวา เอวมาหํสุ “กาโฬ สมโณ โคตโม”ติ.\csromanspacer{} เอกจฺเจ มนุสฺสา เอวมาหํสุ “น กโฬ สมโณ โคตโม, สาโม สมโณ โคตโม”ติ.\csromanspacer{} เอกจฺเจ มนุสฺสา เอวมาหํสุ “น กาโฬ สมโณ โคตโม, นปิ สาโม, มงฺคุรจฺฉวิ สมโณ โคตโม”ติ.\csromanspacer{} ยาวสฺสุ เม ภารทฺวาช ตาว ปริสุทฺโธ ฉวิวณฺโณ ปริโยทาโต อุปหโต โหติ ตาเยวปฺปาหารตาย.}

\proseitem{482}{ตสฺส มยฺหํ ภารทฺวาช เอตทโหสิ “เย โข เกจิ อตีตมทฺธานํ สมณา วา พฺราหฺมณา วา โอปกฺกมิกา ทุกฺขา ติพฺพา ขรา กฏุกา เวทนา \csromanfolio{436}เวทยิํสุ, เอตาวปรมํ นยิโต ภิยฺโย.\csromanspacer{} เยปิ หิ เกจิ อนาคตมทฺธานํ สมณา วา พฺราหฺมณา วา โอปกฺกมิกา ทุกฺขา ติพฺพา ขรา กฏุกา เวทนา เวทยิสฺสนฺติ, เอตาวปรมํ นยิโต ภิยฺโย.\csromanspacer{} เยปิ หิ เกจิ เอตรหิ สมณา วา พฺราหฺมณา วา โอปกฺกมิกา ทุกฺขา ติพฺพา ขรา กฏุกา เวทนา เวทยนฺติ, เอตาวปรมํ นยิโต ภิยฺโย.\csromanspacer{} น โข ปนาหํ อิมาย กฏุกาย ทุกฺกรการิกาย อธิคจฺฉามิ อุตฺตริ มนุสฺสธมฺมา อลมริยญาณทสฺสนวิเสสํ.\csromanspacer{} สิยา นุ โข อญฺโญ มคฺโค โพธายา”ติ.\csromanspacer{} ตสฺส มยฺหํ ภารทฺวาช เอตทโหสิ “อภิชานามิ โข ปนาหํ ปิตุ อ กมฺมนฺเต สีตาย ชมฺพุจฺฉายาย นิสินฺโน วิวิจฺเจว กาเมหิ วิวิจฺจ อกุสเลหิ ธมฺเมหิ สวิตกฺกํ สวิจารํ วิเวกชํ ปีติสุขํ ปฐมํ ฌานํ อุปสมฺปชฺช วิหริตา.\csromanspacer{} สิยา นุ โข เอโส มคฺโค โพธายา”ติ.\csromanspacer{} ตสฺส มยฺหํ ภารทฺวาช สตานุสาริ วิญฺญาณํ อโหสิ “เอเสว มคฺโค โพธายา”ติ.\csromanspacer{} ตสฺส มยฺหํ ภารทฺวาช เอตทโหสิ “กิํ นุ โข อหํ ตสฺส สุขสฺส ภายามิ, ยํ ตํ สุขํ อญฺญเตฺรว กาเมหิ อญฺญตฺร อกุสเลหิ ธมฺเมหี”ติ.\csromanspacer{} ตสฺส มยฺหํ ภารทฺวาช เอตทโหสิ “น โข อหํ ตสฺส สุขสฺส ภายามิ, ยํ ตํ สุขํ อญฺญเตฺรว กาเมหิ อญฺญตฺร อกุสเลหิ ธมฺเมหี”ติ.}

\proseitem{483}{ตสฺส มยฺหํ ภารทฺวาช เอตทโหสิ “น โข ตํ สุกรํ สุขํ อธิคนฺตุํ เอวํ อธิมตฺตกสิมานํ ปตฺตกาเยน, ยํนูนาหํ โอฬาริกํ อาหารํ อาหาเรยฺยํ โอทนกุมฺมาสนฺ”ติ.\csromanspacer{} โส โข อหํ ภารทฺวาช โอฬาริกํ อาหารํ อาหาเรสิํ โอทนกุมฺมาสํ.\csromanspacer{} เตน โข ปน มํ ภารทฺวาช สมเยน ปญฺจวคฺคิยา ภิกฺขู ปจฺจุปฏฺฐิตา โหนฺติ “ยํ โข สมโณ โคตโม ธมฺมํ อธิคมิสฺสติ, ตํ โน อาโรเจสฺสตี”ติ.\csromanspacer{} ยโต โข อหํ ภารทฺวาช โอฬาริกํ อาหารํ อาหาเรสิํ โอทนกุมฺมาสํ, อถ เม เต ปญฺจวคฺคิยา ภิกฺขู นิพฺพิชฺช ปกฺกมิํสุ “พาหุลฺลิโก สมโณ โคตโม ปธานวิพฺภนฺโต อาวตฺโต พาหุลฺลายา”ติ.}

\prose{โส โข อหํ ภารทฺวาช โอฬาริกํ อาหารํ อาหาเรตฺวา พลํ คเหตฺวา วิวิจฺเจว กาเมหิ ฯเปฯ ปฐมํ ฌานํ อุปสมฺปชฺช วิหาสิํ.\csromanspacer{} วิตกฺกวิจารานํ วูปสมา อชฺฌตฺตํ สมฺปสาทนํ เจตโส เอโกทิภาวํ อวิตกฺกํ}

\csromanheader{2}{10. สงฺคารวสุตฺต (100)}
\prose{\csromanfolio{437}อวิจารํ สมาธิชํ ปีติสุขํ ทุติยํ ฌานํ.\csromanspacer{} ตติยํ ฌานํ.\csromanspacer{} จตุตฺถํ ฌานํ อุปสมฺปชฺช วิหาสิํ.}

\prose{โส เอวํ สมาหิเต จิตฺเต ปริสุทฺเธ ปริโยทาเต อนงฺคเณ วิคตูปกฺกิเลเส มุทุภูเต กมฺมนิเย ฐิเต อาเนญฺชปฺปตฺเต ปุพฺเพนิวาสานุสฺสติญาณาย จิตฺตํ อภินินฺนาเมสิํ.\csromanspacer{} โส อเนกวิหิตํ ปุพฺเพนิวาสํ อนุสฺสรามิ.\csromanspacer{} เสยฺยถิทํ, เอกมฺปิ ชาติํ เทฺวปิ ชาติโย ฯเปฯ อิติ สาการํ ส-อุทฺเทสํ อเนกวิหิตํ ปุพฺเพนิวาสํ อนุสฺสรามิ.\csromanspacer{} อยํ โข เม ภารทฺวาช รตฺติยา ปฐเม ยาเม ปฐมา วิชฺชา อธิคตา, อวิชฺชา วิหตา, วิชฺชา อุปฺปนฺนา, ตโม วิหโต, อาโลโก อุปฺปนฺโน, ยถา ตํ อปฺปมตฺตสฺส อาตาปิโน ปหิตตฺตสฺส วิหรโต.}

\proseitem{484}{โส เอวํ สมาหิเต จิตฺเต ปริสุทฺเธ ปริโยทาเต อนงฺคเณ วิคตูปกฺกิเลเส มุทุภูเต กมฺมนิเย ฐิเต อาเนญฺชปฺปตฺเต สตฺตานํ จุตูปปาตญาณาย จิตฺตํ อภินินฺนาเมสิํ.\csromanspacer{} โส ทิพฺเพน จกฺขุนา วิสุทฺเธน อติกฺกนฺตมานุสเกน สตฺเต ปสฺสามิ จวมาเน อุปปชฺชมาเน หีเน ปณีเต สุวณฺเณ ทุพฺพณฺเณ สุคเต ทุคฺคเต, ยถากมฺมูปเค สตฺเต ปชานามิ ฯเปฯ.\csromanspacer{} อยํ โข เม ภารทฺวาช รตฺติยา มชฺฌิเม ยาเม ทุติยา วิชฺชา อธิคตา, อวิชฺชา วิหตา, วิชฺชา อุปฺปนฺนา, ตโม วิหโต, อาโลโก อุปฺปนฺโน, ยถา ตํ อปฺปมตฺตสฺส อาตาปิโน ปหิตตฺถสฺส วิหรโต.}

\prose{โส เอวํ สมาหิเต จิตฺเต ปริสุทฺเธ ปริโยทาเต อนงฺคเณ วิคตูปกฺกิเลเส มุทุภูเต กมฺมนิเย ฐิเต อาเนญฺชปฺปตฺเต อาสวานํ ขยญาณาย จิตฺตํ อภินินฺนาเมสิํ.\csromanspacer{} โส “อิทํ ทุกฺขนฺ”ติ ยถาภูตํ อพฺภญฺญาสิํ, “อยํ ทุกฺขสมุทโย”ติ ยถาภูตํ อพฺภญฺญาสิํ, “อยํ ทุกฺขนิโรโธ”ติ ยถาภูตํ อพฺภญฺญาสิํ, “อยํ ทุกฺขนิโรธคามินี ปฏิปทา”ติ ยถาภูตํ อพฺภญฺญาสิํ. “อิเม อาสวา”ติ ยถาภูตํ อพฺภญฺญาสิํ, “อยํ อาสวสมุทโย”ติ ยถาภูตํ อพฺภญฺญาสิํ, “อยํ อาสวนิโรโธ”ติ ยถาภูตํ อพฺภญฺญาสิํ, “อยํ อาสวนิโรธคามินี ปฏิปทา”ติ ยถาภูตํ อพฺภญฺญาสิํ.\csromanspacer{} ตสฺส เม เอวํ ชานโต เอวํ ปสฺสโต กามาสวาปิ จิตฺตํ วิมุจฺจิตฺถ, ภวาสวาปิ จิตฺตํ วิมุจฺจิตฺถ, \csromanfolio{438}อวิชฺชาสวาปิ จิตฺตํ วิมุจฺจิตฺถ, วิมุตฺตสฺมิํ “วิมุตฺตมฺ”อิติ ญาณํ อโหสิ, “ขีณา ชาติ, วุสิตํ พฺรหฺมจริยํ, กตํ กรณียํ, นาปรํ อิตฺถตฺตายา”ติ อพฺภญฺญาสิํ.\csromanspacer{} อยํ โข เม ภารทฺวาช รตฺติยา ปจฺฉิเม ยาเม ตติยา วิชฺชา อธิคตา, อวิชฺชา วิหตา, วิชฺชา อุปฺปนฺนา, ตโม วิหโต, อาโลโก อุปฺปนฺโน, ยถา ตํ อปฺปมตฺตสฺส อาตาปิโน ปหิตตฺตสฺส วิหรโตติ.}

\proseitem{485}{เอวํ วุตฺเต สงฺคารโว มาณโว ภควนฺตํ เอตทโวจ “อฏฺฐิตวตํ\footnote{อฏฺฐิต วต (สี, สฺยา, กํ, อิ)} โภโต โคตมสฺส ปธานํ อโหสิ, สปฺปุริสวตํ\footnote{สปฺปุริส วต (สี, สฺยา, กํ, อิ)} โภโต โคตมสฺส ปธานํ อโหสิ, ยถา ตํ อรหโต สมฺมาสมฺพุทฺธสฺส.\csromanspacer{} กิํ นุ โข โภ โคตม อตฺถิ เทวา”ติ\footnote{อธิเทวาติ (ก) เอวํ สพฺเพสุ ‘อตฺถิ เทวา’ติปเทสุ.}.\csromanspacer{} ฐานโส เมตํ\footnote{โข ปเนตํ (สฺยา, กํ, ก)} ภารทฺวาช วิทิตํ ยทิทํ อธิเทวาติ\footnote{อตฺถิ เทวาติ (สี, สฺยา, กํ, อิ), อติเทวาติ (?) เอวํ สพฺเพสุ’อธิเทวา’ติปเทสุ.}. “กิํ นุ โข โภ โคตม อตฺถิ เทวา”ติ ปุฏฺโฐ สมาโน “ฐานโส เมตํ ภารทฺวาช วิทิตํ ยทิทํ อธิเทวา”ติ วเทสิ.\csromanspacer{} นนุ โภ โคตม เอวํ สนฺเต ตุจฺฉา มุสา โหตีติ.\csromanspacer{} อตฺถิ เทวาติ ภารทฺวาช ปุฏฺโฐ สมาโน อตฺถิ เทวาติ โย วเทยฺย, ฐานโส เม วิทิตาติ\footnote{ฐานโส วิทิตา เม วิทิตาติ (สี, สฺยา, กํ, อิ), ฐานโส เม วิทิตา อติเทวาติ (?)} โย วเทยฺย, อถ เขฺวตฺถ วินา ปุริเสน เอกํเสน นิฏฺฐํ คนฺตพฺพํ\footnote{คนฺตุํ (ก), คนฺตุํ วา (สฺยา, กํ)} ยทิทํ อตฺถิ เทวาติ.\csromanspacer{} กิสฺส ปน เม ภวํ โคตโม อาทิเกเนว น พฺยากาสีติ\footnote{โคตโม อาทิเกเนว พฺยากาสีติ (ก), โคตโม อตฺถิ เทวาติ น พฺยากาสีติ (?)}.\csromanspacer{} อุจฺเจน สมฺมตํ โข เอตํ ภารทฺวาช โลกสฺมิํ ยทิทํ อตฺถิ เทวาติ.}

\prose{เอวํ วุตฺเต สงฺคารโว มาณโว ภควนฺตํ เอตทโวจ “อภิกฺกนฺตํ โภ โคตม, อภิกฺกนฺตํ โภ โคตม, เสยฺยถาปิ โภ โคตม นิกฺกุชฺชิตํ วา อุกฺกุชฺเชยฺย, ปฏิจฺฉนฺนํ วา วิวเรยฺย, มูฬฺหสฺส วา มคฺคํ อาจิกฺเขยฺย, อนฺธกาเร วา เตลปชฺโชตํ ธาเรยฺย ‘จกฺขุมนฺโต รูปานิ ทกฺขนฺตี’ติ.\csromanspacer{} เอวเมวํ โภตา โคตเมน อเนกปริยาเยน ธมฺโม ปกาสิโต,}

\csromanheader{2}{10. สงฺคารวสุตฺต (100)}
\csromanmatikaanchor{mtk.60}
\csromantocmarkref{subsection}{อุทฺทานคาถาโย}{mtk.60}
\bookmark[dest={mtk.60},level=2]{อุทฺทานคาถาโย}
\prose{\csromanfolio{439}เอสาหํ ภวนฺตํ โคตมํ สรณํ คจฺฉามิ ธมฺมญฺจ ภิกฺขุสํฆญฺจ, อุปากสํ มํ ภวํ โคตโม ธาเรตุ อชฺชตคฺเค ปาณุเปตํ สรณํ คตนฺ”ติ.}

\nitthitam{สงฺคารวสุตฺตํ นิฏฺฐิตํ ทสมํ.}
\nitthitam{\textbf{พฺราหฺมณวคฺโค นิฏฺฐิโต ปญฺจโม.}}
\titlehead{\textbf{ตสฺสุทฺทานํ}}
\csromangathameasure{พฺรหฺมายุ เสลสฺสลายโน, โฆฏมุโข จ พฺราหฺมโณ. \\ จงฺกี เอสุ ธนญฺชานิ, วาเสฏฺโฐ สุภคารโวติ.}
\csromangathagroup{\csromangathastanza{พฺรหฺมายุ เสลสฺสลายโน, โฆฏมุโข จ พฺราหฺมโณ. \\ จงฺกี เอสุ ธนญฺชานิ, วาเสฏฺโฐ สุภคารโวติ.}}

\csromanheader{2}{\textbf{อิทํ วคฺคานมุทฺทานํ}}
\csromangathameasure{วคฺโค คหปติ ภิกฺขุ, ปริพฺพาชกนามโก. \\ ราชวคฺโค พฺราหฺมโณติ, ปญฺจ มชฺฌิม-อาคเม. \\ \textbf{มชฺฌิมปณฺณาสกํ สมตฺตํ.}}
\csromangathagroup{\csromangathastanza{วคฺโค คหปติ ภิกฺขุ, ปริพฺพาชกนามโก. \\ ราชวคฺโค พฺราหฺมโณติ, ปญฺจ มชฺฌิม-อาคเม. \\ \textbf{มชฺฌิมปณฺณาสกํ สมตฺตํ.}}}

\orphannote{ปสฺส ม 1 จูฬหตฺถิปโทปเม (238) ปิฏฺเฐ.}

\orphannote{( ) สี-สฺยา-กํ-อิ-โปตฺถเกสุ นตฺถิ.}

\orphannote{นวจฺฉทนํ (สี)}

\orphannote{“มรเณนปิ เต ฯเปฯ ปพฺพชฺชายา”ติ วากฺยทฺวยํ สี-สฺยา-กํ-อิ-โปตฺถเกสุ ทุติยฏฺฐาเน เยว ทิสฺสติ, ปาราชิกปาฬิยํ ปน ปฐมฏฺฐาเน เยว ทิสฺสติ. ตสฺมา อิธ ทุติยฏฺฐาเน ปุนาคตํ อธิกํ วิย ทิสฺสติ.}

\orphannote{( ) เอตฺถนฺตเร ปาโฐ สี-สฺยา-กํ-อิ-โปตฺถเกสุ นตฺถิ.}

\orphannote{ภวํ สาริปุตฺโต อาหาติ, กตมํ สาริปุตฺโต อาห พฺรหฺมโลโกติ. (ก)}

